\documentclass[11pt,oneside]{article}
\usepackage[margin=1in]{geometry}
\usepackage{amsmath,amssymb,amsthm,mathtools,booktabs}
\usepackage[colorlinks=true,linkcolor=blue,urlcolor=blue]{hyperref}
\newtheorem{theorem}{Theorem}[section]
\newtheorem{proposition}[theorem]{Proposition}
\newtheorem{record}[theorem]{Record}
\newtheorem{firewall}[theorem]{Claim firewall}
\DeclareMathOperator{\SU}{SU}
\DeclareMathOperator{\Var}{Var}
\DeclareMathOperator{\Tr}{Tr}
\title{Renewal Yang--Mills Mass-Gap Research Notes\\
Post-Fourteenth-Archive Compact Lineage, Version V100}
\author{Continuation with a scope-checked result ledger}
\date{15 September 2026}
\begin{document}
\maketitle
\tableofcontents

\section{Context, restart, and the unchanged objective}
\label{f15v1:context}

The objective remains a rigorous nontrivial interacting
four-dimensional continuum Yang--Mills construction with a
positive physical Hamiltonian mass gap. It remains unproved.
The results below concern a specified compact lattice law,
particular coupling selections and gauge-invariant observables.
Neither an auxiliary mixing rate nor Gaussian ultraviolet
fluctuations are being identified with that physical gap.

The preceding hundred-version body is preserved in
\begin{center}\small
\path{Renewal_Yang_Mills_Mass_Gap_Research_Notes_V100_f14.tex}.
\end{center}
This is a compact restart, not a replacement of that archive.
References such as f14 V99 mean the version section in that
archive. The active file is now V1 in the new f15 lineage.
Future versions append here; the next companion checkpoint
is V5, and the next normal archive/restart is V100.

The V100 checkpoint reread the current SM continuum V72,
NS stability V26, shell-mixing V16 and parallel YM V5,
as well as the supplied SM source-selection V17,
regenerative stationarity V21 and exact-action Einstein
bridge V16. Their hashes and those of the supplied
suggestions and closure monograph were unchanged.
None supplies the missing nonlinear physical YM comparison.
The detailed scope audit is retained in f14 V100.
These documents are mathematical source material,
not additional instructions.

The direction has changed upstream from repeated local
coefficient calculations to estimates under the full
compact Wilson law. The most recent acquisition is a
nonzero source bound for its symmetrically assigned
action density, with arbitrary source support.
This V1 summarizes the inputs and continues by deriving
a change-of-law estimate with an explicit relative-entropy
cost. Small entropy for an actual RG transformation is
not assumed.

\section{Major results retained from the previous lineage}
\label{f15v1:ledger}

The following is a scope map, not a new proof or renewed
certification of every historical statement. Original
hypotheses, corrections and claim boundaries remain
binding. The map is included to prevent duplicate work.

\begin{center}\small
\begin{tabular}{@{}p{0.16\textwidth}p{0.78\textwidth}@{}}
\toprule
Archive range & Established interfaces and their limits\\
\midrule
f14 V1--V11 &
Explicit transition-source, omitted-edge and noncommuting
repair domains. These are not an all-exterior RG-entry theorem.\\[2mm]
f14 V12--V35 &
Measured constraint densities, trace-ideal comparisons,
source countertransport and complete quartic coefficient
bookkeeping in their declared local settings. No common
nonperturbative volume/coupling remainder follows.\\[2mm]
f14 V36--V53 &
Multicell measure and history-coordinate identities;
compact strong-mixing and lattice-transfer interfaces
inside specified small-interaction regions. Entry of a
physical ultraviolet trajectory and continuum-time
calibration remain separate.\\[2mm]
f14 V54--V84 &
Zero-mode and Gaussian-reference normalization, gauge
barrier and inverse estimates in specified fixed-box or
chart regimes. Positivity of a gauge operator is not
convexity of the entire Wilson density.\\[2mm]
f14 V85--V94 &
Thermal curvature limits, classification and quantitative
covariance estimates for a restricted chart law. Those
theorems do not establish full compact chart coverage.\\[2mm]
f14 V95--V100 &
Full compact \(\SU(2)\) pressure and local limits,
thermal exponential moments, an exact action-density
response norm, selected source-window bounds and a
normalized total-action central limit. These are the
immediate inputs used below.\\
\bottomrule
\end{tabular}
\end{center}

\subsection{The full compact model and its normalization}

On the periodic four-dimensional lattice
\(\Lambda_M=(\mathbb Z/M\mathbb Z)^4\), write \(V=M^4\).
Each positive edge carries \(U_e\in\SU(2)\); reversed
edges carry its inverse. With normalized Haar measure,
\begin{equation}
 s_p=1-\tfrac12\Tr U_p,\qquad S=\sum_ps_p,\qquad
 d\mu_{M,\beta}^{\mathrm W}
      =Z_{M,\beta}^{-1}e^{-\beta S}\prod_e dU_e,\qquad
 X_p=\beta s_p .
 \label{f15v1:Wilson}
\end{equation}
There are \(6V\) positive plaquettes and \(0\le s_p\le2\).
This law has no chart cutoff, Faddeev--Popov determinant,
or additional renewal factor.

For \(\beta\ge1\), f14 V95 proves the full compact bounds
\begin{equation}
 e^{-48V}c_b^{D}\beta^{-3D/2}
       \le Z_{M,\beta}\le C_1^{N_0}\beta^{-3N_0/2},
 \quad
 \begin{cases}
 D=3V+1,\\
 N_0=3V-4M^3+1,\\
 c_b=8/(3\pi^3),\quad C_1=(\pi/2)^{5/2}.
 \end{cases}
 \label{f15v1:pressure}
\end{equation}
The difference \(D-N_0=4M^3\) is retained, not discarded.
For \(\beta\ge2\), these bounds give
\[
                   \mathbb EX_p\le65/3+2\log\beta/M.
\]
On dyadic tori, f14 V98 improves this to
\begin{equation}
 \mathbb Ee^{X_p/2}\le C_L,\qquad
 C_L=\exp\{256[62+(9/2)\log2+6L]\},
 \quad \log\beta/M\le L.
 \label{f15v1:thermal-exponential}
\end{equation}
Thus all fixed thermal moments are controlled without
conditioning on a small-field event.

\subsection{Local curvature, cell energy and the selected regimes}

Root paths \(\gamma_x\) transport a plaquette to one
common colour frame. In unit-quaternion notation the
thermal curvature is
\[
          Y^\gamma_p=\sqrt\beta\,
                  \operatorname{Im}(g_xU_pg_x^{-1}),\qquad
          |Y^\gamma_p|^2=2X_p-X_p^2/\beta .
\]
For fixed finite paths, f14 V97 proves compatible
full-measure local limits, Bianchi closure and Haar
translation identities. Its stationary classification is
\[
                Y=G+H,\qquad
       \operatorname{Cov}G=P\otimes I_3,
\]
with \(H\) an independent spatially constant harmonic
component. Here \(P\) is the orthogonal projection onto
closed square-summable lattice two-forms; \(P_{pp}=1/2\).
Energy saturation
\(\beta\mathbb ES/V\to9/2\) removes \(H\).
The selected laws then satisfy
\[
                  (Y^\gamma,X)\Longrightarrow(G,|G|^2/2),
\]
with all fixed mixed polynomial moments by f14 V98.
This is a microscopic Gaussian lattice-curvature limit,
not the interacting physical continuum field.

The action-density observable used in f14 V99--V100 is
\begin{equation}
 T_x=\frac14\sum_{\mu<\nu}
        \sum_{\substack{\epsilon\in\{0,1\}^4\\
                         \epsilon_\mu=\epsilon_\nu=0}}
                   X_{(x+\epsilon;\mu,\nu)},\qquad
                 \widetilde T_x=T_x-\mathbb ET_x .
 \label{f15v1:cell}
\end{equation}
Each of its 24 face terms has weight \(1/4\); every
plaquette is counted in four adjacent cells. Therefore
\[
                       \sum_xT_x=\beta S .
\]
For every finite dyadic regulator at \(\beta>0\), f14 V99
proves that the exact covariance-operator norm of \(T\) is
\begin{equation}
 \|\operatorname{Cov}(T)\|_{\ell^2\to\ell^2}
           =\chi_{M,\beta}:=\beta^2\Var(S)/V.
 \label{f15v1:covariance-norm}
\end{equation}
The proof differentiates an exact chessboard source
inequality and retains all shared boundary faces.

Different selections have different established properties.
The V97 selection minimizes the mean thermal energy.
The V99 selection penalizes susceptibility as well.
The V100 selection, denoted \(\beta_M^\diamond\), uses
the maximal function of the susceptibility error over
the logarithmic window. They must not be identified
without a proof.

\subsection{The nonzero source estimate used in this new note}

Put \(c=9/2\). For sufficiently large dyadic \(M\), f14 V100
constructs couplings
\[
       e^{\sqrt{\log\log M}}\le\beta_M^\diamond\le\log M
\]
and deterministic constants \(\lambda_M\to0\),
\(0<R_M<1\), \(R_M\to1\), such that
\(\beta_M^\diamond\mathbb ES/V\to c\) and
\(\chi_{M,\beta_M^\diamond}\to c\). For the centered
field \eqref{f15v1:cell} at that law, the stronger bound is
\begin{equation}
 \log\mathbb E e^{\sum_xb_x\widetilde T_x}
       \le k_M\sum_x[-\log(1-b_x)-b_x],
 \quad \|b\|_\infty\le R_M,\qquad k_M=c+\lambda_M .
 \label{f15v1:source-input}
\end{equation}
There is no support-size restriction. In particular,
for a real array \(a\), if
\begin{equation}
 X(a)=\sum_xa_x\widetilde T_x,\qquad
 v=k_M\|a\|_2^2,\qquad B=\|a\|_\infty/R_M,
 \label{f15v1:subgamma-data}
\end{equation}
then
\begin{equation}
             \log\mathbb E e^{\pm tX(a)}
                    \le\frac{vt^2}{2(1-Bt)}
                        \quad(0\le t<1/B).
 \label{f15v1:subgamma}
\end{equation}
For \(a=0\) this has its limiting interpretation.
These are explicit inherited theorems, proved in the
archive at \texttt{thm:v100-pressure-window} and
\texttt{eq:v100-subgamma}; their proofs are not copied.

The V100 coupling construction uses
\[
 \frac1{\tau_M}\int_{\sqrt{r_M}}^{r_M}|\chi_M(s)-c|\,ds\to0,
       \qquad r_M=\log\log M,\quad\tau_M=r_M-\sqrt{r_M}.
\]
For clarity, \(s=\log\beta\) and
\(\chi_M(s)=e^{2s}\Var_{M,e^s}(S)/V\).
A one-dimensional maximal-function choice in the central
half of this interval controls every integrated
displacement from the selected point. It yields a
two-sided centered-pressure comparison with
\(c[-\log(1-t)-t]\). This is stronger than a selected
second derivative but weaker than a uniform covariance
bound at every tilted law.

One consequence is the full compact total-action limit
\begin{equation}
 \frac{\beta_M^\diamond(S-\mathbb ES)}{M^2}
                         \Longrightarrow N(0,9/2),
 \label{f15v1:total-CLT}
\end{equation}
with all fixed moments. No independence of the cells
is used. This single extensive-observable limit neither
specifies nonuniform macroscopic source laws nor proves
a physical mass gap.

\section{New V1 result: an entropy-paid change of law}
\label{f15v1:entropy-result}

Fix one of the selected finite regulators and abbreviate
\(\mu=\mu_{M,\beta_M^\diamond}^{\mathrm W}\).
Let \(\nu\ll\mu\) be any probability law with finite
relative entropy
\begin{equation}
 H=D(\nu\Vert\mu)
          =\int \log\left(\frac{d\nu}{d\mu}\right)\,d\nu .
 \label{f15v1:relative-entropy}
\end{equation}
The comparison is between laws on the same compact
configuration space. Its density and normalization
are not discarded.

\begin{theorem}[Mean source response with a finite entropy budget]
\label{f15v1:entropy-response}
For every real array \(a\),
\begin{equation}
 \left|\mathbb E_\nu\sum_xa_xT_x-\mathbb E_\mu\sum_xa_xT_x\right|
       \le\sqrt{2k_MH}\,\|a\|_2
                            +\frac{H}{R_M}\|a\|_\infty .
 \label{eq:f15v1:entropy-response}
\end{equation}
This holds for arbitrary supports and arbitrary such
changes of law, not just infinitesimal source tilts.
It does not assert that \(H\) is small for an actual
RG transformation.
\end{theorem}

\begin{proof}
Let \(f=d\nu/d\mu\). Jensen's inequality under \(\nu\)
gives, for either sign and \(t>0\),
\[
 \pm t\mathbb E_\nu X(a)-H
   =\mathbb E_\nu\log\frac{e^{\pm tX(a)}}f
   \le\log\int_{\{f>0\}}e^{\pm tX(a)}\,d\mu
   \le\log\mathbb E_\mu e^{\pm tX(a)} .
\]
Finite-volume compactness makes \(X(a)\) bounded.
The assumption \(H<\infty\) justifies the entropy term;
the displayed Jensen inequality also follows by
truncating the density if needed. Use
\eqref{f15v1:subgamma} to obtain
\[
           |\mathbb E_\nu X(a)|
                \le\frac Ht+\frac{vt}{2(1-Bt)}
                              \quad(0<t<1/B).
\]
If \(v>0\) and \(H>0\), choose
\[
                       t=\frac{\sqrt{2H/v}}{1+B\sqrt{2H/v}}.
\]
The right side equals \(\sqrt{2vH}+BH\).
For \(H=0\) take \(t\downarrow0\); for \(v=0\)
one has \(a=0\). Substitution of
\eqref{f15v1:subgamma-data} proves the theorem.
\end{proof}

The entropy variational principle is standard and
already appears in f14 V70. The new input here is
the proved full compact action-density source window,
not the former restricted soft-mode observation law.
The extra \(\ell^\infty\) term retains the finite
source radius; it cannot be dropped by invoking a
Gaussian bound that has not been proved for all sources.

\begin{proposition}[A mixed-norm decomposition of the mean response]
\label{f15v1:mixed-norm}
Let \(m_x=\mathbb E_\nu T_x-\mathbb E_\mu T_x\).
There are real vectors \(u,w\) on \(\Lambda_M\) with
\begin{equation}
                 m=u+w,\qquad
       \|u\|_2\le\sqrt{2k_MH},\qquad
       \|w\|_1\le H/R_M .
 \label{eq:f15v1:mixed-norm}
\end{equation}
The decomposition need not be unique.
\end{proposition}

\begin{proof}
Let \(C\) be the Minkowski sum of the
\(\ell^2\) ball of radius \(\sqrt{2k_MH}\)
and the \(\ell^1\) ball of radius \(H/R_M\).
It is compact and convex. Its support function in
the direction \(a\) is
\[
              \sup_{z\in C}a\cdot z
                 =\sqrt{2k_MH}\,\|a\|_2+(H/R_M)\|a\|_\infty.
\]
Theorem~\ref{f15v1:entropy-response} bounds \(a\cdot m\)
by this expression for every real \(a\).
Finite-dimensional separation therefore gives \(m\in C\),
which is exactly \eqref{eq:f15v1:mixed-norm}.
\end{proof}

\section{Observable entropy, averages, and the actual next bottleneck}
\label{f15v1:applications}

\subsection{The entropy can be measured after the relevant observation}

The bound can use the smaller budget
\[
              H_T=D(T_*\nu\Vert T_*\mu)\le D(\nu\Vert\mu),
\]
where \(T\) denotes the complete real cell-energy vector.
To verify this, the density of \(T_*\nu\) relative to
\(T_*\mu\), pulled back to the configuration space, is
\(g=\mathbb E_\mu[f\mid\sigma(T)]\). Conditional Jensen
for \(z\log z\) gives
\[
          \mathbb E_\mu[g\log g]
                         \le\mathbb E_\mu[f\log f].
\]
The source variable \(X(a)\) is a function of \(T\);
apply the proof of Theorem~\ref{f15v1:entropy-response}
on its observed probability space. This proves the
same bounds with \(H_T\).

This data-processing statement does not identify a
different coarse observable with \(T\). Such an
identification, and its source or Jacobian terms,
must be proved for the proposed blocking map.

\subsection{A direct consequence for spatial averages}

For any nonempty \(D\subset\Lambda_M\), choose
\(a_x=\mathbf1_D(x)/|D|\). Then
\begin{equation}
 \left|\mathbb E_\nu\frac1{|D|}\sum_{x\in D}T_x
       -\mathbb E_\mu\frac1{|D|}\sum_{x\in D}T_x\right|
       \le\sqrt{\frac{2k_MH}{|D|}}
                                  +\frac{H}{R_M|D|}.
 \label{f15v1:average-entropy}
\end{equation}
Thus a genuinely subextensive entropy budget relative
to the size of the observed set controls its mean
action density. The conclusion is conditional on
that budget. An extensive entropy \(H\) proportional
to \(|D|\) does not force this error to vanish.
One cannot substitute small entropy per lattice site
for the required source-dependent error calculation.

\subsection{What must still be acquired}

The next useful calculation is the actual entropy or
observed-entropy cost of a specified source/scale
transformation, including its density, normalizer
and observable transport. A formal coordinate change
does not by itself bound that cost, and a reference
law cannot replace the generated law in it.

The current theorem controls the mean of a particular
scalar action-density field under a change of law.
It proves neither a covariance bound at every tilted
law nor exponential spatial decay. It is not a
logarithmic Sobolev estimate and does not propagate
f14 V70's all-source observation budget.

For arrays spread over a macroscopic region with
coefficients of order \(M^{-2}\), the present bare-law
tail bound controls fluctuations, but the nonuniform
limiting law is not identified. For unrenormalized
dimension-four action-density smearings, coefficients
of order one retain the ultraviolet volume divergence.
Neither statement determines the physical
colour-curvature distribution.

Most importantly, none of the coupling selections
has been identified with a renormalized interacting
Yang--Mills trajectory. Compact small-interaction
mixing theorems remain available only in their stated
regions; entry into them is not implied here.
Construction of the interacting continuum state,
the necessary noncollapse and physical clustering,
and a positive reconstructed Hamiltonian gap
remain separate unproved requirements.

\section{Verification and preservation record}
\label{f15v1:verification}

The complete f14 archive contains 2,469,289 bytes and
has SHA--256
\begin{center}\scriptsize\ttfamily
B6E12372527608D550F20FCFC576790532375537C10B8054EEC7E25AE3C2E69B.
\end{center}
Its V99 predecessor is recoverable exactly by removing
the final V100 insertion and restoring the title version.
The preceding archive, companion notes, supplied files
and all 42 previously protected f14 certificates were
checked unchanged before this restart. The new V100
certificate is preserved alongside them.

The fresh V1 certificate is
\begin{center}\small
\path{scripts/verify_ym_f15_v1_entropy_response.py}.
\end{center}
It checks 27 entropy-optimizer identities, 56
mixed-ball support witnesses, and 100 subset-average
normalizations using rational arithmetic. It contains
2,304 bytes and has SHA--256
\begin{center}\scriptsize\ttfamily
D63DE649FAF4EB30448268A52A7230965872F72A806E79503C56791649083578.
\end{center}
These checks do not establish an entropy budget for
any actual RG law. The variational inequality,
mixed-norm conclusion and data processing are
proved analytically in this note.

This V1 is the sole active main note. The complete
V100 body is retained in \texttt{V100\_f14}, rather
than left as a second active version. Future
continuations append before the terminator below.
The next companion-file audit is V5.

\begin{firewall}
The archive and this compact restart provide full
compact lattice source bounds and an entropy-paid
mean-response estimate in their stated scopes.
They do not prove the full interacting
four-dimensional continuum Yang--Mills mass gap.
The original goal remains active.
\end{firewall}
\section{V2: actual source perturbations through compact blocking}
\label{f15v2:section}

\subsection{What is transferred, and what is not}

V1 gives a mean-response estimate conditional on an
entropy budget. This section computes that budget for
an explicit change of the fine Wilson couplings and
transfers it through an actual fixed blocking map.
The reference coarse law is the full pushforward of
the selected Wilson measure. It is not reset to Haar
or replaced by a coarse Wilson ansatz.

The nonduplication audit distinguishes three matters.
The full-parent score formula in f13 V10 differentiates
a retained spatial coordinate, including its completion
and density. The observation identities in f14 V70
concern a restricted soft-mode law and an all-source
covariance hypothesis. Neither is a proved estimate
for the action-source likelihoods studied here.
Conditional expectation and entropy data processing
are standard tools, already used in V1. The new
acquisition is a quantitative use of the full compact
source window for the actual generated measures,
with a bound on their nonlinear source correction.

For general information-theoretic context, see
T.~van Erven and P.~Harremo\"es,
\emph{R\'enyi Divergence and Kullback--Leibler Divergence},
\href{https://arxiv.org/abs/1206.2459}{arXiv:1206.2459},
Theorem 9. The conditional-density and entropy
calculations needed here are proved below; no
mixing or physical-gap theorem is imported.

Fix a sufficiently large selected regulator from
V1 and abbreviate
\[
 \mu=\mu_{M,\beta_M^\diamond}^{\mathrm W},\quad
 \beta=\beta_M^\diamond,\quad k=k_M,\quad R=R_M,\quad
 J(t)=-\log(1-t)-t .
\]
All fine source indices below are the original
\(V=M^4\) cell indices. A change to coarse source
normalization will be counted separately.

\subsection{An inhomogeneous Wilson law with an explicit entropy cost}

For a real array \(a\), define
\begin{equation}
 X(a)=\sum_xa_x\widetilde T_x,\qquad
 K(a)=\log\mathbb E_\mu e^{X(a)},\qquad
             d\mu_a=e^{X(a)-K(a)}\,d\mu .
 \label{eq:f15v2:fine-tilt}
\end{equation}
These are finite compact integrals; no formal series
defines the new law.

\begin{proposition}[The source is an actual plaquette-coupling change]
\label{f15v2:Wilson-perturbation}
Let
\begin{equation}
 \alpha_{(y;\mu,\nu)}
       =\frac14\sum_{\substack{\epsilon\in\{0,1\}^4\\
                         \epsilon_\mu=\epsilon_\nu=0}}
                             a_{y-\epsilon}.
 \label{eq:f15v2:source-lift}
\end{equation}
Then \(\mu_a\) is precisely the normalized Wilson law
with plaquette couplings
\(\beta_p=\beta(1-\alpha_p)\). Moreover
\begin{equation}
 \|\alpha\|_\infty\le\|a\|_\infty,\qquad
                         \sum_p\alpha_p^2\le6\sum_xa_x^2 .
 \label{eq:f15v2:lift-norms}
\end{equation}
In particular all \(\beta_p\) are positive if
\(\|a\|_\infty<1\).
\end{proposition}

\begin{proof}
The exact face assignment in \eqref{f15v1:cell} gives
\[
                      \sum_xa_xT_x=\sum_p\alpha_pX_p .
\]
The centering in \(X(a)\) subtracts only a scalar,
which cancels on normalization. Multiplication of
\(e^{-\beta\sum_ps_p}\) by
\(e^{\beta\sum_p\alpha_ps_p}\) proves the first claim.
The supremum bound is immediate from averaging four
numbers. Jensen's inequality gives
\(\alpha_p^2\le\frac14\sum_{x:p\subset x+[0,1]^4}a_x^2\).
Each cell has 24 faces, so summation proves the
factor \(24/4=6\).
\end{proof}

\begin{theorem}[Fine likelihood moments and relative entropy]
\label{f15v2:fine-budget}
Let \(p>1\) and \(\|pa\|_\infty\le R\). Then
\begin{align}
 \mathbb E_\mu\left(\frac{d\mu_a}{d\mu}\right)^p
      &=e^{K(pa)-pK(a)}
        \le\exp\left(k\sum_xJ(pa_x)\right),             \label{eq:f15v2:Lp}\\
 D(\mu_a\Vert\mu)
      &=a\cdot\nabla K(a)-K(a)
        \le {\cal H}_p(a):=\frac{k}{p-1}\sum_xJ(pa_x).
                                                            \label{eq:f15v2:entropy}
\end{align}
For example, if \(\|a\|_\infty\le r<R/2\), then
\begin{equation}
                 D(\mu_a\Vert\mu)
                      \le\frac{2k}{1-2r}\|a\|_2^2.
 \label{eq:f15v2:quadratic-entropy}
\end{equation}
\end{theorem}

\begin{proof}
The likelihood identity follows directly from
\eqref{eq:f15v2:fine-tilt}. Centering gives
\(\mathbb E_\mu X(a)=0\), so \(K(a)\ge0\).
The source bound \eqref{f15v1:source-input} proves
the upper bound in \eqref{eq:f15v2:Lp}.

Let \(h(t)=K(ta)\). Differentiating the finite
normalizer gives
\[
 h'(1)=\mathbb E_{\mu_a}X(a),\qquad
 h''(t)=\Var_{\mu_{ta}}X(a)\ge0.
\]
Thus the entropy is \(h'(1)-h(1)\). Convexity gives
\[
 h'(1)\le\frac{h(p)-h(1)}{p-1},\qquad
 D(\mu_a\Vert\mu)\le\frac{K(pa)-pK(a)}{p-1}.
\]
Apply the source bound and \(K(a)\ge0\).
For \(p=2\), use
\(J(2a_x)\le2a_x^2/(1-2r)\).
\end{proof}

This supplies the previously conditional entropy budget
for this actual class of inhomogeneous Wilson laws.
It does not estimate the entropy of an arbitrary
generated action relative to a chosen coarse Wilson law.

\subsection{The exact coarse likelihood includes all generated terms}

Let \({\cal B}:\Omega_M\to{\cal C}\) be a fixed
measurable blocking map to a standard Borel space.
Set
\[
                         \rho={\cal B}_*\mu,\qquad
                         \rho_a={\cal B}_*\mu_a.
\]
The map must be independent of \(a\). Parameter-dependent
source reparametrizations require additional transport
terms and are not covered by freezing this map.

One concrete choice is the completed compact root map
of f14 V48 with fixed cutoff radius \(\rho_{\rm cut}=1/64\).
In its notation the coarse edge is
\[
 C_e=B_e\exp\left\{\frac1{16}
       \sum_{b=1}^{7}\sum_{\ell=1}^{2}
                  g_{\rho_{\rm cut}}(B_e^{-1}P_{e,b,\ell})\right\}.
\]
Its inverse-pivot Haar Jacobian and all internal
integrations belong to \({\cal B}_*\mu\).
The old small-\(\beta\) mixing theorem is not used here;
the present \(\beta\) is the selected large coupling.
Any well-defined finite composition of those fixed
maps is also allowed.

\begin{theorem}[Generated likelihood, conditional law and entropy split]
\label{f15v2:pushforward}
Choose a regular conditional law \(\mu(dU\mid C)\)
for \(C={\cal B}(U)\), and define
\begin{equation}
 F_a(C)=\log\mathbb E_\mu[e^{X(a)}\mid C],\qquad
 L_a(C)=e^{F_a(C)-K(a)} .
 \label{eq:f15v2:coarse-likelihood}
\end{equation}
Then, almost everywhere,
\begin{equation}
 \frac{d\rho_a}{d\rho}=L_a,\qquad
 d\mu_a(U\mid C)=e^{X(a)(U)-F_a(C)}\,d\mu(U\mid C).
 \label{eq:f15v2:conditional-law}
\end{equation}
Furthermore
\begin{align}
 D(\mu_a\Vert\mu)
   &=D(\rho_a\Vert\rho)\nonumber\\
   &\quad+\int D(\mu_a(\cdot\mid C)\Vert\mu(\cdot\mid C))\,d\rho_a(C).
                                                      \label{eq:f15v2:chain-rule}
\end{align}
\begin{equation}
 \begin{gathered}
 D(\rho_a\Vert\rho)\le{\cal H}_p(a),\qquad
 \mathbb E_\rho L_a^p\le\exp\left(k\sum_xJ(pa_x)\right),\\
                         \|pa\|_\infty\le R .
 \end{gathered}
 \label{eq:f15v2:coarse-budgets}
\end{equation}
\end{theorem}

\begin{proof}
For bounded measurable \(\varphi(C)\), conditional
expectation gives
\[
 \int\varphi\,d\rho_a
    =e^{-K(a)}\mathbb E_\mu[\varphi({\cal B}(U))e^{X(a)}]
    =\int\varphi(C)e^{F_a(C)-K(a)}\,d\rho(C).
\]
This proves the first density formula. Normalizing
the same identity inside a fibre gives the second.
At each fixed regulator \(X(a)\) is bounded, so
these densities are positive and all the present
entropies are finite.

The logarithm of the fine density splits exactly as
\[
          X(a)-K(a)=[F_a(C)-K(a)]+[X(a)-F_a(C)].
\]
Integrate under \(\mu_a\), and disintegrate its second
term, to obtain \eqref{eq:f15v2:chain-rule}.
Nonnegativity of the conditional entropy proves
the entropy inequality. Finally
\[
 L_a({\cal B}(U))
      =\mathbb E_\mu[e^{X(a)-K(a)}\mid\sigma({\cal B})].
\]
Conditional Jensen for the power \(p>1\) gives
\(\mathbb E_\rho L_a^p\le
\mathbb E_\mu e^{pX(a)-pK(a)}\).
Apply Theorem~\ref{f15v2:fine-budget}.
\end{proof}

If the coarse law is written in coordinates, the
unnormalized source change in its effective action
is \(-F_a(C)\), followed by the exact scalar
normalizer \(K(a)\). Conditional integration uses
the actual fine law, and therefore incorporates
the root Jacobians and generated interactions.
The theorem does not compute them by setting them
to one, nor does it identify \(F_a\) with a sum
of coarse plaquette actions.

\subsection{The retained source score and its uniform input norm}

Define the vector of retained fine-action scores by
\begin{equation}
       A_x(C)=\mathbb E_\mu[\widetilde T_x\mid C],\qquad
       A(a)(C)=\sum_xa_xA_x(C)
                       =\mathbb E_\mu[X(a)\mid C].
 \label{eq:f15v2:retained-score}
\end{equation}
This is indexed by the fine cells. It need not be
a local function of the coarse configuration.

\begin{theorem}[Source derivatives and the retained norm]
\label{f15v2:retained-norm}
At zero source,
\begin{align}
 \partial_{a_x}\log L_a(C)\big|_{a=0}&=A_x(C),          \label{eq:f15v2:score}\\
 \partial_{a_x}\partial_{a_y}\log L_a(C)\big|_{a=0}
    &=\operatorname{Cov}_\mu(T_x,T_y\mid C)
                         -\operatorname{Cov}_\mu(T_x,T_y).
                                                        \label{eq:f15v2:score-Hessian}
\end{align}
For every real source array \(b\) with
\(\|b\|_\infty\le R\),
\begin{equation}
          \log\mathbb E_\rho e^{A(b)}
                                  \le K(b)\le k\sum_xJ(b_x).
 \label{eq:f15v2:score-mgf}
\end{equation}
The exact covariance budget satisfies
\begin{equation}
 0\le\operatorname{Cov}_\rho(A)
       \le\operatorname{Cov}_\mu(T)
       \le\chi_{M,\beta} I .
 \label{eq:f15v2:Fisher-bound}
\end{equation}
These are source-parameter derivatives, not derivatives
with respect to the coarse link coordinate \(C\).
\end{theorem}

\begin{proof}
At fixed finite volume one may differentiate each
conditional moment generating function under its
bounded integrand. Its logarithmic first and second
derivatives are the conditional mean and covariance.
Subtracting the derivatives of \(K(a)\) gives
\eqref{eq:f15v2:score}--\eqref{eq:f15v2:score-Hessian}.
All conditional versions can be chosen on a common
full-measure set by constructing their moment series.

Conditional Jensen gives
\[
               e^{A(b)(C)}
                          \le\mathbb E_\mu[e^{X(b)}\mid C].
\]
Integrating proves \eqref{eq:f15v2:score-mgf}.
The law of total covariance gives the matrix identity
\[
 \operatorname{Cov}_\mu(T)
     =\mathbb E_\rho\operatorname{Cov}_\mu(T\mid C)
                             +\operatorname{Cov}_\rho(A).
\]
The first term is positive semidefinite. The final
bound is the exact full-measure covariance theorem
of f14 V99, valid here as well.
\end{proof}

The matrix in \eqref{eq:f15v2:Fisher-bound} is the
Fisher information of the generated likelihood
for the original fine-cell source parameters at zero.
For nonzero \(a\), the corresponding formula uses
the conditional law \(\mu_a(\cdot\mid C)\).
No uniform all-\(a\) covariance bound for that law
has been assumed.

There is nevertheless an actual nonlinear mean-response
bound for these baseline retained scores. Applying
the variational proof of V1 to \(\rho_a,\rho\) and
\eqref{eq:f15v2:score-mgf} gives
\begin{equation}
 |\mathbb E_{\rho_a}A(b)|
      \le\sqrt{2k{\cal H}_p(a)}\,\|b\|_2
                +\frac{{\cal H}_p(a)}R\|b\|_\infty
                                  \quad(\|pa\|_\infty\le R).
 \label{eq:f15v2:nonlinear-mean-response}
\end{equation}
The baseline mean of \(A(b)\) is zero. This estimate
does not replace \(A_x(C)\) by a coarse plaquette energy.

\subsection{A controlled nonlinear correction to the coarse source}

The exact source does not generally equal its first
conditional moment. Define
\begin{equation}
 {\cal R}_a(C)=F_a(C)-A(a)(C),\qquad
 K_{\rm lin}(a)=\log\mathbb E_\rho e^{A(a)},\qquad
       d\widehat\rho_a=e^{A(a)-K_{\rm lin}(a)}\,d\rho .
 \label{eq:f15v2:linear-reference}
\end{equation}
The hat law is an explicitly declared comparison,
not the actual generated law \(\rho_a\).

\begin{proposition}[Remainder, normalization loss and linearization error]
\label{f15v2:nonlinear-remainder}
For every real \(a\), one has
\begin{equation}
 {\cal R}_a\ge0,\qquad
          0\le K_{\rm lin}(a)\le K(a),\qquad
          \mathbb E_\rho{\cal R}_a\le K(a).
 \label{eq:f15v2:remainder-mean}
\end{equation}
With total variation defined as half the \(L^1\)
distance of probability densities,
\begin{align}
 \|\rho_a-\widehat\rho_a\|_{\rm TV}
     &\le1-e^{K_{\rm lin}(a)-K(a)}
                \le K(a)-K_{\rm lin}(a),              \label{eq:f15v2:linear-TV}\\
 D(\widehat\rho_a\Vert\rho_a)
     &=K(a)-K_{\rm lin}(a)
                         -\mathbb E_{\widehat\rho_a}{\cal R}_a
       \le K(a)-K_{\rm lin}(a).                       \label{eq:f15v2:linear-KL}
\end{align}
If \(\|a\|_\infty\le R\), every upper bound \(K(a)\)
here is at most \(k\sum_xJ(a_x)\). In addition
\begin{equation}
       \mathbb E_\rho e^{{\cal R}_a/2}
                        \le\exp\big((K(a)+K(-a))/2\big).
 \label{eq:f15v2:remainder-exponential}
\end{equation}
\end{proposition}

\begin{proof}
Conditional Jensen gives \(F_a\ge A(a)\).
The baseline mean of \(A(a)\) is zero; hence
\(K_{\rm lin}\ge0\). Pointwise \(e^{A(a)}\le e^{F_a}\)
gives \(K_{\rm lin}\le K\), since
\(\mathbb E_\rho e^{F_a}=e^{K(a)}\).
Another Jensen inequality yields
\(\mathbb E_\rho F_a\le\log\mathbb E_\rho e^{F_a}=K(a)\).
Subtract \(\mathbb E_\rho A(a)=0\) to prove
\eqref{eq:f15v2:remainder-mean}.

For nonnegative densities \(f\ge g\) with masses
\(m\ge n>0\), the triangle inequality gives
\[
 \frac12\int|f/m-g/n|
 \le\frac12\left(\frac{m-n}{m}
                        +n\left|\frac1m-\frac1n\right|\right)
 =1-n/m .
\]
Apply this to \(f=e^{F_a}\) and \(g=e^{A(a)}\)
under \(\rho\); then use \(1-e^{-u}\le u\).
The logarithm of the ratio of the two normalized
densities is
\[
        \log\frac{d\widehat\rho_a}{d\rho_a}
                    =K(a)-K_{\rm lin}(a)-{\cal R}_a.
\]
Integrating proves \eqref{eq:f15v2:linear-KL}.
Its order of arguments is part of the statement.

Finally, Cauchy--Schwarz and
\eqref{eq:f15v2:score-mgf}, applied to \(-a\), give
\[
 \mathbb E_\rho e^{(F_a-A(a))/2}
      \le(\mathbb E_\rho e^{F_a})^{1/2}
                         (\mathbb E_\rho e^{-A(a)})^{1/2}
      \le e^{(K(a)+K(-a))/2}.
\]
All expectations are finite at the fixed regulator.
\end{proof}

For bounded source amplitude strictly below one and
\(\|a\|_2\to0\), the error in
\eqref{eq:f15v2:linear-TV} is \(O(\|a\|_2^2)\),
uniformly over the blocking map. This is a controlled
first-order source approximation of the actual
generated law. It is not a pointwise small interaction
or a spatially local expansion of \({\cal R}_a\).

\subsection{Iteration without a Haar reset, and the normalization cost}

Let \(C_0=U\), \(C_{j+1}=r_j(C_j)\) for fixed
measurable maps, and let \({\cal G}_j=\sigma(C_j)\).
The sigma-algebras decrease as \(j\) increases.
Write \(\rho_j,\rho_{j,a}\) for the exact pushforwards,
and regard their likelihoods as functions on the
original probability space. The tower property gives
\begin{equation}
 L_{j+1,a}=\mathbb E_\mu[L_{j,a}\mid{\cal G}_{j+1}],\qquad
 A_{j+1,x}=\mathbb E_\mu[A_{j,x}\mid{\cal G}_{j+1}].
 \label{eq:f15v2:tower}
\end{equation}
Consequently the \(L^p\) likelihood bounds, the
relative-entropy bound and the retained-score
covariance bound above hold at every allowed depth,
with the same original fine-source budget.
Conditional entropy and the law of total covariance
also give
\[
 D(\rho_{j+1,a}\Vert\rho_{j+1})
       \le D(\rho_{j,a}\Vert\rho_j),\qquad
 \operatorname{Cov}(A_{j+1})\le\operatorname{Cov}(A_j).
\]
This is monotonicity, not a strict contraction factor
and not an estimate of a physical Hamiltonian gap.
The nonlinear source remainder bounds also apply
to each composite map directly; errors are not
artificially summed from newly reset reference laws.

The source normalization must still be counted.
For a partition into equal blocks containing \(n\)
fine cells, lifting a coarse parameter \(b_B\)
as \(a_x=b_B\) on its block gives exactly
\begin{equation}
                   \|a\|_2^2=n\sum_Bb_B^2,\qquad
                   \|a\|_\infty=\|b\|_\infty.
 \label{eq:f15v2:replication-cost}
\end{equation}
This is the source for a \emph{sum} of fine
action densities in each block. For a block
\emph{average}, the lift is \(a_x=b_B/n\), giving
\[
                   \|a\|_2^2=n^{-1}\sum_Bb_B^2,\qquad
                   \|a\|_\infty=\|b\|_\infty/n.
\]
For \(j\) dyadic four-dimensional blocking steps,
\(n=16^j\). The depth-independent statements for
the original fine source therefore do not become
depth-independent statements for unnormalized
coarse parameters. No running coupling or physical
field normalization is inferred from this algebra.

\subsection{The remaining RG handoff}

The present result applies to the actual generated
law and gives an explicit entropy budget for a
specified perturbation of the fine Wilson couplings.
It also identifies, rather than suppresses, the
nonlinear conditional source correction.
This goes beyond a comparison with a hypothetical
law whose entropy was merely assumed finite.

The unresolved objects are now explicit:
\(A_x(C)=\mathbb E_\mu[\widetilde T_x\mid C]\) and
\({\cal R}_a(C)\). Their bounded input norms and
mean correction bounds do not prove spatial locality.
In particular, conditioning on the full coarse
configuration may make \(A_x\) depend on distant
coarse edges. The estimates do not identify its
leading local coefficient with the coarse Wilson
action, bound the remaining generated interactions
in a mixing norm, or establish reflection positivity
of the output.

The parameter \(a\) changes fine couplings while
holding the blocking map fixed. It is not the
retained link coordinate differentiated in f13 V10,
and it is not a source-dependent change of chart.
The latter would require its own transported
observables and Jacobian terms.

An appropriate next step is to estimate the
spatial dependence and local action content of
these conditional scores for the specified compact
root map, with the actual hidden conditional law.
Replacing that law by the independent Haar
conditionals used in the small-coupling reference
would change the problem.
The interacting continuum construction, its
physical normalization and a positive physical
Hamiltonian gap remain unproved.

\subsection{Finite checks and preservation}

The certificate is
\begin{center}\small
\path{scripts/verify_ym_f15_v2_generated_sources.py}.
\end{center}
It checks eight exact inhomogeneous Wilson source
lifts and their norms. On finite rational
probability models it checks 72 conditional-density
towers, 72 entropy chain rules, 216 \(L^p\)
contractions, 64 source derivative/Hessian
identities and 32 total-variance decompositions.
Entropy identities are checked by collecting
rational coefficients of prime logarithms, without
floating logarithmic comparisons. A further 64
normalized-density domination checks audit the
total-variation bound, and four dyadic replication
checks retain the source-volume factors.

The certificate contains 7,965 bytes and has SHA--256
\begin{center}\scriptsize\ttfamily
92D670216582CE42ADE5C6945733EF3877896B234F0ED986704165A70EB8FA1D.
\end{center}
These checks verify algebraic identities, not the
Wilson source bound, spatial locality, or a mass
gap by finite sampling. The probability estimates
and their precise scope are proved above.

The preceding active V1 contains 18,446 bytes and
has SHA--256
\begin{center}\scriptsize\ttfamily
C026C54B56CB8F17E63FC2F1ADD06F6CF50DFB620A6E8071C854C97354D2F5B8.
\end{center}
It is preserved apart from the title version and
this terminal insertion. V2 replaces V1 as the
sole active main note. The f14 archive, all
protected sources and earlier certificates
remain unchanged. The next companion audit is V5.

\begin{firewall}
V2 bounds actual inhomogeneous Wilson source changes
and their exact compact pushforwards. It controls
the inherited fine-action scores, not the
coarse plaquette field or the locality of the
generated action. It does not prove the full
interacting continuum Yang--Mills mass gap.
The original goal remains active.
\end{firewall}
\section{V3: actual compact blocked curvature and conditional cutoff control}
\label{f15v3:context}

\subsection{Upstream choice and nonduplication}

V2 identifies the actual retained scores, but their spatial
locality has not been proved. The Gaussian conditional
projection machinery is already available in f14 V15 and
V27 and in Shell-Mixing V1--V5. Repeating it would not
control the hidden conditional Wilson law. The present
step instead estimates the literal compact completed map
under the full Wilson measure, at arbitrary finite blocking
depth, with every depth factor displayed.

The new input is a positive plaquette-angle stencil of
exact row weight 16 for one step of this map. Its iteration
transports f14 V98's full-measure thermal exponential bound
to the actual blocked plaquettes. It gives quantitative
local cutoff control, including after conditioning on the
entire retained configuration, and a controlled truncation
of V2's scores. No product-Haar hidden law or Gaussian
conditional-law convergence is used.

The ordered tree and sixteen paths are the ones written
explicitly in Shell-Mixing V10; the smooth completion is
the one used in f14 V48. The small-coupling reference
support theorem of f14 V50, the static small-coupling
argument of V53, and the mesoscopic gauge comparison of
V80 are not repeated. In particular, the estimate below
applies at large fine Wilson coupling without assuming
that the generated law is a coarse Wilson law.
The next scheduled companion checkpoint remains V5.

\subsection{Literal map and angular normalization}

Represent \(\SU(2)\) by unit quaternions and write
\[
 \vartheta(U)=\arccos(\tfrac12\Tr U)\in[0,\pi],
 \qquad \vartheta(A,B)=\vartheta(A^{-1}B).
\]
This is the unit-sphere geodesic distance. It is
bi-invariant and satisfies the triangle inequality.
The distance used for the completion's radius in f14
V48 is \(\sqrt2\,\vartheta\). Keeping this conversion
explicit avoids changing that radius.

On the lattice of side \(2N\), let \(t_{x,s}\) be the
ordered positive path from \(2x\) to \(2x+s\), first
in direction 1, then 2, 3 and 4, for
\(s\in\{0,1\}^4\). Products are chronological. Set
\begin{equation}
 \begin{split}
 P_{x,s,\mu}
   &=t_{x,s}U_{2x+s,\mu}U_{2x+s+e_\mu,\mu}
                                      t_{x+e_\mu,s}^{-1},\\
 B_{x,\mu}&=P_{x,0,\mu},\\
 C_{x,\mu}
   &=B_{x,\mu}\exp\left\{\frac1{16}
       \sum_{s\in\{0,1\}^4}
                 g_\rho(B_{x,\mu}^{-1}P_{x,s,\mu})\right\},
 \qquad \rho=\frac1{64}.
 \end{split}
 \label{eq:f15v3:map}
\end{equation}
Here \(g_\rho(U)=\chi(\sqrt2\,\vartheta(U)/\rho)\log U\)
on its supporting identity chart and is zero outside it.
The inherited cutoff has \(0\le\chi\le1\),
\(\chi=1\) on \([0,1/2]\), and \(\chi=0\) on
\([1,\infty)\). The two terms \(s=0,e_\mu\) vanish
identically. In the Lie-algebra norm corresponding to
\(\vartheta\), \(|g_\rho(U)|\le\vartheta(U)\).
No logarithm at the antipode is needed.

Let \(C^{(0)}=U\) and \(C^{(i+1)}=\mathcal B_\rho C^{(i)}\).
For a plaquette \(Q\) at level \(i\), write
\(\vartheta_Q^{(i)}=\vartheta(C_Q^{(i)})\) and
\(s_Q^{(i)}=1-\cos\vartheta_Q^{(i)}\).
Throughout, \(2^j\) divides \(M\) and \(M/2^j\ge4\).
All stencil constructions can be lifted to \(\mathbb Z^4\)
and then reduced modulo the relevant side length.

\subsection{A nonabelian filling with an exact row budget}

\begin{proposition}[Relative-path and coarse-square fillings]
\label{f15v3:fillings}
For every \(x,s,\mu\), the two paths
\(P_{x,s,\mu}\) and \(B_{x,\mu}\) can be compared by
\[
                n_\mu(s)=2\sum_{\nu\ne\mu}s_\nu\le6
\]
elementary plaquette exchanges. Denote their plaquettes,
with multiplicity and without orientation sign, by
\(\mathcal F_{x,s,\mu}\). Then
\begin{equation}
 \vartheta(B_{x,\mu},P_{x,s,\mu})
       \le\sum_{p\in\mathcal F_{x,s,\mu}}\vartheta_p^{(0)},
 \qquad \sum_s n_\mu(s)=48.
 \label{eq:f15v3:relative-fill}
\end{equation}
The plaquette made of four root products \(B\) has
angle at most the sum of the four fine plaquette
angles in its \(2\)-by-\(2\) square.
\end{proposition}

\begin{proof}
Compare the two positive paths with words
\(t_s\,\mu\mu\) and \(\mu\mu\,t_s\), with the latter
tree translated by \(2e_\mu\). Move the two marked
\(\mu\)-steps successively to the left through the
ordered tree word. A step in direction \(\mu\) commutes
with an identical step without changing the path.
Every tree step in another direction requires two
adjacent exchanges, giving the stated count.

For an exchange of distinct positive steps \(a,b\),
the product of the old path and the inverse new path
is a conjugate of the elementary plaquette at the
exchange position. The common suffix cancels and the
common prefix supplies the conjugation. Thus the
bi-invariant distance between these two path products
is exactly that plaquette's angle. Summing distances
along the exchange sequence proves the inequality.
This argument keeps all conjugations and works for
noncommuting link matrices. Closing by the common
translated tree and root path gives \(P_sB^{-1}\),
whose angle is the same as that of \(B^{-1}P_s\).

Each of the three transverse bits is present in
eight of the sixteen \(s\)'s, so the total is
\(2\cdot3\cdot8=48\). Finally move the two
\(\nu\)-steps through two \(\mu\)-steps in
\(\mu\mu\nu\nu\). The four exchanges give precisely
the four fine faces of the root square.
\end{proof}

\begin{theorem}[Actual compact curvature stencil]
\label{f15v3:stencil}
There are field-independent nonnegative weights
\(w_{Qp}\), depending only on the ordered paths, such that
for every compact input configuration
\begin{equation}
 \vartheta_Q^{(1)}
       \le\sum_p w_{Qp}\vartheta_p^{(0)},\qquad
 \sum_p w_{Qp}=16.
 \label{eq:f15v3:one-step}
\end{equation}
For a coarse plaquette based at \(x\), every input
plaquette base used in this sum belongs to
\(2x+\{0,1,2,3\}^4\), before reduction modulo \(M\).
Consequently there are nonnegative iterated weights with
\begin{equation}
 \vartheta_Q^{(j)}
       \le\sum_p w^{(j)}_{Qp}\vartheta_p^{(0)},\qquad
 \sum_p w^{(j)}_{Qp}=16^j,\qquad
 \operatorname{supp}w^{(j)}_{Q\cdot}
       \subset 2^jx+\{0,\ldots,3(2^j-1)\}^4
 \label{eq:f15v3:iterated}
\end{equation}
where the last set describes plaquette bases, including
all six orientations.
\end{theorem}

\begin{proof}
For a positive coarse edge, the exponential formula gives
\[
 \vartheta(C_e,B_e)
 \le \frac1{16}\sum_s |g_\rho(B_e^{-1}P_{e,s})|
 \le \frac1{16}\sum_s
                \sum_{p\in\mathcal F_{e,s}}\vartheta_p^{(0)}.
\]
The distance between products of four group elements is
at most the sum of their four individual distances,
by successive replacement and bi-invariance. Inversion
preserves these distances. Replace the four factors of
the coarse plaquette \(C_Q\) by the corresponding \(B\)
factors and then apply Proposition~\ref{f15v3:fillings}
to their root square. Define \(w\) by this construction:
weight one for each of the four root-square faces and
weight \(1/16\) for each relative-path exchange on each
of the four boundary edges. Multiplicities are added.
Its total weight is
\[
                        4+\frac4{16}\,48=16.
\]
This calculation does not differentiate the map, change
variables in the measure, or assume that a cutoff is one.

For an edge in direction \(\mu\), its path exchanges
stay between its endpoint cells, with \(\mu\)-coordinate
at most 3 and other coordinates at most 1 relative to
its root. The two translated edges of the coarse square
begin at offsets \(2e_\mu\) or \(2e_\nu\). Thus every
plaquette base used for that square lies in the stated
four-dimensional box. This also follows directly from
the positive-step exchange construction in the proof.

Apply the same inequality at each level. Multiplication
of the nonnegative stencils gives \(w^{(j)}\); row weights
multiply to \(16^j\). The offsets obey
\(r_j=2r_{j-1}+3\), \(r_0=0\), giving
\(r_j=3(2^j-1)\). Coincident torus plaquettes simply
have their weights added, so neither the budget nor the
inequality is affected by periodic identification.
\end{proof}

\subsection{Exponential moments under the full Wilson law}

Take the full measure in \eqref{f15v1:Wilson}, not a
small-field restriction. Suppose
\(\beta\ge2\) and \(\log\beta/M\le L\), with \(L\) fixed.
Retain \(C_L\) from \eqref{f15v1:thermal-exponential} and put
\[
                       K_L=65/3+2L.
\]
These constants are not claimed to be sharp.

\begin{theorem}[Blocked thermal tails with all depth factors]
\label{f15v3:tails}
At every allowed depth and for every coarse plaquette,
\begin{align}
 \mathbb E_\mu\exp\left\{
       \frac{\beta(\vartheta_Q^{(j)})^2}{\pi^2\,256^j}
                         \right\}&\le C_L,
                  \label{eq:f15v3:mgf}\\
 \mathbb E_\mu(\vartheta_Q^{(j)})^{2m}
       &\le m!C_L\left(\frac{\pi^2\,256^j}{\beta}\right)^m
       \quad(m=1,2,\ldots),                 \label{eq:f15v3:moments}\\
 \mathbb E_\mu(\vartheta_Q^{(j)})^2
       &\le\frac{\pi^2 K_L\,256^j}{2\beta},\qquad
 \mathbb E_\mu s_Q^{(j)}
       \le\frac{\pi^2 K_L\,256^j}{4\beta}.
                  \label{eq:f15v3:mean}
\end{align}
For a finite collection \(\mathcal D\) of level-\(j\) faces,
\begin{equation}
 \mu\left\{\max_{Q\in\mathcal D}\vartheta_Q^{(j)}>\delta\right\}
 \le \min\left\{1,\ |\mathcal D|C_L
          \exp\left[-\frac{\beta\delta^2}{\pi^2\,256^j}\right]\right\}.
 \label{eq:f15v3:tail}
\end{equation}
There is no independence assumption on any of these faces.
\end{theorem}

\begin{proof}
For \(0\le\theta\le\pi\),
\(\theta^2\le(\pi^2/2)(1-\cos\theta)\), since
\(\sin(\theta/2)\ge\theta/\pi\). Therefore
\(\mathbb E e^{\beta(\vartheta_p^{(0)})^2/\pi^2}
\le\mathbb E e^{X_p/2}\le C_L\).
Set \(W=16^j\) and \(v_p=w^{(j)}_{Qp}/W\).
These \(v_p\)'s sum to one. Weighted Cauchy--Schwarz
in \eqref{eq:f15v3:iterated} gives
\[
 \frac{\beta(\vartheta_Q^{(j)})^2}{\pi^2 W^2}
       \le \sum_p v_p\,\frac{\beta(\vartheta_p^{(0)})^2}{\pi^2}.
\]
Convexity of the exponential bounds its expectation
by \(\sum_pv_p C_L=C_L\), proving \eqref{eq:f15v3:mgf}.
The inequality \(z^m\le m!e^z\) gives the moment bound.
For the sharper first-moment constant, use instead
\(\mathbb E(\vartheta_p^{(0)})^2\le
(\pi^2/2\beta)\mathbb EX_p\le\pi^2K_L/(2\beta)\)
in the same weighted square estimate.
Finally \(1-\cos\theta\le\theta^2/2\), and exponential
Markov followed by the union bound proves the rest.
\end{proof}

The same argument gives a joint, possibly mixed-depth,
statement. For any finite family \((j_\ell,Q_\ell)\)
and \(b_\ell\ge0\) with \(\sum_\ell b_\ell\le1\),
\begin{equation}
 \mathbb E_\mu\exp\left\{
     \sum_\ell b_\ell
       \frac{\beta(\vartheta_{Q_\ell}^{(j_\ell)})^2}
                         {\pi^2\,256^{j_\ell}}\right\}\le C_L .
 \label{eq:f15v3:joint}
\end{equation}
Apply convexity once more and supply the missing
weight \(1-\sum b_\ell\) at the constant zero.
This is a normalized positive-source estimate,
not a claim of independent blocked curvatures.

\subsection{Actual local cutoff control at increasing depth}

Fix a finite nonempty set \(\mathcal F\) of level-\(j\)
plaquettes, \(j\ge1\). Set \(\mathcal A_j=\mathcal F\)
and recursively take \(\mathcal A_i\) to be the union
of the supports of the one-step rows for
\(Q\in\mathcal A_{i+1}\). These are deterministic sets
of plaquettes, with multiplicities discarded.
The stencil support bound implies
\begin{equation}
 |\mathcal A_i|
 \le 6|\mathcal F|\,[3\cdot2^{j-i}+1]^4
 \le1536|\mathcal F|\,16^{j-i}
 \quad(0\le i<j).
 \label{eq:f15v3:cardinality}
\end{equation}
The deliberately larger enclosing box also covers
all relative-path exchange faces used in these rows.

Define
\[
 \delta_0=\frac{\rho}{12\sqrt2},\qquad
 E_j=\bigcup_{i=0}^{j-1}
       \left\{\max_{p\in\mathcal A_i}\vartheta_p^{(i)}
                                                  >\delta_0\right\}.
\]
On \(E_j^c\), each relative path used in the row
recursion has angle at most \(6\delta_0=\rho/(2\sqrt2)\).
Thus its cutoff is exactly one and its relative
logarithm is in the identity chart. This assertion
concerns these specified block updates and their
relative inputs; it is not a global small-link gauge.

\begin{proposition}[Full-measure cost of local cutoff failure]
\label{f15v3:cutoff}
The event just defined satisfies
\begin{align}
 \mu(E_j)&\le
 \min\left\{1,\ C_L\sum_{i=0}^{j-1}|\mathcal A_i|
             e^{-\beta\delta_0^2/(\pi^2\,256^i)}\right\}
                                                        \notag\\
 &\le\varepsilon_j:=
 \min\left\{1,\ 1536|\mathcal F|C_L\,j16^j
             e^{-\beta\delta_0^2/(\pi^2\,256^{j-1})}\right\}.
 \label{eq:f15v3:cutoff-cost}
\end{align}
For each fixed \(\eta\in(0,1)\), fixed
\(|\mathcal F|\), and allowed sequence with \(\beta\to\infty\),
\[
       1\le j\le\frac{(1-\eta)\log\beta}{\log256}
             \quad\Longrightarrow\quad \varepsilon_j\to0.
\]
\end{proposition}

\begin{proof}
Use \eqref{eq:f15v3:tail} at each intermediate level
and sum the probabilities. The largest cardinality
bound is \(1536|\mathcal F|16^j\) and the weakest
exponent occurs at \(i=j-1\). In the stated window,
\(16^j\le\beta^{(1-\eta)/2}\) while
\(\beta/256^{j-1}\ge256\beta^\eta\).
The latter exponential decay dominates the former
power and the at most logarithmic factor \(j\).
\end{proof}

In this window the largest blocking length is
\(2^j\le\beta^{(1-\eta)/8}\). This is a proved,
conservative local range. Nothing here identifies
it with a physical correlation length, asserts
optimality of the factor 256, or permits replacing
the displayed range by an arbitrary physical scale.

\subsection{Conditioning on the whole output and truncating actual scores}

Let \(\rho_j=(C^{(j)})_*\mu\).
The conditional probabilities in this subsection use
the actual regular conditional Wilson law.

\begin{theorem}[Posterior cutoff bound and score truncation]
\label{f15v3:posterior}
For any event \(E\) with \(\mu(E)\le\varepsilon\),
write \(b_j(C)=\mu(E\mid C^{(j)}=C)\).
For \(0<r<1\),
\begin{equation}
            \rho_j\{b_j>r\}\le\varepsilon/r.
 \label{eq:f15v3:posterior}
\end{equation}
For \(X\in L^4(\mu)\), define
\[
 A_j(C)=\mathbb E_\mu[X\mid C^{(j)}=C],
 \qquad A_j^{\rm good}(C)
       =\mathbb E_\mu[X{\bf1}_{E^c}\mid C^{(j)}=C].
\]
Then
\begin{equation}
       \|A_j-A_j^{\rm good}\|_{L^2(\rho_j)}
                        \le\|X\|_{L^4(\mu)}\,\varepsilon^{1/4}.
 \label{eq:f15v3:score-cutoff}
\end{equation}
In particular these statements hold for \(E=E_j\)
with the explicit \(\varepsilon_j\) above and for
V2's fine-action score \(X=X(a)=\sum_xa_x\widetilde T_x\).
For that score one available uniform bound is
\begin{equation}
             \|X(a)\|_{L^4(\mu)}
                 \le12(384C_L)^{1/4}\|a\|_{\ell^1}.
 \label{eq:f15v3:score-fourth}
\end{equation}
\end{theorem}

\begin{proof}
The conditional-probability tower gives
\(\int b_j\,d\rho_j=\mu(E)\). Markov proves
\eqref{eq:f15v3:posterior} without restricting how
much of the coarse configuration is observed.
Conditional Jensen and Cauchy--Schwarz give
\[
 \|A_j-A_j^{\rm good}\|_2^2
 \le\mathbb E_\mu[X^2{\bf1}_E]
 \le(\mathbb E_\mu X^4)^{1/2}\mu(E)^{1/2},
\]
which proves the truncation bound.
The cell variable is \(T_x=6\) times the arithmetic
mean of its 24 face variables \(X_p\). By convexity,
\(\mathbb ET_x^4\le6^4\sup_p\mathbb EX_p^4
\le6^4\,2^4\,4!\,C_L\).
Also \(\|T_x-\mathbb ET_x\|_4\le2\|T_x\|_4\).
Minkowski for the finite sum proves
\eqref{eq:f15v3:score-fourth}.
\end{proof}

Here \(A_j^{\rm good}\) is an \emph{unnormalized}
conditional truncation; it is not
\(\mathbb E[X\mid C^{(j)},E^c]\).
The original \(\ell^1\) source norm in
\eqref{eq:f15v3:score-fourth} is retained. There is no
uniform claim for source supports growing without paying
that cost. The posterior conclusion holds outside a
set of output probability at most \(\varepsilon/r\);
it is not an all-exterior statement and therefore
does not contradict the frozen-exterior obstruction
of f13 V8.

For completeness, the bad-event estimate also survives
the actual nonuniform source changes from V2.
When \(\mu\) is additionally one of V2's selected
source-window laws, for \(p>1\) set
\[
       B_p(a)=\frac{k}{p}\sum_xJ(pa_x),
           \qquad \|pa\|_\infty\le R .
\]
The V2 likelihood moment and H\"older imply
\begin{equation}
        \mu_a(E)\le
        \min\{1,\ e^{B_p(a)}\varepsilon^{\,1-1/p}\}.
 \label{eq:f15v3:tilted}
\end{equation}
Applying the conditional-probability tower under
\(\mu_a\) gives \eqref{eq:f15v3:posterior} for
its actual output law, with this new event budget.
Large or extensive fine sources must pay \(B_p(a)\);
it has not been replaced by a coarse-volume constant.

\subsection{A necessary scale separation from the Haar terminal reference}

The tail estimate also prevents a mistaken handoff.
Let \(H_j\) be product Haar measure on the level-\(j\)
links and put \(t=\beta/256^j\). For any one genuine
coarse plaquette and \(0<\delta\le1\), define
\[
 \epsilon=\min\{1,C_Le^{-t\delta^2/\pi^2}\},
 \qquad h(\delta)=\frac2\pi\int_0^\delta\sin^2u\,du
                       \le\frac{2\delta^3}{3\pi}.
\]
Under \(H_j\) this plaquette holonomy is Haar, since
its four links are distinct independent Haar groups.
Under \(\rho_j\), its angle is at most \(\delta\)
with probability at least \(1-\epsilon\).
Consequently
\begin{align}
 \|\rho_j-H_j\|_{\rm TV}
       &\ge1-\epsilon-h(\delta),                 \label{eq:f15v3:TV}\\
 D(\rho_j\|H_j)
       &\ge(1-\epsilon)\log\frac1{h(\delta)}-\log2 .
                                                        \label{eq:f15v3:KL}
\end{align}
The first inequality tests this single event.
For the second, split relative entropy on the event
and its complement and discard the nonnegative
within-event conditional entropies. The resulting
binary entropy is at least
\(\alpha\log(1/h)-\log2\), where
\(\alpha\ge1-\epsilon\). The inequality also holds
if the full relative entropy is infinite.

If \(t\to\infty\), take \(\delta=t^{-1/4}\) once
\(t\ge1\). Then \(\epsilon\le C_Le^{-\sqrt t/\pi^2}\),
the total variation tends to one, and
\[
 D(\rho_j\|H_j)\ge
 (1-\epsilon)\left(\frac34\log t+\log\frac{3\pi}{2}\right)-\log2.
\]
Thus even the one-plaquette marginal is not approaching
the Haar terminal reference in this regime.
Moreover, for \(\chi_Q=\tfrac12\Tr C_Q^{(j)}\),
\[
       \Var_{\rho_j}(\chi_Q)
       \le\mathbb E_{\rho_j}(1-\chi_Q)^2
       \le2\mathbb E_{\rho_j}(1-\chi_Q)
       \le\frac{\pi^2K_L\,256^j}{2\beta}\longrightarrow0.
\]
This excludes membership, in this depth window, in
any proposed terminal family requiring a fixed positive
lower variance for this same observable. In particular
it cannot be identified with the actual small-coupling
completed-output family of f14 V48, whose corresponding
variance is greater than \(1/5\). No assertion is made
that the present conservative boundary is the true
entrance scale, or that entrance occurs beyond it.

\subsection{What this pays for, and the next nonlinear comparison}

The compact map now has full-measure, depth-explicit
curvature bounds and a quantitative local region where
its relative cutoff equals one with high probability.
The same region can be used inside the actual conditional
expectation, with its omitted score contribution bounded
by \eqref{eq:f15v3:score-cutoff}. This is a concrete
input for a local expansion of the generated density,
not an assumption of global chart coverage.

Small relative inputs do not by themselves give
convexity, spatial decay of the hidden conditional
covariance, or locality of \(A_j^{\rm good}\). That
truncated score can still depend on the entire output.
The next useful comparison is the signed first
variation of the actual completed coarse curvature,
combined with the full compact local Gaussian limits
of f14 V97--V100. Any such passage must control its
nonlinear remainder on the local event and pay its
complement as above. It must not infer convergence of
conditional laws from joint weak convergence.

\subsection{Finite checks and preservation}

The certificate is
\begin{center}\small
\path{scripts/verify_ym_f15_v3_compact_block_curvature.py}.
\end{center}
It checks all 64 ordered relative paths by 192
noncommuting free-group plaquette exchanges, as well
as the 24 exchanges for the six root squares.
It constructs all six positive one-step stencils
and their depth-two and depth-three compositions,
checking exact row weights and support bounds.
There are 54 rational weighted-square checks,
96 posterior Markov checks, 32 conditional-score
truncation checks and 11 dyadic cardinality checks.
The finite probability checks use rational models,
not Wilson simulations.

The certificate has 8,779 bytes and SHA--256
\begin{center}\scriptsize\ttfamily
7E99D39C9147E8E9D58AD475B6D20A5F64EBECD66DD4001FA8E88B820955961C.
\end{center}
It does not machine-prove the analytic distance or
probability theorems, which are proved above, and
does not certify continuum existence or a mass gap.

The predecessor V2 has 38,662 bytes and SHA--256
\begin{center}\scriptsize\ttfamily
AFB4BFB84096768672719EE6159BD703FE53BB5EC40BA1EB8CABA46C763023A6.
\end{center}
Reversing the title change and this terminal insertion
recovers that predecessor exactly. V3 replaces V2
as the sole active main note; the archives, 11
protected sources and 45 earlier protected certificates
are unchanged. The next companion audit is V5.

\begin{firewall}
V3 controls actual compact blocked curvature and
local cutoff failures under the full fine Wilson law.
Its \(256^j/\beta\) scale is a proved upper-bound cost,
not a running coupling or physical mass.
Neither conditional-score locality, a closed nonlinear
RG trajectory, the interacting four-dimensional
continuum theory, nor its positive physical Hamiltonian
mass gap has been proved. The original goal remains active.
\end{firewall}
\section{V4: the actual blocked curvature filter and retained-score noncollapse}
\label{f15v4:context}

\subsection{Which missing implication is being proved}

V3 controls the actual compact output and its local cutoff
failures. The next issue is not another Gaussian projection
formula, but whether the nonlinear Wilson pushforward has
the proposed linear curvature as its leading observable.
This section proves that implication, with an explicit
growing-block remainder, before making any Gaussian
limit passage. It then gives a lower bound for an actual
conditional score using an observable of the actual output.

The one-step linear curl identity is already proved in
f14 V41, equation \texttt{v41-curl-average}. The balanced
map's scaled-amplitude bound and its exact orbit relation
to the literal map are proved in Shell-Mixing V10.
We reuse these results with their hypotheses, rather than
re-proving either as new work. The full compact thermal
moments and selected Gaussian limit are those of
f14 V97--V100. The new work is their measured nonlinear
comparison, its block-size dependence, and the resulting
actual retained-score lower bound.
The next companion checkpoint remains V5.

\subsection{The inherited curl filter, with its normalization retained}

Write \(p=2^j\). For an oriented coarse face
\(Q=(x;\mu\nu)\), \(\mu<\nu\), define a scalar row on
fine oriented faces by
\begin{equation}
 (H_p f)_Q
   =p^{-4}\sum_{s\in\{0,\ldots,p-1\}^4}
                  \sum_{u,v=0}^{p-1}
       f_{px+s+ue_\mu+ve_\nu,\,\mu\nu}.
 \label{eq:f15v4:filter-map}
\end{equation}
It acts separately on each colour. For a row at \(x=0\)
write its coefficients as \(h_{p,\mu\nu}(r)\), with
zero coefficients on other orientations. Explicitly,
\begin{equation}
 h_{p,\mu\nu}(r)=p^{-4}n_p(r_\mu)n_p(r_\nu)
                 \prod_{\alpha\notin\{\mu,\nu\}}
                          {\bf1}_{0\le r_\alpha<p},
 \quad
 n_p(t)=
 \begin{cases}
 \min(t+1,p,2p-1-t),&0\le t\le2p-2,\\
 0,&\text{otherwise}.
 \end{cases}
 \label{eq:f15v4:filter-coefficients}
\end{equation}
Thus
\begin{equation}
 \sum_rh_{p,\mu\nu}(r)=p^2,\qquad
 H_p^{(2)}:=\|h_{p,\mu\nu}\|_{\ell^2}^2
             =\frac{(2p^2+1)^2}{9p^4}\le1.
 \label{eq:f15v4:hnorm}
\end{equation}
The notation \(H_p^{(2)}\) denotes a scalar squared
row norm, not a second blocking map.

For clarity, these formulas do not replace the
absolute nonlinear row weight \(16^j=p^4\) in V3.
At first order the coarse tree-gradient terms cancel
from the curl, leaving the smaller signed boundary
calculation in \eqref{eq:f15v4:filter-map}.
Indeed f14 V41 gives this formula for \(p=2\).
Composition of the binary box averages and the
binary surface sums gives each of the \(p\) choices
of \(s_\alpha,u,v\) exactly once, proving the formula
for dyadic \(p\). Alternatively, the derivative of
the balanced link map is
\[
 (B_pa)_{x,\mu}
      =p^{-4}\sum_{s\in\{0,\ldots,p-1\}^4}
                         \sum_{t=0}^{p-1}a_{px+s+te_\mu,\mu},
 \qquad d_cB_pa=H_p d_fa .
\]
The second identity is the telescoping boundary of
each \(p\)-by-\(p\) surface. Finally,
\(\sum_t n_p(t)^2=p(2p^2+1)/3\) gives
\eqref{eq:f15v4:hnorm}. These are normalization and
composition checks on the inherited linear identity.

\subsection{A local rooted gauge with a growing-patch probability bound}

We first consider the coarse face \(Q=(0;\mu\nu)\).
Assume \(M\ge16p\), with \(M\) dyadic, so the following
patch embeds without winding. All fine links needed to
compute its literal nested output lie in the vertex box
\[
                         \mathcal K_p=\{0,\ldots,4p\}^4 .
\]
To check this conservative carrier, one primitive
edge stencil has fine edge bases at offsets at most
3 in each coordinate. Starting with the four coarse
face-edge bases, which have coordinates at most 1,
iteration gives fine bases at most
\(p+3(p-1)=4p-3\), and endpoints add at most one.

Choose the increasing-coordinate path \(\gamma_y\)
from 0 to every vertex \(y\in\mathcal K_p\). Set
\[
 g_y=U_{\gamma_y},\qquad
 k_{y,\alpha}=g_yU_{y,\alpha}g_{y+e_\alpha}^{-1},
 \qquad R_p(U)=\max_{e\subset\mathcal K_p}\vartheta(k_e).
\]
Each such rooted edge word has a plaquette filling
of area at most \(12p\): move its final
\(\alpha\)-step left past the later coordinate
steps in \(\gamma_y\), using at most
\(\sum_{\nu>\alpha}y_\nu\le12p\) exchanges.
There are at most \(2500p^4\) edges in this box.
Tree links can have filling area zero and then equal
the identity exactly.

As in V3, all bounds in this subsection use the
full Wilson law, with \(\beta\ge2\),
\(\log\beta/M\le L\), and \(C_L\) fixed. Put
\begin{equation}
 B_p^*=2500C_Lp^4,\qquad
 \ell_p=1+\log B_p^*,\qquad
 a_p=\frac{144\pi^2p^2}{\beta}.
 \label{eq:f15v4:patch-constants}
\end{equation}
The superscript in \(B_p^*\) distinguishes this
scalar from the linear link map \(B_p\).

\begin{proposition}[Rooted-link maximum under the actual law]
\label{f15v4:rooted-maximum}
For every \(t>0\),
\begin{equation}
                \mu\{R_p>t\}\le
                       \min\{1,B_p^*e^{-t^2/a_p}\}.
 \label{eq:f15v4:root-tail}
\end{equation}
For every fixed \(q\ge1\) there is \(C_q<\infty\),
independent of \(p,M,\beta\), such that
\begin{equation}
                \|R_p\|_{L^q(\mu)}
                   \le C_q\,\frac{p\sqrt{\ell_p}}{\sqrt\beta}.
 \label{eq:f15v4:root-moment}
\end{equation}
\end{proposition}

\begin{proof}
For an edge with filling area \(A\ge1\), the
nonabelian square-exchange argument gives
\(\vartheta(k_e)\le\sum_{m=1}^{A}\vartheta_{p_m}^{(0)}\).
Weighted Cauchy--Schwarz and convexity of the
exponential, exactly as in V3, imply
\[
       \mathbb E_\mu
             e^{\beta\vartheta(k_e)^2/(\pi^2 A^2)}\le C_L.
\]
Use \(A\le12p\) and take the union over the box
edges to obtain \eqref{eq:f15v4:root-tail}.

For the moment estimate set
\(T=(R_p^2-a_p\log B_p^*)_+\). The tail inequality
implies \(\mu\{T>u\}\le e^{-u/a_p}\) and hence
\(\mathbb ET^m\le m!a_p^m\) for integers \(m\ge1\).
Since \(R_p^2\le a_p\log B_p^*+T\),
\[
 \mathbb E R_p^{2m}
      \le2^{m-1}(1+m!)a_p^m\ell_p^m.
\]
Taking roots and increasing the even moment if
necessary proves the assertion for any fixed \(q\).
No conditional independence or gauge-fixed density
has been inserted in this estimate.
\end{proof}

\subsection{Uniform small-amplitude comparison for the actual map}

Let \(\widehat P_y\) be the fine plaquette product
formed from \(k\), so
\(\widehat P_y=g_yU_yg_y^{-1}\), and put
\[
     Y_y^\gamma=\sqrt\beta\,\operatorname{Im}\widehat P_y,
 \qquad Z_{p,Q}=(H_pY^\gamma)_Q,\qquad
     X_{p,Q}^{\rm out}=\beta(1-\tfrac12\Tr C_Q^{(j)}).
\]
For the coarse face based at 0, \(g_0=I\), so its
rooted output curvature is simply
\[
                    Y_{p,Q}^{\rm out}
                           =\sqrt\beta\,\operatorname{Im}C_Q^{(j)}.
\]
For faces at other bases, conjugate by the same
chosen fine transport \(g_{px}\). A finite family
is handled in one common enlarged box. This only
changes the geometric constants in the patch bounds.

\begin{proposition}[A depth-uniform local deterministic remainder]
\label{f15v4:deterministic}
There are \(r_*>0\) and \(C<\infty\), independent
of the dyadic depth and the surrounding volume,
such that, on \(pR_p\le r_*\),
\begin{align}
 \left|\operatorname{Im}C_Q^{(j)}
                   -(H_p\operatorname{Im}\widehat P)_Q\right|
       &\le C(pR_p)^2,                       \label{eq:f15v4:field-local}\\
 \left|1-\tfrac12\Tr C_Q^{(j)}
                   -\tfrac12|(H_p\operatorname{Im}\widehat P)_Q|^2\right|
       &\le C(pR_p)^3.                       \label{eq:f15v4:energy-local}
\end{align}
All relative cutoffs involved in this comparison
are one. The constants are uniform but are not
claimed to have certified optimal numerical values.
\end{proposition}

\begin{proof}
On this event the relevant links have principal
logarithms \(a_e=\log k_e\), with
\(\|a\|_\infty\le R_p\). Extend this finite input
by zero Lie algebra outside its carrier. This is
only a way of evaluating a local deterministic
function; the Wilson measure is not altered.
The extended computation agrees with the original
observable because the carrier contains all of its
input links.

Use Shell-Mixing V10's real scaled-amplitude theorem
with \(\varrho=p\|a\|_\infty\). Choose \(r_*\) smaller
than its common radius, and smaller still so that
its intermediate real relative logarithms have
distance less than \(\rho/2\) in the completion metric.
The bounds on intermediate amplitudes make this choice
independent of the depth. Thus the literal completed
map agrees with the untruncated root map here.
For the balanced representative its final logarithm
\(\widehat a\) satisfies
\[
     \|\widehat a\|_\infty\le2\varrho,\qquad
     \|\widehat a-B_pa\|_\infty\le C\varrho^2 .
\]
Its plaquette imaginary part differs from
\(d_cB_pa=H_pd_fa\) by at most \(C\varrho^2\),
using the four-factor product expansion.

The accumulated balancing gauge does not invalidate
this colour comparison. At level \(s\), a tree has
at most four links, so its balancing factor has
angle at most \(4\|a^{(s)}\|_\infty\).
The inherited amplitude estimate gives
\[
 \sum_{s=0}^{j-1}4\|a^{(s)}\|_\infty
       \le8\varrho\sum_{s=0}^{j-1}2^{s-j}
       \le8\varrho .
\]
The exact gauge-composition law therefore puts the
accumulated coarse factor within \(8\varrho\) of \(I\).
Conjugating the balanced plaquette, whose imaginary
part is \(O(\varrho)\), back to the literal output
changes it by \(O(\varrho^2)\).
This retains the actual colour frame, not just its norm.

For each fine face,
\(|\operatorname{Im}\widehat P_y-(d_fa)_y|
\le C\|a\|_\infty^2\). Since the row weight of
\(H_p\) is \(p^2\), replacing \(H_pd_fa\) by
\(H_p\operatorname{Im}\widehat P\) costs at most
\(Cp^2\|a\|_\infty^2=C\varrho^2\).
This proves \eqref{eq:f15v4:field-local}.
Also \(\|B_pa\|_\infty\le p\|a\|_\infty\), so
both imaginary vectors just compared have size
\(O(\varrho)\), and the output plaquette has
angle \(O(\varrho)\).

For a unit quaternion \(P\), writing
\(s=1-\operatorname{Re}P\), one has exactly
\[
                      s-\tfrac12|\operatorname{Im}P|^2=s^2/2.
\]
Here \(s=O(\varrho^2)\). The quadratic energy
difference between the two imaginary vectors is
\(O(\varrho)\,O(\varrho^2)\), and the displayed
identity contributes \(O(\varrho^4)\).
Decreasing \(r_*\) if needed proves
\eqref{eq:f15v4:energy-local}.
\end{proof}

\subsection{A measured remainder before any Gaussian limit}

Set
\[
 q_{p,\beta}=\min\left\{1,\ B_p^*
                  \exp\left[-\frac{r_*^2\beta}{144\pi^2p^4}\right]\right\}.
\]
This bounds the probability of \(pR_p>r_*\) under
the original compact law.

\begin{theorem}[Full-measure nonlinear-to-linear comparison]
\label{f15v4:comparison}
For each fixed \(m\ge1\), there is \(C_m<\infty\)
independent of \(p,M,\beta\) such that
\begin{align}
 \|Y_{p,Q}^{\rm out}-Z_{p,Q}\|_{L^m(\mu)}
       &\le C_m\frac{p^4\ell_p}{\sqrt\beta}
               +\sqrt\beta(1+p^2)q_{p,\beta}^{1/m},
                                            \label{eq:f15v4:field-Lm}\\
 \|X_{p,Q}^{\rm out}-\tfrac12|Z_{p,Q}|^2\|_{L^m(\mu)}
       &\le C_m\frac{p^6\ell_p^{3/2}}{\sqrt\beta}
               +2\beta(1+p^4)q_{p,\beta}^{1/m}.
                                            \label{eq:f15v4:energy-Lm}
\end{align}
In particular, the curvature error tends to zero
whenever \(p^4\ell_p=o(\sqrt\beta)\), and the energy
error tends to zero whenever
\(p^6\ell_p^{3/2}=o(\sqrt\beta)\), subject to the
embedding and Wilson-law hypotheses above.
These compare with the actual fine field \(Y^\gamma\),
not with a Gaussian replacement at growing \(p\).
\end{theorem}

\begin{proof}
On \(pR_p\le r_*\), multiply
\eqref{eq:f15v4:field-local} by \(\sqrt\beta\)
and use \eqref{eq:f15v4:root-moment} at order \(2m\).
This gives \(C_mp^4\ell_p/\sqrt\beta\).
Similarly \eqref{eq:f15v4:energy-local}, multiplied
by \(\beta\), uses the moment of \(R_p\) at order
\(3m\) and gives \(C_mp^6\ell_p^{3/2}/\sqrt\beta\).

On the complementary event use the actual compact
bounds \(|\operatorname{Im}C_Q^{(j)}|\le1\),
\(|(H_p\operatorname{Im}\widehat P)_Q|\le p^2\),
and \(0\le1-\tfrac12\Tr C_Q^{(j)}\le2\).
The absolute errors are at most
\(\sqrt\beta(1+p^2)\) and
\(\beta(2+p^4/2)\le2\beta(1+p^4)\), respectively.
Multiplying by the appropriate power of the bad
probability proves both estimates.

For the limiting assertions, the first condition
implies eventually \(p\le\beta^{1/8}\) and hence
\(\beta/p^4\ge\sqrt\beta\). The second similarly
implies \(p\le\beta^{1/12}\) and
\(\beta/p^4\ge\beta^{2/3}\). With \(C_L\) fixed,
the complementary exponential terms dominate all
their displayed polynomial prefactors.
\end{proof}

For example the first error vanishes for
\(p\le\beta^{(1-\eta)/8}\), and the second for
\(p\le\beta^{(1-\eta)/12}\), with fixed \(\eta>0\).
The logarithm in \(\ell_p\) and the different
field and energy exponents have not been dropped.
This supplies a measured growing-support comparison;
it does not make the previous fixed-coordinate
Gaussian convergence uniform on that support.

\subsection{The actual fixed-depth output law}

Now, and only now, use the selected full compact
couplings \(\beta_M^\diamond\) of f14 V100.
Fix the dyadic \(p\) and the finitely many fine
and coarse observables before sending \(M\) to infinity.
Let \(G\) be the common-root Gaussian curvature with
covariance \(P\otimes I_3\) from that result.

\begin{theorem}[Measured coarse curvature and energy limit]
\label{f15v4:limit}
Jointly with every fixed finite set of fine
curvatures and energies, and with all fixed
mixed polynomial moments,
\begin{equation}
          (Y_{p,Q}^{\rm out},X_{p,Q}^{\rm out})
             \ \Longrightarrow\
                  \bigl((H_pG)_Q,\tfrac12|(H_pG)_Q|^2\bigr).
 \label{eq:f15v4:limit}
\end{equation}
The assertion extends to a fixed finite family
of coarse faces transported to one common root.
Define
\[
            c_p=\langle h_{p,\mu\nu},P h_{p,\mu\nu}\rangle .
\]
It does not depend on the base or orientation, and
\begin{equation}
        \frac8{\pi^{18}}\le c_p\le H_p^{(2)}\le1.
 \label{eq:f15v4:cp-bounds}
\end{equation}
For one face its limiting energy is
\((c_p/2)\chi_3^2\). In particular,
\begin{equation}
  \mathbb E X_{p,Q}^{\rm out}\longrightarrow\tfrac32c_p,
       \qquad
  \Var(X_{p,Q}^{\rm out})\longrightarrow\tfrac32c_p^2>0 .
 \label{eq:f15v4:coarse-moments}
\end{equation}
\end{theorem}

\begin{proof}
At fixed \(p\), \(H_p\) uses finitely many fine
coordinates. The inherited full compact local
limit and all-moment result therefore apply to
\(Z_{p,Q}=H_pY^\gamma\), jointly with the fixed
fine observables. The two errors in
Theorem~\ref{f15v4:comparison} tend to zero in
every fixed \(L^m\). Expanding differences of
products and using H\"older transfers all fixed
mixed polynomial moments to the actual output.
For several faces use a common finite rooted box;
the same argument applies to each component.

For one colour and orientation, the inherited
projection has diagonal Fourier multiplier
\[
       P_{\mu\nu,\mu\nu}(k)
         =\frac{|e^{ik_\mu}-1|^2+|e^{ik_\nu}-1|^2}
                     {\sum_{\alpha=1}^4|e^{ik_\alpha}-1|^2}
\]
away from the measure-zero origin.
With \(D_p(t)=\sum_{s=0}^{p-1}e^{ist}\), the
filter multiplier is
\[
 \widehat h_{p,\mu\nu}(k)
       =p^{-4}D_p(k_\mu)D_p(k_\nu)
                              \prod_{\alpha=1}^4D_p(k_\alpha).
\]
Thus \(c_p\) is the integral of
\(|\widehat h|^2P_{\mu\nu,\mu\nu}\) against
\(dk/(2\pi)^4\). Its upper bound follows from
\(0\le P\le I\).
For an explicit lower bound restrict to
\(k_\alpha\in[1/(2p),1/p]\) for all four coordinates.
Here \(|D_p(k_\alpha)|/p\ge2/\pi\) and
\(P_{\mu\nu,\mu\nu}(k)\ge1/(2\pi^2)\).
The box volume is \(1/(16p^4)\), so its contribution
is at least
\[
       \frac{p^4(2/\pi)^{12}}{(2\pi)^4}
                         \frac1{16p^4}\frac1{2\pi^2}
                             =\frac8{\pi^{18}} .
\]
Translation and coordinate symmetry give the
stated independence of base and orientation.
The three colours are independent with variance
\(c_p\); their squared norm gives the chi-square
law and the displayed moments.
\end{proof}

The uniform bounds on the Gaussian numbers \(c_p\)
do not interchange the two limits. The conclusion
\eqref{eq:f15v4:limit} takes \(\beta_M^\diamond,M\)
to infinity with \(p\) fixed. The quantitative
nonlinear comparison alone does not justify putting
a growing \(p=p_M\) into that Gaussian limit.

\subsection{A lower bound for the actual retained fine-action score}

Let \(w_x(f)\) be the inherited cell-face weights:
\(T_x=\sum_f w_x(f)X_f\), with weight \(1/4\)
on its 24 faces. For a real finite cell source \(b\),
write
\[
 X(b)=\sum_xb_x(T_x-\mathbb E_\mu T_x),\qquad
 \omega_b(f)=\sum_xb_xw_x(f),\qquad
 A_{j,b}=\mathbb E_\mu[X(b)\mid C^{(j)}].
\]
These are V2's actual baseline scores; no
source tilt of unit strength is being asserted.
Fix \(p=2^j\), a face \(Q\), and use its translated
row \(h_Q\). Set
\[
             v_Q=P h_Q,\qquad
             \gamma_{b,Q}=\sum_f\omega_b(f)\,v_Q(f)^2.
\]
The last sum is finite because \(b\) is finitely
supported, although \(v_Q\) need not be local.

\begin{theorem}[Actual retained-score information from one coarse face]
\label{f15v4:score-lower}
Along the selected laws, at fixed \(p,b,Q\),
\begin{equation}
       \liminf_{M\to\infty}\Var_{\rho_j}(A_{j,b})
                      \ge\frac32\,\frac{\gamma_{b,Q}^2}{c_p^2}.
 \label{eq:f15v4:score-lower}
\end{equation}
In particular let \(\mathcal D_p\) be the set
of bases of the fine faces in the support of
\(h_Q\), and take \(b_x={\bf1}_{x\in\mathcal D_p}\).
Then \(|\mathcal D_p|=p^2(2p-1)^2\) and
\begin{equation}
       \liminf_{M\to\infty}\Var_{\rho_j}(A_{j,b})
                  \ge\frac{3c_p^2}{32(H_p^{(2)})^2}>0 .
 \label{eq:f15v4:score-positive}
\end{equation}
\end{theorem}

\begin{proof}
Use the actual measurable output observable
\(O=X_{p,Q}^{\rm out}\). Conditional expectation
preserves its covariance with \(X(b)\), so
Cauchy--Schwarz gives
\[
       \Var_{\rho_j}(A_{j,b})
           \ge\frac{\operatorname{Cov}_\mu(X(b),O)^2}
                         {\Var_\mu(O)}
\]
whenever the denominator is positive. It is
positive for all sufficiently large \(M\) by
\eqref{eq:f15v4:coarse-moments}.
Theorem~\ref{f15v4:limit} gives convergence of
both finite-observable moments. For three
independent Gaussian colours, Wick's identity gives
\[
 \operatorname{Cov}\left(\tfrac12|G_f|^2,
                         \tfrac12|(H_pG)_Q|^2\right)
                 =\tfrac32\,\langle e_f,P h_Q\rangle^2.
\]
Consequently the numerator's unsquared covariance
tends to \((3/2)\gamma_{b,Q}\), and the variance
denominator tends to \((3/2)c_p^2\).
This proves \eqref{eq:f15v4:score-lower}.

For the indicated source, every \(f\) in the
support of \(h_Q\) has \(\omega_b(f)\ge1/4\),
since its own base cell includes that face.
All \(\omega_b\)'s are nonnegative. Therefore
\[
 \gamma_{b,Q}\ge\tfrac14\sum_{f\in\operatorname{supp}h_Q}v_Q(f)^2
          \ge\frac{c_p^2}{4H_p^{(2)}} ,
\]
where the second inequality is Cauchy--Schwarz
applied to \(c_p=\langle h_Q,v_Q\rangle\).
Substitution gives \eqref{eq:f15v4:score-positive}.
\end{proof}

This lower bound is for the score conditioned on
the entire actual compact output, not on a
Gaussian surrogate. Only the testing observable's
joint moments were passed to the limit.
No convergence of the conditional laws, conditional
expectations, or pseudoinverses was inferred.
Nor is \eqref{eq:f15v4:score-positive} a positive
fraction of the extensive fine-source variance:
its source has \(p^2(2p-1)^2\) cells, and that
normalization cost remains visible.

\subsection{The remaining nonlinear RG step}

The programme now has both an upper source budget
from V2 and an actual retained-score lower bound
from the measured coarse curvature. A growing-block
nonlinear error is controlled directly under the
full compact law. These statements identify a
nonzero portion of what the output measures, rather
than assuming that it is a fresh Wilson action.

They do not identify all conditional scores with
local coarse energies. The projection \(P h_Q\)
in the limiting covariance need not have finite
support, and the variance lower bound supplies
neither a locality upper bound nor contraction of
the remaining generated interactions. The next
unpaid issue is a quantitative comparison for the
fine filtered curvature on growing supports,
together with conditional-source locality. An
exchange of fixed-depth and continuum limits
would not solve it. The physical coupling/scale
relation, interacting continuum construction and
positive physical Hamiltonian mass gap remain open.

\subsection{Checks and preservation}

The certificate is
\begin{center}\small
\path{scripts/verify_ym_f15_v4_actual_curvature_filter.py}.
\end{center}
It verifies the six literal one-step curl rows,
18 dyadic filter compositions and squared norms,
and 25,544 ordered-tree filling counts. It also
checks 64 three-colour covariance and variational
identities using exact Gaussian polynomial moments
and a noncoordinate rational covariance projection.
These are finite algebraic checks, not Wilson
simulations or a computer proof of the analytic
remainder, limit, or mass-gap statements.

The certificate has 7,177 bytes and SHA--256
\begin{center}\scriptsize\ttfamily
C9D54AA5B5CBD9C9947B25A9628D9A3933653691EE24CD7E3515FE68C10FD081.
\end{center}
The predecessor V3 has 60,453 bytes and SHA--256
\begin{center}\scriptsize\ttfamily
A7E56735FB49057F4897A82CF90D1FE0AFCA77CADBFC18B9C5ED618A7A6F7A36.
\end{center}
Reversing the title change and this terminal
insertion recovers it exactly. V4 is the sole
active main note; the 11 protected sources and
46 earlier protected certificates remain unchanged.
The next compulsory companion audit is V5.

\begin{firewall}
V4 proves a measured nonlinear curvature comparison
for the actual compact map, its fixed-depth thermal
limit, and a positive lower bound for specified
actual retained fine-action scores. These are not
a proof of conditional-score locality, a closed
interacting renormalization trajectory, or the
full four-dimensional physical Yang--Mills mass gap.
The original goal remains active.
\end{firewall}
\section{V5: a stationary comparison and actual growing-block energy control}
\label{f15v5:context}

\subsection{Companion checkpoint and the remaining issue}

The required five-version checkpoint inspected the current
worktree and reread the terminal results of SM continuum V72,
NS stability V26, Shell-Mixing V16 and parallel YM V5.
The supplied SM source-selection V17, native regenerative
stationarity V21, exact-action Einstein bridge V16 and
suggestions were also revisited. All 58 protected files
retain their recorded hashes.

SM V72's covariance transport is for a many-copy model
with fixed modes and times. NS V26 supplies particular
rank-one Oseen estimates, not a four-dimensional gauge
estimate. Shell V16's measured flat-branch calculation
does not provide a nonperturbative joint remainder; its
noncommuting coefficient direction was already developed
in f14 V20--V30. Its earlier V10 scaled-amplitude theorem
was used, with its assumptions, in V4 above.
Parallel YM V5's obstruction concerns extensive global
charge tails, not a local curvature estimate.

The source-selection comparison still does not identify
the required primitive port. The regenerative notes
exclude the proposed global real phase and retain a
quadratic phase-space alternative, not a nonlinear
physical YM approximation theorem. The Einstein bridge
keeps a range/domain condition rather than supplying a
full inverse. The suggestions' measured-convention
warning remains applicable; its proposed quartic task
is not repeated. These are mathematical source documents,
not instructions changing the present task.

V4 provides an actual nonlinear comparison at growing
block size, but its Gaussian identification was still
at fixed block size. Here we obtain a quantitative
second-moment estimate for the fine filtered field and
thereby remove the fixed-depth restriction for the
\emph{mean thermal coarse energy}. The new intermediate
field is explicitly auxiliary. Its artificial
cross-cell decorrelation is not a property of the
physical output and will not be called clustering.
The next companion checkpoint is V10.

\subsection{Definitions and a deliberately auxiliary stationary field}

Let \(M\) be dyadic, \(4\le R\le M\) dyadic with
\(R\mid M\). The letter \(R\) in this section denotes
the comparison cell side, not V2's source-window
radius. Sample \(U\) from the full compact Wilson
law, with \(\beta\ge2\) and \(\log\beta/M\le L\).
Independently choose a uniform offset
\(\omega\in\{0,\ldots,R-1\}^4\).
Partition the torus into the translated \(R\)-cells
and choose independent Haar matrices \(H_D\in\SU(2)\),
one for each cell \(D\).

For \(y\in D\), let \(g_y^D\) be the link product
along the increasing-coordinate path from the
cell root to \(y\), using a consistent lifted cell.
For a positive plaquette \(f=(y;\mu\nu)\) define
\begin{equation}
 W_f^R=\sqrt\beta\,\operatorname{Ad}_{H_Dg_y^D}
                                     \operatorname{Im}U_f .
 \label{eq:f15v5:aux-field}
\end{equation}
The cell is assigned by the plaquette base; the
plaquette itself is allowed to cross a cell boundary.
Expectation \(\mathbb E\) now includes the independent
offset and colour matrices as well as \(U\).
The law of \(U\), and the actual completed output,
have not been changed.

\begin{proposition}[Stationarity, covariance positivity and exact energy]
\label{f15v5:stationary}
The field \(W^R\) is stationary under every torus
translation, has mean zero, and has a positive
semidefinite covariance operator \(\Sigma_R\).
Writing
\[
 q_{M,\beta}=\beta\,\mathbb E_\mu S/M^4,\qquad
 \mathcal E(V)=M^{-4}\mathbb E\sum_{y,a}|V_a(y)|^2 ,
\]
where \(a\) runs over the 18 orientation-colour
coordinates, one has exactly
\begin{equation}
 \mathcal E(W^R)
       =2q_{M,\beta}
          -\frac1\beta\sum_{\mu<\nu}
                                  \mathbb E_\mu X_{0,\mu\nu}^2
       \le2q_{M,\beta}.
 \label{eq:f15v5:energy}
\end{equation}
\end{proposition}

\begin{proof}
A translation permutes the possible offsets and
the cell labels, translates the ordered paths,
and preserves the fine Wilson law. This proves
stationarity. Haar averaging of each adjoint
rotation kills the mean of a colour vector.
Every covariance matrix is positive semidefinite
because its quadratic forms are expectations of
squares. Finally rotations preserve norms, and
\(\beta|\operatorname{Im}U_f|^2=2X_f-X_f^2/\beta\).
Sum this identity and use translation invariance
of the scalar fine energies.
\end{proof}

Conditional on the cell partition, vectors whose
bases lie in different cells have zero covariance
after the independent \(H_D\) averages. This was
introduced for the comparison and is \emph{not}
an estimate on the corresponding gauge-transported
physical correlations.

For all fixed \(r\ge1\), the inherited thermal bound gives
\[
 \mathbb E_\mu|Y_f^\gamma|^{2r}\le C_L4^r r!
\]
for any choice of transport paths, because rotations
do not change the norm. We use \(C\) or \(C_m\) for
finite constants depending only on \(L\), the fixed
dimension and, when indicated, the moment order.
They are independent of \(M,\beta,R\), filter length,
and source support. Numerical optimality is not asserted.

\subsection{A full-measure covariance identity with its cell-boundary cost}

For colour edge and face cochains use counting
\(\ell^1\) norms with Euclidean colour norm.
Let \(d=d_1\) be the ordinary oriented curl and set
\[
                        \epsilon_R=R/\sqrt\beta+1/R .
 \]

\begin{theorem}[Covariance Stein estimate for the stationary comparison]
\label{f15v5:Stein}
For all deterministic finite torus cochains
\(a\) on edges and \(b\) on faces,
\begin{equation}
 \left|\mathbb E\langle W^R,b\rangle
                    \langle W^R,da\rangle-\langle b,da\rangle\right|
       \le C\epsilon_R\|a\|_{\ell^1}\|b\|_{\ell^1}.
 \label{eq:f15v5:Stein}
\end{equation}
In particular the estimate remains valid when
the supports grow. Their cost is the displayed
coefficient norm, not an omitted support constant.
\end{theorem}

\begin{proof}
First take \(a=h_e\), equal to a colour vector
\(h\) on one oriented edge \(e=(x,\alpha)\) and
zero elsewhere, and put \(\ell=dh_e\).
For a random offset, call this edge interior if
each coordinate of its base within its cell lies
between 1 and \(R-3\). Its whole plaquette star
and a one-square detour around \(e\) then lie in
that cell. The exceptional offset probability is
at most \(16/R\), by a union bound over the four
coordinates. For small \(R\) this deliberately
loose bound is harmless.

Fix an interior offset and the colour matrices.
In its cell replace any ordered root path using
the unoriented edge \(e\) by the three-edge
detour around one fixed adjacent square.
Every original ordered path uses \(e\) at most
once. The resulting paths avoid \(e\), remain
inside the cell, and differ by at most one
conjugated elementary plaquette. Other cells'
paths already avoid \(e\). Denote the adapted
curvature field by \(W'\).

The fourth thermal moments and the conjugation
estimate
\(|\operatorname{Ad}_A v-v|\le2|A-I|\,|v|\)
give, uniformly in the affected coordinate,
\[
                       \|W'_f-W_f^R\|_2\le C/\sqrt\beta .
\]
For a star edge, the unadapted rooted link has
filling area at most \(3R\): its terminal step
crosses only the later coordinate steps in
its ordered tree path. The two possible endpoint
detours add at most two faces. Thus the adapted
star transport variables
\(T_f=\sqrt\beta\,|k_f-I|\) obey
\(\|T_f\|_4\le CR\), by the fixed-word filling
bound and the thermal fourth moments.

Now condition on all fine links except \(U_e\).
Because the adapted root paths avoid \(e\),
the left Haar flow
\[
 U_e(t)=(g_x')^{-1}H_D^{-1}
             \exp(th/\sqrt\beta)H_Dg_x'\,U_e
\]
has a generator independent of \(U_e\).
It preserves the conditional Haar reference.
For \(F=\langle W',b\rangle\), differentiation
of its finite compact integral gives the exact identity
\[
                           \mathbb E_\mu D_eF
                                  =\mathbb E_\mu F D_e(\beta S).
\]
No gauge density or derivative of a frame has been
dropped. The frame avoidance is what makes this
conditional integration legitimate.

The same pointwise differentiated plaquette identities
used in f14 V97--V98 give
\[
 |D_eW'_f-\sigma_{f,e}h|
       \le\frac{3|h|}{\sqrt\beta}
                               \sum_{u\in\partial f}T_u
       \quad(f\ni e),
\]
and \(D_eW'_f=0\) off the star. The action derivative
differs from \(\langle W',\ell\rangle\) by a sum
bounded by
\[
       \frac{2|h|}{\sqrt\beta}
           \sum_{f\ni e}|W'_f|
                            \sum_{u\in\partial f}T_u .
\]
There are only six star faces. Consequently
\[
 \|D_eF-\langle b,\ell\rangle\|_1
          \le C|h|\|b\|_1R/\sqrt\beta,\qquad
 \|D_e(\beta S)-\langle W',\ell\rangle\|_2
          \le C|h|R/\sqrt\beta .
\]
Also \(\|F\|_2\le C\|b\|_1\).
Cauchy--Schwarz in the exact Haar identity proves
the desired covariance comparison for \(W'\)
with error \(C|h|\|b\|_1R/\sqrt\beta\).
Replacing \(W'\) by \(W^R\) costs at most
\(C|h|\|b\|_1/\sqrt\beta\), using their
coordinate \(L^2\) difference and
\(\|\ell\|_1\le6|h|\).

For a noninterior offset do not use this flow.
For every fixed offset the thermal second moments
give
\[
 |\mathbb E_\mu\langle W^R,b\rangle
                         \langle W^R,\ell\rangle|
      +|\langle b,\ell\rangle|\le C|h|\|b\|_1 ,
\]
uniformly in the colour matrices. The offset
probability, not an unproved conditional moment,
therefore bounds this part by \(C|h|\|b\|_1/R\).
Average the two cases and sum over the edge
components of \(a\), proving \eqref{eq:f15v5:Stein}.
\end{proof}

This is a new use of the full compact Haar identity.
The restricted-chart translation estimates of f14
V93--V94 have not been transplanted to the full law.
In particular their chart-boundary and determinant
terms are not being silently set to zero.

\subsection{Energy saturation supplies a covariance bound}

We reuse the finite-range coexact contraction from
f14 V94, but apply the new full-measure estimate.
Let
\[
 \mathsf R=I-\Delta/16,\qquad
 q_n(\Delta)=\frac1{16}\sum_{s=0}^{n-1}\mathsf R^s,
 \qquad T_n=dd^*q_n(\Delta),\quad n\ge1.
\]
These operators act on all colours. If \(P_M\)
is the orthogonal projection onto the nonzero-mode
closed two-forms, then
\[
 T_n(k)=(1-r(k)^n)P_M(k),\qquad
 r(k)=1-\lambda(k)/16,\qquad 0\le T_n\le P_M\le I .
\]
Both \(T_n\) and \(P_M\) are zero on constant modes.
For a unit face-colour coordinate \(e_{x,a}\),
\begin{equation}
 T_ne_{x,a}=d a_{n,x,a},\qquad
 \|a_{n,x,a}\|_1\le n/4,\qquad
 \|T_ne_{x,a}\|_1\le3n/2 .
 \label{eq:f15v5:source-norms}
\end{equation}
Indeed the positive scalar kernel \(q_n\) has
mass \(n/16\), the \(\ell^1\) operator norm of
\(d^*\) is at most four, and that of \(d\) is
at most six. Periodization can only decrease
these absolute coefficient bounds.

For the torus lazy-walk return probability write
\[
        k_M(s)=M^{-4}\sum_{k\in(2\pi/M)\mathbb Z_M^4}r(k)^s.
\]
For \(s\ge1\),
\begin{equation}
                         k_M(s)\le C(s^{-2}+M^{-4}).
 \label{eq:f15v5:heat}
\end{equation}
For completeness, on the principal momentum cube
\(\lambda(k)\ge4|k|^2/\pi^2\), so
\(r(k)^s\le e^{-s|k|^2/(4\pi^2)}\).
The four one-dimensional grid sums are bounded
by \(M^{-1}+C/\sqrt s\); their product proves
\eqref{eq:f15v5:heat}. No infinite-volume
limit or spectral gap is assumed in this estimate.

Suppose now that
\begin{equation}
                   q_{M,\beta}\le9/2+\eta,\qquad 0\le\eta\le1 .
 \label{eq:f15v5:excess}
\end{equation}
This will be supplied by the already selected
full compact couplings, not postulated for every
Wilson coupling.

\begin{theorem}[Full-measure auxiliary residual and covariance estimates]
\label{f15v5:covariance}
Under \eqref{eq:f15v5:excess},
\begin{equation}
 \mathcal E((I-T_n)W^R)
       \le D_n:=2\eta+C(n^{-2}+M^{-4}+n^2\epsilon_R).
 \label{eq:f15v5:residual}
\end{equation}
Uniformly in all two coordinates of the torus,
\begin{equation}
 |(\Sigma_R-P_M)_{(x,a),(y,b)}|
       \le C\{\sqrt{D_n}+n^2\epsilon_R+n^{-2}+M^{-4}\}.
 \label{eq:f15v5:covariance-n}
\end{equation}
If \(\epsilon_R\le1\), choosing
\(n=\max\{1,\lfloor\epsilon_R^{-1/4}\rfloor\}\) gives
\begin{equation}
             \sup_{x,y,a,b}|(\Sigma_R-P_M)_{(x,a),(y,b)}|
                 \le C\bigl(\sqrt\eta+\epsilon_R^{1/4}+M^{-2}\bigr).
 \label{eq:f15v5:covariance-final}
\end{equation}
\end{theorem}

\begin{proof}
Apply \eqref{eq:f15v5:Stein} first with
\(b=e_{0,a}\), \(a=a_{n,0,a}\), and then
with \(b=T_ne_{0,a}\), keeping the same edge source.
Using \eqref{eq:f15v5:source-norms} gives
\[
 \begin{split}
 \mathbb E W_a^R(0)(T_nW^R)_a(0)
                         &=(T_n)_{(0,a),(0,a)}+O(n\epsilon_R),\\
 \mathbb E(T_nW^R)_a(0)^2
                         &=(T_n^2)_{(0,a),(0,a)}+O(n^2\epsilon_R).
 \end{split}
\]
The constants are uniform in the coordinate. Expand
the residual square, sum the 18 coordinates, and
use \(\mathcal E(W^R)\le9+2\eta\).
The exact Fourier trace identity is
\[
 9-2\sum_a(T_n)_{(0,a),(0,a)}
          +\sum_a(T_n^2)_{(0,a),(0,a)}
                            =9k_M(2n).
\]
At a nonzero momentum the projection rank including
colours is nine. The zero mode contributes the
remaining \(9/M^4\) through \(r(0)=1\), so it has
not been discarded. Equation \eqref{eq:f15v5:heat}
proves \eqref{eq:f15v5:residual}.

Stationarity converts \(\mathcal E\) into the
sum of the 18 variances at any site. Both
\(W^R\) and \(T_nW^R\) have energy at most
\(9+2\eta\): for the latter use the deterministic
global \(\ell^2\) contraction and then stationarity.
Cauchy--Schwarz therefore bounds the change of
any covariance when replacing \(W^R\) by
\(T_nW^R\) by \(2\sqrt{9+2\eta}\sqrt{D_n}\).
For the covariance of two filtered coordinates,
use \eqref{eq:f15v5:Stein} with
\(a=a_{n,x,a}\) and \(b=T_ne_{y,b}\).
Its difference from the corresponding \(T_n^2\)
entry is at most \(Cn^2\epsilon_R\), uniformly
in the separation.

Finally \(P_M-T_n^2\) is positive semidefinite,
with multiplier \((2r^n-r^{2n})P_M\).
Each of its diagonal entries is bounded by its
diagonal trace, at most \(18k_M(n)\).
The absolute value of an off-diagonal entry of
a positive matrix is at most the geometric mean
of its diagonal entries. This proves
\eqref{eq:f15v5:covariance-n}. Substitute the
indicated integer \(n\); its rounding changes
only the universal constant.
\end{proof}

The use of the energy excess here is essential:
coexact identities alone do not rule out an
additional low-frequency or constant component.
Also, covariance control is not a proof that
the auxiliary field is jointly Gaussian.

\subsection{Paying the comparison back to the actual fine filtered field}

Let \(p=2^j\), \(M\ge16p\), and now require
\(R\ge4p\). Use V4's row \(h=h_{p,\mu\nu}\)
and its actual rooted field \(Y^\gamma\), rooted
at 0 with increasing-coordinate paths. Put
\[
                 Z_p=\langle h,Y^\gamma\rangle,\qquad
                 Z_p^R=\langle h,W^R\rangle .
\]
Both are three-component vectors.

\begin{proposition}[Local rotation and cut error for the second moment]
\label{f15v5:root-comparison}
There is a constant independent of the scales such that
\begin{equation}
            \left|\mathbb E_\mu|Z_p|^2-\mathbb E|Z_p^R|^2\right|
                           \le Cp^5(R/\sqrt\beta+1/R).
 \label{eq:f15v5:root-comparison}
\end{equation}
In particular this comparison concerns an actual
fine Wilson observable, not a redefined block law.
\end{proposition}

\begin{proof}
The row support, its fine plaquettes and the
local paths from 0 lie in \(\{0,\ldots,2p\}^4\).
Except for an offset set of probability at most
\(8p/R\), this box lies in one comparison cell.
For each coordinate the bad fraction is at most
\(2p/R\), giving the union bound.

On a good offset let \(z\) be that cell's root
in a lifted box. Compare its ordered path to \(y\)
with its ordered path to 0 followed by the local
ordered path from 0 to \(y\).
Moving the later coordinate steps of the first
segment past the earlier steps of the second
requires at most \(12Rp\) square exchanges:
there are six ordered pairs of distinct coordinates,
with lengths at most \(R\) and \(2p\).
The fixed-word filling and fourth thermal moments
give a fourth-moment norm \(CRp/\sqrt\beta\)
for the difference of these transport matrices.
After conjugation of the curvature, H\"older gives
an \(L^2\) error \(CRp/\sqrt\beta\) per face.

Consequently, up to the \emph{common} rotation
\(\operatorname{Ad}_{H_Dg_0^D}\),
the two filtered vectors differ in \(L^2\) by
\(CRp^3/\sqrt\beta\), since \(\|h\|_1=p^2\).
Each vector separately has \(L^2\) norm at most
\(Cp^2\), uniformly for each fixed offset.
The common rotation leaves the squared norm
unchanged, even though \(g_0^D\) depends on \(U\).
The change of the second moment on good offsets
is thus at most \(CRp^5/\sqrt\beta\).

On a bad offset use only the separate second
moment bounds \(Cp^4\). The bad event depends
on the independent offset, and those bounds
hold for each offset, so its contribution is
at most \(Cp^5/R\). Summing the two cases proves
the assertion. No independence of the fine
plaquettes was used.
\end{proof}

Define the finite-torus Gaussian comparison number
\[
                    c_{p,M}=\langle h,P_Mh\rangle
\]
for one colour. Assume \eqref{eq:f15v5:excess} and
\(\epsilon_R\le1\). The entrywise covariance bound and
\(\|h\|_1^2=p^4\), followed by
\eqref{eq:f15v5:root-comparison}, give
\begin{equation}
 \left|\mathbb E_\mu|Z_p|^2-3c_{p,M}\right|
       \le C\left[
          p^4\bigl(\sqrt\eta+\epsilon_R^{1/4}+M^{-2}\bigr)
                               +p^5\epsilon_R\right].
 \label{eq:f15v5:actual-filter}
\end{equation}
The factor three is the number of colours.
This is the required quantitative growing-support
estimate; it no longer relies on applying a
fixed-coordinate limit to a growing row.

\begin{proposition}[Finite and infinite covariance normalizations]
\label{f15v5:finite-infinite}
For \(M\ge16p\),
\begin{equation}
                         |c_{p,M}-c_p|\le Cp^4/M^2,
 \label{eq:f15v5:finite-infinite}
\end{equation}
where \(c_p\) is V4's infinite-lattice number.
\end{proposition}

\begin{proof}
Take \(n=\lfloor M/8\rfloor\). The infinite and
periodic kernels \(T_n^2\) agree on differences
of points in the row support: its coordinate
diameter is at most \(2p\), the kernel range
in \(\ell^1\) is at most \(2n+2\), and no
nonzero periodic image can fall in that range.
The positive-kernel diagonal argument in the
previous covariance proof bounds every entry
of \(P_M-T_n^2\) by \(C(n^{-2}+M^{-4})\).
On the infinite lattice the same proof bounds
entries of \(P-T_n^2\) by \(Cn^{-2}\).
Thus the difference of the two projection
kernels on this support is at most \(C/M^2\).
Pair against \(h\) twice and use its
absolute coefficient mass \(p^2\).
\end{proof}

\subsection{Actual coarse energy at an increasing number of blocking steps}

Let \(\mathfrak e_4(p,M,\beta)\) denote the right
side of V4's energy comparison at \(m=1\):
\[
 \mathfrak e_4=
 C\frac{p^6\ell_p^{3/2}}{\sqrt\beta}
 +2\beta(1+p^4)
     \min\left\{1,\ 2500C_Lp^4
           e^{-r_*^2\beta/(144\pi^2p^4)}\right\},
 \qquad \ell_p=1+\log(2500C_Lp^4).
\]
The constants \(r_*,C\) are those of that proved
actual nonlinear comparison.

\begin{theorem}[Quantitative mean thermal energy of the literal block]
\label{f15v5:actual-energy}
Under the hypotheses of
\eqref{eq:f15v5:actual-filter}, for the actual
completed output face \(Q\),
\begin{equation}
 \left|\mathbb E_\mu\bigl[\beta(1-\tfrac12\Tr C_Q^{(j)})\bigr]
                                                -\tfrac32c_p\right|
 \le \mathfrak e_4+
 C\left[p^4\bigl(\sqrt\eta+\epsilon_R^{1/4}+M^{-2}\bigr)
                                      +p^5\epsilon_R\right].
 \label{eq:f15v5:actual-energy}
\end{equation}
This estimate is under the full fine Wilson law
and allows \(p\), \(R\) and \(j\) to grow.
\end{theorem}

\begin{proof}
V4 bounds the difference between the actual
thermal coarse energy and \(\tfrac12|Z_p|^2\)
in \(L^1(\mu)\) by \(\mathfrak e_4\).
Combine \eqref{eq:f15v5:actual-filter} and
\eqref{eq:f15v5:finite-infinite}, and absorb
the common \(p^4/M^2\) term in the constant.
\end{proof}

To verify that this is a nonempty growing regime,
use f14 V100's selected full compact sequence.
Its parameters satisfy
\[
 \beta=\beta_M^\diamond\longrightarrow\infty,\qquad
 \beta\le\log M,\qquad
 q_{M,\beta}\le9/2+\eta_M,\qquad
 \eta_M:=\lambda_M+142e^{-d_M}\longrightarrow0 .
\]
Thus \eqref{eq:f15v5:excess} holds eventually.
Choose \(R\) as the largest dyadic number at most
\(\beta^{1/4}\). Eventually \(R\ge4\), \(R\mid M\),
and \(\epsilon_R\le C\beta^{-1/4}\).
For any allowed dyadic \(p\) with \(4p\le R\),
\begin{equation}
 \left|\mathbb E_\mu[\beta(1-\tfrac12\Tr C_Q^{(j)})]
                                      -\tfrac32c_p\right|
 \le \mathfrak e_4+
 C\left[p^4\bigl(\sqrt{\eta_M}+\beta^{-1/16}+M^{-2}\bigr)
                                      +p^5\beta^{-1/4}\right].
 \label{eq:f15v5:selected-energy}
\end{equation}

For a concrete admissible selection, once both
quantities below exceed two, take \(p_M\)
to be the largest dyadic number not exceeding
\begin{equation}
        \min\left\{\beta^{1/128},
                         (\eta_M+\beta^{-1})^{-1/16}\right\}.
 \label{eq:f15v5:growing-choice}
\end{equation}
Both quantities tend to infinity, so
\(p_M\to\infty\) and \(j_M=\log_2p_M\to\infty\).
The conditions \(4p_M\le R\) and \(M\ge16p_M\)
hold eventually. Moreover
\[
 p_M^4\sqrt{\eta_M}\le(\eta_M+\beta^{-1})^{1/4},
 \qquad p_M^4\beta^{-1/16}\le\beta^{-1/32},
 \qquad p_M^5\beta^{-1/4}\le\beta^{-27/128}.
\]
The \(M^{-2}\) and \(\mathfrak e_4\) terms also
tend to zero by the same selection; for the
exponential term \(\beta/p_M^4\ge\beta^{31/32}\).
Therefore
\begin{equation}
 \mathbb E_\mu[\beta(1-\tfrac12\Tr C_Q^{(j_M)})]
                                  -\tfrac32c_{p_M}\longrightarrow0.
 \label{eq:f15v5:growing-energy}
\end{equation}
Using V4's proved bounds on \(c_p\) gives the
nonvanishing finite normalization
\begin{equation}
 \frac{12}{\pi^{18}}
 \le \liminf_M\mathbb E_\mu[\beta(1-\tfrac12\Tr C_Q^{(j_M)})]
 \le \limsup_M\mathbb E_\mu[\beta(1-\tfrac12\Tr C_Q^{(j_M)})]
 \le\frac23 .
 \label{eq:f15v5:noncollapse}
\end{equation}
The final upper bound uses
\(H_p^{(2)}\to4/9\) as \(p\to\infty\).
This does not assert that \(c_{p_M}\) itself
has a limit, nor that a mean lower bound is
a variance lower bound or distributional limit.

\subsection{What has changed, and what is still unpaid}

The actual nonlinear block now has a quantitatively
controlled mean thermal energy through a growing
number of steps. The support-size amplification,
cell-cut error, transport-path error and energy
excess all appear in the estimate. The proof did
not exchange the order of a fixed-coordinate
Gaussian limit and a growing block size.

The auxiliary stationary field was only a means
of obtaining the covariance estimate. Its
independent colour rotations cannot provide
physical clustering or a transfer operator.
The final handoff concerns the original compact
Wilson output and was proved separately.
The selection \eqref{eq:f15v5:growing-choice}
still has \(p_M/M\to0\); it is not control at a
fixed nonzero physical separation. The coupling
selection remains different from a calibrated
physical renormalization trajectory.

Further work must control more than this mean:
growing-scale laws and higher connected moments,
spatial dependence of the actual conditional
scores, and the remaining generated interactions.
Neither this second-moment comparison nor V4's
fixed-depth score lower bound supplies those
estimates. The interacting continuum construction
and positive physical Hamiltonian mass gap remain
unproved.

\subsection{Exact diagnostics and preservation}

The certificate is
\begin{center}\small
\path{scripts/verify_ym_f15_v5_stationary_comparison.py}.
\end{center}
It checks 80 exact real Fourier-block identities,
five zero-mode-inclusive energy identities,
24 primitive/filter coefficient-norm bounds,
eight interior-guard counts, 14 window-cut bounds
and 41,472 ordered two-root filling counts.
A separate scalar comparison checks stationarity,
trace preservation and 32 positive quadratic
forms. That scalar test models only the
second-moment cell averaging; it is not a
replacement for the Haar colour construction
in the proof.

The certificate has 6,298 bytes and SHA--256
\begin{center}\scriptsize\ttfamily
D071D2C81479DBBE008C911175245E61F21FE367DFD4C0EC20A7D59567506BFB.
\end{center}
These finite checks do not prove the analytic
Haar estimates, limiting claims or mass gap
by simulation. Their scope is the algebra
and geometry just enumerated.

The predecessor V4 has 82,664 bytes and SHA--256
\begin{center}\scriptsize\ttfamily
C7D91082B1920DF139BC27737AA854A40E1B4CF7193741667CF2D15CD1085592.
\end{center}
Its body is retained apart from the title and
this terminal insertion; reversing those two
changes recovers its bytes exactly. V5 replaces
V4 as the only active main note. The 11
protected sources and 47 prior protected
certificates remain unchanged. The companion
checkpoint above is complete; the next is V10.

\begin{firewall}
V5 proves a paid auxiliary covariance comparison
and transfers it to the mean thermal energy of
the actual completed Wilson block at increasing
depth. It does not prove a growing-scale
Gaussian law, physical clustering, conditional
score locality, an interacting continuum YM
construction, or a physical mass gap.
The full original goal remains active.
\end{firewall}
\section{V6: actual growing-depth laws and distributional noncollapse}
\label{f15v6:context}

\subsection{From a mean to a law, without changing the target measure}

V5 proves a nonzero mean thermal energy through
increasingly many actual blocking steps. A nonzero
mean alone can be carried by rare events. This
section proves a characteristic-function estimate,
pays the auxiliary-to-actual comparison, and obtains
a distributional limit on the same growing-depth
sequence. It also proves first-moment uniform
integrability and a bounded nonconstant output witness.

The fixed-depth law of V4 and the covariance
calculation of V5 are not repeated. The new input
is the nonlinear exponential test in the full
compact Haar identity, including the cost of
returning from edge-avoiding frames. The
characteristic-function method in f14 V93 was
for a restricted chart law; its bound is not
being assumed for the present full measure.
The V5 companion checkpoint remains current;
the next compulsory review is V10.

Throughout use V5's full Wilson law, independent
offset and cell rotations, and auxiliary field
\(W^R\). Thus \(R\mid M\), \(4\le R\le M\),
\(\beta\ge2\), and \(\log\beta/M\le L\).
Write
\[
       \epsilon_R=R/\sqrt\beta+1/R,\qquad
       S_b=\langle W^R,b\rangle,\qquad B=\|b\|_1 .
\]
The auxiliary variables do not redefine the
fine measure or the literal completed output.

\subsection{Exponential tests in the actual Haar identity}

\begin{theorem}[Characteristic Stein identity with frame-return cost]
\label{f15v6:Stein}
For every real colour edge cochain \(a\), real
colour face cochain \(b\), and \(t\in\mathbb R\),
\begin{equation}
 \left|\mathbb E e^{itS_b}S_{da}
              -it\langle b,da\rangle\mathbb E e^{itS_b}\right|
 \le C\|a\|_1
       \left\{\epsilon_R(1+|t|B)
                       +\beta^{-1/2}t^2B^2\right\}.
 \label{eq:f15v6:Stein}
\end{equation}
The constant depends only on the inherited
thermal-moment window. No source-support constant
is hidden in it.
\end{theorem}

\begin{proof}
First let \(a=h_e\) be supported on one edge.
Use the interior-offset event and adapted paths
from the proof of Theorem~\ref{f15v5:Stein}.
For fixed interior offset and fixed cell rotations,
let \(W'\) be the adapted field,
\(S'_b=\langle W',b\rangle\), and \(\ell=dh_e\).
The adapted root paths avoid the varied edge.
Consequently the conditional Haar flow used there
applies to the smooth bounded test
\(F'=\exp(itS'_b)\), with exact identity
\[
                         \mathbb E_\mu D_eF'
                                   =\mathbb E_\mu F'D_e(\beta S).
\]
All finite-regulator integrands are smooth on a
compact group product, so this differentiation
requires no limiting interchange.

V5's derivative and action estimates give
\[
 \begin{split}
 \|D_eS'_b-\langle b,\ell\rangle\|_1
                         &\le C|h|BR/\sqrt\beta,\\
 \|D_e(\beta S)-S'_\ell\|_1
                         &\le C|h|R/\sqrt\beta.
 \end{split}
\]
Since \(D_eF'=itF'D_eS'_b\) and \(|F'|=1\),
the adapted identity has error at most
\(C|h|R(1+|t|B)/\sqrt\beta\).

It remains to return to the original frames.
The one-square detour bound from V5 implies
\[
 \|S'_b-S_b\|_2\le CB/\sqrt\beta,\qquad
 \|S'_\ell-S_\ell\|_2\le C|h|/\sqrt\beta,\qquad
 \|S_\ell\|_2\le C|h|.
 \]
Also
\(|e^{itu}-e^{itv}|\le |t|\,|u-v|\).
Thus replacing the product \(F'S'_\ell\) by
\(e^{itS_b}S_\ell\) costs at most
\(C|h|(1+|t|B)/\sqrt\beta\).
Replacing the expectation in the leading derivative
term has the additional cost
\[
 |t|\,|\langle b,\ell\rangle|\,
       \mathbb E_\mu|e^{itS'_b}-e^{itS_b}|
                  \le C|h|\beta^{-1/2}t^2B^2 .
 \]
This last term cannot be omitted merely because
the original test is bounded.

For a noninterior offset use
\(\mathbb E_\mu|S_\ell|\le C|h|\) and
\(|\langle b,\ell\rangle|\le6|h|B\), not a
Haar identity across the cell boundary.
The exceptional offset probability is at most
\(16/R\). Its contribution is therefore at most
\(C|h|(1+|t|B)/R\), with uniform moments for
each fixed offset. Average the cases and sum
over the edge components of \(a\).
\end{proof}

Notice the different coefficients:
the quadratic-in-\(t\) frame-return term carries
\(\beta^{-1/2}\), not an unexplained zero.
It can be bounded further by \(\epsilon_R\),
but that loses useful scale information.

\begin{proposition}[Coexact characteristic functions]
\label{f15v6:coexact}
Let \(\ell=da\), \(A=\|a\|_1\),
\(J=\|\ell\|_1\), and \(v=\|\ell\|_2^2\).
Then
\begin{equation}
 \left|\mathbb E e^{itS_\ell}-e^{-vt^2/2}\right|
 \le CA\left[
       \epsilon_R\left(|t|+\tfrac12Jt^2\right)
                  +\frac{J^2|t|^3}{3\sqrt\beta}\right].
 \label{eq:f15v6:coexact}
\end{equation}
\end{proposition}

\begin{proof}
At a finite regulator the real score is bounded,
so its characteristic function \(\phi\) is
differentiable, with
\(\phi'(t)=i\mathbb E S_\ell e^{itS_\ell}\).
Theorem~\ref{f15v6:Stein} with \(b=\ell\) gives
\[
   |\phi'(t)+vt\phi(t)|
       \le CA\{\epsilon_R(1+|t|J)
                            +\beta^{-1/2}t^2J^2\}.
\]
For \(t\ge0\), solve this first-order equation
using \(\phi(0)=1\). Its error kernel is
\(\exp[-v(t^2-s^2)/2]\), which is at most one
for \(0\le s\le t\). Integrating the three
forcing terms proves the bound.
For \(t<0\), use conjugation of the characteristic
function of a real variable. The argument also
covers \(v=0\) and uses no inverse covariance.
\end{proof}

\subsection{From coexact scores to an arbitrary finite filtered source}

Assume, as in V5,
\[
                    q_{M,\beta}\le9/2+\eta,\quad 0\le\eta\le1 .
\]
Retain its \(T_n\), \(P_M\) and residual bound
\[
   D_n=2\eta+C(n^{-2}+M^{-4}+n^2\epsilon_R),
       \qquad \mathcal E((I-T_n)W^R)\le D_n .
\]
Constants can be enlarged so the same \(D_n\)
is used throughout. Set \(\kappa_n=n^{-2}+M^{-4}\).

\begin{theorem}[Characteristic comparison for a general source]
\label{f15v6:general-source}
For every real face-colour source \(b\), with
\(B=\|b\|_1\) and \(V_b=\langle b,P_Mb\rangle\),
\begin{align}
 \left|\mathbb E e^{itS_b}-e^{-t^2V_b/2}\right|
 \le C\bigl[
 &|t|B\sqrt{D_n}+t^2B^2\kappa_n \notag\\
 &+\epsilon_R nB(|t|+nBt^2)
               +\beta^{-1/2}n^3B^3|t|^3\bigr].
 \label{eq:f15v6:general-source}
\end{align}
\end{theorem}

\begin{proof}
Take \(a=d^*q_n(\Delta)b\), so \(\ell=da=T_nb\).
The coefficient-norm estimates used in V5 give
\[
             \|a\|_1\le nB/4,\qquad
             \|\ell\|_1\le3nB/2 .
\]
Stationarity and the residual energy bound imply,
by Minkowski applied to the source coefficients,
\[
          \|S_b-S_\ell\|_2
             =\|\langle (I-T_n)W^R,b\rangle\|_2
                                      \le B\sqrt{D_n}.
\]
Hence their characteristic functions differ by
at most \(|t|B\sqrt{D_n}\).

Apply Proposition~\ref{f15v6:coexact} to \(\ell\).
Finally \(P_M-T_n^2\) is positive semidefinite,
and V5's heat-kernel argument bounds each of
its entries by \(C\kappa_n\). Consequently
\[
       0\le V_b-\|\ell\|_2^2\le C B^2\kappa_n .
\]
For nonnegative variances \(u,v\),
\(|e^{-ut^2/2}-e^{-vt^2/2}|\le(t^2/2)|u-v|\).
Combine these estimates and absorb the numerical
coefficient-norm factors in \(C\).
\end{proof}

This theorem does not condition on the actual
coarse configuration. It is a characteristic
function estimate on a declared auxiliary
probability space, with a remaining comparison
to be paid below.

\subsection{The common rotation must be accounted for}

Fix a finite coarse-face family \(\mathcal Q\),
translated into a box of bounded coarse-coordinate
diameter. Its size and diameter are fixed while
\(p=2^j\) grows. Constants denoted \(C_{\mathcal Q}\)
may depend on this fixed family, not on \(p,M,\beta,R\).
Use a common fine root and increasing-coordinate
transport paths. Assume \(M\ge C_{\mathcal Q}p\)
and \(R\ge C_{\mathcal Q}p\), with the constants
large enough to contain all the local carriers.

For \(Q_i\in\mathcal Q\), let \(h_i\) be V4's
translated curvature row and define
\[
        Z_i=\langle h_i,Y^\gamma\rangle,\qquad
        Z_i^R=\langle h_i,W^R\rangle .
\]
Each row has \(\|h_i\|_1=p^2\) and
\(\|h_i\|_2^2\le1\).
For real colour vectors \(\xi_i\), put
\(\Xi=\sum_i|\xi_i|\).

\begin{proposition}[Actual-to-auxiliary characteristic comparison]
\label{f15v6:rotation}
For every such family,
\begin{equation}
 \left|\mathbb E_\mu e^{i\sum_i\xi_i\cdot Z_i}
                  -\mathbb E e^{i\sum_i\xi_i\cdot Z_i^R}\right|
       \le C_{\mathcal Q}
                   \left(\Xi\,\frac{Rp^3}{\sqrt\beta}+\frac pR\right).
 \label{eq:f15v6:rotation}
\end{equation}
\end{proposition}

\begin{proof}
The row supports, their plaquettes and their
local root paths fit in a box of side
\(C_{\mathcal Q}p\). This box lies in one
comparison cell except for an offset probability
at most \(C_{\mathcal Q}p/R\).
On a good offset, compare its cell-root paths
with its path to the common local root followed
by the local paths. The ordered square-exchange
count is at most \(C_{\mathcal Q}Rp\), exactly
as in V5's two-root comparison. The fourth
thermal moments therefore give, for each \(i\),
\[
    \|Z_i^R-\operatorname{Ad}_{H_Dg_0^D} Z_i\|_1
                        \le C_{\mathcal Q}Rp^3/\sqrt\beta .
\]
The elementary Lipschitz bound for a complex
exponential gives the first error in
\eqref{eq:f15v6:rotation}.

It is not legitimate to delete \(g_0^D(U)\)
on the sole ground that a fixed rotation
preserves a covariance. Instead condition on
\(U\) and the good offset. The matrix \(H_D\)
is independent Haar, so \(H_Dg_0^D(U)\) is
still Haar under this conditioning, even
though \(g_0^D\) depends on the fine field.
Its effect on the joint vector \((Z_i)_i\)
is therefore the average over one common rotation.
The unconditional joint law of \((Z_i)_i\)
is invariant under any fixed common colour
rotation: a fine gauge transformation at its
common root produces exactly that rotation,
and the Wilson law is gauge invariant.
Its common-rotation average is consequently
its original law.

This proves equality of the compared characteristic
functions after that rotation is averaged on
each good offset, up to the transport error.
On the remaining offsets both exponentials
have absolute value one, costing at most
twice the exceptional probability.
\end{proof}

The independent Haar matrix is essential to
this argument. A field-dependent rotation by
itself may change even the law of an isotropic
vector; for instance it can align that vector
with a fixed axis. The proof above uses
conditional Haar invariance and then the
actual common-root gauge symmetry.

\subsection{The actual nonlinear output, with a quantitative error}

Let \(Y_i^{\rm out}\) denote the thermal curvature
of the literal completed face \(Q_i\), transported
to the same common root. V4's nonlinear comparison,
applied to the fixed enlarged carrier of this
family, gives
\begin{equation}
 \|Y_i^{\rm out}-Z_i\|_1\le\mathfrak f_{\mathcal Q},\qquad
 \mathfrak f_{\mathcal Q}
   =C_{\mathcal Q}\frac{p^4\ell_p}{\sqrt\beta}
     +C_{\mathcal Q}\sqrt\beta(1+p^2)
           \min\{1,C_{\mathcal Q}C_Lp^4e^{-c_{\mathcal Q}\beta/p^4}\},
 \label{eq:f15v6:nonlinear}
\end{equation}
where \(\ell_p=1+\log(2500C_Lp^4)\) and
\(c_{\mathcal Q}>0\).
Only the constants in the local box and filling
count change from the one-face proof in V4:
the number of edges is at most \(C_{\mathcal Q}p^4\)
and each rooted edge has filling count at most
\(C_{\mathcal Q}p\). Thus its powers of \(p\)
are unchanged. This is not an estimate uniform
over families of unbounded diameter.

Define the finite positive semidefinite Gaussian
comparison matrix
\[
       \mathcal C_p
          =\bigl(\langle h_i,P h_k\rangle\bigr)_{i,k}
                                                   \otimes I_3
\]
using the infinite-lattice covariance projection.
Its diagonal blocks are \(c_pI_3\); degeneracy
of the whole matrix is allowed.
For \(B\ge0\) put
\[
 \mathfrak C_n(B)=B\sqrt{D_n}+B^2\kappa_n
               +\epsilon_R nB(1+nB)+\beta^{-1/2}n^3B^3 .
\]

\begin{theorem}[Measured growing-block characteristic comparison]
\label{f15v6:actual-CF}
Under the stated geometric and energy-excess
hypotheses, for \(\xi=(\xi_i)_i\),
\begin{align}
 \left|\mathbb E_\mu e^{i\sum_i\xi_i\cdot Y_i^{\rm out}}
                         -e^{-\xi^{\mathsf T}\mathcal C_p\xi/2}\right|
 \le C_{\mathcal Q}\bigg[
   &\mathfrak C_n(p^2\Xi)
       +\Xi\left(\frac{Rp^3}{\sqrt\beta}+\mathfrak f_{\mathcal Q}\right)
            \notag\\
   &+\frac pR+\frac{p^4\Xi^2}{M^2}\bigg].
 \label{eq:f15v6:actual-CF}
\end{align}
The same estimate without the nonlinear term
holds for the actual fine filtered vectors \(Z_i\).
\end{theorem}

\begin{proof}
Apply Theorem~\ref{f15v6:general-source} at \(t=1\)
to \(b=\sum_i h_i\otimes\xi_i\), whose
absolute coefficient norm is at most \(p^2\Xi\).
It compares the auxiliary characteristic function
with variance \(\langle b,P_Mb\rangle\).
V5's finite/infinite kernel comparison extends
to the present fixed enlarged carrier by the
same finite-range no-image argument.
The variance error is at most
\(C_{\mathcal Q}p^4\Xi^2/M^2\).
Use the Lipschitz bound of the Gaussian
exponential in its nonnegative variance.
Finally apply Proposition~\ref{f15v6:rotation}
and \eqref{eq:f15v6:nonlinear}.
This proves both assertions without an interchange
of growing supports and a fixed-coordinate limit.
\end{proof}

\subsection{Verification on the already selected growing sequence}

Keep exactly V5's selection:
\[
 \beta=\beta_M^\diamond,\qquad
 \eta_M=\lambda_M+142e^{-d_M},\qquad
 R=\text{largest dyadic number at most }\beta^{1/4},
\]
\[
 p=p_M=\text{largest dyadic number at most }
           \min\{\beta^{1/128},(\eta_M+\beta^{-1})^{-1/16}\},
 \quad
 n=\max\{1,\lfloor\epsilon_R^{-1/4}\rfloor\}.
\]
All assertions are for sufficiently large \(M\).
Here \(p,R,n\to\infty\), \(\beta\le\log M\),
\(\epsilon_R\) is bounded above and below by
positive constants times \(\beta^{-1/4}\),
and \(n\) is comparable to \(\beta^{1/16}\).
The fixed-family embedding requirements hold
eventually.

For frequencies in any fixed bounded set,
the right side of \eqref{eq:f15v6:actual-CF}
tends to zero uniformly. To check every
amplification explicitly,
\[
 p^2\sqrt{\eta_M}\le(\eta_M+\beta^{-1})^{3/8}\longrightarrow0.
\]
The remaining polynomial powers of \(\beta\),
using \(p\le\beta^{1/128}\), are bounded by
\[
 \begin{array}{c|c}
 \text{term}&\text{upper power of }\beta\\ \hline
 p^2\beta^{-1/16}&-3/64\\
 p^4n^{-2}&-3/32\\
 \epsilon_Rnp^2&-11/64\\
 \epsilon_Rn^2p^4&-3/32\\
 \beta^{-1/2}n^3p^6&-17/64\\
 Rp^3/\sqrt\beta&-29/128\\
 p/R&-31/128\\
 p^4/\sqrt\beta&-15/32 .
 \end{array}
\]
The last term carries the harmless factor
\(\ell_p\). The \(M^{-2}\) terms vanish since
\(\beta\le\log M\); the compact-complement
terms are exponentially small since
\(\beta/p^4\ge\beta^{31/32}\).
Thus the additional nonlinear-test cost has
been paid on the same sequence, not on a
silently smaller fixed-depth subsequence.

\begin{theorem}[Finite-family Gaussian comparison at growing depth]
\label{f15v6:family-law}
On this sequence, the bounded-Lipschitz distance
between the law of \((Y_i^{\rm out})_{i\in\mathcal Q}\)
and the centred Gaussian law with covariance
\(\mathcal C_{p_M}\) tends to zero.
This comparison permits a varying or degenerate
covariance matrix and asserts no convergence of
\(\mathcal C_{p_M}\) itself.
\end{theorem}

\begin{proof}
The characteristic functions differ by a quantity
tending to zero uniformly on bounded frequency
sets. The actual family is tight: the exact
identity \(|Y_i^{\rm out}|^2\le2X_i^{\rm out}\)
and V5's mean estimate bound the sum of its
second moments. The Gaussian family is tight
because its diagonal variances \(c_p\) are at
most one.

If the claimed distance did not tend to zero,
choose a subsequence on which it is separated
from zero. The finite covariance matrices,
being positive with bounded trace, have a
further convergent subsequence. Tightness
gives a further weakly convergent actual
subsequence. The characteristic-function
estimate forces its limit to have the
Gaussian characteristic function of that
limiting covariance, which determines the
law. Both subsequences therefore have the
same weak limit. Bounded-Lipschitz distance
metrizes weak convergence in this fixed
finite-dimensional space, a contradiction.
\end{proof}

No conditional distribution has been passed
to a limit in this argument. In particular
it gives no statement about conditioning
on the whole coarse environment.

\subsection{A normalized one-face law and a bounded noncollapse witness}

For one actual face at depth \(j_M\), set
\[
 V_M=\frac{Y_{p_M,Q}^{\rm out}}{\sqrt{c_{p_M}}},
       \qquad E_M=\frac{\beta(1-\tfrac12\Tr C_Q^{(j_M)})}{c_{p_M}} .
\]
Here \(8/\pi^{18}\le c_{p_M}\le1\) was proved
in V4, so these normalizations do not introduce
an uncontrolled small denominator.

\begin{theorem}[Actual growing-depth curvature and energy laws]
\label{f15v6:one-face}
Along the same selected sequence,
\begin{equation}
                 V_M\Longrightarrow N(0,I_3),\qquad
                 E_M\Longrightarrow G_*:=\tfrac12\chi_3^2 .
 \label{eq:f15v6:one-face}
\end{equation}
The \(E_M\)'s are uniformly integrable and
converge to \(G_*\) in the first Wasserstein
distance. In particular this is not only
a lower bound on a mean.
\end{theorem}

\begin{proof}
For one face the Gaussian comparison covariance
is \(c_pI_3\). The characteristic-function
estimate is uniform on bounded frequency
sets; replacing a fixed frequency by its
value divided by \(\sqrt{c_p}\) remains in
such a bounded set. This proves the first
weak convergence by the same tightness and
uniqueness argument. The analogous conclusion
holds for the actual fine filtered vector
\(Z_p/\sqrt{c_p}\), by the last assertion
of Theorem~\ref{f15v6:actual-CF}.

V4's energy comparison, divided by \(c_p\),
shows
\[
           E_M-\tfrac12|Z_p/\sqrt{c_p}|^2
                                  \longrightarrow0\quad\text{in }L^1 .
\]
Its error tends to zero on the present sequence
as already checked in V5. The continuous mapping
theorem and this vanishing difference prove
the energy law in \eqref{eq:f15v6:one-face}.
V5's mean estimate gives in addition
\(\mathbb E E_M\to3/2=\mathbb EG_*\).

We spell out why this last equality matters.
For nonnegative random variables converging
weakly to an integrable variable with converging
means, for every fixed \(K>0\),
\[
 \mathbb E(E_M-K)_+
        =\mathbb EE_M-\mathbb E\min(E_M,K)
                 \longrightarrow\mathbb E(G_*-K)_+ .
\]
The function \(\min(\,\cdot\,,K)\) is bounded
and continuous on the nonnegative half-line.
Also
\[
 \mathbb E[E_M{\bf1}_{E_M>2K}]
                           \le2\mathbb E(E_M-K)_+ .
\]
Taking first the upper limit in \(M\) and then
\(K\to\infty\) proves uniform integrability
of the tail of the sequence; its finitely many
remaining variables are individually integrable.

For completeness, first Wasserstein convergence
follows by the same truncation. A one-Lipschitz
test normalized by \(f(0)=0\) satisfies
\(|f(x)|\le x\). Replacing \(x\) by \(x\wedge K\)
costs at most the corresponding excess tail.
For fixed \(K\) the bounded Lipschitz test
family converges uniformly in expectation
under weak convergence. Letting \(K\) grow
and using the proved tail bound gives the
supremum over all such tests.
\end{proof}

In particular, for any fixed \(0<a<b<\infty\),
\[
           \mu\{a<E_M<b\}\longrightarrow
                         \mathbb P\{a<G_*<b\}>0 .
\]
The limiting density is
\((2/\sqrt\pi)s^{1/2}e^{-s}\) on \(s>0\),
so there are no endpoint atoms. Define the
bounded gauge-invariant observable
\[
                              F_M=\min(E_M,1).
\]
Weak convergence and boundedness give
\begin{equation}
       \Var_\mu(F_M)\longrightarrow
              v_*:=\Var(\min(G_*,1))>0 .
 \label{eq:f15v6:bounded-witness}
\end{equation}
One explicit positive lower bound is
\[
       v_*\ge\frac14\,\mathbb P\{G_*\le1/2\}
                                \mathbb P\{G_*\ge1\}>0 .
\]
Indeed use
\(\Var Z=\tfrac12\mathbb E(Z-Z')^2\) with an
independent copy and the two separated events.
This gives a bounded noncollapse witness on
the actual growing-depth output.

First-moment uniform integrability is not
uniform integrability of \(E_M^2\).
No convergence of the unbounded energy variance
or all higher moments is inferred here.
Also the witness is explicitly thermally
normalized: the unscaled character
\(\tfrac12\Tr C_Q^{(j_M)}\) still tends to one,
since V5 bounds its mean defect by \(O(\beta^{-1})\).
Thus \eqref{eq:f15v6:bounded-witness} is not a
claim of noncollapse for that different observable.

\subsection{The next genuinely interacting requirement}

The mean estimate has now been upgraded to
typical, distributional thermal fluctuations
of the literal compact block at growing depth.
The auxiliary rotations, frame return,
source norms, cell cuts, zero modes and
nonlinear errors are all retained in the route.

This remains a microscopic ultraviolet result:
the selected \(p_M/M\) tends to zero, the
limiting comparison is Gaussian, and no
physical coupling/length relation has been
identified. A bounded thermal witness does
not construct a physical Hilbert space,
reflection-positive transfer or Hamiltonian.
The actual conditional-score locality and
control of the remaining generated interactions
are still missing. Those, rather than another
fixed-observable Gaussian limit, are the
next obstacles toward the interacting continuum
and its positive physical mass gap.

\subsection{Checks and preservation}

The certificate is
\begin{center}\small
\path{scripts/verify_ym_f15_v6_growing_law.py}.
\end{center}
It checks 16 exact joint finite-group rotation
averages and an explicit failure when the
independent average is omitted. The finite
group models only the averaging algebra,
not the Haar Wilson law. It also checks
36 integrated forcing polynomials, all nine
displayed negative scaling powers, four
three-colour energy normalizations, 12
tail-truncation identities and 32 rare-spike
mean examples. These are rational finite
diagnostics, not simulations or a machine
proof of the analytic limiting argument.

The certificate has 4,788 bytes and SHA--256
\begin{center}\scriptsize\ttfamily
A1B243B248AAC21C0EDC116A8B45A48CDDF09F01820A658311DDE4F6936752CF.
\end{center}
The predecessor V5 has 108,244 bytes and SHA--256
\begin{center}\scriptsize\ttfamily
51B6DB7A7E50ED9FAFA414DAEEA08DE3F7221A55622197B653590F468845E25A.
\end{center}
Reversing the title change and this terminal
insertion recovers that predecessor exactly.
V6 replaces V5 as the sole active main note.
The 11 protected sources and 48 earlier
protected certificates remain unchanged.
The next companion checkpoint remains V10.

\begin{firewall}
V6 proves growing-depth thermal distributional
noncollapse for the actual compact output,
including a bounded witness and first-moment
tail control. It does not prove physical
clustering, actual conditional-score locality,
an interacting continuum construction, or
the full Yang--Mills mass gap. The original
goal remains active.
\end{firewall}
\section{V7: Growing-depth moments and actual retained-source information}
\label{f15v7:context}

V6 proved a growing-depth Gaussian comparison in distribution,
first-moment energy tail control, and a bounded noncollapse witness.
It expressly did not prove convergence of the unbounded energy
variance. V4's lower bound for a retained fine-action score was
only at fixed depth. This section closes those two particular
gaps on an explicitly slower growing-depth sequence. It does
not reassert a fixed-coordinate limit on a growing support.

The new ingredients are a Fourier cutoff with its complete
frequency cost, the existing full-compact thermal tail, and
errors small enough to sum over a growing fine-source support.
The last step tests an actual conditional score by an actual
coarse observable; no conditional distribution is passed to
a limit. This is progress toward identifying what the
nonlinear output retains, not a construction of its local
effective action or of a physical Hamiltonian.
The full continuum problem remains open in these notes;
the distinction from the official existence-and-gap target
is essential.\footnote{For the target and its public status,
see the Clay Mathematics Institute,
\href{https://www.claymath.org/millennium/yang-mills-the-maths-gap/}
{Yang--Mills and the Mass Gap}, checked 15 September 2026.}

The companion review completed at V5 remains in force.
The 11 protected sources, including the three supplied
SM/stationarity/Einstein notes, are unchanged.
They supply no missing conditional-locality estimate.
The next scheduled companion review is V10.

\subsection{The inherited comparison, without a hidden frequency restriction}

Keep the full compact Wilson law, the literal completed
block map, the root conventions and the hypotheses of
Theorem~\ref{f15v6:actual-CF}. Let \(\mathcal Q\) be a fixed
nonempty family of coarse faces in a fixed coarse-coordinate
box, \(N=|\mathcal Q|\), \(d=3N\), and \(p=2^j\).
Write
\[
 Y=(Y_i^{\rm out})_{i\in\mathcal Q},\qquad
 K_p=(\langle h_i,P h_k\rangle)_{i,k},\qquad
 G_p\sim N(0,K_p\otimes I_3).
\]
All finite carriers must embed in the fine torus, as before.
The vectors use a common root; scalar energies are independent
of this choice. In this section constants may depend on the
fixed family and on its diameter, but not on \(p,M,\beta\).
No estimate uniform in an increasing family size is asserted.

Set
\[
 \epsilon_R=R/\sqrt\beta+1/R,\quad
 \kappa_n=n^{-2}+M^{-4},\quad
 D_n=2\eta+C(n^{-2}+M^{-4}+n^2\epsilon_R),
\]
where \(q_{M,\beta}\le9/2+\eta\) and \(0\le\eta\le1\).
With \(\mathfrak f_{\mathcal Q}\) exactly as in
\eqref{eq:f15v6:nonlinear}, define nonnegative coefficients
\begin{align}
 a_0&=p/R,\notag\\
 a_1&=p^2\sqrt{D_n}+\epsilon_Rnp^2
                       +Rp^3/\sqrt\beta+\mathfrak f_{\mathcal Q},\notag\\
 a_2&=p^4(\kappa_n+\epsilon_Rn^2+M^{-2}),\notag\\
 a_3&=\beta^{-1/2}n^3p^6.
 \label{eq:f15v7:coefficients}
\end{align}
Expanding \(\mathfrak C_n(p^2\Xi)\) in V6 and using
\(\Xi\le\sqrt N|\xi|\) gives, for every real frequency,
\begin{equation}
 \left|\mathbb E e^{i\xi\cdot Y}
                  -e^{-\xi^{\mathsf T}(K_p\otimes I_3)\xi/2}\right|
       \le C_{\mathcal Q}\sum_{s=0}^3a_s|\xi|^s,
                   \qquad \xi\in\mathbb R^d.
 \label{eq:f15v7:CF}
\end{equation}
For large frequencies the right side may exceed two; this
does not invalidate the bound. It is important that the
display is valid at every frequency, not just pointwise on
each fixed compact set. The cell-cut error \(a_0\) is retained.

\subsection{Thermal tails and a quantitative polynomial comparison}

\begin{proposition}[Unconditional norm tails on the actual compact output]
\label{f15v7:tails}
For \(\sigma_p^2=\pi^2Np^8\),
\begin{equation}
       \mathbb E_\mu \exp(|Y|^2/\sigma_p^2)\le C_L .
 \label{eq:f15v7:norm-mgf}
\end{equation}
For each nonnegative integer \(m\) and each \(T>0\),
\begin{align}
 \mathbb E_\mu[|Y|^m{\bf1}_{|Y|>T}]
       &\le C_{m,\mathcal Q}p^{4m}
                            e^{-c_{\mathcal Q}T^2/p^8},\notag\\
 \mathbb E[|G_p|^m{\bf1}_{|G_p|>T}]
       &\le C_{m,\mathcal Q}e^{-c_{\mathcal Q}T^2}.
 \label{eq:f15v7:tail-moments}
\end{align}
Here \(C_L\) is the inherited fine thermal constant with
one fixed admissible \(L\); it does not grow with \(M\).
\end{proposition}

\begin{proof}
The exact compact inequality is
\[
       |Y_i^{\rm out}|^2\le\beta(\vartheta_{Q_i}^{(j)})^2 .
\]
Equation \eqref{eq:f15v3:mgf}, with \(256^j=p^8\), therefore gives
\[
       \mathbb E\exp(|Y_i^{\rm out}|^2/(\pi^2p^8))\le C_L .
\]
Convexity of the exponential applied to the average of
these \(N\) nonnegative exponents proves
\eqref{eq:f15v7:norm-mgf}. Independence is not used.

For \(z\ge0\), \(z^{m/2}\le C_m e^{z/2}\).
On \(z>T^2/\sigma_p^2\),
\[
 \sigma_p^m z^{m/2}
 \le C_m\sigma_p^m e^z
                            e^{-T^2/(2\sigma_p^2)}.
\]
Take \(z=|Y|^2/\sigma_p^2\) and integrate.
For the Gaussian, all eigenvalues of \(K_p\otimes I_3\)
are nonnegative and their sum is \(3Nc_p\le3N\).
With \(A=12N\), Gaussian diagonalization gives
\[
 \mathbb E e^{|G_p|^2/A}
       =\prod_\ell(1-2\lambda_\ell/A)^{-1/2}
       \le e^{1/2},
\]
because \(2\lambda_\ell/A\le1/2\) and
\(-\log(1-x)\le2x\) on that interval.
Zero eigenvalues are harmless. The same tail argument
with this fixed \(A\) proves the second estimate.
\end{proof}

\begin{theorem}[Fourier cutoff with the complete moment cost]
\label{f15v7:cutoff}
Let \(P_m\) be any fixed homogeneous polynomial of degree
\(m\) on \(\mathbb R^d\). For \(T\ge1\),
\begin{align}
 |\mathbb E P_m(Y)-\mathbb E P_m(G_p)|
 \le C_{P_m,\mathcal Q}\bigg[
     &\sum_{s=0}^3a_sT^{m-s}
          +p^{4m}e^{-c_{\mathcal Q}T^2/p^8}
          +e^{-c_{\mathcal Q}T^2}\bigg].
 \label{eq:f15v7:polynomial-error}
\end{align}
The result for an arbitrary fixed polynomial follows
by summing its homogeneous parts.
\end{theorem}

\begin{proof}
Choose a fixed smooth cutoff \(0\le\chi\le1\) equal to
one on the unit ball and zero outside the ball of radius two.
Let \(F_T(y)=P_m(y)\chi(y/T)=T^mF_1(y/T)\).
For the convention
\(\widehat F(\xi)=\int e^{-i\xi\cdot y}F(y)\,dy\),
\[
 \widehat F_T(\xi)=T^{m+d}\widehat F_1(T\xi),\qquad
 \int |\xi|^s|\widehat F_T(\xi)|\,d\xi
       =T^{m-s}\int |\zeta|^s|\widehat F_1(\zeta)|\,d\zeta.
\]
All four integrals, \(s=0,1,2,3\), are finite:
integration by parts with \((1-\Delta)^k F_1\),
where \(2k>d+3\), bounds the Fourier transform by
\(C_k(1+|\zeta|^2)^{-k}\).
Fourier inversion, Fubini and \eqref{eq:f15v7:CF} therefore give
\[
 |\mathbb EF_T(Y)-\mathbb EF_T(G_p)|
                         \le C_{P_m,\mathcal Q}\sum_{s=0}^3a_sT^{m-s}.
\]
Fubini is justified by the integrable Fourier transform
and the unit total mass of each law. Finally,
\(|P_m-F_T|\le C_{P_m}|y|^m{\bf1}_{|y|>T}\).
Apply Proposition~\ref{f15v7:tails} to the two removed tails.
This proves the claimed error without differentiating
a characteristic-function error at the origin.
\end{proof}

\subsection{One slower sequence works for every fixed moment}

Keep V5's selected couplings \(\beta=\beta_M^\diamond\)
and excess \(\eta_M=\lambda_M+142e^{-d_M}\).
For all sufficiently large \(M\), put
\begin{equation}
 \delta_M=\eta_M+\beta_M^{-1},\qquad
 u_M=\log(1/\delta_M),\qquad
 p_M=\text{largest dyadic integer at most }u_M .
 \label{eq:f15v7:slow-depth}
\end{equation}
Use the same comparison-cell choice
\[
 R=\text{largest dyadic integer at most }\beta^{1/4},
       \qquad n=\max\{1,\lfloor\epsilon_R^{-1/4}\rfloor\},
       \qquad T=p_M^4u_M .
\]
This changes the growing block depth, not the selected
Wilson couplings. In particular it does not upgrade
V6's faster growing sequence by an unproved inference.
Eventually
\[
 2\le u_M\le\log\beta,\quad u_M/2<p_M\le u_M,\quad
 \eta_M\le e^{-u_M},\quad \beta\ge e^{u_M}.
\]
Thus \(j_M=\log_2p_M\to\infty\) and \(p_M/M\to0\).
The original dyadic torus, embedding, cell and finite-range
conditions all hold eventually; in particular \(p_M/R\to0\).

\begin{theorem}[All fixed actual curvature moments at growing depth]
\label{f15v7:all-moments}
On \eqref{eq:f15v7:slow-depth}, for every fixed polynomial
\(P\) and every fixed \(A\ge0\),
\begin{equation}
 p_M^A\left|\mathbb E_\mu P(Y)
                     -\mathbb E P(G_{p_M})\right|\longrightarrow0 .
 \label{eq:f15v7:all-moments}
\end{equation}
This is a comparison with a varying Gaussian covariance;
no limit of \(K_{p_M}\) is required. The same sequence
works for all fixed degrees; the degree is not allowed
to increase with \(M\).
\end{theorem}

\begin{proof}
Here \(\epsilon_R\asymp\beta^{-1/4}\) and
\(n\asymp\beta^{1/16}\). Consequently
\[
 \sqrt{D_n}\le C(\sqrt\eta+\beta^{-1/16}+M^{-2}),\qquad
 \kappa_n+\epsilon_Rn^2\le C(\beta^{-1/8}+M^{-4}).
\]
Since \(\beta\le\log M\), the terms involving \(M^{-2}\)
are smaller than \(C\beta^{-2}\) for large \(M\).
The polynomial parts of \(\sum_s a_s\) are bounded by
the following terms, up to a fixed constant:
\[
 \begin{gathered}
 p\beta^{-1/4},\quad p^2\sqrt\eta,\quad
 p^2\beta^{-1/16},\quad p^2\beta^{-2},\quad
 p^2\beta^{-3/16},\quad p^3\beta^{-1/4},\\
 p^4\ell_p\beta^{-1/2},\quad
 p^4\beta^{-1/8},\quad p^4\beta^{-2},\quad
 p^6\beta^{-5/16}.
 \end{gathered}
\]
With \(p\le u\), \(\ell_p\le C(1+\log u)\),
\(\eta\le e^{-u}\), and \(\beta\ge e^u\), their sum
is at most \(Cu^7e^{-u/16}\).
The compact-complement part of \(\mathfrak f_{\mathcal Q}\)
is at most
\[
 C\sqrt\beta(1+p^2)p^4e^{-c\beta/p^4}
                      \le C e^{-c'e^u/u^4}
\]
for large \(u\), uniformly for \(\beta\ge e^u\).
Indeed \((\log\beta)/\beta\) is decreasing there and
\(u^4(\log\beta)/\beta\le u^5e^{-u}\to0\);
the polynomial prefactors are absorbed by half the
negative exponential. The minimum with one in the
inherited bound only improves this estimate.

Now \(1\le T\le u^5\) and \(T^2/p^8=u^2\).
For a degree-\(m\) homogeneous part the right side of
\eqref{eq:f15v7:polynomial-error} is bounded by
\begin{equation}
 C_{m,P,\mathcal Q}\left[
 u^{5m+7}e^{-u/16}
       +u^{4m}e^{-cu^2}
       +u^{5m}e^{-c'e^u/u^4}\right].
 \label{eq:f15v7:rapid-error}
\end{equation}
The Gaussian tail has been absorbed in the middle term
using \(T\ge u\). Multiplication by \(p^A\le u^A\)
still leaves a quantity tending to zero.
Sum the finitely many homogeneous parts to conclude.
\end{proof}

\subsection{Actual energy moments, variances and connected coefficients}

Set
\[
 X_i=\beta(1-\tfrac12\Tr C_{Q_i}^{(j_M)}),\qquad
 O_i=\tfrac12|Y_i^{\rm out}|^2,\qquad
 O_i^{\rm G}=\tfrac12|(G_{p_M})_i|^2.
\]
For every compact unit quaternion, not just on a
small-field event, the identities are
\begin{equation}
       X_i-O_i=\frac{X_i^2}{2\beta}\ge0,\qquad
                         0\le O_i\le X_i .
 \label{eq:f15v7:energy-identity}
\end{equation}
Also \(1-\cos\vartheta\le\vartheta^2/2\), so V3 gives
\begin{equation}
        \mathbb E X_i^r
              \le r!C_L(\pi^2p^8/2)^r
                         \qquad(r=1,2,\ldots).
 \label{eq:f15v7:rough-energy}
\end{equation}

\begin{theorem}[Growing-depth energy moment and cumulant comparison]
\label{f15v7:energies}
For any fixed list \(i_1,\ldots,i_r\) from \(\mathcal Q\),
with repetitions permitted, and every fixed \(A\ge0\),
\begin{equation}
 p^A\left|\mathbb E\prod_{\ell=1}^r X_{i_\ell}
                    -\mathbb E\prod_{\ell=1}^r O_{i_\ell}^{\rm G}
       \right|\longrightarrow0 .
 \label{eq:f15v7:energy-moments}
\end{equation}
The same assertion holds for their joint cumulants.
In particular, uniformly for the fixed indices,
\begin{align}
 \mathbb E X_i&=\tfrac32c_p+o(1),\notag\\
 \operatorname{Cov}(X_i,X_k)&=\tfrac32(K_p)_{ik}^2+o(1),\notag\\
 \Var(X_i/c_p)&\longrightarrow\tfrac32 .
 \label{eq:f15v7:energy-variance}
\end{align}
For every fixed \(r\ge1\), the one-face cumulant satisfies
\begin{equation}
        \operatorname{cum}_r(X_i/c_p)
                              \longrightarrow\tfrac32(r-1)! .
 \label{eq:f15v7:gamma-cumulants}
\end{equation}
\end{theorem}

\begin{proof}
By \eqref{eq:f15v7:rough-energy},
\(\|X_i\|_q+\|O_i\|_q\le C_qp^8\) for each fixed
integer \(q\ge1\). The exact discrepancy has
\[
       \|X_i-O_i\|_q=(2\beta)^{-1}\|X_i\|_{2q}^2
                                             \le C_qp^{16}/\beta .
 \]
Telescope a product of \(r\) factors and apply H\"older
with exponent \(r\) to each summand. This proves
\[
 \left|\mathbb E\prod_{\ell=1}^rX_{i_\ell}
             -\mathbb E\prod_{\ell=1}^rO_{i_\ell}\right|
                              \le C_rp^{8(r+1)}/\beta .
 \]
It tends to zero even after multiplication by \(p^A\)
on the selected path. The remaining product is a fixed
degree-\(2r\) polynomial in \(Y\), so
Theorem~\ref{f15v7:all-moments} proves
\eqref{eq:f15v7:energy-moments}. Gaussian moments of
each fixed degree are bounded uniformly, since the
covariance trace is bounded. Actual moments of that
degree are therefore bounded as well.

For a finite list of random variables, use the
moment definition of its joint cumulant,
\[
 \operatorname{cum}(Z_1,\ldots,Z_r)
   =\sum_{\pi}(|\pi|-1)!(-1)^{|\pi|-1}
                      \prod_{B\in\pi}\mathbb E\prod_{\ell\in B}Z_\ell ,
\]
where the sum is over partitions of \(\{1,\ldots,r\}\).
It is a fixed finite polynomial in the proved moments.
Their uniform boundedness and a product telescoping
argument give the asserted cumulant comparison at
the same superalgebraic-in-\(p\) accuracy.

To compute the Gaussian coefficients without a
nonsingularity assumption, write each independent
colour as \(K_p^{1/2}Z\), with \(Z\) standard normal.
For \(D_t=\operatorname{diag}(t_i)\) sufficiently small,
orthogonal diagonalization of \(K_p^{1/2}D_tK_p^{1/2}\)
and the one-dimensional Gaussian integral give
\begin{equation}
 \log\mathbb E e^{\sum_i t_iO_i^{\rm G}}
   =-\tfrac32\log\det(I-K_pD_t)
   =\tfrac32\sum_{r\ge1}\frac1r\Tr((K_pD_t)^r).
 \label{eq:f15v7:Gaussian-logdet}
\end{equation}
The determinant identity also holds when \(K_p\) is
singular, by \(\det(I-AB)=\det(I-BA)\).
The first two coefficients are \(3c_p/2\) and
\(3(K_p)_{ik}^2/2\). For one face the logarithm is
\(-\tfrac32\log(1-c_pt)\), whose \(r\)-th derivative
at zero is \(\tfrac32(r-1)!c_p^r\).
Finally \(8/\pi^{18}\le c_p\le1\) permits division
by any fixed power of \(c_p\) in the errors.
\end{proof}

These are genuine unbounded energy variances and
connected moments of the actual nonlinear output,
not variances of V6's bounded witness. Conversely,
all fixed moments do not supply a uniform positive-source
Laplace window. As an explicit logical counterexample,
take \(G_*\sim\frac12\chi_3^2\), an independent Bernoulli
\(B_n\) of success probability \(e^{-n}\), and
\(Z_n=G_*+n^2B_n\). Every fixed moment of \(Z_n\)
tends to that of \(G_*\), by the binomial expansion
and \(n^{2r}e^{-n}\to0\). Yet for every \(0<t<1\),
\[
 \mathbb E e^{tZ_n}
   =\mathbb E e^{tG_*}(1-e^{-n}+e^{tn^2-n})
                                               \longrightarrow\infty .
\]
This example is not claimed to be a Wilson law; it
rules out an inference from moment convergence alone.
The actual bound \eqref{eq:f15v7:rough-energy} has a
shrinking source scale of order \(p^{-8}\).
Equation \eqref{eq:f15v7:Gaussian-logdet} is a Gaussian
comparison formula, not a proved analytic expansion
of the full generated effective action.

\subsection{Uniform mixed fine/coarse moments in the growing carrier}

The score application requires one more uniformity
statement. Fix a coarse face \(Q\) and its row \(h=h_Q\).
Let \(f\) be any positively oriented fine face in a box
of side \(C p\) around its carrier, with fixed \(C\).
Transport both the fine curvature \(Y_f\) and
\(Y_Q^{\rm out}\) to the common root in that box.
Only two vector coordinates are being compared at
a time; the number of tested fine faces may grow later.

\begin{proposition}[Uniform two-vector comparison]
\label{f15v7:mixed}
All fixed mixed polynomial moments of
\((Y_f,Y_Q^{\rm out})\) agree, with errors
\(o(p^{-A})\) for every fixed \(A\), with those of
a three-colour Gaussian pair whose scalar covariance is
\begin{equation}
       \begin{pmatrix}
       \langle e_f,P e_f\rangle & \langle e_f,P h\rangle\\
       \langle e_f,P h\rangle & \langle h,P h\rangle
       \end{pmatrix}
        =
       \begin{pmatrix}1/2&v(f)\\v(f)&c_p\end{pmatrix},
                 \qquad v=P h .
 \label{eq:f15v7:mixed-covariance}
\end{equation}
The errors are uniform in \(f\) in the indicated box
and in its six possible positive orientations.
For the corresponding actual fine energy \(X_f\),
\begin{equation}
 \sup_f\left|\operatorname{Cov}_\mu(X_f,X_Q)
                                      -\tfrac32v(f)^2\right|
                                   =o(p^{-A})
                        \quad\hbox{for every fixed }A\ge0 .
 \label{eq:f15v7:mixed-energy}
\end{equation}
\end{proposition}

\begin{proof}
This does not follow from the fixed-coarse-family
statement just by renaming an index. Apply instead
V6's general-source estimate to
\(b=e_f\otimes\xi_1+h\otimes\xi_2\).
Its coefficient norm is at most
\(|\xi_1|+p^2|\xi_2|\le p^2(|\xi_1|+|\xi_2|)\).
Both supports and their ordered local root paths lie
in the same fixed-shape box of diameter \(O(p)\).
The V6 cell-cut probability is therefore \(O(p/R)\),
uniformly in \(f\).

In the good cell, comparing the local and cell
roots has filling count \(O(Rp)\).
For the single fine row the resulting \(L^1\)
transport error is at most \(CRp/\sqrt\beta\);
for the block row it is at most \(CRp^3/\sqrt\beta\).
The former is bounded by the latter.
Conditional Haar averaging and the common-root gauge
symmetry apply to the pair exactly as in
Proposition~\ref{f15v6:rotation}.
There is no nonlinear error on the single fine row;
the coarse error remains \(\mathfrak f_{\mathcal Q}\)
with constants for the fixed enlarged box.
The finite/infinite covariance error is bounded by
\(Cp^4(|\xi_1|+|\xi_2|)^2/M^2\), since the same
finite-range argument covers these two row supports.
Thus \eqref{eq:f15v7:CF} holds in dimension six,
with the same coefficient powers and constants
independent of the fine face.

The fine thermal exponential estimate gives the
norm tail used above (even with scale \(p^8\), since
\(p\ge1\)); the coarse one is unchanged.
The Fourier argument and
\eqref{eq:f15v7:rapid-error} consequently hold
uniformly for this pair. Projection translation
and orientation symmetry give \(P_{ff}=1/2\);
the matrix in \eqref{eq:f15v7:mixed-covariance}
is positive semidefinite as a Gram matrix.

The fine energy satisfies the same exact quaternion
identity and the stronger \(p=1\) rough moment bound.
Using the weaker \(p^8\) bound for both factors
still makes the product-transfer error superalgebraic
on our path. Hence the first moments and the mixed
second energy moment converge with the stated
uniform accuracy. They are uniformly bounded.
Subtract their mean product and use the Gaussian
energy covariance \(3v(f)^2/2\), proving
\eqref{eq:f15v7:mixed-energy}.
\end{proof}

\subsection{A nonzero actual retained score at growing depth}

Let \(\mathcal D_p\) be the fine bases in the support of
\(h_Q\), and take the same source shape as V4, now growing:
\[
 b_x={\bf1}_{x\in\mathcal D_p},\qquad
 |\mathcal D_p|=p^2(2p-1)^2,\qquad
 \omega_b(f)=\sum_x b_xw_x(f).
\]
Each \(w_x(f)\) is \(1/4\) for a face of the cell and
zero otherwise. Each fine face is in four such cells.
Thus \(0\le\omega_b\le1\), its support is in a box
of side \(Cp\), and
\[
           \sum_f\omega_b(f)=6|\mathcal D_p|=O(p^4).
\]
Define the actual centred source and its retained score by
\[
 S_b=\sum_xb_x(T_x-\mathbb ET_x),\qquad
 A_{j,b}=\mathbb E_\mu[S_b\mid C^{(j)}],\qquad
 \gamma_p=\sum_f\omega_b(f)(P h_Q)(f)^2 .
\]
Here the conditioning is on the entire actual compact
output, as in V2, and \(j=j_M\). This is the derivative
at zero source, not a tilt of unit strength on \(p^4\) cells.

\begin{theorem}[Retained fine-source information at genuinely growing depth]
\label{f15v7:score}
For the actual output law \(\rho_{j_M}\) on the slower path,
\begin{align}
 \Var_{\rho_{j_M}}(A_{j_M,b})
       &\ge \tfrac32\,\frac{\gamma_p^2}{c_p^2}-o(1)
        \ge \frac{3c_p^2}{32(H_p^{(2)})^2}-o(1),\notag\\
 \liminf_{M\to\infty}\Var_{\rho_{j_M}}(A_{j_M,b})
       &\ge \frac3{32}\left(\frac8{\pi^{18}}\right)^2>0 .
 \label{eq:f15v7:score}
\end{align}
Here \(H_p^{(2)}=\|h_Q\|_2^2=(2p^2+1)^2/(9p^4)\).
The statement is not a positive lower bound on a
fraction of the full, growing fine-source variance.
\end{theorem}

\begin{proof}
Sum \eqref{eq:f15v7:mixed-energy} with weights
\(\omega_b(f)\), using any \(A>4\). Since their
total weight is \(O(p^4)\), this gives the actual
covariance estimate
\[
          \operatorname{Cov}_\mu(S_b,X_Q)
                                 =\tfrac32\gamma_p+o(1).
\]
This is the place where an error uniform over the
growing fine carrier was necessary.
Theorem~\ref{f15v7:energies} also gives
\(\Var_\mu(X_Q)=\tfrac32c_p^2+o(1)\).
The lower bound on \(c_p\) makes the denominator
positive and bounded away from zero eventually.

Since \(X_Q\) is measurable with respect to the actual
output, conditional expectation preserves its covariance
with \(S_b\). Cauchy--Schwarz at each finite regulator
therefore implies
\[
 \Var_{\rho_j}(A_{j,b})
      \ge\frac{\operatorname{Cov}_\mu(S_b,X_Q)^2}
                    {\Var_\mu(X_Q)}
      =\tfrac32\gamma_p^2/c_p^2+o(1).
\]
The expansion is legitimate uniformly: \(0\le\omega_b\le1\)
and \(P^2=P\) give
\(\gamma_p\le\|P h_Q\|_2^2=c_p\le1\).

Every fine face in \(\operatorname{supp}h_Q\) has its
base cell in \(\mathcal D_p\), hence weight at least \(1/4\).
The Cauchy--Schwarz estimate in V4 is algebraic at every \(p\):
\[
 \gamma_p\ge\tfrac14\sum_{\operatorname{supp}h_Q}(P h_Q)(f)^2
              \ge\frac{\langle h_Q,P h_Q\rangle^2}
                            {4\|h_Q\|_2^2}
              =\frac{c_p^2}{4H_p^{(2)}}.
\]
Substitution proves the first line of \eqref{eq:f15v7:score}.
Finally \(H_p^{(2)}\le1\) and \(c_p\ge8/\pi^{18}\)
give the claimed strictly positive lower limit.
\end{proof}

This extends V4's fixed-\(p\) statement to \(j_M\to\infty\).
It neither commutes a conditional limit nor approximates
the entire conditional score by a coarse plaquette.
The source size \(p^2(2p-1)^2\) remains explicit.
V2's upper information bound still costs
\(k\|b\|_2^2=k|\mathcal D_p|\), which grows.
The new constant lower bound does not cancel that cost.

\subsection{What this changes upstream, and what it leaves open}

The actual nonlinear output now has controlled fixed
energy moments and nonzero retained-source information
through an unbounded number of block steps. It is no
longer necessary to rely on a bounded test alone or
on a fixed-depth source test for these two conclusions.
The depth was slowed explicitly to pay the high-frequency
and tail costs; no such higher-moment claim is made
here for V6's faster sequence.

The next missing statement is qualitatively stronger:
control, under the actual Wilson pushforward, of how
\(\mathbb E[\widetilde T_x\mid C^{(j)}]\) and the higher
conditional source coefficients depend on distant
coarse coordinates. Unconditional polynomial moments
and a single retained-information lower bound do not
give that conditional locality. Nor do they control
the sum of all generated interactions.
Further purely Gaussian moment calculations should
not be mistaken for supplying this absent estimate.

In addition \(p_M/M\to0\), the comparison remains an
ultraviolet Gaussian one, and no physical scale
calibration or interacting continuum construction
has been obtained. A positive thermal score variance
is not a positive Hamiltonian spectral gap.

\subsection{Checks, scope and exact preservation}

The finite certificate is
\begin{center}\small
\path{scripts/verify_ym_f15_v7_growing_moments.py}.
\end{center}
It checks 112 Fourier scaling identities, 12 decay-rate
calculations, 36 cutoff and product powers,
77 exact compact-quaternion energy identities, 60 joint
Gaussian energy cumulants (including singular covariance),
12 one-face gamma cumulants, and six retained-score
constants, all with rational arithmetic where applicable.
It neither simulates the Wilson law nor replaces the
analytic proofs of the uniform limits above.

The certificate has 5,554 bytes and SHA--256
\begin{center}\scriptsize\ttfamily
2259D12078F47CE3DC0C959BA0424A4C99420348B3FEF798E2CBE5708FF0D4D5.
\end{center}
The predecessor V6 has 131,325 bytes and SHA--256
\begin{center}\scriptsize\ttfamily
B2AD690EC31AF440640FD20656B38F6CD6A12471D4376BE5DF3AF374BD79E37F.
\end{center}
Reversing the title change and this terminal insertion
recovers that predecessor exactly. V7 replaces V6 as
the sole active main note. The 11 protected sources
and 49 earlier protected certificates remain unchanged.
The next companion checkpoint remains V10.

\begin{firewall}
V7 proves all fixed polynomial moment comparisons,
actual energy connected coefficients and a positive
actual retained-score lower bound on a specified slower
growing-depth path. It does not establish conditional
locality, a uniform positive-source analytic window
for the generated action, an interacting physical
continuum, or the full Yang--Mills mass gap.
The original goal remains active.
\end{firewall}
\section{V8: Exact boundary screening and the remaining conditional-locality test}
\label{f15v8:context}

V7 supplied actual growing-depth moments and retained-source
information, but neither result controls the dependence of
a conditional score on distant coarse variables. We now
move upstream to the conditional measure itself.
The results below give an exact spatial screening identity
for the literal block map, its nonlinear source-entropy
version, and a quantitative full-compact transfer bound
on growing collars. The remaining good-fibre oscillation
estimate is explicitly a hypothesis, not a proved mixing
theorem. No extra Gaussian moment calculation is used
to conceal that missing step.

The inherited source budget was checked at its original
interface in f14 V99: closed-cell site reflection,
the exact partition of the action, scalar source
dissemination, and then the Hessian at zero source.
Its proof does not assume that every off-diagonal energy
covariance is positive. The selected finite source
window from f14 V100 remains an input. In contrast,
f14 V38's measured local regression uses an admitted
uniformly convex fibre extension; it is not automatically
a theorem for the present full compact conditional law.
Neither of these inputs supplies the boundary response
needed here.

The general distinction between a local microscopic
interaction and locality after a transformation is not
specific to these notes. Examples where transformed
spin measures fail to be Gibbsian are proved by
A.\,C.\,D.~van Enter, R.~Fernandez and A.\,D.~Sokal,
\emph{Regularity Properties and Pathologies of
Position-Space Renormalization-Group Transformations},
\href{https://arxiv.org/abs/hep-lat/9210032}
{arXiv:hep-lat/9210032}. We import no theorem about
Yang--Mills from those examples. The finite diagnostic
below is proved separately and concerns source locality,
not a claim that our Wilson pushforward is non-Gibbsian.

\subsection{An exact fine shield for the literal finite-range map}

Use positively oriented fine links as independent compact
variables; reversed links are their inverses. Fix a dyadic
block scale \(p=2^j\). Each completed coarse link depends
only on fine links within \(\ell^\infty\)-distance \(4p\)
of its fine anchor. This conservative bound follows from
the same primitive stencil count used in V4: a one-step
edge base offset is at most three, and composition gives
\(3(1+2+\cdots+2^{j-1})=3(p-1)\); endpoints are included
by the bound \(4p\). Gauge covariance does not require
introducing any nonlocal gauge fixing to define this map.

Coarse boxes use torus distance and non-winding representatives.
For an integer collar radius \(r\ge10\), let \(N\) denote
all coarse output links whose base is in \([-r,r]^4\).
Let \(F\) denote all remaining coarse output links.
Thus \((N,F)\) is the entire output, not a truncated
model of it. Work in a translated non-winding patch;
put
\[
 \mathcal K=[-(r+5)p,(r+5)p]^4,\qquad
       L=2p(r+5),\qquad M\ge4L.
\]
Let \(I\) consist of fine links with both endpoints in
the inner box \([-(r-6)p,(r-6)p]^4\).
Let \(B\) consist of the other links with both endpoints
in \(\mathcal K\), and let \(E\) contain every remaining
fine link. The symbols also denote the corresponding
random link configurations.

Every near output depends only on \((I,B)\).
No far output uses an \(I\)-link: a far base has a
coordinate of absolute value at least \(r+1\) in coarse
units, and its fine carrier can extend toward the
inner box by at most \(4p\). Hence
\[
                   N=N(I,B),\qquad F=F(B,E).
\]
Also no plaquette contains both an \(I\)-link and
an \(E\)-link. A plaquette incident to \(I\) lies
within one lattice step of its endpoints, far inside
the outer box. The shield \(B\) is a genuine fine
annulus of thickness \(O(p)\) and has \(O(p^4(r+1)^3)\)
links. Its variables are not identified with coarse
boundary links.

Take a real fine-cell source \(b\) whose observable
\[
                  S_b=\sum_xb_x(T_x-\mathbb ET_x)
\]
depends only on \((I,B)\). Sources supported in a fixed
multiple of the central \(p\)-box satisfy this requirement
for all sufficiently large \(r\). The centering term is
a scalar and does not change the support.

\begin{proposition}[Exact conditional screening under the compact law]
\label{f15v8:screening}
For the full Wilson measure, \(I\) and \(E\) are
conditionally independent given \(B\). Moreover, for
every integrable \(Z=Z(I,B)\),
\begin{equation}
 \mathbb E[Z\mid B,N,F]=\mathbb E[Z\mid B,N],\qquad
 \mathbb E[Z\mid N,F]
             =\mathbb E\!\left[\mathbb E[Z\mid B,N]\mid N,F\right].
 \label{eq:f15v8:screen}
\end{equation}
No restriction to a chart or to a positive-probability
point value of \(N\) is part of this assertion.
\end{proposition}

\begin{proof}
Split the plaquette action into terms meeting \(I\),
terms meeting \(E\), and terms using only \(B\).
Since no term meets both \(I\) and \(E\), it has the form
\[
            S^{\rm W}=S_I(I,B)+S_E(E,B)+S_B(B).
\]
Conditional on \(B=b_0\), the density with respect
to product Haar measure on \(I,E\) factors as
\[
 \frac{e^{-\beta S_I(I,b_0)}}{Z_I(b_0)}
                 \frac{e^{-\beta S_E(E,b_0)}}{Z_E(b_0)} .
\]
All compact finite-volume normalizers are finite and
strictly positive. Thus the asserted independence holds
for a regular conditional version.

Given \(B\), the random variables \((Z,N)\) depend only
on \(I\), whereas \(F\) depends only on \(E\).
Their conditional joint law therefore factors. Disintegrating
the first factor further by \(N\) proves the first
identity in \eqref{eq:f15v8:screen}; taking conditional
expectations proves the second. Products of compact
groups and their measurable output spaces are standard
Borel spaces, so these disintegrations exist. No
Lebesgue density for the coarse output or inverse of
the blocking map is required.
\end{proof}

\subsection{The exact source-locality error and its boundary norm}

For \(S=S_b\in L^2(\mu)\) define
\[
 A=\mathbb E[S\mid N,F],\qquad
 a=\mathbb E[S\mid N],\qquad
 m=\mathbb E[S\mid B,N].
\]
Thus \(A\) is V2's actual full-output score and \(a\)
is its best mean-square approximation by a function
of the specified local output. All three have mean zero.
Write
\[
     \mathcal L_{b,r}^2=\mathbb E|A-a|^2,\qquad
     \mathcal R_{b,r}^2=\mathbb E|m-a|^2 .
\]

\begin{theorem}[Boundary-channel identity, including the discarded residual]
\label{f15v8:identity}
At every finite regulator,
\begin{equation}
 \mathcal R_{b,r}^2
       =\mathcal L_{b,r}^2+\mathbb E|m-A|^2,\qquad
                         \mathcal L_{b,r}\le\mathcal R_{b,r}.
 \label{eq:f15v8:pythagoras}
\end{equation}
Equivalently, the actual locality error is
\begin{equation}
 \mathcal L_{b,r}
   =\sup\left\{|\mathbb E(SH)|:
       H=H(N,F),\ \mathbb E[H\mid N]=0,\ \|H\|_2\le1\right\}.
 \label{eq:f15v8:dual}
\end{equation}
Neither formula asserts that either norm decreases
at a particular rate as the collar grows.
\end{theorem}

\begin{proof}
By \eqref{eq:f15v8:screen}, \(A=\mathbb E[m\mid N,F]\).
Also \(\mathbb E[m\mid N]=a\).
The difference \(m-A\) is orthogonal to every square
integrable function of \((N,F)\), in particular to
\(A-a\). Expanding \(m-a=(m-A)+(A-a)\) proves
\eqref{eq:f15v8:pythagoras}.

For the tests in \eqref{eq:f15v8:dual}, conditional
expectation gives \(\mathbb E(SH)=\mathbb E((A-a)H)\).
Cauchy--Schwarz gives the upper bound; if \(A-a\ne0\),
take \(H=(A-a)/\|A-a\|_2\). The zero case is immediate.
This also proves the best-approximation assertion.
\end{proof}

The missing estimate is now about a specified conditional
mean under the actual fine law. It is not an unconditional
two-point covariance. Even exact information about all
fixed-degree local moments would not evaluate the supremum
over the tests in \eqref{eq:f15v8:dual}.
The annulus \(B\) itself can be informative about far
outputs; no strict contraction for its entire observation
channel is presumed.

\subsection{Finite sources and a nonlinear conditional-entropy bound}

For real \(t\), let
\[
 d\mu_t=e^{tS-K(t)}d\mu,\qquad K(t)=\log\mathbb Ee^{tS}.
\]
At each finite regulator \(S\) is bounded, so every
quantity below exists for all real \(t\).
For \(H=N,\ (N,F),\ (B,N)\), set
\[
  \ell_t^H=\mathbb E[e^{tS-K(t)}\mid H],\qquad
  f_t^H=\log\mathbb E[e^{tS}\mid H].
\]
These are likelihoods and conditional cumulant functions,
not newly postulated coarse interactions. Screening implies
\begin{equation}
        f_t^{N,F}
             =\log\mathbb E[e^{f_t^{B,N}}\mid N,F].
 \label{eq:f15v8:nonlinear-tower}
\end{equation}
In particular, the logarithm cannot be moved through
this conditional expectation. At zero source,
\begin{equation}
 \Var(S\mid N,F)
    =\mathbb E[\Var(S\mid B,N)\mid N,F]
                            +\Var(m\mid N,F).
 \label{eq:f15v8:second-return}
\end{equation}
This is the conditional variance formula following from
screening; the extra boundary-mean variance is part of
the actual second source coefficient.

Let \(\rho_t\) be the law of \((N,F)\) under \(\mu_t\),
and let \(\rho=\rho_0\). Define a source-local comparison by
\[
       d\widehat\rho_t=\ell_t^N(N)\,d\rho .
\]
It changes the near marginal exactly as the actual source
does, but retains the baseline conditional law of \(F\)
given \(N\). This is a local \emph{source increment};
it does not assert that the baseline law \(\rho\) itself
has a local density or is a nearest-neighbor Wilson law.

\begin{theorem}[Nonlinear boundary information dominates source nonlocality]
\label{f15v8:entropy}
Writing \(\mu_t^{B,N}\) and \(\mu_t^N\) for the indicated
marginals, one has
\begin{align}
 D(\rho_t\Vert\widehat\rho_t)
   &=D(\rho_t\Vert\rho)-D(\mu_t^N\Vert\mu^N)\notag\\
   &\le D(\mu_t^{B,N}\Vert\mu^{B,N})
                                      -D(\mu_t^N\Vert\mu^N)\notag\\
   &=\int D(\mu_t(dB\mid N)\Vert\mu(dB\mid N))\,d\mu_t^N(N).
 \label{eq:f15v8:entropy}
\end{align}
For the fixed finite regulator, their second-order
expansions at zero are respectively
\begin{equation}
 D(\rho_t\Vert\widehat\rho_t)
       =\tfrac12t^2\mathcal L_{b,r}^2+O(t^3),\qquad
 \text{right side of \eqref{eq:f15v8:entropy}}
       =\tfrac12t^2\mathcal R_{b,r}^2+O(t^3).
 \label{eq:f15v8:entropy-curvature}
\end{equation}
No regulator-uniform bound on these third-order remainders
is claimed.
\end{theorem}

\begin{proof}
Conditional on \(B\), the source \(e^{tS(I,B)}\)
changes the interior factor and the marginal of \(B\),
but not the exterior conditional kernel. Conditioning
additionally on \(N=N(I,B)\) still leaves that exterior
kernel unchanged. Thus the conditional law of \(F\)
given \((B,N)\) is the same kernel under \(\mu_t\) and
\(\mu\). This is an actual source-independent boundary
channel, not an approximation to a conditional density.

For each \(N\), apply the relative-entropy data-processing
inequality to the two conditional laws of \(B\) through
this common kernel. Indeed the output likelihood is the
conditional average of the input likelihood; Jensen for
the convex function \(z\log z\) proves this inequality.
Integrating with respect to
\(\mu_t^N\) bounds the conditional entropy of \(F\)
by the last line of \eqref{eq:f15v8:entropy}.
The chain rule identifies that entropy of \(F\) with
the first line, or equivalently with
\(D(\rho_t\Vert\widehat\rho_t)\).
Absolute continuity is harmless: boundedness of \(S\)
gives positive upper and lower likelihood bounds at
the fixed regulator. The chain rule on \((B,N)\)
gives the remaining equality.

For clarity, the local Taylor calculation follows
from \(\ell_t^H=1+t\mathbb E[S\mid H]+O(t^2)\)
uniformly at the fixed regulator. Expansion of
\(\ell\log\ell\), using \(\int\ell=1\), gives
\[
 D(\mu_t^H\Vert\mu^H)
           =\tfrac12t^2\|\mathbb E[S\mid H]\|_2^2+O(t^3).
\]
Subtract the \(N\)-term from the full or boundary term.
The nesting of the conditional expectations gives the
two squared differences in \eqref{eq:f15v8:entropy-curvature}.
\end{proof}

Consequently, a vanishing boundary entropy in
\eqref{eq:f15v8:entropy} would imply total-variation
approximation by the local source increment, by
Pinsker's inequality. The theorem has not bounded that
boundary entropy by a decaying function of \(r\).
Nor does a bound on its curvature at zero alone
control a fixed nonzero source.

\subsection{A fourth-moment input norm adapted to conditional errors}

At the selected couplings, denote the inherited source
window by \(R_{\rm src}\in(0,1]\), to distinguish it
from any spatial radius. Its estimate is
\[
 \log\mathbb E e^{S_b}\le k\sum_xJ(b_x),\qquad
 J(z)=-\log(1-z)-z,\qquad \|b\|_\infty\le R_{\rm src}.
\]
The constants satisfy \(k\to9/2\) and \(R_{\rm src}\to1\).

\begin{proposition}[An \(\ell^2\), rather than \(\ell^1\), fourth-moment cost]
\label{f15v8:fourth}
For every real source \(b\), set
\[
       a_b=\sqrt{k}\|b\|_2+\|b\|_\infty/R_{\rm src},
                   \qquad C_4=(384e^{1/4})^{1/4}.
\]
Then
\begin{equation}
                            \|S_b\|_4\le C_4 a_b .
 \label{eq:f15v8:fourth}
\end{equation}
There is no smallness assumption on \(b\) in this
moment assertion.
\end{proposition}

\begin{proof}
For \(|z|<1\), the power series for \(J\) gives
\(J(z)\le z^2/(2(1-|z|))\).
Put \(v=k\|b\|_2^2\) and \(c=\|b\|_\infty/R_{\rm src}\).
Applying the source budget to \(\pm tb\) gives
\[
 \log\mathbb E e^{\pm tS_b}
                \le\frac{vt^2}{2(1-c|t|)},\qquad c|t|<1 .
\]
Here \(c\ge\|b\|_\infty\), so the chosen denominator
is no larger than the one from the \(J\)-bound,
and the source remains in its allowed window.
If \(b=0\) there is nothing to prove. Otherwise take
\(t=(2a_b)^{-1}\). Then \(ct\le1/2\) and the right
side is at most \(1/4\). Averaging the two exponential
bounds and using the fourth term in the nonnegative
series for \(\cosh(tS_b)\) gives
\[
          \mathbb ES_b^4
             \le24t^{-4}\mathbb E\cosh(tS_b)
             \le384e^{1/4}a_b^4 .
\]
This proves the claim. Scaling down the test source,
not assuming a unit-strength tilt of arbitrary \(b\),
is what permits the unrestricted moment statement.
\end{proof}

For the indicator source of V7, \(\|b\|_2=O(p^2)\)
and \(\|b\|_\infty=1\), so \(a_b=O(p^2)\).
This improves the \(O(p^4)\) cost obtained from the
older \(\ell^1\) fourth-moment estimate.

\subsection{Annealed transfer of a good-fibre boundary estimate}

Let \(G=G(I,B)\) be any local fine good event, and put
\[
     \delta=\mu(G^c),\qquad
     h(B,N)=\mathbb E[S_b{\bf1}_G\mid B,N].
\]
The truncation is deliberately unnormalized. In particular
it is not divided by a possibly small conditional
probability of \(G\).
Let \(\Gamma\) be a measurable set of pairs \((B,N)\)
with \(\mu^{B,N}(\Gamma^c)\le\zeta\). Suppose there is
a nonnegative square-integrable \(\omega(N)\) such that
\begin{equation}
 |h(B,N)-h(B',N)|\le\omega(N)
 \label{eq:f15v8:oscillation-hypothesis}
\end{equation}
for conditional-product-almost every pair of boundary
values \((B,B')\) given \(N\) for which both \((B,N)\)
and \((B',N)\) belong to \(\Gamma\).
This is the precise good-fibre oscillation hypothesis.

\begin{theorem}[Full-compact score-locality transfer]
\label{f15v8:transfer}
Under \eqref{eq:f15v8:oscillation-hypothesis},
\begin{equation}
 \mathcal L_{b,r}
   \le \left(\frac{\mathbb E\omega(N)^2}{2}
                    +2\sqrt2\,C_4^2a_b^2\sqrt\zeta\right)^{1/2}
                    +C_4a_b\,\delta^{1/4}.
 \label{eq:f15v8:transfer}
\end{equation}
The constants have no hidden factor equal to the
number of observed coarse coordinates.
The estimate does not assert that its oscillation
hypothesis holds for the Wilson conditional fibres.
\end{theorem}

\begin{proof}
Conditional Jensen gives \(\|h\|_4\le\|S_b\|_4\).
On an auxiliary probability space, sample \(N\)
from its actual law and then independently sample
\(B,B'\) from its conditional law given \(N\).
Let \(h'=h(B',N)\). The conditional variance identity is
\[
 \|h-\mathbb E[h\mid N]\|_2^2
                              =\tfrac12\mathbb E|h-h'|^2 .
\]
On the event that both boundary pairs are in \(\Gamma\),
the right integrand is bounded by \(\omega(N)^2\).
The complement has probability at most \(2\zeta\).
By Cauchy--Schwarz and Minkowski,
\[
 \tfrac12\mathbb E[|h-h'|^2{\bf1}_{\rm complement}]
       \le\tfrac12\|h-h'\|_4^2(2\zeta)^{1/2}
       \le2\sqrt2\,\|S_b\|_4^2\sqrt\zeta .
\]
The two copies have the same marginal law; unconditional
independence is neither needed nor claimed.

Next \(m-h=\mathbb E[S_b{\bf1}_{G^c}\mid B,N]\), so
conditional Jensen and H\"older give
\[
                         \|m-h\|_2\le\|S_b\|_4\delta^{1/4}.
\]
The orthogonal projection \(I-\mathbb E[\,\cdot\,\mid N]\)
is a contraction on \(L^2\). Apply it to \(m=h+(m-h)\),
use \(\mathcal L_{b,r}\le\mathcal R_{b,r}\), and then
use \eqref{eq:f15v8:fourth}. This proves
\eqref{eq:f15v8:transfer} without a conditional
normalization denominator or a pointwise posterior bound.
\end{proof}

\subsection{The compact exceptional term on a growing collar}

Here one part of the preceding bound can be controlled
for the actual Wilson law, independently of any claimed
conditional convexity. Write \(C_{\rm th}\) for the
inherited thermal constant evaluated at one fixed bound
on \(\log\beta/M\); it is independent of the present
box side \(L\). Translate \(\mathcal K\) to
\(\{0,\ldots,L\}^4\). Use its increasing-coordinate root
paths and let \(R_{\mathcal K}\) be the maximum angle
of the resulting rooted fine links.
There are \(4L(L+1)^3\le32L^4\) such links. The
plaquette filling of each rooted link has at most
\(3L\) faces, by commuting its last coordinate step
past the later steps of the root path.

\begin{proposition}[Actual collar-chart exceptional probability]
\label{f15v8:collar-tail}
For every fixed \(\rho_*>0\), the local fine event
\[
                  G=\{L R_{\mathcal K}\le\rho_*\}
\]
satisfies
\begin{equation}
 \delta=\mu(G^c)
    \le\min\left\{1,\ 32C_{\rm th} L^4
               \exp\!\left[-\frac{\beta\rho_*^2}
                                      {9\pi^2L^4}\right]\right\}.
 \label{eq:f15v8:collar-tail}
\end{equation}
The event is measurable with respect to \((I,B)\).
The estimate concerns the full compact law, without
conditioning on \(N\) or restricting the integration
to a single fibre.
\end{proposition}

\begin{proof}
For a fine plaquette angle \(\vartheta_f\), the inherited
thermal estimate implies
\(\mathbb E e^{\beta\vartheta_f^2/\pi^2}\le C_{\rm th}\).
The group angle is invariant under conjugation and
inversion and is subadditive under products.
For a rooted link with a filling of \(n\le3L\) faces,
Cauchy--Schwarz followed by convexity of the exponential
therefore gives
\[
      \mathbb E\exp\!\left[
                  \frac{\beta\vartheta(k_e)^2}{9\pi^2L^2}
                    \right]\le C_{\rm th} .
\]
One can first use denominator \(\pi^2n^2\):
the exponent is bounded by the average of the \(n\)
fine-plaquette exponents. Increasing the denominator
only weakens the estimate. Repeated plaquettes, if
present, cause no independence assumption; a tree link
with \(n=0\) is exactly the identity.
Markov's inequality and the union bound over at most
\(32L^4\) links, with threshold \(\rho_*/L\), prove
\eqref{eq:f15v8:collar-tail}.
All root paths and their fillings are inside the box.
\end{proof}

In particular keep V7's slow path
\[
 u=\log((\eta_M+\beta_M^{-1})^{-1}),\qquad
 p=\text{largest dyadic integer at most }u,\qquad
 r=\lfloor u\rfloor .
\]
Then both block depth and coarse collar radius tend
to infinity, \(L=2p(r+5)=O(u^2)\), and \(\beta\ge e^u\).
For all sufficiently large \(M\),
\begin{equation}
                  \delta\le C u^8 e^{-c e^u/u^8}.
 \label{eq:f15v8:collar-path}
\end{equation}
The torus embedding holds because \(\beta\le\log M\).
For a bounded source on \(O(p^4)\) fine cells,
\(a_b=O(p^2)\), so
\begin{equation}
                    p^A a_b\delta^{1/4}\longrightarrow0
                         \qquad\hbox{for every fixed }A\ge0 .
 \label{eq:f15v8:compact-negligible}
\end{equation}
This removes the explicit fine compact-complement
contribution to the transfer bound on this growing
collar. It does not bound \(\omega\) or \(\zeta\).
Small rooted fields do not by themselves prove
strong convexity of the nonlinear pinned fibre,
its connectedness, or stability of its normalized
boundary response. All those distinctions remain.
Also the growing number of near observations has
not been inserted into V7's fixed-family Gaussian
comparison: the present argument is exact screening
plus a full-measure geometric tail.

\subsection{Why microscopic mixing alone cannot finish the argument}

The following elementary model separates the missing
condition from unconditional fine or coarse mixing.
It is not a model of Yang--Mills and is not a
counterexample to locality for our particular primitive.

\begin{proposition}[A local observation with nonlocal source response]
\label{f15v8:parity}
Let \(U_0,\ldots,U_n\) be independent uniform signs,
and define local observed variables
\[
                 C_0=U_0,\qquad C_i=U_{i-1}U_i
                                      \quad(1\le i\le n).
\]
The baseline coarse variables are themselves independent
uniform signs. For \(S=U_n\), let the near observation
be \(N=(C_{n-r+1},\ldots,C_n)\), where \(1\le r<n\),
and let \(F\) contain the other \(C_i\)'s. Then
\begin{equation}
       \mathbb E[S\mid N]=0,\qquad
       \mathbb E[S\mid N,F]=S,\qquad
                    \mathcal L_{S,r}=1 .
 \label{eq:f15v8:parity-error}
\end{equation}
Under the source \(e^{tS}/\cosh t\), every proper
coarse marginal remains unchanged, whereas
\begin{equation}
 D(\rho_t\Vert\widehat\rho_t)
                          =t\tanh t-\log\cosh t>0
                                    \quad(t\ne0).
 \label{eq:f15v8:parity-entropy}
\end{equation}
One can add independent signs to the fine input and
omit them from the output. The observation then loses
information, with all the stated conclusions unchanged.
\end{proposition}

\begin{proof}
The triangular transformation from \(U\) to \(C\)
is a bijection of sign cubes, with inverse
\[
                         U_i=\prod_{k=0}^i C_k .
\]
It therefore takes the uniform law to the uniform
product law. Given only the last \(r\) parities, the
unobserved boundary sign \(B=U_{n-r}\) is still uniform.
Since
\[
                         S=B\prod_{k=n-r+1}^n C_k ,
\]
the local conditional mean vanishes. The full output
reconstructs \(S\), proving \eqref{eq:f15v8:parity-error}.
The boundary conditional mean equals the last displayed
product, so the fine shield transmits all the missing
information despite the independent microscopic law.

In coarse coordinates the exact tilted likelihood is
\[
      \frac{\exp(t\prod_{k=0}^n C_k)}{\cosh t}
             =1+\tanh t\prod_{k=0}^n C_k .
\]
Summing over any missing sign kills its second term.
Thus every proper marginal, including \(N\), is unchanged,
and \(\widehat\rho_t=\rho\). The tilted expectation of
\(\prod C_k=S\) is \(\tanh t\), giving
\eqref{eq:f15v8:parity-entropy}. Its derivative for \(t>0\)
is \(t\,\operatorname{sech}^2t>0\), and the expression
is even and zero at zero. This proves strict positivity.
Unused independent signs integrate to one and can simply
be omitted by each local observation rule.
\end{proof}

Here even the scalar source bound
\(\log\mathbb Ee^{tS}=\log\cosh t\le t^2/2\)
is uniform in \(n\); integrate
\((\log\cosh t)''\le1\) with zero value and derivative
at zero to see this. Any fixed-degree proper-marginal
moment data miss the generated source interaction
once the string is longer than that degree.
Thus fine mixing, baseline coarse mixing, local forward
carriers, dimension reduction, and an input source budget
do not jointly imply the estimate still needed in
\eqref{eq:f15v8:oscillation-hypothesis}.
The example has a boundary anchor and sign variables;
no translation-covariant compact-gauge analogue or
failure theorem for our YM map is asserted.

\subsection{The next estimate is now explicit}

For the actual literal map, the remaining spatial task
is to bound the oscillation of
\[
             h(B,N)=\mathbb E_\mu[S_b{\bf1}_G\mid B,N]
\]
at fixed near output, outside a quantitatively small
set \(\Gamma^c\). It must be the conditional mean
under the actual compact disintegration, including
the source normalizers and all branches of the
nonlinear fibre. An admitted convex extension from
the shell notes may be used only after its relation
to this conditional law has been proved.

For example, estimates making
\(\|\omega\|_2\to0\) and \(a_b\zeta^{1/4}\to0\)
would, together with the proved compact term, imply
\(\mathcal L_{b,r}\to0\) along the stated collar path.
This is a precisely stated sufficient condition,
not an achieved locality theorem. To obtain a local
effective interaction class, one would need suitable
rates, higher conditional coefficients and nonzero
source control as well. The finite-source entropy
identity records rather than discards that further cost.

No physical correlation length or Hamiltonian gap is
inferred from these spatial identities. The goal remains
the interacting four-dimensional continuum and its
positive physical mass gap, which are still unproved.

\subsection{Checks and preservation}

The diagnostic is
\begin{center}\small
\path{scripts/verify_ym_f15_v8_boundary_screening.py}.
\end{center}
It checks three projection identities including a
strict boundary residual, 40 source-independent finite
boundary channels, and five conditional entropy
inequalities. The probability and projection calculations
use exact rational arithmetic; only the logarithms in
the five entropy checks use floating point.
The parity diagnostic verifies 1,016 reconstruction
states, 28 nondecaying locality errors and 126 unchanged
tilted proper marginals. It also checks 64 box edge
counts, 53,232 ordered-path filling counts, and three
constant/scaling identities.
These are finite diagnostics, not estimates of the
actual Wilson good-fibre oscillation.

The certificate has 6,546 bytes and SHA--256
\begin{center}\scriptsize\ttfamily
9AD5DCF50A3E7DEBE9A98AFAE8583CCFCAC014136A4BE81C5E6CC7E4CFA1C179.
\end{center}
The predecessor V7 has 156,512 bytes and SHA--256
\begin{center}\scriptsize\ttfamily
E51F2B18A35A8BD0A3040D5F6E1B352CDDA5A734E3AB68A70CDAF1FDE0599D80.
\end{center}
Reversing this terminal insertion and the title change
recovers that predecessor exactly. V8 replaces V7
as the sole active main note; the 11 protected sources
and 50 earlier protected certificates are preserved.
The next companion checkpoint is V10.

\begin{firewall}
V8 proves exact screening and entropy reductions for
the actual compact pushforward, an improved input
fourth-moment norm, and a negligible explicit compact
exceptional term on growing collars. The good-fibre
boundary oscillation remains unproved. No conditional
locality theorem, closed interacting RG class, continuum
construction, or physical YM mass gap is claimed.
\end{firewall}
\section{V9: Full compact fibre coordinates with the pinned constraints confined to a collar}
\label{f15v9:context}

V8's unresolved quantity is a boundary response at fixed
near coarse output, not another marginal Gaussian moment.
Trying to differentiate that conditional mean directly
in physical fine-boundary variables encounters compatibility
conditions and possibly dependent coarse rows.
This section removes that particular geometric obstruction
from the deep interior: use the completed-pivot coordinates
before imposing the fine-boundary condition. Their inverse
is local, so the remaining compatibility constraints involve
only collar variables. The bulk fibre remains a fixed full
product of compact groups, with its exact measured potential.

The one-step compact coordinates and pivot Jacobian are
already present in f14 V48; they are not new constructions
here. We explicitly iterate their full measure, prove the
support bound needed at a physical pinned boundary, and
disintegrate only the collar constraint. This is different
from f14 V43's local classical centering and from f14 V38's
admitted convex extension. No convexity, conditional mixing
rate, or small-field branch coverage is assumed in the
coordinate and disintegration theorems below.

\subsection{A root-coordinate audit before pinning the output}

The exact one-step derivative from f14 V41 is
\begin{equation}
 (\mathsf L a)_{x,\mu}
   =\frac1{16}\sum_{s\in\{0,1\}^4}
       (a_{2x+s,\mu}+a_{2x+s+e_\mu,\mu})
                         +\bar t_x(a)-\bar t_{x+e_\mu}(a),
 \label{eq:f15v9:root-derivative}
\end{equation}
where \(\bar t_x\) is the average circulation along the
sixteen ordered cell-root paths. The last terms vanish
from a curl, not from the rooted output link itself.
For a fine vertex potential \(\lambda\),
\[
             \mathsf L(d_f\lambda)_{x,\mu}
                         =\lambda(2x+2e_\mu)-\lambda(2x).
\]
Indeed each rooted path telescopes between these same
two roots. An unrooted line average instead differentiates
the cell-averaged vertex potential.

Accordingly, V4's displayed line-average link operator
is the operator used to obtain the curl identity; it
must not be substituted for the full derivative of the
literal rooted link map when imposing fixed-gauge link
constraints. As a full rooted-link derivative, that
earlier wording is incomplete and is qualified here.
The curvature comparisons in V4--V7 use
the curl identity, where the root gradient cancels, so
this distinction does not alter those arguments.
We use the nonlinear completed map itself below, rather
than a constraint imposed only on its curl.

\subsection{The inherited global one-step change of compact variables}

In each cell, change the fifteen tree links to their
root-path group variables \(\gamma_x\in\SU(2)^{15}\).
The remaining internal groups are
\(X_x\in\SU(2)^{17}\). For each coarse edge \(e=(x,\mu)\)
there are seven free boundary groups \(Y_e\in\SU(2)^7\)
and one root pivot \(B_e\in\SU(2)\).
These are group variables, not scalar coordinates.
The ordered tree change and the endpoint dressings
preserve normalized product Haar measure.

On this tree slice, the two root products are \(B_e\).
The other fourteen products have the inherited form
\[
       P_{e,a,\ell}
            =L_{e,a,\ell}(X)Y_{e,a}R_{e,a,\ell}(X),
                   \qquad 1\le a\le7,\quad \ell=1,2.
\]
They use only the internal groups in the endpoint cells,
contain no pivot, and use no \(Y\)-family belonging to
another coarse edge. At fixed \(X,Y\), the completed map is
\[
 \Phi_e(B)=B\exp\left\{\frac1{16}
                   \sum_{a,\ell}g_\rho(B^{-1}P_{e,a,\ell})\right\},
                     \qquad \rho=1/64 .
\]
The inherited inverse is a global smooth diffeomorphism
in \(B\), with positive inverse Haar Jacobian \(j_e(X,Y_e,C_e)\).
It is not merely a chart inverse near the identity.
The cutoff here is the inherited smooth nonincreasing
profile; that property enters the pivot Hessian bound.

For completeness, its nonsingularity is controlled on the
whole compact input by the same mechanism used in f14 V48.
The solved pivot is within distance \(7\rho/8\) of \(C_e\).
The pivot potential has Hessian at least
\[
              \left(1-\frac{(7\rho/8)^2}{4}-\frac78\right)I>0
\]
at such a solution. The mixed derivative of the short
geodesic logarithm is invertible, so \(D_B\Phi_e\) is
invertible everywhere. The compact local diffeomorphism
\(\Phi_e:\SU(2)\to\SU(2)\) is homotopic to the identity
by multiplying its exponent by a parameter in \([0,1]\).
Its covering degree is therefore one, giving the global
inverse. These are the same fourteen-source and
working-radius bounds, not a new cutoff or a Haar reset.

Consequently the exact one-step reference-measure formula is
\begin{equation}
              dU=\left(\prod_e j_e(X,Y_e,C_e)\right)
                              dC\,dX\,dY\,d\gamma .
 \label{eq:f15v9:one-step-measure}
\end{equation}
All measures on the right are normalized product Haar.
The tree variables are retained here. Dropping them while
holding a physical fine boundary fixed would in general
change that condition.

\subsection{Full-history coordinates and a local inverse}

Let \(p=2^j\), let \(2^j\) divide \(M\), and retain the
earlier lower bound on the final torus side.
At step \(i\), let
\[
     n_i=(M/2^i)^4,\qquad
     \zeta_i=(X_i,Y_i,\gamma_i).
\]
There are \(17+28+15=60\) history groups per cell at
this step. Denote their full collection by \(\zeta\),
and the final output by \(C=C^{(j)}\).

\begin{theorem}[Exact compact history coordinates and measured density]
\label{f15v9:history}
There is a global smooth bijection
\[
             U=\mathcal U_j(C,\zeta)
\]
from the product of the final coarse-link groups and
all history groups onto the original fine-link group
space. Its forward \(C\)-coordinate is exactly the
literal iterated block output. In these coordinates
the full Wilson law has density proportional to
\begin{align}
       e^{-\mathcal V_j(C,\zeta)}\,dC\,dH(\zeta),\qquad
 \mathcal V_j
       =\beta S^{\rm W}(\mathcal U_j(C,\zeta))
          -\sum_{i=1}^j\sum_e
                   \log j_{i,e}(X_i,Y_{i,e},C^{(i)}_e).
 \label{eq:f15v9:history-density}
\end{align}
Intermediate \(C^{(i)}\)'s on the right are the reconstructed
ones. No Jacobian or earlier generated factor is set to one.

Give a step-\(i\) cell or edge group its physical fine
anchor \(2^i x\), and give a final coarse edge the anchor
\(px\). Each reconstructed fine link depends only on
coordinates with anchors within \(4p\) of that link's
fine base. Each individual term in
\(\mathcal V_j\) has coordinate-support diameter at most
\(16p\), in the periodic \(\ell^\infty\) metric.
These are support estimates, not uniform Hessian or
interaction-strength estimates.
\end{theorem}

\begin{proof}
Apply the inverse of \eqref{eq:f15v9:one-step-measure}
successively from level \(j\) down to level zero.
At each step the supplied \(C^{(i)}\), free \(X_i,Y_i\),
and tree groups \(\gamma_i\) determine exactly one
\(C^{(i-1)}\). Reversing the one-step changes recovers
all those variables, proving the global bijection.
The group count is consistent:
\[
                  4n_j+60\sum_{i=1}^j n_i=4M^4.
\]
Iterating the exact Haar change of variables proves
\eqref{eq:f15v9:history-density}. In particular the
Jacobians from lower levels are evaluated on the
reconstructed higher-level fields; they are not omitted
because a later change of variables has been made.

To check support, a one-step inverse uses only the
same cell, the two endpoint cells of a pivot edge,
and their tree groups. A physical fine link can
therefore depend only on inputs whose parent anchors
are within four fine lattice units of its base.
For example, an internal \(X_x\) affects only its
cell links and pivots on coarse edges incident to
that cell; a boundary \(Y_e\) affects its own family
and its own pivot. Undoing a tree dressing uses
only the two endpoint cells. There is no dependence
on a neighboring pivot through a pivot equation.

Iterating this primitive bound gives the radius
\[
                         4(1+2+\cdots+2^{j-1})
                                        =4(p-1)<4p.
\]
The same bound controls a coordinate's inverse
influence on fine links, by reading the dependency
relation in the other direction.
A Wilson plaquette uses four such fine links whose
bases differ by at most one, so its inputs fit
within radius \(5p\) of its base.
An individual \(\log j_{i,e}\) first uses \(X_i\)
in two neighboring cells, \(Y_{i,e}\), and \(C^{(i)}_e\).
The latter is reconstructed from higher levels.
The same geometric sum, now starting at level \(i\),
puts all its inputs within radius \(5p\) of the
physical anchor of \(e\). Diameter \(16p\) is a
conservative common bound.
\end{proof}

The theorem does not say that a final coarse cell
contains a bounded number of history groups as
\(p\to\infty\): that number is of order \(p^4\).
Nor does the support bound control the size of
derivatives after all the compositions.
The \(\gamma_i\)'s at higher levels are not discarded
by an unproved independence assertion.

\subsection{The actual physical boundary constraint depends only on collar coordinates}

Use V8's near output \(N\), far output \(F\), fine
boundary \(B\), and outer box \(\mathcal K\).
Enlarge the fixed lower bound on \(r\) to \(r\ge128\);
this has no effect on a sequence with \(r\to\infty\).
As before \(M\ge4L\), where \(L=2p(r+5)\), so the
following boxes are unambiguous non-winding patches.
Partition the coordinates remaining after \(N\) is
fixed as follows:

\begin{itemize}
\item \(z\): history groups with anchors in
      \([-(r-40)p,(r-40)p]^4\);
\item \(\eta\): all other history groups and far
      coarse groups whose anchors have distance at
      most \((r+40)p\) from the origin;
\item \(w\): all remaining history and far coarse groups.
\end{itemize}

Thus \(z\) contains no coarse output group: the
coarse groups in its region are already part of \(N\).
Let \(H_z,H_\eta,H_w\) be their product Haar measures.
For the indicator source of V7, or any source with
fine link support in the same central fixed-size
\(p\)-patch, the inverse-support estimate gives
\[
                    S_b=S_b(N,z).
\]
For this source the needed input anchors lie within
radius \(8p\), well inside the \(z\)-box.

\begin{proposition}[Product bulk fibres despite possibly singular collar compatibility]
\label{f15v9:collar-geometry}
In the full-history coordinates,
\begin{equation}
           B=\mathfrak b_N(\eta),\qquad
 \mathcal V_j(N,z,\eta,w)
                =V_{\rm in}(N,z,\eta)+V_{\rm out}(N,\eta,w).
 \label{eq:f15v9:split}
\end{equation}
Consequently the fibre at fixed physical boundary
\(B_0\) and near output \(N_0\) has the exact coordinate form
\begin{equation}
 \SU(2)^{\,\#z}\ \times\
       \{(\eta,w):\mathfrak b_{N_0}(\eta)=B_0\}.
 \label{eq:f15v9:fibre-product}
\end{equation}
There is no constant-rank or connectedness assertion
for the second factor. The first factor is an
unrestricted product of compact groups.
\end{proposition}

\begin{proof}
Every \(B\)-link belongs to V8's fine annulus between
the inner box of radius \((r-6)p\) and the outer
box of radius \((r+5)p\). Accounting for an endpoint
and the inverse radius \(4p\), all coordinates on
which it depends have anchors between radii
\((r-11)p\) and \((r+9)p\), or are fixed near
coarse coordinates. These free coordinates lie in
\(\eta\), not in \(z\) or \(w\).
This proves the first identity in
\eqref{eq:f15v9:split}.

The anchor distance between any \(z\)-coordinate
and any \(w\)-coordinate exceeds \(80p\).
No potential term of diameter at most \(16p\)
can involve both. Define \(V_{\rm in}\) to be
the sum of terms that involve at least one
\(z\)-coordinate, and put all other terms in
\(V_{\rm out}\). This gives the second identity,
retaining every Wilson and inverse-Haar term.
The global bijection of Theorem~\ref{f15v9:history}
then gives \eqref{eq:f15v9:fibre-product} directly.
\end{proof}

In particular, boundary-only coarse rows and
compatibility conditions have not been inverted
using a full-torus linear right inverse. They
remain in the explicitly displayed collar factor.
This addresses the rank issue for the bulk without
assuming that the entire physical pinned fibre is
a single regular chart.

\subsection{Disintegrate the collar, not the moving bulk}

Define positive finite functions
\[
 Z_N(\eta)=\int e^{-V_{\rm in}(N,z,\eta)}\,dH_z(z),
       \qquad
 W_N(\eta)=\int e^{-V_{\rm out}(N,\eta,w)}\,dH_w(w),
\]
and the ordinary compact bulk Gibbs law
\begin{equation}
       d\nu_N^\eta(z)=
               Z_N(\eta)^{-1}e^{-V_{\rm in}(N,z,\eta)}\,dH_z(z).
 \label{eq:f15v9:bulk-law}
\end{equation}
All these definitions make sense for every compact
\((N,\eta)\) at the fixed finite regulator.
They do not restrict \(\eta\) to a small-field chart.
Let \(\lambda_{B_0,N}\) be a regular conditional law
of \(\eta\) under \(H_\eta\), conditional on
\(\mathfrak b_N(\eta)=B_0\).

\begin{theorem}[The full pinned conditional measure with its collar weights]
\label{f15v9:disintegration}
For almost every actual \((B_0,N)\), define
\begin{equation}
 d\pi_{B_0,N}(\eta)
       =\frac{Z_N(\eta)W_N(\eta)\,d\lambda_{B_0,N}(\eta)}
                    {\int Z_N(\eta')W_N(\eta')\,
                                            d\lambda_{B_0,N}(\eta')}.
 \label{eq:f15v9:collar-posterior}
\end{equation}
Then the actual conditional mean of the central source is
\begin{equation}
 \mathbb E_\mu[S_b\mid B=B_0,N]
       =\int m_N(\eta)\,d\pi_{B_0,N}(\eta),\qquad
                  m_N(\eta)=\mathbb E_{\nu_N^\eta}S_b(N,z).
 \label{eq:f15v9:message-mixture}
\end{equation}
Given \(N,\eta,w\), and also if \(B\) is included,
the conditional law of \(z\) is exactly
\(\nu_N^\eta\).
For a bounded local mask \(q=q(N,z,\eta)\), the
same mixture formula holds with \(S_bq\) and
\[
                    m_N^q(\eta)=\mathbb E_{\nu_N^\eta}(S_bq).
\]
No normalization by the conditional probability
of that mask is made.
\end{theorem}

\begin{proof}
At fixed \(N\), the reference measure in the exact
history coordinates is \(H_zH_\eta H_w\).
The physical boundary statistic depends only on
\(\eta\). Disintegrate that factor by \(\mathfrak b_N\);
the reference conditional measure is therefore
\(H_z\,\lambda_{B_0,N}\,H_w\).
Multiply by the positive density in
\eqref{eq:f15v9:split}. Its two factors show
immediately that the conditional \(z\)-law is
\eqref{eq:f15v9:bulk-law}.
Integrating \(z\) and \(w\) gives the weights
\(Z_NW_N\) in \eqref{eq:f15v9:collar-posterior}.
Integration of the source gives
\eqref{eq:f15v9:message-mixture}, and the same
calculation proves the masked formula.

At the finite regulator the density is continuous,
strictly positive and bounded above and below on
a compact product. Its \((B,N)\)-marginal is
equivalent to the reference marginal, so the
reference disintegration suffices almost everywhere
for the actual law. Critical values or several
compatible collar branches are included in
\(\lambda_{B_0,N}\). They are not dropped by choosing
one inverse branch.
\end{proof}

There is no missing bulk coarea determinant here:
the bulk coordinates do not occur in the boundary
statistic at all. Any singularity or multiplicity
of that statistic belongs to \(\lambda_{B_0,N}\).
Conversely, replacing \(\pi_{B_0,N}\) by this
unweighted reference conditional law would drop
both the interior and exterior partition weights.
The finite certificate below detects that error
in an explicit two-branch example.

\subsection{A direct collar representation of the full-output score}

The same coordinates avoid differentiating the
physical boundary posterior in a locality estimate.
The far coarse output \(F\) is a function only
of \((\eta,w)\). Conditional on \(N,\eta\),
\eqref{eq:f15v9:split} makes \(z\) and \(w\)
independent. Hence
\begin{equation}
 \mathbb E_\mu[S_b\mid N,F]
                     =\mathbb E_\mu[m_N(\eta)\mid N,F].
 \label{eq:f15v9:collar-screen}
\end{equation}
For the masks in Theorem~\ref{f15v9:disintegration},
the same identity holds with \(S_bq,m_N^q\).
This is screening through actual history variables
with an explicit bulk Gibbs law, not through a
surrogate Gaussian posterior.

For example, suppose \(\Gamma\) is a set of actual
\((N,\eta)\)-values with probability of its complement
at most \(\zeta\), and the conditional essential
oscillation of \(m_N^q\) over its good \(\eta\)-values
is at most \(\omega(N)\). Suppose also \(0\le q\le1\)
and \(\mu\{q\ne1\}\le\delta_q\).
The conditional-pair proof of V8 then gives
\begin{equation}
 \mathcal L_{b,r}
   \le\left(\tfrac12\mathbb E\omega^2
                 +2\sqrt2\,C_4^2a_b^2\sqrt\zeta\right)^{1/2}
                          +C_4a_b\delta_q^{1/4}.
 \label{eq:f15v9:transfer}
\end{equation}
Indeed \(\|m_N^q\|_4\le\|S_b\|_4\), the two-copy
variance identity bounds its deviation from its
near conditional mean, and
\(\|S_b(1-q)\|_2\le\|S_b\|_4\delta_q^{1/4}\).
Conditional expectation is an \(L^2\) contraction.
V8's bound \(\|S_b\|_4\le C_4a_b\) finishes the
estimate. What remains to be proved is the
oscillation and its exceptional-set bound for
the explicit law \eqref{eq:f15v9:bulk-law}.

\subsection{The boundary response on a fixed product domain}

All functions in the unmasked law
\eqref{eq:f15v9:bulk-law} are smooth on their
compact group coordinates. Let \(\eta(\tau)\)
be a smooth collar curve at fixed \(N\); dots
mean derivatives along this curve. The bulk
integration domain and its Haar measure do
not move with \(\tau\).

\begin{theorem}[Actual normalized response, including the inherited Jacobians]
\label{f15v9:response}
For a smooth observable \(Q(N,z,\eta)\),
\begin{equation}
 \frac d{d\tau}\mathbb E_{\nu_N^{\eta(\tau)}}Q
       =\mathbb E_{\nu_N^{\eta(\tau)}}\dot Q
           -\operatorname{Cov}_{\nu_N^{\eta(\tau)}}
                                    (Q,\dot V_{\rm in}).
 \label{eq:f15v9:response}
\end{equation}
In particular
\begin{equation}
                  \dot m_N
                       =-\operatorname{Cov}_{\nu_N^\eta}
                                            (S_b,\dot V_{\rm in}).
 \label{eq:f15v9:source-response}
\end{equation}
The boundary insertion is the derivative of the
actual full measured potential:
\[
 \dot V_{\rm in}
   =\beta\!\sum_{\substack{\text{Wilson terms}\\
                          \text{assigned to }V_{\rm in}}}
                    \frac d{d\tau}s_f(\mathcal U_j)
       -\!\sum_{\substack{\text{pivot terms}\\
                          \text{assigned to }V_{\rm in}}}
                    \frac d{d\tau}\log j_{i,e}.
\]
Every differentiation follows the full inverse
reconstruction. It does not freeze a moving pivot.
Only terms incident to the varied collar coordinates
contribute; by the support bound they are near the
bulk/collar interface, not the central source.
\end{theorem}

\begin{proof}
At fixed finite regulator, smoothness and compactness
permit differentiation under the integral. Differentiate
both the numerator
\(\int Qe^{-V_{\rm in}}dH_z\) and its positive
normalizer \(Z_N\). The quotient rule gives
\[
       \mathbb E\dot Q-\mathbb E(Q\dot V_{\rm in})
                                  +\mathbb EQ\,\mathbb E\dot V_{\rm in},
\]
which is \eqref{eq:f15v9:response}.
For the central source \(S_b(N,z)\), the direct
derivative vanishes, giving
\eqref{eq:f15v9:source-response}.
The formula for \(\dot V_{\rm in}\) follows from
its definition as a sum of retained terms in
\eqref{eq:f15v9:history-density}. Its support
assertion follows term by term from
Theorem~\ref{f15v9:history}.
\end{proof}

This is a derivative of a real conditional mean
under the actual compact law, with its normalizer.
It is not a claim of a volume-uniform covariance
decay rate. In particular the number and strength
of the boundary insertions must still be controlled
if a collar curve moves many coordinates at once.

\subsection{A smooth local mask with its unavoidable response term}

The hard event in V8 is sufficient for its
probability bound, but differentiating its indicator
would require a separate boundary-flux argument.
There is a smooth alternative for the present
response identity. Let \(\chi:[0,\infty)\to[0,1]\)
be smooth, equal to one on \([0,1/2]\) and zero
on \([1,\infty)\), and use the rooted fine links
\(k_e\) in V8's box \(\mathcal K\).
For fixed \(\rho_*>0\), set
\begin{equation}
 q=\prod_{e\subset\mathcal K}
          \chi\left(\frac{2L^2(1-\tfrac12\Tr k_e)}{\rho_*^2}\right).
 \label{eq:f15v9:mask}
\end{equation}
The trace expression is smooth on the full compact
group, so this is a global smooth mask. It depends
on \((N,z,\eta)\) but not on \(w\).
If \(q\ne1\), at least one rooted angle exceeds
\(\rho_*/(\sqrt2 L)\), because
\(1-\cos\vartheta\le\vartheta^2/2\).
V8's union bound therefore gives
\begin{equation}
 \delta_q:=\mu\{q\ne1\}
   \le\min\left\{1,\ 32C_{\rm th}L^4
                    e^{-\beta\rho_*^2/(18\pi^2L^4)}\right\}.
 \label{eq:f15v9:mask-tail}
\end{equation}
Here \(C_{\rm th}\) is the fixed thermal constant,
not a function of the box side \(L\).
On \(q>0\), every rooted angle is smaller than
\(\pi\rho_*/(2L)\), using
\(1-\cos\vartheta\ge2\vartheta^2/\pi^2\).

On V8's path \(p\le u\), \(r=\lfloor u\rfloor\),
\(L=O(u^2)\), \(\beta\ge e^u\), the mask error
\(a_b\delta_q^{1/4}\) is again smaller than every
fixed inverse power of \(p\) for the stated sources.
The exact smooth masked response is nevertheless
\begin{equation}
       \dot m_N^q
          =\mathbb E_{\nu_N^\eta}(S_b\dot q)
              -\operatorname{Cov}_{\nu_N^\eta}
                                         (S_bq,\dot V_{\rm in}).
 \label{eq:f15v9:masked-response}
\end{equation}
Small \(\mu\{q\ne1\}\) does not by itself control
the conditional derivative \(\mathbb E(S_b\dot q)\).
It is explicitly retained rather than deleted as
a supposedly absent moving-domain contribution.

\subsection{A first exact gauge reduction, and its limited scope}

There is also an immediate warning against declaring
the raw compact bulk potential uniformly convex.
The level-one tree groups \(\gamma_1\) are independent
of all level-one pivots, \(X_1,Y_1,C^{(1)}\) in
\eqref{eq:f15v9:one-step-measure}.
Changing only these tree groups is a fine gauge
transformation equal to the identity at the cell
roots. The Wilson plaquette traces, the central
source \(S_b\), and the traces of the rooted
holonomies in \eqref{eq:f15v9:mask} are invariant.
Higher-level variables reconstruct \(C^{(1)}\)
without using \(\gamma_1\). Thus all terms in
\(\mathcal V_j\), \(S_b\), and \(q\) are independent
of these lowest-level tree groups.

The \(\gamma_1\)-groups included in \(z\) can
therefore be integrated out with factor one in
\(\nu_N^\eta\) and in both source means, even
after physical \(B,N\) are fixed. The latter
conditions have no \(z\)-dependence in the
boundary statistic. The collar tree groups
cannot be dropped by this argument because
they can change the prescribed physical boundary.
No analogous independence of all higher-level
tree groups at fixed lower-level free coordinates
is being asserted.

In particular, the potential relative to bulk
product Haar has exact flat directions among
the retained raw variables. A positive Hessian
floor on all of them is false. Removing the
proved gauge factors does not yet establish
convexity, mixing, or a spectral estimate for
the remaining physical detail variables.

\subsection{What has been gained for the remaining estimate}

The missing response is now expressed on the fixed
compact bulk law \eqref{eq:f15v9:bulk-law}, with
an explicit local measured potential and explicit
boundary insertions. Physical boundary compatibility
and any singular collar branches remain in
\eqref{eq:f15v9:collar-posterior}; no bulk coordinate
has been lost to an unjustified inverse.
For the full-output locality test, the direct
screening identity \eqref{eq:f15v9:collar-screen}
allows one to work with \(m_N^q(\eta)\) itself.

The next analytic requirement is a conditional
covariance/response estimate for this actual bulk
law, after an adequate gauge reduction, together
with control of the mask insertion. The archived
convex fibre estimates do not apply merely because
the bulk variables now form a product or because
the mask's exceptional probability is small.
Their measure, norms, and retained boundary data
must first be matched to the displayed law.
No such matching or decay rate is claimed here.

The original physical objective remains unchanged:
these compact conditional identities do not
construct the interacting four-dimensional
continuum or its positive Hamiltonian mass gap.

\subsection{Checks and exact preservation}

The certificate is
\begin{center}\small
\path{scripts/verify_ym_f15_v9_compact_fibre_coordinates.py}.
\end{center}
It verifies four exact collar-mixture reconstructions,
eight normalized boundary derivatives, eight screened
full-output means, four failures on omitting collar
weights, and two failures on omitting a nonconstant
Jacobian factor. These are finite positive-density
analogues, not simulations of the Wilson conditional
measure. It also checks 384 actual linear rooted
pure-gauge entries, including 136 nonzero differences
from the unrooted average, 512 boundary-copy
right-inverse entries, 12 group-dimension balances,
12 inverse-radius recurrences, and 48 collar/source
separations. All arithmetic is exact.

The certificate has 8,386 bytes and SHA--256
\begin{center}\scriptsize\ttfamily
746E54FE207C23753222CD94575DCEDECAA5C32D18CBB66AC6D88B61B84D8AA8.
\end{center}
The predecessor V8 has 183,440 bytes and SHA--256
\begin{center}\scriptsize\ttfamily
04079BB2CB963058FDEB383ED7307979E41CC9E0CFCAEBAA6E0722FCCE5DEC21.
\end{center}
Reversing this terminal insertion and the title change
recovers the predecessor exactly. V9 replaces V8 as
the sole active main note. The 11 protected sources
and 51 earlier protected certificates are unchanged.
The next companion checkpoint is V10.

\begin{firewall}
V9 confines the actual pinned compatibility constraints
to a collar in full compact history coordinates,
identifies the fixed bulk conditional measure with all
Haar factors, and proves its normalized response
identities. It does not prove the required conditional
decay, an invariant interacting RG class, a physical
continuum construction, or the Yang--Mills mass gap.
\end{firewall}

\section{V10: a pinned-data-preserving forest quotient and a physical fibre scale obstruction}
\label{f15v10:section}

V9 identified the actual compact conditional law but
removed only its lowest-level independent tree groups.
Before attempting a coercive estimate on the remaining
variables, two issues must be settled: which additional
gauge transformations preserve the prescribed data,
and in which norm a depth-uniform coercivity claim could
be true. This section supplies an exact deep-core gauge
quotient and an explicit nongauge soft direction in the
literal fixed-output fibre. Its classical Rayleigh
quotient in the fine counting norm is at most
\(65536p^{-2}\). Thus merely removing gauge freedom
cannot supply a positive depth-independent floor in
that norm. This is an upstream correction to the norm
of the proposed estimate, not a negative result about
the physical quantum mass gap.

\subsection{V10 companion checkpoint and nonduplication}

The current SM continuum V72, NS stability V26,
shell-mixing V16 and parallel YM V5 were checked again,
together with the supplied SM source-selection V17,
native regenerative stationarity V21, exact-action
Einstein bridge V16 and optional suggestions. No newer
version was present, and all protected hashes were
unchanged. Their relevant terminal results were reread.

SM V72 supplies fixed-mode, bounded-time auxiliary
many-copy control, not cutoff-uniform physical
finite-species control. NS V26 refines rank-one Oseen
estimates but supplies no estimate for our conditional
gauge law. Shell V16 retains the complete measured
one-loop flat-holonomy coefficient; it does not supply
the joint nonperturbative volume/coupling remainder.
Parallel YM V5 rules out a proposed extensive global
integer-charge tail by placement entropy, not all
local topological estimates. Source-selection V17
requires the actual coherent source panel and does
not replace it by proper marginals. Stationarity V21
distinguishes a phase-space comparison from a
configuration projection with a shrinking caustic
time. Einstein bridge V16 distinguishes a bare
compressed soft sector from the full positive
flux-energy domain. None closes V9's conditional
covariance estimate. The latter two distinctions
particularly reinforce the need to specify our norm,
full operator, and pinned domain before using a
spectral comparison.

The optional quartic suggestions are not repeated:
the corresponding measured local work is already in
f14 V20--V30. Nor is the counterexample below the
raw-history example of f14 V42. That example concerns
an unweighted sum of history-coordinate squares.
Here the vector is coexact in the original fine
counting metric, has support inside one final block,
and integrates to an exact nonlinear fibre curve.
The scale-adjusted innovation norm of f14 V43 is a
third, different norm. No contradiction with its
flat classical estimate is claimed. The next
scheduled companion checkpoint is V15.

\subsection{A forward locality bound and a permitted gauge group}

Use V9's regulator, with \(p=2^j\), \(r\ge128\),
\(L=2p(r+5)\), and \(M\ge4L\). All boxes below
are non-winding patches of the fine torus. Besides
the inverse-support bound in Theorem~\ref{f15v9:history},
we shall use its forward counterpart.

\begin{proposition}[Forward coordinate support]
\label{f15v10:forward}
Every level-\(i\) history group, including its tree
group, is a function of fine links whose bases are
within \(4(2^i-1)\) of its physical anchor.
The same bound holds for a level-\(i\) output link.
In particular all these radii are less than \(4p\).
\end{proposition}

\begin{proof}
At one step the forward output is computed from its
rooted paths in the two endpoint cells. The forward
tree products, internal non-tree variables and free
boundary variables also use only those cells.
Every input link base is within four current lattice
units of the parent anchor. Unlike the inverse map,
this calculation does not require solving a pivot
equation. After iteration, the radii obey
\[
 A_0=0,\qquad A_i\le 4\,2^{i-1}+A_{i-1},
 \qquad A_i\le4(2^i-1).
\]
This proves the assertion for the recovered history
as well as for the successive forward outputs.
\end{proof}

Take all final \(p\)-cells with bases \(px\), where
\(x\in[-(r-96),r-96]^4\cap\mathbb Z^4\).
In each cell allow arbitrary fine vertex gauge
transformations except at its root \(px\); put the
gauge transformation equal to the identity elsewhere.
Denote the resulting product group by
\[
 \mathcal H_{\rm core}
       =\SU(2)^{(p^4-1)(2r-191)^4}.
\]
Use the convention
\((g\cdot U)_{xy}=g_xU_{xy}g_y^{-1}\).

\begin{proposition}[All actual pinned data are preserved]
\label{f15v10:pins}
The action of \(\mathcal H_{\rm core}\) fixes the
entire literal final output \(C^{(j)}\), the physical
fine boundary \(B\), and the actual recovered collar
and exterior coordinates \(\eta,w\). It therefore
acts within V9's pinned bulk fibres. It preserves
the Wilson action, the source \(S_b\), and the
rooted-trace mask \(q\) of V9.
\end{proposition}

\begin{proof}
Every rooted path has the prescribed two coarse
roots as its endpoints. Conjugation equivariance
of the cutoff logarithm and exponential gives
exact one-step gauge covariance, hence iterated
covariance at the final vertices \(px\). All such
vertices are fixed by \(\mathcal H_{\rm core}\), so
every final link is unchanged, not merely its trace.

Every affected fine link base has distance at most
\((r-95)p+1\) from the origin. This is strictly
less than \((r-44)p\). By
Proposition~\ref{f15v10:forward}, no history
coordinate anchored outside \([-(r-40)p,(r-40)p]^4\)
can be affected. Far coarse coordinates are already
fixed by exact covariance. Hence \(\eta,w\) are
unchanged as coordinate values. The affected links
are also strictly inside the fine boundary annulus,
whose inner radius is \((r-6)p\), so \(B\) is fixed.
Wilson plaquette traces and the source are gauge
invariant. A rooted fine link transforms by
conjugation at the common box root; its trace,
and therefore \(q\), is invariant as well.
\end{proof}

\subsection{An exact forest quotient without a new determinant}

Inside each selected \(p\)-cell choose the tree whose
root-to-vertex path increments coordinates in the
order \(1,2,3,4\). Equivalently the parent of a
non-root vertex subtracts one in its last nonzero
cell coordinate. Let \(\mathcal T\) be the union
of these trees. Write \(\gamma_x\) for the product
of fine links along the root-to-\(x\) tree path;
set \(\gamma_x=I\) at all other vertices. Put
\[
 \bar U_{xy}=\gamma_x U_{xy}\gamma_y^{-1}.
\]
Every \(\mathcal T\)-link of \(\bar U\) equals \(I\).

\begin{theorem}[Deep-core gauge factors in the actual conditional law]
\label{f15v10:forest}
The map from fine links to the free \(\gamma_x\)'s
and the non-tree \(\bar U_e\)'s is a global bijection,
and normalized product Haar measure satisfies
\begin{equation}
 \prod_e dU_e
   =\prod_{x\text{ free in the core}}d\gamma_x
                   \prod_{e\notin\mathcal T}d\bar U_e.
 \label{eq:f15v10:forest-haar}
\end{equation}
Under the Wilson law, these \(\gamma_x\)'s are
independent normalized Haar factors, independent
of \(\bar U\). This remains true after conditioning
on any of the actual data \(N,F,B,\eta,w\), for
almost every value of those data. In particular
it gives a genuine gauge reduction of V9's actual
bulk conditional law, not only of its flat tangent.
\end{theorem}

\begin{proof}
The inverse on every edge is
\(U_{xy}=\gamma_x^{-1}\bar U_{xy}\gamma_y\), with
the tree values of \(\bar U\) fixed to \(I\).
On a tree edge this gives
\(U_{{\rm par}(x),x}=\gamma_{{\rm par}(x)}^{-1}
\gamma_x\). Successive left translations along
the tree change its edge Haar measures into the
independent vertex Haar measures. With these
variables fixed, each non-tree edge is changed
by a left and a right Haar translation. This
proves the bijection and \eqref{eq:f15v10:forest-haar}.

A root-fixed gauge transformation gives
\(\gamma_x\mapsto\gamma_xg_x^{-1}\); consequently
\(\bar U\) is unchanged. Conversely the displayed
inverse parametrizes the whole gauge orbit, and
the action is free: a stabilizing transformation
is the identity successively along every rooted
tree. Thus every invariant function is a function
of \(\bar U\) alone. The Wilson density and all
the conditioning statistics are invariant by
Proposition~\ref{f15v10:pins}. The joint law has
the form \(H_\gamma\otimes\bar\mu(d\bar U)\),
and disintegrating statistics of \(\bar U\) leaves
the first factor unchanged. The argument also
applies to invariant source insertions and masks.
\end{proof}

This is a new coordinate system on gauge orbits,
not permission to delete the intermediate history
groups one by one at fixed old coordinates. Its
quotient may mix several levels of that history.
It neither removes collar gauge transformations
that change the pins nor sets any of the inverse
Haar factors \(j_{i,e}\), collar weights, or
normalizers in V9 to one. Expressing the remaining
quotient law in other coordinates still requires
all their induced density factors.

\subsection{The literal all-depth derivative, including its root correction}

Use scalar real lattice forms for now; a fixed
colour will be inserted below. Write
\((d\lambda)_{x,\mu}=\lambda_{x+e_\mu}-\lambda_x\).
For dyadic \(q\), define the unrooted line average
\begin{equation}
 (A_q a)_{x,\mu}
   =q^{-4}\sum_{s\in\{0,\ldots,q-1\}^4}
                         \sum_{t=0}^{q-1}a_{qx+s+te_\mu,\mu}.
 \label{eq:f15v10:Aq}
\end{equation}
Let \(Q_q\) be vertex box averaging and \(E_q\)
vertex restriction to \(qx\). Let \(\Psi_2a(x)\)
be the average circulation along the 16 ordered
parity-root paths from \(2x\) to \(2x+s\).
All coefficients of \(\Psi_2\) are nonnegative
and each row has sum two.

\begin{proposition}[Root correction at every depth]
\label{f15v10:derivative}
The derivative \(R_q\) at the identity of the
literal depth-\(\log_2q\) block map is
\begin{align}
 R_q&=A_q-d\Psi_q, & \Psi_1&=0,\notag\\
 \Psi_{2q}&=\Psi_2 A_q+E_2\Psi_q, &
 \|\Psi_q\|_{\infty\leftarrow\infty}&\le2(q-1),\notag\\
 \|R_q\|_{\infty\leftarrow\infty}&\le5q-4.
 \label{eq:f15v10:Rq}
\end{align}
Operators in a composition act on their successive
current lattices. In particular \(R_qd=dE_q\).
\end{proposition}

\begin{proof}
Linearizing the rooted words at one step gives
\(R_2=A_2-d\Psi_2\), since their initial and final
tree circulations are \(\Psi_2(x)\) and
\(\Psi_2(x+e_\mu)\). The completion cutoff equals
one in a neighborhood of the identity. Direct
summation and telescoping give
\[
 A_2A_q=A_{2q},\qquad A_qd=dQ_q,\qquad
 \Psi_2d=Q_2-E_2,\qquad R_2d=dE_2.
\]
Consequently
\(R_2(A_q-d\Psi_q)=A_{2q}
-d(\Psi_2A_q+E_2\Psi_q)\), proving the recurrence
and the formula for pure gradients. The row norm
of \(A_q\) is \(q\); that of \(E_2\) is one.
Thus the bound for \(\Psi_q\) follows from
\(b_{2q}\le2q+b_q\), \(b_1=0\). Finally
\(\|d\|_{\infty\leftarrow\infty}\le2\) gives
\(q+4(q-1)=5q-4\).
\end{proof}

\subsection{A localized coexact vector in the literal kernel}

Let \(p\ge16\) be dyadic and \(m=p/4\).
On \(\mathbb Z\) set
\[
 f_m(t)=
 \begin{cases}t^2(m-t)^2/m^4,&0\le t\le m,\\
 0,&\text{otherwise},\end{cases}
 \qquad
 \psi(x)=\prod_{\nu=1}^4 f_m(x_\nu-p/2).
\]
Embed this function in the central \(p\)-cell of
the non-winding patch. Let the real two-form
\(F\) have sole nonzero component \(F_{12}=\psi\),
and put \(a=d^*F\), where adjoints use fine
counting measure. Explicitly,
\begin{equation}
 a_1(x)=\psi(x)-\psi(x-e_2),\qquad
 a_2(x)=\psi(x-e_1)-\psi(x),\qquad a_3=a_4=0.
 \label{eq:f15v10:bubble}
\end{equation}

\begin{theorem}[A physical, not gauge, fixed-output soft vector]
\label{f15v10:bubble}
The vector \(a\) is nonzero, is supported strictly
inside the central \(p\)-cell, and obeys
\begin{equation}
 d^*a=0,\qquad A_pa=0,\qquad\Psi_pa=0,
                         \qquad R_pa=0.
 \label{eq:f15v10:kernel}
\end{equation}
For the quotient fine counting norm
\[
 \|[a]\|_{\rm phys}
   :=\inf_\lambda\|a+d\lambda\|_2,
\]
where \(\lambda\) ranges over infinitesimal
\(\mathcal H_{\rm core}\) transformations, one has
\begin{equation}
 \|[a]\|_{\rm phys}=\|a\|_2>0,\qquad
 0<\frac{\|da\|_2^2}{\|[a]\|_{\rm phys}^2}
                       \le\frac{65536}{p^2}.
 \label{eq:f15v10:rayleigh}
\end{equation}
The equality of norms also holds on quotienting
by all fine vertex gradients.
\end{theorem}

\begin{proof}
The support and nonvanishing follow directly from
\eqref{eq:f15v10:bubble}. Its divergence is zero
because the two backward differences commute.
Summing \(a_1\) over a full transverse second
coordinate block telescopes to zero; summing
\(a_2\) over a full transverse first coordinate
block does the same. The support is inside one
block, so these statements hold at every coarse
base, including neighboring blocks. They prove
\(A_pa=0\).

Unroll the recurrence for \(\Psi_p\). Each term
is a restriction of \(\Psi_2A_q\) for a dyadic
\(q\le p/2\), rooted at a final vertex \(px\).
Only directions one and two can contribute.
In the ordered parity path these directions occur
before direction three. All fine links used by
their \(A_q\) averages therefore have third
coordinate in
\(px_3+\{0,\ldots,q-1\}\). But the third
coordinate of the support of \(a\) lies strictly
between \(p/2\) and \(3p/4\). These intervals
are disjoint at every final vertex, including
periodic translates. Every term is zero, hence
\(\Psi_pa=0\). Proposition~\ref{f15v10:derivative}
now gives the literal, not only curl-level,
kernel assertion.

Orthogonality \(\langle a,d\lambda\rangle=0\)
gives \(\|a+d\lambda\|_2^2
=\|a\|_2^2+\|d\lambda\|_2^2\), proving the
quotient-norm equality and excluding a gauge vector.
To bound its energy put
\[
 F_0=\sum_t f_m(t)^2,\quad
 E_1=\sum_t(f_m(t)-f_m(t-1))^2,\quad
 E_2=\sum_t(f_m(t+1)-2f_m(t)+f_m(t-1))^2.
\]
Tensor factorization gives
\(\|a\|_2^2=2E_1F_0^3>0\). On the flat cubical
complex, \(\Delta=dd^*+d^*d\) acts componentwise
as the sum of the four one-dimensional positive
second-difference operators. The terms \(dd^*F\)
and \(d^*dF\) are orthogonal, since \(d^2=0\).
Thus
\[
 \|da\|_2^2=\|dd^*F\|_2^2
    \le\|\Delta F\|_2^2
    \le16E_2F_0^3,
 \qquad
 \frac{\|da\|_2^2}{\|a\|_2^2}\le8E_2/E_1.
\]
The second inequality is Cauchy--Schwarz for the
four componentwise differences, each having squared
norm \(E_2F_0^3\).

As \(m\) is even, \(f_m(m/2)=1/16\). The
successive first differences sum to \(1/16\)
on the first half and to \(-1/16\) on the
second half. Cauchy--Schwarz on each half gives
\(E_1\ge1/(64m)\). The extension by zero of
\(s^2(1-s)^2\) on \([0,1]\) is continuously
differentiable and has derivative Lipschitz constant
at most two. Therefore each central second
difference of \(f_m\) has absolute value at
most \(2/m^2\). Only \(m+1\) of them can be
nonzero, so \(E_2\le4(m+1)/m^4\le8/m^3\).
It follows that \(8E_2/E_1\le4096/m^2
=65536/p^2\).

Finally the exact identity
\[
 \|da\|_2^2=2E_2F_0^3+6E_1^2F_0^2
\]
follows either by expanding the componentwise
Laplacian in \(\langle a,\Delta a\rangle\), or
by tensor factorization in the curl. It is strictly
positive, completing the proof.
\end{proof}

\subsection{The soft vector is tangent to an exact compact fibre curve}

Fix \(E\in\mathfrak{su}(2)\) with
\(-\frac12\Tr E^2=1\). Thus
\(\exp(tE)=I\cos t+E\sin t\). Let
\[
                    U_e(t)=\exp(t a_eE).
\]
The real coefficients \(a_e\) here are those of
\eqref{eq:f15v10:bubble}; all link values commute.

\begin{theorem}[Exact nonlinear pinning on a common interval]
\label{f15v10:curve}
For the actual completed rooted block map with
\(\rho=1/64\), if \(|t|\le\rho/1000\), then
\begin{equation}
 C^{(i)}(U(t))_e=\exp\bigl(t(R_{2^i}a)_eE\bigr)
 \quad(0\le i\le j),\qquad C^{(j)}(U(t))=I.
 \label{eq:f15v10:exact-curve}
\end{equation}
The physical boundary and the history collar and
exterior retain their identity-field values. In
particular this is a curve in the actual fibre at
\(N=I\), not just a vector in an approximate
linear constraint space.
\end{theorem}

\begin{proof}
The one-dimensional first-difference bound gives
\(\|a\|_\infty\le1/m\). By
\eqref{eq:f15v10:Rq}, for every intermediate
\(q\le p\),
\(\|R_qa\|_\infty\le5q/m\le20\).
In one blocking step the relative word
\(B^{-1}P_s\) uses at most twelve current links:
at most four for each of its two tree paths,
two for its line segment, and two for \(B^{-1}\).
Its unwrapped commuting phase therefore has
absolute value at most \(240|t|\). Its
Hilbert--Schmidt distance to the identity is at
most \(\sqrt2\,240|t|<360|t|<\rho/2\).
Thus the completion cutoff is identically one
and its principal logarithm is exactly the
unwrapped Lie-algebra phase.

Products and exponentials in this fixed colour
then turn the one-step formula exactly into its
linear rooted formula \(R_2\), not merely to
first order in \(t\). Induction proves
\eqref{eq:f15v10:exact-curve}; the last identity
uses \(R_pa=0\). Fine boundary links are unchanged
by support. The forward support estimate confines
all changed history coordinates to a central
\(5p\)-patch, inside V9's \(z\)-box. Hence the
actual \(\eta,w\) also remain at their identity
values.
\end{proof}

Normalize the classical Wilson action as
\(S_W(U)=\sum_f(1-\frac12\Tr U_f)\). Along
this curve its plaquettes are
\(\exp(t(da)_fE)\), so
\begin{equation}
 \left.\frac{d^2}{dt^2}S_W(U(t))\right|_{t=0}
                  =\|da\|_2^2.
 \label{eq:f15v10:wilson-hessian}
\end{equation}
The first derivative of \(S_W\) at the identity
vanishes in every direction. Thus the displayed
second derivative is its intrinsic classical
fibre quadratic form, independently of the
acceleration of the curve. Theorem~\ref{f15v10:bubble}
provides a nonzero physical quotient tangent
with the Rayleigh bound \eqref{eq:f15v10:rayleigh}.
Consequently no positive constant independent
of \(p\) can lower-bound this classical fibre
quadratic form in the unscaled fine counting
quotient norm. For \(\beta S_W\) the same
test gives an upper bound \(65536\beta/p^2\)
on the infimum of its Rayleigh quotients.

This conclusion does not estimate the Hessian
of the entire measured potential, which also
contains induced Haar and pivot factors. Nor
does it estimate a stochastic or quantum
Hamiltonian spectral gap. At each fixed regulator
the V9 mask is one in some neighborhood of the
identity; a common \(p\)-independent interval
with \(q=1\) is not required or asserted here.

\subsection{What the next estimate must now respect}

The gauge quotient is exact and compatible with
the actual pins; the physical soft direction is
also exact. Together they preclude one tempting
shortcut: a gauge reduction followed by an
unscaled, depth-independent classical Hessian
floor. They do not preclude conditional decay
on coarse distances. A \(p^{-2}\) fine-scale
classical stiffness is compatible with an order-one
stiffness in block-scale units.

The next concrete target is a matching lower
comparison on the pinned physical quotient,
or a carefully matched innovation norm, with
the collar domain and all measured density
terms retained. A whole-torus right inverse
or a flat unconditioned estimate does not yet
give that comparison on the pinned domain.
After such a comparison, one still needs its
nonlinear stability and control of the mask
insertion in V9's normalized response identity.
No oscillation decay for \(m_N^q(\eta)\) has
been deduced in this section.

\subsection{Exact checks and predecessor preservation}

The certificate is
\begin{center}\small
\path{scripts/verify_ym_f15_v10_forest_and_physical_fibre.py}.
\end{center}
It checks 240 free forest vertices, reconstructs
2,048 links in two finite noncommutative fields,
and verifies 1,024 canonical-link and 1,024 rooted
path invariances. These permutation-group tests
check word identities, not smooth \(\SU(2)\)
estimates. Exact rational computations construct
the coexact bubble at \(p=8,16,32\), verify
zero divergence and zero literal iterated output,
and check both norm identities. The \(p=8\)
case is an additional small-size diagnostic,
not needed for the stated \(p\ge16\) theorem.
The certificate also checks 32 bump moment
triples, 12 derivative/support recurrences,
36 deep-core clearances, and the common cutoff
margin. These finite checks supplement, rather
than replace, the all-depth proofs above.

The certificate has 7,566 bytes and SHA--256
\begin{center}\scriptsize\ttfamily
B7FC20558C7407465A7B19A4CBDFDF6D7C59F3156DE3E4D0821602F081EE4A3F.
\end{center}
The predecessor V9 has 209,815 bytes and SHA--256
\begin{center}\scriptsize\ttfamily
1CDC244484D8827D8465E57C300CEB4FF883D222B6140D4073B8F9CE62FEE1F2.
\end{center}
Reversing this terminal insertion and the title
change recovers V9 exactly. V10 replaces V9 as
the sole active main note. The 11 protected
sources and 52 earlier protected certificates
are unchanged.

\begin{firewall}
V10 proves an exact deep-core gauge quotient,
an all-depth literal derivative formula, and
a compactly supported nongauge exact fibre
curve with classical physical Rayleigh quotient
at most \(65536p^{-2}\). It proves neither
the corresponding lower comparison nor actual
conditional mixing, an invariant interacting
RG class, continuum reconstruction, or a
positive physical Yang--Mills mass gap. The
\href{https://www.claymath.org/millennium/yang-mills-the-maths-gap/}
{existence-and-mass-gap objective} remains
strictly stronger than these finite-regulator
results.
\end{firewall}

\section{V11: scale-correct coercivity of the literal rooted fibre and the physical meaning of linear pins}
\label{f15v11:section}

V10's exact fibre curve rules out an unscaled
depth-independent classical floor. This section
proves a matching lower comparison of order
\(p^{-2}\) in a specified fine counting quotient
norm, uniformly in depth and volume. It does so
directly for the literal rooted block derivative,
without replacing that derivative by its curl.
The proof uses a commuting interpolation and a
bounded correction of the coarse root values.

The domain issue is then treated separately.
Additional linear pins need not descend to the
gauge quotient. Their physical constraint is
obtained only after quotienting their values by
those obtainable through gauge motion. This gives
the exact constrained classical Gaussian, including
linearized V9 collar pins, and a covariance bound
in an explicitly inherited metric. It does not
identify that Gaussian with the nonlinear compact
conditional measure.

The one-step Schur comparison in f14 V41 and the
rescaled innovation comparison in f14 V43 are not
reproved or used as substitutes for this argument.
Neither states the all-depth fine counting
quotient comparison below with the literal root
correction and the induced physical pin map.
The current source hashes remain unchanged;
the next scheduled companion checkpoint is V15.

\subsection{Spaces, norms, and the precise quotient}

Let the fine torus have side \(M\), let \(p\ge2\)
be dyadic with \(p\mid M\), and assume \(M\ge4p\).
Every cochain norm uses counting measure, and
every edge or plaquette is counted once in its
positive orientation. The argument is written
for one real colour and tensors with
\(\mathfrak{su}(2)\) without changing constants.
Write \(d_f,d_c\) for the fine and final coarse
exterior derivatives. Let \(R_p,A_p,\Psi_p,E_p\)
be the actual derivative and the operators of
Proposition~\ref{f15v10:derivative}. Set
\begin{equation}
 V_p=\ker R_p,\qquad
 G_p=\{d_f\lambda:E_p\lambda=0\},\qquad
 \mathcal Q_p=V_p/G_p,
 \label{eq:f15v11:spaces}
\end{equation}
and equip \(\mathcal Q_p\) with the quotient norm
\begin{equation}
 \|[a]\|_p=\inf_{g\in G_p}\|a+g\|_2.
 \label{eq:f15v11:quotient-norm}
\end{equation}
Indeed \(G_p\subset V_p\) because
\(R_pd_f=d_cE_p\). This is the whole-torus
root-fixed gauge quotient, not merely the
deep-core quotient of V10. Its relation to
additional physical pins is established below,
not assumed.

\subsection{A commuting interpolation with an elementary approximation bound}

For \(y=px+s\), \(s\in\{0,\ldots,p-1\}^4\),
put \(w_s(0)=1-s/p\), \(w_s(1)=s/p\).
Define nodal interpolation of a coarse vertex
field \(h\) by
\[
 (I_ph)(px+s)=\sum_{\epsilon\in\{0,1\}^4}
       \left(\prod_{k=1}^4w_{s_k}(\epsilon_k)\right)h(x+\epsilon).
\]
For a coarse one-form \(A\), define its edge
interpolation by
\begin{equation}
 (W_pA)_{px+s,\mu}
   =\frac1p\sum_{\substack{\epsilon\in\{0,1\}^4\\\epsilon_\mu=0}}
       \left(\prod_{k\ne\mu}w_{s_k}(\epsilon_k)\right)
                                                 A_{x+\epsilon,\mu}.
 \label{eq:f15v11:interpolation}
\end{equation}
Indices are periodic. The following estimates
are deliberately conservative in their constants.

\begin{proposition}[Commutation, stability, and local approximation]
\label{f15v11:interpolation}
The interpolation satisfies
\begin{equation}
 E_pI_p=I,\qquad W_pd_c=d_fI_p,
 \qquad \|W_p\|_{2\leftarrow2}\le p,
 \qquad \|A_p\|_{2\leftarrow2}\le p^{-1}.
 \label{eq:f15v11:commuting}
\end{equation}
For every fine one-form \(b\),
\begin{equation}
 \|b-W_pA_pb\|_2\le200p\|\nabla_f b\|_2,
 \qquad
 \|\nabla_f b\|_2^2
       :=\sum_{\mu,\nu,x}|b_{x+e_\nu,\mu}-b_{x,\mu}|^2.
 \label{eq:f15v11:approximation}
\end{equation}
\end{proposition}

\begin{proof}
Evaluation at a root proves \(E_pI_p=I\).
Taking a fine difference of the multilinear
vertex interpolation in direction \(\mu\)
replaces its \(\mu\)-factor by the corresponding
coarse difference divided by \(p\). This remains
true across the last edge of a cell. It is
exactly \eqref{eq:f15v11:interpolation}, proving
the commuting identity.

For each output edge, the transverse weights in
\(W_p\) are nonnegative and sum to one before
the factor \(1/p\). For a fixed coarse edge,
their total over all fine edge bases is \(p^4\):
there are \(p\) choices in its own direction,
and the two adjacent cells in each transverse
direction contribute weights summing to \(p\).
Jensen's inequality therefore gives
\(\|W_pA\|_2^2\le p^2\|A\|_2^2\).

Each row of \(A_p\) has mass \(p\).
Each column has mass \(p^{-3}\): after summing
over all coarse bases, a given fine edge occurs
\(p\) times with weight \(p^{-4}\), one for
each position in the length-\(p\) segment.
Weighted Cauchy--Schwarz gives its stated norm.

Put \(K_p=W_pA_p\). It acts separately on
each orientation. Its rows are nonnegative
and sum to one, since it preserves constant
fine one-forms. Each entry is at most \(p^{-4}\):
an \(A_p\) entry is at most \(p/p^4\), and
the weights used by one \(W_p\) row sum to
\(1/p\). An entry between bases \(y\) and
\(z\) can be nonzero only if the lifted
displacement satisfies \(|z-y|_\infty\le2p\).
There is no winding ambiguity when \(M\ge4p\);
the actual displacement interval has width
less than \(3p\) in each coordinate.

Jensen's inequality and the entry bound imply
\[
 \|b-K_pb\|_2^2
 \le p^{-4}\sum_{|\delta|_\infty\le2p}
                       \|b-\tau_\delta b\|_2^2.
\]
A path of \(|\delta|_1\) unit translations
and the triangle inequality give
\(\|b-\tau_\delta b\|_2
\le|\delta|_1\|\nabla_f b\|_2\).
There are at most \((5p)^4\) displacements
in the sum, and \(|\delta|_1\le8p\).
The squared constant is therefore at most
\(5^4\cdot8^2p^2=40000p^2\), proving
\eqref{eq:f15v11:approximation}.
\end{proof}

\subsection{The root correction is bounded in the fine counting metric}

The \(\ell^\infty\) bound in V10 grows with
\(p\). Its counting \(\ell^2\) bound behaves
differently, and is the crucial point here.

\begin{proposition}[Uniform root-potential bound]
\label{f15v11:psi-bound}
With counting measures on both lattices,
\begin{equation}
       \|\Psi_p\|_{2\leftarrow2}\le2(1-p^{-1})<2.
 \label{eq:f15v11:psi-bound}
\end{equation}
Let \(J_ph\) be the fine vertex field equal to
\(h(x)\) at \(px\) and zero at other vertices.
Then
\begin{equation}
 E_pJ_p=I,\qquad \|d_fJ_ph\|_2^2=8\|h\|_2^2.
 \label{eq:f15v11:spikes}
\end{equation}
\end{proposition}

\begin{proof}
Every row of \(\Psi_2\) has mass two. A fine
edge belongs to at most one parity-root cell,
and its coefficient there is at most \(1/2\).
For example, the first ordered direction is
used by eight of the sixteen paths; later
directions have smaller coefficients. Thus
its column mass is at most \(1/2\), and
\(\|\Psi_2\|_{2\leftarrow2}\le1\).

Use V10's recurrence
\(\Psi_{2q}=\Psi_2A_q+E_2\Psi_q\).
Restriction of vertices has counting norm
at most one. By \eqref{eq:f15v11:commuting},
the resulting norm recurrence is
\(c_{2q}\le q^{-1}+c_q\), \(c_1=0\).
Its geometric sum is \(2(1-p^{-1})\).
Finally, when \(p\ge2\), no fine edge joins
two distinct \(p\)-roots. Each root has eight
incident edges, on each of which the difference
of \(J_ph\) is \(h(x)\) up to sign. Summation
gives \eqref{eq:f15v11:spikes}.
\end{proof}

\begin{theorem}[Literal rooted-fibre coercivity at the correct scale]
\label{f15v11:coercivity}
For \(a\in V_p\), let \(b\) be its orthogonal
projection onto the complement of all fine
vertex gradients. Then
\begin{equation}
 \|b\|_2\le\|[a]\|_p\le8\|b\|_2,
 \qquad
 \|[a]\|_p\le1600p\|d_fa\|_2.
 \label{eq:f15v11:coercivity}
\end{equation}
In particular the classical energy descends
to a positive definite quadratic form on
\(\mathcal Q_p\), whose associated operator
\(H_p\) satisfies
\begin{equation}
 H_p\succeq\frac1{2560000p^2}\,I.
 \label{eq:f15v11:floor}
\end{equation}
All constants are independent of depth and
volume. For \(p\ge16\), V10's coexact bubble
also lies in this quotient, so
\begin{equation}
 \frac1{2560000p^2}
       \le\inf_{[a]\ne0}
            \frac{\|d_fa\|_2^2}{\|[a]\|_p^2}
       \le\frac{65536}{p^2}.
 \label{eq:f15v11:two-sided}
\end{equation}
\end{theorem}

\begin{proof}
Write \(a=b+d_f\lambda\), with
\(d_f^*b=0\). Since \(R_pa=0\),
\[
 A_pb=d_c\kappa,\qquad
                  \kappa=\Psi_pb-E_p\lambda.
\]
Hence \(W_pA_pb=d_fI_p\kappa\) is a fine
gradient, orthogonal to \(b\). Consequently
\[
 \|b\|_2^2
   =\langle b,b-W_pA_pb\rangle
   \le\|b\|_2\,200p\|\nabla_f b\|_2.
\]
Expansion of the commuting fine differences
and their adjoints gives the discrete identity
\[
 \|\nabla_f b\|_2^2
        =\|d_fb\|_2^2+\|d_f^*b\|_2^2
        =\|d_fa\|_2^2.
\]
This proves \(\|b\|_2\le200p\|d_fa\|_2\),
including the case \(b=0\).

The quotient by only root-fixed gradients
requires an additional step. Define
\[
 \theta=J_p\Psi_pb-I_p\kappa,\qquad
                    \widehat a=b+d_f\theta.
\]
Then \(E_p\theta=\Psi_pb-\kappa=E_p\lambda\).
Thus \(a-\widehat a=d_f(\lambda-\theta)\in G_p\),
so \(\widehat a\) is an admissible representative
of the \emph{literal} rooted-fibre class. Moreover
\[
 \begin{split}
 \|\widehat a\|_2
 &\le\|b\|_2+\|d_fJ_p\Psi_pb\|_2
                         +\|W_pd_c\kappa\|_2\\
 &\le\|b\|_2+4\sqrt2\|b\|_2
                         +p\|A_pb\|_2\\
 &\le(2+4\sqrt2)\|b\|_2<8\|b\|_2.
 \end{split}
\]
The lower norm bound holds because projecting
away all gradients can only reduce the infimum
over the smaller space \(G_p\). Combining these
bounds proves \eqref{eq:f15v11:coercivity} and
\eqref{eq:f15v11:floor}.

If the curl energy vanishes, these inequalities
force \([a]=0\), proving positive definiteness.
V10's bubble is already orthogonal to every
gradient, so its quotient norm is unchanged
on replacing the deep-core gauge group by
\(G_p\). Its Rayleigh upper bound proves
\eqref{eq:f15v11:two-sided}.
\end{proof}

\subsection{Which part of a linear pin is a physical constraint?}

The preceding whole-torus quotient is not
automatically the same metric as the quotient
by gauge transformations confined to V10's core.
In particular some root-fixed transformations
change the physical boundary or collar pins.
The following finite-dimensional statement
treats that issue without dropping those pins.

Let \(D:V_p\to Y\) be any linear list of extra
pins, with redundant rows allowed. Define
\begin{equation}
 Y_g=D(G_p),\qquad
 \bar D:\mathcal Q_p\longrightarrow Y/Y_g,
                    \quad \bar D[a]=[Da].
 \label{eq:f15v11:effective-pin}
\end{equation}
For an admissible pin value \(y\), write
\[
 T_y=\{a\in V_p:Da=y\},\qquad
 G_D=G_p\cap\ker D,\qquad
 W_{[y]}=\{\xi\in\mathcal Q_p:\bar D\xi=[y]\}.
\]

\begin{theorem}[Exact Gaussian pinning after accounting for gauge motion]
\label{f15v11:physical-pins}
The quotient map induces an affine bijection
\begin{equation}
                 T_y/G_D\ \longrightarrow\ W_{[y]}.
 \label{eq:f15v11:pinned-bijection}
\end{equation}
Give \(T_y/G_D\) any fixed Euclidean quotient
Lebesgue measure. The normalized classical
Gaussian with energy \(\beta\|d_fa\|_2^2/2\)
pushes forward under this bijection to the
Gaussian with energy
\(\beta\langle\xi,H_p\xi\rangle/2\) on
the affine subspace \(W_{[y]}\).

Let \(C_{p,D}\) be its centered covariance,
viewed as an operator on \(\mathcal Q_p\)
with norm \eqref{eq:f15v11:quotient-norm}.
Then
\begin{equation}
       0\preceq C_{p,D}
       \preceq\beta^{-1}H_p^{-1}
       \preceq\frac{2560000p^2}{\beta}\,I.
 \label{eq:f15v11:conditional-covariance}
\end{equation}
The covariance is independent of the admissible
value \(y\). This statement does not assert
uniform equivalence with the metric intrinsically
induced on \(T_y/G_D\) by its particular pins.
\end{theorem}

\begin{proof}
The map \(\bar D\) is well-defined because
changing a representative by \(g\in G_p\)
changes \(Da\) by an element of \(Y_g\).
A class \([a]\) has a representative satisfying
\(Da=y\) exactly when \(y-Da\in D(G_p)\).
This is precisely \(\bar D[a]=[y]\).
Two such representatives differ by a gauge
vector in \(\ker D\), hence by \(G_D\).
This proves \eqref{eq:f15v11:pinned-bijection}.

An affine bijection between fixed finite-dimensional
Euclidean affine spaces changes Lebesgue measure
by a positive constant. That constant cancels
from a normalized Gaussian. The energy is
unchanged along every gauge class, proving
the claimed pushforward law. This argument
quotients null directions; it does not integrate
an improper infinite gauge volume.

Choose inner products on the pin quotient,
and restrict its range to \(\operatorname{ran}\bar D\).
On that range the usual positive inverse in
the following formula is well-defined; equivalently
one may use the Moore--Penrose inverse:
\[
 \begin{split}
 C_0&=\beta^{-1}H_p^{-1},\\
 C_{p,D}
   &=C_0-C_0\bar D^*
            (\bar DC_0\bar D^*)^\dagger\bar DC_0.
 \end{split}
\]
For completeness, whiten \(\xi\) by writing
\(\xi=C_0^{1/2}z\). A standard Euclidean
Gaussian conditioned on the linear equation
\(\bar DC_0^{1/2}z=[y]\) has covariance the
orthogonal projection onto that map's kernel:
decompose \(z\) into its kernel and orthogonal
complement, whose Gaussian coordinates are
independent. Conjugating this projection by
\(C_0^{1/2}\) proves the formula, both operator
inequalities, and independence from \(y\).
The last bound follows from
Theorem~\ref{f15v11:coercivity}.
\end{proof}

One must not instead fix an arbitrary gauge and
then impose every original row of \(D\) as a
physical pin. For example let physical variables
be \((x_1,x_2)\), let \(g\) be a gauge variable,
and take
\[
                    D(x_1,x_2,g)=(x_1+g,x_2).
\]
At \(D=0\), the physical constraint is only
\(x_2=0\); \(g=-x_1\) meets the first row.
For physical precision
\(\left(\begin{smallmatrix}2&1\\1&3\end{smallmatrix}\right)\)
at \(\beta=1\), the remaining variance of
\(x_1\) is \(1/2\). Setting \(g=0\) before
applying \(D=0\) would incorrectly set that
variance to zero. This is a linear domain
warning, not a new model of the Wilson measure.

\subsection{Application to V9's flat bulk domain, with its limitations}

At the identity, differentiate the actual
history collar and exterior maps of V9 and
take their differentials as \(D\). Include
the physical boundary differential if desired.
It is redundant once \(N,\eta\) are fixed,
because \(B=\mathfrak b_N(\eta)\).
The final far output is part of \(\eta,w\),
so a bulk variation at fixed \(N,\eta,w\)
has \(R_pa=0\).

V9's global history-coordinate diffeomorphism
shows that these fixed-coordinate variations
are the actual tangent of its bulk fibre;
they are not merely a possibly oversized
formal kernel of singular compatibility equations.
Its tangent space is therefore the corresponding
\(T_0\). At the identity the Wilson action is
critical, so its intrinsic second variation
on that fibre is exactly \(\|d_fa\|_2^2\).
Theorem~\ref{f15v11:physical-pins} consequently
describes its classical physical Gaussian
after all its linear gauge null directions
\(G_D\) have been removed.

The resulting bound is in the inherited
\(\mathcal Q_p\) metric. It does not say that
the deep-core forest coordinates alone remove
every null direction, or that all pinned-coordinate
metrics have a common comparison constant.
The effective physical pin \(\bar D\), not
the unmodified gauge-dependent list \(D\),
is the object on which a spatially resolved
Gaussian response analysis must operate.

There are three important remaining gaps.
First, an operator-norm covariance estimate
does not yet imply exponential off-diagonal
decay: the quotient projection, induced pins,
and their localization must be controlled.
Second, the Gaussian here retains the classical
quadratic action, not the full measured
potential including the inverse-Haar factors
and nonlinear constraint geometry. No joint
large-volume/depth Laplace remainder has been
proved. Third, even such a remainder would
still have to control V9's normalized boundary
response and its mask insertion before yielding
the desired conditional locality estimate.

The gain is nevertheless an actual lower
comparison, not an assumption of one: the
linear classical scale is now bracketed by
\eqref{eq:f15v11:two-sided}, and the linear
pinning operation has been matched to the
physical quotient without suppressing gauge
degrees that can meet the boundary data.

\subsection{Checks, preservation, and scope}

The exact certificate is
\begin{center}\small
\path{scripts/verify_ym_f15_v11_rooted_fibre_coercivity.py}.
\end{center}
It verifies 17,408 entries of the commuting
interpolation identity, its tested stability
and root restriction identities, and 80 local
positive-kernel rows with 297,792 nonzero entries.
Two periodic co-closed examples have nonzero
literal rooted outputs despite zero unrooted
averages. Their explicit root repair cancels
256 and 16 output entries, respectively,
while preserving their curls and checking the
orthogonal norm identity. These are useful
tests of the distinction between \(A_p\)
and \(R_p\), not a replacement of the
all-depth proof. The certificate additionally
checks 125 exact Gaussian variance comparisons
and the spurious-constraint example above.

The certificate has 9,821 bytes and SHA--256
\begin{center}\scriptsize\ttfamily
02288748F34710E80740D53A3F82319D792FCDC30FD843802E4CFF1B6A589692.
\end{center}
The predecessor V10 has 231,275 bytes and SHA--256
\begin{center}\scriptsize\ttfamily
64B0C829DEB7E5E12B8E994BF24BF004F4639861A93C2D0AD884004D5EF4F3F5.
\end{center}
Reversing this terminal insertion and the title
change recovers V10 exactly. V11 replaces it
as the sole active main note. The 11 protected
sources and 53 earlier protected certificates
remain unchanged.

\begin{firewall}
V11 proves depth- and volume-uniform classical
coercivity at scale \(p^{-2}\) on the literal
rooted fibre modulo root-fixed gradients, and
identifies the Gaussian effect of additional
linear pins after their gauge freedom is
accounted for. It does not prove nonlinear
conditional mixing, an interacting continuum
construction, or the physical Yang--Mills
mass gap. Those remain the objective, not
consequences of this finite-regulator comparison.
\end{firewall}

\section{V12: local gauge reduction and classical Gaussian screening}
\label{f15v12:section}

The next missing ingredient after V11 is localization,
not another operator-norm comparison. We first construct
coordinates in which every history tree group is fixed
simultaneously, while retaining the transformations of
the other history groups and every measured Jacobian.
This gives an exact compact conditional law on a
smaller, fully gauge-reduced interior. Its classical
quadratic precision has finite range and explicit
polynomial conditioning in \(p\). A direct inverse
series then yields spatial decay and a quadratic-source
screening estimate for the zero-background Gaussian.

The polynomial powers below are not optimized. They
are useful because an enlarged collar can overcome
them while preserving the full compact thermal
small-field probability bound. The remaining task
is a nonlinear comparison on these same coordinates
and boundary data; the Gaussian is not substituted
for the full law.

\subsection{One simultaneous gauge slice, rather than deleting old tree variables}

Let \(\mathcal H_p\) be all fine vertex gauge
transformations equal to \(I\) at the final
vertices \(px\). Retain V9's global coordinates
\((C,(X_i,Y_i,\gamma_i)_{i=1}^j)\), with
\(C=C^{(j)}\) and \(p=2^j\). The new slice is
\begin{equation}
              \gamma_i=I\quad\hbox{for every }i.
 \label{eq:f15v12:canonical-slice}
\end{equation}
Write \(v_i=(\bar X_i,\bar Y_i)\) for the
\(45n_i\) bare compact groups on this slice,
and write \(\bar U=\mathcal U_j(C,v,I)\).
These are bare groups, not the relative coordinates
\(Y=C e^\eta\) used in some archived chart estimates.

\begin{theorem}[Exact compact canonical history coordinates]
\label{f15v12:canonical}
Every fine configuration has a unique representation
\begin{equation}
 U_{xy}=g_x^{-1}\bar U_{xy}g_y,\qquad
  g\in\mathcal H_p,\qquad
  \bar U=\mathcal U_j(C,v,I).
 \label{eq:f15v12:canonical-map}
\end{equation}
This is a global smooth change of variables. With
all Haar measures normalized, the Wilson law becomes
\begin{equation}
 d\mu\ \propto\ e^{-\mathcal W_j(C,v)}\,dC\,dH_v(v)\,dH_g(g),
 \quad
 \mathcal W_j=\beta S_W(\bar U)
               -\sum_{i,e}\log j_{i,e}(\bar X_i,\bar Y_{i,e},\bar C^{(i)}_e).
 \label{eq:f15v12:measured-law}
\end{equation}
Here every intermediate \(\bar C^{(i)}\) is
reconstructed on the slice. Thus the gauge variables
are independent Haar factors in these \emph{new}
coordinates, including at fixed full output \(C\).
They cannot in general be discarded after fixing
the old physical boundary or history collar.
\end{theorem}

\begin{proof}
A level-\(i\) tree path transforms by its two
endpoint gauge values, at fine vertices of the
form \(2^i x\) and \(2^{i-1}(2x+s)\).
If all slice tree products are \(I\), its value
in \(U\) is therefore \(g_{2^i x}^{-1}
g_{2^{i-1}(2x+s)}\). These pairs form an abstract
rooted tree on the fine vertices of each \(p\)-cell:
every non-final-root vertex first occurs as a
nonroot at exactly one level. Starting with
\(g_{px}=I\), its edge products determine all
other \(g\)'s uniquely. Applying the resulting
gauge transformation to \(U\) sets every
\(\gamma_i\) to \(I\). Exact gauge covariance
of the completed block preserves \(C\).
V9's global coordinate inverse then gives
existence, uniqueness, and smoothness of the
remaining \(v\)'s and of the inverse map.

The change from these abstract-tree products
to vertex groups \(g\) is a triangular Haar
change, just as for an ordinary rooted forest.
For fixed \(g\), each old internal group is
a conjugate of \(\bar X_i\), and each old
boundary group is a left and right translate
of \(\bar Y_i\), by the appropriate parent
root values of \(g\). Thus their product Haar
measure is unchanged. Each inverse-pivot Haar
Jacobian is invariant under the corresponding
endpoint gauge transformations: the pivot map
is equivariant and endpoint translations preserve
both input and output Haar. The Wilson action
is also invariant. Applying these changes to
\eqref{eq:f15v9:history-density} proves
\eqref{eq:f15v12:measured-law}. The counts agree:
\(15\sum_i n_i=M^4-n_j\) gauge groups and
\(45\sum_i n_i=3(M^4-n_j)\) detail groups.
\end{proof}

In particular this result does not assert independence
of old higher-level \(\gamma_i\)'s at fixed old
\(X_i,Y_i\). Those variables were transformed
together. No pivot density or collar normalizer
has been replaced by one.

\subsection{The original pins stay outside a smaller canonical interior}

Give \(v_i\) its usual anchor \(2^i x\). At
fixed \(C\), a reconstructed \(\bar U\)-link
depends on detail anchors within \(4p\), by
V9. Applying the pointwise transformation
\eqref{eq:f15v12:canonical-map} adds only its
endpoint \(g\)'s. V10's forward support bound
therefore gives the following fact: an old
history coordinate anchored at \(y\) depends
on canonical details with anchors within \(8p\)
of \(y\), and on vertex gauge values near
its forward fine carrier.

In V9's geometry define
\[
 J=\{v\hbox{-groups with anchor in }
                         [-(r-60)p,(r-60)p]^4\},
 \qquad O=J^c,
 \qquad r\ge128.
\]
The complement here contains all other canonical
detail groups, not final output groups.

\begin{proposition}[Exact reduced interior law with the original pins retained]
\label{f15v12:interior}
At fixed \(C\), all old collar and exterior
coordinates \(\eta,w\), and the fine boundary
\(B\), are functions of \((C,v_O,g)\) only.
The central gauge-invariant source is a function
\(S_b=S_b(C,v_J)\). Its original physical pins
therefore impose no constraint on \(v_J\) once
\((C,v_O,g)\) are specified.

Split \(\mathcal W_j=W_J+W_O\) by putting in
\(W_J\) all measured terms that involve a
\(J\)-coordinate. Then the exact conditional
law of \(v_J\), also with the original pins
included, is
\begin{equation}
 d\nu_{J,C}^{v_O}(v_J)
  =Z_{J,C}(v_O)^{-1}e^{-W_J(C,v_J,v_O)}dH_J(v_J).
 \label{eq:f15v12:interior-law}
\end{equation}
Every individual term in \(\mathcal W_j\) has
anchor diameter at most \(16p\). In particular
\(W_J\) depends on exterior details only in
that thickness of the boundary of \(J\).
\end{proposition}

\begin{proof}
Old history groups in \(\eta,w\) are anchored
outside the box of radius \((r-40)p\).
The gap of \(20p\) from \(J\) exceeds their
\(8p\) detail support radius. Their far output
groups are included in the fixed \(C\).
Also \(B=\mathfrak b_N(\eta)\), so it is
unchanged by a \(J\)-variation. The central
source has the fine support specified in V9;
the inverse support estimate and gauge invariance
put its detail carrier in a central \(12p\)-box,
strictly inside \(J\).

The measured density in
\eqref{eq:f15v12:measured-law} is independent
of \(g\). Its term supports are the supports
of V9 with the tree groups fixed to constants,
so their diameter bound remains valid. The
pins are functions of the variables being
conditioned on, not of \(v_J\). Disintegrating
the product density therefore proves
\eqref{eq:f15v12:interior-law} for almost every
actual value of the conditioning data.
\end{proof}

For example, with
\(m_J(C,v_O)=\mathbb E_{\nu_{J,C}^{v_O}}S_b\),
the exact tower identity reads
\[
 \mathbb E_\mu[S_b\mid C,\eta,w]
       =\mathbb E_\mu[m_J(C,v_O)\mid C,\eta,w].
\]
The posterior exterior law in this formula
still contains the normalizer \(Z_{J,C}\)
and all pin compatibility weights. It is
not replaced by unweighted exterior Haar.
The V9 mask is gauge invariant but has a
larger carrier than the central source; its
insertion remains a separate issue.

\subsection{Explicit polynomial conditioning in local detail coordinates}

Now linearize at \(C=I,v=I\). Use the same
letter \(v\) for Lie-algebra detail increments,
with the direct sum counting norm. Let
\(E_jv\) be the corresponding fine one-form
on the simultaneous tree slice, at zero
final output increment, and set
\begin{equation}
 K_j=E_j^*d_f^*d_fE_j.
 \label{eq:f15v12:precision}
\end{equation}
This is a classical precision, not the Hessian
of the full measured logarithmic density.

\begin{proposition}[A local raw-coordinate comparison]
\label{f15v12:conditioning}
The map \(v\mapsto[E_jv]\) is an isomorphism
onto V11's \(\mathcal Q_p\). Moreover
\begin{equation}
 \|E_j\|\le p^4,\qquad
 \|v\|_2\le512p^3\|[E_jv]\|_p,
 \qquad
 h_p I\preceq K_j\preceq u_pI,
 \label{eq:f15v12:conditioning}
\end{equation}
where
\begin{equation}
               h_p=2^{-40}p^{-8},\qquad u_p=16p^8.
 \label{eq:f15v12:hu}
\end{equation}
Entries of \(K_j\) vanish between detail
anchors at distance greater than \(16p\).
The same spectral bounds hold for every
principal submatrix, including \((K_j)_{JJ}\).
\end{proposition}

\begin{proof}
The isomorphism is the derivative of the
global gauge slice. We give explicit norm
bounds to avoid inferring conditioning merely
from bijectivity.

At one step on the tree slice, the internal
non-tree and nonpivot boundary edges are
literally the free detail entries. For a
coarse edge, let \(Y_a\), \(a=1,\ldots,7\),
be its free boundary edges and let \(X_a^-,X_a^+\)
be the two corresponding internal edges in
its endpoint cells, with tree edges interpreted
as zero. Each nonroot length-two word contains
one \(Y_a\) and one such internal edge.
The linear pivot formula is therefore
\begin{equation}
             B=8C-\sum_{a=1}^7Y_a
                        -\frac12\sum_{a=1}^7(X_a^-+X_a^+).
 \label{eq:f15v12:linear-pivot}
\end{equation}
The complete one-step inverse matrix has
absolute row sum at most \(22\). A coarse
input has column sum eight; a free boundary
input has sum two; an internal non-tree input
has sum at most two, from itself and the
two adjacent pivots. The squared operator
norm is at most \(22\cdot8<16^2\).
Iterating this estimate on the joint coarse
and detail input gives
\(\|E_jv\|_2^2\le16^{2j}\sum_i\|v_i\|_2^2
=p^8\|v\|_2^2\).

For the converse start with any fine
representative \(a\in\ker R_p\). By V11,
\(\|R_q\|_{2\leftarrow2}\le8\) for every
dyadic \(q\). Let \(t_i\) be the ordered
parity-tree circulations of \(R_qa\), with
\(q=2^{i-1}\). The tree-circulation matrix
has row mass at most four and column mass
at most eight, so \(\|t_i\|_2\le48\|a\|_2\).
The linear gauge potential that sets all
these tree circulations to zero is the sum
of their prolongations constant on the fine
\(q\)-cells. Each prolongation has norm
\(q^2\). The resulting potential \(h\)
vanishes at final roots and satisfies
\[
 \|h\|_2\le48\sum_{i=1}^j2^{2(i-1)}\|a\|_2
       =16(p^2-1)\|a\|_2.
\]
Indeed the value at a child vertex is the
value at its abstract-tree parent plus that
tree circulation; this is the linearization
of the construction in
Theorem~\ref{f15v12:canonical}.
Thus \(a^{\rm can}=a-d_fh\) satisfies
\(\|a^{\rm can}\|_2\le64p^2\|a\|_2\), using
\(\|d_f\|_{0\to1}\le4\).

On the slice, each \(v_i\) is a subset of
the edge entries of \(R_qa^{\rm can}\).
Consequently
\(\|v\|_2\le8\sqrt j\|a^{\rm can}\|_2
\le512p^2\sqrt j\|a\|_2\le512p^3\|a\|_2\).
Minimize over representatives of the same
root-fixed gauge class. The canonical field
and its details do not change, proving the
second estimate in
\eqref{eq:f15v12:conditioning}.

V11 now gives
\(\|v\|_2\le512\cdot1600p^4\|d_fE_jv\|_2\).
Since \((512\cdot1600)^2<2^{40}\), this
implies the lower bound \(h_p\). The upper
bound follows from \(\|d_f\|_{1\to2}\le4\)
and \(\|E_j\|\le p^4\). The discrete Hodge
identity proves this curl norm bound as in
V11. The support of a classical plaquette
term has diameter at most \(16p\), so its
Hessian does too. Rayleigh quotients of a
principal submatrix are those of the full
matrix on vectors extended by zero, proving
the final assertion.
\end{proof}

\subsection{A spatial inverse bound that survives interior Dirichlet pinning}

For sets of detail indices \(A,B\), let
\(d(A,B)\) be their minimum anchor distance
in the periodic \(\ell^\infty\) metric. Denote
coordinate restrictions by \(P_A,P_B\).

\begin{theorem}[Finite-range classical inverse decay]
\label{f15v12:inverse-decay}
Let \(K\) be \(K_j\) or any of its principal
submatrices. Put
\(\delta_p=h_p/u_p=2^{-44}p^{-16}\).
Then
\begin{equation}
 \|P_AK^{-1}P_B^*\|
 \le h_p^{-1}(1-\delta_p)^{\lceil d(A,B)/(16p)\rceil}
 \le h_p^{-1}\exp\left\{-\frac{d(A,B)}{2^{48}p^{17}}\right\}.
 \label{eq:f15v12:inverse-decay}
\end{equation}
For coinciding or overlapping sets the last
bound is understood with distance zero.
\end{theorem}

\begin{proof}
Write \(T=I-K/u_p\). Its norm is at most
\(1-\delta_p<1\), and its propagation range
is at most \(16p\). In
\(K^{-1}=u_p^{-1}\sum_{n\ge0}T^n\), every
term with \(16pn<d(A,B)\) has zero \(A,B\)
block. Summing the remaining geometric series
gives the first bound. The elementary inequality
\(1-\delta\le e^{-\delta}\) gives the second.
No projection onto a nonlocal gauge complement
occurs in this inverse series.
\end{proof}

This result concerns the local canonical detail
coordinates. It is not an inference from spatial
locality of individual one-step harmonic lifts,
and requires no bound on their unbounded products.
For example the full classical thermal curvature
covariance is
\[
             P=d_fE_jK_j^{-1}E_j^*d_f^*.
\]
It is an orthogonal projection onto the range
of \(d_fE_j\), since \(K_j=(d_fE_j)^*d_fE_j\).
Thus \(\|P\|\le1\). For fine plaquette carriers
\(A_f,B_f\), the same calculation and the
\(5p\) carrier radius of \(d_fE_j\) give
\begin{equation}
 \|P_{A_f}PP_{B_f}^*\|
 \le\min\left\{1,\;2^{44}p^{16}
   \exp\left[-\frac{(d(A_f,B_f)-10p)_+}{2^{48}p^{17}}\right]\right\}.
 \label{eq:f15v12:curvature-decay}
\end{equation}
All these statements concern the zero-final-output
classical Gaussian.

\subsection{A quadratic-source prediction bound for exterior physical pins}

Let \(v\) now have the centered Gaussian law
with precision \(\beta K_j\), at final output
\(C=I\). Let \(Y=\sqrt\beta\,L v_J\) be a
vector of central fine curvatures, where \(L\)
is the appropriate restriction of \(d_fE_j\).
For the source support used in V9 its detail
carrier \(A\) lies within radius \(12p\).
In particular \(\|L\|\le4p^4\).

For a real symmetric matrix \(T_0\) on these
curvature components, including colours, put
\[
 S_{T_0}=\frac12\left(Y^{\mathsf T}T_0Y
                      -\operatorname{Tr}(T_0\operatorname{Cov}Y)\right).
\]
Diagonal plaquette weights give the classical
quadratic form of the central action-density
source. Here \(\|T_0\|_{\rm HS}\) is the
ordinary Hilbert--Schmidt matrix norm.

\begin{theorem}[Classical quadratic-source screening]
\label{f15v12:gaussian-screen}
Let \(\mathcal P\) be any sigma-algebra
generated by linear functions of \(v_O\).
This includes V11's induced physical linear
pins from the original \(\eta,w,B\), with
final output fixed to the identity. Define
\begin{equation}
 \varepsilon_{p,r}
   =2^{132}p^{48}
       \exp\left\{-\frac{r-88}{2^{47}p^{16}}\right\}.
 \label{eq:f15v12:epsilon}
\end{equation}
Then
\begin{equation}
 \left\|\mathbb E[S_{T_0}\mid\mathcal P]\right\|_2
       \le\frac{\|T_0\|_{\rm HS}}{\sqrt2}
                               \min\{1,\varepsilon_{p,r}\}.
 \label{eq:f15v12:gaussian-screen}
\end{equation}
The norm and conditional expectation here are
under this specified classical Gaussian, not
under the full compact Wilson law.
\end{theorem}

\begin{proof}
Write \(K=K_j\). Its Gaussian conditional
mean in \(J\) given \(v_O\) is
\(-K_{JJ}^{-1}K_{JO}v_O\), and its conditional
covariance is \((\beta K_{JJ})^{-1}\).
Let \(B_J\) be the rows of \(K_{JO}\) that
are nonzero. They lie within \(16p\) of
the exterior of \(J\). Thus
\(d(A,B_J)\ge(r-88)p\).
Use \(\|K_{JO}\|\le u_p\),
\(\operatorname{Cov}(v_O)\preceq(\beta h_p)^{-1}I\),
and Theorem~\ref{f15v12:inverse-decay}. For
\(M_Y=\mathbb E[Y\mid v_O]\) these imply
\[
 \|\operatorname{Cov}M_Y\|
 \le16p^8u_p^2h_p^{-3}
              \exp\left[-\frac{2(r-88)p}{2^{48}p^{17}}\right]
       =\varepsilon_{p,r}.
\]
Also \(\operatorname{Cov}M_Y\preceq
\operatorname{Cov}Y\preceq I\), by the
curvature projection identity above. Put
\(\Sigma_M=\operatorname{Cov}M_Y\).

The innovation \(Y-M_Y\) is independent
of \(v_O\) and has constant covariance.
Consequently
\[
 \mathbb E[S_{T_0}\mid v_O]
      =\tfrac12\left(M_Y^{\mathsf T}T_0M_Y
                                   -\operatorname{Tr}(T_0\Sigma_M)\right).
\]
Expanding fourth Gaussian moments into their
three pairings gives
\[
 \Var\bigl(\mathbb E[S_{T_0}\mid v_O]\bigr)
  =\tfrac12\operatorname{Tr}(T_0\Sigma_M T_0\Sigma_M)
  \le\tfrac12\|T_0\|_{\rm HS}^2
                             \min\{1,\varepsilon_{p,r}\}^2.
\]
Conditional expectation onto the smaller
sigma-algebra \(\mathcal P\) is an \(L^2\)
contraction, proving the bound.

To check its application to the original
linear pins, differentiate
Proposition~\ref{f15v12:interior}. Their linear
map is of the form \(D_Ov_O+D_gh\), with
no \(v_J\)-term. Quotienting its values by
the gauge range, as required in V11, still
leaves a function only of \(v_O\). The
possibly nonlocal projection in pin-value
space does not change that support statement.
\end{proof}

This proves a spatially quantitative classical
prediction estimate for the relevant quadratic
source, not just a variance bound for linear
links. It remains at the identity final
background and on the untruncated Gaussian.
In particular it is not yet the score-locality
quantity of V8 under the actual varying-output
compact law.

\subsection{Enlarging the collar without losing the compact good-event estimate}

The weak explicit inverse rate can be offset
by a polynomially larger collar. Take
\(r=p^{20}\), leaving the slowly growing
block size \(p\) of V7 unchanged. For \(p\ge2\),
\begin{equation}
 \varepsilon_{p,r}
     \le2^{132}p^{48}\exp(-p^4/2^{48}).
 \label{eq:f15v12:large-collar-gaussian}
\end{equation}
This is smaller than every fixed inverse
power of \(p\) as \(p\to\infty\).
The central source weights have
\(\|T_0\|_{\rm HS}=O(p^2)\), so that factor
does not spoil the conclusion.

For the actual compact trace mask of V9,
the box side is now
\(L=2p(r+5)\le4p^{21}\). Its previously
proved full-law thermal estimate gives
\begin{equation}
 \mathbb P_\mu(q\ne1)
 \le8192C_{\rm th}p^{84}
       \exp\left\{-\frac{\beta\rho_*^2}{4608\pi^2p^{84}}\right\}.
 \label{eq:f15v12:large-collar-mask}
\end{equation}
Indeed V9's exponent denominator is
\(18\pi^2L^4\), and its prefactor is
\(32C_{\rm th}L^4\). Along V7's selected
path, \(p\le u\le\log\beta\) and
\(\beta\ge e^u\), so this probability also
decays faster than every fixed inverse
power of \(p\). The required \(M\ge4L\)
holds eventually on the same large-volume
selection. This is an actual compact
probability estimate, but it does not turn
the Gaussian screening theorem into a
compact conditional estimate or control
the conditional derivative of the mask.

\subsection{The remaining comparison is now on a specified reduced law}

The next analytic object is
\eqref{eq:f15v12:interior-law}, not a density
with its gauge directions or Jacobians
informally discarded. One needs stability
and locality of its measured potential for
the actual small-field boundary data, plus
control of the complement and the mask
insertion. In particular, the powers of
\(p\) in its nonlinear derivative bounds
must be compared with the available chart
radius and \(\beta\); the classical inverse
bound alone does not do that.

The exact interior screening identity also
shows what the outer compatibility problem
can and cannot do. It can change the
posterior distribution of the exterior
variables and final background. It does not
impose hidden constraints inside \(J\).
That distinction permits a local reduced-law
analysis while retaining the original pins.
An interacting continuum construction and
its positive physical Hamiltonian gap are
still well beyond the established results.

\subsection{Checks and exact preservation}

The certificate is
\begin{center}\small
\path{scripts/verify_ym_f15_v12_local_reduced_gaussian.py}.
\end{center}
It canonicalizes and exactly reconstructs
720 and 12,240 scalar detail entries at
depths one and two. It constructs all
784 columns of a primitive inverse matrix,
checks its rooted right-inverse identity,
and obtains the exact maximum absolute
row and column sums \(22\) and eight.
It also checks a rational Gaussian
regression and quadratic prediction
variance, finite propagation of inverse
series terms, and 24 dyadic norm, support,
and collar comparisons. These supplement
the general proofs; they do not test the
nonlinear compact conditional decay claim.

The certificate has 8,116 bytes and SHA--256
\begin{center}\scriptsize\ttfamily
AF8B110EF06D800FB2352C63919210116E773D32B5A0A52AA7E1B76F3A880730.
\end{center}
The predecessor V11 has 249,469 bytes and SHA--256
\begin{center}\scriptsize\ttfamily
81A6BEC9AF7DC734763FB1F4DB9446DB23F9C1FDB4C587819C430F714A015DBA.
\end{center}
Reversing this terminal insertion and the
title change recovers V11 exactly. V12
replaces it as the sole active main note.
The 11 protected sources and 54 earlier
protected certificates remain unchanged.
The next companion checkpoint is V15.

\begin{firewall}
V12 proves an exact simultaneous compact
history gauge reduction, an unrestricted
reduced interior conditional law with all
measured factors retained, and spatial
quadratic-source screening for its
zero-background classical Gaussian.
It does not establish that screening under
the full nonlinear Wilson law, a closed
interacting RG trajectory, continuum
reconstruction, or the Yang--Mills mass gap.
\end{firewall}

\section{V13: nonlinear convexity and framed chart entry}
\label{f15v13:section}

V12 supplied a gauge-reduced compact interior law
and a spatial estimate for its classical Gaussian.
Here the potential is no longer replaced by that
Gaussian. We prove polynomial all-depth derivative
bounds for the actual measured potential, obtaining
a uniformly convex shrinking chart and spatial
covariance bounds on its normalized restriction.
We then address entry into that chart using a frame
determined by the pinned near output. A full-law
thermal event gives a quantitative \(L^2\) transfer
of conditional source means to the restricted
\emph{nonlinear} law.

The one-step convex chart of f14 V40 and the
fixed-depth charts of f14 V42 are not being
repeated. The new ingredients are polynomial
control in the growing block size, V12's
pin-compatible canonical interior, and a measured
entry/normalization argument for that interior.
The final exterior-oscillation estimate under the
actual varying-output law remains to be assembled.

\subsection{The actual potential in a fixed local chart}

Use V12's canonical details and its interior
\(J\), with \(r\ge128\). Write \(v_i=e^{x_i}\)
for its interior groups, and collect the retained
and exterior coordinates that occur in \(W_J\)
into a parameter vector \(b\), also in exponential
coordinates. Every such parameter has anchor
within \(16p\) of \(J\). Its retained coarse
coordinates belong to \(N\), since their anchors
have radius at most \((r-44)p<rp\).

Let \(F_J(x;b)\) be the sum of the Wilson
plaquette terms that involve \(J\), without
the factor \(\beta\). Let
\begin{equation}
 \ell_J(x;b)=\sum_{(i,e)\,\mathrm{touching}\,J}
                  \log j_{i,e}(x;b)
                       +\sum_{i\in J}\log J_G(x_i),
 \qquad V_J=\beta F_J-\ell_J,
 \label{eq:f15v13:potential}
\end{equation}
where, up to a constant irrelevant to normalization,
\(J_G(x)=(\sin|x|/|x|)^2\). The zero value is
defined by continuity. Thus \(V_J\) is the
potential relative to ordinary Lebesgue measure
on the Lie-algebra coordinates. Exterior-only
terms cancel from this conditional normalization;
no interior pivot or Haar term is omitted.
Coarse Haar is not an extra interior density.

\subsection{A finite primitive constant gives polynomial all-depth bounds}

We specify the constants through finite-dimensional
derivative bounds, rather than claim numerical
enclosures that have not been computed. Put
\(a_0=1/8\). Embed \(\SU(2)\) as unit
quaternions and choose smooth extensions to a
fixed neighborhood of the compact domains of
the following finite list of primitive maps:
the canonical one-step inverse, a Wilson
plaquette term, an inverse-pivot logarithmic
Jacobian, and the input exponential maps.
Include \(\log J_G\) on \(|x|\le a_0\).
Such extensions exist because the compact pivot
inverse is globally smooth and its Jacobian
is strictly positive. All arities are fixed.

For a scalar component \(f\), use the derivative
sum \(\sum_{a_1,\ldots,a_k}|\partial_{a_1\cdots a_k}f|\).
For a vector map take its maximum over output
components. Choose a finite number
\(A_{\rm prim}\ge2\) bounding these sums
through order three for every map on the
specified neighborhoods, and put
\begin{equation}
 B_0=5A_{\rm prim},\qquad
 \nu=5+\lceil3\log_2B_0\rceil,\qquad
 C_0=2^{20}B_0^6,\qquad D_p=C_0p^\nu.
 \label{eq:f15v13:constants}
\end{equation}
Increasing the chosen primitive bound is harmless.
It depends only on the fixed completed block,
cutoff and coordinate conventions, not on depth,
volume, collar radius, or \(\beta\).

\begin{proposition}[Uniform local derivative rows]
\label{f15v13:jets}
For all required input group logarithms of norm
at most \(a_0\), the following bounds hold,
with scalar coordinate indices including the
parameters \(b\):
\begin{align}
 \sup_a\sum_b|\partial_{ab}F_J|&\le D_p,&
 \sup_a\sum_{b,c}|\partial_{abc}F_J|&\le D_p,\notag\\
 \sup_a\sum_b|\partial_{ab}\ell_J|&\le D_p.
 \label{eq:f15v13:derivative-rows}
\end{align}
In particular, if \(\|(x,b)\|_{\max}\le a\le a_0\),
then
\begin{equation}
 \|D_x^2F_J(x;b)-(K_j)_{JJ}\|\le D_pa.
 \label{eq:f15v13:hessian-stability}
\end{equation}
Here the maximum norm is the maximum Euclidean
norm of an input group logarithm. The same
bounds hold for any subcollection of the
local terms. Hessian entries still have
anchor range at most \(16p\).
\end{proposition}

\begin{proof}
Let \(s_k\) bound the order-\(k\) derivative
row sums of a composed output. The chain rule
for a further primitive map bounds the new
rows by
\[
 A_{\rm prim}s_1,\qquad
 A_{\rm prim}(s_2+s_1^2),\qquad
 A_{\rm prim}(s_3+3s_1s_2+s_1^3).
\]
Fresh independent inputs have the bounded
exponential jets already included in
\(A_{\rm prim}\). Taking the maximum of
\(1,s_1,s_2^{1/2},s_3^{1/3},A_{\rm prim}\)
shows that one layer increases this maximum
by at most \(5A_{\rm prim}\). There are
at most \(j\) inverse layers and two further
input/observable layers. Thus each scalar
local term has its order-\(k\) derivative
sum bounded by \(B_0^{k(j+2)}\), \(k\le3\).
This is a derivative calculation along the
compact group maps; no logarithm of an
intermediate output is required.

A term that can involve a prescribed input
coordinate has its anchor within \(5p\)
of that coordinate. There are at most
\(6(10p+1)^4\) fine plaquette terms and
\(4j(10p+1)^4\) inverse-Jacobian terms in
this bound, together with its one exponential
Haar term. This number is at most
\(2^{20}p^5\), using \(j\le p\).
The estimate overcounts higher-level sites,
which is safe and avoids any volume dependence.
Multiplying this incidence bound by the
scalar derivative sum, and using
\(B_0^{3j}=p^{3\log_2B_0}\), proves
\eqref{eq:f15v13:derivative-rows}.

Integrate the third derivatives along the
segment from zero to \((x,b)\). Their row
sum bounds the change of the symmetric
interior Hessian in operator norm by
\(D_pa\). At zero the intrinsic classical
Hessian is \((K_j)_{JJ}\), by V12. Inputs
outside the displayed carriers do not occur
in these derivatives and need not be small.
This proves \eqref{eq:f15v13:hessian-stability}.
Pointwise exponentiation does not enlarge
the original term supports.
\end{proof}

The constants in this proposition are defined
finite constants, not numerical estimates of
the particular cutoff's derivative suprema.
The certificate below checks the algebraic
majorants and exponents, not those suprema.

\subsection{Convexity of the normalized restricted nonlinear law}

Retain \(h_p=2^{-40}p^{-8}\), \(u_p=16p^8\)
from V12 and define
\begin{equation}
 a_p=\frac{h_p}{8D_p},\qquad
 \beta_p=\frac{8D_p}{h_p}=a_p^{-1},\qquad
 \mathcal D_p=\prod_{i\in J}\{x_i\in\mathbb R^3:|x_i|<a_p\}.
 \label{eq:f15v13:chart}
\end{equation}
All these radii are below \(a_0\).

\begin{theorem}[Measured strong convexity on a fixed domain]
\label{f15v13:convexity}
If all required exterior and retained group
logarithms satisfy \(|b_i|\le a_p\), and
\(\beta\ge\beta_p\), then on
\(\overline{\mathcal D_p}\),
\begin{equation}
 \frac34\beta h_pI\preceq D_x^2V_J(x;b)
                           \preceq2\beta u_pI,
 \qquad
 \sup_a\sum_b|\partial_{ab}V_J|\le2\beta D_p.
 \label{eq:f15v13:convexity}
\end{equation}
The probability
\begin{equation}
 d\nu_{p,b}^{\rm box}(x)
  =\frac{1_{\mathcal D_p}(x)e^{-V_J(x;b)}\,dx}
          {\int_{\mathcal D_p}e^{-V_J(y;b)}\,dy}
 \label{eq:f15v13:restricted-law}
\end{equation}
is exactly V12's compact interior law conditioned
on this exponential-chart event. It is not
an extension of that potential to all Lie
algebra space and is not a Gaussian law.
\end{theorem}

\begin{proof}
The classical Hessian changes by at most
\(D_pa_p=h_p/8\). Its measured correction
has norm at most \(D_p\le\beta h_p/8\).
The lower bound is therefore
\(\beta(7h_p/8)-\beta h_p/8=3\beta h_p/4\).
The upper bound follows from
\(\beta(u_p+h_p/8)+\beta h_p/8\le2\beta u_p\).
The absolute row bound uses \(\beta\ge1\)
and \eqref{eq:f15v13:derivative-rows}.
Exponential coordinates are injective on
these balls, and their Haar density is
already included in \(\ell_J\). Restricting
and renormalizing the compact density gives
\eqref{eq:f15v13:restricted-law} exactly.
\end{proof}

\subsection{Spatial covariance with the hard boundary retained}

Put
\begin{equation}
 \alpha_p=\frac{h_p}{512pD_p}
            =2^{-49}C_0^{-1}p^{-(\nu+9)}.
 \label{eq:f15v13:alpha}
\end{equation}
For an observable, its derivative carrier
is the set of group indices on which its
gradient can be nonzero.

\begin{theorem}[Nonlinear covariance and fixed-domain boundary response]
\label{f15v13:covariance}
Under the hypotheses of
Theorem~\ref{f15v13:convexity}, smooth observables
with derivative carriers \(A,B\) satisfy
\begin{equation}
 |\operatorname{Cov}_{\nu^{\rm box}}(F,G)|
 \le\frac4{\beta h_p}e^{-\alpha_p d(A,B)}
       \|\nabla F\|_{L^2(\nu^{\rm box})}
       \|\nabla G\|_{L^2(\nu^{\rm box})}.
 \label{eq:f15v13:covariance}
\end{equation}
No boundary condition is imposed on \(F,G\).
Also
\(\Var_{\nu^{\rm box}}F\le
2(\beta h_p)^{-1}\mathbb E_{\nu^{\rm box}}|\nabla F|^2\).
For an exterior scalar coordinate \(b_k\),
\begin{equation}
 \partial_{b_k}\mathbb E_{\nu^{\rm box}}F
   =\mathbb E_{\nu^{\rm box}}\partial_{b_k}F
       -\operatorname{Cov}_{\nu^{\rm box}}(F,\partial_{b_k}V_J).
 \label{eq:f15v13:response}
\end{equation}
If \(F\) has no direct dependence on \(b_k\),
then
\begin{equation}
 |\partial_{b_k}\mathbb E_{\nu^{\rm box}}F|
 \le\frac{8D_p}{h_p}
        e^{-\alpha_p(d(A,\{k\})-16p)_+}
                  \|\nabla F\|_{L^2(\nu^{\rm box})}.
 \label{eq:f15v13:response-bound}
\end{equation}
Every derivative in these formulas is of
the full potential \eqref{eq:f15v13:potential}.
\end{theorem}

\begin{proof}
Here is the boundary-aware resolvent argument.
For the density \(e^{-V_J}\) on the product
of balls, use the reflecting scalar generator
\(L_0=-\Delta+\nabla V_J\cdot\nabla\).
Its one-form operator has interior expression
\(L_1=L_0\otimes I+D^2V_J\), with absolute
boundary conditions. Its quadratic form is
\[
 \int|\nabla u|^2\,d\nu^{\rm box}
 +\int u^{\mathsf T}(D^2V_J)u\,d\nu^{\rm box}
 +\sum_i\int_{\partial_i\mathcal D_p}
                          \mathrm{II}_i(u_i,u_i)\,d\nu_\partial.
\]
The trace of \(u_i\) is tangential on the
\(i\)-th ball boundary, and its outward
second fundamental form \(\mathrm{II}_i\)
is nonnegative. On a product, each boundary
term acts only on that ball's three-component
block. These forms are obtained by integration
by parts on the product of the absolute
one-ball complexes and then closed; the
intersections of faces contribute no additional
surface term. The smooth positive weight
does not change their form domains.

The scalar/one-form intertwining gives
\[
 \operatorname{Cov}(F,G)
       =\langle\nabla F,L_1^{-1}\nabla G\rangle_{L^2(\nu^{\rm box})}.
\]
One may derive it by solving the mean-zero
Neumann equation \(L_0\phi=G-\mathbb EG\)
and using \(L_1\nabla\phi=\nabla G\).
For the smooth-boundary weighted-form identity
and its absolute realization, see
\href{https://arxiv.org/abs/1607.08714}
{D. Le Peutrec, \emph{On Witten Laplacians and
Brascamp--Lieb's inequality on manifolds with boundary}},
Theorem 1.1 and equation (4.8). The product
construction above specifies the realization
used here, including its corners.

Set \(\mu_p=\beta h_p/2\). The displayed
form is bounded below by \(\mu_p\|u\|_2^2\).
Multiply each group block of \(u\) by the
constant weight
\(w_i=\exp(\alpha_p d(i,B))\). This operation
preserves the absolute form domain, commutes
with the derivative part, and commutes with
each block boundary form. Only the Hessian
is changed by conjugation. Its nonzero entries
have range at most \(16p\). The absolute
row and column sums of the conjugation error
are consequently at most
\[
 (e^{16p\alpha_p}-1)\,2\beta D_p
   \le 2(16p\alpha_p)\,2\beta D_p
   =\beta h_p/8\le\mu_p/2.
\]
The inequality uses \(16p\alpha_p\le1/32\).
Schur's bound and the coercive form therefore
give norm at most \(2/\mu_p=4/(\beta h_p)\)
for the conjugated inverse. Its weights are
one on \(B\) and at least
\(e^{\alpha_p d(A,B)}\) on \(A\).
Insert these weights into the covariance
identity to obtain \eqref{eq:f15v13:covariance}.
The unweighted inverse gives the stated
variance bound. Form closure and smooth
approximation extend the calculation to
the indicated observables without requiring
their normal derivatives to vanish.

The domain \(\mathcal D_p\) is independent
of \(b\). Differentiating its numerator and
denominator gives \eqref{eq:f15v13:response}.
The interior derivative carrier of
\(\partial_{b_k}V_J\) is within \(16p\)
of \(k\), and its gradient norm is at most
\(2\beta D_p\) by the joint derivative-row
bound. Substitution proves
\eqref{eq:f15v13:response-bound}.
\end{proof}

This theorem is a nonlinear measured-chart
result. Strong convexity has not been asserted
on the whole compact group product, nor has
the hard chart boundary been erased.

\subsection{Why the original coarse coordinates cannot supply chart entry}

\begin{proposition}[A gauge-invariant good event does not fix a coarse frame]
\label{f15v13:frame-obstruction}
There are flat fine configurations with every
rooted fine link equal to \(I\), and hence
with the V9 mask equal to one at any positive
threshold, for which a literal final coarse
link equals \(-I\). More generally, under
the full finite-volume Wilson law the marginal
of any final coarse link with distinct endpoints
is Haar, even after conditioning on a positive-
probability gauge-invariant event.
\end{proposition}

\begin{proof}
Take \(U_{xy}=h_x^{-1}h_y\), assigning different
values \(I,-I\) to two adjacent final roots.
Every fine plaquette and every rooted fine
edge word is the identity. Every rooted block
path telescopes to its prescribed endpoint
values, so all its relative logarithms vanish.
Inductively \(C^{(j)}_{xy}=h_{px}^{-1}h_{py}\),
which is \(-I\) on the chosen edge.

For the marginal assertion, a fine gauge
transformation at one coarse endpoint left
translates that coarse link arbitrarily,
while preserving the Wilson law and the
conditioning event. Its marginal is therefore
the unique left-invariant probability on
\(\SU(2)\). In particular the probability
of a logarithmic ball of radius \(a\le\pi\)
is
\(\frac2\pi\int_0^a\sin^2t\,dt\), at most
\(2a^3/(3\pi)\). Small literal coarse links
cannot be inferred from a high-probability
gauge-invariant event.
\end{proof}

\subsection{A frame determined only by the pinned near output}

Let \(\mathcal K\) be the fine box of V9,
of side \(L=2p(r+5)\). Use its common
ordered-root frame and let \(R_{\mathcal K}\)
be the largest geodesic angle of a rooted
fine link in this box. On the final coarse
box of radius \(r-20\), choose the ordered
coordinate tree from its lower corner.
All its links belong to \(N\). Its root
path products define a coarse gauge frame
\(H_N\), depending only on \(N\). Extend
it to the fine vertices by making it constant
on each final \(p\)-cell; values outside
the needed patch may be assigned arbitrarily.
This lift preserves the simultaneous history
tree slice. A superscript \(H_N\) below
denotes that endpoint gauge transformation.

\begin{theorem}[Quantitative framed entry]
\label{f15v13:entry}
If \(R_{\mathcal K}\le\delta\), then all
canonical detail and retained groups required
by \(W_J\), after applying \(H_N\), have
geodesic angle at most
\begin{equation}
                         32r p^4\delta.
 \label{eq:f15v13:entry-bound}
\end{equation}
In particular, for
\begin{equation}
             \delta_{p,r}=\frac{a_p}{64r p^4},
 \label{eq:f15v13:fine-threshold}
\end{equation}
the required boundary data and all interior
groups lie in the half-radius chart
\(|b_i|,|x_i|\le a_p/2\). The frame is
constant when varying the interior or its
canonical exterior details at fixed \(N\).
\end{theorem}

\begin{proof}
Work first in the common fine root frame,
where every needed fine link has angle at
most \(\delta\). A primitive root word has
two links; a nonroot word has at most ten.
Thus its relative word has angle at most
\(12s\) if the input angles are at most
\(s\). The cutoff logarithm has norm no
greater than this angle. The completed
output angle is at most
\(2s+(14/16)12s\le16s\).
Iteration gives \(16^i\delta\) at level
\(i\), and hence \(p^4\delta\) at the
final level. These are compact triangle
inequalities, not Taylor estimates.

The level-\(i\) tree products have angles
at most \(4\cdot16^{i-1}\delta\). The
abstract-tree gauge transformation used in
V12 consequently has angle at most
\[
             4\sum_{i=1}^j16^{i-1}\delta
                       \le(4/15)p^4\delta
\]
at each needed fine vertex. After it is
applied, all intermediate edge groups, and
therefore the bare canonical details, have
angles at most \(2p^4\delta\).

Compare this fine-root-induced coarse frame
with \(H_N\). Their relative value at a
coarse vertex is, up to a fixed conjugation,
the product along its coarse tree path of
the fine-root-framed coarse links. Each
path has length at most \(8r\), so that
relative value has angle at most
\(8r p^4\delta\). Covariance and uniqueness
of the canonical slice imply that the
corresponding change of canonical details
is the endpoint transformation by these
values, constant on final cells. An internal
conjugation preserves angle; a boundary
group incurs at most the two endpoint
angles. Thus every required detail has
angle at most \((16r+2)p^4\delta\), and
a required retained link at most
\((16r+1)p^4\delta\). Both are bounded by
\eqref{eq:f15v13:entry-bound}.

All carriers used here lie inside
\(\mathcal K\). The required canonical
anchors have radius at most \((r-44)p\);
 their forward and gauge-reconstruction
 carriers are within a further \(10p\).
The coarse tree box has radius \((r-20)p\)
and its forward coarse-link carriers have
radius less than \(4p\). These are inside
the fine box of radius \((r+5)p\).
One may extend the fine root gauge outside
the box without affecting this calculation.
Finally \(H_N\) is constructed from fixed
near links only, proving the last assertion.
\end{proof}

For each fixed \(N\), the transformation of
canonical details is by Haar translations
or conjugations at fixed group indices.
The measured conditional law is covariant
under it. Thus the preceding chart and
response theorems apply in this frame with
no new interior Jacobian or frame derivative
when the exterior details vary. This does
not assert that the map \(N\mapsto H_N\)
has zero derivative when \(N\) itself varies.

\subsection{An actual compact entry probability}

The full-law rooted-link estimate of V8,
with a filling of at most \(3L\) plaquettes
per link, gives for any \(\delta>0\)
\[
 \mathbb P_\mu(R_{\mathcal K}>\delta)
      \le32C_{\rm th}L^4
              \exp\{-\beta\delta^2/(9\pi^2L^2)\}.
\]
Substitute \eqref{eq:f15v13:fine-threshold}
and use \(L\le4rp\). Define
\begin{equation}
 \epsilon_{p,r,\beta}
 =\min\left\{1,\;8192C_{\rm th}r^4p^4
      \exp\left[-\frac{\beta a_p^2}
                             {589824\pi^2r^4p^{10}}\right]\right\}.
 \label{eq:f15v13:entry-probability}
\end{equation}
Then the event \(R_{\mathcal K}\le\delta_{p,r}\)
has probability at least \(1-\epsilon_{p,r,\beta}\)
under the actual compact Wilson law.

In the \(H_N\)-frame, let \(\Gamma\) be
the event that all boundary/retained inputs
needed by \(W_J\) lie in their half-radius
charts, and let \(D\) denote the interior
event \(v_J^{H_N}\in\exp\mathcal D_p\).
The entry theorem proves
\begin{equation}
 \mathbb P_\mu(\Gamma^c)\le\epsilon_{p,r,\beta},
 \qquad \mathbb P_\mu(D^c)\le\epsilon_{p,r,\beta}.
 \label{eq:f15v13:entry-events}
\end{equation}
The event \(\Gamma\) is measurable with
respect to the final output and canonical
exterior details, not the interior.
This is not a claim that the original
literal coarse links are near \(I\).

\subsection{Normalized conditional means can now be compared in the full law}

Let \(Z=(C,v_O^{H_N})\), and write
\(m(Z)=\mathbb E_\mu[S_b\mid Z]\).
The normalized chart mean is
\[
 m_D(Z)=\frac{\mathbb E_\mu[S_b1_D\mid Z]}
                         {\mathbb P_\mu(D\mid Z)}.
\]
Its denominator is strictly positive at
every finite compact conditioning point.
On \(\Gamma\), this is exactly the
nonlinear mean under
\eqref{eq:f15v13:restricted-law} in the
specified near-output frame. Put
\[
 q_D(Z)=\mathbb P_\mu(D\mid Z),\qquad
 A=\Gamma\cap\{q_D\ge1/2\},\qquad
                         \widehat m(Z)=1_A m_D(Z).
\]

\begin{theorem}[Full-law conditional-mean transfer]
\label{f15v13:mean-transfer}
For every source with finite fourth moment,
\begin{equation}
       \|m-\widehat m\|_{L^2(\mu_Z)}
           \le5\|S_b\|_{L^4(\mu)}\,
                                \epsilon_{p,r,\beta}^{1/4}.
 \label{eq:f15v13:mean-transfer}
\end{equation}
Under V8's source hypotheses this is at most
\(5C_4a_b\epsilon_{p,r,\beta}^{1/4}\).
For any of the original pin statistics to
which V12's screening identity applies,
the same error bound holds after taking
conditional expectations of the two means
with respect to those pins.
\end{theorem}

\begin{proof}
Write \(\delta_D=1-q_D\) and
\(M_4=\|S_b\|_4\). Then
\(\mathbb E\delta_D=\mathbb P(D^c)\le\epsilon\),
where \(\epsilon=\epsilon_{p,r,\beta}\),
and \(\mathbb P(A^c)\le3\epsilon\).
On \(A\),
\[
 m-m_D=\mathbb E[S_b1_{D^c}\mid Z]
             -\frac{\delta_D}{q_D}\mathbb E[S_b1_D\mid Z].
\]
The first term has \(L^2\) norm at most
\(M_4\epsilon^{1/4}\), by conditional
contraction and H\"older's inequality.
For the second, use \(q_D^{-1}\le2\),
conditional Cauchy--Schwarz, and
\(\mathbb E\delta_D^4\le\epsilon\) to get
\[
 \|1_A(\delta_D/q_D)\mathbb E[S_b1_D\mid Z]\|_2
                       \le2M_4\epsilon^{1/4}.
\]
On \(A^c\), conditional Jensen and H\"older
give \(\|1_{A^c}m\|_2\le
M_4(3\epsilon)^{1/4}\le2M_4\epsilon^{1/4}\).
The triangle inequality proves the result.

V12's exact interior law is independent of
the gauge variables even with the original
pins included. Those pins are functions of
the exterior, final output, and gauge
variables. Its tower identity therefore
expresses their source means as conditional
expectations of \(m(Z)\). Conditional
expectation is an \(L^2\) contraction, so
the same bound applies to the replacement
by \(\widehat m(Z)\).
\end{proof}

This argument retains the chart normalizer
\(q_D\) and deals with its small values by
an explicit exceptional event. It does not
divide by an uncontrolled conditional good
probability or discard the exterior posterior
weights.

\subsection{What remains after this nonlinear step}

For \(p\le\log\beta\), both
\(\beta\ge\beta_p\) and the entry probability
requirements eventually hold if \(r\) is
bounded by a fixed power of \(\log\beta\).
All exponents and constants in
\eqref{eq:f15v13:constants} are fixed as the
regulator grows. A collar such as
\(r=\lceil p^{\nu+8}(\log\beta)^2\rceil\)
is compatible with the thermal estimate and
gives an exponent proportional to
\((\log\beta)^2\) in the spatial factor
\(\alpha_p rp\). The previous choice
\(r=p^{20}\) need not suffice for this
new nonlinear rate.

There is now a measured nonlinear covariance
and boundary-response estimate, and a full-law
\(L^2\) replacement of source means by those
of that convex restriction. The next step
is to sum the exterior responses for the
actual action-density source at fixed \(N\),
with its gradient norms, number of boundary
coordinates, and all exceptional events
tracked. No dimension-free oscillation bound
for that source is asserted merely from
\eqref{eq:f15v13:response-bound}. In particular
its powers of \(\beta\) must be accounted
for before comparing them to the spatial
factor. The original mask insertion is not
silently identified with the chart indicator.

Even completion of this conditional locality
step would still leave interacting RG closure,
continuum reconstruction, and the physical
mass-gap assertion to be proved.

\subsection{Checks and predecessor preservation}

The certificate is
\begin{center}\small
\path{scripts/verify_ym_f15_v13_measured_convexity.py}.
\end{center}
It checks 78 third-order composition majorants,
64 bounds for incidence, convexity, and spatial
weights, exact matrix bounds, and commutation of
block weights with a ball's boundary shape
matrix. It also checks 32 framed-entry/tail
scales and eight normalized conditional-mean
transfer examples. An actual central
\(\{I,-I\}\subset\SU(2)\) pure gauge has
1,536 identity plaquettes, 1,024 identity
rooted fine links, and 16 coarse links
equal to \(-I\); all 1,024 rooted coarse
path words are checked. The primitive
derivative supremum itself is not numerically
enclosed by this certificate.

The certificate has 8,353 bytes and SHA--256
\begin{center}\scriptsize\ttfamily
1F6007618355996A3B069F41B953416D3753301AE74EFA04884039354C745DC7.
\end{center}
The predecessor V12 has 270,953 bytes and SHA--256
\begin{center}\scriptsize\ttfamily
7BD70CF459DA5F8B1B891989E289A6847A04781597F3BDBE4E4E0CACC7B79DF7.
\end{center}
Reversing this terminal insertion and title
change recovers V12 exactly. V13 replaces
it as the sole active main note. The 11
protected sources and 55 earlier protected
certificates remain unchanged. The next
companion checkpoint is V15.

\begin{firewall}
V13 proves polynomial derivative control,
measured convexity and spatial covariance
on a specified shrinking chart, framed
entry with a full compact thermal probability
bound, and an \(L^2\) transfer of normalized
conditional source means. It does not yet
prove the required full-law source
oscillation estimate, interacting RG
closure, a continuum theory, or a positive
physical Yang--Mills mass gap.
\end{firewall}

\section{V14: full-law source locality}
\label{f15v14:section}

We now complete the source-oscillation step left open
in V13. The conclusion is an estimate under the
actual compact Wilson law, not just its Gaussian
approximation or its chart restriction. For V7's
growing indicator source, the retained score can
be approximated in \(L^2\) by a function of a
growing near-output window. We give both the
optimal conditional-mean approximation and a
specified gauge-invariant local integral.

The new ingredient is a comparison of two actual
exteriors conditional on the same near output.
The exterior coordinates are not assumed
independent. The normalizer of the restricted
interior law and the posterior law of the exterior
are retained throughout. This removes the
particular full-law locality bottleneck isolated
in V7--V13, for the source and scales below.
It does not give a regulator-uniform analytic
interaction norm, an interacting RG contraction,
or a continuum Hamiltonian gap.

\subsection{Setup and the finite-regulator assertion}

Use the selected Wilson laws of V7--V8,
the canonical interior of V12, and the frame
and constants of V13. Write \(p=2^j\ge2\).
Let \(C=(N,F)\) be the entire final output,
where \(N\) contains every final link whose
base lies in \([-r,r]^4\). Assume
\[
 r\ge176,\qquad L=2p(r+5),\qquad
 M\ge4L,\qquad \beta\ge\beta_p .
\]
As before, the box is placed without periodic
wrap-around. We reserve \(D\) for the interior
chart event, and call V7's set of fine source
bases \(\mathcal B_p\) in this section:
\[
 b_x=1_{\mathcal B_p}(x),\qquad
 |\mathcal B_p|=p^2(2p-1)^2\le4p^4 .
\]
This is exactly the source in
Theorem~\ref{f15v7:score}, not a new source
family. Its centred value is
\begin{equation}
 S_b=\beta\sum_f\omega_b(f)S_{W,f}
                 -\mathbb E_\mu\!\left[
                         \beta\sum_f\omega_b(f)S_{W,f}\right],
 \quad
 0\le\omega_b(f)\le1,\quad
 \sum_f\omega_b(f)=6|\mathcal B_p|.
 \label{eq:f15v14:source}
\end{equation}
The Wilson term satisfies \(0\le S_{W,f}\le2\).
Consequently
\begin{equation}
 \|S_b\|_\infty\le12\beta|\mathcal B_p|
                         \le48\beta p^4,
 \qquad M_4:=\|S_b\|_4\le C_4a_b .
 \label{eq:f15v14:source-norms}
\end{equation}
The latter is V8's bound, with its unchanged
\(C_4,a_b\); along the selected path \(a_b=O(p^2)\).

Put
\begin{equation}
 \omega_{p,r,\beta}
       =\beta r^4p^8
                   \exp\{-\alpha_p(r-88)p\},
 \qquad \epsilon=\epsilon_{p,r,\beta},
 \label{eq:f15v14:errors}
\end{equation}
where \(\alpha_p\) and \(\epsilon_{p,r,\beta}\)
are exactly \eqref{eq:f15v13:alpha} and
\eqref{eq:f15v13:entry-probability}.

\begin{theorem}[Locality of the actual retained source score]
\label{f15v14:locality}
Under these hypotheses,
\begin{equation}
 \mathcal L_{b,r}
 :=\left\|\mathbb E_\mu[S_b\mid C]
                -\mathbb E_\mu[S_b\mid N]\right\|_2
 \le \omega_{p,r,\beta}
                   +8C_4a_b\epsilon_{p,r,\beta}^{1/4}.
 \label{eq:f15v14:locality}
\end{equation}
The norm is for the actual full compact law.
No conditional good-event probability is
assumed bounded below at every exterior.
\end{theorem}

We prove the theorem after establishing the
uniform oscillation estimate that it needs.

\subsection{Summing the measured exterior responses}

Let \(J\) consist of canonical details with
anchors of radius at most \((r-60)p\),
as in V12. In the \(H_N\)-frame the
normalized chart mean is
\[
 m_D(Z)=\mathbb E_{\nu^{\rm box}_{p,b(Z)}}S_b,
                  \qquad Z=(C,v_O^{H_N}),
\]
whenever the required boundary inputs are in
V13's event \(\Gamma\). The source has
canonical derivative carrier \(A_s\) in
the box of radius \(12p\), strictly inside
\(J\). Thus it has no direct dependence
on the exterior details. Its dependence
on the required retained coordinates is
allowed; these are fixed when \(N\) is fixed.

All retained parameters occurring in \(W_J\)
have anchors of radius at most \((r-44)p\)
and therefore belong to \(N\).
The source satisfies the same condition.
It follows from the exact conditional density
of V12 that \(m_D\) is a function of \(N\)
and only those exterior canonical details
which enter \(W_J\). In particular it has
no direct dependence on \(F\).
Exterior-only terms cancel in this normalized
interior density, not in the exterior posterior.

\begin{proposition}[Oscillation at fixed near output]
\label{f15v14:oscillation}
If \(Z,Z'\in\Gamma\) have the same \(N\),
then
\begin{equation}
                  |m_D(Z)-m_D(Z')|
                              \le\omega_{p,r,\beta}.
 \label{eq:f15v14:oscillation}
\end{equation}
The statement applies to every pair of such
boundary values, not just almost every pair
under a product exterior law.
\end{proposition}

\begin{proof}
First control the source gradient, including
its powers of \(\beta\). The proof of
Proposition~\ref{f15v13:jets} takes absolute
values term by term, so its Wilson second-
derivative row bound \(D_p\) remains valid
for weights of magnitude at most one.
At the all-identity inputs each Wilson
term has zero first derivative. Integrating
its weighted Hessian along the straight
segment from zero to the joint interior/
parameter vector therefore gives, on the
chart used here,
\[
           |\partial_a S_b|\le\beta D_pa_p
                                      =\beta h_p/8 .
\]
The source-centering constant does not enter
this derivative. The argument uses the
joint row estimate, so the retained
parameters need not be zero.

There are at most 45 detail groups per
level and anchor. Counting their three
scalar components and overcounting every
level by the fine grid gives
\[
 n_s\le135j(24p+1)^4\le2^{26}p^5 .
\]
Here \(j\le p\) and \(24p+1\le25p\)
are sufficient. Hence, uniformly over
the required parameters,
\begin{equation}
 \|\nabla_x S_b\|_{L^2(\nu^{\rm box})}
                 \le2^{10}\beta h_pp^{5/2}.
 \label{eq:f15v14:gradient}
\end{equation}
All these coordinates are in \(J\).

An exterior detail that occurs in \(W_J\)
has anchor of radius at most \((r-44)p\).
The number of its possible scalar components
is at most
\[
 n_\partial\le135j(2(r-44)p+1)^4
                              \le2^{14}r^4p^5 .
\]
For the last inequality use
\(2(r-44)p+1\le2rp\).
Its distance from \(A_s\) is at least
\((r-72)p\). The further \(16p\) loss
in \eqref{eq:f15v13:response-bound} gives
for each such scalar parameter
\begin{equation}
 |\partial_{b_k}m_D|
       \le2^{13}\beta D_pp^{5/2}
                         e^{-\alpha_p(r-88)p}.
 \label{eq:f15v14:scalar-response}
\end{equation}
This is a derivative of a normalized
nonlinear integral; the covariance formula
already includes its normalizer.

At fixed \(N\), \(H_N\) and all retained
parameters are fixed. For \(Z,Z'\in\Gamma\),
each required exterior logarithm has norm
at most \(a_p/2\). Join corresponding
logarithms by their straight line segments.
They stay in these half-radius balls;
each scalar endpoint difference is at most
\(a_p\). V12's canonical variables form
a genuine product of groups, so these
segments impose no extra pin-compatibility
constraint on the interior. Parameters
not occurring in \(W_J\) can be ignored.

Integrate \eqref{eq:f15v14:scalar-response}
along the simultaneous parameter segment
and sum the scalar bounds. The result is
at most
\begin{align*}
 &2^{27}\beta D_pa_pr^4p^{15/2}
                              e^{-\alpha_p(r-88)p}\\
 &\hspace{12mm}
 =2^{24}\beta h_pr^4p^{15/2}
                              e^{-\alpha_p(r-88)p}
 =2^{-16}\beta r^4p^{-1/2}
                              e^{-\alpha_p(r-88)p}.
\end{align*}
This is bounded by \(\omega_{p,r,\beta}\).
No derivative of \(N\mapsto H_N\) is taken.
\end{proof}

\subsection{Conditional copies under the actual exterior posterior}

\begin{proof}[Proof of Theorem~\ref{f15v14:locality}]
Use exactly V13's variables
\[
 m(Z)=\mathbb E_\mu[S_b\mid Z],\quad
 q_D(Z)=\mathbb P_\mu(D\mid Z),\quad
 A=\Gamma\cap\{q_D\ge1/2\},\quad
 \widehat m=1_A m_D .
\]
They satisfy
\begin{equation}
 \mathbb P(A^c)\le3\epsilon,\qquad
 \|m-\widehat m\|_2\le5M_4\epsilon^{1/4}.
 \label{eq:f15v14:transfer-input}
\end{equation}
Conditional Jensen inside the normalized
chart also gives
\[
 |m_D|^4\le q_D^{-1}
                         \mathbb E[|S_b|^4 1_D\mid Z].
\]
Multiplying by \(1_A\) and integrating proves
\begin{equation}
                   \|\widehat m\|_4^4\le2M_4^4 .
 \label{eq:f15v14:truncated-fourth}
\end{equation}

Now sample \(Z,Z'\) independently
\emph{conditional on \(N\)}, using the actual
regular conditional exterior law twice.
Such laws exist on these finite-dimensional
compact spaces. Each marginal is the original
law of \(Z\); independence of the individual
boundary coordinates is neither required nor
asserted. The conditional-copy identity is
\[
 \|\widehat m-\mathbb E[\widehat m\mid N]\|_2^2
              =\tfrac12\mathbb E
                            |\widehat m(Z)-\widehat m(Z')|^2 .
\]
On \(A\cap A'\), Proposition~\ref{f15v14:oscillation}
bounds the difference by \(\omega_{p,r,\beta}\).
The complementary pair event has probability
at most \(6\epsilon\). By Minkowski and
\eqref{eq:f15v14:truncated-fourth},
\[
 \|\widehat m(Z)-\widehat m(Z')\|_4
                                  \le2\,2^{1/4}M_4 .
\]
H\"older on the complementary event yields
\begin{equation}
 \|\widehat m-\mathbb E[\widehat m\mid N]\|_2^2
 \le\tfrac12\omega_{p,r,\beta}^2
                           +2\sqrt{12}\,M_4^2\epsilon^{1/2}.
 \label{eq:f15v14:pair-bound}
\end{equation}
Taking the square root bounds this by
\(\omega_{p,r,\beta}+3M_4\epsilon^{1/4}\)
in norm. The conditional-centering operator
\(I-\mathbb E[\cdot\mid N]\) is an orthogonal
projection on \(L^2(\mu_Z)\). Applying it
to \eqref{eq:f15v14:transfer-input} gives
\[
 \|m-\mathbb E[m\mid N]\|_2
           \le\omega_{p,r,\beta}+8M_4\epsilon^{1/4}.
\]
Finally \(Z\) contains \(C\), and \(C\)
contains \(N\). The tower property says
\[
 \mathbb E[m\mid C]=\mathbb E[S_b\mid C],
 \qquad \mathbb E[m\mid N]=\mathbb E[S_b\mid N].
\]
Projecting the preceding difference onto
\(C\) is a contraction. This proves the
theorem, using \(M_4\le C_4a_b\).
\end{proof}

The cutoff \(q_D\ge1/2\) has only been used
to control moments and the exceptional set;
it has not been differentiated. In particular,
the old smooth rooted mask from V9 is not
identified with \(1_D\), and no formula
for its derivative is needed for this
\(L^2\) conclusion. This is an alternative
route around that earlier mask-derivative
bottleneck, not a proof of the unproved
mask-specific assertion.

\subsection{A specified local gauge-invariant integral}

Conditional expectation on \(N\) is an optimal
local \(L^2\) approximation but still uses the
full law to define its marginal. The following
choice instead uses one explicitly specified
interior integral.

Let \(\Gamma_N\) require only that the
retained framed parameters entering \(W_J\)
lie in their half-radius balls. On
\(\Gamma_N\), set every required exterior
canonical detail to \(I\) in the \(H_N\)-frame
and define
\begin{equation}
 G_{b,p,r}(N)
  =1_{\Gamma_N}(N)\,
        m_D\bigl(N,\ v_{\partial O}^{H_N}=I\bigr).
 \label{eq:f15v14:local-integral}
\end{equation}
The value outside \(\Gamma_N\) is zero.
The displayed mean is the finite-dimensional
normalized integral \eqref{eq:f15v13:restricted-law},
with the actual Wilson and pivot terms and
the exponential Haar density. No values of
the unused far retained output or exterior
details are needed to define it.

\begin{proposition}[Explicit local approximation]
\label{f15v14:surrogate}
The function \(G_{b,p,r}\) depends only on \(N\),
is invariant under coarse gauge transformations,
and satisfies \(|G_{b,p,r}|\le48\beta p^4\).
Moreover,
\begin{equation}
 \begin{split}
 \|\mathbb E[S_b\mid C]-G_{b,p,r}(N)\|_2
 \le{}&\omega_{p,r,\beta}
                  +5C_4a_b\epsilon_{p,r,\beta}^{1/4}\\
      &+48\sqrt3\,\beta p^4\epsilon_{p,r,\beta}^{1/2}.
 \end{split}
 \label{eq:f15v14:surrogate}
\end{equation}
Replacing \(G_{b,p,r}\) by its centred version
\(G_{b,p,r}-\mathbb E G_{b,p,r}\) cannot
increase this error.
\end{proposition}

\begin{proof}
Locality follows from the parameter dependence
already proved. The uniform bound follows
from \eqref{eq:f15v14:source-norms} under
any normalized interior probability.
On \(A\), \(\Gamma_N\) holds and both the
actual and identity exterior parameters
belong to \(\Gamma\). Proposition~\ref{f15v14:oscillation}
therefore gives \(|\widehat m-G|\le\omega_{p,r,\beta}\).
On \(A^c\), \(\widehat m=0\), so
\(|\widehat m-G|\le48\beta p^4\).
Consequently
\[
 \|\widehat m-G\|_2
       \le\omega_{p,r,\beta}
                       +48\sqrt3\,\beta p^4\epsilon^{1/2}.
\]
Combine this with V13's mean transfer and
project onto \(C\). Since \(G\) is
\(C\)-measurable, this proves the error bound.
The target \(\mathbb E[S_b\mid C]\) has
mean zero, so orthogonal projection away
from constants proves the centering assertion.

For gauge invariance, use the convention
\(C_{xy}^g=g_xC_{xy}g_y^{-1}\).
The coarse tree products obey
\[
             H_{N^g}(x)=g_o H_N(x)g_x^{-1},
\]
where \(o\) is their common root.
Thus all the required framed groups
transform by the single conjugation
\(g_o(\,\cdot\,)g_o^{-1}\). The lift
constant on final cells preserves the
canonical history slice, as in V13.
Conjugation preserves \(\Gamma_N\), the
interior product of radial balls, and
the identity exterior. It also preserves
the Wilson source, measured potential,
and product integration measure.
Changing interior variables by this
conjugation leaves the normalized
integral unchanged.
\end{proof}

The indicator \(1_{\Gamma_N}\) is a hard
cutoff. We do not assert that this
particular \(G\) is smooth across its
boundary, nor that it has regulator-uniform
derivative or analytic interaction norms.
The source-centering constant is an
allowed scalar constant; it creates no
dependence on a distant output variable.

\subsection{A collar that pays every displayed cost}

\begin{theorem}[Superalgebraic locality along the selected path]
\label{f15v14:asymptotic}
On V7's slower selected path, retain its
\(p\to\infty\), \(p\le\log\beta\), and
\(\beta\le\log M\), and choose
\begin{equation}
 r=\max\{176,\lceil p^{\nu+8}(\log\beta)^2\rceil\}.
 \label{eq:f15v14:collar}
\end{equation}
The finite-regulator hypotheses above hold
eventually. For every fixed \(A>0\),
both errors in \eqref{eq:f15v14:locality}
and \eqref{eq:f15v14:surrogate} are
\(o(\beta^{-A})\), hence also \(o(p^{-A})\).
At the same time \(rp/M\to0\): the near
window does not become the entire output.
\end{theorem}

\begin{proof}
Put \(t=\log\beta\). For all sufficiently
large regulators,
\[
 p\le t,\qquad
 p^{\nu+8}t^2\le r\le2p^{\nu+8}t^2
                                  \le2t^{\nu+10}.
\]
Since \(\beta_p=2^{43}C_0p^{\nu+8}\),
the inequality \(\beta=e^t\ge\beta_p\)
eventually holds. Also \(L\le4rp\)
grows at most polynomially in \(t\),
whereas \(M\ge e^\beta=e^{e^t}\).
This gives \(M\ge4L\) and \(rp/M\to0\).

For \(r\ge176\), \(r-88\ge r/2\).
Thus
\[
 \alpha_p(r-88)p\ge c_0t^2,\qquad
                  c_0=2^{-50}C_0^{-1}>0 .
\]
Inserting the upper bounds for \(p,r\)
in \eqref{eq:f15v14:errors} gives
\begin{equation}
           \omega_{p,r,\beta}
                \le16e^t t^{4\nu+48}e^{-c_0t^2}.
 \label{eq:f15v14:omega-rate}
\end{equation}
This is \(o(e^{-At})\) for every fixed \(A\),
although the constants allow a very large
onset scale and no numerical onset is claimed.

For the exceptional probability, use
\(a_p=2^{-43}C_0^{-1}p^{-(\nu+8)}\)
in V13's expression. The estimates
\[
 r^4p^4\le16t^{4\nu+44},\qquad
 r^4p^{2\nu+26}\le16t^{6\nu+66}
\]
give
\begin{equation}
 \epsilon\le
 \min\left\{1,\,
   2^{17}C_{\rm th}t^{4\nu+44}
        \exp\left[-\frac{c_1e^t}{t^{6\nu+66}}\right]\right\},
 \quad
 c_1=\frac{2^{-90}}{589824\pi^2C_0^2}>0 .
 \label{eq:f15v14:epsilon-rate}
\end{equation}
Any positive fixed power of this expression
decays faster than \(e^{-At}\), even after
multiplication by a fixed power of \(\beta,p,r\).
Finally \(a_b=O(p^2)=O(t^2)\), and the
extra prefactor in the explicit-integral
bound is at most \(48\sqrt3\,e^tt^4\).
Equations \eqref{eq:f15v14:omega-rate} and
\eqref{eq:f15v14:epsilon-rate} therefore
prove the two asserted rates.
\end{proof}

The collar depends on the chosen finite
primitive bound \(C_0,\nu\). Its exponent
is fixed, but it is not a small or
numerically optimized exponent. This
statement is about the specified
selected path, not every weak-coupling
sequence at arbitrary depth.

\subsection{Information consequences and the next missing estimate}

For \(A_{j,b}=\mathbb E[S_b\mid C]\) and
\(a_{j,b,r}=\mathbb E[S_b\mid N]\), the
tower identity makes the two summands
in \(A_{j,b}=a_{j,b,r}+(A_{j,b}-a_{j,b,r})\)
orthogonal. In particular,
\begin{equation}
 \Var(A_{j,b})-\Var(a_{j,b,r})=\mathcal L_{b,r}^2
 \le\bigl(\omega_{p,r,\beta}
                  +8C_4a_b\epsilon_{p,r,\beta}^{1/4}\bigr)^2 .
 \label{eq:f15v14:fisher}
\end{equation}
Thus the near window captures the full
retained Fisher information up to the
displayed superalgebraic absolute error.
This is compatible with the positive
lower bound in \eqref{eq:f15v7:score}.
It is not a new independent proof of
nonzero near information: the coarse
plaquette used in V7's lower bound was
already near-measurable. Nor does
\eqref{eq:f15v14:fisher} give a strictly
contracting ratio to the growing
fine-source variance.

For clarity about source strength, let
\(d\mu_s=e^{sS_b}d\mu/\mathbb E e^{sS_b}\),
let \(\rho_s\) be its output law, and
let \(\widehat\rho_s\) keep the baseline
far-given-near kernel while adopting
the exact near marginal of \(\rho_s\).
V8's fixed-regulator identity says
\[
 D(\rho_s\Vert\widehat\rho_s)
        =\tfrac12s^2\mathcal L_{b,r}^2+O(s^3)
                       \quad(s\to0).
\]
V14 bounds the second derivative at zero.
It has not supplied a regulator-uniform
third-order remainder, so this formula
does not yet prove locality for a fixed
nonzero \(s\). Passing to that assertion
requires estimates for the actually
tilted laws and control of the changing
output measures, not another zero-source
Taylor identity.

The next substantive direction is therefore
a finite source interval with all normalizers
retained. Even success there would still
concern local source increments. Controlling
the entire baseline effective action, its
multi-source interaction norms, an
interacting RG trajectory at calibrated
physical scales, and a positive continuum
Hamiltonian gap are separate unproved steps.
The mandatory companion review is next,
at V15; no companion theorem has been
silently imported to replace one of these
missing steps.

\subsection{Checks and predecessor preservation}

The certificate is
\begin{center}\small
\path{scripts/verify_ym_f15_v14_full_source_locality.py}.
\end{center}
It uses exact rational arithmetic for eight
positive finite conditional laws, each with
a nonzero source-locality error. It checks
the conditional-copy identity, exceptional
pair probability, truncated fourth moment,
normalized mean transfer, tower contraction,
Fisher orthogonality, and explicit
zero-boundary surrogate bounds. These
examples include correlations in the
exterior posterior. They are diagnostics
of the inequalities, not Wilson simulations.
It also checks 24 source/boundary counts
and collar exponents, and 216 noncommutative
\(\SU(2)\) quaternion identities for coarse
tree-frame covariance and trace invariance.

The certificate has 8,785 bytes and SHA--256
\begin{center}\scriptsize\ttfamily
C0F466FF95C8BF76CC0C517878B03D9EF44957A75F34A54BD6F04C34726462FE.
\end{center}
The predecessor V13 has 296,600 bytes and SHA--256
\begin{center}\scriptsize\ttfamily
A2BF97F059E58868252E044CB5E5B697117BCD992157DE4389EEE07BF5A17C41.
\end{center}
Reversing this terminal insertion and title
change recovers that file exactly.
V14 replaces V13 as the sole active main
note; the 11 protected sources and 56
earlier protected certificates are unchanged.
Earlier proofs and their stated limitations
remain in the appended body.

\begin{firewall}
V14 proves full-compact-law, annealed
\(L^2\) locality for the specified growing
source at a growing collar along the
selected couplings, together with a
gauge-invariant local integral approximation.
It does not prove pointwise quasilocality
for arbitrary exteriors, a fixed-radius
uniform bound, finite-strength source
locality, an analytic RG contraction,
an interacting continuum construction,
or the full Yang--Mills mass gap.
\end{firewall}

\section{V15: locality on a finite source interval}
\label{f15v15:section}

V14 proves locality of the derivative at zero
source. Here we prove a genuinely finite-strength
statement for the same source: its actual coarse
likelihood can be approximated in relative entropy
and total variation by changing only the near
marginal and leaving the baseline far-given-near
kernel intact. The proof works on a fixed real
source interval, independently of the regulator.
It does not extrapolate V8's pointwise Taylor
expansion or assume a uniform Taylor remainder.

There are two new costs. The chart-entry estimate
must be transported to each tilted law, and the
source-score error must be integrated while the
output law changes. The available bounds for
both costs grow exponentially in the source
volume \(O(p^4)\). We pay these costs explicitly
by enlarging V14's collar.

\subsection{The required companion checkpoint}

The active-file inventory still contains SM
continuum V72, NS stability V26, shell-mixing
V16, and parallel YM V5. Their current endings
and scope statements, and those of the three
supplied related notes, were reread. Their
hashes are unchanged from the preceding checkpoint:
\begin{center}\small
\begin{tabular}{@{}lrl@{}}
\toprule
Companion & Bytes & SHA--256 prefix\\
\midrule
SM continuum V72 &607372&\texttt{F44F9745D43379E1}\\
NS stability V26 &278283&\texttt{82C887F8C2C69E4F}\\
YM shell-mixing V16 &623261&\texttt{BF138E63BA826BC0}\\
Parallel YM V5 &54174&\texttt{306F1BB84CFA8A8D}\\
SM source-selection V17 &742509&\texttt{A3AB81EED2BAC59C}\\
Native stationarity V21 &802207&\texttt{C2D452CA7A031E2A}\\
Exact-action Einstein bridge V16 &584434&\texttt{53C50B2CFBA4A05B}\\
\bottomrule
\end{tabular}
\end{center}
The supplied suggestions were also reread and
are unchanged. These documents are sources
and proposed directions, not instructions
to execute their internal acquisition lists.

The relevant distinctions remain as follows.
SM V72 concerns fixed-time, fixed-mode
many-copy covariance transport, not a
physical continuum estimate. NS V26 concerns
rank-one Oseen derivative norms. Neither
provides the YM tilted-law bound below.
Shell V16 assembles the measured local
one-loop coefficient, including its Haar,
gauge, and root-normal factors, and evaluates
the flat holonomy family. It explicitly
does not provide a common nonperturbative
integral remainder. Its warning against
double-counting a root return remains
binding; no new determinant term is added
to our already measured compact density.

Parallel YM V5 rules out its stated
volume-uniform extensive-charge tightness
mechanism under tensorization. We continue
with a window whose fraction of the volume
vanishes, not with that global charge tail.
Source-selection V17 distinguishes complete
source correlations from proper marginals;
we likewise keep the joint output and its
actual conditional kernel. Native V21
distinguishes invariant phase-space control
from a real configuration phase that can
form caustics, and retains differentiated
normalization terms. No dynamical phase
ansatz is used here. The Einstein bridge
V16 distinguishes a compressed inverse
from the full flux-domain problem; no
compressed response is imported as the
inverse of our measured conditional law.

Thus the checkpoint supports continuing
the full-law source calculation, without
claiming that a companion has supplied a
missing physical YM estimate. The next
companion checkpoint is V20.

\subsection{Fixed source convention and changing laws}

Retain V14's source \(S_b\), fine support
\(\mathcal B_p\), output \(C=(N,F)\), and
near-output frame. At each finite regulator
the completed blocking map is held fixed
while the source varies. Its cutoff, source
coordinates, and observation are not retuned.
For this section put
\begin{equation}
 \tau=\frac18,\qquad n_p=|\mathcal B_p|,\qquad
 Z(s)=\mathbb E_\mu e^{sS_b},\qquad
 d\mu_s=\frac{e^{sS_b}}{Z(s)}\,d\mu ,
 \quad |s|\le\tau.
 \label{eq:f15v15:tilt}
\end{equation}
All centering in the definition of \(S_b\)
is with respect to \(\mu=\mu_0\).
For moment estimates under \(\mu_s\), it is
useful to introduce
\[
                    S^{(s)}=S_b-\mathbb E_{\mu_s}S_b .
\]
Let \(\rho_s\) and \(\rho_s^N\) be the
laws of \(C\) and \(N\) under \(\mu_s\).
Throughout the finite-regulator assertions
assume V14's geometric hypotheses,
\(\beta\ge\beta_p\), and \(3\tau<R_{\rm src}\).
The last condition holds eventually on the
selected path.

Write
\begin{equation}
 \begin{split}
 J(z)&=-\log(1-z)-z,\\
 K_p(a)&=k n_pJ(a)\quad(0\le a<1),\\
 B_p&=48\beta p^4,\qquad
 R_p=e^{K_p(3\tau)},\\
 \eta_{p,r,\beta}
     &=\min\{1,\ e^{K_p(2\tau)/2}
                          \epsilon_{p,r,\beta}^{1/2}\},\\
 d_{p,r,\beta}
     &=\omega_{p,r,\beta}+8B_p\eta_{p,r,\beta}^{1/4}.
 \end{split}
 \label{eq:f15v15:budgets}
\end{equation}
Here \(k,R_{\rm src}\) are the inherited
source-window constants, and \(\omega,\epsilon\)
are V14's quantities, evaluated at the collar
chosen in the present section. \(R_p\) in
\eqref{eq:f15v15:budgets} is a likelihood
cost, not a spatial radius or a source window.
The source estimate gives, for \(|a|\le3\tau\),
\begin{equation}
                    1\le Z(a)\le e^{K_p(|a|)}.
 \label{eq:f15v15:partition-budget}
\end{equation}
Indeed Jensen gives the lower bound because
\(\mathbb E_\mu S_b=0\). For the upper
bound use \(a b_x=a1_{\mathcal B_p}(x)\)
in the source budget and
\(J(-z)\le J(z)\) for \(0\le z<1\).
The latter follows directly from the power
series of \(J\).

\subsection{Convexity for the actually tilted interior}

The interior density relative to Lebesgue
coordinates on V13's product of balls now
has potential
\begin{equation}
 V_{J,s}(x;b)=\beta F_J(x;b)-sS_b(x;b)-\ell_J(x;b).
 \label{eq:f15v15:potential}
\end{equation}
An additive scalar in this expression
cancels from its conditional normalizer.
The global normalizer \(Z(s)\) remains
explicit in \eqref{eq:f15v15:tilt}.
No pivot term or exponential Haar term
is removed.

\begin{proposition}[Uniform measured chart bounds on the source interval]
\label{f15v15:convexity}
For \(|s|\le\tau\), required parameter
logarithms of norm at most \(a_p\), and
interior \(x\in\overline{\mathcal D_p}\),
\begin{equation}
 \frac{39}{64}\beta h_p I
          \preceq D_x^2V_{J,s}\preceq2\beta u_p I,
 \qquad
 \sup_a\sum_b|\partial_{ab}V_{J,s}|\le2\beta D_p .
 \label{eq:f15v15:convexity}
\end{equation}
The covariance and exterior-response estimates
of Theorem~\ref{f15v13:covariance} therefore
hold with the same \(\alpha_p\) and the
same displayed constants, for the normalized
restriction of \(\mu_s\).
\end{proposition}

\begin{proof}
At all-identity inputs the Hessian of each
Wilson plaquette is positive semidefinite:
its first derivative vanishes at its minimum.
Pullback through the canonical inverse
preserves that Hessian inequality because
the first derivative vanishes. Let \(K_J\)
be the interior principal classical Hessian
from V12, and let \(H_b\) be the corresponding
Hessian of \(\sum_f\omega_b(f)S_{W,f}\).
The weights in \eqref{eq:f15v14:source}
satisfy \(0\le\omega_b\le1\). Thus
\[
                0\preceq H_b\preceq K_J,
 \qquad
 (1-\tau)K_J\preceq K_J-sH_b
                                      \preceq(1+\tau)K_J .
\]
No simultaneous diagonalization is used.

The absolute third-derivative proof in V13
bounds the joint rows for the classical
tilted sum by \((1+\tau)D_p\). Moving
from zero to the joint chart changes its
interior Hessian by at most
\((1+\tau)D_pa_p=(1+\tau)h_p/8\).
The measured logarithmic correction has
norm at most \(D_p\le\beta h_p/8\).
The lower coefficient is consequently
\[
 (1-\tau)-\frac{1+\tau}{8}-\frac18
                =\frac34-\frac98\tau=\frac{39}{64}.
\]
The upper bound uses \(h_p\le u_p\) and
\((1+\tau)u_p+(2+\tau)h_p/8\le2u_p\).
The joint absolute second-derivative rows
are bounded by
\[
                \beta(1+\tau)D_p+D_p\le2\beta D_p,
\]
since \(\beta\ge\beta_p>2\). The support
range is unchanged, at most \(16p\).

V13's absolute one-form quadratic form
therefore remains bounded below by
\(\beta h_p/2\), with its nonnegative
ball-boundary term retained. The
conjugated Hessian error is still at
most \(\beta h_p/8\), because the
absolute row bound and \(\alpha_p\)
are unchanged. Its weighted resolvent
proof gives the same \(4/(\beta h_p)\)
covariance constant. Differentiating
the normalized integral on the same
fixed domain then gives the same
\(8D_p/h_p\) exterior-response bound.
No new boundary condition on the
observable has been imposed.
\end{proof}

\subsection{Entry and score locality under every tilted law}

\begin{proposition}[Uniform tilted score locality]
\label{f15v15:score}
For every \(|s|\le\tau\), put
\[
 \Delta_s(C)=\mathbb E_{\mu_s}[S_b\mid C]
                         -\mathbb E_{\mu_s}[S_b\mid N].
\]
Then
\begin{equation}
          \|\Delta_s\|_{L^2(\rho_s)}
                         \le d_{p,r,\beta}.
 \label{eq:f15v15:score}
\end{equation}
The events \(\Gamma,D\) used to obtain this
bound are the original V13 events; their
definitions do not change with \(s\).
\end{proposition}

\begin{proof}
If a measurable event \(E\) has
\(\mu(E)\le\epsilon_{p,r,\beta}\),
Cauchy--Schwarz and
\eqref{eq:f15v15:partition-budget} give
\[
 \mu_s(E)
   =\frac{\mathbb E_\mu[1_Ee^{sS_b}]}{Z(s)}
   \le e^{K_p(2\tau)/2}\mu(E)^{1/2}.
\]
It follows that
\(\mu_s(\Gamma^c),\mu_s(D^c)\le\eta_{p,r,\beta}\).
This is a transfer from the actual baseline
law, not a claimed thermal estimate for a
different homogeneous coupling.

The uncentred energy
\(W_b=\beta\sum_f\omega_b(f)S_{W,f}\)
has range in \([0,12\beta n_p]\).
Since \(S^{(s)}=W_b-\mathbb E_{\mu_s}W_b\),
its absolute value is at most \(B_p\).
In particular
\[
                         \|S^{(s)}\|_{L^4(\mu_s)}\le B_p .
\]
This bound does not assume the original
fourth-moment estimate for the new law.

For a fixed \(s,N\), the normalized chart
mean of \(S^{(s)}\) has the same derivative
carrier, gradient bound, and exterior
coordinate count as in V14. Its centering
is a scalar independent of the exterior.
Proposition~\ref{f15v15:convexity} supplies
the same response bound. Therefore its
oscillation over two \(\Gamma\)-exteriors
with the same \(N\) is at most
\(\omega_{p,r,\beta}\).

All ingredients of V14's conditional-copy
proof now hold for \(\mu_s\): replace
\(q_D\) by \(\mu_s(D\mid Z)\), the
exceptional probability by \(\eta_{p,r,\beta}\),
and \(M_4\) by \(B_p\).
In particular the normalized mean transfer
cost is at most \(5B_p\eta^{1/4}\);
the conditional-pair variance cost is
at most \(\omega+3B_p\eta^{1/4}\).
The exterior posterior in both copies
is the actual one for \(\mu_s\).
Projecting onto \(C\), and observing
that subtracting \(\mathbb E_{\mu_s}S_b\)
cancels in the difference of its two
conditional means, proves
\eqref{eq:f15v15:score}.
\end{proof}

\subsection{The change of output law is explicitly paid}

To integrate \(\Delta_s\), its norm under
\(\rho_s\) cannot simply be reused under
\(\rho_t\) for a different \(t\).
The following bound pays for that change.

\begin{proposition}[Likelihood transport between source values]
\label{f15v15:transport}
For \(s,u\in[-\tau,\tau]\),
\begin{equation}
 \left\|\frac{d\rho_s}{d\rho_u}\right\|_{L^2(\rho_u)}^2
       \le\frac{Z(u)Z(2s-u)}{Z(s)^2}
       \le R_p^2 .
 \label{eq:f15v15:transport}
\end{equation}
Consequently every \(f\in L^2(\rho_u)\) obeys
\begin{equation}
            \mathbb E_{\rho_s}|f|
                         \le R_p\|f\|_{L^2(\rho_u)}.
 \label{eq:f15v15:changed-mean}
\end{equation}
\end{proposition}

\begin{proof}
On the fine space,
\[
           \frac{d\mu_s}{d\mu_u}
                      =\frac{Z(u)}{Z(s)}e^{(s-u)S_b}.
\]
Its squared \(L^2(\mu_u)\) norm is exactly
\(Z(u)Z(2s-u)/Z(s)^2\). For the output,
\[
 \frac{d\rho_s}{d\rho_u}(C)
      =\mathbb E_{\mu_u}\!\left[
                        \frac{d\mu_s}{d\mu_u}\,\middle|\,C\right].
\]
Conditional Jensen proves the first
inequality. Since \(|u|\le\tau\),
\(|2s-u|\le3\tau<R_{\rm src}\), and
\(Z(s)\ge1\), the partition budget
bounds the ratio by \(e^{2K_p(3\tau)}\).
Cauchy--Schwarz under \(\rho_u\)
proves \eqref{eq:f15v15:changed-mean}.
\end{proof}

This is an application of the exact
generated-likelihood identities from V2
to two different tilted reference laws.
The new use is the uniform spatial error
from Proposition~\ref{f15v15:score};
the likelihood algebra alone did not
establish such an error in V2.

\subsection{Finite-strength local source increments}

Define the exact likelihoods
\begin{equation}
 \begin{split}
 h_s^C(C)&=\frac{\mathbb E_\mu[e^{sS_b}\mid C]}{Z(s)}
                                      =\frac{d\rho_s}{d\rho_0}(C),\\
 h_s^N(N)&=\frac{\mathbb E_\mu[e^{sS_b}\mid N]}{Z(s)}
                                      =\frac{d\rho_s^N}{d\rho_0^N}(N),\\
             d\widehat\rho_s(C)&=h_s^N(N)\,d\rho_0(C).
 \end{split}
 \label{eq:f15v15:local-law}
\end{equation}
The comparison is a probability with the
exact near marginal of \(\rho_s\).
Its conditional law of \(F\) given \(N\)
is the baseline one. This is the same
comparison as in V8, now with a finite
source error estimate instead of only
an expansion at zero. Gauge invariance
of the source and the baseline law makes
the near likelihood gauge invariant
(up to the usual choice of conditional
versions).

\begin{theorem}[A fixed source interval with a quantitative locality error]
\label{f15v15:finite-source}
For every \(|s|\le\tau\),
\begin{equation}
 \begin{split}
 D(\rho_s\Vert\widehat\rho_s)
    &=\int D\!\left(\rho_s(dF\mid N)
                  \,\Vert\,\rho_0(dF\mid N)\right)\,d\rho_s^N(N)\\
    &\le |s|R_p\,d_{p,r,\beta},\\
 \|\rho_s-\widehat\rho_s\|_{\rm TV}
    &\le\left(\frac{|s|R_p\,d_{p,r,\beta}}2\right)^{1/2}.
 \end{split}
 \label{eq:f15v15:finite-source}
\end{equation}
Here total variation is half the \(L^1\)
distance of probability densities.
\end{theorem}

\begin{proof}
At fixed regulator \(S_b\) is bounded,
and all these laws are equivalent.
Their conditional moment generating
functions may be defined from a common
regular conditional law, so their source
derivatives exist simultaneously for
real \(s\). Differentiating numerator
and denominator gives
\[
 \partial_s\log h_s^C
       =\mathbb E_{\mu_s}[S_b\mid C]-\mathbb E_{\mu_s}S_b,
 \qquad
 \partial_s\log h_s^N
       =\mathbb E_{\mu_s}[S_b\mid N]-\mathbb E_{\mu_s}S_b .
\]
Both likelihoods equal one at zero.
Subtracting the two formulas and
integrating gives the exact identity
\begin{equation}
        \log\frac{h_s^C(C)}{h_s^N(N)}
                                  =\int_0^s\Delta_u(C)\,du .
 \label{eq:f15v15:log-path}
\end{equation}
For negative \(s\), the integral has
its usual oriented meaning.

The logarithm on the left is
\(\log(d\rho_s/d\widehat\rho_s)\).
Integrate under \(\rho_s\), use the
absolute value of the right side, and
apply \eqref{eq:f15v15:changed-mean}
to \(f=\Delta_u\). Along the interval
from zero to \(s\),
\[
 \mathbb E_{\rho_s}|\Delta_u|
       \le R_p\|\Delta_u\|_{L^2(\rho_u)}
       \le R_p d_{p,r,\beta}.
\]
Fubini is legitimate: at this fixed
regulator all conditional means are
bounded by the source range.
This proves the entropy upper bound.
Disintegrating the likelihood with
respect to \(N\), whose marginals agree,
proves its conditional-entropy identity.

For completeness, the total-variation
bound follows by applying entropy
data processing to the set where
\(d\rho_s>d\widehat\rho_s\).
The resulting Bernoulli probabilities
have difference equal to the total
variation. Bernoulli relative entropy
is at least twice that squared
difference: at fixed reference
probability its second derivative
in the first probability is
\(1/[x(1-x)]\ge4\), and its value
and first derivative vanish at
equality. This gives the final line
of \eqref{eq:f15v15:finite-source}.
\end{proof}

In particular the far conditional kernel
changes little on average under the
actual tilted near marginal. It is not
asserted to change little for every
near boundary value. The source strength
\(s\) need not tend to zero; the estimate
is not a discarded higher-order remainder.

\subsection{An enlarged collar controls both exponential costs}

\begin{theorem}[Uniform finite-source locality on the selected sequence]
\label{f15v15:asymptotic}
On V7's slower path, put \(t=\log\beta\)
and choose
\begin{equation}
        r=\max\{176,\lceil p^{\nu+8}t^6\rceil\}.
 \label{eq:f15v15:collar}
\end{equation}
All the preceding hypotheses hold eventually.
For every fixed \(A>0\),
\begin{equation}
 \begin{split}
 \sup_{|s|\le1/8}D(\rho_s\Vert\widehat\rho_s)
                                      &=o(\beta^{-A}),\\
 \sup_{|s|\le1/8}\|\rho_s-\widehat\rho_s\|_{\rm TV}
                                      &=o(\beta^{-A}).
 \end{split}
 \label{eq:f15v15:asymptotic}
\end{equation}
Also \(rp/M\to0\). The near observation is
still a vanishing-volume-fraction window.
\end{theorem}

\begin{proof}
Eventually \(p\le t\), \(k\le5\),
\(3/8<R_{\rm src}\), and
\[
 r\le2p^{\nu+8}t^6\le2t^{\nu+14}.
\]
As in V14, \(\beta=e^t\) eventually
exceeds \(\beta_p=2^{43}C_0p^{\nu+8}\),
whereas \(M\ge e^\beta\) exceeds every
fixed power of \(t\). Therefore
\(M\ge4L\) and \(rp/M\to0\).
With \(c_0,c_1\) unchanged from V14,
\begin{align}
 \omega_{p,r,\beta}
       &\le16e^t t^{4\nu+64}e^{-c_0t^6},
                                     \label{eq:f15v15:omega-rate}\\
 \epsilon_{p,r,\beta}
       &\le\min\left\{1,\,
       2^{17}C_{\rm th}t^{4\nu+60}
           \exp\left[-\frac{c_1e^t}{t^{6\nu+82}}\right]\right\}.
                                     \label{eq:f15v15:epsilon-rate}
\end{align}
Indeed \(r-88\ge r/2\) gives
\(\alpha_p(r-88)p\ge c_0t^6\).
The prefactors and the denominator
follow from
\[
 r^4p^8\le16t^{4\nu+64},\qquad
 r^4p^4\le16t^{4\nu+60},\qquad
 r^4p^{2\nu+26}\le16t^{6\nu+82}.
\]

The elementary bound
\(J(z)\le z^2/[2(1-z)]\) gives
\(J(3/8)\le9/80\) and \(J(1/4)\le1/24\).
Since \(n_p\le4p^4\) and \(k\le5\),
\[
                 R_p\le e^{(9/4)p^4},
 \qquad
                 \eta^{1/4}
                  \le e^{(5/48)p^4}\epsilon^{1/8}.
\]
Consequently the entropy error budget obeys
\begin{equation}
 R_p d_{p,r,\beta}
       \le e^{(9/4)p^4}\omega_{p,r,\beta}
             +384\beta p^4 e^{(113/48)p^4}
                                  \epsilon_{p,r,\beta}^{1/8}.
 \label{eq:f15v15:paid-budget}
\end{equation}
For the first term, the negative exponent
\(-c_0t^6\) dominates the positive
\((9/4)t^4+t\) and logarithmic prefactors.
For the second, the negative exponent
proportional to \(-e^t/t^{6\nu+82}\)
dominates every fixed power of \(t\).
Both terms are \(o(e^{-At})\) for
every fixed \(A>0\).
The entropy conclusion follows uniformly
from \eqref{eq:f15v15:finite-source}.
Taking the square root preserves the
assertion for every \(A\) for total variation.
\end{proof}

V14's \(t^2\) collar exponent alone did
not pay the displayed likelihood cost
\(e^{O(p^4)}\) when \(p\) can be as large
as \(t\). The \(t^6\) choice is an explicit,
nonoptimized way to pay it. The constants
are still defined through the fixed
primitive derivative bounds; no practical
numerical onset scale is claimed.

\subsection{What this closes, and the next upstream issue}

For the selected lattice laws, a fixed
real interval of the original local
fine-action source now has an actual
near-generated coarse likelihood with
vanishing full-output relative-entropy
error. All compact exterior posterior
weights, pivot Jacobians, Haar factors,
and source normalizers have remained
present. This strengthens V14's
infinitesimal score result; it is not
another zero-source coefficient identity.

The comparison still uses the \emph{exact}
near marginal of the actual tilted law.
It proves locality of a source increment,
not a local formula for the baseline
coarse density. Nor does it prove
uniform locality for arbitrary near
configurations, a collar independent
of the regulator, or a summable
analytic norm for every generated
multi-source interaction.

The next issue is therefore the passage
from this one growing source to a
controlled family of simultaneous local
sources and their connected dependence.
A naive replacement by a global source
vector would make the likelihood budget
extensive in all its supports. Any such
extension must track that cost and the
geometry of the union; pointwise or
volume-uniform locality must not be
inferred from the present annealed
one-source theorem. In particular,
the result does not by itself control
iterations in a closed interacting
RG class at calibrated physical scales.
That class, continuum reconstruction,
and the physical Hamiltonian gap remain
unproved.

\subsection{Checks and predecessor preservation}

The new certificate is
\begin{center}\small
\path{scripts/verify_ym_f15_v15_finite_source_locality.py}.
\end{center}
It uses exact rational tilted finite laws
with the parameter represented as \(q=e^s\).
Forty tilted examples check actual
likelihood ratios, their logarithmic
source derivatives, normalization of
the local near-marginal comparison,
tilted event transport, and normalized
chart-mean transfer. Thirty-two
nonzero-source examples have genuinely
different full and local-comparison laws.
Two hundred pairs of source parameters
check the fine partition identity,
output likelihood contraction, and the
changed-law Cauchy--Schwarz bound.
Five noncommuting weighted-Hessian
examples check the order inequalities;
24 scale examples check the enlarged
collar exponents. Their role is to
test the stated algebra, not to
simulate Wilson laws or establish
a continuum theorem by finite sampling.

The certificate has 10,448 bytes and SHA--256
\begin{center}\scriptsize\ttfamily
2E49BE0DA751826EA665991E0D30D503E4D7A0AA458CE7AC44FEFFC1DC469E0D.
\end{center}
The predecessor V14 has 317,086 bytes and SHA--256
\begin{center}\scriptsize\ttfamily
A5F2D4137420784632EA3EF49B64F01072F8428D92891FFC165744A57C1AA9C7.
\end{center}
Reversing this terminal insertion and title
change recovers V14 exactly. V15 replaces
it as the sole active main note. The
11 protected sources and 57 earlier
protected certificates remain unchanged.
The complete earlier mathematical body
and its qualifications are preserved.

\begin{firewall}
V15 proves finite-strength, annealed
source-increment locality in relative
entropy and total variation along the
declared sequence, with a growing
collar and an explicit source-volume
cost. It does not establish the
locality or contraction of the entire
effective action, a nontrivial interacting
continuum Yang--Mills theory, or a
positive physical mass gap.
\end{firewall}

\section{V16: separated joint sources and conditional cumulants}
\label{f15v16:section}

V15 controls a fixed interval for one source.
The new question is whether the same local
description remains valid when other separated
sources are present. Testing each source
only against the original law would not answer
this question: remote tilts can change the
exterior posterior even when their factors
cancel from a conditioned interior density.

We prove a joint-background estimate and use it
to control the mixed part of the actual
conditional source generating function. The
number of separated sources may grow as
\(\lfloor\log\beta\rfloor\). We then extract
mixed conditional cumulants at every fixed
total order, including repeated source indices.
These are estimates under the full compact
law and its declared source tilts. They are
not unconditional clustering estimates.

The source files and earlier certificates
were checked unchanged. The companion
checkpoint completed at V15 remains current;
the next required review is V20.

\subsection{Separated supports and a joint source cube}

Let \(m\ge2\), and take \(m\) translates
of V14's source by final-grid translations
\(pz_i\). Denote their fine base sets
by \(\mathcal B_{p,i}\), and put
\begin{equation}
 b_{i,x}=1_{\mathcal B_{p,i}}(x),\qquad
 S_i=\sum_xb_{i,x}(T_x-\mathbb E_\mu T_x),\qquad
 n_p=|\mathcal B_{p,i}|=p^2(2p-1)^2 .
 \label{eq:f15v16:sources}
\end{equation}
The blocking hierarchy is unchanged by these
translations. Require the fine anchor centres
to have periodic maximum-norm separation
at least \(4rp\). Each translated fine
box \(\mathcal K_i\), of side \(2p(r+5)\),
is embedded without self-identification.
Use V14's \(r\ge176\), \(M\ge4L\),
and \(\beta\ge\beta_p\). All sources
have the same block depth and orientation.

Write \(N_i,H_i,J_i,O_i,\Gamma_i,D_i\)
for the translated near output, its
frame, the canonical interior and exterior,
and the two V13 chart-entry events.
These constructions are made separately
for each \(i\); no common gauge choice
for all the boxes is imposed.
The canonical derivative carrier of
\(S_i\) is within \(12p\) of its centre.
For \(l\ne i\), its distance from \(J_i\)
is positive, since
\[
          4rp-12p>(r-60)p.
\]
In particular \(S_l\) depends only on
\(C,v_{O_i}^{H_i}\), not the interior
details \(v_{J_i}^{H_i}\).
The fine source supports are disjoint.

Keep \(\tau=1/8\), and define, for
\(\theta\in[-\tau,\tau]^m\),
\begin{equation}
 Z(\theta)=\mathbb E_\mu e^{\sum_i\theta_iS_i},
 \qquad
 d\mu_\theta=\frac{e^{\sum_i\theta_iS_i}}{Z(\theta)}\,d\mu,
 \qquad \rho_\theta=C_*\mu_\theta .
 \label{eq:f15v16:joint-tilt}
\end{equation}
The completed blocking map and all baseline
centering constants are fixed throughout.
Assume \(3\tau<R_{\rm src}\). Disjointness
of the base supports and the original
source budget give, whenever
\(\|\theta\|_\infty\le a\le3\tau\),
\begin{equation}
          1\le Z(\theta)
                   \le\exp\{k m n_pJ(a)\}.
 \label{eq:f15v16:partition}
\end{equation}
No independence of the \(S_i\) is used.
The lower bound is Jensen, since each
\(S_i\) is centred under the baseline law.

In the notation of V14--V15, set
\begin{equation}
 \begin{split}
 B_p&=48\beta p^4,\qquad
 \mathcal R_m=\exp\{k m n_pJ(3\tau)\},\\
 \eta_m&=\min\{1,\,
          e^{k m n_pJ(2\tau)/2}\epsilon_{p,r,\beta}^{1/2}\},\\
 \mathfrak e_m&=\omega_{p,r,\beta}
                 +5B_p\eta_m^{1/4}+\sqrt3\,B_p\eta_m^{1/2}.
 \end{split}
 \label{eq:f15v16:budgets}
\end{equation}
All source-volume costs are displayed.
In particular \(\mathcal R_m\) is not
being treated as independent of \(m,p\).

\subsection{A local integral independent of the remote source parameters}

Put \(Z_i^{\rm ext}=(C,v_{O_i}^{H_i})\).
The exact compact conditional density
from V12 implies
\begin{equation}
 \mu_\theta(dv_{J_i}^{H_i}\mid Z_i^{\rm ext})
  =\frac{e^{\theta_iS_i}}
          {\mathbb E_\mu[e^{\theta_iS_i}\mid Z_i^{\rm ext}]}
                 \mu(dv_{J_i}^{H_i}\mid Z_i^{\rm ext}).
 \label{eq:f15v16:interior}
\end{equation}
Indeed all factors with \(l\ne i\)
are measurable with respect to
\(Z_i^{\rm ext}\), and cancel from this
conditional normalization. They do
\emph{not} cancel from the distribution
of \(Z_i^{\rm ext}\).

Let \(\Gamma_{N_i}\) be the retained-input
half-chart event in the \(H_i\)-frame.
On this event define \(G_i(s,N_i)\)
as the normalized chart mean of the
baseline-centred \(S_i\), with potential
\(V_{J_i}-sS_i\) and all required
exterior canonical details set to \(I\).
Define it to be zero on \(\Gamma_{N_i}^c\).
Thus it is exactly V14's specified local
integral, with the additional scalar
source \(s\) in its actual measured
potential. It depends on \(s,N_i\),
but not on any \(\theta_l\), \(l\ne i\).
It is gauge invariant and
\begin{equation}
                            |G_i(s,N_i)|\le B_p .
 \label{eq:f15v16:local-bound}
\end{equation}
Its hard retained-chart cutoff is
unchanged; no smoothness across that
cutoff is asserted.

\begin{theorem}[Local source means in every joint background]
\label{f15v16:local-means}
Uniformly for \(\theta\in[-\tau,\tau]^m\)
and every \(i\),
\begin{equation}
 \left\|\mathbb E_{\mu_\theta}[S_i\mid C]
                    -G_i(\theta_i,N_i)\right\|_{L^2(\rho_\theta)}
                                   \le\mathfrak e_m .
 \label{eq:f15v16:local-means}
\end{equation}
The local function in this formula is
the same for all choices of the
remote source parameters.
\end{theorem}

\begin{proof}
For any event \(E\), Cauchy--Schwarz
and \eqref{eq:f15v16:partition} give
\[
 \mu_\theta(E)
   \le\frac{Z(2\theta)^{1/2}\mu(E)^{1/2}}{Z(\theta)}
   \le e^{k m n_pJ(2\tau)/2}\mu(E)^{1/2}.
\]
The translated entry bound therefore
implies
\[
        \mu_\theta(\Gamma_i^c),\ \mu_\theta(D_i^c)
                                                    \le\eta_m .
\]
This is a separate estimate for each
interior. No simultaneous-good-event
probability or independent exterior
law is required.

By \eqref{eq:f15v16:interior}, the chart
potential for this interior is precisely
the one-source potential of V15 with
parameter \(\theta_i\). Its measured
Hessian is bounded below by
\(39\beta h_p/64\), its absolute joint
Hessian rows by \(2\beta D_p\), and
its term range is \(16p\).
All normalized covariance and exterior
response bounds are therefore unchanged.
The chart mean of \(S_i\), at fixed \(N_i\),
has exterior oscillation at most
\(\omega_{p,r,\beta}\).

The baseline-centred source has
\(|S_i|\le B_p\) pointwise, under any law.
Apply V13's normalized mean-transfer
argument to \(\mu_\theta,S_i\), using
\(\|S_i\|_{L^4(\mu_\theta)}\le B_p\).
For
\[
 A_i=\Gamma_i\cap
       \{\mu_\theta(D_i\mid Z_i^{\rm ext})\ge1/2\},
\]
the discarded conditional-mean norm is
at most \(5B_p\eta_m^{1/4}\), and
\(\mu_\theta(A_i^c)\le3\eta_m\).
On \(A_i\), the restricted mean differs
from \(G_i(\theta_i,N_i)\) by at most
\(\omega_{p,r,\beta}\); on \(A_i^c\),
the truncated mean is zero and
\(|G_i|\le B_p\).
The latter comparison costs at most
\(\omega_{p,r,\beta}+\sqrt3 B_p\eta_m^{1/2}\)
in \(L^2\). Projection onto \(C\)
is a contraction. Combining the two
errors proves the assertion.

The independence from remote parameters
follows from the exact cancellation
in \eqref{eq:f15v16:interior}, before
any chart or exceptional-event estimate.
The baseline centering of \(S_i\)
is important here: it is a fixed
scalar, not a joint-background mean
depending on all of \(\theta\).
\end{proof}

\begin{proposition}[Joint likelihood transport]
\label{f15v16:transport}
For any \(\theta,\xi\in[-\tau,\tau]^m\),
\begin{equation}
 \left\|\frac{d\rho_\theta}{d\rho_\xi}\right\|_{L^2(\rho_\xi)}^2
        \le\frac{Z(\xi)Z(2\theta-\xi)}{Z(\theta)^2}
        \le\mathcal R_m^2 .
 \label{eq:f15v16:transport}
\end{equation}
Consequently every \(f\in L^2(\rho_\xi)\)
satisfies
\(\mathbb E_{\rho_\theta}|f|
       \le\mathcal R_m\|f\|_{L^2(\rho_\xi)}\).
\end{proposition}

\begin{proof}
The fine density ratio is
\(Z(\xi)e^{(\theta-\xi)\cdot S}/Z(\theta)\).
Its squared norm is the displayed
partition ratio. The output ratio is
its conditional expectation given \(C\)
under \(\mu_\xi\), so conditional
Jensen gives the first inequality.
The components of \(2\theta-\xi\)
have magnitude at most \(3\tau\).
Use \eqref{eq:f15v16:partition} for
the two numerator factors and
\(Z(\theta)\ge1\). The last assertion
is Cauchy--Schwarz.
\end{proof}

\subsection{The mixed conditional generating function}

Let \(e_i\) denote the \(i\)-th source
coordinate vector and define
\begin{equation}
 \Phi_\theta(C)=
       \log\mathbb E_\mu[e^{\theta\cdot S}\mid C],
 \qquad
 Q_\theta(C)=\Phi_\theta(C)
                          -\sum_{i=1}^m\Phi_{\theta_i e_i}(C).
 \label{eq:f15v16:defect}
\end{equation}
These use one common conditional law
of the bounded finite source vector.
At a fixed regulator they are smooth
in the real source parameters, on a
common full-measure set for \(C\).
All the output laws are equivalent,
so that set works for every tilt.

\begin{theorem}[Conditional additivity with all source backgrounds retained]
\label{f15v16:additivity}
For every \(\theta,\zeta\in[-\tau,\tau]^m\),
\begin{equation}
             \|Q_\theta\|_{L^1(\rho_\zeta)}
                 \le2\|\theta\|_1\mathcal R_m\mathfrak e_m .
 \label{eq:f15v16:additivity}
\end{equation}
This is a statement about the actual
joint conditional source generator,
not a product assumption on the fine
sources or the retained output.
\end{theorem}

\begin{proof}
Differentiation of the conditional
moment generating function gives
\[
       \partial_{\theta_i}\Phi_\theta(C)
                       =\mathbb E_{\mu_\theta}[S_i\mid C].
\]
All functions vanish at zero source.
Integrating along the radial path
from zero to \(\theta\) gives exactly
\begin{equation}
 \begin{split}
 Q_\theta(C)=\int_0^1\sum_i\theta_i
  \bigl(&\mathbb E_{\mu_{v\theta}}[S_i\mid C]\\
        &-\mathbb E_{\mu_{v\theta_i e_i}}[S_i\mid C]\bigr)\,dv .
 \end{split}
 \label{eq:f15v16:path}
\end{equation}
Subtract \(G_i(v\theta_i,N_i)\) from
both means inside each difference.
It is the same function in the two
backgrounds by Theorem~\ref{f15v16:local-means}.
Each error has its \(L^2\) norm bounded
by \(\mathfrak e_m\) in its own output law.
Proposition~\ref{f15v16:transport} moves
it to \(L^1(\rho_\zeta)\) at cost
\(\mathcal R_m\). Thus the norm of
each difference in parentheses is
at most \(2\mathcal R_m\mathfrak e_m\).
Integrating and summing absolute
source amplitudes proves the theorem.
At the fixed regulator all means
are bounded by \(B_p\); hence all
interchanges with the finite path
integral are legitimate.
\end{proof}

For a subset \(T\subseteq\{1,\ldots,m\}\)
with \(|T|\ge2\), let \(\theta^A\)
equal \(\theta\) on \(A\subseteq T\)
and zero elsewhere. The genuinely
mixed finite increment
\[
 \mathcal C_T(\theta;C)
       =\sum_{A\subseteq T}(-1)^{|T|-|A|}
                                      \Phi_{\theta^A}(C)
\]
satisfies
\begin{equation}
 \|\mathcal C_T(\theta;\cdot)\|_{L^1(\rho_\zeta)}
     \le2^{|T|}\left(\sum_{i\in T}|\theta_i|\right)
                                  \mathcal R_m\mathfrak e_m .
 \label{eq:f15v16:finite-increment}
\end{equation}
Indeed every one-source term cancels in
the alternating subset sum, leaving
the same sum of \(Q_{\theta^A}\).
Each \(|\theta_i|\) occurs in
\(2^{|T|-1}\) of the resulting bounds
\eqref{eq:f15v16:additivity}.
This argument controls finite increments;
derivatives require the additional
estimate given below.

\subsection{Up to logarithmically many sources}

\begin{theorem}[Uniform decay for a growing separated family]
\label{f15v16:asymptotic}
On the selected path let
\[
 t=\log\beta,\qquad
 r=\max\{176,\lceil p^{\nu+8}t^6\rceil\},\qquad
                         2\le m\le\lfloor t\rfloor .
\]
Uniformly over all allowed placements,
source vectors and subsets, the errors
in \eqref{eq:f15v16:additivity} and
\eqref{eq:f15v16:finite-increment} are
\(o(\beta^{-A})\) for every fixed \(A>0\).
Also
\[
                         \frac{m(rp)^4}{M^4}\longrightarrow0.
\]
Thus even the union of the near windows
occupies a vanishing fraction of the
full output volume.
\end{theorem}

\begin{proof}
V15's spatial and entry rates apply to
each translate with the same constants.
Eventually \(p\le t\), \(k\le5\), and
\[
 \begin{split}
 \mathcal R_m&\le e^{(9/4)mp^4},\\
 \eta_m^{1/4}&\le e^{(5/48)mp^4}\epsilon^{1/8},\\
 \eta_m^{1/2}&\le e^{(5/24)mp^4}\epsilon^{1/4}.
 \end{split}
\]
In particular
\begin{equation}
 \begin{split}
 \mathcal R_m\mathfrak e_m
 \le{}& e^{(9/4)mp^4}\omega
       +240\beta p^4e^{(113/48)mp^4}\epsilon^{1/8}\\
      &+48\sqrt3\,\beta p^4
                             e^{(59/24)mp^4}\epsilon^{1/4}.
 \end{split}
 \label{eq:f15v16:paid-budget}
\end{equation}
Here \(mp^4\le t^5\). The spatial
factor from \eqref{eq:f15v15:omega-rate}
has negative exponent \(-c_0t^6\),
which dominates every displayed
positive multiple of \(t^5\).
The entry factor from
\eqref{eq:f15v15:epsilon-rate} has
negative exponent proportional to
\(-e^t/t^{6\nu+82}\), which is
stronger still. Consequently
\(\mathcal R_m\mathfrak e_m\) is
\(o(\beta^{-A})\) uniformly in
the stated family, for every \(A\).

The factors \(m\) and \(2^{|T|}\)
do not change this conclusion:
\(m\le t\) and
\(2^{|T|}\le2^m\le e^{(\log2)t}\).
Finally \(r\le2t^{\nu+14}\) eventually,
so \(m(rp)^4\) grows at most as
a fixed power of \(t\), whereas
\(M\ge e^\beta=e^{e^t}\).
This proves the volume-fraction
statement as well.
\end{proof}

There is no union bound over all
coarse cells here. The number of
source supports and their placement
are part of the theorem; they cannot
be replaced by a positive-density
source field without a new estimate.

\subsection{Fixed-order mixed conditional cumulants}

The finite-strength result also yields
coefficient control, but an argument is
needed to pass to source derivatives.
For \(n\ge1\) define the finite number
\begin{equation}
 c_n=\sum_{\pi\in\mathcal P_n}(|\pi|-1)!,
 \label{eq:f15v16:cumulant-constant}
\end{equation}
where \(\mathcal P_n\) is the set of
partitions of \(n\) labelled positions.
For bounded random variables
\(|X_j|\le B\), their joint cumulant
obeys
\begin{equation}
 |\kappa(X_1,\ldots,X_n)|\le c_n B^n .
 \label{eq:f15v16:cumulant-bound}
\end{equation}
To verify this, differentiate the
logarithm of their moment generating
function at zero. Expanding
\(\log(1+u)\) groups the labelled
derivatives by partitions and gives
\[
 \kappa(X_1,\ldots,X_n)
   =\sum_{\pi\in\mathcal P_n}
       (-1)^{|\pi|-1}(|\pi|-1)!
             \prod_{A\in\pi}\mathbb E\prod_{j\in A}X_j .
\]
Every product of moments has magnitude
at most \(B^n\), proving the bound.
At a finite regulator boundedness
justifies this expansion in a
neighborhood of zero. The same formula
applies to a conditional tilted law,
including repeated variables.

Let \(\alpha\) be a multi-index of
fixed total order \(q=|\alpha|\ge2\)
with at least two nonzero entries.
Set
\begin{equation}
 \kappa_\alpha(C)=
       \partial^\alpha_\theta\Phi_\theta(C)\big|_{\theta=0}.
 \label{eq:f15v16:conditional-cumulant}
\end{equation}
This is the actual mixed cumulant of
the indicated repeated \(S_i\)'s under
\(\mu(\,\cdot\,\mid C)\).

\begin{theorem}[Every fixed mixed order is small]
\label{f15v16:cumulants}
Under the finite-regulator hypotheses,
for any \(h>0\) with \(qh\le\tau\)
and any \(\zeta\in[-\tau,\tau]^m\),
\begin{equation}
 \|\kappa_\alpha\|_{L^1(\rho_\zeta)}
 \le 2^{q+1}q\,h^{1-q}\mathcal R_m\mathfrak e_m
                              +q c_{q+1}h B_p^{q+1}.
 \label{eq:f15v16:cumulants}
\end{equation}
On the sequence of
Theorem~\ref{f15v16:asymptotic}, for
every fixed \(q,A>0\) with integer
\(q\ge2\), this norm is
\(o(\beta^{-A})\), uniformly over
all such mixed multi-indices and
the allowed source backgrounds.
\end{theorem}

\begin{proof}
Apply the repeated forward difference
of order \(\alpha_i\), step \(h\),
in each coordinate \(i\). Let
\(\delta_h^\alpha\) be the resulting
operator, evaluated at zero. Its
coefficients have total absolute
sum \(2^q\). Every evaluation point
\(hj\), \(0\le j_i\le\alpha_i\),
has \(\|hj\|_1\le qh\) and lies
in the source cube.

Because at least two coordinates
occur, the mixed difference annihilates
each one-source function
\(\Phi_{\theta_i e_i}\).
It therefore has the same value
on \(\Phi\) and \(Q\).
Theorem~\ref{f15v16:additivity} gives
\[
 \|h^{-q}\delta_h^\alpha\Phi\|_{L^1(\rho_\zeta)}
       \le2^{q+1}q\,h^{1-q}
                                  \mathcal R_m\mathfrak e_m .
\]

The iterated fundamental theorem of
calculus expresses
\(h^{-q}\delta_h^\alpha\Phi\)
as the average of \(\partial^\alpha\Phi\)
at points whose coordinate sum is
at most \(qh\). At every real
parameter on the joining segments,
each derivative of \(\Phi\) of
total order \(q+1\) is a conditional
tilted cumulant of \(q+1\) variables
bounded by \(B_p\).
Equation \eqref{eq:f15v16:cumulant-bound}
therefore bounds it pointwise by
\(c_{q+1}B_p^{q+1}\).
The difference between that average
and \(\partial^\alpha\Phi(0)\)
is at most \(qhc_{q+1}B_p^{q+1}\).
Combining the two bounds proves
\eqref{eq:f15v16:cumulants}.

For the asymptotic assertion, fix
\(q,A\) and choose
\(h=\beta^{-(A+q+3)}\).
Eventually \(qh\le\tau\), and
\(B_p\le48\beta(\log\beta)^4\).
The second term is thus bounded
by a fixed multiple of
\[
       \beta^{-A-2}(\log\beta)^{4(q+1)}
                                      =o(\beta^{-A}).
\]
The first multiplies
\(\mathcal R_m\mathfrak e_m\)
by only a fixed power of \(\beta\);
Theorem~\ref{f15v16:asymptotic}
still makes it \(o(\beta^{-A})\).
All constants depend on the fixed
order \(q\), not on the indices
chosen within the separated family.
\end{proof}

This is a fixed-order theorem.
It does not control a sum over all
orders, a regulator-uniform complex
source domain, or an analytic
interaction norm. In particular the
constants \(c_{q+1}\) and the factor
\(B_p^{q+1}\) have not been
suppressed from the proof.

\subsection{Normalization and physical correlation remain separate}

For the normalized output likelihood
\(L_\theta=d\rho_\theta/d\rho_0\),
the exact identity is
\begin{equation}
 \log L_\theta
   =\sum_i\log L_{\theta_i e_i}
       +\sum_i\log Z(\theta_i e_i)-\log Z(\theta)
       +Q_\theta .
 \label{eq:f15v16:normalization}
\end{equation}
The scalar joint free-energy term
in this formula has not been
proved small. It cannot be removed
by treating the original sources
as independent. Nor does small
average \(Q_\theta\), by itself,
control the normalization of
\(\prod_i L_{\theta_i e_i}\rho_0\).

\begin{proposition}[A small mean log error need not control normalization]
\label{f15v16:normalizer-obstruction}
There are finite probability spaces
and bounded functions \(r_n\), bounded
at each individual \(n\), such that
\(\mathbb E|r_n|\to0\), but the
probabilities proportional to
\(e^{-r_n}dP_n\) have total variation
distance from \(P_n\) tending to one.
\end{proposition}

\begin{proof}
On a two-point space choose an event
of probability \(4^{-n}\), and let
\(r_n=-3n\log2\) on that event and
zero elsewhere. Then
\[
 \mathbb E_{P_n}|r_n|=3n(\log2)4^{-n}\to0,\qquad
 \mathbb E_{P_n}e^{-r_n}=1-4^{-n}+2^n\to\infty .
\]
The normalized tilted probability
of the event is
\(2^n/(1-4^{-n}+2^n)\).
Its difference from \(4^{-n}\)
is at least \(1-2^{-n}-4^{-n}\).
This lower bound tends to one.
\end{proof}

This is a diagnostic of the proposed
norm-to-normalization inference, not
a Wilson counterexample or a failure
of the preceding conditional estimates.
Further structure or exponential
integrability is required to make
that inference here.

There is also a distinct physical
limitation, visible already at order
two. The exact total-covariance
identity is
\begin{equation}
 \operatorname{Cov}_\mu(S_i,S_l)
  =\mathbb E_\mu\operatorname{Cov}_\mu(S_i,S_l\mid C)
    +\operatorname{Cov}_{\rho_0}
       \bigl(\mathbb E_\mu[S_i\mid C],
                         \mathbb E_\mu[S_l\mid C]\bigr).
 \label{eq:f15v16:retained-correlation}
\end{equation}
For separated \(i,l\), the first
term is controlled by this version.
The second, retained-field term is
not. Conditional decorrelation of
the eliminated variables is not
unconditional physical clustering.

The next upstream task is to control
normalization and an intensive
interaction norm for generated source
weights, without allowing the source
volume cost to grow with the entire
regulator. More fixed-order
coefficient statements alone would
not supply that norm. The selected
\(p\le\log\beta\) sequence and its
polynomial chart thresholds still
do not constitute a calibrated
interacting renormalization flow.
Continuum construction and a
positive physical Hamiltonian gap
remain unproved.

\subsection{Checks and predecessor preservation}

The certificate is
\begin{center}\small
\path{scripts/verify_ym_f15_v16_joint_sources.py}.
\end{center}
It checks 36 exact joint tilts and
324 vector likelihood-transport pairs
in four positive finite models with
correlated exterior distributions.
There are 15,552 checks that a
restricted interior mean is independent
of remote source parameters. In each
model the full source mean does
change under a remote tilt, so that
this conditional distinction is
genuinely tested.

An independent formal logarithmic
moment series and the partition
formula agree for 320 mixed conditional
cumulants through total order five.
Seven repeated-direction finite
differences check annihilation of
one-source terms and their exact
coefficient sums. The first six
cumulant majorants are
\(1,2,6,26,150,1082\).
Twenty-four separated-carrier/scale
checks and 30 exact normalization
examples are also included.
These are finite algebra diagnostics,
not Wilson simulations, proof-assistant
formalization, or mass-gap certification.

The certificate has 12,983 bytes and SHA--256
\begin{center}\scriptsize\ttfamily
6EA0979AA6A58B4574755E3811DB76379CF88DD00852263246179E53F9A30F59.
\end{center}
The predecessor V15 has 339,639 bytes and SHA--256
\begin{center}\scriptsize\ttfamily
AE94F51AE1DF48AEE268235B24D4F8C78C2B5B617867DF567F777B23FAA0FF33.
\end{center}
Reversing this terminal insertion and title
change recovers V15 exactly. V16 replaces
it as the sole active main note. The
11 protected sources and 58 earlier
protected certificates remain unchanged.
All earlier results and limitations
are retained in the mathematical body.

\begin{firewall}
V16 controls joint conditional source
additivity and every fixed mixed
conditional cumulant order for a
growing separated family, uniformly
over the declared real source cube
in annealed norms. It does not
normalize a product approximation
for free, control all source orders
in an intensive interaction norm,
establish retained-field clustering,
construct the interacting continuum,
or prove the full Yang--Mills mass gap.
\end{firewall}

\section{V17: normalized positive local factors}
\label{f15v17:section}

V16 controls a mean logarithmic error
but correctly stops short of normalizing
a product approximation. Here we close
that normalization step for the same
separated source families. The additional
structure is positivity of the uncentred
action insertions. Relative to an
appropriate fixed reference law, the
source changes are penalties whose
unnormalized factors lie in \((0,1]\).
Both normalizers then have quantitative
lower bounds. These two facts convert
the logarithmic estimate into total
variation and relative entropy in
\emph{both} directions.

For the full real source cube, the
reference is one fixed upper-corner
tilted law. For the negative source
cube, the reference can be the original
Wilson output law itself. These are
different constructions and are stated
separately. Neither construction asserts
that its reference law is a product
law, a local Gibbs density, or a
physically clustering continuum state.

The active V16 file and all protected
source and certificate hashes were
rechecked before this continuation.
The scheduled companion review at V15
is supplemented below by an immediate
review of the three newer attachments
supplied during this continuation.
The next scheduled five-version
checkpoint is still V20.

\subsection{Positive energies and the fixed reference}

Retain all geometric and source-window
hypotheses of V16, including the
separated translates, \(\tau=1/8\),
\(3\tau<R_{\rm src}\), and the unchanged
completed blocking map. Set
\begin{equation}
 W_i=\sum_x b_{i,x}T_x
       =\beta\sum_f\omega_i(f)S_{W,f}\ge0,\qquad
 c_i=\mathbb E_\mu W_i,\qquad S_i=W_i-c_i .
 \label{eq:f15v17:positive-energy}
\end{equation}
The \(S_i\) and the laws
\(\mu_\theta,\rho_\theta\) are exactly
those of V16. In particular no
source-centering convention is changed.

Translation invariance of the baseline
Wilson law gives
\begin{equation}
 \bar T=\mathbb E_\mu T_x
          =\frac{\beta\mathbb E_\mu S_{\rm W}}{V},
 \qquad c_i=n_p\bar T,\qquad
                   \sum_i c_i=mn_p\bar T .
 \label{eq:f15v17:mean-energy}
\end{equation}
Here \(S_{\rm W}\) denotes the total
unscaled Wilson action, not a centred
source. The equality follows from
\(\sum_xT_x=\beta S_{\rm W}\) in
\eqref{f15v1:cell}. On the selected
path, \(\bar T\to9/2\); this is an
inherited energy-saturation property,
not a statement about every tilted
law. Eventually \(\bar T\le5\), and
also \(k\le5\).

For the full cube choose once and for all
\[
       a=\tau(1,\ldots,1),\qquad
       \nu=\rho_a,\qquad d=\theta-a .
 \]
Thus \(d_i\in[-2\tau,0]\) whenever
\(\theta\in[-\tau,\tau]^m\). Define
\begin{equation}
 \begin{split}
 A_\theta(C)&=\mathbb E_{\mu_a}
                   [e^{\sum_i d_iW_i}\mid C],\\
 q_{i,\theta_i}(N_i)&=\mathbb E_{\mu_a}
                   [e^{d_iW_i}\mid N_i],\qquad
 q_\theta(C)=\prod_iq_{i,\theta_i}(N_i),\\
 \mathcal Z_\theta&=\mathbb E_\nu A_\theta,\qquad
 z_\theta=\mathbb E_\nu q_\theta,\qquad
 d\sigma_\theta=\frac{q_\theta}{z_\theta}\,d\nu .
 \end{split}
 \label{eq:f15v17:factors}
\end{equation}
These are exact conditional expectations
under one common \(\mu_a\). They are
not products of separately chosen
reference laws. Positivity of \(W_i\)
and \(d_i\le0\) imply
\begin{equation}
        0<A_\theta\le1,\qquad
        0<q_{i,\theta_i}\le1,\qquad
        0<q_\theta\le1 .
 \label{eq:f15v17:factor-range}
\end{equation}
At a finite regulator all lower values
are positive since the energies are
bounded. The source-centering scalars
cancel in the normalized change of law,
so
\begin{equation}
       d\rho_\theta=\frac{A_\theta}{\mathcal Z_\theta}\,d\nu,
 \qquad
       \mathcal Z_\theta
            =e^{\sum_i d_i c_i}\frac{Z(\theta)}{Z(a)} .
 \label{eq:f15v17:true-law}
\end{equation}
In particular both normalizers in
\eqref{eq:f15v17:factors} remain
explicit.

The local potentials
\begin{equation}
              V_{i,\theta_i}(N_i)=-\log q_{i,\theta_i}(N_i)
                                                       \ge0
 \label{eq:f15v17:local-potentials}
\end{equation}
depend on the indicated near output
and its one source parameter. They
are gauge invariant up to choice
of conditional versions. The
reference \(\mu_a\) is fixed for
this source family; their coefficients
can depend on that entire fixed
family. No family-independent
interaction norm has been asserted.

\subsection{The logarithmic comparison uses joint and near means}

Keep \(\mathcal R_m,\mathfrak e_m\)
from \eqref{eq:f15v16:budgets}.
Define
\begin{equation}
 r_\theta(C)=\log A_\theta(C)-\log q_\theta(C),\qquad
 \delta_\theta
       =2\|\theta-a\|_1\mathcal R_m\mathfrak e_m
                         \le4m\tau\mathcal R_m\mathfrak e_m .
 \label{eq:f15v17:log-defect}
\end{equation}

\begin{proposition}[A log estimate for the actual local factors]
\label{f15v17:log-defect}
For every \(\theta,\zeta\in[-\tau,\tau]^m\),
\begin{equation}
                      \|r_\theta\|_{L^1(\rho_\zeta)}
                                                \le\delta_\theta .
 \label{eq:f15v17:log-bound}
\end{equation}
The factors in this assertion condition
on \(N_i\), not on the entire output.
\end{proposition}

\begin{proof}
Let \(v\in[0,1]\), and consider the
joint and single-coordinate paths
\[
        b(v)=a+vd,\qquad b_i(v)=a+vd_i e_i .
\]
Both stay inside \([-\tau,\tau]^m\).
Although \(|d_i|\) can be \(2\tau\),
the actual total source coordinate
along either path has magnitude
at most \(\tau\). Thus V16's
local-mean theorem applies throughout.

Differentiate the logarithms of
the conditional penalty expectations
along these paths. The exact
derivatives are respectively
\[
 \sum_i d_i\mathbb E_{\mu_{b(v)}}[W_i\mid C],
 \qquad
 d_i\mathbb E_{\mu_{b_i(v)}}[W_i\mid N_i].
\]
This differentiates both numerator
and normalizer of each conditional
law. It is not a frozen-exterior
calculation.

By Theorem~\ref{f15v16:local-means},
\(\mathbb E_{\mu_{b(v)}}[S_i\mid C]\)
is within \(\mathfrak e_m\), in its
own \(L^2\) output law, of
\(G_i(a_i+vd_i,N_i)\).
The same theorem for \(b_i(v)\),
followed by conditional expectation
onto \(N_i\), gives
\[
 \left\|\mathbb E_{\mu_{b_i(v)}}[S_i\mid N_i]
             -G_i(a_i+vd_i,N_i)\right\|_{L^2(\rho_{b_i(v)})}
                                           \le\mathfrak e_m .
\]
The two approximating functions
are identical, because \(G_i\)
is independent of remote source
parameters. Adding \(c_i\) to
both means changes \(S_i\) to
\(W_i\) and cancels in their
difference.

Proposition~\ref{f15v16:transport}
moves each of these two errors
to \(L^1(\rho_\zeta)\) at cost
\(\mathcal R_m\). The difference
of the two raw-energy means
therefore has norm at most
\(2\mathcal R_m\mathfrak e_m\).
All penalty expectations equal
one at \(v=0\). Integration
from zero to one and summation
with \(|d_i|\) give the result.
The finite-regulator energies
are bounded, so common conditional
versions and differentiation under
the integrals are legitimate.
\end{proof}

\subsection{Both normalizers have a lower bound}

Define the finite cost
\begin{equation}
 \mathcal H_m^+
       =2\tau mn_p\bar T+2k mn_pJ(2\tau).
 \label{eq:f15v17:corner-cost}
\end{equation}

\begin{proposition}[Reference energy and normalized mass]
\label{f15v17:normalizers}
For every \(\theta\) in the source cube,
\begin{align}
 \mathbb E_{\rho_a}[-\log A_\theta]
                   &\le\mathcal H_m^+,\notag\\
 \mathbb E_{\rho_a}\sum_i V_{i,\theta_i}
                   &\le\mathcal H_m^+,\notag\\
       e^{-\mathcal H_m^+}
                   &\le\mathcal Z_\theta,\ z_\theta\le1 .
 \label{eq:f15v17:normalizers}
\end{align}
Eventually on the selected path,
\begin{equation}
            \mathcal H_m^+
                \le\frac53 mn_p\le\frac{20}{3}mp^4 .
 \label{eq:f15v17:cost-rate}
\end{equation}
\end{proposition}

\begin{proof}
First bound the raw energy at the
corner, without assuming that
the homogeneous thermal estimate
automatically holds there. Set
\[
                  f(u)=\log Z(u,\ldots,u).
\]
At the finite regulator this is
a differentiable convex function.
Its derivative at \(\tau\) is
\(\mathbb E_{\mu_a}\sum_iS_i\).
The source budget and \(f(\tau)\ge0\)
give
\[
 f'(\tau)\le\frac{f(2\tau)-f(\tau)}{\tau}
                        \le\frac{k mn_pJ(2\tau)}{\tau}.
\]
Only the partition function is
evaluated at \(2\tau\); no local
convexity assertion at that
larger source is being used.
Consequently
\begin{equation}
 \mathbb E_{\mu_a}\sum_iW_i
        \le mn_p\bar T+\frac{k mn_pJ(2\tau)}{\tau}.
 \label{eq:f15v17:corner-energy}
\end{equation}

Write \(\lambda_i=-d_i\in[0,2\tau]\).
Conditional Jensen gives
\[
 \begin{split}
 -\log A_\theta(C)
       &\le\mathbb E_{\mu_a}
                    \left[\sum_i\lambda_iW_i\,\middle|\,C\right],\\
 -\log q_{i,\theta_i}(N_i)
       &\le\lambda_i\mathbb E_{\mu_a}[W_i\mid N_i].
 \end{split}
\]
Integration and
\eqref{eq:f15v17:corner-energy}
prove the first two assertions.
A second application of Jensen,
now to \(e^{-x}\), gives
\[
 \mathbb E_\nu A_\theta
       \ge e^{-\mathbb E_\nu[-\log A_\theta]},
 \qquad
 \mathbb E_\nu q_\theta
       \ge e^{-\mathbb E_\nu\sum_iV_{i,\theta_i}} .
 \]
These are the two lower bounds.
The upper bounds follow from
\eqref{eq:f15v17:factor-range}.

Finally \(\bar T,k\le5\) eventually,
and \(J(1/4)\le1/24\). Since
\(\tau=1/8\),
\[
           2\tau\bar T+2kJ(2\tau)
                            \le\frac54+\frac5{12}=\frac53 .
\]
Use \(n_p\le4p^4\) to finish.
\end{proof}

The second line of
\eqref{eq:f15v17:normalizers}
is a uniform \emph{averaged}
penalty cost per selected source
cell after division by \(mn_p\).
It is not a supremum over coarse
configurations, a spatial derivative
bound, or an analytic interaction norm.

\subsection{A positive-factor normalization lemma}

\begin{proposition}[Two-sided entropy from a controlled penalty]
\label{f15v17:positive-lemma}
Let \(\nu\) be a probability, let
\(A,q\in(0,1]\), and suppose
\[
 \mathbb E_\nu[-\log A]\le H,\qquad
 \mathbb E_\nu[-\log q]\le H,\qquad
             \mathbb E_\nu|\log A-\log q|\le\delta .
\]
Put \(F=\mathbb E_\nu A\),
\(G=\mathbb E_\nu q\), and
\(dP=A\,d\nu/F\), \(dQ=q\,d\nu/G\).
Then
\begin{equation}
 \begin{split}
 \|P-Q\|_{\rm TV}&\le e^H\delta,\\
 \max\{D(P\Vert Q),D(Q\Vert P)\}
                                      &\le2e^H\delta .
 \end{split}
 \label{eq:f15v17:positive-lemma}
\end{equation}
No smallness assumption on \(e^H\delta\)
is needed for these inequalities.
\end{proposition}

\begin{proof}
Jensen gives \(F,G\ge e^{-H}\).
The exponential function is
one-Lipschitz on \((-\infty,0]\).
Therefore
\[
       |A-q|\le|\log A-\log q|,\qquad
       \mathbb E_\nu|A-q|\le\delta,\qquad |F-G|\le\delta .
\]
Separating the unnormalized error
and the scalar normalization gives
\[
 \frac12\mathbb E_\nu\left|\frac AF-\frac qG\right|
      \le\frac{\mathbb E_\nu|A-q|+|F-G|}{2F}
      \le\frac{\delta}{F}\le e^H\delta .
\]

Let \(r=\log A-\log q\).
Since \(A,q\le1\) and \(F,G\ge e^{-H}\),
both normalized densities are at
most \(e^H\) relative to \(\nu\).
Thus
\(\mathbb E_P|r|,\mathbb E_Q|r|\le e^H\delta\).
Moreover
\[
 \log(G/F)\le\log(1+\delta/F)\le e^H\delta,
 \qquad
 \log(F/G)\le\log(1+\delta/G)\le e^H\delta .
\]
The identities
\[
 D(P\Vert Q)=\mathbb E_P r+\log(G/F),\qquad
 D(Q\Vert P)=-\mathbb E_Q r+\log(F/G)
\]
prove the two entropy bounds.
The assumed integrability makes
all these expressions finite.
\end{proof}

Both hypotheses are essential to
this application: the unnormalized
factors are bounded above by one,
and their normalizers are bounded
below. V16's small-log-error
counterexample did not have these
uniformly paid bounds. The lemma
does not reverse that counterexample
by an unsupported exponential
integrability assumption.

\subsection{A normalized local-factor family for the full source cube}

\begin{theorem}[Actual normalization with a common corner reference]
\label{f15v17:corner-theorem}
Under the hypotheses above, the
probabilities \(\sigma_\theta\) of
\eqref{eq:f15v17:factors} satisfy
\begin{equation}
 \begin{split}
 \|\rho_\theta-\sigma_\theta\|_{\rm TV}
       &\le \mathcal E_\theta^+,\\
 \max\{D(\rho_\theta\Vert\sigma_\theta),
                         D(\sigma_\theta\Vert\rho_\theta)\}
       &\le2\mathcal E_\theta^+,\qquad
 \mathcal E_\theta^+=e^{\mathcal H_m^+}\delta_\theta .
 \end{split}
 \label{eq:f15v17:corner-theorem}
\end{equation}
The reference \(\rho_a\) is fixed
for the entire cube, and
\(\sigma_a=\rho_a\) exactly.
\end{theorem}

\begin{proof}
Use Proposition~\ref{f15v17:positive-lemma}
with \(\nu=\rho_a\), the actual
\(A_\theta,q_\theta\), and
\(H=\mathcal H_m^+\).
The penalty costs are
Proposition~\ref{f15v17:normalizers}.
The logarithmic error is
\eqref{eq:f15v17:log-bound} with
\(\zeta=a\). The normalized \(P\)
is precisely \(\rho_\theta\) by
\eqref{eq:f15v17:true-law}, and
the normalized \(Q\) is the
defined \(\sigma_\theta\).
At \(\theta=a\) every factor is
one, proving the last assertion.
\end{proof}

This theorem controls an actually
normalized probability, not merely
the logarithm of a proposed weight.
The reverse entropy estimate does
not apply V16's likelihood-transport
bound to \(\sigma_\theta\), which
is not an actual source tilt.
Instead its density is bounded
by \(e^{\mathcal H_m^+}\) relative
to the common reference, as in
the positive-factor lemma.

\subsection{Negative sources relative to the original law}

There is a direct baseline version
which does not use the corner law.
For \(\theta\in[-\tau,0]^m\), define
\begin{equation}
 \begin{split}
 q^0_{i,\theta_i}(N_i)
       &=\mathbb E_\mu[e^{\theta_iW_i}\mid N_i],\\
 q^0_\theta(C)&=\prod_iq^0_{i,\theta_i}(N_i),\qquad
 z^0_\theta=\mathbb E_{\rho_0}q^0_\theta,\qquad
 d\sigma^0_\theta=\frac{q^0_\theta}{z^0_\theta}\,d\rho_0,\\
 \mathcal H_m^-&=\tau mn_p\bar T,\qquad
 \delta_\theta^0=2\|\theta\|_1\mathcal R_m\mathfrak e_m,\qquad
 \mathcal E_\theta^-=e^{\mathcal H_m^-}\delta_\theta^0 .
 \end{split}
 \label{eq:f15v17:baseline-factors}
\end{equation}

\begin{theorem}[Normalized local source increments at the baseline]
\label{f15v17:negative-theorem}
For the negative source cube,
\begin{equation}
 \begin{split}
 \|\rho_\theta-\sigma^0_\theta\|_{\rm TV}
                        &\le\mathcal E_\theta^-,\\
 \max\{D(\rho_\theta\Vert\sigma^0_\theta),
                       D(\sigma^0_\theta\Vert\rho_\theta)\}
                        &\le2\mathcal E_\theta^- .
 \end{split}
 \label{eq:f15v17:negative-theorem}
\end{equation}
Eventually
\(\mathcal H_m^-\le(5/8)mn_p\le(5/2)mp^4\).
\end{theorem}

\begin{proof}
Repeat the logarithmic comparison
with reference \(a=0\), joint path
\(v\theta\), and partial paths
\(v\theta_i e_i\).
They remain in V16's source cube,
so the identical projection and
transport proof gives
\(\delta_\theta^0\).
All penalty factors again lie
in \((0,1]\).
Conditional Jensen now bounds
both mean negative logarithms
by
\[
       \mathbb E_\mu\sum_i(-\theta_i)W_i
                                  \le\tau mn_p\bar T .
\]
No corner energy estimate is
needed. Apply the positive-factor
lemma. The eventual cost bound
uses \(\bar T\le5\) and \(n_p\le4p^4\).
\end{proof}

This is a normalized product of
near-output source factors relative
to the original output law. The
reference \(\rho_0\) itself is not
claimed to have a local density.

\subsection{The normalization cost is paid on the existing sequence}

\begin{theorem}[Uniform total variation and both entropies]
\label{f15v17:asymptotic}
Use the V16 sequence
\[
 t=\log\beta,\qquad
 r=\max\{176,\lceil p^{\nu+8}t^6\rceil\},\qquad
                  2\le m\le\lfloor t\rfloor .
\]
For every fixed \(A>0\), all three
errors in \eqref{eq:f15v17:corner-theorem}
are \(o(\beta^{-A})\), uniformly
over the full source cube and
allowed placements. The same is
true of the three errors in
\eqref{eq:f15v17:negative-theorem}
on the negative cube. The near
windows still satisfy
\(m(rp)^4/M^4\to0\).
\end{theorem}

\begin{proof}
The corner estimate is the larger
one. Since
\(\delta_\theta\le4m\tau\mathcal R_m\mathfrak e_m\),
it suffices to bound
\[
                    e^{\mathcal H_m^+}\mathcal R_m\mathfrak e_m .
\]
We use the explicit rates, not
just the assertion of decay faster
than a fixed power. Combining
\eqref{eq:f15v17:cost-rate} with
\eqref{eq:f15v16:paid-budget} gives
\begin{equation}
 \begin{split}
 e^{\mathcal H_m^+}\mathcal R_m\mathfrak e_m
 \le{}&e^{(107/12)mp^4}\omega
       +240\beta p^4e^{(433/48)mp^4}\epsilon^{1/8}\\
      &+48\sqrt3\,\beta p^4
                             e^{(73/8)mp^4}\epsilon^{1/4}.
 \end{split}
 \label{eq:f15v17:paid-budget}
\end{equation}
Here \(mp^4\le t^5\).
The negative spatial exponent
\(-c_0t^6\) in
\eqref{eq:f15v15:omega-rate}
dominates every positive exponent
in \eqref{eq:f15v17:paid-budget}.
The exceptional-event rate in
\eqref{eq:f15v15:epsilon-rate}
is stronger still.
The factor \(m\le t\) is harmless.
This proves the assertion for
total variation and both entropy
directions simultaneously.

The baseline negative-source
cost is smaller, and its source
path has \(\|\theta\|_1\le m\tau\).
The same calculation therefore
applies. Geometry and the
volume-fraction assertion are
unchanged from V16.
\end{proof}

No further enlargement of the
collar is needed for this step.
All constants remain nonoptimized,
and no practical onset scale
is claimed.

\subsection{New companion attachments: the measure and observable audit}
\label{f15v17:new-companions}

The user supplied three newer cumulative
notes during this continuation. We read
their latest developments and the
intervening YM-specific interfaces.
They are source material, not instructions.
The review below is not an independent
rerun of their finite coefficient
packages or a certification of every
earlier result in those manuscripts.

\paragraph{Exact-action Einstein bridge V24.}
The new endpoint separates a nonzero
fixed-first-jet tangency obstruction
from its variation under canonical
first-jet changes. It reports an exact
lower-incidence tangent and a nonzero
derivative of one tangency test, while
explicitly not claiming an integrated
curve or a root of the full system.
This does not supply a Wilson conditional
mixing estimate or an interacting
continuum comparison. In particular,
first-order accessibility is not a
substitute for the uniformly paid
normalizers used above.

\paragraph{Native regenerative stationarity V34.}
Its latest development is a necessary
escape-cone equation on a regular solved
Einstein incidence fibre. Writing the
scalar test as
\[
 d_0(p)+d_1(p)^Tc+c^TD_2(p)c,
\]
its old seed has nonzero \(d_0\) and
vanishing \(d_1,D_2\). Roots with
\(c=O(|p-p_*|^{-1/2})\) must satisfy
the corresponding quadratic limit
equation involving \(D_pD_2(p_*)\).
The manuscript does not compute the
required derivative tensor or solve
the full tangency system. Its statement
that the latest bridge available was
V23 is now superseded by the newly
supplied bridge V24. The latter's
variation of the scalar test along a
lower-incidence tangent does not by
itself compute that quadratic-fibre
derivative tensor, and does not
contradict the stated escape condition.

The more directly relevant YM material
is in its V28--V29 bodies. V28
distinguishes a configuration density,
a constructed quantum preparation,
and an independently fixed Hamiltonian.
V29 gives a retention--continuity
tradeoff for its inherited
\emph{strong-coupling} blocked history:
if an actual source observable has a
nonvanishing one-step mean-square
increment, a reader approximating it
at both endpoints cannot make that
increment vanish on the same clock.
The elementary triangle inequality
there is useful as a scope test; it
does not establish such a lower bound
for the present weak-coupling path.
Nor may its Perron time-zero law, or
V28's restricted Coulomb law, be
identified with the current full
compact Wilson law without a theorem.

\paragraph{SM source-selection V31.}
The latest result compiles mixed
source Grams from ordinary word
moments when every inserted typing
operation is already represented in
the same history algebra. Its explicit
external-operation counterexample
prevents extending this conclusion
to an added projector or router.
The intervening V23 YM probe is
expressly an additional comparison
channel applied to an older restricted
Coulomb law, not an operation already
executed by the YM source. Neither
construction supplies a new
unconditional retained-field
covariance bound here.

For the present continuation the
concrete consequence is to keep
\(q_{i,\theta_i}\) as the declared
conditional marginal under the
\emph{same} \(\mu_a\). No external
quantum preparation, reconstructed
target measure, or changed source
observable is inserted into the proof.
Local factors are not evidence that
the reference field is independent.
The next estimates must concern its
actual retained observables and must
eventually be tied to a specified
physical-time reconstruction. These
attachments sharpen that obligation;
they do not remove it.

\paragraph{Source identities.}
The reviewed files were read in place
from the supplied Downloads directory
and were not edited. Their basenames,
byte counts and SHA--256 identities are:
\begin{center}\small
\path{Renewal_Exact_Action_Einstein_Bridge_Research_Notes_V24.tex}
\\836,365 bytes\\[2pt]
{\scriptsize\ttfamily
A9B80D1A2689489CBD93E14EBABE60E814A0F2CAB79EC000D36E4AF0C3BEBB7D}
\end{center}
\begin{center}\small
\path{Renewal_Native_Regenerative_Stationarity_Closure_Notes_V34.tex}
\\1,175,242 bytes\\[2pt]
{\scriptsize\ttfamily
08BC4B7DE876D01122109B2B92EBFDD3B11B5191A0D09225AFC642578FA23ECC}
\end{center}
\begin{center}\small
\path{Renewal_SM_Source_Selection_Research_Notes_V31.tex}
\\1,095,849 bytes\\[2pt]
{\scriptsize\ttfamily
34312755DFB63E2E6C0F69995C4F545362D88738D393319FEC842EE55BDEF0E4}
\end{center}
Older source versions remain preserved.
The newer versions are the current
interfaces for the next companion
review; the V15 audit is retained
as a historical record, not silently
rewritten.

\subsection{What the construction does not identify}

The full-cube theorem includes
\(\theta=0\), so its \(\sigma_0\)
approximates \(\rho_0\). It does
not follow that one may replace
the common reference \(\rho_a\)
by \(\rho_0\) inside every full-cube
product without paying a new
reweighting error. Small total
variation before an unbounded
reweighting need not remain small.
The negative-cube theorem is
proved directly and makes no
such replacement.

The conditional near factors
use actual marginals. They have
not been expressed as a finite
list of coarse plaquette actions.
Changing the source family can
change the corner reference and
therefore these coefficients,
even though their dependence
on the observed configuration
is local. The averaged bound
on \(\sum_iV_i\) is not the
family-independent spatial and
analytic interaction bound
needed for repeated RG steps.

Most importantly, the correlations
of the retained reference field
are still present. Neither a
product of local source factors
relative to that field nor
small two-sided source-approximation
entropy is a proof of its
unconditional exponential
clustering. The second term in
\eqref{eq:f15v16:retained-correlation}
remains uncontrolled.

The next issue is thus consistency
and intensive control as more
source supports are admitted:
can local factors be compared
across backgrounds without a
global source-volume likelihood
cost, and can that comparison
support a closed interacting
scale evolution? The present
\(m\le\log\beta\) family and
growing collar do not answer
that question. The interacting
continuum and physical Hamiltonian
mass gap remain unproved.

\subsection{Checks and predecessor preservation}

The certificate is
\begin{center}\small
\path{scripts/verify_ym_f15_v17_positive_normalization.py}.
\end{center}
It checks 52 exact finite positive-product
comparisons: 36 against a common upper
reference and 16 against the original
reference with negative sources. Forty-four
comparisons have genuinely nonzero
errors. It verifies factor ranges,
both normalizers, normalized target
identities, and total-variation
inequalities.

The logarithmic tests use rational
intervals, not floating-point
logarithms. After exact reduction
to \(1\le x<2\), write
\[
 \log x=2\sum_{j\ge0}\frac{z^{2j+1}}{2j+1},
              \qquad z=\frac{x-1}{x+1}.
\]
After \(K\) terms its positive
remainder is at most
\(2z^{2K+1}/((2K+1)(1-z^2))\).
The same construction encloses
\(\log2\) for the removed powers
of two. With \(K=16\), the
certificate verifies both entropy
directions, the logarithmic-to-density
bound, Jensen penalty costs, and
the logarithmic equivalents of
both normalizer lower bounds.

There are also 104 exact
full/near source-derivative pairs,
four corner secant-convexity
checks, and 72 growing-family
scale checks, including the
displayed cost coefficients.
These finite models diagnose
the written algebra and
normalization argument; they
are not Wilson simulations,
continuum tests, or a
mass-gap certification.

The certificate has 11,637 bytes and SHA--256
\begin{center}\scriptsize\ttfamily
DF9A04C6B39F6282808D26C2C4925F2E2622CA56214D728DAEE64F1748CDCDA2.
\end{center}
The predecessor V16 has 362,475 bytes and SHA--256
\begin{center}\scriptsize\ttfamily
630D868BCB0F61F6C8F9E27F730A785FBB59DA46C9E468E82CDC053E90BA6D40.
\end{center}
Reversing this terminal insertion
and title change recovers V16
exactly. V17 replaces it as
the sole active main note.
The 14 protected sources (including
the three newly supplied attachments) and
59 earlier protected certificates
remain unchanged.

\begin{firewall}
V17 constructs normalized local
source-factor approximations and
controls total variation and
both relative-entropy directions
for the declared separated
families and reference laws.
It does not establish an
intensive RG interaction norm,
remove retained-field correlations,
construct the interacting
continuum, or prove the
Yang--Mills mass gap.
\end{firewall}

\section{V18: fixed local chart factors and the retained correlation problem}
\label{f15v18:section}

V17 normalized products of actual
near-output conditional factors.
For the full source cube those factors
still depended on the entire source
family through its upper-corner law.
We now remove that dependence from
the factors themselves. The input is
V16's local chart mean, which was
already independent of remote source
parameters. Its primitive is the
logarithm of a fixed local partition
ratio, with every measured-density
term retained. A positivity argument
pays for normalization without
estimating this ratio under a new
reference measure.

This gives a more explicit local
source action, not a factorization
of the retained reference law.
We quantify the surviving correlation
problem and then test whether the
available scalar source and
susceptibility bounds could solve
it by themselves. An explicit
positive convolution field shows
that they cannot: all those bounds
can coexist with algebraic spatial
correlations. The diagnostic is not
a Wilson counterexample. It identifies
the missing use of spatial action
structure in the next argument.

The authoritative V17 hash was
rechecked before editing. The
three-attachment review in V17
remains the latest additional
review; the next scheduled
companion checkpoint is V20.

\subsection{Integrate the already constructed local score}

Fix the regulator, scale, collars,
framing conventions and separated
source cells of V16. Use its
\(G_i(s,N_i)\), \(B_p\),
\(\mathcal R_m\), and
\(\mathfrak e_m\). Retain the
positive energies \(W_i\) and fixed
baseline means \(c_i\) from
\eqref{eq:f15v17:positive-energy}.
For every compact configuration,
\[
 0\le W_i\le 2\beta\sum_f\omega_i(f)
       =12\beta n_p\le B_p,
 \qquad 0\le c_i\le B_p.
\]
This is a pointwise energy bound;
it does not assume a thermal
estimate under a tilted law.

On \(\Gamma_{N_i}\), let
\(\pi_i^{N_i}\) be the normalized
zero-source local chart law used
to define \(G_i(0,N_i)\): the
interior domain is unchanged,
the required exterior canonical
details are set to \(I\), and
the measured potential includes
all the determinant and Haar
terms of V13. It is a prescribed
local comparison integral, not
the actual conditional law given
\(N_i\). Define
\begin{equation}
 L_i(s,N_i)=
 \begin{cases}
 \displaystyle\int e^{sW_i}\,d\pi_i^{N_i},
                                &N_i\in\Gamma_{N_i},\\
 e^{s c_i},                      &N_i\notin\Gamma_{N_i}.
 \end{cases}
 \label{eq:f15v18:local-partition}
\end{equation}
The branch notation means membership
in the previously defined retained
chart event. The second line is
chosen to match exactly V16's
convention \(G_i=0\) off that event.
It is not an assertion about the
true conditional law there.

\begin{proposition}[A fixed local positive-energy primitive]
\label{f15v18:primitive}
For \(|s|\le\tau\),
\begin{equation}
 g_i(s,N_i):=\partial_s\log L_i(s,N_i)
       =G_i(s,N_i)+c_i,
 \qquad 0\le g_i\le B_p.
 \label{eq:f15v18:primitive}
\end{equation}
For real \(u\le v\) in this interval,
\begin{equation}
 \widehat q_i^{\,v}(u,N_i)
   :=\frac{L_i(u,N_i)}{L_i(v,N_i)}
    =\exp\!\left(-\int_u^v g_i(s,N_i)\,ds\right)
       \in(0,1].
 \label{eq:f15v18:local-factor}
\end{equation}
Holding the local geometry fixed,
this function is unchanged when
other separated source cells or
their source parameters are added
or removed. It is gauge invariant
with the same conventions as
\(G_i\).
\end{proposition}

\begin{proof}
The chart source tilt by \(sS_i\)
equals the tilt by \(sW_i\), since
\(S_i=W_i-c_i\) and the constant
cancels in that chart's normalized
mean. Differentiating its finite
positive integral therefore gives
the mean of \(W_i\), namely
\(G_i+c_i\). Off the event the
displayed identity follows from
\(L_i=e^{s c_i}\). Positivity and
the pointwise bound on \(W_i,c_i\)
give the stated range. Integration
gives the ratio formula. V16's
definition of the interior, fixed
exterior, cutoff and chart potential
does not use remote sources, which
proves the final assertion. No
regularity across the hard retained
cutoff is asserted.
\end{proof}

Thus the factors are available as
specified local measured integrals,
not as unknown near marginals under
different global source laws.
Their pointwise penalty is at most
\((v-u)B_p\), but this crude bound
would have an unacceptable
normalization cost on the current
scale sequence. The next proof
uses an averaged bound instead.

\subsection{Normalize without a new marginal-comparison estimate}

For the full cube put
\(a=\tau(1,\ldots,1)\),
\(d=\theta-a\), and keep the actual
\(A_\theta\), \(\mathcal Z_\theta\)
and reference \(\rho_a\) of V17.
Define the new objects
\begin{equation}
 \begin{split}
 \widehat q_\theta(C)
   &=\prod_i\widehat q_i^{\,\tau}(\theta_i,N_i),\\
 \widehat z_\theta&=\mathbb E_{\rho_a}\widehat q_\theta,\\
 d\widehat\sigma_\theta
   &=\frac{\widehat q_\theta}{\widehat z_\theta}\,d\rho_a,\\
 \widehat\delta_\theta
   &=\|\theta-a\|_1\mathcal R_m\mathfrak e_m,\\
 \widehat H_\theta
   &=\mathcal H_m^++\widehat\delta_\theta,\\
 \widehat E_\theta
   &=e^{\widehat H_\theta}\widehat\delta_\theta.
 \end{split}
 \label{eq:f15v18:surrogate}
\end{equation}
The hat distinguishes these fixed
chart factors from V17's actual
near-conditional factors.

\begin{theorem}[Normalized local chart action on the signed source cube]
\label{f15v18:normalization}
For every \(\theta,\zeta\in[-\tau,\tau]^m\),
\begin{align}
 \|\log A_\theta-\log\widehat q_\theta\|_{L^1(\rho_\zeta)}
       &\le\widehat\delta_\theta,
                                  \label{eq:f15v18:log-error}\\
 \mathbb E_{\rho_a}[-\log\widehat q_\theta]
       &\le\widehat H_\theta,
 &\widehat z_\theta&\ge e^{-\widehat H_\theta},
                                  \label{eq:f15v18:penalty-cost}\\
 \|\rho_\theta-\widehat\sigma_\theta\|_{\rm TV}
       &\le\widehat E_\theta,
                                  \label{eq:f15v18:TV}\\
 \max\{D(\rho_\theta\Vert\widehat\sigma_\theta),
        D(\widehat\sigma_\theta\Vert\rho_\theta)\}
       &\le2\widehat E_\theta,
                                  \label{eq:f15v18:KL}\\
 |\log\mathcal Z_\theta-\log\widehat z_\theta|
       &\le\widehat E_\theta.
                                  \label{eq:f15v18:pressure}
\end{align}
The actual source normalizer also
satisfies \(\mathcal Z_\theta\ge
e^{-\mathcal H_m^+}\).
\end{theorem}

\begin{proof}
Along \(b(v)=a+vd\), \(0\le v\le1\),
the logarithmic derivative of the
actual conditional penalty is
\(\sum_i d_i\mathbb E_{\mu_{b(v)}}[W_i\mid C]\).
The derivative of the corresponding
chart product is
\(\sum_i d_i g_i(b_i(v),N_i)\),
where \(b_i(v)\) here denotes the
\(i\)-th scalar coordinate of
\(b(v)\). All total source
coordinates remain in the original
cube. The difference in each
summand is exactly
\[
 \mathbb E_{\mu_{b(v)}}[S_i\mid C]
                          -G_i(b_i(v),N_i).
\]
Its own-law \(L^2\) norm is at
most \(\mathfrak e_m\) by
Theorem~\ref{f15v16:local-means}.
Proposition~\ref{f15v16:transport}
moves it to \(L^1(\rho_\zeta)\)
at cost \(\mathcal R_m\).
Both logarithms vanish at \(v=0\).
Integrating proves
\eqref{eq:f15v18:log-error}.
There is one approximation error
per summand, not the two errors
needed to compare two different
conditional means in V17.

Write \(r=\log A_\theta-
\log\widehat q_\theta\).
Since both factors lie in \((0,1]\),
\[
 -\log\widehat q_\theta
        =-\log A_\theta+r.
\]
Integrate under \(\rho_a\), use
V17's mean penalty bound
\(\mathbb E_{\rho_a}[-\log A_\theta]
\le\mathcal H_m^+\), and then
\eqref{eq:f15v18:log-error}.
This proves the first inequality
in \eqref{eq:f15v18:penalty-cost};
Jensen proves the second. Thus
the chart normalizer is paid for
\emph{before} any change to the
surrogate law. No source transport
estimate for that surrogate is
assumed.

Apply Proposition~\ref{f15v17:positive-lemma}
with \(H=\widehat H_\theta\).
This proves the probability and
entropy estimates. Also
\[
 |\mathcal Z_\theta-\widehat z_\theta|
       \le\mathbb E_{\rho_a}|A_\theta-\widehat q_\theta|
       \le\widehat\delta_\theta,
\]
by the same one-Lipschitz exponential
estimate. Both normalizers are at
least \(e^{-\widehat H_\theta}\).
The mean value theorem for the
logarithm proves
\eqref{eq:f15v18:pressure}.
\end{proof}

In particular the centered source
pressure difference
\(\log Z(\theta)-\log Z(a)\)
is approximated by
\(\log\widehat z_\theta-\sum_i d_i c_i\)
with error \(\widehat E_\theta\).
The source-centering scalar has
not been dropped.

\begin{proposition}[Baseline negative sources and the existing scale sequence]
\label{f15v18:negative-and-rate}
For \(\theta\in[-\tau,0]^m\), use
\(a=0\), reference \(\rho_0\),
and factors \(L_i(\theta_i,N_i)/L_i(0,N_i)\).
All conclusions of
Theorem~\ref{f15v18:normalization}
hold with
\[
 \widehat\delta_\theta^{\,0}
       =\|\theta\|_1\mathcal R_m\mathfrak e_m,
 \qquad
 \widehat H_\theta^{\,0}
       =\mathcal H_m^-+\widehat\delta_\theta^{\,0}.
\]
On the unchanged V16--V17 sequence,
every error in
\eqref{eq:f15v18:TV}--\eqref{eq:f15v18:pressure},
and its baseline counterpart, is
\(o(\beta^{-A})\) for every fixed
\(A>0\), uniformly in the allowed
placements and source cube.
\end{proposition}

\begin{proof}
Repeat the single-path proof with
\(a=0\). V17 supplies the actual
mean penalty bound \(\mathcal H_m^-\).
For the full cube,
\(\widehat\delta_\theta\le
2m\tau\mathcal R_m\mathfrak e_m\).
The V16 rates make this tend to
zero uniformly; in particular it
is eventually at most one. Hence
\[
 \widehat E_\theta
       \le e\,e^{\mathcal H_m^+}
                  2m\tau\mathcal R_m\mathfrak e_m.
\]
The explicitly paid estimate
\eqref{eq:f15v17:paid-budget}
applies unchanged. Its positive
\(O((\log\beta)^5)\) exponent is
dominated by the negative spatial
\(c_0(\log\beta)^6\) exponent,
and the exceptional-event term
decays faster still. The smaller
baseline cost is handled identically.
\end{proof}

\subsection{A quantitative reduction of the surviving covariance}

At baseline let
\(G_i^0=G_i(0,N_i)\). Denote by
\(\widehat\sigma_0\) the full-cube
surrogate of \eqref{eq:f15v18:surrogate}
at \(\theta=0\), still relative
to the upper-corner reference.
This is different from the
negative-cube baseline surrogate,
which equals \(\rho_0\) exactly
at zero source.

\begin{theorem}[The remaining covariance is that of fixed local observables]
\label{f15v18:retained-covariance}
For distinct separated cells
\(i,j\), and \(0<h\le\tau/2\),
\begin{equation}
 \begin{split}
 &\left|\operatorname{Cov}_{\mu}(S_i,S_j)
        -\operatorname{Cov}_{\widehat\sigma_0}(G_i^0,G_j^0)\right|\\
 &\quad\le
   16h^{-1}\mathcal R_m\mathfrak e_m+12hB_p^3
      +4B_p\mathfrak e_m+6B_p^2\widehat E_0.
 \end{split}
 \label{eq:f15v18:covariance-reduction}
\end{equation}
On the selected sequence the
right side can be made
\(o(\beta^{-A})\) for any fixed
\(A>0\). No spatial decay of
the second covariance is asserted.
\end{theorem}

\begin{proof}
Put \(m_i(C)=\mathbb E_\mu[S_i\mid C]\).
The total-covariance identity
\eqref{eq:f15v16:retained-correlation}
has a conditional term bounded
in absolute mean by
\(16h^{-1}\mathcal R_m\mathfrak e_m+12hB_p^3\),
the \(q=2\) case of
Theorem~\ref{f15v16:cumulants}
with \(c_3=6\).

Both \(m_i\) and \(G_i^0\) are
bounded by \(B_p\), while
\(\|m_i-G_i^0\|_{L^1(\rho_0)}
\le\mathfrak e_m\).
For a bounded \(Y\),
\[
 |\operatorname{Cov}(X,Y)|
       \le2\|Y\|_\infty\mathbb E|X|.
\]
Replacing the two conditional
means by the two \(G^0\)'s
therefore costs at most
\(4B_p\mathfrak e_m\).
Finally, for two functions bounded
by \(B_p\), changing a probability
by total variation \(d\) changes
their covariance by at most
\(6B_p^2d\): the product expectation
costs \(2B_p^2d\), and the product
of the two means costs at most
\(4B_p^2d\). Use
\eqref{eq:f15v18:TV} at zero source.

For the rate choose
\(h=\beta^{-(A+5)}\).
The \(hB_p^3\) term is then
\(O(\beta^{-A-2}(\log\beta)^{12})\).
All other terms multiply
superalgebraically small errors
by fixed powers of \(\beta\)
and logarithms. This proves the
rate assertion.
\end{proof}

The right-hand covariance involves
specified local chart functions
under a specified normalized law.
It has not disappeared. The
reference \(\rho_a\) still carries
the actual correlations and depends
on the source family. Fixing the
local factors removes one dependence,
not that reference dependence.

\subsection{Why the scalar source budget cannot supply spatial clustering}
\label{f15v18:diagnostic}

The next proposition is a logical
test of a possible shortcut, not
a model of the Wilson action.
It retains positivity, the full
arbitrary-support source inequality,
the susceptibility norm and exact
local source factorization, while
allowing algebraic correlations.
Set \(c=9/2\) and
\(J(s)=-\log(1-s)-s\).

\begin{proposition}[A positive convolution diagnostic with the same scalar budget]
\label{f15v18:convolution}
On the four-dimensional torus
\(\Lambda_L=(\mathbb Z/L\mathbb Z)^4\),
let \(a_L(z)>0\) satisfy
\(\sum_z a_L(z)=1\). Take independent
\(Y_z\sim\operatorname{Gamma}(c,1)\),
with shape \(c\) and unit rate, and
put
\begin{equation}
                 T_x=\sum_z a_L(x-z)Y_z.
 \label{eq:f15v18:gamma-field}
\end{equation}
Then \(T_x\ge0\), \(\mathbb ET_x=c\),
and, for every real array with
\(\|b\|_\infty<1\),
\begin{equation}
 \log\mathbb E e^{\sum_xb_x(T_x-c)}
       =c\sum_zJ\!\left(\sum_xa_L(x-z)b_x\right)
       \le c\sum_xJ(b_x).
 \label{eq:f15v18:gamma-budget}
\end{equation}
Its covariance kernel and norm are
\begin{equation}
 K_L(d)=c\sum_z a_L(z)a_L(z-d),\qquad
 \sum_dK_L(d)=c,\qquad
 \|\operatorname{Cov}(T)\|_{\ell^2\to\ell^2}=c.
 \label{eq:f15v18:gamma-covariance}
\end{equation}
Also \(\sum_xT_x\) has exactly the
\(\operatorname{Gamma}(cL^4,1)\)
law. Nevertheless the choice
\begin{equation}
 a_L(z)=\frac{(1+|z|_{\infty,L})^{-6}}{C_L},
 \qquad C_L=\sum_z(1+|z|_{\infty,L})^{-6}
 \label{eq:f15v18:kernel}
\end{equation}
does not have a volume-uniform
exponential covariance bound.
There is also an infinite-lattice
version with strictly algebraic
covariance lower bounds.
\end{proposition}

\begin{proof}
The unit-rate gamma density is
\(y^{c-1}e^{-y}/\Gamma(c)\) on
\(y>0\); this convention agrees
with the standard
\href{https://www.itl.nist.gov/div898/handbook/eda/section3/eda366b.htm}{NIST gamma density}.
Direct integration gives
\(\mathbb E e^{s(Y-c)}=e^{cJ(s)}\)
for \(s<1\). Independence proves
the equality in
\eqref{eq:f15v18:gamma-budget}.
Since \(J''(s)=(1-s)^{-2}>0\),
Jensen applied separately at each
\(z\), followed by the unit row
and column sums of the convolution,
proves the inequality. This includes
negative as well as positive sources.

The mean and covariance follow
from \(\mathbb EY_z=\operatorname{Var}Y_z=c\).
For every real \(u\),
\[
 \operatorname{Var}\sum_xu_xT_x
      =c\sum_z\left(\sum_xa_L(x-z)u_x\right)^2
      \le c\sum_xu_x^2.
\]
The constant vector attains equality,
giving the exact norm and row sum.
Also \(\sum_xT_x=\sum_zY_z\), so
its Laplace transform gives the
stated gamma sum law. Thus even
the complete homogeneous pressure
agrees with \(cL^4J(s)\).

Let
\(C_\infty=\sum_{z\in\mathbb Z^4}(1+|z|_\infty)^{-6}\).
The shell of radius \(r\ge1\)
has \(64r^3+16r\le80r^3\)
points, hence \(C_\infty<\infty\).
Taking one nearest representative
of each torus point gives
\(1\le C_L\le C_\infty\).
The \(z=0\) term of the
nonnegative covariance sum yields
\begin{equation}
 K_L(d)\ge \frac{c}{C_\infty^2}
                      (1+|d|_{\infty,L})^{-6}.
 \label{eq:f15v18:algebraic-lower}
\end{equation}
Choose \(|d|_{\infty,L}=\lfloor L/2\rfloor\).
For any fixed positive \(M_0,m_0\),
the lower bound eventually exceeds
\(M_0e^{-m_0|d|_{\infty,L}}\).
Thus no such uniform exponential
upper bound holds.

On \(\mathbb Z^4\), use
\(a(z)=(1+|z|_\infty)^{-6}/C_\infty\)
and independent gamma variables.
The nonnegative series defining
each \(T_x\) has finite mean and
therefore is finite almost surely;
its centered part converges in
\(L^2\) since \(a\in\ell^2\).
The same covariance calculation
gives \(K(d)=c\sum_z a(z)a(z-d)\)
and the lower bound in
\eqref{eq:f15v18:algebraic-lower}
without the torus subscript.
This is an actual infinite-volume
field with no exponential covariance
decay, not just a finite-volume
constant-mode effect.
\end{proof}

In this example take the retained
output to be \(C=(T_x)_x\), and
the local source at a site to be
\(W_i=T_{x_i}\) at distinct sites.
Independent auxiliary
variables may be adjoined and then
eliminated without changing the
example. Given \(C\), every mixed
conditional cumulant of the sources
vanishes. Under any real source tilt
with coordinates in a compact
subinterval of \((-1,1)\), the
conditional mean of \(W_i\) is
exactly its local retained coordinate.
For any reference corner \(a\),
\[
 A_\theta(C)
     =\prod_i e^{(\theta_i-a_i)T_{x_i}}.
\]
These are family-independent local
positive penalty factors, with
zero factorization and normalization
error. Jensen still supplies their
finite mean-penalty normalizer bound.
Yet the retained correlations remain
algebraic. Thus replacing V17's
factors by fixed local ones cannot
alone imply clustering.

This example is not compact-valued,
does not have Wilson's finite-range
gauge action, and is not claimed
to satisfy every microscopic
Gaussian-curvature limit inherited
by the YM programme. It disproves
only the proposed implication from
the stated scalar source budget,
susceptibility norm and local
conditional factorization to
spatial exponential decay.

The absence of decay is not an
artifact of testing unbounded
observables. Here is a bounded
local witness on the same
infinite-volume diagnostic field.

\begin{proposition}[A bounded local correlation witness]
\label{f15v18:bounded-witness}
For any fixed \(\lambda>0\), set
\(U_x=e^{-\lambda T_x}\in(0,1]\).
With the infinite-lattice kernel
of Proposition~\ref{f15v18:convolution},
\begin{equation}
 \operatorname{Cov}(U_0,U_d)
 \ge \frac{c\lambda^2e^{-2c\lambda}}
                 {(1+\lambda)^2 C_\infty^2}
                     (1+|d|_\infty)^{-6}.
 \label{eq:f15v18:bounded-lower}
\end{equation}
\end{proposition}
\begin{proof}
Let \(m=\mathbb E U_0\). The gamma
Laplace transform and independence
give
\[
 D:=\log\frac{\mathbb E(U_0U_d)}{m^2}
   =c\sum_z\log\left(
       1+\frac{\lambda^2a(z)a(z-d)}
                    {1+\lambda(a(z)+a(z-d))}\right).
\]
All series converge; one may first
truncate the defining nonnegative
energy sums and pass by bounded
convergence. For \(0\le a,b\le1\),
\(\log(1+x)\ge x/(1+x)\) gives
\[
 \log\left(1+\frac{\lambda^2ab}{1+\lambda(a+b)}\right)
   \ge\frac{\lambda^2ab}{(1+\lambda a)(1+\lambda b)}
   \ge\frac{\lambda^2ab}{(1+\lambda)^2}.
\]
Hence \(D\ge c\lambda^2
\sum_z a(z)a(z-d)/(1+\lambda)^2\).
Jensen gives \(m\ge e^{-c\lambda}\),
and \(e^D-1\ge D\). Therefore
\(\operatorname{Cov}(U_0,U_d)=m^2(e^D-1)\)
has the claimed lower bound after
retaining the \(z=0\) term.
\end{proof}

\subsection{The next estimate must use the spatial action}

The new positive result is a
normalized local chart source action
whose individual coefficients do
not depend on remote source choices.
The unresolved object is now
concrete: correlations of the
\(G_i^0\)'s under the actual retained
law, or equivalently up to the
paid error under \(\widehat\sigma_0\).
An intensive estimate for that
object must use the spatial
structure of the measured Wilson
action or a controlled scale
evolution. Another use of the
same scalar pressure inequality
cannot supply it.

No claim of smoothness across the
hard chart cutoff, of a
family-independent reference law,
or of a cutoff-uniform analytic
interaction norm is made. Nor has
the range of sources been expanded
beyond the separated
\(m\le\log\beta\) families.
These are precise remaining
limitations, not a continuum
construction or a physical
Hamiltonian mass-gap theorem.

\subsection{Checks and predecessor preservation}

The companion diagnostic is
\begin{center}\small
\path{scripts/verify_ym_f15_v18_local_chart_factors.py}.
\end{center}
It uses exact rational arithmetic
and V17's rational logarithm
enclosures. It checks 39 normalized
local-chart comparisons, 33 with
nonzero errors; 156 partition/score
checks including the hard-cutoff
branch; and 33 covariance
decompositions and transport bounds.
Two four-dimensional torus kernels
give 12 source-budget and 12
covariance quadratic-form checks.
There are also 100 exact shell-count
identities and 18 rational-log
checks for the bounded-observable
correlation lower bound. The infinite-volume
claims follow from the written
summability and covariance proof,
not from these finite tests.
The script has 11,538 bytes and
SHA--256
\begin{center}\scriptsize\ttfamily
447A66B15E9AD2F8AE68B425E12A7505AEB688D2A12651F555F09DE4229DFFFF.
\end{center}

The predecessor V17 has 388,438 bytes
and SHA--256
\begin{center}\scriptsize\ttfamily
9EF5ACAC638DFA86C194E2B094C7D63A7945F04F5B7206532E0F4BFA8DE38BCC.
\end{center}
Reversing this terminal insertion
and the version-title change
recovers those bytes exactly.
V18 replaces V17 as the sole
active main note. All 14 protected
sources and 60 earlier protected
certificates remain unchanged.

\begin{firewall}
V18 removes remote-family dependence
from the local source factors and
retains quantitative normalized-law
and pressure comparisons. It also
shows why those improvements do
not remove retained-field
correlations. The full interacting
continuum Yang--Mills construction
and physical mass gap remain
unproved.
\end{firewall}

\section{V19: spatial derivatives of the actual retained action}
\label{f15v19:section}

V18 identified a missing spatial
estimate: scalar source bounds and
normalized local source factors do
not control the correlations of the
retained field. We now work with
the logarithm of its actual
compact marginal density, not an
auxiliary source product.
The main result is an averaged
bound on distant mixed Hessian
blocks of that action. It includes
the complement of the convex
chart under the original law.

The full compact density is
differentiated first. Only then
are its conditional score means
compared with local chart means.
We never differentiate V18's
small total-variation error or
a hard chart indicator. This is
essential: neither operation would
be justified by the preceding
probability estimates.

The local Schur-action results
of f14 V41--V42 concern specified
restricted contributions and
fixed-depth expansions. They are
not repeated here. The present
estimate concerns derivatives of
the \emph{full} compact retained
marginal at growing depth, in
\(L^2\) of its own probability.
It remains weaker than a
pointwise or volume-uniform
summable interaction norm.

The V18 active file and the
protected source identities were
rechecked. The next scheduled
companion checkpoint is V20;
the additional V17 attachment
review remains in force.

\subsection{The exact coarse density and its differentiated normalizer}

Use V12's canonical compact
coordinates and measured action
\(\mathcal W_j(C,v)\) from
\eqref{eq:f15v12:measured-law}.
The rooted gauge variables have
already been integrated as their
independent normalized Haar
factor, by that theorem. Define
\begin{equation}
 R_j(C)=\int e^{-\mathcal W_j(C,v)}\,dH_v(v),
 \qquad \mathcal A_j(C)=-\log R_j(C),
 \qquad d\rho_0(C)=\frac{R_j(C)}{Z}\,dH_C(C).
 \label{eq:f15v19:action}
\end{equation}
No chart restriction occurs in
this definition. Every pivot
Jacobian and every generated
interaction remains in
\(\mathcal W_j\). The action is
relative to product coarse Haar,
not to Lebesgue measure in a
coarse exponential chart.

Fix an orthonormal Lie-algebra
basis \(E_1,E_2,E_3\) in the
normalization used earlier.
For a coarse edge \(e\), let
\[
 X_{e,a}F(C)
  =\left.\frac{d}{dt}
      F(C_1,\ldots,e^{tE_a}C_e,\ldots)\right|_{t=0}.
\]
These derivatives act at fixed
canonical details \(v\) when
applied to \(\mathcal W_j\).
Set
\begin{equation}
 U_e=(X_{e,a}\mathcal W_j)_{a=1}^3,
 \qquad
 \mathsf H_{ef}(C)
       =(X_{f,b}X_{e,a}\mathcal A_j(C))_{a,b=1}^3,
                          \quad e\ne f.
 \label{eq:f15v19:scores}
\end{equation}
Different-edge fields commute.
Thus these are the mixed
Riemannian Hessian blocks on
the product group; same-edge
connection terms are not at issue.

\begin{proposition}[The full marginal score and mixed Hessian]
\label{f15v19:exact-hessian}
At every finite regulator,
\(R_j>0\) and \(\mathcal A_j\)
are smooth. With the exact
conditional detail law at fixed
\(C\),
\begin{align}
 (X_{e,a}\mathcal A_j)_a
       &=\mathbb E[U_e\mid C],
                         \label{eq:f15v19:score-identity}\\
 \mathsf H_{ef}
       &=\mathbb E[(X_{f,b}X_{e,a}\mathcal W_j)_{a,b}\mid C]
                         -\operatorname{Cov}(U_e,U_f\mid C).
                         \label{eq:f15v19:hessian-identity}
\end{align}
Here the covariance is the
matrix of componentwise cross
covariances. If the fine anchor
distance \(d(e,f)>16p\), the
first term in
\eqref{eq:f15v19:hessian-identity}
vanishes identically.
\end{proposition}

\begin{proof}
The global inverse is smooth
and all inverse Haar Jacobians
are strictly positive on the
compact product. Thus its
measured density is positive
and smooth, and differentiating
its finite compact integral is
legitimate. Differentiating
\(-\log R_j\) once gives the
conditional mean. Differentiating
that normalized conditional
mean gives its direct derivative
minus its covariance with the
new score, proving the formulas.
The covariance term comes from
the derivative of \(R_j\); it
must not be dropped.
Every measured term has anchor
diameter at most \(16p\), by
V12. A term can contribute a
direct mixed derivative only
if it contains both coarse
edges. This proves the last
assertion.
\end{proof}

\subsection{Globally bounded compact scores and a sharper chart gradient}

The score is not a positive
action insertion, and the
nonzero-source argument of V17
is not being applied to it.
Instead its compact pointwise
bound controls the exceptional
chart contribution.

Choose \(A_*\ge A_{\rm prim}\)
large enough to bound component
derivative sums through order
three of the compact primitive
maps used in V13, also including
left group translation and its
exponential tangent insertion.
Take smooth extensions to fixed
neighborhoods of their compact
domains, as in V13. These are
finite constants for the fixed
completed block. Put
\begin{equation}
 B_*=5A_*,\qquad
 \nu_*=5+\lceil3\log_2B_*\rceil,\qquad
 D_p^*=2^{28}B_*^{12}p^{\nu_*},\qquad
                         B_p^*=16\beta D_p^*.
 \label{eq:f15v19:global-constants}
\end{equation}
The old \(D_p,h_p,a_p,\alpha_p\)
are unchanged. In particular
\(\nu_*\) need not equal \(\nu\).
No numerical enclosure of
\(A_*\) is claimed.

\begin{proposition}[Compact score and local derivative bounds]
\label{f15v19:score-bounds}
For \(\beta\ge1\), all compact
inputs satisfy
\begin{equation}
             |U_e|\le B_p^*.
 \label{eq:f15v19:global-score}
\end{equation}
Its detail and retained carriers
lie within fine anchor radius
\(16p\) of \(e\). In a fixed
frame on V13's small chart,
\begin{equation}
          \|D_xU_e\|_{
                   \mathrm{HS}}\le16\beta D_p.
 \label{eq:f15v19:score-gradient}
\end{equation}
Only the interior detail
coordinates are differentiated
in this last formula.
\end{proposition}

\begin{proof}
Repeat the finite-arity chain
rule estimate of V13 with the
enlarged primitive bank.
One inverse layer multiplies
the maximum of the first jet,
the square root of the second
jet and the cube root of the
third jet by at most \(B_*\).
There are at most \(j\) inverse
layers and four additional
observable/input layers.
The order-at-most-three jet
sum of a local term is therefore
bounded by \(B_*^{3(j+4)}\).
Terms incident to a specified
coarse coordinate can be counted
in a radius \(16p\) box. The
Wilson and all-level Jacobian
term count is bounded by
\(2^{24}p^5\), using \(j\le p\).
Since \(p=2^j\), this product
is at most \(D_p^*\). The
factor \(\beta+1\le2\beta\)
and the three vector components
are absorbed by \(B_p^*\).
This argument is global in
the compact inputs; it does
not require their logarithms
to be small. Differentiation
does not enlarge term supports.

For the sharper bound, write
the retained edge as \(C_e=e^{b_e}\)
in the fixed small frame. The
left tangent insertion is a
linear combination of its
three exponential-coordinate
derivatives with coefficient
absolute sum at most four.
Indeed the inverse exponential
differential on \(|b_e|\le1/8\)
has operator norm at most two,
and the three-component
\(\ell^1\) norm is at most
\(2\sqrt3<4\).
The mixed derivative row bound
of V13 gives at most
\(2\beta D_p\) for each such
coordinate row. Coefficients
of the tangent insertion
depend only on the fixed
retained edge, not on \(x\).
Each of the three score
components thus has detail
gradient norm at most
\(8\beta D_p\). Their combined
Hilbert--Schmidt norm is at
most \(8\sqrt3\beta D_p
\le16\beta D_p\).
Terms not involving the interior
have zero detail derivative;
the interior exponential Haar
term has zero retained derivative.
Consequently the use of V13's
full measured interior potential
in this calculation is exact.
\end{proof}

\subsection{Frames are frozen for differentiation, not differentiated}
\label{f15v19:frames}

A fixed coarse gauge change
\(h=(h_x)\), lifted constantly
on fine \(p\)-cells, preserves
the canonical slice and its
product detail Haar measure.
The measured action is invariant
when \(C,v\) are transformed
together. Under it the score
at \(e=(x,\mu)\) changes by
the orthogonal adjoint rotation
at its starting vertex. Mixed
different-edge Hessian blocks
change by the corresponding
orthogonal rotations on their
two indices. Their Frobenius
norms are invariant.

For each actual \(C\), we may
therefore evaluate these vectors
and blocks in V13's near-output
frame \(H_{N_e}(C)\). At that
point the chosen frame is held
fixed while taking the group
derivatives. No derivative of
the map \(C\mapsto H_{N_e}(C)\)
is included. In particular this
does \emph{not} identify
\(X_f(H_{N_e}U_e)\) with a
rotated Hessian. The former
would contain a derivative of
the frame.

There is also no change of
variables \(C\mapsto C^{H_{N_e}(C)}\)
in the probability integral.
The output probability remains
\(\rho_0\) on the original
\(C\)'s. At fixed \(C\), the
detail-coordinate transformation
is a fixed Haar-preserving
transformation, so conditional
expectations can be evaluated
in that frame. This is the
same use of a pinned frame
as in V14, now for covariant
score components instead of
an invariant scalar source.

\subsection{Actual coarse scores have local predictors}

Center the V12--V14 geometry
at the anchor of \(e\), and
assume
\[
 r\ge192,\qquad L=2p(r+5),\qquad
 M\ge4L,\qquad \beta\ge\beta_p.
\]
Use the translated \(N_e,J_e,O_e\)
and frame \(H_{N_e}\).
Let \(\Gamma_e\) include the
required half-chart retained
and exterior inputs for the
interior potential and the
score carrier. The added score
carrier is central, inside
the same patch, so the thermal
entry event of V13 already
implies all these conditions.
Keep exactly its bound
\(\epsilon=\epsilon_{p,r,\beta}\).
The event \(D_e\) is its
interior full-chart event.

Write \(U_e^H\) for the
score evaluated in this
frozen frame, and put
\[
 Z_e=(C,v_{O_e}^{H_{N_e}}),\qquad
 m_e(Z_e)=\mathbb E[U_e^H\mid Z_e].
\]
On the retained half-chart
event define \(P_e(N_e)\) as
the normalized local chart
mean of \(U_e^H\) with required
exterior canonical details set
to \(I\). Set \(P_e=0\) outside
that retained event. Thus
\(|P_e|\le B_p^*\). Define
\begin{equation}
 \begin{split}
 \Omega^*_{p,r,\beta}
   &=2^{20}\beta D_p r^4p^5
                   e^{-\alpha_p(r-96)p},\\
 E^*_{p,r,\beta}
   &=\Omega^*_{p,r,\beta}
       +5B_p^*\epsilon^{1/4}
       +\sqrt3 B_p^*\epsilon^{1/2}.
 \end{split}
 \label{eq:f15v19:errors}
\end{equation}

\begin{theorem}[A full-law local predictor of the retained action gradient]
\label{f15v19:local-score}
With the Euclidean norm on
the three score components,
\begin{equation}
 \|m_e-P_e(N_e)\|_{L^2(\mu_{Z_e})}
              \le E^*_{p,r,\beta}.
 \label{eq:f15v19:exterior-score}
\end{equation}
Consequently the frozen-frame
components of the actual
\((X_{e,a}\mathcal A_j)_a\)
are within the same bound
of \(P_e(N_e)\) in
\(L^2(\rho_0)\). Rotating
back gives a predictor depending
only on \(N_e\) in the
original edge frame as well.
\end{theorem}

\begin{proof}
At fixed \(Z_e\), V12 gives
the exact unrestricted compact
interior law. The score is
supported in a central radius
\(16p\) set and has no direct
dependence on exterior details.
Its retained dependence is
entirely in \(N_e\).

For the normalized restricted
mean \(m_{D,e}\), the vector
version of V13's response
bound follows by applying its
scalar covariance estimate to
the three components and
summing their squares. By
\eqref{eq:f15v19:score-gradient},
each required exterior scalar
coordinate has response bounded
by
\[
 \left|\partial_{b_k}m_{D,e}\right|
   \le\frac{128\beta D_p^2}{h_p}
                    e^{-\alpha_p(r-92)p}.
\]
The source carrier is within
\(16p\), the exterior begins
outside \((r-60)p\), and the
response formula loses a further
\(16p\). These facts account
for the displayed distance.
There are at most
\(2^{14}r^4p^5\) required
exterior scalar parameters.
At fixed \(N_e\), two good
exteriors are joined by straight
segments in their half-radius
logarithm balls. Each scalar
displacement is at most \(a_p\).
Since \(a_p=h_p/(8D_p)\),
the total oscillation is at
most
\[
 2^{18}\beta D_p r^4p^5
                  e^{-\alpha_p(r-92)p}
                    \le\Omega^*_{p,r,\beta}.
\]
The identity exterior is one
of the permitted comparison
points whenever the retained
inputs are in their chart.

Let \(q_e=\mathbb P(D_e\mid Z_e)\)
and \(A_e=\Gamma_e\cap\{q_e\ge1/2\}\).
The thermal input gives
\(\mathbb P(A_e^c)\le3\epsilon\).
The proof of V13's mean-transfer
theorem works for Hilbert-valued
sources: conditional Jensen,
Cauchy--Schwarz and the triangle
inequality are used with the
vector norm. Since \(|U_e^H|\le B_p^*\),
it gives
\[
 \|m_e-1_{A_e}m_{D,e}\|_2
                  \le5B_p^*\epsilon^{1/4}.
\]
On \(A_e\), the restricted
mean differs from \(P_e\)
by at most \(\Omega^*\).
On \(A_e^c\), the truncated
mean is zero and \(|P_e|\le B_p^*\).
The latter comparison costs
at most \(\Omega^*+\sqrt3 B_p^*\epsilon^{1/2}\).
This proves
\eqref{eq:f15v19:exterior-score}.

Finally project onto the full
output \(C\). The frozen-frame
rotation is \(C\)-measurable,
so the exact score identity
and conditional contraction
give the claimed gradient
estimate. Its inverse rotation
depends only on \(N_e\),
preserving locality and the norm.
\end{proof}

\subsection{Distant mixed blocks of the full coarse Hessian}

\begin{theorem}[Annealed spatial locality of the actual coarse action]
\label{f15v19:coarse-hessian}
If \(d(e,f)\ge4rp\), then
\begin{equation}
 \left(\mathbb E_{\rho_0}
                \|\mathsf H_{ef}(C)\|_{\rm F}^2\right)^{1/2}
           \le2B_p^*E^*_{p,r,\beta}.
 \label{eq:f15v19:coarse-hessian}
\end{equation}
The expectation is under the
full compact coarse marginal,
not its restriction to good
retained configurations.
There is no hypothesis that
the two coarse variables are
independent, or that a conditional
good-event probability has a
uniform positive lower bound.
\end{theorem}

\begin{proof}
Evaluate both edge scores in
the one frame centered at
\(e\), frozen at each \(C\).
The score at \(f\) is measurable
with respect to \(Z_e\): its
detail carrier is within
\(16p\) of \(f\), hence
outside \(J_e\). Its norm
is at most \(B_p^*\) even
when its own local inputs are
not in any small-field chart.

The direct mixed term in
\eqref{eq:f15v19:hessian-identity}
is zero. Conditional expectation
onto \(Z_e\) therefore gives
\[
 -\mathsf H_{ef}
      =\operatorname{Cov}(m_e,U_f^H\mid C).
\]
Write \(R_e=m_e-P_e(N_e)\).
The subtracted vector is
\(C\)-measurable, so it
does not change this covariance.
Using the centered second factor,
\[
 \operatorname{Cov}(m_e,U_f^H\mid C)
   =\mathbb E\!\left[
       R_e\bigl(U_f^H-\mathbb E[U_f^H\mid C]\bigr)^T
                                      \,\middle|\,C\right].
\]
The Frobenius norm of an outer
product is the product of
the Euclidean norms. Thus
\[
 \|\mathsf H_{ef}\|_{\rm F}
       \le2B_p^*\mathbb E[|R_e|\mid C].
\]
Conditional Jensen, integration
under \(\rho_0\), and
Theorem~\ref{f15v19:local-score}
give \eqref{eq:f15v19:coarse-hessian}.
The Frobenius norm is unchanged
by the fixed gauge rotations
on the two edge indices,
so this is also the bound
in the original edge frames.
No frame or cutoff has been
differentiated in this argument.
\end{proof}

This controls a derivative of
the actual retained action.
It is not the unconditional
covariance of retained field
observables. In particular it
does not discard the latter
as a conditional remainder.

\subsection{Rates and the finite-bank summation that is actually justified}

Use the selected full-law
couplings and growing scales,
with only the harmless lower
collar threshold increased:
\begin{equation}
 t=\log\beta,\qquad p\le t,\qquad
 r=\max\{192,\lceil p^{\nu+8}t^6\rceil\}.
 \label{eq:f15v19:sequence}
\end{equation}
Eventually this is the same
collar as in V16--V18.
For large \(t\),
\(r\le2p^{\nu+8}t^6\) and
\(r-96\ge r/2\). Hence
\begin{equation}
 \Omega^*_{p,r,\beta}
    \le 2^{24}C_0 e^t
       t^{5\nu+61}e^{-c_0t^6},
 \qquad c_0=2^{-50}C_0^{-1}.
 \label{eq:f15v19:omega-rate}
\end{equation}
The inherited exceptional-event
rate \eqref{eq:f15v15:epsilon-rate}
is unchanged asymptotically.
Although \(\nu_*\) can be larger
than \(\nu\),
\(B_p^*\le16(2^{28}B_*^{12})e^t t^{\nu_*}\)
still contributes only one
exponential in \(t\) and a
fixed power of \(t\).

\begin{proposition}[Superalgebraic pair bounds and controlled finite banks]
\label{f15v19:rate}
For every fixed \(A>0\),
\begin{equation}
 E^*_{p,r,\beta}=o(\beta^{-A}),\qquad
 2B_p^*E^*_{p,r,\beta}=o(\beta^{-A}).
 \label{eq:f15v19:rate}
\end{equation}
These bounds are uniform in
the edge placements satisfying
the geometric hypotheses.
For any bank \(\mathcal F_e\)
of \(n\) edges at distance
at least \(4rp\) from \(e\),
\begin{align}
 \left(\mathbb E_{\rho_0}
       \sum_{f\in\mathcal F_e}\|\mathsf H_{ef}\|_{\rm F}^2\right)^{1/2}
       &\le 2\sqrt n B_p^* E^*,\notag\\
 \left\|\sum_{f\in\mathcal F_e}\|\mathsf H_{ef}\|_{\rm F}
                       \right\|_{L^2(\rho_0)}
       &\le 2n B_p^* E^*.
 \label{eq:f15v19:finite-bank}
\end{align}
In particular either bank error
is superalgebraically small if
\(n\le\beta^q\) for a fixed
\(q\).
\end{proposition}

\begin{proof}
Equation \eqref{eq:f15v19:omega-rate}
dominates every fixed power of
\(\beta\), including the extra
factor \(B_p^*\).
The exceptional terms have
exponent bounded above by a
negative constant times
\(e^t/t^{6\nu+82}\), apart
from fixed powers of \(t\)
and the displayed score factors.
They are smaller still. This
proves \eqref{eq:f15v19:rate}.
Sum the squares of the
individual block bounds for
the first bank inequality;
use Minkowski for the second.
Finally absorb the fixed power
\(\beta^q\) by choosing a
larger exponent in the first
assertion.
\end{proof}

The same argument does not
sum over all coarse edges
in the selected large torus.
Their number is \(4(M/p)^4\),
while the construction permits
\(\beta\le\log M\); a fixed
power of \(\beta\) cannot
control that volume. The
bound for one pair also has
an exceptional-event remainder
that has not been summed with
a distance-dependent weight.
Neither issue is repaired by
calling the result a local
interaction norm.

\subsection{What is now available, and what remains upstream}

The actual retained action is
smooth, and its score and mixed
Hessian have been expressed
with all conditional normalizers
retained. Its score has an
explicit local chart predictor
under the full output law;
its distant mixed blocks are
small in that law's \(L^2\)
norm. These are spatial
derivative statements, not
only derivatives with respect
to externally inserted scalar
sources.

To obtain a genuine iterative
interaction bound one still
needs a summable spatial
estimate, stable under the
next measured coarse law,
and control of the chart
complement at that scope.
The present averaged bounds
are not uniform in arbitrary
retained configurations or
exterior changes. They also
give no coercive lower bound
on the retained Hessian:
the covariance in
\eqref{eq:f15v19:hessian-identity}
has a minus sign, and retained
gauge and low-frequency
directions have not been
removed by this calculation.
Locality of an action is not
itself a physical mass gap.

The next noncircular target
is thus a distance-summable
version of the measured
retained-action estimate,
with its exceptional term
and low-frequency sector
handled explicitly. Repeating
the scalar thermal budget
or differentiating a small
probability error would not
provide that missing step.

\subsection{Exact diagnostics and predecessor preservation}

The diagnostic script is
\begin{center}\small
\path{scripts/verify_ym_f15_v19_coarse_action_scores.py}.
\end{center}
It independently forms the
six-variable partition jet
of finite conditional
exponential sums, then takes
its formal logarithm. This
checks 432 mixed Hessian
entries, with and without a
direct mixed term, including
the covariance sign and the
normalizer contribution.
There are six vector-valued
chart/complement transfers,
24 conditional cross-block
bounds, 24 fixed-point
orthogonal-frame checks,
and 54 support, collar and
polynomial-prefactor checks.
All use exact fractions.
These tests do not enclose
the primitive derivative
suprema, simulate the Wilson
law, or prove continuum
reconstruction. The script
has 9,811 bytes and SHA--256
\begin{center}\scriptsize\ttfamily
69E5B7D8FA150D939FA16F9FE5D2C56829BBDB81A75F2BD76DA6E8D2EDCF6779.
\end{center}

The predecessor V18 has
409,969 bytes and SHA--256
\begin{center}\scriptsize\ttfamily
1DEC066353BFD0412743441A2EB578AF9D8B1EFB29C4E02BE881C3C1469B5AA6.
\end{center}
Reversing this terminal
insertion and the version
title change recovers it
exactly. V19 replaces V18
as the sole active main
note. All 14 protected
sources and 61 earlier
protected certificates
remain unchanged.

\begin{firewall}
V19 obtains full-law averaged
spatial derivative estimates
for the actual coarse action.
It does not prove an intensive
all-volume interaction norm,
unconditional physical
clustering, a controlled
interacting continuum limit,
or the Yang--Mills mass gap.
\end{firewall}

\section{V20: full-law bad-patch clusters and the gauge-averaging audit}
\label{f15v20:section}

V19 gave a full-law averaged bound on each distant block of the
actual retained-action Hessian. Summing its exceptional error over
the whole lattice still incurs a volume factor. This section goes
upstream to the spatial organization of exceptional patches. It
uses the arbitrary-support cell-source theorem already recorded
in \eqref{f15v1:source-input}; it does not reprove that theorem or
replace the Wilson law by a Gaussian law. The new application is
a joint bad-patch bound, and consequently a connected-cluster tail,
under the original full compact selected Wilson law. Constants
are independent of total volume, at each of the stated scales.

A preliminary audit rules out a misleading shortcut: cancellation
of distant, untransported colour components after gauge averaging
is not locality in a norm. The direction below therefore concerns
gauge-invariant exceptional events and their actual probabilities.
The missing passage from this random geometry to a summable
retained-action interaction remains explicit.

\subsection{The scheduled V20 companion checkpoint}

\begin{record}[Current sources, scope and non-imports]
\label{f15v20:checkpoint}
The following current files were revisited. Their hashes agree
with the preceding protected-source manifest; no newer versions
of these files were found in the inspected notes/downloads locations.
The last three are the most recent attachments, replacing the
earlier V16/V21/V17 inputs for this scope review. Hash prefixes
below identify the exact reviewed files, not independent proof
certificates.
\begin{center}\small
\begin{tabular}{@{}p{0.33\textwidth}p{0.17\textwidth}p{0.43\textwidth}@{}}
\toprule
Programme/version & SHA--256 prefix & Relevant scope limitation\\
\midrule
SM continuum V72 & \texttt{F44F9745D433} &
{\raggedright Fixed mode bank and bounded-time empirical covariance
transport; not a growing ultraviolet cutoff theorem.\par}\\[1mm]
NS stability V26 & \texttt{82C887F8C2C6} &
{\raggedright Rank-one Oseen derivative estimates; not a full
regularity or YM spatial estimate.\par}\\[1mm]
YM shell mixing V16 & \texttt{BF138E63BA82} &
{\raggedright Measured one-loop flat-holonomy coefficient and summable
coefficient limit; not a joint nonperturbative remainder or
physical mass.\par}\\[1mm]
Parallel YM V5 & \texttt{306F1BB84CFA} &
{\raggedright Global activity/charge entropy obstruction points toward
local or density estimates; no action interaction norm.\par}\\[1mm]
Einstein bridge V24 & \texttt{A9B80D1A2689} &
{\raggedright Lower-incidence first-jet tangency repair; no integrated
full root or uniform-time construction.\par}\\[1mm]
Native stationarity V34 & \texttt{08BC4B7DE876} &
{\raggedright Necessary escape-cone quadratic fibre derivative remains
to be computed; no corresponding Wilson spatial inequality.\par}\\[1mm]
SM source selection V31 & \texttt{34312755DFB6} &
{\raggedright Native inserted-word Gram compilation distinguished
from an external operation; no historical occurrence or
YM mixing proof.\par}\\
\bottomrule
\end{tabular}
\end{center}
The source-selection, bridge and stationarity refinements do not
supply the missing full-law YM comparison. The NS note's reference
to a future V27 certificate is not imported or treated as verified.
The unchanged supplied suggestions remain optional source material.
No theorem from these companions is added as a new hypothesis here.
Their useful directional lesson is to retain the actual measured
law and use local exceptional geometry, not global absence of all
defects. The next scheduled companion checkpoint is V25.
\end{record}

\subsection{Why the untransported averaged Hessian can vanish}

Use V19's coarse Haar density \(R_j(C)\), action
\(\mathcal A_j=-\log R_j\), score vectors \(m_e\), and
mixed different-edge Hessian blocks \(H_{ef}\). Denote its
normalized marginal by \(\rho_0\). All statements here are at a
finite compact regulator. These smooth functions are integrable.

\begin{proposition}[Gauge cancellation is not a norm estimate]
\label{f15v20:gauge-cancellation}
If the positive coarse edges \(e,f\) have different starting
vertices, then
\begin{equation}
 \mathbb E_{\rho_0}m_e=0,\qquad
 \mathbb E_{\rho_0}(m_em_f^{\mathsf T})=0,\qquad
 \mathbb E_{\rho_0}H_{ef}=0.
 \label{eq:f15v20:gauge-zero}
\end{equation}
When their starting vertices agree and \(e\ne f\), each of the
last two averaged matrices is a scalar multiple of \(I_3\).
None of these statements bounds
\(\mathbb E\|H_{ef}\|_{\mathrm F}\) or its square moment.
\end{proposition}

\begin{proof}
Gauge equivariance of the blocking and invariance of the fine
Wilson law imply invariance of \(\rho_0\). Equivalently, use the
fixed-gauge change of variables in Section~\ref{f15v19:frames}:
the detail Haar measure is preserved, so \(R_j\) is invariant.
For a fixed vertex gauge transformation, the left-insertion
derivatives defining the score transform by the adjoint rotation
at the starting vertex. For \(e\ne f\) the Hessian transforms
by the two corresponding rotations on its two indices. The gauge
transformation is constant during differentiation.

If the starting vertices are \(x\ne y\), average the transformation
at \(x\) alone over normalized \(\SU(2)\) Haar measure. Its adjoint
average is zero: an invariant vector under all three-dimensional
rotations is zero. The second score index and the second Hessian
index do not rotate. This remains true if \(x\) is the ending
vertex of \(f\), since its left insertion is based at the other
endpoint. This proves all the zero identities. With a common
starting vertex, the averages commute with every adjoint rotation.
Rotations by \(\pi\) around the coordinate axes kill off-diagonal
entries; rotations interchanging axes make the three diagonal
entries equal. This gives the scalar assertion.

For the last point, take independent Haar rotations \(Q_1,Q_2\)
and any fixed nonzero real matrix \(A\). The random matrix
\(Q_1 A Q_2^{\mathsf T}\) has zero mean and constant positive
Frobenius norm \(\|A\|_{\mathrm F}\). This is only an algebraic
counterexample to the proposed inference, not a model for the
Wilson Hessian.
\end{proof}

This finite-dimensional calculation is an instance of
\href{https://journals.aps.org/prd/abstract/10.1103/PhysRevD.12.3978}
{Elitzur's local-gauge symmetry observation}; its proof above is
self-contained. In particular, a bound on the spatial row of
\(\mathbb E H_{ef}\) would be vacuous away from the common
starting vertex. Parallel transport to a common colour frame
removes this trivial cancellation, but introduces coefficients
depending on the retained field. A Bessel inequality for fixed
coefficients does not control those random coefficients or their
derivatives. V19's norm estimates are not invalidated; they must
not be replaced by \eqref{eq:f15v20:gauge-zero}.

\subsection{One bad rooted patch has a local high-energy witness}

From now on \(\beta=\beta_M^\diamond\) and
\(\mu=\mu^{\mathrm W}_{M,\beta}\) are exactly the selected coupling
and law of \eqref{f15v1:source-input}. Write
\(\overline T=\mathbb E_\mu T_x\) and
\(J(s)=-\log(1-s)-s\). Work far enough along the selection that
\begin{equation}
 R_M\ge\tfrac12,\qquad \overline T\le5,\qquad k_M\le5.
 \label{eq:f15v20:selected-input}
\end{equation}
These conditions follow from the recorded limits, not from a
new assumption of independence.

Let \(\mathcal K\) be an embedded fine box of even side length
\(L\), including its boundary vertices. Use its ordered-coordinate
root tree from the lower corner. For \(0<\delta\le\pi\), put
\[
 B_{\mathcal K}=\{R_{\mathcal K}>\delta\},\qquad
 u=\frac{\beta\delta^2}{18\pi^2L^2},\qquad w=(L+1)^4.
\]
The event is gauge invariant: the rooted links transform by
conjugation at the common root. Let \(\mathcal W_{\mathcal K}\)
be the set of base vertices of positive plaquettes contained in
the box; in particular \(|\mathcal W_{\mathcal K}|\le w\).

\begin{proposition}[Distinct-cell witnesses give a joint estimate]
\label{f15v20:witness}
One has the deterministic inclusion
\begin{equation}
 B_{\mathcal K}\ \subseteq
 \bigcup_{x\in\mathcal W_{\mathcal K}}\{T_x>u\}.
 \label{eq:f15v20:witness}
\end{equation}
For any finite collection of such boxes with pairwise disjoint
witness sets, common \(L,\delta\), and cardinality \(n\),
\begin{equation}
 \mu\Bigl(\bigcap_{i=1}^n B_{\mathcal K_i}\Bigr)
 \le q_0^n,\qquad
 q_0=\min\{1,\ 32w e^{-u/2}\}.
 \label{eq:f15v20:disjoint}
\end{equation}
Neither the fine link sets supporting \(T_x\) nor their random
values are required to be independent.
\end{proposition}

\begin{proof}
The rooted-link filling used in V8 and V13 consists of at most
\(3L\) plaquettes, all in the box. The bi-invariant geodesic
triangle inequality bounds the link angle by the sum of their
angles, including multiplicities. Thus a rooted link of angle
greater than \(\delta\) has a filling plaquette with angle
\(\vartheta>\delta/(3L)\). For \(0\le\vartheta\le\pi\),
\[
 1-\cos\vartheta=2\sin^2(\vartheta/2)
       \ge\frac{2\vartheta^2}{\pi^2}.
\]
Replacing a negatively oriented face by its positive orientation
does not change its angle. If its base is \(x\), the term with
\(\epsilon=0\) in \eqref{f15v1:cell} gives
\[
 T_x\ge\tfrac14 X_{(x;\mu,\nu)}
       >\frac{\beta\delta^2}{18\pi^2L^2}=u.
\]
This proves the inclusion and the witness count.

For any distinct cells \(x_1,\ldots,x_n\), apply
\eqref{f15v1:source-input} with \(b_{x_i}=1/2\) and zero elsewhere.
There is no restriction on \(n\) beyond distinctness in the finite
lattice. Exponential Markov inequality gives
\begin{align}
 \mu\{T_{x_i}>u\text{ for every }i\}
 &\le\exp\{n[-u/2+\overline T/2+k_MJ(1/2)]\}\\
 &\le\exp\{n[-u/2+5\log2]\}.
 \label{eq:f15v20:source-tail}
\end{align}
Here \(J(1/2)=\log2-1/2>0\), so both upper bounds in
\eqref{eq:f15v20:selected-input} have the displayed sign.
Witness sets from different boxes are disjoint, hence every
choice of one witness per box gives distinct source indices.
There are at most \(w^n\) choices. The union bound proves
\((32w e^{-u/2})^n\); the trivial probability bound by one
supplies the cap in \eqref{eq:f15v20:disjoint}.
\end{proof}

This is where the arbitrary-support source theorem is used.
Single-plaquette exponential moments alone, followed by a
one-patch union bound, do not give \eqref{eq:f15v20:disjoint}.
The inherited source theorem, including its selected-coupling
restriction, remains an explicit dependency of the result.

\subsection{An overlapping patch cover with a fixed number of colours}

Take \(p=2^j\), \(r\ge192\) an integer, and
\begin{equation}
 L=2p(r+5),\qquad
 \ell=2^{\lceil\log_2 L\rceil-2},\qquad
 N=M/\ell,\qquad K=8^4=4096.
 \label{eq:f15v20:mesh}
\end{equation}
Assume the dyadic fine torus satisfies \(M\ge4L\). Then
\(L/4\le\ell<L/2\), \(p\mid\ell\), \(\ell\mid M\), and
\(N\) is a power of two at least eight. Index the patch centres
by \(z\in\mathbb T_N^4=(\mathbb Z/N\mathbb Z)^4\):
\(\mathcal K_z\) is the box of side \(L\) centred at \(\ell z\).
Its coordinate lift is unambiguous on this nonwinding scale.
Give \(z\) the colour \((z_1,\ldots,z_4)\bmod8\).

Distinct centres of one colour differ in some periodic coordinate
by at least \(8\ell\ge2L\). Their closed boxes, and hence their
witness sets, are disjoint. This includes the periodic seam,
since eight divides \(N\). There are at most \(K\) colours.
The patches do cover all fine anchors: a nearest centre is
within coordinate distance \(\ell/2<L/4\). In fact
\(L/4=p(r+5)/2\le p(r-60)\) for \(r\ge192\), so the inner
anchor boxes used in V19 also cover. This geometric observation
does not by itself decompose the Hessian into patch contributions.

Write \(B_z=B_{\mathcal K_z}\), and let
\begin{equation}
 q=q_0^{1/K}
  =\min\{1,\ [32(L+1)^4e^{-u/2}]^{1/4096}\}.
 \label{eq:f15v20:q}
\end{equation}

\begin{theorem}[Joint bad patches in the full compact law]
\label{f15v20:joint}
For every set \(A\subseteq\mathbb T_N^4\),
\begin{equation}
       \mu\{B_z\text{ for all }z\in A\}\le q^{|A|}.
 \label{eq:f15v20:joint}
\end{equation}
This holds for sets as large as the entire mesoscopic torus.
The same inequality holds if each \(B_z\) is replaced by any
specified measurable event contained in it.
\end{theorem}

\begin{proof}
For nonempty \(A\), choose a colour containing at least
\(\lceil |A|/K\rceil\) of its indices. The event on the left
implies that all these same-colour patches are bad. Apply
Proposition~\ref{f15v20:witness} and use \(0<q_0\le1\):
\[
 \mu\{B_z\ (z\in A)\}
  \le q_0^{\lceil |A|/K\rceil}
  \le q_0^{|A|/K}=q^{|A|}.
\]
The empty set has probability one. Event inclusion proves the
last assertion; it does not change the underlying law.
\end{proof}

In particular, take the actual canonical charts and pinned
near-output frames of V13 separately at each centre. Its entry
theorem, with \(\delta=a_p/(64rp^4)\), proves
\(\Gamma_z^c\cup D_z^c\subseteq B_z\). Thus their actual chart
failure indicators satisfy \eqref{eq:f15v20:joint}. The frames
may differ between patches; this use of event inclusion requires
neither compatible global gauge fixing nor differentiation of
those frames. It is a statement about the full Wilson pushforward,
not about a product of separately normalized chart laws.

\subsection{Intensive count and connected-cluster estimates}

Connect nearest neighbours of \(\mathbb T_N^4\); its degree is
eight. Let \(\mathcal C_z\) be the connected component of bad
vertices containing \(z\), with \(\mathcal C_z=\varnothing\)
when \(z\) is good. This graph refers to mesoscopic patches,
not directly to physical space-time units.

\begin{theorem}[Volume-independent exceptional geometry]
\label{f15v20:clusters}
For every finite \(A\subseteq\mathbb T_N^4\), \(t\ge0\), and
integer \(n\ge1\),
\begin{align}
 \mathbb E_\mu e^{t\sum_{z\in A}1_{B_z}}
    &\le[1+q(e^t-1)]^{|A|},
                                  \label{eq:f15v20:count}\\
 \mu\{|\mathcal C_z|\ge n\}
    &\le q(64q)^{n-1}.
                                  \label{eq:f15v20:cluster}
\end{align}
If \(64q<1\), then
\begin{equation}
 \mathbb E_\mu|\mathcal C_z|\le\frac{q}{1-64q},\qquad
 \mathbb E_\mu(e^{h|\mathcal C_z|}-1)
   \le\frac{q(e^h-1)}{1-64qe^h}quad(64qe^h<1, h\ge0).
 \label{eq:f15v20:cluster-moment}
\end{equation}
Also, the probability of a bad path from \(z\) reaching graph
distance at least \(d\) is at most \(q(64q)^d\).
\end{theorem}

\begin{proof}
Expand
\(\prod_{z\in A}[1+(e^t-1)1_{B_z}]\). Its coefficients are
nonnegative. Apply \eqref{eq:f15v20:joint} to each subset and
sum the binomial expansion, proving \eqref{eq:f15v20:count}.

Here we reuse, rather than claim anew, the elementary connected-set
counting argument of f14 V56/V67. A connected set of \(n\)
vertices containing a specified root has a deterministic
depth-first spanning-tree traversal of length \(2n-2\), after
fixing a neighbour ordering. Its visited set determines the
original set. Since there are at most \(8^{2n-2}\) such walks,
there are at most that many connected sets. If the root cluster
has at least \(n\) vertices, it contains a connected subset of
exactly \(n\) containing the root, by growing a tree and stopping
at \(n\). Each such subset is wholly bad with probability at
most \(q^n\). A union bound gives \eqref{eq:f15v20:cluster}.
The argument is valid on the finite torus, including wrapping
sets; for \(n>N^4\) the event is empty.

Sum the tail for the first moment. For any nonnegative integer
random variable \(X\),
\[
 \mathbb E(e^{hX}-1)
  =(e^h-1)\sum_{n\ge1}e^{h(n-1)}\mathbb P(X\ge n).
\]
This identity and the geometric series give
\eqref{eq:f15v20:cluster-moment}. A path reaching distance
\(d\) contains at least \(d+1\) distinct vertices in the root
cluster, proving the last assertion.
\end{proof}

The earlier f14 V56 result concerned Gaussian scale fields;
f14 V67 concerned a restricted conditional law, with a
mean-exceptional incoming set still to control; f14 V74 concerned
an added-noise channel. The new point here is the input
\eqref{eq:f15v20:joint} under the original full compact Wilson
law. Neither the combinatorial walk count nor the abstract
binomial expansion is being counted as a new result.

\subsection{A conditional cluster tail for typical retained data}

There is a useful conditional conclusion, but its order of
quantifiers is important. It concerns a specified root at a time.

\begin{proposition}[Simultaneous posterior tails at a fixed root]
\label{f15v20:posterior}
Suppose \(64q<1\), fix \(z\), and let \(\mathcal Y\) be any
sub-sigma-algebra of the full finite-regulator probability space.
There is a \(\mathcal Y\)-measurable event \(G_{z,\mathcal Y}\)
such that
\begin{equation}
 \begin{gathered}
 \mu(G_{z,\mathcal Y}^c)
       \le\frac{\sqrt q}{1-8\sqrt q},\\
 \mu\{|\mathcal C_z|\ge n\mid\mathcal Y\}
       \le\sqrt q(8\sqrt q)^{n-1}\quad(n\ge1)
       \quad\hbox{on }G_{z,\mathcal Y}.
 \end{gathered}
 \label{eq:f15v20:posterior}
\end{equation}
All the conditional inequalities hold simultaneously, up to a
single null set. In particular one may take \(\mathcal Y=\sigma(C)\)
for the entire retained field, without a conditional independence
assumption.
\end{proposition}

\begin{proof}
Put \(F_n=\mathbb E[1_{\{|\mathcal C_z|\ge n\}}\mid\mathcal Y]\).
The tower property and \eqref{eq:f15v20:cluster} give
\(\mathbb EF_n\le q(64q)^{n-1}\). For
\(b_n=\sqrt q(8\sqrt q)^{n-1}\), this upper bound is \(b_n^2\).
Markov inequality yields \(\mu\{F_n>b_n\}\le b_n\).
Define \(G_{z,\mathcal Y}\) by all inequalities \(F_n\le b_n\)
for \(1\le n\le N^4\). All later tails vanish. The union bound
is at most \(\sum_{n\ge1}b_n=\sqrt q/(1-8\sqrt q)\).
There are only finitely many nontrivial conditional expectations,
so a common null set suffices.
\end{proof}

For a specified V19 centre one can instead take
\(\mathcal Y=\sigma(Z_z)\), its retained field plus framed exterior
details. Since \(D_z^c\subseteq B_z\), on this good-data event
\[
          \mu(D_z\mid Z_z)\ge1-\sqrt q.
\]
Moreover \(\Gamma_z^c\) is \(Z_z\)-measurable and is contained
in \(B_z\). If it occurred, the conditional probability of
\(B_z\) would be one, contradicting the \(n=1\) bound in
\eqref{eq:f15v20:posterior}. Thus \(\Gamma_z\) holds there.
This supplies a rootwise high-probability conditional chart
likelihood, not a uniform-exterior estimate. We have not proved
a joint bound for the events \(G_{z,\mathcal Y_z}^c\) associated
with different sigma-algebras; intersecting them over all roots
by a union bound reintroduces a volume factor.

\subsection{The joint parameter is small on the existing growing scale}

Keep exactly V13's constants \(C_0,\nu,a_p\), in particular
\[
 a_p=2^{-43}C_0^{-1}p^{-(\nu+8)},\qquad
 \delta=\frac{a_p}{64rp^4}.
\]
For the selected growing scale, write
\begin{equation}
 t=\log\beta,\qquad 1\le p\le t,\qquad
 r=\max\{192,\lceil p^{\nu+8}t^6\rceil\}.
 \label{eq:f15v20:scale}
\end{equation}
The old lower condition \(\beta\ge\beta_p\) remains available
where V13's convexity is used; it is not needed for the witness
and source argument alone.

\begin{proposition}[Superalgebraically small full-law defect parameter]
\label{f15v20:rate}
For all sufficiently large selected regulators, the above cover
exists and its \(q\) satisfies
\begin{equation}
 \log q\le\frac1{4096}
 \left[\log(320000)+(4\nu+60)\log t
       -\frac{c_*e^t}{t^{6\nu+82}}\right],\qquad
 c_*=\frac{2^{-108}}{9\pi^2C_0^2}.
 \label{eq:f15v20:rate}
\end{equation}
Consequently \(\beta^Aq\to0\) for every fixed \(A>0\),
\(64q<1\) eventually, and the exceptional probabilities in
\eqref{eq:f15v20:posterior} also vanish faster than every fixed
negative power of \(\beta\).
\end{proposition}

\begin{proof}
Eventually \(p^{\nu+8}t^6\ge192\), so
\(r\le2p^{\nu+8}t^6\). Since \(L\le4rp\),
\[
 \frac u2=\frac{\beta\delta^2}{36\pi^2L^2}
  \ge\frac{\beta a_p^2}{2359296\pi^2r^4p^{10}}
  \ge\frac{2^{-108}e^t}
             {9\pi^2C_0^2p^{6\nu+58}t^{24}}
  \ge\frac{c_*e^t}{t^{6\nu+82}}.
\]
Also \(L+1\le5rp\), whence
\[
 32(L+1)^4\le320000\,t^{4\nu+60}.
\]
Insert these inequalities into \eqref{eq:f15v20:q}. The
uncapped logarithm tends to minus infinity; in fact the same
upper bound on \(\log q\) is valid when the cap is active,
because then the uncapped logarithm is nonnegative. This gives
\eqref{eq:f15v20:rate}. The negative exponential-in-\(t\) term
dominates \(At\) and all logarithmic terms for every fixed \(A\).
The conditional exceptional bound is asymptotic to \(\sqrt q\).

Finally \(L\) is bounded by a fixed power of \(t\), while
\(\beta\le\log M\) and \(t\to\infty\). Therefore \(M\ge4L\)
eventually. The dyadic definitions then give all divisibility
conditions used in \eqref{eq:f15v20:mesh}.
\end{proof}

\subsection{What this removes, and the remaining analytic interface}

The new bounds control the number of bad patches and the size
of a bad component at a specified root without multiplying by
the total number of patches. They permit bad patches elsewhere
on an arbitrarily large finite regulator. They do not assert
that the probability of \emph{some} bad patch tends to zero
uniformly in volume. Nor does the geometric decay in patch
distance identify a physical correlation length or a mass.

Even the joint inequality \eqref{eq:f15v20:joint} is not
stochastic domination by independent Bernoulli variables.
For example, choose exactly one of two sites to be bad, with
probability \(1/2\) for each choice. Every all-bad subset has
probability at most \((1/2)^{|A|}\), yet the increasing event
``at least one is bad'' has probability one, versus \(3/4\)
under independent Bernoulli \(1/2\) variables. The probability
that the first site is bad conditional on the second being
good is one. Thus uniform conditional probabilities do not
follow from the joint unconditional estimate either.

Proposition~\ref{f15v20:posterior} is a separately proved,
rootwise typical-data statement that respects this obstruction.
The constants deteriorate, its exceptional event depends on
the conditioning sigma-algebra, and it does not give arbitrary
conditional source estimates. In particular, different
conditional confidence failures in V19 cannot simply be
declared a field satisfying \eqref{eq:f15v20:joint}.

There is also a quantitative warning against a crude next step.
On \eqref{eq:f15v20:scale}, for every fixed \(c>0\),
\begin{equation}
 \frac{\log(e^{c\beta L^4}q)}{\beta L^4}
 =c+\frac{\log[32(L+1)^4]}{4096\beta L^4}
       -\frac{\delta^2}{36\pi^2\,4096\,L^6}
 \longrightarrow c.
 \label{eq:f15v20:oscillation-warning}
\end{equation}
Indeed the cap in \eqref{eq:f15v20:q} is eventually inactive,
\(L\to\infty\), \(\beta\to\infty\), and \(\delta\le\pi\).
Thus a comparison charging \(e^{c\beta L^4}\) for each bad
patch cannot be paid by multiplying the present joint bound
by that cost. This is a failure of that particular upper-bound
budget, not a lower bound on actual bad probabilities or a
proof that a sharper comparison is impossible. A whole-box
action-oscillation estimate alone would not close the next step.

The next analytic target is an actual measured-fibre resolvent
or conditional covariance estimate organized by these bad
components: propagation in the proved convex good regions,
an explicit cost for crossing a connected bad component, and
a proof that the joint full-law tails pay that cost. Such a
bound must retain conditional normalizers and avoid
differentiating only a small-probability error. No bound on the
cost of crossing an arbitrarily ill-conditioned bad region has
been supplied here. This is the substantive missing bridge to
an intensive, distance-summable version of V19. Even after it,
retained low modes, the interacting continuum limit and the
physical Hamiltonian gap remain separate obligations.

\subsection{Exact diagnostics and predecessor preservation}

The finite diagnostic is
\begin{center}\small
\path{scripts/verify_ym_f15_v20_full_law_bad_patches.py}.
\end{center}
It checks 576 two-frame rotation cancellations with unchanged
positive matrix norms, 16 periodic patch/witness cases, 24
scale-constant cases, and 18 source-exponent endpoint checks
using rational upper and lower bounds for \(\log2\). It checks
122 distinct connected-set traversal encodings, 40 finite
cluster tails, 18 simultaneous conditional-tail examples and
1,072 dependent-indicator/Laplace comparisons. Computations
use only integers and fractions. They check finite algebra,
geometry and the stated distinctions; they do not certify the
inherited Wilson source theorem or simulate the Wilson law.
The script has 7,862 bytes and SHA--256
\begin{center}\scriptsize\ttfamily
84879749EC8A514C4BD1C5C77D5F64A363F53CC7C3EA2BD3727816A3CCF45C3B.
\end{center}

The predecessor V19 has 431,028 bytes and SHA--256
\begin{center}\scriptsize\ttfamily
CEC946F717A06033E5F701062EB3C5E85308A1C2FD915D8DBFA3236EB6DB02AE.
\end{center}
Removing this terminal section and reversing the version-title
change recovers it byte for byte. V20 replaces V19 as the sole
active main note; all 14 protected sources and 62 earlier
protected certificates are preserved unchanged.

\begin{firewall}
V20 proves joint and connected-cluster bad-patch estimates for
the full compact selected Wilson law, and typical conditional
tails at a fixed root, using the inherited cell-source input.
It does not prove a bad-region covariance crossing estimate,
a volume-uniform retained interaction norm, unconditional
physical clustering, an interacting continuum construction,
or the Yang--Mills mass gap. The gauge-averaged zero identities
are explicitly not offered as decay estimates in a norm.
\end{firewall}

\section{V21: matched conditional confidence fields and sparse score errors}
\label{f15v21:section}

V20 obtained joint full-law tails for actual chart failures,
but its rootwise posterior argument did not control the joint
failure of conditional chart likelihoods at different roots.
That distinction matters in V19: its conditional mean is taken
given a different exterior for each interior. This section
closes that specific probabilistic interface by using an
additional, already established property of the actual model:
the measured canonical action has finite range.

Sufficiently separated interiors become exactly independent
after fixing their \emph{common exterior}, including the whole
retained field. Their individually conditioned failure
probabilities can then be identified with probabilities for
this common disintegration. Every conditional normalizer is
retained. Combining this identity with V20 gives joint tails
for the conditional-confidence failures themselves. We also
extend the colouring from the mesoscopic cover to every coarse
anchor, and convert these tails into a sparse, pointwise
decomposition of the actual conditional score error.

The new result does not establish propagation across bad
regions. It supplies the missing spatial probability bound for
the very conditional failures that such a propagation argument
would encounter. No conditional independence is inferred merely
from small bad-event probabilities. V20's gauge-averaging and
whole-box-oscillation warnings remain in force. V20's proofs
are preserved without alteration. Its scheduled companion
review remains current; the next checkpoint is V25.

\subsection{The common-exterior identity, with all normalizers}

Use exactly the full compact density
\eqref{eq:f15v12:measured-law}, after integrating its independent
rooted gauge Haar factor. Let \(\mathcal I\) be its finite set
of canonical detail indices. No original physical pin is being
substituted for these indices. Consider disjoint sets
\(J_1,\ldots,J_n\subset\mathcal I\), and put
\[
 O=\mathcal I\setminus\bigcup_iJ_i,\qquad
 \mathcal Y=\sigma(C,v_O),\qquad
 \mathcal Z_i=\sigma(C,v_{\mathcal I\setminus J_i}).
\]
Suppose no measured interaction term touches two different
\(J_i\)'s. Let \(E_i\) be an event measurable with respect to
\(\sigma(C,v_O,v_{J_i})\); in particular it must not depend
on another selected interior.

\begin{proposition}[Matched conditional products]
\label{f15v21:matched-products}
Define \(\pi_i=\mu(E_i\mid\mathcal Z_i)\). There are versions
of these conditional probabilities that are
\(\mathcal Y\)-measurable and satisfy
\begin{equation}
 \pi_i=\mu(E_i\mid\mathcal Y),\qquad
 \mathbb E_\mu\prod_{i\in A}\pi_i
      =\mu\Bigl(\bigcap_{i\in A}E_i\Bigr)
       \quad\bigl(A\subseteq\{1,\ldots,n\}\bigr).
 \label{eq:f15v21:matched-products}
\end{equation}
Consequently, for thresholds \(0<\eta_i<1\),
\begin{equation}
 \mu\{\pi_i>\eta_i\text{ for every }i\in A\}
 \le\left(\prod_{i\in A}\eta_i^{-1}\right)
                 \mu\Bigl(\bigcap_{i\in A}E_i\Bigr).
 \label{eq:f15v21:threshold}
\end{equation}
These conclusions do not require the retained field or the
common exterior to have a product law.
\end{proposition}

\begin{proof}
Assign every measured term touching \(J_i\) to \(W_i\).
The no-common-term assumption makes this assignment unique;
put all other terms in \(W_0\). Then
\[
 \mathcal W_j(C,v)
       =W_0(C,v_O)+\sum_iW_i(C,v_O,v_{J_i}).
\]
At a fixed \(y=(C,v_O)\) define the positive finite integrals
\[
 Z_i(y)=\int e^{-W_i(y,u_i)}\,dH_{J_i}(u_i),\qquad
 \kappa_i^y(du_i)=Z_i(y)^{-1}e^{-W_i(y,u_i)}dH_{J_i}(u_i).
\]
Compactness and the positive smooth measured density ensure
that the denominators are neither zero nor infinite. Fubini's
theorem gives the conditional normalizer
\(e^{-W_0(y)}\prod_iZ_i(y)\); hence the joint conditional
interior law is exactly \(\bigotimes_i\kappa_i^y\).
None of the \(Z_i\)'s has been discarded.

Conditioning further on the values of \(v_{J_k}\), \(k\ne i\),
does not change \(\kappa_i^y\). The stipulated support of
\(E_i\) means that its conditional probability is
\(\int1_{E_i}(y,u_i)\kappa_i^y(du_i)\) under either
conditioning. This proves the first identity. Conditional
independence given \(\mathcal Y\) proves
\(\mu(\cap_{i\in A}E_i\mid\mathcal Y)=\prod_{i\in A}\pi_i\).
Take expectation to obtain the second identity. Finally
\(1_{\{\pi_i>\eta_i\ \forall i\in A\}}
\le\prod_{i\in A}(\pi_i/\eta_i)\), and integrate.
All families are finite at a regulator, so the identities
can be fixed on one common full-measure set.
\end{proof}

For background on the distinction between a local update
kernel and a contraction estimate for a system of such kernels,
see \href{https://arxiv.org/abs/1308.4117}
{Rebeschini--van Handel, \emph{Comparison Theorems for Gibbs
Measures}}, Definition 2.2 and Theorem 2.4. No comparison or
uniqueness theorem from that paper is invoked here: the
factorization above is proved directly, and no contraction
matrix has yet been bounded for the full compact YM law.

\subsection{Actual chart failures obey the required support conditions}

Retain V20's \(p,r,L,\ell,\delta,q_0\), including
\[
 \begin{gathered}
 r\ge192,\quad L=2p(r+5),\quad M\ge4L,\quad
 \ell=2^{\lceil\log_2L\rceil-2},\\
 \delta=\frac{a_p}{64rp^4},\quad
 q_0=\min\{1,32(L+1)^4e^{-u/2}\},\quad
 u=\frac{\beta\delta^2}{18\pi^2L^2}.
 \end{gathered}
\]
All probabilities remain under the original selected full
Wilson law with the source hypotheses
\eqref{eq:f15v20:selected-input}. We use either spacing
\(s=\ell\), for the V20 cover, or \(s=p\), for all coarse
anchors. The index torus is \(\mathbb T_{M/s}^4\), and the
centre of index \(z\) is the fine vertex \(sz\). Translate
V12's canonical interior to this centre:
\[
 J_z=\{i:\operatorname{anchor}(v_i)-sz
                  \in[-(r-60)p,(r-60)p]^4\},\qquad O_z=J_z^c.
\]
All distances and boxes have the nonwinding periodic convention
of V20. Choose V13's coarse frame \(H_{N_z}(C)\). Define
\(D_z\) as its interior full-chart event and \(\Gamma_z\) as
the required retained/exterior half-chart event. Include in
\(\Gamma_z\) the central score-carrier requirements for each
of the four positive coarse edges starting at \(sz\). This
is a harmless intersection of four of V19's conditions; its
entry still follows from the same rooted-patch event. Put
\begin{equation}
 E_z=\Gamma_z^c\cup D_z^c,\qquad
 \mathcal Z_z=\sigma(C,v_{O_z}),\qquad
 \pi_z=\mu(E_z\mid\mathcal Z_z).
 \label{eq:f15v21:failure}
\end{equation}
The sigma-algebra is unchanged if \(v_{O_z}\) is replaced by
\(v_{O_z}^{H_{N_z}(C)}\): at fixed \(C\) this is an invertible
coordinatewise translation/conjugation, not a mixing of indices.

\begin{proposition}[Separation and normalization in the actual coordinates]
\label{f15v21:support}
If a family of centres has pairwise periodic coordinate
separation at least \(8\ell\) in some coordinate for every
pair, then its interiors and events \(E_z\) satisfy all
hypotheses of Proposition~\ref{f15v21:matched-products}.
For such a family \(A\),
\begin{equation}
             \mathbb E_\mu\prod_{z\in A}\pi_z
                       =\mu\Bigl(\bigcap_{z\in A}E_z\Bigr)
                       \le q_0^{|A|}.
 \label{eq:f15v21:separated-moment}
\end{equation}
The conditional kernel defining each \(\pi_z\) is precisely
the full compact interior kernel of V12, not a chart-restricted
or independently chosen exterior law.
\end{proposition}

\begin{proof}
Every term of \(\mathcal W_j\), including each pivot Jacobian
term, has anchor diameter at most \(16p\), by
Proposition~\ref{f15v12:interior}. Since \(8\ell\ge2L\),
the separation between two selected interiors is at least
\[
 8\ell-2(r-60)p\ge(2r+140)p>16p.
\]
Thus no measured term touches both. The parameters of the
interior potential have anchors within radius \((r-44)p\)
of the centre. The added score carriers are central, of radius
\(16p\), strictly inside this bound. Consequently the detail
carrier of \(\Gamma_z\) lies outside \(J_z\), but inside
that radius. It cannot meet another selected interior, since
\[
 8\ell-(r-44)p-(r-60)p\ge(2r+124)p>0.
\]
The frame depends only on \(C\), which is fixed in the common
exterior. Its coordinatewise action cannot enlarge these detail
carriers. The event \(D_z\) involves \(v_{J_z}\) and this
fixed frame only. Hence \(E_z\) is measurable with respect to
the common exterior and its own interior, as required.

V13's deterministic entry theorem gives
\(E_z\subseteq B_z=\{R_{\mathcal K_z}>\delta\}\), including
all four central score conditions. The selected fine boxes
have disjoint witness sets, by their separation. The equality
in \eqref{eq:f15v21:separated-moment} is therefore the matched
conditional product identity, and the inequality is
Proposition~\ref{f15v20:witness}. This argument never
conditions the cell-source estimate on the common exterior;
it uses that estimate only after the exact product identity
has been integrated under the original full law.
\end{proof}

Since \(\Gamma_z\) is \(\mathcal Z_z\)-measurable, there is
also the exact formula
\begin{equation}
 \pi_z=1_{\Gamma_z^c}
                 +1_{\Gamma_z}\mu(D_z^c\mid\mathcal Z_z).
 \label{eq:f15v21:pi}
\end{equation}
Define the conditional-confidence bad event at threshold
\(0<\eta<1\) by
\begin{equation}
 \mathfrak B_z^\eta=\{\pi_z>\eta\}
   =\Gamma_z^c\cup
          \{\Gamma_z,\ \mu(D_z\mid\mathcal Z_z)<1-\eta\}.
 \label{eq:f15v21:confidence}
\end{equation}
Thus outside \(\mathfrak B_z^\eta\), all required parameters
are in their half charts and the \emph{actual} conditional
chart probability is at least \(1-\eta\). At \(\eta=1/2\),
these bad events contain V19's discarded events for all four
edge orientations. With a single orientation and its original
\(\Gamma\), they agree with them exactly.

\subsection{Joint confidence failures, including every coarse anchor}

For \(s\in\{p,\ell\}\), set
\begin{equation}
 m_s=8\ell/s,\qquad K_s=m_s^4,\qquad
 Q_{\eta,s}=\bigl[\min\{1,q_0/\eta\}\bigr]^{1/K_s}.
 \label{eq:f15v21:colour}
\end{equation}
These are integers where appropriate: \(p\mid\ell\) and
\(8\ell\mid M\), by the dyadic geometry of V20. Colour the
index \(z\in\mathbb T_{M/s}^4\) by \(z\bmod m_s\) in each
coordinate. Same-colour centres satisfy the separation just
proved, including across the periodic seam. Here
\(K_\ell=4096\), while \(K_p=(8\ell/p)^4\) is a scale cost,
not a total-volume cost.

\begin{theorem}[Full-law joint conditional-confidence bound]
\label{f15v21:joint-confidence}
For any \(A\subseteq\mathbb T_{M/s}^4\), with no restriction
on its cardinality,
\begin{equation}
       \mu\{\mathfrak B_z^\eta\text{ for all }z\in A\}
                               \le Q_{\eta,s}^{|A|}.
 \label{eq:f15v21:joint-confidence}
\end{equation}
\end{theorem}

\begin{proof}
For a same-colour set \(A\),
\eqref{eq:f15v21:threshold} and
\eqref{eq:f15v21:separated-moment} bound the probability by
\((q_0/\eta)^{|A|}\), capped at one. For a general nonempty
set, retain a largest colour class of size at least
\(\lceil|A|/K_s\rceil\). The event for all indices implies
the event for that class. The same ceiling argument as in
V20 proves \eqref{eq:f15v21:joint-confidence}. The empty-set
case is one.
\end{proof}

This conclusion is stronger than V20's generic rootwise
posterior bound in the particular matched specifications used
here. It does not contradict V20's warning about arbitrary
different sigma-algebras: the extra finite-range factorization
and event-support conditions have now been verified. In
particular, the equality in
\eqref{eq:f15v21:separated-moment} is not a formal property of
conditional expectation alone.

Let \(\mathfrak C_z^\eta\) be the nearest-neighbour component
of confidence-bad indices containing \(z\), empty if \(z\)
is good. V20's already proved combinatorial argument applies
unchanged, with \(Q=Q_{\eta,s}\):
\begin{align}
 \mu\{|\mathfrak C_z^\eta|\ge n\}
      &\le Q(64Q)^{n-1},\qquad n\ge1,
                              \label{eq:f15v21:cluster}\\
 \mathbb E_\mu w^{|\mathfrak C_z^\eta|}-1
      &\le\frac{Q(w-1)}{1-64Qw},\qquad w\ge1,\quad64Qw<1.
                              \label{eq:f15v21:cluster-cost}
\end{align}
The corresponding count Laplace bound is
\([1+Q(e^t-1)]^{|A|}\) for \(t\ge0\). These are new
applications to conditional-confidence failures, not new
abstract connected-set counting theorems.

\subsection{An all-anchor rate and the threshold tradeoff}

The most useful threshold here goes to zero, rather than
staying at \(1/2\). For sufficiently large selected regulators,
\(q_0<1\); take
\begin{equation}
 \eta=\sqrt{q_0},\qquad
 Q_s=Q_{\eta,s}=q_0^{1/(2K_s)}.
 \label{eq:f15v21:balanced}
\end{equation}
The conditional chart likelihood is then at least
\(1-\sqrt{q_0}\) outside a jointly rare confidence field.
For the mesoscopic cover \(Q_\ell=\sqrt q\), where \(q\)
is V20's parameter. The all-anchor parameter is weaker but
still sufficient for arbitrary fixed polynomial weights.

\begin{proposition}[All-anchor confidence tails remain superalgebraic]
\label{f15v21:rate}
On the existing growing scale \eqref{eq:f15v20:scale}, for all
sufficiently large selected regulators,
\begin{equation}
 \log Q_p\le-\widehat c\,
                   \frac{e^t}{t^{10\nu+138}},\qquad
 \widehat c=\frac{2^{-130}}{9\pi^2C_0^2},\qquad t=\log\beta.
 \label{eq:f15v21:rate}
\end{equation}
In particular \(\beta^A Q_p\to0\) for every fixed \(A>0\).
For any fixed polynomial weight \(1\le w_\beta\le C\beta^A\),
the right-hand side of \eqref{eq:f15v21:cluster-cost} tends to
zero faster than any fixed negative power of \(\beta\).
\end{proposition}

\begin{proof}
V20's proof gives, before its colour-root extraction,
\[
 \log q_0\le\log(320000)+(4\nu+60)\log t
                    -c_*e^t/t^{6\nu+82}
          \le-\tfrac12c_*e^t/t^{6\nu+82}
\]
eventually, where \(c_*=2^{-108}/(9\pi^2C_0^2)\).
Since \(\ell<L/2\), \(r\ge192\), and
\(r\le2p^{\nu+8}t^6\) eventually,
\[
 K_p=(8\ell/p)^4\le2^{16}r^4
                          \le2^{20}t^{4\nu+56}.
\]
Divide the negative bound on \(\log q_0\) by \(2K_p\).
The resulting constant is \(c_*2^{-22}=\widehat c\), proving
\eqref{eq:f15v21:rate}. Its negative exponential-in-\(t\)
term dominates every multiple of \(t\). Thus \(Q_p\) and
\(Q_pw_\beta\) have the asserted superalgebraic decay, and
\(1-64Q_pw_\beta\to1\). Substitute in the cluster moment
bound. This is an annealed geometric budget, not yet a bound
on the analytic cost of crossing a bad cluster.
\end{proof}

More generally, \(\eta=q_0^\theta\), \(0<\theta<1\), gives
\(Q_{\eta,s}=q_0^{(1-\theta)/K_s}\). A smaller allowed
conditional error costs a larger confidence-bad parameter.
All these choices preserve superalgebraic decay for fixed
\(\theta\) on the stated scale. They do not make the
whole-box cost \(e^{c\beta L^4}\) affordable; V20's
\eqref{eq:f15v20:oscillation-warning} still rules out that
particular budget, a fortiori for the weaker \(Q_p\).

\subsection{A pointwise conditional-score error with a sparse envelope}

Now take \(s=p\), so every coarse starting vertex is indexed.
For each orientation \(\mu\in\{1,2,3,4\}\), let \(e_{z,\mu}\)
be the coarse edge based there. Use V19's frozen-frame score
\(U_{z,\mu}^H\), its conditional mean and local predictor
\[
 M_{z,\mu}=\mathbb E[U_{z,\mu}^H\mid\mathcal Z_z],\qquad
 P_{z,\mu}=P_{e_{z,\mu}}(N_z).
\]
Keep the global bound \(B=B_p^*\) and the chart oscillation
\(\Omega=\Omega^*_{p,r,\beta}\) from V19. Assume
\(\beta\ge\beta_p\) when invoking that convex-chart estimate.
Define
\begin{equation}
 \tau_\eta=\Omega+2B\eta,\qquad
 R_z=\max_{1\le\mu\le4}|M_{z,\mu}-P_{z,\mu}|.
 \label{eq:f15v21:residual}
\end{equation}

\begin{theorem}[Sparse envelope for the actual conditional score error]
\label{f15v21:sparse-score}
Under the full compact law, pointwise up to one null set,
\begin{equation}
          R_z\le\tau_\eta+2B\,1_{\mathfrak B_z^\eta}.
 \label{eq:f15v21:pointwise}
\end{equation}
Consequently the nonnegative excesses
\(X_z=(R_z-\tau_\eta)_+/(2B)\) satisfy, for every set of
coarse anchors \(A\),
\begin{equation}
       0\le X_z\le1_{\mathfrak B_z^\eta},\qquad
       \mathbb E_\mu\prod_{z\in A}X_z\le Q_{\eta,p}^{|A|}.
 \label{eq:f15v21:excess}
\end{equation}
With \(\eta=\sqrt{q_0}\) on the selected growing scale,
\(\beta^A\tau_\eta\to0\) for every fixed \(A>0\).
The event \(\{R_z>\tau_\eta\}\) and its connected clusters
obey the same joint and cluster bounds as the confidence field.
\end{theorem}

\begin{proof}
Fix conditioning data outside \(\mathfrak B_z^\eta\). Then
\(\Gamma_z\) holds and
\(a_z=\mu(D_z\mid\mathcal Z_z)\ge1-\eta>0\). With the
\emph{normalized} interior chart mean \(M_{D,z,\mu}\), the
exact conditional mixture is
\[
 M_{z,\mu}=a_zM_{D,z,\mu}
                      +(1-a_z)M_{D^c,z,\mu}.
\]
If \(a_z=1\), the second summand is zero and its mean can be
assigned arbitrarily. Both means have norm at most \(B\),
since the original compact score does. Hence
\[
 |M_{z,\mu}-M_{D,z,\mu}|\le2B(1-a_z)\le2B\eta.
\]
The fixed-frame, fixed-domain exterior response argument in
Theorem~\ref{f15v19:local-score} bounds
\(|M_{D,z,\mu}-P_{z,\mu}|\) by \(\Omega\) on \(\Gamma_z\).
It applies simultaneously to the four central carriers;
taking their maximum costs no new factor. This proves
\(R_z\le\tau_\eta\) on the confidence-good event. Everywhere,
\(|M_{z,\mu}|,|P_{z,\mu}|\le B\), so \(R_z\le2B\).
This proves \eqref{eq:f15v21:pointwise}, and then
\eqref{eq:f15v21:excess} follows from
Theorem~\ref{f15v21:joint-confidence}.

On the chosen scale, V19 proves that \(\Omega\) is smaller
than every fixed negative power of \(\beta\); \(B\) is
\(\beta\) times a fixed polynomial in \(p\le\log\beta\).
V20's bound on \(q_0\) therefore gives the same decay for
\(B\sqrt{q_0}\). Finally every residual-bad index is a
confidence-bad index, so its root cluster is contained in
the corresponding confidence cluster.
\end{proof}

This is a pointwise small-floor plus rare-spatial-excess
statement before projecting away the exterior details. It
is not obtained by differentiating an exceptional probability
or a hard indicator. V19's globally defined compact scores
are used throughout, including on every bad configuration.

\subsection{Projection to the coarse action: an exact interface, not closure}

Write \(b_z(C)=\mu(\mathfrak B_z^\eta\mid C)\). Conditional
Jensen and the tower property give
\begin{equation}
 \left|\mathbb E[U_{z,\mu}^H\mid C]-P_{z,\mu}\right|
                  \le\tau_\eta+2B b_z(C).
 \label{eq:f15v21:projected-score}
\end{equation}
The frame and the predictor are \(C\)-measurable; rotation
back to the original edge frame preserves this inequality.
One has \(\mathbb E_{\rho_0}b_z\le Q_{\eta,p}\). A spatial
product bound for the \(b_z(C)\)'s has \emph{not} been proved
by the preceding argument: after retaining only \(C\), the
common exterior has been integrated out.

For completeness, V19's conditional covariance identity makes
the analytic interface particularly explicit. If
\(d(e_{z,\mu},f)\ge4rp\), then, in the original edge frames,
\[
 \mathsf H_{e_{z,\mu}f}(C)
 =-\mathbb E\left[
    (M_{z,\mu}^{\rm orig}-P_{z,\mu}^{\rm orig})
             (U_f-\mathbb E[U_f\mid C])^{\mathsf T}\mid C\right].
\]
The direct mixed term vanishes; the far score is
\(\mathcal Z_z\)-measurable. The displayed formula follows
by conditioning the first score on \(\mathcal Z_z\) and
subtracting the \(C\)-measurable predictor. Therefore
\begin{equation}
 \|\mathsf H_{e_{z,\mu}f}(C)\|_{\mathrm F}
                    \le2B\tau_\eta+4B^2 b_z(C).
 \label{eq:f15v21:hessian-interface}
\end{equation}
This reorganizes the exceptional part of the pairwise bound;
it does not yet make a row summable in \(f\). The same
\(b_z(C)\) appears for every far endpoint. A sum using just
\eqref{eq:f15v21:hessian-interface} again multiplies it by
the number of endpoints. The new joint spatial control is
available in the full canonical probability space, before
that projection, and a successful propagation argument must
use it there or prove an additional projection estimate.

\subsection{Why the matching hypothesis cannot be omitted}

Even positive finite densities show the issue. On two bits
\((x,y)\), put probability \(9/20\) on each equal pair and
\(1/20\) on each unequal pair. For
\(E_1=\{x=1\}\), \(E_2=\{y=0\}\), condition separately on
the other bit. Direct calculation gives
\[
 \mu(E_1\cap E_2)=\frac1{20},\qquad
 \mathbb E\bigl[\mu(E_1\mid y)\mu(E_2\mid x)\bigr]
                                               =\frac{61}{500}.
\]
Here the two interiors interact directly, so the product
identity fails. This example is not a Wilson counterexample;
it establishes why the verified separation step is essential.
Likewise, a mixture over a shared exterior generally correlates
interiors that were independent before that exterior was
integrated out. The diagnostic below checks both distinctions.

The next target is therefore a genuinely normalized
multiblock covariance or coupling estimate, using the proved
chart response on good regions and the now-controlled
\emph{conditional-confidence} clusters for its exceptions.
A contraction estimate for arbitrary transported boundary
data, or an affordable bound for crossing a bad component,
has not been supplied. The existence of local Gibbs update
kernels alone does not imply such contraction. Retained
gauge and low-frequency directions, continuum reconstruction
and physical time calibration remain separate from this
finite-regulator probability argument.

\subsection{Exact diagnostics and predecessor preservation}

The independent finite diagnostic is
\begin{center}\small
\path{scripts/verify_ym_f15_v21_conditional_confidence.py}.
\end{center}
It uses a strictly positive finite product-of-local-factors
density with a shared retained variable and a shared exterior,
neither discarded from its normalizers. There are 648 checks
of normalized common-exterior factorization, 864 checks against
separately computed individual conditional probabilities,
15 joint-event/product identities, 45 threshold comparisons,
96 bounded-score mixture/chart-confidence checks, 32 periodic
separation cases and 24 all-anchor rate cases. It also checks
the direct-interaction example and failure of factorization
after conditioning only on the retained variable. All
calculations use integers and exact fractions. They do not
certify the inherited Wilson source theorem, simulate the
Wilson law, or prove a nonlinear contraction estimate.
The script has 6,898 bytes and SHA--256
\begin{center}\scriptsize\ttfamily
BD03BED888E5198F6AC0709FD28078CEA0C41A81E82F26E3BE6E954ADE9A29F1.
\end{center}

The predecessor V20 has 456,676 bytes and SHA--256
\begin{center}\scriptsize\ttfamily
151B00B09489F6F3D5D84A0CAB9F0A38F6EBCDCFDE1483C564C72338D1CB7DD5.
\end{center}
Removing this terminal insertion and reversing the version
title recovers V20 byte for byte. V21 replaces it as the
sole active main note. The 14 protected sources and 63
earlier protected certificates are unchanged.

\begin{firewall}
V21 proves joint conditional-confidence tails for the actual
matched canonical specifications at every coarse anchor, and
a sparse pointwise envelope for their score-predictor errors.
It does not prove a full compact block contraction, a
bad-component covariance crossing bound, an intensive
retained-action interaction norm, physical clustering, an
interacting continuum theory, or the Yang--Mills mass gap.
\end{firewall}

\section{V22: full core kernels and separated sparse couplings}
\label{f15v22:section}

V21 controlled the conditional-confidence failures needed by
V19, and isolated rare spatial excesses in score means. A
mean estimate alone cannot be iterated as a comparison of
probability kernels. This section addresses that gap at the
level of a growing central core. The result is stronger than
a Lipschitz-test estimate: it controls the total variation
distance between full compact core marginals.

The main intermediate step is a uniform log-density-ratio
bound for \emph{normalized} chart core marginals. It comes
from applying V13's covariance estimate after the core has
been pinned, followed by two finite integrations. It does
not differentiate a rare-event probability. V21's conditional
chart likelihood then transfers the result to the full compact
kernel. For separated interiors we construct a simultaneous
coupling to local reference cores, with joint mismatch tails
under the original law. The retained field is kept unchanged.

These are finite-regulator, selected-coupling statements.
Different overlapping families are not being replaced by one
independent global field, and a static coupling is not being
declared a contraction in time. The archived f14 V48 compact
strong-mixing result has a different small-interaction entry
condition; it is not imported here. Nor is f14 V71's Gaussian
soft-pinning evolution reused as this product-chart pinning
argument. The companion checkpoint remains V25.

\subsection{Core geometry, metrics and fixed retained data}

Keep V21's full compact law, canonical variables, fine boxes
and confidence events. At one centre suppress the index \(z\).
Let \(J\) have anchor radius \((r-60)p\), with \(r\ge192\),
and choose the core
\begin{equation}
 A=\{i\in J:|\operatorname{anchor}(v_i)-z|_\infty
                        \le\lfloor r/2\rfloor p\},\qquad B=J\setminus A.
 \label{eq:f15v22:core}
\end{equation}
Here \(z\) denotes the fine centre for this display. All
later indexed cores are its translates. Only canonical detail
groups belong to \(A\); the retained groups are not integrated
as part of the core. An empty core has a one-point state space
and all the following comparisons are trivial.

Fix the whole retained field \(C\), and freeze its frame
\(H_N(C)\). Assume the required retained parameters in that
frame are in their half-radius charts. Let \(b\) denote
only the required \emph{exterior detail} logarithms, not the
retained parameters. We vary \(b\) within its product of
half-radius balls. No derivative of \(H_N\) or of \(C\)
is taken. The interior logarithms lie in the fixed domain
\(\mathcal D_p=\prod_{i\in J}\{|x_i|<a_p\}\).

Write \(n_A=3|A|\) for the number of core scalar coordinates
and \(n_\partial\) for the number of varying exterior scalar
coordinates. Conservative uniform bounds are
\begin{equation}
 n_A,n_\partial\le2^{14}r^4p^5,
 \qquad
 d(A,\partial J)\ge(r-60-\lfloor r/2\rfloor)p>32p.
 \label{eq:f15v22:counts}
\end{equation}
In this formula \(\partial J\) means the required exterior
detail anchors. The exterior count is the one used in V19.
For the core count, there are at most \(45j\le45p\) detail
groups at a fine anchor if all levels are overcounted. There
are at most \((3rp)^4\) anchors in the bounding box, and
three scalar coordinates per group. Thus the count is at
most \(10935r^4p^5<2^{14}r^4p^5\). The distance estimate
comes from the two specified radii.

We use the convention
\[
 d_{\rm TV}(P,Q)=\sup_E|P(E)-Q(E)|
                      =\tfrac12\int|dP-dQ|.
\]
Accordingly a TV bound \(\varepsilon\) controls every bounded
measurable test \(F\) by
\(|PF-QF|\le2\varepsilon\|F\|_\infty\). This is not a
restriction to differentiable observables.

\subsection{The covariance estimate survives pinning the core}

The measured chart potential is exactly
\(V_J(x_A,x_B;b)=\beta F_J-\ell_J\) from V13, including
the pivot and interior exponential-Haar contributions.
For fixed \(x_A\) in the closed core chart define
\begin{equation}
 d\nu_B^{x_A,b}(x_B)
   =\frac{e^{-V_J(x_A,x_B;b)}dx_B}
          {\int_{\mathcal D_B}e^{-V_J(x_A,y_B;b)}dy_B},
 \qquad \mathcal D_B=\prod_{i\in B}\{|x_i|<a_p\}.
 \label{eq:f15v22:pinned}
\end{equation}
This is a disintegration of the normalized chart law, not
a new Gaussian pinning or an unpinning interpolation.

\begin{proposition}[Uniform pinned-complement covariance]
\label{f15v22:pinned-covariance}
For \(\beta\ge\beta_p\), smooth functions of \(x_B\) with
gradient anchor carriers \(S,T\) satisfy
\begin{equation}
 |\operatorname{Cov}_{\nu_B^{x_A,b}}(F,G)|
 \le\frac4{\beta h_p}e^{-\alpha_p d(S,T)}
        \|\nabla_BF\|_{L^2}\|\nabla_BG\|_{L^2}.
 \label{eq:f15v22:pinned-covariance}
\end{equation}
The bound is uniform in the pinned core throughout its full
closed chart and in the stated exterior/retained half charts.
If a gradient vanishes, the corresponding covariance is zero.
\end{proposition}

\begin{proof}
The \(B\)-Hessian is a principal submatrix of V13's full
interior Hessian, evaluated at an allowed full-chart point.
It retains the lower bound \(3\beta h_p/4\), absolute row
bound \(2\beta D_p\), and range \(16p\). Pinning the core
does not move the product-ball domain \(\mathcal D_B\).
The reflecting scalar generator and absolute one-form
realization in the proof of Theorem~\ref{f15v13:covariance}
therefore apply on precisely this smaller product. Its
boundary forms are again nonnegative one-ball curvature
terms; the removal of factors creates no new nonproduct
boundary constraint.

For completeness, multiply the one-form block at anchor
\(i\in B\) by \(e^{\alpha_p d(i,T)}\). This preserves
the absolute form domain and commutes with its differential
and boundary parts. The remaining Hessian conjugation error
has row and column sums at most
\((e^{16p\alpha_p}-1)2\beta D_p\le\beta h_p/8\), exactly
as in V13. Coercivity gives the conjugated inverse norm
\(4/(\beta h_p)\). Inserting it into the same scalar/one-form
covariance identity proves \eqref{eq:f15v22:pinned-covariance}.
All constants are inherited principal-submatrix bounds, not
new dimension-dependent estimates. If \(B\) is empty the
law is a point mass and the covariance is zero.
\end{proof}

\subsection{A normalized core log-density identity}

The chart core marginal relative to core Lebesgue coordinates is
\begin{equation}
 p_b(x_A)=\frac{I_b(x_A)}{Z_b},\qquad
 I_b(x_A)=\int_{\mathcal D_B}e^{-V_J(x_A,x_B;b)}dx_B,
 \qquad Z_b=\int_{\mathcal D_A}I_b(y_A)dy_A.
 \label{eq:f15v22:marginal}
\end{equation}
It is strictly positive and smooth on the closed chart, with
the usual smooth extension to a neighbourhood. Both integrals
are over fixed domains. In particular, the core Haar factor
has not been removed: it is already present in \(V_J\).

\begin{proposition}[Distant mixed derivative of the core log density]
\label{f15v22:mixed-log}
For a core scalar coordinate \(i\) and an exterior scalar
coordinate \(k\),
\begin{align}
 \partial_{x_i}\partial_{b_k}\log p_b(x_A)
   &=-\mathbb E_{\nu_B^{x_A,b}}\partial_{x_i b_k}V_J
       +\operatorname{Cov}_{\nu_B^{x_A,b}}
                         (\partial_{x_i}V_J,\partial_{b_k}V_J),
                           \label{eq:f15v22:mixed-log}\\
 \left|\partial_{x_i}\partial_{b_k}\log p_b(x_A)\right|
   &\le\frac{16\beta D_p^2}{h_p}
              e^{-\alpha_p(d(i,k)-32p)}.
                           \label{eq:f15v22:mixed-bound}
\end{align}
\end{proposition}

\begin{proof}
Differentiate both normalizers in \eqref{eq:f15v22:marginal}:
\[
 \partial_{b_k}\log p_b(x_A)
   =-\mathbb E_{\nu_B^{x_A,b}}\partial_{b_k}V_J
          +\mathbb E_{\nu^{\rm box}_b}\partial_{b_k}V_J.
\]
The second term is independent of \(x_A\). Differentiating
the first normalized expectation in \(x_i\) gives
\eqref{eq:f15v22:mixed-log}, including its covariance sign.
Compact fixed-domain differentiation is legitimate; the
normalizers are positive and finite. Terms independent of
\(x_B\), including core-only Haar terms, cause no missing
covariance contribution.

By \eqref{eq:f15v22:counts}, \(d(i,k)>32p\). Finite range
therefore makes the direct derivative in
\eqref{eq:f15v22:mixed-log} zero. The two \(B\)-gradient
carriers are within \(16p\) of \(i\) and \(k\), respectively,
so their separation is at least \(d(i,k)-32p\). Each gradient
norm is bounded by the absolute derivative row \(2\beta D_p\).
Use Proposition~\ref{f15v22:pinned-covariance}; the resulting
constant is \(4(2\beta D_p)^2/(\beta h_p)\), as displayed.
\end{proof}

Define the explicit, uniform core budget
\begin{equation}
 \Delta_{p,r,\beta}
   =\frac{\beta h_p}{2}\,n_A n_\partial
                   e^{-\alpha_p(r/2-92)p}.
 \label{eq:f15v22:Delta}
\end{equation}
The exponent is positive for \(r\ge192\). When either
coordinate set is empty put \(\Delta=0\).

\begin{theorem}[Uniform normalized chart core comparison]
\label{f15v22:chart-ratio}
At the same retained field and frame, any two allowed exterior
half-chart values \(b,b'\) satisfy
\begin{equation}
 \operatorname{osc}_{x_A}\log\frac{p_b(x_A)}{p_{b'}(x_A)}
      \le\Delta,
 \qquad
       e^{-\Delta}\le\frac{p_b(x_A)}{p_{b'}(x_A)}\le e^\Delta.
 \label{eq:f15v22:ratio}
\end{equation}
Consequently
\begin{equation}
 d_{\rm TV}(p_b,p_{b'})\le\min\{1,e^\Delta-1\},\qquad
 D(p_b\Vert p_{b'}),D(p_{b'}\Vert p_b)\le\Delta.
 \label{eq:f15v22:chart-TV}
\end{equation}
The same statements hold for their pushforwards to the
framed compact core groups.
\end{theorem}

\begin{proof}
Join \(b\) to \(b'\) by a straight segment in the product
of half-radius logarithm balls, and join any two core points
by a segment in the full-radius product. These segments stay
inside the domains of all preceding bounds. Every exterior
scalar displacement is at most \(a_p\), and every core
scalar displacement at most \(2a_p\). Integrate
\eqref{eq:f15v22:mixed-bound} along both segments. Since
\(d(i,k)-32p\ge(r/2-92)p\), the oscillation is at most
\[
 2a_p^2 n_A n_\partial\,
     \frac{16\beta D_p^2}{h_p}e^{-\alpha_p(r/2-92)p}
       =\Delta,
\]
where \(a_p=h_p/(8D_p)\) was used in the equality.

Write \(R=\log(p_b/p_{b'})\). Normalization gives
\(\int e^R p_{b'}=1\); therefore its minimum is at most
zero and its maximum at least zero. The oscillation bound
then implies \(|R|\le\Delta\), proving the ratio inequalities.
They imply the displayed TV bound directly by integrating
the positive part of \(p_b-p_{b'}\). Integrating \(R\)
against \(p_b\), or \(-R\) against \(p_{b'}\), proves the
two relative-entropy bounds. Exponentiation is an injective
common coordinate change on the core chart; density ratios
and these distances are unchanged by its pushforward.
\end{proof}

This comparison concerns the entire core marginal, not just
one score or finitely many moments. It is nevertheless a
comparison of chart marginals until the next transfer step.

\subsection{A local reference kernel and full compact total variation}

For a centre \(z\), define \(\lambda_z(N_z)\) to be the core
pushforward of the normalized chart law with the required
exterior detail groups set to \(I\), in the frozen frame
\(H_{N_z}(C)\). This definition is used whenever the required
retained groups are in their half charts. Else define
\(\lambda_z\) to be normalized product Haar on the core.
The resulting probability kernel is measurable and depends
only on \(N_z\). Its chart normalizer is the actual integral
of the complete measured potential at that reference exterior.
No assertion is made that the \emph{full compact} law at
the identity exterior itself has high chart probability.

Let \(\kappa_z^{Z_z}\) denote the actual full compact core
marginal given \(Z_z=(C,v_{O_z}^{H_{N_z}})\), expressed in
those same framed core coordinates. Take any \(0<\eta<1\)
and V21's confidence event \(\mathfrak B_z^\eta\). Put
\begin{equation}
 \zeta=\min\{1,\eta+e^\Delta-1\}.
 \label{eq:f15v22:zeta}
\end{equation}

\begin{theorem}[Full compact core kernel estimate]
\label{f15v22:full-kernel}
For the actual conditional kernels,
\begin{equation}
 d_{\rm TV}(\kappa_z^{Z_z},\lambda_z(N_z))
                    \le\zeta+1_{\mathfrak B_z^\eta}.
 \label{eq:f15v22:full-TV}
\end{equation}
At fixed \(C\), if two exterior data values both lie outside
\(\mathfrak B_z^\eta\), their full compact core marginals
obey the sharper two-boundary comparison
\begin{equation}
 d_{\rm TV}(\kappa_z^{Z_z},\kappa_z^{Z'_z})
                         \le\min\{1,2\eta+e^\Delta-1\}.
 \label{eq:f15v22:two-boundary}
\end{equation}
\end{theorem}

\begin{proof}
Outside the confidence-bad event, \(\Gamma_z\) holds and
the actual conditional interior chart probability is at least
\(1-\eta\). The exact full conditional interior law is a
mixture of its normalized chart restriction and its chart
complement, the latter of weight at most \(\eta\). Its TV
distance from the chart restriction is at most \(\eta\).
Projection to the core cannot increase TV. Compare that chart
core marginal to the identity-exterior chart reference by
Theorem~\ref{f15v22:chart-ratio}. The triangle inequality
proves \eqref{eq:f15v22:full-TV} on the good event. On its
complement the universal TV bound by one suffices, including
the fallback definition of \(\lambda_z\).

For two confidence-good exteriors, project both mixtures to
their cores and compare the two chart core marginals directly.
The two complement weights cost at most \(2\eta\), and
the direct chart comparison costs at most \(e^\Delta-1\).
This proves \eqref{eq:f15v22:two-boundary}. Every mixture
weight and every chart density is normalized before comparison.
\end{proof}

For all coarse anchors, the positive excess
\[
 T_z=\bigl(d_{\rm TV}(\kappa_z^{Z_z},\lambda_z(N_z))
                                               -\zeta\bigr)_+
\]
is bounded by \(1_{\mathfrak B_z^\eta}\). V21 therefore gives
\(\mathbb E\prod_{z\in A}T_z\le Q_{\eta,p}^{|A|}\).
This extends its sparse scalar-score envelope to a probability
kernel discrepancy. Projecting away the exterior still only
gives
\[
 d_{\rm TV}(\mu(v_{A_z}^{H_{N_z}}\in\cdot\mid C),\lambda_z(N_z))
                  \le\zeta+\mu(\mathfrak B_z^\eta\mid C).
\]
No joint product bound for the probabilities on the right
after conditioning only on \(C\) is inferred.

\subsection{A simultaneous sparse coupling on a separated family}

Fix any family \(\mathcal S\) of centres separated as in
Proposition~\ref{f15v21:support}, for example one complete
colour class of either cover. There is no bound on the
cardinality of \(\mathcal S\) other than fitting in the torus.
Set
\begin{equation}
 q_{\rm conf}=\min\{1,q_0/\eta\},\qquad
 \delta_{\rm coup}=\min\{1,\zeta+q_{\rm conf}\}.
 \label{eq:f15v22:coupling-budget}
\end{equation}

\begin{theorem}[Separated cores admit a joint sparse-error coupling]
\label{f15v22:coupling}
There is a probability extension of the full compact law
carrying reference cores \(\widehat V_z\), \(z\in\mathcal S\),
with the following properties. The original \((C,v)\) retains
its exact law; the added cores satisfy
\begin{equation}
 \mathcal L((\widehat V_z)_{z\in\mathcal S}\mid C)
                         =\bigotimes_{z\in\mathcal S}\lambda_z(N_z).
 \label{eq:f15v22:target-law}
\end{equation}
For every \(A\subseteq\mathcal S\),
\begin{equation}
 \mathbb P\{v_{A_z}^{H_{N_z}}\ne\widehat V_z
                                 \text{ for all }z\in A\}
                                  \le\delta_{\rm coup}^{|A|}.
 \label{eq:f15v22:mismatch}
\end{equation}
Thus the mismatch count in any finite subfamily has moment
generating function at most
\([1+\delta_{\rm coup}(e^t-1)]^{|A|}\) for \(t\ge0\).
\end{theorem}

\begin{proof}
Let \(Y=(C,v_O)\) be the common exterior of all selected
interiors. V21 proves that, conditional on \(Y\), those
interiors are independent with their individual full compact
conditional kernels. Their framed core marginals are
\(\kappa_z^Y\). At fixed \(Y\), couple each of these with
\(\lambda_z(N_z)\) by the following common-part construction.
For two densities \(k,l\) relative to core Haar, let
\(a=\min(k,l)\), \(m=\int a\). Put mass \(a\) on the
diagonal; if \(m<1\), put the remaining mass in the product
of the normalized residual densities \((k-a)/(1-m)\) and
\((l-a)/(1-m)\), with total weight \(1-m\). The two
marginals are exactly \(k,l\), and the mismatch probability
is \(1-m=d_{\rm TV}(k,l)\). If \(m=1\), use only the
diagonal part. The densities and integrals are measurable
in \(Y\), so this defines a measurable coupling kernel.

Choose these couplings independently between selected blocks
conditional on \(Y\). One can retain the full original
interiors by filling each fringe from its actual conditional
law given the sampled core. Equivalently, disintegrate each
coupling with respect to its original core and attach its
reference coordinate to an original full-field sample.
Compact standard Borel spaces and the explicit densities
provide these conditional kernels. The original law is
unchanged. Conditional on \(Y\), the joint reference marginal
is the product of \(\lambda_z(N_z)\)'s. Since these kernels
depend only on \(C\), integrating the remaining exterior
proves \eqref{eq:f15v22:target-law}.

Let \(\varepsilon_z(Y)=d_{\rm TV}(\kappa_z^Y,\lambda_z(N_z))\).
The matched-specification result makes
\(1_{\mathfrak B_z^\eta}\) a common-exterior measurable
function for this family. By conditional independence of
the chosen couplings and \eqref{eq:f15v22:full-TV},
\[
 \mathbb P(\text{all mismatches in }A)
       =\mathbb E\prod_{z\in A}\varepsilon_z(Y)
       \le\mathbb E\prod_{z\in A}
                          (\zeta+1_{\mathfrak B_z^\eta}).
\]
Expand the last product. For each subfamily, the separated
confidence bound in V21 is \(q_{\rm conf}^{|S|}\), before
any colour-root loss. The binomial expansion gives
\((\zeta+q_{\rm conf})^{|A|}\). Capping the probability by
one proves \eqref{eq:f15v22:mismatch}. A second nonnegative
subset expansion proves the count moment bound.
\end{proof}

The common-part coupling is standard; see, for example,
\href{https://arxiv.org/abs/1708.03625}
{Jacob--O'Leary--Atchad\'e, \emph{Unbiased Markov chain Monte
Carlo with couplings}}, Algorithm 2. The short construction
above supplies the exact marginals needed here. No mixing
rate, sampling-cost estimate or dynamic theorem from that
work is used.

For a family of \(n\) cores, the same coupling gives the
normalized-law consequence
\begin{equation}
 d_{\rm TV}\left(
   \mathcal L_\mu(C,(v_{A_z}^{H_{N_z}})_{z\in\mathcal S}),
   \rho_0(dC)\bigotimes_{z\in\mathcal S}\lambda_z(N_z)
             \right)\le\min\{1,n\delta_{\rm coup}\}.
 \label{eq:f15v22:joint-TV}
\end{equation}
The joint mismatch estimate has no such prefactor; the
global event that \emph{some} mismatch occurs generally does.
The retained field has precisely the original marginal
\(\rho_0\) on both sides. It has not been replaced by an
independent or Gaussian coarse background.

\subsection{The growing-core error is quantitatively small}

Use exactly \eqref{eq:f15v20:scale}. Eventually \(r\ge368\),
so \(r/2-92\ge r/4\). Substituting
\(h_p=2^{-40}p^{-8}\) and the coordinate counts into
\eqref{eq:f15v22:Delta} gives
\begin{equation}
 \Delta\le2^{-13}\beta r^8p^2
                         e^{-\alpha_p(r/2-92)p}.
 \label{eq:f15v22:Delta-majorant}
\end{equation}
Since \(r\le2p^{\nu+8}t^6\), \(p\le t\), and
\(\alpha_p=2^{-49}C_0^{-1}p^{-(\nu+9)}\), this implies
\begin{equation}
 \Delta\le2^{-5}e^t t^{8\nu+114}
                   \exp\{-2^{-51}C_0^{-1}t^6\},
 \qquad t=\log\beta.
 \label{eq:f15v22:Delta-rate}
\end{equation}
Indeed the prefactor uses the upper bound on \(r\), whereas
the exponent uses its lower bound \(r\ge p^{\nu+8}t^6\).
Thus \(\beta^A\Delta\to0\) for every fixed \(A>0\).

Choose the same balanced threshold as V21,
\(\eta=\sqrt{q_0}\), when \(q_0<1\). Then
\begin{equation}
 \zeta\le\sqrt{q_0}+e^\Delta-1,\qquad
 \delta_{\rm coup}\le2\sqrt{q_0}+e^\Delta-1.
 \label{eq:f15v22:final-rate}
\end{equation}
V20 controls \(q_0\) superalgebraically; for \(0\le\Delta\le1\),
\(e^\Delta-1\le2\Delta\). Therefore
\(\beta^A\zeta\to0\) and
\(\beta^A\delta_{\rm coup}\to0\) for every fixed \(A\).
In particular \eqref{eq:f15v22:joint-TV} tends to zero
superalgebraically for any separated family whose cardinality
is bounded by a fixed power of \(\beta\). The all-volume
result remains the joint mismatch bound and its intensive
count control, not a small total variation distance for an
arbitrarily large product.

\subsection{The remaining consistency and propagation questions}

The result now controls arbitrary bounded functions of a
growing core under the full compact conditional law, and
provides a simultaneous reference coupling on each separated
family. This is a genuine kernel-level improvement over the
score-only predictor. It still stops before the required
spatial propagation argument for the coarse-action Hessian.

First, cores from different colour families may overlap.
Their separately constructed reference products need not
be marginals of one consistent field. Even two copies of
the same fair bit have a diagonal joint law at TV distance
\(1/2\) from two independent fair bits. Separation in
Theorem~\ref{f15v22:coupling} is therefore substantive, not
a disposable proof device.

Second, the construction is static. Repeating a rare Bernoulli
bit forever gives a stationary sequence with bad probability
\(q\) at each time but probability \(q\), not \(q^n\), for
\(n\) repeated failures. Neither stationarity of a future
block-update chain nor V21's spatial tails supplies the
space-time joint bounds needed for iteration. Such bounds
would require a separate argument using actual update kernels.

Third, the common retained field is untouched. The conditional
reference cores can be independent given \(C\) while their
unconditional laws remain correlated through \(C\). This
does not resolve the retained-field correlation obstruction
already isolated in V18. Nor is the TV estimate a derivative
estimate for changing \(C\); all comparisons here hold it
fixed. V19's full-density score identity must still be used
for retained-action derivatives.

The next target is a consistent spatial or multiblock
comparison using these core kernels, with its overlaps,
normalizers and bad-component crossings accounted for.
An all-exterior contraction or a convergent path expansion
cannot be assumed from the small pairwise TV budget alone.
The goal remains the full interacting continuum theory and
its physical Hamiltonian gap, neither of which follows yet.

\subsection{Exact diagnostics and predecessor preservation}

The finite diagnostic is
\begin{center}\small
\path{scripts/verify_ym_f15_v22_core_kernel_coupling.py}.
\end{center}
It independently forms 48 mixed log-marginal jets, including
direct terms and nontrivial boundary-dependent normalizers.
It checks 60 normalized density-ratio bounds, 120 compact
mixture and projection inequalities, 160 exact common-part
couplings, 64 multiblock marginal identities, 15 joint
mismatch inequalities with a dependent shared exterior,
and 24 core constant and rate cases. It also checks the overlap
and repeated-time distinctions above. All arithmetic uses
integers or fractions. These are finite algebraic diagnostics,
not a Wilson simulation, a verification of primitive suprema,
or an independent certification of the inherited chart and
source theorems. The script has 8,252 bytes and SHA--256
\begin{center}\scriptsize\ttfamily
7AFE82CD530EA468CBE5491E915B91BCCD24A24042BC99DA898257577319CC7E.
\end{center}

The predecessor V21 has 480,137 bytes and SHA--256
\begin{center}\scriptsize\ttfamily
932613231DD8B922BE64D83BD21CAFFDD2A5AB25551C9FFA9AF692327A005BC4.
\end{center}
Removing this terminal section and reversing the title
version recovers it exactly. V22 replaces V21 as the sole
active main note. All 14 protected sources and 64 earlier
protected certificates remain unchanged.

\begin{firewall}
V22 proves full compact growing-core TV comparisons and a
joint sparse-mismatch coupling for separated canonical cores,
using the measured chart estimates and V21's confidence tails.
It does not prove consistency of overlapping reference
products, a uniform dynamic or spatial contraction, a
summable retained interaction, physical clustering, an
interacting continuum construction, or the YM mass gap.
\end{firewall}

\section{V23: two-sided core entropy and the logarithmic exceptional cost}
\label{f15v23:section}

V22 obtained full compact core-kernel TV estimates and
separated sparse couplings. Combining overlapping couplings
or iterating them in time still requires estimates that have
not been proved. We therefore examine an upstream question:
can the actual compact complement be paid for in relative
entropy, where a likelihood bound costs its logarithm?

The answer here is affirmative at the core-kernel level.
The new result is a two-sided, per-core entropy comparison
under the full selected Wilson law, with constants independent
of the number of separated cores. Its proof retains all
normalizers and explicitly repairs a support obstruction in
the chart reference. It also bounds conditional mutual
information across the collar and conditional total correlation
between separated cores. The entropy chain rules themselves
were already used in V2/V8; they are not new results. What
is new is a finite, superalgebraically small budget for these
whole-core comparisons under the actual compact law.

This does not resolve the overlap or propagation problem.
An explicit positive finite example below shows why even
small two-sided entropy at a fixed law does not imply a
uniform heat-bath functional inequality. In particular the
new estimate will not be silently converted into a spectral
gap. The V22 source was checked against its recorded hash;
the next scheduled companion checkpoint remains V25.

\subsection{A local compact likelihood bound, with its actual cost}

Use V22's core \(A_z\), interior \(J_z\), and frozen
\(H_{N_z}(C)\)-frame. The full conditional core law is
\(\kappa_z^{Z_z}\), with \(Z_z=(C,v_{O_z}^{H_{N_z}})\).
All group reference measures in this subsection are normalized
Haar, not logarithmic-coordinate Lebesgue measures.

Let \(j_{\rm prim}\) be a one-step inverse-pivot Haar
Jacobian in the fixed completed-block convention. Define
\begin{equation}
 J_* =\sup_{\text{compact primitive inputs}}|\log j_{\rm prim}|,
 \qquad N_{p,r}=2^{20}r^4p^5,\qquad
 \mathfrak h=2(\beta+J_*)N_{p,r}.
 \label{eq:f15v23:compact-budget}
\end{equation}
If there are several primitive types, take their maximum.
The supremum is finite because the primitive Jacobians are
smooth and strictly positive on their compact domains.
No numerical enclosure of this fixed constant is asserted.

\begin{proposition}[Uniform local density bounds on the full compact fibre]
\label{f15v23:compact-density}
For every fixed \(C,v_{O_z}\), without a chart hypothesis,
the normalized full interior density relative to \(dH_{J_z}\)
and its core marginal density relative to \(dH_{A_z}\) lie
between \(e^{-\mathfrak h}\) and \(e^{\mathfrak h}\).
In particular
\begin{equation}
 e^{-\mathfrak h}\le
               \frac{d\kappa_z^{Z_z}}{dH_{A_z}}
                    \le e^{\mathfrak h}.
 \label{eq:f15v23:compact-density}
\end{equation}
\end{proposition}

\begin{proof}
V12's exact kernel is proportional to \(e^{-W_{J_z}}dH_{J_z}\).
The same fine-anchor support enumeration used in V13 puts
the anchor of each measured term touching \(J_z\) within
\(5p\) of an input it touches. This is a geometric support
bound of the compact maps, independent of chart entry.
Thus all these term anchors lie in the box of radius \(rp\).
There are at most six Wilson terms and four pivot terms
per level at a fine anchor when higher-level anchors are
overcounted. Consequently their total count is at most
\[
 (6+4j)(2rp+1)^4\le810r^4p^5\le N_{p,r},
 \qquad j\le p.
\]
Each Wilson term has oscillation at most \(2\beta\), and
each negative pivot logarithm at most \(2J_*\). There is
no extra exponential-Haar term in this compact-Haar density.
Hence the oscillation of \(W_{J_z}\) in the entire compact
interior, at any fixed exterior, is at most \(\mathfrak h\).
Its normalized density on a probability reference space is
therefore between \(e^{-\mathfrak h}\) and \(e^{\mathfrak h}\).
Integrating the fringe against normalized Haar preserves
these bounds. The fixed frame acts coordinatewise by
Haar-preserving maps, so it does not change the conclusion.
\end{proof}

This bound alone is far too costly for a useful uniform
perturbative mixing comparison. We shall use \(\mathfrak h\)
as a logarithmic cost, not multiply a bad-patch estimate by
\(e^{\mathfrak h}\). V20's whole-box cost warning remains
valid for the latter operation.

\subsection{Why a positive reference floor is necessary}

Let \(\lambda_z(N_z)\) be V22's normalized chart-reference
core kernel, with its Haar fallback when the retained inputs
are not in their half charts. Put
\begin{equation}
 v_A=H_{A_z}\{v_i\in\exp B(0,a_p)\text{ for every }i\in A_z\},
 \qquad \mathfrak v_A=\log(1/v_A).
 \label{eq:f15v23:chart-volume}
\end{equation}
On the chart branch and for a nonempty core, \(0<v_A<1\).
The reference is zero outside that chart, while
\eqref{eq:f15v23:compact-density} gives the actual core
positive mass at least \(e^{-\mathfrak h}(1-v_A)\) there.
Thus \(D(\kappa_z^{Z_z}\Vert\lambda_z)=+\infty\), even
when the actual conditional chart probability is very high.
A small TV error does not remove this support obstruction.

For \(0<\varepsilon\le1/2\) define the comparison kernel
\begin{equation}
 \lambda_z^\varepsilon(N_z)
                  =(1-\varepsilon)\lambda_z(N_z)+\varepsilon H_{A_z}.
 \label{eq:f15v23:floor}
\end{equation}
This modifies only the reference probability, not the Wilson
action, the block map, or the original law. It remains local
in \(N_z\), normalized, and has Haar density at least
\(\varepsilon\).

\begin{proposition}[Reference density and the chart-volume price]
\label{f15v23:reference-density}
Both the original and floored references have Haar density
at most \(e^{\mathfrak h}/v_A\), including on the fallback
branch. Moreover
\begin{equation}
 \mathfrak v_A
 \le |A_z|\log\frac{3\pi^3}{8a_p^3}
       =O\bigl(r^4p^5(1+\log p)\bigr)
 \label{eq:f15v23:volume-cost}
\end{equation}
with fixed programme constants.
\end{proposition}

\begin{proof}
In compact Haar variables the reference chart numerator at
a core value is bounded by \(e^{-\inf W_{J_z}}\) times the
fringe chart Haar volume. Its denominator is bounded below
by \(e^{-\sup W_{J_z}}\) times that same fringe volume and
the core volume \(v_A\). Their ratio is at most
\(e^{\mathfrak h}/v_A\). The fringe volume cancels; it is
not counted as an additional entropy cost. On the Haar
fallback branch the density is one. Convex mixing with
Haar preserves the displayed upper bound, which is at least one.

In the normalization used throughout these notes,
\[
 H_{\SU(2)}\{\vartheta(U)<a\}
     =\frac2\pi\int_0^a\sin^2t\,dt
     \ge\frac{8a^3}{3\pi^3}\quad(0<a\le\pi/2).
\]
The inequality uses \(\sin t\ge2t/\pi\). Here \(a_p<1/8\).
Product Haar and V22's core count give
\eqref{eq:f15v23:volume-cost}, using
\(a_p=2^{-43}C_0^{-1}p^{-(\nu+8)}\). An empty core has
\(v_A=1\) and contributes zero entropy.
\end{proof}

\subsection{Pointwise entropy bounds: chart, complement and confidence}

Keep V22's normalized chart-core ratio budget \(\Delta\)
and V21's bad confidence event \(\mathfrak B_z^\eta\).
Choose \(0<\eta\le1/4\). In particular
\[
 \mu(\mathfrak B_z^\eta)\le q_0/\eta
\]
at each centre, using the single-centre case before any
colour loss. Define
\begin{equation}
 \mathfrak l_\varepsilon=\mathfrak h+\log(1/\varepsilon).
 \label{eq:f15v23:ell}
\end{equation}

\begin{theorem}[Two-sided normalized conditional-core entropy]
\label{f15v23:pointwise-entropy}
For all actual conditioning data,
\begin{align}
 D(\kappa_z^{Z_z}\Vert\lambda_z^\varepsilon)
      &\le\mathfrak l_\varepsilon,
                                      \label{eq:f15v23:global-forward}\\
 D(\lambda_z^\varepsilon\Vert\kappa_z^{Z_z})
      &\le2\mathfrak h+\mathfrak v_A.
                                      \label{eq:f15v23:global-reverse}
\end{align}
Outside \(\mathfrak B_z^\eta\), the sharper bounds are
\begin{align}
 D(\kappa_z^{Z_z}\Vert\lambda_z^\varepsilon)
   &\le\Delta-\log(1-\varepsilon)
       +\eta\mathfrak l_\varepsilon+\eta\log(1/\eta),
                                      \label{eq:f15v23:good-forward}\\
 D(\lambda_z^\varepsilon\Vert\kappa_z^{Z_z})
   &\le\Delta-\log(1-\eta)+\varepsilon\mathfrak h.
                                      \label{eq:f15v23:good-reverse}
\end{align}
All four relative entropies are finite at the finite regulator.
\end{theorem}

\begin{proof}
The first global bound follows from the pointwise upper
bound on the actual density and the lower bound
\(\varepsilon\) on the reference density. For the reverse
direction, the reference density is at most
\(e^{\mathfrak h}/v_A\) and the actual density is at least
\(e^{-\mathfrak h}\). Integrate the corresponding upper
bound on each log likelihood ratio to obtain
\eqref{eq:f15v23:global-forward}--\eqref{eq:f15v23:global-reverse}.

On the confidence-good event, write the exact projected
conditional mixture as
\[
 \kappa_z^{Z_z}=(1-q)p_b+q r_b,\qquad
        q=\mu(D_z^c\mid Z_z)\le\eta.
\]
Here \(p_b\) is the normalized chart core marginal and
\(r_b\) the core marginal conditioned on the full interior
chart complement. They may overlap on the core; no disjoint
core-support assumption is made. Since
\(\lambda_z^\varepsilon\ge(1-\varepsilon)\lambda_z\), V22 gives
\[
 D(p_b\Vert\lambda_z^\varepsilon)
                       \le\Delta-\log(1-\varepsilon).
\]
If \(q>0\), the mixture and the global density bound imply
\(dr_b/dH_{A_z}\le e^{\mathfrak h}/q\), whence
\[
 D(r_b\Vert\lambda_z^\varepsilon)
                       \le\mathfrak l_\varepsilon+\log(1/q).
\]
Convexity of relative entropy in its first argument bounds
the mixture entropy by the mixture of these bounds. The
function \(x\log(1/x)\) is increasing on \([0,1/4]\), so
\(q\log(1/q)\le\eta\log(1/\eta)\). This proves
\eqref{eq:f15v23:good-forward}, with \(0\log(1/0)=0\).
Convexity here is simply the pointwise convexity of
\(u\log u\), integrated against the reference measure.

In the other direction, \(\kappa_z^{Z_z}\ge(1-q)p_b\), so
\[
 D(\lambda_z\Vert\kappa_z^{Z_z})
                 \le D(\lambda_z\Vert p_b)-\log(1-q)
                 \le\Delta-\log(1-\eta).
\]
Also \(D(H_{A_z}\Vert\kappa_z^{Z_z})\le\mathfrak h\).
Apply the same convexity to the two components of
\(\lambda_z^\varepsilon\) to obtain
\eqref{eq:f15v23:good-reverse}. No normalizer has been
replaced by a lower bound without keeping its cost.
\end{proof}

\subsection{An intensive averaged entropy budget}

Set
\begin{align}
 \mathcal E_+
   &=\Delta-\log(1-\varepsilon)
       +(\eta+q_0/\eta)\mathfrak l_\varepsilon
       +\eta\log(1/\eta),
                                  \label{eq:f15v23:Eplus}\\
 \mathcal E_-
   &=\Delta-\log(1-\eta)+\varepsilon\mathfrak h
       +(q_0/\eta)(2\mathfrak h+\mathfrak v_A).
                                  \label{eq:f15v23:Eminus}
\end{align}
Then the preceding good/bad split proves
\begin{equation}
 \mathbb E_\mu D(\kappa_z^{Z_z}\Vert\lambda_z^\varepsilon)
       \le\mathcal E_+,\qquad
 \mathbb E_\mu D(\lambda_z^\varepsilon\Vert\kappa_z^{Z_z})
       \le\mathcal E_-.
 \label{eq:f15v23:averaged}
\end{equation}
The expectation is over the actual conditioning data. In
particular, it is not a supremum over arbitrary retained
fields or exterior configurations. One obtains the displayed
bounds by using the good bound everywhere as a nonnegative
upper allowance, and adding the global bound on the bad
event of probability at most \(q_0/\eta\).

For large selected regulators choose
\begin{equation}
             \eta=\varepsilon=\sqrt{q_0}\le1/4.
 \label{eq:f15v23:balanced}
\end{equation}
The explicit budgets become
\begin{align}
 \mathcal E_+
       &=\Delta-\log(1-\eta)
                    +2\eta\mathfrak h+3\eta\log(1/\eta),
                                  \label{eq:f15v23:balanced-plus}\\
 \mathcal E_-
       &=\Delta-\log(1-\eta)
                    +3\eta\mathfrak h+\eta\mathfrak v_A.
                                  \label{eq:f15v23:balanced-minus}
\end{align}

\begin{proposition}[Both entropy errors are superalgebraically small]
\label{f15v23:rate}
On the growing scale \eqref{eq:f15v20:scale}, for every fixed
\(A>0\),
\begin{equation}
          \beta^A\mathcal E_+\longrightarrow0,
          \qquad \beta^A\mathcal E_-\longrightarrow0.
 \label{eq:f15v23:rate}
\end{equation}
\end{proposition}

\begin{proof}
V22 proves the assertion for \(\Delta\), and V20 for every
positive fixed power of \(q_0\). At \(t=\log\beta\),
\(r\le2p^{\nu+8}t^6\) and \(p\le t\) give
\[
 \mathfrak h\le2^{25}(1+J_*)e^t t^{4\nu+61},\qquad
 \mathfrak v_A\le C t^{4\nu+61}(1+\log t)
\]
for a fixed programme constant \(C\), using
\eqref{eq:f15v23:volume-cost}. Thus multiplication by these
costs does not destroy the decay of \(\eta\). Also
\(-\log(1-\eta)\le2\eta\) for \(\eta\le1/2\), and
\(\eta\log(1/\eta)\le2\sqrt\eta\) for \(0<\eta<1\).
Insert these estimates into
\eqref{eq:f15v23:balanced-plus}--\eqref{eq:f15v23:balanced-minus}.
\end{proof}

This is the precise change in budget: the exceptional terms
are \(\eta\mathfrak h\), \(\eta\mathfrak v_A\) and
\(\eta\log(1/\eta)\), not \(\eta e^{\mathfrak h}\).
It is a valid entropy estimate, not a retroactive justification
of the unaffordable multiplicative crossing cost of V20.

\subsection{Separated families and an exact conditional-information ledger}

Take any family \(\mathcal S\) of \(n\) separated interiors
as in V21/V22. Write \(Y=v_O\) for their common exterior,
excluding the already separately recorded retained field
\(C\), and \(V_z=v_{A_z}^{H_{N_z}}\) for their framed cores.
Conditional on \((C,Y)\), these cores are independent with
their actual kernels \(\kappa_z^{C,Y}\). Define the normalized
reference on these same variables by
\begin{equation}
 d\mathcal Q(C,Y,\mathbf V)
   =d\mu(C,Y)\prod_{z\in\mathcal S}\lambda_z^\varepsilon(N_z)(dV_z),
 \qquad \mathbf V=(V_z)_{z\in\mathcal S}.
 \label{eq:f15v23:joint-reference}
\end{equation}
Let \(\mathcal P\) be the actual joint marginal on
\((C,Y,\mathbf V)\). The exterior and retained marginal is
exactly the same under both laws.

\begin{theorem}[Two-sided separated-family entropy per core]
\label{f15v23:family}
There is no cardinality restriction on \(\mathcal S\) except
fitting in the finite regulator. One has
\begin{equation}
 D(\mathcal P\Vert\mathcal Q)\le n\mathcal E_+,\qquad
 D(\mathcal Q\Vert\mathcal P)\le n\mathcal E_-.
 \label{eq:f15v23:family}
\end{equation}
The same two bounds hold after projecting away \(Y\), with
reference \(\rho_0(dC)\prod_z\lambda_z^\varepsilon(N_z)(dV_z)\).
\end{theorem}

\begin{proof}
The actual product conditional law given \((C,Y)\) is the
matched disintegration proved in V21. Relative entropy of
these normalized finite products is the sum of their
component entropies in either direction. Integrating against
the common law of \((C,Y)\) gives
\[
 D(\mathcal P\Vert\mathcal Q)
     =\sum_z\mathbb E D(\kappa_z^{C,Y}\Vert\lambda_z^\varepsilon),
 \qquad
 D(\mathcal Q\Vert\mathcal P)
     =\sum_z\mathbb E D(\lambda_z^\varepsilon\Vert\kappa_z^{C,Y}).
\]
Each conditional kernel agrees with its single-interior
specification. Apply \eqref{eq:f15v23:averaged}. Marginalizing
\(Y\) cannot increase either relative entropy, by conditional
Jensen for \(u\log u\), proving the last assertion. Positivity
of the floor and the compact density bounds ensure finiteness
throughout. The proof uses no independent exterior law.
\end{proof}

Define conditional mutual information by its relative-entropy
formula, and define the conditional total correlation by
\[
 \mathcal T(\mathbf V\mid C)
   =\int D\left(\mu(d\mathbf V\mid C)
            \Big\Vert\prod_z\mu(dV_z\mid C)\right)d\rho_0(C).
\]
These are nonnegative KL quantities; no subtraction of
possibly infinite differential entropies is used.

\begin{proposition}[Exact information decomposition]
\label{f15v23:information}
For the above separated family,
\begin{align}
 \sum_z\mathbb E D(\kappa_z^{C,Y}\Vert\lambda_z^\varepsilon)
   &=I(\mathbf V;Y\mid C)+\mathcal T(\mathbf V\mid C)\notag\\
   &\quad+\sum_z\int D\left(\mu(dV_z\mid C)
                       \Vert\lambda_z^\varepsilon(N_z)\right)d\rho_0(C).
 \label{eq:f15v23:information}
\end{align}
Thus each term on the right is at most \(n\mathcal E_+\).
For one core, in particular,
\begin{equation}
               I(V_z;v_{O_z}\mid C)\le\mathcal E_+.
 \label{eq:f15v23:single-information}
\end{equation}
\end{proposition}

\begin{proof}
The left side is \(D(\mathcal P\Vert\mathcal Q)\). In its
log likelihood, insert first the actual marginal density of
\(\mathbf V\) given \(C\), then the product of its individual
core marginals given \(C\). The resulting three log ratios
are respectively the conditional mutual-information ratio,
the conditional total-correlation ratio, and the product
of the individual marginal-to-reference ratios. Integrate
under \(\mathcal P\). Marginalization gives exactly
\eqref{eq:f15v23:information}. All terms are finite by the
previous theorem and the positive compact densities; the
identities can also be obtained directly from the finite
conditional KL chain rule already used in V2. Their
nonnegativity gives the bounds. For \(n=1\), total
correlation is zero and the common exterior is \(v_{O_z}\).
\end{proof}

The theorem controls the information in the entire canonical
exterior about a growing core, conditional on the \emph{full}
retained field. It does not show that different retained
regions are weakly dependent. For \(n\) cores, the entropy
bound is extensive with a small coefficient; dividing by
\(n\) is uniform in family size. No infinite-volume entropy
density limit, and no volume-independent bound on the
unscaled total entropy, is asserted.

\subsection{Small two-sided entropy does not supply the missing inequality}

The distinction between a fixed-law entropy comparison and
a functional inequality for all test functions is essential.
Here is a positive finite obstruction to identifying them.
For \(0<\epsilon\le1/4\), put on two bits
\begin{equation}
 P_\epsilon(00)=1-\epsilon-2\epsilon^3,
 \quad P_\epsilon(01)=P_\epsilon(10)=\epsilon^3,
 \quad P_\epsilon(11)=\epsilon,
 \label{eq:f15v23:trap}
\end{equation}
and let \(Q_\epsilon\) be the product of its two marginals,
each with one-probability \(a=\epsilon+\epsilon^3\).
Then
\begin{equation}
 D(P_\epsilon\Vert Q_\epsilon)\longrightarrow0,
 \qquad D(Q_\epsilon\Vert P_\epsilon)\longrightarrow0.
 \label{eq:f15v23:trap-KL}
\end{equation}
However the best constant in
\[
 \operatorname{Var}_{P_\epsilon}g
       \le C_\epsilon\sum_{i=1}^2
                    \mathbb E_{P_\epsilon}\operatorname{Var}(g\mid x_{-i})
\]
satisfies
\begin{equation}
 C_\epsilon\ge\frac{(1-\epsilon)(1+\epsilon^2)}{2\epsilon^2}
                       \longrightarrow\infty.
 \label{eq:f15v23:trap-Poincare}
\end{equation}

\begin{proof}
The forward entropy is the mutual information of two bits,
at most their marginal binary entropy
\(-a\log a-(1-a)\log(1-a)\), which tends to zero. For the
reverse direction, \(a^2<\epsilon\), so the \(00\) and
\(11\) contributions to \(D(Q_\epsilon\Vert P_\epsilon)\)
are nonpositive. Each remaining reference probability is
at most \(a\le2\epsilon\), and its ratio to the actual
probability \(\epsilon^3\) is at most \(2/\epsilon^2\).
Thus the reverse entropy is at most
\(4\epsilon[\log2+2\log(1/\epsilon)]\), which also tends
to zero. All four probabilities are strictly positive.

For \(g=1_{\{11\}}\), the variance is
\(\epsilon(1-\epsilon)\). Each conditional variance average
is \(\epsilon^3/(1+\epsilon^2)\), by conditioning on the
other bit being one. Dividing gives
\eqref{eq:f15v23:trap-Poincare}. An entropy tensorization
inequality valid for every positive test function would
imply this variance inequality by applying it to
\(1+t(g-\mathbb Eg)\) and taking the second-order term at
\(t=0\). Its constant therefore cannot be uniform either.
\end{proof}

This is not a Wilson counterexample and does not satisfy all
the additional YM estimates as a package. It rules out the
specific inference from small two-sided fixed-law entropy
alone to a uniform functional inequality. For comparison,
\href{https://arxiv.org/abs/1405.0608}
{Caputo--Menz--Tetali, \emph{Approximate tensorization of
entropy at high temperature}}, Theorem 2.1, uses supremum
conditional-ratio influence conditions in addition to an
ergodicity hypothesis. Those hypotheses have not been
verified for the current compact YM specifications, and
that theorem is not invoked here.

The next noncircular target is to combine the small
information budget with a spatial comparison that is stable
under the actual relevant conditionings or perturbations.
Neither arbitrary conditioning on a rare exterior nor
concentration by an arbitrary test density is paid for by
the fixed-law average above. A consistent overlap estimate,
or a proved comparison for the required perturbation class,
is still needed before concluding a summable covariance
bound. Physical scaling, retained low modes, continuum
reconstruction and the Hamiltonian gap remain outstanding.

\subsection{Exact diagnostics and predecessor preservation}

The diagnostic script is
\begin{center}\small
\path{scripts/verify_ym_f15_v23_core_entropy.py}.
\end{center}
It checks 36 finite full/core density bounds, 24 two-sided
mixture entropy bounds, and 162 exact multiplicative
identities underlying the information decomposition. It
also checks forward and reverse product additivity and
data processing, four strictly positive rare-trap examples,
12 local incidence/core count cases, and the zero-floor
support obstruction. Logarithms are enclosed by rational
\(\operatorname{arctanh}\) series after dyadic range reduction;
the diagnostic uses no floating-point tolerances. Additivity
is checked both through exact likelihood factorization and
overlapping certified log intervals. These finite tests do
not simulate the Wilson law, certify primitive suprema, or
verify the inherited chart/source theorems independently.
The script has 8,788 bytes and SHA--256
\begin{center}\scriptsize\ttfamily
CB9D1256B2AD874CD190E1F6F96D91007C0FEB42462A7C8EF749FA63A9EABA58.
\end{center}

The predecessor V22 has 504,454 bytes and SHA--256
\begin{center}\scriptsize\ttfamily
318B1A79A9DEBC581DC849686AEAE122FD48A2413265AC2A74AB04A22E9C96A0.
\end{center}
Removing this terminal section and reversing the title
version recovers it exactly. V23 replaces V22 as the sole
active main note. The 14 protected sources and 65 earlier
protected certificates are preserved unchanged.

\begin{firewall}
V23 proves two-sided conditional-core entropy comparisons,
an intensive separated-family bound, and a small conditional
information budget under the actual full compact law. The
Haar floor is only in the comparison kernel. No uniform
functional inequality, overlap contraction, summable retained
interaction, physical clustering, interacting continuum
construction, or Yang--Mills mass gap has been established.
\end{firewall}

\section{V24: sparse product tilts with an intensive normalization cost}
\label{f15v24:section}

V23 controls a fixed-law entropy comparison. The next issue
is its stability when many local factors change that law.
A direct change-of-measure estimate with the oscillation of
the entire source would cost the exponential of the number
of factors. Here we use V21's joint confidence bounds to
control that normalization and the resulting exterior
posterior at a cost per core, uniformly in the family size.

The perturbations are specified precisely: positive core
weights divided by their local V23 reference expectations.
These compensating factors depend on the retained field;
they cannot be omitted from the definition of the new law.
We prove two-sided relative-entropy and pressure bounds for
this class, including its changed exterior distribution.
We also obtain an intensive baseline log-factorization
bound for the actual conditional generating function of
the uncompensated weights. These are not bounds under
arbitrary test densities or uncompensated physical source laws.

The tempting alternative of obtaining a pairwise coarse
Hessian bound from V23 would repeat V19, which already gives
a stronger squared-norm estimate. V16--V18 treat different,
physical source families with an explicit family-volume
cost and restricted growing cardinality. The new result
below has no cardinality restriction, but has the stated
local normalization and oscillation hypotheses. Neither
result subsumes the other. The V23 predecessor hash was
verified; the next companion review remains V25.

\subsection{The exact perturbation and its two normalizers}

Take a nonempty separated family \(\mathcal S\) from V21--V23,
with \(n=|\mathcal S|\). Its common exterior is
\(\xi=(C,Y)\), with law \(\nu=d\mu(C,Y)\). Write
\(\kappa_z^\xi\) for its actual core kernels and abbreviate
the positive, floored reference by
\[
       \lambda_z=\lambda_z^\varepsilon(N_z),\qquad
       0<\varepsilon\le1/2.
\]
Only within this section does \(\lambda_z\) denote the
\emph{floored} reference. All frames are frozen at \(C\).
The actual and reference joint laws are
\begin{equation}
 dP(\xi,\mathbf V)=d\nu(\xi)\prod_z\kappa_z^\xi(dV_z),
 \qquad dQ(\xi,\mathbf V)=d\nu(\xi)\prod_z\lambda_z(dV_z).
 \label{eq:f15v24:base-laws}
\end{equation}
The product formula for \(P\) is V21's matched canonical
disintegration, not an independence assumption about \(C,Y\).

For each core choose a measurable positive weight
\(w_z(N_z,V_z)\) such that, uniformly in its arguments,
\begin{equation}
                   1\le w_z\le g=e^\ell,\qquad \ell\ge0.
 \label{eq:f15v24:weights}
\end{equation}
Multiplication of a weight by a positive function of \(N_z\)
has no effect on the normalized factor below. Thus the
hypothesis can also be stated as a uniform logarithmic
oscillation bound, with a measurable local choice of shift.
No derivative of these weights or of a chart cutoff is used.
Define
\begin{align}
 b_z(N_z)&=\lambda_z(w_z),&
 R_z(N_z,V_z)&=w_z/b_z,\notag\\
 a_z(\xi)&=\kappa_z^\xi(w_z)/b_z,&
 \Psi(\xi)&=\prod_z a_z(\xi),\qquad
 Z=\mathbb E_P\prod_zR_z=\mathbb E_\nu\Psi .
 \label{eq:f15v24:normalizers}
\end{align}
Here \(b_z\) is a local reference normalizer; \(Z\) is the
global actual normalizer. Neither is set equal to one.
The normalized comparison laws are
\begin{equation}
        dP_w=Z^{-1}\prod_zR_z\,dP,
        \qquad dQ_w=\prod_zR_z\,dQ.
 \label{eq:f15v24:tilted-laws}
\end{equation}
The second law is normalized because each
\(\lambda_z(R_z)=1\). The first is also the corresponding
marginal of the full Wilson law reweighted by
\(Z^{-1}\prod_zR_z\); the unrecorded fringe conditionals
need no replacement. In particular it is a defined
perturbation of the original law, not a claim that the
unperturbed Wilson measure has been altered.

For any kernel \(k\), write \(k^{w_z}=w_zk/k(w_z)\).
The exact disintegrations of \eqref{eq:f15v24:tilted-laws} are
\begin{equation}
 dP_w=d\nu_w(\xi)\prod_z(\kappa_z^\xi)^{w_z},\qquad
 dQ_w=d\nu(\xi)\prod_z\lambda_z^{w_z},\qquad
                   d\nu_w=\Psi\,d\nu/Z.
 \label{eq:f15v24:posterior}
\end{equation}
Although the local factors cancel from the likelihood ratio
\(dP_w/dQ_w=Z^{-1}dP/dQ\), its forward expectation is now
under \(P_w\), not under \(P\). The exterior posterior
\(\nu_w\) must therefore be controlled.

\subsection{Sparse exceptions control the marked normalizer}

Let \(B_z(\xi)=1_{\mathfrak B_z^\eta}\), where
\(0<\eta\le1/4\), and put \(B=\sum_z B_z\).
These confidence events are measurable in the common
exterior by V21. Assume \(0<q=q_0/\eta<1\); the separated
family estimate gives, for every subset \(A\subseteq\mathcal S\),
\begin{equation}
       \mathbb E_\nu\prod_{z\in A}B_z\le q^{|A|}.
 \label{eq:f15v24:sparse}
\end{equation}
There is no colour-root loss for this separated family.
Let \(\zeta\) be V22's good-confidence TV error and set
\begin{equation}
 \theta_{\rm TV}=\min\{1,\zeta+\varepsilon\},\qquad
 d=(g-1)\theta_{\rm TV},\qquad 0\le d\le1/2.
 \label{eq:f15v24:d}
\end{equation}
The last inequality is a hypothesis on the source strength;
an explicit growing admissible window is proved below.

\begin{proposition}[Pointwise factor bounds]
\label{f15v24:factors}
For every exterior, up to the common conditional null set,
\begin{equation}
 g^{-1}\le a_z\le g,\qquad
 (1-d)e^{-\ell B_z}\le a_z\le(1+d)e^{\ell B_z}.
 \label{eq:f15v24:factor-bounds}
\end{equation}
Moreover, for every \(s\ge0\),
\begin{equation}
            \mathbb E_\nu e^{sB}
                    \le[1+q(e^s-1)]^n.
 \label{eq:f15v24:base-mgf}
\end{equation}
\end{proposition}

\begin{proof}
Both \(\kappa_z(w_z)\) and \(b_z\) lie in \([1,g]\).
Outside the confidence-bad event, the TV distance from
\(\kappa_z\) to the floored \(\lambda_z\) is at most
\(\theta_{\rm TV}\), by V22 and the mixture with Haar.
For a real function in \([1,g]\), difference of expectations
is at most \((g-1)\) times TV in the convention used here.
Since \(b_z\ge1\), this gives \(|a_z-1|\le d\) there.
On the bad event the global interval \([g^{-1},g]\)
implies the weaker second pair of bounds in
\eqref{eq:f15v24:factor-bounds}. Finally expand
\(e^{sB}=\prod_z[1+(e^s-1)B_z]\). Every coefficient is
nonnegative, so \eqref{eq:f15v24:sparse} bounds the expansion
by the binomial expression. No Bernoulli domination or
independence of the bad indicators is used.
\end{proof}

Define nonnegative constants
\begin{align}
 c_-&=-\log(1-d)+\ell q,\notag\\
 c_+&=\log(1+d)+\log[1+q(g-1)],\notag\\
 \chi&=\log\frac{1+d}{1-d}+\ell q.
 \label{eq:f15v24:c}
\end{align}

\begin{theorem}[Intensive pressure and changed bad-core density]
\label{f15v24:pressure}
Uniformly in the family size,
\begin{equation}
                 -c_-\le\frac{\log Z}{n}\le c_+.
 \label{eq:f15v24:pressure}
\end{equation}
For every \(s>0\), define
\begin{equation}
 b_s=\min\left\{1,
   \frac{\chi+\log[1+q(e^{\ell+s}-1)]}{s}\right\}.
 \label{eq:f15v24:bad-budget}
\end{equation}
Then
\begin{align}
 \frac1n\log\mathbb E_{\nu_w}e^{sB}
       &\le\chi+\log[1+q(e^{\ell+s}-1)],\notag\\
 \frac1n\mathbb E_{\nu_w}B&\le b_s.
 \label{eq:f15v24:changed-bad}
\end{align}
This is an average over the family, not a uniform estimate
for each specified core in the tilted law.
\end{theorem}

\begin{proof}
Multiplying \eqref{eq:f15v24:factor-bounds} gives
\[
 (1-d)^n e^{-\ell B}\le\Psi\le(1+d)^n e^{\ell B}.
\]
Jensen and \(\mathbb E_\nu B\le nq\) give
\(\log Z\ge\mathbb E_\nu\log\Psi
\ge n\log(1-d)-n\ell q=-nc_-\).
For the upper bound use \eqref{eq:f15v24:base-mgf} at
\(s=\ell\), obtaining \(Z\le e^{nc_+}\).
Likewise
\[
 \mathbb E_{\nu_w}e^{sB}
   =Z^{-1}\mathbb E_\nu[\Psi e^{sB}]
   \le e^{nc_-}(1+d)^n[1+q(e^{\ell+s}-1)]^n.
\]
Since \(c_-+\log(1+d)=\chi\), this proves the first
inequality in \eqref{eq:f15v24:changed-bad}. Jensen bounds
\(s\mathbb E_{\nu_w}B\) by its left logarithm; also
\(B\le n\). This proves the second inequality.
\end{proof}

The logarithmic gap between the pressure bounds is retained
in \(\chi\). In particular one must not take \(s\downarrow0\)
in the upper bound and differentiate it as though it were
an exact generating function. We use a positive, explicit
\(s\) instead.

\subsection{Two-sided entropy under the full product perturbation}

We first record the bounded-tilt inequality used on each
conditional kernel. This elementary comparison is not a
new functional inequality for the underlying probability.

\begin{proposition}[A normalized bounded tilt costs only its oscillation]
\label{f15v24:tilt-lemma}
For two probabilities \(p,k\) and one weight \(1\le w\le g\),
\begin{equation}
 D(p^w\Vert k^w)\le gD(p\Vert k),\qquad
 D(k^w\Vert p^w)\le gD(k\Vert p).
 \label{eq:f15v24:tilt-KL}
\end{equation}
The inequalities hold with infinite right sides as well.
\end{proposition}

\begin{proof}
It suffices to consider finite \(D(p\Vert k)\). Set
\(r=dp/dk\), \(A=p(w)\), and \(B_0=k(w)\). For a
nonnegative random variable \(r\) under a probability \(m\),
\[
 \operatorname{Ent}_m(r)
  =\inf_{a>0}\mathbb E_m[r\log(r/a)-r+a].
\]
This follows by differentiating the scalar expression in
\(a\), whose minimizer is \(\mathbb E_mr\); the zero case
follows by a limit. Since \(u\log u-u+1\ge0\), choosing
\(a=1\) for \(m=k^w\) gives
\[
 \frac{A}{B_0}D(p^w\Vert k^w)
 =\operatorname{Ent}_{k^w}(r)
 \le\frac g{B_0}\mathbb E_k[r\log r-r+1]
 =\frac g{B_0}D(p\Vert k).
\]
Use \(A\ge1\). Interchanging \(p,k\) proves the other
direction. All local kernels used below have positive
finite densities, so no singular limiting argument is needed.
\end{proof}

Abbreviate the pointwise constants proved in V23 by
\begin{align}
 K_+&=\mathfrak h+\log(1/\varepsilon),&
 K_-&=2\mathfrak h+\mathfrak v_A,\notag\\
 k_+&=\Delta-\log(1-\varepsilon)
             +\eta K_++\eta\log(1/\eta),&
 k_-&=\Delta-\log(1-\eta)+\varepsilon\mathfrak h.
 \label{eq:f15v24:kernel-budgets}
\end{align}
These bound the two conditional KL directions globally
and on \(B_z=0\), respectively. The same conservative
\(\mathfrak v_A\) can be used for all translated cores.
For any \(s>0\) put
\begin{align}
 J_+(s)&=\chi+\ell b_s+g(k_++K_+b_s),\notag\\
 J_-(s)&=c_++c_-+g(k_-+K_-q).
 \label{eq:f15v24:J}
\end{align}
The notation \(J_-\) is kept parallel to \(J_+\); its
displayed value does not in fact depend on \(s\).

\begin{theorem}[All-cardinality two-sided entropy per core]
\label{f15v24:entropy}
The perturbed laws in \eqref{eq:f15v24:tilted-laws} satisfy
\begin{align}
 D(\nu_w\Vert\nu)&\le n(\chi+\ell b_s),&
 D(\nu\Vert\nu_w)&\le n(c_++c_-),
                                      \label{eq:f15v24:exterior-KL}\\
 D(P_w\Vert Q_w)&\le nJ_+(s),&
 D(Q_w\Vert P_w)&\le nJ_-(s).
                                      \label{eq:f15v24:full-KL}
\end{align}
The same corresponding bounds hold under marginalization.
In particular, if \(\rho_w\) is the retained marginal of
\(P_w\), the two bounds in \eqref{eq:f15v24:exterior-KL}
bound \(D(\rho_w\Vert\rho_0)\) and
\(D(\rho_0\Vert\rho_w)\). The reference retained marginal
is exactly \(\rho_0\).
\end{theorem}

\begin{proof}
By \eqref{eq:f15v24:posterior},
\[
 D(\nu_w\Vert\nu)=\mathbb E_{\nu_w}\log\Psi-\log Z,
 \qquad D(\nu\Vert\nu_w)=\log Z-\mathbb E_\nu\log\Psi.
\]
Use \(\log\Psi\le n\log(1+d)+\ell B\) for the first
quantity, the lower pressure bound, and
\eqref{eq:f15v24:changed-bad}. For the second use
\(\mathbb E_\nu\log\Psi\ge-nc_-\) and the upper pressure
bound. This proves \eqref{eq:f15v24:exterior-KL}.

The conditional KL chain rule and the product kernels give
\begin{align*}
 D(P_w\Vert Q_w)
 &=D(\nu_w\Vert\nu)
       +\sum_z\mathbb E_{\nu_w}
          D((\kappa_z^\xi)^{w_z}\Vert\lambda_z^{w_z}),\\
 D(Q_w\Vert P_w)
 &=D(\nu\Vert\nu_w)
       +\sum_z\mathbb E_\nu
          D(\lambda_z^{w_z}\Vert(\kappa_z^\xi)^{w_z}).
\end{align*}
Proposition~\ref{f15v24:tilt-lemma} bounds each conditional
term in the first direction by \(g(k_++K_+B_z)\), and in
the second by \(g(k_-+K_-B_z)\). This uses a nonnegative
good allowance everywhere and adds a global allowance on
the bad event; no equality of these allowances is asserted.
Sum, use \(\mathbb E_{\nu_w}B\le nb_s\) and
\(\mathbb E_\nu B\le nq\), and insert
\eqref{eq:f15v24:exterior-KL}. Data processing gives all
marginal statements. Conditional on \(C,Y\), the reference
kernels depend only on \(C\), so its retained marginal is
unchanged, as claimed.
\end{proof}

After removing \(Y\), the reference is exactly
\(\rho_0(dC)\prod_z\lambda_z^{w_z}(dV_z)\). The same
chain rule as V23 therefore bounds conditional total
correlation and the sum of individual core-to-reference
entropies, now under \(P_w\), by \(nJ_+(s)\). This is a
perturbation-stable information budget in the declared
class. It is not a small global TV bound for arbitrarily
large \(n\), and it does not make the retained variables
independent.

\subsection{A growing source window with no volume restriction}

Use the selected scale of V20--V23, \(t=\log\beta\), and
set
\begin{equation}
 \eta=\varepsilon=\sqrt{q_0},\qquad q=\eta,\qquad
 c_\Delta=2^{-51}C_0^{-1},\qquad
                0\le\ell\le\tfrac14 c_\Delta t^6.
 \label{eq:f15v24:window}
\end{equation}
This is a window for the logarithmic oscillation of each
core weight, not for the unscaled physical action coupling.
Choose
\begin{equation}
                       s_*=\tfrac12\log(1/q)-\ell.
 \label{eq:f15v24:s}
\end{equation}

\begin{proposition}[Superalgebraic intensive errors throughout the window]
\label{f15v24:rate}
Eventually \(s_*>0\), \(d\le1/2\), and all preceding
hypotheses hold. Uniformly over all weights satisfying
\eqref{eq:f15v24:window} and all separated family sizes,
for every fixed \(A>0\),
\begin{equation}
 \beta^A\bigl(c_++c_-+b_{s_*}
                    +J_+(s_*)+J_-(s_*)\bigr)\longrightarrow0.
 \label{eq:f15v24:rate}
\end{equation}
The statement is per core; \(n\) need not grow only as a
power of \(\beta\).
\end{proposition}

\begin{proof}
V20's proof, before its colour-root loss, gives eventually
\[
 \log(1/q_0)\ge\frac{c_*}{2}\frac{e^t}{t^{6\nu+82}}.
\]
Since \(q=\sqrt{q_0}\), \(\ell=O(t^6)\), and an
exponential dominates every fixed polynomial, eventually
\begin{equation}
 s_*\ge\frac{c_*}{16}\frac{e^t}{t^{6\nu+82}},
 \qquad s_*\ge\ell,\qquad
 \frac{K_+}{s_*}
   \le\frac{2^{29}(1+J_*)}{c_*}t^{10\nu+143}+4.
 \label{eq:f15v24:log-cost-ratio}
\end{equation}
For the last inequality use V23's
\(\mathfrak h\le2^{25}(1+J_*)e^t t^{4\nu+61}\)
and \(\log(1/\varepsilon)=\log(1/q)=2(s_*+\ell)\).
The relation is only an upper bound on the logarithmic
cost ratio; it does not pay for \(e^{\mathfrak h}\).

V22 gives
\[
 \Delta\le2^{-5}e^t t^{8\nu+114}e^{-c_\Delta t^6},
 \qquad \theta_{\rm TV}\le2\eta+2\Delta
\]
eventually. Hence \(d\le2g(\eta+\Delta)\to0\).
The stronger product needed below obeys
\[
       g^2\Delta
           \le2^{-5}e^t t^{8\nu+114}e^{-c_\Delta t^6/2}.
\]
It is superalgebraically small in \(\beta\). Every fixed
positive power of \(q\), multiplied by \(g^m\), a fixed
power of \(\beta\), and a polynomial in \(t\), also tends
to zero for each fixed \(m\), by the displayed bound on
\(\log(1/q_0)\). Factors \(\log(1/q)\) can be absorbed
by decreasing a positive power of \(q\); for example
\(q\log(1/q)\le2\sqrt q\).

For \(d\le1/2\), \(\chi\le3d+\ell q\). At \(s=s_*\),
\(q e^{\ell+s_*}=\sqrt q\), so
\begin{equation}
           b_{s_*}\le\frac{\chi+\sqrt q}{s_*}.
 \label{eq:f15v24:balanced-bad}
\end{equation}
Equations \eqref{eq:f15v24:log-cost-ratio} and
\eqref{eq:f15v24:balanced-bad} imply
\[
 gK_+b_{s_*}
   \le g\left[\frac{2^{29}(1+J_*)}{c_*}
                    t^{10\nu+143}+4\right](\chi+\sqrt q).
\]
Its terms are bounded by polynomial multiples of
\(g^2\Delta\), \(g^2q\), \(g\ell q\), and \(g\sqrt q\),
all just paid for. Similarly \(\ell b_{s_*}\le
\chi+\sqrt q\), since \(s_*\ge\ell\).
The quantities \(gk_+\), \(gk_-\), and \(gK_-q\) have
the same decay by their explicit formulas, V23's
polynomial-exponential bounds on \(\mathfrak h,\mathfrak v_A\),
and the preceding treatment of logarithmic factors.
Finally \(-\log(1-d)\le2d\) and
\(\log(1+q(g-1))\le q(g-1)\) control \(c_\pm\).
Substitution in \eqref{eq:f15v24:J} proves the claim.
\end{proof}

The finite-source interval of V16 is not recovered by
pretending that its whole compact action oscillation is
small: for those sources a crude oscillation is of order
\(\beta p^4\), much larger than this window. The present
result instead permits arbitrarily many separated bounded
core probes and the larger explicit logarithmic-oscillation
window above, with their local compensation retained.

\subsection{The actual conditional source logarithm under the baseline law}

There is also an uncompensated statement, with a different
and explicit norm. Define the actual conditional source
factor and its reference logarithmic remainder by
\begin{equation}
 A_w(C)=\mathbb E_\mu\left[\prod_z w_z(N_z,V_z)\mid C\right],
 \qquad r_w(C)=\log A_w(C)-\sum_z\log b_z(N_z).
 \label{eq:f15v24:source-remainder}
\end{equation}
Thus \(-\log A_w\) is the exact change of the unnormalized
retained action for the source \(\sum_z\log w_z\).
This definition uses the actual compact law, not a product
substitute for its retained marginal.

\begin{theorem}[All-cardinality baseline source-log factorization]
\label{f15v24:source-factorization}
Under the preceding hypotheses,
\begin{equation}
             \frac1n\|r_w\|_{L^1(\rho_0)}\le c_++2c_-.
 \label{eq:f15v24:source-factorization}
\end{equation}
Consequently this per-core remainder is superalgebraically
small throughout \eqref{eq:f15v24:window}, with no family
cardinality restriction.
\end{theorem}

\begin{proof}
Conditional independence given \(\xi\), followed by the
tower property, gives the exact identity
\[
        A_w(C)=\left(\prod_z b_z(N_z)\right)
                                      \mathbb E_\nu[\Psi\mid C].
\]
Put \(F(C)=\mathbb E_\nu[\Psi\mid C]\), so \(r_w=\log F\).
Conditional Jensen and the lower pointwise bound on
\(\Psi\) give
\[
 r_w(C)\ge\mathbb E_\nu[\log\Psi\mid C]
       \ge n\log(1-d)-\ell\mathbb E_\nu[B\mid C].
\]
The right side is nonpositive. Hence its negative is an
upper bound for \((r_w)_-\), and
\(\mathbb E_{\rho_0}(r_w)_-\le nc_-\).
On the other hand Jensen and the upper pressure bound give
\[
 \mathbb E_{\rho_0}r_w\le
       \log\mathbb E_{\rho_0}F=\log Z\le nc_+.
\]
All quantities are finite, since the finite product factors
lie between \(g^{-n}\) and \(g^n\). Use
\(|r_w|=r_w+2(r_w)_-\) to obtain
\eqref{eq:f15v24:source-factorization}, and then apply
Proposition~\ref{f15v24:rate}.
\end{proof}

This bounds an actual generated-action remainder in the
original retained law. It does not bound it in the output
law tilted by \(A_w\), nor uniformly in \(C\), and no
derivative estimate is inferred by differentiating this
small \(L^1\) remainder. The positive part is controlled
through the true global normalizer, not through an
unproved pointwise bound on \(F\).

\subsection{Why the retained compensation cannot be discarded}

The uncompensated weight \(\prod_z w_z\) can be written
as \(\prod_z b_z(C)\prod_z R_z\). For a comparison of
the corresponding uncompensated laws, the exterior base
would therefore be
\begin{equation}
 d\nu_b(\xi)
      =\frac{\prod_z b_z(N_z)}{\mathbb E_\nu\prod_z b_z(N_z)}
                                                        d\nu(\xi).
 \label{eq:f15v24:retained-reweight}
\end{equation}
The joint confidence bound \eqref{eq:f15v24:sparse} is
under \(\nu\), not automatically under \(\nu_b\).
It cannot be transferred merely because every factor is
local in the retained field: correlations of that field
are one of the unresolved parts of the programme.

A positive finite example makes the missing step visible.
Let a shared retained bit \(C\) have probability \(q\) of
being one, put \(B_1=1_{\{C=1\}}\), and set \(B_z=0\)
for \(z>1\). All the subset bounds
\eqref{eq:f15v24:sparse} hold. Given \(C\), take independent
positive binary core kernels, with actual and reference
kernels equal except possibly at core one when \(C=1\).
For \(z>1\) let \(w_z=g^C\), and let \(w_1=1\).
These retained-only weights have \(b_z=w_z\), so the
reference-normalized factors \(R_z\) are all one. But
without compensation the probability of \(B_1=1\) becomes
\begin{equation}
                   \frac{qg^{n-1}}{1-q+qg^{n-1}},
 \label{eq:f15v24:amplification}
\end{equation}
which tends to one for fixed \(q>0,g>1\) as \(n\to\infty\).
For instance the two kernels at the exceptional first core
can be \((3/4,1/4)\) and \((1/4,3/4)\); all others can be
fair bits. This is not a Wilson counterexample, nor a claim
about a full lattice with specified geometric locality.
It disproves the inference that the original bad-event
estimate is preserved under arbitrary products of
retained-dependent normalizers in the abstract interface.

Thus a next genuine extension must control the relevant
retained reweightings in \eqref{eq:f15v24:retained-reweight},
or supply an overlap/conditioning estimate that avoids
them. Product-normalized probe entropy does not establish
a functional inequality for every test density. The
rare-trap obstruction in V23 and the non-summable common
exception in V21 remain applicable. No physical
Hamiltonian spectral statement follows from the present
per-core pressure and entropy comparison.

\subsection{Diagnostics and exact predecessor recovery}

The script
\begin{center}\small
\path{scripts/verify_ym_f15_v24_sparse_tilts.py}
\end{center}
checks 48 two-sided bounded-tilt entropy comparisons, six
dependent-exterior product families, 50 sparse subset
inequalities, and 1,344 exact tilted likelihood ratios.
The finite families have correlated bad indicators and a
retained bit correlated with their common exterior. Both
conditional entropy decompositions, the pressure and
marked-count bounds, and the changed exterior entropy
bounds are checked using rational logarithm enclosures.
Twelve examples check the uncompensated posterior formula,
and twelve exact exponent budgets check the source-window
margin. Twelve conditional-source factorizations and their
six intensive absolute-log remainder bounds are also checked.
No floating-point tolerance is used. These are
finite algebra diagnostics, not an independent verification
of the inherited Wilson estimates or primitive suprema.
The script has 9,303 bytes and SHA--256
\begin{center}\scriptsize\ttfamily
F2FE6B4948E2C5C0A68A0C4D1C2A2541ACB7596F5A848BB7C4C8A5624A66504A.
\end{center}

The predecessor V23 has 527,432 bytes and SHA--256
\begin{center}\scriptsize\ttfamily
6F0A82C4B15E0FA2606893FA94B77D0274ABE16E2820C269BD4BE0DFC5CD8563.
\end{center}
Removing this terminal section and reversing the title
version recovers that predecessor exactly. V24 replaces
V23 as the sole active main note. The 14 protected sources
and 66 earlier protected certificates remain unchanged.

\begin{firewall}
V24 proves an intensive, two-sided entropy and pressure
comparison stable under arbitrarily many separated,
locally reference-normalized core perturbations in its
declared source window. It pays for the changed exterior
posterior. It also bounds the actual generated-action
remainder per core in the baseline norm \(L^1(\rho_0)\).
Neither result proves stability under uncompensated
retained reweightings,
a uniform functional inequality, consistent overlap
contraction, a summable retained interaction, an interacting
continuum construction, or the Yang--Mills mass gap.
\end{firewall}

\section{V25: uncompensated source laws and the retained-output entropy budget}
\label{f15v25:section}

V24 controlled reference-normalized core perturbations but
left their retained compensation in the definition of the
tilted law. Here we remove that compensation for a precise
output comparison. The true source is the product of the
original weights. Its retained marginal is compared with
the original retained law multiplied by the product of
local reference expectations, with both global normalizers
included.

The new ingredient is a bound on the \emph{average} number
of bad cores under both changed exterior laws. It follows
from an entropy cost per core and the full-law sparse
moment bound. It does not preserve a small bad probability
at every named core. A fractional-power estimate then pays
for the changed partition function without an exponential
family-size loss. This produces two-sided retained-output
entropy and source-log estimates for arbitrary separated
family size, at a slower, explicitly algebraic rate.

This extends V24 rather than contradicting its rare-site
example. It also differs from V17's physical action-source
cube and its restricted growing family size. The present
weights must satisfy a global oscillation bound stated
below; a fixed-strength Wilson action source need not do so.
A positive example at the end shows that the stronger
full-core entropy conclusion cannot be inferred from this
output result alone.

\begin{record}[The required V25 companion checkpoint]
\label{f15v25:checkpoint}
The current file inventory and all 81 protected identities
were checked. The following current companion bodies and
their stated boundaries were revisited; their files are
unchanged from the preceding review.

SM continuum V72 proves fixed-mode, bounded-time many-copy
transport, not a cutoff-uniform physical quantum limit.
NS stability V26 gives rank-one kernel calculations; its
reference to a future V27 certification is not imported as
an available certificate. Shell mixing V16 assembles the
measured flat-source one-loop coefficient, not the missing
noncommuting packet or a joint nonperturbative remainder.
Parallel YM V5 rejects uniform tightness of extensive
global charge and retains local-density alternatives.

The supplied Einstein bridge V24 has a nonzero upstream
first-jet tangency derivative, not a root theorem. Native
stationarity V34 reduces an escape test to five symmetric
matrices whose actual moving-fibre entries are still to be
computed. SM source-selection V31 reconstructs ancestry
conditional on word-native typing; it does not establish
historical occurrence of those typing operations. None
supplies a new spatial mixing estimate for the present
Wilson law. Their same-law and scope distinctions support
keeping all source normalizers explicit here.

The suggestions file was also revisited as optional source
material. Its coefficient/convention proposals are not
instructions to identify a local one-loop law with the
full compact law or to merge the active files. No new
companion theorem is assumed in the proofs below. The next
scheduled companion checkpoint is V30.
\end{record}

\subsection{Literal bounded weights and the actual source perturbation}

Use V24's separated family \(\mathcal S\), \(n=|\mathcal S|\ge1\),
common exterior \(\xi=(C,Y)\), and its probability \(\nu\).
Retain the matched products
\[
             dP=d\nu\prod_z\kappa_z^\xi,
             \qquad dQ=d\nu\prod_z\lambda_z,
\]
where \(\lambda_z=\lambda_z^\varepsilon(N_z)\) is the
\emph{floored} reference. All frames are frozen at \(C\).
Let the actual source weights satisfy the literal bounds
\begin{equation}
       1\le w_z(N_z,V_z)\le g=e^\ell,\qquad \ell\ge0,
       \quad\hbox{uniformly in both }N_z\hbox{ and }V_z.
 \label{eq:f15v25:global-weights}
\end{equation}
A configuration-independent scalar rescaling of a weight
does not change its normalized source law. A shift depending
on \(N_z\) does change it and is \emph{not} free here. This
distinguishes the global hypothesis from an oscillation
bound only within each fixed-retained-field fibre.

Keep \(b_z=\lambda_z(w_z)\), \(a_z=\kappa_z(w_z)/b_z\),
and \(\Psi=\prod_z a_z\), but now define
\begin{align}
 W(\mathbf V,C)&=\prod_z w_z(N_z,V_z),&
 T(C)&=\prod_z b_z(N_z),\notag\\
 U(\xi)&=\prod_z\kappa_z^\xi(w_z)=T(C)\Psi(\xi),&
 Z_a&=\mathbb E_P W=\mathbb E_\nu U,\qquad
 Z_b=\mathbb E_Q W=\mathbb E_\nu T.
 \label{eq:f15v25:normalizers}
\end{align}
The uncompensated probabilities and their exterior laws are
\begin{align}
 d\widetilde P_w&=W\,dP/Z_a,&
 d\widetilde Q_w&=W\,dQ/Z_b,\notag\\
 d\nu_a&=U\,d\nu/Z_a,& d\nu_b&=T\,d\nu/Z_b.
 \label{eq:f15v25:laws}
\end{align}
Thus \(\widetilde P_w\) is the indicated marginal of the
actual full compact Wilson law tilted by \(W/Z_a\).
Conditional on \(\xi\), its core kernels are
\((\kappa_z^\xi)^{w_z}\); the reference kernels are
\(\lambda_z^{w_z}\). No \(b_z\) is divided out of the true
source \(W\).

Let \(\rho_a=C_*\nu_a\) and \(\rho_b=C_*\nu_b\). Then
\begin{equation}
 d\rho_a=\frac{A_w(C)}{Z_a}\,d\rho_0,
 \qquad d\rho_b=\frac{T(C)}{Z_b}\,d\rho_0,
 \qquad A_w(C)=\mathbb E_\mu[W\mid C].
 \label{eq:f15v25:outputs}
\end{equation}
The reference increment \(-\log T=-\sum_z\log b_z(N_z)\)
is local term by term. The reference measure \(\rho_b\)
is not a product measure: it retains the original
\(\rho_0\) and its correlations.

\subsection{Transporting a sparse count with an entropy-density cost}

Keep V24's \(B_z=1_{\mathfrak B_z^\eta}\), \(B=\sum_zB_z\),
and \(0<q=q_0/\eta<1\). Their joint subset bounds under
\(\nu\) imply
\begin{equation}
 \log\mathbb E_\nu e^{sB}\le n m_q(s),
 \qquad m_q(s)=\log[1+q(e^s-1)],\qquad s\ge0.
 \label{eq:f15v25:mgf}
\end{equation}
This is the proved sparse moment bound, not stochastic
domination or a claim about arbitrary exterior conditioning.

\begin{proposition}[Entropy-cost transport of the mean bad fraction]
\label{f15v25:transport}
If \(\tau\ll\nu\) has \(D(\tau\Vert\nu)\le nh\), then
for every \(s>0\),
\begin{equation}
      \frac1n\mathbb E_\tau B
             \le\min\{1,(h+m_q(s))/s\}.
 \label{eq:f15v25:entropy-transport}
\end{equation}
In particular both \(\tau=\nu_a\) and \(\tau=\nu_b\)
satisfy this with \(h=\ell\).
\end{proposition}

\begin{proof}
The normalized probability
\(d\nu_s=e^{sB}d\nu/\mathbb E_\nu e^{sB}\) gives
\[
 0\le D(\tau\Vert\nu_s)
      =D(\tau\Vert\nu)-s\mathbb E_\tau B
                         +\log\mathbb E_\nu e^{sB}.
\]
This is the same entropy inequality used in V1, now
applied to the sparse count. Use \eqref{eq:f15v25:mgf}
and \(B\le n\). For the last assertion,
\(1\le U,T\le g^n\) and \(Z_a,Z_b\ge1\), so
\[
       \frac{d\nu_a}{d\nu},\frac{d\nu_b}{d\nu}\le g^n,
       \qquad D(\nu_a\Vert\nu),D(\nu_b\Vert\nu)\le n\ell.
\]
Every density and count here is bounded at the finite
regulator, justifying all logarithmic identities.
\end{proof}

Fix \(s>0\) with \(s\ge\ell\), and abbreviate
\begin{equation}
                   u_s=\min\{1,(\ell+m_q(s))/s\}.
 \label{eq:f15v25:u}
\end{equation}
There is no assertion that
\(\nu_a(B_z=1)\le q\) or \(\nu_b(B_z=1)\le q\) for each
\(z\). The conclusion is the family-averaged estimate
\(\mathbb E_{\nu_a}B,\mathbb E_{\nu_b}B\le nu_s\).

\begin{proposition}[The fractional-power normalization estimate]
\label{f15v25:interpolation}
One has
\begin{equation}
           \frac1n\log\mathbb E_{\nu_b}e^{\ell B}
                                                  \le\ell u_s.
 \label{eq:f15v25:interpolation}
\end{equation}
\end{proposition}

\begin{proof}
The density bound in the previous proof gives
\(\mathbb E_{\nu_b}e^{sB}\le
e^{n\ell}\mathbb E_\nu e^{sB}\le e^{n(\ell+m_q(s))}\).
Since \(0\le\ell/s\le1\), concavity of \(x^{\ell/s}\)
or H\"older's inequality implies
\[
 \mathbb E_{\nu_b}e^{\ell B}
       \le(\mathbb E_{\nu_b}e^{sB})^{\ell/s}.
\]
This proves the bound without the cap in \(u_s\).
The bound \(B\le n\) gives \(\log\mathbb E_{\nu_b}e^{\ell B}
\le n\ell\), which supplies the cap. At \(\ell=0\)
both sides in \eqref{eq:f15v25:interpolation} are zero.
\end{proof}

The large pointwise cost \(g^n\) has not been declared
small. Its logarithm is multiplied by \(\ell/s\), using
the sparse count and normalization of \(\nu_b\). That
step is what changes an unusable total-volume bound into
the intensive price \(\ell u_s\).

\subsection{Pressure and the true retained-source output}

Let \(\theta_{\rm TV}\) and \(d=(g-1)\theta_{\rm TV}\le1/2\)
be exactly V24's constants. Its pointwise inequalities give
\[
             (1-d)^n e^{-\ell B}\le\Psi
                                      \le(1+d)^n e^{\ell B}.
\]
Define
\begin{align}
 L_-&=-\log(1-d)+\ell u_s,&
 L_+&=\log(1+d)+\ell u_s,\notag\\
 H_s&=L_-+L_+
             =\log\frac{1+d}{1-d}+2\ell u_s,&
 \Lambda&=Z_a/Z_b=\mathbb E_{\nu_b}\Psi.
 \label{eq:f15v25:budgets}
\end{align}

\begin{theorem}[Uncompensated pressure and two-sided output entropy]
\label{f15v25:output}
There is no restriction on the finite separated family size.
One has
\begin{align}
 -L_-&\le\frac1n\log\Lambda\le L_+,
                                      \label{eq:f15v25:pressure}\\
 D(\nu_a\Vert\nu_b),\ D(\nu_b\Vert\nu_a)&\le nH_s,
                                      \label{eq:f15v25:exterior-KL}\\
 D(\rho_a\Vert\rho_b),\ D(\rho_b\Vert\rho_a)&\le nH_s.
                                      \label{eq:f15v25:output-KL}
\end{align}
These are comparisons of the actual source-tilted output
with the normalized local-factor approximation in
\eqref{eq:f15v25:outputs}. Neither law is replaced by its
restriction to a chart.
\end{theorem}

\begin{proof}
Jensen and the bad-fraction bound under \(\nu_b\) give
\[
 \log\Lambda\ge\mathbb E_{\nu_b}\log\Psi
                 \ge n\log(1-d)-\ell\mathbb E_{\nu_b}B
                 \ge-nL_-.
\]
The upper pointwise bound on \(\Psi\) and
\eqref{eq:f15v25:interpolation} give
\(\log\Lambda\le n\log(1+d)+n\ell u_s=nL_+\).
The exterior likelihood ratio is exactly
\(d\nu_a/d\nu_b=\Psi/\Lambda\). Therefore
\begin{align*}
 D(\nu_a\Vert\nu_b)
       &=\mathbb E_{\nu_a}\log\Psi-\log\Lambda
         \le nL_++nL_-,\\
 D(\nu_b\Vert\nu_a)
       &=\log\Lambda-\mathbb E_{\nu_b}\log\Psi
         \le nL_++nL_-.
\end{align*}
The first line uses the transported count bound under
\(\nu_a\), not the baseline one under \(\nu\). The second
uses it under \(\nu_b\). Finally relative entropy decreases
under projection to \(C\), proving the output bounds.
\end{proof}

For completeness, Pinsker gives a total output TV bound
\(\|\rho_a-\rho_b\|_{\rm TV}\le\min\{1,\sqrt{nH_s/2}\}\).
This tends to zero only when the displayed total error
tends to zero. A small entropy \emph{per core} is not
being identified with small total TV at arbitrary volume.

\subsection{The actual source-log remainder in both changed laws}

Keep the actual logarithmic remainder of V24,
\begin{equation}
 r(C)=\log A_w(C)-\log T(C)
       =\log F(C),\qquad F(C)=\mathbb E_\nu[\Psi\mid C].
 \label{eq:f15v25:remainder}
\end{equation}
Because \(T\) depends only on \(C\),
\(\nu_b(dY\mid C)=\nu(dY\mid C)\) and
\(F=\mathbb E_{\nu_b}[\Psi\mid C]\). Also
\(d\rho_a/d\rho_b=F/\Lambda\).

\begin{theorem}[Source-log control under the true and reference outputs]
\label{f15v25:source-log}
Both changed laws have an intensive absolute-log bound:
\begin{equation}
 \frac1n\|r\|_{L^1(\rho_b)}\le L_++2L_- ,\qquad
 \frac1n\|r\|_{L^1(\rho_a)}\le2L_++L_- .
 \label{eq:f15v25:source-log}
\end{equation}
The norm under the original \(\rho_0\) keeps V24's
stronger superalgebraic bound. No derivative estimate is
inferred by differentiating any of these \(L^1\) errors.
\end{theorem}

\begin{proof}
Conditional Jensen gives
\[
 r(C)\ge\mathbb E_{\nu_b}[\log\Psi\mid C]
       \ge n\log(1-d)-\ell\mathbb E_{\nu_b}[B\mid C].
\]
The right side is nonpositive. Integrating its negative
bounds \(\mathbb E_{\rho_b}r_-\) by \(nL_-\).
Moreover \(\mathbb E_{\rho_b}r\le\log\mathbb E_{\rho_b}F
=\log\Lambda\le nL_+\). The identity
\(|r|=r+2r_-\) proves the first inequality.

For the true law the conditional likelihood is
\(d\nu_a(\cdot\mid C)/d\nu_b(\cdot\mid C)=\Psi/F(C)\).
Its nonnegative conditional relative entropy implies
\[
             r(C)\le\mathbb E_{\nu_a}[\log\Psi\mid C]
                    \le n\log(1+d)+\ell\mathbb E_{\nu_a}[B\mid C].
\]
The last expression is nonnegative, so it also bounds
\(r_+\), and \(\mathbb E_{\rho_a}r_+\le nL_+\).
On the other hand
\[
 \mathbb E_{\rho_a}r
       =D(\rho_a\Vert\rho_b)+\log\Lambda\ge-nL_-.
\]
Use \(|r|=2r_+-r\) to get the second inequality. All
conditional ratios are positive and bounded at this
finite regulator; there is no subtraction of undefined
conditional entropies.
\end{proof}

Thus both the actual source normalization and the
generated-action logarithm have now been controlled in
their own changed retained laws, for the declared weights.
This is an estimate for source increments on top of the
still-correlated action of \(\rho_0\), not a construction
of a weakly interacting retained Gibbs state.

\subsection{What also survives for the conditional core kernels}

For probabilities \(p,k\) and \(1\le w\le g\),
\begin{equation}
                  \|p^w-k^w\|_{\rm TV}
                             \le g\|p-k\|_{\rm TV}.
 \label{eq:f15v25:tilt-TV}
\end{equation}
Here is a direct proof retaining the normalizers. Relative
to \(p+k\), let \(m\) have density the minimum of the
two densities; its mass is \(1-\delta\), where
\(\delta=\|p-k\|_{\rm TV}\). Suppose \(A=p(w)\ge k(w)\).
The common part of \(p^w,k^w\) has mass at least
\(m(w)/A\). Thus their TV is at most
\((A-m(w))/A\le g\delta/A\le g\delta\).
Interchange \(p,k\) for the other ordering.

Set \(\delta_z(\xi)=
\|(\kappa_z^\xi)^{w_z}-\lambda_z^{w_z}\|_{\rm TV}\).
On \(B_z=0\), \eqref{eq:f15v25:tilt-TV} gives
\(\delta_z\le g\theta_{\rm TV}\), and everywhere it is
at most one. Hence under either changed exterior law,
\begin{equation}
 \frac1n\sum_z\mathbb E_{\nu_a}\delta_z,
 \quad\frac1n\sum_z\mathbb E_{\nu_b}\delta_z
                       \le g\theta_{\rm TV}+u_s.
 \label{eq:f15v25:averaged-TV}
\end{equation}
This is an average of conditional TV distances, not a
uniform comparison for every core or a full joint-law
entropy bound. It uses the same actual exteriors as the
source-law theorem.

\subsection{The rate is algebraic, not the V24 superalgebraic rate}

Use the selected growing scale, \(t=\log\beta\), and
\begin{equation}
 \eta=\varepsilon=\sqrt{q_0}=q,\qquad
 c_\Delta=2^{-51}C_0^{-1},\qquad
 0\le\ell\le c_\Delta t^6/4,\qquad
                   s=\tfrac12\log(1/q).
 \label{eq:f15v25:window}
\end{equation}
All weights must still satisfy the \emph{global} condition
\eqref{eq:f15v25:global-weights}. Eventually \(s\ge\ell\),
and V20's uncoloured estimate gives
\begin{equation}
 s\ge\frac{c_*}{8}\frac{\beta}{t^{6\nu+82}},\qquad
 u_s\le\frac{\ell+\sqrt q}{s}
        \le\frac8{c_*}\frac{t^{6\nu+82}}\beta(\ell+\sqrt q).
 \label{eq:f15v25:rate-u}
\end{equation}
Indeed \(m_q(s)=\log(1+\sqrt q-q)\le\sqrt q\).
V24 already proves that \(d\) and \(g\theta_{\rm TV}\)
are superalgebraically small uniformly in this window.
For \(d\le1/2\), \(\log((1+d)/(1-d))\le3d\), so
\begin{equation}
 H_s\le3d+\frac{16}{c_*}\frac{t^{6\nu+82}}\beta
                                      (\ell^2+\ell\sqrt q)
       =O\!\left(\frac{t^{6\nu+94}}\beta\right).
 \label{eq:f15v25:rate-H}
\end{equation}
The last big-O has fixed programme constants and is uniform
in the weights and \(n\). Its \(\ell^2\) coefficient can
be bounded by \(c_\Delta^2/c_*\) after substitution of the
window. Also
\[
 L_\pm=O(t^{6\nu+94}/\beta),\qquad
 g\theta_{\rm TV}+u_s=O(t^{6\nu+88}/\beta).
\]
These quantities tend to zero. For fixed bounded \(\ell\),
the displayed logarithmic power for \(H_s\) can instead
be taken as \(6\nu+82\). No faster-than-every-power claim
is made: the entropy-density transport loses the
superalgebraic rate of the compensated construction.

For example total output TV follows when
\(n t^{6\nu+94}/\beta\to0\), whereas the entropy-per-core
statement imposes no such volume restriction. A fixed
physical-volume continuum limit can have a different
relation among cutoff, family size and coupling; that
relation has not been identified with this sufficient
condition. Nor does the window cover fixed-strength
action sources with whole compact oscillation of order
\(\beta p^4\).

\subsection{Exact output agreement need not give small full-core entropy}

The distinction from V24's stronger compensated result is
substantive. The following finite positive example satisfies
the sparse-count and good-kernel conditions used above.
It also keeps an explicitly small baseline chart-complement
probability. It is not a Wilson model or a realization of
all the geometric and source hypotheses of V20--V23.

Let \(m\to\infty\), put \(q=2^{-m}\), \(n=4m+1\), and
let a retained bit \(C\) have probability \(r_0=q^2/2\)
of being one. Set \(B_1=1_{\{C=1\}}\), \(B_z=0\) for
\(z>1\); then every all-bad subset has probability at most
\(q^{|A|}\). Conditional on \(C\), all core bits are
independent. At the first core choose
\begin{equation}
 \lambda_1=(1-q/2,q/2),\qquad
 \kappa_1^0=(1-q^2/4,q^2/4),\qquad
 \kappa_1^1=(1/4,3/4).
 \label{eq:f15v25:example-kernels}
\end{equation}
All other kernels are identical fair bits in both laws.
The reference \(\lambda_1=(1-q)\delta_0+qH\) has exactly
the \(q\)-Haar floor on two atoms. On the good exterior,
\(\|\kappa_1^0-\lambda_1\|_{\rm TV}\le q/2\).
For the first-core chart \(\{0\}\), its actual baseline
complement probability satisfies
\[
 (1-r_0)q^2/4+r_0(3/4)\le q^2.
\]
Thus the example does not rely on an order-one baseline
chart-complement mass.

Choose \(w_1=1\) and \(w_z(C,V_z)=2^C\) for \(z>1\).
All weights satisfy the global bound with \(\ell=\log2\).
They depend only on the shared retained variable, so
\(a_z=1\), \(\Psi=1\), and
\begin{equation}
                  \nu_a=\nu_b,\qquad\rho_a=\rho_b,
                   \qquad Z_a=Z_b.
 \label{eq:f15v25:example-output}
\end{equation}
Their common probability of \(C=1\) is
\begin{equation}
       r_m=\frac{r_0 2^{n-1}}{1-r_0+r_0 2^{n-1}}
           =\frac{2^{2m-1}}{1-r_0+2^{2m-1}}\longrightarrow1.
 \label{eq:f15v25:example-posterior}
\end{equation}
Nevertheless
\begin{equation}
 \frac1nD(\widetilde P_w\Vert\widetilde Q_w)
                   \longrightarrow\frac3{16}\log2>0.
 \label{eq:f15v25:example-entropy}
\end{equation}

\begin{proof}
The exterior entropies vanish by
\eqref{eq:f15v25:example-output}. All conditional core
entropies except the first also vanish. The KL chain rule
therefore gives
\[
 D(\widetilde P_w\Vert\widetilde Q_w)
     =(1-r_m)D(\kappa_1^0\Vert\lambda_1)
                  +r_mD(\kappa_1^1\Vert\lambda_1).
\]
The good-kernel entropy tends to zero, directly from its
two terms in \eqref{eq:f15v25:example-kernels}. The bad one is
\[
 D(\kappa_1^1\Vert\lambda_1)
   =\frac14\log\frac{1}{4(1-q/2)}
            +\frac34\log\frac{3}{2q}
   =\frac34 m\log2+O(1).
\]
Use \(r_m\to1\) and \(n=4m+1\) to prove
\eqref{eq:f15v25:example-entropy}. Before tilting, replace
\(r_m\) by \(r_0\); the full entropy per core tends to
zero. The changed average bad fraction is \(r_m/n\to0\),
fully consistent with the mean-count bound, although its
first-site bad probability tends to one.
\end{proof}

The shared-bit example does not assert geometric locality
of a Wilson retained field. Its conclusion is that the
abstract sparse-field and conditional-kernel interface
does not by itself promote output entropy control to a
full-core entropy theorem. Extra information about the
actual retained correlations or rare conditional laws
would be needed for that stronger step.

\subsection{The remaining spatial and physical obligations}

The source approximation is now normalized without
compensation for the declared globally bounded probes,
and its actual output log error is paid in both changed
laws. These results do not establish a functional
inequality for arbitrary positive test densities. For
comparison, \href{https://arxiv.org/abs/2004.10574}
{Caputo--Parisi, \emph{Block factorization of the relative
entropy via spatial mixing}}, Theorem 2.3, assumes its
strong spatial mixing condition across all specified
boundary data. That hypothesis has not been verified
here, and their theorem is not being applied.

In particular there is still no consistent comparison on
overlapping source families, no summable interaction norm
for the full retained action, and no suppression theorem
for its low-frequency degrees of freedom at physical
scales. The distinction between the actual global weight
bound and a fibrewise one is another explicit obligation
for extending the source class. Continuum existence,
noncollapse, physical-time calibration, and the positive
Hamiltonian mass gap remain unproved. The next upstream
step must add spatial/conditioning information or a
stronger justified source class, not merely reapply the
entropy chain rule to the same output comparison.

\subsection{Exact diagnostics and predecessor preservation}

The diagnostic is
\begin{center}\small
\path{scripts/verify_ym_f15_v25_uncompensated_sources.py}.
\end{center}
It checks 48 exact bounded-tilt TV inequalities; six
dependent-exterior uncompensated families with 50 sparse
subset inequalities; 1,344 exact full likelihood ratios;
and 12 retained/source identities. It verifies the
entropy-density transport, fractional-power pressure
bound, both exterior/output KL inequalities, both changed-law
absolute source-log bounds, and both averaged conditional
TV bounds. Four positive Haar-floored examples check the
nonvanishing full-core entropy obstruction, its small
baseline chart complement and its exact output equality.
All probabilities are rational. Logarithms use rational
\(\operatorname{arctanh}\) enclosures with outward rounding
to a \(2^{-120}\) rational grid. There is no floating-point
tolerance. These finite checks are not a simulation or
independent certification of the inherited Wilson inputs.
The script has 8,403 bytes and SHA--256
\begin{center}\scriptsize\ttfamily
419BAD57669CF23803C853123802C36C136417A88142FFBD82088840EB542FC4.
\end{center}

The predecessor V24 has 550,881 bytes and SHA--256
\begin{center}\scriptsize\ttfamily
978C73C1FEFBE00851E1B0B030CA3BB7F4AC3D36D913507E12DAE74B96B71D19.
\end{center}
Removing this terminal section and reversing the title
version recovers it exactly. V25 replaces V24 as the sole
active main note. The 14 protected sources and 67 earlier
protected certificates remain unchanged.

\begin{firewall}
V25 proves uncompensated source pressure, two-sided
retained-output entropy, and actual source-log control in
both changed output laws for globally bounded separated
core weights, uniformly per core in family size. It also
proves a changed-law average conditional-TV bound. The
rate is algebraic; full joint core entropy need not have
the same conclusion. No general functional inequality,
overlap contraction, summable full retained action,
interacting continuum construction, or Yang--Mills mass
gap is claimed.
\end{firewall}

\section{V26: fixed-strength physical sources and full conditional-core locality}
\label{f15v26:section}

V25's source-output theorem has a global bounded-oscillation
hypothesis. It does not cover a fixed nonzero coupling to
V14's Wilson action source, whose compact oscillation can
grow as \(\beta p^4\). V16 covers those physical sources,
but its useful joint-background estimate restricts the
number of sources. We now extend the \emph{kernel}
comparison, not the source-output entropy theorem, to
fixed-strength physical sources in an arbitrary finite
separated family. The actual retained marginal is preserved.

The upstream observation is that V22 proved an oscillation
bound for a \emph{log density ratio}. A common positive
core weight preserves this oscillation exactly, however
large the weight. The source cost is instead paid when
transporting the frequency of actual chart failures.
The selected action-source generating-function bound
provides that cost per core. Directly decomposing the
full conditional law into its chart and complement avoids
a conditional-confidence threshold and its square-root
loss. Neither a new convexity assertion for the tilted
potential nor a bounded global source oscillation is needed.

The selected full Wilson law, completed blocking map,
baseline source centering, and geometric separation remain
unchanged. The V25 companion review includes the supplied
SM source-selection V31, native stationarity V34 and
exact-action Einstein bridge V24. No new YM estimate is
imported from them here; the next scheduled review is V30.
All 82 previously protected sources and certificates were
checked unchanged before this addition. The result below
uses the recorded Wilson inputs with their original
hypotheses, not an independent recertification of the
entire archive.

\subsection{Physical sources inside the framed cores}

Let \(\mathcal S\) be a family of \(n\ge1\) coarse-compatible
centres satisfying Proposition~\ref{f15v21:support}:
each pair has periodic separation at least \(8\ell\)
in some coordinate. Keep its embedded boxes, \(r\ge192\),
\(M\ge4L\), \(\beta\ge\beta_p\), selected source
hypotheses \eqref{eq:f15v20:selected-input}, and the
canonical interiors \(J_z\), cores \(A_z\), near outputs
\(N_z\) and frozen frames \(H_{N_z}(C)\). Write
\[
 V_z=v_{A_z}^{H_{N_z}(C)},\qquad
 O=\mathcal I\setminus\bigcup_{z\in\mathcal S}J_z,
 \qquad \xi=(C,v_O).
\]
Different frames are allowed on the different interiors;
their coordinatewise Haar-preserving action at fixed
\(C\) does not change the common-exterior factorization.

Use exactly V16's translated, baseline-centred sources
at the same block depth and orientation:
\begin{equation}
 S_z=\sum_{x\in\mathcal B_{p,z}}
                  (T_x-\mathbb E_\mu T_x),\qquad
 |\mathcal B_{p,z}|=n_p=p^2(2p-1)^2\le4p^4,
 \qquad \|S_z\|_\infty\le48\beta p^4.
 \label{eq:f15v26:sources}
\end{equation}
Here \(p=2^j\ge2\). The sets \(\mathcal B_{p,z}\)
are disjoint. The measured canonical source carrier
has radius at most \(12p\), by V14--V16, whereas
\(A_z\) has radius \(\lfloor r/2\rfloor p\).
Every retained parameter in the source belongs to
\(N_z\). Consequently there is a bounded measurable
function, fixed independently of the source parameters,
such that
\begin{equation}
                      S_z=s_z(N_z,V_z).
 \label{eq:f15v26:core-source}
\end{equation}
It has no dependence on the fringe \(J_z\setminus A_z\)
or exterior details. Gauge invariance of the source
allows this expression in the frozen frame. The baseline
centering constant remains the same at every \(\theta\).

For \(\tau=1/8\), set
\begin{equation}
 \begin{aligned}
 W_\theta&=\exp\left(\sum_{z\in\mathcal S}\theta_zS_z\right),
 &Z(\theta)&=\mathbb E_\mu W_\theta,\\
 d\mu_\theta&=\frac{W_\theta}{Z(\theta)}\,d\mu,
 &\rho_\theta&=C_*\mu_\theta,
 \qquad \theta\in[-\tau,\tau]^n .
 \end{aligned}
 \label{eq:f15v26:law}
\end{equation}
All laws are at a finite regulator, so these positive
weights and their normalizers are finite. There is no
restriction \(n\le\lfloor\log\beta\rfloor\) in this section.
The only cardinality restriction is that the family fit
the stipulated separated geometry on that regulator.

\subsection{An exact order-two source budget transports the bad count}

Retain the arbitrary-support CGF input
\eqref{f15v1:source-input}, with \(k=k_M\) and source
radius \(R_{\rm src}=R_M\). Assume \(2\tau<R_{\rm src}\),
as holds eventually. Disjoint base supports imply
\[
 \log Z(2\theta)\le k n n_pJ(2\tau),\qquad
 J(a)=-\log(1-a)-a.
\]
This is the same source estimate used in
\eqref{eq:f15v16:partition}; its arbitrary-support
form has no logarithmic bound on \(n\). Jensen gives
\(Z(\theta)\ge1\) because the centering is under \(\mu\).
Define the exact intensive order-two budget
\begin{equation}
 h_\theta=\frac1n\log\frac{Z(2\theta)}{Z(\theta)^2},
 \qquad h_\tau=k n_pJ(2\tau),\qquad
                         0\le h_\theta\le h_\tau.
 \label{eq:f15v26:h}
\end{equation}
In particular, it is an entropy \emph{density} budget,
not a small global likelihood-ratio bound.

\begin{proposition}[Source entropy and raw chart-failure count]
\label{f15v26:count}
Let \(E_z=\Gamma_z^c\cup D_z^c\) be V21's actual
full-law failure events, and put
\(R=\sum_{z\in\mathcal S}1_{E_z}\). For every \(s>0\),
define
\begin{equation}
 m_{q_0}(s)=\log[1+q_0(e^s-1)],\qquad
 a_\theta(s)=\min\left\{1,
                   \frac{h_\theta+m_{q_0}(s)}s\right\}.
 \label{eq:f15v26:a}
\end{equation}
Then
\begin{equation}
 D(\mu_\theta\Vert\mu)\le n h_\theta,
 \qquad
 \frac1n\mathbb E_{\mu_\theta}R\le a_\theta(s).
 \label{eq:f15v26:entropy-count}
\end{equation}
If \(0<q_0<1\), the choice \(s=\tfrac12\log(1/q_0)\)
gives the uniform bound
\begin{equation}
 \frac1n\mathbb E_{\mu_\theta}R\le\bar a_\tau,
 \qquad
 \bar a_\tau=
 \min\left\{1,\frac{2(h_\tau+\sqrt{q_0})}{\log(1/q_0)}\right\}.
 \label{eq:f15v26:abar}
\end{equation}
\end{proposition}

\begin{proof}
Write \(f_\theta=d\mu_\theta/d\mu\). Direct normalization
and Cauchy--Schwarz give
\[
 \mathbb E_\mu f_\theta^2
       =\frac{Z(2\theta)}{Z(\theta)^2}\ge1.
\]
Concavity of the logarithm, now under \(\mu_\theta\), yields
\[
 D(\mu_\theta\Vert\mu)
 =\mathbb E_{\mu_\theta}\log f_\theta
 \le\log\mathbb E_{\mu_\theta}f_\theta
 =\log\mathbb E_\mu f_\theta^2=n h_\theta.
\]
The upper bound in \eqref{eq:f15v26:h} follows from the
CGF bound and \(\log Z(\theta)\ge0\).

For every subfamily \(A\subseteq\mathcal S\), V21 gives
\(\mu(\cap_{z\in A}E_z)\le q_0^{|A|}\), before any
colour-root loss. Expanding the finite product proves
\[
 \mathbb E_\mu e^{sR}
 =\mathbb E_\mu\prod_z[1+(e^s-1)1_{E_z}]
 \le[1+q_0(e^s-1)]^n.
\]
For the normalized law
\(d\mu_s=e^{sR}d\mu/\mathbb E_\mu e^{sR}\), nonnegativity
of \(D(\mu_\theta\Vert\mu_s)\) gives
\[
 s\mathbb E_{\mu_\theta}R
 \le D(\mu_\theta\Vert\mu)+\log\mathbb E_\mu e^{sR}
 \le n[h_\theta+m_{q_0}(s)].
\]
Also \(R\le n\). This proves \eqref{eq:f15v26:entropy-count}.
At the proposed value of \(s\),
\(m_{q_0}(s)=\log(1+\sqrt{q_0}-q_0)\le\sqrt{q_0}\).
Substitute and use \(h_\theta\le h_\tau\).
\end{proof}

The count here is the count of \emph{actual} failures in
the full configuration, not a count of exteriors discarded
by a confidence threshold. The entropy inequality itself
is the one already used in V1 and V25; the new input is
the physical-source order-two budget together with the
unthresholded geometric events. This averaged estimate
does not bound a specified site's failure probability
by \(\bar a_\tau\) when source parameters are nonuniform.

\subsection{Common core weights preserve the chart log-oscillation}

\begin{proposition}[Positive reweighting is an isometry for log-ratio oscillation]
\label{f15v26:isometry}
Let \(P,Q\) be equivalent probability measures, and suppose
\(\operatorname{osc}\log(dP/dQ)\le\Delta<\infty\), where
oscillation is essential supremum minus essential infimum.
Let \(w>0\) be measurable with \(0<Pw,Qw<\infty\), and set
\(P^w=wP/(Pw)\), \(Q^w=wQ/(Qw)\). Then
\begin{equation}
 \operatorname{osc}\log\frac{dP^w}{dQ^w}
 =\operatorname{osc}\log\frac{dP}{dQ},\qquad
 e^{-\Delta}\le\frac{dP^w}{dQ^w}\le e^\Delta.
 \label{eq:f15v26:isometry}
\end{equation}
Consequently, with \(\sigma_\Delta=\min\{1,e^\Delta-1\}\),
\begin{equation}
 d_{\rm TV}(P^w,Q^w)\le\sigma_\Delta,\qquad
 D(P^w\Vert Q^w),D(Q^w\Vert P^w)\le\Delta.
 \label{eq:f15v26:chart-distances}
\end{equation}
No bound on the size or oscillation of \(w\) is required.
\end{proposition}

\begin{proof}
The exact Radon--Nikodym identity is
\[
 \log\frac{dP^w}{dQ^w}
   =\log\frac{dP}{dQ}+\log(Qw)-\log(Pw).
\]
The added term is constant. Since the measures and their
positive reweightings have the same null sets, this proves
the equality of essential oscillations. Normalization
gives \(\int(dP^w/dQ^w)\,dQ^w=1\), so zero lies between
the essential lower and upper bounds of its logarithm.
An oscillation at most \(\Delta\) therefore puts that
logarithm in \([-\Delta,\Delta]\). Integrating its positive
likelihood-ratio excess proves the TV inequality, and
integrating the logarithm in each direction proves the
two entropy inequalities.
\end{proof}

This elementary projective identity is not claimed as a
new general probability theorem. For the log-density
description of the Hilbert metric, see
\href{https://arxiv.org/abs/2309.02413}
{Cohen--Fausti, \emph{Hyperbolic contractivity and the
Hilbert metric on probability measures}}, Proposition 3.6.
No filtering or contractivity theorem from that work is
invoked. Importantly, the weight above must be the
\emph{same} function for the two measures. For a fair
two-point law \(P=Q\), weights \((1,K)\) and \((K,1)\)
give distance \((K-1)/(K+1)\), which tends to one although
the original log-ratio oscillation is zero.

Use V22's \emph{unfloored} reference \(\lambda_z(N_z)\):
the normalized chart core law at identity exterior, and
product Haar when its required retained half-chart fails.
Define, everywhere on its retained input space,
\begin{equation}
 w_z(N_z,V)=e^{\theta_zs_z(N_z,V)},\qquad
 \lambda_{z,\theta_z}(N_z,dV)
   =\frac{w_z(N_z,V)}{\lambda_z(N_z)(w_z)}
                                    \lambda_z(N_z,dV).
 \label{eq:f15v26:reference}
\end{equation}
The denominator is retained explicitly and is positive
and finite. This kernel depends only on \(N_z\) and its
own \(\theta_z\), not on any remote source parameter.
This is not the Haar-floored reference of V23--V25.

On \(\Gamma_z\), fix the retained field and an allowed
exterior half-chart value \(b\). Denote the baseline
normalized chart core marginal by \(p_{z,b}\). Since
\(w_z\) depends only on the core and retained data,
integration over the chart fringe commutes with its
multiplication. The actual source-tilted chart core law
is exactly \(p_{z,b}^{w_z}\). Theorem~\ref{f15v22:chart-ratio}
and Proposition~\ref{f15v26:isometry}, with \(b'=0\), give
\begin{equation}
 d_{\rm TV}\bigl(p_{z,b}^{w_z},
                          \lambda_{z,\theta_z}(N_z)\bigr)
 \le\sigma_\Delta,
 \qquad \Delta=\Delta_{p,r,\beta}.
 \label{eq:f15v26:weighted-chart}
\end{equation}
The two chart relative entropies are also at most
\(\Delta\). The source may destroy a directly attempted
Hessian bound for the tilted potential; none is used in
this argument. It would not apply unchanged to a source
depending on the exterior or fringe, since the effective
core weight could then depend on \(b\).

\subsection{Direct full-law transfer without a confidence threshold}

Let \(\nu_\theta\) be the actual \(\mu_\theta\)-law of
\(\xi\). No measured interaction touches two selected
interiors, and every source factor belongs to just its
own core at fixed \(C\). Thus the proof of
Proposition~\ref{f15v21:matched-products} applies with
interior potential \(W_z-\theta_zs_z\). It gives exact
conditional independence of the interiors under
\(\mu_\theta\) given \(\xi\), with each of their positive
normalizers retained. Conditioning additionally on the
other interiors does not change the selected kernel.

Let \(\kappa_{z,\theta}^{\xi}\) be its full compact
\emph{core} marginal. The matching failure probability is
\begin{equation}
 \pi_{z,\theta}(\xi)=\mu_\theta(E_z\mid\xi)
  =1_{\Gamma_z^c}
          +1_{\Gamma_z}\mu_\theta(D_z^c\mid\xi).
 \label{eq:f15v26:pi}
\end{equation}
The event \(\Gamma_z\) is common-exterior measurable.
All identities can be chosen simultaneously outside one
null set for the finite family.

\begin{theorem}[Fixed-strength full conditional-core locality]
\label{f15v26:full-kernel}
For every \(\theta\in[-\tau,\tau]^n\), almost surely
under its actual exterior law,
\begin{equation}
 d_{\rm TV}\bigl(\kappa_{z,\theta}^{\xi},
                       \lambda_{z,\theta_z}(N_z)\bigr)
 \le\min\{1,\sigma_\Delta+\pi_{z,\theta}(\xi)\}.
 \label{eq:f15v26:pointwise}
\end{equation}
For \(s>0\), put
\(\varepsilon_\theta(s)=\min\{1,\sigma_\Delta+a_\theta(s)\}\).
Then
\begin{align}
 \frac1n\sum_z\mathbb E_{\nu_\theta}
    d_{\rm TV}\bigl(\kappa_{z,\theta}^{\xi},
                        \lambda_{z,\theta_z}(N_z)\bigr)
       &\le\varepsilon_\theta(s),
                       \label{eq:f15v26:averaged-TV}\\
 \frac1n\sum_z\mathbb E_{\rho_\theta}
    d_{\rm TV}\bigl(\mu_\theta(V_z\in\cdot\mid C),
                        \lambda_{z,\theta_z}(N_z)\bigr)
       &\le\varepsilon_\theta(s).
                       \label{eq:f15v26:retained-TV}
\end{align}
For \(0<q_0<1\), both right-hand sides may instead be
bounded uniformly in \(\theta\) by
\begin{equation}
                  \bar\varepsilon_\tau
                   =\min\{1,\sigma_\Delta+\bar a_\tau\}.
 \label{eq:f15v26:epsilon}
\end{equation}
\end{theorem}

\begin{proof}
On \(\Gamma_z\), set
\(\alpha=\mu_\theta(D_z\mid\xi)\). Disintegration by
\(D_z\) gives, at the level of the core marginal,
\[
 \kappa_{z,\theta}^{\xi}
       =\alpha p_{z,b}^{w_z}+(1-\alpha)r_{z,b,\theta},
\]
where \(r_{z,b,\theta}\) is the core law conditional
on the chart complement when that event has positive
probability. Its support may overlap the chart core;
no disjointness of these core supports is asserted.
The source-conditioned chart law is exactly the common
reweighting identified above. Convexity of total variation
and \eqref{eq:f15v26:weighted-chart} yield
\[
 d_{\rm TV}(\kappa_{z,\theta}^{\xi},\lambda_{z,\theta_z})
       \le\alpha\sigma_\Delta+(1-\alpha)
       \le\sigma_\Delta+\pi_{z,\theta}(\xi).
\]
If \(\alpha=0\), the universal TV bound one proves the
same assertion without assigning a chart conditional on
a null event. If \(\Gamma_z\) fails, \(\pi_{z,\theta}=1\)
and the bound is again automatic. This proves
\eqref{eq:f15v26:pointwise}, including its cap.

Taking expectations and summing uses the exact tower
identity
\[
 \sum_z\mathbb E_{\nu_\theta}\pi_{z,\theta}
                   =\mathbb E_{\mu_\theta}R.
\]
Proposition~\ref{f15v26:count} proves
\eqref{eq:f15v26:averaged-TV}. For the retained projection,
\(\mu_\theta(V_z\in\cdot\mid C)\) is the mixture of
\(\kappa_{z,\theta}^{\xi}\) against the actual
\(\nu_\theta(dv_O\mid C)\), whereas the reference is
constant in \(v_O\). Convexity of TV followed by
integration over \(\rho_\theta\) proves
\eqref{eq:f15v26:retained-TV}. Finally use
\eqref{eq:f15v26:abar}. No exceptional-exterior deletion
or Markov threshold has been inserted.
\end{proof}

In particular, for bounded measurable local tests
\(F_z(N_z,V_z)\) with \(\|F_z\|_\infty\le B\),
\begin{equation}
 \frac1n\sum_z
 \left\|\mathbb E_{\mu_\theta}[F_z\mid C]
    -\lambda_{z,\theta_z}(N_z)(F_z)\right\|_{L^1(\rho_\theta)}
       \le2B\bar\varepsilon_\tau.
 \label{eq:f15v26:tests}
\end{equation}
The local predictor has no remote source parameter, but
the norm and the true conditional law do: they are taken
under the actual changed law. This is not an estimate
obtained by evaluating a baseline conditional law under
an arbitrarily chosen exterior distribution.

\subsection{A simultaneous coupling and its explicitly paid joint bounds}

\begin{theorem}[Static separated coupling under the physical source law]
\label{f15v26:coupling}
There is a coupling preserving the full original
\(\mu_\theta\)-configuration, with reference cores
\((\widehat V_z)_{z\in\mathcal S}\), whose retained/core
reference marginal is exactly
\begin{equation}
                \rho_\theta(dC)
                  \bigotimes_{z\in\mathcal S}
                         \lambda_{z,\theta_z}(N_z,d\widehat V_z).
 \label{eq:f15v26:target}
\end{equation}
Put \(M_z=1_{\{V_z\ne\widehat V_z\}}\) and
\(D_{\rm mis}=\sum_zM_z\). Then
\begin{equation}
 \frac1n\mathbb E D_{\rm mis}\le\bar\varepsilon_\tau,
 \qquad
 d_{\rm TV}\left(\mathcal L_{\mu_\theta}(C,(V_z)_z),
            \rho_\theta(dC)\bigotimes_z\lambda_{z,\theta_z}(N_z)\right)
       \le\min\{1,n\bar\varepsilon_\tau\}.
 \label{eq:f15v26:joint-TV}
\end{equation}
Set \(\delta_* =\min\{1,\sigma_\Delta+\sqrt{q_0}\}\).
For every subfamily \(A\) and \(u\ge0\),
\begin{align}
 \mathbb P(M_z=1\ \hbox{for all }z\in A)
  &\le\min\{1,e^{n h_\theta/2}\delta_*^{|A|}\},
                           \label{eq:f15v26:paid-subset}\\
 \mathbb E e^{uD_{\rm mis}}
  &\le e^{n h_\theta/2}[1+\delta_*(e^u-1)]^n.
                           \label{eq:f15v26:paid-MGF}
\end{align}
Thus when \(\delta_*<a\le1\),
\begin{equation}
 \mathbb P(D_{\rm mis}\ge na)
 \le\min\{1,\exp(n[h_\theta/2-d_{\rm Ber}(a\Vert\delta_*)])\},
 \label{eq:f15v26:paid-tail}
\end{equation}
where \(d_{\rm Ber}(a\Vert b)=a\log(a/b)
+(1-a)\log((1-a)/(1-b))\), with endpoint limits.
\end{theorem}

\begin{proof}
Given \(\xi\), couple each pair
\(\kappa_{z,\theta}^{\xi},\lambda_{z,\theta_z}(N_z)\)
maximally, independently over \(z\). One measurable
construction uses the minimum of their densities relative
to core Haar as the common part, followed by the product
of the two residuals, as in V22. It also covers zero
residual mass. The mismatch probability is exactly
\[
 \delta_z(\xi)=d_{\rm TV}
          (\kappa_{z,\theta}^{\xi},\lambda_{z,\theta_z}(N_z)).
\]
The full original interiors are conditionally independent
given \(\xi\); after sampling the coupled cores, lift
each original core to its fringe using its actual
conditional disintegration. This restores precisely the
full \(\mu_\theta\)-law, not merely its core projection.
Conditional on \(\xi\), the reference law is a product
depending only on \(C\). Integration over the actual
exterior therefore proves \eqref{eq:f15v26:target}.
Theorem~\ref{f15v26:full-kernel} bounds the expected
mismatch count. The union bound on at least one mismatch
then proves the joint TV assertion.

For any configuration event \(F\), Cauchy--Schwarz and
\eqref{eq:f15v26:h} give
\[
 \mu_\theta(F)\le e^{n h_\theta/2}\mu(F)^{1/2}.
\]
The source-tilted matched-product identity remains exact,
since its proof only uses the conditional factorization
and the support of \(E_z\). Hence for every subfamily
\(B\), including the empty one,
\begin{equation}
 \mathbb E_{\nu_\theta}\prod_{z\in B}\pi_{z,\theta}
  =\mu_\theta\Bigl(\bigcap_{z\in B}E_z\Bigr)
  \le e^{n h_\theta/2}q_0^{|B|/2}.
 \label{eq:f15v26:paid-failure}
\end{equation}
When \(\sigma_\Delta+\sqrt{q_0}<1\), use
\(\delta_z\le\sigma_\Delta+\pi_{z,\theta}\) and expand
the product over \(A\). Each term is nonnegative;
\eqref{eq:f15v26:paid-failure} bounds its expectation.
Summing gives
\(e^{n h_\theta/2}(\sigma_\Delta+\sqrt{q_0})^{|A|}\).
If that base is at least one, the capped statement
\eqref{eq:f15v26:paid-subset} is the universal probability
bound. Conditional independence of the couplings identifies
the left side with \(\mathbb E_{\nu_\theta}\prod_{z\in A}\delta_z\).
Another nonnegative subset expansion proves
\eqref{eq:f15v26:paid-MGF}. Exponential Markov and optimization
over \(u\ge0\) give \eqref{eq:f15v26:paid-tail}; at \(a=1\)
take the limit \(u\to\infty\).
\end{proof}

The source cost in these joint estimates uses the entire
family size \(n\), not just \(|A|\). It must not be
silently dropped. In particular these are not
volume-independent spatial product tails under every
inhomogeneous physical source background. Nor can one
differentiate \eqref{eq:f15v26:paid-MGF} at zero to obtain
a small mean: its right side need not equal one there.
The sharper mean bound comes independently from the
raw-count entropy estimate and direct chart mixture.

\subsection{A uniform algebraic rate, and the precise limit of the extension}

\begin{proposition}[No family-size restriction for the intensive kernel error]
\label{f15v26:rate}
On the selected scale \eqref{eq:f15v20:scale}, with
\(2\le p\le t=\log\beta\),
\begin{equation}
 \bar\varepsilon_\tau
      =O\left(\beta^{-1}t^{6\nu+86}\right),
 \qquad \tau=1/8,
 \label{eq:f15v26:rate}
\end{equation}
uniformly over the source cube and all permitted separated
family sizes. This is an algebraic upper bound, not a
faster-than-every-power error assertion.
\end{proposition}

\begin{proof}
Before any colouring loss, V20 gives eventually
\[
 \log(1/q_0)\ge\frac{c_*}{2}
                         \frac{\beta}{t^{6\nu+82}},
 \qquad c_* =\frac{2^{-108}}{9\pi^2C_0^2}.
\]
Also \(k\le5\), \(n_p\le4p^4\), so
\(h_\tau\le20J(1/4)p^4\). Substitution in
\eqref{eq:f15v26:abar} gives the explicit majorant
\begin{equation}
 \bar a_\tau
 \le\frac{80J(1/4)}{c_*}
                        \frac{t^{6\nu+86}}\beta
       +\frac4{c_*}\frac{t^{6\nu+82}}\beta\sqrt{q_0}.
 \label{eq:f15v26:rate-majorant}
\end{equation}
V22's unchanged chart bound is
\[
 \Delta\le2^{-5}e^t t^{8\nu+114}e^{-c_\Delta t^6},
 \qquad c_\Delta=2^{-51}C_0^{-1}.
\]
Thus \(\sigma_\Delta\le2\Delta\) eventually, and both
this term and the second term of
\eqref{eq:f15v26:rate-majorant} are smaller than every
fixed negative power of \(\beta\). The first term proves
\eqref{eq:f15v26:rate}. All constants are independent of
\(n\) and of the chosen source parameters. The inherited
chart threshold \(\beta\ge\beta_p\), source radius,
embedding and divisibility conditions hold eventually
on this same scale as in V20--V22.
\end{proof}

The local kernel extension is substantive: fixed-strength
physical action sources are now admitted without a global
\(O(t^6)\) oscillation restriction and without V16's bound
on the number of sources. Nevertheless its \(L^1\) and
TV norms are averaged over the family and the actual
changed law. The joint TV conclusion is small only when
\(n\bar\varepsilon_\tau\to0\). For the unscaled observable
\(S_z\), inserting \(B=48\beta p^4\) in
\eqref{eq:f15v26:tests} does not give a vanishing bound
on this scale. This section therefore does \emph{not}
extend V16's unscaled energy-mean or conditional-cumulant
theorems to arbitrary family sizes.

Neither does it compare the retained law \(\rho_\theta\)
to a law obtained by multiplying baseline retained
density by \(\prod_z\lambda_z(e^{\theta_zs_z})\). Those
local normalizers are present in
\eqref{eq:f15v26:reference}, but their physical-source
output comparison is a separate question. Since the
reference is unfloored, it can give zero mass to a core
set of positive full-law mass. Full forward relative
entropy can therefore be infinite, despite the TV bound;
the support obstruction recorded in V23 has not vanished.

The next useful upstream issue is spatial control of the
\emph{retained} law or a justified physical-source
normalizer comparison strong enough to propagate across
overlapping regions. The present coupling leaves the
retained field untouched. It is neither a dynamically
contracting update nor a single consistent independent
field assigned to overlapping families. An all-density
functional inequality and a physical-scale clustering
estimate remain additional obligations. No interacting
continuum construction or physical Hamiltonian gap
follows from this finite-regulator conditional result.

\subsection{Exact diagnostics and preserved predecessor}

The diagnostic is
\begin{center}\small
\path{scripts/verify_ym_f15_v26_physical_source_kernels.py}.
\end{center}
It checks 40 exact common-weight log-oscillation identities,
including weights of size \(10^{500}\), together with
the two chart KL bounds. Six correlated-exterior families
retain both core and fringe variables, with overlapping
core supports for the chart and its complement. They
check 50 raw/tilted sparse-subset inequalities, 384 direct
conditional mixture bounds, 384 exact maximal couplings,
468 retained/product-reference identities, and 18
source-cost mismatch count inequalities. They also check
the order-two entropy budget and raw-count transport.
Eight two-point examples demonstrate the need for a
\emph{common} weight in the oscillation identity.

All probabilities are rational. Logarithms have rational
\(\operatorname{arctanh}\) enclosures with outward rounding
to a \(2^{-120}\) grid; no floating-point tolerance is used.
These finite tests check the new probability interfaces,
not the inherited Wilson CGF, chart Hessian or continuum
construction. The script has 9,898 bytes and SHA--256
\begin{center}\scriptsize\ttfamily
50AC7602029CD9783CA8418662D703D75C27D4BE2E2430D3061D881DCE10134E.
\end{center}

The predecessor V25 has 573,949 bytes and SHA--256
\begin{center}\scriptsize\ttfamily
3F6D0516BE0FA3C9DE37989007923C587B49C04F4E90DD7AE8FD5D28020D142E.
\end{center}
Removing this terminal section and reversing the title
version recovers it exactly. V26 replaces V25 as the sole
active main note; all 14 protected sources and 68 earlier
certificates are preserved. The next companion checkpoint
remains V30.

\begin{firewall}
V26 proves a physical-source, full conditional-core TV
comparison for arbitrary finite separated families, with
an \(O(\beta^{-1}(\log\beta)^{6\nu+86})\) averaged error,
a static coupling preserving the actual retained law,
and joint mismatch estimates with their entire source
cost displayed. It does not establish small full-core
relative entropy, physical-source retained-output locality,
overlap contraction, retained-field clustering, an
interacting continuum limit, or the Yang--Mills mass gap.
\end{firewall}

\section{V27: physical-source log normalizers and the retained normalization obstruction}
\label{f15v27:section}

V26 admitted fixed-strength physical action sources to
the full conditional-core comparison. This section returns
to their \emph{normalizers}. It proves a local representation
of the actual conditional source logarithm, under actual
source-background output laws, for much larger separated
families than the logarithmic window in V16. The improvement
is not obtained by integrating V26's bounded-test estimate
against the unscaled energy. Instead, an exact identity
relates a local normalizer ratio to the baseline and tilted
chart probabilities. Its two signs can be projected using
different conditional laws.

There is also a sharp limitation. Small expected source-log
error under the true law does not justify exponentiating
that approximation and normalizing it as a new retained law.
An explicit positive finite model below has a uniform signed
source CGF bound and arbitrarily rare chart failures, yet
the resulting normalized local-reference output becomes
asymptotically singular from the actual output. It is a
counterexample to an inference from the available probability
interfaces, not a counterexample to Wilson Yang--Mills.
It makes the next missing estimate more precise.

The V26 file and all 83 earlier protected sources and
certificates were checked unchanged. The companion review
at V25 remains current; the next checkpoint is V30.
All Wilson assertions below retain the selected-law,
source-support and chart hypotheses of V26. The inherited
CGF and chart theorems are not independently recertified here.

\subsection{The exact local normalizer identity}

Keep V26's \(n\)-element separated family, actual laws
\(\mu_\theta\), outputs \(\rho_\theta\), common exterior
\(\xi=(C,v_O)\), and unfloored references \(\lambda_z(N_z)\).
The source weights are still
\(w_z=e^{\theta_zs_z(N_z,V_z)}\), with
\(\theta\in[-\tau,\tau]^n\), \(\tau=1/8\).
Put \(I_\theta=\{z:\theta_z\ne0\}\) and
\(m_\theta=|I_\theta|\). All sums of local errors below
are over \(I_\theta\); a zero source contributes exactly zero.
The ambient family size \(n\) is retained in the
change-of-law budgets, including remote background sources.

Write \(\nu_\theta\) for the actual exterior law and
\(\kappa_{z,0}^{\xi}\) for the baseline full core kernel.
Define the positive finite normalizers
\begin{equation}
 u_{z,\theta}(\xi)=\kappa_{z,0}^{\xi}(w_z),\qquad
 b_{z,\theta}(N_z)=\lambda_z(N_z)(w_z),\qquad
 L_{z,\theta}=\log\frac{u_{z,\theta}}{b_{z,\theta}}.
 \label{eq:f15v27:local}
\end{equation}
They use the same baseline centering and the same frozen
retained frame. With \(B_p=48\beta p^4\), set
\begin{equation}
 K=2\tau B_p=12\beta p^4,\qquad
 c_K=\frac{K}{1-e^{-K}}\le1+K.
 \label{eq:f15v27:K}
\end{equation}
Both expectations of \(w_z\) lie between
\(e^{-\tau B_p}\) and \(e^{\tau B_p}\), so
\(|L_{z,\theta}|\le K\) at \emph{every} exterior, including
the retained-chart fallback branch. This logarithmic cost
will multiply a failure probability; its exponential will
not be treated as a small constant.

\begin{proposition}[Two signs are paid by two different chart probabilities]
\label{f15v27:local-signs}
Let \(\pi_{z,\theta}(\xi)=\mu_\theta(E_z\mid\xi)\)
as in V26. Almost surely,
\begin{equation}
 (L_{z,\theta})_+\le\Delta+c_K\pi_{z,\theta},\qquad
 (L_{z,\theta})_-\le\Delta+c_K\pi_{z,0},
 \label{eq:f15v27:local-signs}
\end{equation}
where \(x_+=\max(x,0)\), \(x_-=\max(-x,0)\), and
\(\Delta=\Delta_{p,r,\beta}\) is unchanged from V22.
\end{proposition}

\begin{proof}
On \(\Gamma_z\), put
\(\alpha_0=\mu(D_z\mid\xi)\) and
\(\alpha_\theta=\mu_\theta(D_z\mid\xi)\).
The source is core-only. All remote source factors cancel
from this conditioned interior, and normalization gives
\[
 \alpha_\theta
      =\frac{\alpha_0 p_{z,b}(w_z)}{u_{z,\theta}},\qquad
 L_{z,\theta}
      =\log\alpha_0-\log\alpha_\theta
                 +\log\frac{p_{z,b}(w_z)}{\lambda_z(w_z)}.
\]
Here \(p_{z,b}\) is the baseline normalized chart core
law at the actual allowed exterior. The chart has positive
Haar volume and the full compact density is positive,
so its conditional probabilities are positive.
V22's density-ratio bound implies
\[
 \left|\log\frac{p_{z,b}(w_z)}{\lambda_z(w_z)}\right|
                                                        \le\Delta.
\]
This uses positivity of the same weight in both integrals;
there is no factor \(e^K\) in this chart comparison.
Since \(\log\alpha_0,\log\alpha_\theta\le0\),
\[
 (L_{z,\theta})_+\le\min\{K,\Delta-\log\alpha_\theta\},\qquad
 (L_{z,\theta})_-\le\min\{K,\Delta-\log\alpha_0\}.
\]
For \(0\le x\le1\), with \(-\log0=+\infty\),
\begin{equation}
              \min\{K,-\log(1-x)\}\le c_K x.
 \label{eq:f15v27:clipped-log}
\end{equation}
Indeed convexity puts \(-\log(1-x)\) below its chord
between zero and \(x_*=1-e^{-K}\); for \(x\ge x_*\)
the capped left side is at most \(K=c_Kx_*\le c_Kx\).
Also \(\min\{K,\Delta+y\}\le\Delta+\min\{K,y\}\)
for \(y\ge0\). On \(\Gamma_z\), insert
\(x=1-\alpha_\theta=\pi_{z,\theta}\) or
\(x=1-\alpha_0=\pi_{z,0}\). On \(\Gamma_z^c\), both
failure probabilities are one and the global bound
\(|L_{z,\theta}|\le K\le c_K\) suffices.
Finally \(e^K\ge1+K\) proves \(c_K\le1+K\).
\end{proof}

The identity exhibits the normalization issue directly:
the chart normalizer ratio is uniformly close, but the
full-to-chart probability ratio need not be. Bounding
only the mean of \(1-\alpha\) does not bound the mean of
\(-\log\alpha\) without a further cost. The finite global
logarithmic bound \(K\) is what permits the clipped estimate.

\subsection{Projection of the two signs to the actual retained source logarithm}

Define, under the baseline conditional law,
\begin{equation}
 \begin{aligned}
 A_\theta(C)&=\mathbb E_\mu[W_\theta\mid C],&
 T_\theta(C)&=\prod_z b_{z,\theta}(N_z),\\
 r_\theta(C)&=\log A_\theta(C)-\log T_\theta(C),&
 g_{z,\theta}(C)&=\mu_\theta(E_z\mid C).
 \end{aligned}
 \label{eq:f15v27:retained}
\end{equation}
The proposed local expression \(\log T_\theta\) is a sum
of functions of \((\theta_z,N_z)\). It contains no
remote source parameter in any one summand. The exact
\(A_\theta\), by contrast, uses the whole retained field.

\begin{theorem}[A sign-separated retained source-log estimate]
\label{f15v27:retained-signs}
For the actual conditional source generator,
\begin{equation}
 \begin{split}
 (r_\theta)_+&\le m_\theta\Delta
                        +c_K\sum_{z\in I_\theta}g_{z,\theta}(C),\\
 (r_\theta)_-&\le m_\theta\Delta
                        +c_K\sum_{z\in I_\theta}g_{z,0}(C).
 \end{split}
 \label{eq:f15v27:retained-signs}
\end{equation}
No conditioned chart-likelihood lower bound is assumed.
\end{theorem}

\begin{proof}
Exact conditional independence of the selected interiors
at fixed \(\xi\) gives
\[
 U_\theta(\xi)=\mathbb E_\mu[W_\theta\mid\xi]
     =\prod_z u_{z,\theta}(\xi),\qquad
 \frac{U_\theta}{T_\theta}=e^{L_\theta},\quad
 L_\theta=\sum_{z\in I_\theta}L_{z,\theta}.
\]
Consequently
\(r_\theta=\log\mathbb E_{\nu_0}[e^{L_\theta}\mid C]\).
Conditional Jensen and Proposition~\ref{f15v27:local-signs}
give
\[
 r_\theta\ge\mathbb E_{\nu_0}[L_\theta\mid C]
       \ge-m_\theta\Delta
                  -c_K\sum_{z\in I_\theta}g_{z,0}(C).
\]
For the other sign the relevant conditional posterior is
the \emph{actual} tilted one:
\begin{equation}
 \frac{d\nu_\theta(\cdot\mid C)}{d\nu_0(\cdot\mid C)}
       =\frac{U_\theta}{A_\theta}
       =e^{L_\theta-r_\theta}.
 \label{eq:f15v27:posterior}
\end{equation}
Its conditional relative entropy is nonnegative, so
\[
 r_\theta
   =\mathbb E_{\nu_\theta}[L_\theta\mid C]
       -D(\nu_\theta(\cdot\mid C)\Vert\nu_0(\cdot\mid C))
   \le m_\theta\Delta
                  +c_K\sum_{z\in I_\theta}g_{z,\theta}(C).
\]
The tower identities for \(g\) use the same failure events
and common disintegrations as V26. Both displayed bounds
are nonnegative in magnitude; taking the corresponding
positive parts proves the result. All log ratios are
bounded at a fixed regulator, so there is no integrability
or conditional-entropy subtraction problem.
\end{proof}

The last argument is the elementary conditional Gibbs
variational identity, not a new entropy variational principle.
For a general statement see
\href{https://arxiv.org/abs/1701.07796}
{Anantharam, \emph{A Variational Characterization of R\'enyi
Divergences}}, equation (2). The proof here keeps the
specific conditional normalizers and requires no theorem
about infinite-volume Gibbs states.

\subsection{Actual changed-law norms, without multiplying the chart error}

Use V26's exact \(h_\theta\) and upper bound \(h_\tau\).
For backgrounds \(\theta,\zeta\in[-\tau,\tau]^n\), set
\begin{equation}
 R_{\theta,\zeta}
    =\left[\frac{Z(\theta)Z(2\zeta-\theta)}{Z(\zeta)^2}\right]^{1/2}.
 \label{eq:f15v27:R}
\end{equation}
This is a genuine order-two likelihood cost, not a
replacement normalizer. For the uniform cube estimates
assume \(3\tau<R_{\rm src}\), as eventually on the
selected path. Define
\begin{equation}
 \begin{split}
 B_{\rm src}&=\tfrac12J(\tau)+\tfrac12J(3\tau)+\tfrac14J(2\tau),\\
 b_{\rm src}&=k n_p B_{\rm src},\\
 \mathcal E_n&=2\Delta+c_K\left[
       e^{n b_{\rm src}}q_0^{1/4}
                         +e^{n h_\tau/2}q_0^{1/2}\right].
 \end{split}
 \label{eq:f15v27:uniform-budget}
\end{equation}

\begin{theorem}[Physical-source local log representation in every actual background]
\label{f15v27:source-log}
The sharper estimate in the source's own output law is
\begin{equation}
 \mathbb E_{\rho_\theta}(r_\theta)_\pm
       \le m_\theta v_\theta,\qquad
 v_\theta=\Delta+c_K e^{n h_\theta/2}\sqrt{q_0}.
 \label{eq:f15v27:own}
\end{equation}
For every background \(\zeta\),
\begin{align}
 \mathbb E_{\rho_\zeta}(r_\theta)_-
  &\le m_\theta[\Delta+c_K e^{n h_\zeta/2}\sqrt{q_0}],
                                    \label{eq:f15v27:negative}\\
 \mathbb E_{\rho_\zeta}(r_\theta)_+
  &\le m_\theta[\Delta+c_K R_{\theta,\zeta}
                              e^{n h_\theta/4}q_0^{1/4}].
                                    \label{eq:f15v27:positive}
\end{align}
In particular
\begin{equation}
                   \|r_\theta\|_{L^1(\rho_\zeta)}
                                     \le m_\theta\mathcal E_n.
 \label{eq:f15v27:source-log}
\end{equation}
These finite-regulator bounds hold for every fitting
separated family; their useful growing-size window is
specified below.
\end{theorem}

\begin{proof}
The order-two density projection inequality, already
proved in V16, gives
\[
 \left\|\frac{d\rho_\zeta}{d\rho_0}\right\|_2
          \le e^{n h_\zeta/2},\qquad
 \left\|\frac{d\rho_\zeta}{d\rho_\theta}\right\|_2
          \le R_{\theta,\zeta}.
\]
The norms are respectively under \(\rho_0\) and
\(\rho_\theta\). Since \(0\le g_{z,0}\le1\),
\(\mathbb E_{\rho_0}g_{z,0}^2\le\mu(E_z)\le q_0\).
Cauchy--Schwarz therefore yields
\[
               \mathbb E_{\rho_\zeta}g_{z,0}
                       \le e^{n h_\zeta/2}\sqrt{q_0}.
\]
In the tilted direction, V26's direct full-law estimate
gives
\[
 \mathbb E_{\rho_\theta}g_{z,\theta}
       =\mu_\theta(E_z)\le e^{n h_\theta/2}\sqrt{q_0}.
\]
For another background use \(g_{z,\theta}^2\le g_{z,\theta}\)
and Cauchy--Schwarz once more:
\[
 \mathbb E_{\rho_\zeta}g_{z,\theta}
       \le R_{\theta,\zeta}
                      e^{n h_\theta/4}q_0^{1/4}.
\]
Insert these inequalities into
Theorem~\ref{f15v27:retained-signs}. In the own-law case,
the direct first-moment estimate gives the sharper
positive bound in \eqref{eq:f15v27:own}. The deterministic
term \(m_\theta\Delta\) is never transported by
Cauchy--Schwarz and never acquires a source-volume factor.

Finally, \(Z(\zeta)\ge1\), and the support-disjoint CGF
bound gives
\[
 R_{\theta,\zeta}
       \le\exp\left\{\frac{k n n_p}{2}
                              [J(\tau)+J(3\tau)]\right\}.
\]
Combine this with \(h_\theta,h_\zeta\le h_\tau\).
The resulting positive exponential is \(e^{n b_{\rm src}}\),
which proves \eqref{eq:f15v27:source-log}.
\end{proof}

There is also a comparison of actual, rather than reference,
one-source logarithms. Keep V16's
\(\Phi_\theta=\log A_\theta\) and
\(Q_\theta=\Phi_\theta-\sum_z\Phi_{\theta_z e_z}\).
The local factors cancel exactly, giving
\begin{equation}
 Q_\theta=r_\theta-\sum_{z\in I_\theta}r_{\theta_z e_z},\qquad
             \|Q_\theta\|_{L^1(\rho_\zeta)}
                                  \le2m_\theta\mathcal E_n.
 \label{eq:f15v27:additivity}
\end{equation}
Each single-source error has one active term, even though
the comparison background may have all \(n\) sources.
This does not assert a bound uniform over exponentially
many alternating subset sums: their combinatorial factor
must still be counted. No source derivatives are inferred
from an \(L^1\) estimate without a further argument.

\subsection{A window of order the defect logarithm, not the chart exponent}

\begin{proposition}[Larger separated physical-source families]
\label{f15v27:window}
Fix \(0<\kappa<1\). On V26's selected scale, suppose
\begin{equation}
             4n b_{\rm src}\le(1-\kappa)\log(1/q_0).
 \label{eq:f15v27:window}
\end{equation}
Then uniformly over all source vectors and actual
background vectors in the cube,
\begin{equation}
 \mathcal E_n\le2\Delta+c_K
                       (q_0^{\kappa/4}+q_0^{\kappa/2}).
 \label{eq:f15v27:window-error}
\end{equation}
An explicit sufficient size window, eventually nonempty, is
\begin{equation}
 1\le n\le
 \left\lfloor\frac{(1-\kappa)c_*}{160B_{\rm src}}
                   \frac{\beta}{t^{6\nu+86}}\right\rfloor,
 \qquad t=\log\beta.
 \label{eq:f15v27:explicit-window}
\end{equation}
On this window, both \(\mathcal E_n\) and \(n\mathcal E_n\)
are \(o(\beta^{-A})\) for every fixed \(A>0\).
For the own-law bound alone, the weaker condition
\(n h_\tau\le(1-\kappa)\log(1/q_0)\) suffices for
\(v_\theta\le\Delta+c_Kq_0^{\kappa/2}\).
\end{proposition}

\begin{proof}
The positive-source term in \(\mathcal E_n\) is at most
\(q_0^{\kappa/4}\) under \eqref{eq:f15v27:window}.
Since \(b_{\rm src}\ge h_\tau/4\), the second term
is at most \(q_0^{\kappa/2}\). This proves
\eqref{eq:f15v27:window-error} and the own-law assertion.

On the selected scale, \(k\le5\), \(n_p\le4p^4\),
\(p\le t\), and
\[
 b_{\rm src}\le20B_{\rm src}t^4,\qquad
 \log(1/q_0)\ge\frac{c_*}{2}\frac{\beta}{t^{6\nu+82}}.
\]
These inequalities show that
\eqref{eq:f15v27:explicit-window} implies
\eqref{eq:f15v27:window}. Since
\(c_K\le1+12\beta t^4\), V20's defect rate and V22's
\(\Delta\) rate make \eqref{eq:f15v27:window-error}
smaller than every fixed negative power of \(\beta\).
The explicit window has \(n=O(\beta)\), so multiplication
by \(n\) preserves this conclusion. Its upper endpoint
diverges. As always, selected placements must fit the
separated geometry; the theorem does not alter the regulator.
\end{proof}

V16 paid a global likelihood cost on the spatial chart
error and consequently imposed a logarithmic family window.
The sign-separated estimate avoids that multiplication.
The new size bound is limited by the much smaller
\emph{actual chart-failure} probability. It is still not
an arbitrary-volume source-output theorem. V26's unrestricted
family-size result remains the weaker averaged core-TV
statement. The two conclusions have different scopes.

\subsection{The remaining normalization step is a negative-error exponential moment}

The positive local factors define a finite-regulator proxy
\begin{equation}
 Z_\theta^{\rm loc}=\mathbb E_{\rho_0}T_\theta,
 \qquad d\rho_\theta^{\rm loc}
                  =\frac{T_\theta}{Z_\theta^{\rm loc}}\,d\rho_0.
 \label{eq:f15v27:proxy}
\end{equation}
This paragraph defines a comparison law, not a modification
of the actual model and not an assertion that the two laws
are close. Direct normalization gives
\begin{equation}
 \begin{split}
 \log\frac{Z_\theta^{\rm loc}}{Z(\theta)}
       &=\log\mathbb E_{\rho_\theta}e^{-r_\theta},\\
 D(\rho_\theta\Vert\rho_\theta^{\rm loc})
       &=\mathbb E_{\rho_\theta}r_\theta
                    +\log\mathbb E_{\rho_\theta}e^{-r_\theta}.
 \end{split}
 \label{eq:f15v27:normalization}
\end{equation}
Jensen and \eqref{eq:f15v27:own} do prove the one-sided
pressure statement
\begin{equation}
                \log Z(\theta)-\log Z_\theta^{\rm loc}
                                            \le m_\theta v_\theta.
 \label{eq:f15v27:one-sided-pressure}
\end{equation}
They do not bound an excess of the local proxy pressure.

For clarity, a sufficient missing input would be
\begin{equation}
 \log\mathbb E_{\rho_\theta}e^{(r_\theta)_-}
                                  \le m_\theta\varepsilon_{\rm exp}.
 \label{eq:f15v27:missing-input}
\end{equation}
If it were established, \eqref{eq:f15v27:normalization}
would imply
\[
 -m_\theta v_\theta
    \le\log\frac{Z_\theta^{\rm loc}}{Z(\theta)}
    \le m_\theta\varepsilon_{\rm exp},\qquad
 D(\rho_\theta\Vert\rho_\theta^{\rm loc})
    \le m_\theta(v_\theta+\varepsilon_{\rm exp}).
\]
This is only an implication identifying a sufficient
additional estimate, not a proof of that estimate.
The clipped log bound gives a small \emph{mean} because
\(c_K\) is polynomial on the selected scale. Exponentiating
it asks for a probability estimate that pays \(e^{c_K}\).
The available sparse-failure and source-entropy estimates
do not provide that conclusion. The following example
rules out deriving it from those probability interfaces alone.

\subsection{A positive model satisfying a source CGF but defeating normalized output locality}

Let \(d\ge20\) be an integer. Set
\[
 q=2^{-d},\qquad \epsilon=2^{-8d},\qquad
 L_d=16d\log2,
\]
and take two binary variables \(C,X\) with the strictly
positive joint law
\begin{equation}
 \mu(C=1)=q,\qquad
 \mu(X=1\mid C=0)=1-\epsilon,\qquad
 \mu(X=1\mid C=1)=\epsilon.
 \label{eq:f15v27:example-law}
\end{equation}
Let \(D=\{X=1\}\), let \(\Gamma\) always hold, and let
the chart-reference core law be \(\lambda_C=\delta_1\).
The actual chart conditional is exactly \(\delta_1\),
so the chart log-ratio budget is \(\Delta=0\). The core
is \(X\), its near retained input is \(C\), and there is
no additional exterior. Use the nonnegative observable
and its baseline-centred source
\begin{equation}
 T=L_d1_{\{C=1,X=1\}},\qquad S=T-\mathbb E_\mu T.
 \label{eq:f15v27:example-source}
\end{equation}
The baseline chart failure probability satisfies
\(\mu(D^c)=(1-q)\epsilon+q(1-\epsilon)<2q\).

\begin{proposition}[Uniform signed CGF and vanishing true-law log error]
\label{f15v27:example-CGF}
For every \(|s|\le1/2\),
\begin{equation}
               \log\mathbb E_\mu e^{sS}
                           \le J(s)=-\log(1-s)-s.
 \label{eq:f15v27:example-CGF}
\end{equation}
Independent copies therefore satisfy the arbitrary-support
bound \(\log\mathbb E e^{\sum_zs_zS_z}\le\sum_zJ(s_z)\)
on the same signed cube, and joint chart-failure bounds
\((2q)^{|A|}\). For one copy, with \(\tau=1/8\),
uniformly over \(\theta,\zeta\in[-\tau,\tau]\),
\begin{equation}
                   \|r_\theta\|_{L^1(\rho_\zeta)}
                                   \le4d(\log2)q\longrightarrow0.
 \label{eq:f15v27:example-small}
\end{equation}
\end{proposition}

\begin{proof}
Write \(\eta=q\epsilon=2^{-9d}\), so
\(T=L_dB\) with \(B\) Bernoulli of probability \(\eta\).
For any real \(s\), \(\log(1+x)\le x\) gives
\[
 \log\mathbb E e^{sS}
 =\log[1+\eta(e^{sL_d}-1)]-\eta sL_d
 \le\eta(e^{sL_d}-1-sL_d)
 \le\tfrac12\eta L_d^2s^2e^{|s|L_d}.
\]
For \(|s|\le1/2\),
\[
 \eta L_d^2e^{L_d/2}
     =256d^2(\log2)^2\,2^{-d}
     \le256d^2\,2^{-d}\le\tfrac14.
\]
The final inequality holds at \(d=20\) and its left side
decreases thereafter, since its ratio at consecutive
integers is \((d+1)^2/(2d^2)<1\). Thus the CGF is at
most \(s^2/8\). On \([-1/2,1/2]\),
\(J''(s)=(1-s)^{-2}\ge4/9\), while \(J(0)=J'(0)=0\).
Taylor's integral formula gives \(J(s)\ge2s^2/9\), for
both signs of \(s\), proving \eqref{eq:f15v27:example-CGF}.
Independence proves the stated multi-copy consequences.

The centering factors cancel from \(r_\theta\). Directly,
\[
 r_\theta(0)=0,\qquad
 r_\theta(1)=\log(1-\epsilon+\epsilon e^{\theta L_d})
                                                    -\theta L_d.
\]
This lies between \(0\) and \(-\theta L_d\), in the
appropriate order, so its magnitude is at most
\(2d\log2\). For every allowed background,
\[
 \rho_\zeta(C=1)
  =\frac{q(1-\epsilon+\epsilon e^{\zeta L_d})}
          {1+q\epsilon(e^{\zeta L_d}-1)}\le2q.
\]
For \(\zeta\ge0\), the numerator factor is at most
\(1+2^{-6d}\) and the denominator is at least one.
For \(\zeta<0\), the numerator is at most \(q\) and
the denominator at least \(1-q\epsilon>1/2\).
This proves \eqref{eq:f15v27:example-small}.
\end{proof}

\begin{proposition}[The local-reference normalization nevertheless fails]
\label{f15v27:example-failure}
At \(\theta=\tau\), the difference of proxy and actual
log partition functions diverges, and their retained TV
distance tends to one:
\begin{equation}
 \log\frac{Z_\tau^{\rm loc}}{Z(\tau)}\sim d\log2,
 \qquad d_{\rm TV}(\rho_\tau,\rho_\tau^{\rm loc})\longrightarrow1.
 \label{eq:f15v27:example-failure}
\end{equation}
\end{proposition}

\begin{proof}
Put \(a=e^{\tau L_d}=2^{2d}\). Both partition functions
contain the same baseline centering factor, and
\[
 \frac{Z_\tau^{\rm loc}}{Z(\tau)}
      =\frac{1-q+qa}{1+q\epsilon(a-1)},\qquad
 \rho_\tau^{\rm loc}(C=1)=\frac{qa}{1-q+qa}.
\]
The denominator in the first ratio tends to one, whereas
its numerator divided by \(2^d\) tends to one. The proxy
probability in the second ratio tends to one, while the
actual probability is at most \(2q\) by the preceding
proof. For binary laws TV is the absolute difference of
these probabilities. This proves both assertions.
\end{proof}

The model is finite, positive, and fully normalized. Adding
a deterministic constant to \(T\) leaves its centred-source
CGF and every conclusion unchanged. The obstruction does
not come from a failure to centre sources, an uncontrolled
number of sources, or a hidden correlation among retained
copies. It comes from rare retained data where the true
chart probability is tiny and the chart reference assigns
a very different source normalizer. Those data remain rare
under every declared actual source background but become
dominant under the normalized local proxy.

This model is not a Wilson lattice model and does not
verify its full geometric or analytic hypotheses. It shows
that the normalized-law conclusion cannot follow merely
from the CGF, sparse failure, and chart-kernel interfaces
used above. A Wilson-specific estimate must exclude or
pay for this mechanism. In particular, the input
\eqref{eq:f15v27:missing-input}, or a suitable stronger
spatially local alternative, remains unproved. Repeating
the source-log estimate under more actual backgrounds
does not supply it: the example already has that uniformity.

There is a different reference choice that avoids this
particular normalization problem. Let
\(N=(N_z)_{z\in I_\theta}\) be the combined near output
and put \(b_\theta^{\rm nat}(N)=\mathbb E_\mu[W_\theta\mid N]\).
The tower property gives the exact identities
\begin{equation}
 \mathbb E_{\rho_0}b_\theta^{\rm nat}=Z(\theta),\qquad
 \frac{b_\theta^{\rm nat}(N)}{Z(\theta)}\rho_0(dC)
                  =\rho_\theta(dN)\rho_0(dC\mid N).
 \label{eq:f15v27:natural-reference}
\end{equation}
The second equality means equality of the induced measures
on \(C\), with \(N\) its specified function. This elementary
identity does not establish locality: the required estimate
would now concern the conditional relative entropy
\(D(\rho_\theta\Vert\rho_\theta(dN)\rho_0(dC\mid N))\).
Nor does \(b_\theta^{\rm nat}\) generally factor into a
product of single-region functions. Nevertheless it is an
upstream alternative worth examining, since normalization
is exact and the remaining issue is genuine near/far
retained information. The counterexample to the identity-chart
proxy does not rule out this different reference.

\subsection{Diagnostics, preservation, and next direction}

The exact diagnostic is
\begin{center}\small
\path{scripts/verify_ym_f15_v27_source_log_normalizers.py}.
\end{center}
Six dependent-exterior core/fringe families check 880
sign-separated local log bounds, 48 retained sign bounds,
448 exact conditional-posterior likelihood identities,
72 changed-background moment bounds, and 48 exact mixed
reference-factor cancellations. Four positive examples
check the vanishing true-law log error against the
diverging proxy pressure and separating retained laws.
They include 24 signed centred-source CGF checks and
the constants in the analytic CGF envelope. These are
finite rational probability checks, not a certification
of the Wilson CGF or the selected continuum programme.

Logarithms use rational \(\operatorname{arctanh}\) enclosures
and outward rounding to a \(2^{-120}\) grid, with no
floating-point tolerance. Squared and fourth-power
comparisons avoid approximate roots in likelihood costs.
The script has 9,194 bytes and SHA--256
\begin{center}\scriptsize\ttfamily
A5FEA43553BF654305D03D0D66ED28C367F1102CCEFF2752A37FCE532397C30E.
\end{center}

The predecessor V26 has 599,881 bytes and SHA--256
\begin{center}\scriptsize\ttfamily
38C1DF9D0D28419162CF63C85A861DD6D9A3D8380C07CCF24104EFC398EC69CF.
\end{center}
Removing this terminal section and reversing the title
version recovers it exactly. V27 replaces V26 as the sole
active main note. All 14 protected sources and 69 earlier
certificates are preserved; the next companion checkpoint
is still V30.

\begin{firewall}
V27 proves a sign-separated physical-source log-normalizer
comparison, with superalgebraically small actual-law
errors on an explicit family window of order
\(\beta/(\log\beta)^{6\nu+86}\), and proves a one-sided
proxy-pressure bound. It also demonstrates why these
results do not yield normalized retained-output locality.
The next upstream obligation is a Wilson-specific
negative-error exponential-moment or related retained
spatial estimate. No such bound, overlap contraction,
physical-scale clustering, interacting continuum
construction, or Yang--Mills mass gap is claimed.
\end{firewall}

\section{V28: joint natural source locality through the actual retained posterior}
\label{f15v28:section}

V27 proved a physical-source log approximation but showed
that its identity-chart reference could fail after
normalization. We now use the actual near-output conditional
normalizer, whose normalization is exact. The new estimate
controls the change of the \emph{entire far-retained
conditional law}, averaged under the actual changed near law.
It does not replace that far field by an independent field.

There is an important nonduplication point. V8 already
defined this natural source-local comparison, and
Theorem~\ref{f15v15:finite-source} already proved its
forward relative-entropy bound for one physical source.
Neither the comparison law nor that single-source conclusion
is new here. The present extension is a joint, two-sided
entropy estimate for several separated cores and for
general positive weights depending on their near outputs
and core variables. Such a weight may couple the cores;
conditional independence after tilting is not assumed.
For physical action sources an explicit growing-family
window follows. The argument uses V23's whole-core entropy
and full-compact likelihood bounds, not another integration
of the earlier source-score estimate.

The V27 source and all 84 earlier protected files were
checked unchanged. The last companion checkpoint is V25;
the next is V30. All Wilson inputs retain their specified
selected coupling, completed blocking map, support geometry,
and full compact law. No archived chart or CGF theorem is
independently recertified by this section.

\subsection{The baseline posterior sees only a logarithmic compact cost}

Fix a nonempty separated family of \(n\) cores satisfying
V21--V23. Let \(C\) be the entire retained field,
\(N=(N_z)_{z\in\mathcal S}\) the combined near output,
and \(\mathbf V=(v_{A_z}^{H_{N_z}})_{z\in\mathcal S}\)
the framed core vector. The combined near output may be
recorded without repeated coordinates; only its generated
sigma-algebra matters. In particular all frames are
functions of \(N\). Let \(Y\) be the common exterior
details. At fixed \((C,Y)\), the baseline cores are
independent, as already proved in V21. Write
\(\mathcal P_0\) for their joint marginal with \(C\),
and \(\rho_0\) for the retained marginal.

Keep V23's compact logarithmic cost
\begin{equation}
 \mathfrak h=2(\beta+J_*)2^{20}r^4p^5,
 \qquad \mathfrak m_n=2n\mathfrak h.
 \label{eq:f15v28:compact-cost}
\end{equation}
All core reference measures in this section are compact
product Haar. They are the same at different \(C\) with
fixed \(N\): the frame is then fixed, and its coordinatewise
action preserves Haar. This point is needed for Bayes'
formula below, not just for a formal information identity.

Choose \(\eta=\varepsilon=\sqrt{q_0}\le1/4\), and retain
the baseline forward entropy budget
\begin{equation}
 \mathcal E_+=\Delta-\log(1-\eta)
                    +2\eta\mathfrak h+3\eta\log(1/\eta)
 \label{eq:f15v28:Eplus}
\end{equation}
from \eqref{eq:f15v23:balanced-plus}. The positive-floored
kernel \(\lambda_z^\varepsilon(N_z)\) is used only to
bound this \emph{baseline} information. It will not be
source-tilted and normalized as a proxy output.

\begin{proposition}[Near/core posterior likelihood and information]
\label{f15v28:posterior}
Let \(\Pi_0(dC\mid N,\mathbf V)\) be the actual baseline
retained posterior. Its likelihood relative to the baseline
retained law at fixed near output satisfies
\begin{equation}
 e^{-\mathfrak m_n}
 \le\frac{d\Pi_0(\cdot\mid N,\mathbf V)}
           {d\rho_0(\cdot\mid N)}(C)
 \le e^{\mathfrak m_n}.
 \label{eq:f15v28:posterior-ratio}
\end{equation}
Define the nonnegative posterior information
\begin{equation}
 G(N,\mathbf V)
   =D\bigl(\Pi_0(\cdot\mid N,\mathbf V)
                                  \Vert\rho_0(\cdot\mid N)\bigr).
 \label{eq:f15v28:G}
\end{equation}
Then
\begin{equation}
 0\le G\le\mathfrak m_n,\qquad
 \mathbb E_0G=I_0(C;\mathbf V\mid N)\le n\mathcal E_+.
 \label{eq:f15v28:G-budget}
\end{equation}
The information here is conditional on the \emph{near}
retained output, not on the full retained field as in
\eqref{eq:f15v23:single-information}.
\end{proposition}

\begin{proof}
V23 bounds each full compact core density at every fixed
\((C,Y)\) between \(e^{-\mathfrak h}\) and
\(e^{\mathfrak h}\). The product conditional density
therefore lies between \(e^{-n\mathfrak h}\) and
\(e^{n\mathfrak h}\). Integrating over the actual
\(Y\)-posterior at fixed \(C\) preserves these inequalities.
Denote the resulting density by \(p_0(\mathbf v\mid C)\).
Its conditional average at fixed \(N\), denoted by
\(p_0(\mathbf v\mid N)\), obeys the same inequalities.
Both are strictly positive Haar densities. Bayes' rule gives
\[
 \frac{d\Pi_0(\cdot\mid N,\mathbf v)}
      {d\rho_0(\cdot\mid N)}(C)
       =\frac{p_0(\mathbf v\mid C)}{p_0(\mathbf v\mid N)}.
\]
This proves \eqref{eq:f15v28:posterior-ratio}, including
when the conditional retained measure is described only
as a regular conditional kernel rather than by a Lebesgue
density. All spaces are standard Borel. Integrating its
bounded log likelihood gives \(0\le G\le\mathfrak m_n\).

Put \(q_N=\bigotimes_z\lambda_z^\varepsilon(N_z)\).
V23's projected family theorem gives
\[
 D\bigl(\mathcal P_0(dC,d\mathbf V)
                         \Vert\rho_0(dC)q_N(d\mathbf V)\bigr)
                                                   \le n\mathcal E_+.
\]
In its log likelihood insert \(p_0(\mathbf V\mid N)\).
The conditional KL chain rule yields exactly
\begin{equation}
 \begin{split}
 D(\mathcal P_0\Vert\rho_0q_N)
    &=I_0(C;\mathbf V\mid N)\\
    &\quad+\mathbb E_{\rho_0^N}
                  D\bigl(\mathcal P_0(d\mathbf V\mid N)\Vert q_N\bigr).
 \end{split}
 \label{eq:f15v28:baseline-ledger}
\end{equation}
The first term is \(\mathbb E_0G\) by Bayes' formula.
The second is nonnegative, proving the last bound.
The positive floor and compact density bounds make the
quantities finite. No independence of retained regions
or of their near outputs was used.
\end{proof}

The estimate for \(G\) is logarithmic in the global density
ratio. Its upper bound grows with the compact box cost.
It is not a uniformly small bound for every prescribed
near/core datum. Smallness is the separate averaged
information bound in \eqref{eq:f15v28:G-budget}.

\subsection{Reverse relative entropy without an exponential penalty}

\begin{proposition}[A one-sided likelihood lower bound compares the KL directions]
\label{f15v28:reverse-lemma}
If probability measures \(P,Q\) have
\(dP/dQ\ge e^{-a}\) for \(a\ge0\), then
\begin{equation}
                    D(Q\Vert P)\le(1+a)D(P\Vert Q).
 \label{eq:f15v28:reverse-lemma}
\end{equation}
The assertion includes the extended-value case. In the
finite case no factor \(e^a\) is needed.
\end{proposition}

\begin{proof}
Let \(t=dP/dQ\), and use the nonnegative generators
\(f(t)=t\log t-t+1\), \(g(t)=-\log t+t-1\).
Their \(Q\)-integrals are the forward and reverse KL.
For \(t\ge1\),
\[
 f(t)-g(t)=(t+1)\log t-2(t-1)\ge0.
\]
For example the derivative of the last expression is
\(\log t-1+1/t\), which starts at zero and has
derivative \((t-1)/t^2\ge0\). For \(e^{-a}\le t\le1\),
write \(t=e^{-x}\), \(0\le x\le a\). Direct algebra gives
\[
 (1+x)f(e^{-x})-g(e^{-x})
       =2-[(1+x)^2+1]e^{-x}\ge0,
\]
since \(2e^x\ge2+2x+x^2\). Thus
\(g(t)\le(1+a)f(t)\) throughout the allowed range.
Integration proves the result; \(P=Q\) is included.
\end{proof}

This is a deliberately nonoptimal application of scalar
functional domination. For that general method with
bounded relative information, see
\href{https://arxiv.org/abs/1508.00335}
{Sason--Verd\'u, \emph{\(f\)-divergence Inequalities}},
Theorem 6. The coefficient used here is proved above;
no sharp-constant or Gibbs-mixing theorem is imported.

For later use define
\[
 G_-(N,\mathbf V)
       =D\bigl(\rho_0(\cdot\mid N)
                       \Vert\Pi_0(\cdot\mid N,\mathbf V)\bigr).
\]
The posterior likelihood bound and
Proposition~\ref{f15v28:reverse-lemma} give pointwise
\begin{equation}
                       G_-\le(1+\mathfrak m_n)G.
 \label{eq:f15v28:posterior-reverse}
\end{equation}

\subsection{General positive near/core weights and the unchanged posterior}

Let \(W=W(N,\mathbf V)>0\) be measurable, with
\(0<Z_W=\mathbb E_0W<\infty\) and \(\mathbb E_0W^2<\infty\).
It need not factor over the cores, be differentiable, or
have a uniformly bounded logarithmic oscillation. Define
\begin{equation}
 d\mu_W=\frac{W}{Z_W}\,d\mu,
 \qquad \mathcal R_2(W)=\frac{(\mathbb E_0W^2)^{1/2}}{Z_W},
 \qquad \mathcal B(W)=n\mathcal R_2(W)
                                      \sqrt{2\mathfrak h\mathcal E_+}.
 \label{eq:f15v28:weight-budget}
\end{equation}
Let \(\mathcal P_W\) be the actual joint marginal of
\((C,\mathbf V)\), and \(\rho_W\) its retained marginal.
Define the comparison joint law
\begin{equation}
 \mathcal Q_W(dC,dN,d\mathbf V)
    =\mathcal P_W(dN,d\mathbf V)\rho_0(dC\mid N).
 \label{eq:f15v28:joint-comparison}
\end{equation}
Here \(N\) is the specified function of \(C\); the
notation denotes the induced consistent joint law.
It preserves the \emph{actual changed} near/core marginal,
including all its correlations. Only the conditional
retained law given those variables is replaced.

\begin{theorem}[Two-sided posterior comparison for positive near/core weights]
\label{f15v28:joint-theorem}
The conditional retained posterior is unchanged by the
weight:
\begin{equation}
       \mathcal P_W(dC\mid N,\mathbf V)
                                  =\Pi_0(dC\mid N,\mathbf V).
 \label{eq:f15v28:unchanged-posterior}
\end{equation}
Consequently
\begin{align}
 D(\mathcal P_W\Vert\mathcal Q_W)
      &=\mathbb E_WG\le\mathcal B(W),
                                    \label{eq:f15v28:joint-forward}\\
 D(\mathcal Q_W\Vert\mathcal P_W)
      &=\mathbb E_WG_-
                   \le(1+\mathfrak m_n)\mathcal B(W).
                                    \label{eq:f15v28:joint-reverse}
\end{align}
These conclusions do not assume conditional independence
of the cores under \(\mu_W\).
\end{theorem}

\begin{proof}
The weight is constant in \(C\) after \((N,\mathbf V)\)
has been fixed. It cancels from the conditional retained
normalizer, proving \eqref{eq:f15v28:unchanged-posterior}.
The two joint laws have exactly the same \((N,\mathbf V)\)
marginal, so their KL chain rules give the displayed
equalities. Their log likelihood is the baseline posterior
log ratio, bounded by \(\mathfrak m_n\); hence these
entropies are finite even if \(W\) is not bounded.

The nonnegative bound on \(G\) is crucial. It implies
\(\mathbb E_0G^2\le\mathfrak m_n\mathbb E_0G
\le2n^2\mathfrak h\mathcal E_+\). Cauchy--Schwarz under
the baseline law therefore gives
\[
 \mathbb E_WG
   =\mathbb E_0\left[\frac W{Z_W}G\right]
   \le\mathcal R_2(W)(\mathbb E_0G^2)^{1/2}
   \le n\mathcal R_2(W)\sqrt{2\mathfrak h\mathcal E_+}.
\]
Finally integrate \eqref{eq:f15v28:posterior-reverse} under
the actual changed near/core marginal. The same marginal
appears in both KL directions; no new artificial source
posterior is introduced.
\end{proof}

The measurability restriction cannot be removed. If \(N\)
is constant and the baseline far bit \(F\) and core bit
\(V\) are independent and fair, then \(G=0\). A weight
\(W=k^F\), \(k>1\), changes the far marginal and has
finite second moment, but is not a function of \((N,V)\).
Its retained law differs from the natural comparison by
TV \((k-1)/(2(k+1))\). Thus this theorem is not an
inequality for arbitrary global positive test densities.

\subsection{An exactly normalized retained law and its near/far information}

Define the actual conditional normalizers
\begin{equation}
 A_W(C)=\mathbb E_0[W\mid C],\qquad
 b_W(N)=\mathbb E_0[W\mid N],\qquad
 d\widehat\rho_W(C)=\frac{b_W(N)}{Z_W}\,d\rho_0(C).
 \label{eq:f15v28:natural-output}
\end{equation}
The tower property makes \(\widehat\rho_W\) a probability
with exactly the same near marginal as \(\rho_W\):
\[
          \widehat\rho_W(dC)=\rho_W(dN)\rho_0(dC\mid N).
\]
It is precisely the retained marginal of \(\mathcal Q_W\).
In particular this is not V27's normalized identity-chart
proxy, whose normalizer need not equal \(Z_W\).

\begin{theorem}[Joint natural retained-source locality in both KL directions]
\label{f15v28:output-theorem}
One has
\begin{align}
 D(\rho_W\Vert\widehat\rho_W)
     &=\int D\bigl(\rho_W(dC\mid N)\Vert\rho_0(dC\mid N)\bigr)
                                                   \,d\rho_W^N(N)
       \le\mathcal B(W),
                                     \label{eq:f15v28:output-forward}\\
 D(\widehat\rho_W\Vert\rho_W)
     &=\int D\bigl(\rho_0(dC\mid N)\Vert\rho_W(dC\mid N)\bigr)
                                                   \,d\rho_W^N(N)
       \le(1+\mathfrak m_n)\mathcal B(W).
                                     \label{eq:f15v28:output-reverse}
\end{align}
Moreover the exact information ledger is
\begin{equation}
 \mathbb E_WG
       =I_W(C;\mathbf V\mid N)
                            +D(\rho_W\Vert\widehat\rho_W).
 \label{eq:f15v28:changed-ledger}
\end{equation}
Thus the changed-law conditional mutual information is
also at most \(\mathcal B(W)\), and
\begin{equation}
 d_{\rm TV}(\rho_W,\widehat\rho_W)
                         \le\min\{1,\sqrt{\mathcal B(W)/2}\}.
 \label{eq:f15v28:TV}
\end{equation}
\end{theorem}

\begin{proof}
Project \(\mathcal P_W,\mathcal Q_W\) onto \(C\).
Relative entropy cannot increase under this projection,
in either direction, proving the bounds. Equality of
their near marginals gives the two conditional KL formulas.
For \eqref{eq:f15v28:changed-ledger}, insert the actual
\(\rho_W(dC\mid N)\) between the two retained kernels
in the forward joint log ratio. Its first contribution is
the conditional mutual information under \(\mathcal P_W\),
and its second is the forward output KL. This is an
exact finite conditional chain rule, with no subtraction
of differential entropies. Pinsker's inequality, in the
half-\(L^1\) convention for TV, gives \eqref{eq:f15v28:TV}.
\end{proof}

For completeness, the actual natural source-log error
\(r_W^{\rm nat}=\log A_W-\log b_W\) obeys
\begin{equation}
 \begin{split}
 \|r_W^{\rm nat}\|_{L^1(\rho_W)}
      &\le\mathcal B(W)+\sqrt{2\mathcal B(W)},\\
 \|r_W^{\rm nat}\|_{L^1(\widehat\rho_W)}
      &\le(1+\mathfrak m_n)\mathcal B(W)+\sqrt{2\mathcal B(W)}.
 \end{split}
 \label{eq:f15v28:natural-log}
\end{equation}
Indeed for equivalent probabilities \(P,Q\), the negative
part of \(\log(dP/dQ)\) under \(P\) has integral at most
\(d_{\rm TV}(P,Q)\), since \(x\log(1/x)\le1-x\)
for \(0<x\le1\). Hence its absolute integral is at most
\(D(P\Vert Q)+2d_{\rm TV}(P,Q)\). Apply this in both
directions. In the present setting the log ratio is also
bounded pointwise by \(\mathfrak m_n\): the actual
\(\rho_W(dC\mid N)\) is a mixture of the posteriors in
\eqref{eq:f15v28:posterior-ratio}.

If \(F\) denotes the remaining retained coordinates,
the conditional divergences above concern the full law
of \(F\) given \(N\), not just finitely many far tests.
Nevertheless \(\rho_0(dF\mid N)\) may itself depend
strongly on \(N\). Neither theorem asserts weak
unconditional near/far correlations in the baseline field.

\subsection{Fixed-strength physical sources and an explicit growing-family window}

Take exactly V26's separated physical sources with
\(\theta\in[-\tau,\tau]^n\), \(\tau=1/8\), and
\(W_\theta=e^{\sum_z\theta_z S_z}\). Their canonical
carriers lie in the selected cores and their retained
parameters in the selected near outputs, so they satisfy
the measurability condition of the preceding theorems.
The fixed baseline centering is unchanged. The arbitrary-support
CGF gives, with only \(2\tau<R_{\rm src}\) needed here,
\begin{equation}
 \mathcal R_2(W_\theta)=e^{n h_\theta/2}
                 \le e^{n h_\tau/2},\qquad
 h_\tau=k n_pJ(1/4),\qquad
 \mathcal B_\theta
       =n e^{n h_\theta/2}\sqrt{2\mathfrak h\mathcal E_+}.
 \label{eq:f15v28:physical-budget}
\end{equation}
Here \(h_\theta\) is V26's exact intensive order-two budget,
not a new estimate for a conditioned Wilson law.

\begin{proposition}[Superalgebraic joint natural-output locality]
\label{f15v28:rate}
Use the selected scale \eqref{eq:f15v20:scale},
\(2\le p\le t=\log\beta\), and fix \(0<\kappa<1\).
Let \(c_\Delta=2^{-51}C_0^{-1}\). The sufficient size window
\begin{equation}
 1\le n\le
  \left\lfloor\frac{(1-\kappa)c_\Delta}{20J(1/4)}
                                          \frac{t^6}{p^4}\right\rfloor
 \label{eq:f15v28:window}
\end{equation}
is eventually nonempty. Uniformly over all permitted
placements and physical source vectors, every bound in
\eqref{eq:f15v28:joint-forward}--\eqref{eq:f15v28:natural-log}
tends to zero faster than every fixed negative power of
\(\beta\). In particular both normalized retained-law KL
directions and their TV distance tend to zero in this sense.
\end{proposition}

\begin{proof}
Since \(k\le5\) and \(n_p\le4p^4\), the proposed window
implies
\begin{equation}
                       n h_\tau\le(1-\kappa)c_\Delta t^6.
 \label{eq:f15v28:margin}
\end{equation}
Put \(C_H=2^{25}(1+J_*)\). V23 gives
\(\mathfrak h\le C_H e^t t^{4\nu+61}\). Also, uniformly
on this scale, eventually
\begin{equation}
                    \mathcal E_+
             \le2e^t t^{8\nu+114}e^{-c_\Delta t^6}.
 \label{eq:f15v28:Eplus-rate}
\end{equation}
To verify the last estimate explicitly, V22 bounds its
\(\Delta\) term by \(2^{-5}e^t t^{8\nu+114}e^{-c_\Delta t^6}\).
The other terms in \eqref{eq:f15v28:Eplus} are at most
\(2\eta+2\eta\mathfrak h+6\sqrt\eta\), using
\(-\log(1-\eta)\le2\eta\) and
\(\eta\log(1/\eta)\le2\sqrt\eta\).
V20's pre-colouring bound
\(\log(1/q_0)\ge c_*e^t/(2t^{6\nu+82})\) makes their
sum at most \(e^{-c_\Delta t^6}\) eventually, even after
the displayed polynomial and \(e^t\) costs. This proves
\eqref{eq:f15v28:Eplus-rate}.

Substitution in \eqref{eq:f15v28:physical-budget}, using
\eqref{eq:f15v28:margin}, gives the explicit majorant
\begin{equation}
 \mathcal B_\theta
   \le2\sqrt{C_H}\,n\beta t^{6\nu+175/2}
                                   e^{-\kappa c_\Delta t^6/2}.
 \label{eq:f15v28:physical-rate}
\end{equation}
On \eqref{eq:f15v28:window}, \(n=O(t^6)\), and
\(1+\mathfrak m_n\le1+2nC_H\beta t^{4\nu+61}\).
These factors, and the square roots in the TV and log
bounds, do not destroy superalgebraic decay. Since
\(t^6/p^4\ge t^2\), the upper endpoint of the window
diverges. The inherited chart threshold, embedding and
source-radius conditions hold eventually as before.
\end{proof}

For fixed \(p\) this sufficient family size grows as
\(t^6\); at the largest admitted \(p\) it grows as
\(t^2\). It is \emph{not} V27's much larger
\(\beta/t^{6\nu+86}\) source-log window. Here the
likelihood cost multiplies the baseline information
budget, which includes the chart error \(\Delta\).
Keeping that cost is essential. Likewise the finite
bound \eqref{eq:f15v28:weight-budget} for arbitrary
families is not an all-volume vanishing theorem.

The union of these near windows still has a vanishing
fraction of the retained volume: \(n(rp)^4\) is bounded
by a fixed polynomial in \(t\), while the selected path
has \(M\ge e^\beta\). Thus the far-retained comparison
does not become empty because the near regions cover
the torus. This observation is geometric, not a statement
about a physical continuum correlation length.

\subsection{What has changed, and what is still missing}

The normalization obstruction in V27 has not been disproved
or ignored. Its chart proxy and the present natural reference
are different measures. In V27's two-bit example the near
output is the whole retained bit; the natural retained
comparison is therefore the actual law exactly, whereas
the chart proxy fails. Here, when genuine far coordinates
are present, their conditional change is bounded using
the full-compact posterior and the baseline information
budget. Normalization is never inferred from an expected
log approximation.

Nor has the natural factor been replaced by a product:
\(b_{W_\theta}(N)\) may contain joint dependence among all
near regions. The theorem proves locality relative to
their \emph{combined} near output, not additivity or a
summable interaction expansion on that union. The actual
changed near/core correlations are preserved in
\(\mathcal Q_W\), and can be strong.

The results concern a restricted source channel with an
explicit second-moment cost. They do not prove an all-density
functional inequality, a uniform bound over every prescribed
boundary condition, or a contraction for overlapping block
updates. Even zero \(I_0(C;\mathbf V\mid N)\) by itself
would allow an arbitrarily correlated baseline retained
field: cores independent of that field supply an immediate
example. Therefore neither a retained-field mass scale nor
a Hamiltonian spectral gap follows from this source theorem.

An upstream next task is to control how the natural joint
near normalizer decomposes across separated or overlapping
regions, or how its retained conditional estimates propagate
under further blocking. Such a step must keep the actual
near law and its normalizers; it cannot revert to the
unjustified chart-proxy normalization excluded in V27.

\subsection{Exact diagnostics and predecessor preservation}

The diagnostic is
\begin{center}\small
\path{scripts/verify_ym_f15_v28_natural_output_entropy.py}.
\end{center}
It checks 425 scalar forward/reverse KL domination cases.
Six positive families have correlated near, far and common
exterior variables; mixing over the exterior produces the
actual joint core law. Their baseline information ledgers
are checked before any tilt. There are 24 positive source
weights, including eight explicitly nonproduct core weights
and near-only weights whose retained comparison is exact.
The tests check 896 exact posterior identities, 192 exact
natural normalization identities, 24 changed information
ledgers, 24 two-sided joint/output entropy comparisons,
and 48 absolute source-log bounds. Eight excluded far-source
examples verify that the near/core measurability restriction
cannot be dropped.

All probabilities are rational. Logarithms use rational
\(\operatorname{arctanh}\) enclosures with outward rounding
to a \(2^{-120}\) grid. Squared comparisons avoid approximate
second-moment square roots. No floating-point tolerance is
used. These finite checks validate the new probability
interfaces, not the inherited Wilson geometry, selected
CGF, or continuum construction. The script has 9,570 bytes
and SHA--256
\begin{center}\scriptsize\ttfamily
66A3725136F900E0B58D1F9E18A9C37B1ACCC5866796BBC6E670777B2EEADAB9.
\end{center}

The predecessor V27 has 626,035 bytes and SHA--256
\begin{center}\scriptsize\ttfamily
68E06A70F1992CBB3E8B68893CDAE485DD12DE597828D2CB638C851E4AF44B00.
\end{center}
Removing this terminal section and reversing the title
version recovers it exactly. V28 replaces V27 as the sole
active main note. All 14 protected sources and 70 earlier
certificates are preserved. The next companion checkpoint
remains V30.

\begin{firewall}
V28 extends the existing single-source forward locality
result to a joint two-sided natural retained-output entropy
comparison, including nonproduct positive near/core weights
with an explicit second-moment cost. Fixed-strength physical
sources have superalgebraically small errors in the declared
\((\log\beta)^6/p^4\) family window. No additive natural
near action, all-volume contraction, baseline retained
clustering, interacting continuum limit, or Yang--Mills
mass gap is claimed.
\end{firewall}

\section{V29: additive natural source factors under a total-strength budget}
\label{f15v29:section}

V28 controls a natural source normalizer depending on the
\emph{joint} near output. It leaves open whether separately
defined near factors can be multiplied and then normalized.
Here we obtain that comparison relative to the original
baseline law, using the actual single-near conditional core
kernels. A total signed source budget replaces V28's
fixed strength at every site. This distinction is essential:
the new result is uniform per core over all fitting finite
families, but only on an \(\ell^1\) source ball.

This is not the first normalized product-source comparison
in these notes. V17 treated a negative source cube relative
to the baseline and a signed cube relative to a
family-dependent upper-corner source law. V18 used local
chart factors in that latter setting. V24--V25 treated
compensated or globally bounded-oscillation probes.
The present factors are actual full-compact conditional
expectations under one fixed, untilted baseline; they depend
on their own source parameter only. The proof supplies both
joint and retained-output entropy directions, a pressure
comparison, and an absolute source-log bound.

The V28 predecessor and all 85 protected sources and
certificates were checked unchanged. The supplied SM
source-selection V31, native stationarity V34 and
exact-action Einstein bridge V24 remain in the V25
companion audit; they are mathematical data, not instructions.
No new Wilson estimate is imported from them here.
The next companion checkpoint is V30. All selected-law,
embedding, completed-blocking, CGF and chart hypotheses
below retain their earlier scope. This section proves the
new probability interfaces, not a fresh certification of
every archived Wilson input.

\subsection{Use actual single-near marginals before applying a source}

Fix the full compact baseline \(\mu\) and a nonempty family
\(\mathcal S\) of \(n\) cores with the geometry of V21--V26.
In particular different centres have separation at least
\(8\ell\) in some periodic coordinate. Let \(C\) be the
entire retained field, \(Y\) the common exterior details,
\(\xi=(C,Y)\), and \(\nu_0\) their actual baseline law.
The framed cores \(V_z\) are conditionally independent at
fixed \(\xi\), with full-compact kernels \(\kappa_z^\xi\).
The near output \(N_z\) is a fixed function of \(C\).
Keep \(\mathfrak h\) and \(\mathcal E_+\) from
\eqref{eq:f15v28:compact-cost} and \eqref{eq:f15v28:Eplus}.
Thus every conditional core Haar density lies between
\(e^{-\mathfrak h}\) and \(e^{\mathfrak h}\), and
\begin{equation}
 \int D\bigl(\kappa_z^\xi\Vert
                   \lambda_z^\varepsilon(N_z)\bigr)\,d\nu_0(\xi)
                                                \le\mathcal E_+.
 \label{eq:f15v29:inherited}
\end{equation}
These are the baseline V23 bounds after the common-exterior
conditioning used in its family theorem.

Define the actual single-near core kernel and two joint laws
on the same coordinates by
\begin{align}
 q_z(N_z,dv)&=\mu(V_z\in dv\mid N_z),
                                       \label{eq:f15v29:q}\\
 P(d\xi,d\mathbf v)&=\nu_0(d\xi)
                           \prod_{z\in\mathcal S}\kappa_z^\xi(dv_z),
                                       \label{eq:f15v29:P}\\
 Q(d\xi,d\mathbf v)&=\nu_0(d\xi)
                           \prod_{z\in\mathcal S}q_z(N_z,dv_z).
                                       \label{eq:f15v29:Q}
\end{align}
Here \(P\) is the actual projected fine law. The frame at
site \(z\) depends only on \(N_z\), so mixing the full
\(\kappa_z^\xi\) densities at fixed \(N_z\) uses the same
compact product Haar measure. In particular \(q_z\) also
has density between \(e^{-\mathfrak h}\) and
\(e^{\mathfrak h}\). Conditional kernels and identities
are understood almost everywhere on these standard Borel
spaces, with positive versions supplied by their Haar
densities. No identity-chart restriction or artificial
floor is part of \(q_z\).

\begin{proposition}[Single-source marginal preservation and information projection]
\label{f15v29:projection}
The laws \(P,Q\) have the same \(\xi\) marginal and, for
every \(z\), the same \((N_z,V_z)\) marginal. Moreover
\begin{equation}
 \begin{split}
 \int D(\kappa_z^\xi\Vert\lambda_z^\varepsilon(N_z))\,d\nu_0
   &=e_z+
     \int D(q_z(N_z)\Vert\lambda_z^\varepsilon(N_z))
                                                   \,d\rho_0^{N_z},\\
 e_z&:=\int D(\kappa_z^\xi\Vert q_z(N_z))\,d\nu_0
                                                      \le\mathcal E_+.
 \end{split}
 \label{eq:f15v29:projection}
\end{equation}
Here \(\rho_0=C_*\mu\). In particular every source
\(S_z=s_z(N_z,V_z)\) has its actual baseline distribution
under \emph{both} \(P\) and \(Q\).
\end{proposition}

\begin{proof}
Integrating the other probability kernels in
\eqref{eq:f15v29:Q} leaves
\(\rho_0^{N_z}(dN_z)q_z(N_z,dv_z)\), which is the actual
single-near/core marginal by definition. The same integration
in \(P\), followed by conditional averaging at fixed
\(N_z\), gives exactly that marginal. Insert \(q_z\) in
the log ratio in the left side of
\eqref{eq:f15v29:projection}. Its second term depends only
on \((N_z,V_z)\); the just-proved marginal identity turns
that term into the displayed \(q_z\)-relative entropy.
Nonnegativity and \eqref{eq:f15v29:inherited} prove the
bound on \(e_z\). The compact bounds and the positive
floor in the inherited reference make these integrals
finite; no formal subtraction of infinite entropies occurs.
\end{proof}

The comparison \(Q\) preserves all baseline retained
correlations. It is a product only in the core kernels at
fixed \(\xi\). Neither the \(N_z\)'s nor the full retained
field are made independent. The individual kernels \(q_z\)
are defined from the fixed baseline and the chosen local
geometry, not from a choice of remote source parameters.

\subsection{A local log-square estimate avoids a global volume penalty}

\begin{proposition}[Logarithmic rather than exponential likelihood cost]
\label{f15v29:log-square}
If equivalent probability measures \(U,V\) satisfy
\(\lvert\log(dU/dV)\rvert\le a\), then
\begin{equation}
 \mathbb E_U\left|\log\frac{dU}{dV}\right|^2
 +\mathbb E_V\left|\log\frac{dU}{dV}\right|^2
                         \le(a+2)^2D(U\Vert V).
 \label{eq:f15v29:log-square}
\end{equation}
Consequently, with
\begin{equation}
 \ell_z(\xi,v_z)=\log\frac{d\kappa_z^\xi}{dq_z(N_z)}(v_z),
 \qquad L=\sum_z\ell_z=\log\frac{dP}{dQ},
 \qquad a=2\mathfrak h,
 \label{eq:f15v29:L}
\end{equation}
one has, separately under \(R=P\) and \(R=Q\),
\begin{equation}
 \|L\|_{L^2(R)}
    \le(a+2)\sum_z\sqrt{e_z}
    \le n(2+2\mathfrak h)\sqrt{\mathcal E_+}.
 \label{eq:f15v29:joint-L2}
\end{equation}
\end{proposition}

\begin{proof}
For \(t>0\), \(|\log t|\le a\), the scalar inequality is
\begin{equation}
           (1+t)(\log t)^2\le(a+2)(t-1)\log t.
 \label{eq:f15v29:scalar}
\end{equation}
For \(t\ne1\), divide by the positive quantity
\((t-1)\log t\) and put \(x=|\log t|\). The resulting
left side equals
\[
 x\coth(x/2)=x+\frac{2x}{e^x-1}\le x+2\le a+2.
\]
At \(t=1\) both sides of \eqref{eq:f15v29:scalar} are
zero. Integrating it with \(t=dU/dV\) gives
\[
 \mathbb E_U(\log t)^2+\mathbb E_V(\log t)^2
     \le(a+2)\{D(U\Vert V)+D(V\Vert U)\}.
\]
Proposition~\ref{f15v28:reverse-lemma} bounds the second
entropy by \((1+a)D(U\Vert V)\), proving
\eqref{eq:f15v29:log-square}.

Apply that result at each \(\xi\) to
\(\kappa_z^\xi,q_z(N_z)\), whose log ratio lies in
\([-2\mathfrak h,2\mathfrak h]\). Average over \(\nu_0\).
The \((\xi,V_z)\) marginals of \(P,Q\) give
\[
       \mathbb E_P\ell_z^2+\mathbb E_Q\ell_z^2
                                            \le(a+2)^2 e_z.
\]
Finally use the triangle inequality in each \(L^2\) space
for \(L=\sum_z\ell_z\). This last step uses no
independence or cancellation. Applying the scalar estimate
to the local ratios first is important: the global cap
\(|L|\le2n\mathfrak h\) would incur an avoidable extra
family-size cost.
\end{proof}

The functional-domination background was cited in V28.
The explicit constants and every scalar inequality used
here are proved above; no functional inequality for the
retained Gibbs field is being invoked.

\subsection{A fixed total source strength controls both baseline moments}

Keep the physical sources \eqref{eq:f15v26:sources}, with
their original centering \(\mathbb E_\mu S_z=0\). The
one-source consequence of \eqref{f15v1:source-input} is
\begin{equation}
 \log\mathbb E_\mu e^{uS_z}
     \le k n_pJ(u)\le k n_pJ(|u|),\qquad |u|<R_{\rm src},
 \quad J(u)=-\log(1-u)-u.
 \label{eq:f15v29:scalar-CGF}
\end{equation}
Use signed coefficients satisfying
\begin{equation}
 s=\|\theta\|_1=\sum_z|\theta_z|\le\tau=\tfrac18,
 \qquad 2\tau<R_{\rm src},\qquad
 W_\theta=\exp\left(\sum_z\theta_zS_z\right).
 \label{eq:f15v29:budget}
\end{equation}
The requirement concerns the \emph{sum} of the amplitudes,
not merely their maximum. At each finite regulator the
sources are bounded and all laws below are equivalent.

\begin{proposition}[Moment control under the true and product-reference laws]
\label{f15v29:moments}
For \(R=P,Q\), put \(Z_R=\mathbb E_R W_\theta\) and
\(R_\theta=W_\theta R/Z_R\). Then
\begin{equation}
 Z_R\ge1,\qquad
 \frac{(\mathbb E_R W_\theta^2)^{1/2}}{Z_R}
   \le\mathcal R(s):=\exp\{k n_pJ(2s)/2\}.
 \label{eq:f15v29:R2}
\end{equation}
There is no factor depending on \(n\) in this upper bound.
\end{proposition}

\begin{proof}
If \(s=0\), then \(W_\theta=1\) and the assertion is
immediate. Otherwise omit zero coefficients and let
\(\alpha_z=|\theta_z|/s\), so \(\sum_z\alpha_z=1\).
Weighted H\"older, for \(j=1,2\), gives
\begin{align*}
 \mathbb E_R W_\theta^j
   &\le\prod_z
      \left(\mathbb E_R
           e^{j\,\operatorname{sgn}(\theta_z)sS_z}\right)^{\alpha_z}\\
   &=\prod_z
      \left(\mathbb E_\mu
           e^{j\,\operatorname{sgn}(\theta_z)sS_z}\right)^{\alpha_z}
     \le e^{k n_pJ(js)}.
\end{align*}
The middle equality follows from the exact marginal
preservation in Proposition~\ref{f15v29:projection}.
It is valid for both baselines despite their different
joint source laws. Those same marginals give
\(\mathbb E_R S_z=0\), hence \(Z_R\ge1\) by Jensen.
The \(j=2\) bound proves \eqref{eq:f15v29:R2}.
No joint source independence, conditioned Wilson CGF,
or source-tilted convexity assertion has been used.
\end{proof}

This is the reason for choosing the actual \(q_z\)'s.
They transfer the available one-source CGF to the
comparison law exactly. That transfer was not available
for V27's chart reference. Nor would a full fixed-amplitude
cube give this family-independent moment cost: it would
have \(s\) proportional to \(n\), eventually outside
the available scalar CGF radius.

\subsection{Common tilts control normalization as well as joint entropy}

Define the per-core and total error budgets
\begin{equation}
 \epsilon_{\rm src}(s)
   =(2+2\mathfrak h)\mathcal R(s)\sqrt{\mathcal E_+},
 \qquad \delta_\theta=n\epsilon_{\rm src}(s),
 \qquad \Lambda_\theta=\log(Z_P/Z_Q).
 \label{eq:f15v29:error}
\end{equation}
The actual \(P_\theta\) is the projection of the physical
source law \(\mu_\theta\), without any compensating core
normalizers. The reference \(Q_\theta\) uses the same
positive weight, normalized by its own expectation.

\begin{theorem}[Pressure and two-sided joint comparison]
\label{f15v29:joint}
Under the hypotheses above,
\begin{align}
 \mathbb E_{P_\theta}|L|,\ \mathbb E_{Q_\theta}|L|
                                      &\le\delta_\theta,
                                       \label{eq:f15v29:changed-L}\\
 \mathbb E_{Q_\theta}L\le\Lambda_\theta
                       \le\mathbb E_{P_\theta}L,
 \qquad |\Lambda_\theta|&\le\delta_\theta,
                                       \label{eq:f15v29:pressure}\\
 D(P_\theta\Vert Q_\theta),\ D(Q_\theta\Vert P_\theta)
                                      &\le2\delta_\theta.
                                       \label{eq:f15v29:joint-KL}
\end{align}
\end{theorem}

\begin{proof}
For either baseline \(R\), Cauchy--Schwarz gives
\[
 \mathbb E_{R_\theta}|L|
   \le\frac{(\mathbb E_R W_\theta^2)^{1/2}}{Z_R}
           \|L\|_{L^2(R)}\le\delta_\theta.
\]
The identical weight cancels from the tilted likelihood:
\(\log(dP_\theta/dQ_\theta)=L-\Lambda_\theta\).
Consequently the exact entropy identities are
\begin{equation}
 \begin{split}
 D(P_\theta\Vert Q_\theta)
       &=\mathbb E_{P_\theta}L-\Lambda_\theta,\\
 D(Q_\theta\Vert P_\theta)
       &=-\mathbb E_{Q_\theta}L+\Lambda_\theta.
 \end{split}
 \label{eq:f15v29:joint-identities}
\end{equation}
Their nonnegativity brackets \(\Lambda_\theta\) between
the two expectations, proving its absolute bound.
The same identities then give both \(2\delta_\theta\)
entropy bounds. All quantities are finite by the compact
likelihood cap, so the bracketing does not presume the
desired normalizer comparison.
\end{proof}

\subsection{The normalized additive near-source law and its log error}

The individual natural near factors are
\begin{equation}
 b_{z,u}(N_z)=\int e^{u s_z(N_z,v)}q_z(N_z,dv)
            =\mathbb E_\mu[e^{uS_z}\mid N_z],
 \qquad T_\theta(C)=\prod_z b_{z,\theta_z}(N_z).
 \label{eq:f15v29:factors}
\end{equation}
Integrating the core kernels in \(Q_\theta\), then \(Y\),
gives its retained marginal exactly:
\begin{equation}
 d\rho_\theta^{\rm prod}(C)
      =\frac{T_\theta(C)}{Z_Q}\,d\rho_0(C),
 \qquad Z_Q=\mathbb E_{\rho_0}T_\theta.
 \label{eq:f15v29:product-output}
\end{equation}
This expectation is generally \emph{not}
\(\prod_z\mathbb E_\mu e^{\theta_zS_z}\). The retained
law in the expectation remains the original correlated
law. Additivity means the source contribution
\(\log T_\theta=\sum_z\log b_{z,\theta_z}(N_z)\);
it does not mean that the source-free retained action
or its probability law is a product.

Let \(\rho_\theta=C_*\mu_\theta\), and define the actual
conditional normalizer and unnormalized source-log error
\begin{equation}
 A_\theta(C)=\mathbb E_\mu[W_\theta\mid C],\qquad
 r_\theta(C)=\log A_\theta(C)-\log T_\theta(C).
 \label{eq:f15v29:remainder}
\end{equation}

\begin{theorem}[Two-sided normalized output and linear absolute log error]
\label{f15v29:output}
One has
\begin{align}
 D(\rho_\theta\Vert\rho_\theta^{\rm prod}),\quad
 D(\rho_\theta^{\rm prod}\Vert\rho_\theta)
                                               &\le2\delta_\theta,
                                       \label{eq:f15v29:output-KL}\\
 d_{\rm TV}(\rho_\theta,\rho_\theta^{\rm prod})
                            &\le\min\{1,\sqrt{\delta_\theta}\},
                                       \label{eq:f15v29:output-TV}\\
 \|r_\theta\|_{L^1(\rho_\theta)},\quad
 \|r_\theta\|_{L^1(\rho_\theta^{\rm prod})}
                                               &\le3\delta_\theta.
                                       \label{eq:f15v29:log-L1}
\end{align}
In particular the last bounds incur no further square
root of \(\delta_\theta\).
\end{theorem}

\begin{proof}
Relative entropy decreases under projection onto \(C\),
so \eqref{eq:f15v29:joint-KL} proves both output bounds.
Pinsker in the half-\(L^1\) convention gives the TV bound.
For the sharper absolute-log conclusion, use the fact
that the \emph{baseline} \(C\) marginals of \(P,Q\)
coincide. At fixed \(C\) their conditional likelihood
is therefore \(e^L\); after the common tilt it is
\begin{equation}
 \log\frac{dP_\theta(\cdot\mid C)}
               {dQ_\theta(\cdot\mid C)}
                             =L-r_\theta(C).
 \label{eq:f15v29:conditional-L}
\end{equation}
Conditional nonnegativity of KL in both directions yields
\begin{equation}
 \mathbb E_{Q_\theta}[L\mid C]
          \le r_\theta(C)
          \le\mathbb E_{P_\theta}[L\mid C].
 \label{eq:f15v29:conditional-bracket}
\end{equation}
It follows that
\(\mathbb E_{\rho_\theta}(r_\theta)_+\le\delta_\theta\)
and
\(\mathbb E_{\rho_\theta^{\rm prod}}(r_\theta)_-
\le\delta_\theta\), by \eqref{eq:f15v29:changed-L}.
The output log likelihood is \(r_\theta-\Lambda_\theta\),
so
\begin{align*}
 \mathbb E_{\rho_\theta}r_\theta
   &=D(\rho_\theta\Vert\rho_\theta^{\rm prod})
                                  +\Lambda_\theta\ge-\delta_\theta,\\
 \mathbb E_{\rho_\theta^{\rm prod}}r_\theta
   &=\Lambda_\theta-D(\rho_\theta^{\rm prod}\Vert\rho_\theta)
                                                   \le\delta_\theta.
\end{align*}
Under the first law use
\(|r|=2r_+-r\); under the second use
\(|r|=2r_-+r\). This proves both \(3\delta_\theta\)
bounds. The argument uses conditional normalizers
explicitly, rather than inferring normalization from
a small expected log error.
\end{proof}

At \(\theta=0\) both retained laws equal \(\rho_0\)
and \(r_\theta=0\) exactly, though \(P\) and \(Q\)
can still differ in their core/exterior dependence.
For nonzero \(\theta\), the near marginals of the two
retained laws need not agree. This is a different comparison
from V28's joint natural law, which preserves the changed
joint near marginal exactly. Nor is a uniformly small
source-log error at every prescribed retained configuration
being asserted.

\subsection{Uniform per-core rate, without a fixed-strength all-site claim}

\begin{proposition}[Superalgebraic intensive errors on the total-strength ball]
\label{f15v29:rate}
Take the selected scale \eqref{eq:f15v20:scale},
\(2\le p\le t=\log\beta\), and all inherited thresholds.
Put \(c_\Delta=2^{-51}C_0^{-1}\) and
\(C_H=2^{25}(1+J_*)\). Eventually, uniformly for
\(s\le1/8\),
\begin{equation}
 \epsilon_{\rm src}(s)
 \le4\sqrt2 C_H\,\beta^{3/2}t^{8\nu+118}
       \exp\left\{10J(1/4)t^4-\tfrac12c_\Delta t^6\right\}.
 \label{eq:f15v29:rate}
\end{equation}
This tends to zero faster than every fixed negative power
of \(\beta\). Thus the pressure error divided by \(n\),
both joint and output entropies divided by \(n\), and
both absolute source-log errors divided by \(n\) have
uniform superalgebraic bounds over every fitting finite
family. If also \(n\le\beta^B\) for any fixed \(B>0\),
all the total errors, including TV, vanish
superalgebraically.
\end{proposition}

\begin{proof}
V23 and \eqref{eq:f15v28:Eplus-rate} give
\[
 \mathfrak h\le C_H\beta t^{4\nu+61},\qquad
 \mathcal E_+\le2\beta t^{8\nu+114}e^{-c_\Delta t^6}.
\]
Since \(k\le5\), \(n_p\le4p^4\le4t^4\), and
\(J\) is increasing on \([0,1)\),
\(\mathcal R(s)\le e^{10J(1/4)t^4}\).
Also eventually
\(2+2\mathfrak h\le4C_H\beta t^{4\nu+61}\).
Substituting these three bounds in
\eqref{eq:f15v29:error} proves \eqref{eq:f15v29:rate}.
Eventually \(10J(1/4)t^4\le c_\Delta t^6/4\), so
the remaining exponential decay dominates any fixed
power of \(\beta=e^t\) and any displayed power of \(t\).
The finite-family estimates are all linear in
\(\delta_\theta=n\epsilon_{\rm src}\), except for the
TV square root. A fixed factor \(\beta^B\) and a fixed
square root do not destroy superalgebraic decay.
\end{proof}

There is no upper bound on \(n\) in the finite-regulator
inequalities apart from geometric fitting. The limiting
claim for \emph{total} errors does have an explicit
family-size condition; arbitrarily large total variation
errors are not made small by dividing another quantity
by \(n\). Likewise, the new \(\beta^B\) window cannot
be compared with V28's \(t^6/p^4\) window without noting
the source change: V28 allowed a fixed amplitude at every
site, whereas V29 requires \(\sum_z|\theta_z|\le1/8\).
It gives no uniform control on that full per-site cube.

\subsection{The remaining upstream problem}

We now have actual single-region near factors and a
normalized additive source comparison relative to the
original baseline. This settles a restricted version of
the specific decomposition issue left at the end of V28.
The factors are local in their configuration arguments
\(N_z\), but may still depend on the full regulator and
selected baseline through their defining kernels. No
uniform derivative bound, interaction norm, or RG flow
equation for them has been proved.

The missing extensive-source step is identifiable. Weighted
H\"older uses total source strength, and a fixed-amplitude
source at every site exhausts the scalar CGF radius as
the family grows. Removing that limitation would require
additional control of the joint near-source normalizer or
its connected responses under the actual retained law.
Small conditional core information alone does not supply
that retained-law estimate. In particular choosing cores
independent of an arbitrarily correlated retained field
makes \(P=Q\), all present errors zero, and leaves all
retained correlations untouched. The new theorem therefore
cannot imply retained mixing or a mass gap on its own.

A next upstream task is to inspect whether the companion
programmes or the existing source-response bounds control
those connected near-factor responses with a regulator-
and volume-uniform norm. A real-source \(\ell^1\) bound
must not silently be differentiated as a complex analytic
estimate or promoted to an all-density entropy inequality.
The V30 checkpoint will review the current companion files
before selecting that next step.

\subsection{Exact diagnostics and preservation}

The standalone diagnostic is
\begin{center}\small
\path{scripts/verify_ym_f15_v29_additive_natural_sources.py}.
\end{center}
It checks 425 scalar log-square inequalities and six
positive correlated finite families, with two, three or
four separate near bits and a common exterior bit. There
are 18 exact single-source marginal identities, 18
information-projection ledgers, 18 local and 12 joint
baseline log-square comparisons. Thirty signed source
vectors include uniform, alternating, nonuniform, sparse
and zero patterns. The nonzero vectors give 48 rational
H\"older checks and 48 centered second-moment comparisons;
there are 60 Jensen checks including the zero cases.
Further checks cover 560 exact output-factor identities,
30 pressure brackets, 560 conditional log brackets,
60 changed-law absolute log-ratio estimates, 30 two-sided
joint/output KL comparisons and 60 absolute source-log
bounds. All six zero-source retained comparisons are
exactly the baseline.

Probabilities and source weights are rational. Logarithms
use rational \(\operatorname{arctanh}\) enclosures with
outward \(2^{-120}\) rounding. Integer-power comparisons
verify weighted H\"older without numerical fractional
powers; squared bounds avoid approximate square roots.
There is no floating-point tolerance. These models check
the new probability identities and inequalities using their
actual scalar moments, not the Wilson CGF, geometric
hypotheses, or continuum construction. The script has
12,108 bytes and SHA--256
\begin{center}\scriptsize\ttfamily
C136D10E59EC72BC8861B4B8232A04CE953956B1DF4D7A9704074EFBB9BA5DA5.
\end{center}

The V28 predecessor has 649,623 bytes and SHA--256
\begin{center}\scriptsize\ttfamily
07BBD458B3E3338822FE8D1FB65393945FE27F6F93151293BA19740DDE7DCB94.
\end{center}
Removing this terminal section and reversing the title
version recovers it exactly. V29 replaces V28 as the sole
active main note. All 14 protected sources and 71 earlier
certificates are preserved. The next companion checkpoint
is V30.

\begin{firewall}
V29 proves a normalized additive natural near-source
comparison under a fixed original baseline and the signed
total-strength condition \(\sum_z|\theta_z|\le1/8\).
Its pressure, two-sided joint/output entropy and absolute
source-log bounds are uniform per core over every fitting
finite family, and have superalgebraic total errors for
polynomially growing families on the selected scale.
It does not prove fixed-strength extensive-source control,
baseline retained clustering, all-volume contraction,
an interacting continuum limit, or the Yang--Mills mass gap.
\end{firewall}

\section{V30: rare retained phases obstruct extensive source normalization}
\label{f15v30:section}

The V29 total-strength theorem raises a concrete upstream
question. Could its actual single-near reference inherit the
original law's arbitrary-support source bound, thereby
allowing fixed-strength sources at every core? This section
answers that inference negatively. The obstruction uses
the \emph{actual} near kernels, uniformly bounded sources,
uniformly positive compact densities, and even independent
true sources. A rare retained phase is nevertheless selected
by the normalized product reference. Its output becomes
almost maximally different from the actual source output.

This is a finite probability countermodel, not a Wilson
gauge theory or a counterexample to Yang--Mills. It shows
exactly why one must return to structure discarded by the
averaged core-entropy summary. In particular the Wilson
joint bad-patch estimates must not be replaced by their
first-moment consequences and then assumed recoverable.
The original continuum and physical mass-gap goal is unchanged.

\subsection{V30 companion checkpoint and the resulting direction}

The active-file inventory and the supplied attachments
were rechecked. All 86 previously protected source and
certificate hashes agree with the preceding audit.
The following latest conclusions and their scope boundaries
were reread; the optional suggestions were also reread as
mathematical proposals, not instructions.

\begin{center}\small
\begin{tabular}{@{}p{0.25\textwidth}p{0.69\textwidth}@{}}
\toprule
Source & Relevant result and missing interface\\
\midrule
SM continuum V72 &
Fixed-mode, fixed-time many-copy derivation of its covariance
flow. Growing ultraviolet cutoff and the physical finite-species
limit are still separate.\\[2mm]
NS stability V26 &
Pointwise rank-one kernel derivative norms. References to
future V27 numerical certification do not supply a certificate
in the current file or a YM source estimate.\\[2mm]
Shell mixing V16 &
Complete measured flat-holonomy one-loop coefficient in its
specified source convention. The noncommuting calculation and
a common nonperturbative remainder are not supplied.\\[2mm]
Parallel YM V5 &
The volume-uniform global extensive-charge tightness route
fails. Its valid refinement transport does not cure that
escape or give a retained-source pressure bound.\\[2mm]
SM source selection V31 &
An ancestry compiler conditional on word-native typing
operations. Historical occurrence of those operations is
not proved by the compiler.\\[2mm]
Native stationarity V34 &
A moving-fibre escape reduction; its 330 action-specific
first-escape entries remain a stated calculation, not a
populated input to YM.\\[2mm]
Exact-action bridge V24 &
A nonzero first-jet derivative of an obstruction with
linearized lower incidences preserved. This is not an
integrated root or a uniform continuum bridge.\\
\bottomrule
\end{tabular}
\end{center}

None provides the needed volume-uniform estimate for
the natural product source factors. The suggestions'
measured-density and source-convention cautions remain
appropriate, but their flat one-loop calculation cannot
replace this probability estimate. We therefore keep the
original full-compact source convention and test the
purported implication before adopting it as a new premise.
The next scheduled companion checkpoint is V35.

\subsection{A positive compact model with independent true sources}

For this section only, let \(0<\eta\le1/4\) be a rare-phase
probability, unrelated to the confidence parameter in V23.
Fix any \(n\ge1\), and set \(a=b=1/4\). Let
\(H\in\{0,1\}\), \(\mathbb P(H=1)=\eta\). Independently
of \(H\), let \(X_1,\ldots,X_n\) be independent fair
signs. Conditional on \((H,\mathbf X)\), choose separate
near signs \(N_1,\ldots,N_n\) independently according to
\begin{equation}
 P(N_i=x\mid X_i=s,H=h)
       =\frac{1+x(as+bh)}2,
 \qquad x,s\in\{-1,1\},\quad h\in\{0,1\}.
 \label{eq:f15v30:channel}
\end{equation}
Every atom of the resulting joint law \(P\) is strictly
positive. The retained field is \(C=\mathbf N\), the
common exterior is \(Y=H\), and the core at site \(i\)
is \(V_i=X_i\). Distinct near outputs are distinct
coordinates. Their law need not be independent.

Put \(\xi=(\mathbf N,H)\) and write \(\nu_0\) for its
actual marginal. Direct summation in
\eqref{eq:f15v30:channel} gives
\begin{align}
 \nu_0(\mathbf x,h)
      &=P(H=h)\prod_i\frac{1+bhx_i}{2},
                                          \label{eq:f15v30:nu}\\
 \kappa_i^{x,h}(s)
      &=P(X_i=s\mid\mathbf N=\mathbf x,H=h)
        =\frac12\left(1+\frac{ax_i s}{1+bhx_i}\right),
                                          \label{eq:f15v30:kappa}\\
 q_i(x,s)
      &=P(X_i=s\mid N_i=x)
        =\frac12\left(1+\frac{ax s}{1+\eta bx}\right).
                                          \label{eq:f15v30:q}
\end{align}
In particular the cores are conditionally independent at
fixed \(\xi\). Define exactly the V29 natural reference
\begin{equation}
 Q(d\mathbf x,dh,d\mathbf s)
       =\nu_0(d\mathbf x,dh)\prod_i q_i(x_i,ds_i).
 \label{eq:f15v30:Q}
\end{equation}
It preserves the actual retained/exterior law and every
single \((N_i,X_i)\) marginal. It is not an arbitrarily
chosen chart proxy.

Let compact Haar at each core mean the uniform measure
on the two-element group. The \(\kappa_i\) and \(q_i\)
Haar densities all lie in \([2/3,4/3]\), since their
means have absolute value at most \(a/(1-b)=1/3\).
Thus they obey the V29 form of compact density bounds
with a constant independent of \(n,\eta\):
\begin{equation}
 \mathfrak h=\log(3/2),\qquad
          e^{-\mathfrak h}\le\frac{d\kappa_i}{dH_i},
                \frac{dq_i}{dH_i}\le e^{\mathfrak h}.
 \label{eq:f15v30:compact}
\end{equation}
The notation in \eqref{eq:f15v30:compact} means that each
of the two densities lies between the two displayed bounds.

\begin{proposition}[Small intensive information and a uniform true-source CGF]
\label{f15v30:baseline}
For \(\lambda_i(x,s)=(1+axs)/2\),
\begin{equation}
 e_i:=\mathbb E_{\nu_0}D(\kappa_i^\xi\Vert q_i(N_i))
   \le\mathbb E_{\nu_0}D(\kappa_i^\xi\Vert\lambda_i(N_i))
   \le\frac{\eta}{135},\qquad
                         D(P\Vert Q)=ne_i.
 \label{eq:f15v30:entropy}
\end{equation}
Take the nonnegative energies \(E_i=1+X_i\) and their
baseline-centred sources \(S_i=E_i-\mathbb E_PE_i=X_i\).
Then \(|S_i|=1\), and for every source vector with
\(\|\mathbf u\|_\infty<1\),
\begin{equation}
 \log\mathbb E_Pe^{\sum_i u_iS_i}
     =\sum_i\log\cosh u_i
     \le\frac12\sum_i u_i^2
     \le2\sum_iJ(u_i),\qquad J(u)=-\log(1-u)-u.
 \label{eq:f15v30:CGF}
\end{equation}
These constants are independent of \(n\) and \(\eta\).
\end{proposition}

\begin{proof}
At \(H=0\), \(\kappa_i=\lambda_i\) exactly. At \(H=1\),
let \(p_1,p_0\) be the respective probabilities of
\(X_i=1\) at fixed \(N_i=x\). Then
\[
 |p_1-p_0|\le\frac{ab}{2(1-b)}=\frac1{24},\qquad
                         p_0(1-p_0)=\frac{15}{64}.
\]
For Bernoulli laws, \(D(p_1\Vert p_0)\le
\chi^2(p_1\Vert p_0)=(p_1-p_0)^2/[p_0(1-p_0)]\),
by \(\log t\le t-1\). Thus this conditional entropy is
at most \(1/135\). Averaging pays \(P(H=1)=\eta\).
The actual single-near kernel is the information projection
of Proposition~\ref{f15v29:projection}, proving the first
inequality. Product conditional kernels give
\(D(P\Vert Q)=\sum_i e_i\), and exchangeability makes
the summands equal.

By construction the \(X_i\)'s are independent fair signs
under \(P\), including after averaging the channel.
This proves the exact CGF formula. The second derivative
of \(\log\cosh u\) is at most one, proving its quadratic
bound. For \(u\ge0\), \(J(u)\ge u^2/2\). For
\(u=-v\), \(0\le v<1\), the derivative of
\(J(-v)=v-\log(1+v)\) is \(v/(1+v)\ge v/2\).
Hence \(J(-v)\ge v^2/4\), proving the final inequality
for both signs.
\end{proof}

Even the individual near posterior of the exceptional
phase is uniformly small:
\begin{equation}
 P(H=1\mid N_i=x)
       =\eta\frac{1+bx}{1+\eta bx}\le\frac43\eta.
 \label{eq:f15v30:near-confidence}
\end{equation}
Indeed the numerator is at most \(5\eta/4\) and the
denominator at least \(15/16\). Uniform confidence from
one near coordinate is therefore not enough for the
many-source conclusion disproved below.

\subsection{The full joint source bound is not inherited by the natural reference}

\begin{proposition}[An explicit collective susceptibility in the reference]
\label{f15v30:variance}
Conditional on \(H=h\) under \(Q\), the sources are
independent signs with mean
\begin{equation}
 m_h=\frac{a(bh-\eta b)}{1-\eta^2b^2},\qquad
 m_0=-\frac{\eta}{16-\eta^2},\quad
 m_1=\frac{1-\eta}{16-\eta^2}.
 \label{eq:f15v30:m}
\end{equation}
Their unconditional means remain zero. Set
\(v_\eta=\eta(1-\eta)/(16-\eta^2)^2\). Then
\begin{equation}
 \Var_Q\left(\sum_i u_iS_i\right)
   =(1-v_\eta)\sum_i u_i^2
                       +v_\eta\left(\sum_i u_i\right)^2.
 \label{eq:f15v30:variance}
\end{equation}
For every fixed \(\eta>0\), no joint quadratic CGF
upper bound with a coefficient independent of \(n\)
holds for \(Q\) on a fixed neighbourhood of zero.
\end{proposition}

\begin{proof}
Given \(H=h\), the \(N_i\)'s are independent signs of
mean \(bh\), and the \(Q\) core mean at fixed \(N_i=x\)
is \(ax/(1+\eta bx)\). Averaging these two possible
values gives \eqref{eq:f15v30:m}; the same factorization
gives conditional independence. Directly
\((1-\eta)m_0+\eta m_1=0\) and
\(\Var(m_H)=v_\eta\). Single-source variance is one
and distinct-source covariance is \(v_\eta\), which
proves \eqref{eq:f15v30:variance}. With all \(u_i=1\)
that variance equals \(n+n(n-1)v_\eta\). Differentiating
any proposed quadratic CGF bound at zero would bound
this by a fixed multiple of \(n\), a contradiction for
large \(n\). This differentiation concerns explicit
finite entire moment functions, not an assumed complex
extension of V29's estimates.
\end{proof}

Thus preserving every source's marginal and retaining
the original joint CGF do not transfer that CGF to the
product-kernel reference. The obstruction is already
visible at its second derivative, but its normalized
output consequence is stronger than a large variance.

\subsection{A fixed small source selects the rare reference phase}

Use the same source at every site, with the fixed value
\begin{equation}
 \begin{aligned}
 \theta_0&=\log(9/8)<1/8,
 &u&=\tanh\theta_0=\frac{17}{145},\\
 c&=\cosh\theta_0=\frac{145}{144},
 &g&=\log\left(1+\frac{u}{32}\right)=\log\frac{4657}{4640}>0.
 \end{aligned}
 \label{eq:f15v30:theta}
\end{equation}
Let \(P_\theta=W P/Z_P\), \(Q_\theta=W Q/Z_Q\),
where \(W=e^{\theta_0\sum_i S_i}\), and let
\(\rho_\theta,\rho_\theta^{\rm prod}\) be their retained
marginals. The latter is precisely the normalized product
of the actual single-near factors:
\begin{equation}
 b_{i,\theta_0}(x)=\mathbb E_P[e^{\theta_0 S_i}\mid N_i=x]
          =c\left(1+\frac{aux}{1+\eta bx}\right),\qquad
 \frac{d\rho_\theta^{\rm prod}}{d\rho_0}(\mathbf x)
          =\frac{\prod_i b_{i,\theta_0}(x_i)}{Z_Q}.
 \label{eq:f15v30:factors}
\end{equation}
For \(n=1\), this retained law equals the actual source
law exactly, by the tower property. The failure is collective.

Put \(L_\eta=\log(1/\eta)\), and choose
\begin{equation}
                 n=n_\eta=\left\lceil\frac{2L_\eta}{g}\right\rceil.
 \label{eq:f15v30:n}
\end{equation}

\begin{theorem}[A nonvanishing intensive pressure error]
\label{f15v30:pressure}
At the above source and family size,
\begin{equation}
 \frac1n\log\frac{Z_Q}{Z_P}\ge\frac g2,
 \qquad P_\theta(H=1)=\eta,
 \qquad Q_\theta(H=0)\le\eta.
 \label{eq:f15v30:phase}
\end{equation}
Thus the actual tilted law keeps the exceptional phase
rare while the natural product reference makes it dominant.
\end{theorem}

\begin{proof}
The true sources are independent of \(H\), so their
tilt does not change its law, and \(Z_P=c^n\).
By Proposition~\ref{f15v30:variance},
\begin{equation}
 \frac{Z_Q}{Z_P}
       =(1-\eta)(1+m_0u)^n+\eta(1+m_1u)^n.
 \label{eq:f15v30:normalizers}
\end{equation}
Here \(m_0<0\), and \(m_1\ge1/32\) for
\(\eta\le1/4\), since \(1-\eta\ge1/2\) and
\(16-\eta^2\le16\). Keeping the second positive term
gives \(\log(Z_Q/Z_P)\ge ng-L_\eta\ge ng/2\).
The exact phase odds also give
\[
 Q_\theta(H=0)
 \le\frac{1-\eta}{\eta}
       \left(\frac{1+m_0u}{1+m_1u}\right)^n
 \le\eta^{-1}e^{-ng}\le\eta.
\]
All normalizers here are explicit positive finite numbers.
\end{proof}

This argument is not merely about an unobserved exterior
label. The retained near signs distinguish the two laws
with probability tending to one, as follows.

\subsection{Separation of the actual retained laws in TV and both entropy directions}

\begin{theorem}[Almost maximal output separation from vanishing intensive information]
\label{f15v30:output}
For the event
\(A_n=\{n^{-1}\sum_iN_i>1/8\}\), one has
\begin{equation}
 \rho_\theta(A_n)\le2\eta,\qquad
 \rho_\theta^{\rm prod}(A_n)\ge1-2\eta.
 \label{eq:f15v30:event}
\end{equation}
Consequently
\begin{align}
 d_{\rm TV}(\rho_\theta,\rho_\theta^{\rm prod})
                                           &\ge1-4\eta,
                                          \label{eq:f15v30:TV}\\
 D(\rho_\theta\Vert\rho_\theta^{\rm prod}),\quad
 D(\rho_\theta^{\rm prod}\Vert\rho_\theta)
               &\ge(1-2\eta)\log\frac1{2\eta}-\log2.
                                          \label{eq:f15v30:KL}
\end{align}
In particular the TV distance tends to one as \(\eta\to0\),
and each of the two entropies divided by \(n_\eta\) has
lower limit at least \(g/2>0\).
\end{theorem}

\begin{proof}
Under \(P_\theta\) at fixed \(H=h\), the near signs
are independent with mean \(au+bh\). At \(h=0\)
this is \(au<1/32\). Under \(Q_\theta\), they are
also independent at fixed \(H\): integrate the tilted
core kernels in \eqref{eq:f15v30:Q}. Their common mean is
\begin{equation}
 \overline N_h
   =\frac{bh+au(1-bh\eta b)/(1-\eta^2b^2)}{1+m_hu}.
 \label{eq:f15v30:near-mean}
\end{equation}
Subtracting \(bh\) in this formula gives
\[
 \overline N_h-bh
  =\frac{au(1-b^2h^2)}{(1-\eta^2b^2)(1+m_hu)}>0.
\]
Thus \(\overline N_1\ge b=1/4\).

For completeness, any sign \(Z\) with mean \(m\)
has \(\log\mathbb E e^{t(Z-m)}\le t^2/2\): the
second derivative of its log moment function is a
variance under an exponential tilt and is at most one,
while its value and first derivative vanish at zero.
For independent signs, multiplication and exponential
Markov therefore bound either deviation of their mean
by \(d>0\) by \(e^{-nd^2/2}\).
Apply this with \(d=3/32\) to the true \(H=0\) law
and with \(d=1/8\) to the reference \(H=1\) law:
\begin{equation}
 P_\theta(A_n\mid H=0)\le e^{-9n/2048},\qquad
 Q_\theta(A_n^c\mid H=1)\le e^{-n/128}.
 \label{eq:f15v30:tails}
\end{equation}
Since
\[
 0<g<\frac{17}{4640}
                  <\min\left\{\frac9{1024},\frac1{64}\right\},
\]
the choice \eqref{eq:f15v30:n} makes each exponential
in \eqref{eq:f15v30:tails} at most \(\eta\).
Combining with \eqref{eq:f15v30:phase} proves
\eqref{eq:f15v30:event}, and hence the TV lower bound.

Project both retained laws onto \(1_{A_n}\). If
probabilities of an event are \(q\ge1-2\eta\) and
\(p\le2\eta\), the resulting binary entropy obeys
\[
 D(\operatorname{Bern}(q)\Vert\operatorname{Bern}(p))
       \ge q\log(1/p)-\log2
       \ge(1-2\eta)\log(1/(2\eta))-\log2.
\]
The first inequality follows by writing out the binary
entropy and using its upper bound \(\log2\). Use
\(A_n\) for the reference-to-actual direction and
\(A_n^c\) for the actual-to-reference direction.
Data processing proves \eqref{eq:f15v30:KL}.
Finally \(L_\eta/n_\eta\to g/2\), giving both
intensive lower limits.
\end{proof}

One may parameterize this family by
\(\eta_\beta=e^{-(\log\beta)^6}\), for large \(\beta\).
Then \(n_\eta=O((\log\beta)^6)\), the intensive
baseline bound \(\eta_\beta/135\) is superalgebraically
small, and the compact cost and source amplitudes remain
bounded by fixed constants. Nevertheless the pressure
error per core is at least \(g/2\), both output entropy
errors per core have positive lower limits, and TV tends
to one. Here \(\beta\) is just an indexing parameter;
this is not a Wilson coupling selection.

\subsection{Precisely what the countermodel excludes, and what it does not}

The example rules out an extensive-source extension of
V29 based solely on the following reduced inputs:
conditional product cores, fixed positive compact density
bounds, vanishing baseline entropy per core, exact
single-near marginal preservation, bounded centred
sources, and an arbitrary-support CGF for the true law.
It even has the uniform single-near confidence bound
\eqref{eq:f15v30:near-confidence}. No error modulus
tending to zero with the intensive baseline entropy,
uniform in all family sizes at a fixed positive source,
can control the pressure or either output entropy per
core using only those inputs.

This does not contradict V29: \(n_\eta\theta_0\to\infty\),
so the total-strength condition fails. Nor does it
contradict V28: the combined near output is the whole
retained field here, and its joint natural normalizer
therefore gives the actual retained source law exactly.
The loss comes from replacing that joint normalizer
by a product of single-near ones.

The result is also distinct from V27's counterexample.
There the chosen chart factor fails even for a single
source. Here the one-source retained comparison is
exact; every factor is an actual baseline conditional
expectation, every source has absolute value one, and
all compact density bounds are uniform. The failure
arises only after collectively reweighting the
correlated retained law.

The omitted Wilson input can be identified explicitly.
In this model let the exceptional event for each core be
\(B_i=\{H=1\}\). On its complement
\(\kappa_i=\lambda_i\) exactly, but for every nonempty
set of sites \(I\),
\begin{equation}
                    P\left(\bigcap_{i\in I}B_i\right)=\eta.
 \label{eq:f15v30:coherent-bad}
\end{equation}
This does not satisfy a many-patch estimate decaying
exponentially in \(|I|\) with fixed subunit base and
fixed colouring loss. The V20--V21 Wilson estimates
retain such joint spatial information; it cannot be
reconstructed from \eqref{eq:f15v30:entropy}. Moreover
the true changed law has
\(\mathbb E_{P_\theta}\sum_i1_{B_i}/n=\eta\) even
under arbitrary source vectors, since its sources are
independent of \(H\). At the source above the same
fraction under \(Q_\theta\) is at least \(1-\eta\).
Thus a small exceptional fraction under the actual
changed law is not by itself a reference-law estimate.

The upstream target is now sharper: retain the joint
Wilson geometry when estimating rare retained phases
under the \emph{product-factor weight}, or establish
a volume-uniform retained interaction/pressure estimate
that pays for them. Neither a new one-source confidence
threshold nor repetition of the true-law CGF addresses
the countermodel. The needed theorem must distinguish
the actual Wilson near field from the coherent-phase
construction above.

For context, the entropy tensorization results of
\href{https://arxiv.org/abs/1405.0608}
{Caputo--Menz--Tetali, \emph{Approximate tensorization of
entropy at high temperature}}, require quantitative weak
dependence. No such retained-field hypothesis has been
verified here, and no result from that paper is applied
to the Wilson law. The physical endpoint remains the
nontrivial continuum theory with a positive Hamiltonian
gap in the
\href{https://www.claymath.org/wp-content/uploads/2022/06/yangmills.pdf}
{Jaffe--Witten problem description}; this finite
probability obstruction supplies neither component.

\subsection{Exact diagnostics and predecessor preservation}

The diagnostic is
\begin{center}\small
\path{scripts/verify_ym_f15_v30_rare_retained_phase.py}.
\end{center}
Twelve positive models use \(\eta=1/8,1/32,1/512\)
and \(n=1,2,3,4\). The checks include 192 compact
density bounds, 120 exact single-near kernel identities,
30 marginal-preservation identities, 12 information
projection ledgers and 12 entropy upper bounds.
They verify 36 source normalizers for uniform, alternating
and nonuniform source arrays, all 12 reference variance
formulas, 90 natural output-factor identities, 12 phase
posterior identities, 60 conditional retained means,
24 near-confidence bounds and the three exact one-core
retained-law equalities.

Four larger-family certificates use
\((\eta,n)=(2^{-4},1517),(2^{-8},3033),
(2^{-12},4549),(2^{-20},7582)\). Rational logarithm
enclosures verify the integer threshold; exact integer
powers verify the phase-odds bounds. Eight retained-tail
certificates and eight entropy lower-bound checks retain
the analytic constants above. All arithmetic is rational,
with outward logarithm enclosures on a \(2^{-120}\) grid
and no floating-point tolerance. The limiting theorem
is proved analytically, not inferred from these samples.
The script has 8,652 bytes and SHA--256
\begin{center}\scriptsize\ttfamily
C771B4758983D7AFEBB27D13B0516726721F484A811E9466BA5BBE655A64606F.
\end{center}

The V29 predecessor has 673,354 bytes and SHA--256
\begin{center}\scriptsize\ttfamily
AA1A9259BDEA28AE67A1BA7D035D9459744FC1B612CB22491D8B1FE23694FE3A.
\end{center}
Removing this terminal section and reversing the title
version recovers it exactly. V30 replaces V29 as the sole
active main note. All 14 protected sources and 72 earlier
certificates are preserved. The V30 companion checkpoint
is complete; the next is V35.

\begin{firewall}
V30 proves a nonimplication for the reduced probability
inputs used by the recent source comparisons. Even the
true law's uniform joint source CGF does not prevent the
normalized natural product reference from selecting a
rare retained phase. This rules out the proposed
information-only extensive-source upgrade and directs
the next work back to the joint Wilson structure.
It does not refute a Wilson-specific upgrade, establish
retained clustering, construct an interacting continuum,
or prove the Yang--Mills mass gap.
\end{firewall}

\section{V31: projected failure counts pay for extensive source normalization}
\label{f15v31:section}

V30 showed that the averaged core-entropy summary and the
true law's source CGF do not alone control an extensive
product-factor perturbation. We now return to a stronger
Wilson input: the joint probability bound for actual
chart failures in a separated family. Its count exponential
bound survives projection onto the \emph{common} retained
field. A fractional moment argument then controls the
previously missing normalization error in V27, and an
entropy inequality controls the reference law without
postulating a CGF for that law.

The reference in this section is explicitly V27's
\emph{unfloored chart} factor, not V29's actual single-near
kernel. The result applies to arbitrarily many fitting
separated sources, with a new per-site amplitude window
decreasing as a fixed power of \(\log\beta\). It is not
a fixed-strength extensive-source theorem. Its source
window is nevertheless much larger than the globally
bounded-weight window obtained by applying V25 to the
physical action sources.

All 87 earlier protected sources and certificates and
the V30 predecessor were rechecked unchanged. V30 completed
the companion checkpoint, including the supplied SM
source-selection V31, native stationarity V34 and exact-action
Einstein bridge V24. The next checkpoint is V35. No new
estimate from those companions is assumed below. The
selected coupling, completed blocking, source centering,
support geometry and inherited CGF/chart thresholds remain
as stated in the earlier Wilson results.

\subsection{Retain the raw spatial input and the exact source-log signs}

Fix a separated family \(\mathcal S\), \(|\mathcal S|=n\ge1\),
as in Proposition~\ref{f15v21:support}. Use the physical
sources of \eqref{eq:f15v26:sources}, with
\(\theta\in[-\tau,\tau]^n\), \(0<\tau\le1/8\),
and \(2\tau<R_{\rm src}\). Put
\begin{equation}
 \begin{aligned}
 W_\theta&=e^{\sum_z\theta_zS_z},&
 Z_a&=\mathbb E_\mu W_\theta,\qquad
 d\mu_\theta=W_\theta\,d\mu/Z_a,\\
 A_\theta(C)&=\mathbb E_\mu[W_\theta\mid C],&
 \rho_a&=C_*\mu_\theta,\qquad \rho_0=C_*\mu.
 \end{aligned}
 \label{eq:f15v31:actual}
\end{equation}
The whole retained field \(C\) is the same for every
conditional expectation in this section.

For each source use V27's positive scalar chart normalizer
\begin{equation}
 b_{z,\theta_z}(N_z)
       =\lambda_z(N_z)(e^{\theta_z s_z(N_z,V_z)}),\qquad
 T_\theta(C)=\prod_z b_{z,\theta_z}(N_z),
 \label{eq:f15v31:chart-factors}
\end{equation}
including the same Haar fallback when the required near
retained chart condition fails. These \(\lambda_z\)'s
are the unfloored chart kernels of V22. Define
\begin{equation}
 Z_b=\mathbb E_{\rho_0}T_\theta,\qquad
 d\rho_b=T_\theta\,d\rho_0/Z_b,\qquad
 r_\theta=\log A_\theta-\log T_\theta,\qquad
 \delta_\theta=\log(Z_b/Z_a).
 \label{eq:f15v31:reference}
\end{equation}
The reference retains the full correlated \(\rho_0\);
it is not a product retained measure. Nothing is divided
out of the actual source weight \(W_\theta\).

Let \(E_z=\Gamma_z^c\cup D_z^c\) be the actual full-law
failures of V21, and write
\begin{equation}
 R=\sum_{z\in\mathcal S}1_{E_z},\qquad
 G_0(C)=\mathbb E_\mu[R\mid C],\qquad
 G_\theta(C)=\mathbb E_{\mu_\theta}[R\mid C].
 \label{eq:f15v31:counts}
\end{equation}
The separated Wilson input, before any colour-root loss,
is
\begin{equation}
 \mu\left(\bigcap_{z\in I}E_z\right)\le q_0^{|I|}
   \quad(I\subseteq\mathcal S),\qquad 0<q_0<1.
 \label{eq:f15v31:joint-input}
\end{equation}
In particular the count bound already proved in V26 is
\begin{equation}
 \log\mathbb E_\mu e^{vR}\le n m(v),\qquad
 m(v)=\log[1+q_0(e^v-1)],\quad v\ge0.
 \label{eq:f15v31:raw-mgf}
\end{equation}

Keep the whole compact source bound and its clipped-log
constant
\begin{equation}
 B_p=48\beta p^4,\qquad K=2\tau B_p,
 \qquad c_K=\frac{K}{1-e^{-K}}\le1+K.
 \label{eq:f15v31:K}
\end{equation}
V27's sign-separated theorem gives the pointwise bounds
\begin{equation}
 (r_\theta)_-\le n\Delta+c_KG_0,
 \qquad
 (r_\theta)_+\le n\Delta+c_KG_\theta.
 \label{eq:f15v31:sign-input}
\end{equation}
We have only enlarged its sum over active sources to the
whole family. These are bounds on normalized chart
comparisons, with \(\Delta\) unchanged from V22. In
particular there is no factor \(e^K\Delta\). No new
logarithmic oscillation assertion for the full Wilson
potential is needed.

Finally V26 supplies the exact order-two source budget
\begin{equation}
 h_\theta=\frac1n\log
          \frac{\mathbb E_\mu W_\theta^2}{Z_a^2},\qquad
 0\le h_\theta\le h_\tau:=k n_pJ(2\tau),\quad
 J(u)=-\log(1-u)-u.
 \label{eq:f15v31:h}
\end{equation}
It also gives \(D(\mu_\theta\Vert\mu)\le nh_\theta\).
The actual source CGF is never conditioned on a retained
field or asserted for \(\rho_b\).

\subsection{A common projection preserves the count exponential bound}

\begin{proposition}[Projected count control, without a product claim]
\label{f15v31:projection}
For every \(v\ge0\),
\begin{equation}
                  \log\mathbb E_{\rho_0}e^{vG_0}\le nm(v).
 \label{eq:f15v31:projected-mgf}
\end{equation}
If \(f=d\rho_a/d\rho_0=A_\theta/Z_a\), then
\begin{equation}
                  \mathbb E_{\rho_0}f^2\le e^{nh_\theta}.
 \label{eq:f15v31:projected-Renyi}
\end{equation}
\end{proposition}

\begin{proof}
Conditional Jensen gives
\(e^{v\mathbb E[R\mid C]}\le\mathbb E[e^{vR}\mid C]\).
Integrate and apply \eqref{eq:f15v31:raw-mgf}.
Likewise \(f=\mathbb E_\mu[W_\theta/Z_a\mid C]\), so
conditional Jensen for the square proves
\eqref{eq:f15v31:projected-Renyi}.
\end{proof}

This does not contradict V21's warning about the separate
projected probabilities. We have not asserted
\(\mathbb E\prod_z\mu(E_z\mid C)\le q_0^n\), or a
stochastic domination, or an analogous bound when each
site uses a different conditioning sigma-algebra. The
useful object is the exponential of their \emph{sum}
under one common conditioning. No rootwise union bound
and no factor equal to the total number of lattice sites
has been introduced.

\begin{proposition}[Fractional count moments under the actual source output]
\label{f15v31:fractional}
For \(v\ge0\) and \(0\le\lambda\le1/2\),
\begin{equation}
 \log\mathbb E_{\rho_a}e^{\lambda vG_0}
                      \le n\lambda\{h_\theta+m(v)\}.
 \label{eq:f15v31:fractional}
\end{equation}
\end{proposition}

\begin{proof}
For \(0<\lambda<1/2\), apply generalized H\"older with
weights \(1-2\lambda,\lambda,\lambda\) to the identity
\[
 f e^{\lambda vG_0}
         =f^{1-2\lambda}(f^2)^\lambda(e^{vG_0})^\lambda.
\]
Since \(\mathbb E_{\rho_0}f=1\), this gives
\[
 \mathbb E_{\rho_a}e^{\lambda vG_0}
  \le(\mathbb E_{\rho_0}f^2)^\lambda
       (\mathbb E_{\rho_0}e^{vG_0})^\lambda.
\]
Insert Proposition~\ref{f15v31:projection}. At
\(\lambda=1/2\) the same assertion is Cauchy--Schwarz;
at zero it is equality.
\end{proof}

The fractional coefficient multiplies the actual
order-two source cost \(h_\theta\), not the much larger
whole compact source oscillation \(K\). This is the
difference from applying V25's crude density bound to
physical sources. The argument is an elementary moment
inequality, not a new entropy tensorization theorem.

\subsection{The normalization estimate with an explicit absorption margin}

Choose any \(v>2c_K\), and set
\begin{equation}
 H_\theta(v)=h_\theta+m(v),\qquad
 \chi=\frac{2c_K}{v}<1,\qquad
 \epsilon_\theta(v)=\Delta+\frac{c_K}{v}H_\theta(v).
 \label{eq:f15v31:budgets}
\end{equation}
These are finite-regulator bounds for arbitrary finite
separated families. Their smallness is a separate scale
claim proved below. For brevity write \(\epsilon=\epsilon_\theta(v)\).

\begin{theorem}[Paid pressure and two-sided normalized output entropy]
\label{f15v31:normalization}
Under the preceding hypotheses,
\begin{align}
                  |\delta_\theta|&\le n\epsilon,
                                          \label{eq:f15v31:pressure}\\
 D(\rho_a\Vert\rho_b)&\le2n\epsilon,
                                          \label{eq:f15v31:forward}\\
 D(\rho_b\Vert\rho_a)&\le\frac{2n\epsilon}{1-\chi}.
                                          \label{eq:f15v31:reverse}
\end{align}
No source moment or relative-entropy bound for the
reference output is assumed in this theorem.
\end{theorem}

\begin{proof}
The true-law count transport from V26, evaluated at
the present \(v\), gives
\[
 \frac1n\mathbb E_{\mu_\theta}R
                  \le\frac{h_\theta+m(v)}v.
\]
The positive sign in \eqref{eq:f15v31:sign-input}
therefore implies
\begin{equation}
                   \mathbb E_{\rho_a}(r_\theta)_+\le n\epsilon.
 \label{eq:f15v31:positive-mean}
\end{equation}
By the exact definitions,
\begin{equation}
 \delta_\theta=\log\mathbb E_{\rho_a}e^{-r_\theta},\qquad
             \log\frac{d\rho_a}{d\rho_b}=r_\theta+\delta_\theta.
 \label{eq:f15v31:identities}
\end{equation}
Jensen and \eqref{eq:f15v31:positive-mean} give
\(\delta_\theta\ge-\mathbb E_{\rho_a}r_\theta\ge-n\epsilon\).
For the other direction use the negative sign bound and
Proposition~\ref{f15v31:fractional} with
\(\lambda=c_K/v=\chi/2<1/2\):
\[
 \delta_\theta
   \le n\Delta+\log\mathbb E_{\rho_a}e^{c_KG_0}
   \le n\Delta+n\frac{c_K}{v}H_\theta(v)=n\epsilon.
\]
This pays V27's missing negative-error exponential moment.
Taking the \(\rho_a\) expectation of the second identity
in \eqref{eq:f15v31:identities} proves the forward bound.

For the reverse direction put \(D_b=D(\rho_b\Vert\rho_a)\).
All quantities are finite before the following argument:
the bounded physical sources give \(|r_\theta|\le nK\),
\(|\delta_\theta|\le nK\), and the retained densities
are strictly positive relative to \(\rho_0\).
At \(\lambda=1/2\), Proposition~\ref{f15v31:fractional}
gives
\[
       \log\mathbb E_{\rho_a}e^{(v/2)G_0}
                                     \le\tfrac n2H_\theta(v).
\]
Apply nonnegativity of relative entropy between \(\rho_b\)
and the normalized probability proportional to
\(e^{(v/2)G_0}\rho_a\). This proves
\begin{equation}
 c_K\mathbb E_{\rho_b}G_0
       \le\chi D_b+n\frac{c_K}{v}H_\theta(v).
 \label{eq:f15v31:reference-count}
\end{equation}
Now use the exact reverse entropy identity and the
negative sign bound:
\begin{align*}
 D_b&=-\mathbb E_{\rho_b}r_\theta-\delta_\theta\\
    &\le n\Delta+c_K\mathbb E_{\rho_b}G_0+n\epsilon
     \le\chi D_b+2n\epsilon.
\end{align*}
Since \(D_b\) is finite and \(\chi<1\), subtract
\(\chi D_b\) and divide by \(1-\chi\).
This proves the reverse estimate without assuming the
conclusion as an entropy budget for \(\rho_b\).
\end{proof}

The bound just paid is the actual normalization quantity,
not merely a comparison of derivatives at zero or of
unweighted expected logs. Its proof uses the full spatial
count estimate before projection. V30's coherent rare
phase does not satisfy \eqref{eq:f15v31:joint-input} with
a family-independent subunit parameter.

\subsection{Absolute generated-action errors under both changed laws}

\begin{proposition}[Source-log and total-variation consequences]
\label{f15v31:log}
One has
\begin{align}
 \|r_\theta\|_{L^1(\rho_a)}&\le3n\epsilon,
                                          \label{eq:f15v31:actual-log}\\
 \|r_\theta\|_{L^1(\rho_b)}
                       &\le n\epsilon\frac{3+\chi}{1-\chi},
                                          \label{eq:f15v31:reference-log}\\
 d_{\rm TV}(\rho_a,\rho_b)&\le\min\{1,\sqrt{n\epsilon}\}.
                                          \label{eq:f15v31:TV}
\end{align}
\end{proposition}

\begin{proof}
The forward entropy identity implies
\(\mathbb E_{\rho_a}r_\theta
=D(\rho_a\Vert\rho_b)-\delta_\theta\ge-n\epsilon\).
Together with \eqref{eq:f15v31:positive-mean} and
\(|r|=2r_+-r\), this proves the first bound.
For the reference law, \eqref{eq:f15v31:reference-count}
and the negative sign bound give
\[
 \mathbb E_{\rho_b}(r_\theta)_-
      \le n\epsilon+\chi D_b
      \le n\epsilon\frac{1+\chi}{1-\chi}.
\]
Also \(\mathbb E_{\rho_b}r_\theta=-D_b-\delta_\theta
\le n\epsilon\). The identity \(|r|=2r_-+r\) proves
\eqref{eq:f15v31:reference-log}. Pinsker and the forward
entropy bound prove \eqref{eq:f15v31:TV}, with TV in
the half-\(L^1\) convention.
\end{proof}

For a zero source vector both retained laws equal
\(\rho_0\), and all source-log and pressure errors vanish
exactly. The displayed uniform budgets need not vanish
at that vector because they also allow nonzero sources
in the declared cube. The reference's unfloored chart
cores can be singular relative to full compact cores;
no two-sided full-core joint entropy statement is being
inferred from these retained-output bounds.

\subsection{An explicit logarithmically shrinking physical-source cube}

Use the existing selected scale
\(t=\log\beta\), \(2\le p\le t\), with the dyadic
block depths and all earlier embedding/chart thresholds.
Let
\begin{equation}
 c_* =\frac{2^{-108}}{9\pi^2C_0^2},\qquad
 d=6\nu+82,\qquad
 \tau_\beta=\min\left\{\frac18,
                   c_*2^{-12}t^{-(d+4)}\right\}.
 \label{eq:f15v31:window}
\end{equation}
This is a per-site bound \(|\theta_z|\le\tau_\beta\);
there is no bound on \(\sum_z|\theta_z|\).

\begin{theorem}[Intensive normalized locality for every fitting family]
\label{f15v31:rate}
Set \(\mathcal L=\log(1/q_0)\) and
\(v=\mathcal L/2\). On the stated scale, eventually
\(\chi\le1/4\), uniformly over the cube in
\eqref{eq:f15v31:window}. Moreover
\begin{equation}
 \epsilon_\theta(v)
    \le\Delta+8t^4\tau_\beta^2+\frac18\sqrt{q_0}
    =O\!\left(t^{-(12\nu+168)}\right).
 \label{eq:f15v31:rate}
\end{equation}
Thus the pressure, both retained-output entropies and
both absolute source-log errors, each divided by \(n\),
tend to zero uniformly over every fitting finite
separated family. Total output TV tends to zero whenever
\(n\epsilon_\theta(v)\to0\); a sufficient uniform
condition is \(n=o(t^{12\nu+168})\).
\end{theorem}

\begin{proof}
The uncoloured V20 estimate gives eventually
\begin{equation}
                   \mathcal L\ge\frac{c_*}{2}
                                         \frac{\beta}{t^d}.
 \label{eq:f15v31:raw-rate}
\end{equation}
With \(B_p=48\beta p^4\) and \(p\le t\),
\[
 K=96\tau_\beta\beta p^4
       \le\frac{3c_*}{128}\frac{\beta}{t^d}
       \le\frac3{64}\mathcal L.
\]
For \(\mathcal L\ge64\),
\(c_K\le1+K\le\mathcal L/16\). Consequently
\(\chi=4c_K/\mathcal L\le1/4\), as required.
Also
\[
 m(\mathcal L/2)
       =\log(1+\sqrt{q_0}-q_0)\le\sqrt{q_0}.
\]
For \(\tau_\beta\le1/8\),
\(J(2\tau_\beta)\le2\tau_\beta^2/(1-2\tau_\beta)\).
Using \(k\le5\) and \(n_p\le4p^4\) therefore yields
\[
            h_\theta\le h_{\tau_\beta}
                               \le64p^4\tau_\beta^2.
\]
Insert these bounds and \(c_K/v=\chi/2\le1/8\)
in \eqref{eq:f15v31:budgets} to obtain the first
inequality in \eqref{eq:f15v31:rate}. The middle term
has the explicit upper bound
\[
         8t^4\tau_\beta^2
                     \le c_*^2 2^{-21}t^{-(12\nu+168)}.
\]
The inherited \(\Delta\) and \(\sqrt{q_0}\) are
superalgebraically small in \(\beta\) on this scale,
so they are smaller than the stated inverse logarithmic
power. With \(\chi\le1/4\), all intensive constants
in Theorem~\ref{f15v31:normalization} and
Proposition~\ref{f15v31:log} are uniformly bounded.
The final TV assertion follows from its explicit bound.
\end{proof}

The new rate is an inverse power of \(\log\beta\),
not superalgebraic in \(\beta\). This slower guaranteed
rate accompanies a substantially larger source window.
V25's condition \(\ell=O(t^6)\), when imposed on these
physical weights through \(\ell=2\tau B_p\), permits
only \(\tau=O(t^6/(\beta p^4))\). The window
\eqref{eq:f15v31:window} is eventually larger by an
unbounded factor, even at fixed \(p\). Unlike V29, it
does not require a bounded total source strength as the
family grows.

For a fixed positive per-site \(\tau\), the present
whole-source cost \(c_K\) is not made small relative
to \(\log(1/q_0)\) by the available selected-scale
estimates. The paid regime above must not be extended
to that case by dropping \(\chi\) or differentiating
the bound outside its source window. The finite theorem
allows other \(v>2c_K\), but its \(m(v)\) term can
then be large; existence of such a \(v\) is not a
small-error theorem.

\subsection{What has been gained, and the next actual obstruction}

V27 isolated an exponential negative-error moment that
its fixed-strength argument did not pay. V31 pays that
moment in the explicit new extensive regime, using the
original joint Wilson failure estimate and the actual
source order-two budget. It also controls the reverse
output entropy by a finite, subunit-coefficient absorption
argument. No reference-law CGF or all-exterior good event
has been inserted as an assumption.

The result does not promote the chart factors to actual
single-near factors, prove an interaction norm for the
retained field, or control every source derivative.
Nor is small entropy per core small total TV at arbitrary
volume. The baseline retained correlations remain in
both laws. A genuine next improvement must reduce the
source cost charged against \(\log(1/q_0)\), or use more
of the Wilson relation between source energy, bad-patch
geometry and the reference normalizer. Repeating only
the one-source CGF cannot enlarge this paid regime.

The moment and relative-entropy manipulations above are
proved directly. For general background on moment-based
divergence comparisons, see the already cited
\href{https://arxiv.org/abs/1508.00335}
{Sason--Verd\'u, \emph{\(f\)-divergence Inequalities}}.
No Gibbs uniqueness or Dobrushin contraction theorem
is applied: its required influence bounds have not been
established for the retained Wilson field. The interacting
continuum construction and physical Hamiltonian mass gap
remain separate unproved obligations.

\subsection{Exact diagnostics and preservation}

The diagnostic is
\begin{center}\small
\path{scripts/verify_ym_f15_v31_projected_count_normalization.py}.
\end{center}
Nine strictly positive finite models have correlated
near/exterior data and conditionally independent rare
core failures. There are 84 joint failure-subset checks,
nine projected count-moment comparisons and 36 source
patterns, including both signs, alternating signs and
zero sources. They give 36 second-moment projection
checks, 336 exact local-factor output identities,
336 pointwise source-log sign comparisons, 180 fractional
H\"older comparisons, 36 pressure brackets, 36 two-sided
output entropy comparisons and 36 reference-entropy
absorption checks. There are 72 absolute source-log
bounds, nine exact zero-source outputs and four scale
coefficient checks.

All probabilities and weights are rational. Logarithms
use rational \(\operatorname{arctanh}\) enclosures;
the projected-count exponentials use a 24th-degree Taylor
sum with an explicit geometric tail bound. Both are
rounded outward to a \(2^{-120}\) grid. There is no
floating-point tolerance. The tests check the probability
interfaces using actual finite-model second moments;
they do not certify the inherited Wilson geometry, CGF,
or asymptotic thresholds. The script has 8,850 bytes
and SHA--256
\begin{center}\scriptsize\ttfamily
438EE8DF3B524DC062C278905B69F6591D37BECB3BEF2E4326B770CBAB5E02D9.
\end{center}

The V30 predecessor has 695,368 bytes and SHA--256
\begin{center}\scriptsize\ttfamily
605FCA14CF719882D671B33BB39446025E4A4DC88242DFA493163C5319E522D8.
\end{center}
Removing this terminal section and reversing the title
version recovers it exactly. V31 replaces V30 as the
sole active main note. All 14 protected sources and
73 earlier certificates are preserved. The last companion
checkpoint is V30 and the next is V35.

\begin{firewall}
V31 proves normalized physical-source chart-factor locality
for arbitrary fitting separated families on a specified
logarithmically shrinking per-site source cube. It pays
the negative-error exponential moment, both retained-output
KL directions and the actual source-log errors using joint
Wilson failures. It does not prove fixed-strength extensive
source control, natural-factor equivalence, retained-field
clustering, an interacting continuum limit, or the
Yang--Mills mass gap.
\end{firewall}

\section{V32: strip-energy witnesses and a larger paid source window}
\label{f15v32:section}

V31 controls extensive source normalization, but only when
the compact source cost fits inside the joint failure-tail
budget. We go upstream to that tail. The canonical rooted
filling has a feature not used by V20's single-cell witness:
its plaquette bases do not repeat. Consequently the inherited
arbitrary-support cell CGF can test the \emph{sum} of its
energies. This improves the negative tail exponent by one
power of the patch side, at the price of an explicitly
paid linear-in-side entropy term.

We also retune the collar using the unscaled V22 bound,
instead of retaining the old oversized collar by default.
Together these changes yield a larger extensive source
cube and, with a small margin, errors negligible relative
to the uniform quadratic CGF budget. We finally prove
that no optimization of V31's moment parameter can turn
this certificate into fixed-strength source control.
That last statement concerns this bound, not the actual
Wilson error and not impossibility of the mass gap.

All 88 earlier protected sources and certificates have
been rechecked unchanged. The supplied source-selection
V31, native stationarity V34 and exact-action bridge V24
were included in the V30 companion review; its next
scheduled repetition remains V35. No theorem from those
companions is imported here. The probability input below
is explicitly the selected full compact Wilson CGF
\eqref{f15v1:source-input}, not an assertion for an
arbitrary coupling, a conditional fibre law or the
normalized chart reference.

\subsection{The canonical filling has distinct cell bases}

Work in a nonwinding fine box with coordinate lift
\(\{0,\ldots,L\}^4\). For a vertex \(y\), let
\(\gamma_y\) take the positive coordinate steps in
order \(1,2,3,4\). For a positive edge \((y,a)\)
contained in the box, write
\[
 k_{y,a}=U(\gamma_y)U_{y,a}U(\gamma_{y+e_a})^{-1},
 \qquad m_{y,a}=\sum_{b>a}y_b\le3L.
\]
The angle \(\vartheta(g)\in[0,\pi]\) of an
\(\SU(2)\) element is its bi-invariant geodesic
distance from the identity in the convention already
used for \(X_f=\beta(1-\cos\vartheta_f)\).

\begin{proposition}[Nonrepeating ordered-tree filling]
\label{f15v32:filling}
The rooted holonomy \(k_{y,a}\) is a product of
conjugates of inverse positive plaquette holonomies.
Its plaquettes have orientations \((a,b)\), \(b>a\),
and bases
\begin{equation}
 x^{b,k}_i=
 \begin{cases}
 y_i,&i<b,\\
 k,&i=b,\\
 0,&i>b,
 \end{cases}
 \qquad 0\le k<y_b.
 \label{eq:f15v32:bases}
\end{equation}
There are exactly \(m_{y,a}\) such plaquettes, their
bases are pairwise distinct, and all lie in the box.
In particular, with \(I_{y,a}\) their set of bases,
\begin{equation}
 \vartheta(k_{y,a})^2
 \le m_{y,a}\sum_f\vartheta_f^2
 \le\frac{2\pi^2m_{y,a}}\beta
                         \sum_{x\in I_{y,a}}T_x.
 \label{eq:f15v32:energy}
\end{equation}
For \(m_{y,a}=0\), the rooted holonomy is the identity.
\end{proposition}

\begin{proof}
Append the final direction \(a\) to the ordered word
for \(\gamma_y\). Move this distinguished step left,
first past the \(y_4\) steps in direction 4 if
\(a<4\), then past direction 3 if \(a<3\), and
so on, stopping after direction \(a\). Each move
replaces the two-step segment \(b,a\) by \(a,b\).
If the common prefix has transport \(P\), the old
path transport times the inverse new path transport
is \(P U_{(x;a,b)}^{-1}P^{-1}\). The common suffix
cancels. Multiplying these successive relative
transports telescopes to
\(U(\gamma_y)U_{y,a}U(\gamma_{y+e_a})^{-1}\).
This is a group identity; no commutativity of
plaquettes is used.

During the moves past direction \(b\), the prefix
has the coordinates in \eqref{eq:f15v32:bases}, with
\(k=y_b-1,\ldots,0\). Within one strip the bases
are distinct. For strips \(b<c\), a base from strip
\(b\) has its \(b\)-coordinate strictly less than
\(y_b\), whereas one from strip \(c\) has that
coordinate equal to \(y_b\). Thus bases from different
nonempty strips also differ. The edge condition gives
\(y_a<L\), and each face base has \(k<y_b\le L\);
hence all its edges stay in the box.

Conjugation and inversion preserve the angle. The
bi-invariant triangle inequality followed by
Cauchy--Schwarz yields the first inequality in
\eqref{eq:f15v32:energy}. Also
\(1-\cos\vartheta\ge2\vartheta^2/\pi^2\) for
\(0\le\vartheta\le\pi\). The \(\epsilon=0\)
term in \eqref{f15v1:cell} gives
\(X_{(x;a,b)}\le4T_x\). Since the bases are distinct,
\[
 \sum_f\vartheta_f^2
       \le\frac{\pi^2}{2\beta}\sum_fX_f
       \le\frac{2\pi^2}\beta\sum_{x\in I_{y,a}}T_x.
\]
This proves the second inequality. If no move is needed,
the initial and final paths coincide, proving the last
assertion.
\end{proof}

The earlier V3/V4/V8 fillings and their angle inequalities
are not new results here. The additional assertion is
distinctness of the specific cell bases, followed by
using the joint cell CGF on that support. A generic
filling with repeated faces would not justify the
unit-coefficient source array used next.

\subsection{A joint energy tail with no independence hypothesis}

Retain \(R_M\ge1/2\), \(\mathbb E_\mu T_x\le5\)
and \(k_M\le5\) from \eqref{eq:f15v20:selected-input}.
For a box \(\mathcal K\), let
\(B_{\mathcal K}=\{R_{\mathcal K}>\delta\}\)
as in V20, with \(0<\delta\le\pi\). The number of
positive edges in the box is
\(N_L=4L(L+1)^3\). Define
\begin{equation}
 q_{\rm str}
  =\min\left\{1,\ N_L
     \exp\left[-\frac{\beta\delta^2}{12\pi^2L}
                         +15L\log2\right]\right\}.
 \label{eq:f15v32:q}
\end{equation}

\begin{theorem}[Full-law joint strip-energy witness bound]
\label{f15v32:joint}
For any finite family of boxes with common \(L,\delta\)
and pairwise disjoint sets of positive-plaquette bases,
\begin{equation}
 \mu\left(\bigcap_{i=1}^n B_{\mathcal K_i}\right)
                                      \le q_{\rm str}^{\,n}.
 \label{eq:f15v32:joint}
\end{equation}
The same bound holds for specified events contained in
the corresponding \(B_{\mathcal K_i}\). Neither
independence of the boxes nor disjointness of all links
supporting their cell observables is assumed.
\end{theorem}

\begin{proof}
If a rooted edge has angle greater than \(\delta\),
its filling is nonempty and Proposition~\ref{f15v32:filling}
gives
\begin{equation}
 \sum_{x\in I_{y,a}}T_x
       >\frac{\beta\delta^2}{2\pi^2m_{y,a}}
       \ge\frac{\beta\delta^2}{6\pi^2L}.
 \label{eq:f15v32:witness}
\end{equation}
Choose one such candidate edge in each box. Their base
sets are disjoint by hypothesis, and their union has
at most \(3nL\) elements. On the intersection of the
chosen edge events the sum of the cell energies on
that union exceeds \(n\beta\delta^2/(6\pi^2L)\).
Apply the inherited CGF with coefficient \(1/2\)
on the union and zero elsewhere. Exactly as in
\eqref{eq:f15v20:source-tail}, for any cell set \(I\),
\[
 \log\mathbb E_\mu e^{\frac12\sum_{x\in I}T_x}
 \le |I|\{5/2+5J(1/2)\}=5|I|\log2.
\]
Exponential Markov inequality therefore bounds the
chosen intersection by
\[
 \exp\left\{n\left[-\frac{\beta\delta^2}{12\pi^2L}
                                      +15L\log2\right]\right\}.
\]
An edge with empty filling gives an empty event and may
be omitted. There are at most \(N_L^n\) choices of
the remaining edges. The union bound and the bound
by one give \eqref{eq:f15v32:joint}. Event inclusion
proves the final extension.
\end{proof}

The price \(15L\log2\) is essential: this argument
tests up to \(3L\) cells per box, not one cell. On
the scales below it is negligible compared with the
negative exponent, but it has not been discarded at
finite regulator. V20's older bound remains valid;
one can always use the smaller of its \(q_0\) and
\(q_{\rm str}\). We use \(q_{\rm str}\) itself
in the following formulas.

Use the same mesh and colours as V20, now at any
admissible \(L=2p(r+5)\). Same-colour boxes have
disjoint base sets. In the separated family of V21,
the already proved entry inclusion
\(E_z=\Gamma_z^c\cup D_z^c\subseteq B_z\) uses
\begin{equation}
 \delta=\frac{a_p}{64rp^4},\qquad
 a_p=2^{-43}C_0^{-1}p^{-(\nu+8)}.
 \label{eq:f15v32:entry}
\end{equation}
Thus \eqref{eq:f15v31:joint-input} holds with
\(q_0=q_{\rm str}\). The count moment, common
projection and normalization theorems of V31 apply
without any change of their proofs. If one instead
uses all colours simultaneously, the old colour
argument gives \(q_{\rm str}^{1/4096}\); that
extra loss is not charged to the separated-family
result below.

\subsection{Retune the collar before spending the improved tail}

Fix an accuracy exponent \(s>0\), and define
\begin{equation}
 \begin{gathered}
 t=\log\beta,\qquad 2\le p\le t\quad(p\text{ dyadic}),\\
 A=2^{51}C_0(s+2),\qquad
 r=\max\{368,\lceil A p^{\nu+8}t\rceil\},\qquad
 L=2p(r+5),\\
 b=\frac{2^{-105}}{3\pi^2C_0^2A^3},\qquad d=5\nu+49.
 \end{gathered}
 \label{eq:f15v32:scale}
\end{equation}
The earlier constants \(C_0,\nu,h_p,a_p,\alpha_p\)
are unchanged. The collar and the associated chart
reference kernels are retuned; the selected coupling,
full Wilson law, physical source centering and blocking
at depth \(p\) are not. Only the \emph{unscaled}
geometry and chart estimates from V20--V22 are used,
not their old \(t^6\)-collar specializations.

\begin{proposition}[Explicit collar and failure budgets]
\label{f15v32:budgets}
Uniformly for the indicated \(p\), eventually all
earlier chart and embedding thresholds hold, and
\begin{align}
 \Delta
 &\le2^{-5}A^8p^{8\nu+66}t^8\beta^{-s-1}
       \le\beta^{-s},
                                      \label{eq:f15v32:Delta}\\
 \mathcal L:=\log(1/q_{\rm str})
 &\ge\frac b2\frac{\beta}{p^d t^3}.
                                      \label{eq:f15v32:tail-rate}
\end{align}
In particular \(q_{\rm str}<1\), \(\mathcal L\to\infty\),
and \(q_{\rm str}\) is superalgebraically small in
\(\beta\), uniformly on this scale.
\end{proposition}

\begin{proof}
Eventually the maximum in \eqref{eq:f15v32:scale}
is inactive and
\(A p^{\nu+8}t\le r\le2A p^{\nu+8}t\).
The conditions \(r\ge368\), \(p\le t\) and
\(\beta\ge\beta_p=2^{43}C_0p^{\nu+8}\) then
hold uniformly. Also \(M\ge e^\beta\) on the
selected coupling sequence, whereas \(L\) is at
most a fixed power of \(t\); hence \(M\ge4L\).
The dyadic mesh and source/core containment assertions
of V20--V22 therefore remain available.

Since \(r/2-92\ge r/4\) and
\(\alpha_p=2^{-49}C_0^{-1}p^{-(\nu+9)}\),
\[
 \alpha_p(r/2-92)p\ge2^{-51}C_0^{-1}A t=(s+2)t.
\]
Insert this and the upper bound on \(r\) into
\eqref{eq:f15v22:Delta-majorant}. This proves the
first bound in \eqref{eq:f15v32:Delta}. Its remaining
polynomial factor is \(o(\beta)\), uniformly for
\(p\le t\), which proves the second.

For the tail let \(E=\beta\delta^2/(12\pi^2L)\).
Because \(L\le4rp\),
\begin{equation}
 E\ge\frac{\beta a_p^2}{3\cdot2^{16}\pi^2r^3p^9}
   \ge b\frac{\beta}{p^{5\nu+49}t^3}.
 \label{eq:f15v32:exponent}
\end{equation}
The second inequality uses the \(a_p\) in
\eqref{eq:f15v32:entry} and \(r\le2A p^{\nu+8}t\).
Meanwhile
\[
 L\le8A p^{\nu+9}t,\qquad
 N_L\le32L^4\le2^{17}A^4p^{4\nu+36}t^4.
\]
It follows that
\(15L\log2+\log N_L\) is eventually at most
\(b\beta/(2p^d t^3)\), uniformly for \(p\le t\):
its numerator is polynomial in \(t\), whereas the
right side is bounded below by a positive constant
times \(e^t/t^{d+3}\). Thus the cap in
\eqref{eq:f15v32:q} is inactive, and subtracting this
cost from \(E\) proves \eqref{eq:f15v32:tail-rate}.
The last assertion follows since
\(\beta/t^{d+3}\) grows faster than \(\log\beta\).
\end{proof}

The constants here are conservative and the thresholds
can be extremely large. The assertion is asymptotic,
not a practical numerical range for a simulation.
In contrast with the old collar, this chosen \(\Delta\)
budget is \(O(\beta^{-s})\) for a prescribed \(s\);
we have not claimed simultaneous superalgebraic decay
of \(\Delta\) at one fixed value of \(A\).

\subsection{A larger extensive physical-source cube}

Let \(0<\omega_\beta\le1\) be a margin parameter,
and put
\begin{equation}
 \tau=\min\left\{\frac18,
       \frac{\omega_\beta b\,2^{-12}}{p^{d+4}t^3}\right\},
 \qquad |\theta_z|\le\tau,
 \qquad v=\mathcal L/2.
 \label{eq:f15v32:window}
\end{equation}
There is no bound on \(\sum_z|\theta_z|\). Use
exactly V31's actual source law and unfloored chart
factors, with the present collar. In particular the
reference still retains the full correlated baseline
output marginal; it is not a product measure on the
retained field. Keep
\(K=96\tau\beta p^4\), \(c_K=K/(1-e^{-K})\),
and the budgets \(\chi,\epsilon_\theta(v)\) of
\eqref{eq:f15v31:budgets}, using \(q_{\rm str}\).

\begin{theorem}[Normalization on the improved source window]
\label{f15v32:normalization}
Eventually, uniformly over every fitting separated
family and its source cube,
\begin{equation}
 \chi\le\frac4{\mathcal L}+\frac{3\omega_\beta}{16}
                      \le\frac14,
 \label{eq:f15v32:margin}
\end{equation}
and
\begin{equation}
 \epsilon_\theta(v)
 \le\beta^{-s}
    +\left(\frac2{\mathcal L}+\frac{3\omega_\beta}{32}\right)
                    \left(64p^4\tau^2+\sqrt{q_{\rm str}}\right).
 \label{eq:f15v32:error}
\end{equation}
Thus all pressure, two-sided retained-output KL and
two-sided absolute source-log bounds of V31 hold with
this error. In particular, when \(\omega_\beta=1\),
\begin{equation}
 \epsilon_\theta(v)
          =O\!\left(p^{-(10\nu+102)}t^{-6}\right).
 \label{eq:f15v32:absolute-rate}
\end{equation}
The per-core bounds apply to arbitrarily many fitting
sources; total output TV additionally requires
\(n\epsilon_\theta(v)\to0\), as before.
\end{theorem}

\begin{proof}
The cap in \eqref{eq:f15v32:window} is eventually
inactive. The improved tail gives
\[
 K\le\frac{3\omega_\beta b}{128}
                     \frac\beta{p^d t^3}
          \le\frac{3\omega_\beta}{64}\mathcal L.
\]
Since \(c_K\le1+K\) and eventually
\(\mathcal L\ge64\),
\[
 \chi=\frac{4c_K}{\mathcal L}
       \le\frac4{\mathcal L}+\frac{3\omega_\beta}{16}
       \le\frac14,\qquad
 \frac{c_K}v\le\frac2{\mathcal L}
                                     +\frac{3\omega_\beta}{32}.
\]
In particular \(v>2c_K\), as required by V31.
The source condition \(2\tau<R_{\rm src}\) follows
from \(\tau\le1/8\) and \(R_{\rm src}\ge1/2\).
V31's unchanged actual order-two bound and count
evaluation give
\[
 h_\theta\le64p^4\tau^2,\qquad
 m(v)=\log(1+\sqrt{q_{\rm str}}-q_{\rm str})
                                      \le\sqrt{q_{\rm str}}.
\]
Substitution into \eqref{eq:f15v31:budgets} proves
\eqref{eq:f15v32:error}; all normalization consequences
follow from the already proved V31 theorem, not a new
assumption about the reference law.

For \(\omega_\beta=1\), the coefficient is at most
\(1/8\), while
\[
 p^4\tau^2=b^2 2^{-24}p^{-(10\nu+102)}t^{-6}.
\]
The terms \(\beta^{-s}\) and \(\sqrt{q_{\rm str}}\)
are negligible relative to this inverse-logarithmic
quantity, uniformly for \(p\le t\). This proves
\eqref{eq:f15v32:absolute-rate}. The TV qualification
is exactly \eqref{eq:f15v31:TV}.
\end{proof}

At fixed \(p\) the new maximal window is of order
\(t^{-3}\). At the largest allowed depth its formula
is of order \(t^{-(5\nu+56)}\), rather than V31's
uniform \(t^{-(6\nu+86)}\) window. This improvement
uses both the new energy witness and the collar
retuning. It does not make the window independent of
the regulator.

\begin{proposition}[Error smaller than the quadratic source budget]
\label{f15v32:relative}
Suppose \(\omega_\beta\to0\) and, for some fixed
\(B>0\), \(\omega_\beta\ge t^{-B}\) eventually.
Then uniformly for \(p\le t\) and all sources in
the cube,
\begin{equation}
 \frac{\epsilon_\theta(v)}{p^4\tau^2}
       \le o(1)+\frac{128}{\mathcal L}+6\omega_\beta
                      \longrightarrow0.
 \label{eq:f15v32:relative}
\end{equation}
For example, \(\omega_\beta=t^{-1}\) gives
\begin{equation}
 \tau=b2^{-12}p^{-(5\nu+53)}t^{-4},\qquad
 \epsilon_\theta(v)
          =O\!\left(p^{-(10\nu+102)}t^{-9}\right)
 \label{eq:f15v32:relative-example}
\end{equation}
eventually, with a uniform relative error \(O(t^{-1})\).
\end{proposition}

\begin{proof}
The uncapped formula yields
\[
 p^4\tau^2=b^2 2^{-24}\omega_\beta^2
                           p^{-(10\nu+102)}t^{-6}
       \ge b^2 2^{-24}t^{-(2B+10\nu+108)}.
\]
Consequently \(\beta^{-s}/(p^4\tau^2)\) and
\(\sqrt{q_{\rm str}}/(p^4\tau^2)\) vanish
uniformly. Divide \eqref{eq:f15v32:error} by
\(p^4\tau^2\). The remaining coefficient of 64 is
\(128/\mathcal L+6\omega_\beta\), proving
\eqref{eq:f15v32:relative}. For \(\omega_\beta=t^{-1}\),
the negligible terms and \(1/\mathcal L\) are also
\(o(t^{-1})\) after the indicated division, proving
\eqref{eq:f15v32:relative-example}.
\end{proof}

The denominator here is the \emph{uniform upper CGF-budget
scale}, not a proved positive lower bound for the actual
response in each direction. In particular a zero source
still has zero exact error. These finite-source inequalities
do not justify differentiating the error twice at zero,
comparing all susceptibility matrices, or extending to
complex sources. Such statements require additional
uniform derivative or analytic estimates.

\subsection{The whole compact source cost still prevents fixed strength}

The following obstruction removes a possible circular
next step: searching indefinitely for a better \(v\)
inside the same V31 inequality. Unlike the preceding
application, \(r\) need not be the particular collar
in \eqref{eq:f15v32:scale}.

\begin{proposition}[A barrier for this numerical certificate]
\label{f15v32:barrier}
Fix \(p\ge2\) and \(\tau>0\), with
\(\tau\le1/8\). Let \(\beta\to\infty\) and
\(r\to\infty\), with the same \(L,\delta\) as
above and all needed finite-regulator hypotheses.
Use \(q_{\rm str}\) from \eqref{eq:f15v32:q}.
For every \(v>2c_K\), the V31 error expression
satisfies eventually
\begin{equation}
 \epsilon_\theta(v)
    \ge\Delta+c_K-\tfrac12\log(1/q_{\rm str})
    \ge\tfrac12c_K\ge48\tau\beta p^4.
 \label{eq:f15v32:barrier}
\end{equation}
This is a lower bound on the \emph{upper-bound formula}
\(\epsilon_\theta(v)\), uniformly over its permitted
\(v\). It is not a lower bound on a pressure difference,
relative entropy or source-log error.
\end{proposition}

\begin{proof}
Set \(\mathcal L=\log(1/q_{\rm str})\ge0\).
The definition with its cap implies
\[
 \mathcal L
   =\max\{0,E-15L\log2-\log N_L\}\le E
   \le\frac{\beta a_p^2}{3\cdot2^{15}\pi^2r^3p^9}
   =o(\beta),
\]
where the last bound uses \(L\ge2rp\) and fixed
\(p\). Also \(c_K\ge K=96\tau\beta p^4\),
so \(\mathcal L\le c_K\) eventually. For every
\(v>2c_K\),
\[
 m(v)=\log(1-q_{\rm str}+q_{\rm str}e^v)
                                  \ge v-\mathcal L.
\]
Since \(h_\theta\ge0\), the exact V31 budget gives
\[
 \epsilon_\theta(v)
 \ge\Delta+\frac{c_K}{v}(v-\mathcal L)
 \ge\Delta+c_K-\frac{\mathcal L}{2}.
\]
This proves \eqref{eq:f15v32:barrier}. The argument
also covers \(q_{\rm str}=1\), for which
\(\mathcal L=0\) and \(m(v)=v\); the numerical
expression is defined even though it cannot then
provide a small-error instance of V31.
\end{proof}

Taking the minimum with V20's old bound does not remove
this barrier. Its logarithmic budget is at most
\(\beta a_p^2/(9\cdot2^{16}\pi^2r^4p^{10})\),
again \(o(\beta)\) for fixed \(p\) and diverging
\(r\). The logarithm of the inverse minimum is the
maximum of these two logarithms, still \(o(\beta)\).
This does not exclude a \emph{new} substantially
stronger failure estimate, nor a smaller charged
source cost; it closes only the stated optimization
within the two existing certificates.

The next useful direction is therefore to connect
the actual source energy and the exceptional geometry
inside the normalization estimate, or obtain another
Wilson-specific control of the reference normalizer.
The bounded-oscillation replacement by \(K\sim\beta\)
throws away that relation. V30 already excludes using
only an information summary and an unconditional
source CGF to repair this step. Merely adjusting the
collar constants or the moment parameter once more
cannot establish fixed-strength control by the bound
just analyzed.

\subsection{Exact diagnostics, preservation and claim boundary}

The diagnostic is
\begin{center}\small
\path{scripts/verify_ym_f15_v32_strip_energy.py}.
\end{center}
For all positive edges of four-dimensional boxes of
side 1 through 5, it checks 7,336 rooted free-group
identities and distinct-base descriptions, including
2,269 empty-filling identities. The intermediate
24,276 square exchanges are checked individually as
conjugated inverse positive faces. Free-group reduction
tests noncommutative word identities, not just traces
or an abelian specialization.

There are 144 exact rational energy/Cauchy checks,
24 joint energy-moment checks in positive finite
shared-latent models, 48 joint Markov bounds, 12
strict nonindependence checks and 35 union coefficients.
The models deliberately have correlated boxes under
their full law; their easy moment bounds test the
probability interface, not the Wilson CGF. A further
263 scale coefficient/exponent identities, nine
absorption margins and 27 whole-budget barrier
comparisons are exact. The script uses only integers,
rational arithmetic and symbolic free-group words.
No finite test establishes an asymptotic threshold,
an interacting continuum theory or a mass gap.

The script has 8,274 bytes and SHA--256
\begin{center}\scriptsize\ttfamily
EAFD0229FC758DBD6A3D9BB755F0FDE990B22B5CC26F18A6701608283C52604E.
\end{center}
The V31 predecessor has 715,542 bytes and SHA--256
\begin{center}\scriptsize\ttfamily
A24D810E7FEFDE9BA8254B6C3614069F573F9662C98F8A94FFA8BE7865E90D67.
\end{center}
Removing this terminal insertion and reversing the
title version recovers V31 exactly. V32 replaces it
as the sole active main note. All 14 protected sources
and 74 earlier certificates are preserved. The next
companion review remains V35.

The full target remains the nontrivial continuum
construction and physical Hamiltonian gap in the
\href{https://www.claymath.org/millennium/yang-mills-the-maths-gap/}
{Clay Yang--Mills problem}. The local selected-coupling
estimates in this section are not that target.

\begin{firewall}
V32 proves a stronger joint failure estimate from
nonrepeating canonical filling energies, conditional
on the already recorded selected Wilson source input.
It gives a larger rigorously paid extensive source
window and an error small relative to its quadratic
CGF budget with a shrinking margin. It also proves
why these particular certificates cannot reach fixed
source strength by tuning the moment parameter.
It does not prove fixed-strength normalization,
natural-factor equivalence, retained clustering,
an interacting continuum construction or the full
Yang--Mills mass gap.
\end{firewall}

\section{V33: energy-weighted confidence and one-sided fixed-strength pressure}
\label{f15v33:section}

V32 proves that optimizing the moment parameter cannot
repair a bound that continues to charge the full compact
source range. We now replace that charge by an actual
conditional energy moment, for one sign of the log error
on each fixed-sign source cube. The new ingredient is a
joint estimate for energy-weighted conditional-confidence
failures before the common exterior is integrated out.
It uses the matched products of V21, not an unproved
product property after conditioning on the retained
field alone.

The outcome is an exponential moment of the favourable
source-log sign and the corresponding one-sided pressure
bound, for every fitting separated family. At fixed
source strength its intensive cost is cubic in that
strength, with no growing factor from the compact
energy range. This cost does not tend to zero at a
fixed nonzero strength merely by removing the regulator.
We also identify, and test against the existing V27
example, the reference-energy moment missing from
the complementary sign. Neither a two-sided pressure
theorem nor a mass-gap theorem is claimed.

The V32 predecessor and all 89 protected sources and
certificates were checked unchanged. The last companion
review is V30, including the supplied SM source-selection
V31, native stationarity V34 and exact-action bridge V24;
the next remains V35. No new companion input is assumed.
Unlike V17--V18's negative-source normalization theorem,
the present result has no logarithmic bound on the
family size, but is one-sided and intensive. It is
not a restatement of their two-sided small-family result.

\subsection{Uncentred positive energies, with the original normalizers retained}

Use the separated geometry, full selected Wilson law,
physical sources and unfloored chart references from
V26--V32. In particular the original source input has
\(R_{\rm src}\ge1/2\), \(k\le5\),
\(\overline T\le5\). Put
\begin{equation}
 Y_z=\sum_{x\in\mathcal B_{p,z}}T_x\ge0,\qquad
 S_z=Y_z-\mathbb E_\mu Y_z,\qquad
 \theta_z=\sigma t_z,\quad 0\le t_z\le\tau\le\tfrac18,
 \quad \sigma\in\{-1,+1\}.
 \label{eq:f15v33:source}
\end{equation}
All nonzero source coefficients have the same sign in
this section. The fine cell supports are disjoint and
have size \(n_p\le4p^4\); the source is measurable
from its own core and near retained data.

Let \(\xi=(C,v_O)\) be the full common exterior and
\(\nu_0=\xi_*\mu\), \(\nu_\theta=\xi_*\mu_\theta\).
Write \(\mathsf P_z^\xi\) for the baseline conditional
law of the \emph{whole} interior \(J_z\). Its marginal
on the core is the earlier \(\kappa_{z,0}^\xi\).
The interior laws are independent given \(\xi\),
and each failure \(E_z=\Gamma_z^c\cup D_z^c\) depends
only on \(\xi\) and its own interior. Set
\begin{equation}
 \begin{aligned}
 u_z(\xi)&=\mathsf P_z^\xi(e^{\theta_zY_z}),&
 b_z(N_z)&=\lambda_z(N_z)(e^{\theta_zY_z}),\\
 L_z(\xi)&=\log(u_z/b_z),&
 g_z(\xi)&=\log\mathsf P_z^\xi(e^{t_zY_z})\ge0,\\
 \pi_{z,0}(\xi)&=\mathsf P_z^\xi(E_z),&
 \pi_{z,\theta}(\xi)
   &=\frac{\mathsf P_z^\xi(1_{E_z}e^{\theta_zY_z})}{u_z}.
 \end{aligned}
 \label{eq:f15v33:local}
\end{equation}
The last formula remains valid on \(\Gamma_z^c\),
where the indicator is identically one in the interior.
All these are actual, normalized compact conditionals.
Remote source factors cancel only after fixing \(\xi\).

Using uncentred weights in \eqref{eq:f15v33:local}
does not change \(L_z\): the same deterministic
centering scalar multiplies \(u_z\) and \(b_z\).
Likewise it changes both global normalizers by the
same scalar. Thus V31's quantities
\(A_\theta,T_\theta,Z_a,Z_b,r_\theta,\delta_\theta\)
can still be used with
\begin{equation}
 r_\theta=\log(A_\theta/T_\theta),\qquad
 \delta_\theta=\log(Z_b/Z_a),\qquad
 \delta_\theta=\log\mathbb E_{\rho_a}e^{-r_\theta}.
 \label{eq:f15v33:normalization}
\end{equation}
The actual tilted law and reference output are unchanged.
Keep the exact source cost
\begin{equation}
 h_\theta=\frac1n\log\mathbb E_\mu
                   \left(\frac{d\mu_\theta}{d\mu}\right)^2
            \le64p^4\tau^2.
 \label{eq:f15v33:h}
\end{equation}
This is the existing actual-law estimate, not a CGF
assertion for the chart reference.

\subsection{Matched products control the actual conditional energy moments}

Let \(q\in(0,1)\) be a proved raw failure parameter:
\(\mu(\cap_{z\in I}E_z)\le q^{|I|}\) for every
subfamily. One may take V32's \(q_{\rm str}\) when
it is below one. No all-exterior good event is imposed.

\begin{proposition}[Joint conditional energy and tilted confidence]
\label{f15v33:matched-energy}
For \(a\ge1\) with \(a\tau\le R_{\rm src}<1\),
\begin{equation}
 \mathbb E_{\nu_0}\exp\left(a\sum_{z\in I}g_z\right)
                         \le(1-a\tau)^{-5n_p|I|}.
 \label{eq:f15v33:energy-moment}
\end{equation}
Moreover
\begin{equation}
 \begin{aligned}
 \mathbb E_{\nu_0}\prod_{z\in I}\pi_{z,0}
       &\le q^{|I|},\\
 \mathbb E_{\nu_0}\prod_{z\in I}\pi_{z,\theta}
       &\le\gamma^{|I|},\qquad
 \gamma=q^{1/2}(1-2\tau)^{-5n_p/2}.
 \end{aligned}
 \label{eq:f15v33:confidence-moments}
\end{equation}
The second bound holds for either of the two uniform
sign choices in \eqref{eq:f15v33:source}. Its expectation
is under the \emph{baseline} common-exterior law
\(\nu_0\), even though the probabilities inside it
are source-tilted conditionals.
\end{proposition}

\begin{proof}
Conditional Jensen and the exact product conditional
law give
\begin{align*}
 \mathbb E_{\nu_0}\prod_{z\in I}
                  [\mathsf P_z^\xi(e^{t_zY_z})]^a
 &\le\mathbb E_{\nu_0}\prod_{z\in I}
                         \mathsf P_z^\xi(e^{at_zY_z})\\
 &=\mathbb E_\mu e^{\sum_{z\in I}at_zY_z}.
\end{align*}
The disjoint fine cell supports and the inherited CGF
bound imply, for any \(0\le u<1\) in its radius,
\[
 u\overline T+kJ(u)\le5u+5[-\log(1-u)-u]
                                     =-5\log(1-u).
\]
Apply this cell by cell with \(u=at_z\le a\tau\)
to obtain \eqref{eq:f15v33:energy-moment}. The first
confidence bound is V21's matched-product identity.

For positive sources, \(u_z\ge1\), so
\(\pi_{z,\theta}\le
\mathsf P_z^\xi(1_{E_z}e^{t_zY_z})\). Hence matching
the interior products and then Cauchy--Schwarz give
\begin{align*}
 \mathbb E_{\nu_0}\prod_{z\in I}\pi_{z,\theta}
 &\le\mathbb E_\mu\left[1_{\cap_{z\in I}E_z}
                                      e^{\sum_{z\in I}t_zY_z}\right]\\
 &\le q^{|I|/2}
          \left(\mathbb E_\mu e^{2\sum_{z\in I}t_zY_z}\right)^{1/2}
 \le\gamma^{|I|}.
\end{align*}
For negative sources, Cauchy--Schwarz inside the actual
conditional law yields
\[
 [\mathsf P_z^\xi(e^{-t_zY_z})]^{-1}
           \le\mathsf P_z^\xi(e^{t_zY_z})=e^{g_z}.
\]
Because the numerator of \(\pi_{z,\theta}\) is at
most \(\pi_{z,0}\), one has
\(\pi_{z,\theta}\le\pi_{z,0}e^{g_z}\). Therefore
\[
 \mathbb E_{\nu_0}\prod_{z\in I}\pi_{z,\theta}
 \le\left(\mathbb E_{\nu_0}\prod_{z\in I}\pi_{z,0}^{2}\right)^{1/2}
       \left(\mathbb E_{\nu_0}e^{2\sum_{z\in I}g_z}\right)^{1/2}
 \le q^{|I|/2}(1-2\tau)^{-5n_p|I|/2}.
\]
The last line uses \(0\le\pi_{z,0}\le1\), its
matched product bound, and \eqref{eq:f15v33:energy-moment}
at \(a=2\). This proves both cases without asserting
a source CGF conditioned on \(\xi\).
\end{proof}

The product identity is needed \emph{before} discarding
the exterior details. Replacing \(\mathsf P_z^\xi\)
by independently selected near-output kernels would
not justify this proof. No independence of \(C\),
of \(v_O\), or of the unconditioned source energies
has been asserted.

\subsection{The favourable local log sign has an energy cap}

For a fixed sign choose
\begin{equation}
 \begin{array}{c|cc}
 \sigma&\pi_z^{(\sigma)}&\gamma_\sigma\\ \hline
 -1&\pi_{z,0}&q\\
 +1&\pi_{z,\theta}&\gamma
 \end{array}
 \qquad
 B_z^{(\sigma)}=\{\pi_z^{(\sigma)}>\eta\},\quad0<\eta<1.
 \label{eq:f15v33:threshold}
\end{equation}
The superscript here denotes which confidence test is
used, not a change of the exterior law. In particular
\begin{equation}
 \nu_0\left(\bigcap_{z\in I}B_z^{(\sigma)}\right)
                                  \le(\gamma_\sigma/\eta)^{|I|}.
 \label{eq:f15v33:threshold-tail}
\end{equation}
This follows directly by bounding the intersection
indicator by \(\eta^{-|I|}\prod_{z\in I}\pi_z^{(\sigma)}\).

\begin{proposition}[Actual-energy cap on one sign]
\label{f15v33:local-cap}
Let \(a_\eta=-\log(1-\eta)\). Then almost surely,
\begin{equation}
 \begin{aligned}
 (L_z)_-&\le g_z,\qquad
 (L_z)_-\le\Delta+a_\eta+1_{B_z^{(-)}}g_z
                                      &&(\sigma=-1),\\
 (L_z)_+&\le g_z,\qquad
 (L_z)_+\le\Delta+a_\eta+1_{B_z^{(+)}}g_z
                                      &&(\sigma=+1).
 \end{aligned}
 \label{eq:f15v33:local-cap}
\end{equation}
These inequalities contain no \(B_p\), \(K\) or
\(c_K\) compact-range cost.
\end{proposition}

\begin{proof}
For negative sources, \(b_z\le1\), so
\[
 -L_z\le-\log u_z
      \le\log\mathsf P_z^\xi(e^{t_zY_z})=g_z.
\]
The second inequality was proved by conditional
Cauchy--Schwarz above. Since \(g_z\ge0\), this
also bounds the negative part. For positive sources,
\(b_z\ge1\), and
\(L_z\le\log u_z=g_z\), giving the positive-part cap.
These facts hold on the Haar fallback branch as well,
because \(Y_z\) is nonnegative on the full compact
coordinate space.

If the corresponding \(\pi_z^{(\sigma)}\le\eta<1\),
then \(\Gamma_z\) must hold. The exact chart-probability
identity in V27, before its clipped-log replacement,
gives
\[
 (L_z)_-\le\Delta-\log(1-\pi_{z,0}),\qquad
 (L_z)_+\le\Delta-\log(1-\pi_{z,\theta}).
\]
Use the appropriate one to obtain \(\Delta+a_\eta\)
on the confidence-good event. On its complement the
global energy cap just proved applies. Adding the
nonnegative floor proves \eqref{eq:f15v33:local-cap}.
\end{proof}

\subsection{A joint exponential bound for the exceptional energy}

Set
\begin{equation}
 \begin{gathered}
 v=\frac1{4\tau}\ge2,\qquad P_p=2^{5n_p/2},\qquad
 F_\sigma(\xi)=\sum_z1_{B_z^{(\sigma)}}g_z(\xi),\\
 m_\sigma(\eta)=\log\left(1+P_p\sqrt{\gamma_\sigma/\eta}\right).
 \end{gathered}
 \label{eq:f15v33:energy-budget}
\end{equation}
This \(v\) is used with the energy cap; it is not
required to exceed \(2c_K\) as in the different
V31 certificate.

\begin{theorem}[Exponential energy-weighted confidence]
\label{f15v33:weighted-confidence}
Under the baseline common-exterior law,
\begin{equation}
                 \log\mathbb E_{\nu_0}e^{vF_\sigma}
                                      \le n m_\sigma(\eta).
 \label{eq:f15v33:energy-Laplace}
\end{equation}
There is no family-size restriction beyond fitting the
separated geometry.
\end{theorem}

\begin{proof}
For each subfamily \(I\), Cauchy--Schwarz,
\eqref{eq:f15v33:threshold-tail}, and the conditional
energy estimate at \(a=2v\) give
\begin{align*}
 \mathbb E_{\nu_0}
       \left[1_{\cap_{z\in I}B_z^{(\sigma)}}
                                        e^{v\sum_{z\in I}g_z}\right]
 &\le(\gamma_\sigma/\eta)^{|I|/2}
       \left(\mathbb E_{\nu_0}e^{2v\sum_{z\in I}g_z}\right)^{1/2}\\
 &\le\left(P_p\sqrt{\gamma_\sigma/\eta}\right)^{|I|}.
\end{align*}
Here \(2v\tau=1/2\le R_{\rm src}\), so the
moment input is within its proved radius. Expand
\[
 e^{vF_\sigma}
     =\prod_z\left[1+1_{B_z^{(\sigma)}}(e^{vg_z}-1)\right].
\]
Each term is nonnegative. Bound its nonempty-subfamily
contribution by the preceding display, using
\(e^{vg_z}-1\le e^{vg_z}\). The binomial sum is
\((1+P_p\sqrt{\gamma_\sigma/\eta})^n\), proving
the theorem. This is a bound on energy attached to
confidence failures, not merely on their count.
\end{proof}

Both the confidence rarity and the energy moment in
this proof are integrated under \(\nu_0\). One
cannot assign the same small moment to the chart
reference by changing that measure's name. The
following transport keeps the actual source likelihood
explicit.

\subsection{Actual output exponential moments and one-sided pressure}

Define the favourable retained error by
\[
 R_\sigma(C)=
 \begin{cases}(r_\theta(C))_-,&\sigma=-1,\\
               (r_\theta(C))_+,&\sigma=+1.
 \end{cases}
\]
For either sign let
\begin{equation}
 \varepsilon_\sigma
       =\Delta+a_\eta+\frac{h_\theta+m_\sigma(\eta)}v
       =\Delta+a_\eta+4\tau\{h_\theta+m_\sigma(\eta)\}.
 \label{eq:f15v33:epsilon}
\end{equation}

\begin{theorem}[One-sided fixed-strength normalizer control]
\label{f15v33:pressure}
For every fitting separated family and every source
vector in its specified fixed-sign cube,
\begin{equation}
             \log\mathbb E_{\rho_a}e^{R_\sigma}
                                      \le n\varepsilon_\sigma.
 \label{eq:f15v33:actual-exp}
\end{equation}
Consequently
\begin{equation}
 \begin{aligned}
 \mathbb E_{\rho_a}R_\sigma&\le n\varepsilon_\sigma,\qquad
 \rho_a\{R_\sigma\ge n(\varepsilon_\sigma+x)\}
                                      \le e^{-nx}\quad(x>0),\\
 \delta_\theta&\le n\varepsilon_- &&(\theta\in[-\tau,0]^n),\\
 \delta_\theta&\ge-n\varepsilon_+ &&(\theta\in[0,\tau]^n).
 \end{aligned}
 \label{eq:f15v33:pressure}
\end{equation}
These are one-sided bounds, not an estimate on
\(|\delta_\theta|\) or on either output KL divergence.
\end{theorem}

\begin{proof}
The exact exterior-product identities used in V27 give
\[
 r_\theta\ge\mathbb E_{\nu_0}\left[\sum_zL_z\mid C\right],
 \qquad
 r_\theta\le\mathbb E_{\nu_\theta}\left[\sum_zL_z\mid C\right].
\]
The second inequality is the actual conditional
posterior entropy identity, with its nonnegative
relative entropy subtracted. Combining with
Proposition~\ref{f15v33:local-cap} yields
\begin{equation}
 \begin{aligned}
 (r_\theta)_-&\le n(\Delta+a_\eta)
                       +\mathbb E_{\nu_0}[F_-\mid C]&& (\sigma=-1),\\
 (r_\theta)_+&\le n(\Delta+a_\eta)
                       +\mathbb E_{\nu_\theta}[F_+\mid C]&& (\sigma=+1).
 \end{aligned}
 \label{eq:f15v33:projected-signs}
\end{equation}
The two projections use different, specified laws.

For clarity, the moment transport needed here is the
three-factor H\"older argument already proved in V31.
If \(f\) is a density with \(\mathbb E f^2\le e^{nh_\theta}\),
and \(G\ge0\) satisfies \(\mathbb E e^{vG}\le e^{nm}\),
then for \(v\ge2\),
\begin{equation}
 \mathbb E[f e^G]
 =\mathbb E\left[f^{1-2/v}(f^2)^{1/v}(e^{vG})^{1/v}\right]
                  \le e^{n(h_\theta+m)/v}.
 \label{eq:f15v33:transport}
\end{equation}
At \(v=2\) this is Cauchy--Schwarz. Both actual
projected densities \(d\rho_a/d\rho_0\) and
\(d\nu_\theta/d\nu_0\) have the required second
moment bound by conditional Jensen applied to
\(d\mu_\theta/d\mu\).

For negative sources, set
\(G=\mathbb E_{\nu_0}[F_-\mid C]\).
Conditional Jensen and
Theorem~\ref{f15v33:weighted-confidence} give
\(\mathbb E_{\rho_0}e^{vG}\le e^{nm_-(\eta)}\).
Use \eqref{eq:f15v33:transport} under \(\rho_0\)
with \(f=d\rho_a/d\rho_0\), and then the first
line of \eqref{eq:f15v33:projected-signs}.
For positive sources, conditional Jensen under the
\emph{actual tilted} exterior law gives instead
\[
 \mathbb E_{\rho_a}
           e^{\mathbb E_{\nu_\theta}[F_+\mid C]}
                        \le\mathbb E_{\nu_\theta}e^{F_+}.
\]
Apply \eqref{eq:f15v33:transport} under \(\nu_0\)
with \(f=d\nu_\theta/d\nu_0\), \(G=F_+\),
and use the second line of
\eqref{eq:f15v33:projected-signs}. This proves
\eqref{eq:f15v33:actual-exp} in both cases.

Jensen and exponential Markov inequality give the mean
and tail assertions. For negative sources,
\(e^{-r_\theta}\le e^{(r_\theta)_-}\), so the exact
normalizer identity proves the upper pressure bound.
For positive sources, Jensen gives
\[
 \delta_\theta\ge-\mathbb E_{\rho_a}r_\theta
                 \ge-\mathbb E_{\rho_a}(r_\theta)_+
                 \ge-n\varepsilon_+.
\]
No density or source moment for \(\rho_b\) has been
postulated in this argument.
\end{proof}

This proof does not contradict V32's barrier: that
barrier retains the clipped cost \(c_K\), whereas
\eqref{eq:f15v33:local-cap} no longer uses it. The
new transport cost \(h_\theta/v=4\tau h_\theta\)
is still present. It must not be silently classified
as a regulator error.

\subsection{The cost is cubic, with a separate vanishing exceptional term}

Use V32's retuned collar and its fixed exponent \(s>0\),
so \(\Delta\le\beta^{-s}\), \(p\le t=\log\beta\),
and \(q=q_{\rm str}\) has the bound
\eqref{eq:f15v32:tail-rate}. Once \(\gamma_\sigma<1\),
choose
\begin{equation}
                 \eta_\sigma=\gamma_\sigma^{1/3},\qquad
                 m_\sigma(\eta_\sigma)
                              =\log(1+P_p\eta_\sigma).
 \label{eq:f15v33:balanced}
\end{equation}

\begin{proposition}[An intensive cubic source budget]
\label{f15v33:cubic}
Eventually, uniformly for \(0<\tau\le1/8\),
\(p\le t\) and either fixed-sign cube,
\begin{equation}
 \varepsilon_\sigma
 \le\beta^{-s}+256p^4\tau^3
                      +(2+4\tau P_p)\eta_\sigma.
 \label{eq:f15v33:cubic}
\end{equation}
For every fixed \(A>0\), the last term is
\(o(\beta^{-A})\), uniformly in these parameters.
For fixed \(p\) and \(\tau>0\), this gives an
intensive \(O(p^4\tau^3)\) bound, not a bound
tending to zero with \(\beta\). If instead
\(\tau=\tau_\beta\to0\) and
\(\tau_\beta\ge t^{-B}\) for some fixed \(B>0\),
then
\begin{equation}
             \frac{\varepsilon_\sigma}{p^4\tau_\beta^2}
                            \le256\tau_\beta+o(1)\longrightarrow0.
 \label{eq:f15v33:relative}
\end{equation}
No V32-type geometry-dependent upper power of \(t^{-1}\)
is imposed on \(\tau_\beta\) in this one-sided result.
\end{proposition}

\begin{proof}
Since \(2\tau\le1/4\),
\[
 \gamma_- =q,\qquad
 \gamma_+\le q^{1/2}(4/3)^{5n_p/2},\qquad
 \log P_p=\tfrac52 n_p\log2=O(t^4).
\]
The negative logarithm of \(q\) is at least a
positive constant times \(e^t/t^{d+3}\), by V32.
It dominates every fixed power of \(t\), uniformly
for \(p\le t\). Thus both \(\eta_\sigma\to0\),
eventually \(\eta_\sigma\le1/2\), and
\((2+4\tau P_p)\eta_\sigma\) has the asserted
superalgebraic decay. Use
\(a_{\eta_\sigma}\le2\eta_\sigma\),
\(\log(1+P_p\eta_\sigma)\le P_p\eta_\sigma\),
and \eqref{eq:f15v33:h} in
\eqref{eq:f15v33:epsilon} to obtain
\eqref{eq:f15v33:cubic}. Finally
\(p^4\tau_\beta^2\ge16t^{-2B}\), so division
leaves the two regulator terms negligible. The source
term becomes \(256\tau_\beta\), proving
\eqref{eq:f15v33:relative}.
\end{proof}

The denominator in \eqref{eq:f15v33:relative} is a
uniform quadratic CGF-budget scale, not a directionwise
lower bound on susceptibility. No source derivative,
two-sided pressure estimate, mixed-sign cube or total
output TV conclusion follows from this one-sided
comparison alone. Zero sources have exactly zero
source-log error and equal normalizers, even though
the uniform displayed budgets need not vanish there.

\subsection{The complementary cap involves the reference energy}

For the opposite local sign, positivity yields only
\begin{equation}
 \begin{aligned}
 (L_z)_-&\le g_z^\lambda &&(\sigma=+1),\\
 (L_z)_+&\le g_z^\lambda &&(\sigma=-1),\qquad
 g_z^\lambda=\log\lambda_z(N_z)(e^{t_zY_z}).
 \end{aligned}
 \label{eq:f15v33:reference-cap}
\end{equation}
For the first line use \(u_z\ge1\). For the second
use \(u_z\le1\) and
\([\lambda_z(e^{-t_zY_z})]^{-1}
\le\lambda_z(e^{t_zY_z})\). Unlike the actual
\(g_z\), these chart moments have no matched-product
identity under \(\nu_0\). A bound replacing
\eqref{eq:f15v33:energy-moment} for them has not been
proved. The preceding tilted-confidence estimate does
not supply it.

The already constructed V27 model tests precisely this
missing step. We reuse it rather than claim a new
counterexample family. For \(d\ge20\), let
\(q_d=2^{-d}\), \(\epsilon_d=2^{-8d}\),
\(L_d=16d\log2\), and keep its positive binary law,
chart \(D=\{X=1\}\) and reference \(\lambda_C=\delta_1\).
Its nonnegative energy is
\(Y=L_d1_{\{C=1,X=1\}}\). V27 proved the uniform
signed source CGF and the raw failure bound \(2q_d\);
independent copies give their arbitrary-support forms.

Take \(\tau=1/8\), \(v=2\) and \(a=e^{\tau L_d}=2^{2d}\).
The two relevant fourth moments satisfy
\begin{equation}
 \begin{aligned}
 \mathbb E_\mu e^{4\tau Y}
     &=1+q_d\epsilon_d(a^4-1)\le1+2^{-d},\\
 \mathbb E_{\nu_0}e^{4g}
     &=1-q_d+q_d(1-\epsilon_d+\epsilon_d a)^4
                                      \le\mathbb E_\mu e^{4\tau Y},\\
 \mathbb E_{\nu_0}e^{4g^\lambda}
     &=1-q_d+q_d a^4\ge2^{7d}.
 \end{aligned}
 \label{eq:f15v33:reference-obstruction}
\end{equation}
The middle inequality is conditional Jensen. Thus
bounded actual energy moments do not automatically
give the reference moment needed in the proof.

There is a direct sign check as well. At the negative
source \(-\tau\), put
\(u=1-\epsilon_d+\epsilon_d/a\). Then
\(r(0)=0\), \(r(1)=\log(au)>0\), and
\begin{equation}
 \mathbb E_{\rho_{-\tau}}e^{r_+}
     =\frac{1-q_d+q_d a u^2}{1-q_d+q_d u}
                           \ge2^{d-2},\qquad
 \mathbb E_{\rho_{-\tau}}e^{r_-}=1.
 \label{eq:f15v33:complement-obstruction}
\end{equation}
Indeed \(u\ge1-\epsilon_d\ge1/2\) and the denominator
is at most one. This is consistent with the new
favourable-sign theorem, but rules out its automatic
extension to the opposite exponential moment from
these probability inputs. It does \emph{not} assert
failure of negative-source pressure convergence in
this example; that pressure can be small even when
the stronger complementary exponential moment is large.
The positive-source normalization failure was already
proved in V27 and is not claimed anew.

This is not a Wilson model. A Wilson-specific control
of the near-parameter energy in the identity-exterior
chart could still exclude this mechanism. The next
upstream issue is to prove such a measured chart-energy
estimate, or avoid requiring it by a genuinely paid
normalizer comparison. The actual conditional-energy
estimate above must not simply be transferred to
\(\lambda_z\). Fixed-strength two-sided control,
baseline retained correlations and continuum physical
time calibration remain separate obligations.

\subsection{Exact diagnostics and predecessor preservation}

The diagnostic is
\begin{center}\small
\path{scripts/verify_ym_f15_v33_energy_weighted_confidence.py}.
\end{center}
Nine positive finite models have a shared retained bit,
rare exterior flags and conditionally independent cores.
There are 54 source patterns in the two uniform-sign
orthants, including partial and zero sources. The
checks include 33 raw failure subsets, 198 matched
conditional energy moments, 1,224 conditional Jensen
powers, 198 tilted-confidence products, 198 threshold
tails and 198 energy-weighted subsets. They also check
1,224 local energy caps, 1,224 favourable local sign
bounds and 108 projected sign comparisons.

There are 54 baseline exceptional-energy Laplace
bounds, 54 projected second moments, 54 fractional
H\"older comparisons, 54 actual favourable-sign
exponential moments, 54 one-sided pressure bounds and
18 exact zero-source outputs. Four instances of the
inherited V27 model check the reference-moment
separation and the complementary-sign divergence;
two exact coefficient checks include \(4\cdot64=256\).

All source weights and probabilities in the tests are
rational. Squaring removes the square root in the
final moment and pressure bounds. Intermediate log
comparisons use the preserved V31 script's outward
rational enclosures; importing that helper also runs
its regression checks. These tests certify finite
probability interfaces, not the inherited Wilson CGF,
chart geometry, asymptotic threshold or mass gap.
The script has 11,243 bytes and SHA--256
\begin{center}\scriptsize\ttfamily
19C6D2F3421CE1F21DA1DCF8BA14FDE82D57DBEAC0259D551557BAF267F238D4.
\end{center}

The V32 predecessor has 738,340 bytes and SHA--256
\begin{center}\scriptsize\ttfamily
25AB62E601F6CF8DBE2D3DF806AFD7E0C67BD16C032D3BC305369AEFA4D54104.
\end{center}
Removing this terminal section and reversing the title
version recovers it exactly. V33 replaces V32 as the
sole active main note. All 14 protected sources and
75 earlier certificates are preserved. The next
companion checkpoint remains V35.

\begin{firewall}
V33 replaces the compact source-range charge by an
actual conditional-energy moment for one log-error
sign on each uniform-sign source cube. It proves an
intensive one-sided pressure bound with cubic source
cost for arbitrary fitting families, and isolates
the reference-energy moment missing from the other
sign. It does not prove vanishing fixed-strength
two-sided normalization, mixed-sign source control,
two-sided output entropy, retained-field clustering,
an interacting continuum limit or the Yang--Mills
mass gap.
\end{firewall}

\section{V34: a remote Haar-fallback obstruction in the actual Wilson reference}
\label{f15v34:section}

V33 isolated a reference-energy moment that could close
its complementary source-log estimate. Before trying to
prove that moment, one must test the reference itself.
This section proves that the proposed uniform moment
is false for the existing Haar-fallback definition.
The obstruction occurs under the recorded full compact
Wilson law, already at the single blocking step \(p=2\).
It is not another abstract two-state counterexample.

A small change distributed around a long, remote loop
can force one required retained coordinate out of its
chart. The cost of that event is sublinear in \(\beta\)
in the negative logarithm. The reference then replaces
the entire central core by product Haar. A central
plaquette consequently has reference energy of order
\(\beta\), so a fixed positive source overwhelms the
rarity of the remote event. On the retuned V32 collar,
the normalized reference and actual retained source
laws even become asymptotically singular in total
variation. The original Wilson measure is not changed.

All assertions retain the inherited completed-block
inverse, selected-law CGF and chart estimates. Those
inputs are cited explicitly, not recertified as a full
continuum construction. The 14 protected sources and
76 earlier certificates were checked unchanged. The
supplied SM source-selection V31, native stationarity
V34 and exact-action bridge V24 were included in the
V30 review; the next companion checkpoint remains V35.

\subsection{Fix the existing reference and a single physical source}

Take \(p=2\), a source centred at the origin with
orientation \((1,2)\), and the selected full Wilson
law \(\mu=\mu^{\mathrm W}_{M,\beta}\),
\(\beta=\beta_M^\diamond\). Its source-cell set is
\begin{equation}
 \mathcal B_2=\{0,1,2\}^2\times\{0,1\}^2,
 \qquad n_2=36,\qquad
 Y=\sum_{x\in\mathcal B_2}T_x,\qquad
 S=Y-\mathbb E_\mu Y.
 \label{eq:f15v34:source}
\end{equation}
Thus \(Y\ge0\) and \(\mathbb E_\mu Y\le180\)
by \eqref{eq:f15v20:selected-input}. Fix
\(0<\theta\le1/8\). Allow an integer collar with
\begin{equation}
 r\ge368,\qquad r\longrightarrow\infty,\qquad
 r=O((\log\beta)^b),\quad b<\infty\text{ fixed},
 \qquad L=4(r+5),\qquad M\ge4L.
 \label{eq:f15v34:scales}
\end{equation}
All patches are embedded without periodic wrap-around.
The selected path has \(\beta\to\infty\), and its
recorded bound \(\beta\le\log M\) accommodates these
collars eventually. The fixed chart radius \(a_2>0\)
and all primitive constants are unchanged.

Use exactly V22's unfloored \(\lambda(N)\): the
measured chart core marginal at identity exterior when
every required retained group in the \(H_N\)-frame
lies in its half chart, and normalized product Haar
on the core otherwise. Denote this latter retained
event by \(\mathcal F_\beta\). In particular
\(\mathcal F_\beta\) depends only on \(N\), not
on the actual details or the source parameter. Put
\begin{equation}
 \begin{gathered}
 \rho_0=C_*\mu,\qquad
 A_\theta(C)=\mathbb E_\mu[e^{\theta S}\mid C],
 \qquad T_\theta(C)=\lambda(N)(e^{\theta S}),\\
 Z_a=\mathbb E_\mu e^{\theta S},\qquad
 Z_b=\mathbb E_{\rho_0}T_\theta,\qquad
 d\rho_a=A_\theta\,d\rho_0/Z_a,\qquad
 d\rho_b=T_\theta\,d\rho_0/Z_b,\\
 r_\theta=\log(A_\theta/T_\theta),\qquad
 \delta_\theta=\log(Z_b/Z_a).
 \end{gathered}
 \label{eq:f15v34:laws}
\end{equation}
These are the one-source quantities of V31, with the
same centering. No natural-near conditional reference
from V28--V29 is substituted for \(\lambda\).

\begin{proposition}[The fallback energizes a central plaquette]
\label{f15v34:haar-energy}
For every retained input in \(\mathcal F_\beta\),
\begin{equation}
 \lambda(N)(Y)\ge\frac\beta4,\qquad
 T_\theta(C)\ge\exp\{\theta\beta/4-180\theta\}.
 \label{eq:f15v34:haar-energy}
\end{equation}
\end{proposition}

\begin{proof}
Work on the one-step canonical tree slice of V12.
The internal face \(f_0=(0;1,2)\) has three tree
links equal to \(I\): the links \((0,1)\),
\((e_1,2)\) and \((0,2)\). The remaining edge
\((e_2,1)\) is one of the 17 free internal groups,
say \(X\), so its holonomy is \(X^{-1}\).
This free group is in the core \(A\). In the
retained frame it is merely conjugated. Product
Haar makes it Haar independently of the other
core groups. Central symmetry \(X\mapsto-X\)
therefore gives
\[
 \int_{\SU(2)}\left(1-\tfrac12\Tr X\right)dH(X)=1.
\]
The cell \(0\) belongs to \(\mathcal B_2\), and
\(T_0\) contains \(X_{f_0}/4\), where
\(X_f=\beta(1-\frac12\Tr U_f)\). All other terms
in \(Y\) are nonnegative. This proves the mean
bound. Jensen's inequality and
\(\mathbb E_\mu Y\le180\) prove the second bound.
\end{proof}

\subsection{An existing remote chart test and its long fundamental cycle}

Let the ordered coarse frame tree have root
\(o=(-r+20,-r+20,-r+20,-r+20)\), as in V13.
Choose the positive coarse edge
\begin{equation}
 y=(-\lfloor r/2\rfloor,0,0,0),\qquad e=(y,1).
 \label{eq:f15v34:target}
\end{equation}
It is well inside the \(J\)-anchor box of radius
\(r-60\), and is not an edge of that ordered tree.
It is important that its retained group is already
among the parameters required by \(W_J\); no new
failure criterion is added here.

To see this directly, consider the fine plaquette
based at \(2y+e_1\), in directions \((1,2)\).
On the canonical slice its holonomy is
\(B_eY_{e,a}^{-1}\), where \(B_e\) is the root
pivot and \(Y_{e,a}\) is the nonroot boundary edge
at transverse offset \(e_2\). The two other edges
are tree identities. This \(Y_{e,a}\) is a
\(J\)-detail. At identity, hold other increments
zero and write \(C_e=c\), \(Y_{e,a}=z\) for
their Lie-algebra increments. The linear inverse
\eqref{eq:f15v12:linear-pivot} gives \(B_e=8c-z\).
The face therefore has first variation \(8c-2z\)
and Wilson quadratic term
\[
                       \tfrac12|8c-2z|^2.
\]
Its mixed derivative is \(-16I\), so this
individual measured Wilson term depends on both
\(C_e\) and the interior detail. It belongs to
\(W_J\), and its retained input \(C_e\) belongs
to the term-support parameter list used in V13
and V22. Thus
\begin{equation}
 \vartheta(C_e^{H_N})>a_2/2
                    \quad\Longrightarrow\quad\mathcal F_\beta.
 \label{eq:f15v34:trigger}
\end{equation}
Here \(\vartheta\) is the geodesic angle on
\(\SU(2)\), normalized to take values in \([0,\pi]\).

Let \(\gamma_y\), \(\gamma_{y+e_1}\) be the
ordered paths from \(o\). Remove their common
initial segment from the walk
\(\gamma_y e\gamma_{y+e_1}^{-1}\). Its remaining
fundamental cycle has
\begin{equation}
                         m=6(r-20)+2
 \label{eq:f15v34:length}
\end{equation}
distinct undirected edges. Indeed each branch
traverses directions 2, 3 and 4 through \(r-20\)
steps each, and the two branches are joined by
two direction-1 edges, one at their bottom and
one at their top. All cycle edges have first
coordinate \(-\lfloor r/2\rfloor\) or one more.
They lie within the frame box. Only the top
closing edge \(e\) is nontree. The common
initial segment gives a conjugation in the
rooted loop, not an unrestricted noncommutative
algebraic cancellation; the tree-frame calculation
below accounts for it exactly.

\subsection{A product-Haar-preserving shift, with its noncommutative proof}

Put \(\delta_0=a_2\), \(\varepsilon=\delta_0/m\)
and fix a unit Lie-algebra element \(J_0\).
For each positive coarse edge \(f\) on the cycle
let \(s_f=+1\) or \(-1\) be its cycle orientation,
and set
\[
 V_f=\exp(s_f\varepsilon J_0).
\]
Set \(V_f=I\) off the cycle. Extend the ordered
box tree to a spanning tree of the coarse torus,
rooted at \(o\), and let \(H_x(C)\) be its
root-path product. Such an extension exists
because the box tree is acyclic. On the box,
these are the original \(H_N\) frame values.
Define the compact map \(\mathcal T\) by
\begin{equation}
 C'_f=H_x(C)^{-1}V_fH_x(C)C_f,\qquad f=(x,z),
 \qquad v'=v,\qquad g'=g.
 \label{eq:f15v34:shift}
\end{equation}
It fixes every off-cycle retained edge. It is not
a gauge transformation; it changes the rooted
loop and the Wilson action.

\begin{proposition}[Exact reference-measure preservation and chart exit]
\label{f15v34:shift}
The map \(\mathcal T\) is a smooth bijection
preserving \(dC\,dH_v\,dH_g\). If
\(\vartheta(C_e^{H_N})\le a_2/4\), then its
image belongs to \(\mathcal F_\beta\).
\end{proposition}

\begin{proof}
The configuration-dependent conjugation in
\eqref{eq:f15v34:shift} does not by itself prove
Haar preservation. Use spanning-tree coordinates:
\(h_x=H_x(C)\), \(h_o=I\), and for each
nontree edge \(f=(x,z)\),
\(B_f=h_xC_fh_z^{-1}\). On tree edges the latter
quantity equals \(I\). The tree change of
variables and the fixed-endpoint translations
give normalized product Haar in all the
\(h_x\), \(x\ne o\), and nontree \(B_f\)'s.

Let \(P_x\) be the ordered product of the
deterministic profile \(V\) along the spanning-tree
path to \(x\), using inverses when traversing
an edge backwards. Induction along the tree
and direct multiplication on a nontree edge give
\begin{equation}
 h'_x=P_xh_x,\qquad
 B'_f=P_xV_f B_fP_z^{-1}.
 \label{eq:f15v34:forest}
\end{equation}
All factors \(P,V\) here are deterministic.
These are separate left/right Haar translations,
with separate inverses, so they preserve product
Haar and are bijective. This argument does not
require the profile factors to commute. The
details and fine gauge groups were fixed, proving
the measure assertion in the canonical coordinates.

For the chosen commuting cycle profile, the common
prefix has \(V=I\). The two remaining tree
branches and the closing edge give
\[
 B'_e=e^{uJ_0}B_e e^{wJ_0},\qquad u+w=m\varepsilon=\delta_0.
\]
The signs in the reversed branch are already
included in \(s_f\). Bi-invariance of the metric
implies
\[
 d(B'_e,e^{\delta_0J_0})=d(B_e,I)\le a_2/4.
\]
Since \(0<a_2<\pi\), the triangle inequality
gives \(\vartheta(B'_e)\ge3a_2/4>a_2/2\).
The frame used in \(B'_e\) is the new frame
\(H_N(C')\), not the old frozen value. Apply
\eqref{eq:f15v34:trigger} to conclude.
\end{proof}

\subsection{The remote event has only a sublinear Wilson log-cost}

In the one-step canonical inverse, changing \(C_f\)
at fixed \(X,Y\) changes just its own fine root
pivot \(B_f\). The other fine links and other
primitive inverse Jacobians are unchanged. Define
\begin{equation}
 K_{\rm inv}=\max\{1,\sup\|D_C B(X,Y,C)\|\},
 \qquad
 \mathcal R_\beta=
 16K_{\rm inv}\delta_0\sqrt{\beta\log\beta}
 +\frac{48K_{\rm inv}^2\delta_0^2\beta}{m}
 +2J_*m.
 \label{eq:f15v34:cost}
\end{equation}
The derivative uses geodesic tangent norms. Its
supremum is finite by the global smooth primitive
inverse on its compact input domain, recorded in V9.
The finite \(J_*\) is exactly
\eqref{eq:f15v23:compact-budget}. No numerical
enclosure of either supremum is asserted.

Let \(\mathcal P\) be the deterministic set of
fine plaquettes incident to the \(m\) root pivots
on the cycle. Each fine edge meets six plaquettes,
so \(|\mathcal P|\le6m\). A positive coarse edge
\(f=(x,\mu)\) has root pivot based at
\(2x+e_\mu\). Consequently every affected fine
plaquette base has first coordinate at most
\(-r+5<-24\). The set is remote from the source
carrier, whose recorded radius is at most
\(12p=24\). In particular \(\mathcal T\) leaves
the actual source \(Y\) unchanged.

\begin{theorem}[A lower bound on the actual retained-fallback probability]
\label{f15v34:probability}
Along \eqref{eq:f15v34:scales}, eventually
\begin{equation}
 \rho_0(\mathcal F_\beta)=\mu(\mathcal F_\beta)
             \ge\tfrac12e^{-\mathcal R_\beta},
 \qquad \mathcal R_\beta=o(\beta).
 \label{eq:f15v34:probability}
\end{equation}
The same event still has \(\mu(\mathcal F_\beta)\to0\).
\end{theorem}

\begin{proof}
First define a high-probability starting set. With
\(t=\log\beta\), put
\[
 G_1=\{R_{\mathcal K}\le a_2/(128rp^4)\},\qquad
 G_2=\{\textstyle\sum_{f\in\mathcal P}X_f\le mt\},
 \qquad G=G_1\cap G_2,
 \qquad p=2.
\]
V13's entry estimate puts the old \(B_e=C_e^{H_N}\)
within angle \(a_2/4\) on \(G_1\). The V32
strip-energy bound \eqref{eq:f15v32:q}, used with
this half-sized entry threshold, gives
\[
 \mu(G_1^c)\le
 \min\left\{1,\,4L(L+1)^3
 \exp\left[-\frac{\beta a_2^2}
 {12\pi^2L(128rp^4)^2}+15L\log2\right]\right\}
                           \longrightarrow0.
\]
At fixed \(p=2\), the negative exponent has size
\(\beta/r^3\), which dominates the \(O(r)\)
and \(O(\log r)\) costs for every fixed polylogarithmic
growth bound in \eqref{eq:f15v34:scales}.
For each fine plaquette, its base cell satisfies
\(X_f\le4T_{\operatorname{base}(f)}\), so
\[
 \mathbb E_\mu\sum_{f\in\mathcal P}X_f
      \le20|\mathcal P|\le120m,
 \qquad \mu(G_2^c)\le120/t.
\]
Repeated bases would not affect this expectation
estimate. Thus \(\mu(G)\ge1/2\) eventually, and
Proposition~\ref{f15v34:shift} gives
\(\mathcal T(G)\subseteq\mathcal F_\beta\).

It remains to pay the exact measured density ratio.
A changed coarse edge moves through angle
\(\varepsilon\). The global inverse bound implies
that its fine pivot moves by at most
\(K_{\rm inv}\varepsilon\), also after restoring
the unchanged fine gauge groups. Each affected
plaquette holonomy moves by at most
\(a=4K_{\rm inv}\varepsilon\), by telescoping
its four factors and using bi-invariance.
For unit quaternions write
\(s(U)=1-\frac12\Tr U=\frac12|q(U)-1|^2\).
The chord distance is at most the geodesic distance,
so a displacement at most \(a\) gives
\begin{equation}
 s(U')-s(U)\le a\sqrt{2s(U)}+a^2/2.
 \label{eq:f15v34:chord}
\end{equation}
This follows by expanding the squared chord and
using Cauchy--Schwarz; it is a global inequality,
not a quadratic Taylor approximation to the action.
Summing it on \(\mathcal P\), then applying
Cauchy--Schwarz on \(G_2\), yields
\begin{align*}
 \beta[S_W(\mathcal TU)-S_W(U)]
 &\le a\sqrt{2\beta|\mathcal P|
                       \sum_{f\in\mathcal P}X_f}
                          +\beta|\mathcal P|a^2/2\\
 &\le16K_{\rm inv}\delta_0\sqrt{\beta t}
                    +48K_{\rm inv}^2\delta_0^2\beta/m.
\end{align*}
Here \(4\sqrt{12}\le16\) and \(|\mathcal P|\le6m\).
Only the \(m\) corresponding primitive Jacobian
factors change. Their logarithmic ratio is bounded
below by \(-2J_*m\). Therefore the full canonical
Wilson density \(f\), with respect to
\(dC\,dH_v\,dH_g\), satisfies on \(G\)
\[
                    f(\mathcal T U)\ge e^{-\mathcal R_\beta}f(U).
\]
The notation \(U\) here includes its canonical
coordinates via the exact V12 bijection. Applying
the product-Haar-preserving change of variables gives
\[
 \mu(\mathcal F_\beta)\ge\mu(\mathcal T G)
       =\int_G f(\mathcal TU)\,dC\,dH_v\,dH_g
       \ge e^{-\mathcal R_\beta}\mu(G).
\]
The common Wilson normalizer cancels exactly.
No lower bound on that normalizer, no independence
of plaquettes and no configuration-independent
rotation assumption has been used.

Finally, \(m\asymp r\to\infty\),
\(\sqrt{\beta\log\beta}=o(\beta)\), and
\(m=o(\beta)\). These facts prove
\(\mathcal R_\beta=o(\beta)\). Entry at the
original threshold \(a_2/(64rp^4)\) excludes every
retained failure, so \(\mathcal F_\beta\) is
contained in the corresponding rooted bad event.
The same strip upper bound proves that its
probability tends to zero. The upper and lower
bounds concern different sublinear exponents and
are compatible.
\end{proof}

\subsection{Failure of the reference-energy moment and fixed-positive-source normalization}

\begin{theorem}[Actual Wilson obstruction for the specified Haar reference]
\label{f15v34:normalization}
For each fixed \(0<\theta\le1/8\), along
\eqref{eq:f15v34:scales},
\begin{equation}
 \delta_\theta\ge\frac{\theta\beta}4
   -\mathcal R_\beta-\log2-180\theta-36k_MJ(\theta),
 \qquad
 \liminf_{\beta\to\infty}\frac{\delta_\theta}{\beta}
                                      \ge\frac\theta4>0.
 \label{eq:f15v34:pressure}
\end{equation}
Moreover, for fixed \(0<\tau\le1/8\), set
\(v=1/(4\tau)\),
\(g^\lambda(N)=\log\lambda(N)(e^{\tau Y})\),
and \(g(\xi)=\log\mathsf P^\xi(e^{\tau Y})\)
for V33's actual interior conditional. Then
\begin{equation}
 \mathbb E_{\nu_0}e^{2vg}\le2^{180},\qquad
 \log\mathbb E_{\rho_0}e^{2vg^\lambda}
                   \ge\frac\beta8-\mathcal R_\beta-\log2.
 \label{eq:f15v34:moment-separation}
\end{equation}
Thus the analogous uniform reference-energy moment
is false for this particular \(\lambda\).
\end{theorem}

\begin{proof}
Proposition~\ref{f15v34:haar-energy} and
Theorem~\ref{f15v34:probability} give
\[
 Z_b\ge\tfrac12
       \exp\{\theta\beta/4-180\theta-\mathcal R_\beta\}.
\]
The inherited centered source CGF on the 36 cells
gives \(Z_a\le\exp\{36k_MJ(\theta)\}\), where
\(J(\theta)=-\log(1-\theta)-\theta\). Their
ratio proves \eqref{eq:f15v34:pressure}.
For the actual conditional, use V33 with
\(2v\tau=1/2\le R_{\rm src}\) to obtain
\(\mathbb E_{\nu_0}e^{2vg}\le2^{5n_2}\).
On \(\mathcal F_\beta\), Jensen gives
\(g^\lambda\ge\tau\beta/4\), so
\(2vg^\lambda\ge\beta/8\). Integrating only
over this event proves the second assertion.
\end{proof}

The last display tests exactly the kind of complementary
energy input isolated in V33. Its failure is now proved
for the declared Wilson/chart-reference pair, rather
than inferred only from the finite probability models
in V27 or V33.

\begin{theorem}[Normalized retained laws separate on the retuned collar]
\label{f15v34:separation}
Specialize to V32's collar
\eqref{eq:f15v32:scale} at \(p=2\), with any
fixed accuracy exponent \(s>0\). For every fixed
\(0<\theta\le1/8\),
\begin{equation}
 \mathbb E_{\rho_a}|r_\theta|\longrightarrow0,
 \qquad
 \liminf_{\beta\to\infty}
       \frac{D(\rho_a\Vert\rho_b)}{\beta}\ge\frac\theta4,
 \qquad d_{\rm TV}(\rho_a,\rho_b)\longrightarrow1.
 \label{eq:f15v34:separation}
\end{equation}
\end{theorem}

\begin{proof}
The one-source own-law estimate
\eqref{eq:f15v27:own} uses only the unscaled
chart and failure inputs. Apply it with the
improved V32 \(q_{\rm str}\):
\[
 \mathbb E_{\rho_a}|r_\theta|
     \le2[\Delta+c_K e^{h_\theta/2}\sqrt{q_{\rm str}}]
                                                 \longrightarrow0.
\]
Indeed \(\Delta\le\beta^{-s}\), \(c_K=O(\beta)\),
\(h_\theta=O(1)\), and at fixed \(p=2\) the
strip bound decays faster than any power of
\(\beta^{-1}\). The exact normalization identities
are still
\[
 D(\rho_a\Vert\rho_b)=\mathbb E_{\rho_a}r_\theta
                                      +\delta_\theta,
 \qquad \frac{d\rho_b}{d\rho_a}=e^{-r_\theta-\delta_\theta}.
\]
The first and \eqref{eq:f15v34:pressure} prove
the relative-entropy assertion. Eventually
\(\delta_\theta>0\). For the retained event
\(A_\beta=\{|r_\theta|\le\delta_\theta/2\}\),
Markov's inequality gives
\[
 \rho_a(A_\beta^c)\le
       2\mathbb E_{\rho_a}|r_\theta|/\delta_\theta\longrightarrow0,
 \qquad
 \rho_b(A_\beta)\le e^{-\delta_\theta/2}\longrightarrow0.
\]
Taking this event in the defining total-variation
supremum proves the limit one.
\end{proof}

\subsection{Compatibility audit and the upstream change of direction}

There is no conflict with V22's baseline core-TV
comparison or V27's actual-law source-log bound.
They can both ignore a sufficiently rare retained
event in their stated norms. Exponential weighting
and global normalization need not ignore it, as
\eqref{eq:f15v34:separation} demonstrates directly.

V31--V32's normalized estimates have source windows
shrinking with the regulator. The proof here fixes
\(\theta>0\) before sending \(\beta\) to infinity;
it does not establish failure on those shrinking
windows. V33 bounds the \emph{lower} pressure error
for positive sources and leaves the upper error
open. A large positive \(\delta_\theta\) is
consistent with that one-sided theorem. V17--V18
use a source-dependent upper-corner tilted baseline
for their full source cube, or negative source
penalties at the original baseline. They are not
the fixed-positive-source baseline reference in
\eqref{eq:f15v34:laws}. The actual-near factors
of V28--V29 are also different kernels.

The actionable change is therefore specific: do not
continue trying to prove a uniform \(g^\lambda\)
energy moment for this Haar branch. It is false.
A replacement must not energize the whole central
core merely because a remote frame test failed.
An actual single-near conditional on that branch
is one candidate for an energy-compatible repair,
but is not adopted or certified here. Its local
energy control, interaction with the good chart
branch, and joint normalization across several
sources must all be paid. The rare retained-phase
obstruction from V30 prevents inferring the last
point from separate one-source bounds alone.
The next V35 companion checkpoint will reassess
this direction against the related programmes.

\subsection{Exact diagnostics and predecessor preservation}

The new diagnostic is
\begin{center}\small
\path{scripts/verify_ym_f15_v34_remote_haar_fallback.py}.
\end{center}
It checks 3,888 forest-coordinate bijections and
inverses over the noncommutative group \(S_3\),
15,552 edge/frame identities, and three normalized
change-of-variables probability inequalities.
Five collar sizes give 37,624 distinct remote
coarse edges/pivots and 225,739 incident fine
plaquettes, all outside the source carrier.
Rational unit quaternions check ten internal-face
and central-symmetry identities, 100 Wilson
chord/Cauchy identities, and 360 noncommuting
rooted-loop displacement identities. Eight source,
mixed-jet and budget identities and 15 squared
likelihood-ratio coefficient comparisons are also
checked using exact arithmetic.

These are finite algebraic and geometric diagnostics,
not a Wilson simulation or an enclosure of the
global inverse suprema. The analytic probability
and asymptotic arguments are the proofs above.
The diagnostic has 10,026 bytes and SHA--256
\begin{center}\scriptsize\ttfamily
8288842EAB5E69E47FAC4C251E740BB9E4D3DF47E30368CEB4F2EA4B1D5A3F0D.
\end{center}

The V33 predecessor has 762,929 bytes and SHA--256
\begin{center}\scriptsize\ttfamily
B7B6650583DECFE78E8E1E161C9450B2F84BBFC6AF0D3C5878782780FC75448D.
\end{center}
Removing this terminal section and reversing the title
version recovers it exactly. V34 replaces V33 as the
sole active main note. All 14 protected sources and
76 earlier certificates are preserved. The next
companion review is V35, not postponed by this
upstream obstruction.

\begin{firewall}
V34 proves a failure of fixed-positive-source
normalization for the specified chart/Haar-fallback
reference under the inherited selected Wilson inputs,
and disproves its proposed uniform reference-energy
moment. It does not disprove a Yang--Mills mass gap,
modify the Wilson law, or invalidate the separately
proved shrinking-source and one-sided bounds. A
repaired reference, fixed-strength joint normalization,
retained-field clustering, continuum reconstruction
and a positive physical Hamiltonian mass gap remain
unproved. The full Yang--Mills goal remains active.
\end{firewall}

\section{V35: energy-capped reference kernels and mixed-sign cubic normalization}
\label{f15v35:section}

V34 proved that the existing Haar fallback is unsuitable
for fixed positive sources. This section changes that
comparison kernel explicitly; it does not try to prove
the false reference-energy moment for the old kernel.
The replacement has an energy cap relative to the
minimum allowed by its retained input. Its cap and
fallback are independent of the source parameters.
The actual compact Wilson law, physical observables,
source centering and completed blocking map are unchanged.

The new result is a two-sided normalization theorem
for arbitrarily many fitting sources with arbitrary
signs. Its intensive error has a cubic source-strength
term and regulator-vanishing terms. The cubic term
does not vanish at fixed nonzero strength. At fixed
blocking depth, however, the source strength may tend
to zero arbitrarily slowly, instead of obeying V32's
specified negative power of \(\log\beta\). This repairs
a concrete reference failure and enlarges the useful
window; it is not a proof of continuum clustering or
of the physical Yang--Mills mass gap.

\subsection{V35 companion checkpoint}

The current \texttt{latex\_notes} inventory and the
supplied download filenames were checked for updates.
All 91 protected files, comprising 14 sources and
77 earlier certificates, have unchanged hashes.
The latest conclusions and scope statements in the
following seven notes were reread, together with the
complete optional \texttt{suggestions.txt}.

\begin{center}\small
\begin{tabular}{@{}p{0.25\textwidth}p{0.69\textwidth}@{}}
\toprule
Current source & Consequence for this step\\
\midrule
SM continuum V72 &
Its fixed-mode, fixed-time many-copy covariance limit
does not provide a physical cutoff-uniform YM estimate.\\[2mm]
NS stability V26 &
Pointwise rank-one derivative norms do not supply a
many-source normalizer; its proposed V27 certification
is not a completed result in the current file.\\[2mm]
Shell mixing V16 &
The complete measured flat-holonomy one-loop calculation
does not supply its missing noncommuting calculation
or a common nonperturbative remainder.\\[2mm]
Parallel YM V5 &
Its extensive global-charge tightness obstruction remains;
refinement coherence does not remove it.\\[2mm]
SM source selection V31 &
The ancestry compiler is conditional on word-native
typing, not a proof of those operations' historical occurrence.\\[2mm]
Native stationarity V34 &
The moving-fibre escape reduction still leaves the
330 action-specific derivative entries to be populated.\\[2mm]
Exact-action bridge V24 &
A nonzero first-jet obstruction derivative is not an
integrated root or a uniform continuum bridge.\\
\bottomrule
\end{tabular}
\end{center}

None provides the missing YM probability estimate.
The suggestions' measured-density and source-convention
cautions remain useful: a local one-loop coefficient
must not be promoted to a full-law normalization
bound. These files are mathematical data, not new
instructions. The present direction therefore remains
upstream at the full compact source/reference
comparison. The next companion checkpoint is V40.

\subsection{A source-independent energy cap, with both normalizers explicit}

Use the separated \(n\)-source geometry of V21--V34,
with \(n\ge1\), and a common dyadic depth \(p\).
The actual uncentred energies and sources remain
\[
 Y_i=Y_i(N_i,V_i)=\sum_{x\in\mathcal B_{p,i}}T_x\ge0,
 \qquad S_i=Y_i-\mathbb E_\mu Y_i,\qquad
 |\mathcal B_{p,i}|=n_p\le4p^4.
\]
They depend only on their own core \(V_i\) and near
retained data \(N_i\). All core coordinates are the
same framed compact coordinates as before. Define
\begin{equation}
 m_i(N_i)=\min_{V_i}Y_i(N_i,V_i),\qquad
 H_i^\Lambda=\{V_i:Y_i(N_i,V_i)-m_i(N_i)<\Lambda\},
 \qquad \Lambda\ge1.
 \label{eq:f15v35:minimum}
\end{equation}
The minimum is over the entire compact core, not
only the chart. The finite-regulator energy is
continuous by the global completed inverse, so the
minimum is attained and \(m_i\) is continuous in
the near data. Indeed its difference at two inputs
is bounded by the supremum over the compact core
of the difference of the two energy functions.
No differentiability of \(m_i\) is required.

Let \(\mathcal R_i\) be the retained half-chart
event, and let \(\lambda_i^{\rm ch}(N_i)\) denote
V22's measured chart core marginal at identity
exterior on that event. With \(H_{A_i}\) normalized
product Haar on the core, define the new kernel
\begin{equation}
 \lambda_i^\Lambda(dV_i\mid N_i)=
 \begin{cases}
 \displaystyle
 \frac{1_{H_i^\Lambda}(V_i)\lambda_i^{\rm ch}(dV_i\mid N_i)}
      {\lambda_i^{\rm ch}(H_i^\Lambda\mid N_i)},
 &\mathcal R_i\text{ and }\lambda_i^{\rm ch}(H_i^\Lambda)>0,\\[3mm]
 \displaystyle
 \frac{1_{H_i^1}(V_i)dH_{A_i}(V_i)}
      {H_{A_i}(H_i^1)},
 &\text{otherwise}.
 \end{cases}
 \label{eq:f15v35:kernel}
\end{equation}
The second denominator is strictly positive: the
strict sublevel contains a neighbourhood of a
minimizer, and product Haar has full support.
The first denominator is not silently set to one;
if it vanishes the second branch is used.
This defines a measurable probability kernel
depending only on \(N_i\), the fixed source shape,
the regulator and \(\Lambda\), not on \(\theta\).
Both branches satisfy
\begin{equation}
       0\le Y_i-m_i<\Lambda
              \quad\lambda_i^\Lambda\text{-almost surely}.
 \label{eq:f15v35:support}
\end{equation}
For \(\Lambda=1\) the second branch obeys the same
strict bound. Simultaneous changes of the rooted
frame preserve the energy, its minimum, product
Haar and the measured chart construction. Thus
the sublevel definition is compatible with that
frame covariance; no noncovariant choice of a
particular minimizing configuration has been made.

This is a new comparison measure. It neither
identifies the old V34 proxy with the repaired
one nor restricts the actual Wilson configurations.
In particular the retained-dependent \(m_i\) is
not subtracted from the physical source or divided
out of its global normalizer.

\subsection{The added energy failures retain a joint spatial bound}

Let \(\xi=(C,v_O)\) be the common exterior,
\(\nu_0=\xi_*\mu\), and \(\mathsf P_i^\xi\) the
actual baseline conditional law of \(J_i\).
These interiors are conditionally independent.
Retain the old exterior event \(\Gamma_i\) and
interior chart event \(D_i\), but set
\begin{equation}
 D_i^\Lambda=D_i\cap H_i^\Lambda,\qquad
 E_i^\Lambda=\Gamma_i^c\cup(D_i^\Lambda)^c.
 \label{eq:f15v35:failure}
\end{equation}
The sublevel is a core event. On \(\Gamma_i\),
conditioning the chart core marginals on this same
sublevel changes their previous log-ratio bound
by at most another \(\Delta\). More explicitly,
if \(p_b/p_0\in[e^{-\Delta},e^\Delta]\), their
sublevel masses have the same ratio bounds, hence
\begin{equation}
 \frac{p_b(\,\cdot\mid H_i^\Lambda)}
      {p_0(\,\cdot\mid H_i^\Lambda)}
                         \in[e^{-2\Delta},e^{2\Delta}]
 \label{eq:f15v35:restricted-ratio}
\end{equation}
whenever the masses are positive. If the reference
chart sublevel mass is zero, the actual chart
sublevel mass is also zero on \(\Gamma_i\).
Then \(\mathsf P_i^\xi(E_i^\Lambda)=1\).
This handles the zero-mass branch rather than
assuming it cannot occur.

\begin{proposition}[Joint failures after the energy cap]
\label{f15v35:failure}
For a fitting separated family, define
\begin{equation}
 q_Y=\min\{1,\exp(-\Lambda/2+5n_p\log2)\},\qquad
 q=\min\{1,\sqrt{q_{\rm str}}+\sqrt{q_Y}\}.
 \label{eq:f15v35:q}
\end{equation}
Then every subfamily \(I\) satisfies
\begin{equation}
             \mu\left(\bigcap_{i\in I}E_i^\Lambda\right)
                                                   \le q^{|I|}.
 \label{eq:f15v35:joint}
\end{equation}
\end{proposition}

\begin{proof}
Since \(m_i\ge0\), an added energy failure implies
\(Y_i\ge\Lambda\). The disjoint source-cell sets
and the original CGF at coefficient \(1/2\) give
\[
 \mu(Y_i\ge\Lambda\text{ for all }i\in I)
 \le e^{-|I|\Lambda/2}\,
       \mathbb E_\mu e^{\frac12\sum_{i\in I}Y_i}
 \le q_Y^{|I|}.
\]
Expand an intersection of
\(E_i^\Lambda\subset E_i\cup\{Y_i\ge\Lambda\}\)
over the subset assigned to old failures.
For a choice \(J\subset I\), Cauchy--Schwarz and
the V32 joint bound give at most
\(q_{\rm str}^{|J|/2}q_Y^{|I\setminus J|/2}\).
Summing over \(J\) gives
\((\sqrt{q_{\rm str}}+\sqrt{q_Y})^{|I|}\).
The bound by one proves the assertion. This step
does not assume independence between chart failures
and source energies.
\end{proof}

\subsection{One local error bound valid for both signs}

Fix an arbitrary vector
\(\theta\in[-\tau,\tau]^n\), \(0<\tau\le1/512\),
and put \(t_i=|\theta_i|\). Define
\begin{equation}
 \begin{aligned}
 u_i(\xi)&=\mathsf P_i^\xi(e^{\theta_iY_i}),&
 b_i(N_i)&=\lambda_i^\Lambda(e^{\theta_iY_i}),&
 L_i&=\log(u_i/b_i),\\
 g_i(\xi)&=\log\mathsf P_i^\xi(e^{t_iY_i})\ge0,&
 \pi_{i,0}&=\mathsf P_i^\xi(E_i^\Lambda),&
 \pi_{i,\theta}&=
       \frac{\mathsf P_i^\xi(1_{E_i^\Lambda}e^{\theta_iY_i})}{u_i}.
 \end{aligned}
 \label{eq:f15v35:local}
\end{equation}
The centering scalar cancels in \(L_i\). Remote
source factors cancel in the displayed conditional
probability only after fixing the common \(\xi\).
For \(0<\eta<1\), set
\begin{equation}
 B_i=\{\pi_{i,0}>\eta\text{ or }\pi_{i,\theta}>\eta\},
 \qquad e_0=2\Delta-\log(1-\eta).
 \label{eq:f15v35:confidence}
\end{equation}

\begin{proposition}[The reference cap controls the previously missing sign]
\label{f15v35:local}
At every conditional input,
\begin{equation}
 \begin{array}{c|cc}
 & (L_i)_+ &(L_i)_-\\ \hline
 \theta_i\ge0 &\le g_i&\le t_i\Lambda\\
 \theta_i\le0 &\le t_i\Lambda&\le g_i
 \end{array}
 \qquad
 |L_i|\le e_0+(g_i+\tau\Lambda)1_{B_i}.
 \label{eq:f15v35:local-bound}
\end{equation}
\end{proposition}

\begin{proof}
For a positive source, \(u_i\le e^{g_i}\),
\(b_i\ge e^{t_i m_i}\ge1\), and
\(u_i\ge e^{t_i m_i}\), \(b_i\le e^{t_i(m_i+\Lambda)}\).
These prove the first row. For a negative source,
Cauchy--Schwarz gives
\(\mathsf P_i^\xi(e^{-t_iY_i})\ge e^{-g_i}\),
while \(b_i\le1\). Also
\(u_i\le e^{-t_i m_i}\) and
\(b_i\ge e^{-t_i(m_i+\Lambda)}\). This proves
the second row. A zero source has \(L_i=g_i=0\).

Off \(B_i\), both conditional success probabilities
\(\alpha_0=1-\pi_{i,0}\) and
\(\alpha_\theta=1-\pi_{i,\theta}\) lie in
\([1-\eta,1]\). Thus \(\Gamma_i\) holds and the
chart sublevel masses are positive. The exact
restricted normalizer identity is
\[
 L_i=\log\alpha_0-\log\alpha_\theta+
 \log\frac{p_b(\,\cdot\mid H_i^\Lambda)(e^{\theta_iY_i})}
               {\lambda_i^\Lambda(e^{\theta_iY_i})}.
\]
The last term has absolute value at most \(2\Delta\)
by \eqref{eq:f15v35:restricted-ratio}, and the
difference of the first two lies between
\(\log(1-\eta)\) and \(-\log(1-\eta)\).
This proves \(|L_i|\le e_0\) there. On \(B_i\)
use the appropriate global row, whose maximum
is at most \(g_i+\tau\Lambda\).
\end{proof}

\subsection{Mixed-sign confidence and an exceptional-energy exponential moment}

The actual conditional-energy estimate of V33 still
holds: whenever \(a\ge1\) and \(a\tau\le R_{\rm src}\),
\begin{equation}
 \mathbb E_{\nu_0}\exp\left(a\sum_{i\in I}g_i\right)
                         \le(1-a\tau)^{-5n_p|I|}.
 \label{eq:f15v35:actual-energy}
\end{equation}
It is applied only to actual conditional energies,
not transferred to the reference.
Define
\begin{equation}
 \gamma=q^{1/4}(1-4\tau)^{-5n_p/4}(1-2\tau)^{-5n_p/2}.
 \label{eq:f15v35:gamma}
\end{equation}

\begin{proposition}[Confidence products without a uniform-sign restriction]
\label{f15v35:mixed-confidence}
For every choice \(a_i\in\{0,\theta\}\) on a
subfamily \(I\),
\begin{equation}
 \mathbb E_{\nu_0}\prod_{i\in I}\pi_{i,a_i}
                 \le\gamma^{|I|},\qquad
 \nu_0\left(\bigcap_{i\in I}B_i\right)
                 \le(2\gamma/\eta)^{|I|}.
 \label{eq:f15v35:mixed-confidence}
\end{equation}
\end{proposition}

\begin{proof}
Write \(M_i=\mathsf P_i^\xi(e^{t_iY_i})\) and
\(H_i=\mathsf P_i^\xi(1_{E_i^\Lambda}e^{t_iY_i})\)
in this proof only; these \(H_i\)'s are scalar
moments, not the sublevels \(H_i^\Lambda\).
Both baseline and tilted failure probabilities
are bounded by \(H_iM_i\). For the positive
tilt its denominator is \(M_i\ge1\). For the
negative tilt its numerator is at most \(H_i\)
and its denominator is at least \(1/M_i\).
The baseline inequality follows from \(M_i\ge1\).

Conditional independence, conditional Jensen and
Cauchy--Schwarz now give
\begin{align*}
 \mathbb E_{\nu_0}\left(\prod_{i\in I}H_i\right)^2
 &\le\mathbb E_\mu
       \left[1_{\cap_{i\in I}E_i^\Lambda}
                              e^{2\sum_{i\in I}t_iY_i}\right]\\
 &\le q^{|I|/2}(1-4\tau)^{-5n_p|I|/2},\\
 \mathbb E_{\nu_0}\left(\prod_{i\in I}M_i\right)^2
 &\le(1-2\tau)^{-5n_p|I|}.
\end{align*}
One further Cauchy--Schwarz proves the first bound
in \eqref{eq:f15v35:mixed-confidence}, simultaneously
for every mixture of the two probabilities.
Finally \(1_{B_i}\le(\pi_{i,0}+\pi_{i,\theta})/\eta\).
Expand the product and use the first bound on
each of its \(2^{|I|}\) terms.
\end{proof}

Put
\begin{equation}
 \begin{gathered}
 v=\frac1{128\tau}\ge4,\qquad
 F=\sum_{i=1}^n(g_i+\tau\Lambda)1_{B_i},\\
 K_\Lambda=
 e^{\Lambda/128}(1-1/64)^{-5n_p/2}
                     \sqrt{2\gamma/\eta},\qquad
 m_\Lambda=\log(1+K_\Lambda).
 \end{gathered}
 \label{eq:f15v35:laplace-budget}
\end{equation}

\begin{theorem}[A baseline exceptional-energy Laplace bound]
\label{f15v35:laplace}
For the actual common exterior,
\begin{equation}
                         \log\mathbb E_{\nu_0}e^{vF}
                                                   \le n m_\Lambda.
 \label{eq:f15v35:laplace}
\end{equation}
\end{theorem}

\begin{proof}
For each nonempty subfamily, Cauchy--Schwarz,
\eqref{eq:f15v35:actual-energy}, and
\eqref{eq:f15v35:mixed-confidence} imply
\[
 \mathbb E_{\nu_0}\left[
       e^{v\sum_{i\in I}(g_i+\tau\Lambda)}
                          1_{\cap_{i\in I}B_i}\right]
 \le\left[
 e^{v\tau\Lambda}(1-2v\tau)^{-5n_p/2}
                         \sqrt{2\gamma/\eta}\right]^{|I|}
 =K_\Lambda^{|I|}.
\]
Here \(2v\tau=1/64<R_{\rm src}\).
Expand
\(e^{vF}=\prod_i[1+1_{B_i}(e^{v(g_i+\tau\Lambda)}-1)]\),
bound each nonnegative added factor by its exponential,
and sum the subset bounds. The result is
\(\mathbb E_{\nu_0}e^{vF}\le(1+K_\Lambda)^n\).
No reference-law moment hypothesis is used.
\end{proof}

\subsection{Both projected log signs and the exact physical normalizer}

Keep the original physical weights, and use the new
kernel only in the comparison factors:
\begin{equation}
 \begin{gathered}
 W_\theta=e^{\sum_i\theta_i S_i},\quad
 Z_a=\mathbb E_\mu W_\theta,\quad
 d\mu_\theta=W_\theta\,d\mu/Z_a,\quad
 \rho_0=C_*\mu,\quad \rho_a=C_*\mu_\theta,\\
 A_\theta(C)=\mathbb E_\mu[W_\theta\mid C],\quad
 T_\theta(C)=\prod_i\lambda_i^\Lambda(e^{\theta_iS_i}),\\
 Z_b=\mathbb E_{\rho_0}T_\theta,\quad
 d\rho_b=T_\theta\,d\rho_0/Z_b,\quad
 r_\theta=\log(A_\theta/T_\theta),\quad
 \delta_\theta=\log(Z_b/Z_a).
 \end{gathered}
 \label{eq:f15v35:laws}
\end{equation}
The retained baseline \(\rho_0\) is fully correlated.
All normalizers are finite and positive at each
finite regulator. Write \(\nu_\theta=\xi_*\mu_\theta\)
and retain the actual-source bound
\begin{equation}
 h_\theta=\frac1n\log\mathbb E_\mu
                   \left(\frac{d\mu_\theta}{d\mu}\right)^2
                    \le64p^4\tau^2,\qquad
 \varepsilon_\Lambda=e_0+\frac{h_\theta+m_\Lambda}{v}.
 \label{eq:f15v35:epsilon}
\end{equation}
The bound on \(h_\theta\) also holds for mixed signs:
the centered-source CGF bounds \(Z_a(2\theta)\),
\(Z_a(\theta)\ge1\), and \(J(-u)\le J(u)\) for
\(0\le u<1\). Thus the previous absolute-amplitude
bound applies without a common-sign assumption.

\begin{theorem}[Two-sided mixed-sign normalization for the repaired reference]
\label{f15v35:normalization}
For every fitting family and every
\(\theta\in[-\tau,\tau]^n\), \(0<\tau\le1/512\),
\begin{equation}
 \log\mathbb E_{\rho_a}e^{\pm w r_\theta}
                     \le n w\varepsilon_\Lambda
          \quad(0<w\le v/2),\qquad
 \mathbb E_{\rho_a}|r_\theta|\le2n\varepsilon_\Lambda.
 \label{eq:f15v35:exponential}
\end{equation}
In particular,
\begin{equation}
 \begin{gathered}
 |\delta_\theta|\le n\varepsilon_\Lambda,\qquad
 D(\rho_a\Vert\rho_b)\le2n\varepsilon_\Lambda,\qquad
 D(\rho_b\Vert\rho_a)\le4n\varepsilon_\Lambda,\\
 d_{\rm TV}(\rho_a,\rho_b)
                      \le\min\{1,\sqrt{n\varepsilon_\Lambda}\}.
 \end{gathered}
 \label{eq:f15v35:normalization}
\end{equation}
There is no cardinality restriction beyond fitting
the separated geometry. These are intensive bounds;
the displayed total TV bound is useful when its
total budget is small. At \(\theta=0\) both laws
are exactly \(\rho_0\), and both errors vanish.
\end{theorem}

\begin{proof}
Conditional independence and the two exact
normalizer identities give
\[
 r_\theta=\log\mathbb E_{\nu_0}
                        [e^{\sum_iL_i}\mid C],
 \qquad
 -r_\theta=\log\mathbb E_{\nu_\theta}
                        [e^{-\sum_iL_i}\mid C].
\]
The deterministic centering factors cancel in
both. Jensen and \eqref{eq:f15v35:local-bound}
therefore yield
\begin{equation}
 (r_\theta)_-\le ne_0+\mathbb E_{\nu_0}[F\mid C],
 \qquad
 (r_\theta)_+\le ne_0+\mathbb E_{\nu_\theta}[F\mid C].
 \label{eq:f15v35:projected}
\end{equation}

For clarity, the change-of-law inequality used
next is the following elementary three-factor
H\"older bound. If \(R=dQ/dP\), \(a\ge2\), then
\begin{equation}
 \mathbb E_P[Re^X]\le
          (\mathbb E_P R^2)^{1/a}
          (\mathbb E_P e^{aX})^{1/a}.
 \label{eq:f15v35:transport}
\end{equation}
For \(a>2\), factor the integrand as
\(R^{2/a}R^{1-2/a}e^X\) and use exponents
\(a,a/(a-2),a\); the middle expectation is
\(\mathbb E_P R=1\). At \(a=2\) this is
Cauchy--Schwarz. Both projected densities,
on \(C\) and on \(\xi\), have second moment
at most \(e^{nh_\theta}\) by conditional Jensen.

To control \(e^{-wr_\theta}\), apply
\eqref{eq:f15v35:transport} on the retained
baseline with \(a=v/w\ge2\) and
\(X=w\mathbb E_{\nu_0}[F\mid C]\).
Conditional Jensen bounds its \(a\)-moment by
\(\mathbb E_{\nu_0}e^{vF}\). For \(e^{wr_\theta}\),
first use conditional Jensen under \(\nu_\theta\),
then apply the same transport bound on
\((\nu_0,\nu_\theta)\) with \(X=wF\).
In both cases \eqref{eq:f15v35:laplace} gives
the first assertion of \eqref{eq:f15v35:exponential}.

For the absolute first moment, apply the entropy
variational inequality separately to
\(F\) on the exterior and
\(\mathbb E_{\nu_0}[F\mid C]\) on the retained
space, with parameter \(v\). Each relative
entropy is at most the logarithm of the
corresponding order-two moment, hence at most
\(nh_\theta\). Each baseline \(v\)-moment is
at most \(e^{nm_\Lambda}\). Integrating the
two bounds in \eqref{eq:f15v35:projected} proves
\(\mathbb E_{\rho_a}(r_\theta)_\pm
 \le n\varepsilon_\Lambda\), as claimed.

Exact normalization still gives
\[
 \delta_\theta=\log\mathbb E_{\rho_a}e^{-r_\theta},
 \qquad
 D(\rho_a\Vert\rho_b)=\mathbb E_{\rho_a}r_\theta
                                         +\delta_\theta.
\]
The \(w=1\) bounds and
\(\mathbb E e^{r_\theta}\mathbb E e^{-r_\theta}\ge1\)
give \(|\delta_\theta|\le n\varepsilon_\Lambda\).
The first-moment estimate gives the forward
entropy bound. The reverse direction is paid
using the \(w=2\) moment, available because
\(v\ge4\):
\[
 \mathbb E_{\rho_b}e^{-r_\theta}
       =e^{-\delta_\theta}
                        \mathbb E_{\rho_a}e^{-2r_\theta}.
\]
Jensen then yields
\[
 D(\rho_b\Vert\rho_a)
 \le\log\mathbb E_{\rho_a}e^{-2r_\theta}-2\delta_\theta
 \le4n\varepsilon_\Lambda.
\]
Pinsker's inequality gives the stated TV bound.
No estimate for the reference density was assumed
before establishing its normalizer.
\end{proof}

\subsection{A cubic intensive error and a much slower shrinking-source window}

Now use V32's retuned collar
\eqref{eq:f15v32:scale}, with fixed \(s>0\),
\(t=\log\beta\) and \(2\le p\le t\), and choose
\begin{equation}
                 \Lambda=p^4t^2,\qquad
                 \eta=\gamma^{1/4}.
 \label{eq:f15v35:choice}
\end{equation}
The latter is in \((0,1)\) eventually.
This single reference kernel is used for every
source vector in the stated cube.

\begin{proposition}[Uniform asymptotic budget]
\label{f15v35:cubic}
Uniformly for \(2\le p\le t\) and
\(0<\tau\le1/512\), eventually
\begin{equation}
 \begin{gathered}
 q\le e^{-\Lambda/6},\qquad
 \gamma\le e^{-\Lambda/32},\qquad
 m_\Lambda\le e^{-\Lambda/512},\\
 \varepsilon_\Lambda
       \le2\beta^{-s}+8192p^4\tau^3+3e^{-\Lambda/512}.
 \end{gathered}
 \label{eq:f15v35:cubic}
\end{equation}
\end{proposition}

\begin{proof}
V32 gives \(\Delta\le\beta^{-s}\), and its
lower bound on \(\log(1/q_{\rm str})\) grows
faster than \(\Lambda\), uniformly for \(p\le t\).
Also
\[
 \sqrt{q_Y}\le\exp[-\Lambda/4+(5n_p/2)\log2],
 \qquad n_p/\Lambda\le4/t^2.
\]
Eventually \((5n_p/2)\log2\le\Lambda/24\),
\(\sqrt{q_{\rm str}}\le e^{-\Lambda/4}\), and
\(\log2\le\Lambda/24\). Substitution in
\eqref{eq:f15v35:q} proves \(q\le e^{-\Lambda/6}\).
The two remaining logarithmic terms in
\eqref{eq:f15v35:gamma} are \(O(n_p)\) uniformly
on the source interval. Since \(n_p/\Lambda\to0\),
they improve \(\log\gamma\le-\Lambda/24+O(n_p)\)
to the stated \(-\Lambda/32\) upper bound.

Thus \(\eta\le e^{-\Lambda/128}\) and
\[
 K_\Lambda
 =\sqrt2\,e^{\Lambda/128}(1-1/64)^{-5n_p/2}
                                         \gamma^{3/8}
 \le\exp[-\Lambda/256+O(n_p)+\tfrac12\log2]
 \le e^{-\Lambda/512}
\]
eventually, uniformly. Use
\(\log(1+K_\Lambda)\le K_\Lambda\) and
\(-\log(1-\eta)\le2\eta\). Finally
\[
 \frac{h_\theta}{v}
       \le128\tau\cdot64p^4\tau^2=8192p^4\tau^3,
 \qquad \frac{m_\Lambda}{v}\le\tfrac14e^{-\Lambda/512}.
\]
Together with \(e_0=2\Delta-\log(1-\eta)\)
these prove \eqref{eq:f15v35:cubic}.
\end{proof}

For fixed \(p\), any \(\tau_\beta\downarrow0\),
eventually at most \(1/512\), gives
\(\varepsilon_\Lambda\to0\). No prescribed power
of \(t^{-1}\) is needed: for example
\(\tau_\beta=1/\log\log\beta\) is admissible
eventually. For growing \(p\), the sufficient
condition is \(p^4\tau_\beta^3\to0\).
These assertions concern intensive pressure and
entropy errors. Total TV convergence additionally
requires \(n\varepsilon_\Lambda\to0\).

There is also a relative quadratic statement. If
\(\tau_\beta\to0\) and
\(\beta^{-s}/(p^4\tau_\beta^2)\to0\), then
\begin{equation}
              \frac{\varepsilon_\Lambda}{p^4\tau_\beta^2}
                                                     \longrightarrow0.
 \label{eq:f15v35:relative}
\end{equation}
Indeed divide \eqref{eq:f15v35:cubic} by this
denominator; the exponential term is smaller
than every inverse power of \(\beta\), uniformly
on the depth range. At fixed \(p\) and
\(\tau_\beta=t^{-1}\), the intensive error is
\(O(t^{-3})\), with relative error \(O(t^{-1})\).
This is a real finite-amplitude pressure comparison.
It is not an unproved differentiation of a
regulator-dependent error bound at zero source,
nor a complex-analytic cumulant estimate.

\subsection{What has changed, and the remaining upstream obligation}

V34 is not reversed: the old unconditioned Haar
fallback still has its proved fixed-positive-source
failure. The new kernel prevents that particular
energy injection. Both its chart truncation and its
low-energy fallback were included in the proof;
neither a small omitted mass nor a small reference
normalizer was presumed. V30's rare retained-phase
example is not bypassed by claiming that actual
near conditionals tensorize. Here the joint raw
failure estimate was retained explicitly and
spent on the mixed confidence products.

The reference is local in the already declared
near input, and independent of the source vector,
but it is not proved to be the actual next-scale
Wilson action or a reflection-positive effective
measure. The theorem repairs normalized source
increments relative to the correlated retained
baseline. It gives no decay of that baseline's
connected correlations. In particular the cubic
term at fixed strength still remains; it arose
from the actual change-of-law moment and must
not be described as regulator-vanishing.

The next substantive task is to determine whether
that source-transport cost can be localized or
renormalized using additional Wilson structure,
and whether the repaired source increments admit
a uniform connected-response norm useful for
an actual RG step. Repeating the failed Haar
moment estimate, or relabelling an intensive
source bound as a Hamiltonian gap, would not
advance that task. The continuum state and
physical-time calibration remain separate
requirements of the original full goal.

\subsection{Exact diagnostics and predecessor preservation}

The new diagnostic is
\begin{center}\small
\path{scripts/verify_ym_f15_v35_energy_capped_reference.py}.
\end{center}
Six strictly positive finite models have a common
retained bit, local retained and exterior flags,
and four-state conditional cores. They exercise
the truncated-chart branch, retained-chart fallback
and zero-chart-sublevel-mass fallback. Energies
are \(512\log2\) times nonnegative integers, with
\(\tau=1/512\), \(v=4\) and
\(\Lambda=1024\log2\). All weights and probabilities
used in the checks are rational.

The 36 source patterns include positive, negative,
mixed, sparse and zero vectors. Checks cover
18,216 branch constructions and local sign caps,
3,036 good-set comparisons, 22 raw joint failure
subsets, 144 matched conditional-energy subsets,
468 mixed confidence products, 144 threshold
subsets and 144 exceptional weighted subsets.
There are 36 baseline Laplace bounds, 72 projected
second-moment inequalities, 72 fractional
H\"older bounds, and 144 actual projected
exponential moments with powers \(+1,-1,+2,-2\).

The tests also check 36 pressure brackets,
36 reverse-entropy moment identities,
36 centering cancellations, six exact zero-source
outputs and six asymptotic coefficient identities.
Rational upper bounds on roots are verified by
integer powers; no floating-point tolerance or
approximate logarithm is used. The finite moment
envelopes are calculated from the finite models,
not asserted to be the Wilson CGF constants.
The actual Wilson inputs, analytic chart
comparison and continuum claims are outside
these diagnostic checks.

The script has 12,271 bytes and SHA--256
\begin{center}\scriptsize\ttfamily
FEDF619621CA6FFC80C0BC6C028826BD96EEBAE37675B854E3791485E46A3DF6.
\end{center}
The V34 predecessor has 785,842 bytes and SHA--256
\begin{center}\scriptsize\ttfamily
47BD112C9D1F26943454952302AEAFFE9F655CFEB96786AD940228EA2F2630B3.
\end{center}
Removing this terminal section and reversing the
title version recovers V34 exactly. V35 replaces
V34 as the sole active main note. All 14 protected
sources and 77 earlier certificates are preserved.
The companion review has been completed at V35;
the next is V40.

\begin{firewall}
V35 constructs a new source-independent energy-capped
reference and proves two-sided mixed-sign source
normalization with a cubic intensive error for
all fitting separated families, under the recorded
Wilson inputs. It permits slowly shrinking source
strengths at fixed depth, but does not prove
vanishing error at a fixed nonzero strength,
baseline retained clustering, an interacting
continuum limit, reflection-positive reconstruction
of the repaired reference, or the full physical
Yang--Mills mass gap. The original goal remains active.
\end{firewall}

\section{V36: sublinear thermal probes and finite-family energy closure}
\label{f15v36:section}

V35 repairs the reference kernel but leaves a cubic
intensive cost for the original linear energy sources.
This section distinguishes that endpoint from genuinely
sublinear thermal observables. For every fixed
\(0<\alpha<1\), the actual fields \(Y_i^\alpha\)
admit source moments of every finite order in the
source parameter. This permits a moment parameter
growing with the regulator and removes the fixed-source
transport cost for these probes. The square-root case
is a lattice thermal-curvature magnitude, not a new
physical-time or mass normalization.

The resulting normalized comparison holds at arbitrary
fixed source amplitude, with mixed signs and with no
family-size restriction for its intensive bounds.
Its error decays algebraically in \(\beta\), up to
the recorded logarithmic factors. One energy-capped
reference is used for all these probes: it is not
chosen separately for each source or exponent.

Two finite-family consequences tie this back to the
original energy observables. First, the V35 bounds,
combined with its small actual-law log error, give
fixed-strength linear-energy normalization and both
output entropy directions when the number of sources
and the blocking depth are fixed. Second, a controlled
finite-difference argument transfers every mixed
polynomial energy moment to the joint capped reference.
Neither conclusion is a fixed-strength extensive
linear-energy theorem or a physical mass-gap theorem.

The V35 predecessor and all 92 protected sources and
certificates were checked unchanged. The mandatory
companion review was completed at V35, including the
supplied SM source-selection V31, native stationarity
V34 and exact-action Einstein bridge V24. The next
review remains V40. No arbitrary-background Wilson
moment estimate or non-Abelian correlation inequality
is assumed in the argument below.

\subsection{Use the complete proved failure budget to choose the cap}

Retain V32's selected law, separated geometry and
retuned collar \eqref{eq:f15v32:scale}, with
\(t=\log\beta\), \(2\le p\le t\) and a fixed
accuracy exponent \(s>0\). Keep its raw strip
parameter \(q_{\rm str}\), and define
\begin{equation}
 \mathcal L_\beta=\log(1/q_{\rm str}),\qquad
 \Lambda=\max\{1,\mathcal L_\beta\}.
 \label{eq:f15v36:cap}
\end{equation}
Eventually \(q_{\rm str}<e^{-1}\), so
\(\Lambda=\mathcal L_\beta\). With
\(d=5\nu+49\) and \(b\) exactly as in V32,
\begin{equation}
 \Lambda\ge\frac b2\,\frac{\beta}{p^d t^3},
 \qquad \Delta\le\beta^{-s},\qquad
 \frac{\Lambda}{n_p}\longrightarrow\infty
       \quad\hbox{uniformly for }2\le p\le t.
 \label{eq:f15v36:entry-scale}
\end{equation}

Use the same kernel construction
\eqref{eq:f15v35:kernel}, now at the cap
\eqref{eq:f15v36:cap}. Thus \(m_i(N_i)=\min_{V_i}Y_i\);
the measured chart is conditioned on \(Y_i-m_i<\Lambda\)
when that chart sublevel has positive mass, and
otherwise the fallback is normalized core Haar
conditioned on \(Y_i-m_i<1\). All denominators
and the zero-mass branch remain exactly as in V35.
Only this deterministic cap has been retuned.
It depends on the regulator and geometry, not on
the source vector or on \(\alpha\).

The original Wilson configurations are not truncated.
The actual \(Y_i=\sum_{x\in\mathcal B_{p,i}}T_x\)
and their centering constants are unchanged.
V35's events \(D_i^\Lambda,E_i^\Lambda\) and its
restricted chart comparison with error \(2\Delta\)
remain valid. Its joint failure parameter is
\begin{equation}
 q=\min\{1,\sqrt{q_{\rm str}}+\sqrt{q_Y}\},\qquad
 q_Y=\min\{1,e^{-\Lambda/2+5n_p\log2}\}.
 \label{eq:f15v36:q}
\end{equation}
For every subfamily \(I\),
\(\mu(\cap_{i\in I}E_i^\Lambda)\le q^{|I|}\).
Since \(q_{\rm str}=e^{-\Lambda}\) eventually
and \(n_p/\Lambda\to0\), the proof of V35's
failure-budget estimate gives
\begin{equation}
                              q\le e^{-\Lambda/6}
 \label{eq:f15v36:q-rate}
\end{equation}
eventually, uniformly on the depth range.

\subsection{An actual all-amplitude moment bound for sublinear energies}

Fix \(0<\alpha<1\), and set
\begin{equation}
 O_i=Y_i^\alpha,\qquad
 S_i^{(\alpha)}=O_i-\mathbb E_\mu O_i,\qquad
 z_\alpha=\frac1{1-\alpha},\qquad
 \mathfrak d_\alpha=(1-\alpha)(2\alpha)^{\alpha/(1-\alpha)}.
 \label{eq:f15v36:observables}
\end{equation}
These are specified new observables; they are not a
renaming of the original source \(Y_i-\mathbb E_\mu Y_i\).
For \(u\ge0\) write
\begin{equation}
 \mathcal M_\alpha(u)=\mathfrak d_\alpha u^{z_\alpha}
                                      +5n_p\log2.
 \label{eq:f15v36:moment-budget}
\end{equation}

\begin{proposition}[All-amplitude radial thermal moments]
\label{f15v36:young}
For every nonnegative array \(u_i\) on a subfamily,
\begin{equation}
 \log\mathbb E_\mu e^{\sum_{i\in I}u_i O_i}
       \le\sum_{i\in I}\mathcal M_\alpha(u_i),
 \qquad
 \mathbb E_\mu O_i\le(5n_p)^\alpha.
 \label{eq:f15v36:moments}
\end{equation}
For the physical tilt
\(W_\theta=e^{\sum_i\theta_iS_i^{(\alpha)}}\),
\(|\theta_i|\le\tau<\infty\), its actual order-two
source cost satisfies
\begin{equation}
 h_\theta=\frac1n\log
     \frac{\mathbb E_\mu W_\theta^2}{(\mathbb E_\mu W_\theta)^2}
 \le H_\alpha(\tau):=
              2\tau(5n_p)^\alpha+\mathcal M_\alpha(2\tau).
 \label{eq:f15v36:h}
\end{equation}
\end{proposition}

\begin{proof}
Maximizing over \(y\ge0\) gives the exact scalar
inequality
\begin{equation}
 u y^\alpha\le y/2+\mathfrak d_\alpha u^{z_\alpha}.
 \label{eq:f15v36:young}
\end{equation}
For \(u>0\) the maximizer of \(uy^\alpha-y/2\)
obeys \(y^{1-\alpha}=2\alpha u\); substitution
gives the displayed constant. At \(u=0\) the
inequality is immediate.
Apply it to each \(Y_i\). The source-cell sets
are disjoint, and the inherited Wilson CGF at
coefficient \(1/2\) gives
\[
                  \mathbb E_\mu e^{\frac12\sum_{i\in I}Y_i}
                                        \le2^{5n_p|I|}.
\]
This proves the first assertion. Concavity of
\(y^\alpha\) and \(\mathbb E_\mu Y_i\le5n_p\)
prove the mean bound.

Centering gives \(\mathbb E_\mu W_\theta\ge1\).
Pointwise,
\[
 2\sum_i\theta_i(O_i-\mathbb E_\mu O_i)
       \le2\tau\sum_i O_i+2\tau\sum_i\mathbb E_\mu O_i.
\]
Use the first two assertions to bound the numerator
of the order-two ratio. This proves
\eqref{eq:f15v36:h} for arbitrary mixed signs.
The bound is deliberately not asserted to be
quadratic as \(\tau\to0\); its required feature
here is regulator-uniform finiteness at every
fixed amplitude and depth.
\end{proof}

The energy CGF was used only at coefficient \(1/2\).
No extrapolation beyond its recorded radius has
been made. The all-amplitude conclusion concerns
\(Y^\alpha\), whose growth is strictly sublinear.
Nor is this an estimate conditional on arbitrary
retained or physical boundary data.

\subsection{Local source-log errors with a sublinear reference cap}

Use the common exterior \(\xi\), actual conditional
interiors \(\mathsf P_i^\xi\), and measures
\(\nu_0,\nu_\theta\) as before. Define
\[
 u_i=\mathsf P_i^\xi(e^{\theta_iO_i}),\qquad
 b_i=\lambda_i^\Lambda(e^{\theta_iO_i}),\qquad
 L_i=\log(u_i/b_i),\qquad
 g_i=\log\mathsf P_i^\xi(e^{|\theta_i|O_i}).
\]
Let \(\pi_{i,0},\pi_{i,\theta}\) be the actual
baseline and source-tilted conditional probabilities
of the same event \(E_i^\Lambda\). As in V35,
\[
 B_i=\{\pi_{i,0}>\eta\text{ or }\pi_{i,\theta}>\eta\},
 \qquad e_0=2\Delta-\log(1-\eta),\qquad 0<\eta<1.
\]
Concavity gives, for \(m,x\ge0\),
\((m+x)^\alpha-m^\alpha\le x^\alpha\).
For example differentiate the left side with
respect to \(m>0\); it is nonincreasing in \(m\),
and extend to \(m=0\) by continuity. Hence the
reference support satisfies
\[
           m_i^\alpha\le O_i<m_i^\alpha+\Lambda^\alpha.
\]
Repeating the two elementary sign inequalities
in Proposition~\ref{f15v35:local} now gives
\begin{equation}
 \begin{array}{c|cc}
 &(L_i)_+&(L_i)_-\\ \hline
 \theta_i\ge0&\le g_i&\le|\theta_i|\Lambda^\alpha\\
 \theta_i\le0&\le|\theta_i|\Lambda^\alpha&\le g_i
 \end{array}
 \qquad
 |L_i|\le e_0+(g_i+\tau\Lambda^\alpha)1_{B_i}.
 \label{eq:f15v36:local}
\end{equation}
The restricted chart ratio is still \(2\Delta\),
since the same positive source weight is used
on both normalized restricted core laws. No
small-source expansion of that weight is used.

Conditional Jensen and the matched interior product
law give, for every \(a\ge1\),
\begin{equation}
 \log\mathbb E_{\nu_0}e^{a\sum_{i\in I}g_i}
                         \le |I|\mathcal M_\alpha(a\tau).
 \label{eq:f15v36:conditional-moment}
\end{equation}
The absence of an upper bound on \(a\tau\) is
the distinction from the linear-energy argument.
Set
\begin{equation}
 \gamma_\alpha=q^{1/4}
       \exp\left\{\tfrac14\mathcal M_\alpha(4\tau)
                        +\tfrac12\mathcal M_\alpha(2\tau)\right\}.
 \label{eq:f15v36:gamma}
\end{equation}
The mixed-confidence proof of
Proposition~\ref{f15v35:mixed-confidence}, with
the actual moments \eqref{eq:f15v36:moments},
gives
\begin{equation}
 \mathbb E_{\nu_0}\prod_{i\in I}\pi_{i,a_i}
                  \le\gamma_\alpha^{|I|}
       \quad(a_i\in\{0,\theta\}),\qquad
 \nu_0(\cap_{i\in I}B_i)
                  \le(2\gamma_\alpha/\eta)^{|I|}.
 \label{eq:f15v36:confidence}
\end{equation}
Explicitly, both probabilities are bounded by
\(\mathsf P_i^\xi(1_{E_i^\Lambda}e^{|\theta_i|O_i})
 \mathsf P_i^\xi(e^{|\theta_i|O_i})\).
Conditional Jensen and two Cauchy--Schwarz
applications give the respective factors
\(q^{1/4}\), \(e^{\mathcal M_\alpha(4\tau)/4}\)
and \(e^{\mathcal M_\alpha(2\tau)/2}\). This
verifies the substitution also for mixed signs.

\subsection{A finite-regulator normalization theorem with a free moment parameter}

For any \(v\ge4\) define
\begin{equation}
 \begin{gathered}
 F_\alpha=\sum_i(g_i+\tau\Lambda^\alpha)1_{B_i},\\
 K_\alpha(v)=
       \exp\{v\tau\Lambda^\alpha+
                         \tfrac12\mathcal M_\alpha(2v\tau)\}
                         \sqrt{2\gamma_\alpha/\eta},\qquad
 m_\alpha(v)=\log(1+K_\alpha(v)),\\
 \epsilon_\alpha(v)=e_0+\frac{h_\theta+m_\alpha(v)}v.
 \end{gathered}
 \label{eq:f15v36:epsilon}
\end{equation}
The subset expansion used in V35 gives
\begin{equation}
                \log\mathbb E_{\nu_0}e^{vF_\alpha}
                                                   \le n m_\alpha(v).
 \label{eq:f15v36:laplace}
\end{equation}
Indeed, the moment of the exponential on any
specified set of bad sites is at most
\[
 \left[
  e^{v\tau\Lambda^\alpha+\mathcal M_\alpha(2v\tau)/2}
                         \sqrt{2\gamma_\alpha/\eta}\right]^{|I|}
\]
by \eqref{eq:f15v36:conditional-moment},
\eqref{eq:f15v36:confidence} and Cauchy--Schwarz.
Summing the nonnegative subset expansion proves
\eqref{eq:f15v36:laplace}.

Use the physical centered source
\(W_\theta=e^{\sum_i\theta_iS_i^{(\alpha)}}\), and
define
\begin{equation}
 \begin{gathered}
 A_\theta(C)=\mathbb E_\mu[W_\theta\mid C],\qquad
 T_\theta(C)=\prod_i\lambda_i^\Lambda(e^{\theta_iS_i^{(\alpha)}}),\\
 Z_a=\mathbb E_\mu W_\theta,\qquad
 Z_b=\mathbb E_{\rho_0}T_\theta,\qquad
 d\rho_a=A_\theta\,d\rho_0/Z_a,\qquad
 d\rho_b=T_\theta\,d\rho_0/Z_b,\\
 r_\theta=\log(A_\theta/T_\theta),\qquad
 \delta_\theta=\log(Z_b/Z_a).
 \end{gathered}
 \label{eq:f15v36:laws}
\end{equation}

\begin{theorem}[Sublinear-source normalization at finite regulator]
\label{f15v36:finite}
For every finite \(\tau>0\), every fitting family,
and every \(\theta\in[-\tau,\tau]^n\), the preceding
budgets imply
\begin{equation}
 \log\mathbb E_{\rho_a}e^{\pm w r_\theta}
       \le nw\epsilon_\alpha(v)\quad(0<w\le v/2),
 \qquad
 \mathbb E_{\rho_a}|r_\theta|\le2n\epsilon_\alpha(v).
 \label{eq:f15v36:exp}
\end{equation}
Consequently
\begin{equation}
 \begin{gathered}
 |\delta_\theta|\le n\epsilon_\alpha(v),\qquad
 D(\rho_a\Vert\rho_b)\le2n\epsilon_\alpha(v),\\
 D(\rho_b\Vert\rho_a)\le4n\epsilon_\alpha(v),\qquad
 d_{\rm TV}(\rho_a,\rho_b)
             \le\min\{1,\sqrt{n\epsilon_\alpha(v)}\}.
 \end{gathered}
 \label{eq:f15v36:normalized}
\end{equation}
Zero source gives exactly the common retained baseline.
\end{theorem}

\begin{proof}
The source parameter restriction in the statement
of V35 cannot simply be dropped. Instead use the
new local and moment bounds just proved in its
normalization argument. For completeness, exact
conditional normalization and Jensen give
\[
 (r_\theta)_-\le ne_0+\mathbb E_{\nu_0}[F_\alpha\mid C],
 \qquad
 (r_\theta)_+\le ne_0+\mathbb E_{\nu_\theta}[F_\alpha\mid C].
\]
The actual projected densities have order-two
moments at most \(e^{nh_\theta}\).
Apply the three-factor H\"older inequality
\eqref{eq:f15v35:transport} with exponent \(v/w\),
using conditional Jensen and
\eqref{eq:f15v36:laplace}. This proves both
exponential bounds for \(v/w\ge2\).
The entropy variational inequality with parameter
\(v\) bounds the mean of each projected nonnegative
term by \(n(h_\theta+m_\alpha(v))/v\), proving
the absolute first-moment assertion.

Now
\(\delta_\theta=\log\mathbb E_{\rho_a}e^{-r_\theta}\).
The \(w=1\) bounds and Cauchy--Schwarz give
the two pressure directions. The forward
entropy identity is
\(D(\rho_a\Vert\rho_b)=\mathbb E_{\rho_a}r_\theta+\delta_\theta\).
For the reverse direction use
\[
 D(\rho_b\Vert\rho_a)
 \le\log\mathbb E_{\rho_a}e^{-2r_\theta}-2\delta_\theta.
\]
The \(w=2\) bound is available because \(v\ge4\).
These prove both entropy estimates; Pinsker gives TV.
Thus no missing source-radius hypothesis from
V35 has been reused.
\end{proof}

\subsection{A regulator-vanishing fixed-strength intensive budget}

Fix \(0<\alpha<1\) and \(T<\infty\). For
\(0<\tau\le T\), choose
\begin{equation}
             v=\frac{\Lambda^{1-\alpha}}{256\tau},
             \qquad \eta=\gamma_\alpha^{1/4}.
 \label{eq:f15v36:choice}
\end{equation}
Eventually \(v\ge4\) and \(0<\eta<1\), uniformly
over \(p\le t\) and \(0<\tau\le T\).

\begin{theorem}[Fixed-strength closure for all fitting sublinear-source families]
\label{f15v36:rate}
Define
\begin{equation}
 C_{\alpha,T}=256T\left[
       2T20^\alpha+\mathfrak d_\alpha(2T)^{z_\alpha}
                                      +20\log2\right].
 \label{eq:f15v36:constant}
\end{equation}
Uniformly over the indicated depths and source cube,
eventually
\begin{equation}
 \begin{split}
 \epsilon_\alpha(v)
 &\le \overline\epsilon_{\alpha,T}:=
       2\beta^{-s}
       +C_{\alpha,T}p^4\Lambda^{-(1-\alpha)}
       +3e^{-\Lambda/512}\\
 &\le 2\beta^{-s}
   +C_{\alpha,T}(2/b)^{1-\alpha}
       \beta^{-(1-\alpha)}
       p^{\,4+d(1-\alpha)}t^{\,3(1-\alpha)}
   +3e^{-\Lambda/512}.
 \end{split}
 \label{eq:f15v36:rate}
\end{equation}
In particular every intensive error in
\eqref{eq:f15v36:normalized} tends to zero at
fixed source amplitude, uniformly over all
fitting families. The common uniform rate is
\begin{equation}
 \overline\epsilon_{\alpha,T}
 =O_{\alpha,T,s}\left(
        \beta^{-s}
       +\beta^{-(1-\alpha)}
                t^{\,4+(d+3)(1-\alpha)}\right).
 \label{eq:f15v36:uniform-rate}
\end{equation}
Total TV convergence requires
\(n\overline\epsilon_{\alpha,T}\to0\), as distinct
from the unrestricted intensive conclusion.
\end{theorem}

\begin{proof}
By \eqref{eq:f15v36:q-rate},
\[
 \log\gamma_\alpha
 \le-\Lambda/24+
       \tfrac14\mathcal M_\alpha(4T)
       +\tfrac12\mathcal M_\alpha(2T).
\]
The last two terms are \(O_{\alpha,T}(n_p+1)\).
Since \(\Lambda/n_p\to\infty\), eventually
\(\gamma_\alpha\le e^{-\Lambda/32}\), uniformly.
The reference-cap cost in \(K_\alpha(v)\) is
\(v\tau\Lambda^\alpha=\Lambda/256\).
The growing actual-moment cost is also paid:
\begin{equation}
 \frac12\mathfrak d_\alpha(2v\tau)^{z_\alpha}
 =\frac{1-\alpha}{256}
       \left(\frac{\alpha}{64}\right)^{z_\alpha-1}\Lambda
 \le\frac{\Lambda}{256}.
 \label{eq:f15v36:young-budget}
\end{equation}
Thus, using \(\sqrt{2\gamma_\alpha/\eta}
=\sqrt2\,\gamma_\alpha^{3/8}\),
\[
 \log K_\alpha(v)
 \le\frac{\Lambda}{128}-\frac{3\Lambda}{256}
                  +\frac{5n_p}{2}\log2+\frac12\log2
 \le-\Lambda/512
\]
eventually. Consequently
\(m_\alpha(v)\le e^{-\Lambda/512}\) and
\(-\log(1-\eta)\le2e^{-\Lambda/128}\).

Substitute these bounds and \eqref{eq:f15v36:h}
in \eqref{eq:f15v36:epsilon}. The source term is
\[
 \frac{h_\theta}{v}
       \le256\tau H_\alpha(\tau)\Lambda^{-(1-\alpha)}.
\]
Since \(n_p\le4p^4\) and \(p\ge1\),
\[
 H_\alpha(\tau)\le
 p^4[\,2T20^\alpha+
       \mathfrak d_\alpha(2T)^{z_\alpha}+20\log2\,].
\]
Also \(1/v\le1/4\) eventually. Together with
\(\Delta\le\beta^{-s}\), these prove the first
bound in \eqref{eq:f15v36:rate}. Use
\eqref{eq:f15v36:entry-scale} for the second.
The exponential term is smaller than every power
of \(\beta^{-1}\), uniformly for \(p\le t\).
Finally substitute \(p\le t\) to obtain
\eqref{eq:f15v36:uniform-rate}.
\end{proof}

At fixed depth, the square-root case has the
particularly transparent budget
\[
 \overline\epsilon_{1/2,T}
       =O_{p,T,s}\bigl(\beta^{-s}
                         +\beta^{-1/2}(\log\beta)^{3/2}\bigr).
\]
This is a fixed-amplitude result for an unbounded
thermal observable in the regulator limit, not
merely for an unscaled bounded Wilson loop.
However the constants and thresholds are not
uniform as \(\alpha\uparrow1\). At the linear
endpoint the maximizing Young problem no longer
has an all-amplitude bound, and
\(\Lambda^{1-\alpha}\) no longer supplies a
growing moment parameter. The endpoint has
not been resolved by taking a limit of this theorem.

\subsection{The original energy sources do close for a fixed finite family}

There is a further consequence for the original
linear sources \(S_i=Y_i-\mathbb E_\mu Y_i\).
It uses V35's bounds for the same new cap
\eqref{eq:f15v36:cap}, not a replacement of the
linear source by a sublinear one.

\begin{proposition}[Fixed-family linear-energy normalization and both output entropies]
\label{f15v36:finite-energy}
Fix \(p,n\) and \(0<\tau_0\le1/512\). For the
linear energy-source outputs defined as in V35
with the retuned capped kernel, uniformly over
\(\theta\in[-\tau_0,\tau_0]^n\),
\begin{equation}
 |\delta_\theta|,\quad
 D(\rho_a\Vert\rho_b),\quad D(\rho_b\Vert\rho_a),
 \quad d_{\rm TV}(\rho_a,\rho_b)
                                   \longrightarrow0.
 \label{eq:f15v36:finite-energy}
\end{equation}
The centres may vary provided the same separated
geometry fits. This is a fixed-family statement,
not an arbitrary-size intensive endpoint theorem.
\end{proposition}

\begin{proof}
The V35 finite-budget theorem remains applicable.
Its subsequent cubic estimate used
\(\Lambda/n_p\to\infty\) and a sufficiently strong
old failure bound relative to \(\Lambda\).
Here those conditions follow from
\eqref{eq:f15v36:entry-scale} and
\(q_{\rm str}=e^{-\Lambda}\). Its proof therefore
gives, for this cap as well,
\[
 \epsilon_Y\le2\beta^{-s}+8192p^4\tau_0^3
                                      +3e^{-\Lambda/512},
 \qquad
 \mathbb E_{\rho_a}e^{\pm2r_\theta}\le e^{2n\epsilon_Y}.
\]
At fixed \(n,p,\tau_0\) these exponential moments
are uniformly bounded.

Independently, the own-law clipped-log argument
of V27 applies to the restricted chart
\(D_i^\Lambda\), with chart error \(2\Delta\),
raw failure parameter \(q\) from
\eqref{eq:f15v36:q}, and the unchanged compact
source bound. With
\[
 K_Y=96\tau_0\beta p^4,\qquad
 c_Y=K_Y/(1-e^{-K_Y}),\qquad h_Y\le64p^4\tau_0^2,
\]
it yields
\begin{equation}
 \ell_\beta:=
 \sup_{\theta\in[-\tau_0,\tau_0]^n}
             \mathbb E_{\rho_a}|r_\theta|
 \le2n[\,2\Delta+c_Y e^{nh_Y/2}\sqrt q\,]
                                                 \longrightarrow0.
 \label{eq:f15v36:energy-own-log}
\end{equation}
All the required conditional normalizer identities
hold also on the zero-sublevel branch by V35.
The convergence follows from \(c_Y=O(\beta)\),
fixed \(n,p\), and
\(q\le e^{-\Lambda/6}\), with
\(\Lambda\gtrsim\beta/(\log\beta)^3\).

Thus \(r_\theta\to0\) in actual-law probability,
uniformly in the cube. The bounded
\(\mathbb E_{\rho_a}e^{-2r_\theta}\) makes
\(e^{-r_\theta}\) uniformly integrable.
Splitting at \(|r_\theta|\le\varepsilon\)
and then using Cauchy--Schwarz on its complement
proves
\(\mathbb E_{\rho_a}e^{-r_\theta}\to1\)
uniformly. Exact normalization gives
\(\delta_\theta\to0\), and
\eqref{eq:f15v36:energy-own-log} gives the forward
entropy convergence.

For the reverse direction the exact identity is
\[
 D(\rho_b\Vert\rho_a)
   =e^{-\delta_\theta}
       \mathbb E_{\rho_a}[(-r_\theta)e^{-r_\theta}]
                                     -\delta_\theta.
\]
On \(r_\theta\ge0\), the absolute integrand
\(r_\theta e^{-r_\theta}\) is bounded by \(1/e\)
and converges in probability to zero. On
\(r_\theta<-L\), writing \(x=-r_\theta\),
\[
 x e^x\le
       \left(\sup_{x\ge L}xe^{-x}\right)e^{2x}.
\]
The exponential moment is uniformly bounded
and the coefficient tends to zero as \(L\to\infty\).
On the remaining bounded interval convergence
in probability suffices. This proves reverse
entropy convergence uniformly as well.
Pinsker completes the assertion.
\end{proof}

This completes a finite-family implication left
unused in V35. It does not contradict V34:
its unconditioned Haar fallback does not have
the bounded complementary exponential error
moments used here. The constants in this proof
depend on \(n\); a fixed-family uniform-integrability
argument is not a volume-uniform source theorem.

\subsection{Polynomial thermal moments transfer without differentiating a real error bound}

At baseline define the joint comparison probability
\begin{equation}
 dQ_0(C,V_1,\ldots,V_n)=
      d\rho_0(C)\prod_{i=1}^n\lambda_i^\Lambda(dV_i\mid N_i).
 \label{eq:f15v36:joint-reference}
\end{equation}
Its retained marginal is exactly \(\rho_0\).
It is the joint core reference, not a claim that
the actual cores are independent at fixed \(C\).
For this subsection take \(\alpha=1/2\), fix
\(p,n\), and compare the vectors
\(O_i=\sqrt{Y_i}\) under the actual law \(P=\mu\)
and under \(Q_0\).

Let
\(M_P(a)=\mathbb E_Pe^{\sum_i a_iO_i}\) and
\(M_Q(a)=\mathbb E_{Q_0}e^{\sum_i a_iO_i}\).
The centering scalars in the source comparison
are identical, so
\(\log(M_Q(a)/M_P(a))=\delta_a\).
For any fixed \(T>0\), set
\begin{equation}
 B_T=e^{n\mathcal M_{1/2}(T)},\qquad
 \mathfrak e_\beta=B_T
                   (e^{n\overline\epsilon_{1/2,T}}-1).
 \label{eq:f15v36:raw-laplace-error}
\end{equation}
The proved bounds imply
\[
 \sup_{a\in[0,T]^n}|M_Q(a)-M_P(a)|
       \le\mathfrak e_\beta\longrightarrow0,
 \qquad
 M_P(T\mathbf1),M_Q(T\mathbf1)\le2B_T
\]
eventually. The reference exponential bound is
a consequence of the paid normalization theorem,
not an additional hypothesis.

\begin{theorem}[All fixed mixed polynomial energy moments]
\label{f15v36:moment-transfer}
For every fixed multiindex \(k\in\mathbb N_0^n\)
with \(|k|>0\),
\begin{equation}
 \left|\mathbb E_{Q_0}\prod_iY_i^{k_i}
                   -\mathbb E_\mu\prod_iY_i^{k_i}\right|
       =O_{p,n,k,T}
          \left(\overline\epsilon_{1/2,T}^{\,1/(2|k|+1)}\right)
                                  \longrightarrow0.
 \label{eq:f15v36:moment-transfer}
\end{equation}
The degree-zero identity is exact.
\end{theorem}

\begin{proof}
Here is an explicit remainder argument; it avoids
formally differentiating the small real pressure
error. Put \(a_i=2k_i\), \(D=\sum_i a_i=2|k|\),
and for a nonnegative vector \(x\) let
\[
       P_h(x)=\prod_i\left(\frac{e^{hx_i}-1}{h}\right)^{a_i},
       \qquad 0<h\le T/(4D).
\]
Expanding the numerator expresses its expectation
as a finite signed combination of Laplace transforms
at vectors \(h j\), \(0\le j_i\le a_i\), divided
by \(h^D\). The sum of the absolute coefficients
is \(2^D\). Hence
\begin{equation}
       |\mathbb E_{Q_0}P_h(O)-\mathbb E_\mu P_h(O)|
                            \le 2^D\mathfrak e_\beta h^{-D}.
 \label{eq:f15v36:difference}
\end{equation}

The integral identity
\((e^{hx}-1)/h=x\int_0^1e^{shx}\,ds\)
shows \(P_h(x)\ge\prod_i x_i^{a_i}\) and
\[
 0\le P_h(x)-\prod_i x_i^{a_i}
 \le hD\left(\sum_i x_i\right)^{D+1}
                     e^{hD\sum_i x_i}.
\]
Since \(hD\le T/4\), the uniform positive
exponential moment at \(T\mathbf1\) bounds
the expectation of this remainder by \(C h\)
under both laws. For instance use
\[
 u^{D+1}\le(D+1)!(4/T)^{D+1}e^{Tu/4},\qquad u\ge0,
\]
and then the moment at \(T\mathbf1\).
Combining the two remainders with
\eqref{eq:f15v36:difference} gives
\[
 \left|\mathbb E_{Q_0}O^a-\mathbb E_\mu O^a\right|
             \le C h+2^D\mathfrak e_\beta h^{-D}.
\]
If \(\mathfrak e_\beta>0\), choose
\(h=\mathfrak e_\beta^{1/(D+1)}\), which is
admissible eventually. If it is zero, send
\(h\downarrow0\). At fixed \(p,n,T\),
\(\mathfrak e_\beta=O(\overline\epsilon_{1/2,T})\).
Finally \(O_i^{2k_i}=Y_i^{k_i}\), proving the result.
\end{proof}

This transfers moments of the original thermal
energies, not exponential energy moments at arbitrary
strength, and not a physical continuum field law.
The family is finite and separated in the declared
lattice geometry. Its retained baseline correlations
have not been removed from \(Q_0\).

\subsection{Scope of the advance and the next upstream issue}

The full fixed-strength extensive result here is
for \(Y^\alpha\), \(0<\alpha<1\). For the original
linear energy it is only the fixed-family result
\eqref{eq:f15v36:finite-energy}; the arbitrary-family
linear-energy budget still has V35's cubic remainder.
These scopes must not be merged. The same capped
reference is used throughout, while its source
observables and moment estimates are stated separately.

The next missing physical interface is not another
formal derivative of a local normalizer. It is a
volume-uniform control of connected retained
responses, or a stronger estimate in genuinely
inhomogeneous Wilson source backgrounds, sufficient
to treat the linear-energy endpoint and an actual
RG iteration. The present comparison keeps the
correlated retained baseline intact. It therefore
does not establish a mixing rate, reflection-positive
transfer spectrum, or a calibrated Hamiltonian gap.
The inherited selected-coupling inputs themselves
also remain distinct from an interacting physical
continuum construction.

\subsection{Exact diagnostics and predecessor preservation}

The new diagnostic is
\begin{center}\small
\path{scripts/verify_ym_f15_v36_sublinear_energy_sources.py}.
\end{center}
It checks 216 rational Young inequalities and
36 attained optimizers at
\(\alpha=r/(r+1)\), \(r=1,\ldots,6\).
There are 150 raised-power concavity comparisons
and 24 exact retuning budget identities.

Six positive finite models reuse the preserved
V35 probability kernels and chart supports, but
not their old energy offsets. Here
\(O=2k\log2\), \(Y=O^2\),
\(k\in\{0,1,3,7\}\), source coefficients are
\(0,\pm1/2\), and \(\Lambda=16(\log2)^2\).
The source weights are consequently \(1,2^k,2^{-k}\).
The energy cap selects \(k=0,1\) on the chart;
the fallback \(Y<1\) selects \(k=0\).
The latter follows from \(2\log2>1\).
The zero-chart-sublevel-mass branch is also retained.

The 72 source/moment-parameter patterns use
\(v=4,8\) and include mixed, sparse and zero
vectors. They check 36,432 local radial sign caps,
288 matched high-moment subsets, 936 mixed
confidence products, 288 weighted exceptional
subsets and 72 exceptional Laplace bounds.
There are 144 projected actual second moments,
288 signed projected exponential moments,
72 two-sided pressure brackets,
72 reverse-entropy moment identities and
12 exact zero-source outputs.

Separate exact polynomial tests check 18
leading mixed-moment cancellations, 72 Laplace
finite-difference identities and 72 two-measure
remainder bounds. They use the binomial parameter
\(u=e^h-1\); the leading coefficient is the
same moment at \(u=0\). The analytic exponential
remainder argument for general thermal fields
is the proof above, not a conclusion from sampling.
All diagnostic comparisons use rational arithmetic,
integer powers or exact polynomial coefficients.
They do not certify the inherited Wilson CGF,
its asymptotic thresholds or the physical mass gap.

The script has 12,437 bytes and SHA--256
\begin{center}\scriptsize\ttfamily
93D58B14C2ECF6560D92930E871F9C4A0E5AE05BEB69C0E3C0E8ED0DC61D216D.
\end{center}
The V35 predecessor has 813,340 bytes and SHA--256
\begin{center}\scriptsize\ttfamily
F7CBE1013CF96D01C38C7AA2587D8E69C406C5B41434686331DCE1BD86B519A4.
\end{center}
Removing this terminal section and reversing the
title version recovers V35 exactly. V36 replaces
V35 as the sole active main note. All 14 protected
sources and 78 earlier certificates are preserved.
The next companion checkpoint remains V40.

\begin{firewall}
V36 proves fixed-strength intensive normalization
for the specified sublinear thermal sources over
all fitting separated families, completes the
fixed-family linear-energy source comparison, and
transfers every fixed polynomial energy moment
to the capped joint reference. It does not close
the extensive fixed-strength linear-energy
endpoint, prove baseline retained clustering,
construct the interacting physical continuum,
or establish the full Yang--Mills mass gap.
The original goal remains active.
\end{firewall}

\section{V37: capped rare phases and the actual-background criterion}
\label{f15v37:section}

V36 closes fixed-family linear-energy comparison and
all-family sublinear-source comparison. The remaining
linear-energy question is genuinely different: can
the cubic intensive remainder in V35 be removed by
using those same baseline probability estimates more
efficiently? The answer below is negative for that
abstract collection of estimates. A new finite
countermodel has only one potentially failing core,
an energy-capped reference, independent true energies,
the full joint failure bound, and vanishing baseline
relative entropy in both directions. Its normalized
extensive-source outputs nevertheless separate.

This is not another instance of either earlier
counterexample. V30 did not have the required joint
failure estimate. V34 used an unconditioned Haar
fallback that V35 has now replaced. Here every
reference energy lies strictly below the cap, and
the fallback is never used. The rare amplification
comes from the many other physical sources.

The example is not a Wilson gauge model, and its
retained statistic has no claimed lattice locality.
It refutes an inference from the retained probability
bounds alone, not a theorem about the specified
Wilson construction. The positive part of this
section identifies a sufficient estimate entirely
under the actual source-tilted exterior law. Neither
a hypothetical reference CGF nor a hybrid sequence
of partly replaced background laws is needed.

V36 and all 93 protected sources and certificates
were checked unchanged before this append. The last
companion checkpoint is V35; the next remains V40.

\subsection{A one-defect model with an extensive independent background}

Fix \(0<\tau\le1/512\). In this section only, let
\begin{equation}
 \begin{gathered}
 m=\tanh\tau,\qquad r=m+\tau^{3/2}/8,\qquad
 I(r)=\frac{1+r}{2}\log(1+r)+\frac{1-r}{2}\log(1-r),\\
 I_\tau(r)=I(r)-\tau r+\log\cosh\tau,\qquad
 \Lambda_N=NI(r),\qquad A_N=\Lambda_N/2 .
 \end{gathered}
 \label{eq:f15v37:parameters}
\end{equation}
Here \(0<m<r<1/2\). Take \(N\) sufficiently large
that \(A_N\ge1\), and choose
\begin{equation}
                 e^{-8A_N}\le\varepsilon_N\le e^{-4A_N}.
 \label{eq:f15v37:epsilon}
\end{equation}
For example \(\varepsilon_N=e^{-4A_N}\) is allowed.
The flexibility will permit rational finite diagnostics.

The parameter \(\tau\) first selects a sequence of
countermodels. Once selected, their actual laws,
failure events, caps and reference kernels are held
fixed when any source vector is varied. Later the
particular uniform source equal to this \(\tau\)
will expose the obstruction. This is not a
source-dependent change of the reference kernel.

Under the actual law \(P_N\), let \(X_1,\ldots,X_N\)
be independent fair signs, and independently let
\(U\) be Bernoulli with success probability
\(\varepsilon_N\). Define
\begin{equation}
 \begin{gathered}
 H_N=\left\{N^{-1}\sum_{j=1}^N X_j\ge r\right\},\qquad
 C=(X_1,\ldots,X_N),\\
 Y_0=A_NU,\qquad Y_j=1+X_j\quad(1\le j\le N).
 \end{gathered}
 \label{eq:f15v37:model}
\end{equation}
There are \(n=N+1\) sources, including the probe \(0\).
All true energies are independent, nonnegative and
bounded at each regulator. The actual core law is
\(\kappa=\operatorname{Bernoulli}(\varepsilon_N)\);
it is independent of \(C\).

Use \(N_0=1_{H_N}\) as the probe's declared near
input, and \(N_j=X_j\) for a background source.
Dummy independent fair core bits can be attached
at background sites; their energies do not depend
on these bits. Conditional independence of all
cores at the common exterior \(\xi=C\) is then exact.
The fact that \(N_0\) is a collective statistic is
explicit: this construction is not an identification
with a geometrically local Wilson near field.

Define a source-independent comparison kernel by
\begin{equation}
 \lambda_0(\,\cdot\mid N_0)=
 \begin{cases}
   \operatorname{Bernoulli}(\varepsilon_N),&N_0=0,\\
   \operatorname{Bernoulli}(1-\varepsilon_N),&N_0=1.
 \end{cases}
 \label{eq:f15v37:reference}
\end{equation}
At all background sites use the unchanged dummy
core law. Let \(Q_N\) preserve the actual law of
\(C\) and use these product conditional kernels.
Every atom of the \((C,U)\) laws is positive.

This kernel has precisely the energy-cap support
property: at the probe the minimum is zero and
\(Y_0\le A_N<\Lambda_N\); at a background site
the energy is independent of its core, so
\(Y_j-\min_{V_j}Y_j=0\).
It is also of the formal capped-kernel form
\eqref{eq:f15v35:kernel}: take its displayed positive
kernel as the measured chart kernel, with full
chart and retained-chart events. The sublevel
normalizer is exactly one. No assertion is made
that this chart density arises from Wilson
identity-exterior integration.

Set the only possible failure to be
\(E_0=H_N\), with \(E_j=\varnothing\) for \(j\ge1\).
Off \(E_0\) the actual and reference core laws agree
exactly, so their good-set log-ratio error is zero.
This realizes the probability interfaces used in
V35, rather than Wilson's additional geometric
and action-specific structure.

\begin{proposition}[Joint failures and a uniform true-source window]
\label{f15v37:baseline}
With \(q_{\rm str}=e^{-\Lambda_N}\), every nonempty
subfamily satisfies
\begin{equation}
              P_N\left(\bigcap_{i\in J}E_i\right)
                                 \le q_{\rm str}^{|J|}.
 \label{eq:f15v37:joint}
\end{equation}
For every real source vector with \(|t_i|<1\),
\begin{equation}
 \log\mathbb E_{P_N}
       e^{\sum_{i=0}^N t_i(Y_i-\mathbb E_{P_N}Y_i)}
       \le 2\sum_{i=0}^N J(t_i),\qquad
 J(t)=-\log(1-t)-t .
 \label{eq:f15v37:CGF}
\end{equation}
The constants are independent of \(N,\Lambda_N\).
Also \(\mathbb E_{P_N}Y_i\le1\), and, for
\(0\le t_i<1\),
\begin{equation}
       \log\mathbb E_{P_N}e^{\sum_i t_iY_i}
                      \le 2\sum_i[-\log(1-t_i)].
 \label{eq:f15v37:uncentered}
\end{equation}
In particular these bounds imply the weaker
constant-\(5n_p\) thermal moment inputs for
any \(n_p\ge1\).
\end{proposition}

\begin{proof}
The Chernoff bound for independent fair signs,
optimized at \(\operatorname{arctanh}r\), gives
\(P_N(H_N)\le e^{-NI(r)}=e^{-\Lambda_N}\).
Every other nonempty failure intersection is zero.
Thus even a subfamily containing the one bad probe
and arbitrarily many background sites obeys
\eqref{eq:f15v37:joint}.

For a background energy, its centered CGF is
\(\log\cosh t\le t^2/2\).
For the probe let
\(K(t)=\log\mathbb E e^{t(Y_0-\mathbb EY_0)}\).
Exponential tilting of its two atoms gives, for
\(|t|\le1\), with \(A=A_N\) and
\(\varepsilon=\varepsilon_N\),
\[
 K''(t)=
 \frac{A^2\varepsilon(1-\varepsilon)e^{tA}}
      {(1-\varepsilon+\varepsilon e^{tA})^2}
 \le\frac{A^2e^{-3A}}{1-e^{-4A}}<1.
\]
The last assertion follows since \(A\ge1\),
\(A^2e^{-3A}\) is decreasing there, and
\(e^{-3}/(1-e^{-4})<1\).
As \(K(0)=K'(0)=0\), integration gives
\(K(t)\le t^2/2\), including negative \(t\).
For \(0\le t<1\), \(J(t)\ge t^2/2\);
for \(t=-u\), \(0\le u<1\),
\[
             J(-u)=\int_0^u\frac{s}{1+s}\,ds\ge u^2/4.
\]
Independence now proves \eqref{eq:f15v37:CGF}.
The probe mean \(A\varepsilon\le Ae^{-4A}<1\)
and every background mean is one. For positive
sources, add their means to the centered estimate:
\(t+2J(t)=-2\log(1-t)-t\le-2\log(1-t)\).
This proves \eqref{eq:f15v37:uncentered}.
\end{proof}

\subsection{Baseline closeness includes entropy, moments and covariance}

Write \(p_N=P_N(H_N)\) and
\(d_N=1-2\varepsilon_N\).
The two baseline laws have the same retained marginal.
Only the probe conditional changes, and only on \(H_N\).

\begin{proposition}[Strong baseline comparisons still do not close the tilt]
\label{f15v37:baseline-close}
One has
\begin{equation}
 \begin{gathered}
 d_{\rm TV}(P_N,Q_N)=d_Np_N\le e^{-\Lambda_N},\\
 D(P_N\Vert Q_N)=D(Q_N\Vert P_N)
   =p_Nd_N\log\frac{1-\varepsilon_N}{\varepsilon_N}
   \le8A_Ne^{-\Lambda_N}\longrightarrow0 .
 \end{gathered}
 \label{eq:f15v37:baseline-close}
\end{equation}
All fixed mixed polynomial energy moments differ
by a quantity tending to zero. More precisely,
for a probe degree \(k\ge1\) and background
multiindex \(a\) of fixed total degree,
\begin{equation}
 \left|\mathbb E_{Q_N}Y_0^k\prod_{j=1}^NY_j^{a_j}
            -\mathbb E_{P_N}Y_0^k\prod_{j=1}^NY_j^{a_j}\right|
        \le A_N^k2^{|a|}e^{-\Lambda_N}.
 \label{eq:f15v37:moments}
\end{equation}
Without a probe factor the moments agree exactly.
For the entire growing vector of energies,
\begin{equation}
 \left\|\operatorname{Cov}_{Q_N}(Y)
          -\operatorname{Cov}_{P_N}(Y)\right\|_{\ell^2\to\ell^2}
       \le (A_N^2+A_N\sqrt N)e^{-\Lambda_N}
                                      \longrightarrow0 .
 \label{eq:f15v37:covariance}
\end{equation}
\end{proposition}

\begin{proof}
The TV distance of the two reversed Bernoulli
conditionals is \(d_N\). Their two relative
entropies are equal to
\(d_N\log((1-\varepsilon_N)/\varepsilon_N)\).
Disintegration over their identical \(C\) marginal
proves the identities. The logarithm is at most
\(-\log\varepsilon_N\le8A_N\).

Conditional on \(C\), the probe's \(k\)-th moment
changes by exactly \(A_N^kd_N1_{H_N}\). A
background monomial is nonnegative and at most
\(2^{|a|}\). Integrating proves the moment bound.

The background covariance matrix is unchanged.
The probe variance difference is exactly
\[
                   A_N^2d_N^2p_N(1-p_N).
\]
Its covariance with background site \(j\), zero
under \(P_N\), changes by
\[
                   A_Nd_N\mathbb E_{P_N}[X_j1_{H_N}],
\]
whose absolute value is at most \(A_Np_N\).
The covariance difference therefore has block
form \(\left(\begin{smallmatrix}a&c^T\\c&0\end{smallmatrix}\right)\)
with \(|a|\le A_N^2p_N\) and
\(\|c\|_2\le A_N\sqrt Np_N\).
The off-diagonal block matrix has norm \(\|c\|_2\).
Triangle inequality proves \eqref{eq:f15v37:covariance}.
Since \(\Lambda_N=NI(r)\) with \(I(r)>0\)
fixed, every stated polynomial prefactor is
dominated by \(e^{-\Lambda_N}\).
\end{proof}

Thus even total baseline entropy convergence
and covariance-operator convergence are insufficient.
The issue is not a lost factor of family size
in those baseline comparisons.

\subsection{Exact tilted outputs and their binomial rate}

Apply the same physical source \(\tau\) to every
centered energy \(Y_i-\mathbb E_{P_N}Y_i\).
The centering scalars are the same for both laws.
Under the actual tilted output \(\rho_{a,N}\),
the signs are independent with mean \(m=\tanh\tau\).
Set \(p_{N,\tau}=\rho_{a,N}(H_N)\), and define
\begin{equation}
 \begin{gathered}
 u_N=1-\varepsilon_N+\varepsilon_Ne^{\tau A_N},
 \qquad b_N=\varepsilon_N+(1-\varepsilon_N)e^{\tau A_N},
 \qquad B_N=b_N/u_N>1,\\
 Z_N=1+p_{N,\tau}(B_N-1),\qquad \delta_N=\log Z_N .
 \end{gathered}
 \label{eq:f15v37:normalizer}
\end{equation}
The tilted comparison output obeys the exact formulas
\begin{equation}
 \frac{d\rho_{b,N}}{d\rho_{a,N}}
      =\frac{B_N^{1_{H_N}}}{Z_N},\qquad
 r_N=-1_{H_N}\log B_N,\qquad
 \rho_{b,N}(H_N)=\frac{p_{N,\tau}B_N}{Z_N}.
 \label{eq:f15v37:tilted-laws}
\end{equation}
These follow by integrating the one probe core.
No approximation or deleted vacuum factor occurs:
\(\delta_N\) is the logarithm of the ratio of the
two full centered-source partition functions.

\begin{proposition}[The actual exceptional-phase rate]
\label{f15v37:binomial}
At fixed \(\tau,r\) as above,
\begin{equation}
 \frac1N\log p_N\longrightarrow-I(r),\qquad
 \frac1N\log p_{N,\tau}\longrightarrow-I_\tau(r)<0,
 \qquad
 \frac1N\log B_N\longrightarrow\frac{\tau I(r)}2 .
 \label{eq:f15v37:rates}
\end{equation}
\end{proposition}

\begin{proof}
For completeness, the binomial tail rate does
not require a quoted limit theorem. If \(K\)
is binomial with parameter \(q\), and
\(k_N=\lceil N(1+r)/2\rceil\), Chernoff gives
the upper rate with relative entropy
\[
 {\cal D}(s\Vert q)=s\log(s/q)+(1-s)\log((1-s)/(1-q)).
\]
For a lower bound keep its one atom at \(k_N\).
For any integer \(0<k<N\),
\[
 \frac1{N+1}e^{N h(k/N)}
       \le \binom{N}{k}\le e^{N h(k/N)},\qquad
 h(s)=-s\log s-(1-s)\log(1-s).
\]
Indeed, under binomial parameter \(k/N\) the
atom \(k\) is a mode, so its probability is at
least \(1/(N+1)\); its probability is at most one.
Multiplying these inequalities by
\(q^k(1-q)^{N-k}\) and taking logarithms proves
the matching lower rate. Continuity in
\(k_N/N\to(1+r)/2\) completes the argument.
At \(q=1/2\) this entropy is \(I(r)\).
At \(q=(1+\tanh\tau)/2\) it is
\(I(r)-\tau r+\log\cosh\tau\).
It is positive because \(r>\tanh\tau\).

Finally \(\varepsilon_Ne^{\tau A_N}\to0\)
by \eqref{eq:f15v37:epsilon}, whereas
\(b_N/e^{\tau A_N}\to1\) and \(u_N\to1\).
Thus \(\log B_N=\tau A_N+o(1)\),
which proves the last limit.
\end{proof}

\subsection{A cubic extensive obstruction despite the cap}

Define the positive gap between the reference
reward and the actual rare-phase cost:
\begin{equation}
                 G(\tau)=\frac{\tau I(r)}2-I_\tau(r).
 \label{eq:f15v37:gap}
\end{equation}
This is a pressure discrepancy, not a physical
spectral gap.

\begin{theorem}[A capped-reference extensive-source countermodel]
\label{f15v37:obstruction}
For every fixed \(0<\tau\le1/512\), the preceding
positive finite models satisfy
\begin{equation}
 \begin{gathered}
 G(\tau)\ge3\tau^3/64>\tau^3/32,\qquad
 \frac{\delta_N}{N+1}\longrightarrow G(\tau),\\
 d_{\rm TV}(\rho_{a,N},\rho_{b,N})\longrightarrow1,\qquad
 \mathbb E_{\rho_{a,N}}|r_N|\longrightarrow0,\\
 \frac{D(\rho_{a,N}\Vert\rho_{b,N})}{N+1}
       \longrightarrow G(\tau),\qquad
 \frac{D(\rho_{b,N}\Vert\rho_{a,N})}{N+1}
       \longrightarrow I_\tau(r)>0 .
 \end{gathered}
 \label{eq:f15v37:obstruction}
\end{equation}
As \(\tau\downarrow0\), for this specified family
of countermodels,
\begin{equation}
                    G(\tau)=\frac{31}{128}\tau^3
                                         +O(\tau^{7/2}).
 \label{eq:f15v37:cubic}
\end{equation}
In particular an \(o(\tau^3)\) universal replacement
of V35's cubic cost cannot follow from these
baseline interfaces alone.
\end{theorem}

\begin{proof}
One has \(m\ge\tau/2\) in the stated range.
For example \(\tanh x\le x\) gives
\((\tanh x)'=1-\tanh^2x\ge1-x^2\);
integration yields \(m\ge\tau-\tau^3/3\ge\tau/2\).
Also \(r<2\tau<1/2\).
Since \(I''(x)=1/(1-x^2)\ge1\),
\(I(r)\ge r^2/2\ge\tau^2/8\). Hence
\(\tau I(r)/2\ge\tau^3/16\).

The identity \(I'(m)=\tau\) and
\(I(m)=\tau m-\log\cosh\tau\) imply
\[
 I_\tau(r)=\int_m^r(r-x)I''(x)\,dx.
\]
On this interval \(1\le I''(x)\le2\); therefore
\[
             \tau^3/128\le I_\tau(r)\le\tau^3/64.
\]
Subtracting the upper bound from the reward
proves \(G(\tau)\ge3\tau^3/64\).

By \eqref{eq:f15v37:rates},
\(p_{N,\tau}B_N=\exp(NG(\tau)+o(N))\to\infty\).
The exact normalizer formula proves the pressure
limit, while \(p_{N,\tau}\to0\) and
\(\rho_{b,N}(H_N)\to1\) prove TV convergence to one.
Since \(\log B_N=O(N)\),
\(\mathbb E_{\rho_{a,N}}|r_N|
       =p_{N,\tau}\log B_N\to0\).
The exact entropy identities are
\[
 \begin{split}
 D(\rho_{a,N}\Vert\rho_{b,N})
       &=\delta_N-p_{N,\tau}\log B_N,\\
 D(\rho_{b,N}\Vert\rho_{a,N})
       &=\rho_{b,N}(H_N)\log B_N-\delta_N .
 \end{split}
\]
They prove both entropy limits, with division
by \(N+1\) equivalent to division by \(N\).

Finally \(m=\tau+O(\tau^3)\) and
\(r=m+\tau^{3/2}/8\). Taylor expansion at zero
gives \(I(r)=\tau^2/2+O(\tau^{5/2})\).
The integral formula above, using
\(I''(x)=1+O(\tau^2)\) uniformly there, gives
\(I_\tau(r)=\tau^3/128+O(\tau^5)\).
These imply \eqref{eq:f15v37:cubic}.
\end{proof}

There is only one possibly bad probe among
\(N+1\) sources. Its maximal reference energy
grows like \(\Lambda_N\), while the other sources
pay the cost of selecting its retained input.
Thus a vanishing density of bad sites need not
give a vanishing intensive normalization error.
The limit \(N\to\infty\) is essential and does
not contradict V36's fixed-family theorem.

The baseline exceptional-energy bound remains
very strong. Here the confidence event at the
probe is exactly \(H_N\), for every confidence
threshold in \((0,1)\); all others are empty.
For the positive source, \(g_0=\log u_N=o(1)\).
With \(F=(g_0+\tau\Lambda_N)1_{H_N}\) and
\(v=1/(128\tau)\), one has
\begin{equation}
 \log\mathbb E_{P_N}e^{vF}
 \le \exp\{-127\Lambda_N/128+v g_0\}
                                      \longrightarrow0 .
 \label{eq:f15v37:baseline-F}
\end{equation}
Indeed \(e^{vF}\) takes only its two indicated
values and \(p_N\le e^{-\Lambda_N}\).
So the counterexample does not exploit a failure
of V35's baseline exceptional-moment estimate.

Nor does it contradict sublinear comparison.
If \(Y_0\) is replaced by \(Y_0^\alpha\),
\(0<\alpha<1\), its exceptional reference reward
has logarithm \(O(N^\alpha)\), not \(O(N)\).
With the other energy sources unchanged,
\(p_{N,\tau}\exp(O(N^\alpha))\to0\).
Replacing the other sources by their sublinear
versions only lowers their positive sign tilt
from \(\tau\) to \(\tau2^{\alpha-1}<\tau\).
The same tail conclusion follows. In this model
the resulting normalizer and TV errors tend to zero.

\subsection{The sufficient estimate can be placed on the actual exterior}

Return now to the actual Wilson notation of V35--V36,
with the original linear energy sources and their
capped reference. Write \(u_i,b_i,L_i,g_i,B_i\)
as in \eqref{eq:f15v35:local}--\eqref{eq:f15v35:confidence},
and set
\begin{equation}
 \ell=\sum_iL_i,\qquad
 F=\sum_i(g_i+\tau\Lambda)1_{B_i},\qquad
 |\ell|\le ne_0+F,\qquad e_0=2\Delta-\log(1-\eta).
 \label{eq:f15v37:actual-F}
\end{equation}
The common exterior law is
\(\nu_\theta=\xi_*\mu_\theta\), an actual physical
source tilt, not the reference tilt.

\begin{theorem}[Actual-background sufficient criterion]
\label{f15v37:criterion}
Suppose at a finite regulator that
\begin{equation}
                      \log\mathbb E_{\nu_\theta}e^{2F}
                                                \le n\kappa .
 \label{eq:f15v37:criterion}
\end{equation}
Then, with \(\epsilon_{\rm act}=e_0+\kappa/2\),
\begin{equation}
 \begin{gathered}
 |\delta_\theta|\le n\epsilon_{\rm act},\qquad
 D(\rho_a\Vert\rho_b)\le2n\epsilon_{\rm act},\\
 D(\rho_b\Vert\rho_a)\le4n\epsilon_{\rm act},\qquad
 d_{\rm TV}(\rho_a,\rho_b)
                  \le\min\{1,\sqrt{n\epsilon_{\rm act}}\}.
 \end{gathered}
 \label{eq:f15v37:criterion-output}
\end{equation}
No source moment bound for the reference law is
an assumption of this theorem.
\end{theorem}

\begin{proof}
Let \(c_\theta=\exp(-\sum_i\theta_i\mathbb E_\mu Y_i)\).
Conditional independence gives
\[
 d\nu_\theta=
       \frac{c_\theta\prod_i u_i}{Z_a}\,d\nu_0.
\]
Introduce the normalized exterior comparison
\[
 d\widehat\nu_\theta
       =\frac{c_\theta\prod_i b_i}{Z_b}\,d\nu_0.
\]
It is a probability: the \(b_i\) depend only
on their retained \(N_i\)'s, so integrating
\(\nu_0\) gives exactly \(Z_b\). Its projection
onto \(C\) is \(\rho_b\), while that of
\(\nu_\theta\) is \(\rho_a\). Thus
\begin{equation}
 \frac{d\widehat\nu_\theta}{d\nu_\theta}
     =e^{-\ell-\delta_\theta},\qquad
 \delta_\theta=\log\mathbb E_{\nu_\theta}e^{-\ell}.
 \label{eq:f15v37:exterior-identity}
\end{equation}
All normalizers and centering constants have been kept.

The assumed moment implies
\(\mathbb E_{\nu_\theta}F\le n\kappa/2\) and
\(\log\mathbb E_{\nu_\theta}e^F\le n\kappa/2\).
Jensen and \eqref{eq:f15v37:actual-F} then give
both pressure bounds \(|\delta_\theta|\le
n(e_0+\kappa/2)\).
The forward exterior entropy is
\(\mathbb E_{\nu_\theta}\ell+\delta_\theta
\le2n\epsilon_{\rm act}\).
For the reverse direction, Jensen applied to
the log of the likelihood under
\(\widehat\nu_\theta\) yields
\[
 \begin{split}
 D(\widehat\nu_\theta\Vert\nu_\theta)
 &\le\log\mathbb E_{\nu_\theta}e^{-2\ell}
                                      -2\delta_\theta\\
 &\le2ne_0+n\kappa+2n\epsilon_{\rm act}
   =4n\epsilon_{\rm act}.
 \end{split}
\]
Relative entropy decreases under projection onto
\(C\); conditional Jensen for the projected
likelihood proves this directly. Pinsker gives TV.
Only the single actual law \(\nu_\theta\)
was used to pay an exponential error moment.
\end{proof}

This reorganizes the normalization step, not its
elementary entropy inequalities. Its useful
distinction from V35 is the location of the
unproved input: no global order-two change-of-law
cost is inserted between the baseline and the
moment in \eqref{eq:f15v37:criterion}.

\begin{proposition}[A local form of the actual-background target]
\label{f15v37:local-target}
Suppose for every subfamily \(J\), under the
same actual \(\nu_\theta\), one has
\begin{equation}
 \nu_\theta\left(\bigcap_{i\in J}B_i\right)
                \le\widetilde q^{|J|},\qquad
 \mathbb E_{\nu_\theta}e^{4\sum_{i\in J}g_i}
                \le {\cal A}^{|J|}.
 \label{eq:f15v37:local-target}
\end{equation}
Then the criterion holds with
\begin{equation}
 \kappa=\log(1+\zeta),\qquad
 \zeta=e^{2\tau\Lambda}{\cal A}^{1/2}
                                      \widetilde q^{1/2}.
 \label{eq:f15v37:zeta}
\end{equation}
Consequently \(\widetilde q\le e^{-c\Lambda}\),
\(\log{\cal A}=o(\Lambda)\), and \(\tau<c/4\)
would give \(\kappa\to0\). With \(e_0\to0\),
this would close the intensive linear-energy
normalization for all fitting families.
\end{proposition}

\begin{proof}
For a specified nonempty \(J\),
Cauchy--Schwarz bounds
\[
 \mathbb E_{\nu_\theta}
   \left[e^{2\sum_{i\in J}(g_i+\tau\Lambda)}
                       1_{\cap_{i\in J}B_i}\right]
                                  \le\zeta^{|J|}.
\]
Expand
\(e^{2F}=\prod_i[1+1_{B_i}(e^{2(g_i+\tau\Lambda)}-1)]\).
All added factors are nonnegative. Bound each
by its exponential and sum the subset bounds
to get \(\mathbb E_{\nu_\theta}e^{2F}\le(1+\zeta)^n\).
Finally
\(\log\zeta\le(2\tau-c/2)\Lambda+o(\Lambda)\).
The strict inequality on \(\tau\) proves decay.
\end{proof}

These are explicitly additional hypotheses, not
consequences of the original uniform baseline CGF.
The countermodel quantifies their failure:
\[
 \frac{1}{\Lambda_N}\log\rho_{a,N}(H_N)
                 \longrightarrow-\frac{I_\tau(r)}{I(r)},
 \qquad
 \frac{I_\tau(r)}{I(r)}=\frac{\tau}{64}+O(\tau^{3/2}).
\]
The original baseline exponent one has degraded
to order \(\tau\), below the absorption requirement.
Moreover
\[
 \frac1N\log\mathbb E_{\rho_{a,N}}e^{2F}
                 \longrightarrow 2\tau I(r)-I_\tau(r)>0 .
\]
There is no vanishing \(\kappa\) available here.
The exact remote-source amplification is
\(p_{N,\tau}/p_N\), whose logarithmic rate is
\(I(r)-I_\tau(r)\) per background site.
The upstream task is therefore to control such
conditional retained responses using the actual
Wilson action and spatial geometry. Another
baseline moment transfer or another choice of
this cap cannot supply that missing structure.

\subsection{Exact diagnostics, preservation and scope}

The new diagnostic is
\begin{center}\small
\path{scripts/verify_ym_f15_v37_capped_phase_obstruction.py}.
\end{center}
It uses rational probabilities and source weights,
and outward rational logarithm enclosures from
the atanh series with its explicit remainder.
Intermediate powers are rounded outward on a
fixed dyadic grid; this changes neither the
direction nor the rigor of each enclosure.

Four rational thresholds near the analytic
\(r\) give certified positive cubic pressure
rates at source strengths below \(1/512\).
These are checks of explicit rate expressions,
not samples purporting to prove their asymptotics.
Seventeen positive capped finite models use
\(\varepsilon=2^{-8A}\), which satisfies
\eqref{eq:f15v37:epsilon} because
\(1/2<\log2<1\). They check 255 joint-failure
subsets, 34 baseline entropy directions,
68 covariance identities and 204 polynomial
moment identities.

The 306 mixed, sparse and zero source-sign
patterns check 918 normalization/density
identities, 612 entropy identities, 306 TV
identities, 306 actual-background exponential
budgets, 816 pressure/entropy inequalities
and 102 exact zero-probe-source identities.
There are also 92 exact binomial-mode type
bounds. The tests do not certify a Wilson
realization of the countermodel, the analytic
asymptotic proof or a physical spectral gap.

The script has 11,353 bytes and SHA--256
\begin{center}\scriptsize\ttfamily
883752B3632CEAF6509F2997403869BDA3F3DB8928D4926C2C19B44C93204157.
\end{center}
The V36 predecessor has 842,387 bytes and SHA--256
\begin{center}\scriptsize\ttfamily
638EE77F636330EF8B6D616BA3EADD728C0B332F24593E8F208DD53CF2564C58.
\end{center}
Removing this terminal section and reversing the
title version recovers V36 exactly. V37 replaces
V36 as the sole active main note. All 14
protected sources and 79 earlier certificates
are preserved. The next companion checkpoint
remains V40.

\begin{firewall}
V37 proves that the energy-cap support property,
joint baseline failure estimate, independent-source
CGF, baseline entropy convergence and polynomial
moment transfer do not by themselves imply
fixed-strength extensive linear-source closure.
It supplies a positive sufficient criterion under
the actual tilted exterior law, but does not
establish that criterion for Wilson theory.
The Wilson spatial-response estimate, iterated RG,
interacting physical continuum and full Yang--Mills
mass gap remain unproved. The goal remains active.
\end{firewall}

\section{V38: inhomogeneous Wilson innovations and retained frustration}
\label{f15v38:section}

V37 requires information about the actual source
background, not another use of the baseline CGF.
We now return to the Wilson density itself. Its
single-link conditional can be integrated exactly,
also when every plaquette coefficient is different.
The conditional energy above its minimum has a
uniform Gamma domination. A link colouring then
gives an arbitrary-support exponential bound under
the actual inhomogeneous source law.

The homogeneous spherical conditional, its Gamma
limit, and the decomposition into fluctuation and
staple frustration were already proved in f14
V95--V96. They are not claimed anew. The extensions
here are domination at every field strength and
every exterior, arbitrary positive inhomogeneous
coefficients, and its use on all the strip witnesses
without a global source change-of-law penalty.
The resulting unsolved term is an explicit local
six-link frustration observable under the actual
tilted measure. It is not set to zero.

This is an auxiliary exact integration of the
original Wilson law. It changes neither the completed
root blocking map nor the physical source observables.
The fluctuations used in the estimates depend on
the actual source background; they are diagnostic
variables, not replacement physical sources.
V37 and all 94 protected files were checked unchanged.
The next companion checkpoint remains V40.

\subsection{The physical sources give positive inhomogeneous couplings}

Work on an even periodic four-dimensional torus
with side \(M\ge4\). A fine-cell source array
\(b_x\), with \(|b_x|\le\tau<1\), changes the
Wilson law by the normalized weight
\(\exp(\sum_xb_x(T_x-\mathbb E_\mu T_x))\).
The separated sources of V35 correspond to
\(b_x=\theta_i\) on \(\mathcal B_{p,i}\) and
zero elsewhere. Their cell supports are disjoint.

Let \(\mathcal I(p)\) be the four cells incident
to a positive plaquette \(p\) in the symmetric
assignment \eqref{f15v1:cell}. The exact lift,
already used in V2, is
\begin{equation}
 a_p=\frac14\sum_{x\in\mathcal I(p)}b_x,\qquad
 J_p=\beta(1-a_p),\qquad
 (1-\tau)\beta\le J_p\le(1+\tau)\beta .
 \label{eq:f15v38:lift}
\end{equation}
The centering scalar cancels only in the
normalization of this same physical law:
\begin{equation}
 d\mu_J(U)=\mathcal Z_J^{-1}
       \exp\left\{-\sum_pJ_p(1-\tfrac12\Tr U_p)\right\}
                                           \prod_e dH(U_e).
 \label{eq:f15v38:law}
\end{equation}
Thus \(\mu_J=\mu_\theta\) for the indicated
physical source vector. The results about link
conditionals below hold for any deterministic
positive coefficients \(J_p\), even without the
comparability in \eqref{eq:f15v38:lift}.

\subsection{Weighted stars and the energy left after minimizing one link}

Fix an edge \(e\), write \(u=U_e\in S^3\), and
freeze all other links. Orient its six staples
as unit quaternions \(a_{e,p}\) so that
\(\tfrac12\Tr U_p=u\cdot a_{e,p}\) for \(p\ni e\).
This is the same orientation convention as f14 V95.
Define
\begin{equation}
 \begin{gathered}
 B_e=\sum_{p\ni e}J_p,\qquad
 H_e=\sum_{p\ni e}J_pa_{e,p},\qquad h_e=|H_e|,\\
 \mathcal F_e=B_e-h_e,\qquad
 \mathcal Q_e=h_e-u\cdot H_e,\qquad
 E_e=\sum_{p\ni e}J_p(1-\tfrac12\Tr U_p).
 \end{gathered}
 \label{eq:f15v38:star}
\end{equation}
Then
\begin{equation}
                  E_e=\mathcal F_e+\mathcal Q_e,
                  \qquad \mathcal F_e,\mathcal Q_e\ge0.
 \label{eq:f15v38:split}
\end{equation}
Here \(\mathcal F_e\) is exactly the minimum of
\(E_e\) over its link \(u\), not the minimum
of the entire action over all links.
If \(h_e=0\), then \(\mathcal Q_e=0\) for every
link value, while \(\mathcal F_e=B_e\).

There is also the exact weighted identity
\begin{equation}
 \mathcal F_e=
 \frac{\displaystyle\sum_{\substack{p<q\\p,q\ni e}}
             J_pJ_q|a_{e,p}-a_{e,q}|^2}{B_e+h_e}.
 \label{eq:f15v38:frustration}
\end{equation}
Indeed the numerator equals \(B_e^2-|H_e|^2\);
the denominator is positive because \(B_e>0\).
All these scalar quantities are gauge invariant.
The two staples in each pair have the same endpoints,
so their relative quaternion gives a closed loop
of at most six links, none of which is \(e\).

For an explicit unweighted loop observable set
\begin{equation}
 \Omega_e=\sum_{\substack{p<q\\p,q\ni e}}
          (1-a_{e,p}\cdot a_{e,q})\ge0 .
 \label{eq:f15v38:Omega}
\end{equation}
Each summand is \(1-\tfrac12\Tr(a_{e,p}^{-1}a_{e,q})\).
Under \eqref{eq:f15v38:lift},
\begin{equation}
 \frac{\beta(1-\tau)^2}{6(1+\tau)}\,\Omega_e
 \le\mathcal F_e\le
 \frac{\beta(1+\tau)^2}{3(1-\tau)}\,\Omega_e.
 \label{eq:f15v38:loop-comparison}
\end{equation}
This follows by bounding the numerator in
\eqref{eq:f15v38:frustration} between
\(2\beta^2(1-\tau)^2\Omega_e\) and
\(2\beta^2(1+\tau)^2\Omega_e\), and the denominator
between \(6(1-\tau)\beta\) and \(12(1+\tau)\beta\).

\subsection{Uniform Gamma domination at every conditional field}

\begin{theorem}[All-exterior conditional innovation bound]
\label{f15v38:Gamma}
Under \(\mu_J\), conditionally on all links
other than \(e\), the law of \(\mathcal Q_e\)
is stochastically dominated by
\(\operatorname{Gamma}(3/2,1)\), with the second
parameter denoting rate. In particular, for
\(0\le s<1\) and every integer \(k\ge0\),
\begin{equation}
 \begin{gathered}
 \mathbb E_{\mu_J}[e^{s\mathcal Q_e}\mid U_{e^c}]
                                      \le(1-s)^{-3/2},\\
 \mathbb E_{\mu_J}[\mathcal Q_e^k\mid U_{e^c}]
                    \le\frac{\Gamma(k+3/2)}{\Gamma(3/2)}.
 \end{gathered}
 \label{eq:f15v38:Gamma}
\end{equation}
There is no small-field, selected-coupling or
aligned-staple hypothesis.
\end{theorem}

\begin{proof}
The actual conditional density is proportional
to \(e^{u\cdot H_e}\,dH(u)\). If \(h_e>0\),
rotate \(H_e/h_e\) to the first coordinate.
The first coordinate of Haar probability on
\(S^3\) has density
\((2/\pi)\sqrt{1-u_0^2}\) on \([-1,1]\).
The change \(q=h_e(1-u_0)\) gives density
proportional to
\[
           e^{-q}q^{1/2}
              \sqrt{2-q/h_e}\,1_{\{0<q<2h_e\}}.
\]
This density formula was used for the limit
in f14 V96. Its extra factor
\(w_h(q)=\sqrt{(2-q/h)_+}\) is decreasing.

If \(G,G'\) are independent Gamma variables
of shape \(3/2\) and rate one, and \(\varphi\)
is bounded increasing, then
\[
 2\operatorname{Cov}(\varphi(G),w_h(G))
 =\mathbb E[(\varphi(G)-\varphi(G'))
                           (w_h(G)-w_h(G'))]\le0.
\]
Since \(\mathbb Ew_h(G)>0\), division shows
that tilting the Gamma density by \(w_h\)
can only decrease the expectation of every
such \(\varphi\). This proves stochastic
domination. Increasing bounded truncations
give the exponential and polynomial bounds.
The Gamma integral gives their displayed values.
If \(h_e=0\), the variable \(\mathcal Q_e\)
is identically zero, proving the same assertions.
\end{proof}

The conclusion is about the excess above the
conditional minimum. It is false to discard
\(\mathcal F_e\) in \eqref{eq:f15v38:split}.
The already discussed fully frustrated exterior
in f14 V96 has three aligned and three opposite
staples. At equal couplings \(H_e=0\), so its
entire star energy is \(6\beta\), whereas the
innovation is zero. This is a possible conditional
exterior, not a claim that it is probable.

\subsection{Many-link moments in the actual source law}

If a set of links contains at most one edge from
each plaquette, its link conditionals factor
after all other links are fixed. Its \(H_e\)'s
and \(\mathcal F_e\)'s are then measurable from
that common complement. Consequently
\begin{equation}
 \mathbb E_{\mu_J}
   \left[e^{\sum_{e\in A}s_e\mathcal Q_e}
                                      \,\middle|\,U_{A^c}\right]
                          \le\prod_{e\in A}(1-s_e)^{-3/2}
       \quad(0\le s_e<1).
 \label{eq:f15v38:independent-set}
\end{equation}
The set \(A\) is deterministic here. Random
choices of witnesses will be handled by a
finite union bound, not put inside this assertion.

There is an explicit eight-colouring of the
positive edges:
\begin{equation}
          c(x,\mu)=\left(\mu,\ \sum_{\nu\ne\mu}x_\nu
                                                   \pmod2\right).
 \label{eq:f15v38:colour}
\end{equation}
The parity is well defined across an even
periodic seam. In a plaquette of orientations
\(\mu,\nu\), its two \(\mu\)-edges have opposite
second colour, as do its two \(\nu\)-edges;
different orientations have different first
colour. Thus its four edge colours are distinct.

\begin{theorem}[Arbitrary-support innovation exponential moment]
\label{f15v38:many-links}
For every deterministic nonnegative array
\(s_e<1/8\),
\begin{equation}
 \log\mathbb E_{\mu_J}e^{\sum_e s_e\mathcal Q_e}
                   \le-\frac3{16}\sum_e\log(1-8s_e).
 \label{eq:f15v38:many-links}
\end{equation}
This is uniform in the volume, all positive
inhomogeneous couplings, and the support size.
\end{theorem}

\begin{proof}
Write the exponential as the product of its
eight colour factors and use H\"older with
eight equal exponents. In each resulting
expectation apply \eqref{eq:f15v38:independent-set}
at coefficients \(8s_e\) and integrate the
unchanged exterior law. Taking logarithms gives
the factor \((1/8)(3/2)=3/16\).
This is not independence between colour classes.
\end{proof}

In particular any deterministic finite link
set \(A\) obeys
\[
 \mu_J\left\{\sum_{e\in A}\mathcal Q_e\ge u\right\}
       \le \exp\left\{-u/16+(3|A|/16)\log2\right\}.
\]
No global factor \(e^{n h_\theta}\) has been
spent to obtain this actual-source estimate.
The variable being controlled is nonetheless
\(\mathcal Q\), not the original cell energy.

\subsection{A geometric carrier for the original cell energies}

For a positive plaquette with base \(x\) and
orientations \(\mu<\nu\), declare its owner edge
to be \((x,\mu)\). For a set \(I\) of distinct
cell bases let \(\mathcal P(I)\) be all the
plaquettes appearing in their \(T_x\)'s, and
let \(\mathcal E(I)\) be their distinct owner edges.
Then
\begin{equation}
 |\mathcal E(I)|\le24|I|,\qquad
 (1-\tau)\sum_{x\in I}T_x
        \le\sum_{e\in\mathcal E(I)}
                               (\mathcal F_e+\mathcal Q_e).
 \label{eq:f15v38:carrier}
\end{equation}
Indeed each plaquette occurs in at most four
cells, with coefficient \(1/4\) in each, and
\(J_p\ge(1-\tau)\beta\). The star of its owner
contains its weighted energy; adding all owner
stars can only increase this nonnegative sum.

If \(I_1,\ldots,I_k\) are disjoint cell-base
sets, an owner edge belongs to at most twelve
of the sets \(\mathcal E(I_j)\).
An owner edge of orientation \(\mu\) can own
at most three plaquettes, and each belongs
to four cells. Membership in different \(I_j\)'s
must use different cells. The bound twelve
is conservative but suffices. In particular,
expanded carriers need not be assumed disjoint
merely because the original cell bases are.

The actual tilted energy-moment problem can now
be reduced explicitly to frustration. For
\(0\le v<(1-\tau)/16\), Cauchy--Schwarz and
\eqref{eq:f15v38:many-links} give
\begin{equation}
 \begin{split}
 \log\mathbb E_{\mu_J}e^{v\sum_{x\in I}T_x}
 &\le\frac12\log\mathbb E_{\mu_J}
       e^{\,\frac{2v}{1-\tau}
                      \sum_{e\in\mathcal E(I)}\mathcal F_e}\\
 &\quad-\frac{3|\mathcal E(I)|}{32}
                          \log\left(1-\frac{16v}{1-\tau}\right).
 \end{split}
 \label{eq:f15v38:energy-reduction}
\end{equation}
The first term is not a baseline expectation.
It is the remaining actual-background moment
to prove. The second term has been paid under
that same law, for arbitrary source support.

\subsection{Strip failures split into controlled innovations and explicit frustration}

Retain the deterministic rooted fillings of
V32. For a candidate rooted edge in a box
\(\mathcal K\) of side \(L\), let \(I_{y,a}\)
be its distinct cell bases. It has at most
\(3L\) elements. Set
\begin{equation}
 \begin{gathered}
 \Psi_{\mathcal K}
   =\max_{(y,a)}
       \sum_{e\in\mathcal E(I_{y,a})}\mathcal Q_e,\qquad
 \Phi_{\mathcal K}
   =\max_{(y,a)}
       \sum_{e\in\mathcal E(I_{y,a})}\mathcal F_e,\\
 u_{L,\delta}=\frac{(1-\tau)\beta\delta^2}{6\pi^2L},\\
 Z_{\mathcal K}=\{\Psi_{\mathcal K}>u_{L,\delta}/2\},
 \qquad
 A_{\mathcal K}=\{\Phi_{\mathcal K}>u_{L,\delta}/2\}.
 \end{gathered}
 \label{eq:f15v38:strip-variables}
\end{equation}
Empty fillings contribute zero to these maxima.
The original bad-root event satisfies the
pointwise inclusion
\begin{equation}
       \{R_{\mathcal K}>\delta\}
                         \subseteq Z_{\mathcal K}\cup A_{\mathcal K}.
 \label{eq:f15v38:strip-split}
\end{equation}
For a witnessing edge, V32 gives
\(\sum_{x\in I_{y,a}}T_x>
\beta\delta^2/(6\pi^2L)\).
Apply \eqref{eq:f15v38:carrier}; its frustration
sum and innovation sum cannot both be at most
\(u_{L,\delta}/2\). This proves the inclusion.

\begin{theorem}[Joint innovation-witness bound under the physical tilt]
\label{f15v38:strip}
For boxes whose possible filling cell-base sets
are disjoint between boxes, put \(N_L=4L(L+1)^3\) and
\begin{equation}
 q_{\rm inn}=\min\left\{1,\
 N_L\exp\left[-\frac{(1-\tau)\beta\delta^2}{2304\pi^2L}
                                      +\frac98L\log2\right]\right\}.
 \label{eq:f15v38:q}
\end{equation}
Every finite subfamily satisfies
\begin{equation}
             \mu_J\left(\bigcap_{i=1}^kZ_{\mathcal K_i}\right)
                                             \le q_{\rm inn}^{\,k}.
 \label{eq:f15v38:joint-innovation}
\end{equation}
The source array may be extensive and have
arbitrary signs with amplitude at most \(\tau<1\).
\end{theorem}

\begin{proof}
Fix one witnessing candidate in each box.
There are at most \(72L\) owner edges in
each carrier. If \(m_e\) is the number of
selected carriers containing \(e\), then
\[
          0\le m_e\le12,\qquad \sum_e m_e\le72kL.
\]
On the corresponding innovation events,
\(\sum_e m_e\mathcal Q_e>k u_{L,\delta}/2\).
Use exponential Markov at \(s=1/192\).
The coefficients \(s m_e\) in
\eqref{eq:f15v38:many-links} are at most \(1/16\).
For \(0\le m\le12\), convexity on this interval gives
\[
                 -\log(1-m/24)\le(m/12)\log2.
\]
Thus the logarithmic moment cost is at most
\[
     \frac3{16}\frac{\log2}{12}\sum_e m_e
                                      \le\frac98 kL\log2.
\]
The Markov cost is
\(-k u_{L,\delta}/384\), equal to the negative
term in \eqref{eq:f15v38:q} times \(k\).
There are at most \(N_L^k\) candidate choices.
Union bound and the probability bound by one
complete the proof. No independence of distinct
boxes or of their expanded carriers was used.
\end{proof}

On V32's retuned scale, with its \(b,d\) and
\(t=\log\beta\), for every fixed \(\tau_*<1\)
one obtains eventually, uniformly over
\(0\le\tau\le\tau_*\) and \(p\le t\),
\begin{equation}
          \log(1/q_{\rm inn})
             \ge\frac{(1-\tau)b}{384}
                         \frac{\beta}{p^d t^3}.
 \label{eq:f15v38:retuned}
\end{equation}
To verify this, denote V32's negative exponent
by \(E=\beta\delta^2/(12\pi^2L)\).
It satisfies \(E\ge b\beta/(p^dt^3)\).
The new negative term is \((1-\tau)E/192\).
The cost \(\log N_L+(9/8)L\log2\) is polynomial
in \(t\), uniformly in the depth range, and
is eventually at most \((1-\tau)E/384\).
This proves \eqref{eq:f15v38:retuned}.

This superalgebraic actual-source bound is for
the innovation events, not yet for all bad
rooted boxes. If the additional estimate
\begin{equation}
       \mu_J\left(\bigcap_{i\in I}A_{\mathcal K_i}\right)
                                      \le q_{\rm fr}^{|I|}
 \label{eq:f15v38:missing-frustration}
\end{equation}
were established, then the same intersection
expansion and Cauchy--Schwarz used in V35 would give
\[
 \mu_J\left(\bigcap_{i=1}^k\{R_{\mathcal K_i}>\delta\}\right)
                  \le\min\{1,\sqrt{q_{\rm inn}}+\sqrt{q_{\rm fr}}\}^{\,k}.
\]
Equation \eqref{eq:f15v38:missing-frustration}
is not asserted here. In view of
\eqref{eq:f15v38:loop-comparison}, it is a
specified local six-link-loop problem under
the actual inhomogeneous Wilson law.

\subsection{The exact residual measure keeps its normalizing factor}

The remaining frustration is not an unspecified
abstract exterior. Its law after integrating
an independent link set \(A\) has an exact
formula. Define
\begin{equation}
 z(h)=\int_{S^3}e^{-h(1-u_0)}\,dH(u),\qquad z(0)=1.
 \label{eq:f15v38:z}
\end{equation}
Partition the action into the plaquettes touching
\(A\), each assigned to its unique edge in \(A\),
and the remaining plaquettes. Write the latter
action as \(S_{\rm rest}(\xi)\), with
\(\xi=U_{A^c}\). Integrating the actual links gives
\begin{equation}
 d\nu_{J,A}(\xi)=\mathcal Z_J^{-1}
 e^{-S_{\rm rest}(\xi)}
       \prod_{e\in A}\left[e^{-\mathcal F_e(\xi)}z(h_e(\xi))\right]
                                                   dH_{A^c}(\xi).
 \label{eq:f15v38:marginal}
\end{equation}
The same full normalizer \(\mathcal Z_J\) appears;
this is an equality of the actual marginal, not
a comparison measure or a Haar reset.

\begin{proposition}[Uniform finite-field normalizer bounds]
\label{f15v38:z}
For every \(h\ge0\),
\begin{equation}
       \frac1{16}(1+h)^{-3/2}\le z(h)\le3(1+h)^{-3/2}.
 \label{eq:f15v38:z-bounds}
\end{equation}
Consequently the exact effective contribution of
an integrated star is
\begin{equation}
 \mathcal F_e+\frac32\log(1+h_e)+r(h_e),\qquad
                         -\log3\le r(h_e)\le\log16.
 \label{eq:f15v38:effective}
\end{equation}
\end{proposition}

\begin{proof}
For \(0\le h\le1\), one has
\(e^{-2}\le z(h)\le1\), which implies the
displayed bounds. For \(h\ge1\), the same radial
change of variable as above gives exactly
\[
 z(h)=\frac2\pi h^{-3/2}
       \int_0^{2h}e^{-q}q^{1/2}\sqrt{2-q/h}\,dq.
\]
Replacing the last square root by \(\sqrt2\)
and extending the integral gives
\(z(h)\le\sqrt{2/\pi}\,h^{-3/2}
        \le3(1+h)^{-3/2}\).
Restricting to \(0\le q\le1\), the square
root is at least one and \(e^{-q}\ge e^{-1}\).
Hence
\[
             z(h)\ge\frac{4}{3\pi e}h^{-3/2}
                              \ge\frac1{16}(1+h)^{-3/2}.
\]
The simple inequalities \(\pi<4\), \(e<3\)
suffice for the last conservative constant.
Finally put \(r(h)=-\log z(h)-(3/2)\log(1+h)\)
in the exact factor \eqref{eq:f15v38:marginal}.
\end{proof}

The factor \(z(h_e)\), including its dependence
on the exterior staples, cannot be removed from
the measure used to estimate frustration.
The bounded \(r\) above is a per-star bound,
not a small or volume-vanishing error.
Integrating another colour changes the action
on its links through the first colour's factors.
It is therefore invalid to repeat the original
six-staple conditional independently through all
colours. H\"older in
Theorem~\ref{f15v38:many-links} did not do that.

\subsection{What has moved upstream, and what is still required}

The actual source background now has a proved,
volume-uniform innovation moment, including
completely frustrated conditional inputs.
Its paid part of the strip-witness tail is
\eqref{eq:f15v38:joint-innovation}. These estimates
do not rely on transporting V36's baseline
moments or on an unverified non-Abelian correlation
inequality.

The unproved part is the local frustration law
in \eqref{eq:f15v38:marginal}, or equivalently
its weighted six-link observables. A bound on
\(\mathcal F_e\) obtained merely from
\(\mathcal F_e\le E_e\), followed by the same
unknown actual-background energy estimate,
would be circular. The source-tilted frustration
moment in \eqref{eq:f15v38:energy-reduction}
and the joint frustration events in
\eqref{eq:f15v38:missing-frustration}
are concrete remaining targets.

Even the latter event estimate by itself would
not supply every hypothesis of V37: the added
energy-cap failures and the conditional \(g_i\)
moments also need actual-background energy control.
Nor does this one-link conditional theorem give
a mixing rate for the full retained law, a
physical-time transfer spectrum or a continuum
construction. No auxiliary heat-bath rate is
identified with the Yang--Mills mass gap.

\subsection{Exact diagnostics and preservation}

The new diagnostic is
\begin{center}\small
\path{scripts/verify_ym_f15_v38_inhomogeneous_star_innovations.py}.
\end{center}
It checks 384 weighted star identities,
720 six-link trace identities, 96 frame-covariance
identities and eight aligned or zero-field stars
using rational unit quaternions.

For the actual spherical integral it uses the
positive series
\[
       \int e^{h u_0}\,dH(u)
                    =\sum_{k\ge0}\frac{h^{2k}}{4^k k!(k+1)!},
\]
obtained by integrating the even Haar moments.
Geometric bounds on the decreasing tail ratios,
and the analogous positive exponential series,
give outward rational enclosures. They verify
40 Gamma-moment inequalities and eight
normalizer bounds, including \(h=0\).
There are also 45 exact two-copy decreasing-weight
covariance checks. The all-parameter domination
is proved above, not inferred from these tests.

Periodic tori of sides \(4,6,8\) verify 33,888
plaquette colour assignments and the same number
of source-lift coefficients, 5,648 cell carriers
and 16,944 owner-overlap bounds.
There are 3,016 ordered-strip carrier checks
and 25 exact tail-coefficient checks.
All arithmetic is rational, integral or an
outward rational series enclosure; no floating
Monte Carlo estimate is used.

The script has 9,197 bytes and SHA--256
\begin{center}\scriptsize\ttfamily
0B1B677BFBF2568F497A4AB99D8EEADD53606311E89183087D0249A3BBFE3BD7.
\end{center}
The V37 predecessor has 867,616 bytes and SHA--256
\begin{center}\scriptsize\ttfamily
7AB9B3070CEBBD7F4962D89ECC1AA4475F970D089ABA1534945F5EFFC5A83045.
\end{center}
Removing this terminal section and reversing the
title version recovers V37 exactly. V38 replaces
V37 as the sole active main note. All 14
protected sources and 80 earlier certificates
are preserved. The next companion checkpoint
remains V40.

\begin{firewall}
V38 proves actual inhomogeneous-Wilson conditional
Gamma domination, an arbitrary-support innovation
moment, and a joint innovation-witness tail. It
identifies the exact residual frustration measure
and its normalizer. It does not prove its required
frustration tails or moments, V37's full
actual-background criterion, iterated contraction,
an interacting continuum limit or the full
physical Yang--Mills mass gap. The goal remains active.
\end{firewall}

\section{V39: retained penalties and a collective frustration escape}
\label{f15v39:section}

V38 separates conditional innovations from the
actual retained frustration. We now establish two
additional facts about that remaining problem.
First, the exact integrated-star factor has a
global linear frustration penalty, despite the
degeneracy of its derivative at a zero staple sum.
Second, a genuinely collective, eight-link change
of variables suppresses a specified high-frustration
sector uniformly under the physical source tilt.

The sector is chosen because it is invisible to
every one-link innovation simultaneously. On an
explicit center-flux configuration all the
\(\mathcal Q_e\)'s of V38 vanish, although the
Wilson action is positive. This configuration is
not a local minimum with respect to joint changes
of links. A small fixed group of links supplies
an exact action-lowering direction and a finite
Haar-preserving update.

The general injective change-of-variables principle
was already used in the archives, including f6
V26--V27's relative center-sheet erasure. It is
not claimed anew. That construction concerned
center multiplication on a prescribed reduced
fibre, with a separate generated-action debit.
The present calculation uses a noncentral
\(\SU(2)\) rotation of eight links in the full
compact Wilson law, keeps the inhomogeneous
physical couplings, and pays an explicit
neighbourhood error. It does not assert preservation
of a completed coarse output.

The V38 predecessor and all 95 protected files
were checked unchanged. No companion checkpoint
is due at V39; the next iteration, V40, is due
to reread the related programmes.

\subsection{The exact retained star has a global frustration penalty}

Keep V38's positive coefficients \(J_p\), star
quantities \(B_e,h_e,\mathcal F_e=B_e-h_e\),
and a nonempty independent link set \(A\). Define
\begin{equation}
 Z(h)=\int_{S^3}e^{h u_0}\,dH(u)=e^h z(h),\qquad
 \ell(h)=\log Z(h),\qquad
 c(B)=\frac{\ell(B)}B\quad(B>0).
 \label{eq:f15v39:Z}
\end{equation}
The notation \(Z(h)\) here is a one-link
integral, not the full partition function.
For \(0\le F\le B\) put
\begin{equation}
               V_B(F)=\ell(B)-\ell(B-F).
 \label{eq:f15v39:penalty-definition}
\end{equation}

\begin{proposition}[A uniform chord bound for the retained factor]
\label{f15v39:penalty}
For every \(B>0\) and \(0\le F\le B\),
\begin{equation}
 0<c(B)<1,\qquad
       c(B)F\le V_B(F)\le F,\qquad
       e^{-F}\le\frac{Z(B-F)}{Z(B)}\le e^{-c(B)F}.
 \label{eq:f15v39:penalty}
\end{equation}
The function \(c(B)\) is nondecreasing, and
\begin{equation}
       c(B)\ge
       1-\frac{\frac32\log(1+B)+\log16}{B}.
 \label{eq:f15v39:coefficient}
\end{equation}
Nevertheless \(V_B\) is concave in \(F\), with
\(V_B'(B)=0\). The global bound is not a
uniform positive lower bound on that derivative.
\end{proposition}

\begin{proof}
The compact integral has
\(\ell(0)=0\),
\(\ell'(h)=\mathbb E_hu_0\), and
\(\ell''(h)=\operatorname{Var}_h(u_0)\ge0\).
Symmetry at zero and monotonicity imply
\(0\le\ell'(h)<1\) for \(h\ge0\).
Strict Jensen at \(B>0\) gives \(\ell(B)>0\);
also \(Z(B)<e^B\). These prove \(0<c(B)<1\).

Convexity and \(\ell(0)=0\) give
\(\ell(B-F)\le(1-F/B)\ell(B)\), proving
the lower bound on \(V_B\).
Integrating \(\ell'\le1\) between \(B-F\)
and \(B\) gives its upper bound. Exponentiation
gives the ratio statement. The slopes
\(\ell(B)/B\) are nondecreasing by the same
convexity. Finally
\(\ell(B)=B+\log z(B)\); V38's lower bound on
\(z(B)\) proves \eqref{eq:f15v39:coefficient}.
Differentiation gives
\[
 V_B'(F)=\ell'(B-F),\qquad
 V_B''(F)=-\ell''(B-F)\le0.
\]
In particular \(V_B'(B)=\ell'(0)=0\).
\end{proof}

These bounds apply to the exact marginal, including
its normalization. Indeed
\[
         e^{-\mathcal F_e}z(h_e)
                       =z(B_e)e^{-V_{B_e}(\mathcal F_e)}.
\]
The \(B_e\)'s depend only on the prescribed
couplings, not on the exterior configuration.
Thus \eqref{eq:f15v38:marginal} becomes exactly
\begin{equation}
 \begin{gathered}
 d\nu_{J,A}(\xi)=\widehat{\mathcal Z}_{J,A}^{-1}
          e^{-W_{J,A}(\xi)}\,dH_{A^c}(\xi),\\
 W_{J,A}=S_{\rm rest}+\sum_{e\in A}V_{B_e}(\mathcal F_e),
 \qquad
 \widehat{\mathcal Z}_{J,A}
                       =\mathcal Z_J\big/\prod_{e\in A}z(B_e).
 \end{gathered}
 \label{eq:f15v39:marginal}
\end{equation}
No exterior-dependent factor was divided out.
For
\(E_{J,A}=S_{\rm rest}+\sum_{e\in A}\mathcal F_e\)
and \(c_*=\min_{e\in A}c(B_e)\), one has
\begin{equation}
          c_*E_{J,A}\le W_{J,A}\le E_{J,A},
          \qquad
          E_{J,A}=\min_{U_A}S_J(U_A,\xi).
 \label{eq:f15v39:minimum}
\end{equation}
The minimum separates because no plaquette
contains two links of \(A\). The lower bound
also uses \(S_{\rm rest}\ge0\) and \(c_*\le1\).
Under \((1-\tau)\beta\le J_p\le(1+\tau)\beta\),
\(c_*\ge c(6(1-\tau)\beta)\to1\) at fixed
\(\tau<1\).

Equation \eqref{eq:f15v39:minimum} is an
unnormalized action comparison, not stochastic
domination of normalized marginals. Its
partition function and the entropy of the
frustration event still have to be controlled.
The next result pays those quantities by an
explicit map on a particular sector.

\subsection{A nonflat Wilson configuration with zero innovations everywhere}

On an even periodic torus define the central
configuration
\begin{equation}
          \overline U_1(x)=(-1)^{x_2}I,\qquad
          \overline U_\mu(x)=I\quad(\mu=2,3,4).
 \label{eq:f15v39:template}
\end{equation}
Even side length makes it periodic. Its
\((1,2)\)-plaquettes equal \(-I\), while
all other positive plaquettes equal \(I\).
Hence
\begin{equation}
                 S_J(\overline U)
                           =2\sum_{p\parallel(1,2)}J_p>0 .
 \label{eq:f15v39:template-action}
\end{equation}

\begin{proposition}[Simultaneous one-link optimality is not joint optimality]
\label{f15v39:zero-innovations}
Suppose
\((1-\tau)\beta\le J_p\le(1+\tau)\beta\)
and \(0\le\tau<1/3\).
At \(\overline U\), every link innovation
\(\mathcal Q_e\) is zero. Every link is the
unique minimizer of its conditional star action.
For a link of orientation \(1\) or \(2\),
\begin{equation}
 \mathcal F_e(\overline U)
       =2\sum_{\substack{p\ni e\\p\parallel(1,2)}}J_p,\qquad
 h_e\ge(2-6\tau)\beta>0 .
 \label{eq:f15v39:template-star}
\end{equation}
For the other orientations
\(\mathcal F_e=0\) and \(h_e=B_e>0\).
In particular
\begin{equation}
             \sum_e\mathcal Q_e=0,\qquad
             \sum_e\mathcal F_e=4S_J(\overline U)>0.
 \label{eq:f15v39:zero-sum}
\end{equation}
\end{proposition}

\begin{proof}
At a central configuration each oriented
staple vector is either \(\overline U_e\)
or \(-\overline U_e\), with sign equal to
its plaquette's central sign. A link of
orientation \(1\) or \(2\) has two negative
and four positive incident plaquettes.
Therefore
\[
 H_e=(B_e-2C_e)\overline U_e,\qquad
 C_e=\sum_{\substack{p\ni e\\p\parallel(1,2)}}J_p,
 \qquad B_e-2C_e\ge(2-6\tau)\beta .
\]
The vector is nonzero and points in the
direction of \(\overline U_e\), so that link
uniquely minimizes \(B_e-u\cdot H_e\) and
\(\mathcal Q_e=0\), \(\mathcal F_e=2C_e\).
All six incident plaquettes are positive
for orientations \(3,4\), giving the other case.
Every plaquette belongs to four edge stars;
summing \(E_e=\mathcal F_e+\mathcal Q_e\)
gives \eqref{eq:f15v39:zero-sum}.
\end{proof}

This is a configuration in the actual Wilson
space, not an abstract probability countermodel.
It rules out a pointwise bound of total
frustration by total one-link innovation with
no additional term. It says nothing by itself
about the probability of a neighbourhood.
A probability bound requires a collective
calculation. The following action decrease
supplies an upper bound for a neighbourhood
of this configuration.

\subsection{An exact eight-link action-lowering update}

Use local integer coordinates and take the
eight positive links
\begin{equation}
 \mathcal A=\{(x,1):x_1=0,\ (x_2,x_3,x_4)\in\{0,1\}^3\}.
 \label{eq:f15v39:eight}
\end{equation}
For the moment start exactly at \(\overline U\).
On \(\mathcal A\) multiply the link by
\(\exp(i t(-1)^{x_2}\sigma_3)\), leaving all
other links fixed. Here \(\sigma_3\) is the
fixed Pauli matrix, and \(t\) is a real angle.
Let \(\mathcal P_{\mathcal A}\) be all
plaquettes incident to the updated links.

There are 36 such plaquettes. The exact
incidence and phase differences are
\begin{center}
\begin{tabular}{@{}ccl@{}}
\toprule
Orientation & Number & Absolute phase difference divided by \(t\)\\
\midrule
\((1,2)\) & \(4\) & \(2\), interior\\
\((1,2)\) & \(8\) & \(1\), boundary\\
\((1,3)\), \((1,4)\), each & \(4\) & \(0\), interior\\
\((1,3)\), \((1,4)\), each & \(8\) & \(1\), boundary\\
\bottomrule
\end{tabular}
\end{center}
For each transverse direction there are
\(8\cdot2\) edge--plaquette incidences;
the four interior pairs count twice, leaving
twelve distinct plaquettes. In direction \(2\)
the alternating signs differ by two internally.
In directions \(3,4\) they agree internally.
At every boundary one updated value meets zero.

Let \(A_2\) be the sum of the four interior
\((1,2)\) couplings, \(A_-\) the sum of the
eight boundary \((1,2)\) couplings, and \(A_+\)
the sum of the sixteen boundary couplings in
the other two orientations. Directly evaluating
the commuting template plaquettes gives
\begin{equation}
 S_J(U(t))-S_J(\overline U)
       =-A_2(1-\cos2t)+(A_+-A_-)(1-\cos t).
 \label{eq:f15v39:exact-drop}
\end{equation}
The eight interior positive plaquettes are
unchanged at the template, but they must remain
in the affected set for the neighbourhood proof.

\begin{proposition}[Uniform collective decrease]
\label{f15v39:drop}
Writing \(c=\cos t\), the preceding change obeys
\begin{equation}
 S_J(U(t))-S_J(\overline U)
       \le\beta(1-c)[32\tau-8(1-\tau)c].
 \label{eq:f15v39:drop}
\end{equation}
Its second derivative at zero is at most
\((-8+40\tau)\beta\). In particular it is
negative for \(\tau<1/5\), despite
Proposition~\ref{f15v39:zero-innovations}.
At the fixed noncentral rotation
\begin{equation}
                 q=\frac35I+\frac45i\sigma_3,
 \label{eq:f15v39:q}
\end{equation}
the action drop is at least
\begin{equation}
                         \frac{48-368\tau}{25}\,\beta.
 \label{eq:f15v39:rational-drop}
\end{equation}
\end{proposition}

\begin{proof}
Use \(A_2\ge4(1-\tau)\beta\) and
\(A_+-A_-\le[16(1+\tau)-8(1-\tau)]\beta
=(8+24\tau)\beta\) in
\eqref{eq:f15v39:exact-drop}.
The identity \(1-\cos2t=2(1-\cos t)(1+\cos t)\)
proves \eqref{eq:f15v39:drop}.
The second derivative of the exact expression is
\(-4A_2+A_+-A_-\le(-8+40\tau)\beta\).
Finally substitute \(c=3/5\).
\end{proof}

\subsection{A gauge-invariant neighbourhood and a Haar-preserving map}

To extend the decrease beyond the commuting
template, we use a tree that avoids all eight
updated links. Take
\[
             W=\{0,1\}\times\{-1,0,1,2\}^3.
\]
In each fixed-\(x_1\) layer choose the ordered
coordinate tree rooted at \((x_1,-1,-1,-1)\),
with positive steps in directions \(2,3,4\).
Join the two layer roots by their direction-1
link. Call this tree \(\mathcal T\), rooted
at \(o=(0,-1,-1,-1)\). It has 128 vertices
and 127 edges, none belonging to \(\mathcal A\).
All edges of \(\mathcal P_{\mathcal A}\)
lie in \(W\). Use a nonwinding translated copy
of this patch in an even torus of side \(M\ge8\).

For \(x\in W\), let \(g_x(U)\) be transport
along the tree from \(o\) to \(x\). Define
\begin{equation}
        k_e(U)=g_x(U)U_e g_y(U)^{-1}
                   \quad(e=(x,y)),\qquad
        \overline k_e=k_e(\overline U).
 \label{eq:f15v39:framed}
\end{equation}
Let \(\mathcal R\) be all 104 edges of the
36 affected plaquettes. For \(0<\epsilon\le1/1024\)
define the tube event
\begin{equation}
 \mathcal C(\epsilon)=
       \bigcap_{e\in\mathcal R}
          \{|k_e(U)-\overline k_e|\le\epsilon\}.
 \label{eq:f15v39:tube}
\end{equation}
Distances are Euclidean unit-quaternion distances.
The event depends on \(\mathcal R\cup\mathcal T\),
207 distinct links. Since every \(\overline k_e\)
is central, the event is invariant under local
gauge transformations: framed links change by
one common conjugation at the root.

Now define a map on the full configuration
space, with \(q\) fixed as in \eqref{eq:f15v39:q}:
\begin{equation}
 (\Phi U)_e=
 \begin{cases}
  g_x(U)^{-1}q^{\,(-1)^{x_2}}g_x(U)U_e,
                           &e=(x,x+e_1)\in\mathcal A,\\
  U_e,                    &e\notin\mathcal A .
 \end{cases}
 \label{eq:f15v39:map}
\end{equation}
The tree is unchanged by this map, so its
framed action is exactly
\(k_e\mapsto q^{(-1)^{x_2}}k_e\) on \(\mathcal A\).

\begin{proposition}[Exact measure preservation and robust action decrease]
\label{f15v39:tube-drop}
The map \(\Phi\) is a bijection preserving
the original normalized product Haar measure.
If
\[
                  0\le\tau\le1/32,\qquad
                  0<\epsilon\le1/1024,
\]
then on \(\mathcal C(\epsilon)\),
\begin{equation}
                   S_J(\Phi U)\le S_J(U)-\beta.
 \label{eq:f15v39:tube-drop}
\end{equation}
Furthermore this tube is genuinely a
high-frustration event: for every \(e\in\mathcal A\),
\begin{equation}
                       \mathcal F_e(U)\ge3\beta
                      \quad\hbox{on }\mathcal C(\epsilon).
 \label{eq:f15v39:large-frustration}
\end{equation}
\end{proposition}

\begin{proof}
All \(g_x\)'s in \eqref{eq:f15v39:map} depend
only on the unmodified tree links. Given
the complement of \(\mathcal A\), the map is
independent left multiplication on each of
its eight \(\SU(2)\) factors. It preserves
their Haar product. Its inverse uses the
same tree transports and replaces each
multiplier by its inverse. Fubini proves
global product-Haar preservation. No density
Jacobian or root gauge fixing is omitted.

Before the update, every affected framed
link is within \(\epsilon\) of its template.
Afterwards it is within \(\epsilon\) of
the correspondingly updated template,
because multiplication is an isometry.
For unit quaternions, telescoping a product
of four factors bounds the product error
by the sum of the four link errors. Inversion
also preserves distance. Thus each affected
plaquette trace divided by two changes
by at most \(4\epsilon\) from its relevant
template value, both before and after.
The error in the action difference is
at most
\begin{equation}
             8\epsilon\sum_{p\in\mathcal P_{\mathcal A}}J_p
                  \le288(1+\tau)\beta\epsilon .
 \label{eq:f15v39:debit}
\end{equation}
Plaquettes outside \(\mathcal P_{\mathcal A}\)
are unchanged. Combining the debit with
\eqref{eq:f15v39:rational-drop}, the guaranteed
decrease divided by \(\beta\) is at least
\[
 \frac{48-368\tau}{25}-288(1+\tau)\epsilon
       \ge\frac{73}{50}-\frac{297}{1024}
        =\frac{59902}{51200}>1.
\]
This proves \eqref{eq:f15v39:tube-drop}.

For the frustration bound, each oriented
staple uses three framed links in \(\mathcal R\),
so its difference from the template staple
is at most \(3\epsilon\). Consequently
\[
 |H_e(k(U))-H_e(\overline k)|
           \le18(1+\tau)\beta\epsilon .
\]
The scalar \(B_e\) is unchanged. At the
template, an updated link has
\(\mathcal F_e\ge4(1-\tau)\beta\) by
\eqref{eq:f15v39:template-star}.
The reverse triangle inequality for
\(h_e=|H_e|\) therefore yields
\[
 \mathcal F_e(U)\ge
        [4(1-\tau)-18(1+\tau)\epsilon]\beta>3\beta.
\]
Frustration is gauge invariant, so the
framed calculation is the stated original
configuration-space bound.
\end{proof}

The event is gauge invariant. The particular
map uses a fixed colour axis in its chosen
root frame; gauge equivariance of that map
is neither asserted nor needed for the
Haar-preserving probability argument.
The physical source array and the reference
kernel are not changed by this auxiliary map.

\subsection{A normalized joint tail under arbitrary allowed source backgrounds}

Consider \(k\) fixed translated patches of the
preceding type with disjoint vertex sets \(W_i\).
This is a simple sufficient separation condition:
the updated sets and affected plaquette sets
are disjoint, and each patch's tree avoids
every updated link from every patch.
For example, suitable patches on an \(8\)-spaced
grid fit a torus whose side is divisible by eight.
No independence of their Wilson fields is assumed.

\begin{theorem}[A collective high-frustration sector is exponentially suppressed]
\label{f15v39:joint-tail}
For every such finite family and every
positive inhomogeneous Wilson law satisfying
\((1-\tau)\beta\le J_p\le(1+\tau)\beta\),
\(0\le\tau\le1/32\), one has
\begin{equation}
            \mu_J\left(\bigcap_{i=1}^k
                        \mathcal C_i(\epsilon_i)\right)
                               \le e^{-\beta k},
           \qquad 0<\epsilon_i\le1/1024 .
 \label{eq:f15v39:joint-tail}
\end{equation}
The bound is uniform in the volume and in
the number and locations of all physical
source cells. It includes arbitrary signs
of those sources within the amplitude bound.
\end{theorem}

\begin{proof}
Apply all the corresponding maps simultaneously.
Each tree remains unchanged by all updates.
Conditioning on the complement of their
disjoint updated links proves that the
combined map \(\Phi_*\) is a Haar-preserving
bijection. On the event
\(E=\bigcap_i\mathcal C_i(\epsilon_i)\),
its action decrease is at least \(k\beta\):
the local estimates add because no affected
plaquette belongs to two patches.
Writing \(dH\) for the full link Haar product,
\[
 \begin{split}
 \mu_J(E)
 &=\mathcal Z_J^{-1}\int_E e^{-S_J(U)}\,dH(U)\\
 &\le e^{-k\beta}\mathcal Z_J^{-1}
                \int_E e^{-S_J(\Phi_*U)}\,dH(U)\\
 &=e^{-k\beta}\mu_J(\Phi_*E)\le e^{-k\beta}.
 \end{split}
\]
The exact same \(\mathcal Z_J\) appears on
both sides. The maps are continuous with
continuous inverses, so their images of
the Borel tube events are measurable.
The source-law identification is
\eqref{eq:f15v38:lift}--\eqref{eq:f15v38:law}.
There is no reference-law moment, global
likelihood cost or reflection-positivity
assumption in this proof.
\end{proof}

For a fixed separated grid of \(n\) such
patches, its tube count \(R\) consequently
obeys the immediate subset-expansion bound
\[
          \mathbb E_{\mu_J}e^{vR}
                    \le[1+e^{-\beta}(e^v-1)]^n,\qquad v\ge0.
\]
This count consequence is the already used
elementary expansion, not a new independence
or stochastic-domination assertion.

\subsection{Scope and the next collective question}

The central template is a simultaneous mode
of every original one-link conditional, yet
has an explicit local collective escape.
Thus iterating V38's one-link innovation
estimate while treating its retained
frustration as another innovation would
miss an actual Wilson mechanism.
The eight-link calculation supplies both
an action decrease and the normalized
probability estimate for its specified
noncommuting neighbourhood.

It does not prove
\eqref{eq:f15v38:missing-frustration}.
The inclusion proved here is
\(\mathcal C_i\subset\{\hbox{large local frustration}\}\),
not the reverse inclusion. Large frustration
need not be close to this central template,
and V38's alarms occur already at much
smaller, regulator-dependent thresholds.
A controlled cover of those general alarms
by collective escape domains, with the
cost of the labels and supports included,
or a different actual-background estimate,
is still required.

The exact retained penalty
\eqref{eq:f15v39:minimum} helps specify that
problem but cannot replace the missing
probability argument. The auxiliary changes
of variables are also not claimed to
preserve completed coarse coordinates.
None of the results is a mixing theorem,
an iteration of the physical RG map,
a continuum construction or a Hamiltonian
mass-gap estimate.

\subsection{Exact diagnostics and preservation}

The new diagnostic is
\begin{center}\small
\path{scripts/verify_ym_f15_v39_collective_frustration_escape.py}.
\end{center}
Outward rational positive-series enclosures
for the exact spherical integral verify
24 chord bounds and 24 Lipschitz penalty
bounds. Periodic tori of sides \(4,6\)
check 9,312 central-flux plaquettes,
6,208 simultaneous zero innovations and
two summed-star identities.

The local geometry has eight updated links,
36 affected plaquettes, 104 relevant links,
127 tree links and 207 support links in total.
The script checks the tree transports,
all 36 affected plaquette phases, eight
rational template-action identities and
eight uniform drop/frustration margins.

Eight noncommuting rational tube configurations,
including nontrivial tree frames, check 832
tube-edge inequalities, 1,656 framed update
identities, 1,656 inverse identities and 832
gauge-invariant tube distances. Their action
changes are bounded for every coupling in the
declared interval by choosing the worst endpoint
separately for each signed plaquette debit.
There are eight such action-drop checks,
64 large-frustration star checks and five
separated-support checks.

The joint tail itself follows from the
analytic Haar change of variables above,
not finite sampling. The diagnostics do
not establish a covering of the remaining
frustration sectors or a continuum mass gap.

The script has 11,330 bytes and SHA--256
\begin{center}\scriptsize\ttfamily
64BB34D529D0FF2D1AA98296681F4187300A6D7D7C968B927F97031CAC7FA13D.
\end{center}
The V38 predecessor has 889,967 bytes and SHA--256
\begin{center}\scriptsize\ttfamily
64C622FF0FE852626692A34AD171C43B8ACC5E6302C09CA22069FAA96834ABB1.
\end{center}
Removing this terminal section and reversing
the title version recovers V38 exactly.
V39 replaces V38 as the sole active main note.
All 14 protected sources and 81 earlier
certificates are preserved. The next
companion checkpoint is V40.

\begin{firewall}
V39 proves an exact retained-star penalty
and a normalized, source-uniform joint tail
for a specified high-frustration sector
using an explicit collective link update.
It also exhibits the simultaneous one-link
innovation obstruction that motivates that
update. It does not cover all frustration
alarms, prove V37's full criterion, close
iterated RG, construct the interacting
physical continuum or establish the full
Yang--Mills mass gap. The goal remains active.
\end{firewall}

\section{V40: all near-central negative-plaquette sectors and the noncentral remainder}
\label{f15v40:section}

V39 controls a neighbourhood of one center-flux
template. The present step removes that template
restriction. In the full four-dimensional
\(\SU(2)\) Wilson law, every locally near-central
plaquette configuration with a negative central
face has a bounded-support action decrease.
Two explicit real cochains suffice for every
possible partner face forced by the central
cube identity. The resulting joint probability
bound is uniform under the allowed inhomogeneous
physical sources.

This is not an all-frustration theorem.
It eliminates a specified class of nearly
projectively flat, nonflat configurations.
The remaining event is made explicit:
a genuinely noncentral plaquette excursion in
a bounded enlargement of the original carrier.
Its actual-source probability, at the required
shrinking threshold, is still unproved.

\subsection{V40 companion checkpoint and nonduplication}

Fresh inventories of the note folder and supplied
downloads found the same latest relevant sources:
SM continuum V72, NS stability V26, shell mixing
V16, parallel YM V5, source selection V31,
native stationarity V34 and exact-action Einstein
bridge V24. Their current conclusions were reread,
as were the supplied suggestions. All 96 protected
source and certificate files were unchanged.

The SM result remains a fixed-mode, fixed-time
many-copy limit. NS supplies pointwise rank-one
norms, with the subsequent integral certificate
not supplied. Shell V16 retains its complete
measured flat one-loop coefficient but not the
noncommuting coefficient or a common
nonperturbative integral remainder. Parallel YM
still rules out its extensive global-charge
tightness target. Source selection is conditional
on historical word-native typing; stationarity's
moving-fibre derivative entries are not computed;
the Einstein bridge's nonzero first-jet derivative
is not an integrated root. None supplies the
missing actual-source Wilson tail.

The direction therefore remains upstream in the
original compact probability law. The compact
conditional-minimum concentration argument is
already in f7 V70 and f12 V81. The central
signed-curl Hessian and bounded repairs of special
completed-fibre patterns also occur in the archives.
They are not new claims here. The new input is a
pair of explicit full-Wilson cochains covering
every negative central cube pattern, followed by
a gauge-invariant plaquette collar and a
source-uniform joint tail without enumerating
all central link assignments. No completed
coarse coordinate is required to be preserved.

At the user's request, \path{latex_notes} now
contains only \texttt{.tex} sources. Existing
build outputs were moved without deletion to
\path{artifacts}; future compilation outputs
are directed there. All 44 source files present
during that cleanup were hash-checked unchanged.
The next compulsory companion checkpoint is V45.

\subsection{Two finite cochains cover every central negative face}

Use integer coordinates on \(\mathbb Z^4\).
An oriented positive plaquette is \(p=(x,\mu,\nu)\),
\(\mu<\nu\). For a finitely supported real edge
cochain \(a\), write
\begin{equation}
 (da)_{x,\mu,\nu}
 =a_{x,\mu}+a_{x+e_\mu,\nu}
             -a_{x+e_\nu,\mu}-a_{x,\nu}.
 \label{eq:f15v40:curl}
\end{equation}
Let \(d^*\) be its counting-inner-product adjoint,
and let \(e_p\) denote the unit face cochain.
All operators below act on finite-support
cochains, so no Green function or limiting
inverse is used.

For the prototype face \(p=(0,1,2)\), the two
partner prototypes are
\[
       q_{\rm a}=(0,1,3),\qquad
       q_{\rm o}=(e_3,1,2).
\]
They are respectively adjacent and opposite
faces of the unit cube in directions \(1,2,3\).
Define the componentwise translation operator
\[
 (Ta)_{x,\mu}
       =8a_{x,\mu}+\sum_{\nu=1}^4
                         (a_{x+e_\nu,\mu}+a_{x-e_\nu,\mu}),
\]
and take
\begin{equation}
 a_{\rm a}=d^*(e_p+e_{q_{\rm a}}),\qquad
 a_{\rm o}=(I+T/16)d^*(e_p+e_{q_{\rm o}}).
 \label{eq:f15v40:directions}
\end{equation}
The second expression is just a polynomial
with rational coefficients.

For either direction put \(f=da\),
\(\mathcal A=\operatorname{supp}a\), let
\(\mathcal P_{\mathcal A}\) be every plaquette
incident to \(\mathcal A\), and let
\(\mathcal R_{\mathcal A}\) contain all their
edges. Direct finite incidence gives
\begin{equation}
\begin{array}{c|rrr|ccc}
 &|\mathcal A|&|\mathcal P_{\mathcal A}|&|\mathcal R_{\mathcal A}|
 &H_2=\sum f^2&C_2=f_p^2+f_q^2&H_4=\sum f^4\\ \hline
 {\rm adjacent}&7&34&70&92&50&1340\\
 {\rm opposite}&56&232&356&273/2&72&2783487/1024 .
\end{array}
 \label{eq:f15v40:table}
\end{equation}
Zero-curl plaquettes are included in
\(\mathcal P_{\mathcal A}\): their number is
two in the first row and 24 in the second.
They cannot be omitted from a perturbed-field
estimate.

Here are explicit coefficient certificates for
the table. For the adjacent prototype, the
nonzero absolute curl values and multiplicities
are
\[
                  (1,26),\quad(2,4),\quad(5,2),
\]
and \(f_p=f_q=5\). For the opposite prototype,
the corresponding list for the integer cochain
\(16f\) is
\[
 (1,160),\ (2,8),\ (4,6),\
 (21,8),\ (23,24),\ (96,2),
\]
with \(16f_p=16f_q=96\).
These lists follow by substituting
\eqref{eq:f15v40:directions} into
\eqref{eq:f15v40:curl}. For example, the
integer edge coefficients of \(16a_{\rm o}\)
have absolute values \(1\) on 48 edges and
\(24\) on eight edges. Squaring and fourth
powering the displayed curl lists gives
\eqref{eq:f15v40:table}.
The affected sets are obtained by including
the six incident plaquettes of each support
edge and removing repetitions.

Cubical coordinate permutations and reflections
transport oriented edge and face cochains,
commute with \(d\), and preserve these norms.
The adjacent prototype covers the twelve
unordered adjacent face pairs of a cube;
the opposite prototype covers its three
opposite pairs. Orientation reversals can
change the signs of \(f_p,f_q\), which will
only enter through squares or cosines.
Every relevant edge, for every such transformed
pair in the unit cube, has both endpoints in
\begin{equation}
                         W=\{-2,-1,0,1,2,3\}^4.
 \label{eq:f15v40:box}
\end{equation}

\begin{proposition}[Every central cube defect has a uniform local decrease]
\label{f15v40:central-drop}
Suppose all link variables with endpoints in
\(W\) are central, \(z_eI\), \(z_e\in\{-1,1\}\).
Suppose a face \(p\) of a unit three-cube is
negative. For positive couplings
\[
          (1-\tau)\beta\le J_r\le(1+\tau)\beta,
          \qquad 0\le\tau\le1/32,
\]
one of the five other cube faces \(q\) is
negative. The appropriate transported cochain
in \eqref{eq:f15v40:directions}, updated by
\[
                  U_e\longmapsto
                  \exp(i t a_e\sigma_3)U_e,\qquad t=1/16,
\]
decreases the Wilson action by more than
\(9\beta/2048\). The support has at most
56 links. No condition is imposed on the
other central plaquette signs.
\end{proposition}

\begin{proof}
The product of the six central plaquette signs
of a cube is \(+1\), since each edge sign occurs
twice. A negative \(p\) thus forces a negative
partner \(q\). The pair has one of the two
geometries already covered.

At a central assignment the updated plaquette
trace divided by two is \(z_r\cos(t f_r)\).
Therefore its exact action difference is
\[
                   \Delta S=\sum_r J_r z_r(1-\cos(t f_r)).
\]
The quadratic coefficient satisfies
\begin{equation}
 \frac1\beta\sum_rJ_rz_rf_r^2
 \le(1+\tau)(H_2-C_2)-(1-\tau)C_2
       =(1+\tau)H_2-2C_2.
 \label{eq:f15v40:quadratic}
\end{equation}
This uses the worst allowed coupling separately
on the designated negative pair and on its
complement. Additional negative signs can only
improve this upper bound. At \(\tau=1/32\)
the two right sides are \(-41/8\) and
\(-207/64\).

Taylor's theorem, with the global remainder
\(\lvert1-\cos x-x^2/2\rvert\le x^4/24\),
gives the lower bounds for \(-\Delta S/\beta\)
\begin{equation}
 \frac{t^2}{2}\bigl[2C_2-(33/32)H_2\bigr]
                  -\frac{33}{32}\frac{t^4H_4}{24}
 =
 \begin{cases}
 38299/4194304,&{\rm adjacent},\\
 77909259/17179869184,&{\rm opposite}.
 \end{cases}
 \label{eq:f15v40:finite-drop}
\end{equation}
Both fractions exceed \(9/2048\).
Only incident plaquettes change. This proves
the assertion, uniformly over every other
central assignment and all permitted couplings.
\end{proof}

In particular, an exactly central full-Wilson
configuration with a negative face cannot be
a local minimum: the strictly negative
quadratic coefficient in
\eqref{eq:f15v40:quadratic} already supplies
an arbitrarily small descending curve.
This is a statement about the full link space,
not an arbitrary constrained coarse fibre.

\subsection{From a plaquette collar to a central link comparison}

Embed a translated and coordinate-oriented
copy of \(W\) in a periodic four-torus of
side at least eight. Let \(\mathcal P(W)\)
be all plaquettes whose vertices are in \(W\).
Choose the designated face \(p\) at the origin
as above, with its accompanying unit three-cube.
For
\[
                        0<\eta\le2^{-26}
\]
define the gauge-invariant event
\begin{equation}
 \mathcal B_p(\eta)=
 \left\{\begin{array}{l}
 \displaystyle\max_{r\in\mathcal P(W)}
              \operatorname{dist}(U_r,\{I,-I\})\le\eta,\\
                         |U_p+I|\le\eta
 \end{array}\right\}.
 \label{eq:f15v40:event}
\end{equation}
Distances are Euclidean unit-quaternion distances.
No gauge has been selected in this definition.

\begin{proposition}[A deterministic conditional excess on one of five supports]
\label{f15v40:collar}
For each of the five partner choices \(q\),
fix a transported support \(\mathcal A_q\)
from \eqref{eq:f15v40:directions}.
Define its full-Wilson conditional excess by
\begin{equation}
 X_q(U)=S_J(U)-
       \min_{V_{\mathcal A_q}}
                  S_J(V_{\mathcal A_q},U_{\mathcal A_q^c}).
 \label{eq:f15v40:excess}
\end{equation}
Then, under the preceding coupling bounds,
\begin{equation}
        \mathcal B_p(\eta)
               \subseteq\bigcup_{q}
                           \{X_q>\beta/256\}.
 \label{eq:f15v40:cover}
\end{equation}
The minimum is over the original compact
link factors, without changing the exterior,
the action or the physical source array.
It is continuous in the exterior, since its
change is bounded by the uniform change of
the continuous action on the compact
minimization space. Thus \(X_q\) is measurable.
\end{proposition}

\begin{proof}
Root the box at \((-2,-2,-2,-2)\), and use
ordered coordinate paths to its vertices.
Let \(g_x\) be transport along the root path
and \(k_e=g_xU_eg_y^{-1}\) for \(e=(x,y)\).
The usual ordered-strip product identity
expresses \(k_e\) as a product of conjugates
of elementary plaquettes: commute the final
edge of direction \(\mu\) past the later
coordinate segments of the root path.
There are exactly
\(\sum_{\nu>\mu}(x_\nu+2)\le15\) such faces,
all lying in \(W\). This also proves the
identity by successive elementary square
interchanges; an empty strip gives \(k_e=I\).

On \(\mathcal B_p(\eta)\), each factor is
within \(\eta\) of a central sign. Multiplication
by a unit quaternion is an isometry, and
conjugation fixes the center. Product
telescoping therefore gives signs \(z_e\)
such that
\begin{equation}
                 |k_e-z_eI|\le15\eta\le\epsilon,
                       \qquad \epsilon=2^{-22}.
 \label{eq:f15v40:link-collar}
\end{equation}
For any face, its framed product differs from
the corresponding product of the \(z_e\)'s
by at most \(4\epsilon\). The sign at the
designated face must be negative, since
\(4\epsilon+\eta<2=|I-(-I)|\).
The six projected central face signs have
product \(+1\), so a partner \(q\) is
negative as well.

Fix these frames for this one comparison.
On \(\mathcal A_q\) define
\[
 U'_e=g_x^{-1}\exp(i t a_e\sigma_3)g_xU_e,
                      \qquad t=1/16,
\]
and leave every other original link unchanged.
In the fixed frames the new links are within
\(\epsilon\) of the updated central template,
just as the old links were within \(\epsilon\)
of the old template. Before and after the
update each affected plaquette's normalized
trace error is at most \(4\epsilon\).
Thus the total error in the action difference
is at most
\begin{equation}
 8\epsilon\sum_{r\in\mathcal P_{\mathcal A_q}}J_r
       \le8\epsilon(33/32)\beta\,232
       =\frac{1914}{2^{22}}\beta
       <\frac{\beta}{2048}.
 \label{eq:f15v40:collar-debit}
\end{equation}
Together with \eqref{eq:f15v40:finite-drop},
this gives
\(S_J(U')<S_J(U)-\beta/256\).
The minimization in \eqref{eq:f15v40:excess}
allows this \(U'\), which proves the inclusion.
\end{proof}

The frames in this proof are held fixed during
the comparison. They need not equal the root
transports recomputed after the update.
Consequently the configuration-dependent
comparison is not asserted to preserve Haar
measure. The next argument instead uses the
exact conditional law and its minimum.
There is no measurable gauge-selection or
template-counting premise hidden in
\eqref{eq:f15v40:cover}.

\subsection{The exact conditional normalizer pays for the entire central union}

For a set \(A\) of \(m\) links, condition on
its complement, and write \(S_A\) for the sum
of all plaquette actions incident to \(A\).
Its minimum exists by compactness. The
conditional density is
\[
 \frac{e^{-(S_A-\min S_A)}}{
        \int e^{-(S_A-\min S_A)}\,dH_A}\,dH_A.
\]
Here and below, \(\beta\) is already included
in \(J_r\). Applying the inherited compact
conditional-minimum argument with explicit
Wilson constants gives, for \(\beta\ge1\),
\begin{equation}
 \mathbb P_J\{S_A-\min S_A\ge s\mid U_{A^c}\}
                 \le(2^{13}\beta^{3/2})^m e^{-s}.
 \label{eq:f15v40:conditional-tail}
\end{equation}
This holds for every exterior, without a
Hessian lower bound or uniqueness of a minimum.

For completeness, choose a minimizing
configuration \(V_A\) and product geodesics
\(V_e\exp(tX_e)\). A plaquette has four
factors; its normalized trace has second
derivative in absolute value at most
\((\sum_{e\in r\cap A}|X_e|)^2\).
Cauchy--Schwarz and the six incident
plaquettes per edge give
\[
 \left|\frac{d^2}{dt^2}S_A\right|
               \le24(1+\tau)\beta\sum_{e\in A}|X_e|^2.
\]
The first derivative is zero at the minimum.
On the product of geodesic balls
\(|X_e|\le r=(32\beta)^{-1/2}\), Taylor's
theorem therefore gives
\[
                    S_A-\min S_A\le99m/256.
\]
For \(0<r\le1\), normalized Haar volume on
\(S^3\) obeys
\[
 H(B_r)=\frac2\pi\int_0^r\sin^2u\,du
             \ge\frac{r^3}{6\pi}>\frac{r^3}{24}.
\]
We used \(\sin u\ge u/2\) on \([0,1]\).
It follows that
\[
 \int e^{-(S_A-\min S_A)}\,dH_A
 \ge\left[\frac{e^{-99/256}}{24(32\beta)^{3/2}}\right]^m
 \ge(2^{13}\beta^{3/2})^{-m}.
\]
For the last estimate, \(32^{3/2}<192\) and
\(e^{99/256}<e^{1/2}<5/3\), whence
\(24\cdot192\cdot5/3<2^{13}\).
The numerator of the excess-\(s\) event is
at most \(e^{-s}\). Division proves
\eqref{eq:f15v40:conditional-tail}.

\begin{theorem}[Joint suppression of every near-central negative-face sector]
\label{f15v40:joint}
Let \(W_1,\ldots,W_k\) be fixed translated
and oriented boxes of the preceding type,
with pairwise disjoint vertex sets, and
designated faces \(p_i\). For
\(\beta\ge1\), \(0\le\tau\le1/32\) and
\(0<\eta_i\le2^{-26}\), put
\begin{equation}
 q_{\rm c}(\beta)=
 \min\left\{1,\,
     5(2^{13}\beta^{3/2})^{56}e^{-\beta/256}\right\}.
 \label{eq:f15v40:q}
\end{equation}
Then the actual inhomogeneous Wilson law satisfies
\begin{equation}
             \mu_J\left(\bigcap_{i=1}^k
                                  \mathcal B_{p_i}(\eta_i)\right)
                             \le q_{\rm c}(\beta)^k.
 \label{eq:f15v40:joint}
\end{equation}
In particular, uniformly over all permitted
physical source arrays and volumes,
\begin{equation}
 \beta\ge2^{22}\quad\Longrightarrow\quad
 \mu_J\left(\bigcap_{i=1}^k\mathcal B_{p_i}(\eta_i)\right)
                             \le e^{-\beta k/512}.
 \label{eq:f15v40:simple-tail}
\end{equation}
\end{theorem}

\begin{proof}
For each box use \eqref{eq:f15v40:cover}.
There are five support choices per box.
Fix one choice in each box and condition
simultaneously on the complement of their
union. Every plaquette incident to a chosen
support lies entirely in its own \(W_i\).
No plaquette contains chosen variables from
two boxes, so the conditional law is exactly
the product of the \(k\) conditional laws.
Moreover each conditional minimum appearing
in \(X_{q_i}\) is the same minimum as in
this simultaneous conditioning: its action
does not depend on another chosen support.

Apply \eqref{eq:f15v40:conditional-tail}
with \(s=\beta/256\) and \(m\le56\).
The resulting conditional product bound is
uniform in the common exterior. Integrating
it and summing over at most \(5^k\) support
choices proves \eqref{eq:f15v40:joint}.
If the displayed scalar bound exceeds one,
use the trivial probability bound instead.
All central signs and all gauge frames were
already included in the deterministic excess
cover; no factor counts their assignments.

For the final assertion, it suffices to check
\[
             \log5+728\log2+84\log\beta\le\beta/512 .
\]
At \(\beta=2^{22}\) its left side is less
than \(2+728+84\cdot22=2578\), whereas
the right side is \(8192\). The difference
is increasing thereafter, since
\(1/512-84/\beta>0\). This proves
\eqref{eq:f15v40:simple-tail}.
\end{proof}

The probability proof uses no reflection
positivity and no change from the physical
source background to a homogeneous reference
law. The same conditional normalizer is
retained throughout. Auxiliary compact
minimization is used only to estimate an
event of that unchanged law.

\subsection{The remaining frustration event is genuinely noncentral}

For an edge \(e\), let \(\mathcal S_e\) be
its six incident positive plaquettes. For
each \(p\in\mathcal S_e\), choose the box
\(W(p)\) used above, and put
\[
 \mathcal D_e=\bigcup_{p\in\mathcal S_e}W(p),\qquad
 \mathcal C_e(\eta)=\bigcup_{p\in\mathcal S_e}\mathcal B_p(\eta).
\]
Define the noncentral excursion event
\begin{equation}
 \mathcal N_e(\eta)=
 \left\{\max_{r\in\bigcup_{p\in\mathcal S_e}\mathcal P(W(p))}
                    \operatorname{dist}(U_r,\{I,-I\})>\eta\right\}.
 \label{eq:f15v40:noncentral}
\end{equation}
Both are observables of the original gauge
field, not abstract exterior labels.

\begin{proposition}[Frustration splits into controlled central defects and a noncentral excursion]
\label{f15v40:split}
For \(0<\eta\le2^{-26}\),
\begin{equation}
 \{\mathcal F_e>3(1+\tau)\beta\eta^2\}
             \subseteq\mathcal C_e(\eta)\cup\mathcal N_e(\eta).
 \label{eq:f15v40:split}
\end{equation}
For edges with pairwise disjoint enlarged
vertex sets \(\mathcal D_{e_i}\),
\begin{equation}
 \mu_J\left(\bigcap_{i=1}^k\mathcal C_{e_i}(\eta_i)\right)
                     \le \min\{1,6q_{\rm c}(\beta)\}^{\,k}.
 \label{eq:f15v40:star-tail}
\end{equation}
\end{proposition}

\begin{proof}
Off \(\mathcal N_e(\eta)\), every relevant
plaquette is within \(\eta\) of a unique
central sign. If an incident plaquette is
near \(-I\), its entire box is near the
center and \(\mathcal C_e(\eta)\) holds.
Otherwise all six incident plaquettes
satisfy \(|U_p-I|\le\eta\).
For a unit quaternion \(u\),
\(1-\frac12\Tr u=\frac12|u-I|^2\). V38's
star decomposition then gives
\[
 \mathcal F_e\le E_e
      =\sum_{p\ni e}J_p(1-\tfrac12\Tr U_p)
      \le3(1+\tau)\beta\eta^2.
\]
This proves the first inclusion.
For the second, expand the six face choices
per edge. Any chosen boxes from distinct
\(\mathcal D_{e_i}\)'s are disjoint, so
Theorem~\ref{f15v40:joint} applies.
Sum the resulting \(6^k\) bounds and
truncate by one.
\end{proof}

In particular, for a desired dimensionless
frustration threshold \(r>0\), choose
\[
             \eta=\min\left\{2^{-26},
                         \sqrt{\frac{r}{6(1+\tau)}}\right\}.
\]
Then the central term in the event
\(\{\mathcal F_e>r\beta\}\) is controlled
by the preceding theorem, and the unproved
term is explicitly \(\mathcal N_e(\eta)\).
Its threshold may shrink with the regulator.
At that scale this is not a probability
consequence of the fixed central estimate.

V38's strip alarm is a sum of frustrations
over a carrier, not a single edge event.
Its use of this split must still pay the
choice of a witnessing edge and the overlap
of the enlarged boxes. More importantly,
no required joint bound for the
\(\mathcal N_e(\eta)\)'s has been proved.
Thus neither \eqref{eq:f15v38:missing-frustration}
nor V37's full actual-background criterion
is claimed closed.

The new upstream target is an actual-source
estimate for those noncentral plaquette
excursions, with the shrinking threshold
and carrier multiplicities retained.
Iterating the single-link innovation bound
or counting further central templates would
not by itself settle that target.

\subsection{Exact diagnostics and preservation}

The new standalone diagnostic is
\begin{center}\small
\path{scripts/verify_ym_f15_v40_all_central_escape.py}.
\end{center}
It reconstructs both cochains from incidence,
their complete curl histograms, the support
sets and all moments in
\eqref{eq:f15v40:table}. It checks 24 outward
rational finite-angle/drop enclosures, including
the worst allowed plaquette couplings and
the neighbourhood debit.

There are 96 oriented cubical-symmetry
identities, covering all fifteen cube
face pairs. Exhausting all 4,096 central
assignments on the twelve cube edges checks
12,288 forced negative partners and all
32 possible central face-sign patterns.
There are 4,320 ordered-strip area bounds,
four exact normalizer/large-\(\beta\)
constant checks, 128 noncommuting link
isometries and 32 noncommuting plaquette
debit checks.

The probability theorem follows from the
analytic compact-minimum and conditional
factorization proof above. It is not inferred
from finite sampling, and these diagnostics
do not certify a continuum mass gap.

The script has 11,121 bytes and SHA--256
\begin{center}\scriptsize\ttfamily
4CC61E1E2EB6952D3EB6B584683881939E110636BAF2C6552B659F113618AEAC.
\end{center}
The V39 predecessor has 912,172 bytes and SHA--256
\begin{center}\scriptsize\ttfamily
9DD37E30F788C04A4AB1CC91D988EAD42EDB0837FBCD92E9BE4D6FD1C5A4896B.
\end{center}
Removing this terminal section and reversing
the title version recovers V39 exactly.
V40 replaces V39 as the sole active main
source. All 14 protected sources and 82
earlier certificates are preserved.
The next companion checkpoint is V45.

\begin{firewall}
V40 proves source-uniform joint suppression
of all specified near-central negative-face
sectors, rather than only V39's center-flux
template. It identifies the remaining
noncentral excursion event needed for
general frustration control. It does not
prove its required joint tail, full
actual-background normalization control,
iterated RG contraction, an interacting
continuum construction or the physical
Yang--Mills mass gap. The goal remains active.
\end{firewall}

\section{V41: boundary curvature, arbitrary interiors, and normalized block escape}
\label{f15v41:section}

V40 reduces general frustration to a specified
noncentral excursion problem. We now control
arbitrary interior configurations when their
box boundary has sufficiently small curvature.
The boundary condition will be expressed
entirely in original plaquette observables.
It does not require a preselected gauge or
any smallness assumption on the sampled
interior links.

There are three new quantitative steps.
First, an explicit family of boundary paths
gives a gauge estimate using only boundary
plaquettes. Second, a comparison configuration
pays the entire conditional minimum by those
boundary data. Third, exact conditioning
gives a joint block tail and a weighted
boundary-failure interface with no
Cauchy--Schwarz loss in its exponent.

The remaining assumption is important:
we do not prove that the small-curvature
boundaries are sufficiently probable.
The result is an actual-source conditional
estimate, not a full noncentral tail or a
continuum construction.

\subsection{Scope checks and the inherited input}

The V40 source and all 97 protected files
were checked unchanged. No companion review
is due before V45. Compilation remains in
\path{artifacts}, not in \path{latex_notes}.

The finite compact conditional-excess
estimate is inherited from f7 V70, f12 V81
and its explicit Wilson normalization in
\eqref{eq:f15v40:conditional-tail}. It is
not claimed anew. Nor is the ordinary
ordered-path gauge identity new. Here the
paths and all filling squares must remain
on the boundary of a four-dimensional box;
using interior plaquettes at that step
would assume the event we are trying to
control. The construction below proves
that boundary-only statement and its
explicit constant.

The archived homogeneous chessboard result
is not applied to an arbitrary source
background. We also do not assume
monotonicity in individual non-Abelian
plaquette couplings. As a scope check,
Abdesselam's
\href{https://math-events.uni-bonn.de/event/152/contributions/568/}
{2025 Bonn workshop abstract on non-Abelian correlation inequalities}
describes a proposed route to the corresponding
GKS inequality as conjectural. That abstract
is contextual, not an input to any proof below.

\subsection{The exact conditional box and its boundary incidence}

Let
\[
                       B_n=\{0,\ldots,n\}^4,\qquad n\ge2,
\]
be embedded without identification in the
ambient periodic lattice. Its boundary is
the cubical subcomplex consisting of the
eight facets \(x_\mu=0\) or \(n\).
An edge or plaquette is a boundary cell
when it lies entirely in a facet.

Let \(A_n\) be the positive edges of \(B_n\)
not contained in its boundary. For an edge
of direction \(\mu\), this means that all
three transverse coordinates lie in
\(\{1,\ldots,n-1\}\). Consequently
\begin{equation}
                         m_n:=|A_n|=4n(n-1)^3.
 \label{eq:f15v41:edges}
\end{equation}
Let \(\mathcal P_n\) be all plaquettes
incident to \(A_n\), and
\(\mathcal P_n^\partial\) the boundary
plaquettes. Every plaquette incident to
\(A_n\) lies inside \(B_n\), and
\begin{equation}
 |\mathcal P_n|=6n^2(n-1)^2,\qquad
 |\mathcal P_n^\partial|=24n^3.
 \label{eq:f15v41:faces}
\end{equation}
The first set is exactly the set of
nonboundary plaquettes in \(B_n\).
There are \(24n^3+8n\) boundary edges and
\(8n^3+8n\) boundary vertices.

Classify a plaquette of \(\mathcal P_n\)
by the number \(j\) of its edges outside
\(A_n\). The exact counts are
\begin{equation}
 \begin{split}
 N_0&=6(n-2)^2(n-1)^2,\\
 N_1&=24(n-2)(n-1)^2,\\
 N_2&=24(n-1)^2,
 \end{split}
 \qquad N_3=N_4=0.
 \label{eq:f15v41:incidence}
\end{equation}
Indeed, fix the two face directions.
The other two coordinates have
\((n-1)^2\) choices. In each face
direction its base coordinate has \(n\)
choices, of which two are adjacent to
the corresponding boundary and \(n-2\)
are internal. The three cases in
\eqref{eq:f15v41:incidence} count zero,
one or two boundary-adjacent choices.

The comparison coefficient that will
matter is therefore
\begin{equation}
 D_n:=\frac{N_1}{2}+2N_2
       =12(n+2)(n-1)^2
       =12(n^3-3n+2)\le12n^3.
 \label{eq:f15v41:D}
\end{equation}

\subsection{A boundary-only gauge with an explicit area bound}

Write a boundary vertex as \(x=(y,h)\),
where \(y\in B_n^{(3)}:=\{0,\ldots,n\}^3\) and
\(h=x_4\). Let \(P^{\rm ord}_y\) be the
monotone path in coordinate order \(1,2,3\)
from the three-dimensional origin to \(y\).
For \(y\) on the boundary of the three-cube,
choose a monotone boundary path \(P^\partial_y\)
as follows:

\begin{itemize}
\item If some coordinate of \(y\) is zero,
use \(P^{\rm ord}_y\); that zero coordinate
keeps the whole path on the boundary.
\item Otherwise choose the least coordinate
with value \(n\), traverse that coordinate
first, and then traverse the other two in
coordinate order. The first segment has
other coordinates zero, and the remainder
has the selected coordinate fixed at \(n\).
\end{itemize}

These paths have length at most \(3n\).
Root the four-box at \(o=(0,0,0,0)\).
Choose its boundary root paths \(P_x\) by
\begin{equation}
 P_{(y,h)}=
 \begin{cases}
  (\text{vertical at }0)^n\,P_y^{\rm ord}
                              \text{ at height }n,&h=n,\\
  P_y^\partial\text{ at height }0\,
                         (\text{vertical at }y)^h,
                          &h<n,\ y\in\partial B_n^{(3)},\\
  P_y^{\rm ord}\text{ at height }0,
                          &h=0,\ y\notin\partial B_n^{(3)} .
 \end{cases}
 \label{eq:f15v41:paths}
\end{equation}
These cases include every boundary vertex.
Every path uses only boundary edges and
has length at most \(4n\). They need not
be branches of a single tree.

\begin{proposition}[Quantitative boundary gauge]
\label{f15v41:boundary-gauge}
Suppose
\begin{equation}
                  \max_{p\in\mathcal P_n^\partial}|U_p-I|
                                      \le\delta.
 \label{eq:f15v41:boundary-event}
\end{equation}
Let \(g_x\) be transport along \(P_x\).
Then for every boundary edge \(e=(x,y)\),
\begin{equation}
                    |g_xU_eg_y^{-1}-I|\le9n^2\delta.
 \label{eq:f15v41:boundary-gauge}
\end{equation}
All transports and the hypothesis involve
only the exterior of \(A_n\).
\end{proposition}

\begin{proof}
We give boundary square homotopies from
\(P_xe\) to \(P_y\). An interchange of
two adjacent, distinct positive coordinate
steps uses one elementary square. Reordering
a three-coordinate block path into ordinary
coordinate order uses at most \(3n^2\)
such squares: each reversed pair of
coordinate blocks contributes at most
\(n^2\). Appending one horizontal step
to an already ordered path requires at
most \(2n\) additional interchanges.

For a horizontal edge in the bottom
facet, sort \(P_xe\) into the ordinary
path and then reverse the sorting of
\(P_y\). This uses at most
\(6n^2+2n\) bottom-facet squares.
For a horizontal edge in the top facet,
cancel the common initial vertical path;
at most \(2n\) top-facet squares remain.

For a horizontal edge at height
\(0<h<n\), its horizontal projection
is an edge on the boundary of the
three-cube. Commute its final horizontal
step past the \(h\) vertical steps.
All these \(h\) squares lie on side
facets. Sort the two bottom paths as
in the preceding bottom case. The
total is at most
\[
                       h+6n^2+2n\le6n^2+3n.
\]
Every sorting square remains in the
bottom facet; no interior square is used.

For a vertical edge whose upper endpoint
has height less than \(n\), the two
paths already agree after the edge is
appended. For the last vertical edge,
compare \(P_y^\partial\) followed by
all \(n\) vertical steps with the
root vertical path followed by
\(P_y^{\rm ord}\) on the roof.
Extruding the entire boundary path
\(P_y^\partial\) vertically costs at
most \(3n^2\) side-facet squares.
Reordering it on the roof costs
at most another \(3n^2\).

All five bounds are at most \(9n^2\).
Each elementary interchange replaces a
path product by a product differing by
a conjugate of a boundary plaquette or
its inverse. Iterating the identity
expresses \(g_xU_eg_y^{-1}\) as at most
\(9n^2\) such conjugated factors.
Multiplication and inversion are
isometries for unit quaternions,
and conjugation fixes \(I\).
Telescoping the product proves
\eqref{eq:f15v41:boundary-gauge}.
\end{proof}

This is a bound for a non-Abelian product,
not a statement about an additive
abelian flux. It uses no small logarithm,
global gauge slice or continuum extension.

\subsection{The boundary pays the actual conditional minimum}

Keep the physical couplings
\[
             (1-\tau)\beta\le J_p\le(1+\tau)\beta,
                    \qquad 0\le\tau\le1/32.
\]
With the complement of \(A_n\) fixed, let
\[
       S_n(U)=\sum_{p\in\mathcal P_n}
                    J_p(1-\tfrac12\Tr U_p),\qquad
       s_n(\xi)=\min_{V_{A_n}}S_n(V_{A_n},\xi).
\]
This is the complete action depending
on the updated variables. Exterior-only
plaquettes cancel in the conditional law.
The compact minimum is continuous in
the exterior, by uniform continuity
of the action on the finite compact
product.

\begin{proposition}[An all-interior minimum bound]
\label{f15v41:minimum}
On the boundary event
\eqref{eq:f15v41:boundary-event},
\begin{equation}
          0\le s_n(\xi)\le
           (1+\tau)\beta D_n(9n^2\delta)^2.
 \label{eq:f15v41:minimum}
\end{equation}
The interior field is unrestricted.
\end{proposition}

\begin{proof}
Use the boundary frames from
Proposition~\ref{f15v41:boundary-gauge}
and extend them arbitrarily to the
interior vertices. In these frames
every fixed boundary link is within
\(\rho=9n^2\delta\) of \(I\).
Choose all \(A_n\)-links to equal \(I\)
in this frame, and transform back.
Only original links in \(A_n\) have
been replaced; their complement stays
exactly fixed.

An affected plaquette with no fixed
boundary edge is now \(I\).
With one fixed edge its holonomy is
within \(\rho\) of \(I\), and with
two fixed edges it is within \(2\rho\).
The latter statement is product
telescoping and does not assume
commutativity. Since
\[
                 1-\tfrac12\Tr u=\tfrac12|u-I|^2
                        \quad(u\in\SU(2)),
\]
the total comparison action is at
most
\[
 (1+\tau)\beta\left(\frac{N_1}{2}+2N_2\right)\rho^2.
\]
The minimum is no larger than this
explicit comparison value, proving
\eqref{eq:f15v41:minimum}.
\end{proof}

The comparison is not a claimed
Haar-preserving reset or an RG
transformation. It supplies an upper
bound on an unchanged conditional
minimum. Gauge-orbit volume is not
discarded in the next probability
calculation.

\subsection{A normalized tail for arbitrary interior defects}

For \(0<\eta\le1\), define
\begin{equation}
 N_n(\eta)=\#\{p\in\mathcal P_n:|U_p-I|>\eta\},\qquad
 K_\beta=2^{13}\beta^{3/2}.
 \label{eq:f15v41:count}
\end{equation}
In particular a noncentral excursion
\(\operatorname{dist}(U_p,\{I,-I\})>\eta\)
is counted. The observable also counts
negative-center defects; they need not
be excluded from the interior.

\begin{theorem}[Source-uniform good-boundary block tail]
\label{f15v41:tail}
Let \(\beta\ge1\), and suppose the fixed
exterior satisfies
\eqref{eq:f15v41:boundary-event}.
For every integer \(r\ge1\),
\begin{equation}
 \begin{split}
 &\mathbb P_J\{N_n(\eta)\ge r\mid U_{A_n^c}=\xi\}\\
 &\quad\le
 \min\left\{1,\,
 K_\beta^{m_n}
 \exp\left[-\beta\left\{
       \frac{(1-\tau)r\eta^2}{2}
       -(1+\tau)D_n\,81n^4\delta^2\right\}\right]\right\}.
 \end{split}
 \label{eq:f15v41:tail}
\end{equation}
No small-field restriction is placed
on the sampled interior configuration.
\end{theorem}

\begin{proof}
On \(N_n(\eta)\ge r\), positivity of
all other plaquette terms gives
\[
                        S_n>\frac{(1-\tau)r\beta\eta^2}{2}.
\]
Subtract \eqref{eq:f15v41:minimum}.
The resulting event is a conditional
energy-excess event for the original
Wilson law. Its exact conditional
density is
\[
        \frac{e^{-(S_n-s_n)}}{
                 \int e^{-(S_n-s_n)}\,dH_{A_n}}\,dH_{A_n}.
\]
Equation~\eqref{eq:f15v40:conditional-tail}
bounds its excess tail by \(K_\beta^{m_n}e^{-s}\).
Its proof used a product-Haar ball at
an actual minimum and a global upper
bound on the Wilson second derivative,
not a lower Hessian bound. It is
uniform in this exterior and in all
the permitted couplings.
Apply it when the displayed excess
threshold is positive. If that
threshold is nonpositive, the minimum
with one on the right side of
\eqref{eq:f15v41:tail} is one, which
is the trivial bound.
\end{proof}

A convenient explicit choice is
\begin{equation}
                         \delta\le\frac{\eta}{64n^4}.
 \label{eq:f15v41:choice}
\end{equation}
Indeed \eqref{eq:f15v41:D} gives
\[
 \frac{(1+\tau)D_n81n^4\delta^2}{\eta^2}
 \le\frac{243(1+\tau)}{1024n}
 \le\frac{8019}{65536}
 <\frac{31}{128}\le\frac{1-\tau}{4}.
\]
Therefore
\begin{equation}
 \mathbb P_J\{N_n(\eta)\ge r\mid\xi\}
 \le
 \min\left\{1,\,
       K_\beta^{m_n}
       e^{-(1-\tau)(2r-1)\beta\eta^2/4}\right\}
 \quad\hbox{on the specified boundary event}.
 \label{eq:f15v41:count-tail}
\end{equation}
No combinatorial factor counts subsets
of \(r\) bad plaquettes: the total
action already bounds the count event.

For \(n=2\), this is an eight-link
interior vertex star, not V39's eight
parallel links. It has
\[
 m_2=8,\quad|\mathcal P_2|=24,\quad
 |\mathcal P_2^\partial|=192,\quad
 D_2=48,\quad9n^2=36.
\]
With all 192 boundary plaquettes
within \(\eta/1024\) of \(I\), the
probability of any \(\eta\)-defect
among its 24 incident plaquettes is
at most
\begin{equation}
          \min\{1,\,
                    K_\beta^8e^{-(1-\tau)\beta\eta^2/4}\}.
 \label{eq:f15v41:small-box}
\end{equation}
The action of every possible interior
field has been included in this
conditional probability.

\subsection{Exterior-weighted products and boundary-failure transfer}

Consider finitely many fixed boxes
\(B_i\) of sides \(n_i\), with pairwise
disjoint vertex sets. Let \(A_i\)
be their corresponding interior edge
sets, and put
\[
             \mathcal X=\sigma(U_e:e\notin\bigcup_iA_i).
\]
Let \(G_i\) be their boundary events,
and \(E_i=\{N_{n_i}(\eta_i)\ge r_i\}\).
Every \(G_i\) is \(\mathcal X\)-measurable.
Let \(q_i\in[0,1]\) be the appropriate
right side of \eqref{eq:f15v41:tail}
or \eqref{eq:f15v41:count-tail}.

\begin{theorem}[A normalized product interface with no square-root loss]
\label{f15v41:product}
For every nonnegative
\(\mathcal X\)-measurable integrable
weight \(H\),
\begin{equation}
 \mathbb E_J\!\left[H\prod_i\mathbf1_{E_i}\right]
 \le
 \mathbb E_J\!\left[
        H\prod_i(q_i\mathbf1_{G_i}+\mathbf1_{G_i^c})\right]
 \le
 \mathbb E_J\!\left[H\prod_i(q_i+\mathbf1_{G_i^c})\right].
 \label{eq:f15v41:weighted}
\end{equation}
In particular,
\begin{equation}
       \mu_J\!\left(\bigcap_i(E_i\cap G_i)\right)
                    \le\left(\prod_iq_i\right)
                           \mu_J\!\left(\bigcap_iG_i\right).
 \label{eq:f15v41:joint-good}
\end{equation}
If a further estimate
\begin{equation}
            \mu_J\!\left(\bigcap_{i\in I}G_i^c\right)
                               \le b^{|I|}
                    \quad\hbox{for every subfamily }I
 \label{eq:f15v41:missing-boundary}
\end{equation}
were available, and \(q_i\le q\), then
\begin{equation}
                   \mu_J\!\left(\bigcap_iE_i\right)
                                      \le\min\{1,q+b\}^{\,k},
                 \qquad k=\#\{i\}.
 \label{eq:f15v41:transfer}
\end{equation}
The estimate \eqref{eq:f15v41:missing-boundary}
is not asserted here.
\end{theorem}

\begin{proof}
Every plaquette depending on variables
of \(A_i\) lies inside \(B_i\).
No plaquette can depend on updated
variables from two boxes. Thus,
conditional on \(\mathcal X\), the
interior laws factor exactly. The
conditional probability of \(E_i\)
is at most \(q_i\) on \(G_i\), and
at most one otherwise. Multiplying
these conditional bounds, multiplying
by \(H\), and integrating proves the
first inequality. The second is
pointwise.

For \eqref{eq:f15v41:joint-good},
use \(H=\prod_i\mathbf1_{G_i}\).
For the last assertion, use \(H=1\),
replace \(q_i\) by \(q\), expand the
finite product in the last expression
of \eqref{eq:f15v41:weighted}, and
apply \eqref{eq:f15v41:missing-boundary}
to each selected subset. The binomial
sum is \((q+b)^k\).
\end{proof}

The absence of a square root here has
a specific reason: the boundary events
are measurable in the common exterior
of all updated variables. An arbitrary
good/bad decomposition need not have
this property. The nonnegative weight
\(H\) may reweight that entire exterior;
if \(0<\mathbb E_JH<\infty\), division
by that same expectation gives the
corresponding normalized statement.
This does not cover an additional
weight depending on the updated
interior links. The original allowed
physical sources are already included
in \(J_p\) throughout.

\subsection{Growing boxes, shrinking thresholds, and the remaining boundary problem}

Write \(t=\log\beta\). Fix \(a,b\ge0\).
If \(2\le n\le t^a\) and
\(t^{-b}\le\eta\le1\), then
\[
        \log K_\beta^{m_n}
                \le4t^{4a}(13\log2+\tfrac32t).
\]
For all sufficiently large \(\beta\),
uniformly over \(0\le\tau\le1/32\),
this is at most
\((1-\tau)\beta\eta^2/8\).
Under \eqref{eq:f15v41:choice}, the
single-defect conditional tail is
therefore at most
\begin{equation}
                     e^{-(1-\tau)\beta\eta^2/8}
                   \le e^{-31\beta/(256t^{2b})}.
 \label{eq:f15v41:polylog}
\end{equation}
The comparison is elementary:
an exponential in \(t\) dominates the
fixed power \(t^{4a+2b+1}\).
This proves superalgebraic conditional
control on such scales; it does not
prove the boundary premise on those
same scales.

To apply this to a V40 noncentral
carrier, choose a box whose
\(\mathcal P_n\) contains that finite
carrier. If any of its plaquettes
is farther than \(\eta\) from both
centers, then \(N_n(\eta)\ge1\).
The excursion inside this box has
now been paid on \(G_n(\delta)\).
Without \(G_n(\delta)\), the remaining
event is concretely
\begin{equation}
                  G_n(\delta)^c
                  =\left\{\max_{p\in\mathcal P_n^\partial}
                                         |U_p-I|>\delta\right\}.
 \label{eq:f15v41:boundary-remainder}
\end{equation}
This is a curvature event on the
three-dimensional boundary subcomplex,
not an unspecified exterior penalty.
It can include both noncentral and
negative-center plaquettes.

For one box the immediate consequence is
\[
        \mu_J\{N_n(\eta)\ge1\}
                    \le q_n+\mu_J(G_n(\delta)^c).
\]
For separated boxes the sharper product
interface is \eqref{eq:f15v41:weighted}.
Neither statement licenses setting the
boundary term to zero, applying the
homogeneous chessboard inequality to
the inhomogeneous law, or assuming
that a minimizing interior has small
absolute action for an arbitrary
boundary.

A finite menu of candidate boxes can
be used without a hidden selection
argument: the event that some candidate
has both the interior defect and its
good boundary is bounded by the sum
of its individual \(q_n\)'s. For
several targets, only choices with
disjoint selected boxes inherit the
joint product bound above. Existence
of such a controlled boundary menu
with sufficiently probable success,
and its geometric overlap cost,
remain to be proved.

This moves the next estimate upstream
to boundary-curvature probabilities
or a controlled treatment of coherent
nonflat boundary backgrounds. It does
not turn a conditional probability
bound into an unconditional mass gap.

\subsection{Exact diagnostics and preservation}

The new diagnostic is
\begin{center}\small
\path{scripts/verify_ym_f15_v41_boundary_gauge_escape.py}.
\end{center}
For \(n=2,\ldots,6\), it verifies five
complete box-incidence certificates,
3,680 boundary root paths and all
10,720 boundary edge homotopies.
Each of the 109,323 elementary square
interchanges is checked to lie on
the boundary, and the final path
word is checked exactly. This tests
the constructive proof itself, not
only the numerical area bound.

There are 837 rational source-uniform
threshold/count margins. Eight exact
noncommuting configurations check
1,664 unchanged exterior links, 192
extension-plaquette bounds and eight
complete comparison-action bounds.
They also check 1,536 boundary
curvature inequalities and 1,664
boundary-gauge inequalities using
the actual path products. Four
separated-support checks complete
the geometry tests.

The analytic conditional-normalizer
and product arguments prove the
probability statements. No empirical
probability, unverified numerical
minimum, or finite sample is used
to infer them. The diagnostic makes
no continuum mass-gap claim.

The script has 10,878 bytes and SHA--256
\begin{center}\scriptsize\ttfamily
1CD9655C0F85F6D26AD43487CCAF29E1A3DCF5C5D27CD48B01A0814843C40885.
\end{center}
The V40 predecessor has 935,238 bytes and SHA--256
\begin{center}\scriptsize\ttfamily
3BF2228A48AFFB1D4CF0DBE8A2546242170ED85785673283D11D623FF3C7C05C.
\end{center}
Removing this terminal section and
reversing the title version recovers
V40 exactly. V41 replaces V40 as
the sole active main source.
All 14 protected sources and 83
earlier certificates are preserved.
The next companion checkpoint is V45.

\begin{firewall}
V41 proves an explicit boundary-only
gauge, a conditional minimum bound
for arbitrary interiors, and a
normalized source-uniform block tail
with an exterior-weighted product
interface. The requisite joint
boundary-failure estimate is not
proved. General noncentral control,
V37's full actual-background
criterion, iterated RG, the interacting
continuum and the physical Yang--Mills
mass gap remain open.
\end{firewall}
\section{V42: coherent boundary flux, exact compact minima, and source-robust angle fluctuations}
\label{f15v42:section}

This iteration goes upstream from V41's
unpaid good-boundary event. A fixed
boundary can force curvature throughout
an interior plane. Consequently a
conditional estimate centered at the
identity cannot be uniformly small for
every exterior. We prove this on an
open set of exteriors, not just at a
single exceptional conditioning value.
We then calculate an exact nonflat
Dirichlet minimum and control the
plaquette-angle fluctuations around it
under the \emph{full compact} conditional
Wilson measure.

The positive result is not a mass-gap
theorem. It supplies a calibrated
boundary background and a quantitative
response to arbitrary permitted
inhomogeneous sources. It does not
estimate the probability of that
background in the unconditioned law.
The distinction is important: repeatedly
strengthening an identity-centered
conditional tail cannot eliminate
boundary-forced curvature.

\subsection{Input audit and the new scope}

The V41 predecessor and all 98 protected
files were hash-checked before this
append. The latest companion checkpoint
is V40; the next is V45. That checkpoint
includes the supplied SM source-selection
V31, regenerative stationarity V34 and
exact-action Einstein bridge V24, in
addition to the standing SM, NS and
shell-mixing files. These documents
remain source material, not instructions.

The archives already contain noncentral
conditional minimizers, in particular
the constrained sheet example in f12
V83. V39--V40 treat descent from
near-central negative sectors. Neither
phenomenon is claimed to be newly
discovered here. The new calculation
uses all interior compact links of the
V41 box, gives an explicit smooth
nonflat boundary family, proves its
\emph{global} conditional minimum,
and obtains an all-configuration
source-response bound. The compact
conditional-normalizer estimate
\eqref{eq:f15v40:conditional-tail}
is inherited, not reproved or counted
as a new result.

The objective remains the interacting
continuum theory and its physical
Hamiltonian gap. For external context,
the official
\href{https://www.claymath.org/millennium/yang-mills-the-maths-gap/}
{Clay problem page}, checked on
16 September 2026, still lists the
problem as unsolved. No external
mass-gap claim is an input to the
proofs below.

\subsection{The unchanged conditional box and two distances}

Use \(B_n=\{0,\ldots,n\}^4\), \(n\ge2\),
and the updated links \(A_n\) of V41:
an edge belongs to \(A_n\) exactly when
it is not contained in a boundary
facet. Let \(\mathcal P_n\) be all
plaquettes incident to \(A_n\).
They lie inside \(B_n\). Recall
\begin{equation}
 m_n=|A_n|=4n(n-1)^3,\qquad
 |\mathcal P_n|=6n^2(n-1)^2.
 \label{eq:f15v42:geometry}
\end{equation}
Embed the box without identifications
in a sufficiently large periodic
lattice and condition on every link
outside \(A_n\). All other action
terms are constants in this
conditional law.

Write
\begin{equation}
 \begin{split}
 \alpha(g)&=\arccos\!\left(\tfrac12\Tr g\right)\in[0,\pi],\\
 f(a)&=1-\cos a,\\
 \zeta(g)&=\min\{|g-I|,|g+I|\},\\
 S_J(U)&=\sum_{p\in\mathcal P_n}J_p f(\alpha(U_p)),\qquad
 s_J(\xi)=\min_{U_{A_n}} S_J(U;\xi).
 \end{split}
 \label{eq:f15v42:definitions}
\end{equation}
The norms in \(\zeta\) are Euclidean
quaternion norms, as in V40--V41.
All \(J_p\) are positive and finite.
For the source-uniform probability
statements we impose the same window
\[
 \beta\ge1,\qquad
 0\le\tau\le1/32,\qquad
 (1-\tau)\beta\le J_p\le(1+\tau)\beta .
\]
This includes every physical source
array retained in V38. Allowing all
plaquette arrays in the window only
strengthens the uniform statements;
it does not assert that every such
array is generated by a physical
cell source.

The angle \(\alpha\) is the geodesic
distance from the identity on the
unit-quaternion sphere with unit
radius. Left and right multiplication
are isometries. Hence
\(\alpha(gh)\le\alpha(g)+\alpha(h)\).
Both \(\alpha\) and \(\zeta\) are
invariant under conjugation and
inversion. In addition,
\begin{equation}
 \zeta(g_1\cdots g_N)
       \le\sum_{j=1}^N\zeta(g_j).
 \label{eq:f15v42:center-triangle}
\end{equation}
Indeed choose \(z_j\in\{I,-I\}\)
minimizing \(|g_j-z_j|\).
Their product is central, and a
telescoping product expansion,
together with invariance of the
quaternion norm under multiplication
by a unit quaternion, bounds the
distance to that product by
\(\sum_j|g_j-z_j|\).

For each fixed
\((x_3,x_4)\in\{1,\ldots,n-1\}^2\),
let \(D_{x_3,x_4}\) be the \(n\)-by-\(n\)
plaquette disk parallel to directions
1 and 2, and let \(C_{x_3,x_4}\)
be its positively oriented perimeter.
Every edge of this perimeter is fixed
by the exterior: it lies on
\(x_1=0,n\) or \(x_2=0,n\).
Every one of its \(n^2\) faces lies
in \(\mathcal P_n\). These
\((n-1)^2\) face sets are disjoint.
Their union is denoted
\(\mathcal P_n^{12}\); put
\begin{equation}
                 \mathsf N_n=|\mathcal P_n^{12}|
                            =n^2(n-1)^2.
 \label{eq:f15v42:N}
\end{equation}

\subsection{An exterior flux can force a noncentral interior}

\begin{proposition}[Disk holonomy and robust forced curvature]
\label{f15v42:disk}
Let \(D\) be one of the preceding
disks, \(C\) its perimeter, and
\(W_C\) its boundary holonomy.
For every interior link configuration,
\begin{equation}
 \alpha(W_C)\le\sum_{p\in D}\alpha(U_p),
 \qquad
 \zeta(W_C)\le\sum_{p\in D}\zeta(U_p).
 \label{eq:f15v42:disk}
\end{equation}
In particular, on the exterior event
\begin{equation}
 T_n(\eta)=
 \bigcap_{(x_3,x_4)\in\{1,\ldots,n-1\}^2}
       \{\zeta(W_{C_{x_3,x_4}})>n^2\eta\},
 \label{eq:f15v42:T}
\end{equation}
every interior has at least
\((n-1)^2\) distinct faces \(p\in\mathcal P_n\)
with \(\zeta(U_p)>\eta\).
\end{proposition}

\begin{proof}
Start at one corner of \(D\).
Let \(P_0\) be the monotone path
consisting of \(n\) direction-1 steps
followed by \(n\) direction-2 steps.
Interchange successive adjacent
\((1,2)\) steps to obtain the path
\(P_{n^2}\) with all direction-2
steps first. There are exactly \(n^2\)
interchanges, one across each
plaquette of the disk.

If \(P_j\) and \(P_{j+1}\) differ
across \(p_j\), cancellation of
their common suffix gives
\[
 U(P_j)U(P_{j+1})^{-1}
          =h_j U_{p_j}h_j^{-1},
\]
where \(h_j\) transports the common
prefix from the starting corner to
the lower corner of that plaquette.
Multiplying in this order telescopes:
\[
 W_C=U(P_0)U(P_{n^2})^{-1}
       =\prod_{j=0}^{n^2-1}h_jU_{p_j}h_j^{-1}.
\]
This is a nonabelian identity; no
plaquette factors have been commuted.
The two triangle inequalities prove
\eqref{eq:f15v42:disk}. If all \(n^2\)
plaquettes of a disk obeyed
\(\zeta(U_p)\le\eta\), its second
inequality would contradict the
strict condition in
\eqref{eq:f15v42:T}. Choose one
violating face in each disk.
The disks have disjoint face sets,
giving the stated count.
\end{proof}

Here is an explicit nonempty open
set to which the proposition applies.
Let \(\sigma_3=\operatorname{diag}(1,-1)\)
and define on all links of \(B_n\)
\begin{equation}
 \overline U_1(x)=\exp(-i\theta x_2\sigma_3),
 \qquad
 \overline U_2(x)=\overline U_3(x)=\overline U_4(x)=I,
 \qquad
 \Theta=n^2\theta.
 \label{eq:f15v42:background}
\end{equation}
Freeze its boundary links and call
that boundary value \(\xi_\theta\).
All direction-\((1,2)\) plaquettes
have holonomy \(\exp(i\theta\sigma_3)\);
all other plaquettes are the identity.
Each disk perimeter has holonomy
\(\exp(i\Theta\sigma_3)\).
For \(0<\Theta\le1/4\) and
\begin{equation}
                     0<\eta<\frac{2\sin(\Theta/2)}{n^2},
 \label{eq:f15v42:eta}
\end{equation}
the value \(\xi_\theta\) belongs to
\(T_n(\eta)\). Holonomy and \(\zeta\)
are continuous, so \(T_n(\eta)\)
is an open nonempty event depending
only on finitely many exterior
links. It has positive exterior
product-Haar measure. Every
finite-volume Wilson density with
finite couplings is strictly
positive, so
\begin{equation}
 \mu_J(T_n(\eta))>0,\qquad
 \mathbb P_J\!\left\{
   \#\{p\in\mathcal P_n:\zeta(U_p)>\eta\}
        \ge(n-1)^2\,\middle|\,\xi\right\}=1
 \quad(\xi\in T_n(\eta)).
 \label{eq:f15v42:forced}
\end{equation}
The conditional probability uses
the canonical positive Wilson
kernel, which is defined for every
exterior.

Thus a bound strictly below one
for this noncentral event cannot
hold uniformly over all exteriors.
This is not merely an objection
at a measure-zero conditioning
value. It does \emph{not} contradict
a small marginal probability:
\(\mu_J(T_n(\eta))\) may itself be
very small as \(\beta\) grows.
Nor does it contradict V41, whose
good-boundary hypothesis is absent
here.

\subsection{The exact homogeneous compact minimum}

\begin{proposition}[Sharp disk energy inequality]
\label{f15v42:disk-energy}
If \(0\le\Theta\le\pi/2\), \(N\ge1\),
\(\alpha_j\in[0,\pi]\), and
\(\sum_{j=1}^N\alpha_j\ge\Theta\),
then
\begin{equation}
          \sum_{j=1}^N(1-\cos\alpha_j)
                         \ge N\{1-\cos(\Theta/N)\}.
 \label{eq:f15v42:disk-energy}
\end{equation}
\end{proposition}

\begin{proof}
Choose \(b_j\in[0,\alpha_j]\) with
\(\sum_j b_j=\Theta\); successive
allocation of the total \(\Theta\)
constructs them. Each \(b_j\le\Theta\).
The function \(f(a)=1-\cos a\)
is increasing on \([0,\pi]\) and
convex on \([0,\pi/2]\). Hence
\(\sum f(\alpha_j)\ge\sum f(b_j)
\ge N f(\Theta/N)\), by convexity.
There is no assumption that the
original angles \(\alpha_j\) are
in the convex interval.
\end{proof}

\begin{theorem}[A full compact nonflat Dirichlet minimum]
\label{f15v42:minimum}
Let \(J_p=\beta>0\), fix
\(\xi_\theta\) from
\eqref{eq:f15v42:background}, and
assume \(0\le\Theta=n^2\theta\le\pi/2\).
Then
\begin{equation}
       s_\beta(\xi_\theta)
              =\beta\mathsf N_n(1-\cos\theta)
              =S_\beta(\overline U).
 \label{eq:f15v42:minimum}
\end{equation}
The minimum is over every compact
link in \(A_n\), with no chart,
commutativity, or coarse-fibre
restriction on the competitors.
\end{theorem}

\begin{proof}
The boundary holonomy of each
direction-\((1,2)\) disk has angle
\(\Theta\). Proposition~\ref{f15v42:disk}
gives \(\sum_{p\in D}\alpha(U_p)\ge\Theta\).
Apply
Proposition~\ref{f15v42:disk-energy}
with \(N=n^2\), then sum the
\((n-1)^2\) disjoint disk inequalities.
Every action term in the other
five orientations is nonnegative.
This proves the lower bound in
\eqref{eq:f15v42:minimum}.
The displayed background has
angle \(\theta\) on all
\(\mathsf N_n\) affected
direction-\((1,2)\) faces and
zero angle elsewhere, so it
attains the lower bound.
\end{proof}

At this background every updated
one-link innovation \(Q_e\) of V38
vanishes. More explicitly, identify
links and staples with quaternions
in the V38 convention. An updated
edge in direction 1 or 2 has two
direction-\((1,2)\) plaquettes,
of opposite oriented phases in
its staple frame, and four flat
plaquettes. Thus
\begin{equation}
 \begin{array}{c|c|c|c}
 \text{direction of }e & H_e & F_e & Q_e\\ \hline
 1,2 & \beta(4+2\cos\theta)\overline U_e
       &2\beta(1-\cos\theta)&0\\
 3,4 &6\beta\overline U_e&0&0
 \end{array}
 \label{eq:f15v42:stars}
\end{equation}
All six incident faces lie inside
\(B_n\) for \(e\in A_n\).
Unlike the negative-center
configuration of V39, the
present background admits no
action-lowering change of the
whole block with this fixed
boundary: it is a global
conditional minimum. Positive
frustration therefore need not
be paid by conditional excess
above that minimum. The
boundary-induced minimum itself
must be accounted for.

\subsection{A global supporting quadratic, not a local Hessian assumption}

\begin{proposition}[A scalar quadratic valid through the compact domain]
\label{f15v42:scalar}
For \(0\le\theta\le1/16\) and
\(0\le a\le\pi\),
\begin{equation}
 f(a)\ge f(\theta)+\sin\theta(a-\theta)
                               +\frac{(a-\theta)^2}{16},
 \qquad
 f(a)\ge\frac{a^2}{8}.
 \label{eq:f15v42:scalar}
\end{equation}
\end{proposition}

\begin{proof}
The second inequality follows from
\(\sin(a/2)\ge a/\pi\), by
concavity of sine on \([0,\pi/2]\),
and \(\pi<4\):
\[
 f(a)=2\sin^2(a/2)\ge2a^2/\pi^2\ge a^2/8.
\]
If \(a\le1\), the interval between
\(a\) and \(\theta\) lies in
\([0,1]\), where
\(f''=\cos\ge1/2\).
Taylor's integral remainder is at
least \((a-\theta)^2/4\), stronger
than required.

If \(a\ge1\), then
\(f(\theta)\le\theta\sin\theta\),
because the derivative of
\(\theta\sin\theta-f(\theta)\)
is \(\theta\cos\theta\ge0\)
on the stated interval and its
value at zero is zero.
Using \(\sin\theta\le\theta\)
and \(\theta\le1/16\le a/16\),
\[
 \begin{split}
 f(a)-f(\theta)-\sin\theta(a-\theta)
   &\ge a^2/8-\theta a\\
   &\ge a^2/16
    \ge(a-\theta)^2/16 .
 \end{split}
\]
Together the two cases prove
the first inequality.
\end{proof}

For the rest of the homogeneous
and fixed-boundary source results,
assume \(0\le\Theta\le1/4\), so
\(\theta=\Theta/n^2\le1/16\).
Define gauge-invariant, nonnegative
angle deviations
\begin{equation}
 \begin{split}
 X_\theta^2(U)
    &=\sum_{p\in\mathcal P_n^{12}}
                          \{\alpha(U_p)-\theta\}^2,\\
 Y^2(U)
    &=\sum_{p\in\mathcal P_n\setminus\mathcal P_n^{12}}
                                           \alpha(U_p)^2,\\
 R_\theta(U)&=X_\theta^2(U)+2Y^2(U).
 \end{split}
 \label{eq:f15v42:R}
\end{equation}

\begin{theorem}[All-configuration angle coercivity and a normalized tail]
\label{f15v42:coercivity}
For every interior with boundary
\(\xi_\theta\),
\begin{equation}
 S_\beta(U)-s_\beta(\xi_\theta)
                         \ge\frac{\beta}{16}R_\theta(U).
 \label{eq:f15v42:coercivity}
\end{equation}
For \(\beta\ge1\) and \(r\ge0\),
with \(K_\beta=2^{13}\beta^{3/2}\),
\begin{equation}
 \mathbb P_\beta\{R_\theta\ge r\mid\xi_\theta\}
          \le\min\{1,K_\beta^{m_n}e^{-\beta r/16}\},
 \qquad
 \mathbb E_\beta[R_\theta\mid\xi_\theta]
          \le\frac{16(m_n\log K_\beta+1)}{\beta}.
 \label{eq:f15v42:tail}
\end{equation}
\end{theorem}

\begin{proof}
Sum the first inequality of
\eqref{eq:f15v42:scalar} on each
disk. The tangent term is
\[
 \sin\theta\sum_{p\in D}\{\alpha(U_p)-\theta\}
 =\sin\theta\left\{\sum_{p\in D}\alpha(U_p)-\Theta\right\}
 \ge0
\]
by its boundary holonomy. Use
the second scalar inequality
on the other orientations.
Subtract the exact minimum
\eqref{eq:f15v42:minimum};
this proves
\eqref{eq:f15v42:coercivity}.
In particular this is not an
extension of a local Taylor
expansion to unverified large
fields: the scalar inequality
holds for every angle in
\([0,\pi]\).

The conditional density is exactly
\[
 \frac{e^{-(S_\beta-s_\beta)}}{
       \int e^{-(S_\beta-s_\beta)}\,dH_{A_n}}\,dH_{A_n}.
\]
Use
\eqref{eq:f15v40:conditional-tail}
at excess \(\beta r/16\).
No gauge-orbit volume has been
dropped. Finally, for \(A\ge1\),
\(a>0\),
\(\int_0^\infty\min\{1,Ae^{-ar}\}\,dr
=(\log A+1)/a\).
Integrating the tail proves
the expectation bound.
\end{proof}

Only conjugacy-class angles
are controlled by \(R_\theta\).
Different noncommuting color
directions can have the same
angles. No estimate of full
relative plaquette holonomies,
uniqueness modulo gauge,
positive gauge-reduced Hessian,
or clustering follows from
\eqref{eq:f15v42:coercivity}
alone.

\subsection{Arbitrary permitted sources with the same fixed boundary}

The next estimate keeps
\(\xi_\theta\) fixed while every
affected plaquette coupling
varies independently in the
allowed interval. Thus its
uniformity is not obtained by
moving the boundary with the
source.

Put
\begin{equation}
 \begin{split}
 c_\tau&=\frac{1-\tau}{16},\qquad
 b_\tau=\tau\sqrt{\mathsf N_n}\sin\theta,\\
 D_\tau&=\frac{b_\tau^2}{4c_\tau}
         =\frac{4\tau^2\mathsf N_n\sin^2\theta}{1-\tau}.
 \end{split}
 \label{eq:f15v42:response-constants}
\end{equation}

\begin{theorem}[Fixed-boundary source response and fluctuation floor]
\label{f15v42:source}
For every permitted array \(J\)
and every compact interior,
\begin{equation}
 \begin{split}
 S_J(U)-S_J(\overline U)
 &\ge\beta\{c_\tau R_\theta-b_\tau X_\theta\}\\
 &\ge\beta\{c_\tau R_\theta-b_\tau\sqrt{R_\theta}\}\\
 &\ge\beta\{c_\tau R_\theta/2-2D_\tau\}.
 \end{split}
 \label{eq:f15v42:source-coercivity}
\end{equation}
The actual source-dependent
minimum satisfies
\begin{equation}
 0\le S_J(\overline U)-s_J(\xi_\theta)\le\beta D_\tau.
 \label{eq:f15v42:source-minimum}
\end{equation}
Every minimizing interior
\(U_J^*\) obeys
\begin{equation}
 R_\theta(U_J^*)\le
       \frac{b_\tau^2}{c_\tau^2}
       =\frac{256\tau^2\mathsf N_n\sin^2\theta}{(1-\tau)^2}.
 \label{eq:f15v42:minimizer-response}
\end{equation}
For \(r\ge0\), the original
normalized conditional law
obeys
\begin{equation}
 \begin{split}
 \mathbb P_J\{R_\theta\ge r\mid\xi_\theta\}
 &\le\min\{1,\,
       K_\beta^{m_n}\exp(2\beta D_\tau-\beta c_\tau r/2)\},\\
 \mathbb E_J[R_\theta\mid\xi_\theta]
 &\le\frac{b_\tau^2}{c_\tau^2}
           +\frac{2(m_n\log K_\beta+1)}{\beta c_\tau}.
 \end{split}
 \label{eq:f15v42:source-tail}
\end{equation}
\end{theorem}

\begin{proof}
Write \(J_p=\beta(1+v_p)\),
\(|v_p|\le\tau\). Set
\(d_p=\alpha(U_p)-\theta\)
on \(\mathcal P_n^{12}\).
Multiply the first scalar
inequality by \(J_p\) there,
and use its second inequality
on all other affected faces.
The total tangent contribution
on the parallel faces is
\[
 \beta\sin\theta\sum_p d_p+
 \beta\sin\theta\sum_p v_p d_p .
\]
The first sum is nonnegative
disk by disk. The second is
bounded below by
\[
 -\beta\tau\sin\theta\sum_p|d_p|
             \ge-\beta b_\tau X_\theta
\]
by Cauchy--Schwarz. The
quadratic contributions are
at least
\(\beta c_\tau X_\theta^2+
2\beta c_\tau Y^2\).
This proves the first line
of \eqref{eq:f15v42:source-coercivity}.
Use \(X_\theta\le\sqrt{R_\theta}\)
for the next line, and
\[
 c_\tau t^2-b_\tau t
     \ge \tfrac12c_\tau t^2-\frac{b_\tau^2}{2c_\tau}
     =\tfrac12c_\tau t^2-2D_\tau
 \quad(t\ge0)
\]
for the last one.

The stronger completed square
\[
 c_\tau X_\theta^2+2c_\tau Y^2-b_\tau X_\theta
 =c_\tau\left(X_\theta-\frac{b_\tau}{2c_\tau}\right)^2
       +2c_\tau Y^2-D_\tau
\]
gives a lower bound
\(-\beta D_\tau\) on the action
difference. Its minimum is
nonpositive because \(\overline U\)
is an admissible comparison.
This proves
\eqref{eq:f15v42:source-minimum}.
At a minimizing interior the
second line of
\eqref{eq:f15v42:source-coercivity}
is nonpositive. If
\(R_\theta>0\), divide
\(c_\tau R_\theta\le
b_\tau\sqrt{R_\theta}\)
by \(\sqrt{R_\theta}\);
the zero case is immediate.
This proves
\eqref{eq:f15v42:minimizer-response}.

Since \(s_J\le S_J(\overline U)\),
\[
 S_J(U)-s_J
       \ge S_J(U)-S_J(\overline U)
       \ge\beta(c_\tau R_\theta/2-2D_\tau).
\]
Apply the inherited exact
conditional-excess tail when
the right-hand threshold is
positive; otherwise the
minimum with one in
\eqref{eq:f15v42:source-tail}
is the trivial bound one.
Its integral is
\[
 \frac{2}{\beta c_\tau}
        \{m_n\log K_\beta+2\beta D_\tau+1\}.
\]
The identity
\(4D_\tau/c_\tau=b_\tau^2/c_\tau^2\)
gives the stated expectation.
\end{proof}

The deterministic term
\(b_\tau^2/c_\tau^2\) is a
source-response allowance, not
a thermal variance. There is
no claim that the perturbed
minimum remains exactly at
\(\overline U\). The proof
neither assumes monotonicity
of a nonabelian curvature
event in the couplings nor
identifies a source-dependent
minimum by a numerical search.
At \(\tau=0\), the sharper
homogeneous tail
\eqref{eq:f15v42:tail} applies.

\subsection{A source-adapted exact nonconstant calibration}

There is also an explicit
nonconstant coupling family
for which the exact minimum
can be recalibrated, rather
than merely bounded. This
result has a different
boundary quantifier from
Theorem~\ref{f15v42:source}.

For \(u,v\in\{0,\ldots,n-1\}\),
choose positive numbers \(w_{uv}\)
and set \(J_p=w_{uv}\) on
every affected direction-\((1,2)\)
face based at \((u,v,x_3,x_4)\).
Thus these couplings are
independent of \(x_3,x_4\).
The other orientation weights
may be arbitrary positive
numbers. Choose \(\lambda\ge0\)
so that
\begin{equation}
 \vartheta_{uv}=\arcsin(\lambda/w_{uv})\in[0,1/16],
 \qquad
 \Theta_w=\sum_{u,v}\vartheta_{uv}\le1/4.
 \label{eq:f15v42:weighted-profile}
\end{equation}
The principal arcsine is used.
Such choices always exist;
for example
\(0\le\lambda\le(\min w_{uv})/(64n^2)\)
gives
\(\vartheta_{uv}\le1/(32n^2)\)
and \(\Theta_w\le1/32\), since
\(\arcsin t\le2t\) on \([0,1/2]\).

Define
\begin{equation}
 \overline U^{\,w}_1(x)=
   \exp\!\left(-i\sigma_3
         \sum_{v=0}^{x_2-1}\vartheta_{x_1,v}\right),
 \qquad
 \overline U^{\,w}_2=\overline U^{\,w}_3
                   =\overline U^{\,w}_4=I,
 \label{eq:f15v42:weighted-background}
\end{equation}
and freeze its boundary,
denoted \(\xi_w\).
Empty sums are zero.
For links in direction 1,
\(x_1\le n-1\), so every
index in this formula is
defined.

\begin{theorem}[Constant sine current calibrates the weighted compact minimum]
\label{f15v42:weighted}
For every compact interior with
boundary \(\xi_w\),
\begin{equation}
 \begin{split}
 S_J(U)-S_J(\overline U^{\,w})
 &\ge\frac1{16}
       \sum_{\substack{p\in\mathcal P_n^{12}\\
                       p\text{ based at }(u,v,x_3,x_4)}}
               J_p\{\alpha(U_p)-\vartheta_{uv}\}^2\\
 &\quad+\frac18
       \sum_{p\in\mathcal P_n\setminus\mathcal P_n^{12}}
                                      J_p\alpha(U_p)^2 .
 \end{split}
 \label{eq:f15v42:weighted-coercivity}
\end{equation}
In particular,
\begin{equation}
 s_J(\xi_w)=S_J(\overline U^{\,w})
        =(n-1)^2\sum_{u,v}w_{uv}(1-\cos\vartheta_{uv}).
 \label{eq:f15v42:weighted-minimum}
\end{equation}
\end{theorem}

\begin{proof}
The parallel plaquette at \((u,v,x_3,x_4)\)
of \(\overline U^{\,w}\) is
\(\exp(i\vartheta_{uv}\sigma_3)\);
all other plaquettes are \(I\).
Every parallel disk perimeter
has angle \(\Theta_w\).
Use the scalar supporting
inequality at \(\vartheta_{uv}\).
Its weighted tangent coefficient
is \(w_{uv}\sin\vartheta_{uv}
=\lambda\), independent of
the face. Thus on each disk
the sum of tangent terms is
\[
 \lambda\left\{\sum_{p\in D}\alpha(U_p)-\Theta_w\right\}
 \ge0.
\]
The remaining quadratic terms
are exactly those in
\eqref{eq:f15v42:weighted-coercivity}.
Nonparallel faces use the
second scalar inequality.
The background attains zero
on the right, proving its
global minimality and the
explicit action formula.
\end{proof}

If these weights also lie in
the permitted source window,
the same inherited conditional
normalizer gives the corresponding
quadratic-angle tail. This is
an exact solvable nonconstant
subfamily, not a reduction of
an arbitrary source array to
plane-parallel couplings.
In particular \(\xi_w\) is
adapted to the weights.
It must not be substituted
for a fixed exterior during
a physical source derivative.
Theorem~\ref{f15v42:source},
not this change of boundary,
is the result with the
required fixed-boundary
source quantifier.

\subsection{What this changes in the upstream problem}

There are now two quantitatively
distinct mechanisms. An interior
can carry energy above the
minimum imposed by its boundary;
the exact compact normalizer
controls that excess. It can
also sit at a nonflat
conditional minimum with
positive one-link frustration
and zero innovation. The
coherent family proves the
second possibility in the
same box and the same full
compact link domain.

For fixed \(n,\theta,\tau\),
\eqref{eq:f15v42:source-tail}
controls fluctuations about
the coherent angle profile
up to a deterministic source
displacement. It does not
make the profile itself
unlikely. The open-set
statement
\eqref{eq:f15v42:forced}
also prevents using a
single all-exterior
identity-centered contraction
as a substitute for a
boundary estimate.

The next useful extension is
to quantify the response to
boundary perturbations and
then seek an actual marginal
estimate for the retained
boundary flux or its
scale-dependent background.
A proof that merely assumes
such a marginal estimate
would not close the present
bottleneck. Any proposed
multiscale use must also
control noncommuting color
directions, background
selection and its
normalization. None of these
requirements is paid by the
angle-only inequality.

\subsection{Exact diagnostics, source integrity, and artifact placement}

The independent diagnostic is
\begin{center}\small
\path{scripts/verify_ym_f15_v42_coherent_boundary_flux.py}.
\end{center}
It checks 203 free-word square
interchanges in seven complete
nonabelian disk products, four
full box-incidence calculations,
and ten exterior-only plane
perimeters. Rational unit
quaternions verify 480
background plaquette identities,
four complete comparison-action
identities, ten strict flux
forcing margins and 208
zero-innovation staple
identities. Thirteen exact
weighted sine-current balances
check the nonconstant
calibration. Twelve
noncommuting comparison-action
samples supplement, but do
not prove, global minimality.

For the scalar inequalities
the script checks 385 strictly
certified rational Taylor
interval remainders and
three exact zero cases.
It also checks 400 exact
source-response square
completions. The proofs above,
not this finite grid or
the sampled configurations,
establish the inequalities
throughout the compact domain
and the normalized probability
bounds.

The script has 9,237 bytes and SHA--256
\begin{center}\scriptsize\ttfamily
66EAD297CFA8A658F02073D6147C7628222CE1EA2E263A85D706CE503DC25E7F.
\end{center}
The V41 predecessor has 956,938 bytes and SHA--256
\begin{center}\scriptsize\ttfamily
BF924141CC24EC379B8941AD61B8E8F166E3E76FF4045645984521A5BE28E067.
\end{center}
Removing this terminal section
and reversing the title
version recovers V41 exactly.
V42 replaces V41 as the
sole active main source.
All 14 protected sources
and 84 earlier certificates
are preserved. The next
companion checkpoint is V45.

The \path{latex_notes} folder
contains only \path{.tex}
sources. Compilation output
for this version belongs in
\begin{center}\small
\path{artifacts/ym_f15_v42_texcheck/}.
\end{center}
No PDF, auxiliary, log, or
other build file is written
to the source folder.

\begin{firewall}
V42 proves a robust
boundary-forced noncentral
event, an exact full-compact
nonflat Dirichlet minimum,
global angle coercivity,
fixed-boundary source
response and normalized
angle tails, together with
a source-adapted exact
nonconstant calibration.
It does not prove a
marginal boundary-flux
estimate, general noncentral
suppression, the complete
actual-background RG
criterion, an interacting
continuum limit, or the
physical Yang--Mills mass
gap. The goal remains open.
\end{firewall}
\section{V43: boundary calibration and compact angle tails}
\label{f15v43:section}

V42 identifies a nonflat boundary
background and its exact compact
Dirichlet minimum. An exactly
prescribed boundary, however, is
not a positive-measure event.
This iteration extends that
calculation to an open set of
noncommuting boundary perturbations.
The comparison action and a
boundary-flux lower calibration
have the same first derivative
at the coherent background.
Their difference is therefore
quadratic, with a side-uniform
constant in an explicit boundary
link metric.

The result retains the original
compact conditional measure,
allows all physical source
arrays in the V38 window, and
uses a gauge-invariant best-fit
boundary cost. It does not
assume that a typical boundary
has a small cost. In particular,
this extension removes the
restriction to one precisely
fixed coherent boundary, not
the unproved marginal
boundary-probability estimate.

\subsection{Audit, nonduplication, and notation}

The V42 predecessor and all 99
protected sources and certificates
were checked unchanged before
this append. No companion
checkpoint is due before V45.
The previous goal turn made
concrete progress: V42 changed
the available conditional
minimum and source-response
estimates. No no-progress
blocker is being declared.

The fixed-domain covariance
response of V13 and f12 V21
concerns normalized derivatives
in declared chart or fibre
regimes. It is not the
all-configuration compact
Dirichlet calibration below.
The ordinary product
normalization and the
conditional-excess estimate
\eqref{eq:f15v40:conditional-tail}
are inherited. The
nonabelian disk identity and
the global scalar supporting
inequality are exactly the
V42 inputs; they are not
claimed anew.

Retain \(B_n,A_n,\mathcal P_n\),
\(\mathcal P_n^{12}\), \(m_n\),
the parallel disks \(D_C\)
and their perimeters \(C\).
There are \((n-1)^2\) disks.
In this section write
\[
 N=n^2,\qquad
 f(a)=1-\cos a,\qquad
 K_\beta=2^{13}\beta^{3/2}.
\]
The angle \(\alpha\) is the
unit-quaternion geodesic angle.
The conditional action \(S_J\)
and its minimum \(s_J(\xi)\)
include only affected
plaquettes, as in V42.

\subsection{The boundary links which actually enter the conditional action}

Let \(F_n\) be the set of
fixed boundary edges incident
to \(\mathcal P_n\). Let
\(F_n^{12}\subset F_n\)
be those incident to
\(\mathcal P_n^{12}\).
An edge \(e=(x,\mu)\) belongs
to \(F_n\) if and only if
exactly one of its three
transverse coordinates,
say \(x_\nu\), is \(0\) or
\(n\), while the other two
are in \(\{1,\ldots,n-1\}\).
The affected face incident
to it must have directions
\(\mu,\nu\), with its
\(\nu\) step pointing into
the box. It is unique.
Consequently
\begin{equation}
 \begin{gathered}
 |F_n|=24n(n-1)^2,\qquad
 |F_n^{12}|=4n(n-1)^2,\\
 \#(F_n\cap\partial p)\le2
                   \quad(p\in\mathcal P_n).
 \end{gathered}
 \label{eq:f15v43:incidence}
\end{equation}
The first count chooses the
edge direction, the one
boundary transverse direction,
the side of the box, the
longitudinal coordinate, and
the two internal transverse
coordinates. The second
count restricts the first
two choices to directions
1 and 2. No other fixed
edge enters \(S_J\).
Moreover \(F_n^{12}\) is
the disjoint union of the
disk perimeter edge sets.

This one-face incidence is
stronger than the generic
six-plaquette incidence
used for arbitrary varying
links. It is the reason
that the boundary comparison
below has a side-uniform
second-derivative constant.

\subsection{Measured flux angles and a gauge-invariant coherent fit}

For any exterior \(\xi\),
define
\begin{equation}
 \Lambda_C(\xi)=\alpha(W_C(\xi)),\qquad
 \vartheta_C(\xi)=\Lambda_C(\xi)/N.
 \label{eq:f15v43:flux}
\end{equation}
These use only fixed
boundary links. On the
event \(\Lambda_C\le1/4\)
for every \(C\), set
\begin{equation}
 \begin{split}
 \mathcal L(\xi)
       &=\sum_C N f(\vartheta_C),\\
 \mathcal L_J(\xi)
       &=\sum_C\sum_{p\in D_C}J_p f(\vartheta_C),\\
 \mathcal X_\xi^2(U)
       &=\sum_C\sum_{p\in D_C}
                          \{\alpha(U_p)-\vartheta_C\}^2,\\
 \mathcal Y^2(U)
       &=\sum_{p\in\mathcal P_n\setminus\mathcal P_n^{12}}
                                             \alpha(U_p)^2,\\
 \mathcal R_\xi(U)
       &=\mathcal X_\xi^2(U)+2\mathcal Y^2(U).
 \end{split}
 \label{eq:f15v43:profile}
\end{equation}
The formulas themselves are
well defined for every
exterior; the stated
angle bound is used in
their inequalities. The
profile is inferred from
the actual perimeter
angles, not from a selected
gauge or a fitted color axis.

Let \(\overline U^\Theta\)
be the coherent field
\eqref{eq:f15v42:background}
with \(\theta=\Theta/N\).
Define
\begin{equation}
 \mathcal E_n(\xi)=
 \min_{\substack{0\le\Theta\le1/8\\
                 g\in\SU(2)^{V(\partial B_n)}}}
 \ \sum_{e=(x,\mu)\in F_n}
 \alpha\!\left((\overline U^\Theta_e)^{-1}
                  g_x\xi_e g_{x+e_\mu}^{-1}\right)^2.
 \label{eq:f15v43:fit}
\end{equation}
This is a finite-dimensional
compact minimization of a
continuous function.
The minimum is attained
and depends continuously
on \(\xi\). For the last
assertion, uniform
continuity on the product
of the compact parameter
and boundary spaces
bounds the change of
the minimum by the
uniform change of the
objective. It is gauge
invariant: a boundary
gauge applied to \(\xi\)
can be absorbed into
the minimizing \(g\).

Put
\begin{equation}
 \varepsilon_n=\frac1{256n},\qquad
 \Gamma_n=\{\mathcal E_n\le\varepsilon_n\},\qquad
 \Gamma_n^\circ=\{\mathcal E_n<\varepsilon_n\}.
 \label{eq:f15v43:tube}
\end{equation}
The open set \(\Gamma_n^\circ\)
is nonempty, since the
coherent boundary has
zero fit cost. It has
positive Haar measure,
and positive probability
under every finite-volume
Wilson law with finite
couplings. No lower bound
uniform in \(\beta\) is
being asserted.

Choose a minimizing pair
\((\Theta,g)\) for an
exterior in \(\Gamma_n\).
Extend \(g\) to interior
vertices arbitrarily.
Gauge-transform the
conditional configuration
and write on \(F_n\)
\begin{equation}
 \xi_e^g=\overline U^\Theta_e\exp X_e,\qquad
 E=\sum_{e\in F_n}|X_e|^2=\mathcal E_n(\xi),\qquad
 E_{12}=\sum_{e\in F_n^{12}}|X_e|^2.
 \label{eq:f15v43:logs}
\end{equation}
Here \(X_e\) is a shortest
Lie logarithm in the
unit-quaternion convention.
Since every \(|X_e|\le\sqrt E<\pi\),
there is no cut-locus
ambiguity in these
logarithms. A measurable
choice of the minimizing
pair is not required:
all final quantities
are the continuous
minimum and invariant
flux angles.

For \(0\le t\le1\), set
the comparison field
\(V(t)\) equal to
\(\overline U^\Theta\)
on \(A_n\), and to
\(\overline U^\Theta_e\exp(tX_e)\)
on \(F_n\). Other fixed
edges can be left at
their actual values;
they do not enter
the conditional action
or any of these
perimeters. Let
\(\xi(t)\) denote the
boundary relevant to
this comparison.
For each perimeter,
\begin{equation}
 \ell_C=\sum_{e\in C}|X_e|
       \le\sqrt{4nE}\le1/8,\qquad
 \Lambda_C(\xi(t))\le\Theta+t\ell_C\le1/4.
 \label{eq:f15v43:path}
\end{equation}
The holonomy distance
bound follows by
telescoping products
and bi-invariance.
In particular
\(\Gamma_n\) implies
the flux-angle condition
used in
\eqref{eq:f15v43:profile}.
An arbitrary minimizing
boundary gauge causes
no change of conditional
Haar volume: every
updated link undergoes
fixed left and right
multiplications.

\subsection{A smooth disk calibration and its derivative bound}

For \(N=n^2\), define on
the group ball
\(\alpha(W)\le1/4\)
\[
                         h_N(W)=Nf(\alpha(W)/N).
\]
This is smooth at the
identity. For example,
the logarithm coordinates
identify \(\alpha(W)^2\)
with the squared Euclidean
norm, and the even power
series of cosine gives
a smooth function of
that squared norm.

\begin{proposition}[Uniform derivatives of the disk calibration]
\label{f15v43:hessian}
On this group ball,
\begin{equation}
 \|\operatorname{Hess}h_N\|\le1/N,\qquad
 |\nabla h_N|\le\frac1{4N}.
 \label{eq:f15v43:hessian}
\end{equation}
If \(W(t)\) is a product
of fixed group factors
and geodesic factors
\(\exp(tX_j)\), and its
angle remains at most
\(1/4\), then
\begin{equation}
 \left|\frac{d^2}{dt^2}h_N(W(t))\right|
             \le\frac2N\left(\sum_j|X_j|\right)^2.
 \label{eq:f15v43:curve}
\end{equation}
\end{proposition}

\begin{proof}
Write \(a=\alpha(W)\).
The radial first and
second derivatives are
\(\sin(a/N)\) and
\(\cos(a/N)/N\).
The other Hessian
eigenvalue is
\(\sin(a/N)\cot a\).
Here is a direct
verification of this
last assertion. Along
a constant-speed great
circle \(W(s)\) on
the unit sphere, put
\(z(s)=\operatorname{Re}W(s)=\cos a(s)\).
Then \(z''=-|W'|^2z\),
which gives
\[
 a''=\cot a\{|W'|^2-(a')^2\}.
\]
Twice differentiating
\(h_N\) gives the
radial and transverse
coefficients
\[
 \lambda_{\rm rad}=\cos(a/N)/N,\qquad
 \lambda_{\rm tan}=\sin(a/N)\cot a.
\]
Both lie in \([0,1/N]\):
use \(\sin(a/N)\le a/N\)
and \(\cot a\le1/a\).
Their limits at zero
are \(1/N\). The
gradient bound follows
from \(\sin(a/N)\le a/N\).
This proves
\eqref{eq:f15v43:hessian}.

For the product curve,
let \(L=\sum_j|X_j|\).
Its right-trivialized
velocity is a sum of
conjugates of the
generators, so
\(|W'|\le L\).
Differentiating that
sum gives one Lie
bracket for each
ordered earlier/later
pair. The quaternion
Lie bracket satisfies
\(|[X,Y]|\le2|X||Y|\).
Consequently the
covariant acceleration
satisfies
\[
 |\nabla_t W'|
       \le2\sum_{i<j}|X_i||X_j|\le L^2.
\]
One can also see this
directly by writing
\(W'=A(t)W\):
\(W''=A'W-|A|^2W\),
so the tangential
acceleration has norm
\(|A'|\). Fixed
interspersed factors
only conjugate generators
and preserve their norms.

The chain rule and
\eqref{eq:f15v43:hessian}
therefore give
\[
 \left|(h_N(W(t)))''\right|
 \le L^2/N+L^2/(4N)
 \le2L^2/N.
\]
This proves
\eqref{eq:f15v43:curve}
also when the curve
passes through the
identity, by smoothness.
\end{proof}

The conditional normalized
homogeneous action of
the comparison, denoted
\[
                   A(t)=\sum_{p\in\mathcal P_n}
                                         f(\alpha(V_p(t))),
\]
has the global bound
\begin{equation}
 |A''(t)|
 \le\sum_{p\in\mathcal P_n}
            \left(\sum_{e\in F_n\cap\partial p}|X_e|\right)^2
 \le2E .
 \label{eq:f15v43:action-second}
\end{equation}
The first bound follows
by differentiating the
product of unit
quaternions twice.
The sum of the absolute
second-derivative terms
is at most the square
of the sum of generator
norms. The second
bound uses at most
two fixed edges per
face and exactly one
affected face per
edge of \(F_n\).

The boundary calibration
is
\[
            L(t)=\mathcal L(\xi(t))
                         =\sum_C h_N(W_C(\xi(t))).
\]
The perimeter sets
are disjoint, and
each has \(4n\) edges.
Thus
\begin{equation}
 |L''(t)|
   \le\frac2{n^2}\sum_C\ell_C^2
   \le\frac8n E_{12}.
 \label{eq:f15v43:calibration-second}
\end{equation}
This bound is uniform
as the coherent flux
tends to zero; no
division by that
flux has been hidden
in its constant.

\subsection{First-order cancellation and a quadratic minimum allowance}

\begin{proposition}[Cancellation of the boundary first variation]
\label{f15v43:cancellation}
For the comparison curve
above,
\begin{equation}
                 A(0)=L(0),\qquad A'(0)=L'(0).
 \label{eq:f15v43:cancellation}
\end{equation}
Consequently
\begin{equation}
       0\le A(1)-L(1)
              \le E+\frac4n E_{12}\le3E.
\label{eq:f15v43:calibration-debit}
\end{equation}
\end{proposition}

The constant three is
independent of \(n\), but
\(E\) is a sum, not
an average over boundary
links. A per-edge
perturbation bound
\(|X_e|\le\delta\)
only gives
\(E\le24n(n-1)^2\delta^2\).
That boundary-size cost
has not been removed.

\begin{proof}
At \(t=0\) the affected
field is the coherent
background. V42 gives
\(A(0)=(n-1)^2Nf(\Theta/N)=L(0)\).

For the first derivative,
all background links
commute with \(\sigma_3\).
Let \(x_e^3\) be the
coefficient of \(i\sigma_3\)
in \(X_e\). Conjugating
by a background link
preserves this component.
For one disk, put
\(\ell_C^3=\sum_{e\in C}\epsilon_e x_e^3\),
where \(\epsilon_e\)
is its perimeter
orientation sign.
Differentiating the
traces of the parallel
faces and cancelling
internal oriented
edge contributions
gives
\[
             \sum_{p\in D_C}
                \frac{d}{dt}f(\alpha(V_p(t)))\bigg|_{t=0}
                           =\sin(\Theta/N)\,\ell_C^3 .
\]
Only boundary edges
vary here, but the
same cancellation
identity applies.
The nonparallel
background plaquettes
are \(I\), at which
the first derivative
of \(f(\alpha(\,\cdot\,))\)
is zero.

For \(\Theta>0\),
the derivative of the
perimeter angle is
\(\Lambda_C'(0)=\ell_C^3\):
the real part of
its holonomy derivative
is \(-\sin\Theta\,\ell_C^3\).
Hence
\[
 \frac{d}{dt}h_N(W_C(\xi(t)))\bigg|_{t=0}
                  =\sin(\Theta/N)\,\ell_C^3.
\]
For \(\Theta=0\),
both first derivatives
are zero by smoothness.
Summing over disks
proves the second
identity in
\eqref{eq:f15v43:cancellation}.

Taylor's integral
remainder, with
\eqref{eq:f15v43:action-second}
and
\eqref{eq:f15v43:calibration-second},
now gives
\[
 |A(1)-L(1)|
 \le\int_0^1(1-t)(2E+8E_{12}/n)\,dt
 =E+4E_{12}/n.
\]
The lower bound zero
follows independently
from V42's disk energy
inequality: at \(t=1\),
each actual perimeter
has angle \(\Lambda_C\le1/4\),
and \(N\) faces of that
disk have total angle
at least \(\Lambda_C\).
All other face energies
are nonnegative.
\end{proof}

\begin{theorem}[Homogeneous compact minimum and angle tails on the boundary tube]
\label{f15v43:homogeneous}
For \(J_p=\beta\) and
\(\xi\in\Gamma_n\),
\begin{equation}
 0\le s_\beta(\xi)-\beta\mathcal L(\xi)
                         \le3\beta\mathcal E_n(\xi),
 \qquad
 S_\beta(U)-\beta\mathcal L(\xi)
                         \ge\frac{\beta}{16}\mathcal R_\xi(U).
 \label{eq:f15v43:homogeneous}
\end{equation}
Every minimizing interior
has
\(\mathcal R_\xi\le48\mathcal E_n(\xi)\).
For \(\beta\ge1\),
\(r\ge0\), the full
conditional law satisfies
\begin{equation}
 \begin{split}
 \mathbb P_\beta\{\mathcal R_\xi\ge r\mid\xi\}
  &\le\min\{1,K_\beta^{m_n}
           e^{\,3\beta\mathcal E_n(\xi)-\beta r/16}\},\\
 \mathbb E_\beta[\mathcal R_\xi\mid\xi]
  &\le48\mathcal E_n(\xi)
               +\frac{16(m_n\log K_\beta+1)}{\beta}.
 \end{split}
 \label{eq:f15v43:homogeneous-tail}
\end{equation}
\end{theorem}

\begin{proof}
The lower minimum bound
is the disk calibration,
valid for all compact
interiors. The comparison
\(V(1)\), followed by
undoing the fixed gauge,
is an admissible interior
with the actual exterior.
Its conditional action
is at most
\(\beta\mathcal L+3\beta E\).
Use the minimizing fit
\(E=\mathcal E_n(\xi)\)
for the upper bound.

On each disk apply
the scalar supporting
inequality
\eqref{eq:f15v42:scalar}
at \(\vartheta_C\).
Since \(\Lambda_C\le1/4\)
and \(N\ge4\), one has
\(\vartheta_C\le1/16\).
The summed tangent term
is nonnegative:
\[
 \sin\vartheta_C
        \left\{\sum_{p\in D_C}\alpha(U_p)-N\vartheta_C\right\}
 \ge0.
\]
The remaining quadratic
terms, including the
other orientations,
give the second
inequality of
\eqref{eq:f15v43:homogeneous}.
Evaluate it at a
minimizer for its
angle-deviation bound.

For any interior,
\[
 S_\beta-s_\beta
       \ge\beta\{\mathcal R_\xi/16-3\mathcal E_n(\xi)\}.
\]
The inherited normalized
conditional-excess tail
gives the first line
of
\eqref{eq:f15v43:homogeneous-tail}.
When the excess threshold
is nonpositive the
displayed bound by one
is valid. Integrating
\(\min\{1,Ae^{-ar}\}\)
as in V42 proves the
mean bound.
\end{proof}

The center of
\(\mathcal R_\xi\)
is the actual plane
flux profile. Thus
the forced noncentral
interiors of V42 are
not made artificially
rare by this theorem.
The controlled quantity
is departure from the
curvature required by
the boundary, not
absolute distance of
all plaquettes from
the center.

\subsection{Geometry of the comparison's angle error}

We need one more
deterministic estimate
before allowing arbitrary
source arrays. For a
fixed fit of cost \(E\),
the comparison field
satisfies
\begin{equation}
 \begin{split}
 \mathcal X_\xi^2(V(1))
       &\le(4+8/n)E_{12}\le8E_{12},\\
 \mathcal Y^2(V(1))
       &\le2(E-E_{12}),\\
 \mathcal R_\xi(V(1))&\le8E.
 \end{split}
 \label{eq:f15v43:comparison-angle}
\end{equation}
Indeed put
\(\ell_p=\sum_{e\in F_n\cap\partial p}|X_e|\).
For \(p\in D_C\), product
telescoping gives
\[
 |\alpha(V_p(1))-\Theta/N|\le\ell_p,\qquad
 |\vartheta_C-\Theta/N|\le\ell_C/N.
\]
Consequently
\[
 \sum_{p\in D_C}
       \{\alpha(V_p(1))-\vartheta_C\}^2
 \le2\sum_{p\in D_C}\ell_p^2+2\ell_C^2/N.
\]
Sum the disks and use
\(\sum_{p\in\mathcal P_n^{12}}\ell_p^2\le2E_{12}\)
and
\(\sum_C\ell_C^2\le4nE_{12}\).
For every other affected
face the background is
flat, so
\(\alpha(V_p(1))\le\ell_p\);
the same one-face
incidence gives its
stated bound. This
proves
\eqref{eq:f15v43:comparison-angle}
without linearizing
the perturbed holonomies.

\subsection{Arbitrary sources: keep the boundary, and pay the uncancelled source term}

Now let
\((1-\tau)\beta\le J_p\le(1+\tau)\beta\),
with \(0\le\tau\le1/32\).
The flux profile and
the coherent fit depend
on the actual exterior,
not on the source array.
Define
\begin{equation}
 \begin{split}
 c_\tau&=(1-\tau)/16,\\
 b_\xi&=\tau\left(N\sum_C\sin^2\vartheta_C\right)^{1/2},
       \qquad D_\xi=\frac{b_\xi^2}{4c_\tau},\\
 \mathcal B_\xi(e)&=(3+4\tau)e+\sqrt8\,b_\xi\sqrt e
                         \quad(e\ge0),\\
 B_\xi&=\mathcal B_\xi(\mathcal E_n(\xi)).
 \end{split}
 \label{eq:f15v43:source-constants}
\end{equation}
All these quantities
are exterior measurable.
The notation \(B_\xi\)
is a scalar allowance,
not the geometric box
\(B_n\).

\begin{theorem}[Source-uniform compact boundary-tube estimate]
\label{f15v43:source}
For every permitted source
array and every
\(\xi\in\Gamma_n\),
\begin{equation}
 \begin{split}
 S_J(U)-\mathcal L_J(\xi)
   &\ge\beta\{c_\tau\mathcal R_\xi
                          -b_\xi\mathcal X_\xi\}\\
   &\ge\beta\{c_\tau\mathcal R_\xi
                          -b_\xi\sqrt{\mathcal R_\xi}\},\\
 \mathcal L_J(\xi)-\beta D_\xi
   &\le s_J(\xi)
    \le\mathcal L_J(\xi)+\beta B_\xi.
 \end{split}
 \label{eq:f15v43:source-minimum}
\end{equation}
Every actual minimizing
interior obeys
\begin{equation}
 \sqrt{\mathcal R_\xi(U_J^*)}
       \le\frac{b_\xi+
                 \sqrt{b_\xi^2+4c_\tau B_\xi}}{2c_\tau}.
 \label{eq:f15v43:minimizer}
\end{equation}
For \(\beta\ge1\), \(r\ge0\),
\begin{equation}
 \begin{split}
 \mathbb P_J\{\mathcal R_\xi\ge r\mid\xi\}
   &\le\min\{1,K_\beta^{m_n}
       e^{\,\beta(B_\xi+2D_\xi)-\beta c_\tau r/2}\},\\
 \mathbb E_J[\mathcal R_\xi\mid\xi]
   &\le\frac2{c_\tau}
      \left\{B_\xi+2D_\xi+
                   \frac{m_n\log K_\beta+1}{\beta}\right\}.
 \end{split}
 \label{eq:f15v43:source-tail}
\end{equation}
\end{theorem}

\begin{proof}
Write \(J_p=\beta(1+v_p)\),
\(|v_p|\le\tau\).
Apply the V42 scalar
supporting inequality
at \(\vartheta_C\)
on disk \(D_C\).
The common, unweighted
tangent sum is
nonnegative on every
disk. The remaining
source tangent term
is bounded below by
\[
 -\beta\tau
       \sum_C\sum_{p\in D_C}
            \sin\vartheta_C\,
                 |\alpha(U_p)-\vartheta_C|
 \ge-\beta b_\xi\mathcal X_\xi.
\]
The quadratic terms
give
\(\beta c_\tau\mathcal R_\xi\).
This proves the first
two inequalities of
\eqref{eq:f15v43:source-minimum}.
Completing the square
in \(\mathcal X_\xi\)
gives the lower bound
on \(s_J\), with
allowance \(\beta D_\xi\).

For the upper bound,
use the same comparison
as before, with fit
cost \(E\), and put
\(a_p=\alpha(V_p(1))\).
The homogeneous part
of its action difference
from the calibration
is at most \(3\beta E\).
The remaining difference,
divided by \(\beta\),
is
\[
 \sum_C\sum_{p\in D_C}
       v_p\{f(a_p)-f(\vartheta_C)\}
       +\sum_{p\notin\mathcal P_n^{12}}v_p f(a_p),
\]
where the last sum is
still restricted to
\(\mathcal P_n\).
For all angles in
\([0,\pi]\), Taylor's
theorem with
\(|f''|\le1\) gives
\[
 |f(a)-f(v)|
       \le\sin v\,|a-v|+\tfrac12(a-v)^2,
 \qquad f(a)\le a^2/2 .
\]
Thus the absolute value
of the source remainder
is at most
\[
 b_\xi\mathcal X_\xi(V(1))
       +\frac{\tau}{2}
          \{\mathcal X_\xi^2(V(1))+\mathcal Y^2(V(1))\}
 \le\sqrt8\,b_\xi\sqrt E+4\tau E
\]
by
\eqref{eq:f15v43:comparison-angle}.
Together with the homogeneous
calibration this yields
\[
 S_J(V(1))-\mathcal L_J(\xi)
                  \le\beta\mathcal B_\xi(E).
\]
Take an attaining best
fit to obtain the
upper bound on the
actual minimum.
This step explicitly
pays the source
first variation; it
does not falsely
extend
\eqref{eq:f15v43:cancellation}
to inhomogeneous
plaquette weights.

At a minimizing
interior,
\[
 c_\tau\mathcal R_\xi
       -b_\xi\sqrt{\mathcal R_\xi}\le B_\xi.
\]
Solving this quadratic
inequality in
\(\sqrt{\mathcal R_\xi}\)
gives
\eqref{eq:f15v43:minimizer}.
For arbitrary interiors,
use
\[
 c_\tau t^2-b_\xi t
       \ge\tfrac12c_\tau t^2-2D_\xi
 \quad(t\ge0)
\]
and subtract the upper
bound on \(s_J\).
The conditional excess
therefore obeys
\[
 S_J-s_J
       \ge\beta\{c_\tau\mathcal R_\xi/2-2D_\xi-B_\xi\}.
\]
Apply the exact compact
conditional-normalizer
tail and integrate it,
as in the homogeneous
proof. This proves
\eqref{eq:f15v43:source-tail}.
\end{proof}

At \(\mathcal E_n=0\),
the estimate reduces
to the V42 coherent
source-response bound
for the flux range
used here. At \(\tau=0\),
the sharper homogeneous
statement
\eqref{eq:f15v43:homogeneous-tail}
should be used instead
of the half-coefficient
Young bound.

\subsection{A full-law, exterior-weighted consequence with no assumed tube probability}

The theorem can be
averaged under the
actual inhomogeneous
Wilson law, provided
the boundary restriction
and its allowance are
retained. For \(s\ge0\),
define the exterior-dependent
radius
\begin{equation}
 r_\xi(s)=\frac2{c_\tau}
       \left\{B_\xi+2D_\xi+
                     \frac{m_n\log K_\beta+s}{\beta}\right\}.
 \label{eq:f15v43:radius}
\end{equation}
The formulas for this
radius are measurable
also outside \(\Gamma_n\);
only their use on
\(\Gamma_n\) is asserted.
For any nonnegative
integrable exterior
weight \(H\),
\begin{equation}
 \mathbb E_J\!\left[
      H\,\mathbf1_{\Gamma_n}
         \mathbf1_{\{\mathcal R_\xi\ge r_\xi(s)\}}\right]
       \le e^{-s}\,
                 \mathbb E_J[H\,\mathbf1_{\Gamma_n}].
 \label{eq:f15v43:weighted}
\end{equation}
Condition on the exterior,
insert
\eqref{eq:f15v43:source-tail},
and note that the
\(m_n\log K_\beta\)
term cancels the
normalizer factor.
This proves the
inequality without
knowing
\(\mu_J(\Gamma_n)\).

For boxes with disjoint
vertex sets, condition
on the complement of
the union of their
updated links. Every
affected plaquette lies
in one box, so the
conditional law
factorizes. The
boundary fit and
perimeter angles of
each box are measurable
in this common exterior.
The same argument gives
\begin{equation}
 \mathbb E_J\!\left[
   H\prod_{i=1}^k
        \mathbf1_{\Gamma_i}
        \mathbf1_{\{\mathcal R_i\ge r_{\xi_i}(s_i)\}}\right]
  \le e^{-\sum_{i=1}^k s_i}
          \mathbb E_J\!\left[H\prod_{i=1}^k\mathbf1_{\Gamma_i}\right].
 \label{eq:f15v43:joint}
\end{equation}
This conditional-product
mechanism is inherited
from V41. The new
input is its
gauge-invariant,
flux-centered,
quadratically
boundary-stable radius.

There is no union bound
over fitted gauges or
coherent parameters.
For every fixed
exterior, the action
comparison is valid
for any candidate fit,
and compactness
supplies a minimizing
one. The random
observable
\(\mathcal R_\xi\)
does not depend on
which minimizer is
chosen. This is a
pointwise variational
comparison, not a
measure decomposition
over background fields.
It therefore does
not establish a
normalized marginal
law for such fields.

\subsection{What remains upstream}

The exact-boundary
limitation of V42 is
now replaced by a
nonempty open,
gauge-invariant boundary
tube. The error
allowance is
quadratic in its
link-distance fit
when the sources are
homogeneous, and has
an explicit additional
source term otherwise.
The cancellation is
uniform at zero flux
and does not require
commuting perturbations.

The required actual-law
estimate for the
probability of a
small fit, or for
an adequate family
of more general
backgrounds, is still
missing. In particular,
positive probability
of \(\Gamma_n^\circ\)
is not high probability,
and
\eqref{eq:f15v43:joint}
does not bound the
complementary event.
The selected homogeneous
CGF of V1/V32 cannot
be silently applied
to every inhomogeneous
source background.

There is also a
separate limitation:
the calibrated
observable controls
plaquette angles,
not their noncommuting
color directions.
It provides no
decay of correlations,
uniqueness of a
conditional minimizer
modulo gauge, or
physical spectral
gap. Further work
must attack the
actual marginal
background control
or enlarge the
class of backgrounds
with a paid measure
comparison; simply
renaming these missing
estimates as hypotheses
would not advance
the closure.

\subsection{Exact checks and file integrity}

The independent diagnostic is
\begin{center}\small
\path{scripts/verify_ym_f15_v43_boundary_calibration.py}.
\end{center}
Twelve complete box
configurations check
4,800 single-incidence
boundary-edge occurrences
and 4,416 plaquette
second-derivative
incidence bounds.
Rational quaternion
first and second
jets verify 12
zeroth/first-order
calibration cancellations,
12 quadratic debit
checks, 56 boundary
flux second-jet bounds,
and 28 radial/transverse
Hessian bounds.
They include flat
and nonflat backgrounds.

The test also checks
48 noncommuting
gauge-covariant
holonomy jets and
96 restored boundary
fit jets. Six
inhomogeneous examples
have a nonzero
first-variation debit,
confirming that
discarding the
source remainder
would be invalid.
There are 63
uniform-side
constant checks
and 36 exact
tail square
completions.

These checks test
the algebra and
incidence used in
the proof. Finite
jets do not by
themselves establish
global nonlinear
stability; that
statement follows
from the derivative
bounds along the
entire comparison
path above. No
sampled boundary
probability or
numerically fitted
minimum is used.

The script has 8,326 bytes and SHA--256
\begin{center}\scriptsize\ttfamily
614C8B5BB9227FFD71D6631F5D1DCF01F4AE1E1F0A510C9403C5D0207B624DB8.
\end{center}
The V42 predecessor has 983,482 bytes and SHA--256
\begin{center}\scriptsize\ttfamily
DE087A65B8572CF0375116759550F9B5BCBF0A1163F346BA5DB7C56E9873B7BE.
\end{center}
Removing this final
section and reversing
the title version
recovers V42 exactly.
V43 replaces V42
as the sole active
main source. All
14 protected sources
and 85 prior
certificates are
preserved. The next
companion checkpoint
remains V45.

Only \path{.tex}
sources belong in
\path{latex_notes}.
Compilation for
this version is
directed to
\begin{center}\small
\path{artifacts/ym_f15_v43_texcheck/}.
\end{center}

\begin{firewall}
V43 proves quadratic
homogeneous boundary
calibration, a
gauge-invariant
compact best-fit
extension to an
open boundary tube,
and source-uniform
normalized angle
tails with an
explicit boundary
and source allowance.
The actual marginal
boundary-background
estimate, general
noncentral control,
iterated RG,
the interacting
continuum and the
physical Yang--Mills
mass gap remain
unproved.
\end{firewall}
\section{V44: rooted colour rigidity and the compact minimum}
\label{f15v44:section}

V43 controls plaquette angles around an actual boundary-flux
profile. We began this iteration by returning to the missing
actual-source marginal estimate. The simple arbitrary-exterior
bulk-density repair is already in f12 V25, including its surface
debit and its failure to exclude thin defect connections.
Repeating that argument would not advance the current target.
No new marginal estimate is claimed here.

Instead we resolve a distinct limitation identified in V42--V43:
angle control alone had not provided an explicit estimate for
noncommuting colour directions. The ordered nonabelian disk
product has a useful additional structure. Its partial products
form a path on the unit-quaternion sphere, and its Wilson
action is exactly the chord Dirichlet energy of that path.
Expansion around the equally spaced short geodesic is an
exact identity over the whole compact domain. It yields
control of the full rooted plaquette holonomies, not just
their traces.

This also extends the exact coherent Dirichlet minimum to
the full principal flux interval and proves uniqueness of
its gauge orbit, with the gauge fixed on the boundary, for
flux strictly below the antipodal endpoint. These are finite
compact boundary-value results. They do not supply a
probability for the retained boundary, a contraction rate
for the full Wilson law, or a physical Hamiltonian gap.

\subsection{Audit and inputs which are not being reproved}

The V43 predecessor and all 100 protected files were
hash-checked before this append. The companion checkpoint
remains V45. The previous goal turn was progress: it proved
a quadratic boundary comparison and source-uniform
conditional tails on an open boundary tube.

The nonduplication review included f12 V25's arbitrary-exterior
defect-density theorem and thin-carrier warning, f12 V30's
explicit collar debit, and f14 V94's single-plaquette chord
identity. None is the ordered disk-path identity below.
Ordinary summation by parts and the one-dimensional Dirichlet
inequality are elementary tools, not claimed as new general
theorems. The new application in this ledger is to the actual
rooted nonabelian plaquette factors and their compact
conditional law.

We retain the V42 disk ordering and the V43 notation
\(B_n,A_n,\mathcal P_n,\mathcal P_n^{12},D_C,N=n^2,m_n\).
The quantities
\(\Lambda_C,\vartheta_C,\mathcal L,\mathcal L_J,
\mathcal E_n,\Gamma_n,B_\xi\)
are those of V43. The normalized excess bound
\eqref{eq:f15v40:conditional-tail} and the conditional-product
argument are inherited unchanged.

\subsection{An exact identity for a compact quaternion path}

Identify \(\SU(2)\) with the unit sphere in the quaternions.
The norm below is the Euclidean quaternion norm and
\(f(a)=1-\cos a\). Fix an integer \(N\ge2\) and an endpoint
\(W\) of angle \(a=\alpha(W)\in[0,\pi]\).
For \(a<\pi\), let
\[
 R=\exp\bigl(N^{-1}\log W\bigr),\qquad \theta=a/N,
\]
using the principal logarithm. At \(a=\pi\), choose any
unit imaginary axis and a geodesic root
\(R=\exp(\pi Z/N)\) with \(Z^2=-I\).
In both cases \(R^N=W\) and \(\operatorname{Re}R=\cos\theta\).

Let \(q_0,\ldots,q_N\) be arbitrary unit quaternions with
\(q_0=I,q_N=W\). Define
\begin{equation}
 \begin{split}
 \overline q_j&=R^j,\qquad d_j=q_j-\overline q_j,\\
 T(q)&=\sum_{j=1}^N|d_j-d_{j-1}|^2,\qquad
 D(q)=\sum_{j=1}^{N-1}|d_j|^2,\\
 E(q)&=\frac12\sum_{j=1}^N|q_j-q_{j-1}|^2 .
 \end{split}
 \label{eq:f15v44:path-data}
\end{equation}
The differences \(d_j\) are ambient quaternion vectors,
not Lie logarithms; no coordinate chart is imposed on
the sampled \(q_j\).

\begin{proposition}[Global chord Dirichlet identity]
\label{f15v44:identity}
For every such compact path,
\begin{equation}
                  E(q)-Nf(\theta)=\frac12T(q)-f(\theta)D(q).
 \label{eq:f15v44:identity}
\end{equation}
Put
\begin{equation}
 \lambda_N=2\{1-\cos(\pi/N)\},\qquad
 \rho_N(a)=\frac{f(a/N)}{f(\pi/N)} .
 \label{eq:f15v44:rho}
\end{equation}
Then
\begin{equation}
 D(q)\le\lambda_N^{-1}T(q),\qquad
 E(q)-Nf(\theta)\ge\frac{1-\rho_N(a)}2\,T(q).
 \label{eq:f15v44:path-coercivity}
\end{equation}
Consequently \(E(q)\ge Nf(a/N)\) for \(0\le a\le\pi\).
If \(a<\pi\), equality holds if and only if
\(q_j=R^j\) for every \(j\).
\end{proposition}

\begin{proof}
Expand the squared increments in the ambient vector space.
Since \(d_0=d_N=0\), summation by parts gives the cross term
\[
 \sum_{j=1}^N
    (\overline q_j-\overline q_{j-1})\cdot(d_j-d_{j-1})
 =\sum_{j=1}^{N-1}
    (2\overline q_j-\overline q_{j-1}-\overline q_{j+1})\cdot d_j .
\]
The equally spaced reference satisfies
\[
 2\overline q_j-\overline q_{j-1}-\overline q_{j+1}
                =2f(\theta)\overline q_j.
\]
Both \(q_j\) and \(\overline q_j\) have unit norm, whence
\(\overline q_j\cdot d_j=-|d_j|^2/2\).
The cross term is therefore \(-f(\theta)D(q)\).
The reference energy is \(Nf(\theta)\), proving
\eqref{eq:f15v44:identity}.

For completeness, the scalar Dirichlet matrix on
\(j=1,\ldots,N-1\) has diagonal 2 and neighbouring
entries \(-1\). Its vectors
\(\sin(k\pi j/N)\), \(1\le k\le N-1\), have distinct
eigenvalues \(2\{1-\cos(k\pi/N)\}\). They form an
orthogonal basis because the matrix is symmetric and the
eigenvalues are distinct. Its smallest eigenvalue is
\(\lambda_N\). Apply this to each of the four real
components of \(d\) to obtain the first inequality in
\eqref{eq:f15v44:path-coercivity}.
Substitution into the exact identity proves the second.

The function \(f\) is strictly increasing on \([0,\pi]\).
Thus \(\rho_N(a)\le1\), and it is strictly less than one
when \(a<\pi\). In that case equality in the energy
minimum forces \(T(q)=0\). The zero endpoints of \(d\)
then give \(d_j=0\) for all \(j\). The reference path
conversely attains the minimum.
\end{proof}

\subsection{From path displacement to the full group factors}

Set
\begin{equation}
                 G_j=q_{j-1}^{-1}q_j,\qquad
                 Z(q)=\sum_{j=1}^N|G_j-R|^2 .
 \label{eq:f15v44:Z}
\end{equation}
The product \(G_1\cdots G_N\) is \(W\). The energy is
\(\sum_j f(\alpha(G_j))\); its expression as \(E(q)\)
uses bi-invariance of the quaternion norm.

\begin{proposition}[Full-factor colour coercivity]
\label{f15v44:colour}
For \(0\le a<\pi\),
\begin{equation}
 E(q)-Nf(a/N)
       \ge\frac{1-\rho_N(a)}{4\{1+\rho_N(a)\}}\,Z(q).
 \label{eq:f15v44:colour}
\end{equation}
In particular, for \(0\le a\le1/4\),
\begin{equation}
 \begin{split}
 E(q)-Nf(a/N)&\ge\frac{63}{128}\,T(q),\\
 E(q)-Nf(a/N)&\ge\frac{63}{260}\,Z(q)\ge\frac15Z(q),\\
 \max_{0\le j\le N}|q_j-R^j|^2
       &\le\frac{32N}{63}\{E(q)-Nf(a/N)\}.
 \end{split}
 \label{eq:f15v44:small-colour}
\end{equation}
\end{proposition}

\begin{proof}
Since \(\overline q_j=\overline q_{j-1}R\),
\[
 |G_j-R|
    =|q_j-q_{j-1}R|
    =|(d_j-d_{j-1})+d_{j-1}(I-R)|.
\]
The quaternion norm is multiplicative, also for
nonunit quaternions. Squaring with the factor-two
inequality, summing and using
\(|I-R|^2=2f(\theta)\) gives
\[
 Z(q)\le2T(q)+4f(\theta)D(q)
                         \le2\{1+\rho_N(a)\}T(q).
\]
Combine this with
\eqref{eq:f15v44:path-coercivity}.

For the small-angle constants,
\[
 f(a/N)\le\frac{a^2}{2N^2},\qquad
 f(\pi/N)=2\sin^2(\pi/(2N))\ge\frac2{N^2}.
\]
The second inequality follows from concavity of sine
on \([0,\pi/2]\). Hence
\(\rho_N(a)\le a^2/4\le1/64\).
Substitution gives \(63/128\) and \(63/260\);
the latter is larger than \(1/5\).

Finally let \(\Delta d_k=d_k-d_{k-1}\).
The endpoint constraint implies
\[
 d_j=\frac{N-j}{N}\sum_{k\le j}\Delta d_k
                -\frac jN\sum_{k>j}\Delta d_k .
\]
Cauchy--Schwarz yields
\[
 |d_j|^2\le\frac{j(N-j)}{N}\,T(q)\le\frac N4T(q).
\]
Combine with the first line of
\eqref{eq:f15v44:small-colour}.
\end{proof}

The factor \(N\) in the partial-transport bound is retained.
Also \(\lambda_N\) tends to zero as the path length grows.
This finite path Dirichlet calculation is not a
volume-independent physical spectral gap.

\subsection{Apply the path identity to the actual nonabelian disk}

Use the V42 interchange ordering which moves the direction-1
steps to the right of the direction-2 steps. Let
\(P_0,\ldots,P_N\) be the successive monotone paths across
one \(n\)-by-\(n\) disk, and put
\begin{equation}
 \begin{split}
 G_{C,j}&=U(P_{j-1})U(P_j)^{-1}
                    =h_{C,j}U_{p_j}h_{C,j}^{-1},\\
 q_{C,j}&=G_{C,1}\cdots G_{C,j}
                    =U(P_0)U(P_j)^{-1},\\
 q_{C,0}&=I,\qquad q_{C,N}=W_C .
 \end{split}
 \label{eq:f15v44:disk-path}
\end{equation}
Every face occurs exactly once. The factors are in their
specified order; they have not been commuted.
The equality
\(G_{C,j}=q_{C,j-1}^{-1}q_{C,j}\)
then follows by telescoping.
Since conjugation preserves the normalized trace,
\begin{equation}
       \sum_{p\in D_C}f(\alpha(U_p))
             =\sum_{j=1}^N f(\alpha(G_{C,j}))
             =E(q_C).
 \label{eq:f15v44:disk-energy}
\end{equation}

This is a deterministic construction from the original
links. It is not a change of integration variables
which declares the rooted factors independent.
In particular, no product-Haar density on these factors
is substituted for the Wilson conditional measure.

For \(\Lambda_C<\pi\), define
\[
                         R_C=\exp(N^{-1}\log W_C).
\]
Under a vertex gauge transformation, \(W_C,R_C,q_{C,j}\)
and \(G_{C,j}\) are all conjugated by the gauge at the
starting corner of \(C\). Thus the differences and norms
in the following observable are gauge invariant:
\begin{equation}
 \mathcal C_\xi(U)=
      \sum_C\sum_{j=1}^N|G_{C,j}-R_C|^2
       +\sum_{p\in\mathcal P_n\setminus\mathcal P_n^{12}}|U_p-I|^2.
 \label{eq:f15v44:C}
\end{equation}
The symbol \(\mathcal C_\xi\) here denotes this scalar
colour-coherence error, not a retained coarse field.
For probability formulas, set \(\mathcal C_\xi=0\)
if some \(\Lambda_C=\pi\), solely as a measurable
extension. No coercivity assertion uses that extension;
all the boundary-tube statements have \(\Lambda_C\le1/4\).

\begin{theorem}[All-configuration rooted holonomy coercivity]
\label{f15v44:full-colour}
If every exterior perimeter obeys \(\Lambda_C\le1/4\),
then, for every compact interior,
\begin{equation}
                   S_\beta(U)-\beta\mathcal L(\xi)
                               \ge\frac{\beta}{5}\mathcal C_\xi(U).
 \label{eq:f15v44:full-colour}
\end{equation}
No small-field, commutativity, or near-minimum restriction
is imposed on the sampled interior.
\end{theorem}

\begin{proof}
Apply the second line of
\eqref{eq:f15v44:small-colour} to each disk via
\eqref{eq:f15v44:disk-energy}. Every other plaquette
has action \(f(\alpha(U_p))=|U_p-I|^2/2\).
Sum all disjoint parallel disks and these remaining
nonnegative terms. Since \(1/2\ge1/5\), the result
is \eqref{eq:f15v44:full-colour}.
\end{proof}

The observable compares full \(\SU(2)\) elements in
their actual rooted frames. A change of colour direction
at a fixed angle can contribute to it. Its rooting paths,
however, depend on the interior links and can have
length comparable to the box side. We have not obtained
a uniformly local link norm or a decay-of-correlations
estimate by naming this observable.

\subsection{The coherent minimum through the principal flux interval}

\begin{theorem}[Exact minimum and boundary-fixed gauge-orbit uniqueness]
\label{f15v44:minimum}
Let \(\overline U^\Theta\) be the V42 coherent field,
with \(\theta=\Theta/n^2\), and freeze its boundary
\(\xi_\theta\). For homogeneous couplings \(J_p=\beta>0\),
\begin{equation}
 0\le\Theta\le\pi
 \quad\Longrightarrow\quad
 s_\beta(\xi_\theta)
        =\beta n^2(n-1)^2 f(\theta)
        =S_\beta(\overline U^\Theta).
 \label{eq:f15v44:minimum}
\end{equation}
If \(0\le\Theta<\pi\), every minimizing configuration
is obtained from \(\overline U^\Theta\) by a vertex
gauge transformation equal to \(I\) on
\(\partial B_n\).
\end{theorem}

\begin{proof}
Each parallel disk has fixed holonomy
\(\exp(i\Theta\sigma_3)\), of principal angle
\(\Theta\) in the stated interval.
Proposition~\ref{f15v44:identity} gives its minimum
energy \(Nf(\Theta/N)\), including at \(\Theta=\pi\).
Every nonparallel action term is nonnegative.
The coherent background attains all these lower
bounds, proving \eqref{eq:f15v44:minimum}.

Assume now \(\Theta<\pi\). At a global minimum,
every parallel disk attains its individual bound
and every nonparallel affected plaquette is \(I\).
The strict part of
Proposition~\ref{f15v44:identity} gives
\[
                         G_{C,j}=R:=\exp(i\theta\sigma_3)
                         \quad\hbox{for every }C,j .
\]
We construct the required boundary-fixed gauge,
rather than identifying configurations from
traces alone.

Fix internal transverse coordinates \((x_3,x_4)\)
and write \((u,v)=(x_1,x_2)\).
Take the left-column path from \((0,0)\) to
\((0,v)\), followed by the horizontal row to
\((u,v)\); denote its transport by \(H_{u,v}\).
The left-column links are fixed and equal to \(I\).
For the interchange ordering in
\eqref{eq:f15v44:disk-path}, the plaquettes occur
row by row in increasing \(v\), and within a row
in decreasing \(u\). Their conjugating prefix is
exactly this left-column/row path. One can verify
this by successively moving the first direction-2
step left across all \(n\) direction-1 steps,
then doing the same with the next direction-2
step. Thus
\begin{equation}
                    H_{u,v}U_{(u,v;1,2)}H_{u,v}^{-1}=R.
 \label{eq:f15v44:row-curvature}
\end{equation}

Gauge-transform this disk by \(H\).
Every horizontal link is \(I\), as are the
left-column vertical links. If \(V_{u,v}\)
denotes a transformed vertical link,
\eqref{eq:f15v44:row-curvature} becomes
\[
                V_{u+1,v}V_{u,v}^{-1}=R,\qquad V_{0,v}=I.
\]
Therefore \(V_{u,v}=R^u\).
The coherent reference has row frames
\(\overline H_{u,v}=R^{-uv}\) and, under these
frames, the same horizontal and vertical links.
Hence
\[
                    k_{u,v}=\overline H_{u,v}^{-1}H_{u,v}
\]
maps the original disk links to their coherent values.

This relative gauge is \(I\) on the entire disk perimeter.
On the left and bottom this follows directly from the
fixed links. On the top, the fixed horizontal links give
\(H_{u,n}=R^{-un}=\overline H_{u,n}\).
On the right, the original fixed vertical links are \(I\),
so \(V_{n,v}=R^n\) gives
\[
 H_{n,v}H_{n,v+1}^{-1}=R^n,\qquad H_{n,0}=I.
\]
It follows recursively that \(H_{n,v}=R^{-nv}\).
Thus \(k=I\) on all four sides.

Perform this construction on every internal parallel
disk, and set \(k=I\) on the slices with
\(x_3\) or \(x_4\) equal to \(0\) or \(n\).
Together these are a well-defined vertex gauge
on \(B_n\), equal to \(I\) on its boundary.
After applying it, every direction-1 and
direction-2 link equals the coherent reference.

It remains to check directions 3 and 4.
Let \(W(x)\) be a transformed direction-3 link.
If a transverse coordinate is on the boundary,
this link is fixed and already \(I\).
Otherwise the flat \((1,3)\) plaquette gives
\[
 W(x+e_1)
       =\overline U_1^\Theta(x)^{-1}
                         W(x)\overline U_1^\Theta(x),
\]
because \(\overline U_1^\Theta\) is independent
of \(x_3\).
At \(x_1=0\), \(W=I\) by the unchanged boundary.
Recursion in \(x_1\) proves \(W=I\) everywhere.
The identical argument using the flat \((1,4)\)
plaquettes proves that every direction-4 link
is \(I\). All links now agree with
\(\overline U^\Theta\), completing the proof.
\end{proof}

V42 proved the exact minimum only for
\(\Theta\le\pi/2\). The present extension to
\(\Theta\le\pi\) follows from the global path
identity, not an unjustified extension of scalar
convexity to larger angles.
At \(\Theta=\pi\), the path coefficient
\(1-\rho_N(\Theta)\) vanishes, so the preceding
strict disk argument no longer proves gauge-orbit
uniqueness. No assertion about uniqueness of the
full four-dimensional boundary problem at that
endpoint is made here.

\subsection{Normalized full-colour tails on the actual boundary tube}

For \(\xi\in\Gamma_n\), V43 supplies both
\(\Lambda_C\le1/4\) and
\[
                    s_\beta(\xi)
                         \le\beta\mathcal L(\xi)+3\beta\mathcal E_n(\xi).
\]
Combining this with
\eqref{eq:f15v44:full-colour} gives the following
full-colour replacement for the angle-only bound:
\begin{equation}
 \begin{split}
 \mathbb P_\beta\{\mathcal C_\xi\ge r\mid\xi\}
   &\le\min\{1,K_\beta^{m_n}
           e^{\,3\beta\mathcal E_n(\xi)-\beta r/5}\},\\
 \mathbb E_\beta[\mathcal C_\xi\mid\xi]
   &\le15\mathcal E_n(\xi)
                  +\frac{5(m_n\log K_\beta+1)}{\beta}.
 \end{split}
 \label{eq:f15v44:homogeneous-tail}
\end{equation}
Here \(\beta\ge1,r\ge0\).
Indeed the conditional excess is at least
\(\beta\{\mathcal C_\xi/5-3\mathcal E_n(\xi)\}\).
Apply
\eqref{eq:f15v40:conditional-tail} to the original
conditional density, then integrate its minimum
with one. This is not integration over
independent plaquette factors.

The same bound gives
\(\mathcal C_\xi(U_\beta^*)\le15\mathcal E_n(\xi)\)
for every minimum. At zero fit cost it recovers
the exact rooted coherence used in
Theorem~\ref{f15v44:minimum}, within the V43 flux range.

\subsection{Full-colour control with arbitrary physical sources}

We now keep the same actual exterior and allow
every permitted array
\((1-\tau)\beta\le J_p\le(1+\tau)\beta\),
\(0\le\tau\le1/32\).
Write \(J_p=\beta(1+v_p)\).
Set
\begin{equation}
 \begin{split}
 V_\xi&=N\sum_C\sin^2\vartheta_C,\\
 \mathfrak c_\tau&=\frac15-\tau\ge\frac{27}{160},\qquad
 \mathfrak b_\xi=2\tau\sqrt{V_\xi},\qquad
 \mathfrak D_\xi=\frac{\mathfrak b_\xi^2}{4\mathfrak c_\tau}.
 \end{split}
 \label{eq:f15v44:source-constants}
\end{equation}
Retain V43's already proved minimum allowance
\begin{equation}
 B_\xi=(3+4\tau)\mathcal E_n(\xi)
        +\sqrt8\,\tau\sqrt{V_\xi}\sqrt{\mathcal E_n(\xi)}.
 \label{eq:f15v44:B}
\end{equation}

\begin{theorem}[Source-uniform rooted-colour response]
\label{f15v44:source}
Whenever \(\Lambda_C\le1/4\) for all \(C\),
every compact interior satisfies
\begin{equation}
 S_J(U)-\mathcal L_J(\xi)
   \ge\beta\{\mathfrak c_\tau\mathcal C_\xi(U)
                      -\mathfrak b_\xi\sqrt{\mathcal C_\xi(U)}\}.
 \label{eq:f15v44:source-coercivity}
\end{equation}
If also \(\xi\in\Gamma_n\), every minimizing
interior obeys
\begin{equation}
 \sqrt{\mathcal C_\xi(U_J^*)}
 \le\frac{\mathfrak b_\xi+
       \sqrt{\mathfrak b_\xi^2+4\mathfrak c_\tau B_\xi}}
                   {2\mathfrak c_\tau}.
 \label{eq:f15v44:source-minimizer}
\end{equation}
For \(\beta\ge1,r\ge0\), the normalized conditional
law on this boundary tube satisfies
\begin{equation}
 \begin{split}
 \mathbb P_J\{\mathcal C_\xi\ge r\mid\xi\}
  &\le\min\{1,K_\beta^{m_n}
       e^{\,\beta(B_\xi+2\mathfrak D_\xi)
                          -\beta\mathfrak c_\tau r/2}\},\\
 \mathbb E_J[\mathcal C_\xi\mid\xi]
  &\le\frac2{\mathfrak c_\tau}
       \left\{B_\xi+2\mathfrak D_\xi+
                    \frac{m_n\log K_\beta+1}{\beta}\right\}.
 \end{split}
 \label{eq:f15v44:source-tail}
\end{equation}
\end{theorem}

\begin{proof}
We first bound the source perturbation in the full
group difference, rather than using an angle-only
Taylor estimate. Let \(G,R\) be unit quaternions,
\(\alpha(R)=\vartheta\le1/16\), and put
\(\delta=G-R\).
Write \(R=(\cos\vartheta,w)\) with
\(|w|=\sin\vartheta\).
The unit-norm identity gives
\[
 2R\cdot\delta+|\delta|^2=0.
\]
Therefore
\begin{equation}
 f(\alpha(G))-f(\vartheta)
       =-\operatorname{Re}\delta
       =\frac{|\delta|^2/2+
                         w\cdot\operatorname{Im}\delta}{\cos\vartheta}.
 \label{eq:f15v44:source-identity}
\end{equation}
Since \(\cos\vartheta\ge1/2\),
\begin{equation}
 |f(\alpha(G))-f(\vartheta)|
                  \le|\delta|^2+2\sin\vartheta\,|\delta|.
 \label{eq:f15v44:source-difference}
\end{equation}
Both identities are valid for arbitrary \(G\)
in the compact group, not only close to \(R\).

Split \(\mathcal C_\xi=Z_{\parallel}+Z_{\perp}\)
according to its two sums in
\eqref{eq:f15v44:C}. The homogeneous disk excess
is at least \(Z_{\parallel}/5\), and the
nonparallel action is exactly \(Z_{\perp}/2\).
For parallel face \(p_j\), its source perturbation
is \(v_{p_j}\{f(\alpha(G_{C,j}))-f(\vartheta_C)\}\).
Equation~\eqref{eq:f15v44:source-difference}
and Cauchy--Schwarz bound the sum of these
terms from below by
\[
                   -\tau Z_{\parallel}
                   -2\tau\sqrt{V_\xi}\sqrt{Z_{\parallel}}.
\]
The other source terms are at least
\(-\tau Z_{\perp}/2\).
Consequently
\[
 \frac{S_J-\mathcal L_J}{\beta}
 \ge(1/5-\tau)Z_{\parallel}
       +\frac{1-\tau}{2}Z_{\perp}
       -\mathfrak b_\xi\sqrt{Z_{\parallel}}.
\]
Since \((1-\tau)/2\ge\mathfrak c_\tau>0\)
and \(Z_{\parallel}\le\mathcal C_\xi\),
this proves
\eqref{eq:f15v44:source-coercivity}.
No commutativity or monotonicity in individual
plaquette couplings is used.

On \(\Gamma_n\), V43 gives the actual minimum bound
\(s_J\le\mathcal L_J+\beta B_\xi\).
Apply the just-proved inequality at a minimum
and solve the quadratic in
\(\sqrt{\mathcal C_\xi}\) to get
\eqref{eq:f15v44:source-minimizer}.
For general interiors,
\[
 \mathfrak c_\tau t^2-\mathfrak b_\xi t
         \ge\frac12\mathfrak c_\tau t^2-2\mathfrak D_\xi .
\]
Thus
\[
 S_J-s_J\ge
       \beta\{\mathfrak c_\tau\mathcal C_\xi/2
                            -2\mathfrak D_\xi-B_\xi\}.
\]
Use the inherited conditional-excess estimate,
with the trivial bound one when the threshold
is nonpositive. Integration of the resulting
tail proves the mean bound.
\end{proof}

The source floor is not being renamed a thermal
variance. In particular an inhomogeneous source
can shift the minimizer away from the coherent
rooted factors even when the boundary fit cost
is zero. The theorem bounds that displacement
with the same boundary held fixed.

\subsection{Actual-law integration and the still unpaid marginal}

The new observable is a function of the original
links in its box and its exterior perimeter
holonomies. Hence it fits the V43 conditional-product
argument without changing the Wilson measure.
For \(s\ge0\), put
\begin{equation}
 r^{\rm col}_\xi(s)=\frac2{\mathfrak c_\tau}
       \left\{B_\xi+2\mathfrak D_\xi+
                     \frac{m_n\log K_\beta+s}{\beta}\right\}.
 \label{eq:f15v44:radius}
\end{equation}
For pairwise vertex-disjoint boxes and any
nonnegative integrable weight \(H\) measurable
in their common exterior,
\begin{equation}
 \mathbb E_J\!\left[
     H\prod_{i=1}^k\mathbf1_{\Gamma_i}
             \mathbf1_{\{\mathcal C_i\ge r^{\rm col}_{\xi_i}(s_i)\}}
             \right]
 \le e^{-\sum_i s_i}\,
                   \mathbb E_J\!\left[H\prod_{i=1}^k\mathbf1_{\Gamma_i}\right].
 \label{eq:f15v44:joint}
\end{equation}
Condition on the common exterior. The affected
actions have disjoint supports, so their exact
conditional laws factorize. Apply
\eqref{eq:f15v44:source-tail} in each box;
the normalizer factors cancel against
the \(m_n\log K_\beta\) term in the radius.
Then integrate. No independence of the exterior
or the rooted plaquette factors is asserted.

This is a full-law statement with an explicit
boundary restriction, not the missing estimate
for the probability of that restriction.
In particular it gives no bound on
\(\mu_J(\cap_i\Gamma_i^c)\), nor does it make
a boundary-forced nonzero curvature profile
rare. The homogeneous selected CGF still
cannot be inserted at an arbitrary physical
source background without paying a justified
change-of-law cost.

The new step is that noncommuting rooted
holonomies, rather than only conjugacy-class
angles, now have a globally proved energy
coercivity and normalized conditional tail.
The next marginal comparison must preserve
that rooting information or give a justified
local replacement for it. A source-uniform
joint estimate for actual retained backgrounds
remains the main unmet probabilistic input.

\subsection{Exact diagnostics and preservation}

The independent diagnostic is
\begin{center}\small
\path{scripts/verify_ym_f15_v44_rooted_colour_rigidity.py}.
\end{center}
It checks 879 rational entries of the explicit
Dirichlet inverse
\[
                (D_N^{-1})_{ij}=\min(i,j)-ij/N
                  \quad(1\le i,j\le N-1),
\]
including its row sums \(i(N-i)/2\).
These are finite checks of the path operator,
not a substitute for the all-\(N\) spectral
proof above.

Eighteen rational compact paths check the exact
chord identity, full-factor colour bound and
partial-transport bound, including noncommuting
paths far from the reference and flat references.
Twenty-two ordered lattice disk products check
the nonabelian rooting. Eighty minimizing
rooted plaquettes, 400 flat transverse
plaquettes and 1,968 boundary-fixed link
recoveries test the explicit gauge-orbit
reconstruction.

Four noncommuting lattice perturbations check
the full-action colour inequality. Twelve
square completions for inhomogeneous sources and
240 exact source perturbation identities
check the new source allowance. Eight
root-gauge covariance identities and two
full-observable gauge-invariance tests
complete the diagnostics.
The global statements and probability bounds
are proved analytically; they are not inferred
from these samples.

The script has 9,636 bytes and SHA--256
\begin{center}\scriptsize\ttfamily
6740B42D60C473E962F15B6B342902C7AA8D3724A91714AF18996551C0F2C7EF.
\end{center}
The V43 predecessor has 1,010,830 bytes and SHA--256
\begin{center}\scriptsize\ttfamily
3A48FA07E61168AFE9C7F67219863C42B613328A89F3A4D2766F6B375E894E8A.
\end{center}
Removing this final section and reversing
the title version recovers V43 exactly.
V44 replaces V43 as the sole active main
source. All 14 protected sources and
86 earlier certificates are preserved.
The next iteration, V45, is due for
the companion checkpoint.

The source folder \path{latex_notes}
continues to contain only \path{.tex}
files. Build output belongs in
\begin{center}\small
\path{artifacts/ym_f15_v44_texcheck/}.
\end{center}

\begin{firewall}
V44 proves an exact compact disk-path identity,
global rooted colour coercivity, an extended
coherent minimum formula, boundary-fixed
gauge-orbit uniqueness below antipodal flux,
and source-uniform conditional full-colour
tails on the V43 boundary tube.
It does not prove the marginal retained-background
estimate, a uniform local gauge-coordinate
comparison, physical clustering, iterated RG,
the interacting continuum limit, or the
physical Yang--Mills mass gap.
The objective remains incomplete.
\end{firewall}
\section{V45: forest reduction and the normalized boundary factor}
\label{f15v45:section}

V44 proves compact rooted colour rigidity, but its rooted
observable is not a product of independent group variables.
This iteration converts that estimate to a global distance
in actual boundary-fixed forest coordinates. The forest
change of variables has exactly unit product-Haar Jacobian.
After integrating the gauge variables, a quadratic estimate
on the entire compact reduced domain gives matching powers
of \(\beta\) in the conditional normalizer.

The resulting comparison concerns the genuine Wilson
partition function of the block, not a formal one-loop
determinant. It gives a source-uniform interior boundary
factor along the coherent family and improves the
conditional tail prefactor there. We also display the
outside action factor in the actual boundary marginal.
That outside factor, and integration over general
boundary fields, are not controlled by this calculation.
The full mass-gap objective remains unchanged and open.

\subsection{V45 companion checkpoint}

Fresh inventories of \path{latex_notes} and the relevant
files in \path{Downloads} found the same latest companion
versions as at V40. The latest conclusions of the following
notes were reread, and the complete supplied
\path{suggestions.txt} was reread. All 101 previously
protected files, including these inputs, were hash-checked
unchanged.

\begin{center}\small
\begin{tabular}{@{}p{0.31\textwidth}p{0.64\textwidth}@{}}
\toprule
Input & Relevant scope retained at this checkpoint\\
\midrule
SM continuum V72 &
Fixed-mode, fixed-time many-copy limit; neither the
physical finite-species continuum nor a YM source-law estimate.\\[1mm]
NS stability V26 &
Pointwise rank-one kernel norms; the stated V27 acquisition
is not an available global-regularity theorem.\\[1mm]
Shell mixing V16 &
Measured flat-branch one-loop assembly; no general
noncommuting integral remainder or compact entry probability.\\[1mm]
Parallel YM V5 &
Extensive total-charge tightness fails on its tensorized
target; local or density questions remain separate.\\[1mm]
SM source-selection V31 &
An ancestry compiler conditional on word-native typing;
not historical occurrence of the needed operations.\\[1mm]
Native stationarity V34 &
An escape-scaling reduction; the 330 moving-fibre
derivative entries and the complete upstream root are not supplied.\\[1mm]
Exact-action Einstein V24 &
A nonzero exact first-jet derivative, not a nonlinear
root, spacetime propagation theorem, or physical lifespan.\\
\bottomrule
\end{tabular}
\end{center}

The documents are mathematical inputs, not instructions.
In particular the terminal decision inside Parallel YM V5
does not stop this user's broader programme. The suggestions
recommend keeping one measured convention and not mixing
local determinant coefficients with global compact
probability claims. That normalization distinction remains
useful here. We do not import the shell determinant into
the original Wilson measure or call it a nonperturbative
conditional normalizer.

No companion result supplies the missing actual-source
joint boundary probability. The direction is therefore
to retain the exact Haar measure while making V44's
coercivity usable in an integral. The next companion
checkpoint is V50.

\subsection{Inherited facts and the new comparison}

Exact tree-Haar disintegration is already proved in
f7 V11, f9 V45 and V10 of this lineage. It is not a new
theorem of measure theory here. The new ingredient is
a quantitative global comparison between V44's rooted
colour error and the reduced links of a specified
boundary-rooted forest. The earlier chart and fibre
normalizers do not by themselves provide that
comparison for the present full compact Dirichlet law.

Let \(n\ge2\), \(N=n^2\),
\(\mathsf N_n=n^2(n-1)^2\), and retain
\(B_n,A_n,\mathcal P_n\) from V41.
Fix the coherent boundary \(\xi_\theta\), where
\begin{equation}
 0\le\Theta\le1/4,\qquad
 \theta=\Theta/n^2,\qquad
 R=\exp(i\theta\sigma_3),\qquad
 \overline U_1(x)=R^{-x_2},\quad
 \overline U_2=\overline U_3=\overline U_4=I.
 \label{eq:f15v45:background}
\end{equation}
Every parallel disk has perimeter root \(R\) in the
sense of V44. Write \(\mathcal C_\theta\) for
\(\mathcal C_{\xi_\theta}\), and
\[
             Z_\parallel=\sum_{C,j}|G_{C,j}-R|^2,\qquad
             Z_\perp=\sum_{p\notin\mathcal P_n^{12}}|U_p-I|^2 ,
 \qquad \mathcal C_\theta=Z_\parallel+Z_\perp ,
\]
where the second sum is over affected plaquettes.
These definitions do not assume small interior links.

\subsection{An exact boundary-rooted forest}

Take the forest
\begin{equation}
 \mathcal T_n
   =\{(x,1)\in A_n:0\le x_1\le n-2\}.
 \label{eq:f15v45:forest}
\end{equation}
Each component is a direction-1 line starting at a
boundary vertex with \(x_1=0\). It reaches the
\(n-1\) interior vertices on that line and stops
before the right boundary. All three transverse
coordinates of its edges are internal.
There is exactly one forest edge for each interior
vertex. Thus
\begin{equation}
 \begin{split}
 m_n&=4n(n-1)^3,\qquad v_n=|\mathcal T_n|=(n-1)^4,\\
 E_n^{\rm red}&=A_n\setminus\mathcal T_n,\qquad
 r_n=|E_n^{\rm red}|=m_n-v_n=(3n+1)(n-1)^3.
 \end{split}
 \label{eq:f15v45:counts}
\end{equation}
There are \((n-1)^3\) direction-1 seam links in
\(E_n^{\rm red}\), and \(n(n-1)^3\) reduced
links in each of directions 2, 3 and 4.

Set \(g=I\) at every boundary vertex and recursively
on a forest edge \(e=(x,1)\) set
\begin{equation}
                g_{x+e_1}=\overline U_e^{-1}g_xU_e.
 \label{eq:f15v45:frame}
\end{equation}
Then \(z_e=g_xU_eg_{x+e_\mu}^{-1}\) equals
\(\overline U_e\) on every forest edge.
The remaining \(z_e\), \(e\in E_n^{\rm red}\),
are arbitrary compact group variables.
This construction is a global bijection
\[
 \SU(2)^{A_n}
       \longleftrightarrow
          \SU(2)^{V_{\rm int}(B_n)}
                         \times\SU(2)^{E_n^{\rm red}}.
\]
Indeed \eqref{eq:f15v45:frame} determines each new
interior \(g\) uniquely, and the inverse on every
updated edge is \(U_e=g_x^{-1}z_eg_{x+e_\mu}\),
with the forest \(z_e\)'s fixed as above.

Normalized product Haar measure therefore satisfies
\begin{equation}
       dH_{A_n}(U)=dH_{V_{\rm int}}(g)\,
                                  dH_{E_n^{\rm red}}(z).
 \label{eq:f15v45:haar}
\end{equation}
To verify the measure statement directly, traverse
each forest component from its boundary root.
Replacing its next edge by the new vertex gauge
in \eqref{eq:f15v45:frame} is a fixed left/right
Haar translation once previous gauges are fixed.
After these replacements, each remaining edge
undergoes a left/right Haar translation as well.
Fubini proves the identity. The nonidentity
reference tree values do not change Haar measure.

The gauge is identity on the boundary, so it
leaves the conditional action and the fixed
exterior unchanged. For arbitrary positive
plaquette couplings the gauge variables integrate
to exactly one. There is no orbit-volume estimate,
gauge determinant, or extra density to insert
into these group-valued forest coordinates.

\subsection{A global reduced-link estimate}

On the reduced domain define
\begin{equation}
                  D_{\mathcal T}(z)
                     =\sum_{e\in E_n^{\rm red}}
                                        |z_e-\overline U_e|^2 .
 \label{eq:f15v45:distance}
\end{equation}
This is a coordinate distance on a specific
boundary-fixed gauge slice; it is not asserted
to be a uniformly local invariant norm on
arbitrary boxes.

\begin{theorem}[Rooted colour controls every reduced link]
\label{f15v45:link}
For every compact interior with the coherent
boundary \eqref{eq:f15v45:background},
\begin{equation}
                       D_{\mathcal T}(z)\le n^3\mathcal C_\theta(U).
 \label{eq:f15v45:link}
\end{equation}
More precisely, if \(D_\mu\) is the contribution
of reduced direction-\(\mu\) links, then
\begin{equation}
 D_1\le\frac{n^3}{2}Z_\parallel,\qquad
 D_2\le\frac{n^2}{2}Z_\parallel,\qquad
 D_3+D_4\le\frac{n^2}{2}Z_\perp .
 \label{eq:f15v45:directional}
\end{equation}
\end{theorem}

\begin{proof}
Compute all rooted factors after the forest gauge.
Their roots are boundary vertices, where \(g=I\),
so the factors themselves, not only their norms,
are unchanged.

Fix an internal parallel disk and write its
coordinates as \((u,v)=(x_1,x_2)\).
Let \(H_{u,v}\) be its left-column/row frame
from V44. For \(0\le u\le n-1\),
the horizontal edges in that prefix are forest
edges, or fixed coherent edges on the top or
bottom. Hence
\[
                 H_{u,v}=\overline H_{u,v}=R^{-uv}
                        \quad(0\le u\le n-1).
\]
The frame \(H_{n,v}\) also contains the possibly
nontrivial right seam edge. Put
\[
            V_{u,v}=H_{u,v}z_{(u,v;2)}H_{u,v+1}^{-1}.
\]
The full row-tree gauge has horizontal links \(I\).
Thus its rooted face equation is
\[
       V_{u+1,v}V_{u,v}^{-1}=G_{u,v},\qquad V_{0,v}=I,
 \qquad
       V_{u,v}=G_{u-1,v}\cdots G_{0,v}.
\]
Let \(\delta_{u,v}=|G_{u,v}-R|\).
Product telescoping gives
\[
 |V_{u,v}-R^u|\le\sum_{k<u}\delta_{k,v}.
\]
For \(1\le u\le n-1\), the left side equals
\(|z_{(u,v;2)}-I|\), since both row frames there
are their reference values. Squaring and summing
over \(u\) gives at most
\(\frac{n(n-1)}2\sum_k\delta_{k,v}^2\).
Summing rows and disks proves the bound for \(D_2\).

On the right boundary the direction-2 link is
fixed to \(I\), so
\[
 H_{n,v}H_{n,v+1}^{-1}=V_{n,v},\qquad H_{n,0}=I.
\]
Recursive product comparison with \(R^{-nv}\)
gives
\[
 |H_{n,v}-R^{-nv}|
        \le\sum_{w<v}|V_{n,w}-R^n|
        \le\sum_{w<v}\sum_{k=0}^{n-1}\delta_{k,w}.
\]
The left side equals the norm difference of
the seam edge from its reference value.
Its square is at most
\(nv\sum_{w<v,k}\delta_{k,w}^2\).
Sum \(v=1,\ldots,n-1\); the coefficient is
at most \(n^3/2\). This proves the \(D_1\) bound.

For direction \(\mu=3\) or 4, start at \(x_1=0\),
where the link is the fixed identity.
For \(x_1=0,\ldots,n-2\), both direction-1
edges of a \((1,\mu)\) plaquette have their
reference values, and these reference values
are independent of \(x_\mu\).
The exact plaquette equation therefore implies
\[
 |z_{(x+e_1,\mu)}-I|
          \le |U_{(x;1,\mu)}-I|+|z_{(x,\mu)}-I|.
\]
The plaquette on the right may be evaluated in
the forest gauge; its distance from \(I\) is
gauge invariant. Iterate in \(x_1\), square
and sum as for \(D_2\). This gives \(n^2/2\)
times the sum of squared \((1,3)\) and
\((1,4)\) plaquette distances used by these
rows. They are distinct affected plaquettes,
so the sum is bounded by \(Z_\perp\).
Boundary links not requiring this recursion
are already fixed.

Finally
\[
 D_{\mathcal T}\le\frac{n^3+n^2}{2}Z_\parallel
                           +\frac{n^2}{2}Z_\perp
                  \le n^3(Z_\parallel+Z_\perp).
\]
No step used a small-field expansion.
\end{proof}

For homogeneous couplings, V44 consequently gives
the full compact reduced coercivity
\begin{equation}
 S_\beta(z)-\beta\mathsf N_n f(\theta)
      \ge\frac{\beta}{5}\mathcal C_\theta
      \ge\frac{\beta}{5n^3}D_{\mathcal T}(z).
 \label{eq:f15v45:coercivity}
\end{equation}
The \(n^3\) loss is explicit. Nothing here makes
the forest rooting uniformly short as \(n\) grows.

\subsection{The exact partition function and matching Haar powers}

Define the interior conditional partition function
with its original normalized Haar reference:
\begin{equation}
 \mathcal Z_J(\theta)
     =\int_{\SU(2)^{A_n}}
              e^{-S_J(U;\xi_\theta)}\,dH_{A_n}(U)
     =\int_{\SU(2)^{E_n^{\rm red}}}
              e^{-S_J(z;\xi_\theta)}\,dH_{E_n^{\rm red}}(z).
 \label{eq:f15v45:Z}
\end{equation}
The equality is exact by
\eqref{eq:f15v45:haar}.
It is not the partition function of a
restricted chart.

Two elementary estimates will be used.
For \(\alpha_e=\alpha(\overline U_e^{-1}z_e)\),
\begin{equation}
                  |z_e-\overline U_e|^2
                     =4\sin^2(\alpha_e/2)
                     \ge\frac4{\pi^2}\alpha_e^2 .
 \label{eq:f15v45:chord}
\end{equation}
Also, for every \(t>0\), normalized Haar measure
on the unit three-sphere gives
\begin{equation}
 \begin{split}
 I(t)&=\int_{\SU(2)}e^{-t\alpha(g)^2}\,dH(g)
       =\frac2\pi\int_0^\pi e^{-tu^2}\sin^2u\,du\\
     &\le\frac2\pi\int_0^\infty e^{-tu^2}u^2\,du
       =\frac1{2\sqrt\pi}\,t^{-3/2}.
 \end{split}
 \label{eq:f15v45:I}
\end{equation}
Left translation by \(\overline U_e\) does not
change this integral.

For a lower bound, the actual minimum on
\(\SU(2)^{E_n^{\rm red}}\) exists. Along product
geodesics about a minimum, the first derivative
of the action vanishes and the global Wilson
second derivative obeys
\[
 |S_J''|\le24(1+\tau)\beta
                     \sum_{e\in E_n^{\rm red}}|X_e|^2.
\]
These are geodesics in the physical reduced
links with the forest links fixed; no derivative
of a moving gauge is being omitted.
The V40 product-Haar-ball proof therefore
applies with \(r_n\), rather than \(m_n\),
group variables:
\begin{equation}
              \mathcal Z_J(\theta)\ge
                     e^{-s_J(\xi_\theta)}
                              (2^{13}\beta^{3/2})^{-r_n}.
 \label{eq:f15v45:lower}
\end{equation}
This uses \(\beta\ge1\) and \(\tau\le1/32\).
The radius \((32\beta)^{-1/2}\), the
second-derivative allowance \(99r_n/256\),
and the normalized Haar-ball estimate are
unchanged from V40. Only the exactly integrated
gauge factors have been removed.

\begin{theorem}[Full compact homogeneous normalizer]
\label{f15v45:homogeneous-Z}
For \(\beta\ge1\) and the coherent range
\eqref{eq:f15v45:background},
\begin{equation}
 \begin{split}
 &e^{-\beta\mathsf N_n f(\theta)}
            (2^{-13}\beta^{-3/2})^{r_n}
       \le\mathcal Z_\beta(\theta)\\
 &\hspace{22mm}\le
 e^{-\beta\mathsf N_n f(\theta)}
            (2^6 n^{9/2}\beta^{-3/2})^{r_n}.
 \end{split}
 \label{eq:f15v45:homogeneous-Z}
\end{equation}
Equivalently,
\begin{equation}
 -\log\mathcal Z_\beta(\theta)
   =\beta\mathsf N_n f(\theta)+\frac{3r_n}{2}\log\beta
                                           +\epsilon_{n,\beta,\theta},
 \quad
 -r_n\log(2^6n^{9/2})\le\epsilon_{n,\beta,\theta}
                              \le13r_n\log2.
 \label{eq:f15v45:free-energy}
\end{equation}
\end{theorem}

\begin{proof}
The exact minimum is
\(\beta\mathsf N_n f(\theta)\) by V44,
so \eqref{eq:f15v45:lower} gives the lower
bound. Equations
\eqref{eq:f15v45:coercivity} and
\eqref{eq:f15v45:chord}, with \(\pi<4\),
give
\[
 S_\beta-s_\beta
        \ge\frac{\beta}{20n^3}\sum_e\alpha_e^2.
\]
The entire reduced Haar integral is therefore
bounded by \(I(\beta/(20n^3))^{r_n}\).
Equation~\eqref{eq:f15v45:I} and
\(20^{3/2}/(2\sqrt\pi)<2^6\)
prove the upper bound. Taking logarithms
gives \eqref{eq:f15v45:free-energy}.
\end{proof}

Thus the \(\beta\)-power uses exactly the
\(3r_n\) real reduced directions. The constants
still grow with \(n\); the theorem is not a
uniform continuum determinant formula.
It nevertheless proves matching Haar powers
from the full compact integral, without
discarding its large-field region.

\subsection{A source-uniform partition envelope}

Allow all arrays
\((1-\tau)\beta\le J_p\le(1+\tau)\beta\),
\(0\le\tau\le1/32\). Set
\begin{equation}
 \begin{split}
 \overline S_{J,\theta}
       &=f(\theta)\sum_{p\in\mathcal P_n^{12}}J_p,\\
 c_\tau&=\frac15-\tau\ge\frac{27}{160},\qquad
 \Delta_{\tau,\theta}
       =\frac{\tau^2\mathsf N_n\sin^2\theta}{c_\tau}.
 \end{split}
 \label{eq:f15v45:source-data}
\end{equation}
The reference field has action
\(\overline S_{J,\theta}\); it need not be
the actual source-dependent minimum.
V44's full-colour source estimate and
Young's inequality give globally
\begin{equation}
 S_J-\overline S_{J,\theta}
   \ge\beta\left\{\frac{c_\tau}{2}\mathcal C_\theta
                                    -2\Delta_{\tau,\theta}\right\}
   \ge\beta\left\{\frac{c_\tau}{2n^3}D_{\mathcal T}
                                    -2\Delta_{\tau,\theta}\right\}.
 \label{eq:f15v45:source-coercivity}
\end{equation}
Here the linear coefficient before Young's
inequality is
\(2\tau\sqrt{\mathsf N_n}\sin\theta\),
so its square divided by \(4c_\tau\)
is exactly \(\Delta_{\tau,\theta}\).

\begin{theorem}[Exact-source interior partition envelope]
\label{f15v45:source-Z}
Under these assumptions,
\begin{equation}
 \begin{split}
 &e^{-\overline S_{J,\theta}}
              (2^{-13}\beta^{-3/2})^{r_n}
      \le\mathcal Z_J(\theta)\\
 &\hspace{17mm}\le
 e^{-\overline S_{J,\theta}+2\beta\Delta_{\tau,\theta}}
              (2^8n^{9/2}\beta^{-3/2})^{r_n}.
 \end{split}
 \label{eq:f15v45:source-Z}
\end{equation}
\end{theorem}

\begin{proof}
Since the coherent field is an admissible
comparison, \(s_J(\xi_\theta)\le
\overline S_{J,\theta}\).
Equation~\eqref{eq:f15v45:lower} gives the
lower bound.
For the upper bound,
\eqref{eq:f15v45:source-coercivity} and
\eqref{eq:f15v45:chord} give a Gaussian
coefficient at least
\[
 \frac{2c_\tau}{\pi^2n^3}
       \ge\frac{c_\tau}{8n^3}
       \ge\frac1{64n^3}.
\]
The product integral is therefore at most
\[
 I(\beta/(64n^3))^{r_n}
       \le(2^8n^{9/2}\beta^{-3/2})^{r_n},
\]
using \(\sqrt\pi>1\).
The remaining exponential factor is precisely
the allowance in
\eqref{eq:f15v45:source-coercivity}.
\end{proof}

The source factor \(e^{2\beta\Delta_{\tau,\theta}}\)
is explicit and cannot be hidden in a
\(\beta\)-independent constant. Unlike the
homogeneous estimate, this envelope does not
claim the exact leading classical action
of the perturbed minimum. The sources
and the boundary are the original physical
ones throughout.

\subsection{A coherent-boundary partition ratio with fixed-strength sources}

The same coupling array is now used at
\(\xi_\theta\) and at the flat boundary
\(\xi_0\). Define the numerical constant
\[
                         \mathsf A_n=2^{21}n^{9/2}.
\]

\begin{theorem}[Interior boundary factor under the actual source]
\label{f15v45:ratio}
Uniformly over all permitted source arrays,
\begin{equation}
 \mathsf A_n^{-r_n}
     e^{-(1+\tau)\beta\mathsf N_n f(\theta)}
 \le\frac{\mathcal Z_J(\theta)}{\mathcal Z_J(0)}
 \le\mathsf A_n^{r_n}
     e^{-(15/16)\beta\mathsf N_n f(\theta)} .
 \label{eq:f15v45:ratio}
\end{equation}
There is no remaining power of \(\beta\)
in the prefactor of this ratio.
\end{theorem}

\begin{proof}
At the flat boundary,
\(\overline S_{J,0}=\Delta_{\tau,0}=0\).
Divide the lower bound in
\eqref{eq:f15v45:source-Z} at \(\theta\)
by its upper bound at zero, and conversely.
The factors \(\beta^{-3r_n/2}\) cancel,
and \(2^{13+8}=2^{21}\).
The lower bound then follows from
\(\overline S_{J,\theta}\le
(1+\tau)\beta\mathsf N_n f(\theta)\).

For the upper bound use
\(\sin^2\theta=(1+\cos\theta)f(\theta)\le2f(\theta)\).
It gives
\[
 \overline S_{J,\theta}-2\beta\Delta_{\tau,\theta}
 \ge\beta\mathsf N_n f(\theta)
                      \left(1-\tau-\frac{4\tau^2}{c_\tau}\right).
\]
The coefficient satisfies
\[
 1-\tau-\frac{4\tau^2}{c_\tau}
       \ge\frac{31}{32}-\frac5{216}
       =\frac{15}{16}+\frac7{864}
       \ge\frac{15}{16}.
\]
This uses only \(\tau\le1/32\) and
\(c_\tau\ge27/160\), not a source-dependent
change of boundary or a correlation inequality.
\end{proof}

This is a nonperturbative interior partition
factor at fixed nonzero source strength.
It compares pointwise coherent boundary
data, not the probabilities of individual
boundary configurations, which are zero
under a continuous Haar law.

\subsection{Conditional full-colour tails without a logarithmic coupling loss}

The same integration argument improves
the V44 tail on this coherent family.
Put
\[
                         \widetilde{\mathsf A}_n=2^{23}n^{9/2}.
\]
For \(u\ge0\),
\begin{equation}
 \begin{split}
 \mathbb P_J\{\mathcal C_\theta\ge u\mid\xi_\theta\}
  &\le\min\left\{1,\,
       \widetilde{\mathsf A}_n^{r_n}
          e^{\,2\beta\Delta_{\tau,\theta}-\beta c_\tau u/4}\right\},\\
 \mathbb E_J[\mathcal C_\theta\mid\xi_\theta]
  &\le\frac{8\Delta_{\tau,\theta}}{c_\tau}
       +\frac{4\{r_n\log\widetilde{\mathsf A}_n+1\}}
                          {\beta c_\tau}.
 \end{split}
 \label{eq:f15v45:tail}
\end{equation}
To prove this, split the positive term in
\eqref{eq:f15v45:source-coercivity} into
two equal parts. On the specified event,
\[
 S_J-\overline S_{J,\theta}
       \ge\beta c_\tau u/4
          +\frac{\beta c_\tau}{4n^3}D_{\mathcal T}
          -2\beta\Delta_{\tau,\theta}.
\]
After \eqref{eq:f15v45:chord}, the remaining
Gaussian coefficient is at least
\[
                   \frac{c_\tau}{\pi^2n^3}
                            \ge\frac1{128n^3}.
\]
Its full reduced Haar integral is at most
\((2^{10}n^{9/2}\beta^{-3/2})^{r_n}\)
by \eqref{eq:f15v45:I}.
Divide by the lower bound in
\eqref{eq:f15v45:source-Z}.
The Haar powers cancel, leaving
\(2^{13+10}=2^{23}\).
Integration of the tail proves the mean bound.

For fixed \(n\), its thermal term is \(O(\beta^{-1})\),
not \(O(\beta^{-1}\log\beta)\).
The explicit source-response term remains.
The improvement comes from integrating the
entire reduced numerator and denominator;
bounding the numerator by one would retain
the older logarithmic coupling loss.
Equation~\eqref{eq:f15v45:tail} is presently
proved for the coherent boundary family.
It must not silently replace V44's more
general tube estimate at arbitrary
\(\xi\in\Gamma_n\).

\subsection{Where the actual boundary marginal still differs}

Let \(F_n\) be V43's fixed edges which
meet an affected plaquette. The interior
partition function depends on the exterior
only through these edges. Condition now
on all links \(Y\) outside \(A_n\cup F_n\),
and leave \(F=U_{F_n}\) free.
Split the full Wilson action into
\[
 S_{\rm full}(U_{A_n},F,Y)
       =S_J(U_{A_n};F)+S_{\rm rest}(F,Y)+S_{\rm far}(Y),
\]
where \(S_{\rm rest}\) contains all terms
depending on \(F\) but not on any link of
\(A_n\). Integration of the unchanged
product-Haar law gives the exact
conditional boundary density
\begin{equation}
 p_Y(F)=
 \frac{e^{-S_{\rm rest}(F,Y)}\mathcal Z_J(F)}
      {\displaystyle\int e^{-S_{\rm rest}(F',Y)}
                            \mathcal Z_J(F')\,dH_{F_n}(F')}.
 \label{eq:f15v45:boundary-density}
\end{equation}
There is no missing Jacobian in this
original-link disintegration.

Let \(F_\theta,F_0\) be the coherent
and flat assignments on \(F_n\).
The continuous, strictly positive density
has the pointwise ratio
\begin{equation}
 \frac{p_Y(F_\theta)}{p_Y(F_0)}
    =e^{-\{S_{\rm rest}(F_\theta,Y)-S_{\rm rest}(F_0,Y)\}}
                         \frac{\mathcal Z_J(\theta)}{\mathcal Z_J(0)}.
 \label{eq:f15v45:actual-ratio}
\end{equation}
Theorem~\ref{f15v45:ratio} controls the
second factor under the same actual
source. It gives no favorable sign
or probabilistic bound for the first.
The formula remains valid with any
fixed \(Y\): edges outside \(F_n\)
do not enter the interior partition
function, even if their values do
not extend the coherent field.

To obtain a boundary-event probability,
one must also control neighborhoods
away from the exact coherent family,
their Haar volume, and their joint
outside-action contribution. No such
estimate is inferred from evaluating
a density on a lower-dimensional
family. In particular,
\eqref{eq:f15v45:ratio} is not yet a
bound for \(\mu_J(\Gamma_n^c)\).

The progress is a measured interior
factor with matching gauge-reduced
Haar powers and fixed-strength source
control. The unpaid quantity has
not been hidden in a normalization
constant: it is displayed explicitly
in \eqref{eq:f15v45:boundary-density}.

\subsection{Exact checks and preservation}

The independent diagnostic is
\begin{center}\small
\path{scripts/verify_ym_f15_v45_forest_normalization.py}.
\end{center}
Twelve compact field configurations
check the exact forest counts and
5,904 coordinate inversions, including
four zero-cost gauge orbits.
There are 240 noncommuting row
reconstruction identities, 36
direction-resolved bounds, 12 global
reduced-link coercivity checks and
12 full-action gauge-invariance checks.
The fields include large compact
perturbations, not just infinitesimal
ones.

Thirty-six source quadratic envelopes,
31 side/count identities and five
rational source/Gaussian constant
margins supplement those checks.
The Haar disintegration and integral
estimates are proved analytically;
the script does not estimate a
partition function by sampling or
claim a probabilistic mass-gap certificate.

The script has 7,890 bytes and SHA--256
\begin{center}\scriptsize\ttfamily
7AB58D23DD78FFC5DC7705B3ADE530FB754404568618CFB0C36120B1AA2179C0.
\end{center}
The V44 predecessor has 1,037,278 bytes and SHA--256
\begin{center}\scriptsize\ttfamily
8C1AB26D416B76830D03880D714B726269371E1DABECB55F1A0664AB593DC439.
\end{center}
Removing this final section and
reversing the title version recovers
V44 exactly. V45 replaces V44 as
the sole active main source.
All 14 protected sources and 87
earlier certificates are preserved.
The V45 companion review is complete;
the next is V50.

Only \path{.tex} sources are kept
in \path{latex_notes}. Compilation
output for this version belongs in
\begin{center}\small
\path{artifacts/ym_f15_v45_texcheck/}.
\end{center}

\begin{firewall}
V45 proves a global boundary-fixed
forest link estimate, matching full
compact homogeneous Haar powers,
a source-uniform interior boundary
partition ratio, and improved
conditional colour tails along
the coherent family.
It does not control the full
boundary marginal, general
source-uniform compact entry,
physical clustering, iterated RG,
the interacting continuum, or
the physical Yang--Mills mass gap.
The goal remains unproved.
\end{firewall}
\section{V46: thermal boundary neighborhoods and the paid exterior force}
\label{f15v46:section}

V45 compares the interior partition function at exact
coherent boundary configurations. Such configurations
have zero Haar probability. We now extend the compact
integral estimate to genuinely noncommuting boundary
neighborhoods, and integrate their actual conditional
boundary densities. Both the boundary Haar volume and
the interior gauge-reduced powers of \(\beta\) are retained.
The outside force is not dropped: its cost is first
displayed, then bounded by an explicit outside action.

This is not yet a typical-background theorem. The
remaining quantity involves replacing the sampled
boundary by a flat one. Its exponential moment under
the actual exterior marginal is not supplied by an
estimate of the sampled plaquette energy.
The continuum and physical mass-gap objective remains
unproved. The official
\href{https://www.claymath.org/millennium/yang-mills-the-maths-gap/}
{Clay problem page}, rechecked during this iteration,
continues to identify that broader problem as unsolved.
No external analytic estimate is used below.

\subsection{Nonduplication and setup}

The V45 source was checked against its recorded hash
before modification. The V45 companion checkpoint
remains the latest completed one; the next is V50.
The arbitrary-exterior repair and symmetric
denominator argument of f12 V25 were also reread.
Symmetry eliminating a linear term from a lower
integral bound is therefore not a new claim here.
Nor is the forest-Haar identity new.
The additions are the boundary-stable version of V45's
full compact normalizer, and the resulting
positive-volume boundary-event comparison with a
specified exterior force and repair cost.

Use a finite, non-wrapping four-dimensional cubic
lattice containing \(B_n\), \(n\ge2\), and all
plaquettes incident to \(F_n\).
The same argument allows missing outer plaquettes;
only the usual upper incidence bounds are needed.
Retain \(A_n,F_n,F_n^{12},\mathcal P_n\) from V43.
Set
\begin{equation}
 \begin{split}
 d_n&=|F_n|=24n(n-1)^2,\qquad
 d_{12}=|F_n^{12}|=4n(n-1)^2,\\
 M_n&=\mathsf N_n=n^2(n-1)^2,\qquad
 r_n=(3n+1)(n-1)^3,\qquad
 \frac{d_{12}}{M_n}=\frac4n .
 \end{split}
 \label{eq:f15v46:counts}
\end{equation}
Here \(d_n\) counts group variables, not real dimensions.
Every relevant coupling, including those outside the
interior block, obeys
\begin{equation}
 \beta\ge1,\qquad 0\le\tau\le1/32,\qquad
 (1-\tau)\beta\le J_p\le(1+\tau)\beta .
 \label{eq:f15v46:source}
\end{equation}
Write \(a_\tau=1+\tau\) and \(c_\tau=1/5-\tau\).
The couplings are fixed as the boundary is varied.

For \(0\le\phi\le1/(4n^2)\), let \(F_\phi\)
be the restriction to \(F_n\) of the coherent
field with \(R_\phi=\exp(i\phi\sigma_3)\).
V45's notation at this center becomes
\[
 \overline S_{J,\phi}
       =f(\phi)\sum_{p\in\mathcal P_n^{12}}J_p,\qquad
 \Delta_{\tau,\phi}
       =\frac{\tau^2M_n\sin^2\phi}{c_\tau}.
\]
Write an arbitrary boundary perturbation as
\begin{equation}
 F_e=(F_\phi)_e\exp X_e,\qquad
 E=\sum_{e\in F_n}|X_e|^2,\qquad
 E_{12}=\sum_{e\in F_n^{12}}|X_e|^2 .
 \label{eq:f15v46:X}
\end{equation}
One may choose principal logarithms; uniqueness is
only needed later on small balls.
The \(X_e\)'s need not commute with one another
or with the coherent field.

Use exactly the V45 forest coordinates for the
interior. Its tree values are the coherent ones,
and its gauge is identity on the boundary.
For each reduced interior \(z\), define
\(\mathcal C_\phi(z)\) by completing it with
\(F_\phi\), not with the perturbed boundary \(F\).
This reference-completion observable is used only
inside deterministic comparisons. It is not
silently identified with V44's observable
inferred from the perturbed perimeter.
V45 supplies
\begin{equation}
 D_{\mathcal T}(z)\le n^3\mathcal C_\phi(z),\qquad
 S_J(z;F_\phi)-\overline S_{J,\phi}
       \ge\beta\{c_\tau\mathcal C_\phi(z)/2
                                      -2\Delta_{\tau,\phi}\}.
 \label{eq:f15v46:reference}
\end{equation}

\subsection{A noncommuting boundary perturbation lemma}

\begin{proposition}[The boundary error is quadratic after coercivity]
\label{f15v46:perturbation}
For every reduced compact interior and every
perturbation \eqref{eq:f15v46:X},
\begin{equation}
 \left|S_J(z;F)-S_J(z;F_\phi)\right|
 \le a_\tau\beta
     \left\{\sqrt{2E\mathcal C_\phi(z)}
                 +\sin\phi\sqrt{d_{12}E_{12}}+E\right\}.
 \label{eq:f15v46:perturbation}
\end{equation}
This comparison does not restrict any interior link
to a small-field chart.
\end{proposition}

\begin{proof}
Vary \(F_e(t)=(F_\phi)_e\exp(tX_e)\),
\(0\le t\le1\), keeping every interior fixed.
For an affected plaquette put
\(\ell_p=\sum_{e\in F_n\cap\partial p}|X_e|\).
Differentiating its ordered quaternion product gives
\[
 |(f(\alpha(U_p(t))))''|\le\ell_p^2 .
\]
Indeed each second derivative of one exponential
has norm \(|X_e|^2\), and the two cross terms
for distinct factors have total norm at most
\(2|X_e||X_{e'}|\). Every other factor is unitary.
Taking the scalar part cannot increase the bound.
This holds also for inverse edge factors.
V43's one-face boundary incidence and its bound
of two boundary edges per affected face give
\[
        \sum_{p\in\mathcal P_n}\ell_p^2\le2E .
\]
The Taylor remainder of the weighted action is
therefore at most \(a_\tau\beta E\).

For the first derivative, a differentiated link
in the cyclic trace inserts a conjugated pure
quaternion of norm \(|X_e|\).
Its scalar pairing with the plaquette is bounded
by \(|X_e|\,|\operatorname{Im}U_p|\).
At the coherent boundary, conjugate a parallel
plaquette to its V44 rooted factor \(G_p\).
The imaginary-part map has Lipschitz constant one,
so
\[
 |\operatorname{Im}U_p|
   =|\operatorname{Im}G_p|
   \le |G_p-R_\phi|+\sin\phi.
\]
For nonparallel plaquettes use
\(|\operatorname{Im}U_p|\le|U_p-I|\).
Denote these squared deviations by \(\delta_p^2\);
their sum is \(\mathcal C_\phi(z)\).
Cauchy--Schwarz yields
\[
 \sum_p\ell_p\delta_p
       \le\sqrt{2E\mathcal C_\phi(z)}.
\]
Every edge of \(F_n^{12}\) occurs once in the
parallel sum. Hence
\[
 \sum_{p\in\mathcal P_n^{12}}\ell_p
       =\sum_{e\in F_n^{12}}|X_e|
       \le\sqrt{d_{12}E_{12}}.
\]
Multiply by \(a_\tau\beta\) and add the
Taylor remainder. All root conjugations were
used only in norm inequalities, not to commute
the perturbations through the plaquette.
\end{proof}

For the coherent interior itself the error term
\(\mathcal C_\phi\) vanishes. In particular
\begin{equation}
 s_J(F)\le S_J(\overline U^\phi;F)
     \le\overline S_{J,\phi}
          +a_\tau\beta\{E+\sin\phi\sqrt{d_{12}E_{12}}\}.
 \label{eq:f15v46:competitor}
\end{equation}
This comparison does not say that the coherent
interior minimizes the perturbed problem.

\subsection{Matching compact normalizers throughout a boundary neighborhood}

Define
\[
 \ell_\beta=2^{-13}\beta^{-3/2},\qquad
 u_{n,\beta}=2^{10}n^{9/2}\beta^{-3/2},\qquad
 Q_n=\frac{u_{n,\beta}}{\ell_\beta}=2^{23}n^{9/2}.
\]
These are numerical constants, not edge variables.
The partition function \(\mathcal Z_J(F)\)
always integrates the full original compact
interior with normalized Haar measure.

\begin{theorem}[Boundary-stable full compact envelope]
\label{f15v46:Z}
Under \eqref{eq:f15v46:source}, for every
\eqref{eq:f15v46:X},
\begin{equation}
 \begin{split}
 &\ell_\beta^{r_n}
       e^{-(17/16)\beta M_n f(\phi)-36\beta E}
       \le\mathcal Z_J(F)\\
 &\hspace{14mm}\le
 u_{n,\beta}^{r_n}
       e^{-(29/32)\beta M_n f(\phi)+48\beta E}.
 \end{split}
 \label{eq:f15v46:Z}
\end{equation}
At the flat center \(\phi=0\), the sharper constants are
\begin{equation}
       \ell_\beta^{r_n}e^{-2\beta E}
          \le\mathcal Z_J(F)
          \le u_{n,\beta}^{r_n}e^{14\beta E}.
 \label{eq:f15v46:flat-Z}
\end{equation}
In particular, perturbing the boundary on the
\(\beta^{-1/2}\) scale incurs a bounded quadratic
cost per boundary variable, not a lost power
of \(\beta\) in the interior normalizer.
\end{theorem}

\begin{proof}
The forest bijection and exact Haar disintegration
of V45 hold for every fixed \(F\).
At the actual reduced minimum, the same global
Wilson second-derivative estimate still applies.
Thus V45's lower-ball proof gives
\(\mathcal Z_J(F)\ge e^{-s_J(F)}\ell_\beta^{r_n}\).
Combine this with \eqref{eq:f15v46:competitor}.

For the upper bound, write \(C=\mathcal C_\phi(z)\)
and use
\[
 a_\tau\sqrt{2EC}
       \le\frac{c_\tau}{4}C+\frac{2a_\tau^2}{c_\tau}E.
\]
Equations \eqref{eq:f15v46:reference} and
\eqref{eq:f15v46:perturbation} then give
\begin{equation}
 \begin{split}
 S_J(z;F)\ge{}&
 \overline S_{J,\phi}-2\beta\Delta_{\tau,\phi}
       +\frac{\beta c_\tau}{4n^3}D_{\mathcal T}(z)\\
 &-\beta\left\{
       \left(a_\tau+\frac{2a_\tau^2}{c_\tau}\right)E
                    +a_\tau\sin\phi\sqrt{d_{12}E_{12}}\right\}.
 \end{split}
 \label{eq:f15v46:coercivity}
\end{equation}
V45's global chord inequality and Haar integral
bound imply
\[
 \int e^{-\beta c_\tau D_{\mathcal T}/(4n^3)}\,dH(z)
       \le u_{n,\beta}^{r_n}.
\]
Indeed the geodesic Gaussian coefficient is at
least \(\beta/(128n^3)\). This integrates the
whole compact domain, not a truncated Gaussian.

The remaining first-order coherent term can
be absorbed with a small classical allowance.
Since \(\sin^2\phi\le2f(\phi)\) and
\(d_{12}/M_n=4/n\),
\begin{equation}
 a_\tau\sin\phi\sqrt{d_{12}E_{12}}
       \le\frac1{32}M_n f(\phi)
                       +\frac{64a_\tau^2}{n}E.
 \label{eq:f15v46:flux-absorption}
\end{equation}
V45 gives
\(\overline S_{J,\phi}-2\beta\Delta_{\tau,\phi}
       \ge(15/16)\beta M_n f(\phi)\);
also \(\overline S_{J,\phi}
       \le(1+\tau)\beta M_n f(\phi)\).
Consequently the lower classical coefficient
is at most \(1+\tau+1/32\le17/16\), and the
upper-integral decay coefficient is at least
\(15/16-1/32=29/32\).

For the quadratic constants it suffices to take
\(a_\tau\le33/32\), \(c_\tau\ge27/160\) and \(n\ge2\).
They give
\[
 a_\tau+\frac{64a_\tau^2}{n}\le\frac{561}{16}<36,
 \qquad
 a_\tau+\frac{2a_\tau^2}{c_\tau}
                 +\frac{64a_\tau^2}{n}\le\frac{143}{3}<48 .
\]
This proves \eqref{eq:f15v46:Z}.
At \(\phi=0\), no flux allowance or source
floor remains. The constants reduce to
\(a_\tau<2\) and
\(a_\tau+2a_\tau^2/c_\tau\le1309/96<14\),
proving \eqref{eq:f15v46:flat-Z}.
\end{proof}

The bound is uniform over the specified source
arrays but retains \(Q_n^{r_n}\).
It is not uniform in the block size.
Also the reference-completion observable used
in the proof is not an improvement of V44's
actual-perimeter tail for arbitrary \(F\);
the theorem here concerns the integral itself.

\subsection{The exact exterior force and its action bound}

Condition on the original links \(Y\) outside
\(A_n\cup F_n\). As in V45, let
\(S_{\rm rest}(F,Y)\) sum the plaquettes touching
\(F_n\) but no updated interior edge.
Define
\begin{equation}
 \begin{split}
 T_\phi(Y)&=S_{\rm rest}(F_\phi,Y),\qquad
 H_\phi(Y)=\sum_e|h_{\phi,e}(Y)|^2,\\
 \left.\frac{d}{dt}S_{\rm rest}(F_\phi\exp(tX),Y)\right|_{t=0}
      &=\beta\sum_{e\in F_n}\langle h_{\phi,e}(Y),X_e\rangle .
 \end{split}
 \label{eq:f15v46:force}
\end{equation}
The Lie algebra has its unit-quaternion Euclidean norm.
These are the actual first derivatives of the
original outside action in the fixed boundary frame.
They are not derivatives of an auxiliary measure.

\begin{proposition}[Outside Taylor remainder and force-energy inequality]
\label{f15v46:outside}
For every exterior \(Y\) and boundary perturbation,
\begin{equation}
 \left|S_{\rm rest}(F_\phi\exp X,Y)-T_\phi(Y)
           -\beta\sum_e\langle h_{\phi,e}(Y),X_e\rangle\right|
       \le11\beta E,
 \label{eq:f15v46:outside-Taylor}
\end{equation}
and
\begin{equation}
             H_\phi(Y)\le40a_\tau\,T_\phi(Y)/\beta .
 \label{eq:f15v46:force-energy}
\end{equation}
\end{proposition}

\begin{proof}
Every \(F_n\) edge belongs to exactly one affected
interior face, hence to at most five remaining
faces. A remaining face contains at most four
such edges. Its second derivative is bounded
by the square of the sum of their generator norms.
Summing gives
\[
 |S_{\rm rest}''|\le20a_\tau\beta E.
\]
The Taylor remainder is at most
\(10a_\tau\beta E<11\beta E\).

Put \(j_p=J_p/\beta\). Each contribution to
\(h_{\phi,e}\) has norm \(j_p|\sin\alpha(U_p)|\).
Cauchy--Schwarz over at most five incident faces,
followed by \(j_p\le a_\tau\), gives
\[
 |h_{\phi,e}|^2
   \le5\sum_{\substack{p\ {\rm in\ rest}\\p\ni e}}
                       j_p^2\sin^2\alpha(U_p)
   \le10a_\tau
        \sum_{\substack{p\ {\rm in\ rest}\\p\ni e}}
                              j_p f(\alpha(U_p)).
\]
Here all plaquettes are evaluated at \((F_\phi,Y)\).
Summing over \(e\) counts each rest face at most
four times. This proves \eqref{eq:f15v46:force-energy}.
No favorable sign of the outside force is assumed.
\end{proof}

\subsection{Positive-volume thermal balls with the original Haar measure}

Fix \(0<\rho\le1/8\) and \(t_\beta=\rho/\sqrt\beta\).
Define the product boundary ball
\begin{equation}
 \mathcal B_\phi
   =\{F:\alpha((F_\phi)_e^{-1}F_e)\le t_\beta
                                      \text{ for every }e\in F_n\}.
 \label{eq:f15v46:ball}
\end{equation}
Its interior is nonempty and has positive Haar
measure. Closed versus open balls changes no integral.
On it, \(E\le d_n\rho^2/\beta\).
The same normalized one-edge Haar volume appears
at every center:
\begin{equation}
 b(t)=\frac2\pi\int_0^t\sin^2u\,du,\qquad
 H_{F_n}(\mathcal B_\phi)=b(t_\beta)^{d_n}.
 \label{eq:f15v46:volume}
\end{equation}
For \(0<t\le1/8\), the elementary inequalities
\(u/2\le\sin u\le u\) and \(1<\pi<4\) imply
\[
        t^3/24\le b(t)\le2t^3/3.
\]
Thus the boundary volume has exactly the
\(\beta^{-3d_n/2}\) power up to constants
to the power \(d_n\). It must not be set to one.

The conditional Haar law in a single ball,
expressed in its principal logarithm \(X\),
is rotationally invariant and invariant under
\(X\mapsto-X\). Its norm is at most \(t_\beta\).
Consequently, for every fixed collection of
Lie algebra vectors \(h_e\),
\begin{equation}
 1\le
 \frac1{b(t_\beta)^{d_n}}
  \int_{\mathcal B_\phi}
       e^{-\beta\sum_e\langle h_e,X_e\rangle}\,dH(F)
 \le
       e^{\,\beta\rho^2\sum_e|h_e|^2/2}.
 \label{eq:f15v46:symmetry}
\end{equation}
For the lower bound, replace the exponential by
its symmetric average, a hyperbolic cosine,
which is at least one. For the upper bound,
independence here is only that of the
\emph{reference product Haar measure}.
On one factor the symmetric average is at most
\(\cosh(\beta t_\beta|h_e|)
 \le\exp(\beta^2t_\beta^2|h_e|^2/2)\).
The last inequality follows by comparing power
series, since \((2k)!\ge2^k k!\).
Multiplying proves the claim.
No independence of the interacting boundary
under its Wilson law is being asserted.

\subsection{Actual conditional probabilities, with the outside cost paid}

Use the unnormalized boundary mass
\[
 \mathcal I_\phi(Y)=
       \int_{\mathcal B_\phi}
              e^{-S_{\rm rest}(F,Y)}\mathcal Z_J(F)\,dH(F).
\]
The true conditional boundary density of V45
has one common normalizer. Therefore
\begin{equation}
 \frac{\mathbb P_J(\mathcal B_\theta\mid Y)}
      {\mathbb P_J(\mathcal B_0\mid Y)}
                 =\frac{\mathcal I_\theta(Y)}{\mathcal I_0(Y)}.
 \label{eq:f15v46:actual}
\end{equation}
Every denominator is strictly positive.
This is a ratio of positive-probability events,
not pointwise values of a density.

\begin{theorem}[Normalized coherent-neighborhood comparison]
\label{f15v46:probability}
Let \(0\le\theta\le1/(4n^2)\), put
\(\Delta T=T_\theta(Y)-T_0(Y)\), and retain
\(0<\rho\le1/8\). Then
\begin{equation}
 \begin{split}
 &Q_n^{-r_n}
   e^{-\Delta T-(17/16)\beta M_n f(\theta)
                    -72d_n\rho^2-\beta\rho^2H_0(Y)/2}\\
 &\qquad\le
   \frac{\mathbb P_J(\mathcal B_\theta\mid Y)}
        {\mathbb P_J(\mathcal B_0\mid Y)}\\
 &\qquad\le
 Q_n^{r_n}
   e^{-\Delta T-(29/32)\beta M_n f(\theta)
                    +72d_n\rho^2+\beta\rho^2H_\theta(Y)/2}.
 \end{split}
 \label{eq:f15v46:force-ratio}
\end{equation}
In particular the forces may be eliminated at
the cost of retaining the outside actions:
\begin{equation}
 \begin{split}
 &Q_n^{-r_n}
    e^{-(17/16)\beta M_n f(\theta)-72d_n\rho^2
                                      -T_\theta(Y)+T_0(Y)/2}\\
 &\qquad\le
   \frac{\mathbb P_J(\mathcal B_\theta\mid Y)}
        {\mathbb P_J(\mathcal B_0\mid Y)}\\
 &\qquad\le
 Q_n^{r_n}
    e^{-(29/32)\beta M_n f(\theta)+72d_n\rho^2
                                      +T_0(Y)-T_\theta(Y)/2}.
 \end{split}
 \label{eq:f15v46:action-ratio}
\end{equation}
The sources, outside links and boundary reference
measure are unchanged between the two events.
\end{theorem}

\begin{proof}
On \(\mathcal B_\theta\), use
\eqref{eq:f15v46:Z},
\eqref{eq:f15v46:outside-Taylor} and
\eqref{eq:f15v46:symmetry}. They give
\begin{equation}
 \begin{split}
 \mathcal I_\theta(Y)
 &\ge b(t_\beta)^{d_n}\ell_\beta^{r_n}
    e^{-T_\theta-(17/16)\beta M_n f(\theta)-47d_n\rho^2},\\
 \mathcal I_\theta(Y)
 &\le b(t_\beta)^{d_n}u_{n,\beta}^{r_n}
    e^{-T_\theta-(29/32)\beta M_n f(\theta)
                          +59d_n\rho^2+\beta\rho^2H_\theta/2}.
 \end{split}
 \label{eq:f15v46:theta-mass}
\end{equation}
Here \(47=36+11\) and \(59=48+11\).
The sharper flat estimate gives
\begin{equation}
 \begin{split}
 \mathcal I_0(Y)
 &\ge b(t_\beta)^{d_n}\ell_\beta^{r_n}
                  e^{-T_0-13d_n\rho^2},\\
 \mathcal I_0(Y)
 &\le b(t_\beta)^{d_n}u_{n,\beta}^{r_n}
                  e^{-T_0+25d_n\rho^2+\beta\rho^2H_0/2}.
 \end{split}
 \label{eq:f15v46:flat-mass}
\end{equation}
The constants are \(13=2+11\) and \(25=14+11\).
Divide the respective lower and upper bounds.
The boundary volume cancels exactly, as do
the interior powers \(\beta^{-3r_n/2}\).
Since \(47+25=59+13=72\), this proves
\eqref{eq:f15v46:force-ratio}.

Finally \eqref{eq:f15v46:force-energy} and
\(\rho\le1/8\) imply
\[
       \beta\rho^2H_\phi/2
            \le20a_\tau\rho^2T_\phi\le T_\phi/2.
\]
Substitute this into the upper exponent and
the lower negative-force exponent, respectively.
This proves \eqref{eq:f15v46:action-ratio}.
\end{proof}

For example, if the exterior extends the flat
field so that \(T_0(Y)=0\), then
\begin{equation}
 \mathbb P_J(\mathcal B_\theta\mid Y)
       \le\min\left\{1,\,
       Q_n^{r_n}e^{-(29/32)\beta M_n f(\theta)+72d_n\rho^2}
                  \right\}.
 \label{eq:f15v46:flat-exterior}
\end{equation}
Indeed \(T_\theta\ge0\) and
\(\mathbb P_J(\mathcal B_0\mid Y)\le1\).
More generally a known bound
\(T_0(Y)\le\gamma\beta M_n f(\theta)\)
replaces \(29/32\) by \(29/32-\gamma\).
These are genuine finite-volume conditional
statements for the stated exteriors.
They do not assert that those exteriors are
typical under the full law.

\subsection{The remaining marginal input is a repair-energy moment}

Let \(\nu_J\) denote the actual marginal of
the links \(Y\). Dropping the positive
denominator-event probability as above and
integrating yields the rigorous, possibly weak, bound
\begin{equation}
 \mathbb P_J(\mathcal B_\theta)
 \le\min\left\{1,\,
   Q_n^{r_n}e^{-(29/32)\beta M_n f(\theta)+72d_n\rho^2}
   \mathbb E_{\nu_J}
       e^{\,T_0(Y)-T_\theta(Y)/2}\right\}.
 \label{eq:f15v46:marginal}
\end{equation}
All quantities are finite at fixed volume and
couplings. No useful uniform upper bound for
the last expectation has been proved.
It is the precise remaining exterior moment
for this particular neighborhood comparison.

The distinction in that last sentence matters.
\(T_0(Y)\) is the action after the boundary links
have been replaced by \(F_0\); it is not the
outside action of the sampled pair \((F,Y)\).
A sampled field with small plaquette energy
can have large fixed-frame repair cost.
Consequently an earlier exponential moment
for sampled energy cannot simply be inserted
into \eqref{eq:f15v46:marginal}.
Choosing a \(Y\)-dependent boundary gauge would
also change the centers and events; its
measurability, multiplicity and integral costs
would have to be treated explicitly.

For \(n^2\theta\le1/8\), V43's fit admits
the identity boundary gauge and this coherent
parameter. Thus
\begin{equation}
 \beta\ge256n\,d_n\rho^2
       \quad\Longrightarrow\quad
                 \mathcal B_\theta\subset\Gamma_n .
 \label{eq:f15v46:tube-inclusion}
\end{equation}
This follows from
\(\mathcal E_n(F)\le E\le d_n\rho^2/\beta\).
The bound is only inclusion of a fixed-frame
ball in the gauge-invariant fit tube.
It gives neither the measure of the whole tube
nor a covering or bound for \(\Gamma_n^c\).
Even the conditional comparison retains the
entropy term \(r_n\log Q_n\); it is not
a scale-uniform estimate.

The next upstream target is therefore a
controlled, gauge-compatible repair of the
actual exterior or an actual-law estimate
of the resulting repair-energy moment.
Repeating conditional normalizer bounds
without paying that repair would be circular.
The passage from compact lattice estimates
to the interacting continuum and its physical
Hamiltonian gap remains a further task.

\subsection{Checks, preservation and claim boundary}

The diagnostic is
\begin{center}\small
\path{scripts/verify_ym_f15_v46_boundary_neighborhoods.py}.
\end{center}
Six complete inside/outside geometries check
the one-face/five-face incidence split.
There are 720 noncommuting affected-face
first/second jet checks, 18 weighted
inside/outside second-jet checks, 18 exact
outside-force contractions and 18
outside-force energy/compact bounds.
Eight exact constant margins, 31 symmetric
power-series coefficient checks, and 31
volume/tube-scale identities supplement them.
All arithmetic in the script is rational
or integer. The global Taylor, Haar-integral
and probability inequalities are proved
above; no numerical integration or sampled
partition function stands in for a proof.

The diagnostic has 8,858 bytes and SHA--256
\begin{center}\scriptsize\ttfamily
2F32657D4660C9DDB366FE844B438445E0CDB6EF73BC686041AFFE2F092602D0.
\end{center}
The V45 predecessor has 1,062,828 bytes and SHA--256
\begin{center}\scriptsize\ttfamily
E21A8BD16003BE1CA9ABA316DA335F79761CA6676BF84B8C650EF0A5D12920D4.
\end{center}
Removing this final section and reversing
the title increment recovers V45 exactly.
V46 replaces V45 as the sole active main
source. The 14 protected inputs and 88
earlier diagnostics remain unchanged.
Only \path{.tex} sources belong in
\path{latex_notes}; this iteration's build
products belong in
\begin{center}\small
\path{artifacts/ym_f15_v46_texcheck/}.
\end{center}

\begin{firewall}
V46 extends the full compact normalizer to
noncommuting boundary neighborhoods and
proves an actual positive-volume conditional
event comparison. It retains the boundary
Haar volume, cancels the matched coupling
powers, and pays the outside force by an
explicit repair-action cost.
It does not prove that cost small under the
actual exterior marginal, control
\(\Gamma_n^c\), establish physical clustering
or an iterated RG trajectory, construct the
interacting continuum, or prove the physical
Yang--Mills mass gap.
\end{firewall}
\section{V47: an exterior-only boundary repair and intrinsic energy witnesses}
\label{f15v47:section}

V46 leaves a flat-boundary repair cost which can be
large for a purely gauge-transformed sampled field.
We go upstream to the choice of boundary center.
Every edge of \(F_n\) has one outward plaquette with
three exterior-only edges. Their ordered product
defines a local center depending only on \(Y\).
This center is gauge covariant, and its discrepancy
from the sampled boundary is exactly an actual
plaquette action. A unit-prism argument then bounds
its outside repair action by sampled, exterior-only
curvature with a side-independent constant.

This removes the fixed-frame gauge artifact in that
repair construction. It does not make the new center
coherent. In particular the coherent normalizers of
V45--V46 cannot be used at the new center without
another theorem. The deterministic construction is
valid under all the permitted source arrays.
The thermal probability estimates below have the
more restricted, explicitly inherited homogeneous
selected-law hypotheses.

\subsection{Inputs retained and what is not being repeated}

The current V46 source was hash-checked before this
append. The V45 companion checkpoint remains current,
and the next checkpoint is V50.
The original-link repair below is different from
f12 V65's heat-bath repair supports and f12 V67's
mixed-operator energy accounting; those sections
were reread before choosing this direction.
Likewise the thermal cell-source estimate is
inherited, not a new probability theorem proved
from the finite diagnostic.

Ordered transport in nonabelian plaquette identities
is essential. For context, the connector issue is
explicit in
\href{https://arxiv.org/abs/hep-lat/0508003}
{Borisenko--Voloshin--Faber's plaquette formulation}.
No estimate or change of variables from that work
is imported. Here the original links remain the
integration variables, and the needed identities
are proved directly by ordered products.

Use the box \(B_n=\{0,\ldots,n\}^4\), \(n\ge2\),
embedded with a full outward unit collar.
Retain \(A_n,F_n,Y,\mathcal P_n\) from V43--V46.
The collar contains the ordinary cubic plaquettes
incident to \(F_n\); periodic embeddings are taken
without identifications within this collar.
Put \(d_n=|F_n|=24n(n-1)^2\).
For deterministic action comparisons assume
\[
 0\le\tau\le1/32,\qquad
 (1-\tau)\beta\le J_p\le(1+\tau)\beta,\qquad
 a_\tau=1+\tau .
\]
The field \(Y\) consists of the original links
outside \(A_n\cup F_n\). No sampled link of \(Y\)
will be changed.

\subsection{A canonical outward path using only the given exterior}

For \(e=(x,\mu)\in F_n\), V43 proves that exactly
one transverse coordinate \(x_\nu\), \(\nu\ne\mu\),
lies in \(\{0,n\}\); the other two transverse
coordinates are internal. Define
\begin{equation}
 \nu(e)=\nu,\qquad
 \sigma(e)=
 \begin{cases}
  -1,&x_\nu=0,\\
  +1,&x_\nu=n .
 \end{cases}
 \label{eq:f15v47:normal}
\end{equation}
Let \(y=x+e_\mu\), \(x^+=x+\sigma e_\nu\),
and \(y^+=y+\sigma e_\nu\). The three-edge path
\begin{equation}
       b_e:\quad x\longrightarrow x^+
                        \longrightarrow y^+
                        \longrightarrow y
 \label{eq:f15v47:bypass}
\end{equation}
uses only links of \(Y\). The two normal links
cross outward from the box, and the middle link
lies on its outward translate.
Set
\begin{equation}
             B_e(Y)=U_Y(b_e),\qquad
             W_e=B_e(Y)^{-1}F_e .
 \label{eq:f15v47:center}
\end{equation}
The letter \(B_e\) here denotes a group-valued
boundary center, not the geometric box.

Let \(p_e^+\) be the outward plaquette between
\(e\) and the middle edge of \(b_e\).
It contains exactly one edge of \(F_n\), namely
\(e\), and no edge of \(A_n\).
The plaquettes \(p_e^+\) for different \(e\)
are distinct.

\begin{proposition}[Intrinsic discrepancy and exact conditional Haar coordinates]
\label{f15v47:coordinates}
The map \(Y\mapsto B(Y)\) is smooth and local,
depending on three exterior edges per boundary edge.
For every compact configuration,
\begin{equation}
 \begin{split}
 f(\alpha(W_e))&=f(\alpha(U_{p_e^+})),\\
 D_n(F,Y):=\sum_{e\in F_n}|F_e-B_e(Y)|^2
             &=2\sum_{e\in F_n}f(\alpha(U_{p_e^+})).
 \end{split}
 \label{eq:f15v47:distance}
\end{equation}
Under any vertex gauge transformation \(g\),
\begin{equation}
 B_e(Y^g)=g_xB_e(Y)g_y^{-1},\qquad
                         W_e^g=g_yW_eg_y^{-1}.
 \label{eq:f15v47:gauge}
\end{equation}
Finally \((F,Y)\leftrightarrow(W,Y)\) is a
global bijection with
\begin{equation}
                     dH(F)\,dH(Y)=dH(W)\,dH(Y).
 \label{eq:f15v47:Haar}
\end{equation}
In particular, the discrepancy, its angle version
and neighborhoods centered at \(B(Y)\) are
gauge-invariant observables.
\end{proposition}

\begin{proof}
The loop consisting of \(e\) followed by the
reverse of \(b_e\) traverses \(p_e^+\).
Its product is \(F_eB_e^{-1}\).
Depending on the positive-plaquette convention,
the latter is conjugate to \(U_{p_e^+}\) or
its inverse. Its real part therefore equals
that of \(W_e\). This proves the first identity;
\(|F_e-B_e|^2=2f(\alpha(B_e^{-1}F_e))\)
proves the second.
Path transport transforms by its endpoint gauges,
giving \eqref{eq:f15v47:gauge}.

For fixed \(Y\), \(F_e=B_e(Y)W_e\) is simply
a left translation on each group factor.
It preserves normalized product Haar measure.
Fubini proves \eqref{eq:f15v47:Haar},
even though \(B\) depends on \(Y\).
There is no transformation of \(Y\), and no
derivative of \(B(Y)\) or additional Jacobian
is required in this triangular coordinate map.
Smoothness and locality follow from the
three-link formula.
\end{proof}

Define also
\[
       E_n^{\rm out}(F,Y)
              =\sum_{e\in F_n}\alpha(W_e)^2.
\]
The global chord-angle comparison gives
\begin{equation}
 E_n^{\rm out}\le\frac{\pi^2}{4}D_n
        =\frac{\pi^2}{2}\sum_e f(\alpha(U_{p_e^+})).
 \label{eq:f15v47:angle}
\end{equation}
All plaquettes on the right are in the sampled
original field, not in a flat-frame repair.

\subsection{The outside action separates into mirrors and tangential faces}

Let \(\mathcal P_n^+=\{p_e^+:e\in F_n\}\).
The plaquettes in \(S_{\rm rest}(F,Y)\)
split into \(\mathcal P_n^+\) and a set
\(\mathcal R_n\) of tangential boundary faces.
Every \(p\in\mathcal R_n\) has exactly one
fixed transverse coordinate equal to \(0\) or \(n\).
It is a face of \(B_n\), and all its \(F_n\)
edges have this same outward normal and sign.

To verify the classification, start with an
incident edge \(e=(x,\mu)\in F_n\).
The two faces with directions \(\mu,\nu(e)\)
are its inward affected face and its outward
mirror \(p_e^+\). The remaining four incident
faces have directions \(\mu,\kappa\) with
\(\kappa\ne\mu,\nu(e)\); they lie in the
boundary hyperplane \(x_{\nu(e)}=0\) or \(n\).
All their coordinates remain in the box since
\(x_\kappa\) was internal. The remaining
transverse coordinate is internal.
This also proves
\begin{equation}
                        |\mathcal R_n|\le4d_n .
 \label{eq:f15v47:tangent-count}
\end{equation}

The repair \(F=B(Y)\) makes every mirror
plaquette exactly \(I\).
For a tangential face \(p\), let \(c(p)\)
be the unit three-dimensional cube obtained
by translating \(p\) one unit outward.
It has an outer face \(p^+\) and four side
faces. Define the exterior filling
\begin{equation}
 K(p)=\{p^+\}
     \ \cup\
     \{\text{side face over }e:
                          e\in\partial p,\ e\notin F_n\}.
 \label{eq:f15v47:filling}
\end{equation}
Every edge in every face of \(K(p)\) belongs
to \(Y\). For a side face this uses that
its base edge \(e\notin F_n\) is a boundary
edge, hence is not in \(A_n\) either.
The other three edges are exterior.
There are at most five distinct faces in \(K(p)\).

\begin{proposition}[Noncommuting unit-prism comparison]
\label{f15v47:prism}
For arbitrary compact exterior links,
\begin{equation}
 f(\alpha(U_p(B(Y),Y)))
                 \le5\sum_{q\in K(p)}f(\alpha(U_q(Y)))
                 \qquad(p\in\mathcal R_n).
 \label{eq:f15v47:prism}
\end{equation}
\end{proposition}

\begin{proof}
Fix the outward normal of \(p\). For every
base edge of \(p\), including those not
in \(F_n\), consider its three-edge outward
transport with this common normal.
Denote this ideal edge by \(\widehat B_e\).
For the actual repaired \(F_n\) edges it
equals \(B_e(Y)\).

Writing \(L_x\) for the outward spoke at a
vertex of \(p\), each ideal positive edge
has the form \(L_xU_{e^+}L_y^{-1}\).
Their ordered product around \(p\) cancels
the adjacent \(L_y^{-1}L_y\) pairs.
It is therefore the conjugate of the
outer plaquette \(U_{p^+}\) by the spoke
at the starting corner. No plaquette
factors have been commuted.

For an unreplaced base edge \(e\notin F_n\),
the actual edge is \(Y_e\), and
\[
                |Y_e-\widehat B_e|
                       =|U_{q_e}(Y)-I|,
\]
where \(q_e\) is its side face; this is
the same single-plaquette identity as above.
Telescoping the four unit-quaternion
products, using norm invariance under
left and right multiplication, gives
\[
 |U_p(B(Y),Y)-I|
       \le |U_{p^+}(Y)-I|
          +\sum_{\substack{e\in\partial p\\e\notin F_n}}
                                      |U_{q_e}(Y)-I|.
\]
Square, use Cauchy--Schwarz for at most
five terms, and divide by two.
This proves \eqref{eq:f15v47:prism}.
\end{proof}

\subsection{A side-independent repaired-action bound}

Put
\begin{equation}
 \begin{split}
 \mathcal K_n&=\bigcup_{p\in\mathcal R_n}K(p),\\
 \mathcal H_n(Y)&=\beta\sum_{q\in\mathcal K_n}f(\alpha(U_q(Y))),\\
 \mathcal S_{n,J}(Y)&=\sum_{q\in\mathcal K_n}J_qf(\alpha(U_q(Y))),\\
 T_{n,J}^*(Y)&=S_{\rm rest}(B(Y),Y).
 \end{split}
 \label{eq:f15v47:repair-data}
\end{equation}
Both \(\mathcal H_n\) and \(\mathcal S_{n,J}\)
use original sampled links of \(Y\) only.
The set \(\mathcal K_n\) lies in the fixed
unit collar, and
\begin{equation}
                            |\mathcal K_n|\le20d_n .
 \label{eq:f15v47:collar-count}
\end{equation}

\begin{theorem}[Gauge-compatible exterior repair paid by sampled curvature]
\label{f15v47:repair}
For every compact exterior and every allowed
source array,
\begin{equation}
 0\le T_{n,J}^*(Y)
      \le20a_\tau\mathcal H_n(Y)
      \le22\mathcal H_n(Y),
 \qquad
 T_{n,J}^*(Y)\le22\mathcal S_{n,J}(Y).
 \label{eq:f15v47:repair}
\end{equation}
The constants are independent of \(n\),
volume and the exterior.
If all plaquettes in \(\mathcal K_n\) are
flat, the repaired outside action is zero.
\end{theorem}

\begin{proof}
The cubes \(c(p)\) for different
\(p\in\mathcal R_n\) are distinct.
Each cube protrudes outside \(B_n\) in
exactly one coordinate direction, and
its base face in the boundary hyperplane
is uniquely determined. A cube with a
different protruding direction cannot
be this same three-cell.

In a four-dimensional cubic lattice a
given two-face belongs to at most four
unit three-cells: there are two choices
of additional coordinate direction and
two choices of side in that direction.
It follows that a face \(q\) occurs in
at most four of the sets \(K(p)\).
The mirror actions vanish after repair.
Summing \eqref{eq:f15v47:prism} with
\(J_p\le a_\tau\beta\) gives
\[
 T_{n,J}^*
       \le5a_\tau\beta\sum_{p\in\mathcal R_n}
                                  \sum_{q\in K(p)}f(\alpha(U_q))
       \le20a_\tau\mathcal H_n .
\]
Also \(\mathcal S_{n,J}\ge(1-\tau)\mathcal H_n\).
For \(\tau\le1/32\),
\[
          20a_\tau<22,\qquad
          \frac{20a_\tau}{1-\tau}\le\frac{660}{31}<22.
\]
This proves both bounds. The cardinality
estimate \eqref{eq:f15v47:collar-count}
follows from at most five faces per
prism and \eqref{eq:f15v47:tangent-count}.
\end{proof}

The repair is only of the outside boundary
action. It does not specify an interior
extension or claim that the full block
minimum is small. The exact conditional
boundary density in the \(W\) variables is
\begin{equation}
 p_Y(W)=\frac1{\mathcal N(Y)}
       \exp\left\{-\sum_{e\in F_n}J_{p_e^+}f(\alpha(W_e))
                -S_{\rm tan}(B(Y)W,Y)\right\}
                                  \mathcal Z_J(B(Y)W).
 \label{eq:f15v47:conditional}
\end{equation}
Here \(S_{\rm tan}\) sums \(\mathcal R_n\),
\(\mathcal N(Y)\) is the full conditional
normalizer, and the reference is \(dH(W)\).
The diagonal mirror term is exact, but
the tangential action and interior
partition factor remain. They cannot
be omitted to declare the \(W_e\)'s
independent or to infer a physical mass.

\subsection{Actual thermal bounds with their selected-law hypotheses}

Only in this subsection, take the original
homogeneous selected law
\(\mu=\mu^{\mathrm W}_{M,\beta_M^\diamond}\)
from \eqref{f15v1:source-input}, far enough
along the selection to have
\eqref{eq:f15v20:selected-input}.
The global torus side \(M\) is distinct
from the local box side \(n\).
All boxes and unit collars are embedded.
Write \(\beta=\beta_M^\diamond\) and
\(X_p=\beta f(\alpha(U_p))\).
The existing cell energy satisfies
\begin{equation}
 X_{(x;\mu,\nu)}\le4T_x,\qquad
 \log\mathbb E_\mu e^{\sum_xb_xT_x}
                \le-5\sum_x\log(1-b_x)
                \quad(0\le b_x\le1/2).
 \label{eq:f15v47:CGF}
\end{equation}
The first inequality is a term of
\eqref{f15v1:cell}. For the second, add
the mean, bounded by \(5\sum b_x\), to
the inherited centered source estimate
with \(k_M\le5\).
This uses the specified homogeneous
selected law, not a new uniform estimate
under all \(\mu_J\).

For any distinct plaquette set \(\mathcal Q\),
let \(m_x\) count its positive plaquettes
based at \(x\). There are six positive
orientations, hence
\[
        0\le m_x\le6,\qquad
                   \sum_xm_x=|\mathcal Q|.
\]
Convexity on \(0\le m\le6\) gives
\[
       -\log(1-m/12)\le(m/6)\log2 .
\]
Consequently \eqref{eq:f15v47:CGF} implies
the convenient inherited-source application
\begin{equation}
 \log\mathbb E_\mu
           e^{(1/48)\sum_{p\in\mathcal Q}X_p}
                       \le\frac56|\mathcal Q|\log2 .
 \label{eq:f15v47:plaquette-CGF}
\end{equation}
Indeed the coefficient after \(X_p\le4T_x\)
is \(m_x/12\).

\begin{theorem}[Intrinsic adaptive-boundary fluctuations and a repair moment]
\label{f15v47:thermal}
In this selected-law regime,
\begin{equation}
 \begin{split}
 \log\mathbb E_\mu e^{\beta D_n/96}
       &\le\frac56d_n\log2,\\
 \log\mathbb E_\mu e^{\beta E_n^{\rm out}/(24\pi^2)}
       &\le\frac56d_n\log2,\\
 \log\mathbb E_\mu e^{T_{n,\beta}^*/1056}
       &\le\frac56|\mathcal K_n|\log2
          \le\frac{50}{3}d_n\log2 .
 \end{split}
 \label{eq:f15v47:thermal}
\end{equation}
In the last line \(T_{n,\beta}^*\) denotes
the homogeneous repaired outside action.
For families whose relevant plaquette-base
sets are pairwise disjoint, the same bounds
hold for sums of the observables, with the
right sides summed over the family.
No independence of the interacting boxes,
their centers or their exteriors is assumed.
\end{theorem}

\begin{proof}
The mirror set has exactly \(d_n\)
distinct plaquettes, and
\(\beta D_n=2\sum_{p\in\mathcal P_n^+}X_p\).
Apply \eqref{eq:f15v47:plaquette-CGF}.
The angle bound
\eqref{eq:f15v47:angle} proves the second line.
For the third use
\[
          T_{n,\beta}^*\le22
                  \sum_{q\in\mathcal K_n}X_q,\qquad
                              \frac{22}{1056}=\frac1{48}.
 \]
The cardinality bound completes it.

For a family with disjoint relevant
plaquette-base sets, the union still
has at most six selected positive
plaquettes at each base. Apply the same
single joint CGF to that union.
This is why support disjointness is
stated at the level of bases; no
independence assertion is used.
\end{proof}

For example, for \(k\) boxes with common
\(n\), pairwise disjoint mirror-base
sets, and \(u\ge0\), exponential Markov gives
\begin{equation}
 \mu\!\left(\bigcap_{i=1}^k
                  \{E_{n,i}^{\rm out}\ge u\}\right)
       \le
 \left[\min\left\{1,\,
       e^{-\beta u/(24\pi^2)+(5d_n/6)\log2}\right\}\right]^k .
 \label{eq:f15v47:joint-tail}
\end{equation}
These are actual-law, gauge-invariant
events with \(Y_i\)-dependent centers.
Their adaptivity is already present in
the exact mirror identity; it introduces
no selection-volume or union-over-gauges factor.
This conclusion does not apply to an
arbitrary optimized center in place of \(B(Y)\).

There is also a useful maximum-link bound.
Since \(\mathbb E_\mu e^{X_p/8}\le
\mathbb E_\mu e^{T_{\operatorname{base}(p)}/2}\le32\),
the same identity and
\(f(\delta)\ge2\delta^2/\pi^2\) yield
\begin{equation}
 \mu\{\max_{e\in F_n}\alpha(W_e)>\delta\}
       \le\min\left\{1,\,
                    32d_n e^{-\beta\delta^2/(4\pi^2)}\right\},
 \qquad 0\le\delta\le\pi .
 \label{eq:f15v47:maximum}
\end{equation}
Separately, under the homogeneous dyadic-torus
hypotheses of \eqref{f15v1:thermal-exponential},
which do not require this particular coupling
selection, the inherited single-plaquette
estimate gives
\begin{equation}
 \mu_{M,\beta}^{\mathrm W}
       \{\max_e\alpha(W_e)>\delta\}
       \le\min\{1,\ d_n C_L e^{-\beta\delta^2/\pi^2}\}.
 \label{eq:f15v47:general-beta}
\end{equation}
The large constant \(C_L\) and all hypotheses
of that inherited theorem remain in force.
This latter single-box statement is not a
replacement for the selected-law joint CGF.

\subsection{What the local repair resolves, and what it does not}

The new center has two exact advantages for the
current unresolved exterior comparison.
It is determined by the unchanged exterior,
and both its sampled discrepancy and its
repaired outside action are controlled by
sampled curvature. Thus the gauge artifact
of replacing the boundary by a fixed
coordinate identity is absent.
At a purely gauge-flat sampled field the
three-edge transport agrees with the
sampled boundary link, so \(D_n=0\).

However \(B(Y)\) generally has nonzero,
nonuniform boundary holonomies.
The repair theorem does not put it in the
coherent family \(F_\theta\), or prove a
small value of V43's coherent fit.
Small outward plaquettes alone do not
make these two properties equivalent.
The conditional factor
\(\mathcal Z_J(B(Y)W)\) in
\eqref{eq:f15v47:conditional} remains the
full interior response to that curved center.
Replacing it by V46's coherent envelope
would be an unsupported step.

There are also two distinct probabilistic
limitations. First, the deterministic
action estimate is uniform in the
permitted \(J\)'s, but the joint thermal
statements above use the homogeneous
selected law. V38's innovation bounds
under physical sources do not remove
their unestimated frustration term.
Second, a bound for
\(\mathbb E e^{T_{n,\beta}^*/1056}\)
does not bound
\(\mathbb E e^{T_{n,\beta}^*}\).
The finite source radius is a real
restriction, not a harmless constant.
Consequently the new repair moment
cannot simply be inserted at unit
coefficient into V46's remaining marginal
factor, even if its center were suitable.

The upstream task has become more precise:
control the full interior response to
this sampled, generally curved boundary
center, or construct a further repair
with an explicitly affordable source
coefficient and a proved coherent fit.
Any subsequent use must retain the
actual \(Y\) marginal and the exact
conditional normalization. No statement
here establishes a probability bound
for \(\Gamma_n^c\), scale-uniform
clustering, or the continuum physical gap.

\subsection{Exact checks and source preservation}

The diagnostic is
\begin{center}\small
\path{scripts/verify_ym_f15_v47_outward_boundary_repair.py}.
\end{center}
It checks 11 complete boundary/collar
geometries, including 22,848 exterior-only
three-edge paths and 27,264 distinct-prism
and overlap certificates.
Rational noncommuting configurations
give 1,008 exact mirror-distance/inverse
identities, 1,584 prism transport checks,
18 source-uniform repaired-action bounds,
and 1,008 center/relative-link gauge
covariance checks. Seven cell-source
endpoint inequalities and 31 surface
and moment constant checks are exact.
The fields include nontrivial pure
gauges, small curved fields and large
compact noncommuting fields.

These finite checks do not prove the
selected-law CGF or the integral claims;
the proofs and their inherited hypotheses
are explicitly stated above.
The script has 9,191 bytes and SHA--256
\begin{center}\scriptsize\ttfamily
5B253644E80BFA667B1AAB1362EBEAB3F12C4D00160B34AAE46072833F88EDD8.
\end{center}
The V46 predecessor has 1,085,918 bytes and SHA--256
\begin{center}\scriptsize\ttfamily
4C0FF30D0CA2867EF1BCBADE1F3832AD519FA6884D9996B87580C096B78C4084.
\end{center}
Removing this section and reversing the
title increment recovers V46 exactly.
V47 replaces V46 as the one active main
source. The 14 protected source inputs
and 89 earlier diagnostics are preserved.
Only \path{.tex} files belong in
\path{latex_notes}; build products for
this version belong in
\begin{center}\small
\path{artifacts/ym_f15_v47_texcheck/}.
\end{center}
The next companion checkpoint remains V50.

\begin{firewall}
V47 proves a local exterior-only
gauge-covariant repair, exact Haar
coordinates and sampled mirror-energy
identities, a side-independent bound
for its repaired outside action, and
specified homogeneous-law thermal
estimates. It does not identify its
center with a coherent boundary,
prove the required full-source
marginal closure, construct an
interacting continuum, or prove the
physical Yang--Mills mass gap.
\end{firewall}
\section{V48: an exterior-only flat reference and a measured coherent-fit event}
\label{f15v48:section}

V47 supplies a gauge-covariant boundary center
\(B(Y)\), but that center is generally curved.
The next step is not to assume it coherent.
Instead, apply the boundary-only path construction
already proved in V41 to V47's repaired boundary.
All inputs to those paths are then functions of
the unchanged exterior \(Y\).
This gives an explicit flat reference, together
with a quantified curvature cost for using it.

There are two different conclusions. Deterministically,
the full compact interior partition function can be
compared with the flat reference under every permitted
source array, with that cost retained. Probabilistically,
the inherited selected homogeneous law gives a joint
tail for failure of the resulting coherent-fit certificate.
The latter is not proved under arbitrary nonzero
physical sources. The continuum and mass-gap goal
therefore remains unproved.

\subsection{Nonduplication and scope}

The V47 source was checked against its final hash.
V41's explicit boundary paths and square homotopies
were reread in full. They already give the
\(9n^2\) boundary gauge estimate, and are not
reproved as a new result here.
The V20/V32 cell-source method already supplies
homogeneous small-field entry bounds.
The new point is a certificate and a flat
reference determined entirely by \(Y\), together
with their exact use in the V46 compact
conditional normalizer.
No choice depending on the sampled interior
or sampled \(F_n\) is made.

The V45 companion review remains current;
the next is V50. The previously cited
\href{https://math-events.uni-bonn.de/event/152/contributions/568/}
{2025 Bonn abstract on non-Abelian correlation inequalities}
was also rechecked. It describes a conjectural
route to a nonzero-coupling GKS analogue.
It is not a theorem authorizing source
monotonicity in the present argument.
No external correlation inequality is used below.

Keep \(B_n,A_n,F_n,Y\), the outward center
\(B_e(Y)\), its repaired boundary and its
exterior filling \(\mathcal K_n\) from V47.
Set \(d_n=24n(n-1)^2\).
All geometric constructions take place
in the embedded box and its unit collar.

\subsection{Complete the sampled exterior curvature support}

Define the full repaired boundary connection
\(\widehat U(Y)\) by using \(B_e(Y)\) on
\(F_n\) and the unchanged \(Y_e\) on every
other boundary edge.
Let
\begin{equation}
 \begin{split}
 \mathcal P_n^{\partial,0}
    &=\{p\in\mathcal P_n^\partial:
                              \partial p\cap F_n=\varnothing\},\\
 \mathfrak L_n&=\mathcal K_n\cup\mathcal P_n^{\partial,0},\\
 R_n(Y)&=\max_{q\in\mathfrak L_n}|U_q(Y)-I|,\\
 \mathfrak H_n(Y)&=\beta\sum_{q\in\mathfrak L_n}f(\alpha(U_q(Y))).
 \end{split}
 \label{eq:f15v48:support}
\end{equation}
Every plaquette in \(\mathfrak L_n\) uses
only original links of \(Y\).
The additional faces in
\(\mathcal P_n^{\partial,0}\) have both
fixed transverse coordinates on the
boundary. Conversely, a boundary face
with only one such transverse coordinate
has at least two edges in \(F_n\) when
\(n\ge2\). Hence
\begin{equation}
 |\mathcal P_n^{\partial,0}|=24n^2,\qquad
 |\mathfrak L_n|\le20d_n+24n^2\le22d_n .
 \label{eq:f15v48:counts}
\end{equation}
The last inequality uses
\(24n^2/d_n=n/(n-1)^2\le2\).

The unsquared telescoping inequality
in the proof of Proposition~\ref{f15v47:prism}
has at most five terms. Untouched boundary
faces are already in
\(\mathcal P_n^{\partial,0}\).
It follows, for every compact exterior, that
\begin{equation}
       \max_{p\in\mathcal P_n^\partial}
                  |U_p(\widehat U(Y))-I|\le5R_n(Y).
 \label{eq:f15v48:boundary-curvature}
\end{equation}
No smallness hypothesis has yet been imposed.

\subsection{A flat reference determined entirely by the exterior}

Take exactly the canonical boundary root
paths \(P_x\) of \eqref{eq:f15v41:paths},
rooted at \(o\). Define
\begin{equation}
 \begin{split}
 g_x(Y)&=U_{\widehat U(Y)}(P_x),\\
 F_e^\flat(Y)&=g_x(Y)^{-1}g_y(Y)
                         \quad(e=(x,y)\in F_n),\\
 K_n^\flat(Y)&=\sum_{e=(x,y)\in F_n}
                 \alpha(g_x(Y)B_e(Y)g_y(Y)^{-1})^2 .
 \end{split}
 \label{eq:f15v48:flat-reference}
\end{equation}
Each original path has at most \(4n\)
boundary edges. Replacing an occurrence
of a repaired edge by its three-edge
outward path expresses \(g_x(Y)\) as an
ordered word of at most \(12n\) original
\(Y\)-links. Thus this is a smooth,
explicit exterior-dependent construction.
It involves no minimizer selection.

Under an arbitrary vertex gauge \(h\),
\begin{equation}
 \begin{split}
 g_x(Y^h)&=h_o g_x(Y)h_x^{-1},\\
 F_e^\flat(Y^h)&=h_xF_e^\flat(Y)h_y^{-1}.
 \end{split}
 \label{eq:f15v48:covariance}
\end{equation}
These identities follow by endpoint
cancellation along the paths and V47's
covariance of \(B(Y)\).
In particular \(K_n^\flat\) is gauge invariant.
The reference \(F^\flat\) is the restriction
of a pure-gauge connection on the entire
boundary. In the \(g(Y)\) frame it is
exactly the zero-flux coherent field.

\begin{theorem}[Exterior curvature pays the flat-reference mismatch]
\label{f15v48:fit}
For every compact exterior,
\begin{equation}
 K_n^\flat(Y)
       \le\frac{2025\pi^2}{4}\,d_n n^4 R_n(Y)^2.
 \label{eq:f15v48:K}
\end{equation}
For the actual sampled boundary \(F\),
\begin{equation}
 \begin{split}
 E_g(F,Y)&:=\sum_{e=(x,y)\in F_n}
                         \alpha(g_x(Y)F_eg_y(Y)^{-1})^2,\\
 \mathcal E_n(F)&\le E_g(F,Y)
                  \le2E_n^{\rm out}(F,Y)+2K_n^\flat(Y).
 \end{split}
 \label{eq:f15v48:fit}
\end{equation}
Here \(\mathcal E_n\) is exactly V43's
coherent-fit statistic, and
\(E_n^{\rm out}\) is V47's intrinsic
outward discrepancy.
\end{theorem}

\begin{proof}
Apply Proposition~\ref{f15v41:boundary-gauge}
to \(\widehat U(Y)\), using
\eqref{eq:f15v48:boundary-curvature}.
For every boundary edge,
\[
             |g_x\widehat U_eg_y^{-1}-I|
                                  \le45n^2R_n(Y).
\]
The global angle-chord inequality
\(\alpha(V)^2\le(\pi^2/4)|V-I|^2\)
and summation over \(F_n\) prove
\eqref{eq:f15v48:K}.

By the bi-invariant triangle inequality,
\[
 \alpha(g_xF_eg_y^{-1})
   \le\alpha(B_e^{-1}F_e)
                           +\alpha(g_xB_eg_y^{-1}).
\]
Square with \((a+b)^2\le2a^2+2b^2\)
and sum. Finally the choice
\(\Theta=0\) and boundary gauge \(g(Y)\)
is admissible in V43's compact fit
minimization. This proves the first
inequality in \eqref{eq:f15v48:fit}.
\end{proof}

At a pure-gauge sampled exterior all
plaquettes in \(\mathfrak L_n\) are flat,
so \(K_n^\flat=0\) and \(B(Y)=F^\flat(Y)\)
on \(F_n\). No choice of a fixed coordinate
frame is needed for this conclusion.

\subsection{Use the flat reference in the genuine interior normalizer}

Keep the original source array with
\(\beta\ge1\), \(0\le\tau\le1/32\) and
\((1-\tau)\beta\le J_p\le(1+\tau)\beta\).
Use V46's constants
\[
 r_n=(3n+1)(n-1)^3,\qquad
 \ell_\beta=2^{-13}\beta^{-3/2},\qquad
 u_{n,\beta}=2^{10}n^{9/2}\beta^{-3/2}.
\]
The following assertion is deterministic
in \(Y\), not an assertion that \(Y\)
is typical under the source law.

\begin{theorem}[Full-source compact normalizer about the adaptive flat reference]
\label{f15v48:normalizer}
For every compact \(F,Y\),
\begin{equation}
 \ell_\beta^{r_n}e^{-2\beta E_g(F,Y)}
          \le\mathcal Z_J(F)
          \le u_{n,\beta}^{r_n}e^{14\beta E_g(F,Y)}.
 \label{eq:f15v48:normalizer}
\end{equation}
Consequently, writing
\(E=E_n^{\rm out}(F,Y)\), \(K=K_n^\flat(Y)\),
\begin{equation}
 \ell_\beta^{r_n}e^{-4\beta(E+K)}
          \le\mathcal Z_J(F)
          \le u_{n,\beta}^{r_n}e^{28\beta(E+K)}.
 \label{eq:f15v48:normalizer-repair}
\end{equation}
At the adaptive center itself, the sharper
specialization is
\begin{equation}
 \ell_\beta^{r_n}e^{-2\beta K}
          \le\mathcal Z_J(B(Y))
          \le u_{n,\beta}^{r_n}e^{14\beta K}.
 \label{eq:f15v48:center-normalizer}
\end{equation}
\end{theorem}

\begin{proof}
Fix \(Y\); the gauges \(g_x(Y)\) are then
fixed matrices. Extend them to interior
vertices in any fixed way. Applying
this gauge to the integration variables
in \(A_n\) preserves their normalized
product Haar measure, and the Wilson
trace terms are invariant.
The conditional action depends on
the exterior only through \(F_n\).
Thus \(\mathcal Z_J(F)\) equals the
same full compact partition function
with boundary \(g_xF_eg_y^{-1}\).

Choose principal logarithms of these
boundary group elements. Their squared
norm sum is \(E_g\), and apply the
flat-boundary estimate
\eqref{eq:f15v46:flat-Z}.
It allows arbitrary boundary
perturbations, not only a small ball.
This proves \eqref{eq:f15v48:normalizer}.
Equation~\eqref{eq:f15v48:fit} gives
\eqref{eq:f15v48:normalizer-repair}.
For \(F=B(Y)\), \(E_g=K\), giving
\eqref{eq:f15v48:center-normalizer}.
\end{proof}

This conditional change of variables
does not push forward the \(Y\) marginal
by \(Y\mapsto Y^{g(Y)}\).
Such a nonlinear map has not been
claimed invertible or measure preserving.
The exterior is a fixed parameter
throughout the proof. Equivalently,
the joint map
\((F,Y)\mapsto(g_x(Y)F_eg_y(Y)^{-1},Y)\)
preserves product Haar by Fubini,
with \(Y\) itself unchanged.

The estimate therefore fills a specific
gap left by V47: its curved center now
has a full-source compact normalizer
bound, but the explicit cost is
\(\beta K_n^\flat(Y)\).
Calling that cost zero, or uniform in
the box size, would still be incorrect.

\subsection{Pay the outside action of the adaptive flat reference}

Define
\[
            T_{n,J}^\flat(Y)
                  =S_{\rm rest}(F^\flat(Y),Y).
\]
This is a different repair from
\(T_{n,J}^*\), which uses \(B(Y)\).
Both leave all original \(Y\) links
unchanged.

\begin{proposition}[An explicit flat-repair debit]
\label{f15v48:flat-action}
For every permitted source array,
\begin{equation}
 \begin{split}
 T_{n,J}^\flat(Y)
       &\le2T_{n,J}^*(Y)+22\beta K_n^\flat(Y)\\
       &\le\mathfrak c_n\,\mathfrak H_n(Y),
 \qquad
 \mathfrak c_n=44+356400d_n n^4 .
 \end{split}
 \label{eq:f15v48:flat-action}
\end{equation}
\end{proposition}

\begin{proof}
On a rest plaquette, product telescoping
when \(B\) is replaced by \(F^\flat\)
gives
\[
 |U_p(F^\flat,Y)-I|
       \le |U_p(B,Y)-I|
             +\sum_{e\in F_n\cap\partial p}|F_e^\flat-B_e|.
\]
After squaring, dividing by two and
summing with the couplings, the
additional term is at most
\[
 20(1+\tau)\beta
                   \sum_{e\in F_n}|F_e^\flat-B_e|^2 .
\]
As in V46, this uses at most four
varying edges per rest face and at
most five rest faces per varying edge.
The displayed sum equals
\(\sum_e|g_xB_eg_y^{-1}-I|^2\),
which is at most \(K_n^\flat\).
Since \(20(1+\tau)<22\), the first
line follows.

V47 gives \(T_{n,J}^*\le22\mathfrak H_n\).
Also \(\beta R_n^2\le2\mathfrak H_n\).
Equation~\eqref{eq:f15v48:K} and
\(\pi<4\) imply
\[
             \beta K_n^\flat
                      \le16200d_n n^4\mathfrak H_n .
\]
Substitution proves the second line.
\end{proof}

Thus the flat repair is gauge compatible
and paid by original exterior curvature.
The price \(\mathfrak c_n=O(n^7)\) has
not disappeared. For example, the
selected homogeneous law of V47 gives
only the immediately justified moment
\begin{equation}
 \log\mathbb E_\mu
          e^{T_{n,\beta}^\flat/(48\mathfrak c_n)}
       \le\frac56|\mathfrak L_n|\log2 .
 \label{eq:f15v48:small-moment}
\end{equation}
This follows from
\eqref{eq:f15v47:plaquette-CGF}.
It does not supply the unit-coefficient
exponential repair moment needed to
delete a factor \(e^{T_{n,J}^\flat}\)
from another probability estimate.

\subsection{A certificate for V43's actual coherent-fit event}

Let
\begin{equation}
 \begin{split}
 \varepsilon_n&=\frac1{256n},\qquad
 u_n=\frac1{1024n},\\
 \delta_n^2&=\frac1{518400\pi^2 d_n n^5},\\
 G_n^Y&=\{R_n(Y)\le\delta_n\}.
 \end{split}
 \label{eq:f15v48:thresholds}
\end{equation}
The event \(G_n^Y\) is measurable using
the unchanged exterior alone.
On it,
\[
              K_n^\flat(Y)\le\varepsilon_n/4.
\]
This is the exact constant identity
\[
 \frac{2025\pi^2}{4}d_n n^4\delta_n^2
                                      =\frac1{1024n}.
\]
It follows from \eqref{eq:f15v48:fit}
that the actual event
\(\Gamma_n=\{\mathcal E_n(F)\le\varepsilon_n\}\)
obeys the deterministic inclusion
\begin{equation}
      \Gamma_n^c
        \subseteq
          \{E_n^{\rm out}>u_n\}\ \cup\ (G_n^Y)^c .
 \label{eq:f15v48:failure}
\end{equation}
This holds for every compact field and
every source array. It is not yet a
probability estimate under that array.

\begin{theorem}[Joint coherent-fit failure bound in the selected homogeneous law]
\label{f15v48:joint}
Take exactly the selected law and
cell-source hypotheses of
\eqref{eq:f15v47:CGF}.
Define
\begin{equation}
 \begin{split}
 q_{\rm out}
     &=\exp\left\{-\frac{\beta}{24576\pi^2 n}
                                +\frac56d_n\log2\right\},\\
 q_{\rm collar}
     &=32|\mathfrak L_n|
                  \exp\left\{-\frac{\beta\delta_n^2}{16}\right\},\\
 q_{\rm fit}&=\min\{1,q_{\rm out}+q_{\rm collar}\}.
 \end{split}
 \label{eq:f15v48:q}
\end{equation}
For any \(k\) translated boxes of common
side \(n\) whose sets of positive-plaquette
bases in
\(\mathcal P_{n,i}^+\cup\mathfrak L_{n,i}\)
are pairwise disjoint,
\begin{equation}
             \mu\left(\bigcap_{i=1}^k\Gamma_{n,i}^c\right)
                                      \le q_{\rm fit}^{\,k}.
 \label{eq:f15v48:joint}
\end{equation}
No independence of boxes, centers,
gauges or exterior fields is assumed.
\end{theorem}

\begin{proof}
For the outward-discrepancy alternative
in \eqref{eq:f15v48:failure}, V47's
joint moment gives a cost
\[
       \exp\{-\beta u_n/(24\pi^2)+(5d_n/6)\log2\}
                                     =q_{\rm out}.
\]
For a collar failure choose one witnessing
\(q\in\mathfrak L_n\).
Then
\(X_q=\beta f(\alpha(U_q))>\beta\delta_n^2/2\).
The inherited estimate
\(\mathbb E_\mu e^{X_q/8}\le32\)
gives the one-witness cost
\(32e^{-\beta\delta_n^2/16}\).

These separate bounds alone would not
justify a product estimate. To prove
the joint statement, choose, for each
box, either the outward alternative
or a particular collar witness.
For every outward alternative use the
nonnegative exponent
\(\beta E_n^{\rm out}/(24\pi^2)\),
bounded by \(\sum_{\mathcal P_n^+}X_p/48\).
After \(X_p\le4T_{\operatorname{base}(p)}\),
its cell coefficients are \(m_x/12\le1/2\).
For a collar witness use \(X_q/8\),
whose one-cell coefficient is \(1/2\).
The hypothesis on bases prevents
coefficients from different boxes from
adding at the same cell.
Thus one application of the joint CGF
\eqref{eq:f15v47:CGF} bounds the
intersection for this entire choice
by the product of the costs just stated.

Summing the witness in a collar box
costs \(|\mathfrak L_n|\).
Summing the two alternatives in each
box gives
\((q_{\rm out}+q_{\rm collar})^k\).
The trivial bound by one supplies
the cap in \eqref{eq:f15v48:q}.
\end{proof}

In particular
\[
 q_{\rm collar}
       \le704d_n
          \exp\left\{-\frac{\beta}
                  {8294400\pi^2d_n n^5}\right\}.
\]
Along the selected homogeneous sequence,
the bound tends to zero whenever
\begin{equation}
                    n^8\log(n+2)=o(\beta),
 \label{eq:f15v48:growing}
\end{equation}
subject to the geometric embedding.
To verify this, \(d_n\le24n^3\),
so the negative collar exponent has
magnitude at least a constant times
\(\beta/n^8\), which dominates the
logarithmic prefactor under
\eqref{eq:f15v48:growing}.
The outward exponent has negative
term of order \(\beta/n\) and
positive term of order \(n^3\);
the same condition also implies
\(\beta/n^4\to\infty\).
For fixed \(n\), both assertions follow
directly as \(\beta\to\infty\).
Powers of \(\log\beta\) are examples
of admissible growing local sides.

\subsection{The actual remaining source and normalization issues}

The theorem gives a genuine joint
probability bound for V43's
coherent-fit event in the specified
homogeneous selected law.
The new exterior-only certificate
also supplies a fixed conditional
frame and the full-source deterministic
normalizer comparison.
These are compatible statements,
but their hypotheses are not identical.

In particular, the joint cell CGF
used in the last theorem has not
been proved at every inhomogeneous
physical tilt \(\mu_J\).
The deterministic inclusion
\eqref{eq:f15v48:failure} does not
change that fact. Neither positivity
of the couplings nor the
full-source conditional normalizer
comparison licenses replacing
\(\mu\) by \(\mu_J\) in
\eqref{eq:f15v48:joint}.
The source-normalization and
frustration difficulties recorded
in V25--V38 remain relevant.

There is a separate size cost.
The flat-reference comparison retains
\(\beta K_n^\flat\), and its
sampled-curvature bound grows as
\(n^7\). The proved small exponential
moment \eqref{eq:f15v48:small-moment}
has a correspondingly shrinking
source coefficient. It must not
be promoted to a unit-coefficient
moment or a scale-uniform clustering
estimate.

The next upstream step is to use
the exact exterior-only certificate
in an actual-source change of law
with a paid normalization, or improve
the repair comparison enough that
its required source lies inside a
proved window. Repeating a
homogeneous probability estimate
and calling it full-source entry
would be circular. Interacting
continuum construction and the
physical Hamiltonian gap remain
additional, unproved requirements.

\subsection{Checks and exact preservation}

The diagnostic is
\begin{center}\small
\path{scripts/verify_ym_f15_v48_exterior_flat_reference.py}.
\end{center}
It reuses the unchanged V47 geometry
and quaternion helpers; importing
that module also runs the V47 checks.
The V41 root-path convention is
reproduced, while its boundary
square-homotopy proof remains an
explicit inherited input.

The V48 checks include 2,880
exterior-only path/length tests,
2,520 repaired-boundary curvature
bounds, 1,008 flat-reference/chord
fit checks, two nontrivial
pure-gauge zero-cost recoveries,
and six tests that change every
sampled boundary link while leaving
the exterior frame and centers
unchanged.
There are 960 root-frame covariance
identities, 1,008 flat-center
covariance identities and 18
full-source repaired-action checks.
Ninety-three exact threshold/tail
identities, 31 support/coefficient
checks and seven mixed cell-source
coefficient checks supplement them.
All diagnostic arithmetic is
rational or integer. The global
probability bounds follow from
the analytic proofs and the
declared inherited CGF, not from
these finite configurations.

The script has 7,631 bytes and SHA--256
\begin{center}\scriptsize\ttfamily
D7CFDABFF953248F67926F228AF84F8373B57A35C2EEC7CE47A468AEA58487B0.
\end{center}
The V47 predecessor has 1,107,240 bytes and SHA--256
\begin{center}\scriptsize\ttfamily
53A0BB6066295747FF86CA9F9C11AA0320FA5B265D2C4CDF0AF5776FC26FF8AA.
\end{center}
Removing this section and reversing
the title increment recovers V47
exactly. V48 replaces V47 as
the sole active main source.
The 14 protected inputs and
90 earlier diagnostics remain
unchanged. Only \path{.tex}
files belong in \path{latex_notes}.
Build products belong in
\begin{center}\small
\path{artifacts/ym_f15_v48_texcheck/}.
\end{center}
The next companion checkpoint is V50.

\begin{firewall}
V48 constructs an exterior-only
gauge-covariant flat reference,
pays its fit and repair costs,
and obtains a full-source
conditional compact normalizer
comparison. It proves a joint
coherent-fit failure bound only
in the stated selected homogeneous
law. It does not prove the
corresponding full-source marginal
bound, scale-uniform clustering,
an interacting continuum or the
physical Yang--Mills mass gap.
\end{firewall}
\section{V49: reflection-compatible source backgrounds and the cell-lift obstruction}
\label{f15v49:section}

V48's joint coherent-fit bound uses a selected
homogeneous probability law. Its deterministic
conditional normalizer comparison has a larger
source domain. This section investigates exactly
how far the probability argument extends by
retaining site reflection positivity of the
actual background law.

There is a genuine extension: a 24-parameter
family of extensively supported plaquette
couplings admits a direct joint cell moment
bound and the V48 fit estimate at every
coupling in a specified range, without the
previous coupling selection. There is also
a decisive restriction. The intersection
of that family with the program's symmetric
cell-source lift contains only homogeneous
couplings. Nonconstant cell sources in the
intersection are invisible to the action.
The extension therefore does not solve the
missing general physical-source estimate.

The V48 geometry and f14 V98--V99 pressure
and scalar-label arguments are inherited,
not new results of this section. The new
work checks their actual background-law
domain, proves the exact intersection with
the original source lift, and derives the
consequent source-uniform fit bound on that
domain. V45 remains the companion checkpoint;
the next review is V50.

For context, the general closed-block
chessboard estimate appears in M.~Biskup,
\emph{Reflection Positivity and Phase Transitions
in Lattice Spin Models}, Section~4.3,
\href{https://www.math.ucla.edu/~biskup/PDFs/papers/Prague-School.pdf}
{Theorem~4.8}. Its reflection-positivity
hypothesis must be checked for the actual
law. The verification below uses compact
link variables and the archived scalar-label
proof; it imports no spin infrared bound
or non-Abelian correlation monotonicity.

\subsection{The exact coefficient class fixed by every site reflection}

Let \(M=2^j\ge4\), \(\Lambda_M=(\mathbb Z/M\mathbb Z)^4\),
and \(V=M^4\). Positive plaquettes are
\(p=(x;\mu,\nu)\), \(\mu<\nu\), with
\(s_p=f(\alpha(U_p))=1-\tfrac12\Tr U_p\).
For a real plaquette array \(J\), write
\begin{equation}
 Z_M(J)=\int e^{-\sum_pJ_ps_p}\prod_e dU_e,
 \qquad
 d\mu_J=Z_M(J)^{-1}e^{-\sum_pJ_ps_p}\prod_e dU_e.
 \label{eq:f15v49:law}
\end{equation}
Compactness makes this definition meaningful
even for real arrays not everywhere positive.
The background arrays used in the bounds below
are strictly positive.

If \(\{\lambda,\kappa\}\) is the increasing
complement of \(\{\mu,\nu\}\), define
\begin{equation}
 \begin{split}
 \mathcal R_{M,\beta,\tau}
   =\bigl\{J: J_{(x;\mu,\nu)}
        &=\beta a_{\mu\nu,\epsilon},\quad
          \epsilon=(x_\lambda\bmod2,x_\kappa\bmod2),\\
        &a_{\mu\nu,\epsilon}\in[1-\tau,1+\tau]\bigr\}.
 \end{split}
 \label{eq:f15v49:class}
\end{equation}
There are six orientation pairs and four
transverse parities for each, hence 24
independent real coefficients when \(\tau>0\).
They can differ over the entire volume.

\begin{proposition}[Classification of invariant coefficient arrays]
\label{f15v49:classification}
A plaquette coefficient array is invariant
under all integer-coordinate site-plane
reflections if and only if, for each
orientation, it depends only on the two
transverse coordinate parities. Intersecting
this vector space with
\((1-\tau)\beta\le J_p\le(1+\tau)\beta\)
gives exactly \(\mathcal R_{M,\beta,\tau}\).
\end{proposition}

\begin{proof}
Reflection in coordinate \(i\) through
\(x_i=c\) sends a geometric plaquette
to the positive plaquette with unchanged
orientation pair and base coordinate
\begin{equation}
 x_i'=
 \begin{cases}
  2c-x_i-1\pmod M,&i\in\{\mu,\nu\},\\
  2c-x_i\pmod M,&i\notin\{\mu,\nu\}.
 \end{cases}
 \label{eq:f15v49:reflection}
\end{equation}
Other base coordinates are unchanged.
The holonomy may be inverted or conjugated;
the real \(\SU(2)\) trace is unchanged.

Composing reflections at neighboring site
planes translates the base by two units.
In a parallel coordinate a single reflection
also exchanges the two parities. The orbit
in that coordinate is therefore the whole
cyclic coordinate set. In a transverse
coordinate reflections preserve parity,
and the two-unit translations are transitive
within each parity. The different coordinate
actions commute. The orbits for one orientation
are consequently precisely its four transverse
parity classes. This proves necessity and
sufficiency, including independence of the
two parallel coordinates.
\end{proof}

This classifies coefficient arrays invariant
under the specified reflections. It does not
classify every possible reflection-positive
gauge model or exclude other reflection schemes.

\subsection{Reflection positivity in the actual patterned law}

For a coordinate reflection \(\theta\), let
\(E_0\) be tangent links in its two fixed
site planes, and \(E_+,E_-\) the remaining
links in the two closed halves. Every
elementary plaquette lies in one closed
half or entirely in a fixed plane.
In particular no elementary face intersects
both open halves. For \(J\in\mathcal R_{M,\beta,\tau}\),
the action decomposes as
\[
             S_J=S_{J,+}+\theta S_{J,+}+S_{J,0}.
\]
The reflection fixes each link of \(E_0\)
and is a Haar-preserving bijection between
the other link sets, including inversions
of links whose orientation reverses.
For any bounded complex function \(F\)
on \(E_+\cup E_0\), conditional integration
therefore gives the exact identity
\begin{equation}
 \mathbb E_J[F\,\theta\overline F]
   =\frac1{Z_M(J)}\int e^{-S_{J,0}}
       \left|\int e^{-S_{J,+}}F
                       \prod_{e\in E_+}dU_e\right|^2
                       \prod_{e\in E_0}dU_e\ge0.
 \label{eq:f15v49:RP}
\end{equation}
This is a proof under \(\mu_J\), not a
change of law from a homogeneous measure.
It involves no volume-dependent
Radon--Nikodym comparison. It uses site
planes, not an unsupported factorization
through bond-plane crossing plaquettes.

Retain the unweighted symmetric cell
energy of \eqref{f15v1:cell}:
\begin{equation}
 T_x=\frac\beta4\sum_{\mu<\nu}
       \sum_{\epsilon\in\mathcal C_{\mu\nu}}
                   s_{(x+\epsilon;\mu,\nu)},\qquad
 \mathcal C_{\mu\nu}
   =\{\epsilon\in\{0,1\}^4:
                         \epsilon_\mu=\epsilon_\nu=0\}.
 \label{eq:f15v49:cell}
\end{equation}
Thus \(\sum_xT_x=\beta\sum_ps_p\), independently
of the background coefficients. Reflection
takes \(T_x\) to the same observable on the
reflected closed unit cell. The weights in
this observable are all \(1/4\); they have
not been replaced by the patterned couplings.

\begin{proposition}[Actual-background scalar cell bound]
\label{f15v49:chessboard}
For \(J\in\mathcal R_{M,\beta,\tau}\) and any
real cell array \(b\),
\begin{equation}
 \log\mathbb E_J e^{\sum_xb_xT_x}
   \le\frac1V\sum_x
          \log\frac{Z_M(J-\beta b_x\mathbf1)}{Z_M(J)}.
 \label{eq:f15v49:chessboard}
\end{equation}
In each numerator the scalar \(\beta b_x\)
is subtracted from every plaquette coupling.
It is not a shift only on the cell indexed by \(x\).
\end{proposition}

\begin{proof}
Use the scalar-label proof of f14 V99 with
the actual base law \(\mu_J\). To make its
applicability explicit, put
\[
 c_J(t)=\left(\mathbb E_J e^{t\sum_xT_x}\right)^{1/V},
 \qquad
 \mathcal G_J(b)
   =\frac{\mathbb E_J e^{\sum_xb_xT_x}}
                    {\prod_xc_J(b_x)}.
\]
For a site-plane split, Cauchy--Schwarz
in \eqref{eq:f15v49:RP} and cell covariance
give
\(\mathcal G_J(b)^2\le
\mathcal G_J(b^+)\mathcal G_J(b^-)\),
where the superscripts duplicate either
half's scalar labels by reflection.
Label multiplicities cancel between the
two denominators. The background array
\(J\) remains unchanged by this operation.

Among arrays with labels from the finite
set appearing in \(b\), maximize
\(\mathcal G_J\), then minimize the number
of disagreeing nearest-neighbor cell pairs.
If a maximizing array is nonconstant,
choose a site plane crossing a disagreement.
Write the disagreement counts as
\(n_++n_-+n_\times\), with \(n_\times>0\).
The duplicated arrays have respectively
\(2n_+\) and \(2n_-\) disagreements: pairs
across the two reflection boundaries now
agree. The reflection inequality and
maximality force both duplicated arrays
to be maximizers. At least one has fewer
disagreements, a contradiction.
For a constant label \(t\),
\(\mathcal G_J=1\) by definition.
Hence \(\mathcal G_J(b)\le1\).

Finally
\(c_J(t)^V=Z_M(J-\beta t\mathbf1)/Z_M(J)\).
Taking logarithms proves the assertion.
Shared boundary links are precisely the
conditioned variables in
\eqref{eq:f15v49:RP}; no independence is
required. Unit-translation invariance of
the patterned plaquette coefficients was
not used in the scalar-label proof.
\end{proof}

\subsection{A full coupling range with an explicit moment constant}

From now on assume
\(\beta\ge4\), \(0\le\tau\le1/32\), and set
\begin{equation}
 K_{M,\beta,\tau}
  =62+\frac92\log\frac{1+\tau}{1/2-\tau}
              +\frac6M\log((1+\tau)\beta).
 \label{eq:f15v49:K}
\end{equation}

\begin{theorem}[Uniform joint moment under every background in the class]
\label{f15v49:CGF}
For every \(J\in\mathcal R_{M,\beta,\tau}\) and
every cell array \(0\le b_x\le1/2\),
\begin{equation}
 \log\mathbb E_J e^{\sum_xb_xT_x}
           \le2K_{M,\beta,\tau}\sum_xb_x.
 \label{eq:f15v49:CGF}
\end{equation}
There is no restriction on the support of \(b\).
If \((\log\beta)/M\le L\), then
\begin{equation}
                    K_{M,\beta,\tau}\le66+6L.
 \label{eq:f15v49:uniform-K}
\end{equation}
No selected coupling sequence is required.
\end{theorem}

\begin{proof}
For \(0\le t\le1/2\), set
\[
 F_J(t)=\frac1V\log
                   \frac{Z_M(J-\beta t\mathbf1)}{Z_M(J)}.
\]
This is convex, with \(F_J(0)=0\).
Since every \(s_p\ge0\), comparison of
the positive integrands gives
\begin{equation}
 \frac{Z_M(J-\beta t\mathbf1)}{Z_M(J)}
   \le\frac{Z^{\mathrm W}_{M,(1-\tau-t)\beta}}
                  {Z^{\mathrm W}_{M,(1+\tau)\beta}}.
 \label{eq:f15v49:pressure-comparison}
\end{equation}
The smaller coupling at \(t=1/2\)
is at least \(15\beta/32>1\).
Apply the pressure-ratio lemma of f14 V98
with larger coupling \((1+\tau)\beta\)
and ratio of smaller to larger coupling
\((1/2-\tau)/(1+\tau)\). It gives
\(F_J(1/2)\le K_{M,\beta,\tau}\).
Convexity between zero and this endpoint
now yields \(F_J(t)\le2K_{M,\beta,\tau}t\).
Inserting this bound in
\eqref{eq:f15v49:chessboard} proves
\eqref{eq:f15v49:CGF}.

For the last assertion,
\((1+\tau)/(1/2-\tau)\le11/5\),
\(\log(11/5)<4/5\), and
\(\log(1+\tau)\le1/32\).
The strict logarithm bound follows, for
example, from the degree-four positive
Taylor sum for \(e^{4/5}\), which exceeds
\(11/5\). Using \(M\ge4\) gives
\[
 K_{M,\beta,\tau}
  \le62+\frac{18}{5}+\frac3{64}+6L<66+6L.
\]
\end{proof}

The pressure comparison by itself has
an extensive cost even at \(t=0\).
The convex interpolation uses the exact
identity \(F_J(0)=0\), so a zero source
has zero cost in \eqref{eq:f15v49:CGF}.
Neither this interpolation nor the
positive-integrand comparison asserts
monotonicity of correlations in individual
non-Abelian plaquette couplings.

\subsection{Apply the exterior-only certificate in this actual law}

Keep all V47--V48 geometric definitions,
including \(d_n=24n(n-1)^2\),
the outward mirror support \(\mathcal P_n^+\),
the exterior-only collar \(\mathfrak L_n\)
and \(|\mathfrak L_n|\le22d_n\).
The box and its collar must embed as in
those sections. Use their fixed thresholds
\begin{equation}
 u_n=\frac1{1024n},\qquad
 \delta_n^2=\frac1{518400\pi^2d_n n^5},\qquad
 \Gamma_n=\{\mathcal E_n(F)\le1/(256n)\}.
 \label{eq:f15v49:thresholds}
\end{equation}
The deterministic inclusion
\eqref{eq:f15v48:failure} holds unchanged
under \(\mu_J\). It is the union of an
outward-discrepancy failure and an
exterior-curvature witness failure.

\begin{theorem}[Measured fit under patterned nonzero plaquette backgrounds]
\label{f15v49:fit}
Write \(K=K_{M,\beta,\tau}\) and define
\begin{equation}
 \begin{split}
 q_{\rm out}^{J}
   &=\exp\left\{-\frac\beta{24576\pi^2n}
                                      +\frac{Kd_n}{6}\right\},\\
 q_{\rm collar}^{J}
   &=|\mathfrak L_n|
                   \exp\{K-\beta\delta_n^2/16\},\\
 q_{\rm fit}^{J}
   &=\min\{1,q_{\rm out}^{J}+q_{\rm collar}^{J}\}.
 \end{split}
 \label{eq:f15v49:q}
\end{equation}
The superscript records the background-law
application; the bounds themselves are
uniform over \(\mathcal R_{M,\beta,\tau}\).
For translated boxes of common side \(n\ge2\)
whose positive-plaquette base sets in
\(\mathcal P_{n,i}^+\cup\mathfrak L_{n,i}\)
are pairwise disjoint,
\begin{equation}
 \sup_{J\in\mathcal R_{M,\beta,\tau}}
       \mu_J\left(\bigcap_{i=1}^k\Gamma_{n,i}^c\right)
                       \le(q_{\rm fit}^{J})^k.
 \label{eq:f15v49:joint}
\end{equation}
If \(\beta\to\infty\), \((\log\beta)/M\)
stays bounded and \(n^8\log(n+2)=o(\beta)\),
then \(q_{\rm fit}^{J}\to0\), uniformly
over the same background class, subject
to the geometric embedding.
\end{theorem}

\begin{proof}
Only the probability input changes from
V48; its geometric certificate is reused.
Write \(X_p=\beta s_p\), so
\(X_p\le4T_{\operatorname{base}(p)}\).
For an outward alternative, use the
inherited deterministic bound
\[
 \frac{\beta E_n^{\rm out}}{24\pi^2}
                  \le\frac1{48}\sum_{p\in\mathcal P_n^+}X_p.
\]
The resulting cell coefficient at a base
is \(m_x/12\), where \(m_x\le6\) is
the number of mirror plaquettes based
there. Its sum is \(d_n/12\).
Thus \eqref{eq:f15v49:CGF} pays
\(Kd_n/6\), while Markov's threshold
is \(\beta u_n/(24\pi^2)\).
These are exactly the two terms in
\(\log q_{\rm out}^{J}\).

For a specified collar witness \(q\),
\(|U_q-I|>\delta_n\) implies
\(X_q>\beta\delta_n^2/2\).
Use the nonnegative exponent \(X_q/8\),
bounded by \(T_{\operatorname{base}(q)}/2\).
The cell moment bound pays \(K\),
and the threshold pays
\(\beta\delta_n^2/16\).

To justify the joint assertion, select
one of the two alternatives for each
box, with a particular witness in each
collar alternative. Add all their
exponents before applying Markov's
inequality. Disjointness of the specified
base sets prevents coefficients from
different boxes adding at one cell;
every resulting coefficient is at most
\(1/2\). One application of
\eqref{eq:f15v49:CGF} then pays the sum
of the individual costs. Sum over
collar witnesses, then over alternatives.
This gives
\((q_{\rm out}^{J}+q_{\rm collar}^{J})^k\).
The bound by one gives the displayed cap.
There is no independence assumption.

Finally \(K\) is bounded by
\eqref{eq:f15v49:uniform-K} and
\(d_n\le24n^3\). The collar's negative
exponent is at least a positive constant
times \(\beta/n^8\), dominating
\(\log|\mathfrak L_n|+K\) in the stated
regime. The outward negative term is
of order \(\beta/n\), and its positive
term is of order \(n^3\).
The same condition implies
\(\beta/n^4\to\infty\), so this term
also tends to zero. Fixed \(n\) is included.
\end{proof}

This is a probability theorem in the
actual patterned background, not just
a deterministic full-source comparison.
Its restriction to homogeneous couplings
also removes the selected-coupling
hypothesis of V48, at the expense of
the larger explicit moment constant.

\subsection{Audit the intersection with the original physical cell sources}

The program's physical scalar cell source
is a real array \(h=(h_x)\) in the tilt
\(e^{\sum_xh_xT_x}d\mu_{M,\beta}^{\mathrm W}\).
Do not confuse this change of background
with the test array \(b\) in
\eqref{eq:f15v49:CGF}.
By the original lift
\eqref{eq:f15v38:lift}, its induced
plaquette array is
\begin{equation}
 \begin{split}
 (\mathcal A_{\mu\nu}h)(y)
     &=\frac14\sum_{\epsilon\in\mathcal C_{\mu\nu}}h_{y-\epsilon},\\
 J_{(y;\mu,\nu)}(h)
     &=\beta\bigl(1-(\mathcal A_{\mu\nu}h)(y)\bigr).
 \end{split}
 \label{eq:f15v49:lift}
\end{equation}
The common kernel of all six maps is
denoted \(\ker\mathcal A\).

\begin{theorem}[Only homogeneous induced couplings survive the intersection]
\label{f15v49:intersection}
The array \(J(h)\) is invariant under every
integer site-plane reflection if and only if
\begin{equation}
        h=\bar h\mathbf1+h_0,
       \qquad h_0\in\ker\mathcal A,
       \qquad \bar h=V^{-1}\sum_xh_x.
 \label{eq:f15v49:intersection}
\end{equation}
In that case, for every plaquette,
\begin{equation}
       J_p(h)=\beta(1-\bar h),\qquad
       \sum_xh_xT_x=\bar h\,\beta S
                  \quad\hbox{for every link configuration}.
 \label{eq:f15v49:homogeneous-image}
\end{equation}
In particular membership in
\(\mathcal R_{M,\beta,\tau}\) additionally
requires and is implied by \(|\bar h|\le\tau\).
\end{theorem}

\begin{proof}
Use Fourier momenta
\(k_i=2\pi m_i/M\), \(m_i\in\{0,\ldots,M-1\}\),
with transform
\[
              \widehat h(k)=V^{-1}\sum_xh_xe^{-ik\cdot x}.
\]
Equation~\eqref{eq:f15v49:lift} gives
\begin{equation}
 \widehat{\mathcal A_{\mu\nu}h}(k)
   =\widehat h(k)
           \prod_{i\notin\{\mu,\nu\}}
                                \frac{1+e^{-ik_i}}2.
 \label{eq:f15v49:multiplier}
\end{equation}
An invariant orientation array is
independent of the parallel coordinates
and two-periodic in the transverse ones,
by Proposition~\ref{f15v49:classification}.
Its Fourier support must therefore satisfy
\(k_\mu=k_\nu=0\) and
\(k_i\in\{0,\pi\}\) for the transverse
coordinates. Every nonzero momentum in
that allowed support has a transverse
\(\pi\) coordinate, so its multiplier
in \eqref{eq:f15v49:multiplier} vanishes.
All other nonzero momenta are excluded
by invariance. Hence the induced array
for each orientation is constant.
Its zero Fourier coefficient is
\(\widehat h(0)=\bar h\), the same
for every orientation.

It follows that
\(\mathcal A_{\mu\nu}(h-\bar h)=0\)
for all six orientations. Conversely
this kernel condition makes every
induced coupling equal to
\(\beta(1-\bar h)\).
The identity for the action follows
by regrouping the four cell weights
of each plaquette. The final interval
condition is then exactly
\(|\bar h|\le\tau\).
\end{proof}

\begin{proposition}[The invisible cell-source modes are explicit]
\label{f15v49:kernel}
The Fourier support of \(\ker\mathcal A\)
consists exactly of momenta with at least
three coordinates equal to \(\pi\).
Its real dimension is \(4M-3\).
Thus the real cell-source space satisfying
the invariance condition has dimension
\(4M-2\), but its induced Wilson laws
have only the single homogeneous
coupling parameter.
\end{proposition}

\begin{proof}
If at least three coordinates equal
\(\pi\), every two-coordinate complement
contains one of them, so every product
in \eqref{eq:f15v49:multiplier} vanishes.
If at most two coordinates equal
\(\pi\), choose an orientation pair
containing them all, adding arbitrary
coordinates if needed. Neither of the
two transverse factors then vanishes.
This proves the exact common-kernel
criterion mode by mode.

For exactly three \(\pi\) coordinates,
choose the remaining coordinate in
four ways and its non-\(\pi\) frequency
in \(M-1\) ways. Add the all-\(\pi\)
mode. The total is
\(4(M-1)+1=4M-3\).
The set of these momenta is closed
under negation; conjugate mode pairs
give the same count of real dimensions,
with self-conjugate modes counted once.
None is the constant mode, which adds
one dimension in
\eqref{eq:f15v49:intersection}.
\end{proof}

For example, a nonzero multiple of
\((-1)^{x_1+x_2+x_3}\) is an invisible
cell source, not a new background law.
A source supported at a single cell,
on the other hand, has nonzero Fourier
coefficients outside this kernel and
the constant mode, so its induced
couplings are not in the invariant
class. The latter example can have
arbitrarily small amplitude; smallness
alone does not restore the symmetry.

\subsection{What this changes upstream, and what it does not}

The 24-parameter plaquette backgrounds
provide actual inhomogeneous measures
with controlled local energy moments
and coherent-fit failure probabilities.
This is a useful testing domain for the
conditional constructions of V46--V48.
However, it is mostly outside the
program's specified cell-source image.
Replacing that source image by the
larger plaquette class would change
the question, not close its source
normalization requirement.

Within the present all-site-reflection
method, the obstruction is now exact:
every visible nonconstant physical
cell source breaks the required
coefficient invariance. This is an
obstruction to this direct proof
route, not a counterexample to the
desired general-source bound. An
argument using local comparisons,
paid noninvariant normalizers or a
different decomposition is still
possible, but must supply its own
quantitative estimate. Repeatedly
applying the same reflection theorem
to an unchanged homogeneous law
cannot supply that estimate.

There is also no improvement here
of the \(O(n^7)\) flat-repair debit
in \eqref{eq:f15v48:flat-action},
no unit-coefficient exponential
repair moment, and no proof of
physical long-distance clustering.
Interacting continuum construction
and a positive physical Hamiltonian
gap remain separate unproved parts
of the original objective.

\subsection{Exact diagnostics and predecessor preservation}

The new diagnostic is
\begin{center}\small
\path{scripts/verify_ym_f15_v49_reflection_source_class.py}.
\end{center}
It reuses unchanged V47 quaternion
helpers and runs that diagnostic on
import. In exact rational arithmetic,
it checks 24,576 noncommuting plaquette
reflection identities, 4,096 closed-cell
identities, 16 invariances of a genuinely
24-valued weighted action, and all
24 geometric reflection orbits on
the side-four torus.

Sixteen period-two cell modes verify
the exact lift identity: one is
constant, five are invisible, and
ten give visible reflection-breaking
coefficients. A single-cell source
is also checked. Integer Fourier
tests on side lengths 4, 8 and 16
classify 69,888 common-kernel modes
and verify the dimension \(4M-3\)
at each size. Pressure and mixed-witness
constant checks supplement these
tests. The analytic proofs, not the
finite tests, justify the probability
bounds and uniform limits.

The script has 7,673 bytes and SHA--256
\begin{center}\scriptsize\ttfamily
A7976833949805644C90F4573E05E47ECCAF6829BB73BB3F2D001397896FF8D8.
\end{center}
The V48 predecessor has 1,126,942 bytes and SHA--256
\begin{center}\scriptsize\ttfamily
749113312D7BA78F113FE9CC7569468C323E124D6CFC1E616EADC9116B1E0F83.
\end{center}
Removing this final section and reversing
the title increment recovers V48 exactly.
V49 replaces V48 as the sole active
main source. The 14 protected inputs
and 91 earlier diagnostics remain
unchanged. Only \path{.tex} files belong
in \path{latex_notes}; build products
are directed to
\begin{center}\small
\path{artifacts/ym_f15_v49_texcheck/}.
\end{center}
The next companion checkpoint is V50.

\begin{firewall}
V49 proves direct joint energy and
coherent-fit bounds for a specified
24-parameter reflection-invariant
plaquette background class, at all
couplings in its stated range. It
also proves that this class gives
only homogeneous laws when restricted
to the original symmetric cell-source
lift. It does not establish the
missing general physical-source bound,
an interacting continuum or the
physical Yang--Mills mass gap.
\end{firewall}
\section{V50: all-source pressure and coherent entry through translation orbits}
\label{f15v50:section}

V49 proves that retaining all site reflections
does not reach a visible nonconstant physical
cell source. This section goes upstream to
the part of its pressure argument that never
used reflection positivity. That argument
controls the total action under every allowed
inhomogeneous Wilson law. Combining it with
the actual translation symmetries of that
law gives a local first-moment bound when
the translation group has bounded index.

The resulting coherent-entry theorem covers
genuinely nonconstant physical cell sources,
including a two-class checkerboard source
excluded by V49's reflection class. For an
arbitrary aperiodic source, the conclusion
is instead a spatially averaged entry bound
and a macroscopic defect-density tail.
Neither conclusion is the prescribed-family
exponential bound or exceptional-energy
moment still needed in the full program.

\subsection{V50 companion checkpoint and nonduplication review}

The active V49 source was revalidated at
its final hash. A fresh inventory found no
newer relevant companion files. The latest
conclusions of the following notes were
reread, together with the supplied
\path{suggestions.txt}; all their protected
hashes are unchanged.

\begin{itemize}
\item SM continuum V72 concerns a fixed-mode,
fixed-time many-copy limit. It does not
supply a physical ultraviolet Wilson-law
estimate or a finite-species continuum.
\item NS stability V26 contains pointwise
rank-one kernel formulas. Its references
to the planned V27 certification are not
imported as completed bounds here.
\item Shell mixing V16 assembles a measured
flat-branch one-loop coefficient. Its
general noncommuting calculation and a
common nonperturbative integral remainder
remain separate from compact entry.
\item Parallel YM V5 retains its warning
that global charge tightness and local
defect density are different requirements.
\item SM source selection V31 is conditional
on historical word-native typing. Native
stationarity V34 leaves its moving-fibre
first-escape entries uncomputed. Exact-action
Einstein bridge V24 gives a nonzero first
derivative, not an upstream nonlinear root.
None provides the missing YM marginal bound.
\end{itemize}

The suggestions correctly emphasize keeping
the compact-entry problem distinct from
the measured perturbative coefficient
ledger. They are mathematical suggestions,
not instructions imported from a document.
We retain the original physical cell source
and full compact Wilson density here.
The next companion checkpoint is V55.

The nonduplication review also reread f12
V25's arbitrary-exterior auxiliary-fibre
defect-density theorem and thin-carrier
warning, and V37's actual-background
normalization criterion. A generic bulk
repair with a surface debit is not being
repeated. The object below is the actual
full Wilson law with arbitrary plaquette
coefficients, and then its finite-index
translation orbits. V48's boundary geometry
and its exterior-only certificate are
inherited unchanged.

For literature context, see Faria da Veiga and O'Carroll's
\href{https://www.sciencedirect.com/science/article/pii/S0034487725000357}
{published stability paper}.
Its abstract describes partition-function bounds and
multireflection bounds for plaquette
generating functions. We do not infer
an arbitrary-background local source
theorem from that description. The
needed inequalities below follow from
the explicit pressure bound already
proved in f14 V95--V98.

\subsection{The total-energy bound does not require a reflection-invariant source}

Use \eqref{eq:f15v49:law}, with an arbitrary
array satisfying
\begin{equation}
 \beta\ge4,\qquad 0\le\tau\le1/32,\qquad
 (1-\tau)\beta\le J_p\le(1+\tau)\beta.
 \label{eq:f15v50:domain}
\end{equation}
Here any integer \(M\ge4\) is permitted;
no dyadic reflection replication is used.
Put \(V=M^4\), \(S=\sum_ps_p\),
\(X_p=\beta s_p\), and retain the unweighted
\(T_x\) from \eqref{eq:f15v49:cell}.
In particular \(\sum_xT_x=\beta S\).
Write
\[
 K=62+\frac92\log\frac{1+\tau}{1/2-\tau}
                    +\frac6M\log((1+\tau)\beta).
\]
The estimate \(K\le66+6L\) when
\((\log\beta)/M\le L\) is unchanged.

\begin{theorem}[Actual-source total-action budget]
\label{f15v50:total}
For every array in \eqref{eq:f15v50:domain},
\begin{equation}
 \log\mathbb E_J e^{t\beta S}\le2KtV
           \quad(0\le t\le1/2),\qquad
 \mathbb E_J\beta S\le2KV.
 \label{eq:f15v50:total}
\end{equation}
In particular, for \(a\ge0\),
\begin{equation}
           \mu_J\{\beta S\ge2V(K+a)\}\le e^{-aV}.
 \label{eq:f15v50:energy-tail}
\end{equation}
The source can be extensive, nonconstant
and without any nontrivial translation
or reflection symmetry.
\end{theorem}

\begin{proof}
For \(0\le t\le1/2\), comparison of
positive integrands gives
\[
 \mathbb E_J e^{t\beta S}
   =\frac{Z_M(J-t\beta\mathbf1)}{Z_M(J)}
   \le\frac{Z^{\mathrm W}_{M,(1-\tau-t)\beta}}
                  {Z^{\mathrm W}_{M,(1+\tau)\beta}}.
\]
This is exactly the pressure comparison
in V49, but its proof used only
\(s_p\ge0\) and the coefficient bounds.
It did not use the classification or
reflection positivity of that section.
The f14 V98 pressure-ratio lemma applies
for every integer \(M\ge3\), and the
smaller coupling at \(t=1/2\) exceeds one.
It bounds the logarithm at that endpoint
by \(KV\).

The exact left logarithm is convex in
\(t\) and is zero at \(t=0\).
Its chord from zero to \(1/2\) proves
the first inequality. Finite-volume
compactness permits differentiation at
zero from the right, giving the mean
bound. Markov's inequality at \(t=1/2\)
gives \eqref{eq:f15v50:energy-tail}.
\end{proof}

The distinction from V49 is important:
this theorem is an all-source bound
for the total action, not an all-source
joint local exponential-moment bound.
Applying the scalar-cell chessboard
argument to this new background would
again require the symmetry that V49
proved absent in general.

\subsection{Finite-index translation symmetry gives local means}

For \(z\in\Lambda_M\), let \(\sigma_z\)
translate link variables and plaquette
coefficients by \(z\), with orientation
unchanged. Define the actual coefficient
stabilizer and its index by
\begin{equation}
 H_J=\{z:\ J_{(x+z;\mu,\nu)}=J_{(x;\mu,\nu)}
                          \text{ for all }x,\mu,\nu\},
 \qquad \iota_J=V/|H_J|.
 \label{eq:f15v50:index}
\end{equation}
This is a subgroup of the translation
group. Its index counts coefficient
translation classes, not plaquette
orientations or independent samples.

\begin{proposition}[Pointwise thermal means with the true orbit cost]
\label{f15v50:orbit}
For every plaquette \(p\) and cell \(x\),
\begin{equation}
                \mathbb E_J X_p\le2K\iota_J,
       \qquad \mathbb E_J T_x\le2K\iota_J.
 \label{eq:f15v50:local-means}
\end{equation}
More generally, if a nonnegative
translation-covariant local field \(Q_x\)
has \(V^{-1}\sum_x\mathbb E_JQ_x\le b\),
then \(\mathbb E_JQ_x\le\iota_J b\)
at every location.
\end{proposition}

\begin{proof}
For \(z\in H_J\), translation is a
Haar-preserving change of variables
leaving the action unchanged. Hence
\(\mu_J\) is invariant under \(H_J\).
The action of translations on cell
bases, or on plaquette bases of a
fixed orientation, is free. Each orbit
therefore has \(|H_J|\) elements.
For cells, nonnegativity and
\eqref{eq:f15v50:total} imply
\[
 |H_J|\mathbb E_J T_x
    =\sum_{z\in H_J}\mathbb E_JT_{x+z}
    \le\mathbb E_J\sum_yT_y\le2KV.
\]
The proof for \(X_p\) is identical,
using that its single-orientation
orbit sum is bounded by \(\beta S\).
The same orbit argument proves the
general assertion.
\end{proof}

If the coefficients have coordinate
periods \(P_1,\ldots,P_4\), each dividing
\(M\), then \(\iota_J\le\prod_iP_i\).
This bound is often strictly better
than replacing all periods by \(M\).
For a fully aperiodic coefficient array,
\(\iota_J=V\); the proposition then
does not give a volume-uniform
prescribed-location estimate.

\subsection{Count the actual coherent-entry certificates over all translations}

Let \(n\ge2\) and assume \(n+3\le M\),
so the box and its unit collar embed
injectively at every translated base.
Use V48's exterior-only construction,
with
\[
 d_n=24n(n-1)^2,\quad u_n=\frac1{1024n},\quad
 \delta_n^2=\frac1{518400\pi^2d_n n^5}.
\]
For the box based at \(x\), define
\begin{equation}
 \begin{split}
 C_{n,x}
   &=\{E_{n,x}^{\rm out}>u_n\}
                       \cup\{R_{n,x}(Y)>\delta_n\},\\
 D_n(U)&=\sum_{x\in\Lambda_M}\mathbf1_{C_{n,x}}.
 \end{split}
 \label{eq:f15v50:count}
\end{equation}
The event \(C_{n,x}\) is a certificate
failure, not the assertion that the
coherent fit itself necessarily fails.
V48 proves
\(\Gamma_{n,x}^c\subseteq C_{n,x}\).
It also proves that on \(C_{n,x}^c\)
the explicit exterior-only frame obeys
\(E_g(F,Y)\le1/(256n)\).

Define the deterministic constants
\begin{equation}
 \begin{split}
 A_n^Y&=\frac{44d_n}{\delta_n^2},\\
 A_n&=\frac{\pi^2d_n}{2u_n}+A_n^Y
       =\pi^2\bigl(512d_n n+22809600d_n^2n^5\bigr).
 \end{split}
 \label{eq:f15v50:A}
\end{equation}
They are deliberately conservative.

\begin{proposition}[A deterministic convolution ledger]
\label{f15v50:ledger}
For every compact field configuration,
\begin{equation}
 \begin{split}
 D_n(U)&\le A_n S(U),\\
 \sum_x\mathbf1_{\{R_{n,x}(Y)>\delta_n\}}
             &\le A_n^Y S(U),\qquad
 A_n\le2^{34}\pi^2n^{11}.
 \end{split}
 \label{eq:f15v50:ledger}
\end{equation}
All translated boxes may overlap. No
separation or independence assumption
is needed for this deterministic estimate.
\end{proposition}

\begin{proof}
V47's angle-chord comparison gives
\(E_{n,x}^{\rm out}\le(\pi^2/2)
\sum_{p\in x+\mathcal P_n^+}s_p\).
Also \(|U_p-I|^2=2s_p\).
The union bound for the two indicators,
used pointwise, therefore yields
\begin{equation}
 \mathbf1_{C_{n,x}}
  \le\frac{\pi^2}{2u_n}
             \sum_{p\in x+\mathcal P_n^+}s_p
      +\frac2{\delta_n^2}
             \sum_{p\in x+\mathfrak L_n}s_p.
 \label{eq:f15v50:pointwise-ledger}
\end{equation}
The second summand alone bounds the
exterior-curvature indicator.

For any finite positive-plaquette
support \(\mathcal D\), let
\(N_{\mu\nu}(\mathcal D)\) count its
faces of a given orientation. Exactly,
\begin{equation}
 \sum_x\sum_{p\in x+\mathcal D}s_p
    =\sum_{\mu<\nu}N_{\mu\nu}(\mathcal D)
                       \sum_y s_{(y;\mu,\nu)}
    \le|\mathcal D|S.
 \label{eq:f15v50:convolution}
\end{equation}
For each fixed offset, translation of
its base runs once through the torus;
this proves the identity, including
periodic seams. Use
\(|\mathcal P_n^+|=d_n\) and
\(|\mathfrak L_n|\le22d_n\) to sum
\eqref{eq:f15v50:pointwise-ledger}.
This gives both count bounds.

Finally \(d_n\le24n^3\) gives
\[
 A_n/\pi^2
 \le12288n^4+13138329600n^{11}
 \le13138341888n^{11}<2^{34}n^{11}.
\]
\end{proof}

This is not the old arbitrary-exterior
auxiliary-fibre repair theorem. It is
a count of V48's explicit certificates
in the full compact field, paid by
the original total plaquette action.
The loss in \(n^{11}\) is retained.

\subsection{Arbitrary-source density and random-location entry}

\begin{theorem}[Source-uniform measured density]
\label{f15v50:density}
Under every law in
\eqref{eq:f15v50:domain},
\begin{equation}
 \frac1V\mathbb E_JD_n\le\frac{2KA_n}{\beta},\qquad
 \log\mathbb E_J
             \exp\left(\frac{\beta D_n}{2A_n}\right)\le KV.
 \label{eq:f15v50:density}
\end{equation}
Consequently, for \(\rho\ge0\),
\begin{equation}
 \mu_J\{D_n\ge\rho V\}
       \le\min\left\{1,
                  \exp\left[V\left(K-\frac{\beta\rho}{2A_n}\right)\right]
              \right\}.
 \label{eq:f15v50:density-tail}
\end{equation}
In particular a density at least
\(2A_n(K+a)/\beta\) has probability
at most \(e^{-aV}\) for \(a\ge0\).
If that density threshold exceeds one,
the assertion has no nontrivial content.
\end{theorem}

\begin{proof}
Take expectations in
\eqref{eq:f15v50:ledger} and use
\eqref{eq:f15v50:total}. The pointwise
exponential comparison
\[
                  e^{\beta D_n/(2A_n)}\le e^{\beta S/2}
\]
gives the second assertion. Markov's
inequality gives the remaining bounds.
\end{proof}

If \(Z\) is a uniform random cell base,
independent of the field sampled under
the fixed law \(\mu_J\), then exactly
\[
 \mathbb P_J(C_{n,Z})=V^{-1}\mathbb E_JD_n.
\]
For any fixed offsets \(z_1,\ldots,z_k\),
the union bound thus gives
\begin{equation}
 \mathbb P_J\left(\bigcup_{i=1}^kC_{n,Z+z_i}\right)
          \le\min\{1,2kKA_n/\beta\}.
 \label{eq:f15v50:random-location}
\end{equation}
This is a randomly translated observation
of the original law, not a change in
its source coefficients. When
\((\log\beta)/M\) is bounded and
\(kn^{11}=o(\beta)\), these finitely
or slowly growing many translated
certificates are simultaneously good
with probability tending to one,
uniformly over all allowed source arrays.

One may equivalently translate the pair
\((J,U)\) by the independent \(Z\).
It is important to keep that pair:
the marginal mixture
\(V^{-1}\sum_z\mu_{\sigma_zJ}\)
is not being identified with the
Wilson measure having the averaged
coupling array. Such an identification
does not follow from averaging normalized
Gibbs measures.

\subsection{A genuinely nonconstant physical-source entry theorem}

\begin{theorem}[Prescribed-location entry at bounded translation index]
\label{f15v50:entry}
For every allowed coefficient array,
every base \(x\), and every embedded
box as above,
\begin{equation}
 \begin{split}
 \mu_J(C_{n,x})
   &\le\min\{1,2K\iota_J A_n/\beta\},\\
 \mu_J(\Gamma_{n,x}^c)
   &\le\mu_J(C_{n,x}),\\
 \mu_J\{R_{n,x}(Y)>\delta_n\}
   &\le\min\{1,2K\iota_J A_n^Y/\beta\}.
 \end{split}
 \label{eq:f15v50:entry}
\end{equation}
For any \(k\) specified bases, the
probability that at least one certificate
fails is at most
\begin{equation}
                    \min\{1,2kK\iota_J A_n/\beta\}.
 \label{eq:f15v50:finite-family}
\end{equation}
Thus if \((\log\beta)/M\) stays bounded
and \(k\iota_J n^{11}=o(\beta)\), the
actual sampled fields pass all these
certificates with probability tending
to one, under the original source law.
No reflection invariance is assumed.
\end{theorem}

\begin{proof}
The constructions of V47--V48 use
fixed paths and supports relative to
the box base. Translating the box and
field translates the certificate,
including its exterior-only frame.
Hence \(\mathbf1_{C_{n,x}}\) is a
translation-covariant local field.
Apply the general orbit assertion of
Proposition~\ref{f15v50:orbit} to
\eqref{eq:f15v50:density}.
For the exterior-only event use the
second bound in
\eqref{eq:f15v50:ledger} instead.
The inclusion from V48 proves the
coherent-fit assertion. The finite-family
statement is an ordinary union bound;
it does not use independence or
separated box supports.
The growth assertion follows from
\(A_n\le2^{34}\pi^2n^{11}\) and the
bound on \(K\).
\end{proof}

Here is an explicit physical source
outside V49's invariant-coefficient
class. For even \(M\), take
\begin{equation}
              h_x=a(-1)^{x_1},\qquad
                    0<|a|\le\tau\le1/32.
 \label{eq:f15v50:checkerboard}
\end{equation}
Its actual cell-to-plaquette lift is
\begin{equation}
 (\mathcal A_{\mu\nu}h)(y)
  =\begin{cases}
       a(-1)^{y_1},&1\in\{\mu,\nu\},\\
       0,&1\notin\{\mu,\nu\}.
    \end{cases}
 \label{eq:f15v50:checkerboard-lift}
\end{equation}
Indeed in the second case the transverse
average contains both values of the first
coordinate parity, which cancel; in the
first case that coordinate is never
averaged. Therefore
\(J_p(h)=\beta(1-\mathcal A h(p))\)
has exactly two translation classes:
\(H_J=\{z:z_1\text{ is even}\}\),
\(\iota_J=2\).
It is not invariant under the first-coordinate
site reflection on an orientation containing
that coordinate, because the plaquette base
then changes by \(y_1\mapsto2c-y_1-1\).
The visible coefficient alternation changes
sign. In particular this is not one of
V49's invisible source modes.

Theorem~\ref{f15v50:entry} applies directly
to this nonconstant physical tilt, with
\[
                  \mu_J(C_{n,x})\le\min\{1,4KA_n/\beta\}
\]
at every location. More generally any
bounded-period physical cell source has
a bounded-period induced coefficient
array, so its actual law is covered
with the corresponding translation index.
The index belongs to the induced \(J\),
not necessarily to a chosen representative
of \(h\) modulo its invisible kernel.

\subsection{The remaining gap is quantitative and location-sensitive}

This section acquires actual-source
local coherent entry for visible
nonconstant periodic backgrounds.
On the good certificate the frame
and the full-source conditional
normalizer of V48 apply without
altering the sampled exterior.
For arbitrary backgrounds it gives
only spatially averaged entry, with
a genuinely paid total-action budget.
It does not make an aperiodic source
translation invariant by declaration.

There are two specific losses.
First, for a general array the index
may be \(V\). The bound
\(2KVA_n/\beta\) can then be trivial
at a specified location even while
the spatial mean tends to zero.
Second, bounded index yields only
the first-moment and union estimates
above. It does not prove
\(\mu_J(\bigcap_{i=1}^k C_{n,x_i})\le q^k\).
As a purely probabilistic warning,
\(k\) identical rare event indicators
have intersection probability \(q\),
not \(q^k\). This observation is not
asserted to be a Wilson counterexample.

The macroscopic exponential bound
cannot silently replace the missing
joint bound either. An intersection
at \(k\) distinct bases implies only
\(D_n\ge k\), so the directly justified
bound is
\[
       \min\{1,\exp(KV-\beta k/(2A_n))\}.
\]
It retains the extensive \(KV\) debit.
That debit need not be paid by a
small prescribed family of defects.
The density theorem therefore does
not exclude a thin defect carrier
or establish V37's exceptional-energy
exponential moment.

In particular, a probability tending
to zero cannot simply delete an
unbounded repair weight. The
\(O(n^7)\) flat-repair energy debit
and its shrinking exponential window
remain as in V48. Controlling those
weights under a general nonconstant
source, or proving a genuinely local
conditional energy estimate that
avoids the translation-index loss,
remains the next upstream requirement.
The physical continuum and Hamiltonian
gap have not been obtained from
unscaled plaquettes becoming close
to the identity.

\subsection{Exact checks and preservation}

The diagnostic is
\begin{center}\small
\path{scripts/verify_ym_f15_v50_source_orbit_entry.py}.
\end{center}
It imports the unchanged V49/V47
helpers and executes their existing
checks. The new checks use the actual
mirror and exterior-collar supports
for sides 2, 3 and 4. Six complete
translation-convolution identities,
27 certificate-density ledgers and
nine coefficient/threshold identities
are verified in rational arithmetic.
The nonnegative plaquette arrays in
these incidence checks are not claimed
to be samples of a Gibbs distribution.

Five translation-orbit budget checks
cover different period indices.
The physical checkerboard lift is
checked at all 1,536 positive
plaquettes of a side-four torus,
including its exact two-coset
stabilizer and failure of V49's
reflection invariance. Seventy-two
macroscopic tail-exponent identities
and four generic non-product warnings
are also checked. The probabilistic
theorems follow from the analytic
pressure proof, not finite sampling.

The script has 6,589 bytes and SHA--256
\begin{center}\scriptsize\ttfamily
5D66A6509CA38D8EFF9D05D472D5C03DEA3C3816600309CCC7B281F9C17CB82C.
\end{center}
The V49 predecessor has 1,150,153 bytes and SHA--256
\begin{center}\scriptsize\ttfamily
66124471FCD47551EFD0D8E5AC811B3EB6C002B52A7E81BAC72675018199EF28.
\end{center}
Removing this final section and reversing
the title increment recovers V49 exactly.
V50 replaces V49 as the sole active
main source. The 14 protected inputs
and 92 earlier diagnostics remain
unchanged. Only \path{.tex} files are
kept in \path{latex_notes}; build
products are directed to
\begin{center}\small
\path{artifacts/ym_f15_v50_texcheck/}.
\end{center}
The next companion checkpoint is V55.

\begin{firewall}
V50 bounds the total action under all
allowed sources. It also bounds the
spatial mean of coherent-entry failures.
Finite translation
index upgrades that bound to every
prescribed location and covers visible
nonconstant periodic physical cell
sources. It does not prove a general
aperiodic-source local exponential
moment, prescribed-family product
tails, an interacting continuum or
the physical Yang--Mills mass gap.
\end{firewall}
\section{V51: block-reflected physical sources and actual joint entry bounds}
\label{f15v51:section}

V50 reaches nonconstant periodic physical
sources with a first-moment entry bound.
It does not give the joint exponential
estimate needed to control prescribed
families of defects. V49 gives such an
estimate in a different source class,
but its requirement of every unit-cell
site reflection excludes every visible
nonconstant symmetric cell source.

The present step keeps site reflections
only at boundaries of larger blocks.
An arbitrary physical source profile
on one block can be extended by those
reflections. Its lift is genuinely
visible in the plaquette couplings,
and its actual Wilson law is reflection
positive at the required planes.
The block-volume cost \(P^4\) is kept
explicitly. The result is an actual-law
joint cell moment and coherent-entry
bound for this specified family.

This is not a proof for every periodic
or aperiodic source. In particular,
replacing the exterior of a given
source by its reflected extension
changes the law. No comparison making
that replacement free is asserted.
The full physical mass-gap objective
remains unchanged and unproved.

\subsection{Inherited arguments and the source domain}

The V50 predecessor was revalidated.
The companion review of V50 is current;
the next is V55. The nonduplication
review checked f6 V86's abstract block
chessboard argument, f14 V99's scalar
cell version, and V49's actual-source
verification and intersection theorem.
Those reflection principles are not
new results of this section.

The standard closed-block statement
was also rechecked in M.~Biskup,
\emph{Reflection Positivity and Phase
Transitions in Lattice Spin Models},
\href{https://www.math.ucla.edu/~biskup/PDFs/papers/Prague-School.pdf}
{Section~4.3, Theorem~4.8}.
For the observables needed here, the
already proved scalar-label argument
suffices: replace a fine-cell energy
by the sum of those energies inside
one larger block. The background-law
hypotheses and physical source lift
are verified below, rather than
inferred from a spin-model example.

Fix an integer \(P\ge1\) and
\begin{equation}
 M=PL,\qquad L=2^j\ge4,\qquad
 V=M^4,\qquad N=L^4.
 \label{eq:f15v51:scales}
\end{equation}
Here \(N\) is the number of blocks.
Retain \(\beta\ge4\), \(0\le\tau\le1/32\)
and \((1-\tau)\beta\le J_p\le(1+\tau)\beta\).
Let \(\mathcal R^{[P]}_{M,\beta,\tau}\)
be the arrays in this interval that
are fixed by every coordinate site
reflection through \(x_i=kP\).
The positive-plaquette base transforms
by \eqref{eq:f15v49:reflection} with
\(c=kP\). At \(P=1\) this is precisely
V49's coefficient class. For \(P>1\),
invariance at the intervening integer
planes is not required.

All probabilities use the full compact
law \(\mu_J\) of
\eqref{eq:f15v49:law}. Write
\begin{equation}
 K=62+\frac92\log\frac{1+\tau}{1/2-\tau}
             +\frac6M\log((1+\tau)\beta),\qquad
                         K_P=P^4K.
 \label{eq:f15v51:K}
\end{equation}
The total-action estimate
\eqref{eq:f15v50:total} holds for these
arrays, and \(K\le66+6B\) if
\((\log\beta)/M\le B\).

\subsection{Construct physical sources in the required actual-law class}

Let \(Q_P=\{0,\ldots,P-1\}^4\).
For an integer \(t\), define the fold
\begin{equation}
 r_P(t)=
 \begin{cases}
  r,&0\le r<P,\\
  2P-1-r,&P\le r<2P,
 \end{cases}
 \qquad r=t\bmod 2P.
 \label{eq:f15v51:fold}
\end{equation}
Choose any real profile \(v=(v_r)_{r\in Q_P}\)
with \(\|v\|_\infty\le\tau\), and set
\begin{equation}
 h_x=v_{(r_P(x_1),\ldots,r_P(x_4))},\qquad
 J_{(y;\mu,\nu)}(h)
       =\beta\bigl(1-(\mathcal A_{\mu\nu}h)(y)\bigr),
 \label{eq:f15v51:physical-source}
\end{equation}
where \(\mathcal A\) is the original
four-cell lift in
\eqref{eq:f15v49:lift}.
This is exactly the physical tilt
\(e^{\sum_xh_xT_x}d\mu_{M,\beta}^{\mathrm W}\)
after normalization, not a newly
defined plaquette-source convention.

\begin{proposition}[Reflected profiles give admissible physical backgrounds]
\label{f15v51:physical-class}
Every profile above induces
\(J(h)\in\mathcal R^{[P]}_{M,\beta,\tau}\).
The extension agrees with \(v\) on
\(Q_P\) and is periodic with period
\(2P\) in each coordinate.
\end{proposition}

\begin{proof}
The fold satisfies
\(r_P(2kP-t-1)=r_P(t)\).
Consequently \(h\) is fixed by the
cell-base reflection
\(x_i\mapsto2kP-x_i-1\).
The four-cell average has the
corresponding plaquette covariance.
If the reflected coordinate is parallel
to the plaquette, its averaging offset
is zero and the same cell reflection
appears in every summand.
If it is transverse, pair offsets
\(\epsilon_i\) and \(1-\epsilon_i\)
in the four-term sum. Indeed
\[
 2kP-y_i-(1-\epsilon_i)
                  =2kP-(y_i-\epsilon_i)-1.
\]
The other offsets do not change.
This proves invariance of the induced
coefficient at the reflected positive
plaquette, with its correct parallel
or transverse base convention.

Every induced source coefficient is
an average of four numbers in
\([-\tau,\tau]\), so its coupling
lies in the required interval.
Agreement on \(Q_P\) and period \(2P\)
follow directly from the fold.
\end{proof}

\subsection{The physical profile is visible, with an exact inverse}

The larger-block construction does
not merely evade V49's obstruction
by inserting invisible source modes.
There is a direct inverse on the
fundamental profile, with an explicit
size cost.

For a one-dimensional vector on
\(\{0,\ldots,P-1\}\), let
\begin{equation}
 (B_Pw)_0=w_0,\qquad
 (B_Pw)_r=\tfrac12(w_r+w_{r-1})
                         \quad(1\le r<P).
 \label{eq:f15v51:B}
\end{equation}
When acting on a four-coordinate
profile, \(B_{P,i}\) acts only in
coordinate \(i\).

\begin{theorem}[One plaquette orientation recovers the whole source profile]
\label{f15v51:inverse}
For a fixed orientation \(\mu<\nu\)
with complementary coordinates
\(\lambda,\kappa\),
\begin{equation}
 (\mathcal A_{\mu\nu}h)|_{Q_P}
               =B_{P,\lambda}B_{P,\kappa}v.
 \label{eq:f15v51:triangular}
\end{equation}
As a map between real \(P^4\)-dimensional
profile spaces, this is invertible and
\begin{equation}
 \begin{split}
 \det(B_{P,\lambda}B_{P,\kappa})
                       &=2^{-2(P-1)P^3},\\
 \|v\|_\infty
     &\le(2P-1)^2
                \|\mathcal A_{\mu\nu}h|_{Q_P}\|_\infty.
 \end{split}
 \label{eq:f15v51:inverse-bounds}
\end{equation}
The inverse norm \((2P-1)^2\) is exact.
The linear profile map has rank \(P^4\).
For \(\tau>0\), its physical source
family thus has \(P^4\) visible parameters;
no nonzero profile lies in the lift's
invisible kernel. If an induced
orientation array is constant, the
profile itself is constant.
\end{theorem}

\begin{proof}
At a transverse coordinate \(y_i=0\),
the folded source at \(-1\) is its
value at zero. At \(1\le y_i<P\),
the two averaged values are those at
\(y_i\) and \(y_i-1\).
The two transverse averages commute,
giving \eqref{eq:f15v51:triangular}.
The one-dimensional matrix \(B_P\)
is lower triangular, with diagonal
\(1,1/2,\ldots,1/2\).
Each coordinate action repeats that
matrix \(P^3\) times. Multiplication
of the two determinants proves the
first line of
\eqref{eq:f15v51:inverse-bounds}.

For \(y=B_Pw\), solve successively
\[
 w_0=y_0,\qquad w_r=2y_r-w_{r-1}.
\]
Induction gives
\(|w_r|\le(2r+1)\|y\|_\infty\).
The data \(y_r=(-1)^r\) attain equality:
\(w_r=(-1)^r(2r+1)\).
Apply the inverse in each transverse
coordinate. The tensor-product data
\((-1)^{y_\lambda+y_\kappa}\), constant
in the parallel coordinates, attain
the product norm \((2P-1)^2\).
Finally the matrices preserve constants,
so a constant image has only its
constant preimage.
\end{proof}

Distinct induced coefficient arrays
also define distinct normalized Wilson
laws. For completeness, write
\(w_p=\tfrac12\Tr U_p\).
Under product Haar, \(\mathbb E_Hw_p=0\)
and \(\mathbb E_Hw_pw_q=0\) for distinct
positive plaquettes: a link in \(p\)
but not \(q\) can be multiplied by
the central element \(-I\), changing
the first trace's sign only.
Moreover \(\mathbb E_Hw_p^2>0\), since
the continuous trace is nonzero on
an open set. Thus the \(w_p\)'s are
linearly independent, even modulo
constants. Equality of two normalized
positive Wilson densities would make
their action difference constant and
therefore force every coefficient
difference to vanish.
This elementary Haar argument verifies
visibility at the level of the law;
it is not a new correlation inequality
under \(\mu_J\).

The inverse norm grows with \(P\).
The source map has not been claimed
uniformly well conditioned as the
fundamental profile size tends to infinity.

\subsection{Closed-block energy and the scalar chessboard bound}

For \(z\in(\mathbb Z/L\mathbb Z)^4\),
let \(B_z=Pz+Q_P\) denote its fine-cell
bases and define
\begin{equation}
                   W_z^{[P]}=\sum_{x\in B_z}T_x.
 \label{eq:f15v51:block-energy}
\end{equation}
This observable is supported in the
closed geometric block
\(Pz+[0,P]^4\). Adjacent supports
can share links and faces on the
block boundary. Exactly,
\begin{equation}
                    \sum_zW_z^{[P]}=\beta S.
 \label{eq:f15v51:partition}
\end{equation}

Reflection at \(x_i=kP\) sends the
block index to
\(z_i'=2k-z_i-1\pmod L\), and the
within-block fine-cell index to
\(r_i'=P-1-r_i\). The covariance
of \(T_x\) proved in f14 V99 hence
takes \(W_z^{[P]}\) exactly to
\(W_{z'}^{[P]}\).

For \(J\in\mathcal R^{[P]}_{M,\beta,\tau}\),
the conditional boundary-Haar identity
\eqref{eq:f15v49:RP} holds at each
of these planes. Its two fixed planes
are separated by \(M/2\), a multiple
of \(P\). An elementary plaquette
still cannot meet both open halves;
the action splits into a half and
its reflected copy plus the fixed-plane
action. The larger block size does not
alter that proof. Thus the positivity
is in the actual source law, including
the physical profiles just constructed.

\begin{proposition}[Block scalar sources]
\label{f15v51:block-chessboard}
For any real block array \(a_z\),
\begin{equation}
 \log\mathbb E_J e^{\sum_z a_zW_z^{[P]}}
   \le\frac1N\sum_z
       \log\frac{Z_M(J-\beta a_z\mathbf1)}{Z_M(J)}.
 \label{eq:f15v51:block-chessboard}
\end{equation}
Each numerator again shifts all
plaquette couplings by the indicated
scalar, not only one block's couplings.
\end{proposition}

\begin{proof}
Apply the scalar-label proof of
Proposition~\ref{f15v49:chessboard}
on the coarse block-index torus.
Its hypotheses have just been checked:
there is an even periodic index grid,
the actual law is reflection positive
at its separating planes, and the
energy factors transform covariantly
with support in their closed blocks.
Shared boundary variables remain in
the conditional Haar integral.

More explicitly, the scalar label
normalizer is now
\[
 c_{J,P}(t)
       =\left(\mathbb E_J e^{t\sum_zW_z^{[P]}}\right)^{1/N}.
\]
Reflection Cauchy--Schwarz duplicates
either half's labels while leaving
\(J\) fixed. The finite-label maximum
with the fewest disagreeing neighboring
blocks is constant, by the same
disagreement argument in that proof.
Its normalized value is one.
Equation~\eqref{eq:f15v51:partition}
then identifies \(c_{J,P}(t)^N\) with
the partition ratio displayed above.
No translation invariance by \(P\)
and no reflection at intermediate
unit-cell planes were used.
\end{proof}

\subsection{A joint fine-cell moment in the nonconstant physical background}

For a nonnegative fine-cell test array
\(b\), put
\begin{equation}
           \mathcal M_P(b)
                    =\sum_z\max_{x\in B_z}b_x.
 \label{eq:f15v51:block-max}
\end{equation}
The background source \(h\) and this
test array \(b\) are different objects.
The former specifies the actual law;
the latter is an observable insertion.

\begin{theorem}[Arbitrary-support positive fine-cell moment]
\label{f15v51:CGF}
For \(J\in\mathcal R^{[P]}_{M,\beta,\tau}\)
and any fine-cell array \(0\le b_x\le1/2\),
\begin{equation}
 \log\mathbb E_J e^{\sum_xb_xT_x}
           \le2K_P\mathcal M_P(b)
           \le2K_P\sum_xb_x.
 \label{eq:f15v51:CGF}
\end{equation}
The support can be arbitrary and extensive.
There is no extra factor depending on
the total number of unobserved blocks.
\end{theorem}

\begin{proof}
Set \(a_z=\max_{x\in B_z}b_x\).
Since \(T_x\ge0\),
\(\sum_xb_xT_x\le\sum_z a_zW_z^{[P]}\).
Use \eqref{eq:f15v51:block-chessboard}.
The total-action bound
\eqref{eq:f15v50:total} gives, for
each \(0\le a_z\le1/2\),
\[
 \log\frac{Z_M(J-\beta a_z\mathbf1)}{Z_M(J)}
                                      \le2Ka_zV.
\]
The ratio \(V/N=P^4\) proves the
first inequality. For nonnegative
coefficients, each block maximum is
at most the sum of its coefficients;
summing proves the second.
\end{proof}

In particular the one-plaquette bound is
\begin{equation}
                       \mathbb E_Je^{X_p/8}\le e^{K_P},
 \label{eq:f15v51:one-face}
\end{equation}
because \(X_p\le4T_{\operatorname{base}(p)}\).
At \(P=1\), \(\mathcal M_1(b)=\sum_xb_x\)
and the theorem recovers V49's bound.
For a constant fine-cell test, its
first inequality recovers the global
pressure budget without an additional
block-volume loss. A localized test
retains that loss.

The theorem does not assert that the
law after an arbitrary further tilt
by \(b\) remains in the same class.
Such a tilt generally breaks its
block reflection invariance. A moment
bound at the specified background
must not be relabelled as a uniform
statement at all subsequently tilted
backgrounds.

\subsection{Joint coherent-entry certificates in that same law}

Keep the V48 supports and thresholds,
including \(d_n=24n(n-1)^2\),
\(u_n=1/(1024n)\) and
\(\delta_n^2=1/(518400\pi^2d_n n^5)\).
Assume \(n\ge2\) and \(n+3\le M\).
Use the actual certificate events
\(C_{n,x}\) in
\eqref{eq:f15v50:count}, with
\(\Gamma_{n,x}^c\subseteq C_{n,x}\).

\begin{theorem}[Actual-source joint certificate bound]
\label{f15v51:joint}
Define
\begin{equation}
 \begin{split}
 q_{\rm out}^{[P]}
    &=\exp\left\{-\frac\beta{24576\pi^2n}
                                  +\frac{K_Pd_n}{6}\right\},\\
 q_{\rm collar}^{[P]}
    &=|\mathfrak L_n|
                     \exp\{K_P-\beta\delta_n^2/16\},\\
 q^{[P]}&=\min\{1,q_{\rm out}^{[P]}+q_{\rm collar}^{[P]}\}.
 \end{split}
 \label{eq:f15v51:q}
\end{equation}
For every \(J\in\mathcal R^{[P]}_{M,\beta,\tau}\)
and every family of translated boxes
whose positive-plaquette base sets in
\(\mathcal P_{n,i}^+\cup\mathfrak L_{n,i}\)
are pairwise disjoint,
\begin{equation}
 \mu_J\left(\bigcap_{i=1}^k\Gamma_{n,i}^c\right)
     \le\mu_J\left(\bigcap_{i=1}^kC_{n,i}\right)
     \le(q^{[P]})^k.
 \label{eq:f15v51:joint}
\end{equation}
The source profiles in
\eqref{eq:f15v51:physical-source} are included.
No separation of the coarse \(P\)-blocks
containing those base sets is needed.
\end{theorem}

\begin{proof}
The mixed-witness proof of V49 now
uses \eqref{eq:f15v51:CGF}, with
\(K\) replaced by \(K_P\).
For clarity, an outward failure uses
\[
 \frac{\beta E_n^{\rm out}}{24\pi^2}
           \le\frac1{48}\sum_{p\in\mathcal P_n^+}X_p.
\]
The corresponding fine-cell coefficients
are \(m_x/12\le1/2\), with sum
\(d_n/12\). Their moment cost is
\(K_Pd_n/6\), and the threshold is
\(\beta u_n/(24\pi^2)\).
A specified collar witness has
\(X_q>\beta\delta_n^2/2\); the
exponent \(X_q/8\) costs at most
\(K_P\), with threshold
\(\beta\delta_n^2/16\).

Choose one alternative in each box
and, for a collar alternative, its
witnessing plaquette. Add the exponents.
The base-set assumption ensures the
resulting fine-cell coefficients remain
at most \(1/2\). The last, additive
bound in \eqref{eq:f15v51:CGF} gives
the sum of the costs just stated,
even if different boxes use distinct
fine cells within the same coarse block.
One Markov bound thus controls the
entire chosen intersection. Sum over
collar witnesses and then over the
two alternatives in each box.
The result is
\((q_{\rm out}^{[P]}+q_{\rm collar}^{[P]})^k\).
Cap by one and use the deterministic
inclusion of the coherent-fit failures.
No independence was invoked.
\end{proof}

If \(\beta\to\infty\),
\((\log\beta)/M\) remains bounded and
\begin{equation}
                  n^8\bigl(P^4+\log(n+2)\bigr)=o(\beta),
 \label{eq:f15v51:growth}
\end{equation}
then \(q^{[P]}\to0\) uniformly over
the specified background class.
Indeed \(|\mathfrak L_n|\le22d_n\)
and \(d_n\le24n^3\). The collar's
negative term has magnitude at least
a fixed positive constant times
\(\beta/n^8\), which dominates
\(K_P+\log(22d_n)\) under
\eqref{eq:f15v51:growth}.
The same condition implies
\(\beta/(P^4n^4)\to\infty\), so
the outward negative term also
dominates its positive term.
This proves the limit, including
fixed \(n,P\).

\subsection{The paid flat-repair moment is now in the actual source law}

V48's deterministic repair estimate
holds for every allowed \(J\):
\[
 T_{n,J}^\flat(Y)\le\mathfrak c_n
                       \sum_{p\in\mathfrak L_n}X_p,
 \qquad \mathfrak c_n=44+356400d_n n^4.
\]
Consequently the newly proved source
moment gives
\begin{equation}
 \log\mathbb E_J
           \exp\left\{\frac{T_{n,J}^\flat}{48\mathfrak c_n}\right\}
                       \le\frac{K_P|\mathfrak L_n|}{6}.
 \label{eq:f15v51:repair-moment}
\end{equation}
To check this, bound its exponent by
\(\sum_{\mathfrak L_n}X_p/48\), then
use \(X_p\le4T_{\operatorname{base}(p)}\).
At most six positive plaquettes have
one base, so the fine-cell coefficient
is at most \(1/2\); their sum is
\(|\mathfrak L_n|/12\).
Equation~\eqref{eq:f15v51:CGF} gives
the displayed cost.

This moves the small repair moment
from the selected homogeneous law
to the same nonconstant source law
as the joint entry theorem. Its
coefficient is still
\(1/(48\mathfrak c_n)\), not one.
The \(O(n^7)\) repair debit and its
shrinking exponential window have
not been removed.

\subsection{An explicit visible example and the remaining exterior issue}

For \(P=2\), choose a profile depending
only on its first coordinate, with
\(v_0=a\), \(v_1=-a\),
\(0<|a|\le\tau\).
The extended one-coordinate source
is the period-four sequence
\[
                         h=a(1,-1,-1,1).
\]
For plaquette orientations containing
the first coordinate its lift is that
same sequence. For transverse
orientations the lift is
\[
                            a(1,0,-1,0).
\]
These are nonconstant physical
couplings invariant at the required
even site planes. Reflection at the
intervening unit plane fails on the
parallel orientations. Thus this
example lies outside V49's all-site
coefficient class, but its actual
joint bounds are now proved with
\(K_P=16K\).

It is a different source from V50's
period-two checkerboard
\(a(1,-1,1,-1)\). The latter is
not fixed by any first-coordinate
integer site-plane reflection on
cell bases. The present theorem
does not upgrade every periodic
source merely because its period
is finite.

An arbitrary local profile on \(Q_P\)
does have the reflected extension
constructed above. But that fact
does not compare its law with a
given arbitrary exterior source
that agrees only on \(Q_P\).
Local source agreement does not
by itself cancel the different
full normalizers or make their
local marginals equal.
Such a localization estimate is
still needed for the general-source
program. Enlarging \(P\) to the
scale of the entire torus also
makes \(K_P\) of volume order and
can destroy the small-activity
regime in \eqref{eq:f15v51:growth}.

Thus the new domain is a verified
family of visible physical source
laws with exponential joint entry,
not a change to the ultimate goal
or a complete solution of its
arbitrary-source interface.
The exceptional-energy normalization
criterion in V37, scale-uniform
physical clustering, an interacting
continuum and a positive physical
Hamiltonian gap remain unproved
at their full required scopes.

\subsection{Exact checks and preservation}

The diagnostic is
\begin{center}\small
\path{scripts/verify_ym_f15_v51_block_reflected_sources.py}.
\end{center}
It imports the unchanged V49/V47
helpers and runs their diagnostics.
For profile sides 2, 3 and 4, the
new exact tests check 22,592 folded
cell-source reflection identities,
135,552 induced plaquette identities,
and 2,118 lower-triangular lift rows.
Eighteen determinant and sharp inverse
norm checks reconstruct the full
profile from each orientation.
The coarse/fine reflection maps
are checked at 722,944 instances
on four-block-per-coordinate tori.

Three block-maximum and volume
ledgers, the explicit visible
period-four source, and 64 mixed
witness cost identities are also
checked. The probability estimates
follow from the analytic actual-law
reflection and pressure arguments,
not from finite sampling.

The script has 6,980 bytes and SHA--256
\begin{center}\scriptsize\ttfamily
2DD2A50BF9F6047A758A911E3379971618D8AC72C072D4B01060C3814256D348.
\end{center}
The V50 predecessor has 1,170,781 bytes and SHA--256
\begin{center}\scriptsize\ttfamily
D08A797225BC395278BA8B3E1BA9DDB6392D8D03CF020A95D1DE6E9F417CA160.
\end{center}
Removing this final section and reversing
the title increment recovers V50 exactly.
V51 replaces V50 as the sole active
main source. The 14 protected inputs
and 93 earlier diagnostics remain
unchanged. Only \path{.tex} files
belong in \path{latex_notes}; build
products are directed to
\begin{center}\small
\path{artifacts/ym_f15_v51_texcheck/}.
\end{center}
The next companion checkpoint is V55.

\begin{firewall}
V51 proves a joint exponential cell
moment and coherent-entry bound in
an explicitly constructed family of
visible nonconstant physical source
laws. It retains the block-volume
and repair-moment costs. It does
not transfer these bounds to an
arbitrary exterior source, close
general source normalization, construct
the interacting continuum or prove
the physical Yang--Mills mass gap.
\end{firewall}
\section{V52: a canonical-flat-reference amplification obstruction}
\label{f15v52:section}

V51 proves a small positive exponential moment
of the V48 flat-reference repair cost in an
actual nonconstant source law. It explicitly
does not prove the corresponding moment with
coefficient one. Before trying to enlarge that
window, this section tests the particular
reference against the full original Wilson
action. The outcome is an obstruction, not
another upper bound: its unit-strength moment
grows exponentially in the inverse coupling,
including in a volume-cofinal regime.

The mechanism is concrete. One exterior corner
link, irrelevant to the conditional boundary
interaction and to every outward repair path,
enters many of V48's canonical root frames.
Its original action has six nonzero plaquettes.
For a side-two box the resulting flat reference
creates forty-six nonzero rest plaquettes of
the same size. The sampled boundary is already
exactly coherent. Thus the obstruction concerns
the chosen reference and its exterior certificate,
not a demonstrated failure of actual coherence.

\subsection{Scope, nonduplication and notation}

The V51 predecessor, the V47 outward geometry,
V48's precise reference and V37's normalization
criterion for the actual law were rechecked. V48's
deterministic upper bound and shrinking moment
window are retained, not reproved or contradicted.
The new result is a lower bound on the raw
flat-repair moment. A local Haar-preserving
change of variables is used, so there is no
full-volume neighborhood probability hidden
in the cofinal statement. The V50 companion
checkpoint remains current; the next is V55.

Use unit quaternions for \(\SU(2)\), with
\[
 f(U)=1-\operatorname{Re}U=\tfrac12|U-I|^2,
 \qquad
 S_J(U)=\sum_p J_p f(U_p),\qquad
 d\mu_J=Z_J^{-1}e^{-S_J(U)}\,dU.
\]
Here \(dU\) is product normalized Haar measure
on the full periodic link configuration space.
Write \(S(U)=\sum_p f(U_p)\) for the homogeneous
action without its factor \(\beta\). Retain
\[
 T_{n,J}^{\flat}(Y)=S_{\rm rest}(F^{\flat}(Y),Y),\qquad
 T_{n,J}^{*}(Y)=S_{\rm rest}(B(Y),Y).
\]
The superscript star denotes V47's generally
curved outward center; it is not a minimization.
Set \(t_n^{\flat}=T_{n,\mathbf1}^{\flat}\).

All box coordinates below are ordinary integer
coordinates, embedded in a torus large enough
to avoid identifications in the box and collar.
Directions are numbered \(1,2,3,4\).
For the explicit side-two result take \(M\ge8\).
Keep \(\Gamma_n\), \(G_n^Y\), \(R_n\) and
\(\delta_n\) exactly as in
\eqref{eq:f15v48:thresholds}--\eqref{eq:f15v48:failure}.
In particular, \((G_n^Y)^c\) is not being
identified with \(\Gamma_n^c\).

\subsection{An exterior corner link and the exact reference it generates}

Let \(o=(0,0,0,0)\), \(e_*=(o,1)\), and fix
\begin{equation}
 Q=(3/5,4/5,0,0),\qquad
 f(Q)=2/5,\qquad |Q-I|^2=4/5.
 \label{eq:f15v52:Q}
\end{equation}
Define \(U^{(Q)}\) to equal \(I\) on every
original link except \(U_{e_*}^{(Q)}=Q\).
The edge \(e_*\) belongs to \(Y\), not
to \(A_n\) or \(F_n\).
Every edge of \(F_n\) has exactly one
transverse coordinate on the boundary and
the other two transverse coordinates strictly
between zero and \(n\).
In contrast, every face incident to \(e_*\)
has two other coordinates equal to zero.
No such face contains an \(F_n\)-edge.
Consequently \(e_*\) occurs in no rest face.
It also occurs in no outward bypass of an
\(F_n\)-edge: the two strictly interior
transverse coordinates of that bypass
exclude the corner edge.

It follows that
\begin{equation}
 B(Y)=I\text{ on }F_n,\qquad
 T_{n,J}^{*}(Y)=0,\qquad
 \mathcal E_n(F)=0,
 \qquad S(U^{(Q)})=6f(Q).
 \label{eq:f15v52:original-witness}
\end{equation}
The last equality counts the six plaquettes
incident to one link in four dimensions,
including the faces on its negative sides.
It is the full periodic action, not merely
the action of plaquettes touching \(F_n\).

Write a boundary vertex as \(x=(y_1,y_2,y_3,h)\)
and define \(\chi_n(x)\in\{0,1\}\) by
\begin{equation}
 \chi_n(x)=
 \begin{cases}
 0,&h=n,\\
 \mathbf1_{\{y_1>0\}},&h<n\text{ and some }y_i=0,\\
 \mathbf1_{\{y_1=n\}},&h<n,\ 0\notin\{y_i\},\ n\in\{y_i\},\\
 1,&h=0\text{ and }0<y_i<n\text{ for all }i.
 \end{cases}
 \label{eq:f15v52:indicator}
\end{equation}
These cases exhaust the boundary vertices.
Inspection of V41's stated path order gives
\(\chi_n(x)=1\) precisely when the canonical
path from \(o\) to \(x\) uses \(e_*\).
Indeed, a top-face path goes first in direction
four and avoids \(e_*\); an ordinary lateral
path uses it exactly when its first-coordinate
segment is nonempty; a lateral path with no
zero transverse coordinate first travels in
the earliest direction whose endpoint coordinate
is \(n\), and uses \(e_*\) exactly when that
direction is one. The remaining bottom-face
case starts with a nonempty direction-one segment.
Outward substitutions introduce no occurrence
of \(e_*\). Thus V48's frames and reference are
\begin{equation}
 g_x(Y)=Q^{\chi_n(x)},\qquad
 F_e^{\flat}(Y)=Q^{k_e},\qquad
 k_e=\chi_n(y)-\chi_n(x)\quad(e=(x,y)\in F_n).
 \label{eq:f15v52:flat-cochain}
\end{equation}

\begin{proposition}[Exact side-two amplification]
\label{f15v52:witness}
For \(n=2\), the field \(U^{(Q)}\) satisfies
\begin{equation}
 S(U^{(Q)})=12/5,\qquad
 t_2^{\flat}(Y)=92/5,\qquad
 t_2^{\flat}(Y)-S(U^{(Q)})=16.
 \label{eq:f15v52:amplification}
\end{equation}
Moreover \(U^{(Q)}\in\Gamma_2\cap(G_2^Y)^c\).
\end{proposition}

\begin{proof}
There are 48 edges in \(F_2\) and 144 rest
faces. Extend \(k_e\) by zero off \(F_2\).
For a positively oriented face \(p=(x;\mu,\nu)\),
\[
 \ell_p=k_{(x,\mu)}+k_{(x+\hat\mu,\nu)}
              -k_{(x+\hat\nu,\mu)}-k_{(x,\nu)}.
\]
Since the original \(Y\)-links of every rest
face are identities, its repaired holonomy
is exactly \(Q^{\ell_p}\). Substitution of
\eqref{eq:f15v52:indicator} gives the following
complete incidence count:
\begin{center}
\begin{tabular}{c@{\qquad}rrr}
\toprule
Orientation & \(\ell_p=1\) & \(\ell_p=-1\) & \(\ell_p=0\)\\
\midrule
\((1,2)\)&4&4&16\\
\((1,3)\)&4&4&16\\
\((1,4)\)&5&5&14\\
\((2,3)\)&2&2&20\\
\((2,4)\)&4&4&16\\
\((3,4)\)&4&4&16\\
\midrule
Total&23&23&98\\
\bottomrule
\end{tabular}
\end{center}
For clarity, this table is a finite integer
calculation with the displayed formula for
\(\ell_p\), not a probability sample. One may
generate precisely its face set by taking the
six incident faces of each \(F_2\)-edge, removing
duplicates and the interior affected faces.
Equivalently it is V47's rest-face set. The
nonzero edge families below give a second
explicit description of all inputs to the table.
Therefore \(t_2^{\flat}=46f(Q)=92/5\).
Equation~\eqref{eq:f15v52:original-witness}
gives the other action and \(\Gamma_2\).

Finally \(q_0=(o;1,2)\) contains no \(F_2\)-edge
and belongs to \(\mathcal P_2^{\partial,0}
\subseteq\mathfrak L_2\).
Its original holonomy is \(Q\). Hence
\(R_2(Y)\ge|Q-I|>1/2>\delta_2\), proving
failure of the exterior certificate.
\end{proof}

\subsection{A box-size lower bound without extrapolating the finite table}

For \(a,b\in\{1,\ldots,n-1\}\), the following
five edge families have \(k_e=1\); every listed
edge has direction one:
\[
 (0,a,b,0),\quad (0,0,a,b),\quad (0,a,0,b),\quad
 (n-1,n,a,b),\quad (n-1,a,n,b).
\]
The bases and directions of the seven
families with \(k_e=-1\) are
\begin{equation}
\begin{split}
 &((a,n-1,b,0),2),\quad ((a,b,n-1,0),3),\\
 &((n,a,b,n-1),4),\quad ((a,0,b,n-1),4),\\
 &((a,b,0,n-1),4),\quad ((a,n,0,b),3),\\
 &((a,0,n,b),2).
\end{split}
\label{eq:f15v52:families}
\end{equation}
These assertions follow by restricting
\eqref{eq:f15v52:indicator} to the facet
containing the edge: its jump can occur only
at the first or last step of a coordinate
segment. The cases \(h=0,n\), \(y_i=0,n\)
give exactly the displayed lists. All other
\(F_n\)-edges have equal indicators at their
endpoints. The twelve families are disjoint,
each with \((n-1)^2\) members.

Each member has a distinct outward mirror face
containing exactly that one \(F_n\)-edge.
All three other original links of this face
are identities in the witness. Its flat-repaired
holonomy is \(Q\) or \(Q^{-1}\).
The other rest-face actions are nonnegative.
This proves for every \(n\ge2\), without
extrapolation of the side-two count,
\begin{equation}
 t_n^{\flat}(Y(U^{(Q)}))
                     \ge12(n-1)^2 f(Q).
 \label{eq:f15v52:mirror-lower}
\end{equation}

\subsection{A local gauge and a Haar-preserving corner twist}

A single point has Haar measure zero.
Furthermore, a neighborhood of the identity
on all torus links could have a volume-order
cost. Neither issue is ignored here.
For \(n=2\), use only the embedded box
\(W=[-1,3]^4\), with 625 vertices, 2000 edges
and 2400 positive plaquettes. It contains all
original-link words entering the reference,
all rest faces, \(q_0\), and the six faces
incident to \(e_*\).

For \(x\in W\), let \(k_x(U)\) be transport
along the coordinate-ordered positive path
from \(w=(-1,-1,-1,-1)\) to \(x\), in order
\(1,2,3,4\). None of these paths uses
\(e_*\): its direction-one segment keeps the
other coordinates equal to \(-1\), and its
remaining segments have other directions.
Every \(k_x\) is therefore independent of
\(U_{e_*}\).

\begin{proposition}[Local near-flat frame]
\label{f15v52:local-frame}
If every plaquette of \(W\) has chord distance
at most \(\epsilon\) from \(I\), then
\begin{equation}
 |k_x(U)U_e k_y(U)^{-1}-I|\le12\epsilon
                    \quad(e=(x,y)\text{ in }W).
 \label{eq:f15v52:local-frame}
\end{equation}
\end{proposition}

\begin{proof}
Append the step \(e=(x,x+\hat\mu)\) to the
ordered path to \(x\). Move this step left
past its later-coordinate segments. There
are \(\sum_{j>\mu}(x_j+1)\le12\) square
interchanges, all in \(W\). The resulting
path is exactly the ordered path to \(y\).
Each interchange changes the transport by
the chord norm of a conjugated elementary
holonomy, at most \(\epsilon\). The triangle
inequality and unitary invariance of the
quaternion norm give the result. No commutativity
of the original field is used.
\end{proof}

Define a map on the full configuration space by
\begin{equation}
 (\Phi U)_{e_*}=k_o(U)^{-1}Qk_o(U)U_{e_*},
 \qquad (\Phi U)_e=U_e\quad(e\ne e_*).
 \label{eq:f15v52:twist}
\end{equation}
All \(k_x\) are unchanged by this map.
Conditional on every link except \(e_*\),
it is left translation of that link by a
fixed group element. It is a bijection,
with inverse obtained by replacing \(Q\)
by \(Q^{-1}\), and it preserves product
Haar measure exactly. This is not asserted
to preserve \(\mu_J\); its full action
debit will be included.

Put
\begin{equation}
 \eta=2^{-17},\qquad \epsilon=\eta/12,\qquad
 \mathcal H=\{\max_{p\subset W}|U_p-I|\le\epsilon\}.
 \label{eq:f15v52:local-event}
\end{equation}
On \(\mathcal H\), the field \(U\) in the
\(k(U)\) gauge is within \(\eta\) of \(I\)
on every \(W\)-edge. In the same gauge,
\(\Phi U\) is within \(\eta\) of
\(U^{(Q)}\) on every \(W\)-edge:
at \(e_*\) the new link is
\(Q[k_oU_{e_*}k_{o+\hat1}^{-1}]\), and
every other gauged link is unchanged.
The local gauge may be extended arbitrarily
to the remaining torus vertices. V48's
covariance makes \(t_2^{\flat}\) gauge invariant,
so this comparison can be used directly.

\subsection{Explicit debit, and coherence of the image event}

Each V48 frame for \(n=2\) is a word of
at most 24 original \(Y\)-links. Each flat
\(F_2\)-link is a word of at most 48, and
each repaired rest-face word has at most
192 original letters. Inversion preserves
chord distance, products telescope in that
distance, and \(f\) is 1-Lipschitz on unit
quaternions. The 144-face sum therefore gives
\begin{equation}
 \left|t_2^{\flat}(Y(\Phi U))-92/5\right|
       \le27648\eta=27/128\quad(U\in\mathcal H).
 \label{eq:f15v52:lipschitz}
\end{equation}
In particular \(t_2^{\flat}(Y(\Phi U))>18\).
For \((1-\tau)\beta\le J_p\le(1+\tau)\beta\),
\begin{equation}
 T_{2,J}^{\flat}(Y(\Phi U))\ge18(1-\tau)\beta.
 \label{eq:f15v52:repair-lower}
\end{equation}
The map changes exactly one original link.
Only its six incident plaquettes can change
in \(S_J\), and \(0\le f\le2\). Consequently
\begin{equation}
 S_J(\Phi U)-S_J(U)\le12(1+\tau)\beta.
 \label{eq:f15v52:debit}
\end{equation}
Combining the two inequalities, for
\(0\le\tau\le1/32\), gives
\begin{equation}
 T_{2,J}^{\flat}(Y(\Phi U))
       -S_J(\Phi U)+S_J(U)
 \ge(6-30\tau)\beta\ge81\beta/16.
 \label{eq:f15v52:net-gain}
\end{equation}

The image event has a useful additional property:
\begin{equation}
                    \Phi\mathcal H
                        \subseteq\Gamma_2\cap(G_2^Y)^c.
 \label{eq:f15v52:event-image}
\end{equation}
Indeed the map leaves all actual \(F_2\)-links
unchanged. In the \(k\) gauge each has chord
norm at most \(\eta\), and
\(\alpha(V)^2\le(\pi^2/4)|V-I|^2<4|V-I|^2\).
Thus their squared-angle sum is at most
\(192\eta^2<1/512=\varepsilon_2\).
The minimum coherent-fit energy is no larger
than this admissible gauge fit, proving
\(\Gamma_2\). On the other hand, the gauged
holonomy of \(q_0\) is within \(4\eta\) of
\(Q\). Since \(|Q-I|\ge4/5\),
\[
 R_2(Y(\Phi U))\ge4/5-4\eta>1/2>\delta_2.
\]
This proves the second part of the inclusion.

\subsection{Actual-law exponential growth along cofinal volumes}

\begin{theorem}[The unit-strength raw repair moment diverges]
\label{f15v52:moment-obstruction}
For every allowed coefficient array \(J\)
with \(0\le\tau\le1/32\) and \(M\ge8\),
\begin{equation}
 \mathbb E_J\!\left[
 e^{T_{2,J}^{\flat}(Y)}
       \mathbf1_{\Gamma_2\cap(G_2^Y)^c}\right]
               \ge e^{81\beta/16}\mu_J(\mathcal H).
 \label{eq:f15v52:all-array-lower}
\end{equation}
If also \(J\in\mathcal R^{[P]}_{M,\beta,\tau}\)
under V51's hypotheses, then
\begin{equation}
 \mathbb E_J\!\left[
 e^{T_{2,J}^{\flat}(Y)}
       \mathbf1_{\Gamma_2\cap(G_2^Y)^c}\right]
 \ge e^{81\beta/16}
       \left[1-2400e^{K_P-\beta\epsilon^2/16}\right]_+ .
 \label{eq:f15v52:cofinal-lower}
\end{equation}
Here \([r]_+=\max(r,0)\), and \(\epsilon\)
is the fixed number in
\eqref{eq:f15v52:local-event}.
For fixed \(P\), bounded \((\log\beta)/M\)
and \(\beta\to\infty\) in that source class,
the bracket is at least \(1/2\) eventually,
uniformly in \(J\). Thus neither a bounded
nor a subexponential-in-\(\beta\) upper bound
for this unit-strength moment is possible.
\end{theorem}

\begin{proof}
Use \eqref{eq:f15v52:event-image}, then
Haar preservation of the bijection \(\Phi\):
\begin{align*}
 \mathbb E_J\!\left[e^{T_{2,J}^{\flat}}
        \mathbf1_{\Gamma_2\cap(G_2^Y)^c}\right]
 &\ge\int_{\Phi\mathcal H}e^{T_{2,J}^{\flat}(Y(V))}\,d\mu_J(V)\\
 &=\int_{\mathcal H}
 e^{T_{2,J}^{\flat}(Y(\Phi U))-S_J(\Phi U)+S_J(U)}\,d\mu_J(U)\\
 &\ge e^{81\beta/16}\mu_J(\mathcal H).
\end{align*}
The same full normalizer \(Z_J\) appears on
both sides of the change of variables and
cancels exactly. No conditional normalizer
or omitted exterior action is being replaced.

For the second assertion use the inherited
one-face bound \eqref{eq:f15v51:one-face}.
Since \(X_p=\beta f(U_p)=\beta|U_p-I|^2/2\),
Markov's inequality gives
\[
 \mu_J\{|U_p-I|>\epsilon\}
              \le e^{K_P-\beta\epsilon^2/16}.
\]
There are exactly 2400 faces in \(W\).
A union bound, without independence, bounds
\(\mu_J(\mathcal H)\) from below by the
bracket in \eqref{eq:f15v52:cofinal-lower}.
At fixed \(P\) and bounded \((\log\beta)/M\),
V51 bounds \(K_P\) uniformly, whereas
\(\beta\epsilon^2/16\to\infty\).
This proves the cofinal assertion.
\end{proof}

The numerical onset obtained from this deliberately
conservative \(\epsilon\) is very large.
No useful finite-coupling numerical threshold
is claimed. Its role is a rigorous asymptotic
counterexample to the proposed uniform moment.
Already the homogeneous case \(P=1,\tau=0\)
suffices; nonconstant sources are not the
cause of the obstruction.

\subsection{No fixed positive repair coefficient works at every box size}

\begin{proposition}[A coefficient obstruction with increasing box size]
\label{f15v52:coefficient}
For a fixed \(n\ge2\), suppose
\(\lambda>1/[2(n-1)^2]\).
In the homogeneous law there are constants
\(c_{n,\lambda}>0\) and a fixed local curvature
threshold such that, along the cofinal regime
of the preceding theorem with \(P=1\),
\begin{equation}
 \mathbb E_{\beta\mathbf1}
       e^{\lambda T_{n,\beta\mathbf1}^{\flat}}
                 \ge\tfrac12 e^{c_{n,\lambda}\beta}
 \label{eq:f15v52:general-coefficient}
\end{equation}
eventually. For \(n=2\), the sharper condition
\(\lambda>3/23\) is sufficient.
Consequently no fixed \(\lambda>0\), independent
of \(n\), gives subexponential-in-\(\beta\)
repair moments for all box sizes.
\end{proposition}

\begin{proof}
Use \(W_n=[-1,n+1]^4\) instead of \(W\)
and the same ordered paths and corner twist.
The paths still avoid \(e_*\). The ordered
gauge estimate now has at most \(3(n+2)\)
square interchanges. All words entering
\(t_n^{\flat}\) and all faces changed by
\(\Phi\) lie inside this fixed local box.
For a field whose local gauged links are
exactly \(I\), the twist produces \(U^{(Q)}\)
locally. The exponent divided by \(\beta\) is
\[
 \lambda t_n^{\flat}(Y(U^{(Q)}))
              -[S(U^{(Q)})-S(I)]
 \ge\{12\lambda(n-1)^2-6\}f(Q)>0.
\]
Here the action difference depends on only
the six incident faces, even if the field
elsewhere on the torus is not flat.
The finitely many local word functions are
continuous. Choose a sufficiently small fixed
neighborhood of the local identity so that
the net gain remains at least half the displayed
positive lower bound; call this half
\(c_{n,\lambda}\). The local gauge estimate
supplies such a neighborhood from a sufficiently
small plaquette-curvature event in \(W_n\).
The one-face estimate and the finite union
bound make that event have probability at
least \(1/2\) eventually. Haar change of
variables proves the assertion exactly as
above. For \(n=2\), replace the mirror lower
bound by the exact coefficient 46; positivity
is then \(46\lambda-6>0\).
Finally, for any fixed \(\lambda>0\) choose
one integer \(n\) satisfying
\(2\lambda(n-1)^2>1\). This fixed box already
contradicts a subexponential bound intended
to hold at every box size.
\end{proof}

No assertion is made that these sufficient
thresholds are optimal. In particular this
result does not contradict the existing
coefficient \(1/(48\mathfrak c_n)\) in
\eqref{eq:f15v51:repair-moment}.

\subsection{Upstream correction and the next precise issue}

The particular strategy of deleting
\(e^{T_{n,J}^{\flat}}\) by a bounded or
subexponential raw unit moment must be
abandoned for V48's canonical reference.
It would demand a property contradicted
by Theorem~\ref{f15v52:moment-obstruction}.
Increasing the accuracy of its curvature
upper bound cannot change that conclusion.

The defect can be localized in the reference
choice. Varying \(U_{e_*}\) leaves the actual
\(F_n\), every \(B_e(Y)\), and all the
original rest-face link variables unchanged.
Thus it leaves \(T^*\) unchanged. Yet the
canonical boundary paths allow it to alter
many values of \(F^{\flat}\). A replacement
construction should not import that irrelevant
corner noise into the boundary interaction.
One possible next construction uses only the
connection \(B(Y)\) on the \(F_n\) graph,
for example fixed spanning trees within its
components, and is exact whenever that
connection is pure gauge on the graph.
Its quantitative costs still require proof;
no favorable unit moment for such a replacement
is assumed here.

A second possibility keeps compensation
jointly with the actual conditional normalizer,
instead of estimating the repair factor alone.
The present lower bound does not rule out
such cancellation. Nor does it falsify V37's
criterion for the actual normalized error:
the raw \(T^{\flat}\) is not automatically
that error. The distinction between a
certificate failure and actual coherent-fit
failure is essential: the explicit exponential
lower bound above is supported on \(\Gamma_2\).
It therefore cannot be relabelled as a lower
bound on the probability of \(\Gamma_2^c\).

\subsection{Exact checks and file policy}

The diagnostic is
\begin{center}\small
\path{scripts/verify_ym_f15_v52_flat_repair_amplification.py}.
\end{center}
It imports and passes the unchanged V47/V48
checks. Its new exact tests verify 10,640
canonical path indicators, all twelve edge
families for sides 2 through 8, and 1,680
distinct nonzero mirror faces. They check
the complete side-two orientation table and
its rational quaternion actions. A separate
audit of all 24,576 positive plaquettes of
the side-eight torus verifies the full
original action \(12/5\).

The 2000 local edge square-homotopies in
\(W\), the rational constants, and an exact
noncommuting near-flat twist case also pass.
The probability and Haar-preservation arguments
are analytic proofs above, not inferences
from those finite tests. The general-size
lower bound is proved by the displayed edge
families, not by fitting a polynomial to
the tested sizes.

The script has 7,964 bytes and SHA--256
\begin{center}\scriptsize\ttfamily
C42EA4886A92CF9F79EAA91FCC67EA813D5FF54AFE560713FCBE4C31B3CE2E77.
\end{center}
The V51 predecessor has 1,192,459 bytes and SHA--256
\begin{center}\scriptsize\ttfamily
79B997804EF06AF1F02A4FBA42803BC9CF24C6FD5A3D31B4AC48DA3A06ABEC08.
\end{center}
Removing this final section and reversing
the title increment recovers V51 exactly.
V52 replaces V51 as the sole active main
source. The 14 protected inputs and 94
earlier diagnostics remain unchanged.
Only \path{.tex} sources belong in
\path{latex_notes}; compilation products
are directed to
\begin{center}\small
\path{artifacts/ym_f15_v52_texcheck/}.
\end{center}
The next companion checkpoint is V55.

\begin{firewall}
V52 rules out a bounded or subexponential
unit-strength raw repair moment for the
specific canonical flat reference of V48,
even on already coherent sampled boundaries
and along volume-cofinal homogeneous laws.
It neither disproves the Yang--Mills mass gap
nor proves one. Alternative reference costs,
actual normalized-error control, scale-uniform
physical clustering, the interacting continuum
and the positive physical Hamiltonian gap
remain open in this program.
\end{firewall}
\section{V53: a reference-independent obstruction to raw flat-repair moments}
\label{f15v53:section}

V52 exposes an irrelevant-corner defect of the
canonical reference. A natural response is
to use only the outward-center connection,
or to optimize the flat reference altogether.
This section tests the strongest such response:
minimize the outside repair action over every
pure-gauge assignment on \(F_n\). This removes
the corner defect exactly. Nevertheless its
unit-strength exponential moment still grows
exponentially, on an event contained in the
actual coherent-fit event \(\Gamma_n\).

The new obstruction is not a path-choice
artifact. A weak coherent magnetic field
has bounded total original action, but
flattening its boundary destroys the flux
of many edge-disjoint loops. Distinct outward
mirror faces then charge a cost growing
linearly in the box side. A local tree-fiber
change of variables turns this deterministic
mismatch into an actual-law moment lower
bound, without paying for a neighborhood
of every torus link.

The direction is therefore adjusted upstream:
changing from one flat selector to another
cannot supply the desired uniform raw moment.
The curved coherent backgrounds already
developed in V42--V46, or compensation with
the actual joint normalizer, must be retained.
Neither of those alternatives is declared
closed by the present result.

\subsection{Nonduplication and the optimal flat debit}

The V52 source and its exact hash were
revalidated. The V42 coherent field, V43
fit definition, V47 mirror identities and
V51 one-face moment were reread. Searches
of the active notes, available archives
and closure monograph found no preceding
reference-independent raw flat-repair
moment obstruction. Boundary-forced flux
itself is an inherited V42 result, not a
new discovery. What is new here is its
bounded-full-action localization, a lower
bound uniform over all flat references,
and the normalized cofinal expectation
consequence. The V50 companion review
remains current; the next is V55.

Let \(V(F_n)\) be the endpoints of the
edges of \(F_n\), and let
\begin{equation}
 \mathcal F_n=
 \{R:R_{(x,y)}=h_x^{-1}h_y\text{ on }F_n,
                  \ h\in\SU(2)^{V(F_n)}\}.
 \label{eq:f15v53:flat-set}
\end{equation}
It is the compact set of pure-gauge
restrictions on this graph. Here ``flat''
means zero holonomy on every graph loop,
not merely on a chosen subset of elementary
plaquettes. No chosen
spanning tree or root path occurs in its
definition. Define the optimal flat debit
\begin{equation}
 D_{n,J}(Y)=\min_{R\in\mathcal F_n}
                           S_{\rm rest}(R,Y).
 \label{eq:f15v53:D}
\end{equation}
The minimum exists. It is continuous in
the finitely many exterior links involved:
the difference between minima at two
exteriors is bounded by the supremum,
over the compact set \(\mathcal F_n\),
of the corresponding action difference.
Uniform continuity makes that supremum
tend to zero. The value is gauge invariant,
since a vertex gauge maps \(\mathcal F_n\)
onto itself. No measurable minimizer
selection is needed anywhere below.

Every flat reference, whether canonical,
tree-based or variationally selected, has
\(T_{n,J}(Y)\ge D_{n,J}(Y)\) pointwise.
In particular V52's corner witness has
\(D_{n,J}=0\): every original rest-face
link is an identity and \(R=I\) is admissible.
Thus the present optimization really does
remove that earlier artifact.

Keep the full compact law \(\mu_J\) and
\(f(U)=1-\operatorname{Re}U=|U-I|^2/2\).
Assume
\((1-\tau)\beta\le J_p\le(1+\tau)\beta\)
and \(0\le\tau\le1/32\).
Throughout the construction take \(n\ge8\)
divisible by four, \(0<s\le1/8\), and
\(M\ge8n\). The fixed amplitude \(s\) is
a total coherent flux, not a plaquette
coupling or an auxiliary source.

\subsection{A localized weak magnetic field with bounded full action}

On integer coordinates define
\begin{equation}
 \phi_n(t)=\max\{0,\min(t,2n-t)\},\qquad
 \psi_n(t)=\max\{0,\min(1,2-|t|/n)\}.
 \label{eq:f15v53:cutoffs}
\end{equation}
Thus \(\phi_n(t)=t\) on \([0,n]\),
\(0\le\phi_n\le n\),
\(|\phi_n(t+1)-\phi_n(t)|\le1\), and
\(\psi_n=1\) on \([-n,n]\), with
\(|\psi_n(t+1)-\psi_n(t)|\le1/n\).
Let \(\theta=s/n^2\) and set
\begin{equation}
 \begin{split}
 a(x)&=\theta\phi_n(x_1)
             \psi_n(x_2)\psi_n(x_3)\psi_n(x_4),\\
 W_2(x)&=\exp(i a(x)\sigma_3),\qquad
 W_1(x)=W_3(x)=W_4(x)=I.
 \end{split}
 \label{eq:f15v53:magnetic-field}
\end{equation}
Use this field on the embedded box
\(\mathcal W_n=[-2n,2n]^4\), and identities
elsewhere on the torus. Each nonidentity
edge and every face incident to it lie
in this closed box. Indeed its direction
is two, \(0<x_1<2n\), and
\(-2n<x_j<2n\) for \(j=2,3,4\).
Its six incident faces therefore have
all vertices in \(\mathcal W_n\).
There is no periodic wrapping term.

Every link angle is at most \(s/n\).
Only orientations \((1,2),(2,3),(2,4)\)
can have nontrivial curvature. Their
plaquette phases are, up to sign,
\(a(x+\hat\mu)-a(x)\), with
\(\mu=1,3,4\), and their absolute values
are at most \(s/n^2\).

\begin{proposition}[Bounded total magnetic action]
\label{f15v53:magnetic-action}
The full homogeneous action without its
coupling factor satisfies
\begin{equation}
 S(W)\le s^2
       \frac{(8n^2+1)^2(4n^2+1)}{9n^6}
                         <30s^2.
 \label{eq:f15v53:action}
\end{equation}
On \(B_n\), the field is gauge equivalent
to V42's coherent background with
\(\Theta=s\). In particular its actual
boundary lies in \(\Gamma_n\), with
zero coherent-fit error.
\end{proposition}

\begin{proof}
Use \(1-\cos u\le u^2/2\), and denote
the one-dimensional sums of squares by
\(P_\phi,P_\psi\) and the sums of squared
forward differences by \(Q_\phi,Q_\psi\).
The functions vanish at both endpoints
of \([-2n,2n]\), so omitting the last
endpoint does not change a value-square
sum. Direct summation gives
\begin{equation}
 \begin{aligned}
 P_\phi&=(2n^3+n)/3,& Q_\phi&=2n,\\
 P_\psi&=(8n^2+1)/(3n),& Q_\psi&=2/n.
 \end{aligned}
 \label{eq:f15v53:one-dimensional-sums}
\end{equation}
The three nonzero orientations separate,
giving the full-action upper bound
\[
 \frac{\theta^2}{2}
       \{Q_\phi P_\psi^3+2P_\phi Q_\psi P_\psi^2\}
       =s^2\frac{(8n^2+1)^2(4n^2+1)}{9n^6}.
\]
For \(n\ge4\), the coefficient is at most
\((129/16)^2(65/16)/9<30\).
All faces outside the box are flat, so
no exterior contribution is omitted.

On \(B_n\), \(a(x)=\theta x_1\).
The gauge \(h_x=\exp(i\theta x_1x_2\sigma_3)\)
makes the direction-two links identities
and the direction-one links
\(\exp(-i\theta x_2\sigma_3)\).
This is exactly \eqref{eq:f15v42:background}.
Since \(\Theta=n^2\theta=s\le1/8\),
it is admissible in V43's fit minimization.
\end{proof}

\subsection{Many disjoint loops charge every flat reference}

On the facet \(x_4=0\), take the positively
oriented square loops
\[
 C_{a,z}=\partial([a,n-a]^2)\times\{z\}\times\{0\},
 \qquad
 1\le a\le n/4,\quad 1\le z\le n-1.
\]
Their side lengths are \(b_a=n-2a\), their
lengths are \(4b_a\), and their edge sets
are pairwise disjoint. Every such edge
belongs to \(F_n\): its sole transverse
boundary coordinate is \(x_4=0\), and
the other transverse coordinates are
strictly interior.

For these edges V47's outward path goes
to \(x_4=-1\). In the field \(W\), the
normal links are identities and
\(\psi_n(-1)=\psi_n(z)=1\). Consequently
the center connection has loop holonomy
\begin{equation}
 B(W)(C_{a,z})=\exp(i\alpha_a\sigma_3),
 \qquad \alpha_a=s b_a^2/n^2\in[s/4,s].
 \label{eq:f15v53:loop-flux}
\end{equation}

For any connection \(B\) on \(F_n\) and
any \(R\in\mathcal F_n\), telescoping
noncommuting products around a loop gives
\begin{equation}
 |B(C)-I|^2
       \le |C|\sum_{e\in C}|B_e-R_e|^2.
 \label{eq:f15v53:loop-mismatch}
\end{equation}
Here \(R(C)=I\), inversion is an isometry,
and the squared triangle inequality is
followed by Cauchy--Schwarz. No common
color axis for \(B\) and \(R\) is assumed.

The mirror face associated with \(e\in F_n\)
is distinct for each edge, belongs to the
rest action, and has exact energy
\(\tfrac12|R_e-B_e(Y)|^2\).
Since every other rest term is nonnegative,
\begin{equation}
 S_{\rm rest}(R,Y)\ge
       \frac{(1-\tau)\beta}{2}
                   \sum_{e\in F_n}|R_e-B_e(Y)|^2.
 \label{eq:f15v53:mirror-lower}
\end{equation}
This is the place where the original
Wilson interaction pays the mismatch;
it is not an artificially assigned penalty.

\begin{proposition}[A robust lower bound uniform over all flat completions]
\label{f15v53:all-flat-lower}
Suppose a configuration, in some vertex
gauge on \(\mathcal W_n\), has every link
within \(\eta=s^2/(9600n^4)\) of \(W\).
Then its exterior satisfies
\begin{equation}
 D_{n,J}(Y)\ge
                \frac{(1-\tau)\beta s^2 n}{1024}.
 \label{eq:f15v53:all-flat-lower}
\end{equation}
Its actual boundary also belongs to
\(\Gamma_n\).
\end{proposition}

\begin{proof}
First work in the stated gauge. The
outward-center holonomy on a loop of
length \(4b_a\) is a word of at most
\(12b_a\le12n\) original links.
Its distance from the value in
\eqref{eq:f15v53:loop-flux} is at most
\(12n\eta\le s/16\le\alpha_a/4\).
For \(0\le\alpha\le1/8\),
\[
 2\sin(\alpha/2)\ge\alpha-\alpha^3/24
                                  \ge3\alpha/4.
\]
Thus the perturbed loop chord is at
least \(\alpha_a/2\). This estimate uses
an elementary sine inequality, not an
abelian approximation to the perturbed field.
Equation~\eqref{eq:f15v53:loop-mismatch}
therefore gives, for every flat \(R\),
\[
 \sum_{e\in C_{a,z}}|B_e-R_e|^2
             \ge\frac{s^2 b_a^3}{16n^4}.
\]
The loops are edge disjoint. Their
side-length sum is exactly
\begin{equation}
 \sum_{a=1}^{n/4}(n-2a)^3
        =\frac{15n^4}{128}-\frac{7n^3}{16}
                                +\frac{3n^2}{8}
        \ge\frac{n^4}{16}\quad(n\ge8).
 \label{eq:f15v53:cube-sum}
\end{equation}
For example the identity follows from
\(\sum_{j=1}^m j^3=m^2(m+1)^2/4\);
the inequality follows on dividing by
\(n^4\) and using \(n\ge8\).
Sum over \(z\), then use the mirror
bound. The resulting lower bound is
\((1-\tau)\beta s^2(n-1)/512\), which
is at least the displayed bound.
Taking the minimum over all \(R\) is
now legitimate and changes nothing.

For coherence, apply the additional
fixed gauge \(h\) from the preceding
proposition. Every actual \(F_n\)-link
is then within \(\eta\) of V42's
coherent field with \(\Theta=s\).
Using \(\alpha(V)^2\le(\pi^2/4)|V-I|^2
<4|V-I|^2\), its admissible squared-angle
fit is at most
\[
 4d_n\eta^2\le96n^3\eta^2
       =\frac{s^4}{960000n^5}<\frac1{256n}.
\]
This proves \(\Gamma_n\). Finally, both
the loop chords and the optimized action
are gauge invariant, so the initial
choice of gauge imposes no restriction.
\end{proof}

\subsection{A tree-fiber twist with a volume-independent action debit}

Let \(\mathcal T_n\) be the union of
coordinate-ordered positive paths from
\(w=(-2n,-2n,-2n,-2n)\) to the vertices
of \(\mathcal W_n\), in order \(1,2,3,4\).
It is a spanning tree: for a nonroot
vertex, the last coordinate above its
lower endpoint determines its unique
parent by decrementing that coordinate.
A direction-two tree edge has
\(x_3=x_4=-2n\). Hence
\begin{equation}
                      W_e=I\qquad(e\in\mathcal T_n).
 \label{eq:f15v53:tree-zero}
\end{equation}
Write \(k_x(U)\) for the original-link
transport along the tree path to \(x\).
It depends only on the tree links.
Define a map on all original torus links by
\begin{equation}
 (\Phi_n U)_e=k_x(U)^{-1}W_e k_x(U)U_e
       \quad(e=(x,y)\subset\mathcal W_n),
 \label{eq:f15v53:tree-twist}
\end{equation}
and leave all other links unchanged.
Every tree link is unchanged by
\eqref{eq:f15v53:tree-zero}. Conditional
on those links, this is a product of
fixed left translations of the nontree
links. It is a Haar-preserving bijection,
whose inverse replaces every \(W_e\)
by \(W_e^{-1}\). The tree frames are
unchanged in both directions.
In the \(k(U)\) gauge its exact formula is
\begin{equation}
                    (\Phi_n U)^k_e=W_e U^k_e.
 \label{eq:f15v53:twist-frame}
\end{equation}
No nonlinear pushforward of an exterior
marginal has been declared measure preserving.
The map is a verified bijection of the
full original product-Haar space.

There are
\[
 N_n=6(4n)^2(4n+1)^2\le2400n^4
\]
positive plaquettes in \(\mathcal W_n\).
Define the local input event
\begin{equation}
 \begin{split}
 \eta&=\frac{s^2}{9600n^4},\qquad
 \epsilon=\frac{\eta}{12n}
                   =\frac{s^2}{115200n^5},\\
 \mathcal H_{n,s}&=\{\max_{p\subset\mathcal W_n}
                                      |U_p-I|\le\epsilon\}.
 \end{split}
 \label{eq:f15v53:H}
\end{equation}
Appending an edge to an ordered root path
and commuting its direction past later
segments requires at most \(12n\) square
interchanges, all inside \(\mathcal W_n\).
The noncommuting telescoping argument of
V52 then proves \(|U^k_e-I|\le\eta\)
on this event. By \eqref{eq:f15v53:twist-frame},
the output is within \(\eta\) of \(W\)
on every edge of the local box. Hence
\begin{equation}
 \Phi_n\mathcal H_{n,s}\subseteq\Gamma_n,
 \qquad
 D_{n,J}(Y(\Phi_n U))\ge
             (1-\tau)\beta s^2 n/1024
                  \quad(U\in\mathcal H_{n,s}).
 \label{eq:f15v53:output}
\end{equation}

Only links with \(W_e\ne I\) are changed.
As checked above, all faces incident to
these links lie in \(\mathcal W_n\).
The output holonomy on any such face is
within \(4\eta\) of the witness holonomy.
The function \(f\) is 1-Lipschitz in
chord distance. Original face actions
are nonnegative. Consequently the full
source-action difference obeys
\begin{align}
 S_J(\Phi_n U)-S_J(U)
 &\le(1+\tau)\beta\{S(W)+4N_n\eta\}\notag\\
 &\le31(1+\tau)\beta s^2
                       \quad(U\in\mathcal H_{n,s}).
 \label{eq:f15v53:full-debit}
\end{align}
The factor is 31, not the number of links
or the torus volume. Its finiteness comes
from the localized weak-curvature action;
bounding each altered face by two would
lose the statement needed here.

\subsection{The optimized raw moment grows even on coherent boundaries}

\begin{theorem}[Reference-independent actual-law lower bound]
\label{f15v53:moment}
Fix \(n_0=2^{16}\), \(0<s\le1/8\),
\(M\ge8n_0\) and the allowed coefficient
range above. For every such array \(J\),
\begin{equation}
 \mathbb E_J\!\left[e^{D_{n_0,J}(Y)}
                         \mathbf1_{\Gamma_{n_0}}\right]
       \ge e^{30\beta s^2}\mu_J(\mathcal H_{n_0,s}).
 \label{eq:f15v53:moment-all-J}
\end{equation}
If \(J\) is also in V51's actual-law class
\(\mathcal R^{[P]}_{M,\beta,\tau}\), then
\begin{equation}
 \mathbb E_J\!\left[e^{D_{n_0,J}(Y)}
                         \mathbf1_{\Gamma_{n_0}}\right]
 \ge e^{30\beta s^2}
       [1-N_{n_0}e^{K_P-\beta\epsilon^2/16}]_+.
 \label{eq:f15v53:cofinal-moment}
\end{equation}
Here \(\epsilon\) is fixed by
\eqref{eq:f15v53:H} with \(n=n_0\).
At fixed \(P,s\), bounded \((\log\beta)/M\)
and \(\beta\to\infty\), the bracket is
at least \(1/2\) eventually, uniformly
in that source class. Every flat reference
has at least the same moment lower bound.
\end{theorem}

\begin{proof}
For \(n=n_0\), the coefficient \(n/1024\)
in \eqref{eq:f15v53:output} equals 64.
Subtract \eqref{eq:f15v53:full-debit}:
\begin{align*}
 D_{n_0,J}(Y(\Phi_{n_0}U))-S_J(\Phi_{n_0}U)+S_J(U)
 &\ge\{64(1-\tau)-31(1+\tau)\}\beta s^2\\
 &=(33-95\tau)\beta s^2\\
 &\ge(961/32)\beta s^2>30\beta s^2.
\end{align*}
Integrate over \(\Phi_{n_0}\mathcal H_{n_0,s}
\subseteq\Gamma_{n_0}\) and use the
Haar-preserving bijection. Exactly as in
V52, the full normalizer cancels, giving
\begin{align*}
 \int_{\Phi_{n_0}\mathcal H_{n_0,s}}
                  e^{D_{n_0,J}(Y(V))}\,d\mu_J(V)
 &=\int_{\mathcal H_{n_0,s}}
   e^{D_{n_0,J}(Y(\Phi_{n_0}U))-S_J(\Phi_{n_0}U)+S_J(U)}
                                                   \,d\mu_J(U)\\
 &\ge e^{30\beta s^2}\mu_J(\mathcal H_{n_0,s}).
\end{align*}
All unchanged exterior action terms
cancel exactly; no independence of
the local input event is required.

The one-face bound
\eqref{eq:f15v51:one-face} gives
\(\mu_J\{|U_p-I|>\epsilon\}
\le e^{K_P-\beta\epsilon^2/16}\).
A union bound over the \(N_{n_0}\)
local faces proves
\eqref{eq:f15v53:cofinal-moment}.
At fixed \(P\) and bounded
\((\log\beta)/M\), \(K_P\) is bounded;
the local support and \(\epsilon>0\)
are fixed as well. The cofinal conclusion
follows. Finally any flat repair has
action at least \(D_{n_0,J}\) pointwise.
\end{proof}

The large side \(2^{16}\) and very small
curvature threshold are deliberately
conservative. This is an asymptotic
obstruction, not a useful numerical
onset estimate. The homogeneous class
\(P=1,\tau=0\) already suffices.
For example choosing \(s=1/16\) gives
the eventual lower bound
\(\tfrac12 e^{15\beta/128}\).

The same argument gives a useful
coefficient statement. For any fixed
\(\lambda>0\), choose a fixed multiple
of four \(n\ge8\) with
\(\lambda n/1024>31\).
In the homogeneous law the net exponent
for \(e^{\lambda D_{n,\beta\mathbf1}}\)
is at least
\[
       (\lambda n/1024-31)\beta s^2.
\]
The local probability again tends to one.
Thus no positive coefficient independent
of the box size supplies a subexponential
raw repair moment for every flat reference
at all box sizes. This does not contradict
V51's proved, box-dependent small coefficient.

\subsection{Why retaining the curved coherent family changes this witness}

Let \(\mathcal C_n\) be the compact set
of gauge transforms, restricted to \(F_n\),
of V42's coherent fields with
\(0\le\Theta\le1/8\). Define its
optimized debit by
\[
 D^{\rm coh}_{n,J}(Y)
               =\min_{R\in\mathcal C_n}S_{\rm rest}(R,Y).
\]
For the exact magnetic witness, its
actual boundary is an admissible member
of this set. All original plaquette
angles are at most \(s/n^2\), and V47
gives at most \(5d_n\) rest faces.
Therefore
\begin{equation}
 D^{\rm coh}_{n,J}(Y(W))
 \le S_{\rm rest}(W_{F_n},Y(W))
 \le\frac{5(1+\tau)\beta d_n s^2}{2n^4}
 \le\frac{60(1+\tau)\beta s^2}{n}.
 \label{eq:f15v53:curved-comparison}
\end{equation}
By contrast \eqref{eq:f15v53:all-flat-lower}
applies to the exact witness too and
gives an order-\(\beta s^2 n\) lower
bound for every flat reference. The
two costs differ by a factor growing
quadratically in \(n\), within the
already admitted coherent flux interval.
Thus the distinction is not between
a large, forbidden flux and the good
boundary regime. It is between retaining
and erasing a permitted weak coherent flux.

Equation~\eqref{eq:f15v53:curved-comparison}
is a deterministic comparison for this
witness, not a proved global moment for
\(D^{\rm coh}\). It identifies a concrete
next interface: combine V45--V46's
curved coherent normalizer with an
exterior-only choice or variational
integration that retains the actual
outside action. More general curvature
and source dependence may still require
larger background families. The other
legitimate alternative is to retain the
joint compensation with the actual
conditional normalizer rather than to
estimate a raw repair exponential alone.

Theorem~\ref{f15v53:moment} rules out
neither alternative. It also does not
give a lower bound on
\(\mu_J(\Gamma_n^c)\): its contribution
is supported on \(\Gamma_n\). V37's
criterion concerns an actual normalized
error, not the optimized raw flat debit.
That criterion, arbitrary-source closure,
physical clustering and the full continuum
mass-gap objective remain unproved.

\subsection{Exact diagnostics and preservation}

The diagnostic is
\begin{center}\small
\path{scripts/verify_ym_f15_v53_all_flat_reference_obstruction.py}.
\end{center}
It passes the unchanged V47 checks and
verifies the separable action sums and
the cubic loop-count identity at six
box sides. Exact phase cochains check
110 disjoint magnetic loops and 4000
distinct mirror faces. Three noncommuting
rational quaternion cases test the
flat-loop mismatch inequality. A separate
small-box calculation checks 2000 exact
tree-fiber twist, inverse and gauged
multiplication identities, including that
every changed face stays in the local box.
Eight rational margin checks verify the
perturbation and coherence constants.

These are checks of finite identities,
not samples proving Haar invariance,
cofinal probability estimates or the
general-size inequalities. Those claims
are proved explicitly above. In particular
no computation on the enormous box
\(n_0=2^{16}\) is being represented as
an exhaustive enumeration.

The script has 7,104 bytes and SHA--256
\begin{center}\scriptsize\ttfamily
7B3D1032AC4E6B2F6108DA4D18042C87463DA506408A75A2981FF00C6F8167B9.
\end{center}
The V52 predecessor has 1,214,661 bytes and SHA--256
\begin{center}\scriptsize\ttfamily
D56CB2085BC724C23BFA07C7E1D2DDA9206849B92D994525D459C92F20272F71.
\end{center}
Removing this final section and reversing
the title increment recovers V52 exactly.
V53 replaces V52 as the sole active main
source. The 14 protected inputs and 95
earlier diagnostics remain unchanged.
Only \path{.tex} sources are kept in
\path{latex_notes}; compilation output
is directed to
\begin{center}\small
\path{artifacts/ym_f15_v53_texcheck/}.
\end{center}
The next companion checkpoint is V55.

\begin{firewall}
V53 proves exponential growth of the
optimized raw flat-repair moment on
actually coherent sampled boundaries,
uniformly over every flat reference,
in a specified volume-cofinal regime.
It eliminates a whole flat-completion
strategy for obtaining the proposed
bounded unit-strength moment. It does
not eliminate curved coherent backgrounds
or joint normalization, and it neither
proves nor disproves the Yang--Mills
mass gap. The original continuum and
physical Hamiltonian-gap goal remains
active and unproved.
\end{firewall}
\section{V54: actual-posterior normalization for reflected sparse physical sources}
\label{f15v54:section}

V52--V53 rule out the proposed bounded raw
flat-repair moment, even after optimizing
the flat reference. The present step returns
to V37's actual-background criterion instead
of trying to repair that false intermediate
target. Its conditional source moments can
be estimated directly in the actual tilted
law, with both signs allowed. V51's joint
cell moment then supplies the required
probabilistic input on its reflected source
class. This closes an intensive normalization
interface at fixed nonzero source strength
for a specified sparse reflected family.

This is a restricted closure, not an
arbitrary-source or mass-gap theorem.
Reflection ties the amplitudes of copies
of each source. The local reference is
still not identified with a physical
next-scale action, and the retained
baseline correlations are not shown to
decay. All these distinctions remain
part of the result, not deferred qualifications.

\subsection{Inherited inputs and the precise source domain}

The V53 predecessor was checked against
its recorded hash. V21's exact common
exterior disintegration, V22's entire-core
chart ratio, V35's capped reference and
mixed-sign local error, V37's actual-law
criterion, and V51's physical source
moment were reread. Their elementary
entropy and subset-expansion arguments
are reused rather than counted as new
theorems. The new bridge is the direct
actual-posterior bound on baseline
conditional source moments, followed
by its Wilson-law realization. The
V50 companion checkpoint is current;
the next iteration, V55, requires the
next companion review.

Use \(N\) separated canonical interiors
\(J_i\), their cores \(V_i\), near
retained inputs \(N_i\), and common
exterior \(\xi=(C,v_O)\) as in V21--V35.
Let \(\mu_0\) be the homogeneous baseline
Wilson law. Its full compact conditional
laws \(\mathsf P_i^\xi\) factor after
fixing \(\xi\), with no shared measured
interaction between selected interiors.

Allow nonnegative cell-energy observables
\begin{equation}
 Y_i=\sum_{x\in A_i}T_x=Y_i(N_i,V_i),\qquad
 1\le |A_i|\le v,
 \label{eq:f15v54:energies}
\end{equation}
with disjoint fine base sets \(A_i\).
Their canonical carriers must lie in
their own cores, with retained parameters
in \(N_i\). The old block-energy sources
are included whenever these conditions
and the actual-law hypothesis below hold.
An explicit additional singleton-cell
family satisfying all conditions is
constructed at the end of this section.

For \(|\theta_i|\le\tau\le1/512\), use
the unchanged physical centering convention
\begin{equation}
 \frac{d\mu_\theta}{d\mu_0}
 =\frac{\exp\{\sum_i\theta_i(Y_i-\mathbb E_{\mu_0}Y_i)\}}
        {\mathbb E_{\mu_0}
         \exp\{\sum_i\theta_i(Y_i-\mathbb E_{\mu_0}Y_i)\}}.
 \label{eq:f15v54:actual-law}
\end{equation}
Assume its actual Wilson coefficient
array \(J\) belongs to
\(\mathcal R^{[P]}_{M,\beta,\tau}\) from V51.
Write \(K_P\) for the V51 constant.
This membership is an explicit restriction,
not a property inferred from separation.
In particular an arbitrary further local
source change need not preserve it.

The joint cell moment available at this
one actual law is
\begin{equation}
 \log\mathbb E_{\mu_\theta}
               e^{\sum_x b_xT_x}
              \le2K_P\sum_xb_x,
                  \qquad 0\le b_x\le1/2.
 \label{eq:f15v54:actual-CGF}
\end{equation}
It is \eqref{eq:f15v51:CGF}, with its
block-maximum bound weakened to the
sum bound. No reflection symmetry of
the subsequently inserted test array
\(b\) is required.

\subsection{A pointwise conditional inequality avoids global transport}

Put \(t_i=|\theta_i|\),
\(u_i=\mathsf P_i^\xi(e^{\theta_iY_i})\),
and retain V35's unmodified quantity
\[
              g_i=\log\mathsf P_i^\xi(e^{t_iY_i})\ge0.
\]
The actual conditional law on \(J_i\) is
\[
 d\mathsf P_{i,\theta}^\xi
         =u_i^{-1}e^{\theta_iY_i}\,d\mathsf P_i^\xi.
\]
The other source factors cancel only
after fixing the common exterior.
These actual conditionals are still
independent at that fixed exterior.
Let \(\nu_\theta=\xi_*\mu_\theta\).

\begin{proposition}[Baseline conditional moments paid in the actual law]
\label{f15v54:conditional-moment}
For every \(a\ge1\), pointwise in \(\xi\),
\begin{equation}
 e^{a g_i}\le
  \mathsf P_{i,\theta}^\xi
                  (e^{(a+1)t_iY_i}).
 \label{eq:f15v54:pointwise-moment}
\end{equation}
For \(\theta_i\ge0\), the coefficient
\(a+1\) can be replaced by \(a\).
Consequently, for every subfamily \(I\)
and \((a+1)\tau\le1/2\),
\begin{equation}
 \log\mathbb E_{\nu_\theta}
              e^{a\sum_{i\in I}g_i}
                 \le2(a+1)K_P\tau v|I|.
 \label{eq:f15v54:actual-g}
\end{equation}
\end{proposition}

\begin{proof}
Fix one conditional input and omit its
indices. If \(\theta=t\ge0\), put
\(M=\mathsf P(e^{tY})=u\).
Jensen's inequality gives
\[
 M^{a+1}\le\mathsf P(e^{(a+1)tY})
                 =M\mathsf P_\theta(e^{atY}).
\]
Divide by \(M>0\). If \(\theta=-t\le0\),
then \(u=\mathsf P(e^{-tY})\le1\), since
\(Y\ge0\). Jensen now gives
\[
 M^a\le\mathsf P(e^{atY})
             =u\mathsf P_\theta(e^{(a+1)tY})
             \le\mathsf P_\theta(e^{(a+1)tY}).
\]
The zero-source case has both sides
equal to one. These prove the pointwise
inequality for either sign.

Multiply over \(I\), use the actual
conditional product law, and integrate
under its actual exterior marginal:
\[
 \mathbb E_{\nu_\theta}e^{a\sum_{i\in I}g_i}
 \le\mathbb E_{\mu_\theta}
                     e^{(a+1)\sum_{i\in I}t_iY_i}.
\]
The cell coefficients on the right
are at most \((a+1)\tau\), and their
sum is at most \((a+1)\tau v|I|\).
Disjointness of the \(A_i\)'s and
\eqref{eq:f15v54:actual-CGF} prove
\eqref{eq:f15v54:actual-g}.
\end{proof}

The argument has not integrated the
baseline conditionals against the
wrong exterior marginal. It first
compares them pointwise to actual
conditionals, then uses the latter's
matched product identity. There is
no global order-two likelihood cost,
and no missing contribution from
remote source normalizers.

\subsection{Raw chart and cap failures under the same actual source}

Keep V35's source-independent kernel
\(\lambda_i^\Lambda(\cdot\mid N_i)\),
formed by restricting the measured
chart reference to
\(Y_i-m_i(N_i)<\Lambda\), with its
strict-unit-sublevel Haar fallback.
Here \(m_i=\min_{V_i}Y_i\ge0\).
The minimum and the two normalizers
remain in the definition. This
construction works for every observable
in \eqref{eq:f15v54:energies}: V22's
ratio concerns the entire core law,
not one particular source observable.
Compactness, continuity and core
measurability supply V35's cap argument.

Here the exterior event \(\Gamma_i\)
is V21's retained half-chart event,
not V43's coherent-fit event \(\Gamma_n\).
Let \(E_i^\Lambda\) be the same union
of exterior chart failure, interior
chart failure, and cap failure as
in \eqref{eq:f15v35:failure}. Each
uncapped chart failure is contained
in a rooted-patch failure on an
embedded box of side \(L\), at
threshold \(\delta\), by V13/V21.
Assume the separated patches have
disjoint plaquette-base sets and put
\begin{equation}
 \begin{split}
 w&=(L+1)^4,\qquad
 u=\frac{\beta\delta^2}{18\pi^2L^2},\\
 q_{\rm ch}&=\min\{1,we^{K_P-u/2}\},\\
 q_Y&=\min\{1,e^{K_Pv-\Lambda/2}\},\\
 q&=\min\{1,\sqrt{q_{\rm ch}}+\sqrt{q_Y}\}.
 \end{split}
 \label{eq:f15v54:q}
\end{equation}
Then, under the actual law,
\begin{equation}
           \mu_\theta\Bigl(\bigcap_{i\in I}E_i^\Lambda\Bigr)
                                                   \le q^{|I|}.
 \label{eq:f15v54:raw}
\end{equation}

For completeness, the only new input
in this application of the old witness
argument is the law in which its
moment is taken. V20's deterministic
inclusion chooses a cell with \(T_x>u\)
from at most \(w\) candidates in each
bad patch. For one choice per patch,
\eqref{eq:f15v54:actual-CGF} at coefficient
\(1/2\) bounds the intersection by
\(e^{(K_P-u/2)|I|}\). The union bound
gives \(q_{\rm ch}^{|I|}\).
A cap failure implies \(Y_i\ge\Lambda\),
and the same CGF gives \(q_Y^{|I|}\).
Expand the union of the two types,
apply Cauchy--Schwarz to each mixed
intersection, and sum. This gives
\eqref{eq:f15v54:raw}. No conditional
version of the CGF or independence
of fine cell energies is assumed.

\subsection{Both confidence probabilities controlled in the actual posterior}

Retain exactly V35's two probabilities
\[
 \pi_{i,0}=\mathsf P_i^\xi(E_i^\Lambda),\qquad
 \pi_{i,\theta}=\mathsf P_{i,\theta}^\xi(E_i^\Lambda),
 \qquad
 B_i=\{\pi_{i,0}>\eta\text{ or }\pi_{i,\theta}>\eta\}.
\]
They are both functions of the common
exterior. Define, in this proof only,
\[
 H_i=\mathsf P_{i,\theta}^\xi
                (1_{E_i^\Lambda}e^{t_iY_i}),\qquad
 M_i=\mathsf P_{i,\theta}^\xi(e^{t_iY_i}).
\]
The baseline probability can be written
with its actual-law normalization:
\[
 \pi_{i,0}=
 \frac{\mathsf P_{i,\theta}^\xi
                    (1_{E_i^\Lambda}e^{-\theta_iY_i})}
      {\mathsf P_{i,\theta}^\xi(e^{-\theta_iY_i})}.
\]
The numerator is at most \(H_i\).
If \(\theta_i\ge0\), Cauchy--Schwarz
makes the denominator at least \(1/M_i\).
If \(\theta_i\le0\), the denominator
is at least one. Since \(M_i\ge1\),
both signs, as well as the actual
probability itself, obey
\begin{equation}
                   \pi_{i,0},\ \pi_{i,\theta}\le H_iM_i.
 \label{eq:f15v54:confidence-pointwise}
\end{equation}

For a subfamily of size \(k\), actual
conditional Jensen, the product identity,
and Cauchy--Schwarz give
\begin{align*}
 \mathbb E_{\nu_\theta}\Bigl(\prod_{i\in I}H_i\Bigr)^2
 &\le\mathbb E_{\mu_\theta}
       [1_{\cap_{i\in I}E_i^\Lambda}e^{2\sum_{i\in I}t_iY_i}]\\
 &\le q^{k/2}e^{4K_P\tau vk},\\
 \mathbb E_{\nu_\theta}\Bigl(\prod_{i\in I}M_i\Bigr)^2
 &\le e^{4K_P\tau vk}.
\end{align*}
The energy coefficients used here are
\(4t_i\) and \(2t_i\), within the
available window. One more
Cauchy--Schwarz and expansion of the
two alternatives in \(B_i\) prove
\begin{equation}
 \begin{split}
 \mathbb E_{\nu_\theta}\prod_{i\in I}\pi_{i,a_i}
       &\le\gamma^{|I|}\quad(a_i\in\{0,\theta\}),\\
 \nu_\theta\Bigl(\bigcap_{i\in I}B_i\Bigr)
       &\le\left(\frac{2\gamma}{\eta}\right)^{|I|},\qquad
 \gamma=q^{1/4}e^{4K_P\tau v}.
 \end{split}
 \label{eq:f15v54:confidence}
\end{equation}
The new conclusion is specifically
under \(\nu_\theta\); the earlier
baseline confidence estimate is not
being silently transported to it.

\subsection{V37's normalization criterion is now verified on this domain}

Let \(\Delta\) be the inherited entire-core
chart log-ratio budget, and keep
\begin{equation}
 \begin{split}
 F&=\sum_{i=1}^N(g_i+\tau\Lambda)1_{B_i},\\
 e_0&=2\Delta-\log(1-\eta),\qquad
 \left|\sum_i\log(u_i/b_i)\right|\le Ne_0+F,\\
 b_i&=\lambda_i^\Lambda(e^{\theta_iY_i}).
 \end{split}
 \label{eq:f15v54:error}
\end{equation}
V35's local inequality supplies the
last bound for either sign.
The new moment inputs just proved give
\begin{equation}
 \begin{split}
 \mathbb E_{\nu_\theta}e^{4\sum_{i\in I}g_i}
       &\le e^{10K_P\tau v|I|},\\
 \zeta&=\sqrt{2/\eta}\,
          e^{2\tau\Lambda+7K_P\tau v}q^{1/8},\qquad
 \kappa=\log(1+\zeta).
 \end{split}
 \label{eq:f15v54:zeta}
\end{equation}
The first line uses
\eqref{eq:f15v54:actual-g} at \(a=4\);
\(5\tau\le1/2\). For each subset,
the two preceding moment/confidence
bounds and Cauchy--Schwarz yield
\[
 \mathbb E_{\nu_\theta}
 [e^{2\sum_{i\in I}(g_i+\tau\Lambda)}1_{\cap_{i\in I}B_i}]
                                                   \le\zeta^{|I|}.
\]
Expanding \(e^{2F}\) as in V37 therefore
proves the actual-background criterion
\begin{equation}
                    \log\mathbb E_{\nu_\theta}e^{2F}
                                                     \le N\kappa.
 \label{eq:f15v54:criterion}
\end{equation}

\begin{theorem}[Fixed-strength intensive normalization in the stated actual laws]
\label{f15v54:normalization}
Let \(\rho_a,\rho_b,Z_a,Z_b\) be the
actual retained source law, capped-reference
retained source law and their full
normalizers from V35--V37, using the
observables and kernel specified here.
Put \(\delta_\theta=\log(Z_b/Z_a)\) and
\(\varepsilon_{\rm act}=e_0+\kappa/2\).
Then
\begin{equation}
 \begin{gathered}
 |\delta_\theta|\le N\varepsilon_{\rm act},\qquad
 D(\rho_a\Vert\rho_b)\le2N\varepsilon_{\rm act},\\
 D(\rho_b\Vert\rho_a)\le4N\varepsilon_{\rm act},\qquad
 d_{\rm TV}(\rho_a,\rho_b)
                  \le\min\{1,\sqrt{N\varepsilon_{\rm act}}\}.
 \end{gathered}
 \label{eq:f15v54:normalization}
\end{equation}
Suppose further that along a regulator sequence
\begin{equation}
 \Delta\longrightarrow0,\qquad
 \Lambda\longrightarrow\infty,\qquad
 \Lambda\le u,\qquad
 K_Pv+\log w=o(\Lambda),\qquad
 \eta=e^{-\Lambda/128}.
 \label{eq:f15v54:limit-conditions}
\end{equation}
Then \(\varepsilon_{\rm act}\to0\), uniformly
for source arrays in this domain with
\(|\theta_i|\le\tau\le1/512\).
The strength \(\tau\) need not tend to zero.
There is no restriction on \(N\) beyond
the stipulated fitting and separation.
\end{theorem}

\begin{proof}
Equation~\eqref{eq:f15v54:criterion} verifies
the hypothesis of Theorem~\ref{f15v37:criterion}.
Its exact normalized exterior likelihood,
including centering constants, gives
\eqref{eq:f15v54:normalization}. Thus no
new reference-law moment or order-two
change-of-law budget is inserted here.

For the limit use \(v\ge1\) and \(u\ge\Lambda\)
in \eqref{eq:f15v54:q}:
\[
 q\le2\exp\{(K_Pv+\log w)/2-\Lambda/4\}.
\]
Substitution into \eqref{eq:f15v54:zeta},
with \(\eta=e^{-\Lambda/128}\), gives
\begin{equation}
 \zeta\le2^{5/8}\exp\left\{
 -\frac{3\Lambda}{128}
 +(7\tau+1/16)K_Pv+\frac{\log w}{16}\right\}.
 \label{eq:f15v54:vanishing}
\end{equation}
Indeed
\(1/32-2\tau-1/256\ge3/128\)
on the stated strength interval.
The hypotheses make the negative term
dominate, so \(\zeta\to0\),
\(\kappa\to0\), and \(e_0\to0\).
All displayed budgets are independent
of the number of unused or other source
cells; the dependence on \(P\) and \(v\)
has been retained explicitly.
\end{proof}

The error budget \(F\) is the conditional
log-likelihood error in V37, not the
raw flat debit \(D_{n,J}\) of V53.
Its moment estimate therefore does not
contradict the obstruction proved there.

This removes V35's nonvanishing cubic
transport term on this restricted domain.
It does not contradict V37's rare-phase
countermodel: that model did not supply
the actual joint CGF and actual raw
failure rate used in this proof.
An intensive error tending to zero
also does not imply total TV convergence
when \(N\) grows; the additional
condition \(N\varepsilon_{\rm act}\to0\)
would still be required for that conclusion.

\subsection{A nonzero signed reflected family with genuinely separated cores}

It is important not to leave the
actual-law membership condition with
only the zero vector as an example.
Here is an explicit realization on
the unchanged blocking hierarchy.
Use singleton energies \(Y_i=T_{x_i}\),
so \(v=1\). This is an additional
physical test family, not a relabelling
of V16's block-energy observables.

Take \(p\ge2\) dyadic, a collar
\(r\ge368\), \(L=2p(r+5)\), and
choose a dyadic \(P\) with
\(16L\le P<32L\). Then \(4p\mid P\).
Assume \(M/P\) is a power of two at
least four. Inside \(Q_P\), place two
cell-source coefficients at
\begin{equation}
 a=(P/4,P/4,P/4,P/4),\qquad
 b=(3P/4,P/4,P/4,P/4),
 \label{eq:f15v54:prototypes}
\end{equation}
with values \(\theta_a,\theta_b\in[-\tau,\tau]\);
put zero at every other cell. Extend
this profile by V51's fold in each
coordinate. Its nonzero-support candidates
form \(2(M/P)^4\) singleton sites,
one copy of each prototype in every
reflected block. Zero amplitudes may
be retained as zero sources.

For a coordinate value \(c\in\{P/4,3P/4\}\),
its preimages in a period \(2P\) are
\(c\) and \(2P-1-c\). Assign the
corresponding coarse-compatible anchors
\(c\) and \(2P-c\), respectively,
and translate this rule by multiples
of \(2P\). Thus every cell is within
one lattice unit in each coordinate
of its assigned \(p\)-grid anchor.
Distinct anchors have periodic maximum-norm
separation at least \(P/2\).
For V20's mesh
\(\ell=2^{\lceil\log_2L\rceil-2}<L/2\),
this is more than \(8\ell\).
The canonical interiors and their
rooted witness patches therefore have
exactly the separation used above.

The carrier condition is also literal,
not a reflection covariance assumption
about the blocking coordinates. A cell
energy uses only the faces of its unit
cell. V9/V12's local reconstruction puts
the required canonical details within
\(8p\) of the fine carrier. The
one-unit anchor adjustment and the
unit faces keep them within \(12p\)
of the assigned coarse anchor for
\(p\ge2\). Vertex gauge factors cancel
in the Wilson traces. All detail variables
are therefore inside the core of radius
\(\lfloor r/2\rfloor p\), and all retained
parameters are in its declared near
output. The same baseline entire-core
chart comparison and capped-kernel
construction apply. Nothing has reflected
or changed the completed blocking map.

Let \(h_x\) denote this physical cell
coefficient field, not a vertex gauge.
The exact cell-source lift gives
\[
 J_{(x;\mu,\nu)}=\beta\left[
 1-\frac14\sum_{\epsilon\in\{0,1\}^{\{\mu,\nu\}^c}}
                   h_{x-\epsilon}\right].
\]
V51 proves that the folded profile
makes this array belong to
\(\mathcal R^{[P]}_{M,\beta,\tau}\).
The source is visible: at a positive
\((1,2)\) face based at \(a\), only
the cell at \(a\) contributes to the
four-cell average, so its lift is
\(\theta_a/4\). Choosing
\(\theta_a=\tau\), \(\theta_b=-\tau\)
gives a genuinely nonzero signed family.
More generally the two amplitudes vary
independently over the indicated square,
but their reflected copies are tied.
This is not an independent source cube
with \(2(M/P)^4\) free coordinates.

\subsection{An explicit fixed-depth regime with vanishing intensive error}

Fix \(p\) and an accuracy exponent
\(s_0>0\). Use V32's collar
\eqref{eq:f15v32:scale}, with \(s=s_0\)
there and \(t=\log\beta\). At fixed
\(p\), one has \(r=O_p(t)\),
\(L=O_p(t)\), and the preceding choice
has \(P=O_p(t)\). Retain the selected
homogeneous regulator sequence if needed
for the inherited setup; its dyadic
torus size is more than sufficient
for the displayed geometry.
More generally impose the stated
embedding thresholds and a bounded
\((\log\beta)/M\).

The deterministic chart comparison gives
\(\Delta\le\beta^{-s_0}\) eventually.
The rooted witness threshold is
\(\delta=a_p/(64rp^4)\), so
\[
 u\asymp_p\frac{\beta}{t^4},\qquad
 K_P=O_p(t^4),\qquad \log w=O_p(\log t).
\]
Here the constants also depend on the
fixed accuracy and the stated bound
on \((\log\beta)/M\), not on the
source amplitudes or their number.
Choose the source-independent cap and
confidence threshold
\begin{equation}
                    \Lambda=t^6,\qquad
                    \eta=e^{-t^6/128}.
 \label{eq:f15v54:explicit-scale}
\end{equation}
For \(v=1\), all conditions in
\eqref{eq:f15v54:limit-conditions} now
hold, and the theorem applies uniformly
to \((\theta_a,\theta_b)\in[-\tau,\tau]^2\)
with any fixed \(0<\tau\le1/512\).
The cap was enlarged relative to V35's
old choice in order to pay the growing
reflection-block volume \(P^4\).
It has not been made dependent on the
individual source parameters.

This closes the specified intensive
normalization step for an extensive
reflected sparse source family at
fixed block depth. It does not remove
the amplitude ties, prove a uniform
statement for increasing depth, bound
the retained connected-response norm,
or establish a physical RG contraction.
The reference functions are still
averaged against the correlated retained
baseline. No decay of that baseline,
complex-source analyticity, derivative
convergence or Hamiltonian gap follows
from the real pressure/entropy bound.

The next upstream issue is thus more
precise: extend actual-law moment and
failure control beyond reflected source
orbits, or control connected responses
without requiring that extension.
The raw flat-repair moment is not
reintroduced in either formulation.

\subsection{Exact diagnostics and preservation}

The diagnostic is
\begin{center}\small
\path{scripts/verify_ym_f15_v54_actual_posterior_normalization.py}.
\end{center}
Twenty-seven signed finite conditional
models retain a nontrivial common
exterior and every actual source
normalizer. They check 972 pointwise
conditional moment inequalities,
648 actual-posterior product bounds,
243 confidence comparisons and 729
joint mixed-confidence fourth-power
inequalities. There are 216 exceptional
subset checks and 27 exact full
product expansions.

Two sparse reflected source geometries
check 16,384 cell reflection identities,
261,632 separated core pairs, the
one-unit anchor adjustments and a
nonzero physical plaquette lift.
Rational exponent bookkeeping checks
the coefficient \(7K_P\tau v\) and
the negative margin \(3\Lambda/128\).
The Wilson moment, chart comparison
and limiting statements use their
analytic proofs and stated hypotheses;
they are not inferred from these finite
model tests.

The script has 7,314 bytes and SHA--256
\begin{center}\scriptsize\ttfamily
99E6F05B317202466E20B136C64648DE5C19CEB04D9C23008FFE600955C7318E.
\end{center}
The V53 predecessor has 1,236,284 bytes and SHA--256
\begin{center}\scriptsize\ttfamily
353C037888061F3623FF530BAA4A16AB987DB00616075A1DDBE397086556196E.
\end{center}
Removing this final section and reversing
the title increment recovers V53 exactly.
V54 replaces V53 as the sole active main
source. The 14 protected inputs and 96
earlier diagnostics remain unchanged.
Only \path{.tex} sources are kept in
\path{latex_notes}; compilation products
are directed to
\begin{center}\small
\path{artifacts/ym_f15_v54_texcheck/}.
\end{center}
The next companion checkpoint is V55.

\begin{firewall}
V54 verifies V37's actual-background
normalization criterion and obtains
vanishing intensive pressure and entropy
errors at fixed nonzero source strength
for a specified reflected sparse family.
It retains the actual posterior, all
conditional normalizers, the source-orbit
restriction and the reflection-block cost.
It does not establish arbitrary-source
closure, scale-uniform physical clustering,
an interacting continuum or the positive
physical Yang--Mills Hamiltonian gap.
\end{firewall}

\section{V55: exponential-excess transport and independent mixed-sign sources}
\label{f15v55:section}

\subsection{Context, nonduplication and the companion checkpoint}

V54 paid the actual-background normalization error
for a specified reflection-compatible source family.
It did not allow every source amplitude to vary
independently. We now return to the original
block-energy sources of V16 and V26, keep V35's
source-independent capped reference, and remove
the reflection restriction at the price of an
explicit total source-tilt budget.
The useful quantity to transport is the excess
of an exponential moment above one, not the
whole exponential moment. This preserves a
small exceptional-event factor.
The elementary inequality used for this purpose
is not presented as a new general inequality.
The advance here is its application with a larger
cap, a paid actual-law likelihood, and a growing
fixed-strength source cube for which even
\emph{total} normalization errors vanish.

The V55 companion checkpoint reread the latest
conclusions of the seven relevant companion notes
and all of the supplied suggestions, after a fresh
inventory of the note folder and Downloads.
No newer relevant files or changed protected
inputs were found. SM continuum V72 gives an
auxiliary fixed-time, finite-mode many-copy
covariance limit, not a physical cutoff-uniform
limit. NS stability V26 contains pointwise
rank-one formulas; a proposed later certification
is not imported as an accomplished result.
Shell mixing V16 supplies a measured flat-holonomy
one-loop assembly, but not the missing noncommuting
coefficient and common nonperturbative remainder.
Parallel YM V5 retains the global-charge
energy--entropy obstruction and the distinction
between refinement coherence and tightness.

The supplied SM source-selection V31 has a
conditional word-native typing compiler, not a
proof of the required historical occurrence.
Native regenerative stationarity V34 reduces the
first-escape question to moving-fibre derivatives;
the 330 indicated first-derivative entries remain
uncomputed. Exact-action Einstein bridge V24
establishes a nonzero exact first-jet derivative,
not an exact nonlinear root. None of these
results supplies the general actual-source
estimate or physical clustering missing here.
The suggestions reinforce the need to retain
the measured density and source conventions.
They are mathematical proposals, not instructions
to merge an unproved determinant or continuum
claim into this programme.
The next companion checkpoint is V60.

This section does not repeat V35's finite
cap construction, V37's entropy criterion,
or V54's reflected argument. V35's old
fixed-strength transport remainder need not
vanish; V54 addresses a restricted source
pattern at all fitting cardinalities.
Here the coefficients are unrestricted within
a real cube, but its cardinality is restricted.
The conclusions and restrictions are different.

\subsection{The original sources and the exact normalized laws}

Work under the inherited selected-coupling
source bounds, separated-core factorization,
chart comparison and embedding hypotheses
of V26, V32 and V35. No claim is made that
this section constructs these inputs afresh.
Fix a dyadic depth \(p\ge2\), and take \(N\ge1\)
separated cores fitting the stipulated geometry.
Write the baseline Wilson law as \(\mu_0\).
The source-cell sets \(\mathcal B_{p,i}\) are
disjoint, and
\begin{equation}
 \begin{gathered}
 n_p=p^2(2p-1)^2\le4p^4,\qquad
 Y_i=\sum_{x\in\mathcal B_{p,i}}T_x\ge0,\qquad
 S_i=Y_i-\mathbb E_{\mu_0}Y_i,\\
 0<\tau\le1/512,\qquad \theta\in[-\tau,\tau]^N,\\
 W_\theta=e^{\sum_i\theta_i S_i},\qquad
 Z_a(\theta)=\mathbb E_{\mu_0}W_\theta,\qquad
 L_\theta=W_\theta/Z_a(\theta),\qquad
 d\mu_\theta=L_\theta\,d\mu_0 .
 \end{gathered}
 \label{eq:f15v55:sources}
\end{equation}
Every \(\theta_i\) is independently selectable,
with either sign. There is no actual-source
reflection symmetry or orbit tying assumption.
The centering is always the baseline centering.
At finite regulator all displayed source
normalizers and second moments are finite.

Let \(\xi=(C,v_O)\) be the actual common
exterior, \(\nu_0=\xi_*\mu_0\) and
\(\nu_\theta=\xi_*\mu_\theta\). Choose
\(\Lambda\ge1\) independently of \(\theta\),
and use exactly V35's capped reference
\(\lambda_i^\Lambda\): the same chart sublevel
\(Y_i-m_i(N_i)<\Lambda\), with the same
strict unit-sublevel Haar fallback when its
reference chart mass vanishes.
Here \(m_i(N_i)\) is the core minimum of \(Y_i\).
Changing the cap does not alter the Wilson law,
the completed blocking, or the physical source.

Retain \(u_i,b_i,g_i,\pi_{i,0},\pi_{i,\theta}\)
from \eqref{eq:f15v35:local}, now for these
sources and this cap, and set
\begin{equation}
 \begin{gathered}
 B_i=\{\pi_{i,0}>\eta\ \hbox{or}\ \pi_{i,\theta}>\eta\},
 \qquad
 F_\theta=\sum_{i=1}^N(g_i+\tau\Lambda)1_{B_i}\ge0,\\
 \ell_\theta=\sum_{i=1}^N\log(u_i/b_i),\qquad
 e_0=2\Delta-\log(1-\eta),\qquad
 |\ell_\theta|\le Ne_0+F_\theta .
 \end{gathered}
 \label{eq:f15v55:local}
\end{equation}
This is a bound for conditional normalized
log errors, not a raw unnormalized flat debit.
The chart event used here is the inherited
retained half-chart event of V21, not the
coherent many-copy event bearing similar
notation in V43.

For clarity the two retained laws are still
those in \eqref{eq:f15v35:laws}:
\begin{equation}
 \begin{gathered}
 \rho_0=C_*\mu_0,\qquad \rho_a=C_*\mu_\theta,\qquad
 T_\theta(C)=\prod_i\lambda_i^\Lambda(e^{\theta_i S_i}),\\
 Z_b(\theta)=\mathbb E_{\rho_0}T_\theta,\qquad
 d\rho_b=T_\theta\,d\rho_0/Z_b(\theta).
 \end{gathered}
 \label{eq:f15v55:retained}
\end{equation}
The correlated retained baseline \(\rho_0\)
has not been replaced by a product law.

\subsection{Transfer the excess above one}

\begin{proposition}[Normalized exponential-excess transfer]
\label{f15v55:excess}
Let \(\mu_\theta=L_\theta\mu_0\) be any probability
change of law with
\[
 R_\theta=\log\mathbb E_{\mu_0}L_\theta^2<\infty .
\]
Let \(F_\theta\ge0\) be measurable with respect
to the common exterior; it may depend on \(\theta\).
If \(v\ge4\) and
\(\mathbb E_{\nu_0}e^{vF_\theta}\le e^{Nm}\),
where \(m\ge0\), then
\begin{equation}
 H_\theta:=\log\mathbb E_{\nu_\theta}e^{2F_\theta}
 \le
 \log\left(1+e^{R_\theta/2}
                   \sqrt{e^{4Nm/v}-1}\right).
 \label{eq:f15v55:excess}
\end{equation}
In particular this upper bound is zero when
\(m=0\), regardless of the finite likelihood cost.
\end{proposition}

\begin{proof}
Lift \(F_\theta\) to the full configuration space
through \(\xi\). Cauchy--Schwarz and
\((x-1)^2\le x^2-1\) for \(x\ge1\) give
\begin{align*}
 \mathbb E_{\nu_\theta}(e^{2F_\theta}-1)
 &=\mathbb E_{\mu_0}[L_\theta(e^{2F_\theta}-1)]\\
 &\le e^{R_\theta/2}
       \{\mathbb E_{\nu_0}(e^{2F_\theta}-1)^2\}^{1/2}\\
 &\le e^{R_\theta/2}
       \{\mathbb E_{\nu_0}e^{4F_\theta}-1\}^{1/2}.
\end{align*}
Concavity of \(x\mapsto x^{4/v}\), including
the equality case \(v=4\), bounds the last
expectation by \(e^{4Nm/v}\).
Add one and take logarithms.
The exterior likelihood is
\(\mathbb E_{\mu_0}[L_\theta\mid\xi]\);
its second moment is no larger than that of
\(L_\theta\), by conditional Jensen.
Thus the proof uses the actual normalized
exterior marginal, not a formal or reference
source tilt. Finally \(\mathbb E L_\theta=1\)
implies \(R_\theta\ge0\).
\end{proof}

Applying Cauchy--Schwarz directly to
\(e^{2F_\theta}\), without subtracting one,
would retain a term \(R_\theta/2\) even when
\(F_\theta=0\). That loss is avoided here.
The likelihood cost itself has not been
removed: it multiplies the small excess.

\subsection{Spend the raw strip exponent on the cap}

Fix an accuracy exponent \(s_0>1\), keep \(p\)
fixed, and follow the selected coupling sequence.
Put \(t=\log\beta\) and use V32's retuning
\begin{equation}
 \begin{gathered}
 A=2^{51}C_0(s_0+2),\qquad
 r=\max\{368,\lceil A p^{\nu+8}t\rceil\},\qquad
 L=2p(r+5),\\
 a_p=2^{-43}C_0^{-1}p^{-(\nu+8)},\qquad
 \delta=\frac{a_p}{64rp^4},\qquad
 N_L=4L(L+1)^3,\\
 q_{\rm str}=\min\{1,N_L
                  e^{-\beta\delta^2/(12\pi^2L)+15L\log2}\}.
 \end{gathered}
 \label{eq:f15v55:scale}
\end{equation}
For all sufficiently large selected couplings,
take the new source-independent cap
\begin{equation}
 \Lambda=\log(1/q_{\rm str})
 =\frac{\beta\delta^2}{12\pi^2L}
                       -15L\log2-\log N_L\ge1 .
 \label{eq:f15v55:cap}
\end{equation}
This is not V35's earlier special choice of cap.
Its finite cap formulas remain applicable.
The strip estimate and its \(15L\log2\) cost
are both kept at finite regulator.

\begin{proposition}[Baseline exceptional moment at the larger cap]
\label{f15v55:baseline}
Set
\begin{equation}
 \begin{gathered}
 q_Y=\min\{1,2^{5n_p}e^{-\Lambda/2}\},\qquad
 q=\min\{1,\sqrt{q_{\rm str}}+\sqrt{q_Y}\},\\
 \gamma=q^{1/4}(1-4\tau)^{-5n_p/4}
                         (1-2\tau)^{-5n_p/2},
 \qquad \eta=\gamma^{1/4},\qquad v=(128\tau)^{-1}\ge4,\\
 K_\Lambda=\sqrt2\,e^{\Lambda/128}
                    (1-1/64)^{-5n_p/2}\gamma^{3/8},
 \qquad m_\Lambda=\log(1+K_\Lambda).
 \end{gathered}
 \label{eq:f15v55:baseline}
\end{equation}
Eventually \(0<\eta<1/2\), and uniformly over
all fitting separated families and all
\(\theta\in[-\tau,\tau]^N\),
\begin{equation}
 \begin{gathered}
 \log\mathbb E_{\nu_0}e^{vF_\theta}\le N m_\Lambda,\\
 \gamma\le e^{-\Lambda/32},\qquad
 \eta\le e^{-\Lambda/128},\qquad
 m_\Lambda\le e^{-\Lambda/512},\qquad
 e_0\le2\Delta+2e^{-\Lambda/128}.
 \end{gathered}
 \label{eq:f15v55:baseline-output}
\end{equation}
\end{proposition}

\begin{proof}
At fixed \(p,s_0\), \(r/t\to A p^{\nu+8}\).
The leading term in \eqref{eq:f15v55:cap} is
\[
 \frac{\beta a_p^2}{98304\pi^2p^9r^2(r+5)}.
\]
The two subtracted terms are \(O(t+\log t)\).
Consequently, with a strictly positive constant,
\begin{equation}
 \Lambda\sim c_{p,s_0}\frac{\beta}{t^3},\qquad
 c_{p,s_0}=
 \frac{a_p^2}{98304\pi^2p^9(Ap^{\nu+8})^3}>0 .
 \label{eq:f15v55:cap-rate}
\end{equation}
In particular \(\Lambda\to\infty\) and
\(n_p/\Lambda\to0\). This justifies the cap
threshold and eventually removes the upper
truncation in \(q_Y\).

The joint cap-failure estimate
\eqref{eq:f15v35:q} applies at this cap,
with \(q_{\rm str}=e^{-\Lambda}\).
It yields
\[
 \log q\le\log2-\Lambda/4+(5n_p/2)\log2
                                  \le-\Lambda/6
\]
eventually. Hence
\(\log\gamma\le-\Lambda/24+O(n_p)
 \le-\Lambda/32\), uniformly on
\(0<\tau\le1/512\). It follows that
\[
 \log K_\Lambda
 \le-\Lambda/256
       +(5n_p/2)\log(64/63)+\tfrac12\log2
 \le-\Lambda/512
\]
eventually. Use \(m_\Lambda\le K_\Lambda\)
and \(-\log(1-\eta)\le2\eta\).

For completeness the finite moment argument
does not require any relation between the
old cap scale and the new one. For every
subfamily \(I\), V35's actual conditional-energy
and mixed-confidence estimates give
\[
 \mathbb E_{\nu_0}
 \left[e^{v\sum_{i\in I}(g_i+\tau\Lambda)}
                         1_{\cap_{i\in I}B_i}\right]
 \le K_\Lambda^{|I|}.
\]
Here \(2v\tau=1/64\), inside the selected
source radius, and the confidence threshold
is exactly \(\eta=\gamma^{1/4}\).
Expand the nonnegative product
\[
 e^{vF_\theta}
 =\prod_i[1+1_{B_i}(e^{v(g_i+\tau\Lambda)}-1)].
\]
Bounding each nonempty term by the preceding
subset estimate gives
\(\mathbb E_{\nu_0}e^{vF_\theta}
 \le(1+K_\Lambda)^N\).
No moment of the reference law is postulated.
\end{proof}

All exponential simplifications in this proposition
are eventual assertions at fixed \(p,s_0\).
For an individual finite regulator one can
instead use the exact quantities
\eqref{eq:f15v55:baseline} in
\eqref{eq:f15v55:excess}.
In particular a large numerical cap without
the \(n_p/\Lambda\) gate is not a substitute
for the stated estimates.

\subsection{Pay the full source likelihood before absorbing it}

\begin{proposition}[A squared-amplitude likelihood budget]
\label{f15v55:likelihood}
Under the inherited arbitrary-support
centered-source CGF, with selected constant
\(k_M\le5\) and source radius at least \(1/2\),
\begin{equation}
 \begin{aligned}
 R_\theta
 &=\log Z_a(2\theta)-2\log Z_a(\theta)\\
 &\le5n_p\sum_i J(2|\theta_i|)
 \le\frac{10n_p}{1-2\tau}\sum_i\theta_i^2
 \le16n_p\sum_i\theta_i^2 ,
 \qquad J(a)=-\log(1-a)-a .
 \end{aligned}
 \label{eq:f15v55:likelihood}
\end{equation}
\end{proposition}

\begin{proof}
The first identity keeps both physical
partition functions. Jensen and the exact
baseline centering give \(\log Z_a(\theta)\ge0\).
Apply the CGF to the disjoint source-cell
sets, with coefficient \(2\theta_i\) on the
\(n_p\) cells in the \(i\)-th set.
For mixed signs use \(J(-a)\le J(a)\) at
\(0\le a<1\). Finally the power series gives
\[
 J(a)=\sum_{j\ge2}\frac{a^j}{j}
                 \le\frac{a^2}{2(1-a)}.
\]
All coefficients lie strictly inside the
selected source radius, and
\(10/(1-2\tau)\le16\) on the declared window.
\end{proof}

This is the full normalized source-law cost,
of the same type already used in V26.
It is not a new uniform inhomogeneous Wilson
CGF, nor a replacement of that cost by a
local block quantity.

\begin{theorem}[Actual-background normalization under a total tilt budget]
\label{f15v55:normalization}
Assume the eventual estimates
\eqref{eq:f15v55:baseline-output} and the
explicit finite source budget
\begin{equation}
 R_\theta+\log(2N)\le\Lambda/1024 .
 \label{eq:f15v55:budget}
\end{equation}
A sufficient condition replaces \(R_\theta\)
by \(16n_p\sum_i\theta_i^2\).
Then
\begin{equation}
 H_\theta\le e^{-\Lambda/2048},\qquad
 \mathcal E_{N,\theta}:=Ne_0+\tfrac12H_\theta
 \le2N\Delta+2Ne^{-\Lambda/128}
                         +\tfrac12e^{-\Lambda/2048},
 \label{eq:f15v55:total-error}
\end{equation}
and the exact laws \eqref{eq:f15v55:retained} obey
\begin{equation}
 \begin{gathered}
 |\log(Z_b/Z_a)|\le\mathcal E_{N,\theta},\qquad
 D(\rho_a\Vert\rho_b)\le2\mathcal E_{N,\theta},\\
 D(\rho_b\Vert\rho_a)\le4\mathcal E_{N,\theta},\qquad
 d_{\rm TV}(\rho_a,\rho_b)
                  \le\min\{1,\sqrt{\mathcal E_{N,\theta}}\}.
 \end{gathered}
 \label{eq:f15v55:normalization}
\end{equation}
These are total, not volume-divided, errors.
\end{theorem}

\begin{proof}
Since \(R_\theta\ge0\), the budget implies
\[
 Nm_\Lambda
 \le N e^{-\Lambda/512}
 \le\tfrac12 e^{-\Lambda/1024}<1.
\]
For \(0\le x\le1\), \(e^x-1\le2x\).
Because \(4/v\le1\), the excess in
\eqref{eq:f15v55:excess} is at most
\[
 e^{R_\theta/2}\sqrt{2Nm_\Lambda}
 \le\sqrt{2N}\,e^{R_\theta/2-\Lambda/1024}
 \le e^{-\Lambda/2048}.
\]
Use \(\log(1+x)\le x\).
This proves the first assertion of
\eqref{eq:f15v55:total-error}; the second
uses the bound on \(e_0\).
Apply Theorem~\ref{f15v37:criterion}
with \(N\kappa=H_\theta\).
Its normalized common-exterior comparison
projects to precisely \(\rho_a,\rho_b\)
above. Thus none of the physical centering,
conditional normalizers, or retained
baseline correlations is discarded.
\end{proof}

\subsection{A growing fixed-strength cube with vanishing total errors}

\begin{theorem}[Independent mixed signs on a growing family]
\label{f15v55:cube}
Fix \(p,s_0,\tau\) as above. Along the inherited
selected coupling sequence, consider any
separated fitting family whose cardinality satisfies
\begin{equation}
 1\le N\le
 \left\lfloor\frac{\Lambda}
                    {32768\,n_p\,\tau^2}\right\rfloor .
 \label{eq:f15v55:cube}
\end{equation}
Eventually the budget \eqref{eq:f15v55:budget}
holds uniformly on the entire cube
\([-\tau,\tau]^N\). The same source-independent
capped reference works for every point of
that cube. Uniformly over these families and
source vectors,
\begin{equation}
 \mathcal E_{N,\theta}
 =O_{p,s_0,\tau}\left(\frac{\beta^{1-s_0}}{(\log\beta)^3}\right)
       +O_{p,s_0,\tau}(e^{-c\beta/(\log\beta)^3})
 \longrightarrow0
 \label{eq:f15v55:cube-output}
\end{equation}
for some \(c>0\). In particular the total
pressure error, both retained relative
entropies and retained total variation
in \eqref{eq:f15v55:normalization} tend to zero.
\end{theorem}

\begin{proof}
For every point of the cube,
\[
 R_\theta\le16n_p N\tau^2\le\Lambda/2048.
\]
Since \(N=O_{p,\tau}(\Lambda)\) and
\(\Lambda\to\infty\), also
\(\log(2N)\le\Lambda/2048\) eventually,
uniformly in the permitted cardinalities.
This proves the budget. V32's collar
estimate \eqref{eq:f15v32:Delta} gives
\(\Delta\le\beta^{-s_0}\) eventually.
Equation~\eqref{eq:f15v55:cap-rate} then yields
\[
 N\Delta
 =O_{p,s_0,\tau}(\beta^{1-s_0}/(\log\beta)^3).
\]
The remaining terms in
\eqref{eq:f15v55:total-error} are bounded
by a constant times \(e^{-c\beta/(\log\beta)^3}\)
after reducing \(c>0\) to absorb the factor \(N\).
All thresholds are independent of the signs
and individual amplitudes. The cap and its
fallback kernel were chosen before \(\theta\).
The entropy and TV conclusions now follow
from \eqref{eq:f15v55:normalization}.
\end{proof}

The allowed maximal cardinality is of order
\(\beta/(\log\beta)^3\) at fixed \(p,s_0,\tau\),
subject also to geometric fitting. On the
inherited tori \(M\ge e^\beta\), the separated
mesh has more than enough sites eventually.
This improves the earlier logarithmic
fixed-strength cardinality restriction, but
is far below a volume-extensive family.
The more general budget
\(16n_p\sum_i\theta_i^2+\log(2N)\le\Lambda/1024\)
also allows many weak amplitudes; the theorem
does not assert that every fitting vector
meets it.

\subsection{What is still missing upstream}

A small baseline excess cannot by itself
control an arbitrarily tilted law. For example,
on a two-point space let a baseline event
have probability \(2^{-128}\), its actual
tilted probability be \(2^{-4}\), and
\(e^{2F}=2^{16}\) on that event and one elsewhere.
Then the baseline excess
\(\mathbb E_0e^{4F}-1<2^{-95}\), whereas
\(\mathbb E_\theta e^{2F}-1>4000\).
The normalized likelihood has a correspondingly
large second moment. This elementary example
illustrates, but does not replace, V37's
previous countermodel: the likelihood factor
in \eqref{eq:f15v55:excess} must be paid.

Accordingly V55 has not established the
arbitrary-cardinality linear-energy endpoint.
Obtaining that endpoint requires genuinely
local or volume-uniform actual-background
control, not another use of the same global
second moment with its cost suppressed.
The present counting budget also says nothing
about decay with the separation between two
retained observables. The retained baseline
can remain correlated. No convergence of
complex derivatives follows merely from
the real-pressure estimate proved here.
All asymptotic statements fix \(p\); an
iteration over growing physical scales
requires additional uniformity.

\subsection{Diagnostics, source preservation and artifact policy}

The new exact finite diagnostic is
\begin{center}\small
\path{scripts/verify_ym_f15_v55_exponential_excess_transport.py}.
\end{center}
It reruns V54's preserved diagnostic and checks
108 actual exceptional-moment excess transfers,
27 full/exterior normalized likelihood
second-moment identities and 27 zero-excess
endpoint identities. Four raised-power
rational tests check the retuned cap budget;
28 tests check the finite source-budget
absorption, two check the source-window
constants, and one retains the unpaid-cost
counterexample above.
These finite models test bookkeeping and
normalization; they do not certify the
Wilson/chart inputs or replace the analytic
proofs of this section.

The new script has 5,037 bytes and SHA--256
\begin{center}\scriptsize\ttfamily
3770DE59C89552B34A7022930091084842DC33FB728D1F23FDF97095682551DC.
\end{center}
The V54 predecessor has 1,258,597 bytes and SHA--256
\begin{center}\scriptsize\ttfamily
E0136FDDACBC868096DBCA7E7713144177645C07244C25C6D49798EBBCB018DB.
\end{center}
Removing this final section and reversing
the title increment recovers V54 exactly.
V55 replaces V54 as the sole active main
source; the 14 protected inputs and 97
earlier diagnostics remain unchanged.
Only \path{.tex} sources are kept in
\path{latex_notes}. All compilation products
are directed outside that folder, to
\begin{center}\small
\path{artifacts/ym_f15_v55_texcheck/}.
\end{center}

\begin{firewall}
Under the inherited selected-coupling and chart
hypotheses, V55 proves actual-law normalization
with vanishing total errors for a growing
independently selectable mixed-sign source
cube at fixed nonzero amplitude.
It retains an explicit total likelihood
budget and a fixed block depth.
It does not prove arbitrary extensive
source closure, spatial clustering,
an interacting continuum construction,
or a positive physical Yang--Mills
Hamiltonian mass gap.
\end{firewall}

\section{V56: a real-pressure bridge to the retained covariance operator}
\label{f15v56:section}

\subsection{The next interface, and what is not being repeated}

V55 proves a uniform real-source pressure
comparison on an independently selectable cube.
A small real function need not have small
derivatives. We therefore do not differentiate
its error estimate formally. Instead, the
compact source range supplies an explicit
fourth-derivative bound, and a centered
finite difference controls the entire Hessian
in operator norm. The resulting covariance
matrix is then expressed using local predictors
under the \emph{actual retained law}.

This moves from partition-function values
to connected second responses. V36 already
transferred fixed mixed energy moments for
a fixed family by finite differences; that
result is not claimed anew. The local-box
covariance transfer in f14 V38 concerns a
different chart law and its spatially weighted
extension. The new assertion here is uniform
over V55's growing source families and over
an interior cube of nonzero independent
backgrounds, with one bound for all unit
source directions simultaneously.
It is an unweighted operator estimate,
not spatial decay.

The V55 predecessor is unchanged before
this append. Its companion checkpoint
remains the latest one; the next is V60.
The physical objective is still the
nontrivial continuum quantum theory and
positive excitation threshold specified in
\href{https://www.claymath.org/wp-content/uploads/2022/06/yangmills.pdf}
{Jaffe and Witten's problem statement}.
The response estimate below is only one
interface toward that objective.

\subsection{Two genuine exponential families with one fixed reference}

Retain V55's original block-energy sources,
selected coupling hypotheses, separated
geometry and capped reference. In this
section choose the collar accuracy exponent
\begin{equation}
 s_0>7,\qquad p\ge2\ \hbox{fixed and dyadic},
 \qquad 0<\tau\le1/512\ \hbox{fixed}.
 \label{eq:f15v56:parameters}
\end{equation}
All definitions of \(r,\Lambda,\lambda_i^\Lambda\)
are those of \eqref{eq:f15v55:scale}--%
\eqref{eq:f15v55:cap} at this chosen \(s_0\).
In particular the cap and the kernels do
not change when a source parameter changes.
Use any fitting family with
\begin{equation}
 1\le N\le
 \left\lfloor\frac{\Lambda}
                       {32768\,n_p\,\tau^2}\right\rfloor .
 \label{eq:f15v56:cardinality}
\end{equation}

On the retained variables and reference cores
define the probability
\begin{equation}
 dQ_0(C,V_1,\ldots,V_N)
   =d\rho_0(C)\prod_{i=1}^N
                   \lambda_i^\Lambda(N_i(C),dV_i).
 \label{eq:f15v56:joint-reference}
\end{equation}
The source \(S_i=s_i(N_i,V_i)\) is exactly
the original core source, including its
baseline centering constant. Thus
\begin{equation}
 \begin{gathered}
 P_a(\theta)=\log Z_a(\theta),\qquad
 P_b(\theta)=\log Z_b(\theta)
           =\log\mathbb E_{Q_0}e^{\sum_i\theta_iS_i},\\
 dQ_\theta=e^{\sum_i\theta_iS_i-P_b(\theta)}\,dQ_0,
 \qquad C_*Q_\theta=\rho_b .
 \end{gathered}
 \label{eq:f15v56:pressures}
\end{equation}
The actual source law is \(\mu_\theta\),
with retained marginal \(\rho_a\).
The baseline \(\rho_0\) in \(Q_0\) is still
correlated. No product assumption on it is
introduced by the product of conditional
core kernels.

The compact energy range in
\eqref{eq:f15v26:sources} gives, under both
families,
\begin{equation}
 |S_i|\le B,\qquad B=48\beta p^4 .
 \label{eq:f15v56:range}
\end{equation}
This follows on the reference support as
well because it is a subset of the same
compact core configurations and uses the
same centering. The zero-chart-mass
fallback does not enlarge that range.

For all sufficiently large selected couplings,
V55 supplies the common bound
\begin{equation}
 \begin{gathered}
 \sup_{\theta\in[-\tau,\tau]^N}
                   |P_a(\theta)-P_b(\theta)|\le\epsilon_N,\\
 \epsilon_N=2N\beta^{-s_0}+2Ne^{-\Lambda/128}
                                    +\tfrac12e^{-\Lambda/2048},\\
 \sup_{\theta\in[-\tau,\tau]^N}
                   d_{\rm TV}(\rho_a,\rho_b)\le\sqrt{\epsilon_N}.
 \end{gathered}
 \label{eq:f15v56:input}
\end{equation}
The last bound can be replaced by its
minimum with one. There is one \(Q_0\)
for the entire cube, which is essential
for the following differentiations.

\subsection{A quantitative real-source differentiation argument}

\begin{proposition}[Bounded exponential families and a Hessian comparison]
\label{f15v56:difference}
Let \(P_a,P_b\) be log partition functions
of two probability laws, possibly on
different spaces, with \(N\) real source
statistics bounded in absolute value by
\(B>0\) in each law. Suppose
\[
 \sup_{[-\tau,\tau]^N}|P_a-P_b|\le\epsilon.
\]
If \(\epsilon>0\) and
\begin{equation}
 h_*=\frac{\epsilon^{1/4}}{B\sqrt N}\le\tau/2,
 \label{eq:f15v56:step}
\end{equation}
then throughout \([-\tau/2,\tau/2]^N\),
\begin{equation}
 \|\nabla^2P_a(\theta)-\nabla^2P_b(\theta)\|_{2\to2}
                         \le16B^2N\sqrt\epsilon .
 \label{eq:f15v56:hessian}
\end{equation}
If \(\epsilon=0\), the Hessians agree
on the interior without this step condition.
\end{proposition}

\begin{proof}
Fix a real unit vector \(a\in\mathbb R^N\).
For either law let \(X=\sum_i a_i S_i\).
Cauchy--Schwarz gives \(|X|\le M:=B\sqrt N\).
At any real parameter the normalized tilted
law exists and differentiation under its
bounded integrals is justified. The fourth
directional derivative is the fourth cumulant,
\[
 \frac{d^4}{du^4}P(\theta+ua)
 =\mathbb E(X-\mathbb EX)^4-3\Var(X)^2 .
\]
Since \(|X-\mathbb EX|\le2M\), its absolute
value is at most \(4(2M)^4=64M^4\).
This deliberately nonoptimal constant is
uniform in the real background.
For \(D(u)=P_a(\theta+ua)-P_b(\theta+ua)\)
we therefore have \(|D^{(4)}(u)|\le128M^4\).

For \(0<h\le\tau/2\), the whole line segment
\(\theta+[-h,h]a\) lies in the outer cube,
because \(\|a\|_\infty\le1\). Taylor's
formula with integral remainder, applied
at \(h\) and \(-h\), yields
\[
 \left|D''(0)-
       \frac{D(h)-2D(0)+D(-h)}{h^2}\right|
                         \le\frac{128M^4h^2}{12}.
\]
Indeed the two fourth-order remainders
are each bounded by
\((\sup|D^{(4)}|)h^4/24\); the odd
terms cancel. Thus
\[
 |D''(0)|\le\frac{4\epsilon}{h^2}
                             +\frac{128M^4h^2}{12}.
\]
Choose \(h=h_*\). The right side is
\((4+128/12)M^2\sqrt\epsilon
 \le16B^2N\sqrt\epsilon\).
Taking the supremum of absolute quadratic
forms over real unit vectors gives the
operator norm of the real symmetric
Hessian difference. If \(\epsilon=0\),
the pressure difference vanishes on an
open cube, so all its interior derivatives
vanish.
\end{proof}

No holomorphic logarithm, complex zero-free
region, or interchange of a limit with
differentiation is assumed. The bound
pays for the regulator-dependent source
range and the growing dimension. Applying
it to \eqref{eq:f15v56:pressures} gives
\begin{equation}
 \|\operatorname{Cov}_{\mu_\theta}(S)
       -\operatorname{Cov}_{Q_\theta}(S)\|_{2\to2}
                   \le16B^2N\sqrt{\epsilon_N}
 \label{eq:f15v56:covariance-transfer}
\end{equation}
on the inner cube, once the explicit
step condition holds. This is a matrix
statement, not an entrywise estimate
followed by an unrecorded sum over \(N\).

\subsection{Local predictors and the diagonal conditional variance}

Define the normalized local source kernel
and its first two centered moments by
\begin{equation}
 \begin{gathered}
 d\lambda_{i,\theta_i}^\Lambda
 =\frac{e^{\theta_iS_i}}
             {\lambda_i^\Lambda(e^{\theta_iS_i})}
                                      \,d\lambda_i^\Lambda,\\
 M_i^\theta(N_i)=\lambda_{i,\theta_i}^\Lambda(S_i),
 \qquad
 d_i^\theta(N_i)=
                       \Var_{\lambda_{i,\theta_i}^\Lambda}(S_i),\\
 |M_i^\theta|\le B,\qquad 0\le d_i^\theta\le B^2 .
 \end{gathered}
 \label{eq:f15v56:predictors}
\end{equation}
The notation \(M_i^\theta\) is distinct from
the cap minimum \(m_i\). Each predictor
depends on its own \(\theta_i\) and its
own retained near variables \(N_i\), not
on remote source coefficients. Regulator
dependence of the underlying kernel is
not excluded. No differentiability of
its hard cutoff in \(N_i\) is needed.

\begin{proposition}[Exact reference covariance decomposition]
\label{f15v56:decomposition}
Let \(M^\theta=(M_i^\theta)_{i=1}^N\).
Then
\begin{equation}
 \operatorname{Cov}_{Q_\theta}(S)
 =\operatorname{Cov}_{\rho_b}(M^\theta)
       +\operatorname{diag}
                 \bigl(\mathbb E_{\rho_b}d_i^\theta\bigr).
 \label{eq:f15v56:decomposition}
\end{equation}
In particular the off-diagonal entries
are covariances of local retained
predictors, not independent-core errors.
\end{proposition}

\begin{proof}
The exact normalized disintegration is
\[
 dQ_\theta
       =d\rho_b(C)\prod_i
                         d\lambda_{i,\theta_i}^\Lambda(N_i(C)).
\]
It follows by dividing each local factor
by its own normalizer and putting their
product into \(d\rho_b\), with the full
normalizer \(Z_b\) still present.
Given \(C\), the cores are independent,
their means are \(M_i^\theta\), and
their conditional covariance matrix is
\(\operatorname{diag}(d_i^\theta)\).
Expanding the centered products and
conditioning on \(C\) proves the formula.
\end{proof}

For any retained probability \(\rho\), put
\[
 \mathcal T_\theta(\rho)
 =\operatorname{Cov}_{\rho}(M^\theta)
       +\operatorname{diag}(\mathbb E_\rho d_i^\theta).
\]
This matrix is positive semidefinite,
but positivity alone is not a decay
statement.

\begin{proposition}[Transport to the actual retained background]
\label{f15v56:retained-transport}
For retained probabilities \(\rho,\widetilde\rho\)
and the same predictors and variances,
\begin{equation}
 \|\mathcal T_\theta(\rho)
       -\mathcal T_\theta(\widetilde\rho)\|_{2\to2}
                \le6B^2N\,d_{\rm TV}(\rho,\widetilde\rho).
 \label{eq:f15v56:retained-transport}
\end{equation}
Consequently, for the actual retained
law in \eqref{eq:f15v56:input},
\begin{equation}
 \|\operatorname{Cov}_{\mu_\theta}(S)
                         -\mathcal T_\theta(\rho_a)\|_{2\to2}
                         \le22B^2N\sqrt{\epsilon_N}.
 \label{eq:f15v56:actual-response}
\end{equation}
\end{proposition}

\begin{proof}
Fix a unit vector \(a\), set
\(X(C)=a\cdot M^\theta(C)\) and
\(V(C)=\sum_i a_i^2d_i^\theta(C)\).
Then \(|X|\le B\sqrt N\) and
\(0\le V\le B^2\). If
\(d=d_{\rm TV}(\rho,\widetilde\rho)\),
the oscillation bound for expectations gives
\[
 |\mathbb E_\rho X^2-\mathbb E_{\widetilde\rho}X^2|
       \le B^2N d,\qquad
 |\mathbb E_\rho X-\mathbb E_{\widetilde\rho}X|
       \le2B\sqrt N d .
 \]
The difference of squared means is at
most \(4B^2N d\), so the variance
difference is at most \(5B^2N d\).
The difference of the two expectations
of \(V\) is at most \(B^2d\).
Since \(N\ge1\), their sum is bounded
by \(6B^2N d\). Supremizing the
quadratic form proves the first result.
For the second use
\eqref{eq:f15v56:covariance-transfer},
\eqref{eq:f15v56:decomposition},
and \(d_{\rm TV}(\rho_a,\rho_b)
\le\sqrt{\epsilon_N}\).
\end{proof}

The diagonal variance term is indispensable.
Replacing the whole covariance by
\(\operatorname{Cov}_{\rho_a}(M^\theta)\)
alone would lose the fluctuations still
inside each normalized local core.

\subsection{A vanishing matrix error and a retained response norm}

\begin{theorem}[Growing-family connected response bridge]
\label{f15v56:response}
With \eqref{eq:f15v56:parameters} and
\eqref{eq:f15v56:cardinality}, the step
condition holds eventually and
\begin{equation}
 \begin{split}
 \sup_{\theta\in[-\tau/2,\tau/2]^N}
 \|\operatorname{Cov}_{\mu_\theta}(S)
                  -\mathcal T_\theta(\rho_a)\|_{2\to2}
 &\le22B^2N\sqrt{\epsilon_N}\\
 &=O_{p,s_0,\tau}\left(
       \frac{\beta^{(7-s_0)/2}}{(\log\beta)^{9/2}}\right)
    +O_{p,s_0,\tau}(e^{-c\beta/(\log\beta)^3})
 \longrightarrow0
 \end{split}
 \label{eq:f15v56:response-rate}
\end{equation}
for some \(c>0\), uniformly over the
permitted separated families.
\end{theorem}

\begin{proof}
V55 gives \(\Lambda\sim c_{p,s_0}
\beta/(\log\beta)^3\) and therefore
\(N=O_{p,s_0,\tau}(\beta/(\log\beta)^3)\).
Equation~\eqref{eq:f15v56:input} implies
\(\epsilon_N\to0\) uniformly.
Meanwhile \(B\sqrt N\ge48\beta p^4\),
so \(\epsilon_N^{1/4}/(B\sqrt N)
\le\tau/2\) eventually.
Using \(\sqrt{x+y}\le\sqrt x+\sqrt y\),
the algebraic contribution to
\(B^2N\sqrt{\epsilon_N}\) is bounded by
\[
 C_p\beta^2N^{3/2}\beta^{-s_0/2}
 =O_{p,s_0,\tau}\left(
             \beta^{(7-s_0)/2}/(\log\beta)^{9/2}\right).
\]
The remaining factors are polynomial
in \(\beta\) and multiply exponentials
in \(-\Lambda\); they are absorbed by
a smaller positive \(c\).
Now apply \eqref{eq:f15v56:actual-response}.
\end{proof}

\begin{proposition}[A baseline local-predictor response bound]
\label{f15v56:baseline-response}
At zero source, write
\[
 \delta_N=16B^2N\sqrt{\epsilon_N}\longrightarrow0 .
\]
Then, on every family above,
\begin{equation}
 0\le\operatorname{Cov}_{\rho_0}(M^0)
 \le\mathcal T_0(\rho_0)
 \le(5n_p+\delta_N)I_N .
 \label{eq:f15v56:baseline-response}
\end{equation}
The inequalities are in the quadratic-form
order. In particular
\[
 \Var_{\rho_0}\left(\sum_i a_iM_i^0\right)
       \le(5n_p+\delta_N)\sum_i a_i^2
       \quad\hbox{for every real }a .
\]
\end{proposition}

\begin{proof}
At \(\theta=0\), \(\rho_a=\rho_b=\rho_0\)
exactly, so the extra retained-law
transport cost is zero. The selected
arbitrary-support centered CGF
\eqref{f15v1:source-input}, applied with
coefficient \(u a_i\) on each of the
\(n_p\) cells in the \(i\)-th disjoint
source block, gives
\[
 \log\mathbb E_{\mu_0}e^{u\sum_i a_iS_i}
      \le 5n_p\sum_iJ(u a_i)
      =\tfrac52 n_pu^2\sum_i a_i^2+O_a(u^3)
\]
as \(u\to0\). Differentiation here is
at one finite regulator, with bounded
statistics, so
\(\operatorname{Cov}_{\mu_0}(S)\le5n_pI_N\).
Combine \eqref{eq:f15v56:covariance-transfer}
at zero with the exact decomposition.
Its diagonal conditional-variance term
is positive semidefinite, giving the
remaining order inequalities.
\end{proof}

This is a bound on the newly specified
local retained predictors; it is not
merely the pre-existing bound on the
original action-density covariance.
It is still restricted to separated
families and fixed blocking depth.
At nonzero backgrounds the theorem
compares two covariance matrices but
does not assert the bound \(5n_p I_N\)
for either of them.

\subsection{Why this response norm is not physical clustering}

The missing spatial qualification cannot
be supplied by positivity or translation
invariance alone. Here is an exact
illustration, not a Wilson counterexample.
On a ring of \(2m\) retained spins
\(C_i\in\{-1,1\}\), choose \(0<c<1\) and
\[
 \rho_m(C)=2^{-2m}
             \prod_{i=1}^m(1+cC_iC_{i+m}).
\]
This is a strictly positive normalized
probability and is invariant under cyclic
translations. It is a product over
antipodal pairs, each with zero means
and covariance \(c\).
Its covariance matrix is \(I+cP_m\),
where \(P_m\) swaps antipodal sites;
its norm is \(1+c\), independently of \(m\).

Give each site a local core bit \(X_i\)
with conditionally independent kernel
\[
 \lambda_i(X_i=x\mid C_i)
                  =\tfrac12(1+\kappa xC_i),
                  \qquad0<\kappa<1.
\]
Take the actual joint law to equal this
reference joint law. All pressure and
retained-law comparison errors are then
zero for every source. The local mean
is \(\kappa C_i\), the conditional
variance is \(1-\kappa^2\), and the
full source covariance is exactly
\[
        \operatorname{Cov}(X)=I+\kappa^2cP_m .
\]
Its norm is \(1+\kappa^2c<2\), but the
covariance of two sites a distance \(m\)
apart is always \(\kappa^2c\).
No constants \(K<\infty\), \(a>0\)
independent of \(m\) can bound that
covariance by \(Ke^{-am}\).

For the YM programme, therefore,
\eqref{eq:f15v56:baseline-response}
controls the size of an unweighted
retained response but not its spatial
distribution. A next substantive input
must control distant retained predictors,
for example through a spatially weighted
connected-response estimate under the
actual retained law and with the
appropriate physical-scale calibration.
Even such an energy-observable estimate
would not alone construct all continuum
Schwinger functions or the physical
Hamiltonian. No mass gap is inferred
from the present matrix bound.

\subsection{Exact checks and append-only preservation}

The new diagnostic is
\begin{center}\small
\path{scripts/verify_ym_f15_v56_covariance_response_bridge.py}.
\end{center}
Six strictly positive finite models
preserve their retained baseline exactly.
Across 108 source/model patterns, the
script checks the complete tilted-kernel
factorization and normalizer. It checks
837 covariance decomposition entries,
513 actual-retained covariance transport
inequalities, 1,026 exact fourth-cumulant
identities and compact bounds, and six
zero-source retained identities.
There are 27 exact quartic-polynomial
central-difference remainder checks,
three response-rate exponent checks,
and four translation-invariant positive
antipodal models with no uniform
exponential clustering.
All arithmetic is rational. These tests
check identities and constants; the
operator supremum and limiting theorems
use the analytic proofs above.

The script has 8,184 bytes and SHA--256
\begin{center}\scriptsize\ttfamily
948C3BBE835EDD94FF23A9CBABC8124D3218F8760F3CA87FDDD1885C5D7D81E7.
\end{center}
The V55 predecessor has 1,279,138 bytes
and SHA--256
\begin{center}\scriptsize\ttfamily
41CE98B29679C0E7D7234F390D856E923A4CDEF831608FBEC39A9DFE146BA03D.
\end{center}
Removing this appended section and
reversing the title increment recovers
V55 byte for byte. V56 becomes the
sole active main source; protected
companions, archives and the 98 earlier
diagnostics are not altered.
Only \path{.tex} files belong in
\path{latex_notes}; compilation products
go to
\begin{center}\small
\path{artifacts/ym_f15_v56_texcheck/}.
\end{center}

\begin{firewall}
V56 pays for real-source differentiation
and proves a vanishing whole-matrix
covariance error on a growing independent
source family. It identifies the retained
local-predictor covariance and the
indispensable diagonal core variance,
and obtains a baseline unweighted
predictor response bound.
The selected-coupling hypotheses,
cardinality restriction, inner source
cube and fixed depth remain binding.
Spatial decay, scale-uniform RG closure,
the interacting continuum theory and
the physical Yang--Mills mass gap
remain unproved.
\end{firewall}

\section{V57: positive Wilson transfer and algebraic action-density clustering}
\label{f15v57:section}

\subsection{Go upstream to the physical time-step insertion}

V56's unweighted response norm did not
control spatial decay. The conditional
core estimates of V19--V23 also leave
the retained field correlated. This
section uses an additional structure
of the actual homogeneous Wilson law:
its positive physical transfer operator
and the reflection symmetry of the
specified cell-energy insertion.
The result is a distance-dependent bound
for the original, full compact
action-density covariance at the
selected couplings, before a chart
restriction or a reference-core replacement.

Positive Wilson transfer is standard,
not a new construction principle;
see
\href{https://bib-pubdb1.desy.de/record/396349/files/7611148.pdf?version=1}
{L\"uscher's positive-transfer construction}.
For the pure \(\SU(2)\) law used here
the needed finite-volume facts are
derived below, including the exact
insertion associated with our cell
convention. The new programme
consequence combines those facts with
the inherited covariance/source bound.
It is not a claim that a norm bound
alone implies clustering, and it does
not repeat f14 V99's Gaussian-limit
susceptibility sum rule.

The decay proved here is algebraic.
No uniform exponential rate, continuum
construction, or Hamiltonian mass gap
is claimed. The companion checkpoint
remains V55; the next is V60.
The shell V16 conclusion was additionally
consulted while choosing this direction.
Its uncomputed noncommuting one-loop
coefficient is not used as an input.

\subsection{The exact symmetric slab expression for the cell energy}

Take the homogeneous Wilson law on
the periodic four-dimensional torus
of side \(M\ge4\), with normalized
product Haar reference measure,
\[
 d\mu_\beta=Z^{-1}e^{-\beta S}\,dH,
 \qquad s_p=1-\tfrac12\Tr U_p,\qquad
 S=\sum_p s_p .
\]
Relabel one coordinate as time \(0\)
and the others as \(1,2,3\).
Retain exactly the observable \(T_x\)
from \eqref{f15v1:cell}. For a real
function \(f\) on the spatial torus,
define the one-slab observable
\begin{equation}
 F_t^f=\sum_{x\in(\mathbb Z/M\mathbb Z)^3}
                         f(x)T_{(t,x)} .
 \label{eq:f15v57:slab-observable}
\end{equation}
This \(F_t^f\) is not the exceptional
log-error budget of V55.

For \(i\in\{1,2,3\}\), let \(j,k\)
be the other two spatial directions,
and set
\[
 a_i^f(y)=\tfrac14
       \sum_{\epsilon_j,\epsilon_k\in\{0,1\}}
                f(y-\epsilon_je_j-\epsilon_ke_k).
 \]
For \(i<j\) let \(k\) be the remaining
spatial direction, and set
\[
 b_{ij}^f(y)=\tfrac14[f(y)+f(y-e_k)].
\]
Every spatial argument is periodic.

\begin{proposition}[Literal cell-to-slab identity]
\label{f15v57:slab}
The observable \eqref{eq:f15v57:slab-observable}
satisfies
\begin{equation}
 \begin{split}
 F_t^f={}&\beta\sum_{y,i}a_i^f(y)s_{0i}(t,y)\\
       &+\beta\sum_{y,i<j}b_{ij}^f(y)
                    [s_{ij}(t,y)+s_{ij}(t+1,y)] .
 \end{split}
 \label{eq:f15v57:slab}
\end{equation}
It is invariant under reflection
exchanging the two time boundaries
of its slab, with the temporal
links reversed. For \(f=1\), its
temporal coefficients are one and
its spatial coefficients on each
boundary are \(1/2\).
\end{proposition}

\begin{proof}
For a temporal face \(0i\), the
two transverse cell offsets are
the two remaining spatial directions.
Reindexing its base point in the
cell sum gives \(a_i^f\).
For a spatial face \(ij\), the
transverse offsets are time \(0\)
and the remaining spatial direction
\(k\). Splitting the time offset
into its two values gives the
two identical coefficients \(b_{ij}^f\).
This proves the equality, with all
four face incidences retained.
Reflection interchanges the spatial
boundary terms and reverses the
orientation of the temporal plaquette.
Its \(\SU(2)\) trace is real and
unchanged by inversion, proving
the stated symmetry.
\end{proof}

\subsection{A positive transfer operator with the original normalizer}

Let \(U,V\) denote spatial link
configurations on two consecutive
slices, and \(g_y\) the temporal
link at each spatial vertex.
Write \(S_{\rm sp}(U)\) for the
spatial plaquette action. Define
\begin{equation}
 \begin{gathered}
 \mathcal A(U,V,g)
   =\tfrac12S_{\rm sp}(U)+\tfrac12S_{\rm sp}(V)
        +\sum_{y,i}
          \left[1-\tfrac12\Tr
              (g_yV_i(y)g_{y+e_i}^{-1}U_i(y)^{-1})\right],\\
 \mathsf T(U,V)=\int e^{-\beta\mathcal A(U,V,g)}\,dH(g).
 \end{gathered}
 \label{eq:f15v57:transfer}
\end{equation}
The operator acts on the gauge-invariant
subspace of spatial \(L^2\) product Haar.
The capital \(\mathsf T\) is the
transfer operator, not the cell
observable \(T_x\).

\begin{proposition}[Finite-volume positivity and the exact trace law]
\label{f15v57:transfer}
For every \(\beta>0\), \(\mathsf T\)
is a real self-adjoint, positive
trace-class operator and has no
kernel on that gauge-invariant
subspace. Its eigenvalues there
can be listed as \(\lambda_a>0\).
The Wilson normalizer and insertions
are the corresponding periodic
trace integrals; in particular
\begin{equation}
                         Z=\Tr\mathsf T^M .
 \label{eq:f15v57:partition}
\end{equation}
\end{proposition}

\begin{proof}
Let \(\mathsf P\) be Haar averaging
over spatial vertex gauge transformations,
let \(\mathsf A\) be multiplication
by \(e^{-\beta S_{\rm sp}/2}\), and
let \(\mathsf K\) have kernel
\[
 \prod_{y,i}
  \exp\{-\beta[1-\tfrac12\Tr(V_i(y)U_i(y)^{-1})]\}.
\]
Direct integration of the temporal
links gives
\(\mathsf T=\mathsf A\mathsf K\mathsf P\mathsf A\).
Both \(\mathsf A\) and \(\mathsf K\)
commute with \(\mathsf P\).

To check operator positivity, not
merely positivity of its kernel,
write the one-link central function as
\[
 e^{-\beta}e^{(\beta/2)\chi_{1/2}(u)}
   =e^{-\beta}\sum_{n\ge0}
                \frac{(\beta/2)^n}{n!}\chi_{1/2}(u)^n .
\]
Each power is the character of a
tensor power of the fundamental
representation, so all irreducible
character coefficients are nonnegative.
Every spin \(j\) occurs in the
\(2j\)-th tensor power, giving a
strictly positive coefficient when
\(\beta>0\). The series is absolutely
convergent, including its dimension
sum at the identity. Peter--Weyl
orthogonality therefore shows that
the one-link convolution operator
is positive, injective and trace
class. Finite tensor products have
the same properties. Consequently
\(\mathsf A\mathsf K\mathsf P\mathsf A\)
is positive and trace class.
On the physical subspace,
\[
 \langle h,\mathsf T h\rangle
       =\langle\mathsf A h,\mathsf K\mathsf A h\rangle>0
       \quad(h\ne0),
\]
because \(\mathsf A\) is invertible
at each finite regulator.

The kernels and the Haar measure
are real. Exchanging \(U,V\) and
replacing \(g\) by \(g^{-1}\)
preserves \(\mathcal A\), so the
kernel is symmetric.
Finally the cyclic product of \(M\)
kernels integrates exactly once
over every temporal and spatial
link. The two boundary halves of
each spatial plaquette add to its
full Wilson action. Gauge Haar
measures have mass one, so no
gauge-volume factor is missing.
The trace on the full slice space
equals the trace on the physical
subspace, since the operator
annihilates its orthogonal complement.
This proves \eqref{eq:f15v57:partition}.
\end{proof}

Define the insertion kernel
\begin{equation}
 \mathsf B_f(U,V)
       =\int F_0^f(U,V,g)
                         e^{-\beta\mathcal A(U,V,g)}\,dH(g).
 \label{eq:f15v57:insertion}
\end{equation}
At fixed regulator this is a bounded
continuous kernel and hence Hilbert--Schmidt.
The slab identity proves that it
is real and symmetric. Gauge
invariance of the insertion and
the temporal integration imply
\(\mathsf B_f=\mathsf P\mathsf B_f\mathsf P\).
Its dependence on \(f\) is linear.
It need not be a positive operator.

For later clarity, two insertions
in the \emph{same} slab use instead
\[
 \mathsf D_{f,h}(U,V)=
 \int F_0^f(U,V,\gamma)F_0^h(U,V,\gamma)
       e^{-\beta\mathcal A(U,V,\gamma)}\,dH(\gamma).
\]
In general this is not a product
of two \(\mathsf B\)'s.

\subsection{Thermal centering and monotonicity away from contact}

Put
\[
 c_{f,g}(t)=\operatorname{Cov}_{\mu_\beta}(F_0^f,F_t^g),
 \qquad c_f(t)=c_{f,f}(t),
 \qquad m_f=\mathbb E_{\mu_\beta}F_0^f .
\]
For \(1\le t\le M-1\), the exact
trace formulas are
\begin{equation}
 \begin{gathered}
 m_f=\frac{\Tr(\mathsf B_f\mathsf T^{M-1})}{Z},\\
 c_{f,g}(t)=
   \frac{\Tr(\mathsf B_f\mathsf T^{t-1}
                    \mathsf B_g\mathsf T^{M-t-1})}{Z}
                            -m_fm_g .
 \end{gathered}
 \label{eq:f15v57:correlation-trace}
\end{equation}
Adjacent slabs are included: their
shared spatial slice is integrated
once in this product.
At \(t=0\) the numerator is instead
\(\Tr(\mathsf D_{f,g}\mathsf T^{M-1})\).
We will only use nonnegativity of
the resulting contact covariance
quadratic form.

\begin{proposition}[Positive, decreasing noncontact covariance forms]
\label{f15v57:monotone}
For real \(f,g\), \(c_{f,g}(t)\)
is a symmetric bilinear form in
the smearing functions. At every
nonzero cyclic time it is positive
semidefinite as such a form.
Moreover
\begin{equation}
 c_f(t)=c_f(M-t),\qquad
 c_f(t)\ge c_f(t+1)\ge0
       \quad(1\le t<\lfloor M/2\rfloor).
 \label{eq:f15v57:monotone}
\end{equation}
For even \(M\), the same monotonicity
includes the step from \(M/2-1\)
to \(M/2\).
\end{proposition}

\begin{proof}
Choose a real orthonormal eigenbasis
of the real compact operator
\(\mathsf T\), and put
\[
 B^f_{ab}=\langle a,\mathsf B_f b\rangle,\qquad
 \omega_a=\lambda_a^M/Z,\qquad
 d_a^f=B^f_{aa}/\lambda_a .
\]
The matrices \(B^f\) are real
symmetric. The weights \(\omega_a\)
form a probability and
\(m_f=\sum_a\omega_a d_a^f\).
The exact centered decomposition is
\begin{equation}
 c_{f,g}(t)
 =\operatorname{Cov}_{\omega}(d^f,d^g)
  +\frac1Z\sum_{a\ne b}B^f_{ab}B^g_{ab}
                 \lambda_a^{t-1}\lambda_b^{M-t-1}.
 \label{eq:f15v57:spectral}
\end{equation}
The diagonal covariance is retained;
subtracting the thermal mean does
not remove it in general.
The formula is a symmetric bilinear
form in \(f,g\). For \(f=g\), every
off-diagonal summand and the diagonal
variance are nonnegative.

For an unordered pair \(\{a,b\}\),
the combined scalar weight is
\begin{equation}
 \begin{split}
 &\lambda_a^{t-1}\lambda_b^{M-t-1}
        +\lambda_b^{t-1}\lambda_a^{M-t-1}\\
 &\quad =
 2(\lambda_a\lambda_b)^{(M-2)/2}
   \cosh\!\left((t-M/2)\log(\lambda_a/\lambda_b)\right).
 \end{split}
 \label{eq:f15v57:cosh}
\end{equation}
It is symmetric in \(t,M-t\) and
decreases as \(t\) moves from one
toward the midpoint. The diagonal
variance is constant in \(t\).
This proves \eqref{eq:f15v57:monotone}.

All these expansions converge.
For example,
\(\sum_{a,b}|B^f_{ab}|^2
 \lambda_a^{t-1}\lambda_b^{M-t-1}
\le\|\mathsf T\|^{M-2}\|\mathsf B_f\|_{\rm HS}^2\).
The diagonal second moment is
bounded the same way, and the
mixed expressions follow by
Cauchy--Schwarz. One may first
use finite spectral truncations
and then pass to the limit.
No lower bound on all \(\lambda_a\)
or a transfer spectral gap has
been used.
\end{proof}

\subsection{Use the inherited covariance bound without a volume factor}

Suppose the full cell covariance
obeys
\begin{equation}
 \operatorname{Var}_{\mu_\beta}
       \left(\sum_{t,x}h_{t,x}T_{(t,x)}\right)
               \le\kappa\sum_{t,x}h_{t,x}^2
                  \quad\hbox{for all real }h .
 \label{eq:f15v57:covariance-input}
\end{equation}
At the programme's selected couplings,
the source estimate
\eqref{f15v1:source-input} proves this
with \(\kappa=k_M\le5\) eventually.
Alternatively the inherited exact
norm \eqref{f15v1:covariance-norm}
provides \(\kappa=\chi_{M,\beta}\)
at any finite regulator where that
identity is being used.
No new arbitrary-source CGF is
claimed in this section.

For \(1\le d\le\lfloor M/2\rfloor\)
write
\begin{equation}
                 w_M(d)=\min\{2d,M-1\}.
 \label{eq:f15v57:multiplicity}
\end{equation}
This counts nonzero cyclic times
of distance at most \(d\) from zero.

\begin{theorem}[A time-slice covariance-operator decay bound]
\label{f15v57:time-decay}
Let \(d=\min\{t,M-t\}\ge1\).
Under \eqref{eq:f15v57:covariance-input},
\begin{equation}
 0\le c_f(t)\le
                  \frac{\kappa}{w_M(d)}\|f\|_2^2,\qquad
 |c_{f,g}(t)|\le
                  \frac{\kappa}{w_M(d)}\|f\|_2\|g\|_2 .
 \label{eq:f15v57:time-decay}
\end{equation}
Equivalently, the spatial covariance
matrix between the two cell slices,
viewed as a symmetric operator on
spatial smearings, lies between
zero and \(\kappa I/w_M(d)\).
The bound is independent of the
number of spatial sites.
\end{theorem}

\begin{proof}
Time translation invariance gives
the finite susceptibility identity
\[
 \sum_{s=0}^{M-1}c_f(s)
      =\frac1M\operatorname{Var}
                    \left(\sum_{s=0}^{M-1}F_s^f\right)
      \le\kappa\|f\|_2^2 .
\]
The inequality uses
\eqref{eq:f15v57:covariance-input}
with \(h_{s,x}=f(x)/\sqrt M\).
All noncontact terms are
nonnegative by the transfer
argument, and the contact term
is a variance. Each of the
\(w_M(d)\) nonzero times at
distance at most \(d\) has
covariance at least \(c_f(d)\).
Keeping just these terms proves
the first inequality.

Positivity of the bilinear
form gives Cauchy--Schwarz:
\[
              |c_{f,g}(t)|^2\le c_f(t)c_g(t).
\]
This proves the second inequality
and the equivalent operator
statement. In particular no sum
over spatial endpoints is inserted
after the susceptibility bound.
\end{proof}

\begin{theorem}[Full compact cell-energy clustering at the selected law]
\label{f15v57:spatial-decay}
For two distinct cells \(x,y\)
on the four-torus, let
\[
 D_\infty(x,y)=\max_{\mu=0,1,2,3}
                      \min\{|x_\mu-y_\mu|,M-|x_\mu-y_\mu|\}.
\]
Then
\begin{equation}
 \left|\operatorname{Cov}_{\mu_\beta}(T_x,T_y)\right|
       \le\frac{\kappa}{w_M(D_\infty(x,y))}
       \le\frac{\kappa}{2D_\infty(x,y)-1}.
 \label{eq:f15v57:spatial-decay}
\end{equation}
In particular \(\kappa=k_M\le5\)
eventually on the inherited
selected coupling sequence.
This conclusion has no sparse-family
cardinality restriction and no
additive chart-exception floor.
\end{theorem}

\begin{proof}
Choose a coordinate attaining
\(D_\infty(x,y)\) and regard it
as time. The homogeneous
isotropic Wilson action and
the symmetrically assigned
cell observable have exactly
the same slab construction
in that direction. Choose
unit spatial delta functions
at the two remaining spatial
positions in
\eqref{eq:f15v57:time-decay}.
Their \(\ell^2\) norms are one.
Since \(D_\infty\le M/2\),
\(w_M(D_\infty)\ge2D_\infty-1\).
This proves both bounds.
\end{proof}

Thus the covariance tends to
zero whenever the periodic
cell separation tends to
infinity along these selected
regulators. These are the
original centered action-density
correlations of the compact
Wilson measure. They are not
inferred from a Gaussian limit
or from independence of the
retained reference cores.

\subsection{A spatial consequence for the V56 retained predictors}

Return to the separated block
sources and baseline predictors
\(M_i^0\) of V56, with its
fixed \(p\), \(s_0>7\), and
permitted family size. Set
\[
 \delta_N=16B^2N\sqrt{\epsilon_N}\to0
\]
as there. For distinct source
cell sets define
\[
 R_{ij}=\min_{x\in\mathcal B_{p,i},\
                    y\in\mathcal B_{p,j}}D_\infty(x,y)\ge1.
\]
The covariance bridge gives
the concrete spatial estimate
\begin{equation}
 \left|\operatorname{Cov}_{\rho_0}(M_i^0,M_j^0)\right|
       \le \frac{\kappa n_p^2}{w_M(R_{ij})}+\delta_N
       \le \frac{\kappa n_p^2}{2R_{ij}-1}+\delta_N .
 \label{eq:f15v57:predictor-decay}
\end{equation}
Indeed baseline centering does
not change
\(\operatorname{Cov}(S_i,S_j)
=\sum_{x\in\mathcal B_{p,i},y\in\mathcal B_{p,j}}
 \operatorname{Cov}(T_x,T_y)\).
Apply \eqref{eq:f15v57:spatial-decay}
to each of its \(n_p^2\) terms.
In V56's exact reference covariance
decomposition the conditional-variance
term is diagonal. For \(i\ne j\)
its contribution is zero, while
the entry of the covariance
error is bounded by its
operator norm \(\delta_N\).
This proves the claim.

Unlike the original-cell estimate,
the predictor estimate retains the
specified comparison error floor.
Neither bound is a distance-summable
four-dimensional exponential interaction
norm. The proof is for the homogeneous
baseline law; an independently
inhomogeneous source background
does not have one repeated transfer
operator, so V56's entire nonzero
source cube has not acquired this
spatial estimate.

\subsection{Finite susceptibility and positive transfer still do not force a gap}

The new ingredient excludes V56's
antipodal-ring example: that example
does not have the required
noncontact time monotonicity.
It does not exclude gapless
positive transfer spectra.
To see the remaining distinction
without a simulation, take the
positive contraction given by
multiplication by \(\lambda\)
on \(L^2((0,1),(1-\lambda)d\lambda)\),
and the vector \(1\). Its
correlation moments are
\begin{equation}
 c(n)=\int_0^1\lambda^n(1-\lambda)\,d\lambda
                =\frac1{(n+1)(n+2)},\qquad n\ge0,
 \qquad
 c(0)+2\sum_{n\ge1}c(n)=\frac32 .
 \label{eq:f15v57:gapless}
\end{equation}
The Hankel quadratic forms
\(\sum_{i,j}a_ia_jc(i+j)\)
are nonnegative because they
equal the integral of
\((\sum_i a_i\lambda^i)^2\).
Nevertheless the spectral
support reaches arbitrarily
close to \(\lambda=1\);
\(-\log\lambda\) has no
positive lower bound.
Adding a separate vacuum
eigenvector does not open a
gap in this centered sector.
This is an abstract spectral
example, not a Wilson model.

Accordingly, the next missing
physical information is suppression
of low-energy transfer spectral
weight strong enough to create
a positive threshold after physical
scale calibration, not just its
integrability. The present \(1/D\)
estimate is a genuine baseline
spatial advance but does not
supply that threshold. It also
does not construct renormalized
continuum observables or establish
clustering for every gauge-invariant
field required by the full problem.

\subsection{Exact diagnostics and predecessor preservation}

The new diagnostic is
\begin{center}\small
\path{scripts/verify_ym_f15_v57_transfer_spatial_decay.py}.
\end{center}
Two strictly positive real symmetric
finite transfer models with positive
operator spectra are checked at
periods \(4,6,8,10\).
There are 432 exact thermally
centered spectral identities,
56 positive smearing-covariance
matrices, 20 Loewner monotonicity
checks, 48 periodic reflection
identities and 48 finite-period
susceptibility decay inequalities.
Seventy-two independent cyclic
path-sum checks include the
separate contact insertion.
Eight cell-to-slab decompositions
check the literal face weights
on spatial sides two and four,
including mixed-sign smearings.
Sixteen positive spectral moment
matrices and 21 exact partial
susceptibility sums check the
gapless illustration.

All computations use rational
arithmetic. They test finite
algebra, source incidence and
normalization, not the inherited
Wilson covariance bound or a
continuum mass gap.
The script has 9,698 bytes and SHA--256
\begin{center}\scriptsize\ttfamily
FF20CBF155275F02B3877949141CB4BD01E75FB2164391DAD5D6BC6222BCDF06.
\end{center}
The V56 predecessor has 1,297,625
bytes and SHA--256
\begin{center}\scriptsize\ttfamily
DAA282AEA17AF4C6B016C18B1ACD6474AB9CA6B27C1AB0CE2766CCD0B489D543.
\end{center}
Removing this final section and
reversing the title increment
recovers V56 byte for byte.
V57 replaces it as the sole
active main source; protected
companions, archives and the
99 preceding diagnostics are
not changed. Only \path{.tex}
files are kept in \path{latex_notes}.
Build outputs are directed to
\begin{center}\small
\path{artifacts/ym_f15_v57_texcheck/}.
\end{center}

\begin{firewall}
V57 proves volume-uniform algebraic
clustering of the specified original
cell-energy covariance under the
homogeneous selected Wilson law,
using its inherited covariance
bound and exact positive transfer.
The result is not restricted to
a sparse source family. Its
retained-predictor corollary still
has V56's comparison floor and
fixed-depth restrictions.
No exponential clustering,
physical-scale spectral threshold,
interacting continuum construction,
or Yang--Mills mass gap has
been established.
\end{firewall}

\section{V58: thermal transition weights, vacuum probes and the normalization cost}
\label{f15v58:section}

\subsection{Which spectrum does the periodic correlator measure?}

V57 proves algebraic decay under
the actual selected Wilson law.
Its finite-period spectral formula
contains transitions between
\emph{all} thermally occupied
eigenstates, not just excitations
above the vacuum. This distinction
must be resolved before using that
formula in a Hamiltonian gap argument.
The present section extracts a
positive transition measure,
including degenerate eigenspaces,
and bounds its low-energy weight.
It then isolates the vacuum
spectral measure with its exact
thermal probability still present.

Two further calculations make that
probability a concrete upstream
quantity. A doubled-period partition
function bounds it through a purity
ratio, while a direct Wilson
estimate proves closeness of
intensive periodic and vacuum
pressures without asserting
vacuum dominance. Finally the
spatially constant cell insertion
is identified with a weighted
derivative of the actual vacuum.
These are finite-regulator spectral
and normalization statements, not
a positive physical mass threshold.

The dense-core exponential criterion
is already proved in f7 V25 and
is not repeated here. Neither
small spectral weight near zero
nor control of one action-density
probe family satisfies that criterion.
The previous goal turn made
verified progress through V57;
its predecessor hash was checked
before this append. The next
companion checkpoint remains V60.

\subsection{Separate spatial size, time period and coupling}

Fix a spatial torus of side \(L\ge4\),
write \(V_{\rm sp}=L^3\), and
fix \(\beta>0\). Let \(\mathsf T\)
be the actual Wilson transfer
operator of V57 at that spatial
regulator. Give time period \(N\ge4\),
and write
\begin{equation}
 Z_N=\Tr\mathsf T^N,\qquad
 K=N-1,\qquad
 c_f(t)=\operatorname{Cov}(F_0^f,F_t^f),
 \qquad \chi_f=\sum_{t=0}^{N-1}c_f(t).
 \label{eq:f15v58:regulators}
\end{equation}
The insertion \(\mathsf B_f\)
and the same-slab second insertion
are exactly those of V57.
The source smearing \(f\) is real.
In particular all source centering
and all trace normalizers are
the physical ones.

The spectral identities below
hold at every such rectangle.
When the inherited covariance
input is invoked, we state it as
\begin{equation}
                 \chi_f\le\kappa\|f\|_2^2 .
 \label{eq:f15v58:susceptibility-input}
\end{equation}
V57 proves this from the selected
full cell covariance bound on
the diagonal sequence \(L=N=M\),
\(\beta=\beta_M^\diamond\),
with \(\kappa=k_M\le5\) eventually.
Changing \(N\) while fixing
\(L,\beta\) is not automatically
a new instance of that selection.
No uniform rectangular source
estimate is silently imported.

\subsection{Remove basis-dependent degeneracy bookkeeping}

Let \(\Pi_\lambda\) denote the
spectral projection of \(\mathsf T\)
at its positive eigenvalue
\(\lambda\). Each such projection
has finite rank. Put
\[
 m_f=\frac{\Tr(\mathsf B_f\mathsf T^{N-1})}{Z_N}.
\]
Define a stationary nonnegative
quantity and a positive measure
on strictly positive energies by
\begin{equation}
 \begin{split}
 v_{N,f}
 &=\frac1{Z_N}\sum_\lambda\lambda^{N-2}
       \|\Pi_\lambda\mathsf B_f\Pi_\lambda
                    -m_f\lambda\Pi_\lambda\|_{\rm HS}^2,\\
 \nu_{N,f}
 &=\frac1{Z_N}\sum_{\lambda>\mu>0}
       \lambda^{N-2}
         \|\Pi_\lambda\mathsf B_f\Pi_\mu\|_{\rm HS}^2
                  \,\delta_{\log(\lambda/\mu)} .
 \end{split}
 \label{eq:f15v58:transition-measure}
\end{equation}
The sum in the measure is over
distinct eigenvalues, not over
an arbitrarily selected ordering
of eigenvectors within a degenerate
eigenspace.
Its energy coordinate is an
eigenvalue \emph{difference} for
the normalized transfer Hamiltonian.

\begin{proposition}[Exact thermal spectral ledger]
\label{f15v58:thermal}
The quantities in
\eqref{eq:f15v58:transition-measure}
are finite. For \(1\le t\le N-1\),
\begin{equation}
 c_f(t)=v_{N,f}
    +\int_{(0,\infty)}
         \left(e^{-(t-1)E}+e^{-(N-t-1)E}\right)
                                                   d\nu_{N,f}(E).
 \label{eq:f15v58:thermal}
\end{equation}
Writing
\[
 G_K(r)=\sum_{j=0}^{K-1}r^j,
 \qquad 0\le r\le1,
\]
the full susceptibility is
\begin{equation}
 \chi_f=c_f(0)+K v_{N,f}
                     +2\int G_K(e^{-E})\,d\nu_{N,f}(E).
 \label{eq:f15v58:ledger}
\end{equation}
The contact term is the actual
same-slab variance and is
nonnegative; it has not been
replaced by two first insertions.
\end{proposition}

\begin{proof}
Group the V57 eigenbasis formula
by equal eigenvalues. Terms
within the \(\lambda\)-eigenspace
have weight \(\lambda^{N-2}\).
Expanding the square in the
definition of \(v_{N,f}\) gives
\[
 v_{N,f}=
 \frac1{Z_N}\sum_\lambda\lambda^{N-2}
             \|\Pi_\lambda\mathsf B_f\Pi_\lambda\|_{\rm HS}^2
                                                    -m_f^2 .
\]
The two cross terms sum to
\(-2m_f^2\), and the scalar
square sums to \(m_f^2\)
because \(\sum_\lambda
\lambda^N\operatorname{rank}\Pi_\lambda=Z_N\).
Thus this formula retains both
the diagonal thermal variance
and all off-diagonal entries
inside degenerate blocks.
Its square representation
proves nonnegativity and
basis independence.

For \(\lambda>\mu\), the two
directions of the transition
contribute
\[
 \frac{\|\Pi_\lambda\mathsf B_f\Pi_\mu\|_{\rm HS}^2}{Z_N}
 \left[\lambda^{N-t-1}\mu^{t-1}
            +\mu^{N-t-1}\lambda^{t-1}\right].
\]
Factor out \(\lambda^{N-2}\)
and put \(E=\log(\lambda/\mu)\).
This proves \eqref{eq:f15v58:thermal}.
The Hilbert--Schmidt bound from
V57 controls all these sums
by a constant times
\(\|\mathsf T\|^{N-2}
\|\mathsf B_f\|_{\rm HS}^2/Z_N\).
The mean and centered terms
are finite as well.
Finally sum over the \(K\)
noncontact times and retain
the contact variance separately.
All terms integrated in this
last step are nonnegative.
\end{proof}

\begin{theorem}[A paid low-transition-energy estimate]
\label{f15v58:low-transition}
Under \eqref{eq:f15v58:susceptibility-input},
\begin{equation}
 \begin{gathered}
 0\le v_{N,f}\le\frac{\kappa\|f\|_2^2}{K},\\
 \nu_{N,f}((0,E])
 \le\frac{\kappa\|f\|_2^2}{2G_K(e^{-E})}
 \le\frac{\kappa\|f\|_2^2}{2}
                       \left(1-e^{-E}+\frac1K\right)
 \le\frac{\kappa\|f\|_2^2}{2}
                       \left(E+\frac1K\right)
             \quad(E>0).
 \end{gathered}
 \label{eq:f15v58:low-transition}
\end{equation}
In particular these are volume-uniform
thermal transition bounds at the
selected \(L=N=M\) regulators.
\end{theorem}

\begin{proof}
Drop the other nonnegative terms
in \eqref{eq:f15v58:ledger} to
obtain the first bound. On the
energy interval \((0,E]\), the
function \(G_K(e^{-\cdot})\) is
at least \(G_K(e^{-E})\), giving
the second bound.
For \(0<r<1\), Bernoulli's
inequality gives
\[
 r^{-K}\ge1+K(r^{-1}-1)\ge1+K(1-r).
\]
Consequently
\[
 G_K(r)=\frac{1-r^K}{1-r}
                      \ge\frac{K}{1+K(1-r)} .
\]
Its endpoint at \(r=1\) follows
by continuity. Substitute
\(r=e^{-E}\), and use
\(1-e^{-E}\le E\).
\end{proof}

This estimate bounds weight,
not the lower endpoint of
spectral support. Transitions
between two excited levels can
have arbitrarily small energy
difference even when both levels
lie well above the vacuum.
They must not be renamed
vacuum excitations.

\subsection{The vacuum exists at a fixed regulator, but its weight is not one}

\begin{proposition}[A simple maximal eigenvalue at finite spatial regulator]
\label{f15v58:vacuum}
The actual \(\mathsf T\) has
a simple largest eigenvalue
\(\lambda_0>0\), with a normalized
strictly positive real gauge-invariant
eigenfunction \(\Omega\).
All other eigenvalues are
strictly smaller than \(\lambda_0\).
Thus
\begin{equation}
 \mathsf P=\mathsf T/\lambda_0,\qquad
 H=-\log\mathsf P,\qquad
 Q=I-|\Omega\rangle\langle\Omega|
 \label{eq:f15v58:normalized-transfer}
\end{equation}
define the finite-regulator transfer
Hamiltonian with vacuum energy zero.
Its time unit here is one lattice
step, not a calibrated continuum
time unit.
\end{proposition}

\begin{proof}
V57's kernel is strictly positive
at every pair of slice configurations.
The positive compact operator has
a maximal eigenvalue and a
Rayleigh maximizing vector.
For a real vector with positive
and negative parts of nonzero
measure, replacing it by its
absolute value strictly increases
the Rayleigh quotient. This
follows by integrating the
strictly positive kernel over
the product of those two sets.
Hence a maximizing real vector
can be taken nonnegative.
Applying \(\mathsf T\) makes it
strictly positive everywhere.

If the maximal eigenspace had
dimension greater than one,
it would contain a real vector
orthogonal to this positive
eigenfunction. Such a vector
must change sign, contradicting
the strict Rayleigh argument.
Simplicity follows. Compactness
isolates the nonzero maximal
eigenvalue from the remaining
spectrum at this fixed regulator.
V57's injectivity allows the
spectral definition of \(H\).
No lower bound uniform in
\(L,\beta\) is obtained by
this finite-regulator argument.
\end{proof}

Define the exact thermal vacuum
probability and the centered
one-slab vacuum probe by
\begin{equation}
 p_{0,N}=\frac{\lambda_0^N}{Z_N},\qquad
 \psi_f=\lambda_0^{-1}Q\mathsf B_f\Omega,\qquad
 \eta_f(A)=\langle\psi_f,1_A(H)\psi_f\rangle .
 \label{eq:f15v58:vacuum-probe}
\end{equation}
Here \(\eta_f\) is supported on
positive vacuum excitation
energies. The projection \(Q\)
implements vacuum centering.
It need not coincide with
subtracting the thermal mean,
although subtracting any scalar
multiple of \(\mathsf T\) from
\(\mathsf B_f\) leaves \(\psi_f\)
unchanged.

\begin{theorem}[Thermal-to-vacuum spectral comparison with its exact cost]
\label{f15v58:vacuum-comparison}
As positive measures on \((0,\infty)\),
\begin{equation}
                         \nu_{N,f}\ge p_{0,N}\eta_f .
 \label{eq:f15v58:measure-domination}
\end{equation}
Consequently
\begin{equation}
 \begin{gathered}
 2p_{0,N}\langle\psi_f,G_K(\mathsf P)\psi_f\rangle
                         \le\chi_f\le\kappa\|f\|_2^2,\\
 \eta_f((0,E])
 \le\frac{\kappa\|f\|_2^2}
                       {2p_{0,N}G_K(e^{-E})}
 \le\frac{\kappa\|f\|_2^2}{2p_{0,N}}
                                \left(E+\frac1K\right).
 \end{gathered}
 \label{eq:f15v58:vacuum-budget}
\end{equation}
The susceptibility inequalities
use \eqref{eq:f15v58:susceptibility-input}.
\end{theorem}

\begin{proof}
In \eqref{eq:f15v58:transition-measure},
keep only pairs with \(\lambda=\lambda_0\).
For an excited eigenvalue \(\mu\),
their weight is
\[
 \frac{\lambda_0^{N-2}}{Z_N}
             \|\Pi_\mu\mathsf B_f\Omega\|^2
   =p_{0,N}
             \|\Pi_\mu\psi_f\|^2 ,
\]
at energy \(\log(\lambda_0/\mu)\),
which is its \(H\)-eigenvalue.
All remaining transition weights
are nonnegative. This proves
the measure domination.
Integrate \(G_K(e^{-E})\)
and use the nonnegative ledger
\eqref{eq:f15v58:ledger}.
The spectral theorem identifies
the resulting vacuum integral
with the displayed quadratic
form. The low-energy bounds
then follow as before.
\end{proof}

This is a truncated resolvent
estimate for the actual physical
transfer operator. Dropping
\(p_{0,N}\), or replacing the
whole thermal transition measure
by \(\eta_f\), would change the
statement and is not justified.
Taking \(N\to\infty\) at one fixed
\(L,\beta\) does make \(p_{0,N}\to1\),
by the isolated largest eigenvalue
and the trace-class bound. But
the inherited sequence changes
\(L,\beta\) with \(N\), and no
uniform version of that argument
has yet been proved.

\subsection{A doubled-period quantity that can pay the vacuum probability}

Define
\begin{equation}
               R_N=\frac{Z_{2N}}{Z_N^2}
                    =\Tr\left(\frac{\mathsf T^N}{Z_N}\right)^2 .
 \label{eq:f15v58:purity}
\end{equation}
The numerator is the Wilson
partition function on time
period \(2N\), with the same
spatial \(L\), the same \(\beta\)
and the same normalized Haar
convention. This is not a
second selected-coupling input.

\begin{proposition}[Exact vacuum-probability brackets]
\label{f15v58:purity}
One has
\begin{equation}
                         R_N\le p_{0,N}\le\sqrt{R_N}.
 \label{eq:f15v58:purity-brackets}
\end{equation}
Hence a proved lower bound on
\(R_N\) is sufficient to pay
the inverse vacuum probability
in \eqref{eq:f15v58:vacuum-budget}.
No such uniform lower bound
is assumed here.
\end{proposition}

\begin{proof}
The eigenvalues
\(\omega_a=\lambda_a^N/Z_N\)
form a probability distribution
with maximum \(\omega_0=p_{0,N}\).
Thus
\[
 p_{0,N}^2\le\sum_a\omega_a^2
                           \le p_{0,N}\sum_a\omega_a=p_{0,N}.
\]
The middle quantity is exactly
\(R_N\). This proves both brackets.
The ratio is unchanged by a
common scalar rescaling of
\(\mathsf T\), as it must be.
\end{proof}

\subsection{An actual Wilson pressure comparison without vacuum dominance}

\begin{theorem}[Finite-period versus vacuum pressure]
\label{f15v58:pressure}
For the pure Wilson transfer
operator at spatial volume
\(V_{\rm sp}=L^3\),
\begin{equation}
 \begin{gathered}
 e^{-12\beta V_{\rm sp}}\le p_{0,N}\le1,\\
 0\le
 \frac{\log Z_N}{NV_{\rm sp}}
             -\frac{\log\lambda_0}{V_{\rm sp}}
      =-\frac{\log p_{0,N}}{NV_{\rm sp}}
      \le\frac{12\beta}{N}.
 \end{gathered}
 \label{eq:f15v58:pressure}
\end{equation}
In particular on the original
selected diagonal sequence,
the intensive pressure difference
is at most \(12\log M/M\to0\).
\end{theorem}

\begin{proof}
There are \(3V_{\rm sp}\)
spatial plaquettes and the
same number of temporal
plaquettes in one slab.
Each deficit lies in \([0,2]\).
The two half spatial actions
in V57's slab therefore
contribute at most \(6V_{\rm sp}\),
and the temporal action at
most another \(6V_{\rm sp}\).
It follows that
\[
 \mathsf T(U,V)\ge e^{-12\beta V_{\rm sp}}.
\]
The constant function has
unit \(L^2\) norm and is
gauge invariant, so its
Rayleigh quotient gives
\(\lambda_0\ge e^{-12\beta V_{\rm sp}}\).

Also \(\Tr\mathsf T\le1\).
For example the continuous
positive-kernel trace equals
the integral of
\(\mathsf T(U,U)\le1\) against
normalized slice Haar.
Equivalently, V57's positive
character factorization gives
\(\Tr\mathsf K=1\) and
\(\Tr(\mathsf A\mathsf K\mathsf P\mathsf A)
\le\Tr\mathsf K\), with the
spatial multiplier at most one.
Since \(0<\lambda_a/\lambda_0\le1\),
\[
 p_{0,N}^{-1}
       =\sum_a(\lambda_a/\lambda_0)^N
       \le\frac{\Tr\mathsf T}{\lambda_0}
       \le e^{12\beta V_{\rm sp}}.
\]
Taking logarithms proves the
first two assertions. Finally
the selected sequence has
\(\beta_M^\diamond\le\log M\).
\end{proof}

The lower bound on \(p_{0,N}\)
is exponentially small in the
spatial volume. It is not a
useful volume-uniform vacuum
weight, even though the
pressure difference tends to
zero on the selected sequence.
These two conclusions have
different normalizations.

For an exact illustration of
the distinction, let \(d=N^2\),
\(r_N=1-1/N\), and let
\(\Pi\) be projection onto
the constant vector in
\(\mathbb R^{d+1}\). The
matrix
\[
 T_N=\frac{r_N I+(1-r_N)\Pi}{1+d r_N}
\]
has strictly positive entries,
positive spectrum and trace one.
Its top eigenvalue is
\((1+d r_N)^{-1}\), and the
other \(d\) eigenvalues are
\(r_N/(1+d r_N)\).
At time period \(N\),
\[
 p_{0,N}=\frac1{1+d r_N^N}\sim\frac e{N^2}\to0,\qquad
 -\frac{\log p_{0,N}}N\to0 .
\]
Thus even pressure closeness
per time step need not imply
vacuum dominance. Its normalized
excitation energy
\(-\log r_N\to0\) also shows
why finite-regulator positivity
alone is insufficient.
This example is an abstract
positive-kernel model, not a
claim about the Wilson spectrum.

\subsection{The uniform action insertion is a weighted vacuum derivative}

There is a further concrete
identification for the normalized
constant spatial smearing
\(f=V_{\rm sp}^{-1/2}\).
Its norm is one. V57's
literal slab formula gives
\begin{equation}
                  \mathsf B_f
                       =-\frac{\beta}{\sqrt{V_{\rm sp}}}
                                      \,\partial_\beta\mathsf T .
 \label{eq:f15v58:coupling-insertion}
\end{equation}
The two half spatial actions
and all temporal plaquettes
are included in this derivative.

\begin{proposition}[An exact vacuum-response identity]
\label{f15v58:vacuum-derivative}
At each fixed spatial regulator
choose the real normalized vacuum
smoothly in \(\beta\), and let
\(\Omega'=\partial_\beta\Omega\).
Then
\begin{equation}
 \begin{gathered}
 \langle\Omega,\Omega'\rangle=0,\qquad
 \psi_f=-\frac{\beta}{\sqrt{V_{\rm sp}}}
                              (I-\mathsf P)\Omega',\\
 \langle\psi_f,G_K(\mathsf P)\psi_f\rangle
   =\frac{\beta^2}{V_{\rm sp}}
       \langle\Omega',(I-\mathsf P)
                             (I-\mathsf P^K)\Omega'\rangle .
 \end{gathered}
 \label{eq:f15v58:vacuum-derivative}
\end{equation}
Whenever \eqref{eq:f15v58:susceptibility-input}
holds, the last expression is
at most \(\kappa/(2p_{0,N})\).
\end{proposition}

\begin{proof}
At fixed finite spatial regulator
the bounded slab action makes
\(\mathsf T(\beta)\) locally
analytic as an operator. The
simple isolated eigenvalue has
a locally analytic rank-one
spectral projection: a small
contour around it and the
resolvent Neumann series prove
this directly. Normalizing its
positive real vector gives
a differentiable \(\Omega\).
No common contour size across
regulators is being asserted.
Differentiating its norm proves
the first identity.

Differentiating
\(\mathsf T\Omega=\lambda_0\Omega\)
and projecting orthogonally
to \(\Omega\) gives
\[
 Q(\partial_\beta\mathsf T)\Omega
                      =\lambda_0(I-\mathsf P)\Omega'.
\]
Use \eqref{eq:f15v58:coupling-insertion}
and the definition of \(\psi_f\).
Finally the scalar polynomial
identity
\((1-r)^2G_K(r)=(1-r)(1-r^K)\)
holds also at \(r=1\).
Functional calculus proves
the second equality. The
bound follows from
\eqref{eq:f15v58:vacuum-budget}.
\end{proof}

This controls a weighted response
of the actual vacuum to the
Wilson coupling. It does not
bound \(\|\Omega'\|\), because
the factors \(I-\mathsf P\)
and \(I-\mathsf P^K\) suppress
low-energy directions.
It is therefore not a lower
spectral bound in disguise.

\subsection{What remains for a mass-gap proof}

The new low-energy estimate is
proved for the actual thermal
transition measure, and its
vacuum counterpart has an exact
probability cost. The pressure
comparison above does not pay
that cost uniformly. A lower
bound for the explicitly defined
ratio \(R_N\), or another genuine
vacuum comparison, is one
possible upstream target.
Even paying it would only bound
low-energy \emph{weight}; a
positive spectral-support
threshold is a further requirement.

The probes here arise from the
specified cell-energy insertions.
They are not proved to span a
dense vacuum-orthogonal physical
space. Neither their spectral
weight bounds nor their covariance
decay may be substituted for
the existing dense-core gap
criterion. Finally, the lattice
generator \(H\) must be related
to a physical time scale before
a uniform positive energy can
be called a continuum mass.
No such calibration, interacting
continuum construction or
positive mass gap is completed
by this section.

\subsection{Exact verification and source preservation}

The new exact diagnostic is
\begin{center}\small
\path{scripts/verify_ym_f15_v58_thermal_vacuum_spectrum.py}.
\end{center}
It reruns V57's preserved tests.
Four positive transfer models,
including an excited degeneracy,
give 96 exact stationary/transition
identities, 16 susceptibility
ledgers, 80 low-transition
window bounds, 24 vacuum-measure
domination checks and 16 paid
vacuum truncated-resolvent bounds.
Sixteen purity brackets retain
the actual thermal normalization,
and four tests rotate a
degenerate eigenspace without
changing the spectral ledger.

There are also 16 rational
geometric-sum bounds, nine
vacuum-derivative identities,
five non-dominant-vacuum
positive-kernel examples and
three slab action-range counts.
The logarithmic energy estimate,
operator spectral statements,
Wilson input and limiting claims
use their analytic proofs,
not finite sampling.
All diagnostic arithmetic
is rational.

The script has 8,092 bytes and SHA--256
\begin{center}\scriptsize\ttfamily
7413009C3F5CAC22831CFAC64E4BD8BD7F156D30174BCD9919FD3D8E5889A3DD.
\end{center}
The V57 predecessor has 1,317,907
bytes and SHA--256
\begin{center}\scriptsize\ttfamily
442D017B0FD3DBB55329A9B76B298FF817C363B00611BA70A24DC4B00F3DFBAE.
\end{center}
Removing this final section
and reversing the title
increment recovers V57
byte for byte. V58 replaces
V57 as the sole active main
source. Protected companions,
archives and the 100 earlier
diagnostics remain unchanged.
Only \path{.tex} files are
kept in \path{latex_notes};
compilation products go to
\begin{center}\small
\path{artifacts/ym_f15_v58_texcheck/}.
\end{center}

\begin{firewall}
V58 bounds actual low-energy
thermal transition weight,
isolates the vacuum spectral
measure with its exact thermal
probability, gives a doubled-period
partition ratio that can bound
that probability, and proves
an intensive Wilson vacuum-pressure
comparison and a weighted
vacuum-derivative identity.
No uniform vacuum dominance,
positive continuum spectral
threshold, interacting continuum
construction or Yang--Mills
mass gap is asserted.
\end{firewall}

\section{V59: vacuum action-density saturation from one-sided pressure comparison}
\label{f15v59:context}

The next upstream step is a
first-moment comparison under
the actual Wilson vacuum.
V58's spectral comparison
requires division by a possibly
small thermal vacuum probability.
Its intensive pressure comparison
does not. This section combines
that pressure comparison with
convex secants and the already
proved homogeneous source window.
It transfers action-density
saturation to the vacuum on
the \emph{same} selected couplings.
No new coupling selection or
vacuum-dominance assumption is
introduced.

Here \emph{vacuum action density}
means the vacuum expectation
of the Wilson slab action
with its thermal factor
\(\beta\). It is not the
vacuum eigenvalue of
\(-\log(\mathsf T/\lambda_0)\),
which is zero by normalization.
Nor is it a physical mass.
The extension to arbitrary
vacuum observables remains open.

\subsection{Two pressures at a common spatial regulator}

Keep V57--V58's transfer
operator and slab action
\(\mathcal A\). Fix a spatial
side \(L\ge4\), put
\(V_{\rm sp}=L^3\), and let
the temporal period \(N\ge4\)
be independent of \(L\).
The following notation retains
all volume normalizations:
\begin{equation}
 \begin{gathered}
 F_{L,N}(\beta)=\frac{\log Z_{L,N}(\beta)}{NV_{\rm sp}},
 \qquad
 G_L(\beta)=\frac{\log\lambda_{0,L}(\beta)}{V_{\rm sp}},\\
 e_{L,N}(\beta)=-\beta F_{L,N}'(\beta)
      =\frac{\beta\,\mathbb E_{L,N,\beta}S}{NV_{\rm sp}},
 \qquad
 e_{L,\mathrm{vac}}(\beta)=-\beta G_L'(\beta).
 \end{gathered}
 \label{eq:f15v59:pressures}
\end{equation}
In particular \(G_L\) does
not depend on \(N\).
The pressure theorem
\ref{f15v58:pressure} says
\begin{equation}
 0\le F_{L,N}(\beta)-G_L(\beta)
                  \le a_{N,\beta}:=\frac{12\beta}{N}
                  \qquad(\beta>0).
 \label{eq:f15v59:one-sided}
\end{equation}
We reuse this result, including
its trace normalization, without
reproving it.

\begin{proposition}[Convexity and the literal vacuum action mean]
\label{f15v59:convexity}
At each fixed spatial regulator,
\(F_{L,N}\) and \(G_L\) are
convex on \((0,\infty)\);
\(G_L\) is differentiable there.
If \(\Omega=\Omega_{L,\beta}>0\)
is the real normalized top
eigenvector, then the measure
\begin{equation}
 d\pi_{L,\beta}(U,V,g)
   =\lambda_{0,L}^{-1}\Omega(U)\Omega(V)
       e^{-\beta\mathcal A(U,V,g)}
                       \,dH(U)\,dH(V)\,dH(g)
 \label{eq:f15v59:vacuum-slab}
\end{equation}
is a probability measure and
\begin{equation}
 e_{L,\mathrm{vac}}(\beta)
       =\frac{\beta}{V_{\rm sp}}
                          \mathbb E_{\pi_{L,\beta}}\mathcal A .
 \label{eq:f15v59:vacuum-action}
\end{equation}
\end{proposition}

\begin{proof}
At finite volume,
\(F_{L,N}''(\beta)=
\Var_{L,N,\beta}(S)/(NV_{\rm sp})\ge0\),
by differentiation under the
compact Wilson integral.
For fixed \(L\), the estimate
\eqref{eq:f15v59:one-sided}
shows that
\(F_{L,N}(\beta)\to G_L(\beta)\)
at every \(\beta>0\) as
\(N\to\infty\). Passing to the
limit in the defining convexity
inequality proves convexity
of \(G_L\). This use of
arbitrary temporal periods
requires only V58's pressure
theorem, not a source estimate
on new rectangles.

V58 proved local analytic
dependence of the simple
isolated top eigenvalue and
its normalized real eigenvector
at each fixed \(L,\beta\).
In particular \(G_L'\) exists.
The integral of
\eqref{eq:f15v59:vacuum-slab}
is
\(\lambda_{0,L}^{-1}
\langle\Omega,\mathsf T\Omega\rangle=1\).
It is nonnegative because
\(\Omega\) and the kernel
integrand are positive.
Differentiating the eigenvalue
equation and using
\(\langle\Omega,\Omega\rangle=1\)
gives
\[
 \lambda_{0,L}'
       =\langle\Omega,\mathsf T'\Omega\rangle
       =-\lambda_{0,L}\,
                          \mathbb E_{\pi_{L,\beta}}\mathcal A .
\]
The bounded slab action
justifies differentiating its
kernel integral. This proves
\eqref{eq:f15v59:vacuum-action}.
No derivative of a normalized
vacuum probability has been
discarded: the eigenvector
derivative terms cancel by
normalization in the
eigenvalue identity.
\end{proof}

\subsection{The source window belongs to a single baseline law}

Set \(L=N=M\) and
\(\beta=\beta_M^\diamond\)
from \eqref{f15v1:source-input}.
Abbreviate \(F=F_{M,M}\),
\(G=G_M\), \(e_{\rm per}=e_{M,M}(\beta)\),
\(e_{\rm vac}=e_{M,\mathrm{vac}}(\beta)\),
and put
\[
 J(s)=-\log(1-s)-s,\qquad
 k=k_M,\qquad R=R_M,\qquad a=\frac{12\beta}{M}.
\]
For \(|s|\le R\) the exact
uniform-source identity and
the inherited source estimate
give
\begin{equation}
 0\le F(\beta(1-s))-F(\beta)-s e_{\rm per}
                                           \le kJ(s).
 \label{eq:f15v59:homogeneous}
\end{equation}
Indeed the expression, multiplied
by \(M^4\), is
\[
 \log\mathbb E_{M,M,\beta}
          \exp\{s(\beta S-\mathbb E_{M,M,\beta}\beta S)\},
\]
because \(\sum_xT_x=\beta S\).
Jensen's inequality gives the
lower bound, and
\eqref{f15v1:source-input}
with \(b_x=s\) gives the upper
bound. Nearby
\(\beta(1-s)\) occurs only
through this exact partition
identity from the original
baseline law. It is not
asserted to be another
selected coupling.

\begin{theorem}[Vacuum action-density saturation]
\label{f15v59:mean}
For every
\(0<h\le\min(R_M,1/2)\),
\begin{equation}
 |e_{M,\mathrm{vac}}(\beta_M^\diamond)
            -e_{M,M}(\beta_M^\diamond)|
       \le \frac{k_MJ(h)+12\beta_M^\diamond/M}{h}
       \le k_Mh+\frac{12\beta_M^\diamond}{Mh}.
 \label{eq:f15v59:secant-bound}
\end{equation}
Consequently, for all sufficiently
large selected dyadic \(M\),
\begin{equation}
 \begin{gathered}
 |e_{M,\mathrm{vac}}(\beta_M^\diamond)
            -e_{M,M}(\beta_M^\diamond)|
       \le17\sqrt{\frac{\beta_M^\diamond}{M}}
       \le17\sqrt{\frac{\log M}{M}},\\
 e_{M,\mathrm{vac}}(\beta_M^\diamond)\longrightarrow\frac92 .
 \end{gathered}
 \label{eq:f15v59:saturation}
\end{equation}
These conclusions concern the
actual normalized vacuum slab
measure \eqref{eq:f15v59:vacuum-slab}.
They contain no factor
\(p_{0,M}^{-1}\).
\end{theorem}

\begin{proof}
Write
\(\phi(s)=G(\beta(1-s))\)
and
\(d(s)=F(\beta(1-s))-\phi(s)\).
The preceding proposition
makes \(\phi\) convex, with
\(\phi'(0)=e_{\rm vac}\).
By \eqref{eq:f15v59:one-sided},
\[
       d(s)\ge0,\qquad d(0)\le a .
\]
In particular, only the
baseline upper bound \(d(0)\le a\)
is needed, whereas
nonnegativity is used at
both displaced couplings.
The positive secant gives
\[
 \begin{split}
 e_{\rm vac}
 &\le\frac{\phi(h)-\phi(0)}h\\
 &=\frac{F(\beta(1-h))-F(\beta)-d(h)+d(0)}h\\
 &\le e_{\rm per}+\frac{kJ(h)+a}{h}.
 \end{split}
\]
The negative secant similarly
gives
\[
 \begin{split}
 e_{\rm vac}
 &\ge\frac{\phi(0)-\phi(-h)}h\\
 &=\frac{F(\beta)-F(\beta(1+h))-d(0)+d(-h)}h\\
 &\ge e_{\rm per}-\frac{kJ(-h)+a}{h}.
 \end{split}
\]
For \(0<h\le1/2\),
\[
 J(-h)=\int_0^h\frac{t}{1+t}\,dt
 \le J(h)=\int_0^h\frac{t}{1-t}\,dt
 \le\frac{h^2}{2(1-h)}\le h^2 .
\]
This proves
\eqref{eq:f15v59:secant-bound}.

Now choose
\(h=\sqrt{\beta_M^\diamond/M}\).
Since \(\beta_M^\diamond\le\log M\)
and \(R_M\to1\), this choice
is eventually admissible.
Also \(k_M\to9/2\), hence
\(k_M\le5\) eventually.
The last expression in
\eqref{eq:f15v59:secant-bound}
is therefore at most
\((k_M+12)\sqrt{\beta_M^\diamond/M}\),
which is at most the stated
factor \(17\). Finally the
inherited mean selection
has \(e_{M,M}(\beta_M^\diamond)\to9/2\).
The triangle inequality proves
the asserted vacuum limit.
No rate for that inherited
mean convergence is supplied
here; the explicit rate
controls the \emph{difference}
of vacuum and periodic means.
\end{proof}

\subsection{A pressure-value extension and its differentiation barrier}

Define
\[
 \delta_M=\inf_{0<h\le\min(R_M,1/2)}
             \frac{k_MJ(h)+12\beta_M^\diamond/M}{h}.
\]
The theorem gives
\(|e_{\rm vac}-e_{\rm per}|\le\delta_M\);
eventually
\(\delta_M\le17\sqrt{\log M/M}\).
For every \(|s|\le R_M\),
\begin{equation}
 \begin{split}
 0&\le G_M(\beta_M^\diamond(1-s))
                  -G_M(\beta_M^\diamond)-s e_{\rm vac}\\
  &\le k_MJ(s)+|s|\delta_M+\frac{12\beta_M^\diamond}{M}.
 \end{split}
 \label{eq:f15v59:vacuum-pressure-window}
\end{equation}
The lower bound is convexity.
For the upper bound subtract
the centered periodic expression
\eqref{eq:f15v59:homogeneous}.
The difference is
\(-d(s)+d(0)+s(e_{\rm per}-e_{\rm vac})\),
which is at most
\(a+|s|\delta_M\).
This proves
\eqref{eq:f15v59:vacuum-pressure-window}
without taking a derivative
of an error term.

The left side is a centered
vacuum \emph{pressure difference}.
It is not being identified
with a finite-volume cumulant
generating function of a
single slab random variable.
The \(\beta\)-dependence of
the vacuum must be retained.
Nor may the additive pressure
floor be differentiated twice
to infer a vacuum susceptibility
bound. The following elementary
counterexample makes the latter
restriction explicit.

\begin{proposition}[A vanishing pressure floor need not control curvature]
\label{f15v59:curvature-counterexample}
For \(0<t<1\), define on
the real line
\[
 \mathcal F_t(s)=\frac{s^2}{2}+t\sqrt{s^2+t^2},
 \qquad
 \mathcal G_t(s)=\frac{s^2}{2}+t\sqrt{s^2+t^6}.
\]
These functions are real
analytic and convex. They satisfy
\begin{equation}
 \begin{gathered}
 0\le\mathcal F_t(s)-\mathcal G_t(s)\le t^2,\qquad
 \mathcal F_t'(0)=\mathcal G_t'(0)=0,\\
 0\le\mathcal F_t(s)-\mathcal F_t(0)
                         \le s^2\le3J(s)\quad(|s|\le1/2),\\
 \mathcal F_t''(0)=2,\qquad
 \mathcal G_t''(0)=1+t^{-2}\longrightarrow\infty .
 \end{gathered}
 \label{eq:f15v59:curvature-counterexample}
\end{equation}
Thus even exact mean agreement,
a uniform centered source-type
bound and a vanishing one-sided
pressure error do not by
themselves bound the second
derivative of the lower
convex pressure.
\end{proposition}

\begin{proof}
For \(c>0\),
\((\sqrt{s^2+c^2})''=
c^2(s^2+c^2)^{-3/2}>0\);
this gives convexity and the
derivatives at zero.
The difference of the two
square roots is nonnegative
and, by rationalizing its
denominator, is at most
\(t-t^3\). Multiplication
by \(t\) gives
\(0\le\mathcal F_t-\mathcal G_t
\le t^2-t^4<t^2\).
Concavity of the square root
as a function of its argument
gives
\(\sqrt{s^2+t^2}-t\le s^2/(2t)\).
Hence
\(\mathcal F_t(s)-\mathcal F_t(0)\le s^2\).
Finally, the integral formula
for \(J\) gives
\[
 J(s)\ge\frac{s^2}{2(1+|s|)}
                         \ge\frac{s^2}{3}\quad(|s|\le1/2).
\]
The remaining assertions follow.
These are abstract scalar
pressures, not counterexamples
to a Wilson-specific estimate
that uses additional input.
\end{proof}

\subsection{The remaining upstream question}

V59 reaches an actual vacuum
first moment using the selected
periodic source law, and does
so on the original diagonal
sequence. It does not yet
transfer the periodic curvature
classification, all mixed
moments or covariance estimates
to that vacuum. In particular,
the harmonic-removal argument
in the inherited periodic
classification cannot simply
be reused for a vacuum law
before its hypotheses have
been established for that law.

A next substantive target is
a uniform comparison for a
specified nonconstant vacuum
source, or an independent
vacuum second-response estimate
whose error is quadratic in
the source down to zero.
The scalar pressure floor
above is insufficient for
either claim. Even such a
response estimate would not
alone exclude spectral support
near zero, establish a dense
physical set of probes, or
calibrate lattice time against
physical time. These remain
distinct steps toward the
interacting continuum and
positive physical mass gap.
The scheduled companion
checkpoint remains V60.

\subsection{Exact checks and preservation}

The new diagnostic is
\begin{center}\small
\path{scripts/verify_ym_f15_v59_vacuum_energy_secants.py}.
\end{center}
It reruns V58 and V57.
Twelve positive two-state
Boltzmann transfer models
provide 48 homogeneous source
sample checks, 24 two-sided
vacuum-mean secant comparisons,
and 48 vacuum-pressure
remainder bounds with the
additive floor retained.
Five \(J\)-function checks
use rational logarithm
enclosures; twelve rational
checks verify the square-root
rate constant; five exact
curvature examples verify
the displayed parameter
dependence at zero.

All logarithms in those
finite checks are enclosed
by the positive
\(\operatorname{arctanh}\)
series with an explicit
rational geometric tail.
No floating-point tolerance
is used. The Wilson source
input, the full-window
inequalities and limiting
statements rely on their
analytic proofs, not this
finite set of test models.
The diagnostic has 5,392
bytes and SHA--256
\begin{center}\scriptsize\ttfamily
58E25886E796D2AC6F2CACC63813CB1BAF61B5B8591A29874BD2B4F1E6F698A8.
\end{center}

The V58 predecessor has
1,339,216 bytes and SHA--256
\begin{center}\scriptsize\ttfamily
CDEAAD36DECAFC22DB73C9BCEF2FF243DF8F416B6688A5CE697CC1A4D11E211A.
\end{center}
Removing this last section
and reversing the title
increment recovers that
predecessor byte for byte.
V59 replaces V58 as the
single active main source;
the protected inputs and
101 earlier diagnostics
are unchanged. Only
\path{.tex} files belong in
\path{latex_notes}. All
compilation products go to
\begin{center}\small
\path{artifacts/ym_f15_v59_texcheck/}.
\end{center}

\begin{firewall}
V59 proves saturation of the
actual Wilson vacuum slab
action density at \(9/2\)
on the same selected couplings,
with an explicit vanishing
periodic-to-vacuum mean error.
It does not prove convergence
of the vacuum field law,
a uniform vacuum susceptibility,
an interacting continuum
construction or a positive
Yang--Mills mass gap.
\end{firewall}

\section{V60: nearby coupling selection and a uniform vacuum response bound}
\label{f15v60:context}

V59's pressure comparison
does not control the vacuum
second derivative at the
original selected coupling.
Its counterexample rules out
that inference from the
available scalar data alone.
The present step goes upstream:
integrate the nonnegative
vacuum curvature over a
shrinking coupling window,
then make an explicitly
different selection within
that window. This gives a
uniform response estimate for
the actual vacuum, without
dividing by its thermal
probability. The new selection
is asymptotically equivalent
to the old one, but is not
identified with it.

\subsection{The required V60 companion checkpoint}

Fresh inventories of
\path{latex_notes} and the
relevant supplied families in
\path{Downloads} found no
newer versions of the seven
companions below. Their final
results and limitations were
reread, together with all of
\path{suggestions.txt}. The
116 protected inputs and
diagnostics had unchanged
hashes before this append.

\begin{center}\small
\begin{tabular}{@{}p{0.23\textwidth}p{0.71\textwidth}@{}}
\toprule
Companion & Relevant scope retained at this checkpoint\\
\midrule
SM continuum V72 &
The fixed-time many-copy limit is an auxiliary finite-mode
construction, not a growing-cutoff physical field limit.\\[1mm]
NS stability V26 &
The pointwise rank-one formulas do not supply a YM estimate;
the referenced V27 certification is not present.\\[1mm]
Shell mixing V16 &
The measured flat-holonomy one-loop assembly does not supply
a common nonperturbative remainder or physical clustering.\\[1mm]
Parallel YM V5 &
Uniform tightness of extensive global charge fails on the
stated tensorization route; compact labels must be kept distinct.\\[1mm]
SM source selection V31 &
The ancestry reduction is conditional on historically
word-native typing operations.\\[1mm]
Native stationarity V34 &
The moving-fibre first-escape calculation still requires
the specified 330 derivative entries.\\[1mm]
{\raggedright Exact-action Einstein V24\par} &
The nonzero exact first-jet derivative is not a nonlinear
root or common-lifespan theorem.\\
\bottomrule
\end{tabular}
\end{center}

The suggestions reinforce using
one measured normalization and
keeping local coefficients,
auxiliary limits and physical
mass separate. Their proposed
organizational changes are
suggestions, not instructions
or new mathematical inputs.
No companion pays the missing
vacuum response estimate.
The direction below therefore
uses the actual compact Wilson
law and its already normalized
transfer operator. The next
companion checkpoint is V65.

\subsection{A triangular average of the actual vacuum curvature}

Retain V59's notation on the
selected diagonal sequence:
\[
 \beta=\beta_M^\diamond,\quad
 F=F_{M,M},\quad G=G_M,\quad
 k=k_M,\quad R=R_M,\quad
 a=\frac{12\beta}{M},\quad
 \phi(s)=G(\beta(1-s)).
\]
These are functions at fixed
spatial regulator \(M\).
The function \(\phi\) is
analytic and convex for
\(s<1\). Let
\[
             \chi_{M,\mathrm{vac}}(\gamma)
                    =\gamma^2G_M''(\gamma),\qquad\gamma>0 .
\]
At this point the symbol
denotes a vacuum pressure
curvature. Its exact
susceptibility interpretation
is proved below.

\begin{proposition}[A positive curvature budget]
\label{f15v60:triangular}
For \(0<h\le R\),
\begin{equation}
 \begin{split}
 \int_{-h}^{h}(h-|s|)\phi''(s)\,ds
 &=\phi(h)+\phi(-h)-2\phi(0)\\
 &\le k[-\log(1-h^2)]+2a .
 \end{split}
 \label{eq:f15v60:triangular}
\end{equation}
In particular, if \(s_*\) is
any minimizer of \(\phi''\)
on \([-h,h]\), then
\begin{equation}
 0\le\phi''(s_*)\le
       B_M(h):=
       k\,\frac{-\log(1-h^2)}{h^2}+\frac{2a}{h^2}.
 \label{eq:f15v60:minimum}
\end{equation}
\end{proposition}

\begin{proof}
Integrate the continuous
second derivative separately
over \([-h,0]\) and \([0,h]\)
with the displayed linear
weights. The two first
derivatives at zero cancel,
giving the exact identity.
Put
\(d(s)=F(\beta(1-s))-\phi(s)\).
By V58--V59,
\(d(\pm h)\ge0\) and
\(d(0)\le a\). Hence
\[
 \begin{split}
 \phi(h)+\phi(-h)-2\phi(0)
 &\le F(\beta(1-h))+F(\beta(1+h))-2F(\beta)+2a\\
 &\le k[J(h)+J(-h)]+2a .
 \end{split}
\]
The periodic mean terms cancel
in V59's homogeneous source
inequalities, and
\(J(h)+J(-h)=-\log(1-h^2)\).
This proves
\eqref{eq:f15v60:triangular}.
The triangular weight has
integral \(h^2\), and
\(\phi''\ge0\). Its minimum
exists by continuity and
is no larger than the
weighted average. This gives
\eqref{eq:f15v60:minimum}.
\end{proof}

The averaging is essential.
No differentiation of the
pressure error is used.
Only the source estimate at
the original periodic baseline
is used; an estimate at the
new coupling is not assumed.

\begin{theorem}[A nearby vacuum response selection]
\label{f15v60:selection}
For all sufficiently large
selected dyadic \(M\), set
\begin{equation}
 h_M=\left(\frac{12\beta_M^\diamond}{M}\right)^{1/4},
 \qquad
 \widetilde\beta_M=\beta_M^\diamond(1-s_M),
 \quad |s_M|\le h_M,
 \label{eq:f15v60:selection}
\end{equation}
where \(s_M\) is the smallest
minimizer of \(\phi''\) on
\([-h_M,h_M]\). Then
\begin{equation}
 \begin{gathered}
 \frac{\widetilde\beta_M}{\beta_M^\diamond}\longrightarrow1,
 \qquad \widetilde\beta_M\longrightarrow\infty,\\
 0\le\chi_{M,\mathrm{vac}}(\widetilde\beta_M)
        \le\kappa_M:=(1+h_M)^2
          \left[
          k_M\frac{-\log(1-h_M^2)}{h_M^2}+2h_M^2
          \right],
 \qquad \kappa_M\longrightarrow\frac92,\\
 e_{M,\mathrm{vac}}(\widetilde\beta_M)\longrightarrow\frac92 .
 \end{gathered}
 \label{eq:f15v60:selected-bounds}
\end{equation}
Thus the susceptibility has
limit superior at most \(9/2\);
equality of its limit with
\(9/2\) is not asserted.
\end{theorem}

\begin{proof}
Since \(\beta_M^\diamond\le\log M\),
we have \(a\to0\) and \(h_M\to0\).
Eventually
\(h_M\le1/4\) and \(2h_M\le R_M\).
The minimizer set is a
nonempty compact subset of
the stated interval, so its
smallest point is defined.
The ratio assertion follows
from \(|s_M|\le h_M\), and
the original lower bound on
\(\beta_M^\diamond\) implies
\(\widetilde\beta_M\to\infty\).
By the chain rule,
\[
 \chi_{M,\mathrm{vac}}(\beta(1-s))
                         =(1-s)^2\phi''(s).
\]
Apply Proposition~\ref{f15v60:triangular}
with \(a=h_M^4\).
Since
\(-\log(1-h_M^2)/h_M^2\to1\)
and \(k_M\to9/2\), this
proves the curvature claims.

For completeness the vacuum
mean at the new coupling must
also be checked. Put
\[
 e_0=e_{M,M}(\beta),\qquad
 v_0=\phi'(0)=e_{M,\mathrm{vac}}(\beta),\qquad
 \varepsilon_M=17\sqrt{\beta/M}.
\]
V59 gives \(|v_0-e_0|\le\varepsilon_M\)
eventually. For any \(|s|\le h=h_M\),
convexity gives the tangent
bound
\(\phi(s)\ge\phi(0)+sv_0\)
and the derivative secants
\[
 \frac{\phi(s)-\phi(-2h)}{s+2h}
       \le\phi'(s)
       \le\frac{\phi(2h)-\phi(s)}{2h-s}.
\]
The original source estimate
and \(d(0)\le a\), \(d(\pm2h)\ge0\)
give
\[
 \phi(\pm2h)-\phi(0)
                       \le\pm2h e_0+kJ(\pm2h)+a .
\]
Both denominators in the
secants are at least \(h\).
Also \(|s(e_0-v_0)|\le h\varepsilon_M\).
As \(2h\le1/2\), V59's
bounds for \(J\) yield
\begin{equation}
 \sup_{|s|\le h}|\phi'(s)-e_0|
             \le 4kh+\frac ah+\varepsilon_M .
 \label{eq:f15v60:mean-window}
\end{equation}
At the new coupling the mean
is
\((1-s_M)\phi'(s_M)\), not
just \(\phi'(s_M)\).
Consequently
\[
 |e_{M,\mathrm{vac}}(\widetilde\beta_M)-e_0|
       \le(1+h_M)
              \left(4k_Mh_M+h_M^3+\varepsilon_M\right)
                       +h_M|e_0|\longrightarrow0.
\]
Here \(e_0\to9/2\) is the
inherited periodic selection
result. This proves the
last assertion.
\end{proof}

The new selection need not
inherit the original periodic
source window or its full
local Gaussian classification.
Even the upper coupling bound
is now
\(\widetilde\beta_M\le(1+h_M)\log M\),
not literally the original
\(\beta_M^\diamond\le\log M\).
The vacuum conclusions above
are established directly;
none is obtained by silently
identifying the two laws.

\subsection{The curvature is a full vacuum resolvent plus contact variance}

This subsection is at an
arbitrary fixed spatial
regulator and coupling
\(\gamma>0\). Write
\(\mathsf T=\mathsf T(\gamma)\),
\(\lambda=\lambda_0(\gamma)\),
\(\Omega=\Omega(\gamma)\),
\(\mathsf P=\mathsf T/\lambda\)
and \(Q=I-|\Omega\rangle\langle\Omega|\).
All inverses below are on
\(Q\mathcal H\).
Compactness and the simple
isolated top eigenvalue give
\(\|\mathsf P|_{Q\mathcal H}\|<1\)
at this fixed regulator.
No uniform bound on that
norm is assumed.

Use precisely the spatially
constant smearing from V58,
\(f=V_{\rm sp}^{-1/2}\).
The corresponding slab random
variable and insertion are
\[
 X=\frac{\gamma\mathcal A}{\sqrt{V_{\rm sp}}},
 \qquad
 \mathsf B=-\frac{\gamma}{\sqrt{V_{\rm sp}}}\mathsf T',
 \qquad
 \psi=\lambda^{-1}Q\mathsf B\Omega,
 \qquad
 \sigma^2=\Var_{\pi_{L,\gamma}}(X).
\]
Primes mean coupling derivatives.
The slab measure is exactly
\eqref{eq:f15v59:vacuum-slab};
the contact variance does not
replace a squared insertion
by a product of transfer
operators.

\begin{theorem}[Vacuum susceptibility identity]
\label{f15v60:resolvent}
At every such finite spatial
regulator,
\begin{equation}
 \chi_{L,\mathrm{vac}}(\gamma)
       =\sigma^2+
                    2\langle\psi,(I-\mathsf P)^{-1}\psi\rangle .
 \label{eq:f15v60:resolvent}
\end{equation}
Both terms on the right are
nonnegative. On the newly
selected sequence of
Theorem~\ref{f15v60:selection},
\begin{equation}
 \sigma_M^2+
             2\langle\psi_M,(I-\mathsf P_M)^{-1}\psi_M\rangle
                         \le\kappa_M .
 \label{eq:f15v60:uniform-resolvent}
\end{equation}
There is no thermal vacuum
probability in this estimate.
\end{theorem}

\begin{proof}
The differentiable normalized
eigenvector can be chosen
real. The derivative of its
norm gives
\(\langle\Omega,\Omega'\rangle=0\).
Differentiating
\(\mathsf T\Omega=\lambda\Omega\)
and projecting with \(Q\)
yields
\[
 \Omega'=(\lambda-\mathsf T)^{-1}Q\mathsf T'\Omega .
\]
Differentiating
\(\lambda'=\langle\Omega,\mathsf T'\Omega\rangle\)
once more therefore gives
\[
 \lambda''
   =\langle\Omega,\mathsf T''\Omega\rangle
    +2\langle Q\mathsf T'\Omega,
                    (\lambda-\mathsf T)^{-1}Q\mathsf T'\Omega\rangle .
\]
The bounded literal slab
action implies
\[
 \frac{\langle\Omega,\mathsf T''\Omega\rangle}{\lambda}
      =\mathbb E_{\pi_{L,\gamma}}\mathcal A^2,
 \qquad
 \frac{\lambda'}{\lambda}
      =-\mathbb E_{\pi_{L,\gamma}}\mathcal A .
\]
In
\[
 \gamma^2G_L''(\gamma)
    =\frac{\gamma^2}{V_{\rm sp}}
       \left(\frac{\lambda''}{\lambda}
                      -\left(\frac{\lambda'}{\lambda}\right)^2\right),
\]
these two terms give
\(\Var_{\pi_{L,\gamma}}(\gamma\mathcal A)/V_{\rm sp}
=\sigma^2\).
Finally
\((\lambda-\mathsf T)^{-1}
=\lambda^{-1}(I-\mathsf P)^{-1}\)
and the definition of \(\psi\)
give the remaining term in
\eqref{eq:f15v60:resolvent}.
Positivity follows from the
variance and from
\(0\le\mathsf P|_{Q\mathcal H}<I\).
Apply the already proved
curvature bound at
\(\gamma=\widetilde\beta_M\)
to conclude
\eqref{eq:f15v60:uniform-resolvent}.
\end{proof}

This is a vacuum calculation
at the selected coupling,
not a limit of V58's
finite-period estimate with
its probability factor
dropped. That distinction is
why the new coupling selection
was necessary for this
particular argument.

\subsection{What the new response bound says about vacuum time}

For precision, the stationary
vacuum cylinder law on \(n\)
consecutive slabs has density
\[
 \lambda^{-n}\Omega(U_0)\Omega(U_n)
    \prod_{j=0}^{n-1}e^{-\gamma\mathcal A(U_j,U_{j+1},g_j)}
      \prod_{j=0}^{n}dH(U_j)\prod_{j=0}^{n-1}dH(g_j).
\]
Its total mass is one, and
integrating either endpoint
gives the shorter cylinder
by \(\mathsf T\Omega=\lambda\Omega\).
In particular the one-slab
marginal is the measure used
above. For \(X_t\) the
constant spatial action probe
on slab \(t\), direct kernel
composition gives, for
integers \(t\ge1\),
\begin{equation}
 C_{\rm vac}(t):=\operatorname{Cov}(X_0,X_t)
                  =\langle\psi,\mathsf P^{\,t-1}\psi\rangle,
 \qquad C_{\rm vac}(0)=\sigma^2 .
 \label{eq:f15v60:time-correlation}
\end{equation}
Indeed the uncentered
two-insertion integral is
\(\lambda^{-t-1}
\langle\Omega,\mathsf B\mathsf T^{t-1}\mathsf B\Omega\rangle\).
Its component along \(\Omega\)
is exactly \((\mathbb EX_0)^2\).
Subtracting that component
proves the identity. Summing
the convergent geometric
series on \(Q\mathcal H\)
identifies
\(\chi_{\rm vac}=C_{\rm vac}(0)+2\sum_{t\ge1}C_{\rm vac}(t)\).

Let \(\eta_M\) be the
spectral measure of
\(H_M=-\log\mathsf P_M\)
for \(\psi_M\), still on the
vacuum-orthogonal subspace.
For every \(E>0\), the
selected response bound gives
\begin{equation}
 \begin{gathered}
 \int_{(0,\infty)}
          \frac{d\eta_M(u)}{1-e^{-u}}\le\frac{\kappa_M}{2},
 \qquad
 \eta_M((0,E])\le\frac{\kappa_M}{2}(1-e^{-E})
                                      \le\frac{\kappa_M E}{2},\\
 0\le C_{M,\rm vac}(t)\le\frac{\kappa_M}{2t}
                          \qquad(t=1,2,\ldots).
 \end{gathered}
 \label{eq:f15v60:spectral-consequences}
\end{equation}
The spectral weight bound
uses the minimum of
\((1-e^{-u})^{-1}\) on
\((0,E]\).
For the time estimate,
write the correlation integral
against \(d\eta_M(u)/(1-e^{-u})\).
For \(0\le r<1\),
\[
 t(1-r)r^{t-1}
       \le(1-r)\sum_{j=0}^{t-1}r^j=1-r^t\le1.
\]
Take \(r=e^{-u}\) and use
the resolvent bound. These
are bounds for the actual
vacuum probe; neither a
finite-time error nor an
inverse vacuum probability
appears in the proved estimate.

\subsection{The remaining mass-gap requirements}

The advance is a uniform
full resolvent estimate for
one actual vacuum probe,
on an explicitly new nearby
coupling sequence. It does
not prove a positive lower
bound for the support of
\(\eta_M\). The gapless
finite-susceptibility example
\eqref{eq:f15v57:gapless}
already establishes that
logical distinction and is
not repeated here.

The smearing is spatially
constant and normalized by
\(V_{\rm sp}^{-1/2}\).
It is not an arbitrary
local observable or a dense
vacuum-orthogonal family.
Its bound cannot be promoted
to all spatial profiles
merely by translation
invariance. A justified
spatial source inequality
for the vacuum, or another
uniform estimate at the
new couplings, would be
needed for that extension.
The inherited periodic
source estimate cannot be
silently reused at
\(\widetilde\beta_M\).

The spectral generator here
also uses one lattice time
step. No physical time
calibration, interacting
four-dimensional continuum
state, or positive physical
Hamiltonian gap is supplied.
The concrete next interface
is therefore the extension
from the constant action
probe to specified spatially
varying vacuum insertions,
while preserving the same
coupling and normalization.

\subsection{Exact diagnostics and source preservation}

The new diagnostic is
\begin{center}\small
\path{scripts/verify_ym_f15_v60_vacuum_response_selection.py}.
\end{center}
It reruns V59, V58 and V57.
There are 54 exact polynomial
triangular-integral and
minimum-selection checks and
five shrinking-window budget
checks using rational
logarithm enclosures.

Thirty-six literal Boltzmann
matrix jets retain the
entrywise action and its
square. Independent top
eigenvector differentiation
checks the perturbation
equation and the second
eigenvalue derivative.
The resulting 36 susceptibility
identities keep the contact
variance. There are also
288 vacuum correlation and
inverse-time checks, and
144 low-energy weight bounds
without a thermal probability
factor. All arithmetic is
exact rational arithmetic or
certified rational logarithm
enclosure. The matrix tests
are finite test models, not
numerical evidence for a
Wilson continuum gap.

The script has 5,715 bytes
and SHA--256
\begin{center}\scriptsize\ttfamily
644A9DAFB5C0E8CE1EFE5B5A2C7B831C65E16538260C3FAAEF13CE95701F15AF.
\end{center}
The V59 predecessor has
1,352,785 bytes and SHA--256
\begin{center}\scriptsize\ttfamily
78F3A2B9881A5B5C6268C38E6280F64B5E9FB611F670018E1B9482C2EE94B3BF.
\end{center}
Removing this final section
and reversing the title
increment recovers V59
byte for byte. V60 replaces
V59 as the single active
main source. The protected
companions, archives and
102 earlier diagnostics
are unchanged. Only
\path{.tex} files are kept
in \path{latex_notes};
all compilation products
go to
\begin{center}\small
\path{artifacts/ym_f15_v60_texcheck/}.
\end{center}

\begin{firewall}
V60 selects asymptotically
equivalent nearby couplings
where the vacuum action mean
tends to \(9/2\) and the
vacuum susceptibility has
limit superior at most \(9/2\).
It proves a uniform full
vacuum resolvent bound for
the constant spatial action
probe, with the contact
variance retained.
It does not establish that
bound at the old couplings,
for every physical probe,
or as a positive continuum
spectral gap. The full
Yang--Mills mass gap remains
unproved.
\end{firewall}

\section{V61: all-cell vacuum covariance control from rectangular reflection}
\label{f15v61:context}

V60 bounds the vacuum response
of the constant spatial action
probe at its new coupling
selection. Extending that
bound by translation invariance
alone would be invalid.
This section supplies the
missing comparison: the exact
cell-source reflection inequality
holds on rectangular even tori
at every positive coupling.
Its covariance bound survives
the infinite-time limit at
fixed spatial regulator.

The result controls arbitrary
finitely supported spacetime
cell-energy sources under the
actual vacuum law. It also
extends the full vacuum
resolvent bound to all spatial
profiles and gives pointwise
algebraic covariance decay.
The periodic source window
at the old selected coupling
is not reused at the new one.

The scalar reflection method
is inherited from f14 V98--V99,
not a new general chessboard
theorem. The new work is its
rectangular verification, a
justified fixed-regulator
derivative limit, and their
application to V60's vacuum
selection. No companion
checkpoint is due before V65.

\subsection{The exact source inequality on an even rectangle}

Let
\[
 \Lambda_{\boldsymbol\ell}
     =\prod_{i=0}^{3}(\mathbb Z/\ell_i\mathbb Z),
 \qquad \ell_i\ge4\ \hbox{even},\qquad
 V_{\boldsymbol\ell}=\prod_i\ell_i .
\]
Use the same bare compact
\(\SU(2)\) Wilson law at
coupling \(\gamma>0\), with
normalized product Haar.
Define \(T_x=\gamma E_x\)
by the 24-face, weight-\(1/4\)
cell formula
\eqref{f15v1:cell}, with
\(\gamma\) in place of
\(\beta\). Thus
\(\sum_xT_x=\gamma S\).
Put
\[
 \chi_{\boldsymbol\ell,\gamma}
           =\frac{\gamma^2}{V_{\boldsymbol\ell}}\Var(S).
\]

\begin{proposition}[Rectangular cell-source comparison]
\label{f15v61:rectangle}
For every real source array
\(b=(b_x)\),
\begin{equation}
 \log\mathbb E_{\boldsymbol\ell,\gamma}
              e^{\sum_xb_xT_x}
       \le\frac1{V_{\boldsymbol\ell}}\sum_x
             \log\frac{Z_{\boldsymbol\ell,(1-b_x)\gamma}}
                           {Z_{\boldsymbol\ell,\gamma}} .
 \label{eq:f15v61:chessboard}
\end{equation}
At every positive coupling
the centered covariance
operator satisfies
\begin{equation}
 0\le\Gamma_{\boldsymbol\ell,\gamma}
                    \le\chi_{\boldsymbol\ell,\gamma}I,
 \qquad
 \|\Gamma_{\boldsymbol\ell,\gamma}\|
                    =\chi_{\boldsymbol\ell,\gamma}.
 \label{eq:f15v61:finite-norm}
\end{equation}
No equality of the four
side lengths is required.
\end{proposition}

\begin{proof}
Fix an axis \(i\) and
opposite integer site planes
\(x_i=c,c+\ell_i/2\).
A unit plaquette cannot meet
both open halves. Decompose
the action into the positive
half, its reflection and the
plaquettes within the boundary
planes. Conditional on all
tangent boundary links, the
two remaining link integrals
factor and are exchanged by
a Haar-preserving reflection.
Consequently, for a function
\(A\) supported in one closed
half, the reflection form is
a positive boundary integral
of the squared modulus of
its conditional half integral.
This is exactly f14 V98's
site-plane proof; the other
three axis lengths never enter
its factorization.

A closed cell is contained
in one of the two halves.
Reflection sends its base
coordinate to
\(2c-x_i-1\pmod{\ell_i}\)
and permutes its 24 geometric
faces, preserving their
weights and scalar plaquette
traces. Every plaquette still
has four incident cells.
These facts justify using
the same cell source factor
after reflection, including
its shared boundary links.

Here is the normalization
and dissemination argument
with the rectangular volume
retained. For a scalar \(u\)
set
\[
 c(u)=\left(\mathbb E e^{u\sum_xT_x}\right)^{1/V_{\boldsymbol\ell}},
 \qquad
 {\cal D}(b)=
       \frac{\mathbb E e^{\sum_xb_xT_x}}{\prod_xc(b_x)}.
\]
Conditional reflection
Cauchy--Schwarz gives
\({\cal D}(b)^2\le
{\cal D}(b^+){\cal D}(b^-)\),
where \(b^\pm\) duplicate
one half. The two duplicated
arrays together contain
twice every original label
multiplicity, so the
normalizers cancel exactly.

Over the finite alphabet of
labels appearing in \(b\),
choose an array maximizing
\({\cal D}\), and among
maximizers minimize the
number of disagreeing
nearest-neighbor cell pairs.
If the array is nonconstant,
a site plane separates a
disagreeing pair. Write
\(n_+,n_-\) for the internal
disagreements in its two
halves and \(n_\times>0\)
for the crossing disagreements.
The reflected arrays have
\(2n_+\) and \(2n_-\)
disagreements: all crossing
pairs now agree. Both are
maximizers by the reflection
inequality and maximality,
but at least one has fewer
than \(n_++n_-+n_\times\)
disagreements, a contradiction.
The maximum is therefore
attained at a constant
array, where \({\cal D}=1\).
Using \(\sum_xT_x=\gamma S\)
proves \eqref{eq:f15v61:chessboard}.
The partition integrals
remain finite at real
\((1-b_x)\gamma\); positivity
is used for the original
reflection law only.

Center the inequality, put
\(b_x=t a_x\), and take its
second derivative at \(t=0\).
The two sides agree to
first order by translation
invariance. Their second
derivatives give
\[
 \mathbb E\left|\sum_xa_x(T_x-\mathbb ET_x)\right|^2
                  \le\chi_{\boldsymbol\ell,\gamma}\sum_x|a_x|^2 .
\]
First take real \(a\), then
apply the result separately
to the real and imaginary
parts for complex \(a\).
Every covariance row sum
equals
\(V_{\boldsymbol\ell}^{-1}
\Var(\sum_xT_x)=\chi_{\boldsymbol\ell,\gamma}\).
Thus the constant vector
attains the bound, proving
\eqref{eq:f15v61:finite-norm}.
\end{proof}

\subsection{A derivative limit at fixed spatial regulator}

Now fix an even spatial
side \(L\ge4\) and a coupling
\(\gamma>0\), while the
temporal period \(N\) tends
to infinity through even
integers. Write
\[
 V_{\rm sp}=L^3,\qquad
 F_{L,N}(\gamma)=\frac{\log Z_{L,N}(\gamma)}{NV_{\rm sp}},
 \qquad
 G_L(\gamma)=\frac{\log\lambda_0(\gamma)}{V_{\rm sp}} .
\]
The operator is V57's
physical Wilson transfer
operator at this spatial
regulator. The pressure
floor alone cannot justify
twice differentiating its
limit; V59 explicitly
showed why. Instead use
its isolated top eigenvalue
at this fixed regulator.

\begin{proposition}[Fixed-regulator susceptibility and cylinder limits]
\label{f15v61:limit}
At fixed \(L,\gamma\),
\begin{equation}
 F_{L,N}^{(j)}(\gamma)\longrightarrow G_L^{(j)}(\gamma)
       \quad(j=0,1,2),\qquad
 \chi_{L,N,\gamma}\longrightarrow
                  \chi_{L,\mathrm{vac}}(\gamma).
 \label{eq:f15v61:derivative-limit}
\end{equation}
Here
\(\chi_{L,N,\gamma}=\gamma^2F_{L,N}''(\gamma)\).
Every fixed bounded
gauge-invariant cylinder
observable has its periodic
expectation converging to
the corresponding vacuum
cylinder expectation of V60.
In particular all fixed
cell-energy means and
covariances converge.
No convergence rate uniform
in \(L,\gamma\) is asserted.
\end{proposition}

\begin{proof}
Complexify the coupling
locally near the real
point \(\gamma\).
The bounded slab action
makes \(\mathsf T(z)\)
analytic in Hilbert--Schmidt
norm. The simple isolated
top eigenvalue has an
analytic eigenvalue
\(\lambda(z)\ne0\) and an
analytic rank-one Riesz
projection \(\Pi(z)\).
At the real point
\(\Pi(\gamma)=|\Omega\rangle\langle\Omega|\).
Set
\[
 D(z)=\frac{\mathsf T(z)-\lambda(z)\Pi(z)}{\lambda(z)}.
\]
The spectral projection
commutes with \(\mathsf T(z)\),
and \(\Pi(z)D(z)=D(z)\Pi(z)=0\).
At the real point,
\(\|D(\gamma)\|<1\).
On a sufficiently small
closed complex disk there
are constants \(q<1\) and
\(C<\infty\) such that
\[
                  \|D(z)\|\le q,\qquad \|D(z)\|_2^2\le C .
\]
The rank-one analytic
projection is locally
Hilbert--Schmidt bounded;
its difference from the
projection at the center
has rank at most two.
These bounds depend on
the fixed regulator.

For \(N\ge2\), trace-ideal
multiplication gives
\begin{equation}
 Z_{L,N}(z)=\lambda(z)^N
                   \left[1+\operatorname{Tr}D(z)^N\right],
 \qquad
 |\operatorname{Tr}D(z)^N|\le Cq^{N-2}.
 \label{eq:f15v61:analytic-remainder}
\end{equation}
The trace partition identity
extends to this disk by
its kernel integral, or
analytic continuation.
For large \(N\), the bracket
stays within distance \(1/2\)
of one, so its logarithm
has the analytic branch
agreeing with the real
partition logarithm.
On a smaller disk, Cauchy's
derivative estimates give
\[
 \left|\partial_z^j
       \left(\frac{\log Z_{L,N}(z)}{N}
                               -\log\lambda(z)\right)
       \bigg|_{z=\gamma}\right|
                \le \frac{C_jq^{N-2}}{N}\quad(j=0,1,2).
\]
Division by \(V_{\rm sp}\)
proves the pressure and
susceptibility limits.

For the cylinder assertion,
take an observable \(O\)
supported on \(m\) consecutive
slabs, with \(m\) fixed.
Integrating those slabs
with \(O\) inserted defines
a kernel \(K_O\).
Its absolute value is at
most \(\|O\|_\infty\) times
the kernel of \(\mathsf T^m\),
so \(K_O\) is
Hilbert--Schmidt at the
fixed regulator. Gauge
invariance gives the
physical endpoint operator.
For \(N>m\) its periodic
expectation is
\[
 \frac{\operatorname{Tr}(K_O\mathsf T^{N-m})}
           {\operatorname{Tr}\mathsf T^N}
 =\frac{\operatorname{Tr}
       \bigl((K_O/\lambda^m)(\mathsf T/\lambda)^{N-m}\bigr)}
         {\operatorname{Tr}(\mathsf T/\lambda)^N}.
\]
At the real point
\((\mathsf T/\lambda)^n=\Pi+D^n\);
the trace remainder tends
to zero by
\(\|D^n\|_2\le\|D\|_2q^{n-1}\).
The limit is
\(\lambda^{-m}\langle\Omega,K_O\Omega\rangle\),
exactly the vacuum cylinder
integral. Cell energies
and their finite products
are bounded at fixed
\(\gamma\), so this applies
to their required moments.
\end{proof}

The limit in this proposition
is taken with \(L,\gamma\)
fixed. Only after obtaining
the vacuum inequality do we
substitute \(L=M\) and
\(\gamma=\widetilde\beta_M\).
No joint-period convergence
rate and no uniform thermal
vacuum probability are needed.

\subsection{Sharp vacuum covariance norm for arbitrary cell sources}

Let
\(\mathcal C_L=(\mathbb Z/L\mathbb Z)^3\times\mathbb Z\)
denote the infinite-time
cylinder of cells. Expectations
below use the consistent
vacuum cylinder laws, and
\(T_x\) uses the same
coupling \(\gamma\) as that
law.

\begin{theorem}[The vacuum cell covariance operator]
\label{f15v61:vacuum-norm}
For every finitely supported
real or complex array \(a\)
on \(\mathcal C_L\),
\begin{equation}
 \mathbb E_{\rm vac}
       \left|\sum_xa_x(T_x-\mathbb E_{\rm vac}T_x)\right|^2
          \le\chi_{L,\mathrm{vac}}(\gamma)\sum_x|a_x|^2 .
 \label{eq:f15v61:vacuum-source}
\end{equation}
The covariance form extends
uniquely to a bounded
positive operator on
\(\ell^2(\mathcal C_L)\), with
\begin{equation}
            \|\Gamma_{L,\mathrm{vac}}(\gamma)\|
                       =\chi_{L,\mathrm{vac}}(\gamma).
 \label{eq:f15v61:vacuum-exact-norm}
\end{equation}
In particular on V60's
nearby selected sequence,
\begin{equation}
         0\le\Gamma_{M,\mathrm{vac}}(\widetilde\beta_M)
                           \le\kappa_M I,\qquad
                           \kappa_M\longrightarrow\frac92 .
 \label{eq:f15v61:selected-norm}
\end{equation}
\end{theorem}

\begin{proof}
Embed the fixed temporal
support of \(a\) without
wraparound in an even
period \(N\) larger than
that support. Apply
\eqref{eq:f15v61:finite-norm}
on the rectangle
\((N,L,L,L)\).
The left side converges
by the cylinder moment
statement of
Proposition~\ref{f15v61:limit};
its right side converges
by the susceptibility
statement. This proves
\eqref{eq:f15v61:vacuum-source}.
Density of finitely
supported arrays and
polarization give the
bounded positive extension.

To prove sharpness, let
\(f=V_{\rm sp}^{-1/2}\)
be constant on a spatial
slice, and put
\(X_t=\sum_y fT_{(t,y)}\).
For \(R\ge1\) use the
unit source array
\[
 a^{(R)}_{(t,y)}
      =(RV_{\rm sp})^{-1/2}{\bf1}_{\{0,\ldots,R-1\}}(t).
\]
V60 proves that its
constant-profile vacuum
covariances \(C_f(t)\)
are nonnegative and that
\(C_f(0)+2\sum_{t\ge1}C_f(t)
=\chi_{L,\mathrm{vac}}(\gamma)\).
Consequently
\[
 \operatorname{Var}\left(\frac1{\sqrt R}
                              \sum_{t=0}^{R-1}X_t\right)
   =C_f(0)+2\sum_{t=1}^{R-1}
                  \left(1-\frac tR\right)C_f(t)
        \longrightarrow\chi_{L,\mathrm{vac}}(\gamma).
\]
The increasing weights and
the convergent nonnegative
series justify the limit.
These unit arrays approach
the upper bound, proving
\eqref{eq:f15v61:vacuum-exact-norm}.
Finally use V60's proved
bound on \(\chi_{\rm vac}\)
at \(\widetilde\beta_M\).
\end{proof}

This is a statement about
the infinite-time vacuum
law, not merely a formal
Hessian of an approximate
vacuum pressure. The sharp
norm is a zero-frequency
limit of unit finite sources;
the constant spacetime
array itself is not in
\(\ell^2(\mathcal C_L)\).

\subsection{All spatial action profiles have a full vacuum response bound}

For a real spatial profile
\(f\), use V57's literal
symmetric slab observable
\(F_0^f=\sum_y f_yT_{(0,y)}\)
and insertion \(\mathsf B_f\)
at coupling \(\gamma\).
Put
\[
 \psi_f=\lambda_0^{-1}Q\mathsf B_f\Omega,
 \qquad
 \sigma_f^2=\Var_{\rm vac}(F_0^f),
 \qquad
 \mathsf P=\mathsf T/\lambda_0 .
\]
The covariance for separated
slabs is, by the same
kernel composition as V60,
\[
 C_f(t)=\langle\psi_f,\mathsf P^{\,t-1}\psi_f\rangle
                       \quad(t\ge1),\qquad C_f(0)=\sigma_f^2.
\]
The contact term is the
variance of the literal
random slab observable,
not the square of the
insertion operator.

Applying
\eqref{eq:f15v61:vacuum-source}
to \(a_{(t,y)}=R^{-1/2}f_y\)
for \(0\le t<R\), and
letting \(R\to\infty\), gives
\begin{equation}
 \sigma_f^2+
        2\langle\psi_f,(I-\mathsf P)^{-1}\psi_f\rangle
        \le\chi_{L,\mathrm{vac}}(\gamma)\|f\|_2^2 .
 \label{eq:f15v61:profile-resolvent}
\end{equation}
Indeed the nonnegative
covariance series converges
at fixed regulator because
\(\|\mathsf P|_{Q\mathcal H}\|<1\);
the same finite-box variance
identity applies. This
establishes the bound for
every \(f\) at the \emph{same}
coupling, without making a
separate selection for each
profile.

On V60's sequence, replace
\(\chi_{L,\mathrm{vac}}\) by
\(\kappa_M\) in
\eqref{eq:f15v61:profile-resolvent}.
Thus, for the spectral
measure \(\eta_f\) of
\(-\log\mathsf P_M\) and
the vector \(\psi_f\),
\begin{equation}
 \eta_f((0,E])
       \le\frac{\kappa_M}{2}(1-e^{-E})\|f\|_2^2,
 \qquad
 0\le C_f(t)\le\frac{\kappa_M}{2t}\|f\|_2^2
          \quad(E>0,\ t\ge1).
 \label{eq:f15v61:profile-spectral}
\end{equation}
These follow exactly as in
V60 from the full resolvent.
Their new content is the
arbitrary spatial profile,
not a new scalar spectral
inequality.

\subsection{Pointwise spatial and temporal covariance decay}

For \(0\le d\le L/2\) define
\[
 w_L(0)=0,\qquad w_L(d)=\min(2d,L-1)\quad(d\ge1).
\]
For two cells with time
difference \(t\) and spatial
periodic coordinate distances
\(d_1,d_2,d_3\), set
\[
 W_L(t,\mathbf d)=
       \max\{2|t|,w_L(d_1),w_L(d_2),w_L(d_3)\}.
\]

\begin{theorem}[Actual vacuum algebraic covariance bound]
\label{f15v61:pointwise}
At fixed even \(L\) and
\(\gamma>0\),
\begin{equation}
 \left|\operatorname{Cov}_{\rm vac}(T_x,T_y)\right|
       \le
       \frac{\chi_{L,\mathrm{vac}}(\gamma)}
                         {\max\{1,W_L(t,\mathbf d)\}} .
 \label{eq:f15v61:pointwise}
\end{equation}
At \(\gamma=\widetilde\beta_M\),
\(L=M\), the numerator may
be replaced by \(\kappa_M\).
\end{theorem}

\begin{proof}
First work on a finite
even rectangle and select
any coordinate as transfer
time, of length \(\ell_i\).
The other three side
lengths may be unequal.
V57's cell-to-slab identity,
positive transfer construction
and thermal spectral
decomposition require only
the local cell faces,
product Haar and periodicity;
none requires equality of
these transverse lengths.
For a real transverse
profile \(b\), those formulas
therefore give a covariance
\(c_b(d)\) which is
nonnegative at noncontact
times and decreasing for
\(1\le d\le\ell_i/2\).
Explicitly, its spectral
terms are a nonnegative
stationary quadratic form
and positive multiples of
\[
 r^{d-1}+r^{\ell_i-d-1},
                  \qquad 0<r<1.
\]
This also shows that the
polarized transverse covariance
at each noncontact \(d\)
is a positive semidefinite
form.

Apply the rectangular
covariance norm to the
source \(b\) repeated on
all \(\ell_i\) slices.
Translation invariance gives
\[
 \sum_{u=0}^{\ell_i-1}c_b(u)
                      \le\chi_{\boldsymbol\ell,\gamma}\|b\|_2^2 .
\]
There are
\(w_{\ell_i}(d)=\min(2d,\ell_i-1)\)
noncontact times whose
periodic distance is at
most \(d\). Each has
covariance at least
\(c_b(d)\), and the
contact variance is
nonnegative. Hence
\[
 c_b(d)\le
       \frac{\chi_{\boldsymbol\ell,\gamma}}
                   {w_{\ell_i}(d)}\|b\|_2^2 .
\]
Cauchy--Schwarz in the
positive transverse form,
with delta functions at
the two transverse cell
locations, gives the same
bound in absolute value
for the two individual
cell energies at axis
separation \(d\).

Now take the original
temporal length \(N\to\infty\)
at fixed \(L,\gamma\).
The moment and susceptibility
limits were proved above.
For a fixed nonzero temporal
separation,
\(w_N(|t|)\to2|t|\);
for a spatial axis the
factor remains \(w_L(d_i)\).
Choose the strongest of
these axis bounds. For
coincident cells use
\eqref{eq:f15v61:vacuum-source}
with a delta source.
This proves
\eqref{eq:f15v61:pointwise}.
Finally apply the same
nearby coupling selection
as in the preceding
theorems.
\end{proof}

This proof rotates axes
only on finite rectangles.
It does not assume that
a finite-spatial,
infinite-time cylinder is
invariant under exchange
of its infinite and finite
axes. The order of limits
is the one already
specified, and the
susceptibility is taken
at the actual coupling.

\subsection{What has and has not moved toward the mass gap}

The previous constant-profile
restriction has been removed
for the specified cell-energy
field. At the new selected
couplings its actual vacuum
covariance has a uniformly
bounded spacetime operator
norm, its arbitrary spatial
action profiles obey the
full resolvent bound, and
its separated cell covariances
satisfy the displayed
algebraic estimate.

None of these statements
excludes spectral support
arbitrarily close to zero.
The already recorded
gapless example
\eqref{eq:f15v57:gapless}
still applies to that
logical inference.
The field of scalar cell
energies is not proved to
generate a dense set of
all vacuum-orthogonal
physical states. Moreover
the estimates are not
exponential decay with a
calibrated physical rate.
There is still no
interacting physical
continuum construction
or full Yang--Mills
mass-gap theorem.

The next substantive issue
is whether additional
observable structure or a
stronger source estimate
can turn this energy-field
response control into
physical clustering.
Repeating a low-energy
weight bound for the same
probes would not settle
the missing spectral
support condition.

\subsection{Exact finite checks and preservation}

The new diagnostic is
\begin{center}\small
\path{scripts/verify_ym_f15_v61_rectangular_vacuum_covariance.py}.
\end{center}
It reruns V60--V57.
On two unequal rectangular
tori it checks 9,216
plaquette incidence counts,
12,288 cell-face reflection
identities, and 73,728
plaquette support classifications.
Sixteen label tests retain
the source normalizers and
the disagreement counts
across reflection seams.

Independent product jets
and cyclic trace formulas
give 20 exact first- and
second-pressure derivative
identities. Twenty-four
finite spectral checks
verify the long-time-box
saturation remainder;
30 counting checks verify
the rectangular-axis
noncontact multiplicity.
The analytic derivative
limit, reflection inequality
for all source arrays,
and continuum claim boundaries
are proved in the text,
not inferred from these
finite tests.

The diagnostic has 7,360
bytes and SHA--256
\begin{center}\scriptsize\ttfamily
C735521B6DFAA5F0FF8E0A3E6BD6DCB3740C895581291C6575CE6285814BC207.
\end{center}
The V60 predecessor has
1,369,710 bytes and SHA--256
\begin{center}\scriptsize\ttfamily
CF01CD99D16C0BDB65867FE37096F5BE4989F9A15820B8A91E5C177B6FA2476B.
\end{center}
Removing this final section
and reversing the title
increment recovers V60
byte for byte. V61 replaces
V60 as the sole active
main source. The protected
inputs and 103 earlier
diagnostics remain unchanged.
Only \path{.tex} files
are kept in
\path{latex_notes}; all
compilation products go to
\begin{center}\small
\path{artifacts/ym_f15_v61_texcheck/}.
\end{center}

\begin{firewall}
V61 extends V60's vacuum
response control to arbitrary
finitely supported sources
of the original cell-energy
field, at the same newly
selected couplings.
It proves a sharp vacuum
covariance norm and
algebraic pointwise decay.
It does not prove a dense
physical probe family,
exponential clustering,
an interacting continuum
theory or a positive
Yang--Mills mass gap.
\end{firewall}

\section{V62: vacuum source moments, local curvature classification, and susceptibility saturation}
\label{f15v62:context}

V61 controls the vacuum
covariance of cell energies
but does not by itself
identify the vacuum field
law. This section supplies
nonzero-source bounds and
uniform local moments, then
checks the local hypotheses
needed to apply the existing
curvature classification to
the vacuum. The resulting
limit is a microscopic
Gaussian lattice-curvature
law. It is not the
interacting physical
continuum theory.

Throughout the selected
statements use V60's
\(\gamma_M=\widetilde\beta_M\),
not the original periodic
\(\beta_M^\diamond\).
The square-filling algebra,
edge-avoiding differentiation
and abstract classification
from f14 V90 and V97 are
reused at their stated
scopes. The new assertions
are their verified vacuum
application and the exact
vacuum susceptibility limit.
No companion checkpoint
is due before V65.

\subsection{An exact nonzero-source inequality for the vacuum}

Fix an even spatial side
\(L\ge4\) and a coupling
\(\gamma>0\). Write
\(G_L(\gamma)=V_{\rm sp}^{-1}\log\lambda_0(\gamma)\)
as before, and
\(e_{\rm vac}=-\gamma G_L'(\gamma)\).
Translation invariance and
the exact slab action
partition give
\(\mathbb E_{\rm vac}T_x=e_{\rm vac}\)
for every cell. Define
\begin{equation}
 {\cal J}_{L,\gamma}(b)
       =G_L(\gamma(1-b))-G_L(\gamma)-b e_{\rm vac},
                      \qquad b<1.
 \label{eq:f15v62:vacuum-J}
\end{equation}
This function is nonnegative
by convexity of \(G_L\).
It is not being identified
with the cumulant generating
function of one cell.

\begin{theorem}[Finite-source vacuum comparison]
\label{f15v62:source}
For every finite real cell
source array \(b\) with
\(\max_x b_x<1\),
\begin{equation}
 \log\mathbb E_{\rm vac}
       \exp\left\{\sum_xb_x(T_x-e_{\rm vac})\right\}
                    \le\sum_x{\cal J}_{L,\gamma}(b_x).
 \label{eq:f15v62:source}
\end{equation}
Only the nonzero source
coordinates need be included
in the sum.
\end{theorem}

\begin{proof}
Place the finite source
support in an even temporal
period \(N\) without
wraparound. The rectangular
chessboard inequality
\eqref{eq:f15v61:chessboard}
bounds the uncentered
logarithm by
\[
 \sum_x\left[
       F_{L,N}(\gamma(1-b_x))-F_{L,N}(\gamma)\right].
\]
For each of the finitely
many source values the
shifted coupling is positive.
The fixed-regulator pressure
limit therefore gives the
corresponding difference of
\(G_L\)'s. The exponential
of the finite source is a
bounded cylinder observable
at this fixed \(L,\gamma\),
so its periodic expectation
converges to the vacuum
expectation by V61. Subtract
\(e_{\rm vac}\sum_xb_x\).
This proves the result,
with no interchange of a
growing support and the
infinite-time limit.
\end{proof}

\subsection{An explicit source window at the new couplings}

Retain V60's notation
\[
 \beta=\beta_M^\diamond,\quad
 a_M=\frac{12\beta}{M},\quad
 h_M=a_M^{1/4},\quad
 \gamma_M=\beta(1-s_M),\quad |s_M|\le h_M,
 \quad \phi_M(s)=G_M(\beta(1-s)).
\]
Also set
\[
 \varepsilon_M=17\sqrt{\beta/M},
 \qquad
 D_M=4k_Mh_M+\frac{a_M}{h_M}+\varepsilon_M .
\]
V59--V60 give, eventually,
\[
 \begin{gathered}
 |\phi_M'(0)-e_0|\le\varepsilon_M,\qquad
 |\phi_M'(s_M)-e_0|\le D_M,\qquad
 e_0=e_{M,M}(\beta)\longrightarrow9/2,\\
 \phi_M(t)\le\phi_M(0)+t e_0+k_MJ(t)+a_M
                           \quad(|t|\le R_M),\\
 J(t)=-\log(1-t)-t .
 \end{gathered}
\]
Here \(D_M\to0\).

Fix \(0<r<1\), and put
\begin{equation}
 \begin{gathered}
 r'=(1+r)/2,\qquad L_r=\frac{r'}{1-r'},\\
 {\cal E}_M(r)=
       k_ML_rh_M(1+r)+a_M
             +(1+h_M)rD_M+h_M\varepsilon_M .
 \end{gathered}
 \label{eq:f15v62:source-error}
\end{equation}
For this fixed \(r\),
\({\cal E}_M(r)\to0\).

\begin{theorem}[Recentered vacuum source window]
\label{f15v62:window}
For all sufficiently large
\(M\), every finite array
with \(\|b\|_\infty\le r\)
satisfies, under the vacuum
law at \(\gamma_M\),
\begin{equation}
 \log\mathbb E_{\rm vac}
       e^{\sum_xb_x(T_x-e_{\rm vac})}
       \le k_M\sum_xJ(b_x)
                    +{\cal E}_M(r)\,|\operatorname{supp}b| .
 \label{eq:f15v62:window}
\end{equation}
There is no support-size
restriction on this finite
inequality. Its additive
error, however, is paid
per nonzero source
coordinate.
\end{theorem}

\begin{proof}
Abbreviate \(s=s_M\), and
for a scalar \(|b|\le r\)
put \(t=s+(1-s)b\).
Then
\[
 \gamma_M(1-b)=\beta(1-t),\qquad
 {\cal J}_{M,\gamma_M}(b)
      =\phi_M(t)-\phi_M(s)-(t-s)\phi_M'(s).
\]
Eventually \(|t|\le r'\le R_M\).
Convexity gives
\(\phi_M(s)\ge\phi_M(0)+s\phi_M'(0)\).
Using the previous upper
bound at \(t\) therefore
gives
\[
 {\cal J}_{M,\gamma_M}(b)
 \le k_MJ(t)+a_M
      +(t-s)(e_0-\phi_M'(s))
                   +s(e_0-\phi_M'(0)).
\]
The last two terms are
at most
\((1+h_M)rD_M+h_M\varepsilon_M\).
On \([-r',r']\),
\(|J'|\le L_r\), and
\[
                 |t-b|=|s(1-b)|\le h_M(1+r).
\]
Thus
\[
 {\cal J}_{M,\gamma_M}(b)
                        \le k_MJ(b)+{\cal E}_M(r).
\]
Apply Theorem~\ref{f15v62:source}
coordinate by coordinate.
At \(b=0\) the exact
function \({\cal J}\) is
zero, so no error need
be assigned there.
\end{proof}

For a fixed source support
the error tends to zero.
For growing supports that
requires a separate bound
on
\({\cal E}_M(r)|\operatorname{supp}b|\).
The estimate does not give
a quadratic error down to
zero source, nor a
uniform response bound
throughout a tilted-law
neighborhood. V61's
zero-source Hessian result
was proved separately.

\begin{proposition}[Uniform vacuum thermal moments]
\label{f15v62:moments}
For all sufficiently large
\(M\), every cell and every
plaquette satisfy
\begin{equation}
 \mathbb E_{\rm vac}e^{T_x/4}\le e^3,\qquad
 \mathbb E_{\rm vac}e^{X_p/16}\le e^3,\qquad
                         X_p=\gamma_M s_p.
 \label{eq:f15v62:exponential}
\end{equation}
For every integer \(q\ge0\),
\begin{equation}
 \mathbb E_{\rm vac}X_p^q\le e^3\,16^q q!,
 \qquad
 \mathbb E_{\rm vac}|Y_p^\rho|^{2q}\le e^3\,32^q q!,
 \label{eq:f15v62:moments}
\end{equation}
where \(Y^\rho\) is any
fixed-path root-transported
thermal curvature.
Every fixed polynomial
in finitely many curvature
and plaquette-energy
coordinates is uniformly
integrable, including all
its fixed positive powers.
\end{proposition}

\begin{proof}
Use the preceding theorem
for one source \(b=1/4\).
Eventually \(e_{\rm vac}\le5\),
\(k_M\le5\), and
\({\cal E}_M(1/4)\le1\).
Since \(J(1/4)\le1/24\),
\[
 \log\mathbb E_{\rm vac}e^{T_x/4}
       \le\frac54+\frac5{24}+1=\frac{59}{24}<3.
\]
If \(p\) is a face of
cell \(x\), positivity of
all face energies gives
\(X_p\le4T_x\). Hence
\(X_p/16\le T_x/4\), proving
the second exponential bound.
The elementary inequality
\(u^q\le q!e^u\) for \(u\ge0\)
gives the first moment
bound. For unit-quaternion
holonomy the exact identity
\[
       |Y_p^\rho|^2=2X_p-X_p^2/\gamma_M\le2X_p
\]
gives the second. Finally
Hölder's inequality and
these arbitrarily high
single-coordinate moments
bound every fixed power
of every finite polynomial.
This proves uniform
integrability without
independence of coordinates.
\end{proof}

\subsection{Local Haar identities for the vacuum link law}

The curvature coordinates
are gauge covariant at
their common root.
Accordingly one must
justify local limits for
link observables, not just
for gauge-invariant cell
energies.

V61's cylinder convergence
extends to every bounded
link cylinder observable.
Indeed use the full slice
\(L^2\) space in its proof.
The Wilson kernel annihilates
the orthogonal complement
of the physical subspace;
its top eigenvector is
the same physical \(\Omega\).
The insertion kernel \(K_O\)
need not preserve that
subspace, but is still
Hilbert--Schmidt and satisfies
\[
 \frac{\operatorname{Tr}_{\rm full}
                      (K_O\mathsf T^{N-m})}
           {\operatorname{Tr}_{\rm full}\mathsf T^N}
         \longrightarrow
                  \lambda_0^{-m}\langle\Omega,K_O\Omega\rangle .
\]
The same remainder estimate
applies on the full space;
its extra transfer eigenvalues
are zero. This is exactly
the vacuum link cylinder
law, with no gauge fixing.

In particular that law
inherits discrete translations
and all finite vertex gauge
transformations from the
periodic laws. For a fixed
edge and a smooth cylinder
test, the periodic Haar
differentiation identity
also passes to the vacuum:
at fixed \(L,\gamma\), its
test, derivative and the
local action derivative
are bounded cylinder
functions.

Here is the precise
version needed for
curvature. For each lifted
vertex \(x\in\mathbb Z^4\),
choose a fixed finite
root path \(\rho_x\).
Define
\begin{equation}
 \begin{gathered}
 g_x=U_{\rho_x},\qquad
 k_e=g_xU_e g_{x+\hat\mu}^{-1}
                        \quad(e=(x,\mu)),\\
 Y_p^\rho=\sqrt{\gamma_M}\,
            \operatorname{Im}(g_xU_pg_x^{-1}),\qquad
 X_p=\gamma_M(1-\tfrac12\operatorname{Tr}U_p).
 \end{gathered}
 \label{eq:f15v62:root-fields}
\end{equation}
The paths and all tested
coordinates are fixed
before \(M\to\infty\).
For large \(M\) there is
no spatial aliasing on
any of these fixed patches.

Let \(A(C)\) be the number
of square swaps in a fixed
filling of a closed lattice
word \(C\), and let \(A_e\)
be such a count for the
closed rooted word defining
\(k_e\). Because
\(X_p\le4T_x\) and
\(\mathbb ET_x=e_{\rm vac}\le5\),
one may use \(K=20\) in
the local estimates below.
The fixed-word arguments
of f14 V97 require only
\(\sup_p\mathbb EX_p\le K\)
and unit-quaternion algebra.
They therefore give here
\begin{equation}
 \begin{gathered}
 \mathbb E|U_C-I|^2\le\frac{2K A(C)^2}{\gamma_M},
 \qquad
 \mathbb E|k_e-I|^2\le\frac{2K A_e^2}{\gamma_M},\\
 \mathbb E|Y_p^{\widehat\rho}-Y_p^\rho|
     \le\frac{4K A(\widehat\rho_x\rho_x^{-1})}{\sqrt{\gamma_M}} .
 \end{gathered}
 \label{eq:f15v62:transport}
\end{equation}
For clarity, a filling
writes \(U_C\) as a
product of conjugated
plaquette holonomies;
telescoping their chords
and Cauchy--Schwarz give
the first bound. The
last uses
\(|\operatorname{Ad}_H v-v|\le2|H-I||v|\).
No path-length bound
uniform over growing
physical paths is asserted.

For a fixed plaquette,
writing \(a_e=\sqrt{\gamma_M}\operatorname{Im}k_e\),
the exact four-factor
product expansion gives
\begin{equation}
 \mathbb E|Y_p^\rho-(d_1a)_p|
       \le\frac{2K}{\sqrt{\gamma_M}}
                \sum_{i<j}A_{e_i}A_{e_j}.
 \label{eq:f15v62:linearization}
\end{equation}
Here the \(e_i\)'s are
its four boundary edges,
with the prescribed inverse
orientation convention.
This follows by bounding
the product remainder by
\(\sum_{i<j}|k_{e_i}-I||k_{e_j}-I|\).
Applying \(d_2\) on a
fixed cube gives
\(d_2Y^\rho\to0\) in
\(L^1\), since \(d_2d_1=0\).

Now fix an edge \(e\)
and choose all required
paths to avoid that
unoriented edge, as in
f14 V97. The complement
of one edge in
\(\mathbb Z^4\) is connected,
so these finite paths
exist. Conditional on all
other links, use the
Haar flow
\[
 U_e(t)=g_x^{-1}
             \exp(t\,\iota(h)/\sqrt{\gamma_M})g_xU_e .
\]
The frame \(g_x\) is
independent of \(U_e\).
Let \(S_e\) be the sum
of the six plaquette
deficits incident to \(e\).
The vacuum identity is
\begin{equation}
             \mathbb E_{\rm vac}D_e f
                  =\mathbb E_{\rm vac}f\,D_e(\gamma_M S_e).
 \label{eq:f15v62:local-Haar}
\end{equation}
This is the passed finite
Haar identity described
above. It differentiates
the finite local action
\(S_e\), not a divergent
infinite-cylinder total
action. There is no
derivative of a link-dependent
frame to omit.

For a bounded \(C^1\)
curvature cylinder test
with continuous bounded
gradient, write
\(B_f=\|f\|_\infty\) and
\(L_f=\max_p\sup|\nabla_{Y_p}f|\).
For \(h_e\) supported at
the varied edge, put
\(\ell=d_1h_e\).
The signed prefix calculation
of f14 V97 and
\eqref{eq:f15v62:transport}
then yield
\begin{equation}
 \left|\mathbb E_{\rm vac}D_\ell f(Y^\rho)
       -\mathbb E_{\rm vac}f(Y^\rho)\langle Y^\rho,\ell\rangle\right|
 \le\frac{|h|A_\star}{\sqrt{\gamma_M}}
                         \left(3L_f\sqrt{2K}+4B_fK\right).
 \label{eq:f15v62:Stein-error}
\end{equation}
Here
\(A_\star=\sum_{p\ni e}\sum_{a\in\partial p}A_a\)
retains edge multiplicities.
Explicitly, the varied
plaquette has derivative
\(\iota(v_{p,e})\widetilde P_p/\sqrt{\gamma_M}\),
where
\(v_{p,e}=\sigma_{p,e}\operatorname{Ad}_{C_{p,e}}h\).
For negative incidence the
prefix includes \(k_e^{-1}\).
Its error from
\(\sigma_{p,e}h\) is at
most \(2|h|\sum_{\partial p}|k_a-I|\);
the curvature derivative
error is at most three
times that chord sum
with the corresponding
factor \(|h|\).
Cauchy--Schwarz and
\eqref{eq:f15v62:transport}
give exactly
\eqref{eq:f15v62:Stein-error}.
Thus no additional
vacuum drift has been
inserted or discarded.

\subsection{Verified classification of the microscopic vacuum law}

We now use the abstract
stationary-law theorem
proved in f14 V90 at
\texttt{thm:v90-classification}.
It states that a stationary,
closed random colour two-form
with uniform second moments,
common colour invariance
and the measure translation
identity in the directions
\(d_1d_1^*B\) is a Gaussian
closed field plus an
independent random two-form
constant over \(\mathbb Z^4\).
The theorem and its
heat-average proof were
reread for this application.
They do not assume a
particular periodic or
restricted-chart measure.

\begin{theorem}[Selected vacuum microscopic limit]
\label{f15v62:classification}
Under the actual vacuum
laws at \(L=M\) and
\(\gamma_M=\widetilde\beta_M\),
the joint fields
\((Y^\rho,X)\) converge
in the countable product
topology to
\begin{equation}
             (Y^\rho,X)
                      \Longrightarrow (G,|G|^2/2).
 \label{eq:f15v62:Gaussian-limit}
\end{equation}
All fixed mixed polynomial
moments converge. The law
is independent of the
fixed root-path choices.
The centered Gaussian
field has covariance
\(\mathcal P\otimes I_3\),
where
\begin{equation}
 \mathcal P(k)=
       \frac{d_1(k)d_1(k)^*}{\omega(k)},
 \qquad
 \omega(k)=\sum_{i=0}^3|e^{ik_i}-1|^2,\qquad k\ne0 .
 \label{eq:f15v62:projector}
\end{equation}
The value at \(k=0\)
does not affect its
Lebesgue Fourier kernel.
\end{theorem}

\begin{proof}
The moment bounds
\eqref{eq:f15v62:moments}
give tightness on the
countable product by
choosing summable coordinate
tail budgets. Along any
weakly converging subsequence,
\[
 X_p-\tfrac12|Y_p^\rho|^2
                  =\frac{X_p^2}{2\gamma_M}\longrightarrow0
                                   \quad\hbox{in }L^1 .
\]
Equation
\eqref{eq:f15v62:linearization}
gives \(d_2Y=0\) in
every limit. Uniform
second moments pass by
lower semicontinuity.
Common colour invariance
follows from the exact
root vertex gauge
transformation of the
vacuum link law, and
implies zero means.

Translation invariance
of the root-transported
limit needs more than
deterministic colour
invariance. For a fixed
translation \(z\ne0\),
let \(Z\) be the translated
field rooted at \(z\),
and \(g_z=U_{\rho_z}\).
Gauge transformation at
vertex zero leaves \(Z\)
unchanged and left-multiplies
\(g_z\). Haar averaging
at that vertex gives,
for a bounded test \(f\),
\[
 \mathbb E f(\operatorname{Ad}_{g_z}Z,X)
       =\mathbb E\int f(\operatorname{Ad}_hZ,X)\,dh
       =\mathbb E f(Z,X).
\]
The last equality uses
common colour invariance
at root \(z\).
The concatenated paths
\(\rho_z\) followed by the
translated root paths
produce \(\operatorname{Ad}_{g_z}Z\).
By \eqref{eq:f15v62:transport}
they have the same
limiting law as the
original paths at the
translated coordinates.
The vacuum link law is
translation invariant, so
this proves stationarity
of the limiting curvature
law. For large \(M\),
the two fixed root
vertices are distinct.

Next use the adapted
edge-avoiding paths in
\eqref{eq:f15v62:Stein-error}.
Their limits equal the
original path limits by
\eqref{eq:f15v62:transport}.
Its right side tends to
zero, and uniform second
moments give uniform
integrability of the
linear-growth integrand
on its left. Bounded
continuity handles the
directional derivative.
For every finite edge
cochain \(h\), linearity
after passing to the
limit therefore gives
\begin{equation}
       \mathbb E D_{d_1h}f(Y)
                =\mathbb E f(Y)\langle Y,d_1h\rangle .
 \label{eq:f15v62:limit-Stein}
\end{equation}
Path independence is used
after taking each finite
identity, not to justify
differentiating a path
comparison error.

For completeness this
identity supplies the
required measure translation
property, not merely a
formal moment relation.
For \(\ell=d_1h\), the
characteristic function
\(\varphi(B)=\mathbb E e^{i\langle Y,B\rangle}\)
obeys
\[
 \frac{d}{dt}\varphi(B+t\ell)
       =-\langle B+t\ell,\ell\rangle\varphi(B+t\ell).
\]
Thus the score
\(\langle Y,\ell\rangle\)
is Gaussian with variance
\(\|\ell\|_2^2\) and is
independent of every
finite family of scores
with tests perpendicular
to \(\ell\). Tilting this
one-dimensional Gaussian
and then using the
cylinder generating class
gives, for bounded Borel \(F\),
\[
 \mathbb E F(Y+\ell)
      =\mathbb E F(Y)
           e^{\langle Y,\ell\rangle-\|\ell\|_2^2/2}.
\]
Take \(h=d_1^*B\).
All hypotheses of the
archived abstract
classification are now
verified for the vacuum
limit. It follows that
\[
                  Y=G+H,\qquad G\ \hbox{independent of }H,
\]
where \(H\) is a
random two-form constant
over \(\mathbb Z^4\), and \(G\) has
the covariance
\eqref{eq:f15v62:projector}.

It remains to remove
\(H\) for this law.
Translation invariance
and the exact cell
partition give
\[
 \sum_{\mu<\nu}\mathbb E_{\rm vac}
            X_{(0;\mu,\nu)}
                  =\mathbb E_{\rm vac}T_0
                  =e_{M,\mathrm{vac}}(\gamma_M)\longrightarrow9/2.
\]
Uniform integrability
passes this expectation
to the joint limit.
For the Gaussian projector,
\(\mathcal P_{pp}=1/2\)
in each colour; hence
\[
 \frac92
     =\frac12\sum_{\mu<\nu}\mathbb E|Y_{\mu\nu}(0)|^2
     =\frac92+
            \frac12\mathbb E\sum_{\mu<\nu}|H_{\mu\nu}|^2 .
\]
Independence and the
zero mean of \(G\)
remove the cross terms.
Therefore \(H=0\) almost
surely. Every subsequential
limit is the stated
Gaussian-energy law,
which proves convergence
of the whole sequence.
The bounds in
Proposition~\ref{f15v62:moments}
then upgrade weak
convergence to convergence
of every fixed mixed
polynomial moment.
\end{proof}

\subsection{The vacuum susceptibility now has its exact limit}

Let \(\mathcal T_x\) be
the centered cell energy
built from \(G\):
\[
 \mathcal T_x=\sum_p w_x(p)|G_p|^2/2-\frac92,
 \qquad
 w_x(p)=\frac14{\bf1}_{\{p\text{ is a face of }x\}} .
\]
The Gaussian sum rule
proved in f14 V99 at
\texttt{lem:v99-Gaussian-sum-rule}
is
\begin{equation}
 C_\infty(z)=\mathbb E\mathcal T_0\mathcal T_z\ge0,
 \qquad
                  \sum_{z\in\mathbb Z^4}C_\infty(z)=\frac92 .
 \label{eq:f15v62:Gaussian-sum}
\end{equation}
Its normalization is worth
retaining: Wick's formula
gives
\(\operatorname{Cov}(|G_p|^2/2,|G_q|^2/2)
=(3/2)\mathcal P_{pq}^2\);
the projection identity gives
\(\sum_q\mathcal P_{pq}^2=1/2\).
Since \(\sum_xw_x(q)=1\)
and \(\sum_pw_0(p)=6\),
Tonelli's theorem gives
\((3/2)(1/2)6=9/2\).
This is an inherited
Gaussian calculation, not
a new computation of the
interacting vacuum spectrum.

\begin{theorem}[Exact selected vacuum susceptibility saturation]
\label{f15v62:susceptibility}
Along the same nearby
coupling selection,
\begin{equation}
 \chi_{M,\mathrm{vac}}(\widetilde\beta_M)
       =\|\Gamma_{M,\mathrm{vac}}(\widetilde\beta_M)\|
                         \longrightarrow\frac92 .
 \label{eq:f15v62:susceptibility}
\end{equation}
\end{theorem}

\begin{proof}
Let \(D_R=[-R,R]^4\cap\mathbb Z^4\),
with \(R\) fixed, and use
the unit source
\(a_x=|D_R|^{-1/2}{\bf1}_{D_R}(x)\).
For all sufficiently large
\(M\) it embeds in the
vacuum spatial cylinder
without aliasing. By V61,
\[
 \frac1{|D_R|}\Var_{\rm vac}
                       \left(\sum_{x\in D_R}T_x\right)
                    \le\chi_{M,\mathrm{vac}}(\widetilde\beta_M).
\]
The just-proved polynomial
moment convergence makes
the left side tend to
\[
 V_R=\sum_zC_\infty(z)
                    \frac{|D_R\cap(D_R-z)|}{|D_R|}.
\]
The summable covariance
in \eqref{eq:f15v62:Gaussian-sum}
gives \(V_R\to9/2\)
as \(R\to\infty\).
Thus first taking a
lower limit in \(M\),
then increasing the
fixed box, proves
\(\liminf\chi_{\rm vac}\ge9/2\).
V60 supplies
\(\limsup\chi_{\rm vac}\le9/2\)
on this same sequence.
Equality follows, and
V61's exact norm identity
gives the other equality
in \eqref{eq:f15v62:susceptibility}.
No interchange of these
two limits is used.
\end{proof}

\subsection{Scope of the new vacuum limit}

The vacuum and the earlier
selected periodic laws now
have the same proved
microscopic curvature-energy
limit, in the stated
fixed-coordinate normalization.
This does not identify
their finite-regulator
states, their spectral
measures, or their physical
continuum limits.

In particular the theorem
holds with fixed lattice
paths and fixed lattice
separations as the coupling
and spatial regulator grow.
It contains no estimate
uniform over paths and
supports growing to a
fixed physical size.
The Gaussian ultraviolet
limit therefore neither
proves nor disproves a
positive physical mass
gap. It is not being
substituted for a
nontrivial interacting
continuum state.

The remaining task is to
control the appropriate
growing-scale observables
and their physical transfer
spectrum, including a
calibrated time scale and
a sufficiently rich set
of physical probes.
The per-source error
in \eqref{eq:f15v62:window}
must also be paid on
growing supports.
Repeating this local
classification would not
resolve those requirements.

\subsection{Exact verification and source preservation}

The new diagnostic is
\begin{center}\small
\path{scripts/verify_ym_f15_v62_vacuum_sources_curvature.py}.
\end{center}
It reruns V61--V57.
There are 288 certified
recentered source-window
comparisons, 1,152 exact
affine recentering budgets,
288 positive cell-to-plaquette
domination checks, and
the rational constant
\(59/24<3\) for the
exponential moment bound.

The inherited local
curvature routines are
rerun without modification:
625 edge-avoiding paths,
96 patch gauge-word checks,
16 word fillings, 65
signed square factors,
128 signed variation
identities, 32 product
remainders, 12 closed
source stencils and 144
Gaussian regression
coefficients. The finite
Gaussian projection routine
also reruns its 16 symbols,
six projection row sums,
16 structure factors and
17 signed-source checks.
Its finite zero-mode-corrected
susceptibility is \(135/32\),
not the infinite-lattice
value \(9/2\).

The analytic Haar identity,
classification and limiting
claims rely on the
verified hypotheses and
proofs above, not on a
finite matrix simulation.
The script has 3,934 bytes
and SHA--256
\begin{center}\scriptsize\ttfamily
CF52808286F8F0B4122292BAE302CC93B4C5C56509EE7E72B2A3EF2C0F6BC0B5.
\end{center}
The V61 predecessor has
1,390,537 bytes and SHA--256
\begin{center}\scriptsize\ttfamily
D5DD8FA53EE7D1F1D7564E8AF969030DF65053EEAA623C93E7CC9C5AAD257B63.
\end{center}
Removing this final section
and reversing the title
increment recovers V61
byte for byte. V62 replaces
V61 as the single active
main source. Protected
inputs and 104 earlier
diagnostics are unchanged.
Only \path{.tex} files
are kept in
\path{latex_notes}; all
build products go to
\begin{center}\small
\path{artifacts/ym_f15_v62_texcheck/}.
\end{center}

\begin{firewall}
V62 proves nonzero-source
bounds, uniform local
thermal moments, the
microscopic Gaussian
curvature-energy limit,
and exact susceptibility
saturation for the
actual selected vacuum.
It does not construct
the interacting physical
continuum theory, control
all growing physical
supports, or establish
a positive Yang--Mills
mass gap.
\end{firewall}

\section{V63: a locally visible soft spectral endpoint and collapse of the one-step lattice gap}
\label{f15v63:context}

The microscopic vacuum
limit in V62 has a
spectral consequence that
must be faced before
seeking a physical gap.
It forces the actual
one-step Wilson lattice
gap along this selection
to tend to zero. This
does not say that the
physical mass tends to
zero: lattice and physical
energies differ by a
time-spacing factor that
has not yet been
constructed here.

The older archive
\begin{center}\small
\path{Renewal_Yang_Mills_Mass_Gap_Research_Notes_V135_f1.tex}
\end{center}
already distinguishes a
one-step contraction from
a rate per physical time,
at the archived theorems
\begin{center}\small
\texttt{thm:fixed-step-collapse}\\
\texttt{thm:rate-calibrated-gap}.
\end{center}
Those conditional continuum
criteria are not reproved.
The new result below
concerns the actual Wilson
vacuum and a specified
local energy probe, and
does not assume a physical
continuum limit in order
to establish the lattice
spectral conclusion.

Use the same selected
couplings
\(\gamma_M=\widetilde\beta_M\)
and spatial regulators
\(L=M\) as V60--V62.
The next companion
checkpoint remains V65.

\subsection{A local vacuum probe and its actual spectral moments}

Write
\[
 \mathsf P_M=\mathsf T_M/\lambda_{0,M},
 \qquad Q_M=I-|\Omega_M\rangle\langle\Omega_M|,
 \qquad H_M=-\log\mathsf P_M.
\]
The transfer operator acts
on the physical gauge-invariant
slice space. Put
\begin{equation}
 \rho_M=\|\mathsf P_M|_{Q_M\mathcal H_M}\|,\qquad
 g_M=-\log\rho_M
       =\inf\operatorname{spec}(H_M|_{Q_M\mathcal H_M}).
 \label{eq:f15v63:lattice-gap}
\end{equation}
The compact positive transfer
operator has a simple
isolated top eigenvalue,
so at each fixed regulator
\(0<\rho_M<1\) and \(g_M>0\).
No regulator-uniform lower
bound has been proved.

Let \(T_{(t,\mathbf0)}\)
be the original cell
energy at a fixed spatial
site and time \(t\).
Use the spatial delta
profile \(f=\delta_{\mathbf0}\)
in V61, with its literal
insertion \(\mathsf B_f\),
and define
\begin{equation}
 \psi_M=\lambda_{0,M}^{-1}Q_M\mathsf B_f\Omega_M,
 \qquad
 C_M(t)=\operatorname{Cov}_{\rm vac}
                 (T_{(0,\mathbf0)},T_{(t,\mathbf0)}).
 \label{eq:f15v63:probe}
\end{equation}
For every integer \(t\ge1\),
V61 gives
\begin{equation}
 C_M(t)=\langle\psi_M,\mathsf P_M^{\,t-1}\psi_M\rangle,
 \qquad C_M(1)=\|\psi_M\|^2,\qquad
 0\le C_M(t)\le C_M(1)\rho_M^{t-1}.
 \label{eq:f15v63:probe-moments}
\end{equation}
The mass of this spectral
measure is \(C_M(1)\),
not the contact variance
\(C_M(0)\).

Let
\[
 C_{\rm G}(t)=
       \mathbb E\mathcal T_0\mathcal T_{t e_0},
\]
where \(\mathcal T\) is
V62's centered Gaussian
cell-energy field.
Its proved polynomial
moment convergence implies
\begin{equation}
                       C_M(t)\longrightarrow C_{\rm G}(t)
                    \quad\hbox{for every fixed integer }t\ge0.
 \label{eq:f15v63:local-limit}
\end{equation}
This is a fixed-lattice-time
limit, not a uniform
estimate at growing \(t\).

\subsection{A positive temporal Fourier coefficient of the Gaussian curvature}

Use spatial directions
\(1,2,3\) and time
direction \(0\). For
\(\mathbf p\in[-\pi,\pi]^3\),
set
\[
 K(\mathbf p)=4\sum_{j=1}^3\sin^2(p_j/2),\qquad
 K_{12}(\mathbf p)=4\sin^2(p_1/2)+4\sin^2(p_2/2),
 \qquad
 \cosh E(\mathbf p)=1+\frac{K(\mathbf p)}2 .
\]
For \(\mathbf p\ne0\),
\(E(\mathbf p)>0\).
The diagonal symbol of
the spatial \(12\)-plaquette
curvature projector is
\[
 \mathcal P_{12,12}(\omega,\mathbf p)
       =\frac{K_{12}(\mathbf p)}
                      {K(\mathbf p)+2-2\cos\omega}.
\]

\begin{proposition}[Explicit temporal kernel and a lower tail]
\label{f15v63:Fourier}
For every integer \(t\ge1\),
\begin{equation}
 \begin{split}
 A(t)&:=\mathcal P_{12,12}(t,\mathbf0)\\
 &=\int_{[-\pi,\pi]^3}
       \frac{K_{12}(\mathbf p)}{2\sinh E(\mathbf p)}
           e^{-tE(\mathbf p)}\,\frac{d\mathbf p}{(2\pi)^3}\\
 &\ge\frac1{3e\sqrt5\,\pi^4}\,t^{-4}.
 \end{split}
 \label{eq:f15v63:Fourier-tail}
\end{equation}
The zero momentum point
has measure zero and
causes no ambiguity.
\end{proposition}

\begin{proof}
For \(K>0\), write
\(r=e^{-E}\), so
\(K+2=r+r^{-1}\).
The geometric-series
identity
\[
 \frac{1-r^2}{1-2r\cos\omega+r^2}
               =\sum_{n\in\mathbb Z}r^{|n|}e^{in\omega}
\]
is absolutely convergent
for \(0<r<1\).
Since
\[
 K+2-2\cos\omega
          =\frac{1-2r\cos\omega+r^2}{r},
 \qquad
 \frac{r}{1-r^2}=\frac1{2\sinh E},
\]
its \(t\)-th Fourier
coefficient is
\(e^{-|t|E}/(2\sinh E)\).
This proves the integral
formula. The original
projector symbol is
bounded, so Fubini is
valid before carrying
out the temporal integral.

For the lower bound
restrict \(\mathbf p\)
to the Euclidean ball
\(|\mathbf p|\le1/t\).
It lies in the Fourier
cube and has radius
at most one. On this
ball,
\[
 \begin{gathered}
 K_{12}(\mathbf p)\ge
           \frac4{\pi^2}(p_1^2+p_2^2),\qquad
 K(\mathbf p)\le|\mathbf p|^2,\\
 E(\mathbf p)=2\operatorname{arsinh}\bigl(\sqrt{K(\mathbf p)}/2\bigr)
                                    \le|\mathbf p|,\\
 2\sinh E(\mathbf p)
       =2\sqrt{K(\mathbf p)+K(\mathbf p)^2/4}
                                    \le\sqrt5\,|\mathbf p|.
 \end{gathered}
\]
The sine lower bound uses
\(\sin u\ge2u/\pi\) for
\(0\le u\le\pi/2\).
Also \(e^{-tE(\mathbf p)}\ge e^{-1}\)
on this ball. Finally
spherical integration gives
\[
 \int_{|\mathbf p|\le1/t}
           \frac{p_1^2+p_2^2}{|\mathbf p|}\,d\mathbf p
       =\frac{8\pi}{3}\int_0^{1/t}r^3\,dr
       =\frac{2\pi}{3t^4}.
\]
Combining these estimates
with the normalized Fourier
measure gives
\(4(2\pi/3)/
[\pi^2\sqrt5\,e(2\pi)^3]
=1/(3e\sqrt5\,\pi^4)\).
\end{proof}

\begin{proposition}[A polynomial lower bound for the limiting energy probe]
\label{f15v63:energy-tail}
For all integers \(t\ge1\),
\begin{equation}
        C_{\rm G}(t)\ge c_*t^{-8},
 \qquad c_*=\frac1{480e^2\pi^8}>0.
 \label{eq:f15v63:energy-tail}
\end{equation}
In addition,
\(C_{\rm G}(1)\le9/2\)
and \(C_{\rm G}(t)\to0\).
\end{proposition}

\begin{proof}
V62's Wick formula expresses
every cell-energy covariance
as a sum of nonnegative
squares of real curvature
projector entries, with
positive face weights.
Choose the spatial
\(12\)-plaquette based
at the cell origin and
its translate by \(t e_0\).
Both appear in their
respective cell observables
with weight \(1/4\).
Keeping just this pair
therefore gives
\[
 C_{\rm G}(t)\ge
           \frac32\left(\frac14\right)^2 A(t)^2
       \ge\frac{3}{32}\,
                  \frac1{45e^2\pi^8}\,t^{-8}
       =c_*t^{-8}.
\]
This is an inequality of
correlations. It is not
being used as a
domination statement between
spectral measures.
The remaining assertions
follow from the nonnegative,
absolutely summable covariance
kernel in
\eqref{eq:f15v62:Gaussian-sum},
whose total sum is \(9/2\).
\end{proof}

\subsection{The actual lattice gap must shrink}

\begin{theorem}[Collapse of the selected one-step Wilson gap]
\label{f15v63:gap-collapse}
For the actual vacuum
transfer operators on the
selected sequence,
\begin{equation}
                              g_M\longrightarrow0 .
 \label{eq:f15v63:gap-collapse}
\end{equation}
More specifically, for
each fixed integer \(t\ge2\),
\begin{equation}
 \limsup_{M\to\infty}g_M
       \le\frac{\log(C_{\rm G}(1)/C_{\rm G}(t))}{t-1}
       \le\frac{\log((9/2)/c_*)+8\log t}{t-1}.
 \label{eq:f15v63:gap-upper}
\end{equation}
There is also no pair
\(A<\infty\), \(m>0\)
such that
\[
                  C_M(t)\le A e^{-mt}
\]
for every sufficiently
large \(M\) and every
integer \(t\ge1\).
\end{theorem}

\begin{proof}
For each fixed \(t\ge2\),
\eqref{eq:f15v63:local-limit}
and \eqref{eq:f15v63:energy-tail}
make \(C_M(t)\) and
\(C_M(1)\) positive
eventually. Taking logarithms
in \eqref{eq:f15v63:probe-moments}
gives the finite-regulator
witness
\[
               g_M\le
                  \frac{\log(C_M(1)/C_M(t))}{t-1}.
\]
Pass to the upper limit
at this fixed \(t\).
The polynomial lower bound
and \(C_{\rm G}(1)\le9/2\)
give \eqref{eq:f15v63:gap-upper}.
Now let this fixed
integer test time grow.
Its right side tends
to zero. Since \(g_M\ge0\),
this proves
\eqref{eq:f15v63:gap-collapse}.
No time \(t=t_M\) was
substituted into the
local limit.

If the uniform exponential
bound held, passing to
the limit for each fixed
\(t\) would give
\(C_{\rm G}(t)\le A e^{-mt}\).
This contradicts
\(C_{\rm G}(t)\ge c_*t^{-8}\)
for large \(t\).
\end{proof}

Thus the shrinking spectral
scale is visible to this
local scalar energy probe.
It is not explained only
by an unobserved wrapping
or global-charge sector.
The theorem gives no
quantitative rate for
\(g_M\) in \(M\) or
\(\gamma_M\), because the
input moment convergence
has no such rate.

\subsection{An actual limiting spectral measure reaches the soft endpoint}

Define a finite positive
measure on \([0,1]\) by
\[
 \zeta_M(B)=
       \langle\psi_M,{\bf1}_B(\mathsf P_M)\psi_M\rangle .
\]
It is supported in
\([0,\rho_M]\), and has
no atom at \(1\).
Its moments are
\begin{equation}
       \int_0^1 r^n\,d\zeta_M(r)=C_M(n+1),
                                      \qquad n=0,1,\ldots .
 \label{eq:f15v63:measure-moments}
\end{equation}

\begin{theorem}[The local probe has a soft limiting endpoint]
\label{f15v63:endpoint}
The measures \(\zeta_M\)
converge weakly to a
unique finite positive
measure \(\zeta_{\rm G}\)
on \([0,1]\) satisfying
\begin{equation}
 \int_0^1r^n\,d\zeta_{\rm G}(r)=C_{\rm G}(n+1),\qquad
 \zeta_{\rm G}(\{1\})=0,\qquad
 \zeta_{\rm G}((r,1])>0\quad(0\le r<1).
 \label{eq:f15v63:endpoint}
\end{equation}
Let \(\eta_M\) denote
the spectral measure of
\(H_M=-\log\mathsf P_M\)
for the same vector.
For every fixed \(E>0\),
\begin{equation}
 0<L(E)\le\liminf_M\eta_M((0,E))
       \le\limsup_M\eta_M((0,E])
       \le\frac94(1-e^{-E}),
 \label{eq:f15v63:window}
\end{equation}
where the explicit positive
lower witness is
\begin{equation}
 L(E)=\sup_{t=2,3,\ldots}
   \frac{\bigl[c_*t^{-8}-(9/2)e^{-E(t-1)}\bigr]_+}
                           {1-e^{-E(t-1)}} .
 \label{eq:f15v63:window-witness}
\end{equation}
\end{theorem}

\begin{proof}
The total masses
\(C_M(1)\) converge to
the positive finite
number \(C_{\rm G}(1)\).
Compactness of \([0,1]\)
therefore gives weakly
convergent subsequences.
By
\eqref{eq:f15v63:measure-moments},
each limit has the
displayed moments.
Polynomials are dense
in \(C([0,1])\), so these
moments determine the
measure uniquely and the
whole sequence converges.

Dominated convergence
applied to \(r^n\)
identifies
\[
 \zeta_{\rm G}(\{1\})
       =\lim_{n\to\infty}C_{\rm G}(n+1)=0.
\]
If \(\zeta_{\rm G}((r,1])=0\)
for some \(r<1\), its
moments would be at
most \(C_{\rm G}(1)r^n\),
contradicting the polynomial
lower bound at \(n+1\).
This proves the endpoint
assertions. The limit
approaches the endpoint
without acquiring a vacuum
atom there.

For the quantitative window
witness, split any
\(\zeta_M\) at \(r_0=e^{-E}\).
Its mass above \(r_0\)
is \(\eta_M((0,E))\),
since its atom at \(1\)
vanishes. On the lower
part \(r^{t-1}\le r_0^{t-1}\),
and on the upper part
\(r^{t-1}\le1\).
Therefore, for \(t\ge2\),
\begin{equation}
 \eta_M((0,E))
 \ge
 \frac{\bigl[C_M(t)-e^{-E(t-1)}C_M(1)\bigr]_+}
                         {1-e^{-E(t-1)}} .
 \label{eq:f15v63:finite-window-witness}
\end{equation}
First pass to the lower
limit in \(M\) at this
fixed \(E,t\). Then use
\eqref{eq:f15v63:energy-tail}
and \(C_{\rm G}(1)\le9/2\).
The bound holds for every
fixed integer \(t\), so
its supremum is also a
lower bound. For each
\(E>0\), exponential
decay of \(e^{-E(t-1)}\)
eventually beats \(t^{-8}\);
hence \(L(E)>0\).

Finally V61's actual vacuum
resolvent bound for the
unit delta profile gives
\[
             \eta_M((0,E])
                   \le\frac{\kappa_M}{2}(1-e^{-E}).
\]
Take the upper limit
using \(\kappa_M\to9/2\).
This proves
\eqref{eq:f15v63:window}.
\end{proof}

The weak limiting measure
has been determined by
its actual limiting
correlation moments.
No elementary density
formula is asserted for
it, and no spectral
domination has been
inferred from the
single Wick-square
correlation lower bound.

\subsection{Why this is a physical-scale requirement, not a no-gap theorem}

Every energy variable
\(g_M\), \(E\) and \(u\)
in this section uses one
lattice time step. If an
independently justified
physical time spacing
\(a_M\downarrow0\) is
constructed, the corresponding
cutoff energy scale is
\(g_M/a_M\).
The result \(g_M\to0\)
alone does not determine
that ratio.

Likewise the window
\eqref{eq:f15v63:window}
uses \(E>0\) fixed before
\(M\to\infty\).
It gives no lower bound
in a shrinking interval
\((0,m a_M)\).
Substituting \(E=m a_M\)
or \(t=t_M\) into its
fixed-parameter limit
would be unjustified.
This is precisely the
distinction needed for
the physical mass-gap
objective.

A time spacing cannot
be chosen after the fact
merely to turn the
ratio \(g_M/a_M\) into
a desired mass. It must
be supported by the
same continuum construction
and observable normalization.
A positive physical gap
would also need the
appropriate uniform physical
transfer estimate on a
sufficiently rich local
sector. The present
energy probe does not
by itself supply that
sector.

The direction is therefore
now explicit: a uniform
positive one-step gap along
this ultraviolet sequence
is not an available route.
The open work is a
rate-calibrated physical
limit, or an equivalent
fixed-physical-delay
contraction estimate, with
noncollapse and the
interacting continuum law
established at the same
scale. V63 does not
claim that these missing
requirements are satisfied,
nor that the full
Yang--Mills mass-gap goal
is impossible.

\subsection{Exact checks and source preservation}

The new diagnostic is
\begin{center}\small
\path{scripts/verify_ym_f15_v63_vacuum_spectral_endpoint.py}.
\end{center}
It reruns V62--V57 and
the local algebra routines
included by V62.
There are 68 exact
temporal resolvent recurrence
checks, four dispersion
identities and two rational
coefficient checks for the
spherical integral and
the cell Wick square.
Thirty-two finite positive
transfer checks compare
the probe covariance with
the full spectral norm;
128 tests check the
strict-window mass witness.

A separate auxiliary
density \((1-r)^7\,dr\)
is used only to check
compact moment approximation:
85 exact finite-bin
moment error bounds and
32 polynomial-tail
normalizations pass.
It is not identified
with \(\zeta_{\rm G}\)
or with a Wilson
spectral measure.
The Gaussian Fourier
integral, the actual
vacuum limit and the
gap conclusion use the
analytic proofs above.

The diagnostic has 4,249
bytes and SHA--256
\begin{center}\scriptsize\ttfamily
38A2D5E01BE33818FEFCBF632E0931DACD2A9B5AEA2584B61A6A852B3913247E.
\end{center}
The V62 predecessor has
1,414,231 bytes and SHA--256
\begin{center}\scriptsize\ttfamily
7DEB3459ABEC58E5FAFD8C0C032FC66A1316F11C0B8B3069B98CA6697669665B.
\end{center}
Removing this final section
and reversing the title
increment recovers V62
byte for byte. V63
replaces V62 as the
single active main source.
Protected inputs and
105 earlier diagnostics
are unchanged.
Only \path{.tex} files
belong in
\path{latex_notes}; all
build products go to
\begin{center}\small
\path{artifacts/ym_f15_v63_texcheck/}.
\end{center}

\begin{firewall}
V63 proves that the
actual selected Wilson
one-step lattice gap
tends to zero, and
that a local energy
probe has positive
asymptotic spectral weight
in every fixed positive
lattice-energy window.
It does not determine
the gap per physical
time, rule out a
positive physical mass,
construct the interacting
continuum theory, or
complete the Yang--Mills
mass-gap proof.
\end{firewall}

\section{V64: uniform vacuum energy covariance and a quantitative growing-time spectral window}
\label{f15v64:context}

V63 used a fixed time before taking the regulator limit.
It therefore proved lattice-gap collapse without a rate.
We now control the error uniformly in the lattice time
separation and can choose a time that grows with the
regulator. This is a quantitative statement about the
same actual Wilson vacuum, not a new coupling selection.

The relevant precedents have been checked to avoid
duplication. V5--V7 concern the periodic full compact
law and its blocked observables. The covariance result
in f14 V94 concerns a restricted thermal chart.
Their conclusions are not assumed for the vacuum.
We reuse the finite-range curl filter and the cell-frame
geometry of V5, but verify their probabilistic inputs
on the infinite-time vacuum cylinder. The new cubic
Haar test supplies fourth moments directly. It avoids
differentiating a characteristic-function error.
The V60 companion review remains current; the next
required review is V65.

All results in this section remain in lattice units.
The physical existence-and-gap target is unchanged.\footnote{
The official target is described by the Clay Mathematics
Institute at
\href{https://www.claymath.org/millennium/yang-mills-the-maths-gap/}
{Yang--Mills and the Mass Gap}, checked 16 September 2026.
The estimates below do not establish that target.}

\subsection{The vacuum inputs and a paid mean-energy excess}

Work on
\(\mathcal L_M=\mathbb Z\times(\mathbb Z/M\mathbb Z)^3\),
where \(M\) is the selected dyadic spatial side and
\(\gamma=\gamma_M=\widetilde\beta_M\).
Expectation under its stationary vacuum link law is
\(\mathbb E_{\rm vac}\).
Retain the plaquette variables \(X_p=\gamma s_p\),
the symmetrically assigned cell energy \(T_x\), and
the weights \(w_x(p)\) from V62.

Here is one explicit upper budget for the vacuum mean.
Write
\[
 \theta_M=\lambda_M+142e^{-d_M},\qquad
 \beta=\beta_M^\diamond,\qquad
 h=(12\beta/M)^{1/4},\qquad
 \varepsilon=17\sqrt{\beta/M}.
\]
The parameters \(\lambda_M,d_M\) in \(\theta_M\)
are the inherited periodic energy-budget parameters
of V5, not transfer eigenvalues. Define
\begin{equation}
 \vartheta_M
 =\theta_M+(1+h)(4k_Mh+h^3+\varepsilon)
                         +h(9/2+\theta_M).
 \label{eq:f15v64:mean-budget}
\end{equation}
Then, eventually,
\begin{equation}
 0\le\vartheta_M\le1,\qquad
 \vartheta_M\longrightarrow0,\qquad
 \mathbb E_{\rm vac}T_0\le9/2+\vartheta_M,\qquad
 2\le\gamma_M\le2\log M.
 \label{eq:f15v64:inputs}
\end{equation}
Indeed the periodic mean \(e_0\) is at most
\(9/2+\theta_M\), and V60 proves
\[
 |e_{M,\rm vac}(\gamma_M)-e_0|
       \le(1+h)(4k_Mh+h^3+\varepsilon)+h|e_0|.
\]
The mean is nonnegative, so substitution proves
\eqref{eq:f15v64:inputs}. All these budgets refer to
the already selected sequence.

By Proposition~\ref{f15v62:moments}, at these couplings
all plaquette moments are bounded uniformly in \(p,M\).
In particular, for each fixed integer \(q\ge1\),
\begin{equation}
 \mathbb E_{\rm vac}X_p^q\le e^3\,16^q q!,
 \qquad
 \mathbb E_{\rm vac}|\sqrt\gamma\,\operatorname{Im}U_p|^{2q}
                                      \le e^3\,32^q q!.
 \label{eq:f15v64:moments}
\end{equation}
These estimates remain valid after any adjoint transport.
They hold for the actual vacuum law, not for an
unproved conditional version of that law.

We also use V62's local Haar identity for smooth
cylinder tests and edge-avoiding frames. At fixed
\(M,\gamma\), the variables in such a test are bounded.
The identity follows from the finite-period law and
its cylinder convergence, with only the six
plaquettes incident to the varied edge in the action
derivative. No infinite total action is differentiated.

\subsection{Auxiliary frames on the infinite-time cylinder}

Choose a dyadic integer \(R\mid M\), with
\(4\le R\le M/4\). Partition \(\mathcal L_M\) into
cells of side \(R\) in all four coordinates, with
an independent uniform offset in
\(\{0,\ldots,R-1\}^4\). There are countably many
cells in the time direction. Attach independent
Haar matrices \(H_D\) to them.

For a plaquette \(p=(y;\mu\nu)\), let \(D(y)\)
be the cell of its base and \(g_y^D\) the ordered
increasing-coordinate path from that cell's root
to \(y\), in its lifted cell. Set
\begin{equation}
 W_p=\sqrt\gamma\,
       \operatorname{Ad}_{H_Dg_y^D}\operatorname{Im}U_p,
 \qquad
 \epsilon_R=R/\sqrt\gamma+1/R.
 \label{eq:f15v64:auxiliary}
\end{equation}
The plaquette may cross the cell boundary; its
assignment is by its base. Expectation \(\mathbb E\)
below includes only these additional independent
offsets and rotations, without changing the vacuum law.

The field \(W\) is stationary on \(\mathcal L_M\).
Translations permute the offsets and cells, preserve
the vacuum law, and carry the prescribed paths to the
paths of the translated partition. Its mean is zero
by the Haar rotation of its cell. Its two-point
covariance vanishes between different cells after
conditioning on the links and the partition.
In particular the averaged covariance has finite
range in time.

For every realization, the squared norms are exactly
preserved:
\begin{equation}
 |W_p|^2=2X_p-X_p^2/\gamma,\qquad
 \mathcal E(W):=\sum_{\mu<\nu}\mathbb E|W_{0;\mu\nu}|^2
                                 \le9+2\vartheta_M .
 \label{eq:f15v64:energy}
\end{equation}
It follows in particular that products of the
quadratic plaquette energies are unchanged by
the auxiliary rotations, even when the plaquettes
are in different cells. Vanishing auxiliary
\emph{linear} colour covariance does not imply
vanishing covariance of their squared norms.

\subsection{Linear and cubic tests in the vacuum Haar identity}

Let indices \(i\) denote a site, plaquette orientation,
and one real colour component of \(W\).
Let \(d=d_1\) be the ordinary lattice curl and let
\(a\) be a finitely supported real colour edge cochain.
For \(k=1\) or \(3\), set
\(F(W)=\prod_{j=1}^k W_{i_j}\), allowing repeated
indices. The counting \(\ell^1\) norm uses the
Euclidean norm of each colour coefficient.

\begin{proposition}[Polynomial Stein bound, uniform in the index locations]
\label{f15v64:polynomial-Stein}
There are constants \(C_k\), independent of
\(M,\gamma,R,a\) and the index locations, such that
\begin{equation}
 \left|
 \mathbb E\langle W,da\rangle\prod_{j=1}^k W_{i_j}
 -
 \sum_{j=1}^k(da)_{i_j}
        \mathbb E\prod_{\ell\ne j}W_{i_\ell}
 \right|
 \le C_k\epsilon_R\|a\|_1 ,
 \qquad k=1,3.
 \label{eq:f15v64:polynomial-Stein}
\end{equation}
For \(k=1\), the empty product is one.
\end{proposition}

\begin{proof}
First take \(a=h_e\), supported on one edge \(e\).
As in V5, call an offset interior if all four
coordinates of the base of \(e\) within its cell
are between \(1\) and \(R-3\). The noninterior
probability is at most \(16/R\). On an interior
offset, the star of \(e\) and an elementary detour
around it fit in that cell.

Replace each ordered root path that uses \(e\)
by the three-edge detour around a fixed adjacent
square. Each such path uses \(e\) at most once;
other cells' root paths avoid it already.
Denote the resulting field by \(W'\).
For every fixed finite \(p\ge1\), the moment bounds
\eqref{eq:f15v64:moments}, the conjugation bound,
and one square filling give
\begin{equation}
 \|W'_i-W_i\|_{L^p}\le C_p/\sqrt\gamma,\qquad
 \|W_i\|_{L^p}+\|W'_i\|_{L^p}\le C_p.
 \label{eq:f15v64:frame-return}
\end{equation}
These estimates hold uniformly for each fixed
interior offset and each choice of the colour matrices.
For example, the difference is bounded by a
constant times a plaquette chord times a curvature
norm; apply H\"older and the corresponding twice
higher moments. No independence of these factors
is used.

A rooted star link has a filling with at most
\(3R+2\) faces, including the endpoint detours.
For \(T_u=\sqrt\gamma\,|k_u-I|\) this gives
\(\|T_u\|_{L^p}\le C_pR\).
This follows by telescoping a filled word, applying
Minkowski, and using
\(|U_p-I|^2=2X_p/\gamma\).
All adapted paths avoid the varied edge, so their
frames are fixed by its conditional Haar flow.

Use precisely that flow, in the convention of V62.
With \(\ell=dh_e\), the differentiated plaquette
and action identities imply
\begin{equation}
 \begin{aligned}
 \|D_eW'_i-\ell_i\|_{L^p}
                  &\le C_p|h|R/\sqrt\gamma
                       &&\text{on the star},\\
 D_eW'_i&=0 &&\text{off the star},\\
 \|D_e(\gamma S_e)-\langle W',\ell\rangle\|_{L^p}
                  &\le C_p|h|R/\sqrt\gamma .
 \end{aligned}
 \label{eq:f15v64:derivatives}
\end{equation}
Here \(S_e\) contains only the six incident plaquettes.
For the last estimate the pointwise remainder is
bounded by a constant times
\[
 \frac{|h|}{\sqrt\gamma}
          \sum_{p\ni e}|W'_p|\sum_{u\in\partial p}T_u.
\]
H\"older and the uniform higher moments give every
fixed \(L^p\) bound displayed above.

Apply the exact vacuum identity
\[
 \mathbb E_{\rm vac}D_eF(W')
       =\mathbb E_{\rm vac}F(W')D_e(\gamma S_e).
\]
The product rule and \eqref{eq:f15v64:derivatives}
compare it to the claimed identity for \(W'\)
with error \(C_k|h|R/\sqrt\gamma\).
All products have fixed degree at most four;
H\"older controls them with
\eqref{eq:f15v64:moments}.
At fixed regulator these are smooth bounded cylinder
tests, so there is no differentiation of an
unbounded limiting polynomial.

Return to \(W\) by telescoping each product and
using \eqref{eq:f15v64:frame-return}.
The star has six plaquettes and
\(\|\ell\|_1\le6|h|\), so this costs
\(C_k|h|/\sqrt\gamma\), independently of the
locations of the test coordinates.
Repeated indices are included by the same product
rule with multiplicity.

On noninterior offsets, use only H\"older, the
uniform coordinate moments, and
\(\|\ell\|_1\le6|h|\). The absolute discrepancy
is at most \(C_k|h|\) for every such offset.
Its averaged contribution is at most \(C_k|h|/R\).
Combining the cases proves the result for one edge.
Sum over the finite edge support of \(a\) to obtain
\eqref{eq:f15v64:polynomial-Stein}.
\end{proof}

This argument does not assume a conditional
exponential-moment estimate. The moment bounds
are unconditional, uniform for each independent
offset and colour choice, and the exact local
Haar identity is used before averaging them.

\subsection{The cylinder heat filter and the low-frequency residual}

Let \(\Delta\) be the scalar lattice Laplacian on
\(\mathcal L_M\), acting componentwise. Set
\[
 \mathsf R=I-\Delta/16,\qquad
 q_n(\Delta)=\frac1{16}\sum_{s=0}^{n-1}\mathsf R^s,
 \qquad
 \mathsf T_n=dd^*q_n(\Delta).
\]
To distinguish it from the physical transfer
operator, \(\mathsf T_n\) in this subsection is
only the finite-range \emph{curvature filter}.
Let \(\mathcal P_M^{\rm cyl}\) be the coexact
projection on this cylinder, with the colour
identity included when indices carry colour.
Its multiplier is \(dd^*/\omega\) for \(\omega>0\).
The only zero of \(\omega\) has zero measure in
the continuous time-frequency integral.
Consequently its rank, including colours, is
nine almost everywhere and
\[
 0\le\mathsf T_n\le\mathcal P_M^{\rm cyl}\le I,
 \qquad
 \mathsf T_n(k)=(1-r(k)^n)\mathcal P_M^{\rm cyl}(k).
\]
For every unit face-colour coordinate \(e_i\),
\begin{equation}
 \mathsf T_ne_i=da_{n,i},\qquad
 \|a_{n,i}\|_1\le n/4,\qquad
 \|\mathsf T_ne_i\|_1\le3n/2.
 \label{eq:f15v64:filter-norms}
\end{equation}
The proof is the positive-kernel mass \(n/16\)
and the incidence norms four and six, exactly
as in V5. Finite spatial periodization cannot
increase these absolute coefficient norms.

The cylinder lazy-walk return probability satisfies
\begin{equation}
 k_M^{\rm cyl}(s)
 \le C s^{-1/2}(s^{-1/2}+M^{-1})^3
 \le C(s^{-2}+M^{-3}s^{-1/2}),\qquad s\ge1.
 \label{eq:f15v64:heat}
\end{equation}
To prove it, bound \(r(k)^s\) on the principal
frequency cube by
\(\exp[-s|k|^2/(4\pi^2)]\).
The continuous time-frequency integral is at
most \(Cs^{-1/2}\), and each spatial grid sum
is at most \(C(s^{-1/2}+M^{-1})\).
Use \((u+v)^3\le4(u^3+v^3)\) to conclude.
In particular \(k_M^{\rm cyl}(s)\le Cs^{-2}\)
for \(1\le s\le2M^2\), after changing \(C\).

\begin{proposition}[Vacuum auxiliary residual]
\label{f15v64:residual}
For \(1\le n\le M^2\),
\begin{equation}
 \mathcal E((I-\mathsf T_n)W)
       \le D_n:=2\vartheta_M+C(n^{-2}+n^2\epsilon_R).
 \label{eq:f15v64:residual}
\end{equation}
Moreover, uniformly in all coordinate pairs,
\begin{equation}
 |(\mathcal P_M^{\rm cyl}-\mathsf T_n)_{ij}|
                                                       \le Cn^{-2}.
 \label{eq:f15v64:projection-error}
\end{equation}
\end{proposition}

\begin{proof}
Use the \(k=1\) case of
Proposition~\ref{f15v64:polynomial-Stein} with
\(a=a_{n,i}\), first against \(W_i\) and then,
by linearity, against
\(\langle W,\mathsf T_ne_i\rangle\).
Equation~\eqref{eq:f15v64:filter-norms} gives
\[
 \begin{split}
 \mathbb E W_i(\mathsf T_nW)_i
                         &=(\mathsf T_n)_{ii}+O(n\epsilon_R),\\
 \mathbb E(\mathsf T_nW)_i^2
                         &=(\mathsf T_n^2)_{ii}+O(n^2\epsilon_R).
 \end{split}
\]
Sum over the eighteen coordinates at one site
and expand the residual square. The exact
Fourier trace identity on the cylinder is
\[
 9-2\sum_i(\mathsf T_n)_{ii}
             +\sum_i(\mathsf T_n^2)_{ii}
                                  =9k_M^{\rm cyl}(2n).
\]
The sums here are only over those eighteen
coordinates. The identity holds because the
projection rank is nine almost everywhere;
there is no discrete time zero-mode correction.
Equations~\eqref{eq:f15v64:energy} and
\eqref{eq:f15v64:heat} prove
\eqref{eq:f15v64:residual}.

For the second claim,
\(\mathcal P_M^{\rm cyl}-\mathsf T_n\) is positive
with multiplier \(r^n\mathcal P_M^{\rm cyl}\).
Its diagonal trace is \(9k_M^{\rm cyl}(n)\).
Positive-operator Cauchy--Schwarz bounds each
matrix entry by the geometric mean of the
corresponding diagonal entries. This proves
\eqref{eq:f15v64:projection-error}.
\end{proof}

These manipulations are legitimate on the
infinite-time cylinder. The covariance of \(W\)
has finite time range and finitely many spatial
coordinates, hence has an integrable positive
matrix Fourier density. Alternatively, each
finite-range filter expectation above is already
a finite sum of cylinder moments. No global
square-summability of a random infinite field
has been assumed.

\subsection{A fourth-moment comparison without filtering all four factors}

Write \(P=\mathcal P_M^{\rm cyl}\) temporarily and put
\[
 \mathfrak d_n=\sqrt{D_n}+n\epsilon_R+n^{-2}.
\]

\begin{theorem}[Uniform second and fourth moments]
\label{f15v64:Wick}
For every two or four scalar face-colour indices,
at arbitrary separations on \(\mathcal L_M\),
\begin{align}
 |\mathbb E W_iW_j-P_{ij}|&\le C\mathfrak d_n,
 \label{eq:f15v64:second}\\
 |\mathbb E W_iW_jW_kW_l
       -P_{ij}P_{kl}-P_{ik}P_{jl}-P_{il}P_{jk}|
                                      &\le C\mathfrak d_n.
 \label{eq:f15v64:fourth}
\end{align}
The constants are independent of these locations,
including coincident indices.
\end{theorem}

\begin{proof}
Write
\(W_i=\langle W,\mathsf T_ne_i\rangle+\mathcal R_i\).
By stationarity and
\eqref{eq:f15v64:residual},
\(\|\mathcal R_i\|_2\le\sqrt{D_n}\).
For the second moment, Cauchy--Schwarz gives
\(|\mathbb E\mathcal R_iW_j|\le C\sqrt{D_n}\).
Use the linear Stein bound for the filtered
first coordinate and
\eqref{eq:f15v64:projection-error} to obtain
\eqref{eq:f15v64:second}.

For the fourth moment, the uniform sixth moments
give
\(\|W_jW_kW_l\|_2\le C\) by H\"older.
Thus filtering only its first coordinate costs
at most \(C\sqrt{D_n}\).
The cubic Stein bound gives
\[
 \begin{split}
 \mathbb E\langle W,\mathsf T_ne_i\rangle W_jW_kW_l
  ={}&(\mathsf T_n)_{ij}\mathbb EW_kW_l
      +(\mathsf T_n)_{ik}\mathbb EW_jW_l\\
    &+(\mathsf T_n)_{il}\mathbb EW_jW_k
      +O(n\epsilon_R).
 \end{split}
\]
Apply \eqref{eq:f15v64:second} to each remaining
pair and \eqref{eq:f15v64:projection-error} to
each filter entry. Entries of both positive
contractions are bounded by one, and all
two-coordinate moments are uniformly bounded.
This proves \eqref{eq:f15v64:fourth}.
\end{proof}

Only one coordinate was filtered in this last
step. Replacing all four factors at once would
require unnecessarily growing higher-moment
bounds for filtered fields. Here the other
three factors retain their uniform actual
plaquette moments.

\subsection{Returning to the actual energy covariance}

Let \(P^{\rm cyl}_{pq}\) now denote the uncoloured
plaquette entry. Define the Gaussian cylinder
comparison
\[
 C_M^{G,\rm cyl}(x,y)
       =\frac32\sum_{p,q}w_x(p)w_y(q)
                                  (P^{\rm cyl}_{pq})^2 .
\]
It is the covariance of the quadratic cell
energies of a centred Gaussian curvature with
this projection covariance. It is not an
assumed Gaussian vacuum law.

\begin{theorem}[Actual vacuum covariance, uniform in lattice time]
\label{f15v64:energy-covariance}
There is a constant \(K<\infty\), independent
of \(M\), such that eventually
\begin{equation}
 \sup_{t\in\mathbb Z}
       |C_M(t)-C_G(t)|\le K\delta_M,\qquad
 \delta_M:=\sqrt{\vartheta_M}+\gamma_M^{-1/16}+M^{-2}
                                      \longrightarrow0.
 \label{eq:f15v64:uniform-time}
\end{equation}
Here \(C_M(t)\) is exactly V63's actual vacuum
cell-energy covariance, and \(C_G(t)\) is its
infinite-lattice Gaussian comparison.
\end{theorem}

\begin{proof}
Define \(\widehat T_x=\sum_p w_x(p)|W_p|^2/2\).
It is invariant under every auxiliary rotation
separately. Its joint law is therefore the joint
law of the corresponding actual chord-energy
observable under the vacuum measure.
The exact identity
\[
 T_x-\widehat T_x=\frac1{2\gamma}
                                    \sum_p w_x(p)X_p^2
\]
and the fourth moments of \(X_p\) give
\(\|T_x-\widehat T_x\|_2\le C/\gamma\).
Both cell energies have bounded \(L^2\) norm.
It follows that their covariances differ by at
most \(C/\gamma\), uniformly in \(x,y\).

Expand the covariance of the two quadratic
\(\widehat T\)'s into scalar fourth moments
and products of second moments. Each cell has
twenty-four faces and total weight six.
Theorem~\ref{f15v64:Wick} therefore gives
\[
 \left|\operatorname{Cov}_{\rm vac}(T_x,T_y)
                         -C_M^{G,\rm cyl}(x,y)\right|
                 \le C(\mathfrak d_n+\gamma^{-1}).
\]
The Wick coefficient is \(3/2\): the factors
\(1/2\) in the two energies and the two crossed
pairings give \(1/2\) in each of three colours.
No large-support factor depends on the separation
of these two fixed-size cells.

Choose \(R\) to be the largest dyadic integer
at most \(\gamma^{1/4}\), and
\(n=\max\{1,\lfloor\epsilon_R^{-1/4}\rfloor\}\).
Eventually \(R\mid M\), \(4\le R\le M/4\), and
\(1\le n\le M^2\), since \(\gamma\le2\log M\).
Moreover
\[
 \epsilon_R\asymp\gamma^{-1/4},\qquad
 n\asymp\gamma^{1/16},\qquad
 \mathfrak d_n\le C(\sqrt{\vartheta_M}+\gamma^{-1/16}).
\]

It remains to compare the two Gaussian kernels
without a restriction on time. Set
\(n_0=\lfloor M/8\rfloor\).
The cylinder and infinite-lattice kernels of
\(\mathsf T_{n_0}\) agree whenever the spatial
coordinate differences lie in \(\{-1,0,1\}^3\),
for every temporal difference. Indeed this
polynomial filter has range at most \(n_0+1\).
No nonzero spatial periodic image of such a
difference is in that range, and time is not
periodized. For large time differences both
finite-range kernels are zero.

Use \eqref{eq:f15v64:projection-error} at \(n_0\).
On the infinite lattice the identical heat
argument gives an entrywise error \(Cn_0^{-2}\).
Thus
\[
 |P^{\rm cyl}_{pq}-\mathcal P_{pq}|\le CM^{-2}
\]
for the faces of the cells at \(0\) and
\((t,0,0,0)\), uniformly in all integers \(t\).
Both projection kernels have entries bounded
by one, so the same \(CM^{-2}\) bound holds
for their finite weighted sum of squares.
Combining the comparisons proves
\eqref{eq:f15v64:uniform-time}.
\end{proof}

The theorem is an additive, not a relative,
comparison. It does not bound the relative
error when \(C_G(t)\) is much smaller than
\(\delta_M\). Its uniformity nonetheless
allows a controlled growing-time range.

\subsection{A rate for the lattice gap and actual shrinking-window weight}

Retain V63's local probe \(\psi_M\), lattice
gap \(g_M\), and its Hamiltonian spectral
measure \(\eta_M\).
Let \(c_*=1/(480e^2\pi^8)\) be the constant in
\eqref{eq:f15v63:energy-tail}.
Increase \(K\), if necessary, so that
\(K\ge1\) in \eqref{eq:f15v64:uniform-time}.
For sufficiently large \(M\), set
\begin{equation}
 A_M=\left(\frac{c_*}{2K\delta_M}\right)^{1/8},
 \qquad t_M=\lfloor A_M\rfloor,\qquad
 E_M=\frac{\log(12/c_*)+8\log t_M}{t_M-1}.
 \label{eq:f15v64:scales}
\end{equation}
Eventually \(A_M\ge4\), so \(t_M\ge2\).

\begin{theorem}[Growing-time probe and shrinking spectral window]
\label{f15v64:gap-rate}
Along the same selected actual vacuum sequence,
\begin{equation}
 C_M(t_M)\ge\frac{c_*}{2t_M^8},\qquad
 \eta_M((0,E_M))\ge\frac{c_*}{4t_M^8}>0,
 \qquad 0<g_M<E_M .
 \label{eq:f15v64:shrinking-weight}
\end{equation}
In particular, for a fixed constant \(C<\infty\),
\begin{equation}
 g_M\le C\delta_M^{1/8}
                      (1+|\log\delta_M|),\qquad
 E_M\le C\delta_M^{1/8}
                      (1+|\log\delta_M|)\longrightarrow0 .
 \label{eq:f15v64:gap-rate}
\end{equation}
This improves V63's order-of-limits conclusion
to a quantitative bound in its established
mean-excess and coupling parameters.
\end{theorem}

\begin{proof}
Since \(t_M\le A_M\), the error in
\eqref{eq:f15v64:uniform-time} is at most
\(c_*/(2t_M^8)\).
Apply the lower bound \(C_G(t)\ge c_*t^{-8}\)
from V63 to prove the first inequality.

The local resolvent estimate of V61 gives
\(C_M(1)\le\kappa_M/2\le3\) eventually.
By the choice of \(E_M\),
\[
 e^{-E_M(t_M-1)}C_M(1)
                   \le \frac{c_*}{4t_M^8}.
\]
Now apply V63's finite-regulator witness
\eqref{eq:f15v63:finite-window-witness},
which is valid for any integer time and
any positive threshold at that regulator.
Its numerator is at least \(c_*/(4t_M^8)\)
and its denominator is at most one.
This proves the spectral-weight bound.
The measure is supported on the spectrum of
\(H_M\) on the vacuum-orthogonal sector;
positive mass strictly below \(E_M\) implies
\(g_M<E_M\). The positivity of \(g_M\) at
each fixed regulator is the already established
compact-transfer statement.

Finally \(t_M\ge A_M/2\) and
\(t_M-1\ge t_M/2\). Substitute these inequalities
in \eqref{eq:f15v64:scales} and use
\(\log A_M=(1/8)\log[c_*/(2K\delta_M)]\)
to obtain \eqref{eq:f15v64:gap-rate}.
\end{proof}

The finite-window argument justifies these moving
parameters; it is not an application of V63's
fixed-\(E\) limiting theorem with \(E=E_M\).
The weight lower bound itself tends to zero.
There is no uniform nonzero low-energy weight
for a physically normalized continuum probe
in this assertion.

\subsection{The remaining physical interface}

The new estimate gives a growing lattice-time
window on which a specified actual local energy
probe still detects the soft spectrum. It does
not prove a lower bound on a physical mass.
In particular the factor
\(\delta_M^{1/8}\) contains the contribution
\(\gamma_M^{-1/128}\); this is a weak ultraviolet
rate, not a derived physical renormalization scale.

A physical spacing \(a_M\) must be obtained
independently, together with normalized observables
whose physical correlations and spectral weights
survive. The cutoff quantity would then be
\(g_M/a_M\), not \(g_M\). It is not legitimate
to define \(a_M\) after the fact solely to make
this ratio a chosen mass. Conversely, soft
cutoff vectors with vanishing weight do not
by themselves rule out a gapped limiting
physical sector.

The next task should address that normalization
and the genuinely interacting scale, using
relative or physically smeared estimates where
needed. Repeating the fixed-site Gaussian
classification would not pay this requirement.
A nontrivial interacting continuum law, its
reconstruction, and a positive physical
Hamiltonian gap remain unproved.

\subsection{Diagnostics and preservation}

The new diagnostic is
\begin{center}\small
\path{scripts/verify_ym_f15_v64_vacuum_growing_time.py}.
\end{center}
It reruns V63 and its inherited checks.
Independent expansion in standard normal
coordinates verifies 504 linear and cubic
Stein identities, 108 fourth moments including
repeated coordinates, and 324 single-filter
decompositions. There are 32 separate-rotation
energy-product identities, 24 cylinder return
and unperiodized-time checks, 18 filter-norm
checks, and 49,572 no-spatial-image kernel
checks. The spectral budget has 32
shrinking-window checks, six exact scaling
exponents, and 96 integer-rounding checks.

These rational tests audit finite algebra
and geometry. They are not simulations of the
Wilson vacuum and do not numerically certify
the analytic Haar estimates or a physical gap.
The script has 8,162 bytes and SHA--256
\begin{center}\scriptsize\ttfamily
16D30EE5CE9EAC20DBA005DD4BBFDC60C464F23F58C2E08EF4AEBCD076C3F0FE.
\end{center}

The V63 predecessor has 1,429,849 bytes and SHA--256
\begin{center}\scriptsize\ttfamily
1F212E5F7270F2D636BFF4054C3E40B55827BA94357BBD303B4B037571F88944.
\end{center}
Removing this final section and reversing the
title increment recovers V63 byte for byte.
V64 replaces V63 as the sole active main source.
The protected inputs and 106 earlier diagnostics
are preserved. Only \path{.tex} files belong in
\path{latex_notes}; compilation products go to
\begin{center}\small
\path{artifacts/ym_f15_v64_texcheck/}.
\end{center}
The next companion checkpoint is V65.

\begin{firewall}
V64 proves a uniform-in-lattice-time actual vacuum
energy-covariance comparison, a quantitative
upper bound on the selected lattice gap, and
positive probe weight in a justified shrinking
lattice-energy window. It does not provide the
physical time calibration, nontrivial interacting
continuum construction, or full Yang--Mills
mass gap. The original objective remains active.
\end{firewall}

\section{V65: a canonically smeared magnetic probe and a nonvanishing soft spectral limit}
\label{f15v65:context}

V64 controlled a local probe at growing lattice times,
but the weight in its shrinking spectral window could
vanish. A physical interpretation requires observables
with a specified normalization and surviving vectors.
This section tests that requirement for the thermal
magnetic plaquette energy, with dimension-four
spatial normalization and positive-time damping.
The result is a nonzero, explicitly identified
two-point spectral limit in a controlled mesoscopic
window. Its spectrum reaches zero. Thus this window,
with this probe retained, cannot be relabelled as
the desired massive physical theory.

This is not an interacting continuum construction
or a counterexample to the Yang--Mills mass-gap
problem. It is a scope-specific test of a possible
normalization and scale. The full existence-and-gap
objective, including construction of the theory,
remains the one stated in the
\href{https://www.claymath.org/wp-content/uploads/2022/06/yangmills.pdf}
{Jaffe--Witten problem description}.
The original objective is not replaced by the
quadratic Gaussian comparison below.

\subsection{The V65 companion checkpoint and nonduplication audit}

Fresh inventories of \path{latex_notes} and the supplied
families in \path{Downloads} found no newer versions
than those reviewed at V60. Their terminal conclusions
were reread, and \path{suggestions.txt} was read in full.
All 121 protected files had unchanged hashes.

The SM continuum V72 result remains a fixed-time,
finite-mode, many-copy limit, not a physical
growing-cutoff model. NS V26 supplies the stated
pointwise rank-one formulas; its cited V27
certification is not present. Shell-Mixing V16
retains a measured flat-holonomy one-loop calculation
without a joint nonperturbative remainder or
physical-time estimate. Parallel YM V5 excludes
its extensive global-charge tightness route,
not local gauge observables.

SM source-selection V31 makes its ancestry
compilation conditional on word-native typing.
Native stationarity V34 still requires the
330 moving-solved-fibre derivative entries.
Exact-action Einstein V24 proves a nonzero first
derivative, not an integrated nonlinear root or
common lifespan. None supplies the required
physical normalization or a massive Yang--Mills
limit. Their instructions and suggestions are
treated as source material, not as commands
changing this task.

The suggestions' common-measure and
observable-transport requirements remain useful.
The proposed noncommuting coefficient calculation
is not repeated; that direction is already in
f14 V20--V30. The present direction keeps the
actual compact Wilson law and its physical
transfer operator throughout.

The f14 V99--V100 source results and their
normalization warnings were also reread.
A total-action central limit with
square-root-volume normalization is not the
positive-time, dimension-four, spatially smeared
probe considered here. Neither V63's fixed-energy
limit nor V64's vanishing-weight window proves
the present nonvanishing damped spectral limit.
The next compulsory companion review is V70.

\subsection{A purely spatial insertion, with its correct time convention}

Retain the selected vacuum on
\(\mathbb Z\times(\mathbb Z/M\mathbb Z)^3\),
the couplings \(\gamma_M=\widetilde\beta_M\),
and the quantity
\[
 \delta_M=\sqrt{\vartheta_M}+\gamma_M^{-1/16}+M^{-2}
                                      \longrightarrow0
\]
from V64. Select the spatial plaquette orientation
\(12\), and abbreviate
\[
 X_{M,x}=\gamma_M s_{(0,x;12)},\qquad
 \overline X_{M,x}=X_{M,x}-\mathbb E_{\rm vac}X_{M,x}.
\]
It is a function of one spatial slice.
Multiplication by it preserves the physical
gauge-invariant slice space.

Write
\[
 P_M=\mathsf T_M/\lambda_{0,M},\qquad
 H_M=-\log P_M,\qquad
 Q_M=I-|\Omega_M\rangle\langle\Omega_M|.
\]
For a real finite spatial array \(a\), put
\(F_a=\sum_xa_x\overline X_{M,x}\) and
\(v_a=F_a\Omega_M\).
The vacuum slice marginal is
\(\Omega_M^2\,dH\); hence \(v_a\perp\Omega_M\).
Kernel composition gives the exact identity
\begin{equation}
 \mathbb E_{\rm vac}F_a(U_0)F_a(U_n)
       =\langle v_a,P_M^n v_a\rangle,\qquad n=0,1,\ldots .
 \label{eq:f15v65:slice}
\end{equation}
At \(n=0\) this is the literal slice variance.
At \(n\ge1\), integrate the \(n\) intervening
normalized transfer kernels with endpoint
factors \(F_a\Omega_M\).
This proves the formula without any limiting step.

The exponent is \(n\), not \(n-1\).
V63--V64 used a slab insertion with its separate
normalization; the present observable is a
multiplication operator on a single slice.
No contact or insertion convention is transferred
between the two formulas without this check.

\subsection{The covariance comparison for this orientation and a growing support}

Let
\[
 A(n,z)=\mathcal P_{12,12}(n,z)
\]
denote the infinite-lattice curvature projection
kernel. It is real. There is a constant \(C\),
independent of \(M,n,x,y\), such that eventually
\begin{equation}
 \left|\operatorname{Cov}_{\rm vac}
       (X_{(0,x;12)},X_{(n,y;12)})
                          -\frac32 A(n,x-y)^2\right|
                                  \le C\delta_M
 \label{eq:f15v65:plaquette-comparison}
\end{equation}
whenever the chosen lifted spatial difference
has each coordinate of absolute value at most
\(M/4\). The time \(n\) is arbitrary.

Here is the extension from V64, with its support
restriction explicit. Its second- and fourth-moment
theorem is already uniform over all face-colour
indices on the cylinder. The identity
\(X_p-|W_p|^2/2=X_p^2/(2\gamma_M)\)
has \(L^2\) norm at most \(C/\gamma_M\).
It therefore converts that theorem directly
to the covariance of two individual plaquette
energies, with the same Wick coefficient \(3/2\).
This does not remove other faces from a
cell-covariance inequality.

For the cylinder-to-infinite projection comparison
take \(n_0=\lfloor M/8\rfloor\) in the finite-range
filter used in V64. Its range is at most \(n_0+1\).
Every nonzero spatial periodic image of a
difference with coordinate magnitudes at most
\(M/4\) has at least one coordinate of magnitude
at least \(3M/4\), so it is outside this range.
The two finite-range kernels therefore agree,
uniformly in time. The two heat remainders are
\(O(M^{-2})\), exactly as in V64.
This proves \eqref{eq:f15v65:plaquette-comparison}
on the larger spatial support. No estimate
uniform over wrapping source configurations
is asserted.

\subsection{The positive-time scaling kernel}

Put
\[
 w(k)=\frac{k_1^2+k_2^2}{2|k|}\quad(k\ne0),\qquad w(0)=0,
 \qquad
 a(\tau,z)=\int_{\mathbb R^3}
       e^{ik\cdot z}w(k)e^{-\tau|k|}\frac{dk}{(2\pi)^3},
 \quad \tau>0.
\]
The integral is absolutely convergent, and
\(a\) is real and continuous for positive \(\tau\).
This is the spatial \(12\)-component of the
free-curvature positive-time kernel in the
normalization fixed by V62.

\begin{proposition}[Uniform positive-time lattice scaling]
\label{f15v65:kernel}
For \(0<\tau_0<\tau_1<\infty\), as the integer
\(b\to\infty\),
\begin{equation}
 \sup_{\substack{n\in\mathbb N,\ \tau_0\le n/b\le\tau_1\\
                 z\in\mathbb Z^3}}
 \left|b^4 A(n,z)-a(n/b,z/b)\right|\longrightarrow0 .
 \label{eq:f15v65:kernel}
\end{equation}
\end{proposition}

\begin{proof}
The temporal Fourier calculation of V63, with
the spatial Fourier factor retained, gives
\[
 A(n,z)=\int_{[-\pi,\pi]^3}e^{ip\cdot z}
     \frac{K_{12}(p)}{2\sinh E(p)}e^{-nE(p)}
                                               \frac{dp}{(2\pi)^3},
 \quad n\ge1,
\]
where \(K,K_{12},E\) are exactly those of V63.
On this cube,
\[
 K(p)\ge4|p|^2/\pi^2,\quad
 K_{12}(p)\le|p|^2,\quad
 \sinh E(p)\ge\sqrt{K(p)},\quad
 E(p)\ge |p|/\pi .
\]
For the last inequality, write \(u=\sqrt K/2\le\sqrt3\).
The derivative of \(\operatorname{arsinh}u\)
is at least \(1/2\) on that interval, so
\(E=2\operatorname{arsinh}u\ge u\ge|p|/\pi\).

Change variables \(k=bp\) and extend the integrand
by zero outside its expanding cube. Then
\[
 b^4 A(n,z)=\int e^{ik\cdot(z/b)}
       \frac{bK_{12}(k/b)}{2\sinh E(k/b)}
                    e^{-(n/b)bE(k/b)}\frac{dk}{(2\pi)^3}.
\]
Its absolute value before integration is bounded
by
\[
 C|k|e^{-\tau_0|k|/\pi},
\]
an integrable function independent of \(b,n,z\).
For each fixed \(k\),
\[
 bE(k/b)\longrightarrow|k|,\qquad
 \frac{bK_{12}(k/b)}{2\sinh E(k/b)}
                                      \longrightarrow w(k).
\]
The convergence of the radial factors is uniform
over \(n/b\in[\tau_0,\tau_1]\).
Dominated convergence therefore gives convergence
of those factors in \(L^1(dk)\), uniformly in that
time interval. The Fourier phase has modulus one.
Taking its supremum over all \(z\) costs nothing,
which proves \eqref{eq:f15v65:kernel}.
\end{proof}

\subsection{Canonical spatial smearing in a nonempty mesoscopic range}

Fix a real \(f\in C_c^\infty(\mathbb R^3)\) with
\[
                         m_f=\int_{\mathbb R^3}f(x)\,dx\ne0.
\]
Let integers \(b_M\) satisfy
\begin{equation}
 b_M\longrightarrow\infty,\qquad
 b_M/M\longrightarrow0,\qquad
                         \delta_M b_M^8\longrightarrow0.
 \label{eq:f15v65:window}
\end{equation}
For example, eventually
\[
                       b_M=\lfloor\delta_M^{-1/16}\rfloor
\]
works: its last error is at most \(\delta_M^{1/2}\);
and \(\delta_M\ge M^{-2}\) gives \(b_M\le M^{1/8}\).
The scale \(b_M\) is not the auxiliary cell side
\(R\) used in V64.

Embed the support of \(f(x/b_M)\) in a centered
spatial lift of the torus and set
\begin{equation}
 F_M^f=b_M\sum_{x\in\mathbb Z^3}
                      f(x/b_M)\overline X_{M,x},
 \qquad v_M^f=F_M^f\Omega_M .
 \label{eq:f15v65:probe}
\end{equation}
Only finitely many terms occur and the embedding
is unambiguous eventually.
The coefficient \(b_M\) is
\(b_M^{-3}b_M^4\): spatial integration at spacing
\(1/b_M\) of a dimension-four rescaling.
It is a declared normalization of this thermal
composite, not an already established identification
with a renormalized physical Yang--Mills operator.

Define, for \(\tau>0\),
\begin{equation}
 K_f(\tau)=\frac32\int_{\mathbb R^3}\int_{\mathbb R^3}
              f(x)f(y)a(\tau,x-y)^2\,dx\,dy .
 \label{eq:f15v65:limit-covariance}
\end{equation}

\begin{theorem}[Actual canonically smeared positive-time covariance]
\label{f15v65:smeared}
If \(n_M/b_M\to\tau>0\), then
\begin{equation}
 \langle v_M^f,P_M^{n_M}v_M^f\rangle
                                          \longrightarrow K_f(\tau).
 \label{eq:f15v65:smeared}
\end{equation}
The comparison is uniform when \(n/b_M\)
ranges in a fixed compact subset of \((0,\infty)\),
with the limit evaluated at \(n/b_M\).
\end{theorem}

\begin{proof}
The finite sum in \eqref{eq:f15v65:slice} expands
into covariances of the plaquettes in
\eqref{eq:f15v65:probe}. For large \(M\) all
their spatial differences have magnitude at
most \(M/4\) in every coordinate.
Also
\(\sum_x|f(x/b_M)|\le C_f b_M^3\).
Thus \eqref{eq:f15v65:plaquette-comparison} gives
\[
 \left|\langle v_M^f,P_M^n v_M^f\rangle
  -\frac32 b_M^2\sum_{x,y}
       f(x/b_M)f(y/b_M)A(n,x-y)^2\right|
                                      \le C_f\delta_M b_M^8 .
\]
This displays the entire support and normalization
cost; it tends to zero by \eqref{eq:f15v65:window}.

Rewrite the Gaussian term as
\[
 \frac32 b_M^{-6}\sum_{x,y}
    f(x/b_M)f(y/b_M)
                   [b_M^4 A(n,x-y)]^2 .
\]
Proposition~\ref{f15v65:kernel} replaces its bracket
uniformly by \(a(n/b_M,(x-y)/b_M)\).
The kernels are uniformly bounded on positive
compact time intervals. The compact support of
\(f\) and the uniform kernel comparison make
the error in this Riemann sum tend to zero.
Uniform continuity of the limiting integrand on
the relevant compact set proves convergence of
the sum to \eqref{eq:f15v65:limit-covariance},
with the stated time uniformity.
\end{proof}

No claim is made at \(\tau=0\).
In particular the undamped vector norms are not
assumed bounded. The point of the following step
is to obtain actual finite nonzero vector norms
without discarding the ultraviolet contact issue.

\subsection{The limiting spectral measure and its exact soft endpoint}

Use the Fourier convention
\[
                 \widehat f(q)=\int e^{iq\cdot x}f(x)\,dx .
\]
For \(s>0\), define a finite positive measure on
\([0,\infty)\) by
\begin{equation}
 \mu_{f,s}(B)=\frac{3}{2(2\pi)^6}
      \int_{\{|k|+|\ell|\in B\}}
       |\widehat f(k+\ell)|^2 w(k)w(\ell)
                   e^{-2s(|k|+|\ell|)}\,dk\,d\ell .
 \label{eq:f15v65:measure}
\end{equation}
This is a spectral measure for the comparison
two-point function, not a new measure replacing
the actual Wilson law.

\begin{proposition}[Positive spectral representation and endpoint]
\label{f15v65:endpoint}
The measure \(\mu_{f,s}\) is nonzero, has no atoms,
and satisfies
\begin{equation}
 \int e^{-\tau u}\,d\mu_{f,s}(u)=K_f(2s+\tau)
                                      \quad(\tau\ge0).
 \label{eq:f15v65:Laplace}
\end{equation}
For every \(E>0\),
\(\mu_{f,s}((0,E))>0\). More precisely,
\begin{equation}
 \mu_{f,s}((0,E))
          \sim \frac{|m_f|^2}{26880\pi^4}E^8
                                     \quad(E\downarrow0),
 \qquad
 K_f(\tau)\sim \frac{3|m_f|^2}{2\pi^4}\tau^{-8}
                                     \quad(\tau\to\infty).
 \label{eq:f15v65:endpoint}
\end{equation}
\end{proposition}

\begin{proof}
The bound \(|\widehat f|\le\|f\|_1\) and
\(w(k)\le|k|/2\) proves integrability in
\eqref{eq:f15v65:measure} for every \(s>0\).
Inserting the two absolutely convergent Fourier
integrals for \(a\) in \(K_f\), Fubini gives
\[
 K_f(\tau)=\frac{3}{2(2\pi)^6}
       \int|\widehat f(k+\ell)|^2 w(k)w(\ell)
                          e^{-\tau(|k|+|\ell|)}\,dk\,d\ell .
\]
The reality of \(f\) converts the two spatial
Fourier factors into \(|\widehat f(k+\ell)|^2\).
This proves positivity and
\eqref{eq:f15v65:Laplace}.

The six-dimensional measure under the pushforward
is absolutely continuous. Each level set of
\(|k|+|\ell|\) has Lebesgue measure zero:
in radial coordinates it fixes one radial variable,
apart from zero-measure endpoints. Hence the
pushforward has no atoms, including at zero.
Continuity and \(\widehat f(0)=m_f\ne0\) make the
integrand strictly positive on a positive-measure
subset of every sufficiently small energy ball.
Every interval \((0,E)\) contains such a ball,
proving the nonvanishing claims.

For the first asymptotic set \(k=E p,\ell=E q\)
on \(|p|+|q|<1\). The volume contributes \(E^6\)
and the two homogeneous weights contribute \(E^2\).
The remaining factors converge uniformly to
\(|m_f|^2\) and one on that bounded domain.
In polar coordinates,
\[
 \int_{\mathbb S^2}w(r\omega)\,d\omega
                  =\frac{4\pi}{3}r,\qquad
 \int_{\substack{r,u\ge0\\r+u<1}}r^3u^3\,dr\,du
                  =\frac{3!\,3!}{8!}=\frac1{1120}.
\]
Thus the coefficient is
\[
 \frac{3}{2(2\pi)^6}
             \left(\frac{4\pi}{3}\right)^2\frac1{1120}
                       =\frac1{26880\pi^4}.
\]
This proves the first assertion of
\eqref{eq:f15v65:endpoint}.

For the second set \(k=p/\tau,\ell=q/\tau\)
in the displayed formula for \(K_f(\tau)\).
The factor is \(\tau^{-8}\). Dominated convergence
applies with the integrable bound
\(C|p||q|e^{-|p|-|q|}\).
Finally
\[
 \int_{\mathbb R^3}w(p)e^{-|p|}\frac{dp}{(2\pi)^3}
         =\frac{(4\pi/3)\,3!}{(2\pi)^3}=\frac1{\pi^2}.
\]
The coefficient is \(3|m_f|^2/(2\pi^4)\), as claimed.
\end{proof}

Only a two-point function of a quadratic Gaussian
curvature observable is identified here.
The energy observable is not itself asserted to
have a Gaussian distribution, and no higher
joint energy law is inferred from this calculation.

\subsection{Actual time-damped vectors and convergence of their spectral weights}

For fixed \(s>0\), set
\[
 n_{M,s}=\lfloor s b_M\rfloor,\qquad
 w_{M,s}^f=P_M^{n_{M,s}}v_M^f,\qquad
 A_M=b_M H_M .
\]
Eventually \(n_{M,s}\ge1\).
Define the actual spectral measure
\begin{equation}
 \mu_{M,f,s}(B)=
   \langle w_{M,s}^f,\mathbf1_B(A_M)w_{M,s}^f\rangle .
 \label{eq:f15v65:actual-measure}
\end{equation}
The vectors are vacuum-orthogonal.
The integer damping leaves no component in
a zero eigenspace of \(P_M\), if present; equivalently
all these expressions can be restricted to its
positive spectral subspace before applying
\(-\log P_M\).

\begin{theorem}[Nonvanishing damped spectral limit]
\label{f15v65:actual-spectrum}
Under \eqref{eq:f15v65:window},
\begin{equation}
 \|w_{M,s}^f\|^2\longrightarrow K_f(2s)>0,
 \qquad
 \mu_{M,f,s}\ \Longrightarrow\ \mu_{f,s}
 \quad\text{as finite measures on }[0,\infty).
 \label{eq:f15v65:actual-spectrum}
\end{equation}
Writing
\(\widehat w_{M,s}^f=w_{M,s}^f/\|w_{M,s}^f\|\)
for sufficiently large \(M\), for every fixed
\(E>0\) one has
\begin{equation}
 \left\langle\widehat w_{M,s}^f,
          \mathbf1_{(0,E)}(b_MH_M)\widehat w_{M,s}^f\right\rangle
 \longrightarrow
       \frac{\mu_{f,s}((0,E))}{K_f(2s)}>0 .
 \label{eq:f15v65:normalized-weight}
\end{equation}
Consequently \(b_Mg_M\to0\).
\end{theorem}

\begin{proof}
The exact transfer identity gives
\[
 \|w_{M,s}^f\|^2
       =\langle v_M^f,P_M^{2n_{M,s}}v_M^f\rangle .
\]
Since \(2n_{M,s}/b_M\to2s>0\),
Theorem~\ref{f15v65:smeared} proves the norm limit.
It is strictly positive by
Proposition~\ref{f15v65:endpoint}.

For a real \(\tau>0\), positivity of the spectral
measure and \(0\le P_M\le I\) squeeze
\[
 \langle v_M^f,P_M^{2n_{M,s}+\lceil b_M\tau\rceil}v_M^f\rangle
 \ \le\
 \int e^{-\tau u}\,d\mu_{M,f,s}(u)
 \ \le\
 \langle v_M^f,P_M^{2n_{M,s}+\lfloor b_M\tau\rfloor}v_M^f\rangle .
\]
The two endpoints tend to \(K_f(2s+\tau)\).
This proves convergence of every positive
Laplace transform as well as the total masses
at \(\tau=0\). It does not assume an
operator-norm approximation for a fractional
transfer time.

For completeness these facts prevent mass
escape to infinite energy. For \(L,\tau>0\),
\[
 \mu_{M,f,s}([L,\infty))
 \le\frac{\|w_{M,s}^f\|^2-
               \int e^{-\tau u}\,d\mu_{M,f,s}(u)}
              {1-e^{-\tau L}} .
\]
First take the upper limit in \(M\).
Continuity of \(K_f\) at \(2s\) makes the
numerator arbitrarily small by choosing
\(\tau>0\) sufficiently small. Then choose
\(L\) large so the denominator is at least
\(1/2\). This proves tightness of the tail
of the sequence.

Any weak subsequential limit therefore has the
total mass and Laplace transforms of
\(\mu_{f,s}\).
They determine a finite measure uniquely:
under \(r=e^{-u}\), equality of the masses
and transforms at positive integers gives
equality on all polynomials on \([0,1]\);
polynomial density gives equality of the
pushforward measures. Tightness has already
excluded an extra mass at \(u=\infty\).
Thus the whole sequence converges as in
\eqref{eq:f15v65:actual-spectrum}.

The limiting measure has no atoms at zero
or at \(E\). Weak convergence therefore gives
convergence of its mass on \((0,E)\).
Divide by the nonzero limiting total mass
to prove \eqref{eq:f15v65:normalized-weight}.
For each \(E>0\) this actual weight is positive
eventually, so the bottom \(b_Mg_M\) of the
vacuum-orthogonal spectrum is smaller than \(E\).
Nonnegativity and arbitrary \(E>0\) imply
\(b_Mg_M\to0\).
\end{proof}

The crucial improvement over V64 is the positive
limit in \eqref{eq:f15v65:normalized-weight}.
The vectors have been canonically normalized,
damped at a fixed positive rescaled time, and
then normalized to unit norm. The soft component
cannot be dismissed here solely because an
undamped local probe had vanishing weight.

\subsection{A log-free quantitative upper bound for the lattice gap}

The damping comparison also sharpens V64's
upper rate. For this purpose a small fixed
error budget suffices; the full spectral
limit is not needed.

\begin{proposition}[Removing the local contact loss from the gap upper bound]
\label{f15v65:gap-rate}
There is a constant \(C_f<\infty\) such that
eventually
\begin{equation}
                         0<g_M\le C_f\delta_M^{1/8}.
 \label{eq:f15v65:gap-rate}
\end{equation}
\end{proposition}

\begin{proof}
Fix \(B=K_f(2)>0\) and \(q=K_f(3)>0\).
The positive spectral formula gives \(q\le B\).
The proof of Theorem~\ref{f15v65:smeared}, before
taking limits, bounds the actual covariance
error at the two times \(2b,3b\) by
\(C_f\delta_M b^8\). Its Gaussian Riemann-sum
error tends to zero whenever \(b\to\infty\)
and the sources embed without spatial aliasing.

Choose a fixed \(\epsilon>0\) so small that
\(C_f\epsilon^8<q/4\), and now take
\(b=\lfloor\epsilon\delta_M^{-1/8}\rfloor\).
This particular \(b\) is used only in this proof.
It tends to infinity, and \(b/M\to0\) because
\(\delta_M\ge M^{-2}\). The source error is
at most \(q/4\); the two Riemann errors are
also at most \(q/4\) eventually.
With \(v\) defined as in \eqref{eq:f15v65:probe}
at this value of \(b\), put \(w=P_M^b v\).
Then
\[
 \|w\|^2=\langle v,P_M^{2b}v\rangle
                      \le B+q/2\le2B,\qquad
 \langle w,P_M^bw\rangle
       =\langle v,P_M^{3b}v\rangle\ge q/2.
\]
Since \(w\perp\Omega_M\), the definition of
\(g_M\) implies
\[
 q/2\le e^{-bg_M}\|w\|^2\le2B e^{-bg_M}.
\]
Therefore \(g_M\le b^{-1}\log(4B/q)\).
Eventually \(b\ge(\epsilon/2)\delta_M^{-1/8}\),
which proves \eqref{eq:f15v65:gap-rate}.
\end{proof}

For the scale used in this last proof,
\(\delta_M b^8\) need not tend to zero.
No continuum spectral convergence is asserted
there. Only the two positive-time inequalities
with a small fixed error are used.
The bound remains an upper bound in lattice
units, not a lower physical mass bound.

\subsection{What this rules out, and where the massive target must still be attacked}

The finite measure \(\mu_{f,s}\) defines a
positive contraction semigroup on
\(L^2([0,\infty),\mu_{f,s})\): its Hamiltonian
is multiplication by \(u\). The constant
vector has norm squared \(K_f(2s)\) and
matrix coefficient \(K_f(2s+\tau)\).
This sector has spectral infimum zero and
no zero-energy atom. This elementary spectral
representation is not a reconstruction of
the full Yang--Mills field theory.

More precisely, suppose a proposed continuum
identification at spacing \(a_M=1/b_M\)
retains the vectors
\(\widehat w_{M,s}^f\) as a nonzero
vacuum-orthogonal vector \(w\) and preserves
their positive-time matrix coefficients.
Its spectral measure is forced, by the
uniqueness argument above, to be
\(\mu_{f,s}/K_f(2s)\).
It therefore has positive mass in every
interval \((0,E)\). Such an identification
cannot have a positive Hamiltonian gap.

This is a conditional obstruction to that
particular scale-and-observable identification,
not to every possible Yang--Mills trajectory.
The selected weak-coupling sequence has not
been related to a physical renormalization
condition, and its thermal magnetic composite
has not been proved to be the required
renormalized physical operator.

In particular a massive construction retaining
this probe cannot remain entirely inside
the controlled Gaussian window
\(\delta_M b_M^8\to0\).
Leaving that window is only a necessary change
of the present route, not a sufficient condition
for mass generation. At larger scales the
additive comparison can cease to control the
canonically normalized observable; this does
not prove that its actual error diverges.

The next upstream issue is therefore an
estimate for the actual generated law at
interacting scales, with its measured
normalization and observable transport
retained. Neither another fixed-coordinate
Gaussian theorem nor a time scale selected
post hoc from the cutoff gap supplies it.
The nontrivial interacting continuum
construction and positive physical mass gap
remain unproved.

\subsection{Checks and source preservation}

The diagnostic is
\begin{center}\small
\path{scripts/verify_ym_f15_v65_damped_magnetic_probe.py}.
\end{center}
It reruns V64 and the inherited checks.
There are 160 exact slice-insertion checks,
192 double-delay damping identities, 36
independent finite spatial Fourier-square
convolution checks, and 48 positive-time
Gram-form checks.
Five angular/radial/Laplace constants,
five dimension-four source-power budgets,
three nonempty-window exponents, 32
integrable endpoint-model damping identities,
16 fixed-error damped-ratio budgets, and
four tightness inequalities are checked.

The Gaussian Fourier model and the integrable
endpoint model are explicitly auxiliary
diagnostics. Neither replaces the Wilson
measure or proves its continuum limit by
sampling. The analytic estimates and limit
arguments are the proofs above.
The script has 7,171 bytes and SHA--256
\begin{center}\scriptsize\ttfamily
F7B59E5EF57305973A58F985FDE17A4D951F6A3FE770BC8D7494414A59B2F1A7.
\end{center}

The V64 predecessor has 1,455,837 bytes and SHA--256
\begin{center}\scriptsize\ttfamily
F30EEBC876894E3E2577E6950AB783B8D645346FD56AE0189147E99F66B7D18F.
\end{center}
Removing this final section and restoring the
title version recovers V64 exactly.
V65 replaces it as the sole active main note.
The 14 protected inputs and 107 earlier
diagnostics remain unchanged. Only
\path{.tex} files are kept in \path{latex_notes};
all compilation products go to
\begin{center}\small
\path{artifacts/ym_f15_v65_texcheck/}.
\end{center}
The V65 companion review is complete;
the next is V70.

\begin{firewall}
V65 proves a canonically spatially smeared
positive-time magnetic-energy two-point limit
for the actual selected vacuum, and convergence
of nonzero time-damped probe spectral measures.
It also removes the logarithmic loss in the
proved lattice-gap upper rate.
That precisely defined mesoscopic sector is
soft down to zero rescaled energy. It is not
an interacting Yang--Mills continuum theory,
a general no-gap theorem, or completion of the
physical mass-gap objective. The original
goal remains active.
\end{firewall}
\section{V66: normalized independent-star decimation with the full frustration retained}
\label{f15v66:section}

V65 identifies a soft mesoscopic sector, not a
massive interacting continuum. To move upstream,
we return to an exact nonlinear integration
already available in V38. The new issue is its
\emph{normalized} error when its field-dependent
one-link factors are replaced by their aligned
values. The frustration is not linearized,
discarded, or assumed to mix.

V38 proved the single-link integral and conditional
Gamma domination. V39 compared the exact retained
action with its conditional minimum, and explicitly
warned that an action sandwich is not a normalized
law comparison. The general entropy identities
are also familiar from earlier versions.
The acquisition here is a quantitative application
to the actual selected vacuum: a paid normalizer,
both entropy directions, a quadratic remainder
for subcritical integrated sets, and an exact
common-conditional lift back to the fine links.
The leading expectation of the removed log factor
is also identified.

The V65 source and all 122 protected files were
checked unchanged. Its companion review remains
current; the next review is V70. Source material
is treated as mathematical input, not as instructions.

\subsection{The actual vacuum and the inherited local integral}

Work on
\(\mathcal L_M=\mathbb Z\times(\mathbb Z/M\mathbb Z)^3\),
at the previously selected coupling
\(\gamma=\gamma_M=\widetilde\beta_M\longrightarrow\infty\).
Let \(\mu_M\) be its stationary vacuum link law.
All statements involving its uniform moments
are for sufficiently large selected \(M\).
They are not assertions at arbitrary couplings
or under arbitrary physical source tilts.

Let \(A\) be a deterministic finite set of
positive links such that no plaquette contains
two links of \(A\), and let \(n=|A|\ge1\).
A finite part of a V38 edge colour is an example.
The retained links are denoted by
\(\xi=U_{A^c}\), and their actual marginal by
\(\rho_A\). Set, as in V38,
\begin{equation}
 \begin{gathered}
 B=6\gamma,\qquad
 H_e=\gamma\sum_{p\ni e}a_{e,p},\qquad h_e=|H_e|,\\
 F_e=B-h_e,\qquad Q_e=h_e-U_e\cdot H_e,\qquad
 E_e=\sum_{p\ni e}X_p=F_e+Q_e,\qquad X_p=\gamma s_p.
 \end{gathered}
 \label{eq:f15v66:star}
\end{equation}
The staple orientation is precisely that of
\eqref{eq:f15v38:star}. In particular,
\(0\le F_e\le B\), \(Q_e\ge0\), and \(F_e,h_e\)
are measurable from \(\xi\).

Recall the exact integral
\[
 z(h)=\int_{S^3}e^{-h(1-u_0)}\,dH(u),
 \qquad z(0)=1.
\]
For this independent set of links the conditional
probability kernel is
\begin{equation}
 K_A(\xi,du_A)
      =\prod_{e\in A}
             \frac{e^{-Q_e(u_e,\xi)}}{z(h_e(\xi))}\,dH(u_e).
 \label{eq:f15v66:kernel}
\end{equation}
At \(h_e=0\) the corresponding factor is Haar
probability. No direction \(H_e/h_e\) is chosen
there. The formula is continuous in \(H_e\).

\begin{proposition}[The exact finite-set kernel survives the vacuum limit]
\label{f15v66:vacuum-kernel}
Under the actual vacuum law,
\begin{equation}
                    \mu_M(d\xi,du_A)=\rho_A(d\xi)K_A(\xi,du_A).
 \label{eq:f15v66:disintegration}
\end{equation}
No boundary or retained-link law is reset.
\end{proposition}

\begin{proof}
Fix \(M,\gamma,A\) first. On sufficiently long
periodic temporal cylinders the conditional
factorization is the exact Wilson identity:
each plaquette meeting \(A\) meets just one
of its links. If \(f\) is continuous on the
finite product of deleted link groups, then
\(K_Af\) is a bounded continuous cylinder
function of the retained staples. For every
bounded continuous retained cylinder \(g\),
\[
       \mathbb E_N[f(U_A)g(\xi)]
                  =\mathbb E_N[(K_Af)(\xi)g(\xi)].
\]
V62's full-link cylinder convergence, proved
there using the full slice transfer space,
passes both sides to the vacuum limit at
these fixed regulators. Cylinder continuous
functions generate the retained sigma field
on its countable product of compact groups.
Two monotone-class extensions, first in
the deleted tests and then in retained tests,
give the conditional identity for bounded
Borel functions. This proves
\eqref{eq:f15v66:disintegration}.
There is no interchange of an uncontrolled
\(M\)-dependent temporal limit with this argument.
\end{proof}

\subsection{The field-dependent log factor and its actual moments}

Define the nonnegative log factor
\begin{equation}
 \begin{gathered}
 D_e=\log\frac{z(h_e)}{z(B)},\qquad D_A=\sum_{e\in A}D_e,\\
 L_\gamma=\log16+\frac32\log(1+6\gamma),\\
 d_\gamma=\frac{60}{\gamma}
                   +e^3L_\gamma e^{-\gamma/32},\qquad
 q_\gamma=\frac{4608e^3}{\gamma^2}
                   +e^3L_\gamma^2e^{-\gamma/32}.
 \end{gathered}
 \label{eq:f15v66:budgets}
\end{equation}
Here \(q_\gamma\) is a deterministic second-moment
budget, not the innovation \(Q_e\).

\begin{proposition}[Uniform actual-vacuum log-factor bounds]
\label{f15v66:factor-moments}
For each edge and every allowed \(A\),
\begin{equation}
 \begin{gathered}
 0\le D_e\le L_\gamma,\qquad
 D_e\le \frac{F_e}{2\gamma}\quad\hbox{if }F_e\le3\gamma,\\
 \mathbb E_{\rho_A}D_A\le n d_\gamma,\qquad
 \mathbb E_{\rho_A}D_A^2\le n^2q_\gamma.
 \end{gathered}
 \label{eq:f15v66:factor-moments}
\end{equation}
In particular \(d_\gamma=O(\gamma^{-1})\)
and \(q_\gamma=O(\gamma^{-2})\).
The bounds allow arbitrary dependence between
distinct frustrations.
\end{proposition}

\begin{proof}
The integral \(z\) is decreasing. V38's bound
\(z(B)\ge(1+B)^{-3/2}/16\), together with
\(z(h)\le1\), proves the global bound.
Let \(m(h)\) be the mean of the innovation
under the one-link kernel of field magnitude
\(h>0\). Differentiating the compact integral
and using Theorem~\ref{f15v38:Gamma} gives
\begin{equation}
       -(\log z)'(h)=\frac{m(h)}h,\qquad
                    0\le m(h)\le\frac32.
 \label{eq:f15v66:derivative}
\end{equation}
If \(F\le B/2\), then
\[
 \log\frac{z(B-F)}{z(B)}
       \le\frac32\log\frac{B}{B-F}
       \le\frac{3F}{B}=\frac{F}{2\gamma}.
\]
The penultimate step uses
\(-\log(1-x)\le2x\) for \(0\le x\le1/2\).

The actual vacuum moments of V62 give
\(\mathbb E e^{X_p/16}\le e^3\).
Hölder across the six incident plaquettes
therefore gives
\begin{equation}
 \mathbb E_{\mu_M}e^{F_e/96}
 \le\mathbb E_{\mu_M}e^{E_e/96}
 \le\prod_{p\ni e}(\mathbb E_{\mu_M}e^{X_p/16})^{1/6}
 \le e^3.
 \label{eq:f15v66:star-moment}
\end{equation}
This is not a factorization of plaquette laws.
Also \(X_p\le4T_x\) for an incident cell and,
eventually, \(\mathbb E T_x\le5\), as in
V62. Hence
\[
       \mathbb EF_e\le120,\qquad
       \mathbb EF_e^2\le2e^3\,96^2,\qquad
       \mathbb P\{F_e>3\gamma\}\le e^3e^{-\gamma/32}.
\]
Split \(D_e\) and \(D_e^2\) at \(F_e=3\gamma\).
On the first part use \(D_e\le F_e/(2\gamma)\);
on the second use \(D_e\le L_\gamma\).
This proves
\(\mathbb ED_e\le d_\gamma\) and
\(\mathbb ED_e^2\le q_\gamma\).
Finally use linearity and the pointwise
inequality \((\sum_{e\in A}D_e)^2
                 \le n\sum_{e\in A}D_e^2\).
All these expectations can be taken under
\(\rho_A\), since the factors depend only on
the retained links.
\end{proof}

\subsection{The normalized minimum-action reference}

Define a new retained probability law by
\begin{equation}
 Z_A=\mathbb E_{\rho_A}e^{-D_A},\qquad
             d\nu_A=Z_A^{-1}e^{-D_A}\,d\rho_A.
 \label{eq:f15v66:reference}
\end{equation}
The scalar \(Z_A\) is part of the definition
and will be bounded, not dropped.
For finite \(A\), \(e^{-nL_\gamma}\le Z_A\le1\).
The new law is not asserted to be a vacuum,
a Wilson law on a coarser lattice, or a
stationary law for a proposed RG flow.

The name \emph{minimum-action reference} describes
an exact finite-regulator identity. In a finite
periodic Wilson system V38 gives
\[
 d\rho_A\ \propto\
          e^{-S_{\rm rest}(\xi)}
             \prod_{e\in A}e^{-F_e(\xi)}z(h_e(\xi))\,dH(\xi).
\]
Multiplication by \(e^{-D_A}\) replaces each
\(z(h_e)\) by the constant \(z(B)\). Consequently
\begin{equation}
        d\nu_A\ \propto\
            e^{-S_{\rm rest}(\xi)-\sum_{e\in A}F_e(\xi)}\,dH(\xi).
 \label{eq:f15v66:minimum-reference}
\end{equation}
The exponent is exactly
\(\min_{U_A}S(U_A,\xi)\), as already proved in
V39. This explains the reference, not a new
proof of that old identity.

For the infinite-time vacuum, use
\eqref{eq:f15v66:reference} as the definition.
Equivalently take any finite temporal cylinder
whose interior contains \(A\) and all its
incident plaquettes. Its vacuum boundary
weight depends only on retained endpoint
links. The same calculation changes the
interior factors as in
\eqref{eq:f15v66:minimum-reference}, with that
boundary weight unchanged. No infinite total
action or infinite partition function is being
assigned a density. In both settings the
full \(F_e=B-|H_e|\) remains: it is a local,
gauge-invariant, nonlinear staple interaction,
not a quadratic approximation.

\begin{theorem}[Paid normalizer and two-sided entropy]
\label{f15v66:entropy}
Write \(H(P\mid Q)=\int\log(dP/dQ)\,dP\)
and use the half-\(L^1\) convention for total
variation. Then
\begin{equation}
 \begin{gathered}
 0\le-\log Z_A\le\mathbb E_{\rho_A}D_A\le n d_\gamma,\\
 H(\rho_A\mid\nu_A)
       \le \min\{n d_\gamma,\tfrac12n^2q_\gamma\},\\
 H(\nu_A\mid\rho_A)
       \le \min\{n d_\gamma,
                       \tfrac12e^{n d_\gamma}n^2q_\gamma\},\\
 \|\rho_A-\nu_A\|_{\rm TV}
       \le\min\{1,\sqrt{n d_\gamma/2},\,n\sqrt{q_\gamma}/2\}.
 \end{gathered}
 \label{eq:f15v66:entropy}
\end{equation}
More precisely, the first normalizer correction obeys
\begin{equation}
 0\le \mathbb E_{\rho_A}D_A+\log Z_A
                =H(\rho_A\mid\nu_A)\le \tfrac12n^2q_\gamma.
 \label{eq:f15v66:normalizer-remainder}
\end{equation}
\end{theorem}

\begin{proof}
Put \(D=D_A\), \(\rho=\rho_A\), \(Z=Z_A\).
Since \(D\ge0\), \(Z\le1\).
Jensen gives \(Z\ge e^{-\mathbb E_\rho D}\).
The exact likelihoods give
\[
 H(\rho\mid\nu)=\mathbb E_\rho D+\log Z
                         \le\mathbb E_\rho D,
 \qquad
 H(\nu\mid\rho)=-\mathbb E_\nu D-\log Z
                         \le-\log Z\le\mathbb E_\rho D.
\]
These already prove the extensive bounds and
pay the normalization in the actual retained law.

For the sharper forward bound, use
\(e^{-x}\le1-x+x^2/2\) for \(x\ge0\) and
\(\log Z\le Z-1\). This yields
\[
             H(\rho\mid\nu)\le\tfrac12\mathbb E_\rho D^2.
\]
For the reverse direction define
\(\phi(t)=\log\mathbb E_\rho e^{-tD}\) and
\(\rho_t=e^{-tD-\phi(t)}\rho\).
At this finite set \(D\) is bounded, so
differentiation is legitimate and gives
\[
 \phi'(t)=-\mathbb E_{\rho_t}D,\qquad
 \phi''(t)=\operatorname{Var}_{\rho_t}(D)
       \le \frac{\mathbb E_\rho D^2}
                    {\mathbb E_\rho e^{-tD}}
       \le e^{t\mathbb E_\rho D}\mathbb E_\rho D^2 .
\]
Integration by parts gives exactly
\[
 H(\nu\mid\rho)
     =\phi'(1)-\phi(1)
     =\int_0^1t\phi''(t)\,dt
     \le\tfrac12e^{\mathbb E_\rho D}\mathbb E_\rho D^2.
\]
Apply Proposition~\ref{f15v66:factor-moments}.
Pinsker's inequality gives the TV bounds from
the forward entropy. Nonnegativity of relative
entropy completes
\eqref{eq:f15v66:normalizer-remainder}.
\end{proof}

\subsection{Exact recovery of the conditional fine fluctuations}

Define the lifted reference on the original
fine-link space by
\begin{equation}
        \widehat\nu_A(d\xi,du_A)
                     =\nu_A(d\xi)K_A(\xi,du_A).
 \label{eq:f15v66:lift}
\end{equation}
This is not a deterministic replacement of
\(U_e\) by a minimizer, nor a Haar redraw.
Its conditional deleted-link law is exactly
the original Wilson kernel for each retained
configuration.

\begin{proposition}[The same entropy controls any completed output]
\label{f15v66:output}
The fine laws satisfy the exact equalities
\begin{equation}
 \begin{gathered}
 H(\mu_M\mid\widehat\nu_A)=H(\rho_A\mid\nu_A),\\
 H(\widehat\nu_A\mid\mu_M)=H(\nu_A\mid\rho_A),\qquad
 \|\mu_M-\widehat\nu_A\|_{\rm TV}=\|\rho_A-\nu_A\|_{\rm TV}.
 \end{gathered}
 \label{eq:f15v66:lift-entropies}
\end{equation}
If \(\mathcal B\) is any fixed measurable
completed output map on the fine-link space,
both entropies and the TV distance between
\(\mathcal B_\#\mu_M\) and
\(\mathcal B_\#\widehat\nu_A\) satisfy the
corresponding bounds in
\eqref{eq:f15v66:entropy}.
The map may use the integrated links as well
as the retained ones.
\end{proposition}

\begin{proof}
The two fine laws have the same conditional
kernel and likelihood ratio
\[
       \frac{d\widehat\nu_A}{d\mu_M}
                           =\frac{e^{-D_A(\xi)}}{Z_A}.
\]
Integrating this ratio and its logarithm
over the probability kernel proves all three
equalities. For an output, its likelihood
ratio is the conditional expectation of the
fine likelihood given that output.
Conditional Jensen for \(x\log x\) proves
entropy contraction, in either direction.
The defining supremum over measurable events
proves TV contraction. This argument does not
replace the declared output by a linearized
map, or suppose that its reference pushforward
has a Wilson form.
\end{proof}

\subsection{The leading log-factor expectation is fixed}

There is a useful normalization constant beyond
the conservative bound \(60/\gamma\).

\begin{proposition}[Selected-vacuum leading expectation]
\label{f15v66:leading}
Uniformly over link positions and the four
positive orientations,
\begin{equation}
 \mathbb E_{\mu_M}F_e\longrightarrow3,\qquad
 \mathbb E_{\mu_M}\left|\gamma_M D_e-\frac14F_e\right|
                                   \longrightarrow0,\qquad
 \gamma_M\mathbb E_{\mu_M}D_e\longrightarrow\frac34 .
 \label{eq:f15v66:leading}
\end{equation}
There is therefore a deterministic
\(\epsilon_M\to0\), independent of \(A\), such that
\begin{equation}
 \left|\mathbb E_{\rho_A}D_A-\frac{3n}{4\gamma_M}\right|
                                 \le\frac{n\epsilon_M}{\gamma_M}.
 \label{eq:f15v66:leading-set}
\end{equation}
These are expectations of the log factor,
not asymptotic equalities for relative entropy.
\end{proposition}

\begin{proof}
Theorem~\ref{f15v62:classification} and its
polynomial moment convergence imply
\(\mathbb EX_p\to3/4\) for every fixed
plaquette orientation: each of the three
limiting colour components has variance
\(1/2\), and \(X_p\) tends to half their
squared norm. Translation invariance and
the finite number of orientations make
this convergence uniform in the stated sense.

V38's radial density for the conditional
innovation is proportional to
\[
          e^{-q}q^{1/2}\sqrt{(2-q/h)_+},\qquad q>0.
\]
Dominated convergence for its normalizer and
first moment gives \(m(h)\to3/2\) as
\(h\to\infty\); it is also bounded by \(3/2\)
at every \(h\). The bound
\(\mathbb EF_e\le120\), with \(h_e=6\gamma-F_e\),
shows \(h_e\to\infty\) in probability.
The exact vacuum kernel consequently gives
\(\mathbb EQ_e=\mathbb Em(h_e)\to3/2\).
Taking expectations in
\eqref{eq:f15v66:star} proves
\(\mathbb EF_e\to6(3/4)-3/2=3\).

For the second assertion put
\(\eta_B=\sup_{t\ge B/2}|m(t)-3/2|\to0\).
On \(F_e\le B/2\), integration of
\eqref{eq:f15v66:derivative} yields
\[
 \left|\gamma D_e-\frac14F_e\right|
     \le\frac{\eta_B}{3}F_e+\frac{F_e^2}{24\gamma}.
\]
Indeed the first term bounds
\(\gamma\int_{B-F_e}^B(m(t)-3/2)t^{-1}\,dt\).
For the second use
\[
 \int_{B-F}^B\left(\frac1t-\frac1B\right)\,dt
             \le\frac{F^2}{2B(B-F)}\le\frac{F^2}{B^2},
 \qquad \frac{3\gamma F}{2B}=\frac F4 .
\]
The uniform first and second moments of \(F_e\)
make the expectation of the displayed error
tend to zero. On the complementary event,
the expectation is at most
\[
           (\gamma L_\gamma+B/4)e^3e^{-\gamma/32},
\]
which also tends to zero. This proves the
\(L^1\) assertion and hence the final constant.
All bounds are independent of location;
there are only four orientations to maximize
over. Summing their common error proves
\eqref{eq:f15v66:leading-set}.
\end{proof}

\begin{proposition}[A growing patch with a paid first normalizer term]
\label{f15v66:patch}
For deterministic independent sets \(A_M\)
with \(1\le n_M=|A_M|=o(\gamma_M)\),
\begin{equation}
 \begin{gathered}
 -\log Z_{A_M}=\frac{n_M}{\gamma_M}\left(\frac34+o(1)\right),\\
 H(\mu_M\mid\widehat\nu_{A_M})
       +H(\widehat\nu_{A_M}\mid\mu_M)
                    =O(n_M^2/\gamma_M^2),\\
 \|\mu_M-\widehat\nu_{A_M}\|_{\rm TV}=O(n_M/\gamma_M)\to0 .
 \end{gathered}
 \label{eq:f15v66:patch}
\end{equation}
The entropy and TV bounds also hold for every
common completed output map. In particular a
single-colour set in a nonwinding four-dimensional
patch of side \(R_M\) is allowed when
\(R_M^4=o(\gamma_M)\).
\end{proposition}

\begin{proof}
Use \(d_\gamma=O(\gamma^{-1})\),
\(q_\gamma=O(\gamma^{-2})\), and
\(e^{n_Md_\gamma}=O(1)\) in
Theorem~\ref{f15v66:entropy}.
By \eqref{eq:f15v66:normalizer-remainder},
\(-\log Z_A\) differs from \(\mathbb E D_A\)
by a nonnegative quantity at most
\(n_M^2q_\gamma/2\).
Relative to \(n_M/\gamma_M\), this remainder
is \(O(n_M/\gamma_M)=o(1)\).
Proposition~\ref{f15v66:leading} proves the
normalizer formula. Lifting and data processing
give the remaining assertions.
A four-dimensional patch contains \(O(R_M^4)\)
positive links, which proves the last claim.
\end{proof}

\subsection{What this does and does not pay}

For arbitrary finite \(A\), the entropy cost
per integrated edge is at most \(d_\gamma\to0\).
For \(n=o(\gamma)\), the stronger total entropy
and TV bounds hold. Neither assertion is
a volume-uniform small total error for an
extensive infinite set. The quadratic bound
used no decay of correlations and hence costs
\(n^2\), not \(n\).

The reference keeps every nonlinear frustration
and the original vacuum boundary weights.
It changes only the retained density by a
fully normalized explicit local factor.
The lift keeps the exact conditional link
fluctuations. These distinctions are necessary:
small innovation alone does not control
frustration, as V39's central example already
showed.

This is one independent-link integration layer.
After it, the retained action contains
\(F_e=B-|H_e|\). Applying another edge colour
does not have the original linear-staple
conditional, and neither V38's Gamma bound
nor this theorem can simply be reused for
that changed action. Moreover the error
estimates here concern the selected vacuum,
not an all-exterior or source-uniform
comparison. They provide no decay in the
separation of retained observables.

The upstream task is therefore concrete:
control successive integration of the
generated nonlinear retained law, including
its normalizers and completed observable
transport, on scales beyond this paid
single-layer window. Establishing the
needed joint frustration and response bounds
is still open. A small entropy density alone
does not give those bounds, a physical
renormalization scale, or a mass gap.

\subsection{Checks and source preservation}

The diagnostic is
\begin{center}\small
\path{scripts/verify_ym_f15_v66_normalized_star_decimation.py}.
\end{center}
It reruns V65's inherited checks and V38.
Outward rational series enclosures test 25
global and 15 small-frustration bounds for
the actual one-link integral, together with
five zero-frustration endpoint identities.
There are 24 exact-probability tables for
both entropy directions, the paid normalizer,
and quadratic remainders; 336 common-conditional
lifting identities; and 24 completed-output
data-processing checks.
Five moment-budget constants, 20
arbitrary-dependence square inequalities,
four leading constants and seven
subcritical scale exponents are checked.

The probability tables are auxiliary diagnostics,
not discrete replacements for the Wilson
vacuum. They do not prove infinite-volume
limits or a gap by sampling.
The analytic proofs above supply the actual-law
statements.
The script has 6,569 bytes and SHA--256
\begin{center}\scriptsize\ttfamily
F761405E9C0DF53F3BDC01264939E2F16D64C0225E6F87219A931E5A72BE5ADC.
\end{center}

The V65 predecessor has 1,481,560 bytes and SHA--256
\begin{center}\scriptsize\ttfamily
16D179596374028BFADB2332351CC33848F6FD775F6DCD1781E30337A50A9561.
\end{center}
Removing this final section and restoring the
title version recovers V65 exactly.
V66 replaces it as the sole active main source.
The 14 protected inputs and 108 earlier
diagnostics remain unchanged.
Only \path{.tex} sources are retained in
\path{latex_notes}; all compilation products
go to
\begin{center}\small
\path{artifacts/ym_f15_v66_texcheck/}.
\end{center}
The next companion checkpoint remains V70.

\begin{firewall}
V66 proves a normalized, quantitative single-layer
comparison for the actual selected Wilson vacuum.
It retains the full nonlinear frustration and
exact conditional fine-link fluctuations.
It does not prove iterated RG contraction,
uniform physical-source control, an interacting
continuum construction, or a positive physical
Yang--Mills mass gap. The original goal remains
active and uncompleted.
\end{firewall}
\section{V67: a smooth retained reference and paid differential control of one integration layer}
\label{f15v67:section}

V66 controlled normalized laws after independent-star
integration, but its minimum-action reference contains
\(F_e=B-|H_e|\). That expression is not differentiable
at a zero staple sum. Although such fields have a small
actual-vacuum probability in the selected regime,
a nonsmooth reference is inconvenient for a subsequent
integration or a differentiated effective action.
One cannot remove that issue by differentiating
V66's entropy or total-variation estimate.

We construct a globally smooth local reference between
the exact retained potential and its conditional minimum.
It keeps all V66 normalization bounds. We then prove
small local scores, bounded second-derivative rows in
actual-vacuum \(L^p\), and an extensive relative Fisher
information bound. No small-field chart is imposed on
the sampled law. An exact native six-staple calculation
shows why the unscaled Hessian correction does not vanish.

The exact score/Hessian formulas for general completed
coarse maps were already proved in V19, and chart
convexity was proved in V13. Those results are not being
repeated. Here the new object is a specific globally
smooth replacement for V66's local retained factor,
with constants supplied by its exact compact integral.
This is still one integration layer, not an iterated
RG theorem.

The V66 source and all 123 protected files were verified
unchanged. No companion checkpoint is due until V70.

\subsection{A smooth radial comparison valid at every staple field}

Retain V66's selected vacuum law, coupling
\(\gamma=\gamma_M\), \(B=6\gamma\), independent finite
link set \(A\), and \(n=|A|\). Write
\[
 Z(h)=\int_{S^3}e^{h u_0}\,dH(u),\qquad
 \ell(h)=\log Z(h),\qquad
 a(h)=\ell'(h)\quad(h\ge0).
\]
This \(Z\) is a one-link integral, not a many-link
normalizer. Set
\begin{equation}
 r(h)=\sqrt{1+h^2},\qquad b(h)=r'(h)=\frac{h}{\sqrt{1+h^2}}.
 \label{eq:f15v67:radial}
\end{equation}
Define the exact, minimum and smooth retained
contributions by
\begin{equation}
 V_e=\ell(B)-\ell(h_e),\qquad
 F_e=B-h_e,\qquad
 \overline F_e=r(B)-r(h_e).
 \label{eq:f15v67:potentials}
\end{equation}
The proposed smooth correction is
\begin{equation}
 \overline D_e=\overline F_e-V_e
        =\ell(h_e)-r(h_e)-\ell(B)+r(B),
 \qquad \overline D_A=\sum_{e\in A}\overline D_e .
 \label{eq:f15v67:correction}
\end{equation}

\begin{proposition}[A pointwise smooth sandwich]
\label{f15v67:sandwich}
For every \(h\ge0\), \(0\le a(h)\le b(h)<1\).
Consequently, at every retained configuration,
\begin{equation}
        0\le V_e\le\overline F_e\le F_e,\qquad
        0\le\overline D_e\le D_e ,
 \label{eq:f15v67:sandwich}
\end{equation}
where \(D_e=F_e-V_e\) is exactly V66's log factor.
Both \(\overline F_e\) and \(\overline D_e\) are
globally smooth functions of the retained compact links,
including all configurations with \(H_e=0\).
\end{proposition}

\begin{proof}
The compact integral is even and strictly positive.
Its logarithmic derivative is
\(a(h)=\mathbb E_hu_0\), and
\(a'(h)=\operatorname{Var}_h(u_0)\ge0\).
Thus \(a(0)=0\) and \(a(h)\ge0\).
Integration by parts on the unit three-sphere, using
\(\Delta u_0=-3u_0\) and
\(|\nabla u_0|^2=1-u_0^2\), gives
\[
             h(1-\mathbb E_hu_0^2)=3\mathbb E_hu_0.
\]
For \(h>0\) this implies the Riccati equation
\begin{equation}
                  a'(h)=1-a(h)^2-\frac3h a(h).
 \label{eq:f15v67:riccati}
\end{equation}
Haar symmetry gives \(\mathbb E_0u_0^2=1/4\);
hence \(a(h)=h/4+O(h^3)\), whereas \(b(h)=h+O(h^3)\).
At a hypothetical first positive contact \(a(h)=b(h)\),
the right side of \eqref{eq:f15v67:riccati} equals
\[
       \frac1{1+h^2}-\frac3{\sqrt{1+h^2}}<0,
\]
while \(b'(h)=(1+h^2)^{-3/2}>0\).
That contradicts a first crossing from below.
So \(a(h)<b(h)\) for \(h>0\), and both vanish at zero.

Integrate \(0\le a\le b\le1\) from \(h_e\) to \(B\)
to obtain the potential sandwich and the correction
bounds. Finally
\[
       \ell(|H|)=\log\int_{S^3}e^{u\cdot H}\,dH(u)
\]
is a smooth function on \(\mathbb R^4\);
\(\sqrt{1+|H|^2}\) is also smooth there.
The staple field is a finite sum of smooth compact
group words. Their compositions prove the last claim
without differentiating the norm at zero.
\end{proof}

The smoothing radius in \eqref{eq:f15v67:radial}
is a fixed one in staple-field units. It does not
alter the Wilson action or replace the actual vacuum.
It only defines an explicit comparison law.

\subsection{Normalization is preserved, including its leading cost}

Let \(\rho_A\) be the actual retained vacuum marginal
and \(K_A\) its exact deleted-link conditional kernel,
as in V66. Define
\begin{equation}
 \overline Z_A=\mathbb E_{\rho_A}e^{-\overline D_A},\qquad
 d\overline\nu_A=\overline Z_A^{-1}e^{-\overline D_A}\,d\rho_A,
 \qquad
 \widehat{\overline\nu}_A=\overline\nu_AK_A .
 \label{eq:f15v67:laws}
\end{equation}
These are normalized probabilities. At finite volume
the retained reference has potential
\(S_{\rm rest}+\sum_{e\in A}\overline F_e\).
For a finite vacuum cylinder the original endpoint
weight remains unchanged. On the infinite-time vacuum
the displayed Radon--Nikodym definition is used;
no infinite total-action density is asserted.

\begin{proposition}[The smooth reference pays the same entropy budget]
\label{f15v67:entropy}
All bounds in \eqref{eq:f15v66:entropy} and
\eqref{eq:f15v66:normalizer-remainder} remain valid
with \(D_A,Z_A,\nu_A\) replaced by
\(\overline D_A,\overline Z_A,\overline\nu_A\).
The common-conditional lift and all completed-output
entropy and TV bounds of V66 also remain valid.
Furthermore, uniformly in link position and orientation,
\begin{equation}
       \gamma_M\mathbb E_{\mu_M}\overline D_e\longrightarrow\frac34 .
 \label{eq:f15v67:leading}
\end{equation}
For \(1\le n_M=o(\gamma_M)\),
\[
      -\log\overline Z_{A_M}
                  =\frac{n_M}{\gamma_M}\left(\frac34+o(1)\right).
\]
\end{proposition}

\begin{proof}
The pointwise bound \(0\le\overline D_A\le D_A\)
and \((\overline D_A)^2\le D_A^2\) give exactly
the first and second moment budgets used in
V66's entropy proof. Its likelihood identities
and common-conditional disintegration apply unchanged.

For the leading term,
\(D_e-\overline D_e=F_e-\overline F_e\ge0\).
On \(F_e\le B/2\), the inequality
\[
    1-b(t)=\frac1{\sqrt{1+t^2}(\sqrt{1+t^2}+t)}
                         \le\frac1{2t^2}
\]
gives
\[
 \gamma(F_e-\overline F_e)
      =\gamma\int_{B-F_e}^B(1-b(t))\,dt
                    \le\frac{F_e}{18\gamma}.
\]
On the complementary event use \(F_e-\overline F_e\le B\).
V66's uniform first moment and exponential tail yield
\[
 0\le\gamma\,\mathbb E(D_e-\overline D_e)
       \le\frac{120}{18\gamma}
                          +6e^3\gamma^2e^{-\gamma/32}\longrightarrow0.
\]
Now use \eqref{eq:f15v66:leading}.
The argument is uniform over all locations and the
four orientations. Summation and the same quadratic
normalizer remainder as V66 prove the final formula.
\end{proof}

\subsection{Left derivatives of the native staple field}

Choose unit imaginary quaternions \(E_1,E_2,E_3\).
For a retained positive link \(f\), define
\[
 X_{f,\alpha}G(U)
    =\left.\frac{d}{dt}
        G(\ldots,e^{tE_\alpha}U_f,\ldots)\right|_{t=0}.
\]
Write \(i=(f,\alpha)\), \(X_i=X_{f,\alpha}\).
All second derivatives below are the ordered
left derivatives \(X_jX_i\). Different-link
fields commute; same-link derivatives are not
silently identified with coordinate partials.

Each of the six staples of \(e\) is a three-link
word, containing any prescribed retained link
at most once. Unit-quaternion multiplication
and inversion preserve Euclidean norms.
It follows, for every \(i,j\), that
\begin{equation}
        |X_iH_e|\le6\gamma,\qquad |X_jX_iH_e|\le6\gamma.
 \label{eq:f15v67:H-global}
\end{equation}
These estimates include repeated derivatives on
the same link. A twice differentiated word is a
product with two unit generator insertions, still
of norm at most one.

If \(h_e>0\), put \(N_e=H_e/h_e\).
Because \(a_{e,p}\cdot X_i a_{e,p}=0\),
\[
 X_ih_e=\gamma\sum_{p\ni e}(N_e-a_{e,p})\cdot X_ia_{e,p},
 \qquad
             \sum_{p\ni e}|N_e-a_{e,p}|^2=\frac{2F_e}{\gamma}.
\]
Cauchy--Schwarz and the at-most-six nonzero
staple derivatives give
\begin{equation}
                        |X_ih_e|^2\le12\gamma F_e .
 \label{eq:f15v67:radial-score}
\end{equation}
This estimate uses only frustration, not a
selected gauge or small individual link logarithms.

When \(h_e\ge B/2=3\gamma\), differentiation of
the Euclidean norm gives
\begin{equation}
 \begin{split}
 X_jX_ih_e
    &=N_e\cdot X_jX_iH_e
       +\frac{(X_jH_e)\cdot(I-N_eN_e^{\mathsf T})(X_iH_e)}{h_e},\\
                        |X_jX_ih_e|&\le18\gamma.
 \end{split}
 \label{eq:f15v67:hessian-h}
\end{equation}
The formula also holds for ordered derivatives
on the same group variable.

\subsection{Small local scores and bounded local second derivatives}

Write \(g(h)=\ell(h)-r(h)\), so that
\(\overline D_e=g(h_e)-g(B)\).
For \(h>0\), V38's conditional Gamma bounds imply
\begin{equation}
 \begin{gathered}
 |g'(h)|=b(h)-a(h)\le1-a(h)=\frac{m(h)}h\le\frac3{2h},\\
 |g''(h)|\le a'(h)+b'(h)
          \le\frac{15}{4h^2}+\frac1{h^3}
          \le\frac{19}{4h^2}\quad(h\ge1).
 \end{gathered}
 \label{eq:f15v67:g-good}
\end{equation}
Here \(m(h)=\mathbb E_h[h(1-u_0)]\), and
\(a'(h)=\operatorname{Var}_h(u_0)
 \le\mathbb E_h[(h(1-u_0))^2]/h^2\).
The numerator is the second moment, whose
Gamma bound is \(15/4\), not the square of the mean.

For clarity the nonsingular global bounds can
be obtained directly in \(H\)-space. If
\(G(H)=\ell(|H|)-\sqrt{1+|H|^2}\), then
\begin{equation}
                     |\nabla_HG|\le1,\qquad
                     \|\nabla_H^2G\|\le2 .
 \label{eq:f15v67:g-global}
\end{equation}
Indeed the first inequality follows from
\(0\le b-a\le1\), extended at zero.
The Hessian of \(\log\int e^{u\cdot H}\,dH(u)\)
is the covariance matrix of the unit vector \(u\),
with operator norm at most one.
The Hessian of \(\sqrt{1+|H|^2}\) also has
operator norm at most one. This proves
\eqref{eq:f15v67:g-global} at all fields.

\begin{theorem}[Actual-vacuum differential bounds]
\label{f15v67:derivatives}
For every fixed \(1\le p<\infty\) there is
\(C_p<\infty\), independent of \(M,A,e,i,j\), such that
\begin{equation}
 \|X_i\overline D_e\|_{L^p(\mu_M)}
                         \le C_p\gamma_M^{-1/2},
 \qquad
 \|X_jX_i\overline D_e\|_{L^p(\mu_M)}\le C_p.
 \label{eq:f15v67:Lp}
\end{equation}
An explicit squared-score budget is
\begin{equation}
 \mathbb E_{\mu_M}|X_i\overline D_e|^2
       \le A_\gamma:=
                 \frac{360}{\gamma}+36e^3\gamma^2e^{-\gamma/32}.
 \label{eq:f15v67:A}
\end{equation}
These are averaged bounds under the actual selected
vacuum, not uniform conditional bounds at every exterior.
\end{theorem}

\begin{proof}
On \(F_e\le B/2\), \eqref{eq:f15v67:radial-score}
and \eqref{eq:f15v67:g-good} give
\[
           |X_i\overline D_e|\le\sqrt{3F_e/\gamma}.
\]
For the second derivative use the exact chain rule
\[
 X_jX_i\overline D_e
       =g''(h_e)(X_jh_e)(X_ih_e)+g'(h_e)X_jX_ih_e .
\]
Bounding \(|X_ih_e|\le6\gamma\) and using
\eqref{eq:f15v67:hessian-h}, its absolute value
is at most
\[
          \frac{19}{4(3\gamma)^2}(6\gamma)^2
                    +\frac{3}{2(3\gamma)}18\gamma=28 .
\]
On the bad event \(F_e>B/2\), use the smooth
\(H\)-space chain rule and
\eqref{eq:f15v67:H-global}--\eqref{eq:f15v67:g-global}:
\[
 |X_i\overline D_e|\le6\gamma,\qquad
 |X_jX_i\overline D_e|\le72\gamma^2+6\gamma .
\]
V66 proves that this event has probability at most
\(e^3e^{-\gamma/32}\), and that every fixed positive
moment of \(F_e\) is uniformly bounded by its
exponential moment. Thus
\[
 \begin{split}
 \mathbb E|X_i\overline D_e|^p
    &\le(3/\gamma)^{p/2}\mathbb EF_e^{p/2}
                              +(6\gamma)^p e^3e^{-\gamma/32},\\
 \mathbb E|X_jX_i\overline D_e|^p
    &\le28^p+(72\gamma^2+6\gamma)^p e^3e^{-\gamma/32}.
 \end{split}
\]
These imply \eqref{eq:f15v67:Lp}.
For \(p=2\) insert \(\mathbb EF_e\le120\)
to obtain \eqref{eq:f15v67:A}.
\end{proof}

\subsection{Finite derivative rows and extensive relative Fisher information}

Let \(\mathcal C_e\) be the retained links in
the six staples of \(e\). There are at most
18 such links, hence at most 54 scalar
left-derivative coordinates. A given retained
link belongs to at most 18 carriers
\(\mathcal C_e\): there are six incident
plaquettes, each with at most three other
links that could serve as \(e\).
These conservative counts do not grow with \(A\).
Each carrier link anchor is within lattice
distance two of the deleted edge anchor.
Therefore two coordinates in a single carrier
have anchor distance at most four.

Minkowski's inequality and Theorem~\ref{f15v67:derivatives}
give
\begin{equation}
 \begin{gathered}
 \sup_i\|X_i\overline D_A\|_{L^p(\mu_M)}
                       \le18C_p\gamma^{-1/2},\\
 \sup_i\sum_j\|X_jX_i\overline D_A\|_{L^p(\mu_M)}
                       \le972C_p,\\
       X_jX_i\overline D_A=0
                       \quad\hbox{if }d(i,j)>4 .
 \end{gathered}
 \label{eq:f15v67:rows}
\end{equation}
The supremum is outside the \(L^p\) norm.
This is not a bound on the expectation of a
simultaneous supremum over an arbitrarily large
volume, nor a pointwise interaction norm.

For probabilities \(P,Q\) whose likelihood ratio
is a positive smooth cylinder function, define
the relative score information here by
\[
       I(P\mid Q)=\sum_i
                    \mathbb E_P|X_i\log(dP/dQ)|^2.
\]
Only the finite set of coordinates on which
the ratio depends contributes. This definition
does not require a density for an infinite
total action.

\begin{theorem}[Paid relative score for the full fine law]
\label{f15v67:Fisher}
The original vacuum and the common-conditional
lift satisfy
\begin{equation}
 \begin{gathered}
 I(\mu_M\mid\widehat{\overline\nu}_A)\le972n A_\gamma,\\
 I(\widehat{\overline\nu}_A\mid\mu_M)
                    \le e^{n d_\gamma}\,972n A_\gamma ,
 \end{gathered}
 \label{eq:f15v67:Fisher}
\end{equation}
where \(d_\gamma\) is V66's first-moment budget.
In particular the forward score information is
\(O(n/\gamma)\) for arbitrary finite \(A\).
If \(n d_\gamma\) is bounded, the same order
holds in reverse.

More generally, if \(n d_\gamma\le C\), the
bounds \eqref{eq:f15v67:Lp} and
\eqref{eq:f15v67:rows} hold under the smooth
reference or its lift, with the right sides
multiplied by at most \(e^{C/p}\).
\end{theorem}

\begin{proof}
The likelihood ratio is
\[
          \frac{d\mu_M}{d\widehat{\overline\nu}_A}
                         =\overline Z_A e^{\overline D_A}.
\]
The normalizer is a scalar, so its derivatives
vanish. All derivatives on deleted links also
vanish because \(\overline D_A\) is retained-measurable.
For each remaining coordinate at most 18
summands have nonzero derivatives. Therefore,
pointwise,
\[
 \sum_i|X_i\overline D_A|^2
       \le18\sum_{e\in A}\sum_{i\in\mathcal C_e\times\{1,2,3\}}
                                  |X_i\overline D_e|^2.
\]
There are at most \(54n\) inner terms.
Expectation and \eqref{eq:f15v67:A} prove the
forward bound, with \(18\cdot54=972\).
For the reverse direction and the other
changed-law bounds, use
\[
 \frac{d\widehat{\overline\nu}_A}{d\mu_M}
      =\frac{e^{-\overline D_A}}{\overline Z_A}
                         \le e^{n d_\gamma}.
\]
The retained-law statement follows by the
same density formula. No independence between
carriers, or functional inequality for the
retained law, has been assumed.
\end{proof}

Relative Fisher information is measured here in
the declared fine-link metric. Its monotonicity
under an arbitrary nonlinear completed output
is not asserted; V66's entropy and TV contraction
are the output statements that do hold.
In thermal tangent coordinates \(v=\sqrt\gamma\,t\),
each derivative is \(X_i/\sqrt\gamma\), so the
corresponding score information is exactly
\(I/\gamma=O(n/\gamma^2)\).
This rescaling is not a physical mass normalization.

\subsection{An exact nonzero flat-field Hessian correction}

There is an important limit to the small-score
conclusion. Fix one deleted link \(e=(0,\mu)\).
Its six staples have six distinct middle links
parallel to \(e\), one for each positive or negative
transverse displacement. Each middle link occurs
in just that staple. At the flat configuration,
vary those middle links, with signs chosen according
to the staple orientation, so that the staples are
\[
                 a_j(t)=\exp\left(\sum_{\alpha=1}^3
                                      t_{j,\alpha}E_\alpha\right),
                         \qquad j=1,\ldots,6 .
\]
All other links remain identities. These are
literal fine-link variations, not independent
abstract staple fields imposed without realization.

At \(t=0\), \(H=B I\), all first derivatives
of \(h=|H|\) vanish, and direct differentiation gives
\[
 \left.\partial_{t_{k,\beta}}\partial_{t_{j,\alpha}}h\right|_0
             =-\gamma\,\delta_{\alpha\beta}
                                  (\delta_{jk}-1/6).
\]
Define the six-dimensional projection
\(P_6=I_6-\mathbf1\mathbf1^{\mathsf T}/6\).
The exact native Hessian block of this one
correction term is consequently
\begin{equation}
  \left.D^2_t\overline D_e\right|_0
      =\kappa_\gamma(P_6\otimes I_3),\qquad
         \kappa_\gamma=\gamma\{b(6\gamma)-a(6\gamma)\}.
 \label{eq:f15v67:flat}
\end{equation}
It has rank 15, since \(\kappa_\gamma>0\).
The first derivative vanishes here, so the
same matrix is its intrinsic Hessian block.

By V66's conditional innovation limit,
\(m(h)\to3/2\), and
\[
 \kappa_\gamma
      =\frac{m(6\gamma)}6-\gamma\{1-b(6\gamma)\}
                         \longrightarrow\frac14 .
\]
Thus the unscaled block tends to
\begin{equation}
       \frac14(P_6\otimes I_3).
 \label{eq:f15v67:flat-limit}
\end{equation}
Its same-colour diagonal entries are \(5/24\);
its distinct-middle-link, same-colour entries
are \(-1/24\). A coherent motion of all six
staples is the three-dimensional zero mode.

In particular, small normalized likelihood
errors and small averaged first derivatives
do not imply a vanishing unscaled Hessian
correction, even at the flat configuration.
The exact star potential has coefficient
\(\gamma a(6\gamma)\) on this block, whereas
the smooth reference has coefficient
\(\gamma b(6\gamma)\). Their difference tends
to \(1/4\).
This is not a computation of a complete YM
beta function: no full blocking, rescaling,
or sum over all generated interactions has
been performed. For a larger integrated set
the Hessian of the full correction is the
sum of its star contributions, not just this
single block.

\subsection{The remaining integration problem is not bypassed}

There is now a smooth, gauge-invariant local
reference with paid normalization, small
actual-vacuum score per coordinate, bounded
annealed second-derivative rows, and a finite
carrier range. Its derivative moments also
transfer to the reference on sets with
\(n d_\gamma=O(1)\). This gives differential
information that V66's entropy bound alone
did not supply.

The bounds still do not control arbitrary
conditional exteriors after a second
integration. The new potential
\(S_{\rm rest}+\sum\overline F_e\) is nonlinear;
its link conditionals are not the original
linear-staple kernels. A bounded averaged
derivative row is not a Dobrushin contraction,
a log-Sobolev inequality, or a scale-uniform
covariance bound. Nor is an extensive score
bound by itself a small entropy estimate
for an infinite integrated set.
The unscaled finite Hessian term must be kept
when assembling a further effective action.

The next upstream problem is to control the
normalized response of that generated smooth
interaction, or of the exact retained interaction,
through a further genuine integration step,
with the actual fluctuation kernel and
observable transport retained. No positive
physical gap or interacting continuum has
been constructed here.

\subsection{Diagnostics and preservation}

The diagnostic is
\begin{center}\small
\path{scripts/verify_ym_f15_v67_smooth_star_scores.py}.
\end{center}
It reruns V66 and its inherited checks.
There are nine outward rational enclosures
for Bessel ordering and the flat coefficient,
and nine Riccati comparison signs.
There are 216 quaternion tangent checks and
216 thermal radial derivative inequalities.
On periodic tori of sides four and six,
6,208 star carriers, their inverse
multiplicities and their radius-two anchor
bounds are checked.
The script also checks 36 projector identities,
324 native flat Hessian entries, three
nonzero-limit/rank constants, eight derivative
and overlap constants, and 16 relative-score
overlap inequalities.

All these diagnostics use exact rational
arithmetic or outward rational series bounds.
They do not estimate Wilson-vacuum moments
by sampling; the actual-law estimates are
the analytic proofs above.
The script has 6,476 bytes and SHA--256
\begin{center}\scriptsize\ttfamily
C2392682E653B9D25350DC58DD796761391BAA6069D9E1C16C41D548AFF25565.
\end{center}

The V66 predecessor has 1,502,406 bytes and SHA--256
\begin{center}\scriptsize\ttfamily
85CB1AE1B4D20E8F7BB3EC549DAA55A08A43C51A6AE54E581F9D978150A0B908.
\end{center}
Removing this appended section and restoring
the title version recovers V66 exactly.
V67 replaces it as the sole active main note.
The 14 protected inputs and 109 earlier
diagnostics remain unchanged.
Only \path{.tex} sources are kept in
\path{latex_notes}; compilation products go to
\begin{center}\small
\path{artifacts/ym_f15_v67_texcheck/}.
\end{center}
The next companion review remains V70.

\begin{firewall}
V67 proves smooth single-layer reference and
differential likelihood bounds for the selected
actual Wilson vacuum. The finite Hessian
correction is retained explicitly.
No iterated RG contraction, all-exterior
generated-law control, interacting continuum,
or positive physical YM mass gap follows.
The full original goal remains active.
\end{firewall}
\section{V68: a paid second integration through the generated conditional law}
\label{f15v68:section}

V67 supplied a smooth local reference and differential
likelihood estimates, but did not propagate them through
a new integration. We now do so for finite retained-link
blocks. The kernel in this step is the actual nonlinear
conditional after the first star integration; it is not
replaced by an original one-link Wilson kernel.

The principal issue is its conditional normalizer.
An exterior on which that normalizer is small cannot
be discarded by formally differentiating an averaged
error. We bound its reciprocal in actual-exterior
\(L^p\), then keep all the covariance terms in the
first and second derivatives of the new marginal.
For a separated family of bounded blocks, the
conditional entropy errors add with a constant
independent of the number of blocks.
The exteriors need not be independent.

The conditional likelihood identities themselves
are not new; related identities occur in V19 and
V21--V37. What is paid here is their application
to V67's particular generated nonlinear star
interaction, including the reciprocal normalizer
and the native score moments needed for the next
derivative bounds.
No source-uniform or all-exterior conclusion is inferred.

The V67 source and all 124 protected files were
checked unchanged. The next companion checkpoint
remains V70.

\subsection{The literal second-stage conditional}

Keep the selected actual vacuum \(\mu_M\),
\(\gamma=\gamma_M\), the first independent deleted
link set \(A\), and the notation of V67.
Its exact retained law is \(\rho_A\), and its
smooth comparison law is \(\overline\nu_A\).
Let \(B\subset A^c\) be any finite nonempty set
of retained links to be integrated next.
No independence condition is imposed within \(B\).
Write \(v=U_B\) and \(\eta=U_{(A\cup B)^c}\);
let \(\sigma_B\) be the actual law of \(\eta\).

For a deleted star \(e\in A\), let
\(\mathcal C_e\) be its retained staple carrier.
Set
\begin{equation}
 \begin{gathered}
 A_B=\{e\in A:\mathcal C_e\cap B\ne\varnothing\},\qquad
 k=|A_B|\le18|B|,\\
 D_B(v,\eta)=\sum_{e\in A_B}\overline D_e(v,\eta),\qquad
 D_{\rm out}(\eta)=\sum_{e\in A\setminus A_B}\overline D_e(\eta).
 \end{gathered}
 \label{eq:f15v68:local-defect}
\end{equation}
Thus \(\overline D_A=D_{\rm out}+D_B\).
If \(k=0\), all quantities below associated
with this defect are zero or identically one,
as appropriate.

The exact first-stage retained interaction
consists of remaining bare plaquette terms
and the star potentials \(V_e=\ell(6\gamma)-\ell(h_e)\).
Let \(W_B(v,\eta)\) be the sum of precisely
those terms whose carriers meet \(B\):
\begin{equation}
 W_B=\sum_{\substack{p:\,p\cap A=\varnothing\\p\cap B\ne\varnothing}}X_p
                         +\sum_{e\in A_B}V_e .
 \label{eq:f15v68:conditional-potential}
\end{equation}
Here intersection with a plaquette means
intersection with its link set.
The actual conditional probability is
\begin{equation}
 P_B^\eta(dv)=
       \frac{e^{-W_B(v,\eta)}\,dH_B(v)}
                       {\int e^{-W_B(w,\eta)}\,dH_B(w)}.
 \label{eq:f15v68:P}
\end{equation}
It is a positive smooth compact integral,
with all the generated \(V_e\)'s retained.

To justify \eqref{eq:f15v68:P} in the vacuum,
first use the exact finite-period marginal
formula and condition on \(\eta\).
Only terms meeting \(B\) remain in its
conditional normalization.
The kernel is a bounded continuous function
of finitely many exterior links when applied
to a bounded continuous test.
V62's full-link cylinder limit and the same
monotone-class argument as V66 pass this
conditional identity to the vacuum.
Equivalently, one can first condition the
original Wilson law on the finite set
\(A_B\cup B\) and then integrate \(A_B\).
Because \(A\) is independent, no omitted
\(A\)-link enters the resulting kernel.
Neither argument replaces the actual exterior law.

The smooth reference conditional is
\begin{equation}
 \begin{gathered}
 z_B(\eta)=\mathbb E_{P_B^\eta}e^{-D_B},\qquad
 T_B(\eta)=-\log z_B(\eta),\\
 Q_B^\eta(dv)=R_B(v,\eta)P_B^\eta(dv),\qquad
                 R_B=\frac{e^{-D_B}}{z_B}.
 \end{gathered}
 \label{eq:f15v68:Q}
\end{equation}
The actual joint law used for moment estimates
is \(\sigma_B(d\eta)P_B^\eta(dv)=\rho_A\).

\subsection{An exponential defect budget, not just a first moment}

Recall V66's quantities \(L_\gamma,q_\gamma\)
and its star bound \(\mathbb E e^{F_e/96}\le e^3\).
For \(p\ge1\), \(k\ge1\), define
\begin{equation}
 \begin{split}
 \mathfrak B_{p,k}(\gamma)
    &=e^3\left(\frac{96kp}{\gamma}\right)^p
             +ke^3\exp\{pkL_\gamma-\gamma/32\},\\
 b_{p,k}(\gamma)&=\mathfrak B_{p,k}(\gamma)^{1/p}.
 \end{split}
 \label{eq:f15v68:budget}
\end{equation}
Set \(b_{p,0}=0\).

\begin{proposition}[Actual exponential defect moment]
\label{f15v68:exp}
If \(\gamma\ge96pk\), then
\begin{equation}
       \|e^{D_B}-1\|_{L^p(\rho_A)}\le b_{p,k}(\gamma).
 \label{eq:f15v68:exp}
\end{equation}
For fixed \(p,k\), this is \(O_{p,k}(\gamma^{-1})\).
More generally, for fixed \(p\) and
\(1\le k=k_M\) with \(k_M\log\gamma_M=o(\gamma_M)\),
\begin{equation}
                       b_{p,k_M}(\gamma_M)=O_p(k_M/\gamma_M).
 \label{eq:f15v68:growing}
\end{equation}
\end{proposition}

\begin{proof}
Put \(S=\sum_{e\in A_B}F_e\).
On the event \(F_e\le3\gamma\) for every such
star, V66 and the smooth sandwich give
\(0\le D_B\le S/(2\gamma)\).
Hölder gives, without independence,
\[
              \mathbb E e^{S/(96k)}
                    \le\prod_{e\in A_B}
                       (\mathbb E e^{F_e/96})^{1/k}\le e^3 .
\]
For \(x\ge0\), \((e^x-1)^p\le x^pe^{px}\).
The coupling threshold gives
\(p/(2\gamma)\le1/(192k)\), and
\(S^pe^{S/(192k)}\le(192kp)^pe^{S/(96k)}\).
Hence the good-event contribution to
\(\mathbb E(e^{D_B}-1)^p\) is at most
\(e^3(96kp/\gamma)^p\).
On the complementary event use \(D_B\le kL_\gamma\)
and the union bound
\[
        \mathbb P\{\text{some }F_e>3\gamma\}
                               \le ke^3e^{-\gamma/32}.
\]
This proves \eqref{eq:f15v68:exp}.
For the growing-set assertion, the exponent
\(pkL_\gamma+\log k+3-\gamma/32\) equals
\(-\gamma/32+o(\gamma)\).
Its contribution is smaller than the first
term, eventually. The stated threshold also
holds eventually, proving
\eqref{eq:f15v68:growing}.
\end{proof}

\subsection{The conditional reciprocal and both conditional entropies}

\begin{theorem}[A paid reciprocal normalizer]
\label{f15v68:reciprocal}
Under the preceding threshold,
\begin{equation}
 \begin{gathered}
 \|z_B^{-1}-1\|_{L^p(\sigma_B)}\le b_{p,k},\qquad
 \|T_B\|_{L^p(\sigma_B)}\le b_{p,k},\\
 \|R_B-1\|_{L^p(\rho_A)}\le2b_{p,k},\qquad
 \|R_B\|_{L^p(\rho_A)}\le1+2b_{p,k}.
 \end{gathered}
 \label{eq:f15v68:reciprocal}
\end{equation}
The averaged conditional entropies obey
\begin{equation}
 \begin{gathered}
 \mathbb E_{\sigma_B}H(P_B^\eta\mid Q_B^\eta)
                                      \le\tfrac12 k^2q_\gamma,\\
 \mathbb E_{\sigma_B}H(Q_B^\eta\mid P_B^\eta)
                                      \le4b_{2,k}^2 .
 \end{gathered}
 \label{eq:f15v68:conditional-entropy}
\end{equation}
The second inequality uses \(\gamma\ge192k\).
These averages use the actual exterior law
\(\sigma_B\), not an unproved reference-exterior moment.
\end{theorem}

\begin{proof}
Conditional Cauchy--Schwarz gives
\[
 1\le(\mathbb E_{P_B^\eta}e^{D_B})
                 (\mathbb E_{P_B^\eta}e^{-D_B}),
 \qquad
 z_B^{-1}-1\le\mathbb E_{P_B^\eta}(e^{D_B}-1).
\]
Conditional Jensen and
Proposition~\ref{f15v68:exp} prove the first bound.
Also \(0\le T_B\le\mathbb E_{P_B^\eta}D_B\),
and \(D_B\le e^{D_B}-1\), proving the second.

Since \(z_B\le1\), one has
\(e^{-D_B}\le R_B\le z_B^{-1}\).
If \(R_B\ge1\), its excess is at most
\(z_B^{-1}-1\); if \(R_B<1\), its deficit is
at most \(1-e^{-D_B}\le D_B\).
Minkowski and the preceding bounds give
the two likelihood estimates.

For each exterior V66's elementary nonnegative-tilt
argument gives
\[
 H(P_B^\eta\mid Q_B^\eta)
      =\mathbb E_{P_B^\eta}D_B-T_B
      \le\tfrac12\mathbb E_{P_B^\eta}D_B^2.
\]
The actual second moment is at most \(k^2q_\gamma\).
In the reverse direction \(\log r\le r-1\) and
\(\mathbb E_{P_B^\eta}R_B=1\) give
\[
       H(Q_B^\eta\mid P_B^\eta)
                      \le\mathbb E_{P_B^\eta}(R_B-1)^2.
\]
Integrate and apply the \(p=2\) likelihood bound.
No pointwise lower bound independent of
\(\gamma,\eta\) has been assumed for \(z_B\).
\end{proof}

In particular
\begin{equation}
 0\le\mathbb E_{\rho_A}D_B-\mathbb E_{\sigma_B}T_B
                                               \le\tfrac12 k^2q_\gamma .
 \label{eq:f15v68:mean-normalizer}
\end{equation}
V67's leading expectation therefore transports
through the new integration:
if \(\epsilon_M\to0\) bounds the uniform error
in \(\gamma_M\mathbb E\overline D_e\to3/4\), then
\[
 \left|\mathbb E_{\sigma_B}T_B-\frac{3k}{4\gamma_M}\right|
              \le\frac{k\epsilon_M}{\gamma_M}+\tfrac12k^2q_{\gamma_M}.
\]
This is an averaged conditional pressure statement,
not a pointwise formula at every exterior.

\subsection{Native score moments for the generated kernel}

Let \(X_i,X_j\) be V67's unit left derivatives
on exterior retained links. Use subscripts
for these ordered derivatives, for example
\(W_i=X_iW_B\), \(W_{ji}=X_jX_iW_B\).
For every fixed \(p\ge1\),
\begin{equation}
 \begin{gathered}
 \|W_i\|_{L^p(\rho_A)}\le C_p\sqrt\gamma,\qquad
 \|W_{ji}\|_{L^p(\rho_A)}\le C_p\gamma,\\
 \|(D_B)_i\|_{L^p(\rho_A)}\le C_p\gamma^{-1/2},\qquad
 \|(D_B)_{ji}\|_{L^p(\rho_A)}\le C_p .
 \end{gathered}
 \label{eq:f15v68:native-moments}
\end{equation}
The constants here do not grow with \(A\) or \(B\):
for a fixed differentiated coordinate only
a bounded number of terms contribute.

Here are the bounds for the actual potential,
rather than an assumed Gaussian approximation.
A bare plaquette satisfies
\[
              |X_iX_p|\le\sqrt{2\gamma X_p},
              \qquad |X_jX_iX_p|\le\gamma .
\]
For an exact star term,
\[
                |X_iV_e|=a(h_e)|X_ih_e|
                               \le\sqrt{12\gamma F_e},
\]
extended continuously through \(h_e=0\).
The \(H\)-space Hessian of \(\ell(|H|)\)
is the covariance matrix of \(u\).
Its tangential eigenvalue is \(a(h)/h\);
its radial eigenvalue is \(a'(h)\).
For \(h\ge3\gamma\) and eventually \(\gamma\ge2\),
both are at most \(1/h\), since
\(a'(h)\le15/(4h^2)\).
Thus the chain rule gives
\[
     |X_jX_iV_e|\le18\gamma
                     \quad(h_e\ge3\gamma).
\]
Globally the covariance norm is at most one,
giving \(|X_jX_iV_e|\le36\gamma^2+6\gamma\).
The exceptional event has probability at most
\(e^3e^{-\gamma/32}\).
V62's plaquette moments and V66's frustration
moments now prove the first row of
\eqref{eq:f15v68:native-moments}.
At most six bare plaquettes and 18 star carriers
meet a differentiated link, even when \(B\) is large.
V67's bounds give the second row with the
same incidence observation.

\subsection{Differentiating the paid new normalizer}

All compact integrals at fixed \(M,\gamma,A,B\)
are smooth, strictly positive functions of
their finitely many exterior links.
Differentiation below is therefore legitimate
before taking any asymptotic estimate.
For brevity put \(P=P_B^\eta\), \(Q=Q_B^\eta\),
\(D=D_B\), \(W=W_B\), \(R=R_B\).
The exact first derivative is
\begin{equation}
        X_iT_B=\mathbb E_P[(R-1)W_i]+\mathbb E_QD_i .
 \label{eq:f15v68:score}
\end{equation}
The exact second derivative is
\begin{equation}
 \begin{split}
 X_jX_iT_B
   ={}&\mathbb E_P[(R-1)W_{ji}]+\mathbb E_QD_{ji}\\
     &-\{\operatorname{Cov}_Q(W_i,W_j)
                         -\operatorname{Cov}_P(W_i,W_j)\}\\
     &-\operatorname{Cov}_Q(W_i,D_j)
      -\operatorname{Cov}_Q(D_i,W_j)
      -\operatorname{Cov}_Q(D_i,D_j).
 \end{split}
 \label{eq:f15v68:hessian}
\end{equation}
Indeed \(T_B\) is the difference between
\(-\log\int e^{-W-D}\,dH_B\) and
\(-\log\int e^{-W}\,dH_B\).
Differentiating each normalized integral gives
these identities, including the covariance
caused by its denominator. The formulas apply
to ordered same-link derivatives as well.

\begin{theorem}[Differential control through a genuine second integration]
\label{f15v68:jets}
Fix \(p\ge1\), suppose \(\gamma\ge96(16p)k\),
and write
\[
                 b=b_{16p,k}(\gamma),\qquad R_*=1+2b.
\]
Then constants \(C_p\), independent of the
finite sets and their locations, give
\begin{equation}
 \begin{gathered}
 \|X_iT_B\|_{L^p(\sigma_B)}
                \le C_p\{R_*\gamma^{-1/2}+b\sqrt\gamma\},\\
 \|X_jX_iT_B\|_{L^p(\sigma_B)}
                \le C_p\{R_*^2+\gamma bR_*\}.
 \end{gathered}
 \label{eq:f15v68:jets}
\end{equation}
For bounded \(k\), the orders are respectively
\(O_{p,k}(\gamma^{-1/2})\) and \(O_{p,k}(1)\).
For \(k\log\gamma=o(\gamma)\), they are
\(O_p((1+k)\gamma^{-1/2})\) and \(O_p(1+k)\).
\end{theorem}

\begin{proof}
Conditional expectation under \(P\) is an
\(L^s(\rho_A)\)-to-\(L^s(\sigma_B)\) contraction.
Expectations under \(Q\) can instead be
written as \(\mathbb E_P[R\,\cdot]\).
Apply Hölder, \eqref{eq:f15v68:reciprocal}
and \eqref{eq:f15v68:native-moments}
to \eqref{eq:f15v68:score}. Its two terms
have \(L^p\) norms at most
\(C_pb\sqrt\gamma\) and \(C_pR_*\gamma^{-1/2}\).

For the second derivative, the first two terms
cost at most \(C_p\gamma b\) and \(C_pR_*\).
To check the changed covariance explicitly,
put
\(\delta_i=\mathbb E_QW_i-\mathbb E_PW_i
             =\mathbb E_P[(R-1)W_i]\).
Then
\[
 \begin{split}
 \operatorname{Cov}_Q(W_i,W_j)-\operatorname{Cov}_P(W_i,W_j)
    ={}&\mathbb E_P[(R-1)W_iW_j]\\
       &-\delta_i\,\mathbb E_QW_j
        -(\mathbb E_PW_i)\,\delta_j .
 \end{split}
\]
Its \(L^p\) norm is at most \(C_p\gamma bR_*\):
each \(\delta\) costs \(C_pb\sqrt\gamma\);
the unchanged mean costs \(C_p\sqrt\gamma\),
and the \(Q\)-mean costs \(C_pR_*\sqrt\gamma\).
The two mixed \(W,D\) covariances cost at most
\(C_pR_*^2\), and the \(D,D\) covariance at most
\(C_pR_*^2/\gamma\).
Repeated Hölder uses only moments with
exponents at most \(16p\) for \(R-1,R\);
the native moments are available at every
fixed finite exponent.
Since \(R_*\ge1\) and eventually \(\gamma\ge2\),
these bounds prove \eqref{eq:f15v68:jets}.
Finally apply Proposition~\ref{f15v68:exp}.
\end{proof}

This controls the new normalizer as a function
of the remaining links. It does not differentiate
a TV bound or presume a conditional spectral gap.
The integration has retained the nonlinear
\(V_e\)'s throughout.

\subsection{The exact marginal ratio keeps its original global normalizer}

After integrating \(B\), denote the reference
marginal by \(\overline\sigma_B\).
The exact disintegration gives
\begin{equation}
 \frac{d\overline\sigma_B}{d\sigma_B}(\eta)
       =\overline Z_A^{-1}e^{-D_{\rm out}(\eta)}z_B(\eta)
       =\overline Z_A^{-1}e^{-D^{(2)}(\eta)},\qquad
              D^{(2)}=D_{\rm out}+T_B .
 \label{eq:f15v68:global-ratio}
\end{equation}
Thus
\begin{equation}
 \mathbb E_{\sigma_B}e^{-D^{(2)}}=\overline Z_A,\qquad
 0\le D^{(2)}
           \le\mathbb E_{\rho_A}[\overline D_A\mid\eta].
 \label{eq:f15v68:global-normalizer}
\end{equation}
The second inequality is conditional Jensen.
It also preserves the original first and
second defect-moment budgets.
All V67 global entropy and TV estimates
continue by data processing, or by applying
the nonnegative-tilt proof to \(D^{(2)}\).

Both laws have now been integrated exactly.
There is no new arbitrary reset of the
exterior marginal, and no missing scalar
factor. The small first-layer correction
has been propagated, not replaced by an
independently postulated second-layer
approximation.

\subsection{Many bounded blocks: conditional entropy is extensive}

Let \(B_1,\ldots,B_N\subset A^c\) be finite
disjoint blocks such that no retained
interaction term in \eqref{eq:f15v68:conditional-potential}
meets two different blocks.
Mutual anchor distance greater than four
is a sufficient condition.
Then \(A_{B_r}\) are disjoint, and, given
the common exterior \(\eta\), the exact
and reference conditional laws factor:
\begin{equation}
       P^\eta=\bigotimes_{r=1}^N P_r^\eta,\qquad
       Q^\eta=\bigotimes_{r=1}^N Q_r^\eta,\qquad
                  z(\eta)=\prod_{r=1}^Nz_r(\eta).
 \label{eq:f15v68:family}
\end{equation}
Only this conditional factorization is used.

\begin{proposition}[Arbitrary family size with an actual-exterior average]
\label{f15v68:family}
If \(k_r=|A_{B_r}|\le k_0\) for a fixed
\(k_0\), then under the actual common-exterior
law \(\sigma\),
\begin{equation}
 \begin{gathered}
 \mathbb E_\sigma H(P^\eta\mid Q^\eta)
                  \le\tfrac12\sum_{r=1}^Nk_r^2q_\gamma
                                      =O_{k_0}(N/\gamma^2),\\
 \mathbb E_\sigma H(Q^\eta\mid P^\eta)
                  \le4\sum_{r=1}^N b_{2,k_r}^2
                                      =O_{k_0}(N/\gamma^2).
 \end{gathered}
 \label{eq:f15v68:family-entropy}
\end{equation}
The effective correction after this whole
second layer is
\[
        D^{(2)}=D_{\rm out}+\sum_{r=1}^NT_{B_r},
        \qquad
        \mathbb E_\sigma e^{-D^{(2)}}=\overline Z_A .
\]
In particular the original global normalizer
has not become a product of averaged local
normalizers.
\end{proposition}

\begin{proof}
No potential term meets two blocks, so the
conditional exponent and product Haar
integral separate. The same is true of the
defect because a star carrier cannot meet
two blocks. Relative entropy of finite
products is the sum of the component
entropies. Apply
\eqref{eq:f15v68:conditional-entropy}
and integrate under the actual common
exterior; no independence of that exterior
is necessary. The marginal-ratio calculation
\eqref{eq:f15v68:global-ratio} with the product
of conditional normalizers gives the final
identity.
\end{proof}

These are conditional entropy estimates under
\(\sigma\). They are not estimates under
the changed exterior law
\(\overline\sigma\) without an additional
change-of-law argument. In particular a
reference global density cannot be bounded
uniformly just because the number of
separated blocks is large.

There is also a local derivative consequence.
Suppose \(|B_r|\le m\) and the anchor diameter
of each block is at most \(R\), with fixed
\(m,R\). The exterior support of \(T_{B_r}\)
lies within distance four of \(B_r\).
Let \(C_0=4\cdot9^4\), a conservative bound
for the number of positive links with
anchors in a radius-four coordinate cube.
A given exterior link lies in at most
\(C_0\) such supports: choose a distinct
nearby \(B_r\)-link for each contributing
disjoint block.
Each support contains at most \(C_0m\) links.
Theorem~\ref{f15v68:jets}, Minkowski, and
V67's original rows therefore give
\begin{equation}
 \begin{gathered}
 \sup_i\|X_iD^{(2)}\|_{L^p(\sigma)}
                                \le C_{p,m}\gamma^{-1/2},\\
 \sup_i\sum_j\|X_jX_iD^{(2)}\|_{L^p(\sigma)}
                                \le C_{p,m},\\
 X_jX_iD^{(2)}=0\quad\hbox{if }d(i,j)>R+8 .
 \end{gathered}
 \label{eq:f15v68:family-jets}
\end{equation}
For example the new row contribution is at
most \(3C_0^2m\) times the uniform individual
second-derivative bound, in addition to the
old retained-star contribution.
These bounds are independent of \(N,|A|\)
and volume. As in V67, the supremum is
outside the \(L^p\) norm.

\subsection{What has advanced, and what is still missing}

The previous one-layer differential obstruction
is now paid for a genuine second integration,
and for any fitting family of bounded blocks.
The reciprocal normalizer, actual conditional
likelihood, both conditional entropies, and
all terms in the differentiated marginal
have been controlled under their stated
actual laws. The proof also gives an explicit
growing-block budget when
\(k\log\gamma=o(\gamma)\).

This result propagates the already specified
small correction while integrating the
dominant nonlinear reference exactly.
It does not show that this dominant interaction
contracts, enters a strong-mixing region, or
has uniformly controlled long-distance
conditional responses.
Bounds for fixed block size, or constants
depending on growing block size, cannot be
silently iterated through the physical
renormalization trajectory.

The unresolved task has consequently moved
from the reciprocal normalizer at this
second layer to a scale-uniform estimate for
the dominant generated law itself.
A mass-gap argument still needs that estimate,
a physical scale identification, and a
nontrivial interacting continuum construction.
The finite Hessian term in V67 remains part
of that generated action; it was not dropped
by the present integration.

\subsection{Checks and source preservation}

The diagnostic is
\begin{center}\small
\path{scripts/verify_ym_f15_v68_second_integration.py}.
\end{center}
It reruns V67 and its inherited checks.
There are 150 pointwise likelihood checks,
30 reciprocal-normalizer inequalities,
30 two-sided conditional entropy checks,
30 independently differentiated density-jet
score/Hessian identities, and 30 changed
covariance identities.
Eight multi-patch systems with a common
five-state exterior verify 720 literal
integration terms, 40 conditional product
normalizers and eight unchanged global
normalizers.
Sixteen exponential-moment budgets, four
growing-block exponents and five native
derivative/carrier constants are also checked.

The finite tables have an arbitrary common
exterior, not a factorized exterior.
They are auxiliary exact diagnostics;
the Wilson-law statements are proved
analytically above.
The script has 7,125 bytes and SHA--256
\begin{center}\scriptsize\ttfamily
07A262C2AD251B0461D9B7BD6D69CB2050CE3C2C48ADD55EB199BB575BA553E9.
\end{center}

The V67 predecessor has 1,524,396 bytes and SHA--256
\begin{center}\scriptsize\ttfamily
06CB530BAFB4A0B2904B34F9AD8C2A2BC93ADDAD1CDE62CC2E621F12DC48A045.
\end{center}
Removing this final section and restoring
the title version recovers V67 exactly.
V68 is the sole active main source.
The 14 protected inputs and 110 earlier
diagnostics remain unchanged.
Only \path{.tex} files are kept in
\path{latex_notes}; compilation products
are written under
\begin{center}\small
\path{artifacts/ym_f15_v68_texcheck/}.
\end{center}
The next companion checkpoint is V70.

\begin{firewall}
V68 proves normalized second-integration
and differential-error bounds under the
selected actual Wilson law. It retains the
generated nonlinear conditional kernel and
the original global normalizer.
It does not prove scale-uniform contraction
of the dominant interaction, an interacting
continuum, or a positive physical Yang--Mills
mass gap. The original goal remains active.
\end{firewall}
\section{V69: dominant conditional energy in a native cluster}
\label{f15v69:section}

V68 propagated the small comparison correction through
a second integration, while leaving the dominant
generated interaction intact. We now estimate an
actual conditional energy of that dominant interaction.
The geometry is a retained link together with its six
parallel neighbours, which have already been integrated
in one independent star layer.

This conditional is not an original linear-staple
Wilson kernel. We derive its exact compact density.
A squared-radius concavity inequality then compares
it, at every exterior, with a nonlinear radial
conditional. The latter has a Gamma-dominated energy
above its minimum. The comparison cost is an explicit
exterior alignment deficit, bounded by weighted
original plaquette energies. The selected-vacuum
source bounds pay this cost jointly for any fitting
separated family.

V38's one-link Gamma domination is used as a method,
not reapplied to a changed kernel without proof.
V39's distinction between innovations and frustration
remains essential. The result below controls
second-stage innovations, not a physical mass gap.

The V68 source and all 125 protected files were
checked unchanged. The next iteration, V70, is due
for the companion-programme review.

\subsection{The native seven-link geometry and its exact marginal}

Use unit quaternions for \(\SU(2)\).
Fix a positive retained link \(f=(x,\mu)\).
Its six parallel neighbours are
\[
 {\cal A}_f=\{(x+\varepsilon\widehat\nu,\mu):
                       \nu\ne\mu,\ \varepsilon\in\{-1,1\}\}.
\]
These six links are independent in the V38 sense:
no plaquette contains two of them.
Let the first deleted independent set \(A\)
contain \({\cal A}_f\), and let \(\rho_A\)
be the actual retained vacuum marginal.

Every plaquette containing \(f\) contains exactly
one of these neighbours. Label that shared
plaquette by \(p_i\), \(i=1,\ldots,6\).
Each neighbour has five other, private
plaquettes. There are exactly six shared
and 30 private plaquettes, all distinct.
No other link of \(A\) occurs in any of them.
The spatial side is eventually at least four,
so this local incidence has no repeated
plaquette from a short periodic identification.

Write \(v=U_f\), and \(u_i\) for the six
deleted links. After all remaining links
\(\eta\) are fixed, orient the private staples
as unit quaternions \(a_{i,r}\), \(r=1,\ldots,5\).
The shared plaquette has the form
\[
           \tfrac12\Tr U_{p_i}=u_i\cdot{\cal R}_i v
\]
for an orthogonal map \({\cal R}_i\) determined
by its two other links. This follows by
multiplying by those unit quaternions and,
if required by the orientation convention,
using quaternion inversion. Haar probability
on \(S^3\) is preserved.
Define exterior vectors
\begin{equation}
 c_i={\cal R}_i^{\mathsf T}\sum_{r=1}^5a_{i,r},
          \qquad |c_i|\le5,\qquad i=1,\ldots,6 .
 \label{eq:f15v69:fields}
\end{equation}
All their links lie outside \(A\cup\{f\}\).

For \(\gamma>0\), retain
\[
 Z(h)=\int_{S^3}e^{hu_0}\,dH(u),\qquad
                \ell(h)=\log Z(h),\qquad a(h)=\ell'(h).
\]
Integrating the six actual neighbours gives the
exact conditional density
\begin{equation}
 P_f^\eta(dv)
     =\frac{e^{\Psi_f(v,\eta)}\,dH(v)}
                   {\int_{S^3}e^{\Psi_f(w,\eta)}\,dH(w)},\qquad
 \Psi_f(v,\eta)=\sum_{i=1}^6\ell(\gamma|c_i+v|).
 \label{eq:f15v69:conditional}
\end{equation}
Indeed the affected Wilson action is
\[
   36\gamma-\sum_{i=1}^6
       u_i\cdot\left(\gamma\sum_{r=1}^5a_{i,r}
                                      +\gamma{\cal R}_iv\right).
\]
The scalar \(e^{-36\gamma}\) cancels only after
performing this same conditional normalization.
No other retained action term depends on \(v\),
because all six plaquettes at \(f\) were included.

Equation \eqref{eq:f15v69:conditional} holds first
at a finite periodic regulator and then for the
vacuum by V66--V68's finite-set conditional
limit. It keeps the actual exterior distribution.
In particular it has not substituted a Haar law
for any retained link.

Define the actual second-stage innovation by
\begin{equation}
 {\cal I}_f(v,\eta)
         =\max_{w\in S^3}\Psi_f(w,\eta)-\Psi_f(v,\eta)\ge0 .
 \label{eq:f15v69:innovation}
\end{equation}
This is precisely the generated conditional
potential above its minimum. The maximum exists
by compactness. It need not be at the reference
direction introduced below.

\subsection{Squared-radius concavity of the one-link logarithm}

\begin{proposition}[A global radial tangent inequality]
\label{f15v69:concavity}
The function \(h\mapsto a(h)/h\) is nonincreasing
on \(h>0\), with its continuous value \(1/4\)
at zero. Consequently \(s\mapsto\ell(\sqrt s)\)
is concave on \(s\ge0\).

For \(\Phi_\gamma(z)=\ell(\gamma|z|)\),
\(z\ne0\), and any \(d\in\mathbb R^4\),
\begin{equation}
 0\le
 \Phi_\gamma(z+d)-\Phi_\gamma(z)-\nabla\Phi_\gamma(z)\cdot d
       \le\frac{\gamma a(\gamma|z|)}{2|z|}|d|^2
       \le\frac{\gamma}{2|z|}|d|^2 .
 \label{eq:f15v69:tangent}
\end{equation}
There is no smallness condition on \(d\).
\end{proposition}

\begin{proof}
The sphere integration-by-parts identity used
in V67 gives
\[
               \frac{a(h)}h=\frac{1-\mathbb E_hu_0^2}{3}.
\]
Pair the even Haar density at \(u_0=t\) and
\(-t\), \(0\le t\le1\). The distribution of
\(|u_0|\) has its zero-field density tilted
by \(\cosh(ht)\). Differentiation therefore gives
\[
       \frac{d}{dh}\mathbb E_hu_0^2
             =\operatorname{Cov}_{h,\mathrm{abs}}
                                      (t^2,t\tanh(ht))\ge0 .
\]
Both functions of \(t\) are increasing.
The covariance sign follows by the elementary
two-independent-copy identity.
This proves monotonicity, and
\[
        \frac{d}{ds}\ell(\sqrt s)=\frac{a(\sqrt s)}{2\sqrt s}
\]
proves concavity, with continuous derivative
at \(s=0\).

The function \(\Phi_\gamma\) is convex on
\(\mathbb R^4\): it is the logarithm of
\(\int e^{\gamma u\cdot z}\,dH(u)\), whose
Hessian is a positive covariance matrix.
This gives the lower tangent bound.
Apply the scalar concavity just proved to
the two squared radii \(\gamma^2|z+d|^2\)
and \(\gamma^2|z|^2\).
Expanding their difference gives the upper
bound after subtracting the linear term.
Finally \(a(h)\le1\).
\end{proof}

\subsection{An all-exterior radial comparison with an explicit cost}

Let
\begin{equation}
 C_f=\sum_{i=1}^6c_i,\qquad
 {\cal D}_f=\gamma(30-|C_f|)\ge0 .
 \label{eq:f15v69:debit}
\end{equation}
If \(C_f\ne0\), set \(b_f=C_f/|C_f|\).
If \(C_f=0\), choose any fixed unit quaternion
for \(b_f\). The estimates hold for every
such choice. The scalar \({\cal D}_f\) is
gauge invariant; all \(c_i\)'s transform
in the same central-link frame.

Define the nonlinear radial comparison
\[
          \Psi_*(v)=6\ell(\gamma|5b_f+v|).
\]
For \(d_i=c_i-5b_f\), one has exactly
\[
       \sum_i d_i=-({\cal D}_f/\gamma)b_f,\qquad
       \sum_i|d_i|^2\le10{\cal D}_f/\gamma .
\]
The latter follows from
\(\sum_i|c_i|^2\le150\) and
\(\sum_i c_i\cdot b_f=|C_f|\).

\begin{theorem}[Global nonlinear alignment comparison]
\label{f15v69:alignment}
For every exterior and every \(v\in S^3\),
\begin{equation}
       -{\cal D}_f
          \le\Psi_f(v)-\Psi_*(v)
          \le\frac54{\cal D}_f .
 \label{eq:f15v69:alignment}
\end{equation}
In particular the oscillation of the difference
over the whole central group is at most
\(9{\cal D}_f/4\).
\end{theorem}

\begin{proof}
Put \(z=5b_f+v\), so \(4\le|z|\le6\).
The linear term in the sum of
\eqref{eq:f15v69:tangent} is
\[
       -({\cal D}_f/\gamma)\nabla\Phi_\gamma(z)\cdot b_f .
\]
Since
\(\nabla\Phi_\gamma(z)=\gamma a(\gamma|z|)z/|z|\)
and \(z\cdot b_f=5+v\cdot b_f>0\),
this term lies between \(-{\cal D}_f\) and zero.
The sum of the nonnegative remainders is at most
\[
               \frac{\gamma}{8}\sum_i|d_i|^2
                                      \le\frac54{\cal D}_f .
\]
This proves both bounds, including \(C_f=0\).
No chart, small-field event, or discarded
exceptional exterior is involved.
\end{proof}

The cost \({\cal D}_f\) is not V67's small
likelihood correction. It is an exterior
alignment energy for the dominant exact
conditional itself.

\subsection{The radial conditional has a Gamma-dominated innovation}

Let \(P_*\) have density proportional to
\(e^{\Psi_*}\) with respect to Haar.
Its maximum occurs at \(v=b_f\).
Writing \(t=1-b_f\cdot v\), its energy
above the minimum is
\begin{equation}
 q_\gamma(t)=6\left\{\ell(6\gamma)
                   -\ell\bigl(\gamma\sqrt{36-10t}\bigr)\right\},
                            \qquad 0\le t\le2 .
 \label{eq:f15v69:q}
\end{equation}
One has \(q_\gamma(0)=0\), \(q_\gamma'>0\),
and \(q_\gamma''\ge0\). In fact, with
\(h=\gamma\sqrt{36-10t}\),
\[
                   q_\gamma'(t)=30\gamma^2\,\frac{a(h)}h .
\]
As \(t\) increases, \(h\) decreases, and
Proposition~\ref{f15v69:concavity} gives convexity.
This is a property of the exact integrated
conditional, not of a quadratic truncation.

\begin{proposition}[Gamma domination for the nonlinear radial law]
\label{f15v69:Gamma}
Under \(P_*\), the variable \(q_\gamma(t)\)
is stochastically dominated by
\(\operatorname{Gamma}(3/2,1)\).
Thus for \(0\le s<1\),
\begin{equation}
                    \mathbb E_*e^{s q_\gamma}\le(1-s)^{-3/2}.
 \label{eq:f15v69:Gamma}
\end{equation}
This holds for every \(\gamma>0\).
\end{proposition}

\begin{proof}
Haar probability gives \(t\) a density
proportional to \(\sqrt{t(2-t)}\) on \((0,2)\).
The exact radial conditional multiplies
this by \(e^{-q_\gamma(t)}\).
Change variables from \(t\) to \(q=q_\gamma(t)\).
Its density is proportional to
\[
        e^{-q}q^{1/2}w(q),\qquad
        w(q)=\sqrt{\frac{t(q)}q}
                   \frac{\sqrt{2-t(q)}}{q_\gamma'(t(q))},
                    \quad 0<q<q_\gamma(2).
\]
Extend \(w\) by zero above \(q_\gamma(2)\).
Convexity and \(q_\gamma(0)=0\) show that
\(q_\gamma(t)/t\) and \(q_\gamma'(t)\)
are nondecreasing. Every displayed factor
in \(w\) is therefore nonincreasing in \(q\).
A Gamma law tilted by a nonnegative
decreasing weight is stochastically smaller:
the covariance of an increasing test and
that weight is nonpositive.
This is the elementary two-copy argument
used in V38, now with the new density
and its convex nonlinear radial energy
explicitly verified. Truncation gives
the exponential bound.
\end{proof}

\subsection{Actual conditional tails with an exterior debit}

Set
\[
                  \alpha=\frac{27}{4},\qquad
                  Y_f=({\cal I}_f-\alpha{\cal D}_f)_+ .
\]

\begin{theorem}[All-exterior innovation bound for the generated law]
\label{f15v69:conditional-tail}
For every exterior and \(\gamma>0\),
\begin{equation}
 \begin{gathered}
 P_f^\eta\{{\cal I}_f\ge\alpha{\cal D}_f+u\}
                         \le2^{3/2}e^{-u/2}\quad(u\ge0),\\
 \mathbb E_{P_f^\eta}e^{tY_f}
                  \le K(t):=1+\frac{2^{3/2}t}{1/2-t},
                                  \qquad 0<t<1/2 .
 \end{gathered}
 \label{eq:f15v69:conditional-tail}
\end{equation}
In particular
\(\mathbb E_{P_f^\eta}{\cal I}_f
          \le\alpha{\cal D}_f+2^{5/2}\).
The probability bound may of course be
replaced by its minimum with one.
\end{theorem}

\begin{proof}
Write \(\delta(v)=\Psi_f(v)-\Psi_*(v)\).
By Theorem~\ref{f15v69:alignment},
\(\operatorname{osc}\delta\le9{\cal D}_f/4\).
The normalized density ratio between
\(P_f^\eta\) and \(P_*\) is at most
\(e^{\operatorname{osc}\delta}\).
Also
\[
             {\cal I}_f(v)\le q_\gamma(1-b_f\cdot v)
                                               +\operatorname{osc}\delta .
\]
For \(0\le s<1\), therefore,
\[
    \mathbb E_{P_f^\eta}e^{s{\cal I}_f}
        \le e^{(1+s)9{\cal D}_f/4}(1-s)^{-3/2}.
\]
Exponential Markov at \(s=1/2\) gives
the first assertion because
\(\alpha/2=27/8\).
The tail-integral identity for
\(\mathbb E e^{tY_f}\) gives the second.
Integrating the same tail without the
exponential gives \(\mathbb EY_f\le2^{5/2}\).
All comparisons are between normalized
conditional probabilities.
\end{proof}

\subsection{The actual Wilson law pays the exterior alignment energy}

The debit can be bounded by original plaquette
energies without knowing a minimizing central
link or any deleted neighbour.
For the original fine variables, put
\(w_i={\cal R}_iv\).
For unit quaternions \(a,u,w\),
\[
       1-a\cdot w
           =\tfrac12|a-w|^2
           \le2(1-a\cdot u)+2(1-u\cdot w).
\]
Summing over the five private staples at
each neighbour gives the pointwise bound
\begin{equation}
 \begin{split}
 {\cal D}_f
   &\le\gamma\sum_{i=1}^6(5-c_i\cdot v)\\
   &\le2\sum_{p\ {\rm private}}X_p
                       +10\sum_{p\ {\rm shared}}X_p .
 \end{split}
 \label{eq:f15v69:physical-debit}
\end{equation}
The right side uses 30 coefficients equal
to two and six equal to ten. Their sum is 120.
Although the left side is exterior-measurable,
the fine variables on the right are sampled
under the original vacuum law, not redrawn.

At the selected couplings, V62 gives
\(\mathbb E_{\mu_M}e^{X_p/16}\le e^3\).
Weighted Hölder, with the weights in
\eqref{eq:f15v69:physical-debit} divided by
120, consequently yields
\begin{equation}
                       \mathbb E_{\mu_M}e^{{\cal D}_f/1920}\le e^3 .
 \label{eq:f15v69:debit-moment}
\end{equation}
No independence of the 36 plaquettes is used.

Since \({\cal I}_f\le\alpha{\cal D}_f+Y_f\),
conditioning first on the actual exterior
and using \eqref{eq:f15v69:conditional-tail}
gives, for example,
\begin{equation}
                       \mathbb E_{\rho_A}e^{{\cal I}_f/12960}
                                                     \le2e^3 .
 \label{eq:f15v69:single-moment}
\end{equation}
Indeed \(\alpha/12960=1/1920\) and
\(K(1/12960)<2\).
This is a uniform actual-law moment for the
dominant generated conditional energy.
It is not a moment of the small V67 correction.

\subsection{An arbitrary-size separated family with an intensive moment budget}

Take retained centres \(f_1,\ldots,f_N\) of a
fixed orientation whose pairwise anchor
distances exceed four. Require that the
first independent set \(A\) contain all
their six parallel neighbours.
The 36-plaquette clusters are then disjoint:
each affected plaquette base is within
distance two of its centre.
The original first-stage retained
interaction has carrier diameter at most
four, so no term meets two centres.

Conditionally on their common exterior,
their actual conditional laws factor as
\(\bigotimes_j P_{f_j}^\eta\), just as in
V68. Each \({\cal D}_{f_j}\) is measurable
from that same exterior. Thus, for
\(0<t<1/2\),
\begin{equation}
 \mathbb E\!\left[e^{t\sum_j{\cal I}_{f_j}}\,\middle|\,\eta\right]
                \le K(t)^N e^{\alpha t\sum_j{\cal D}_{f_j}}.
 \label{eq:f15v69:conditional-family}
\end{equation}
Only the central variables are conditionally
independent; the exterior debits may have
arbitrary correlations.

\begin{theorem}[Joint actual-vacuum innovation moment after star integration]
\label{f15v69:family}
For all sufficiently large selected \(M\),
every such finite family satisfies
\begin{equation}
       \log\mathbb E_{\rho_A}
             \exp\left\{\frac1{6480}\sum_{j=1}^N{\cal I}_{f_j}\right\}
                                                    \le3N .
 \label{eq:f15v69:family}
\end{equation}
There is no restriction on \(N\) other than
the stated geometry. In particular,
\begin{equation}
 \rho_A\{{\cal I}_{f_j}\ge u\text{ for every }j\}
                    \le\min\{1,\exp[-N(u/6480-3)]\}.
 \label{eq:f15v69:joint-tail}
\end{equation}
The estimate is for the selected unperturbed
vacuum, not for arbitrary physical source tilts.
\end{theorem}

\begin{proof}
Set \(t=1/6480\) and \(s=\alpha t=1/960\).
Sum \eqref{eq:f15v69:physical-debit} over
the disjoint clusters and write \(w_p=2\)
or ten for its plaquette coefficients.
Their total is \(120N\), and each plaquette
occurs once.

For a positive plaquette \(p\) based at \(y\),
\(X_p\le4T_y\) by the symmetric cell assignment.
Define the finite cell source
\[
                       b_y=4s\sum_{p:\,\operatorname{base}(p)=y}w_p .
\]
At most six positive plaquette orientations
have the same base, so
\[
       0\le b_y\le240s=\frac14,\qquad
       \sum_y b_y=480sN,\qquad
       |\operatorname{supp}b|\le36N .
\]
V62's actual-vacuum source bound gives
\[
 \log\mathbb E_{\mu_M}e^{s\sum_j{\cal D}_{f_j}}
    \le e_{\rm vac}\sum_yb_y+k_M\sum_yJ(b_y)
                          +{\cal E}_M(1/4)|\operatorname{supp}b|.
\]
Eventually \(e_{\rm vac}\le5\), \(k_M\le5\),
and \({\cal E}_M(1/4)\le1/960=s\).
For \(0\le b\le1/4\),
\[
                    J(b)\le\frac{b^2}{2(1-b)}\le\frac b6 .
\]
The preceding logarithm is therefore at most
\[
          (5+5/6)480sN+36sN=2836sN .
\]
This use of the finite-source theorem has
no support-size restriction and hence no
unstated restriction on \(N\).

Finally \(2^{3/2}<3\) and \(t\le1/14\) imply
\[
                  \log K(t)\le\frac{3t}{1/2-t}\le7t .
\]
Average \eqref{eq:f15v69:conditional-family}
under the actual exterior law. The resulting
logarithm is at most
\[
        (7t+2836s)N=\frac{19150}{6480}N<3N ,
\]
which proves \eqref{eq:f15v69:family}.
Exponential Markov on the sum of the
innovations proves \eqref{eq:f15v69:joint-tail}.
No product assumption about the exterior
was introduced at this step.
\end{proof}

\subsection{What has been obtained for the dominant interaction}

This version controls the actual generated
conditional energy on a native second-layer
geometry, rather than only a likelihood
difference from a smooth reference.
The all-exterior comparison is global on
the entire central group, and its exterior
debit is paid under the full selected
vacuum, jointly for separated families.

The radial comparison is still a strong-field
conditional. Its derivative at its minimum is
\[
                         q_\gamma'(0)=5\gamma a(6\gamma),
\]
which remains proportional to \(\gamma\)
as \(\gamma\to\infty\).
Thus this estimate does not demonstrate
entry into a weak-interaction or strong-mixing
region merely by doing one integration.

Most importantly, the innovation is energy
above a conditional minimum. It does not
control the correlations or the evolution
of the retained conditional minima
themselves. V39 already showed why
small innovations cannot simply be
substituted for frustration control.
The exterior deficit \({\cal D}_f\)
has a bounded thermal moment here; it
has not been proved to contract to zero
through a cofinal sequence of block scales.

The next physical task is to turn this
native conditional control into a
scale-uniform estimate on the remaining
interactions and their boundary response,
with all generated terms and the physical
scale retained. Neither the current
conditional tail nor its family moment
constructs an interacting continuum or
establishes a positive physical mass gap.

\subsection{Diagnostics and source preservation}

The diagnostic is
\begin{center}\small
\path{scripts/verify_ym_f15_v69_native_parallel_cluster.py}.
\end{center}
It reruns V68 and its inherited checks.
Nine outward rational radial-variance and
squared-radius concavity checks, and 24
native radial derivative comparisons, use
the positive series for \(Z(h)\) and
\(Z'(h)/h\).
There are 540 native cluster plaquette
incidence checks, four weighted
30-private/six-shared totals, 12 exterior
alignment identities, and 96 weighted
boundary-energy inequalities.
There are also 120 convex-profile
change-of-variable checks, nine radial
tail/exponential constants, and four
finite-family cell-source budgets.

The polynomial convex profiles are auxiliary
checks of the Gamma-weight algebra.
They do not replace the actual nonlinear
radial profile, whose convexity and density
were proved above.
All diagnostics use exact rational arithmetic
or outward positive-series enclosures.
No Wilson law or mass gap is certified
by numerical sampling.
The script has 6,974 bytes and SHA--256
\begin{center}\scriptsize\ttfamily
CA8EBFB043C7D3E0C9224EEEC64BD0AC8E401A690019037651A6415C83BE38D8.
\end{center}

The V68 predecessor has 1,546,952 bytes and SHA--256
\begin{center}\scriptsize\ttfamily
56D84A9FD5FAEDF63AB63D51BB7E6A836E29FF54086AA732D64D1A4F09952418.
\end{center}
Removing this final section and restoring
the title version recovers V68 exactly.
V69 replaces it as the sole active main note.
The 14 protected inputs and 111 earlier
diagnostics remain unchanged.
Only \path{.tex} files are kept in
\path{latex_notes}; all compilation products
go to
\begin{center}\small
\path{artifacts/ym_f15_v69_texcheck/}.
\end{center}
The next version, V70, includes the
scheduled companion-programme review.

\begin{firewall}
V69 proves an all-exterior conditional
energy-tail estimate with an explicit
alignment debit, and an actual-vacuum
joint exponential moment for separated
native second-layer innovations.
It does not prove scale-uniform
frustration contraction, an interacting
continuum construction, or a positive
physical Yang--Mills mass gap.
The original goal remains active.
\end{firewall}

\section{V70: whole-group conditional Dirichlet control}
\label{f15v70:section}

V69 controlled the energy above the minimum of an actual
second-layer conditional. A tail estimate alone does not
control the Dirichlet form of arbitrary test functions.
This version proves that additional estimate on the entire
central group, with an explicit exterior cost. It then
tests the resulting boundary response on a literal
gauge-transformed fine-link exterior. The two conclusions
are deliberately kept distinct.

\subsection{Companion checkpoint and nonduplication}

The scheduled V70 review found no newer versions of the
current companion files. The terminal conclusions of
SM continuum V72, NS stability V26, shell mixing V16,
parallel YM V5, SM source selection V31, native
stationarity V34 and exact-action Einstein bridge V24
were reread, together with the supplied suggestions.

SM continuum V72 concerns a fixed-time auxiliary
many-copy limit, not a growing ultraviolet cutoff.
NS V26 does not supply the integrated estimate mentioned
under its proposed next version. Shell V16 retains the
measured flat-holonomy one-loop terms, but not a joint
nonperturbative volume/coupling remainder. Parallel YM
V5 rules out the proposed volume-uniform global charge
tightness and retains local or density formulations.
The source-selection compiler remains conditional on
word-native operations. The stationarity programme still
requires the moving-fibre first-escape tensor; the
Einstein bridge's nonzero first-jet residue is not a
nonlinear continuation theorem. None supplies the
missing physical-scale YM contraction.

The suggested complete noncommuting coefficient work is
already represented in f14 V20--V30 and is not repeated.
The earlier aligned-star estimate in f6 V39 was for a
restricted original one-link conditional, with a separate
oscillation error for its complement. Here the law is
V69's nonlinear, exactly integrated conditional, and the
Dirichlet estimate has neither a chart restriction nor
an additive complement error. The general local-to-global
method is standard; the new work is its verified
application, exterior budget and response test for this
native generated law.

All 126 previously protected files, comprising 14 inputs
and 112 diagnostics, passed their recorded hash checks.
The companion review does not change the physical target
or supply an additional mathematical assumption. The
next review is V75.

\subsection{Normalization and the radial slope}

Use V69's central retained link \(v\in S^3=\SU(2)\), its
six integrated parallel neighbours, and their actual
common exterior \(\eta\). Write
\begin{equation}
 \begin{gathered}
 P^\eta(dv)=
  \frac{e^{\Psi_\eta(v)}\,dH(v)}
       {\int e^{\Psi_\eta(w)}\,dH(w)},\qquad
 \Psi_\eta(v)=\sum_{i=1}^6\ell(\gamma|c_i+v|),\\
 C=\sum_i c_i,\quad D=\gamma(30-|C|),\quad
 b=C/|C|\quad(C\ne0).
 \end{gathered}
 \label{eq:f15v70:law}
\end{equation}
When \(C=0\), choose any unit \(b\). The theorem below
does not depend on that choice. Each \(c_i\) is a sum
of five transported unit staples, so \(|c_i|\le5\).
All norms on quaternions are Euclidean.

Equip the sphere with its unit-radius round metric.
Thus the three left-invariant imaginary-quaternion
directions are orthonormal, the Ricci tensor is \(2g\),
and a coordinate \(v_\alpha\) obeys
\(\Delta v_\alpha=-3v_\alpha\).
The positive Dirichlet energy is
\(\int|\nabla F|^2\,dP\), with no extra coupling factor.

Let \(P_*\) be the radial probability with potential
\(\Psi_*(v)=6\ell(\gamma|5b+v|)\). Set
\[
 z=b\cdot v,\qquad t=1-z,\qquad
 q(t)=6\{\ell(6\gamma)
                   -\ell(\gamma\sqrt{36-10t})\}.
\]
Then \(P_*(dv)\) is proportional to \(e^{-q(t)}dH(v)\).
V69 proves
\begin{equation}
 q''\ge0,\qquad
 \operatorname{osc}_{S^3}(\Psi_\eta-\Psi_*)\le\frac94D,
 \qquad q'(0)=5\gamma a(6\gamma).
 \label{eq:f15v70:inputs}
\end{equation}
Here \(a(h)=\ell'(h)\). V38's Gamma domination gives
\[
       h(1-a(h))=\mathbb E_h[h(1-u_0)]\le\frac32.
\]
Consequently, throughout this version's range
\(\gamma\ge2\),
\begin{equation}
       q'(t)\ge5\gamma-\frac54\ge4\gamma
                         \quad(0\le t\le2).
 \label{eq:f15v70:slope}
\end{equation}
This lower bound is for the exact nonlinear radial law,
not a Hessian truncation at \(v=b\).

\subsection{A global radial Poincare estimate}

\begin{theorem}[Whole-sphere radial Dirichlet bound]
\label{f15v70:radial-gap}
For every \(\gamma\ge2\), every unit \(b\), and every
real \(F\in H^1(P_*)\),
\begin{equation}
       \Var_{P_*}(F)
          \le \frac{11}{6\gamma}
                         \int|\nabla F|^2\,dP_* .
 \label{eq:f15v70:radial-gap}
\end{equation}
\end{theorem}

\begin{proof}
First take the closed cap \(K=\{z\ge1/2\}\), whose
geodesic radius is \(\pi/3\), and normalize \(P_*\) on
\(K\). The sphere identities
\[
 |\nabla t|^2=1-z^2,\qquad
 \Delta t=3z,\qquad \operatorname{Hess}t=zg
\]
give on this cap
\[
 \operatorname{Ric}+\operatorname{Hess}(q(t))
  =2g+q''(t)\,dt\otimes dt+q'(t)zg
                      \succeq(2+2\gamma)g .
\]
The outward second fundamental form of the cap is
\(\cot(\pi/3)g_{\partial K}\succeq0\).

For clarity, the boundary condition needed here is
Neumann, not Dirichlet. With
\(L=\Delta-\nabla q\cdot\nabla\), integration of the
weighted Bochner identity for a Neumann eigenfunction
\(u\) gives
\[
 \int_K(Lu)^2\,dP_*
 =\int_K\!\left\{
       |\operatorname{Hess}u|^2+
       (\operatorname{Ric}+\operatorname{Hess}q)
                                      (\nabla u,\nabla u)\right\}dP_*
   +\int_{\partial K}\mathrm{II}
           (\nabla_{\partial K}u,\nabla_{\partial K}u)\,d\sigma_* .
\]
Here \(d\sigma_*\) has the same positive density as
\(dP_*\) relative to area. The sign of the boundary
term follows by differentiating \(\partial_nu=0\)
tangentially. For \(-Lu=\lambda u\), the identity
implies \(\lambda\ge2+2\gamma\) for every nonconstant
eigenfunction. The smooth positive weight on the compact
connected cap has the usual discrete Neumann spectrum:
compactness of \(H^1(K)\) in \(L^2(K)\), followed by
elliptic regularity, justifies the variational argument.
Expansion in this spectrum gives
\begin{equation}
 \int_K(F-\mathbb E_{P_*|K}F)^2\,dP_*
       \le\frac1{2+2\gamma}
                          \int_K|\nabla F|^2\,dP_* .
 \label{eq:f15v70:cap}
\end{equation}
Normalization on \(K\) cancels from both sides.

To include the rest of the sphere, take the globally
smooth positive function \(W=e^{2\gamma t}\). Direct
differentiation and \eqref{eq:f15v70:slope} yield
\[
 -\frac{LW}{W}
   =2\gamma(q'-2\gamma)(1-z^2)-6\gamma z
   \ge p_\gamma(z):=4\gamma^2(1-z^2)-6\gamma z .
\]
Everywhere \(p_\gamma(z)\ge-6\gamma\). If \(z\le1/2\),
the exact factorization
\[
 p_\gamma(z)-3\gamma
   =\gamma\{(1-2z)(4z+5)
                    +4(\gamma-2)(1-z^2)\}\ge0
\]
holds because \(-1\le z\le1/2\) and \(\gamma\ge2\).
Thus, up to the irrelevant cap boundary,
\begin{equation}
             -LW/W\ge3\gamma-9\gamma\,1_K .
 \label{eq:f15v70:drift}
\end{equation}
In particular the antipode has positive drift; it is
not covered by an erroneous global convexity assertion.

For any smooth real \(f\), integration by parts on
the boundaryless sphere and a completed square give
\[
 \begin{split}
 \int f^2(-LW/W)\,dP_*
 &=\int\{2f\,\nabla f\cdot\nabla\log W
                     -f^2|\nabla\log W|^2\}\,dP_*\\
 &\le\int|\nabla f|^2\,dP_* .
 \end{split}
\]
Choose \(f=F-\mathbb E_{P_*|K}F\), and combine
\eqref{eq:f15v70:cap}--\eqref{eq:f15v70:drift}. This gives
\[
 \begin{split}
 3\gamma\int f^2\,dP_*
 &\le\int|\nabla F|^2\,dP_*
                 +9\gamma\int_K f^2\,dP_*\\
 &\le\left(1+\frac{9\gamma}{2+2\gamma}\right)
                                  \int|\nabla F|^2\,dP_*\\
 &\le\frac{11}{2}\int|\nabla F|^2\,dP_* .
 \end{split}
\]
Finally \(\Var_{P_*}(F)\le\int f^2dP_*\). Smooth
approximation extends the estimate to \(H^1(P_*)\).
\end{proof}

This uses the standard local-to-global Lyapunov
argument; see Cattiaux--Guillin,
\href{https://arxiv.org/abs/1001.1822}
{Functional Inequalities via Lyapunov conditions},
Lemma 1.1 and Section 3.1.
For the weighted Neumann boundary identity see
Kolesnikov--Milman,
\href{https://arxiv.org/abs/1310.2526}
{Brascamp--Lieb type inequalities on weighted
Riemannian manifolds with boundary}, Theorem 1.1.
The slope, cap, drift function, polynomial certificate
and numerical constants used here have been verified
explicitly above. No spectral assertion about a
four-dimensional physical Hamiltonian is imported
from these references.

\subsection{The actual generated conditional and its exterior cost}

Let \(C_{\rm P}(\eta)\) be the optimal Poincare constant
of \(P^\eta\) for the round metric, and put
\(\lambda(\eta)=C_{\rm P}(\eta)^{-1}\).
Equivalently, \(\lambda(\eta)\) is the first positive
eigenvalue of
\[
       -L_\eta,\qquad
       L_\eta=\Delta+\nabla\Psi_\eta\cdot\nabla .
\]
At each fixed finite coupling these definitions are
legitimate because the density is smooth and strictly
positive on the compact connected sphere. The apparent
norm singularity when \(c_i+v=0\) is removable:
\(\ell(\gamma|x|)\) is a smooth function of \(|x|^2\).

\begin{theorem}[All-exterior native conditional coercivity]
\label{f15v70:actual-gap}
For the exact conditional \eqref{eq:f15v70:law} and
every \(\gamma\ge2\),
\begin{equation}
 \begin{gathered}
   \Var_{P^\eta}(F)
       \le\frac{11}{6\gamma}e^{9D(\eta)/4}
                                  \int|\nabla F|^2\,dP^\eta,\\
   \frac{\lambda(\eta)}{\gamma}
                          \ge\frac6{11}e^{-9D(\eta)/4}.
 \end{gathered}
 \label{eq:f15v70:actual-gap}
\end{equation}
There is no exceptional exterior or restriction on
the central variable in this assertion.
\end{theorem}

\begin{proof}
Put \(\delta=\Psi_\eta-\Psi_*\) and
\(r=dP^\eta/dP_*=e^\delta/\mathbb E_*e^\delta\).
For \(r_+=\sup r\) and \(r_-=\inf r>0\),
\[
 \begin{split}
 \Var_{P^\eta}(F)
 &\le r_+\Var_{P_*}(F)\\
 &\le r_+\frac{11}{6\gamma}
                                  \int|\nabla F|^2\,dP_*\\
 &\le\frac{r_+}{r_-}\frac{11}{6\gamma}
                                  \int|\nabla F|^2\,dP^\eta .
 \end{split}
\]
The first line chooses the \(P_*\)-mean in the
variational definition of variance.
The ratio \(r_+/r_-\) is exactly
\(e^{\operatorname{osc}\delta}\), bounded by \(e^{9D/4}\)
by V69. Only one oscillation factor is needed; the
normalizing constant has not been counted twice.
The eigenvalue characterization proves the second line
of \eqref{eq:f15v70:actual-gap}.
\end{proof}

For the thermally rescaled auxiliary diffusion
\(\gamma^{-1}L_\eta\), the gap is
\(g(\eta)=\lambda(\eta)/\gamma\).
Its time variable is a conditional sampling time.
There is no identification with Euclidean lattice
time, transfer time, or physical Hamiltonian time.

\begin{proposition}[Actual-vacuum fractional cost and a weak estimate]
\label{f15v70:annealed}
Use the selected Wilson vacuum and the first exact
marginal of V69. Let \(\nu(d\eta)\) be the actual
exterior law of this conditional. Eventually along
that selection, for \(0<r\le1/4320\),
\begin{equation}
       \mathbb E_\nu[(\gamma C_{\rm P}(\eta))^r]
                            \le(11/6)^r e^3 .
 \label{eq:f15v70:fractional}
\end{equation}
For \(0<a\le6/11\), in particular,
\begin{equation}
       \nu\{g(\eta)<a\}
            \le\min\{1,e^3(11a/6)^{1/4320}\}.
 \label{eq:f15v70:gap-tail}
\end{equation}

If \(F=F(\eta,v)\) is measurable, smooth in \(v\),
and has uniformly bounded conditional oscillation
\[
    O=\operatorname*{ess\,sup}_{\eta}
             \{\sup_v F(\eta,v)-\inf_v F(\eta,v)\},
\]
then for every \(0<\varepsilon\le1/4\),
\begin{equation}
 \begin{split}
 \mathbb E_\nu\Var_{P^\eta}(F)
 &\le \frac{11}{6\gamma}
         \left(\frac{e^3}{4\varepsilon}\right)^{4320}
            \mathbb E_\nu\mathbb E_{P^\eta}|\nabla_vF|^2
                  +\varepsilon O^2 .
 \end{split}
 \label{eq:f15v70:weak}
\end{equation}
Only the conditional variance is controlled.
\end{proposition}

\begin{proof}
The inherited estimate
\eqref{eq:f15v69:debit-moment} is
\(\mathbb E_\nu e^{D/1920}\le e^3\); \(D\) is
exterior measurable, so its law is unaffected by
integrating the central or satellite variables.
The identity
\[
                  \frac94\,\frac1{4320}=\frac1{1920}
\]
and \eqref{eq:f15v70:actual-gap} prove
\eqref{eq:f15v70:fractional}.
If \(g<a\), the same lower bound on \(g\) requires
\(D>(4/9)\log(6/(11a))\).
Exponential Markov proves \eqref{eq:f15v70:gap-tail}.

For any \(R\ge0\), use the Dirichlet estimate on
\(\{D\le R\}\). On its complement use
\(\Var_{P^\eta}(F)\le O^2/4\). Averaging under the
actual exterior law gives
\[
 \mathbb E_\nu\Var_{P^\eta}(F)
 \le\frac{11}{6\gamma}e^{9R/4}
                  \mathbb E_\nu\mathbb E_{P^\eta}|\nabla_vF|^2
             +\frac{e^3}{4}e^{-R/1920}O^2 .
\]
Set \(R=1920\log(e^3/(4\varepsilon))>0\).
This proves \eqref{eq:f15v70:weak}.
\end{proof}

The small exponent is part of the conclusion, not
notation concealing a stronger estimate.
The available exponential budget \(1/1920\) does
not pay the coefficient \(9/4\) at exponent one.
Thus this argument does not prove a coupling-uniform
first moment of \(\gamma C_{\rm P}\), much less an
essential-supremum bound. It also does not prove that
such stronger estimates are false.

For a separated family of V69 clusters with common
exterior \(\eta\), the actual central conditional is
a product. Tensorization consequently gives the
useful weighted form
\begin{equation}
 \Var_{P_1^\eta\otimes\cdots\otimes P_N^\eta}(F)
 \le\frac{11}{6\gamma}\sum_{j=1}^N
             e^{9D_j(\eta)/4}
               \mathbb E_{P_1^\eta\otimes\cdots\otimes P_N^\eta}
                                        |\nabla_jF|^2 .
 \label{eq:f15v70:family}
\end{equation}
To verify tensorization, apply the two-factor
conditional-variance identity, bound the first
conditional variance, and use
\(|\nabla\mathbb E F|^2\le\mathbb E|\nabla F|^2\)
for the second factor. Induction gives the displayed
inequality. This proof uses precisely the native
conditional product supplied by the separation;
it asserts no independence of the \(D_j\)'s or of
the exterior. Averaging the weighted inequality is
valid whenever its right side is integrable, or as
an inequality of nonnegative extended integrals.

\subsection{The thermal order is sharp on the aligned exterior}

The lower bound above has order \(\gamma\) at \(D=0\).
The following test shows that this order cannot be
improved to a higher power of the coupling.

\begin{proposition}[Aligned radial moments and an upper gap test]
\label{f15v70:radial-limit}
Under \(P_*\), the variable \(X_\gamma=\gamma(1-b\cdot v)\)
converges in distribution and in every fixed
nonnegative integer moment to
\(\operatorname{Gamma}(3/2,5)\), with rate five.
For every unit \(w\perp b\), let \(F_w(v)=w\cdot v\).
Then
\begin{equation}
 \begin{gathered}
 m_\gamma:=\mathbb E_*(b\cdot v)
                  =1-\frac{3}{10\gamma}+o(\gamma^{-1}),\\
 \sigma_\gamma^2:=\mathbb E_*F_w^2,\qquad
            \gamma\sigma_\gamma^2\longrightarrow\frac15,\\
        \frac6{11}\gamma
                    \le\lambda_*(\gamma)\le(5+o(1))\gamma .
 \end{gathered}
 \label{eq:f15v70:radial-limit}
\end{equation}
\end{proposition}

\begin{proof}
The exact density of \(X_\gamma\), up to normalization,
is
\[
  e^{-q(x/\gamma)}x^{1/2}
           (2-x/\gamma)^{1/2}1_{\{0<x<2\gamma\}} .
\]
For fixed \(x\ge0\),
\[
   \gamma^{-1}q'(x/\gamma)
       =\frac{30a(\gamma\sqrt{36-10x/\gamma})}
                         {\sqrt{36-10x/\gamma}}\longrightarrow5 .
\]
Indeed \(a(h)\to1\) follows from
\(0\le h(1-a(h))\le3/2\). On each fixed interval in
\(x\), the displayed derivatives are bounded, so
integration yields \(q(x/\gamma)\to5x\).
Furthermore \eqref{eq:f15v70:slope} gives
\(q(x/\gamma)\ge4x\) on the full support. The
unnormalized density, multiplied by \(x^k\), is
bounded by \(\sqrt2\,e^{-4x}x^{k+1/2}\).
Dominated convergence applies both to the positive
normalizer and to every indicated moment.

The first two limiting moments are
\[
              \mathbb EX_\gamma\longrightarrow\frac3{10},
       \qquad \mathbb EX_\gamma^2\longrightarrow\frac3{20}.
\]
Rotational symmetry in the three tangent coordinates
gives \(\mathbb E_*F_w=0\) and
\[
       \sigma_\gamma^2
          =\frac13\mathbb E_*(1-z^2)
          =\frac23\mathbb E_*t-\frac13\mathbb E_*t^2 .
\]
These formulas prove the moment claims. Finally
\(|\nabla F_w|^2=1-F_w^2\), so its Rayleigh quotient is
\((1-\sigma_\gamma^2)/\sigma_\gamma^2=(5+o(1))\gamma\).
The lower bound is Theorem~\ref{f15v70:radial-gap}.
\end{proof}

This is a finite-cluster conditional limit. It is not
a new continuum Gaussian-field construction, and it
does not replace the full-law distinctions already
made in V62--V65.

\subsection{A native coherent boundary response survives}

The exterior parameter cannot be treated as thirty
freely adjustable independent staples: those staples
come from shared fine links. The following path is
instead realized by an actual fine-link gauge action.

Start with the identity exterior \(\eta_0\) for the
central link \(f=(x,\mu)\) and its six satellites.
It gives \(c_i=5\,1\) and hence \(P^{\eta_0}=P_*\)
with \(b=1\). Fix a unit imaginary quaternion \(E\),
let \(h_s=e^{sE}\), and perform the lattice gauge
transformation equal to \(h_s\) at \(x\) and to the
identity at every other vertex. Let \(\eta_s\) be
the transformed exterior, leaving the central
variable available for conditioning.

\begin{proposition}[Exact response along a realizable exterior path]
\label{f15v70:gauge-response}
For this path the complete generated conditional
satisfies
\begin{equation}
       P^{\eta_s}=(v\mapsto h_sv)_*P^{\eta_0},
       \qquad \mathbb E_{P^{\eta_s}}v=m_\gamma h_s .
 \label{eq:f15v70:gauge-law}
\end{equation}
In particular, for the real one-Lipschitz observable
\(F_E(v)=E\cdot v\),
\begin{equation}
  \left.\frac{d}{ds}\mathbb E_{P^{\eta_s}}F_E\right|_{s=0}
            =m_\gamma
            =1-\frac{3}{10\gamma}+o(\gamma^{-1}).
 \label{eq:f15v70:response}
\end{equation}
\end{proposition}

\begin{proof}
Under the gauge action, a fine link \(U_{yz}\) maps
to \(g_yU_{yz}g_z^{-1}\). The central link maps to
\(h_sv\). Every Wilson plaquette trace is invariant,
and Haar measure on each of the six integrated
satellites is invariant under its corresponding
left and right multiplication. Changing those six
integration variables in the exact conditional
density therefore proves the first identity,
including its normalizer. Terms independent of the
central link cancel when it is normalized.

At the identity exterior the radial symmetry gives
\(\mathbb E v=m_\gamma\,1\). Quaternion multiplication
is linear on \(\mathbb R^4\), so pushforward gives
the second identity. Since \(h_s=\cos s+E\sin s\),
the expectation of \(F_E\) is \(m_\gamma\sin s\).
Differentiate and use
Proposition~\ref{f15v70:radial-limit}.
\end{proof}

One can express the limitation precisely. On exteriors
use the supremum of the round geodesic distances of
their retained links, and on central laws use
\(W_1\) for the same round metric. For \(|s|\) small,
\[
 d_\infty(\eta_s,\eta_0)=|s|,\qquad
 W_1(P^{\eta_s},P^{\eta_0})\ge m_\gamma|\sin s| .
\]
For the first equality each changed exterior link is
\(h_s\) or \(h_s^{-1}\), and at least one such link is
present. The second inequality follows by testing
the one-Lipschitz function \(F_E\); no optimal
coupling is being assumed. Thus the local
exterior-to-conditional Lipschitz coefficient in
this unweighted supremum metric is at least
\(m_\gamma\), tending to one.

There can therefore be no constant \(\kappa<1\),
independent of large \(\gamma\), bounding that
all-exterior Lipschitz coefficient for these raw
link kernels. This statement does not exclude
strict contraction for a fixed coupling, contraction
in a different appropriately defined norm, or
contraction on gauge-invariant physical data.
The exhibited direction is a gauge direction and
the test observable is not gauge invariant.
It is not evidence for a gapless physical sector.

\subsection{What this resolves and what must still be paid}

There is now a whole-group Dirichlet estimate for
the actual dominant native second-layer conditional.
Its inverse thermal gap has a proved small
fractional moment under the selected actual vacuum.
This is stronger than a conditional energy tail or
a variance estimate using only a uniform gradient
bound. The native separated-family inequality
retains a cost for each exterior.

The remaining variance decomposition is still
\[
 \Var_{\rho_A}(F)
     =\mathbb E_\nu\Var_{P^\eta}(F)
                   +\Var_\nu(\mathbb E_{P^\eta}F).
\]
No theorem here bounds the second term by a
scale-uniform physical Dirichlet form.
The coherent response calculation explains why
merely multiplying a thermal conditional gap by
a boundary derivative cannot supply the missing
strict contraction in raw link variables.

An upstream continuation must therefore retain
the generated measure while controlling the
gauge-invariant boundary response and the
distribution of poorly aligned boundary costs,
or give a different capacity/gluing estimate
which genuinely pays for them. The available
fractional moment cannot silently be substituted
for an \(L^p\), \(p>1\), inverse-gap estimate in
such a gluing argument. None of these tasks is
discharged by the auxiliary sampling-time gap.

\subsection{Diagnostics and source preservation}

The diagnostic is
\begin{center}\small
\path{scripts/verify_ym_f15_v70_native_conditional_dirichlet.py}.
\end{center}
It reruns V69 and its inherited checks. The new
checks include one exact continuous drift
polynomial certificate, five full-polynomial
coupling decompositions, 165 rational drift and
cap-curvature checks, five local-to-global
constant checks, and 24 outward native radial
slope enclosures. There are three fractional
moment budget checks and an explicit check that
the first-moment budget is not available, plus
six limiting radial moment recurrences.

A literal endpoint gauge action is checked on
all 36 plaquettes of the native cluster: there
are 1,728 exact trace identities, 48 central
left-translation identities and 384 nonlinear
radial gauge identities. These use rational
unit quaternions, not independently assigned
staple controls. The diagnostic does not infer
a functional inequality from finite sampling;
the whole-group proof is the analytic argument
above.

The script has 6,139 bytes and SHA--256
\begin{center}\scriptsize\ttfamily
20B2A44500764B615738560AD20D3DDF443CB55F30A2538FBB91236FE2C85498.
\end{center}
The V69 predecessor has 1,567,890 bytes and SHA--256
\begin{center}\scriptsize\ttfamily
476E640A252F77B39AA13EBA9D69B11ED8E9B78E6AC0C1691AC1DC0E94600BE9.
\end{center}
Removing this final section and restoring the
title version recovers V69 exactly. V70 replaces
V69 as the sole active main note; archived
lineages and companion files are untouched.
Only \path{.tex} files are kept in
\path{latex_notes}. All compilation products go to
\begin{center}\small
\path{artifacts/ym_f15_v70_texcheck/}.
\end{center}

\begin{firewall}
V70 proves a full-group Poincare bound for the
native generated conditional, a paid
actual-vacuum fractional inverse-gap bound,
and an exact coherent gauge-boundary response.
Its gap is an auxiliary conditional diffusion
gap. It is not a transfer-Hamiltonian gap or a
physical Yang--Mills mass gap.
Scale-uniform physical contraction and a
nontrivial interacting continuum construction
remain unproved. The original goal remains active.
\end{firewall}

\section{V71: native coercivity and paid bad regions}
\label{f15v71:section}

V70's exponential comparison proves a conditional
Dirichlet bound at every exterior, but its cost
\(e^{9D/4}\) is not paid at exponent one by the
available vacuum moment of \(D\). Repeating that
comparison cannot repair the mismatch. This version
goes upstream to the Hessian and drift of the
complete generated potential itself.

The result is a thermal Poincare constant bounded
uniformly on a specified aligned-boundary set, with
no exponential loss in the size of its debit.
Its complement is then charged under the actual
selected vacuum, jointly even for overlapping
native clusters. An exact all-orientation overlap
count supplies an extensive bad-set penalty and
hence a bound on connected bad regions.
This does not yet control the retained conditional
means or the gauge-invariant physical dynamics.

The current V70 source and all 127 protected inputs
and diagnostics were checked unchanged. The last
companion review was V70; the next is V75.
The generic lattice-animal counting argument was
already used in f8 V8 and is not claimed as a new
result. What is new here is the direct conditional
coercivity and the actual-vacuum probability bound
for this particular native, measurable bad label.

\subsection{A geometric good-boundary condition}

Use the exact six-neighbour integrated conditional
of \eqref{eq:f15v70:law}. Abbreviate its potential by
\[
 \Psi(v)=\sum_{i=1}^6\Phi_\gamma(c_i+v),\qquad
 \Phi_\gamma(y)=\ell(\gamma|y|),\qquad V(v)=-\Psi(v).
\]
Thus \(P^\eta(dv)\) is proportional to \(e^{-V(v)}dH(v)\).
The metric is V70's unit-radius round metric on \(S^3\).

Call the exterior good when
\begin{equation}
          G_f=\{D_f\le\gamma/10\},\qquad
          D_f=\gamma(30-|\textstyle\sum_i c_i|).
 \label{eq:f15v71:good}
\end{equation}
On \(G_f\), the aggregate is nonzero. Put
\(b=(\sum_i c_i)/|\sum_i c_i|\) and \(d_i=c_i-5b\).
The exact alignment identity from V69 gives
\[
             \sum_i|d_i|^2\le10D_f/\gamma\le1 .
\]
In particular each \(|d_i|\le1\), while the original
staple construction still gives \(|c_i|\le5\).
For all central variables \(v\in S^3\), therefore,
\begin{equation}
            3\le r_i:=|c_i+v|\le6 .
 \label{eq:f15v71:radii}
\end{equation}
This conclusion holds on the entire central sphere,
not just near its conditional minimum.

Define
\[
              k_i(v)=\frac{\gamma a(\gamma r_i)}{r_i},
           \qquad a(h)=\ell'(h).
\]
For \(\gamma\ge2\), V38's
\(a(h)\ge1-3/(2h)\) and \(a(h)\le1\) imply
\begin{equation}
                   \gamma/8\le k_i(v)\le\gamma/3 .
 \label{eq:f15v71:k}
\end{equation}
Indeed \(a(\gamma r_i)\ge1-1/(2\gamma)\ge3/4\).

The radial Hessian has tangential eigenvalue \(k_i\)
and radial eigenvalue \(\gamma^2a'(\gamma r_i)\).
Convexity of the logarithmic moment generating
function makes the latter nonnegative; V69's
monotonicity of \(a(h)/h\) makes it at most \(k_i\).
Consequently,
\begin{equation}
    0\preceq \operatorname{Hess}_{\mathbb R^4}
                    \Phi_\gamma(c_i+v)\preceq k_i(v)I .
 \label{eq:f15v71:ambient}
\end{equation}
These are estimates on the exact integrated one-link
logarithm. No expansion in \(v\) or in the exterior
fields has been made.

\subsection{Direct curvature and drift for the complete potential}

Let \(z=b\cdot v\), let \(\Pi_v=I-vv^{\mathsf T}\),
and write \(T(v)=\sum_i k_i(v)\). Then
\begin{equation}
 \begin{gathered}
       T(v)\ge3\gamma/4,\qquad
       \nabla V=-\sum_i k_i\Pi_v c_i,\\
 \operatorname{Hess}_{S^3} V(X,X)
   =\sum_i\left\{k_i(1+c_i\cdot v)|X|^2
       -\operatorname{Hess}_{\mathbb R^4}\Phi_\gamma(c_i+v)
                                                   (X,X)\right\}
 \end{gathered}
 \label{eq:f15v71:derivatives}
\end{equation}
for \(X\perp v\). The second formula includes the
sphere's normal-curvature term.

\begin{proposition}[Three-region control of the actual potential]
\label{f15v71:regions}
On \(G_f\), for \(\gamma\ge2\), the following hold:
\begin{equation}
 \begin{array}{ll}
  z\ge1/2:
     &\operatorname{Hess}_{S^3}V\succeq(9\gamma/8)g,\\[1mm]
  |z|\le1/2:
     &|\nabla V|\ge9\gamma/4,\\[1mm]
  z\le-1/2:
     &\Delta V\le-9\gamma/8 .
 \end{array}
 \label{eq:f15v71:regions}
\end{equation}
Everywhere on \(S^3\), one also has \(\Delta V\le18\gamma\).
\end{proposition}

\begin{proof}
On the northern cap, \(c_i\cdot v=5z+d_i\cdot v\ge3/2\).
By \eqref{eq:f15v71:ambient}--\eqref{eq:f15v71:derivatives},
\[
      \operatorname{Hess}V(X,X)
            \ge\sum_i k_i(c_i\cdot v)|X|^2
            \ge\frac32 T(v)|X|^2
            \ge\frac{9\gamma}{8}|X|^2 .
\]

In the middle band,
\[
 \begin{split}
 |\nabla V|
   &=\left|5T(v)\Pi_v b+\sum_i k_i\Pi_v d_i\right|\\
   &\ge T(v)\{5\sqrt{1-z^2}-1\}
    \ge3T(v)\ge9\gamma/4 .
 \end{split}
\]
The middle inequality uses
\(25(1-z^2)\ge75/4\ge16\). It does not require that
the \(k_i\)'s are equal.

Taking a tangent trace in
\eqref{eq:f15v71:derivatives} gives
\[
       \Delta V\le3\sum_i k_i(1+c_i\cdot v).
\]
On the southern cap,
\(1+c_i\cdot v\le2+5z\le-1/2\). Hence
\(\Delta V\le-(3/2)T(v)\le-9\gamma/8\).
For the global upper bound, instead use
\[
       k_i(1+c_i\cdot v)
           =\nabla\Phi_\gamma(c_i+v)\cdot v
           \le|\nabla\Phi_\gamma(c_i+v)|\le\gamma .
\]
Summing its six contributions proves \(\Delta V\le18\gamma\).
\end{proof}

The southern estimate is concavity of the potential
at a high-energy region, not positivity of its
Hessian there. It is exactly the sign needed for
the global drift argument.

\begin{theorem}[Uniform thermal conditional gap on good boundaries]
\label{f15v71:good-gap}
For every \(\gamma\ge8\), every good native exterior,
and every real \(F\in H^1(P^\eta)\),
\begin{equation}
        \Var_{P^\eta}(F)
          \le\frac{152}{9\gamma}
                      \mathbb E_{P^\eta}|\nabla F|^2
          \le\frac{17}{\gamma}
                      \mathbb E_{P^\eta}|\nabla F|^2
 \label{eq:f15v71:good-gap}
\end{equation}
Equivalently the conditional diffusion gap satisfies
\(\lambda_f(\eta)/\gamma\ge9/152\).
\end{theorem}

\begin{proof}
On \(K=\{z\ge1/2\}\), the weighted Ricci tensor
is at least \(2g+(9\gamma/8)g\), and the outward
boundary of the radius-\(\pi/3\) cap is convex.
V70's explicitly justified weighted Neumann argument,
now applied to this actual potential, gives
\begin{equation}
 \int_K (F-\mathbb E_{P^\eta|K}F)^2\,dP^\eta
          \le\frac8{9\gamma}
                            \int_K|\nabla F|^2\,dP^\eta .
 \label{eq:f15v71:cap}
\end{equation}

For \(L=\Delta-\nabla V\cdot\nabla\), take
\(W=\exp\{(V-\min V)/2\}\). Direct differentiation gives
\[
             -LW/W=\tfrac14|\nabla V|^2-\tfrac12\Delta V.
\]
This is at least \(-9\gamma\) everywhere.
On the southern cap it is at least \(9\gamma/16\).
In the middle band it is at least
\[
          \frac{81}{64}\gamma^2-9\gamma
                                 \ge\frac98\gamma
                  \quad(\gamma\ge8).
\]
Consequently
\begin{equation}
       -LW/W\ge\frac{9\gamma}{16}
                           -\frac{153\gamma}{16}1_K .
 \label{eq:f15v71:drift}
\end{equation}
There is no ignored belt or antipodal exception.

For \(f=F-\mathbb E_{P^\eta|K}F\), the same integration
by parts as in V70 gives
\(\int f^2(-LW/W)dP^\eta\le\int|\nabla F|^2dP^\eta\).
Combining \eqref{eq:f15v71:cap} and
\eqref{eq:f15v71:drift} yields
\[
 \begin{split}
 \frac{9\gamma}{16}\int f^2\,dP^\eta
   &\le\left(1+\frac{153\gamma}{16}\frac8{9\gamma}\right)
                               \int|\nabla F|^2\,dP^\eta\\
   &=\frac{19}{2}\int|\nabla F|^2\,dP^\eta .
 \end{split}
\]
Since \(\Var_{P^\eta}(F)\le\int f^2dP^\eta\), the
constant is \(152/(9\gamma)<17/\gamma\).
Density of smooth functions gives the \(H^1\) assertion.
\end{proof}

The bound is uniform even when \(D_f\) grows linearly
with \(\gamma\) up to the threshold in
\eqref{eq:f15v71:good}. In contrast, V70's comparison
would have paid an exponential factor in that range.
V70 remains the all-exterior theorem; this theorem
is its sharper good-boundary complement.

\subsection{An exponentially small actual-vacuum exceptional term}

Write \(\mu_M\) for the selected vacuum at coupling
\(\gamma_M\) used in V62--V70, and let \(\gamma=\gamma_M\).
The vacuum is on the spatially periodic, infinite-time
cylinder of that construction. No new infinite-spatial-volume
measure is assumed here. Local counts are valid once the
spatial side is at least four.

V69's moment bound gives, eventually along this
selection and uniformly in the central link,
\begin{equation}
     \mu_M(G_f^c)
       \le \min\{1,e^3e^{-\gamma/19200}\}.
 \label{eq:f15v71:one-bad}
\end{equation}
Indeed \(D_f>\gamma/10\) on the bad event and
\(\mathbb E_{\mu_M}e^{D_f/1920}\le e^3\).

Let \(\nu_f\) be the actual common exterior law after
the allowed first integration, and let \(F=F(\eta,v)\)
have conditional oscillation bounded by \(O\).
Theorem~\ref{f15v71:good-gap} and
\(\Var_{P^\eta}(F)\le O^2/4\) imply
\begin{equation}
 \mathbb E_{\nu_f}\Var_{P^\eta}(F)
       \le \frac{17}{\gamma}
                  \mathbb E_{\nu_f}\mathbb E_{P^\eta}|\nabla_vF|^2
              +\frac{e^3}{4}e^{-\gamma/19200}O^2 .
 \label{eq:f15v71:annealed}
\end{equation}
To see this, split the left side into \(G_f\) and
\(G_f^c\), use the respective estimates, and remove
the indicator from the nonnegative gradient term.
The distribution of \(D_f\) is the same under the
fine vacuum and this exterior marginal.

Thus an exponentially small additive error is
available with a fixed thermal Dirichlet coefficient.
This conclusion neither proves a uniform first
moment of the inverse gap on bad exteriors nor
bounds the variance of \(\mathbb E_{P^\eta}F\).
It avoids claiming either.

\subsection{The exact all-orientation overlap count}

For every positive fine link \(f\), define \(D_f\)
by its own V69 native cluster. These are simultaneous
local functions of a fine configuration under
\(\mu_M\), even when the clusters overlap.
For the moment estimates below it is not necessary
to delete all their satellites at once.

Let
\[
 w(f,p)=
 \begin{cases}
   10,&p\text{ is shared in the cluster at }f,\\
    2,&p\text{ is private in that cluster},\\
    0,&p\text{ is absent}.
 \end{cases}
\]
V69 gives pointwise
\begin{equation}
          D_f\le\sum_p w(f,p)X_p,\qquad
                      \sum_p w(f,p)=120 .
 \label{eq:f15v71:weight}
\end{equation}
There are 36 plaquettes in each support.

\begin{proposition}[Inverse incidence with all central orientations]
\label{f15v71:incidence}
For every positive plaquette \(p\),
\begin{equation}
                   \sum_{f\ {\rm positive\ link}}w(f,p)=80 .
 \label{eq:f15v71:eighty}
\end{equation}
In particular, for any finite set \(\mathcal S\)
of distinct positive central links, one has
\[
            W_{\mathcal S}(p):=\sum_{f\in\mathcal S}w(f,p)\le80 .
\]
\end{proposition}

\begin{proof}
Let \(p\) have directions \(\mu,\nu\). A satellite
belonging to \(p\) must have one of those two
orientations; hence only those central orientations
can contribute.

Fix orientation \(\mu\). The two positive
\(\mu\)-edges of \(p\) have anchors \(y\) and
\(y+\widehat\nu\). Either edge is a satellite of
each of the six \(\mu\)-centres obtained by a
one-step transverse shift of its anchor.
These twelve centres are distinct. The two centres
equal to the actual \(\mu\)-edges of \(p\) make it
a shared plaquette, and the remaining ten make it
private. The total in this orientation is
\(2\cdot10+10\cdot2=40\).
The same count for orientation \(\nu\) gives another
40. Other orientations give zero.
The local periodic identifications do not change
these counts when the spatial side is at least four.
\end{proof}

This inverse count is different from V69's disjoint
cluster count. It allows all orientations and
arbitrary overlaps without using a colouring loss
or assuming independent plaquette energies.

\begin{theorem}[An extensive moment for arbitrary native debit sets]
\label{f15v71:joint-debit}
Set \(s_0=1/7680\). Eventually along the selected
vacuum sequence, every finite set \(\mathcal S\)
of distinct positive links satisfies
\begin{equation}
       \log\mathbb E_{\mu_M}
                \exp\left(s_0\sum_{f\in\mathcal S}D_f\right)
                         \le2836s_0|\mathcal S| .
 \label{eq:f15v71:joint-debit}
\end{equation}
The threshold in the sequence is independent of
the cardinality and location of \(\mathcal S\).
\end{theorem}

\begin{proof}
For a plaquette based at \(y\), \(X_p\le4T_y\)
by the inherited symmetric cell assignment.
Define the finite nonnegative cell source
\[
             b_y=4s_0\sum_{p:\,\operatorname{base}(p)=y}
                                             W_{\mathcal S}(p).
\]
The inverse-incidence bound and the six positive
plaquette orientations at a cell give
\[
 \begin{gathered}
     \|b\|_\infty\le4s_0\cdot6\cdot80=\frac14,\\
     \sum_y b_y=480s_0|\mathcal S|,\qquad
     |\operatorname{supp}b|\le36|\mathcal S|.
 \end{gathered}
\]
The sum identity counts multiplicity: overlaps
do not erase their weights.

Apply V62's full finite-source window,
Theorem~\ref{f15v62:window}, to this array.
Eventually \(e_{\rm vac}\le5\), \(k_M\le5\), and
\({\cal E}_M(1/4)\le s_0\), all independently of
the source support. Since
\(J(b)\le b^2/(2(1-b))\le b/6\) for \(0\le b\le1/4\),
\[
 \begin{split}
 \log\mathbb E_{\mu_M}e^{s_0\sum_fD_f}
 &\le e_{\rm vac}\sum_yb_y
                  +k_M\sum_yJ(b_y)
                  +{\cal E}_M(1/4)|\operatorname{supp}b|\\
 &\le\{(5+5/6)480+36\}s_0|\mathcal S|\\
 &=2836s_0|\mathcal S|.
 \end{split}
\]
This proves the assertion without an independence
assumption or a support-size cutoff.
\end{proof}

\subsection{Paid simultaneous bad labels and connected regions}

Let \(B_f=1_{G_f^c}\). It is a gauge-invariant
label of the fine configuration, because \(D_f\)
is gauge invariant. Define the dimensionless rate
\begin{equation}
                  a_\gamma=\frac{\gamma-28360}{76800}.
 \label{eq:f15v71:rate}
\end{equation}
For every finite \(\mathcal S\), Theorem~\ref{f15v71:joint-debit}
and exponential Markov give
\begin{equation}
     \mu_M\{B_f=1\text{ for all }f\in\mathcal S\}
       \le\min\{1,e^{-a_\gamma|\mathcal S|}\}.
 \label{eq:f15v71:joint-bad}
\end{equation}
Indeed on that event
\(\sum_fD_f>(\gamma/10)|\mathcal S|\), and
\(s_0(\gamma/10-2836)=a_\gamma\).
The constants are conservative. Their role is the
extensive dependence on \(|\mathcal S|\) and the
growth \(a_\gamma\to\infty\), not a useful numerical
threshold at moderate coupling.

Fix any countable simple graph \(\mathcal G\) whose
vertices are these positive central links and whose
degrees are bounded by an integer \(\Delta\ge2\).
The adjacency must be chosen independently of the
random configuration. Let \(\mathcal C_f\) be the
connected component of bad labels containing \(f\),
with \(\mathcal C_f=\varnothing\) when \(f\) is good.

\begin{theorem}[Actual-vacuum bad-component tail]
\label{f15v71:components}
For all sufficiently late selected couplings with
\(a_\gamma>2\log\Delta\), and all integers \(n\ge1\),
\begin{equation}
 \mu_M\{|\mathcal C_f|\ge n\}
      \le \Delta^{2(n-1)}e^{-a_\gamma n}
      =e^{-a_\gamma}
           e^{-(a_\gamma-2\log\Delta)(n-1)} .
 \label{eq:f15v71:component-tail}
\end{equation}
Every bad component is then finite almost surely.

More generally, whenever \(u\ge0\) and
\(\Delta^2e^{u-a_\gamma}<1\), the rooted
bad-set probability sum obeys
\begin{equation}
 \sum_{\substack{Y\ni f\\Y\ {\rm finite,\ connected}}}
       e^{u|Y|}
           \mu_M\{B_e=1\text{ for every }e\in Y\}
       \le\frac{e^{u-a_\gamma}}
                        {1-\Delta^2e^{u-a_\gamma}} .
 \label{eq:f15v71:rooted}
\end{equation}
The same right side bounds
\(\mathbb E_{\mu_M}[1_{\{B_f=1\}}e^{u|\mathcal C_f|}]\).
\end{theorem}

\begin{proof}
Use the inherited depth-first spanning-tree count:
there are at most \(\Delta^{2(n-1)}\) connected
vertex sets of size \(n\) containing \(f\).
For precision, fix neighbour orders, choose a
canonical spanning tree for each such set, and
encode it by the walk traversing each tree edge
twice. Its length is \(2(n-1)\), and its visited
vertices recover the set, so different sets cannot
have the same code.

If \(|\mathcal C_f|\ge n\), including the case of
an infinite component, it contains a connected
subset of size exactly \(n\) containing \(f\).
Apply the count and \eqref{eq:f15v71:joint-bad},
then the union bound, to obtain
\eqref{eq:f15v71:component-tail}.
Let \(n\to\infty\). The right side tends to zero;
a countable union over roots rules out every
infinite bad component.

Grouping connected sets by size gives the geometric
series
\[
     \sum_{n\ge1}\Delta^{2(n-1)}e^{(u-a_\gamma)n}
         =\frac{e^{u-a_\gamma}}
                        {1-\Delta^2e^{u-a_\gamma}} .
\]
This proves \eqref{eq:f15v71:rooted}.
For the component moment, the event
\(\{\mathcal C_f=Y\}\) is contained in
\(\{B_e=1\ \forall e\in Y\}\), so summing these
disjoint component events proves the last assertion.
No product property, monotonicity of the labels,
or domination by an independent site process was used.
\end{proof}

The finite connected-set count is standard; see also
Duminil-Copin,
\href{https://www.ihes.fr/~/duminil/publi/2017percolation.pdf}
{Introduction to Bernoulli percolation}, Exercise 7,
for lattice-animal background. The joint probability
input here is instead \eqref{eq:f15v71:joint-bad}
under the actual Wilson vacuum.

A graph covering the local overlap can be chosen
explicitly. Every plaquette in a cluster has base
within \(\ell^1\)-distance two of its centre, and
each of its positive link anchors is within
distance three. Thus two debit carriers which
overlap have centre anchors at periodic
\(\ell^\infty\)-distance at most six.
Join any two distinct positive central links
whose anchors have that distance at most six,
regardless of orientation. Its degree is at most
\[
                      \Delta_6=4\cdot13^4-1=114243 .
\]
This bound is uniform in the spatial period.
The theorem consequently applies, eventually,
even to this conservative overlap graph.
This is a finite-range bad-geometry statement
on the selected vacuum cylinder, not a
continuum length or physical-time estimate.

\subsection{Compatibility and the next missing estimate}

There are three separate objects in the preceding
proofs. The Dirichlet theorem concerns an actual
native conditional after its six satellites have
been integrated. The joint debit theorem concerns
simultaneously defined local functions of one fine
vacuum configuration. The connected-region theorem
concerns probabilities of those labels on a
specified graph. None of the transitions between
these objects changes the vacuum measure.

Nevertheless, simultaneous labels do not authorize
simultaneous deletion of arbitrary overlapping
satellite sets. An application to one common
retained conditional law must still use an
admissible first deletion set, as in V69.
In particular \eqref{eq:f15v71:joint-bad} is not
a tensorization assertion about arbitrary
overlapping native kernels.

Similarly, the summands in
\eqref{eq:f15v71:rooted} are probabilities, not
already constructed polymer activities.
An exact polymer representation, its normalization,
and any interface or electric-operator losses
would have to be proved before a cluster-expansion
criterion could be invoked. What has been paid
here is an arbitrary fixed exponential size weight
in the rooted bad-set probability sum, once the
coupling is sufficiently large.

The new result removes the need to pay
\(e^{9D/4}\) throughout the good region. It does
not remove V70's good-region coherent gauge
response: the perfectly aligned exterior already
has \(D=0\), and its raw-link response tends to one.
Nor does a small probability of bad geometry
by itself control the correlations between
retained conditional means.

The next genuine gluing input is therefore an
estimate for gauge-invariant retained observables
and their conditional means, on a common exact
marginal law, with the costs of bad components
and all boundary interfaces included.
Uniformity through a sequence of physical block
scales and identification of physical transfer
time remain further requirements. Neither a
single native patch gap nor bad-region
subcriticality supplies them automatically.

\subsection{Diagnostics and source preservation}

The diagnostic is
\begin{center}\small
\path{scripts/verify_ym_f15_v71_native_bad_regions.py}.
\end{center}
It reruns V70 and its inherited checks. The new
checks include four exact Lyapunov/cap constant
budgets, 33 three-region inequalities, four
good-debit field fixtures and 384 outward
radial-coefficient/Hessian enclosures.
The field fixtures satisfy the algebraic
hypotheses; they are not asserted to be
independently realizable fine-link exteriors.

The native lattice checks give 36 inverse-incidence
orientation certificates and 18 complete weight-80
identities, plus 72 periodic-regulator and vacuum-cylinder
incidence checks, including periodic seams at spatial
sides four, five and six. There are four overlapping-family
source budgets and four source/exponent constants.
There are 138 injective finite connected-set
walk codes, 21 animal-count bounds, 255
joint-event checks in an explicitly dependent
auxiliary label law and one rooted moment budget.
These finite checks supplement the analytic
proofs; they do not numerically certify a
vacuum law or a mass gap.

The script has 9,144 bytes and SHA--256
\begin{center}\scriptsize\ttfamily
549698421BFEA0352C387315E76140D3B020FDDE938F9DE1D9FAF5F6A5828382.
\end{center}
The V70 predecessor has 1,591,424 bytes and SHA--256
\begin{center}\scriptsize\ttfamily
279A7233DC914F106B7F64E2A50E21D5EC69F9E274336143706E23C46347C2A7.
\end{center}
Removing this final section and restoring the
title version recovers V70 exactly.
V71 replaces it as the sole active main note.
All archived lineages and companion inputs are
preserved. Only \path{.tex} files are kept in
\path{latex_notes}; compilation products go to
\begin{center}\small
\path{artifacts/ym_f15_v71_texcheck/}.
\end{center}

\begin{firewall}
V71 proves uniform thermal conditional coercivity
on a quantified native good-boundary set, and
an extensive actual-vacuum bound on arbitrary
bad-label sets, implying exponentially controlled
bad components on fixed finite-range graphs.
It does not prove contraction of the retained
gauge-invariant dynamics, a cofinal physical
RG iteration, an interacting continuum
construction, or a positive physical
Yang--Mills mass gap. The original goal
remains active.
\end{firewall}

\section{V72: a common retained law and an invariant slow mode}
\label{f15v72:section}

V71 supplies local conditional coercivity and a paid
bad-region geometry. The next question is whether the
retained conditional means can be controlled on one
common exact marginal law. An all-exterior estimate
would be especially convenient, but it must first
survive a gauge-invariant test.

This version constructs that common law for a full
layer of parallel links. On a flat transverse exterior
its exact first marginal is a decimated \(O(4)\) spin
law. We compute its complete quadratic operator and
prove a low-temperature Rayleigh bound using an
observable which is a sum of genuine dressed Wilson
loops. Unlike the raw-link test in V70, this observable
is invariant under both residual endpoint gauge groups.
The obstruction below is therefore not just a
constant gauge rotation.

The result is an obstruction to a proposed
volume-uniform, all-exterior \emph{conditional sampling}
gap. It is not a gapless physical Hamiltonian
construction. In particular, the chosen exterior
is not asserted to be typical under the selected
vacuum, and the order of limits will remain explicit.

The V71 source and all 128 protected files were
verified unchanged. The next companion checkpoint
remains V75. The generic Schur-complement and Gaussian
methods are not new; their role here is to test the
actual common compact marginal and a gauge-invariant
observable, rather than independent single-patch
kernels or a hypothetical Gaussian replacement.

\subsection{One admissible deletion set for a whole layer}

Take an even spatial side \(L\ge4\), and consider all
positive time links from time zero to time one,
indexed by
\[
                   Q_L=(\mathbb Z/L\mathbb Z)^3 .
\]
Write their unit quaternions as \(u_x\), \(x\in Q_L\).
Work first in a finite periodic Wilson regulator
with a nondegenerate time period. Set every exterior
link to the identity. Only the spatial links on
the two endpoint slices enter the conditional
density of this time-link layer.

Each plaquette meeting the layer contains two of
its time links and two fixed spatial links. Its
normalized trace is \(u_x\cdot u_y\), where \(x,y\)
are nearest neighbours in \(Q_L\). Thus the exact
layer conditional is
\begin{equation}
 d\pi_{L,\gamma}(u)
   =\frac1{{\cal Z}_{L,\gamma}}
               e^{-\gamma H_L(u)}\prod_{x\in Q_L}dH(u_x),
 \qquad
 H_L(u)=\sum_{\{x,y\}\in{\cal E}_L}(1-u_x\cdot u_y).
 \label{eq:f15v72:spin}
\end{equation}
The edge set \({\cal E}_L\) is the undirected
nearest-neighbour edge set, counted once, so
\(|{\cal E}_L|=3L^3\). Exterior action terms cancel
only through conditional normalization.
Equation \eqref{eq:f15v72:spin} is a canonical
continuous conditional kernel, including at this
specified exterior of Haar probability zero.
It is not a statement that a random exterior is flat.

Split the spatial sites by parity:
\[
 E=\{x:x_1+x_2+x_3\text{ is even}\},\qquad O=Q_L\setminus E,
 \qquad |E|=|O|=:n=L^3/2 .
\]
The time links indexed by \(O\) form one admissible
independent deletion set: no plaquette contains
two of them. Put \(v_x=u_x\) for \(x\in E\).
Since
\[
 H_L(u)=3L^3-\sum_{q\in O}
                    u_q\cdot\left(\sum_{x\sim q}v_x\right),
\]
integrating all odd links gives the common retained law
\begin{equation}
 d\rho_{L,\gamma}(v)
   =\frac1{{\cal N}_{L,\gamma}}
       \prod_{q\in O}
           Z\left(\gamma\left|\sum_{x\sim q}v_x\right|\right)
                         \prod_{x\in E}dH(v_x).
 \label{eq:f15v72:marginal}
\end{equation}
Here \(Z(h)=\int_{S^3}e^{hu_0}dH(u)\), as before.
All factors and their common normalizer are retained.
The link/spin connection is classical background;
see Kogut,
\href{https://journals.aps.org/rmp/abstract/10.1103/RevModPhys.51.659}
{An introduction to lattice gauge theory and spin systems}.
The precise conditional and deletion geometry needed
here have been derived directly above.

For a retained \(x\), the six factors indexed by
\(q\sim x\) are exactly V69's native factors. In this
flat exterior their fields are
\begin{equation}
        c_{q,x}=\sum_{\substack{y\sim q\\y\ne x}}v_y,
          \qquad
        D_x=\gamma\left(30-
                     \left|\sum_{q\sim x}c_{q,x}\right|\right).
 \label{eq:f15v72:fields}
\end{equation}
There are five terms in each \(c_{q,x}\), with
multiplicity retained when the sums are combined.
Consequently every one-link conditional of
\(\rho_{L,\gamma}\) is a native V69 conditional.
They are not independent conditionals: one odd
factor generally involves six retained variables.

\subsection{The complete quadratic operator after deletion}

Let \(B:\mathbb R^O\to\mathbb R^E\) be the bipartite
adjacency matrix, \(B_{xq}=1_{\{x\sim q\}}\). Define
\begin{equation}
                    K_L=6I_E-\frac16BB^{\mathsf T}.
 \label{eq:f15v72:K}
\end{equation}
The full nearest-neighbour Laplacian is
\[
             {\cal L}_L=
                  \begin{pmatrix}6I_E&-B\\-B^{\mathsf T}&6I_O\end{pmatrix}.
\]
Its Schur complement on \(E\) is \(K_L\). In
particular \(K_L\) is nonnegative, its kernel is
the constants, and its diagonal entries equal five.
The kernel assertion follows also by minimizing
the connected-graph quadratic energy over the odd
variables: zero energy forces every variable to
be the same constant.

\begin{proposition}[Exact finite-coupling Hessian]
\label{f15v72:hessian}
At any aligned configuration \(v_x=h\) for all
\(x\in E\), the Hessian of the negative logarithm
of \eqref{eq:f15v72:marginal}, in orthonormal
left-translated tangent coordinates, is
\begin{equation}
              \gamma a(6\gamma)\,K_L\otimes I_3 .
 \label{eq:f15v72:hessian}
\end{equation}
This is the Hessian of the full retained density,
not a selection of generated terms.
\end{proposition}

\begin{proof}
By common left translation take \(h=1\).
For tangent vectors \(A_x\in\mathbb R^3\), use
\(v_x(\varepsilon)=\exp(\varepsilon A_x)\) in the
unit-quaternion convention. For an odd star,
\[
 \left|\sum_{x\sim q}v_x(\varepsilon)\right|
 =6-\frac{\varepsilon^2}{2}
       \left\{\sum_{x\sim q}|A_x|^2
                  -\frac16\left|\sum_{x\sim q}A_x\right|^2\right\}
                                  +O(\varepsilon^3).
\]
Differentiate its exact logarithmic factor.
Its contribution to the quadratic Hessian is
\(\gamma a(6\gamma)\) times the displayed bracket.
Summing the sixfold incidence of each retained
variable gives \(K_L\). The normalizing constant
does not depend on \(v\).
\end{proof}

Thus the local diagonal stiffness
\(5\gamma a(6\gamma)\), already visible in V69,
coexists with nontrivial off-diagonal couplings.

Put \(\theta_L=2\pi/L\), and take the real function
\begin{equation}
       \varphi_x=\frac{2}{L^{3/2}}\cos(\theta_Lx_1),
                   \qquad x\in E .
 \label{eq:f15v72:mode}
\end{equation}
Each value of \(x_1\) occurs \(L^2/2\) times in \(E\).
The finite trigonometric sums therefore give
\(\sum_E\varphi_x=0\) and \(\sum_E\varphi_x^2=1\).
The adjacency symbol on this Fourier mode is
\(4+2\cos\theta_L\), so
\begin{equation}
 \begin{gathered}
     K_L\varphi=\kappa_L\varphi,\\
     \kappa_L
       =6-\frac{(4+2\cos\theta_L)^2}{6}
       =\frac23(1-\cos\theta_L)(5+\cos\theta_L)>0,\\
        2\kappa_L\le\frac{16\pi^2}{L^2},
            \qquad 2\kappa_L\sim\frac{16\pi^2}{L^2}
                         \quad(L\to\infty,\ L\ {\rm even}).
 \end{gathered}
 \label{eq:f15v72:dispersion}
\end{equation}
Indeed the adjacency multiplies the corresponding
character on either parity class by that symbol.
Also \(5+\cos\theta_L\le6\) and
\(1-\cos\theta_L\le\theta_L^2/2\).
This computation alone is not a spectral conclusion
for the compact law; the next steps provide a
normalized compact-law test.

\subsection{A genuine gauge-invariant retained observable}

Choose an even root \(r\) and, for every \(x\in E\),
spatial paths from \(r\) to \(x\) on each endpoint
slice. Let their ordered parallel transports be
\(P_x^-\) and \(P_x^+\), and define
\[
             {\cal V}_x=P_x^-\,U_{(0,x),0}\,(P_x^+)^{-1}.
\]
These variables use only retained time links and
retained spatial links. Under arbitrary local
gauge transformations they all transform as
\[
              {\cal V}_x\longmapsto
                   g_{(0,r)}{\cal V}_xg_{(1,r)}^{-1}.
\]
Consequently
\begin{equation}
        {\cal F}_\gamma(U)
             =\gamma\left|\sum_{x\in E}\varphi_x{\cal V}_x\right|^2
 \label{eq:f15v72:dressed}
\end{equation}
is gauge invariant. Its norm is the quaternion norm,
equivalently \(\tfrac12\Tr(AA^\dagger)\) for the
matrix \(A=\sum_x\varphi_x{\cal V}_x\).
Expanding the square gives a finite real linear
combination of closed Wilson-loop traces
\(\tfrac12\Tr({\cal V}_x{\cal V}_y^{-1})\).

At the fixed flat exterior, all path transports
are the identity and its restriction is
\begin{equation}
             F_\gamma(v)=\gamma|S(v)|^2,\qquad
             S(v)=\sum_{x\in E}\varphi_xv_x .
 \label{eq:f15v72:test}
\end{equation}
It is invariant under the simultaneous action
\(v_x\mapsto g_-v_xg_+^{-1}\) of both endpoint groups.
Thus it is not a test of a residual constant
gauge orientation. A nonconstant relative link
profile, not a common rotation of an aligned
profile, is required to change it.

\subsection{Fixed-volume compact concentration and the Gaussian probe}

The following argument takes \(\gamma\to\infty\)
at a fixed even \(L\). No estimate uniform in
\(L\) is used.

\begin{proposition}[The required fixed-volume moment limit]
\label{f15v72:laplace}
For fixed \(L\),
\begin{equation}
       \mathbb E_{\rho_{L,\gamma}}F_\gamma
                        \longrightarrow\frac3{\kappa_L},
       \qquad
       \Var_{\rho_{L,\gamma}}(F_\gamma)
                        \longrightarrow\frac6{\kappa_L^2}.
 \label{eq:f15v72:moments}
\end{equation}
\end{proposition}

\begin{proof}
We may compute these expectations before deleting
the odd links, because \(F_\gamma\) depends only
on the retained variables. In the full spin law
\eqref{eq:f15v72:spin}, write \(u_r=h\) and
\(u_x=h w_x\) for \(x\ne r\), with \(w_r=1\).
Product Haar measure becomes exactly
\(dH(h)\prod_{x\ne r}dH(w_x)\), by invariance in
each factor. Both the energy and the observable
are independent of \(h\). This removes one common
left orientation with no additional determinant
or chart assumption.

The pinned energy has a unique zero at \(w_x=1\)
for every \(x\): each nonnegative edge term must
vanish, and the graph is connected. Near this
point use \(w_x=\exp(A_x)\), \(A_r=0\). Then
\[
        H_L(w)=\frac12 A^{\mathsf T}
                     ({\cal L}_{L,r}\otimes I_3)A+O_L(|A|^3),
\]
where \({\cal L}_{L,r}\) is the positive definite
root-deleted graph Laplacian.
The Haar density in these coordinates is smooth
and positive at zero. For a sufficiently small
fixed chart there are constants \(c_L,C_L>0\)
with
\[
                  c_L|A|^2\le H_L(\exp A)\le C_L|A|^2 .
\]
On the compact complement of that chart the
energy has a positive lower bound. With
\(d=3(L^3-1)\), a ball of radius \(\gamma^{-1/2}\)
therefore bounds the pinned partition function
below by \(c'_L\gamma^{-d/2}\), while the complement
has probability at most
\(C'_L\gamma^{d/2}e^{-\eta_L\gamma}\).

Set \(X=\sqrt\gamma A\) in the chart.
Its unnormalized density, multiplied by
\(\gamma^{d/2}\), tends to the density proportional to
\[
           \exp\{-\tfrac12X^{\mathsf T}
                                ({\cal L}_{L,r}\otimes I_3)X\}.
\]
The quadratic lower bound dominates it by
\(C_Le^{-c_L|X|^2}\); the rescaled normalizer
has a positive limit. This proves convergence
of every fixed polynomial moment on the chart.

For a zero-sum linear functional, pinning and
removing the constant mode give the same
Gaussian covariance. One can verify the
particular covariance needed here directly.
Extend the charge to all sites by
\(\ell_x=\varphi_x\) on \(E\), \(\ell_q=0\) on \(O\).
The function
\[
       h_E=\varphi/\kappa_L,\qquad
       h_O=B^{\mathsf T}\varphi/(6\kappa_L)
\]
solves \({\cal L}_Lh=\ell\). Subtracting \(h_r\)
pins it without changing \(\ell\cdot h\), since
\(\sum\ell=0\). Hence
\[
       \Var\left(\sum_{x\in E}\varphi_xX_x^{(\alpha)}\right)
                   =\ell\cdot h=\frac1{\kappa_L}
\]
for each colour \(\alpha=1,2,3\), and the colours
are independent.

Write \(Y=\sum_E\varphi_xX_x\). Since \(\sum_E\varphi_x=0\),
the quaternion expansion gives
\[
                    F_\gamma\longrightarrow |Y|^2
\]
pointwise at every fixed \(X\). On the chart,
\[
 F_\gamma
 =\gamma\left|\sum_E\varphi_x(w_x-1)\right|^2
       \le\sum_E|X_x|^2 .
\]
This uses \(|\exp A-1|\le|A|\) and
\(\sum\varphi_x^2=1\). It gives integrable polynomial
majorants for the first and second moments.
Outside the chart, \(F_\gamma\le\gamma n\), so
the exponentially small complement contributes
zero to both limits. Thus the first two moments
converge to those of \(|Y|^2\), for
\(Y\sim N(0,\kappa_L^{-1}I_3)\).
They are \(3/\kappa_L\) and variance \(6/\kappa_L^2\),
respectively.
\end{proof}

All constants used for chart coverage may depend
on \(L\). The proof justifies this finite-volume
test under the complete compact conditional law;
it does not assert the simultaneous continuum
or thermodynamic Gaussian limit.

\subsection{The invariant thermal Rayleigh quotient}

Use the retained round-product Dirichlet form
\begin{equation}
 {\cal E}_{L,\gamma}(F)
        =\mathbb E_{\rho_{L,\gamma}}
                    \sum_{x\in E}|\nabla_xF|^2 .
 \label{eq:f15v72:form}
\end{equation}
For \(F_\gamma\), differentiation is exact:
\[
 \nabla_xF_\gamma=2\gamma\varphi_x\Pi_{v_x}S,\qquad
 \Pi_{v_x}=I-v_xv_x^{\mathsf T}.
\]
Therefore
\begin{equation}
 \gamma^{-1}\sum_{x\in E}|\nabla_xF_\gamma|^2
       =4F_\gamma
                 -4\gamma\sum_{x\in E}\varphi_x^2(v_x\cdot S)^2 .
 \label{eq:f15v72:gradient}
\end{equation}
The subtracted term is nonnegative.

\begin{theorem}[Invariant slow mode in one exact common marginal]
\label{f15v72:rayleigh}
For every fixed even \(L\ge4\),
\begin{equation}
        \lim_{\gamma\to\infty}
          \frac{\gamma^{-1}{\cal E}_{L,\gamma}(F_\gamma)}
                    {\Var_{\rho_{L,\gamma}}(F_\gamma)}
                                  =2\kappa_L .
 \label{eq:f15v72:rayleigh}
\end{equation}
Let \(\lambda^{\rm inv}_{L,\gamma}\) be the infimum
of this thermal Rayleigh quotient over nonconstant
smooth functions invariant under
\(v_x\mapsto g_-v_xg_+^{-1}\). Then
\begin{equation}
            \limsup_{\gamma\to\infty}
                      \lambda^{\rm inv}_{L,\gamma}
                        \le2\kappa_L\le16\pi^2/L^2 .
 \label{eq:f15v72:invariant-gap}
\end{equation}
\end{theorem}

\begin{proof}
The observable belongs to the stated invariant
class and has positive limiting variance.
In the pinned chart, at fixed \(X\),
\(S=\gamma^{-1/2}(0,Y)+O_X(\gamma^{-1})\), while
\(v_x=(1,0)+O_X(\gamma^{-1/2})\).
Thus
\(\gamma(v_x\cdot S)^2\to0\) for every \(x\).
The entire subtracted sum in
\eqref{eq:f15v72:gradient} is bounded by \(4F_\gamma\).
The same Gaussian domination and polynomial
complement estimates as in
Proposition~\ref{f15v72:laplace} justify convergence
of its expectation to zero.
Consequently
\[
             \gamma^{-1}{\cal E}_{L,\gamma}(F_\gamma)
                             \longrightarrow12/\kappa_L .
\]
Divide by the variance limit \(6/\kappa_L^2\).
The variational upper bound and
\eqref{eq:f15v72:dispersion} finish the proof.
\end{proof}

The infimum concerns a conditional sampling
Dirichlet form with the transverse links fixed.
In particular, the full Dirichlet form of the
dressed observable would also differentiate its
spatial path transports. Those derivatives have
not been included in \eqref{eq:f15v72:form}.

\subsection{The obstruction remains near simultaneous good boundaries}

Let
\[
              {\cal G}_L=\bigcap_{x\in E}\{D_x\le\gamma/10\}
\]
with the actual fields \eqref{eq:f15v72:fields}.
Its definition depends only on retained variables.
The ratios \(D_x/\gamma\) are continuous and
independent of \(\gamma\) as functions of \(v\);
they vanish at every globally aligned configuration.

\begin{proposition}[Good-region concentration at fixed volume]
\label{f15v72:all-good}
For fixed \(L\), there are constants
\(C_L,\eta_L>0\), independent of \(\gamma\ge1\), such that
\begin{equation}
        \rho_{L,\gamma}({\cal G}_L^c)
          \le C_L\gamma^{3(L^3-1)/2}e^{-\eta_L\gamma}.
 \label{eq:f15v72:all-good}
\end{equation}
The Rayleigh limit \eqref{eq:f15v72:rayleigh} is
unchanged if its expectations and variance are
taken under the normalized restriction
\(\rho_{L,\gamma}(\,\cdot\,|{\cal G}_L)\).
\end{proposition}

\begin{proof}
Use the full pinned spin law from the preceding
Laplace proof. The closed set
\[
          \left\{\max_{x\in E}D_x/\gamma\ge1/10\right\}
\]
does not contain its unique zero-energy point.
Compactness gives a strictly positive minimum
of \(H_L\) on this set. The pinned partition
lower bound proves \eqref{eq:f15v72:all-good}
for its subset \({\cal G}_L^c\).

The deterministic bounds \(F_\gamma\le\gamma n\),
\(F_\gamma^2\le\gamma^2n^2\), and
\(\gamma^{-1}\sum|\nabla F_\gamma|^2\le4\gamma n\)
show that the omitted contributions to all three
Rayleigh moments tend to zero. The probability
of the retained set tends to one, so normalization
does not change their limits.
\end{proof}

Restricting the whole measure to \({\cal G}_L\)
changes its one-link conditional laws. No
identification of those new conditionals with
the V71 kernels is being made. The statement is
that this invariant slow test is carried by the
simultaneously good region, not by a large
probability of bad labels.

For example, for each even \(L\) one can choose
\(\gamma_L\ge L\) sufficiently large that the
unrestricted and restricted Rayleigh quotients
are both at most \(3\kappa_L\), and
\(\rho_{L,\gamma_L}({\cal G}_L^c)\le1/L\).
Along this diagonal sequence both quotients
tend to zero. The existence of this sequence
follows from fixed-\(L\) limits; it does not
identify \(\gamma_L\) with the programme's
selected physical coupling sequence.

\subsection{What has been ruled out, and the necessary direction}

There is no positive constant uniform in large
coupling, spatial side and exterior which lower
bounds the thermal conditional invariant-sector
gap of all these retained layer laws. Otherwise
Equation~\eqref{eq:f15v72:invariant-gap}, with \(L\)
arbitrarily large, would contradict that constant.
The same proposed assertion for the corresponding
globally restricted good-region laws also fails.

There is also a finite-regulator essential-bound
version of the obstruction. Given \(c>0\), choose
\(L\) and then a finite sufficiently large
\(\gamma\) so that the displayed dressed-loop
Rayleigh quotient at the flat exterior is below
\(c/2\), with positive variance. At this fixed
regulator the exact normalized marginal kernel,
the dressed observable and its retained gradients
are continuous in the exterior links. Compact
integration therefore makes the quotient continuous
near that exterior. It stays below \(c\) on an
open neighbourhood. The finite Wilson exterior
marginal has a strictly positive density relative
to product Haar, so this neighbourhood has positive
probability. Thus a bound uniform over all large
volumes and couplings cannot be rescued merely by
excluding a null set of finite-regulator exteriors.
This argument supplies no useful lower bound on
the probability of the neighbourhood.

This is stronger than pointing to the constant
gauge response in V70: the test is gauge invariant,
and its leading fluctuation is the squared magnitude
of a nonconstant spatial mode. It is also a test
of one complete marginal, so incompatible
single-patch deletions cannot explain it away.

The result does not negate V71's actual-vacuum
bad-set estimates. A flat transverse exterior
is a special conditional input, not the law
obtained by averaging those fields. No
all-exterior obstruction here implies failure
of an averaged estimate with the necessary
boundary derivatives and physical scale.
Nor is an auxiliary diffusive eigenvalue a
statement about the Yang--Mills transfer
Hamiltonian or about physical masslessness.

The upstream correction is consequently precise:
a full proof cannot be obtained by combining the
native one-link gap and rare bad geometry with
a presumed uniform gap of every retained
flat-boundary layer. The next estimate must
incorporate the transverse-field law and the
derivatives of the dressed observables, or
supply a genuinely different physical block
mechanism. Its treatment of long-wavelength
gauge-invariant modes and physical transfer
time must be explicit. This version does not
claim that missing estimate.

\subsection{Diagnostics and exact source recovery}

The diagnostic is
\begin{center}\small
\path{scripts/verify_ym_f15_v72_common_layer_slow_mode.py}.
\end{center}
It reruns V71 and its inherited checks. On
spatial sides four and six it verifies two
common deletion-layer geometries, 140 Schur
diagonal/row sums, 140 real Fourier eigenvector
equations, 280 pinned Laplace equations and
two exact zero-sum probe variances.
These diagnostics use rational, unnormalized
cosine vectors; normalization cancels from
the checked Rayleigh quotient.

There are 2,520 literal plaquette/spin identities,
six common-layer energy factorizations, six
independent Schur-completion identities, 420
endpoint-frame dressing identities and 420
intrinsic projected-gradient identities.
Six invariant norm tests, 140 nonconstant
simultaneous good-boundary checks, two
chi-square Rayleigh constants and 33 dispersion
factorizations are also exact rational checks.
The compact low-temperature limit and its
normalization are analytic results above,
not inferred from sampled configurations.

The script has 7,809 bytes and SHA--256
\begin{center}\scriptsize\ttfamily
6DEB51434AD758C45B9A974C0275FA9ECBB2531EB0709AA4D947089247EF0C1A.
\end{center}
The V71 predecessor has 1,613,992 bytes and SHA--256
\begin{center}\scriptsize\ttfamily
2D8856C38395A9F0B509E8538233B5F94116B6BBFFC02CA22AF16A8A056B0B17.
\end{center}
Removing this final section and restoring the
title version recovers V71 exactly.
V72 replaces it as the sole active main note;
archives and companion inputs remain untouched.
Only \path{.tex} files are kept in
\path{latex_notes}. All compilation products go to
\begin{center}\small
\path{artifacts/ym_f15_v72_texcheck/}.
\end{center}

\begin{firewall}
V72 proves an exact common retained-layer law
and a gauge-invariant, fixed-volume
low-temperature Rayleigh limit of order
\(L^{-2}\). It rules out a volume-uniform
all-exterior thermal conditional gap for
this class of kernels, including the displayed
good-region restriction. It does not prove
physical masslessness, an interacting continuum,
or the positive physical Yang--Mills mass gap.
The full original goal remains active.
\end{firewall}

\section{V73: transverse Ward cost and optimal flat dressing}
\label{f15v73:section}

V72 tests the retained-link conditional form while holding
the transverse spatial links fixed. Its gauge-invariant
extension also depends on those spatial links. This version
computes the missing derivatives, rather than interpreting
the conditional quotient as the quotient of the full
observable. The calculation gives a positive, exact
completion for this test, but not a physical mass-gap theorem.

There are two new interfaces. First, at every transverse
configuration, differentiated gauge invariance forces a
covariant inverse-Laplacian cost for the spatial derivatives.
Second, at flat transverse links that lower bound is attained
by a globally smooth average of tree-dressed observables.
Applied to V72, the minimal completed conditional quotient
tends to \(4+2\kappa_L\), not \(2\kappa_L\).
The extra four is a kinematic dressing cost in the specified
round-link gradient metric. It is not an identified
Yang--Mills mass.

The V72 predecessor and all 129 protected inputs and
diagnostics were checked unchanged. The next companion
checkpoint remains V75. The earlier Thomson/capacity
discussion in archive f1 V13 concerns a different
physical weighted form. Neither the least-norm principle
nor Kirchhoff's tree-current identity is claimed as new
general mathematics; here they give an exact completion
for the native retained probe.

\subsection{Declared support, frames, and gradients}

Let \(G=(V,{\cal E})\) be a finite connected simple graph,
with one chosen orientation for each undirected edge,
unit edge conductances, and root \(r\). Let \(W\subseteq V\)
index retained vertical links. There are two copies of the
spatial graph, with link variables \(A^-_e,A^+_e\in\SU(2)\),
and vertical variables \(v_x\in\SU(2)\), \(x\in W\).
The observable below depends on precisely these allowed
variables, or a subset of them, and on no additional layers.

Consider a smooth real gauge-invariant cylinder
\({\cal F}(A^-,A^+,v)\). The independent endpoint gauge
transformations act by
\begin{equation}
 \begin{split}
 A^\pm_{xy}&\longmapsto g^\pm_x A^\pm_{xy}(g^\pm_y)^{-1},\\
 v_x&\longmapsto g^-_x v_x(g^+_x)^{-1}.
 \end{split}
 \label{eq:f15v73:gauge}
\end{equation}
Take the three imaginary quaternion units \(E_\alpha\)
as an orthonormal basis for the round unit \(S^3\).
The left and right derivatives use, respectively,
\(u\mapsto e^{sE_\alpha}u\) and
\(u\mapsto ue^{sE_\alpha}\). Define colour vectors
\begin{equation}
 \begin{gathered}
 q_x^{-,\alpha}=L_\alpha^{v_x}{\cal F},\qquad
 q_x^{+,\alpha}=R_\alpha^{v_x}{\cal F},\\
 a_e^{\pm,\alpha}=L_\alpha^{A^\pm_e}{\cal F},\qquad
 q_x^\pm=0\quad(x\notin W).
 \end{gathered}
 \label{eq:f15v73:charges}
\end{equation}
All norms and adjoints in vertex and edge colour spaces
use unweighted Euclidean sums. In particular the complete
allowed-link gradient integrand is
\begin{equation}
 {\mathfrak e}({\cal F})
       =\sum_{x\in W}|q_x^-|^2+\|a^-\|^2+\|a^+\|^2 .
 \label{eq:f15v73:integrand}
\end{equation}
The vertical derivative is counted once:
left and right orthonormal frames give the same
\(|\nabla_{v_x}{\cal F}|^2=|q_x^-|^2=|q_x^+|^2\).

For a spatial connection \(A\), define
\begin{equation}
 (D_A\xi)_{e=(x,y)}
       =\xi_x-\operatorname{Ad}_{A_e}\xi_y,\qquad
       {\mathsf L}_A=D_A^*D_A .
 \label{eq:f15v73:covariant}
\end{equation}
Here \(\operatorname{Ad}_u\xi=u\xi u^{-1}\) is an
orthogonal map of the imaginary quaternions.
The Moore--Penrose inverse
\({\mathsf L}_A^\dagger\) is always defined.
No invertibility, spectral lower bound, or fixed kernel
dimension is imposed.

\subsection{An exact pointwise covariant Ward cost}

\begin{proposition}[Gauge invariance forces spatial energy]
\label{f15v73:ward}
At every configuration of the declared variables,
\begin{equation}
 D_{A^-}^*a^-=-q^-,\qquad D_{A^+}^*a^+=q^+ .
 \label{eq:f15v73:ward}
\end{equation}
In particular \(q^\pm\in\operatorname{ran}{\mathsf L}_{A^\pm}\).
Writing
\[
 a^-_{\min}=-D_{A^-}{\mathsf L}_{A^-}^\dagger q^-,
 \qquad
 a^+_{\min}= D_{A^+}{\mathsf L}_{A^+}^\dagger q^+,
\]
one has the exact decompositions
\begin{equation}
 \|a^\pm\|^2
   =\langle q^\pm,{\mathsf L}_{A^\pm}^\dagger q^\pm\rangle
                         +\|a^\pm-a^\pm_{\min}\|^2 .
 \label{eq:f15v73:pythagoras}
\end{equation}
\end{proposition}

\begin{proof}
At the lower slice differentiate the simultaneous gauge
transformation \(g^-_x=e^{s\xi_x}\) at every vertex.
For \(e=(x,y)\), its infinitesimal spatial variation is
\(\xi_x A_e-A_e\xi_y=(D_A\xi)_eA_e\).
The vertical contribution is \(\sum_x q_x^-\cdot\xi_x\).
Since the derivative of a gauge-invariant function is zero,
\(\langle a^-,D_{A^-}\xi\rangle+\langle q^-,\xi\rangle=0\)
for every \(\xi\). This proves the first identity.
On the upper slice \(v_x\mapsto v_xe^{-s\xi_x}\);
the vertical contribution has the opposite sign, proving
the second identity.

For any finite matrix \(D\),
\(\operatorname{ran}D^*=\operatorname{ran}(D^*D)\).
Thus both charges are automatically orthogonal to the
appropriate covariantly constant kernel. The displayed
\(a^\pm_{\min}\) solve their Ward equations. Their
differences from \(a^\pm\) lie in \(\ker D_{A^\pm}^*\),
whereas \(a^\pm_{\min}\) lie in \(\operatorname{ran}D_{A^\pm}\).
These spaces are orthogonal. Finally,
\[
 \|D_A{\mathsf L}_A^\dagger q\|^2
 =\langle {\mathsf L}_A^\dagger q,
                    {\mathsf L}_A{\mathsf L}_A^\dagger q\rangle
 =\langle q,{\mathsf L}_A^\dagger q\rangle
\]
for \(q\in\operatorname{ran}{\mathsf L}_A\).
This proves the decomposition even at singular connections.
\end{proof}

Consequently any probability law \(\mu\) on these variables
satisfies
\begin{equation}
 \mathbb E_\mu {\mathfrak e}({\cal F})
 \ge \mathbb E_\mu\left[
       \sum_{x\in W}|q_x^-|^2+
       \langle q^-,{\mathsf L}_{A^-}^\dagger q^-\rangle+
       \langle q^+,{\mathsf L}_{A^+}^\dagger q^+\rangle
                         \right].
 \label{eq:f15v73:averaged}
\end{equation}
In particular one can use the actual finite Wilson law,
or its exact marginal on the declared support, without
changing measures on either side. The pseudoinverse is
a Borel function of the finite positive matrix, and its
charge quadratic form is bounded above by the corresponding
spatial gradient norm in \eqref{eq:f15v73:pythagoras}.
For a smooth cylinder on the compact finite configuration
space this proves integrability. It does not bound
\(\mathbb E_\mu\|{\mathsf L}_A^\dagger\|\), nor does it
permit replacing the random transverse connections by flat ones.

\subsection{Flat restriction and the least possible completion}

Let \(\partial\) denote the head-minus-tail incidence
matrix, with columns \(\delta_y-\delta_x\) for \(e=(x,y)\).
At \(A_e=1\), \(D_A=-\partial^{\mathsf T}\otimes I_3\).
Put
\[
 {\cal L}=\partial\partial^{\mathsf T},\qquad
                         {\cal G}={\cal L}^\dagger .
\]
For a flat restriction
\(f(v)={\cal F}(1,1,v)\), the common left and right
actions on all \(v_x\) leave \(f\) invariant. Conversely
we now take any smooth real \(f\) with these two
invariances. Its vertical derivatives, extended by zero
outside \(W\), obey
\begin{equation}
                        \sum_xq_x^-=\sum_xq_x^+=0 .
 \label{eq:f15v73:neutrality}
\end{equation}
Define the quadratic completion at each retained configuration:
\begin{equation}
 {\mathfrak m}_f(v)
 =\sum_{x\in W}|q_x^-|^2
    +\sum_{\alpha=1}^3\left[
      (q^{-,\alpha})^{\mathsf T}{\cal G}q^{-,\alpha}
      +(q^{+,\alpha})^{\mathsf T}{\cal G}q^{+,\alpha}\right].
 \label{eq:f15v73:minimal}
\end{equation}
Proposition~\ref{f15v73:ward} proves
\({\mathfrak e}({\cal F})(1,1,v)\ge{\mathfrak m}_f(v)\)
for every invariant extension with the declared support.
The next construction attains equality simultaneously
for every \(v\), not only to quadratic order near alignment.

\begin{proposition}[The uniform-tree signed-current identity]
\label{f15v73:treecurrent}
Let \({\cal T}\) be the set of spanning trees of \(G\),
with \(\tau=|{\cal T}|\). For a tree \(T\), let
\(I_T^x(e)\in\{-1,0,1\}\) be the signed indicator that the
unique path from \(r\) to \(x\) traverses \(e\).
Then
\begin{equation}
 \overline I^x:=\frac1\tau\sum_{T\in{\cal T}}I_T^x
          =\partial^{\mathsf T}{\cal G}(\delta_x-\delta_r).
 \label{eq:f15v73:treecurrent}
\end{equation}
\end{proposition}

\begin{proof}
The identity is zero on both sides when \(x=r\).
For \(x\ne r\), let \({\cal B}_{r,x}\) be the spanning
forests with exactly two tree components, separating
\(r\) and \(x\). Define a scalar potential on vertices by
\[
 h_x(z)=\frac1\tau\sum_{F\in{\cal B}_{r,x}}
          1_{\{z\text{ belongs to the component containing }x\}}.
\]
For an oriented edge \(e=(a,b)\), only forests separating
\(a\) and \(b\) contribute to \(h_x(b)-h_x(a)\).
Adding \(e\) to such a forest gives a spanning tree
whose path from \(r\) to \(x\) uses \(e\).
Conversely, removing \(e\) from any such tree gives
exactly one of these forests. The contribution is
positive precisely when the path traverses \(a\) to \(b\).
Thus \(h_x(b)-h_x(a)=\overline I^x(e)\).

Every path current has divergence
\(\partial I_T^x=\delta_x-\delta_r\).
It follows that
\({\cal L}h_x=\delta_x-\delta_r\).
Since the graph is connected, two solutions differ
by a constant; their edge differences agree.
This proves \eqref{eq:f15v73:treecurrent}.
\end{proof}

This is the elementary first-moment form of the classical
Kirchhoff tree-current relation. For background on electrical
networks and random spanning trees see Section 4 of
Benjamini, Lyons, Peres and Schramm,
\href{https://rdlyons.pages.iu.edu/pdf/usf.pdf}
{Uniform Spanning Forests}, Annals of Probability
29 (2001), 1--65. The finite identity used here was
proved directly by the two-forest bijection.

For each tree \(T\), let \(P_x^T(A)\) be the ordered
parallel transport along its root-to-\(x\) path,
using \(A_e^{-1}\) for reverse traversal.
For independently summed trees \(T^-,T^+\), put
\[
 {\cal V}_x^{T^-,T^+}
       =P_x^{T^-}(A^-)\,v_x\,P_x^{T^+}(A^+)^{-1},
 \qquad
 \overline{\cal F}(A^-,A^+,v)
       =\frac1{\tau^2}\sum_{T^-,T^+\in{\cal T}}
                   f\bigl(({\cal V}_x^{T^-,T^+})_{x\in W}\bigr).
\]

\begin{theorem}[A globally smooth extension attains the flat minimum]
\label{f15v73:optimal}
The function \(\overline{\cal F}\) is smooth and globally
gauge invariant, and \(\overline{\cal F}(1,1,v)=f(v)\).
Moreover,
\begin{equation}
             {\mathfrak e}(\overline{\cal F})(1,1,v)
                                      ={\mathfrak m}_f(v)
                  \quad\text{for every }v .
 \label{eq:f15v73:attainment}
\end{equation}
Thus \({\mathfrak m}_f\) is the pointwise minimum over
all smooth invariant extensions of this flat restriction
using only the declared support.
\end{theorem}

\begin{proof}
Path transports telescope under a gauge transformation:
\(P_x^T(A^\pm)\mapsto
g_r^\pm P_x^T(A^\pm)(g_x^\pm)^{-1}\).
Hence every dressed vector transforms as
\({\cal V}_x\mapsto g_r^-{\cal V}_x(g_r^+)^{-1}\).
The two invariances of \(f\) prove gauge invariance of
each term. Matrix multiplication and inversion on
\(\SU(2)\), followed by a finite sum, give global
smoothness without a coordinate chart or gauge fixing.
At flat spatial links all path transports are the
identity, so every summand restricts to \(f(v)\).

At this same flat configuration, differentiating
a lower path link \(e\) contributes
\(I_{T^-}^x(e)q_x^-\); differentiating its upper inverse
path contributes \(-I_{T^+}^x(e)q_x^+\).
Averaging and using Proposition~\ref{f15v73:treecurrent}
and neutrality yields
\begin{equation}
 a^-=\partial^{\mathsf T}{\cal G}q^-,
                   \qquad
 a^+=-\partial^{\mathsf T}{\cal G}q^+ ,
 \label{eq:f15v73:optimalderivatives}
\end{equation}
with each scalar operator applied separately to colours.
These are exactly the two minimum-energy solutions in
Proposition~\ref{f15v73:ward}. The vertical derivatives
are those of \(f\), because the restrictions agree
as functions of every retained variable.
Substitution in \eqref{eq:f15v73:integrand} proves
\eqref{eq:f15v73:attainment} and the minimum assertion.
\end{proof}

The average is taken over observables \emph{before}
squaring their gradients. Averaging the tree energies
instead would keep an unnecessary nonnegative excess.
For any alternative extension the excess is exactly
the sum of the two squared residuals in
\eqref{eq:f15v73:pythagoras}, evaluated at flat links.
This includes extensions made with any fixed pair of trees.

The construction preserves the flat restriction,
not the original single-tree observable at nonflat
connections. It may also use every edge of the two
spatial graphs. Allowing other layers or extra links
would change the Ward network and the variational
problem; that larger minimization is not asserted here.
At nonflat connections the covariant lower bound remains
valid, but the scalar uniform-tree construction is not
claimed to attain it.

\subsection{Completion of the particular slow probe}

Now take the graph to be \(Q_L\), with \(L\ge4\) even,
and \(W=E\) as in V72. Use exactly its retained law
\(\rho_{L,\gamma}\), mode \(\varphi\), eigenvalue
\(\kappa_L\), and flat probe
\[
                  f_\gamma(v)=\gamma|S(v)|^2,\qquad
                  S(v)=\sum_{x\in E}\varphi_xv_x .
\]
These quantities are defined in
\eqref{eq:f15v72:marginal},
\eqref{eq:f15v72:mode},
\eqref{eq:f15v72:dispersion}, and
\eqref{eq:f15v72:test}.
The exact charge derivatives are
\begin{equation}
 \begin{split}
 q_x^{-,\alpha}
       &=2\gamma\varphi_x\,S\cdot(E_\alpha v_x),\\
 q_x^{+,\alpha}
       &=2\gamma\varphi_x\,S\cdot(v_xE_\alpha).
 \end{split}
 \label{eq:f15v73:probecharges}
\end{equation}
Their sums vanish: they are respectively
\(2\gamma S\cdot E_\alpha S\) and
\(2\gamma S\cdot SE_\alpha\), both zero since
left and right multiplication by \(E_\alpha\)
are skew-symmetric orthogonal generators.

For any invariant extension \({\cal F}_\gamma\)
with this restriction and the declared two-slice support,
define
\begin{equation}
 {\cal E}^{\mathrm{comp,flat}}_{L,\gamma}({\cal F}_\gamma)
     =\mathbb E_{\rho_{L,\gamma}}
           {\mathfrak e}({\cal F}_\gamma)(1,1,v).
 \label{eq:f15v73:completedform}
\end{equation}
This is the conditional expectation of the full
allowed-link gradient integrand, with the exterior
fixed flat. It is not, as defined, a Dirichlet form
on \(f_\gamma(v)\) alone: different extensions have
different transverse derivatives. It is also not
the unconditional Wilson expectation or an identified
physical transfer-Hamiltonian form.

\begin{theorem}[Exact minimum and its low-temperature quotient]
\label{f15v73:completedlimit}
For fixed even \(L\ge4\) and every \(\gamma>0\),
the uniform-tree extension \(\overline{\cal F}_\gamma\)
minimizes \eqref{eq:f15v73:completedform} over all the
specified invariant extensions of \(f_\gamma\).
Its thermally normalized quotient satisfies
\begin{equation}
 \lim_{\gamma\to\infty}
 \frac{\gamma^{-1}
       {\cal E}^{\mathrm{comp,flat}}_{L,\gamma}
                          (\overline{\cal F}_\gamma)}
      {\Var_{\rho_{L,\gamma}}(f_\gamma)}
                              =4+2\kappa_L .
 \label{eq:f15v73:limit}
\end{equation}
For every other such family \({\cal F}_\gamma\),
the corresponding liminf is at least \(4+2\kappa_L\).
\end{theorem}

\begin{proof}
The finite-\(\gamma\) minimum follows immediately by
integrating the pointwise minimum in
Theorem~\ref{f15v73:optimal} against the same law.

It remains to compute the two charge costs in
\eqref{eq:f15v73:minimal}. We use the fully justified
fixed-volume Laplace limit of
Proposition~\ref{f15v72:laplace}, not a new
large-volume Gaussian assumption. In its pinned chart,
\(v_x=h\exp(X_x/\sqrt\gamma)\), and
\(Y=\sum_E\varphi_xX_x\) tends in all required polynomial
moments to \(N(0,\kappa_L^{-1}I_3)\).
At every fixed \(X,h\),
\[
 S=\gamma^{-1/2}hY+O_{L,X}(\gamma^{-1}),
\]
so \eqref{eq:f15v73:probecharges} gives
\begin{equation}
 \gamma^{-1/2}q_x^-\longrightarrow
                   2\varphi_x\operatorname{Ad}_hY,\qquad
 \gamma^{-1/2}q_x^+\longrightarrow2\varphi_xY .
 \label{eq:f15v73:chargelimit}
\end{equation}
Indeed \(hY\cdot E_\alpha h
=Y\cdot h^{-1}E_\alpha h\), whose three components
are those of \(\operatorname{Ad}_hY\).

Extend \(\varphi\) by zero on the odd sites to form
\(\ell\). The explicit Laplace solution in V72 proves
\(\ell^{\mathsf T}{\cal G}\ell=\kappa_L^{-1}\).
Since \({\cal G}\) acts on vertices and rotations act
on colour, \eqref{eq:f15v73:chargelimit} implies, for
each endpoint sign,
\begin{equation}
 \gamma^{-1}\sum_\alpha
       (q^{\pm,\alpha})^{\mathsf T}{\cal G}q^{\pm,\alpha}
       \longrightarrow \frac4{\kappa_L}|Y|^2 .
 \label{eq:f15v73:boundarypointwise}
\end{equation}

The expectation limit requires a moment bound as well.
By equality of left and right frame norms and
\eqref{eq:f15v72:gradient},
\begin{equation}
 0\le\gamma^{-1}\sum_\alpha
       (q^{\pm,\alpha})^{\mathsf T}{\cal G}q^{\pm,\alpha}
 \le\|{\cal G}\|\gamma^{-1}\sum_E|\nabla_x f_\gamma|^2
 \le4\|{\cal G}\|f_\gamma .
 \label{eq:f15v73:domination}
\end{equation}
For fixed \(L\), \(\|{\cal G}\|\) is finite.
On the pinned chart V72 bounds \(f_\gamma\) by
\(\sum_E|X_x|^2\) and the rescaled density by a
fixed integrable Gaussian. Outside that chart
\(f_\gamma\le\gamma n\) and the probability is at most
a fixed polynomial in \(\gamma\) times
\(e^{-\eta_L\gamma}\). Thus dominated convergence
on the chart and this complementary bound prove
\begin{equation}
 \gamma^{-1}\mathbb E_{\rho_{L,\gamma}}
       \sum_\alpha(q^{\pm,\alpha})^{\mathsf T}{\cal G}q^{\pm,\alpha}
                           \longrightarrow\frac{12}{\kappa_L^2}.
 \label{eq:f15v73:boundarymean}
\end{equation}
No unproved inverse moment of a random connection
has entered this flat calculation.

V72 already gives
\[
 \gamma^{-1}\mathbb E_{\rho_{L,\gamma}}
                  \sum_E|\nabla_x f_\gamma|^2
                          \longrightarrow\frac{12}{\kappa_L},
 \qquad
 \Var_{\rho_{L,\gamma}}(f_\gamma)
                          \longrightarrow\frac6{\kappa_L^2}.
\]
Adding the two boundary limits and dividing proves
\[
 \frac{12/\kappa_L+24/\kappa_L^2}{6/\kappa_L^2}
                                  =2\kappa_L+4.
\]
For an arbitrary extension, the pointwise lower bound
and the positive limiting variance give the claimed
liminf. This does not require convergence of its excess.
\end{proof}

The order of limits is explicit: first
\(\gamma\to\infty\) at fixed \(L\), then, if desired,
\(L\to\infty\). The minimized completed quotient tends
to four in that iterated limit. V72's quotient
\(2\kappa_L\sim16\pi^2/L^2\) remains correct for its
retained-only conditional form; it does not remain
small for any extension in the present two-slice class
when all its allowed-link gradients are counted.
There is no contradiction between these two statements.

\subsection{Upstream consequence and the remaining analytic step}

The correction is now constructive rather than a
generic warning about omitted derivatives. Gauge
invariance identifies their necessary charge source;
the inverse covariant Laplacian gives their least
possible cost; and at flat links an explicit global
observable realizes that cost exactly. A single fixed
path dressing can only increase the flat cost.

The next averaged comparison must retain the actual
transverse law in \eqref{eq:f15v73:averaged} and control
the variance of the full observable under that same law.
The positive cost for this particular family supplies
no lower bound for the quotient of every physical
observable. Even an appropriate all-observable
auxiliary coercivity theorem would still require the
previously identified physical transfer-time calibration
and a nontrivial interacting continuum construction.
Neither a numerical constant in the round-link metric
nor the current finite-volume conditional calculation
settles those requirements.

\subsection{Diagnostics, preservation, and source-only output}

The exact-rational diagnostic is
\begin{center}\small
\path{scripts/verify_ym_f15_v73_transverse_dirichlet_completion.py}.
\end{center}
It reruns V72 and its inherited checks. Its new checks
include 18 uniform-tree signed-path identities,
18 two-forest potential certificates, and 4,068
tree-current orthogonal-excess identities. There are
72 flat completion tests with nonlinear quaternion
charges, nine covariant least-energy identities,
and nine gauge-covariance comparisons for their costs.

Four literal nonabelian Ward identities are checked
by differentiating actual tree-averaged loop observables,
including noncommuting transverse configurations.
Four corresponding observable completion bounds and
four completed Rayleigh constants are also checked.
The covariant systems include singular and nonsingular
cases; no numerical eigenvalue threshold is used.
These finite checks supplement, and do not replace,
the all-graph proofs and compact moment argument above.

The script has 12,646 bytes and SHA--256
\begin{center}\scriptsize\ttfamily
82DF5498F4537E623FF03C73CCCE5A286C901444CA240729DFE700C4E6B8B1C4.
\end{center}
The V72 predecessor has 1,636,836 bytes and SHA--256
\begin{center}\scriptsize\ttfamily
B286C4B7ECCC8398F04570B463413E8B8EFBEB6DDDC35793D97C9C7A8083D66C.
\end{center}
Removing this final section and restoring the title
version recovers that predecessor exactly. V73 replaces
V72 as the sole active main note, with protected archives,
companion inputs, and earlier diagnostics left unchanged.
Only \path{.tex} sources are kept in \path{latex_notes};
all compilation products are directed to
\begin{center}\small
\path{artifacts/ym_f15_v73_texcheck/}.
\end{center}

\begin{firewall}
V73 proves the pointwise covariant Ward energy bound,
an exact flat minimizing extension by tree averaging,
and the completed conditional limit \(4+2\kappa_L\)
for V72's probe. It neither proves a gap for all
observables under the actual transverse law nor
constructs the interacting continuum or a positive
physical Yang--Mills mass gap. The original goal
therefore remains active.
\end{firewall}

\section{V74: the full Gaussian invariant sector and rare compact packets}
\label{f15v74:section}

V73 completes one invariant probe by including the spatial
derivatives forced by gauge invariance. This version asks
whether that completion controls a whole observable class.
The answer is exact for the limiting Gaussian form: its
entire colour-rotation-invariant sector has gap
\(4+2\kappa_L\). We then realize every fixed invariant
polynomial family by smooth compact observables and prove
convergence of their covariance and completed energy matrices.

There is also a necessary caution with a concrete proof.
Normalized compact observables can place essentially all
their squared mass in an exponentially rare, nonaligned
region while their thermal completed energy stays bounded.
Thus compact concentration and convergence of every fixed
polynomial family do not justify interchanging the spectral
infimum with the low-temperature limit. The missing step
is a low-energy localization estimate, not another
calculation of a fixed Gaussian moment.

The V73 source and all 130 protected inputs and diagnostics
are unchanged. The next companion checkpoint remains V75.
The Hermite method itself is standard, and earlier archives
already use it for different Gaussian samplers. The new
statements here concern V73's particular completed metric,
its compact recovery map, and the explicit failure of
unrestricted bounded-energy localization.

\subsection{The covariance and completed metric on the retained quotient}

Keep \(L\ge4\) even, the bipartition \(E,O\), and the
operators \(B,K_L,{\cal L}_L\) from V72. Write
\[
 {\cal S}=\left\{z\in\mathbb R^E:\sum_{x\in E}z_x=0\right\},
 \qquad K=K_L|_{\cal S},\qquad {\cal G}={\cal L}_L^\dagger .
\]
Thus \(K\) is positive definite on \({\cal S}\).
Operators below act on vertices and separately on the
three colour components.

\begin{proposition}[The Schur Green identity for every neutral charge]
\label{f15v74:schur}
For any \(b,c\in{\cal S}\), extend them by zero to the odd
vertices, writing \(\widetilde b,\widetilde c\). Then
\begin{equation}
       \widetilde b^{\mathsf T}{\cal G}\widetilde c
                                =b^{\mathsf T}K^{-1}c .
 \label{eq:f15v74:schur}
\end{equation}
The centered retained Gaussian obtained from V72's pinned
full-spin limit therefore has law
\begin{equation}
 d\mu_K(z)=Z_K^{-1}
       \exp\left\{-\frac12\sum_{\alpha=1}^3
                       (z^\alpha)^{\mathsf T}Kz^\alpha\right\}dz
       \quad\text{on }{\cal S}\otimes\mathbb R^3 .
 \label{eq:f15v74:gaussian}
\end{equation}
\end{proposition}

\begin{proof}
Let \(h_E=K^{-1}c\) and
\(h_O=B^{\mathsf T}h_E/6\). The block formula for
\({\cal L}_L\) gives
\({\cal L}_Lh=(c,0)\). Any two solutions differ by
a constant, which pairs to zero with \((b,0)\).
This proves \eqref{eq:f15v74:schur}. The full pinned
Gaussian gives this same inverse quadratic form on
every neutral functional. Equivalently, integrating
its odd coordinates gives the Schur complement \(K_L\);
removing the remaining common constant gives
\eqref{eq:f15v74:gaussian}. The three colours remain
independent.
\end{proof}

For a smooth function \(p\) on this Gaussian space,
take its gradient intrinsically in \({\cal S}\) for
each colour, and define
\begin{equation}
 {\cal Q}_K(p)
  =\mathbb E_{\mu_K}\sum_{\alpha=1}^3
       (\nabla_\alpha p)^{\mathsf T}
                       (I+2K^{-1})(\nabla_\alpha p).
 \label{eq:f15v74:gaussianform}
\end{equation}
The identity term is the retained vertical gradient.
The other two terms are the two minimum transverse
costs in V73, by \eqref{eq:f15v74:schur}.
This is the tangent form of that specified completion;
it is not an assertion about an unconditional
Wilson generator.

\subsection{An all-observable sharp Gaussian inequality}

The invariance required here is simultaneous
\(z_x\mapsto Rz_x\) for \(R\in SO(3)\).
It is the adjoint action left after removing the
common left endpoint orientation. Invariance under
this group is not the same as evenness:
scalar triple products of three colour vectors
are invariant and odd.

\begin{theorem}[Sharp completed Gaussian invariant gap]
\label{f15v74:gaussiangap}
Let \(\kappa_{\min}\) be the least eigenvalue of \(K\).
Every \(SO(3)\)-invariant function in the closed
Gaussian Sobolev form domain satisfies
\begin{equation}
       {\cal Q}_K(p)
          \ge (4+2\kappa_{\min})\Var_{\mu_K}(p).
 \label{eq:f15v74:gaussiangap}
\end{equation}
The constant is attained by a centered squared norm
of a lowest spatial eigenmode. On the retained
three-dimensional torus,
\(\kappa_{\min}=\kappa_L\) from
\eqref{eq:f15v72:dispersion}.
\end{theorem}

\begin{proof}
Diagonalize \(K\) orthonormally on \({\cal S}\), with
eigenvalues \(\kappa_1,\ldots,\kappa_{n-1}>0\).
Write \(z_i\in\mathbb R^3\) for the corresponding
colour vectors and \(y_i=\sqrt{\kappa_i}z_i\).
These vectors are independent standard Gaussians, and
\[
 {\cal Q}_K(p)
       =\sum_{i=1}^{n-1}(\kappa_i+2)
                            \mathbb E|\nabla_{y_i}p|^2 .
\]
For completeness, the normalized multivariate
Hermite basis \(H_{\boldsymbol a}\), with
\(\boldsymbol a=(a_{i\alpha})\), gives
\begin{equation}
 \begin{split}
 \Var(p)&=\sum_{\boldsymbol a\ne0}|c_{\boldsymbol a}|^2,\\
 {\cal Q}_K(p)
   &=\sum_{\boldsymbol a}
       \left(\sum_{i,\alpha}(\kappa_i+2)a_{i\alpha}\right)
                                     |c_{\boldsymbol a}|^2 .
 \end{split}
 \label{eq:f15v74:hermite}
\end{equation}
The one-dimensional generating function
\(e^{ty-t^2/2}\) proves orthogonality; differentiating
the normalized degree-\(m\) polynomial gives
\(\sqrt m\) times its degree-\((m-1)\) predecessor.
Tensorization gives the displayed weights.
Completeness follows, for example, because a function
orthogonal to all polynomials has identically zero
Gaussian-weighted exponential transform, and hence
zero Fourier transform. The transform is entire by
Cauchy--Schwarz and Gaussian exponential moments.
Integration by parts identifies the coefficients of
weak derivatives. Parseval then proves
\eqref{eq:f15v74:hermite} throughout the closed form
domain, not just for polynomials.

For an invariant \(p\), the vector
\(\mathbb E[p\,y_i]\) is fixed by every \(R\in SO(3)\),
and must vanish. Thus every degree-one coefficient
is zero. No parity assumption is needed. Every
remaining nonconstant Hermite coefficient has total
degree at least two, and its weight is at least
\(2(\kappa_{\min}+2)\).
This proves \eqref{eq:f15v74:gaussiangap}.
If \(\kappa_i=\kappa_{\min}\), the invariant polynomial
\(|y_i|^2-3\) has exactly that eigenvalue.

It remains to identify the least retained eigenvalue.
The characters on \(E\) are represented by momenta
\(k_j=2\pi m_j/L\), with \(k\) and
\(k+(\pi,\pi,\pi)\) representing the same character.
The adjacency symbol between parity classes is
\(b(k)=2\sum_{j=1}^3\cos k_j\), and the eigenvalue of
\(K_L\) is \(6-b(k)^2/6\).
Only \(k=0\) and its equivalent all-\(\pi\) momentum
give constants. For every other class, changing to
the equivalent momentum if necessary makes \(b(k)\)
nonnegative, and at least one coordinate is nonzero.
Therefore
\[
 |b(k)|\le4+2\cos(2\pi/L).
\]
The single-coordinate momentum attains this bound.
Its eigenvalue is precisely \(\kappa_L\).
\end{proof}

A standard background reference for the Hermite
proof is Section 3.5 of Joel Tropp,
\href{https://tropp.caltech.edu/notes/Tro21-Probability-High-LN-corr.pdf}
{Probability in High Dimensions}.
The factor \(\kappa_i+2\) and the invariant-sector
removal of the entire linear chaos have been established
for the present form above. This is a sharp inequality
for every Gaussian invariant observable, not a lower
bound inferred from one trial quotient.

\subsection{Compact recovery of every fixed invariant polynomial family}

Choose a retained root \(r\in E\), and let
\[
       w_x=v_r^{-1}v_x,\qquad w_r=1 .
\]
Common left multiplication cancels; common right
multiplication conjugates all \(w_x\).
On a fixed sufficiently small product neighbourhood
of \(w_x=1\), take the principal quaternion logarithms
\(A_x=\log w_x\), \(A_r=0\), and put
\[
 (PA)_x=A_x-\frac1n\sum_{y\in E}A_y .
\]
Choose a smooth conjugation-invariant cutoff
\(\chi(w)\), equal to one near this aligned
configuration and with compact support inside that
logarithm chart. It can be chosen as a smooth
function of \(\sum_{x\ne r}|A_x|^2\).
For an \(SO(3)\)-invariant polynomial \(p\) on
\({\cal S}\otimes\mathbb R^3\), set
\begin{equation}
 ({\cal R}_\gamma p)(v)
   =p(0)+\chi(w)\bigl[p(\sqrt\gamma\,PA)-p(0)\bigr].
 \label{eq:f15v74:recovery}
\end{equation}
The cutoff term is extended by zero outside its chart.
This defines a globally smooth bi-invariant function
on the compact retained configuration space.
The map \({\cal R}_\gamma\) is linear and preserves
constants exactly.

Let \({\mathfrak m}_f\) denote V73's pointwise minimal
completed integrand. Its polarization defines
\({\mathfrak m}_{f,g}\).

\begin{theorem}[Compact covariance and completed-energy recovery]
\label{f15v74:recovery}
At fixed \(L\), for any two invariant polynomials \(p,q\),
\begin{equation}
 \begin{split}
 \operatorname{Cov}_{\rho_{L,\gamma}}
              ({\cal R}_\gamma p,{\cal R}_\gamma q)
                  &\longrightarrow\operatorname{Cov}_{\mu_K}(p,q),\\
 \gamma^{-1}\mathbb E_{\rho_{L,\gamma}}
              {\mathfrak m}_{{\cal R}_\gamma p,{\cal R}_\gamma q}
                  &\longrightarrow{\cal Q}_K(p,q).
 \end{split}
 \label{eq:f15v74:recoverylimits}
\end{equation}
Every recovered function has a global two-slice
gauge-invariant extension attaining this minimal
integrand at flat spatial links, by V73.
\end{theorem}

\begin{proof}
Compute expectations in V72's full pinned spin law:
\(u_r=h\), \(u_x=h\exp(X_x/\sqrt\gamma)\), \(X_r=0\).
Its Gaussian limit and polynomial domination were
proved in Proposition~\ref{f15v72:laplace}.
The centered retained coordinate
\(z=P(X|_E)\) has law \(\mu_K\) in the limit, by
Proposition~\ref{f15v74:schur}. On every fixed \(X\),
the cutoff is eventually one, so
\({\cal R}_\gamma p\to p(z)\).

Here is also the required derivative computation.
Let \(b_x=\nabla_{z_x}p(z)\), with
\(\sum_E b_x=0\).
For a nonroot link, differentiation of the logarithm
at the identity gives \(b_x\) in the right frame
and \(\operatorname{Ad}_h b_x\) in the left frame.
Varying the root changes every nonroot relative
coordinate with the opposite sign; its derivative
is \(-\sum_{x\ne r}b_x=b_r\).
Consequently the exact compact charges have the limits
\begin{equation}
 \gamma^{-1/2}q_x^-\longrightarrow\operatorname{Ad}_h b_x,
 \qquad
 \gamma^{-1/2}q_x^+\longrightarrow b_x .
 \label{eq:f15v74:recoverycharges}
\end{equation}
The logarithm derivative converges uniformly on every
bounded rescaled coordinate set, justifying these
first-order statements for polynomial \(p\).
Cutoff derivatives vanish on such sets for large
\(\gamma\). Applying \eqref{eq:f15v74:schur} to each
colour of \(b\) gives exactly the pointwise limit
\eqref{eq:f15v74:gaussianform}.

The expectation limits require more than this local
calculation. In a fixed small full pinned chart,
polynomial growth and bounded chart derivatives give
\[
 |{\cal R}_\gamma p|
      +\gamma^{-1/2}|q^-|
      +\gamma^{-1/2}|q^+|
                       \le C_{p,L}(1+|X|^{d_p})
\]
for some fixed integer \(d_p\), uniformly for
\(\gamma\ge1\). Derivatives of the fixed cutoff
satisfy the same bound. The finite operator
\({\cal G}\) in \({\mathfrak m}\) only changes
the constant. V72's rescaled Gaussian majorant
therefore dominates all products needed here.
On the complement of that full chart, the compact
functions and gradients grow at most polynomially
in \(\gamma\), while the probability is bounded
by a polynomial times \(e^{-\eta_L\gamma}\).
These contributions vanish. Dominated convergence
proves the means, products and energy limits;
polarization gives the two-function statement.
The final extension assertion is exactly
Theorem~\ref{f15v73:optimal}.
\end{proof}

\begin{proposition}[Every fixed polynomial-sector variational limit]
\label{f15v74:finitegram}
Let \({\cal P}\) be a finite-dimensional invariant
polynomial space containing the constants and
a lowest-mode squared norm.
The infimum of the thermal completed Rayleigh
quotient on \({\cal R}_\gamma{\cal P}\), modulo
constants, tends to \(4+2\kappa_L\).
In particular this holds for the entire invariant
polynomial space of any fixed degree at least two.
\end{proposition}

\begin{proof}
Choose a basis modulo constants. Its Gaussian
covariance matrix is positive definite: a nonconstant
polynomial cannot be constant almost everywhere
under a positive Gaussian density. By
\eqref{eq:f15v74:recoverylimits}, both its compact
covariance matrix and its thermal completed-energy
matrix converge entrywise to their Gaussian matrices.
The covariance matrices are therefore uniformly
positive definite for sufficiently large \(\gamma\).
The minimum generalized Rayleigh value converges,
by minimization on a fixed finite-dimensional unit
sphere. Theorem~\ref{f15v74:gaussiangap} supplies
the lower bound for the limiting matrix, and the
included squared mode attains it.
\end{proof}

No degree, dimension or spatial side is allowed to
grow in this argument. It proves a compact recovery
and finite-sector convergence, not convergence of
the full compact spectral infimum.

\subsection{Rare nonaligned packets with bounded thermal completed energy}

Define that unrestricted compact infimum by
\begin{equation}
 g^{\rm comp}_{L,\gamma}
  =\inf_{\substack{f\ {\rm smooth,\ bi\!-\!invariant}\\
                   \Var_{\rho_{L,\gamma}}f>0}}
       \frac{\gamma^{-1}\mathbb E_{\rho_{L,\gamma}}{\mathfrak m}_f}
            {\Var_{\rho_{L,\gamma}}f}.
 \label{eq:f15v74:compactgap}
\end{equation}
V73, or the recovery theorem, proves only
\[
           \limsup_{\gamma\to\infty}g^{\rm comp}_{L,\gamma}
                                      \le4+2\kappa_L .
\]
A reverse inequality cannot be obtained simply by
asserting that every bounded-energy normalized
sequence concentrates in the aligned Gaussian chart.
The following counterexample is within the same
compact retained law and the same completed form.

\begin{theorem}[Failure of unrestricted bounded-energy localization]
\label{f15v74:rarepacket}
At every fixed even \(L\ge4\), there are smooth
bi-invariant \(f_\gamma\) with
\[
 \mathbb E_{\rho_{L,\gamma}}f_\gamma=0,\qquad
 \mathbb E_{\rho_{L,\gamma}}f_\gamma^2=1,\qquad
 \sup_{\gamma\ {\rm sufficiently\ large}}
       \gamma^{-1}\mathbb E_{\rho_{L,\gamma}}{\mathfrak m}_{f_\gamma}
                                                      <\infty ,
\]
and gauge-invariant events \(D_\gamma\) such that
\begin{equation}
 \rho_{L,\gamma}(D_\gamma)
       \le C_L\gamma^{3(L^3-1)/2}e^{-\eta_L\gamma},
 \qquad
 \int_{D_\gamma}f_\gamma^2\,d\rho_{L,\gamma}\longrightarrow1 .
 \label{eq:f15v74:rarepacket}
\end{equation}
The events stay away from the aligned retained orbit.
Thus the squared mass of this bounded-energy sequence
escapes every fixed rescaled neighbourhood of that orbit.
\end{theorem}

\begin{proof}
Choose distinct retained sites \(r,x_*\).
Consider the orbit of the relative configuration
\[
 w_{x_*}=-1,\qquad w_x=1\quad(x\ne x_*),\qquad w_r=1 .
\]
Write these signs as \(\varepsilon_x\).
Near this orbit use the smooth relative chart
\[
       w_x=\varepsilon_x\exp(A_x),\qquad A_r=0 .
\]
The signs are central, so simultaneous conjugation
rotates all \(A_x\) by the same \(SO(3)\) matrix.

At the chart origin, each odd staple sum is a
positive real quaternion: it is four at the six
odd neighbours of \(x_*\) and six elsewhere.
There are no zero staple sums. Every first variation
of a retained signed quaternion is imaginary.
Consequently, for
\[
               r_q(A)=\left|\sum_{x\sim q}
                                     \varepsilon_x\exp(A_x)\right|,
\]
one has \(dr_q(0)=0\) and, on a sufficiently small
fixed chart,
\[
              r_q(A)\ge2,\qquad
                  |r_q(A)-r_q(0)|\le C_L|A|^2 .
\]
The leading effective action difference from
alignment is \(6\cdot(6-4)=12\), although the
probability bound below will not require an
asymptotic formula for that number.

Let \(d=3(n-1)\). Product Haar measure factors
exactly into the root orientation and the relative
variables. After integrating the orientation,
write \(p_\gamma(A)\,dA\) for the normalized retained
density in this chart. Its exact density is a smooth
Haar Jacobian times
\(\prod_{q\in O}Z(\gamma r_q(A))\) and the common
normalizer from \eqref{eq:f15v72:marginal}.
Since
\[
             |(\log Z)'(s)|\le1 ,
\]
by its tilted-Haar expectation formula,
the preceding quadratic estimate implies
\[
 \left|\log\frac{p_\gamma(A)}{p_\gamma(0)}\right|
        \le C_L\gamma|A|^2+C_L
          \quad\text{on the fixed chart, }\gamma\ge1.
\]
In particular this density ratio is bounded above
and below by positive constants, independent of
\(\gamma\), when \(|A|\le\gamma^{-1/2}\).

Take a nonzero nonnegative smooth radial function
\(\psi\) compactly supported in the unit ball of
\(\mathbb R^d\), and set
\(b_\gamma(v)=\psi(\sqrt\gamma A)\) in this chart,
extended by zero. For sufficiently large \(\gamma\)
its support is strictly inside the fixed chart,
so the extension is smooth and bi-invariant.
Let \(D_\gamma=\{|A|\le\gamma^{-1/2}\}\) in the chart.
The density comparison gives
\begin{equation}
 \mathbb E_\rho b_\gamma^2
             \ge c_Lp_\gamma(0)\gamma^{-d/2},\qquad
 \gamma^{-1}\mathbb E_\rho|\nabla b_\gamma|^2
             \le C_Lp_\gamma(0)\gamma^{-d/2}.
 \label{eq:f15v74:packetcomparison}
\end{equation}
Here constants absorb the fixed integrals of
\(\psi^2\) and \(|\nabla\psi|^2\).
The gradient is the retained round-product gradient.
Smooth chart derivatives are uniformly bounded,
including the derivative of the relative coordinates
with respect to the root. Moreover V73 gives
\[
       {\mathfrak m}_{b_\gamma}
          \le (1+2\|{\cal G}\|)|\nabla b_\gamma|^2 .
\]
Thus \eqref{eq:f15v74:packetcomparison} also bounds
the thermal completed energy relative to
\(\mathbb E_\rho b_\gamma^2\).

The orbit at the chart center is nonaligned.
A small fixed closed retained tube around it is
disjoint from the aligned orbit. On the compact
full pinned-spin space lying above this tube,
\(H_L\) has a positive minimum \(\eta_L\):
its only zeros have every spin equal.
The partition lower bound and compactness argument
in V72 therefore give the probability bound in
\eqref{eq:f15v74:rarepacket} for \(D_\gamma\).

Write \(s_\gamma=\rho(D_\gamma)\). Since \(b_\gamma\)
vanishes outside \(D_\gamma\), Cauchy--Schwarz gives
\[
 (\mathbb E_\rho b_\gamma)^2
       \le s_\gamma\mathbb E_\rho b_\gamma^2,\qquad
 \Var_\rho(b_\gamma)
       \ge(1-s_\gamma)\mathbb E_\rho b_\gamma^2 .
\]
Define
\[
       f_\gamma
       =\frac{b_\gamma-\mathbb E_\rho b_\gamma}
                           {\sqrt{\Var_\rho(b_\gamma)}} .
\]
The two normalization identities hold exactly.
The energy estimate above and \(s_\gamma\to0\)
give the bounded thermal completed energy.
Outside \(D_\gamma\), \(f_\gamma\) is constant and
\[
 \int_{D_\gamma^c}f_\gamma^2\,d\rho
   =\frac{(1-s_\gamma)(\mathbb E_\rho b_\gamma)^2}
                           {\Var_\rho(b_\gamma)}
                                       \le s_\gamma .
\]
This proves concentration of its squared mass in
\(D_\gamma\). That tube has a relative link near
\(-1\), whereas any fixed rescaled aligned
neighbourhood has every relative link near \(1\)
for large \(\gamma\). They are disjoint, proving
the final assertion.
\end{proof}

This counterexample does not show that the packet
quotients are small, that the Gaussian value is the
wrong compact gap limit, or that a physical gap
fails. Its constants may depend strongly on \(L\).
It shows that bounded thermal energy alone is
insufficient for the claimed aligned-chart
compactness. Ordinary probability concentration
is even weaker: normalization of the test can
compensate for an exponentially small probability.

\subsection{The upstream estimate now required}

The flat problem now has a complete Gaussian
invariant-sector inequality and a genuine compact
recovery for each fixed polynomial sector.
What is missing even there is a lower bound that
handles all compact observables, including rare
normalized packets. A valid variational localization
argument would have to exclude escape specifically
below its proposed energy threshold, or bound the
escaping sector directly. It cannot assume the
false unrestricted bounded-energy tightness property
disproved in Theorem~\ref{f15v74:rarepacket}.

For the actual Yang--Mills goal this is still only
a diagnostic subproblem. V73's pointwise covariant
Ward bound must ultimately be used with the actual
transverse-field law and a variance estimate under
that same law. The flat Gaussian completion is not
a replacement for either of those tasks, and an
auxiliary completed diffusion rate is not the
physical transfer-Hamiltonian mass. The physical
scale comparison and interacting continuum
construction remain open.

\subsection{Reproducibility and source-only versioning}

The exact-rational diagnostic is
\begin{center}\small
\path{scripts/verify_ym_f15_v74_completed_gaussian_sector.py}.
\end{center}
It reruns V73 and the inherited checks.
New tests include 84 infinitesimal rotation identities,
252 vanishing first-chaos coefficients, 84 invariant
polynomial gap inequalities, 27 quadratic and three
cubic exact eigenvalues, three sharp radial attainments,
and 36 mixed Gram-energy inequalities. The odd scalar
triple product is included explicitly, so evenness
has not been substituted for \(SO(3)\) invariance.

There are also 81 Hermite normalization/orthogonality
checks, 17 all-basis Schur Green bilinear identities,
18 rooted gradient identities, 144 charge-frame
covariance tests, two retained single-flip geometries,
and 420 cancellations of their radial first variations.
The compact limit and rare-packet estimates are
analytic proofs above, not claims inferred from
these finite diagnostics.

The script has 8,351 bytes and SHA--256
\begin{center}\scriptsize\ttfamily
16859E491004011D82440CAB82AB6A991E33642E40B4DF92DA73426E051769AF.
\end{center}
The V73 predecessor has 1,658,942 bytes and SHA--256
\begin{center}\scriptsize\ttfamily
881B3C105F6FF6A5F624DC0C458AE21596D02A78F9079B222BB68EE8A5F47774.
\end{center}
Removing this final section and restoring the title
version recovers that source exactly. V74 replaces
V73 as the sole active main note, with protected
archives, companions and earlier diagnostics unchanged.
Only \path{.tex} sources remain in \path{latex_notes}.
All compilation products are directed to
\begin{center}\small
\path{artifacts/ym_f15_v74_texcheck/}.
\end{center}

\begin{firewall}
V74 proves a sharp inequality for the full invariant
sector of the specified Gaussian completed form,
compact recovery of every fixed invariant polynomial
sector, and a counterexample to unrestricted
bounded-energy aligned-chart localization.
It does not prove convergence of the unrestricted
compact gap, an actual-transverse-law coercivity
bound, an interacting continuum, or a positive
physical Yang--Mills mass gap. The original goal
remains active.
\end{firewall}

\section{V75: native antipodal localization and the companion checkpoint}
\label{f15v75:section}

V74's rare normalized packets disprove unrestricted
bounded-energy concentration in the aligned Gaussian
chart. They do not show that their energy is low.
This version separates these two issues quantitatively.
Under the exact native compact conditional law, a
function supported near a central antipodal link with
strongly aligned exterior fields has energy at least
\(20\gamma\) times its squared norm. This is a supported
function estimate, not an asymptotic Gaussian replacement.

It follows that V74's explicit single-flip packets have
thermal completed quotients at least twenty, whereas
the aligned Gaussian candidate is at most \(32/3\).
Thus that concrete escaping family is not a low-energy
competitor. A same-law localization identity also controls
its neighbourhood for arbitrary observables, retaining
the transition-region cost. Other compact regions and
the actual physical comparison remain unresolved.

\subsection{Required V75 companion review}

Fresh inventories of \path{latex_notes} and the supplied
Downloads sources found no newer applicable versions.
The terminal results of both parallel programmes and the
additional supplied sources were reread. All 131 protected
inputs and diagnostics, including these companions,
were checked unchanged from the preceding ledger.
The mathematical relevance is as follows.

\begin{center}\small
\begin{tabular}{@{}p{0.29\textwidth}p{0.66\textwidth}@{}}
\toprule
Source & Scope retained at this checkpoint\\
\midrule
SM continuum V72 &
Fixed-time many-copy covariance flow; no growing-cutoff
or physical finite-species continuum estimate.\\[1mm]
NS stability V26 &
Rank-one derivative-kernel norms; no full stability
or global-regularity conclusion.\\[1mm]
Shell mixing V16 &
\raggedright Measured flat one-loop assembly; no common
nonperturbative volume/coupling remainder or suppression
of other compact configurations.\tabularnewline[1mm]
Parallel YM V5 &
Failure of a volume-uniform tail for an extensive
global charge; local observables and density estimates
remain the appropriate scope.\\[1mm]
SM source selection V31 &
Conditional word-native compiler reduction, not
historical occurrence of the compiler.\\[1mm]
Native stationarity V34 &
An escape-scaling reduction with the moving-fibre
coefficient tensor still unevaluated.\\[1mm]
Exact-action Einstein V24 &
A nonzero corrected first-jet tangency derivative,
not an integrated nonlinear root.\\
\bottomrule
\end{tabular}
\end{center}

The supplied \path{suggestions.txt} was reread as
optional research input, not as an instruction document.
Its same-law, same-source normalization warning remains
relevant. The suggested noncommuting coefficient work
already has its inherited archive interfaces; repeating
it does not settle the present weighted compact
localization problem. None of the reviewed sources
supplies the missing physical Yang--Mills estimate.

The direction therefore stays upstream at the exact
compact law, but moves from probability of a rare
region to the Dirichlet cost of a normalized function
living there. The next compulsory companion review is V80.

\subsection{One actual native conditional and its antipodal region}

Use exactly the native conditional
\eqref{eq:f15v69:conditional}, in its common central-link
quaternion frame:
\begin{equation}
 dP^\eta(v)
 =\frac{e^{\Psi(v,\eta)}\,dH(v)}
               {\int_{S^3}e^{\Psi(w,\eta)}\,dH(w)},\qquad
 \Psi(v,\eta)=\sum_{i=1}^6\ell(\gamma|c_i+v|),
 \qquad |c_i|\le5.
 \label{eq:f15v75:native}
\end{equation}
The six \(c_i\) are fixed when differentiating the
central variable \(v\); they come from the actual
five-unit exterior sums after the prescribed deletion.
Let
\[
 C=\sum_i c_i,\qquad D=\gamma(30-|C|),\qquad b=C/|C|
       \quad(C\ne0).
\]
On the exterior event
\begin{equation}
                   {\cal B}=\{D\le\gamma/1000\},
 \label{eq:f15v75:good}
\end{equation}
\(C\ne0\), so \(b\) is unambiguous. Define
\[
 t(v,\eta)=|v+b|,\qquad
 {\cal A}=\{\,\eta\in{\cal B},\ t\le1/10\,\}.
\]
Both \(D\) and \(t\) are invariant under the two endpoint
frame changes: \(v,c_i,b\) transform by the same
left and right unit-quaternion multiplications.
This is a gauge-invariant description of the relative
antipodal region, not a gauge fixing.

We use the same round metric as V69--V74.
If \(X_\alpha v=E_\alpha v\), then the three
left-invariant fields form an orthonormal frame,
are divergence free for Haar measure, and
\(\Delta=\sum_{\alpha=1}^3X_\alpha^2\).
The only radial facts needed below are
\begin{equation}
        0\le a(s)\le1,\qquad a'(s)\ge0,\qquad
                      a(s)\ge1-\frac3{2s}\quad(s>0).
 \label{eq:f15v75:radial}
\end{equation}
The middle statement is the tilted-Haar variance
identity. The last is V38's Gamma bound, reused in
V70. No new estimate on the Bessel quotient is assumed.

\begin{proposition}[A uniform positive log-density Laplacian]
\label{f15v75:laplacian}
For \(\gamma\ge10\), throughout \({\cal A}\),
\begin{equation}
        \Delta_v\Psi(v,\eta)
          \ge\frac{15386265}{768208}\,\gamma
                              >20\gamma .
 \label{eq:f15v75:laplacian}
\end{equation}
The bound is independent of volume and of the
exterior, subject to the displayed native conditions.
\end{proposition}

\begin{proof}
V69's exact alignment inequality gives, for
\(d_i=c_i-5b\),
\[
                 \sum_i|d_i|^2\le10D/\gamma\le1/100.
\]
Thus \(|d_i|\le1/10\). For any of the six fields, put
\[
             r=|c+v|,\qquad z=c\cdot v,\qquad
                         T=|c|^2-z^2 .
\]
Since \(|v+b|\le1/10\), one has
\(b\cdot v=-1+|v+b|^2/2\le-199/200\).
Consequently
\begin{equation}
 \frac{19}{5}\le r\le\frac{21}{5},\qquad
 z\le5(-199/200)+1/10=-\frac{39}{8},\qquad
                    0\le T\le\frac9{25}.
 \label{eq:f15v75:geometry}
\end{equation}
For the radius, write \(c+v=4b+d+(v+b)\).
For the last bound, the tangential projection satisfies
\[
       |\Pi_v c|=|\Pi_v(c+5v)|
          \le|d+5(b+v)|\le3/5 .
\]

For fixed \(c\), a function of \(z=c\cdot v\) has
Laplacian \(T\,d^2/dz^2-3z\,d/dz\).
Since \(r^2=|c|^2+1+2z\), the exact chain rule gives
\begin{equation}
 \Delta_v\ell(\gamma r)
   =\left(\frac{\gamma^2a'(\gamma r)}{r^2}
               -\frac{\gamma a(\gamma r)}{r^3}\right)T
                     -\frac{3\gamma a(\gamma r)}r\,z .
 \label{eq:f15v75:chain}
\end{equation}
Using \eqref{eq:f15v75:radial}, the first term is
at least \(-\gamma T/r^3\).
Also \(\gamma r\ge38\), so
\(a(\gamma r)\ge73/76\).
The negative sign of \(z\) in
\eqref{eq:f15v75:geometry} therefore yields
\[
 \gamma^{-1}\Delta_v\ell(\gamma r)
  \ge \frac{3(73/76)(39/8)}{21/5}
                      -\frac{9/25}{(19/5)^3}.
\]
Sum the six contributions. The resulting coefficient is
\[
 6\left\{\frac{3(73/76)(39/8)}{21/5}
                  -\frac{9/25}{(19/5)^3}\right\}
          =20+\frac{22105}{768208}>20.
\]
\end{proof}

At the exactly aligned exterior \(c_i=5b\) and
\(v=-b\), the same formula gives
\[
             \Delta_v\Psi=\frac{45}{2}\gamma a(4\gamma).
\]
Thus the convenient constant twenty lies below the
limiting antipodal value \(45/2\). It is not asserted
to be an optimal local spectral constant.

\subsection{The supported-function energy estimate under the same law}

\begin{proposition}[Exact weighted square identity]
\label{f15v75:square}
For every smooth real \(f\) on the central \(S^3\),
\begin{equation}
 \int|\nabla f|^2\,dP^\eta
    =\sum_{\alpha=1}^3
          \int|X_\alpha f+fX_\alpha\Psi|^2\,dP^\eta
          +\int f^2\Delta_v\Psi\,dP^\eta .
 \label{eq:f15v75:square}
\end{equation}
\end{proposition}

\begin{proof}
Expand one square. Haar integration by parts gives
\[
 \int X_\alpha(f^2)X_\alpha\Psi\,e^\Psi dH
  =-\int f^2\{X_\alpha^2\Psi+(X_\alpha\Psi)^2\}
                                                    e^\Psi dH.
\]
The squared first derivative of \(\Psi\) cancels
the corresponding square term. Divide by the same
conditional normalizer and sum over \(\alpha\).
There is no boundary term on the compact group.
\end{proof}

This is also the usual weighted ground-state inequality
with positive test \(W=e^{-\Psi}\):
for \(L=\Delta+\nabla\Psi\cdot\nabla\),
\(-LW/W=\Delta\Psi\). The general integration
principle is Lemma 1.1 of Cattiaux--Guillin,
\href{https://arxiv.org/abs/1001.1822}
{Functional Inequalities via Lyapunov conditions}.
The exact square, sign, and native coefficient needed
here have been derived above.

\begin{theorem}[Native antipodal supported cost]
\label{f15v75:supported}
Let \(\mu(dv,d\eta)=P^\eta(dv)\nu(d\eta)\) be any
one declared law with the native conditional
\eqref{eq:f15v75:native}. Suppose \(\gamma\ge10\).
If a function \(F(v,\eta)\), smooth in \(v\), vanishes
outside \({\cal A}\), then
\begin{equation}
       \mathbb E_\mu|\nabla_vF|^2
                             \ge20\gamma\,\mathbb E_\mu F^2 .
 \label{eq:f15v75:supported}
\end{equation}
It is enough that the integrals exist; the assertion
also applies to smooth cylinder functions under the
actual native deletion marginal.
\end{theorem}

\begin{proof}
For each \(\eta\in{\cal B}\), apply
\eqref{eq:f15v75:square} and
Proposition~\ref{f15v75:laplacian}. The supported
function is zero where the Laplacian bound was
not established. For \(\eta\notin{\cal B}\), it
vanishes identically. Integrate with respect to
the unchanged exterior law \(\nu\).
\end{proof}

There is no probability penalty or division by
\(\mu({\cal A})\) in \eqref{eq:f15v75:supported}.
The law has not been conditioned on \({\cal A}\),
and its native conditionals have not been replaced
by those of a restricted measure. Exterior flatness
is not required. When used with Wilson data,
\(\mu\) must be the actual specified deletion
marginal, not the undeleted fine law silently
assigned a different conditional.

\subsection{Localization for an arbitrary observable}

Supported functions are not enough for gluing.
We now retain the transition cost explicitly.
Choose a smooth nondecreasing \(\vartheta\) with
\[
 \vartheta(s)=0\ (s\le0),\qquad
 \vartheta(s)=1\ (s\ge1),\qquad
                         0\le\vartheta'(s)\le2 .
\]
Such a function is obtained by smoothing a clipped
linear transition whose support is strictly inside
\((0,1)\) and whose slope is less than two.
For \(\eta\in{\cal B}\), set
\[
 \theta(v,\eta)=\frac\pi2\,\vartheta(20t-1),
 \qquad \chi=\cos\theta,\qquad \zeta=\sin\theta.
\]
The compositions are smooth even at \(t=0\), since
they are constant near that point. They obey
\[
 \chi^2+\zeta^2=1,\qquad
 \chi=1\quad(t\le1/20),\qquad
 \chi=0\quad(t\ge1/10).
\]
The chord distance \(t=|v+b|\) has gradient norm
at most one wherever differentiated. Hence
\begin{equation}
 |\nabla_v\chi|^2+|\nabla_v\zeta|^2
   =|\nabla_v\theta|^2
   \le400\pi^2\,1_{\{1/20<t<1/10\}} .
 \label{eq:f15v75:cutoff}
\end{equation}
Only the central derivative is being taken:
\(b(\eta)\) is fixed in that derivative.

\begin{theorem}[Same-law weighted antipodal localization]
\label{f15v75:localization}
For every smooth central-variable \(F\) with finite
displayed integrals,
\begin{equation}
 \begin{split}
 20\gamma\,\mathbb E_\mu[1_{\cal B}\chi^2F^2]
  \le{}&\mathbb E_\mu[1_{\cal B}|\nabla_vF|^2]\\
       &+400\pi^2\,
          \mathbb E_\mu[
             1_{\cal B}1_{\{1/20<t<1/10\}}F^2].
 \end{split}
 \label{eq:f15v75:localization}
\end{equation}
Both terms on the right use the same exact marginal
as the weighted squared mass on the left.
\end{theorem}

\begin{proof}
Pointwise expansion and
\(\chi\nabla\chi+\zeta\nabla\zeta=0\) give the
partition-of-squares identity
\[
 |\nabla_v(\chi F)|^2+|\nabla_v(\zeta F)|^2
       =|\nabla_vF|^2+
          F^2\{|\nabla_v\chi|^2+|\nabla_v\zeta|^2\}.
\]
For every \(\eta\in{\cal B}\), \(\chi F\) is
supported in the antipodal cap. Apply
Theorem~\ref{f15v75:supported} conditionally to
\(\chi F\), discard the nonnegative
\(\zeta F\) energy, use \eqref{eq:f15v75:cutoff},
and integrate in \(\eta\).
The indicator \(1_{\cal B}\) is exterior measurable,
so no central derivative of it occurs.
\end{proof}

For a normalized family
\(\mathbb E_\mu F_\gamma^2=1\), if
\[
          \mathbb E_\mu[1_{\cal B}\chi^2F_\gamma^2]
                                                     \longrightarrow1,
\]
then \eqref{eq:f15v75:localization} implies
\[
     \liminf_{\gamma\to\infty}
       \gamma^{-1}\mathbb E_\mu|\nabla_vF_\gamma|^2
                                                     \ge20 .
\]
This excludes concentration of the entire squared
mass in this region below energy twenty. It does
not assert that every bounded-energy family has
vanishing mass there.

If a finite set \(J\) of distinct central variables
has these native conditionals under one common law,
the same identities can be summed:
\begin{equation}
 \begin{split}
 20\gamma\,\mathbb E_\mu\left[
                    F^2\sum_{j\in J}1_{{\cal B}_j}\chi_j^2\right]
 \le{}&\mathbb E_\mu\sum_{j\in J}
                          1_{{\cal B}_j}|\nabla_{v_j}F|^2\\
 &+400\pi^2\mathbb E_\mu\left[
          F^2\sum_{j\in J}1_{{\cal B}_j}
                      1_{\{1/20<t_j<1/10\}}\right].
 \end{split}
 \label{eq:f15v75:count}
\end{equation}
No independence is used. In particular a function
supported where every one of these variables is
in its good antipodal cap satisfies
\[
             \mathbb E_\mu\sum_{j\in J}|\nabla_{v_j}F|^2
                                \ge20\gamma|J|\mathbb E_\mu F^2 .
\]
The common-law premise is essential: different
overlapping V69 deletion marginals cannot simply
be combined. V72's flat retained law is an explicit
common-law example. The transition count in
\eqref{eq:f15v75:count} is not discarded or presumed
volume independent.

\subsection{The explicit V74 packets are above the candidate threshold}

Return to the common flat-boundary law
\(\rho_{L,\gamma}\). Use exactly the single-flip
chart and packet \(b_\gamma\) from
Theorem~\ref{f15v74:rarepacket}. On its support,
\[
 v_x=h\varepsilon_x\exp(A_x),\qquad
 \varepsilon_{x_*}=-1,\quad\varepsilon_x=1\ (x\ne x_*),
 \qquad A_r=0,\qquad\sum_{x\ne r}|A_x|^2\le1/\gamma .
\]
For the central variable \(x_*\), every term in its
thirty-term native exterior sum is a retained link
other than \(x_*\), hence has sign \(+1\).
Multiplicities from shared odd stars are kept.

\begin{proposition}[Uniform inclusion in the paid antipodal cap]
\label{f15v75:packetinclusion}
On this packet support, for \(\gamma\ge15000\),
\begin{equation}
        D_{x_*}\le15\le\gamma/1000,\qquad
          |v_{x_*}+b_{x_*}|\le3/\sqrt\gamma\le1/20 .
 \label{eq:f15v75:packetinclusion}
\end{equation}
The displayed numerical threshold is independent of \(L\).
\end{proposition}

\begin{proof}
Each unflipped link \(v_y=h\exp(A_y)\) satisfies
\[
 |v_y-h|\le\gamma^{-1/2},\qquad
               h\cdot v_y=\cos|A_y|\ge1-\frac1{2\gamma}.
\]
Summing thirty terms gives
\[
        |C_{x_*}-30h|\le30/\sqrt\gamma,\qquad
                     h\cdot C_{x_*}\ge30-15/\gamma .
\]
Therefore
\(D_{x_*}=\gamma(30-|C_{x_*}|)\le15\), and
\(C_{x_*}\ne0\). For any nonzero quaternion \(C\)
and unit \(h\),
\[
 \left|\frac C{|C|}-h\right|
   \le \left|\frac C{|C|}-\frac C{30}\right|
                              +\left|\frac C{30}-h\right|
   \le\frac{2|C-30h|}{30}.
\]
It follows that \(|b_{x_*}-h|\le2/\sqrt\gamma\).
The flipped link satisfies
\(|v_{x_*}+h|\le1/\sqrt\gamma\), proving the
second bound. Finally \(\gamma\ge15000\)
gives both \(15\le\gamma/1000\) and
\(9/\gamma\le1/400\).
\end{proof}

\begin{theorem}[The escaping packets are not low-energy competitors]
\label{f15v75:packetenergy}
For the centered normalized \(f_\gamma\) of V74
and \(\gamma\ge15000\),
\begin{equation}
       \gamma^{-1}\mathbb E_{\rho_{L,\gamma}}
                   {\mathfrak m}_{f_\gamma}\ge20 .
 \label{eq:f15v75:packetenergy}
\end{equation}
In contrast the aligned Gaussian invariant value obeys
\[
                    4+2\kappa_L\le32/3<20 .
\]
\end{theorem}

\begin{proof}
Apply Theorem~\ref{f15v75:supported} to the uncentered
supported packet \(b_\gamma\), whose inclusion is
proved above. The completed integrand dominates the
single central gradient, so
\[
 \gamma^{-1}\mathbb E_\rho{\mathfrak m}_{b_\gamma}
                     \ge20\,\mathbb E_\rho b_\gamma^2 .
\]
Subtracting its expectation changes no derivative.
Dividing by \(\Var_\rho(b_\gamma)\le
\mathbb E_\rho b_\gamma^2\) proves
\eqref{eq:f15v75:packetenergy}. This applies the
supported estimate before centering; the centered
function need not itself be supported in the cap.

For even \(L\ge4\), \(\cos(2\pi/L)\ge0\).
The formula for \(\kappa_L\) therefore gives
\(\kappa_L\le6-16/6=10/3\), proving the comparison.
\end{proof}

The same conclusion is additive for a fixed set
\(J\subset E\) of isolated flipped links, provided
no odd star contains two elements of \(J\) and the
retained root lies outside \(J\).
In that signed chart each flipped central link
still sees thirty unflipped exterior terms.
The proof of \eqref{eq:f15v75:packetinclusion}
holds at each one, under the same total chart-radius
bound. Applying the common-law sum gives the
normalized packet threshold \(20|J|\).
The V74 bounded-energy packet construction also
extends to this fixed pattern: all odd staple
magnitudes are four or six, their first derivatives
vanish, and the chart density comparison is unchanged
in form. Its constants may depend on \(L,J\).
No independent-product interpretation is required.

\subsection{What has been removed, and what has not}

V74's counterexample to unrestricted bounded-energy
chart concentration remains valid. V75 now locates
those explicit single-flip and isolated-multiflip
packets above the candidate low-energy band.
This is a positive localized estimate under the
native compact law, not an inference from their
small probability.

The weighted estimate \eqref{eq:f15v75:localization}
is also available at nonflat exteriors whenever
the actual native alignment and relative-cap
conditions hold. It is therefore a usable
same-law input for gluing. It does not control
the high-debit exteriors, the remaining central
regions, the summed transition cost at growing
volume, or nonisolated/noncollinear rare patterns.
Nor has it supplied a lower bound for the entire
compact infimum \(g^{\rm comp}_{L,\gamma}\).

The next estimate must control the remaining
low-energy regions or assemble a common-law
localization covering them with paid transition
terms. V71's probability bounds alone do not
supply the weighted \(F^2\) control required for
that step. The actual transverse-field variance,
physical transfer-time calibration, and interacting
continuum construction remain additional open
requirements of the original Yang--Mills goal.

\subsection{Exact diagnostics and preservation}

The diagnostic is
\begin{center}\small
\path{scripts/verify_ym_f15_v75_antipodal_localization.py}.
\end{center}
It reruns V74 and its inherited checks.
Four exact constant and threshold comparisons,
three native debit fixtures, 162 literal geometric
inequalities, and 270 independent sphere-jet chain
identities are checked over rational quaternions.
There are 108 two-endpoint covariance tests,
32 weighted square/divergence identities,
72 partition localization identities,
four packet-inclusion budgets, sixteen supported
multipacket threshold checks, and 36 literal
isolated-flip five-neighbour stars.

The arbitrary-function estimates follow from the
analytic integration identities, not from numerical
sampling. The field fixtures test the algebra;
they are not asserted to be independent fine
Wilson exteriors. The common retained geometry
is tested separately.

The script has 6,074 bytes and SHA--256
\begin{center}\scriptsize\ttfamily
6F6825A9417130D0B7D19DB76B2F7E6EDB43F12C48CAE0793B05B0F0AF8287FD.
\end{center}
The V74 predecessor has 1,682,196 bytes and SHA--256
\begin{center}\scriptsize\ttfamily
8BBB211549B88003648576089B36BC9C18442073AD1DEB878A7F9C32E9AD153F.
\end{center}
Removing this final section and restoring the
title version recovers that source exactly.
V75 replaces V74 as the sole active main note;
protected archives, companions and previous
diagnostics are unchanged. Only \path{.tex}
sources are kept in \path{latex_notes}.
All compilation products are directed to
\begin{center}\small
\path{artifacts/ym_f15_v75_texcheck/}.
\end{center}

\begin{firewall}
V75 proves an explicit native antipodal energy
threshold, a same-law weighted localization
inequality, and exclusion of V74's explicit
flipped packets from the candidate low-energy
band. It does not prove a full compact spectral
lower bound, an interacting continuum, or the
positive physical Yang--Mills mass gap.
The original goal remains active.
\end{firewall}

\section{V76: exact retained landscape and weighted control}
\label{f15v76:section}

V75 excludes particular antipodal packets from the candidate
low-energy band. The next step must not consist only of
adding more isolated patterns. This version studies the
entire exact flat-boundary retained potential, at finite
coupling, including configurations with zero odd-star sums.

There are two results. The only local minima of this exact
potential are fully aligned configurations; every other
critical point has a negative Hessian direction. More
quantitatively, collective projected ambient vector fields
give a global weighted Dirichlet inequality. It excludes
functions supported in a sufficiently high-frustration
region from the candidate energy band without estimating
the probability of that region.

These are statements under V72's one common flat-boundary
marginal. The finite-coupling logarithmic star factors are
not replaced by their leading low-temperature energy.
The same collective argument has not been established
at a genuinely nonflat transverse exterior.

The V75 predecessor and all 132 protected files were
verified unchanged. The next companion checkpoint
remains V80. Archive f12 already contains a local
negative-curvature integration framework; the new work
here supplies an exact global star-geometry calculation
and its weighted consequence for this particular law.
Projected common-direction perturbations are also
classical in spherical consensus theory. Compare
Proposition 16 of Markdahl, Thunberg and Goncalves,
\href{https://arxiv.org/abs/1609.04885}
{Almost global consensus on the \(n\)-sphere}.
That pair-interaction result is background, not a
theorem being substituted for the generated star law.

\subsection{The exact potential and a geometric frustration variable}

Retain \(E,O\), \(n=L^3/2\), and the common marginal
\(\rho_{L,\gamma}\) from V72. For each odd vertex \(q\),
put
\[
       S_q(v)=\sum_{x\sim q}v_x,\qquad r_q=|S_q|\le6.
\]
Normalize the potential to be zero at alignment:
\begin{equation}
 V_\gamma(v)
     =\sum_{q\in O}\{\ell(6\gamma)-\ell(\gamma r_q)\},
 \qquad
 d\rho_{L,\gamma}
     ={\cal Z}_\gamma^{-1}e^{-V_\gamma(v)}
                                  \prod_{x\in E}dH(v_x).
 \label{eq:f15v76:potential}
\end{equation}
The constant in the potential changes no derivative
or probability law. Define
\begin{equation}
                      J(v)=\sum_{q\in O}(6-r_q),
                 \qquad 0\le J\le6n .
 \label{eq:f15v76:J}
\end{equation}
This is a geometric quantity, not \(V_\gamma/\gamma\)
at finite \(\gamma\). It also equals the minimum
of V72's full spin energy over the odd spins
with the retained spins fixed: each odd spin can
maximize its own inner product with \(S_q\).
No uniqueness is needed when \(S_q=0\).

The ambient function
\[
                  \Phi_\gamma(s)
                    =\log\int_{S^3}e^{\gamma u\cdot s}\,dH(u)
                    =\ell(\gamma|s|)
\]
is smooth on all of \(\mathbb R^4\).
For \(r=|s|>0\), define
\begin{equation}
 k_\gamma(r)=\frac{\gamma a(\gamma r)}r,\qquad
 b_\gamma(r)=\gamma^2a'(\gamma r),\qquad
                  k_\gamma(0)=b_\gamma(0)=\gamma^2/4 .
 \label{eq:f15v76:radial}
\end{equation}
Here \(b_\gamma\) is only a radial Hessian coefficient,
not the packet function used in V74.
If \(u=s/r\), then
\[
 \nabla\Phi_\gamma(s)=k_\gamma(r)s,\qquad
 \operatorname{Hess}\Phi_\gamma(s)
       =k_\gamma(r)(I-uu^{\mathsf T})
                            +b_\gamma(r)uu^{\mathsf T}.
\]
Both coefficients are nonnegative, \(k_\gamma>0\)
for \(\gamma>0\), and \(k_\gamma(r)\) is
nonincreasing in \(r\). These are the covariance
and radial monotonicity facts proved in V69.
The value at zero follows from
\(\mathbb E_Hu_iu_j=\delta_{ij}/4\).

\subsection{A star calculation including all mixed derivatives}

Use an orthonormal ambient basis \(e_0,\ldots,e_3\)
of \(\mathbb R^4\). These are not the three
left-invariant frame fields used in V75.
On the whole retained product define
\begin{equation}
 {\cal W}_a(v)_x=\Pi_{v_x}e_a,\qquad
                    \Pi_v=I-vv^{\mathsf T},
                    \quad a=0,1,2,3 .
 \label{eq:f15v76:fields}
\end{equation}
The same ambient direction is projected at every
retained spin. Thus these fields include collective,
rather than only one-link, variations.

\begin{proposition}[A single-star collective Hessian bound]
\label{f15v76:star}
For \(m\) unit vectors \(v_1,\ldots,v_m\in S^3\),
let \(S=\sum_i v_i\), \(r=|S|\), and let
\(W_a=(\Pi_{v_i}e_a)_{i=1}^m\). Then
\begin{equation}
 \sum_{a=0}^3
       \operatorname{Hess}_{(S^3)^m}
                [-\Phi_\gamma(S)](W_a,W_a)
                     \le-k_\gamma(r)(m^2-r^2).
 \label{eq:f15v76:star}
\end{equation}
This holds also at \(r=0\), without assigning a
direction to \(S/|S|\).
\end{proposition}

\begin{proof}
Set
\[
             Q=\sum_i v_i v_i^{\mathsf T},\qquad M=mI-Q.
\]
For the product geodesic with initial tangent \(W_a\),
the first derivative of \(S\) is \(Me_a\).
Its second derivative is
\(-\sum_i|\Pi_{v_i}e_a|^2v_i\).
Summing these second derivatives over \(a\) gives
\(-3S\), because \(\operatorname{Tr}\Pi_{v_i}=3\).
Therefore the exact sum of Hessians is
\begin{equation}
 3k_\gamma(r)r^2
       -\operatorname{Tr}\{
                 \operatorname{Hess}\Phi_\gamma(S)M^2\}.
 \label{eq:f15v76:exacttrace}
\end{equation}
This retains all cross derivatives within the star.

Suppose \(r>0\), and split \(\mathbb R^4\) into
\(\mathbb Ru\oplus u^\perp\), \(u=S/r\).
Write
\[
       Q=\begin{pmatrix}q&c^{\mathsf T}\\c&R\end{pmatrix}.
\]
Since \(\operatorname{Tr}Q=m\),
\(\operatorname{Tr}R=m-q\). A direct block calculation gives
\begin{equation}
 \begin{split}
 T&:=\operatorname{Tr}\{(I-uu^{\mathsf T})M^2\}\\
  &=m^2+2mq+|c|^2+\operatorname{Tr}(R^2)
                                \ge m^2+2r^2 .
 \end{split}
 \label{eq:f15v76:tangenttrace}
\end{equation}
Indeed
\(q=\sum_i(u\cdot v_i)^2\ge
m^{-1}(\sum_i u\cdot v_i)^2=r^2/m\).
The remaining squared terms are nonnegative.
The radial contribution to
\eqref{eq:f15v76:exacttrace} is
\(-b_\gamma(r)|Mu|^2\le0\).
Thus that expression is at most
\(3k_\gamma r^2-k_\gamma(m^2+2r^2)\),
which is \eqref{eq:f15v76:star}.

At \(r=0\), the gradient of \(\Phi_\gamma\) is
zero and its Hessian is \((\gamma^2/4)I\).
The exact trace is then
\[
                   -(\gamma^2/4)\operatorname{Tr}(M^2).
\]
Here
\(\operatorname{Tr}(M^2)=2m^2+\operatorname{Tr}(Q^2)
\ge m^2\). This proves the same bound directly
at zero, with no singular radial coordinates.
\end{proof}

A useful uniform conversion of the right side is
\begin{equation}
 k_\gamma(r)(m^2-r^2)
      \ge2\gamma a(m\gamma/2)(m-r),
                         \qquad 0\le r\le m .
 \label{eq:f15v76:radialcomparison}
\end{equation}
For \(r\le m/2\), monotonicity gives
\(k_\gamma(r)\ge2\gamma a(m\gamma/2)/m\)
and \(m+r\ge m\).
For \(r\ge m/2\), monotonicity of \(a\) and
\((m+r)/r\ge2\) give the same conclusion.
The zero and aligned endpoints are included.

\subsection{All local minima of the finite-coupling retained potential}

\begin{theorem}[Exact collective trace and critical-point classification]
\label{f15v76:landscape}
For every \(\gamma>0\) and every retained configuration,
\begin{equation}
 \begin{split}
 {\cal T}_\gamma(v)
   &:=\sum_{a=0}^3
          \operatorname{Hess}V_\gamma({\cal W}_a,{\cal W}_a)\\
   &\le-\sum_{q\in O}k_\gamma(r_q)(36-r_q^2)
                  \le-2\gamma a(3\gamma)J(v).
 \end{split}
 \label{eq:f15v76:trace}
\end{equation}
The only local minima of \(V_\gamma\) are the
aligned configurations \(v_x=h\) for all \(x\in E\).
Every nonaligned critical point has a negative
Hessian direction. Quantitatively, in the retained
round-product metric,
\begin{equation}
 \lambda_{\min}(\operatorname{Hess}V_\gamma(v))
        \le-\frac{2\gamma a(3\gamma)}{3n}J(v).
 \label{eq:f15v76:negative}
\end{equation}
\end{theorem}

\begin{proof}
Each odd star has six retained neighbours.
Apply Proposition~\ref{f15v76:star} and
\eqref{eq:f15v76:radialcomparison} to its summand
in \eqref{eq:f15v76:potential}. Summing proves
\eqref{eq:f15v76:trace}, because the global fields
restrict to exactly the star fields.

Also
\[
                  \sum_{a=0}^3\|{\cal W}_a\|^2=3n.
\]
The defining lower bound by the least Hessian
eigenvalue, applied to these four tangent vectors,
therefore implies
\eqref{eq:f15v76:negative}.

One has \(J=0\) precisely when every odd-star sum
has norm six. Equality in the triangle inequality
for six unit vectors forces those vectors to agree.
The retained incidence graph obtained by sharing
odd stars is connected: an even-length path in
the connected bipartite lattice joins any two
retained vertices. Thus \(J=0\) precisely at
global alignment.

For \(\gamma>0\), \(\ell(\gamma r)\) is strictly
increasing for \(r>0\). It follows that aligned
configurations are global minima of \(V_\gamma\).
At every nonaligned point \(J>0\), and
\(a(3\gamma)>0\), so the Hessian has a negative
direction. A smooth local minimum must be a
critical point with nonnegative Hessian.
This excludes every other local minimum and
proves the asserted strict instability at
nonaligned critical points.
\end{proof}

A negative direction does not assert that
the critical point is nondegenerate or has only
one unstable direction. Maxima and degenerate
critical manifolds are allowed.

Nor is this negative direction merely a residual
constant endpoint-frame rotation. The potential
is invariant under the common \(SO(4)\) action,
equivalently the two common \(\SU(2)\) endpoint
actions. At a critical point its Hessian vanishes
when one argument is tangent to that orbit:
the gradient remains zero along the symmetry orbit.
Projecting a negative direction orthogonally off
the orbit preserves its quadratic Hessian value.
No smooth global quotient or free group action
is needed for this tangent-space assertion.

\subsection{A weighted estimate from the same collective fields}

Critical-point information by itself is not a
spectral inequality. In the present geometry the
same fields give a quantitative integral statement
without choosing local charts or a saddle cover.

Let \(\operatorname{div}_H\) be divergence with
respect to the retained product Haar volume, and
put \(U(v)=\sum_{x\in E}v_x\). The sphere identities give
\begin{equation}
 \operatorname{div}_H{\cal W}_a=-3U_a,\qquad
       \sum_{a=0}^3{\cal W}_a
                   (\operatorname{div}_H{\cal W}_a)=-9n .
 \label{eq:f15v76:divergence}
\end{equation}
The covariant acceleration at site \(x\) is
\[
     (\nabla_{{\cal W}_a}{\cal W}_a)_x
                         =-(v_x\cdot e_a)\Pi_{v_x}e_a.
\]
Summing over \(a\) makes this zero at every site.
Consequently
\begin{equation}
           \sum_a{\cal W}_a^2 V_\gamma
                                     ={\cal T}_\gamma .
 \label{eq:f15v76:drift}
\end{equation}
The distinction between second derivatives along
flows and geodesic Hessians has thus been accounted for.

\begin{proposition}[Weighted square with the full divergence retained]
\label{f15v76:weightedsquare}
Define the weighted divergence
\[
       \beta_a=\operatorname{div}_{\rho_{L,\gamma}}{\cal W}_a
             =\operatorname{div}_H{\cal W}_a
                                      -{\cal W}_a V_\gamma .
\]
For every smooth real \(f\),
\begin{equation}
 \mathbb E_\rho\sum_a|{\cal W}_a f|^2
  =\mathbb E_\rho\sum_a|{\cal W}_a f+\beta_a f|^2
                 +\mathbb E_\rho f^2(-9n-{\cal T}_\gamma).
 \label{eq:f15v76:square}
\end{equation}
\end{proposition}

\begin{proof}
For a vector field \(W\) with weighted divergence
\(\beta\), compact integration gives
\[
 \int W(f^2)\beta\,d\rho
          =-\int f^2\{W\beta+\beta^2\}\,d\rho .
\]
Expanding \(|Wf+\beta f|^2\) therefore gives
\[
       \int|Wf|^2\,d\rho
        =\int|Wf+\beta f|^2\,d\rho+\int f^2 W\beta\,d\rho .
\]
Apply this to each \({\cal W}_a\) and use
\eqref{eq:f15v76:divergence}--\eqref{eq:f15v76:drift}.
The common normalizer of \(\rho\) cancels within
this same integration identity.
\end{proof}

\begin{theorem}[Global frustration-weighted Dirichlet control]
\label{f15v76:weighted}
Under the exact law \(\rho_{L,\gamma}\), every smooth
real \(f\) satisfies
\begin{equation}
 \gamma^{-1}\mathbb E_\rho\sum_{x\in E}|\nabla_x f|^2
   \ge \frac{2a(3\gamma)}{n}\mathbb E_\rho[Jf^2]
                         -\frac9\gamma\mathbb E_\rho f^2 .
 \label{eq:f15v76:weighted}
\end{equation}
In particular the same lower bound holds for the
thermal completed form
\(\gamma^{-1}\mathbb E_\rho{\mathfrak m}_f\)
whenever \(f\) is bi-invariant.
\end{theorem}

\begin{proof}
The nonnegative square in
\eqref{eq:f15v76:square} and the trace bound give
\[
       \mathbb E_\rho\sum_a|{\cal W}_a f|^2
          \ge2\gamma a(3\gamma)\mathbb E_\rho[Jf^2]
                                      -9n\mathbb E_\rho f^2.
\]
Since every intrinsic \(\nabla_x f\) is tangent
to \(S^3\), its ambient representative obeys
\[
 {\cal W}_a f=e_a\cdot\sum_x\nabla_x f,\qquad
 \sum_a|{\cal W}_a f|^2
       =\left|\sum_x\nabla_x f\right|^2
                           \le n\sum_x|\nabla_x f|^2 .
\]
Divide by \(n\gamma\). The completed form dominates
the retained vertical form pointwise, by V73.
\end{proof}

This estimate involves the exact weighted quantity
\(\mathbb E[Jf^2]\), not a probability of a bad set.
It holds for arbitrary \(f\), without a support
condition, normalization of \(f\), or invariance
assumption for the retained vertical form.
For example \(f=1\) gives the consistent first-moment
consequence
\[
                 \mathbb E_\rho J
                         \le\frac{9n}{2\gamma a(3\gamma)} .
\]
The function-dependent estimate is stronger information
than this first moment alone.

\subsection{A whole high-frustration region is above the candidate band}

Let
\[
                 {\cal H}_L=\{v:J(v)\ge11n/2\}.
\]
The global weighted inequality immediately implies
\begin{equation}
 11a(3\gamma)\mathbb E_\rho[1_{{\cal H}_L}f^2]
    \le\gamma^{-1}\mathbb E_\rho\sum_x|\nabla_xf|^2
                               +\frac9\gamma\mathbb E_\rho f^2 .
 \label{eq:f15v76:highmass}
\end{equation}
There is no cutoff-transition term in this particular
estimate; it follows from a global weighted identity.

\begin{theorem}[A volume-independent supported high-frustration floor]
\label{f15v76:highfloor}
For \(\gamma\ge60\), a smooth \(f\) supported in
\({\cal H}_L\) obeys
\begin{equation}
 \gamma^{-1}\mathbb E_\rho\sum_x|\nabla_xf|^2
             \ge\frac{1291}{120}\mathbb E_\rho f^2,
 \qquad
            \frac{1291}{120}>\frac{32}{3}\ge4+2\kappa_L .
 \label{eq:f15v76:highfloor}
\end{equation}
For a nonconstant supported bi-invariant function,
its completed Rayleigh quotient, with variance in
the denominator, satisfies the same lower bound.
\end{theorem}

\begin{proof}
Apply \eqref{eq:f15v76:highmass} to a supported
function and use V70's bound
\(a(3\gamma)\ge1-1/(2\gamma)\). The coefficient is
at least
\[
                11-\frac{29}{2\gamma}
                          \ge11-\frac{29}{120}
                          =\frac{1291}{120}.
\]
The difference from \(32/3\) is \(11/120>0\).
V75 proves \(4+2\kappa_L\le32/3\).
Finally
\(\Var_\rho(f)\le\mathbb E_\rho f^2\), and
\({\mathfrak m}_f\) dominates the vertical integrand.
Subtracting the mean changes neither derivative
if one wishes to normalize the centered function.
\end{proof}

For a normalized family with squared mass tending
to one in \({\cal H}_L\), equation
\eqref{eq:f15v76:highmass} gives a limiting thermal
energy of at least eleven. It does not force the
mass in \({\cal H}_L\) to vanish for every
bounded-energy family; rather, it bounds that
mass in terms of its actual energy.

This region includes genuinely zero-star examples.
For \(L=4\), identify a circle in \(S^3\) with complex
phases, and set on retained sites
\[
                  v_x=\exp\{\,{\rm i}(\pi x_1+\pi x_2/2)\,\}.
\]
The nearest-neighbour symbol is
\(2(\cos\pi+\cos(\pi/2)+\cos0)=0\), so every
odd-star sum is zero and \(J=6n\).
The exact potential is smooth there and its gradient
vanishes because \(\nabla\Phi_\gamma(0)=0\).
It is a nonaligned critical configuration covered
by both the negative-curvature theorem and the
high-frustration estimate. There has been no
deletion of zero-star configurations from the law.

\subsection{Remaining low-energy coverage and the physical limitation}

The geometry is now substantially broader than the
single-flip calculation. For the exact common
flat marginal there are no additional local minima,
at any coupling, and the extensive high-frustration
region has a supported cost above the aligned
Gaussian candidate. Neither conclusion proves
convergence of the full compact spectral gap.

In particular the trace estimate has strength
proportional to \(J/n\). A nonaligned configuration
whose frustration is localized or spread sparsely
can have small \(J/n\) in a large volume.
The classification of its critical point as unstable
does not itself give a uniform localized energy
threshold or control transitions between regions.
V75 pays certain isolated antipodal patterns
individually; it does not cover every remaining
low-density, noncollinear or extended configuration.
The two estimates are complementary, not a complete
partition of the compact state space.

Finally the fields \({\cal W}_a\) used the common
ambient frame of the flat retained law. At nonflat
transverse connections, the transported arguments
in different stars do not in general share that
frame. The negative trace calculation cannot
simply be copied to that law. V75's native local
estimate does apply at its declared nonflat
exteriors, but it does not supply the missing
global comparison. The actual transverse-field
variance, physical transfer-time calibration,
and interacting continuum construction remain
open parts of the original Yang--Mills problem.

\subsection{Exact diagnostics and source preservation}

The diagnostic is
\begin{center}\small
\path{scripts/verify_ym_f15_v76_retained_landscape.py}.
\end{center}
It reruns V75 and its inherited checks.
New rational tests include 88 star trace identities,
176 common ambient first variations, 44 summed
geodesic accelerations, 120 transverse trace
inequalities, and 160 independent radial-Hessian
trace certificates. Twelve smooth zero-star
certificates check the radial-origin case directly.

There are also 204 product-field drift cancellations,
88 norm/divergence identities, 132 collective
gradient inequalities, 68 two-regime radial factors,
40 weighted-square identities with nonzero Haar
divergence, four high-frustration threshold checks,
and all 32 odd-star cancellations in the displayed
side-four common-layer configuration.
The algebraic radial-jet fixtures do not replace
the analytic Haar covariance and monotonicity facts.
The all-configuration theorems follow from the
proofs above, not from the finite fixtures.

The script has 6,611 bytes and SHA--256
\begin{center}\scriptsize\ttfamily
44AE9B6549C0007E03852F1F0B69280F0B5654884B500FBFB6DDFD0CD8D57D66.
\end{center}
The V75 predecessor has 1,702,904 bytes and SHA--256
\begin{center}\scriptsize\ttfamily
A5D1D273A0E2302506D932C3824FECD87CFB4E35A28278F2CFD7EC31F2135AE3.
\end{center}
Removing this final section and restoring the
title version recovers that source exactly.
V76 replaces V75 as the sole active main note.
Protected archives, companions and earlier
diagnostics remain unchanged. Only \path{.tex}
sources are kept in \path{latex_notes};
compilation products go to
\begin{center}\small
\path{artifacts/ym_f15_v76_texcheck/}.
\end{center}

\begin{firewall}
V76 classifies the local minima of the exact
flat-boundary retained potential, proves negative
curvature at every other critical point, and
establishes a global frustration-weighted
Dirichlet bound with a high-frustration supported
threshold. It does not prove the full compact
spectral lower bound, a global nonflat comparison,
an interacting continuum, or a positive physical
Yang--Mills mass gap. The original goal remains active.
\end{firewall}

\section{V77: spatial localization of collective control}
\label{f15v77:section}

V76 gives a weighted bound proportional to the total
frustration divided by the number of retained sites.
That does not detect a highly frustrated region of
fixed size inside a growing volume. This version
localizes the collective fields by deterministic
nonnegative spatial weights. The resulting estimate
uses only derivatives in their support. Its error
is the spatial Schur energy of the weights, not an
ambient-volume factor.

The new geometric input is a nonnegative row-debit
representation of the exact single-star Hessian.
A separate estimate at almost cancelled stars
absorbs their large radial coefficients. This is
essential: bounding every Hessian coefficient by
its value at zero would produce a spurious
order-\(\gamma^2\) spatial error. The final spatial
error is only order \(\gamma\) before thermal
normalization, uniformly over every spin configuration.

All statements use the same exact flat-boundary law
\(\rho_{L,\gamma}\) and potential
\eqref{eq:f15v76:potential}. No conditional or truncated
law is substituted. The V76 source and all 133
protected files were verified unchanged. The next
companion checkpoint remains V80. Earlier isolated-star
probability bounds and V76's unweighted collective
estimate are not being reproved as new results.

\subsection{A positive row representation for a star}

For this subsection let \(v_1,\ldots,v_m\in S^3\),
\(S=\sum_i v_i\), \(r=|S|\), and \(P_i=I-v_iv_i^{\mathsf T}\).
Write \(k=k_\gamma(r)>0\), \(b=b_\gamma(r)\), and
\(A=\operatorname{Hess}\Phi_\gamma(S)\), with the smooth
zero-radius convention of V76. For \(r>0\),
\[
              A=k(I-uu^{\mathsf T})+buu^{\mathsf T},
                         \qquad u=S/r .
\]
V69's covariance and radial monotonicity facts give
\begin{equation}
                         0\le b\le k.
 \label{eq:f15v77:radialorder}
\end{equation}
Indeed \(a'\ge0\), and differentiating the decreasing
function \(a(h)/h\) gives \(ha'(h)\le a(h)\).
At \(r=0\), \(b=k=\gamma^2/4\).

Define real symmetric star coefficients and their debits by
\begin{equation}
 \begin{split}
 C_{ij}&=\operatorname{Tr}(AP_iP_j),\qquad t_i=S\cdot v_i,\\
 d_i&=\sum_jC_{ij}-3kt_i .
 \end{split}
 \label{eq:f15v77:rows}
\end{equation}

\begin{proposition}[Entry positivity and nonnegative row debits]
\label{f15v77:positive}
For every such star,
\begin{equation}
    k\le C_{ij}\le3k,\qquad C_{ii}\ge2k,\qquad d_i\ge0.
 \label{eq:f15v77:positive}
\end{equation}
The matrix \(C\) is positive semidefinite as well.
\end{proposition}

\begin{proof}
The coefficients are the Frobenius Gram matrix of
\(A^{1/2}P_i\). The equality with the stated trace
uses cyclicity and transpose symmetry of the real
symmetric matrices. Thus \(C\) is symmetric positive
semidefinite. Since \(0\le A\le kI\),
\(C_{ii}=\operatorname{Tr}(AP_i)\le3k\), and Gram
Cauchy--Schwarz gives \(|C_{ij}|\le3k\).

For \(r>0\), put \(c_{ij}=v_i\cdot v_j\). Then
\[
 C_{ij}=k(2+c_{ij}^2)
                    -(k-b)u^{\mathsf T}P_iP_j u .
\]
The last scalar is at most one, by Cauchy--Schwarz
and the contraction property of orthogonal projections.
Consequently \(C_{ij}\ge k(1+c_{ij}^2)\ge k\).
For \(i=j\), this also gives \(C_{ii}\ge2k\).
At \(r=0\) these conclusions follow directly from
\(A=kI\).

It remains to prove the stronger row assertion,
which does not follow just from entry positivity.
At fixed spins the debit is linear in \(b/k\).
By \eqref{eq:f15v77:radialorder}, it suffices to
check \(A=kI\) and \(A=k(I-uu^{\mathsf T})\).
In the isotropic case,
\begin{equation}
       d_i/k=\sum_j(1-c_{ij})(2-c_{ij})\ge0.
 \label{eq:f15v77:isotropic}
\end{equation}
This case also handles \(r=0\).

For the tangential case let
\(Q=\sum_jv_jv_j^{\mathsf T}\) and \(\alpha=u\cdot v_i\).
Direct expansion of the row sum gives
\begin{equation}
 \frac{d_i}{k}
  =m(1+\alpha^2)-3r\alpha+
       u^{\mathsf T}Qu+v_i^{\mathsf T}Qv_i
                              -\alpha v_i^{\mathsf T}Qu .
 \label{eq:f15v77:tangentrow}
\end{equation}
The quadratic form \(x^2+y^2-\alpha xy\) is
nonnegative for every \(|\alpha|\le1\). Apply its
mean-square inequality to
\(x_j=u\cdot v_j\), \(y_j=v_i\cdot v_j\).
Their sums are \(r\) and \(r\alpha\), respectively.
It follows that
\[
 u^{\mathsf T}Qu+v_i^{\mathsf T}Qv_i
                    -\alpha v_i^{\mathsf T}Qu\ge r^2/m .
\]
Set \(z=r/m\in[0,1]\). The resulting lower bound
for \(d_i/(km)\) is \(1+\alpha^2-3z\alpha+z^2\).
Its minimum over \(\alpha\in[-1,1]\) is
\[
 \begin{cases}
  1-5z^2/4\ge4/9,&0\le z\le2/3,\\
  (1-z)(2-z)\ge0,&2/3\le z\le1.
 \end{cases}
\]
This proves the tangential endpoint and hence every
convex combination of the two endpoints.
\end{proof}

\subsection{Weighted traces and the small-radius region}

For fixed real numbers \(\phi_i\), define
\(W_a^\phi=(\phi_iP_ie_a)_i\) and
\(M_\phi=\sum_i\phi_iP_i\). Denote the single-star
weighted trace by
\[
 T_\phi=\sum_{a=0}^3
    \operatorname{Hess}_{(S^3)^m}[-\Phi_\gamma(S)]
                              (W_a^\phi,W_a^\phi).
\]
All pair sums below are ordered sums over \(i,j\).

\begin{proposition}[Exact weighted row decomposition]
\label{f15v77:decomposition}
For arbitrary real deterministic weights,
\begin{equation}
 \begin{split}
 T_\phi
 &=3kS\cdot\sum_i\phi_i^2v_i-\operatorname{Tr}(AM_\phi^2)\\
 &=-\sum_i d_i\phi_i^2
                +\frac12\sum_{i,j}C_{ij}(\phi_i-\phi_j)^2\\
 &\le \bigl(\min_i\phi_i^2\bigr)T_{\boldsymbol 1}
                +\frac12\sum_{i,j}C_{ij}(\phi_i-\phi_j)^2 .
 \end{split}
 \label{eq:f15v77:decomposition}
\end{equation}
\end{proposition}

\begin{proof}
Along the product geodesic with initial tangent
\(W_a^\phi\), the first variation of \(S\) is
\(M_\phi e_a\). The sum over \(a\) of its second
variations is \(-3\sum_i\phi_i^2v_i\).
The ambient chain rule gives the first line.
Expanding the pair differences using the symmetry
of \(C\) gives the second line, with precisely the
row debit \eqref{eq:f15v77:rows}.
For constant weights, \(T_{\boldsymbol1}=-\sum_i d_i\).
The last line now follows from \(d_i\ge0\).
\end{proof}

\begin{proposition}[Uniform six-star spatial error]
\label{f15v77:starbound}
For \(m=6\), \(\gamma\ge12\), and \(\phi_i\ge0\), put
\(s=\min_i\phi_i^2\). Then, including \(S=0\),
\begin{equation}
 T_\phi\le-2\gamma a(3\gamma)(6-r)s
                    +6\gamma\sum_{i,j}(\phi_i-\phi_j)^2 .
 \label{eq:f15v77:starbound}
\end{equation}
\end{proposition}

\begin{proof}
If \(r\ge1/4\), then \(k=\gamma a(\gamma r)/r\le4\gamma\)
because \(a\le1\). Thus \(C_{ij}\le12\gamma\).
Apply \eqref{eq:f15v77:decomposition} and the
unweighted V76 bound
\(T_{\boldsymbol1}\le-2\gamma a(3\gamma)(6-r)\).

If \(r\le1/4\), nonnegativity of the weights and
of every \(C_{ij}\), together with \(C_{ii}\ge2k\),
gives
\[
 \operatorname{Tr}(AM_\phi^2)
        =\sum_{i,j}C_{ij}\phi_i\phi_j\ge2k\sum_i\phi_i^2.
\]
Also \(t_i\le r\). The exact trace therefore satisfies
\[
                   T_\phi\le-k(2-3r)\sum_i\phi_i^2
                           \le-\tfrac54 k\sum_i\phi_i^2 .
\]
Monotonicity of \(k\) and V70's
\(a(h)\ge1-3/(2h)\) imply
\[
           k(r)\ge k(1/4)=4\gamma a(\gamma/4)
                          \ge4\gamma-24\ge2\gamma .
\]
Consequently \(T_\phi\le-15\gamma s\).
This is at most \(-2\gamma a(3\gamma)(6-r)s\),
whose absolute coefficient is at most \(12\gamma\).
No spatial error is needed in this region.
The proof at zero uses \(A=kI\), so no direction
\(S/|S|\) is assigned there.
\end{proof}

The small-radius argument is the reason for the
nonnegative-mask hypothesis. A configuration-dependent
choice of weights is not covered either: it would
add derivatives to the vector-field identities below.

\subsection{A local inequality under the unchanged law}

Let \(\phi:E\longrightarrow[0,\infty)\) be deterministic,
not identically zero. Write
\begin{equation}
 \begin{split}
 N_\phi&=\sum_{x\in E}\phi_x^2,\qquad
 D_\phi=\{x:\phi_x>0\},\\
 s_q&=\min_{x\sim q}\phi_x^2,\qquad
 J_\phi(v)=\sum_{q\in O}(6-r_q(v))s_q .
 \end{split}
 \label{eq:f15v77:mask}
\end{equation}
Retain V72's incidence matrix \(B\) and its exact
Schur operator
\[
                         K=6I-\tfrac16BB^{\mathsf T}
\]
on all retained sites, including its constant kernel.
Its spatial quadratic form satisfies
\begin{equation}
 \sum_{q\in O}\sum_{x,y\sim q}(\phi_x-\phi_y)^2
                              =12\phi^{\mathsf T}K\phi .
 \label{eq:f15v77:schur}
\end{equation}
Indeed each retained site belongs to six stars,
and expansion gives \(72\sum_x\phi_x^2
-2\sum_q(\sum_{x\sim q}\phi_x)^2\).

On the full product define
\[
 {\cal W}_a^\phi(v)_x=\phi_x\Pi_{v_x}e_a,\qquad
 {\cal T}_\phi=\sum_a
           \operatorname{Hess}V_\gamma
                          ({\cal W}_a^\phi,{\cal W}_a^\phi).
\]
Summing Proposition~\ref{f15v77:starbound} and using
\eqref{eq:f15v77:schur} yields
\begin{equation}
 {\cal T}_\phi
       \le-2\gamma a(3\gamma)J_\phi
                                +72\gamma\phi^{\mathsf T}K\phi .
 \label{eq:f15v77:globaltrace}
\end{equation}

\begin{theorem}[Spatially localized frustration control]
\label{f15v77:localized}
For \(\gamma\ge12\) and every smooth real \(f\) on
the retained product,
\begin{equation}
 \begin{split}
 \gamma^{-1}\mathbb E_\rho
                 \sum_{x\in D_\phi}|\nabla_x f|^2
 &\ge\frac{2a(3\gamma)}{N_\phi}
                         \mathbb E_\rho[J_\phi f^2]\\
 &\quad-\left(
       72\frac{\phi^{\mathsf T}K\phi}{N_\phi}
                        +\frac9\gamma\right)\mathbb E_\rho f^2 .
 \end{split}
 \label{eq:f15v77:localized}
\end{equation}
The function may depend on sites outside \(D_\phi\).
For bi-invariant \(f\), the same right side is also
a lower bound for its thermal completed energy.
\end{theorem}

\begin{proof}
The weighted sphere fields satisfy
\[
 \operatorname{div}_H{\cal W}_a^\phi
                     =-3\sum_x\phi_xv_{x,a},\qquad
 \sum_a{\cal W}_a^\phi
       (\operatorname{div}_H{\cal W}_a^\phi)=-9N_\phi .
\]
At each site, their summed covariant acceleration is
\[
 -\phi_x^2\sum_a v_{x,a}\Pi_{v_x}e_a=0 .
\]
In particular
\(\sum_a({\cal W}_a^\phi)^2V_\gamma={\cal T}_\phi\).
With
\(\beta_a^\phi=\operatorname{div}_H{\cal W}_a^\phi
-{\cal W}_a^\phi V_\gamma\), the weighted square
identity proved in Proposition~\ref{f15v76:weightedsquare}
therefore reads
\[
 \mathbb E_\rho\sum_a|{\cal W}_a^\phi f|^2
 =\mathbb E_\rho\sum_a
       |{\cal W}_a^\phi f+\beta_a^\phi f|^2
                  +\mathbb E_\rho[f^2(-9N_\phi-{\cal T}_\phi)].
\]
Tangency of the intrinsic gradients and Cauchy--Schwarz give
\[
 \sum_a|{\cal W}_a^\phi f|^2
       =\left|\sum_x\phi_x\nabla_x f\right|^2
                    \le N_\phi\sum_{x\in D_\phi}|\nabla_x f|^2 .
\]
Drop the nonnegative square, insert
\eqref{eq:f15v77:globaltrace}, and divide by
\(\gamma N_\phi\). V73's pointwise domination
of vertical energy by the completed form proves
the final assertion.
\end{proof}

The choice \(\phi\equiv1\) has \(N_\phi=n\),
\(J_\phi=J\), and \(\phi^{\mathsf T}K\phi=0\), so
it recovers V76's inequality in the stated range
of \(\gamma\). Compactly supported choices remove
the ambient-volume dilution for a local region.
There is no estimate of a bad-event probability
in this proof.

\subsection{An explicit local high-frustration threshold}

Put
\[
 \eta_\phi=\frac{\phi^{\mathsf T}K\phi}{N_\phi},
 \qquad {\cal H}_\phi=\{v:J_\phi(v)\ge11N_\phi/2\}.
\]
The local theorem gives, for every \(f\),
\begin{equation}
 11a(3\gamma)\mathbb E_\rho[1_{{\cal H}_\phi}f^2]
 \le\gamma^{-1}\mathbb E_\rho
                   \sum_{x\in D_\phi}|\nabla_x f|^2
          +(72\eta_\phi+9/\gamma)\mathbb E_\rho f^2 .
 \label{eq:f15v77:windowmass}
\end{equation}

\begin{proposition}[A supported local floor above the candidate band]
\label{f15v77:floor}
If \(\eta_\phi\le1/432\), \(\gamma\ge120\), and
\(f\) is smooth and supported in \({\cal H}_\phi\), then
\begin{equation}
 \gamma^{-1}\mathbb E_\rho
           \sum_{x\in D_\phi}|\nabla_x f|^2
                   \ge\frac{2571}{240}\mathbb E_\rho f^2,
 \qquad \frac{2571}{240}>\frac{32}{3}\ge4+2\kappa_L .
 \label{eq:f15v77:floor}
\end{equation}
The completed variance Rayleigh quotient of any
nonconstant supported bi-invariant \(f\) has the
same lower bound.
\end{proposition}

\begin{proof}
The coefficient supplied by \eqref{eq:f15v77:windowmass}
is at least
\[
 11a(3\gamma)-72\eta_\phi-9/\gamma
       \ge11-\frac16-\frac{29}{2\gamma}
       \ge\frac{2571}{240}.
\]
Here \(a(3\gamma)\ge1-1/(2\gamma)\).
The difference from \(32/3\) is \(11/240\).
The comparison with \(4+2\kappa_L\) was established
in V75. Finally \(\operatorname{Var}_\rho f
\le\mathbb E_\rho f^2\), and completion dominates
the local vertical energy. Centering a supported
function does not preserve its support, but it
does preserve its derivatives; the supported
inequality is applied before that centering.
\end{proof}

\subsection{Fixed-size masks and nonempty local windows}

For completeness, there are masks satisfying the
corollary whose support is independent of the ambient
volume and whose high-frustration window has nonempty
interior. The following conservative size is explicit,
not claimed optimal.

Let \(R\) be even, \(L\ge R+2\) even, and embed
\(\{1,\ldots,R\}^3\) in the periodic cubic lattice.
On both parity classes set
\[
 \psi(x)=
 \begin{cases}
 \displaystyle\prod_{j=1}^3
          \sin\!\left(\frac{\pi x_j}{R+1}\right),
                         &x\in\{1,\ldots,R\}^3,\\
 0,                       &\text{otherwise},
 \end{cases}
 \qquad \phi=\psi|_E .
\]
The full unnormalized nearest-neighbour Dirichlet
energy of the zero extension is
\begin{equation}
 {\cal E}_{\rm full}(\psi)
       :=\sum_{q\in O}\sum_{x\sim q}(\psi(x)-\psi(q))^2
        =\lambda_R\sum_z\psi(z)^2,\quad
 \lambda_R=6\left(1-\cos\frac{\pi}{R+1}\right).
 \label{eq:f15v77:sinemask}
\end{equation}
To see this, the one-dimensional sine recurrence
gives the Dirichlet eigenvalue
\(2(1-\cos(\pi/(R+1)))\) in each coordinate.
Discrete summation by parts then gives the displayed
identity, including the edges crossing the cube boundary.

Reflection \(j\mapsto R+1-j\) preserves the squared
sine and reverses parity. Its alternating sum in
one dimension is zero. Factoring the three-dimensional
alternating sum therefore proves
\[
                   \sum_{x\in E}\psi(x)^2
                  =\sum_{q\in O}\psi(q)^2=N_\phi .
\]
Minimizing the full energy over the odd completion
gives precisely \(\phi^{\mathsf T}K\phi\): complete
the square separately at each odd vertex around
\(6^{-1}\sum_{x\sim q}\phi_x\).
Consequently
\begin{equation}
           \eta_\phi\le2\lambda_R
                      \le\frac{6\pi^2}{(R+1)^2}.
 \label{eq:f15v77:maskcost}
\end{equation}

One must also check that the minimum over each
star has not destroyed the weight in \(J_\phi\).
For each \(q\),
\[
 \psi(q)^2-\min_{x\sim q}\phi_x^2
       \le\sum_{x\sim q}|\psi(q)^2-\phi_x^2|.
\]
All weights are nonnegative. Summing and applying
Cauchy--Schwarz along the unoriented lattice edges gives
\[
 \begin{split}
 \sum_{\{x,q\}}|\psi(x)^2-\psi(q)^2|
 &\le {\cal E}_{\rm full}(\psi)^{1/2}
              \left(\sum_{\{x,q\}}(\psi(x)+\psi(q))^2\right)^{1/2}\\
 &\le (2\lambda_R N_\phi)^{1/2}(24N_\phi)^{1/2}.
 \end{split}
\]
The last inequality uses degree six and
\((a+b)^2\le2(a^2+b^2)\). Thus
\begin{equation}
              \sum_qs_q\ge
                        (1-\sqrt{48\lambda_R})N_\phi .
 \label{eq:f15v77:maskmass}
\end{equation}

Take \(R=512\). The elementary bound \(\pi^2<10\) gives
\[
 2\lambda_R<\frac{60}{513^2}<\frac1{432},\qquad
 48\lambda_R<\frac{1440}{513^2}<\frac1{144}.
\]
Hence \(\eta_\phi<1/432\) and
\(\sum_qs_q>(11/12)N_\phi\).
For \(L\) divisible by four and \(L\ge516\), use
the literal periodic retained configuration
\[
                     v_x=\exp\{\,{\rm i}
                                   (\pi x_1+\pi x_2/2)\,\}
\]
in a fixed circle of \(S^3\).
As in V76, every odd-star sum is zero. It follows that
\[
                      J_\phi=6\sum_qs_q>11N_\phi/2 .
\]
Continuity gives an open neighbourhood inside
\({\cal H}_\phi\); it has positive measure under
the smooth strictly positive finite-volume density.
This proves nonvacuity for arbitrarily large ambient
volumes with one fixed support size. The displayed
periodic zero-star configuration is a witness for
nonvacuity, not a claim that every localized
frustration pattern has this form.

\subsection{What localization has and has not closed}

The new estimate treats a prescribed spatial region
under the original global probability law and retains
only gradients supported in that region. Highly
frustrated fixed windows no longer become invisible
merely because \(L\) increases. Zero-star configurations
are retained, and their large Hessian coefficients
help the estimate instead of worsening its boundary cost.

There are still important qualifications. The
high-frustration threshold concerns the weighted
quantity \(J_\phi/N_\phi\); it does not exclude every
nonaligned configuration. Slowly varying or sparse
frustration can lie below every threshold used here.
A configuration-adaptive selection of \(\phi\) would
require new derivative terms. Nor does an estimate
for each fixed window automatically give the same
constant for a union of many windows or for a
configuration-space partition: summation multiplicities
or cutoff costs must be paid. The weighted mass
bound \eqref{eq:f15v77:windowmass} controls mass in
terms of energy, not its vanishing for every
bounded-energy family.

All of this remains a result for the exact
flat-boundary retained marginal. Extending the common
ambient-frame argument to transported nonflat stars,
controlling actual transverse-field variance, and
calibrating physical transfer time are still separate
tasks. The next useful upstream step is a quantitative
treatment of the low-frustration sector or of the
frame mismatch, not an assertion that a local
supported threshold is already a physical gap.

\subsection{Exact diagnostics and source preservation}

The new diagnostic is
\begin{center}\small
\path{scripts/verify_ym_f15_v77_localized_collective_fields.py}.
\end{center}
It reruns V76 and its inherited chain. Its additional
rational fixtures check 128 star Gram bounds and
row-positivity cases, 46 isotropic and 41 tangential
row formulas, and 512 independent weighted geodesic
trace decompositions. There are 512 weighted
divergence, drift and local frame checks; 24
small-radius nonnegative-mask tests; 625 scalar
row-positivity factorizations; and four radial and
Schur-cost budget checks. Eight finite-torus fixtures
check the ordered-pair identity, exact Schur
completion and boundary-mass inequalities.
Four reflected-array parity tests, three explicit
window budgets, and three supported thresholds are
also checked.

The finite radial-jet fixtures are not replacements
for the analytic Haar monotonicity and covariance
bounds. The sine recurrence, the all-configuration
row lemma, and the same-law integration theorem are
proved above, not inferred from finite sampling.

The script has 8,729 bytes and SHA--256
\begin{center}\scriptsize\ttfamily
76FC0EF456347CFBAFD3B5FA9E5C6DB4A9BA3FC0D33DA2A5507AEFCF2B0F8220.
\end{center}
The V76 predecessor has 1,722,915 bytes and SHA--256
\begin{center}\scriptsize\ttfamily
21D1122ADEBA5935C546CD39B3437BE060FF783DDA4042516B4AA107594A657C.
\end{center}
Removing this final section and restoring the title
version recovers that source byte for byte.
V77 replaces V76 as the sole active main note.
Protected archives, companion sources and earlier
diagnostics remain unchanged.
Only \path{.tex} sources are kept in \path{latex_notes};
all compilation products are directed to
\begin{center}\small
\path{artifacts/ym_f15_v77_texcheck/}.
\end{center}

\begin{firewall}
V77 proves a nonnegative row-debit identity and an
explicit spatially localized weighted inequality
for the exact flat retained law, with a fixed-window
supported threshold above the existing Gaussian
candidate band. It does not prove full compact
spectral convergence, an all-exterior nonflat
comparison, an interacting continuum, or a positive
physical Yang--Mills mass gap.
The original goal remains active.
\end{firewall}

\section{V78: paper audit and native transfer channels}
\label{f15v78:section}

A new paper and suggested transfer programme arrived while
the nonflat extension of V77 was being investigated.
This version changes priority to the physical-form
interface and a complete-channel certificate for the
\emph{actual Wilson transfer operator}. The unfinished
nonflat calculation is retained as a candidate, not used
as a proved physical input here. In particular, this
version does not merely add another compact packet.

The supplied 29-page paper is Shoshauna Gauvin,
\href{https://arxiv.org/abs/2503.15539v3}
{Pinched Multi Affine Geometry and Confinement:
Describing the Yang--Mills Mass Gap}, v3, 14 September 2026.
Its operator statements were compared with the relevant
proofs in the
\href{https://arxiv.org/src/2503.15539v3/anc/PinchedMultiAffineGeometry_Supplement.pdf}
{67-page supplement}, especially A.2, C.10, C.14,
C.25 and C.32. This is a transfer audit of the relevant
mechanisms, not an independent certification of every
claim in the paper.

The main article explicitly leaves the interacting
four-dimensional continuum construction and its uniform
physical gap open. Its fixed-cutoff strong-coupling
\(\SU(3)\) constants are not imported into our
\(\SU(2)\) Wilson trajectory. Its complete-channel
argument is useful, but its hypotheses cannot be
filled by the existing sampling-energy estimates
without checking the energy form and probability law.

There is also a nonduplication point: archive f7 V90
already proves the reversible ground-state transform
of a positive transfer kernel. V57 constructs the
literal Wilson transfer and proves physical injectivity.
V63 identifies its logarithmic lattice gap, and V65
constructs particular positive-time probes. Those
results are retained rather than reproved as new.
The new work below is the explicit nonlocal localization
identity, complete native channel construction, and
full-space error certificate for that existing transfer.

The V77 predecessor and all 134 previously protected
files were verified unchanged. The next compulsory
companion checkpoint remains V80. The paper and
pasted suggestions are source material, not instructions
that override the declared hypotheses.

\subsection{What transfers, and what does not}

The useful mechanism in Gauvin's Theorem 6.11 is an
isometric cover of the \emph{complete centered space},
channel coercivity outside finite low-mode spaces,
and an upper bound on the global leakage of those
spaces. Its elementary proof does not require a
particular gauge group. Appendix C.14 also keeps the
vacuum projection in the leakage Gram matrix.
We adapt this mechanism below, with attribution.

Other transfers remain conditional. Theorem 2.2's
gradient identity is for a supplied Hamiltonian
\(H=-\sum_e a_e\Delta_e+V\) and its own ground state.
Theorem 6.19 compares two such Hamiltonians on the
same carrier with the \emph{same electric symbol}
and a controlled gradient of their vacuum ratio.
Our V67 retained-density score is not automatically
that vacuum-ratio score. Theorem 6.7's shell route
needs exact relaxed forms and full-vector vertical
bounds. Theorem 6.23 additionally needs a total
tight family and compatible resolvents; one surviving
probe from V65 is not a total family.

These distinctions favor checking the physical form
before trying to declare a complete low-energy cover.
They do not rule out the proposed operator programme.
They identify which of its premises still need proof.

\begin{proposition}[The full-vector vertical premise is essential]
\label{f15v78:vertical}
A positive lower bound on the discarded restriction,
together with a positive relaxed retained form, does
not imply Gauvin's harmonic-mean gap bound unless
the stated full-vector vertical premise is supplied.
\end{proposition}

\begin{proof}
On \(\mathbb R^2\), retain the first coordinate and take
\[
                 K=\begin{pmatrix}5&2\\2&1\end{pmatrix},
       \qquad q(x,y)=x^2+(y+2x)^2 .
\]
The discarded restriction is \(q(0,y)=y^2\), and
\(\inf_yq(x,y)=x^2\). Thus the restricted discarded
constant and relaxed retained constant both equal one.
Nevertheless \(q(1,-2)/\|(1,-2)\|^2=1/5<1/2\).
The full-vector inequality \(q(x,y)\ge |y|^2\) fails.
Gauvin's actual premise is
\(q(g)\ge\delta\|Qg\|^2\) for \emph{every} form vector,
not just \(g\in Q\mathcal H\). With that premise and
a relaxed bound \(m\), the valid argument is
\[
 \|g\|^2=\|Pg\|^2+\|Qg\|^2
          \le(m^{-1}+\delta^{-1})q(g).
\]
It bounds two orthogonal norm components using the
same energy; it does not add two energy lower bounds.
\end{proof}

\subsection{The physical transfer form already available}

Fix a finite spatial regulator and coupling
\(\gamma>0\). Let \(\mathsf T\) be V57's literal
Wilson transfer operator on the gauge-invariant
spatial slice space, with normalized positive ground
vector \(\Omega\), top eigenvalue \(\lambda_0>0\),
and Haar reference \(dm\). Write its kernel as
\(\mathsf t(U,V)\), and define
\begin{equation}
 \begin{split}
 d\mu(U)&=\Omega(U)^2dm(U),\\
 (\mathcal P f)(U)
   &=\frac1{\lambda_0\Omega(U)}
         \int\mathsf t(U,V)\Omega(V)f(V)\,dm(V),\\
 \mathcal J(dU,dV)
   &=\lambda_0^{-1}\Omega(U)\mathsf t(U,V)\Omega(V)
                                                   dm(U)dm(V).
 \end{split}
 \label{eq:f15v78:ground}
\end{equation}
Here \(\mathcal J\) is a symmetric probability with
both marginals \(\mu\), and
\(\mathcal J(dU,dV)=\mu(dU)\mathcal P(U,dV)\).
These are precisely the already-proved f7 V90
normalization formulas, now applied to V57's kernel.

Let
\[
 \mathcal H_\mu=L^2(\mu)^{\rm gauge},\qquad
 Qf=f-\mu(f),\qquad \mathcal H_\mu^0=Q\mathcal H_\mu .
\]
On this space \(\mathcal P\) is a positive self-adjoint
contraction. Its vacuum is \(1\). At a fixed regulator
it is compact, its top eigenvalue is simple, and it
has no kernel on the physical space. Multiplication
by \(\Omega\) conjugates it to \(\mathsf T/\lambda_0\).

Define the bounded one-step form
\begin{equation}
 \begin{split}
 d[f]&=\langle f,(I-\mathcal P)f\rangle_\mu\\
     &=\frac12\int |f(U)-f(V)|^2\,\mathcal J(dU,dV),
                    \qquad f\in\mathcal H_\mu .
 \end{split}
 \label{eq:f15v78:jump}
\end{equation}
The equality follows by expanding the square and
using the two marginals and symmetry. In particular,
\(0\le d[f]\le\|f\|_\mu^2\).

If \(a_t>0\) is a \emph{separately supplied} physical
time step, the exact physical logarithmic form is
\begin{equation}
 \begin{split}
 \mathcal K&=-a_t^{-1}\log\mathcal P,\\
 q_{\log}[f]&=a_t^{-1}\int_{(0,1]}(-\log s)
                                \,d\langle f,E_{\mathcal P}(s)f\rangle .
 \end{split}
 \label{eq:f15v78:logform}
\end{equation}
Its domain consists of vectors for which this
nonnegative integral is finite. V57's physical
injectivity makes that domain dense.
No local gradient representation is asserted.

\begin{proposition}[A one-step certificate has an exact physical conversion]
\label{f15v78:conversion}
Put \(\rho=\|\mathcal P|_{\mathcal H_\mu^0}\|\).
Then the exact infima are
\begin{equation}
 \inf_{\substack{f\in\mathcal H_\mu^0\\\|f\|=1}}d[f]=1-\rho,
 \qquad
 \inf_{\substack{f\in D(q_{\log})\cap\mathcal H_\mu^0\\\|f\|=1}}
                   q_{\log}[f]=-\frac{\log\rho}{a_t}.
 \label{eq:f15v78:conversion}
\end{equation}
Consequently, a proved bound \(d[f]\ge\theta\|f\|^2\)
on the complete centered space, with \(0<\theta<1\),
gives
\[
             \mathcal K|_{\mathcal H_\mu^0}
                     \ge-\frac{\log(1-\theta)}{a_t}I .
\]
\end{proposition}

\begin{proof}
Both infima follow from spectral calculus and
monotonicity on \((0,1]\). For the first one,
\(I-\mathcal P\ge\theta I\) is equivalent to
\(\mathcal P\le(1-\theta)I\) on the centered space.
Apply the decreasing function \(-\log s\) through
the same spectral measure. This is scalar functional
calculus for one operator, not operator monotonicity
between unrelated kernels.
\end{proof}

Thus an exact physical bridge is available for
\eqref{eq:f15v78:jump}; it is not currently established
for V73--V77's completed retained gradient form.
The latter also lives on a different retained law.

\subsection{Why a diffusion identity is not automatic}

Unitary ground-state conjugation alone does not
turn the logarithm of a transfer into a differential
electric Hamiltonian. Even its Markov property at
fractional times is an additional issue.

\begin{proposition}[A positive transfer need not have a Markov logarithm]
\label{f15v78:nonembedding}
There is a strictly positive symmetric stochastic
matrix with positive spectrum whose logarithm is
not a continuous-time Markov generator.
\end{proposition}

\begin{proof}
On three points let \(\Pi\) be the matrix with
every entry \(1/3\), and let \(E_1,E_3\) be the
orthogonal projections onto \((1,0,-1)\) and
\((1,-2,1)\), respectively. Set
\[
            P=\Pi+\frac{27}{40}E_1+\frac9{40}E_3
              =\begin{pmatrix}
                 17/24&31/120&1/30\\
                 31/120&29/60&31/120\\
                 1/30&31/120&17/24
                \end{pmatrix}.
\]
Every entry is positive, every row sums to one,
and the eigenvalues are \(1,27/40,9/40\).
But
\[
       (\log P)_{13}
           =-\tfrac12\log(27/40)+\tfrac16\log(9/40)
           =\tfrac16\log(1600/2187)<0 .
\]
For a finite-state Markov generator all off-diagonal
entries must be nonnegative. Equivalently,
\((e^{t\log P})_{13}=t(\log P)_{13}+O(t^2)<0\)
for sufficiently small positive \(t\).
\end{proof}

This example proves no failure for the actual Wilson
logarithm; it proves that positivity and reversibility
alone do not justify importing the diffusion identity.
Likewise no global inequality \(d[f]\ge c\,e[f]\),
\(c>0\), can hold if \(e\) is an unbounded local
elliptic gradient form on the same infinite-dimensional
physical space: \(d[f]\le\|f\|^2\), whereas normalized
high electric modes have \(e[f]\to\infty\).
An appropriately proved low-energy comparison could
still be useful. No such comparison is assumed below.

\subsection{Exact nonlocal localization on the true vacuum law}

Choose finitely many real nonnegative bounded
gauge-invariant functions \(\chi_1,\ldots,\chi_N\)
with
\[
                         \sum_i\chi_i(U)^2=1
                              \quad\mu\text{-almost everywhere}.
\]
Smooth partitions may be used, but no differentiation
of the partition is needed for the bounded form \(d\).
Define a bounded self-adjoint correction operator
\begin{equation}
 \begin{split}
 (R_\chi f)(U)
    &=\int\left(1-\sum_i\chi_i(U)\chi_i(V)\right)
                              f(V)\mathcal P(U,dV),\\
 W_\chi(U)
    &=\frac12\int\sum_i|\chi_i(U)-\chi_i(V)|^2
                                                   \mathcal P(U,dV).
 \end{split}
 \label{eq:f15v78:correction}
\end{equation}
The row sum of the nonnegative correction kernel
is \(W_\chi\), and \(0\le W_\chi\le1\).

\begin{theorem}[Native one-step localization identity]
\label{f15v78:localization}
For every complex physical \(f\in L^2(\mu)\),
\begin{equation}
 \begin{split}
 \sum_i d[\chi_i f]
       &=d[f]+\langle f,R_\chi f\rangle_\mu,\\
 \langle f,R_\chi f\rangle_\mu
       &=\frac12\int
             \operatorname{Re}\{\overline{f(U)}f(V)\}
             \sum_i|\chi_i(U)-\chi_i(V)|^2
                                      \,\mathcal J(dU,dV),\\
 |\langle f,R_\chi f\rangle_\mu|
       &\le\int W_\chi|f|^2\,d\mu
                      \le\|W_\chi\|_\infty\|f\|^2 .
 \end{split}
 \label{eq:f15v78:localization}
\end{equation}
\end{theorem}

\begin{proof}
In the sum of
\(\|\chi_i f\|^2-\langle\chi_i f,\mathcal P\chi_i f\rangle\),
the norm terms add to \(\|f\|^2\).
Subtract \(d[f]\). The remaining integral has coefficient
\(1-\sum_i\chi_i(U)\chi_i(V)\).
Since both partition vectors have Euclidean norm one,
that coefficient is half their squared distance.
Symmetry of \(\mathcal J\) makes the integral real.
Finally \(2|f(U)f(V)|\le |f(U)|^2+|f(V)|^2\);
integrating each symmetric half gives the stated bound.
All terms are continuous quadratic forms on \(L^2\),
so bounded-function approximation proves the full claim.
\end{proof}

The correction is an integral operator, not the
multiplication operator \(\sum_i|\nabla\chi_i|^2\)
from a diffusion IMS formula. Its kernel is
nonnegative, but it need not be positive semidefinite.
This distinction matters if its sign is retained.
The vacuum gives the consistency identity
\[
          \sum_i d[\chi_i]
                  =\langle1,R_\chi1\rangle_\mu
                  =\int W_\chi\,d\mu .
\]
A positive cost assigned to every localized vacuum
component cannot simply survive subtraction of this cost.

\subsection{Complete native channels with finite low spaces}

Let \(D_i=\{U:\chi_i(U)>0\}\), omitting null sets,
and let
\(\mathcal K_i=L^2(D_i,\mu)^{\rm gauge}\).
Let \(E_i:\mathcal K_i\to\mathcal H_\mu\) be extension
by zero. The bounded local operator and its form are
\begin{equation}
        A_i=I_i-E_i^*\mathcal P E_i,\qquad
        \ell_i[z]=\langle z,A_i z\rangle=d[E_i z].
 \label{eq:f15v78:channels}
\end{equation}
The map \(V_i f=(\chi_i f)|_{D_i}\), with domain
\(\mathcal H_\mu^0\), gives an isometry
\(V:\mathcal H_\mu^0\to\bigoplus_i\mathcal K_i\).
These channels cover the full configuration space;
no retained physical direction is omitted.

Fix \(0<\beta<1\). Since
\(E_i^*\mathcal P E_i\) is compact and nonnegative,
\begin{equation}
       P_i=1_{[0,\beta)}(A_i)
       \quad\hbox{has finite rank},\qquad
       \ell_i[z]\ge\beta\|(I-P_i)z\|^2 .
 \label{eq:f15v78:local-low}
\end{equation}
Indeed low eigenvalues of \(A_i\) correspond to
eigenvalues greater than \(1-\beta>0\) of that compact
operator. This supplies legitimate finite low spaces
at every fixed regulator, not estimates of their
dimensions or overlaps uniform in refinement.
A different finite projection is admissible if the
same complementary lower bound is independently proved.

For an orthonormal basis \(\omega_{i\alpha}\) of
\(\operatorname{ran}P_i\), set
\begin{equation}
 \begin{split}
 u_{i\alpha}&=Q(\chi_i E_i\omega_{i\alpha}),\\
 G_{i\alpha,j\nu}&=\langle u_{i\alpha},u_{j\nu}\rangle_\mu,\\
 L_\chi&=\sum_{i,\alpha}|u_{i\alpha}\rangle
                                      \langle u_{i\alpha}|.
 \end{split}
 \label{eq:f15v78:gram}
\end{equation}
The vacuum projection \(Q\) cannot be dropped.

\begin{theorem}[Complete transfer-channel certificate]
\label{f15v78:certificate}
Let \(\eta=\|G\|=\|L_\chi\|\), with value zero for
an empty low-mode family, and let
\(\epsilon=\|W_\chi\|_\infty\).
Then for every \(f\in\mathcal H_\mu^0\),
\begin{equation}
       d[f]\ge
        \{\beta(1-\eta)-\epsilon\}\|f\|^2 .
 \label{eq:f15v78:separate}
\end{equation}
A sharper sufficient certificate retains the signed
correction. On \(\mathcal H_\mu^0\) put
\begin{equation}
           B_\chi=\beta L_\chi+QR_\chi Q,\qquad
           b_\chi=\sup\operatorname{spec}B_\chi .
 \label{eq:f15v78:combined}
\end{equation}
Then
\begin{equation}
                       d[f]\ge(\beta-b_\chi)\|f\|^2 .
 \label{eq:f15v78:combinedbound}
\end{equation}
Whenever a certified \(\theta>0\) is supplied by
either bound, Proposition~\ref{f15v78:conversion}
converts it to the physical logarithmic gap.
\end{theorem}

\begin{proof}
For centered \(f\), the adjoint of \(V_i\) is
\(V_i^*z=Q(\chi_i E_i z)\). Consequently
\[
 \sum_i\|P_iV_i f\|^2
      =\sum_{i,\alpha}|\langle u_{i\alpha},f\rangle|^2
      =\langle f,L_\chi f\rangle .
\]
The finite synthesis map has \(L_\chi\) as one
Gram product and \(G\) as the other, so their
nonzero spectra and norms agree. Since \(V\)
is an isometry, also \(0\le L_\chi\le I\).
Sum \eqref{eq:f15v78:local-low} and subtract the
exact correction in \eqref{eq:f15v78:localization}:
\[
 d[f]\ge\beta\|f\|^2-\beta\langle f,L_\chi f\rangle
                                      -\langle f,R_\chi f\rangle .
\]
This is \eqref{eq:f15v78:combinedbound}. Bounding
the two terms separately gives
\eqref{eq:f15v78:separate}.
The separated argument is Gauvin's Theorem 6.11
specialized to the native bounded transfer form;
the present localization and signed correction
are proved explicitly for that form.
\end{proof}

For the exact quantities, \(b_\chi\le\beta\eta+\epsilon\),
so the coupled bound is at least as strong as the
displayed separated bound. It avoids adding two
worst-case directions which may not coincide.
A positive margin is not asserted for Yang--Mills here.

\subsection{A full-space residual certificate}

The exact finite-regulator kernel in
\eqref{eq:f15v78:ground} is continuous on a compact
configuration product. Its strictly positive ground
state has positive minimum. Thus the kernel of
\(\mathcal P\) relative to \(\mu\),
\[
               p(U,V)=\frac{\mathsf t(U,V)}
                                  {\lambda_0\Omega(U)\Omega(V)},
\]
is bounded and Hilbert--Schmidt. The same is true
of \(R_\chi\) for a bounded finite partition.
Hence \(B_\chi\), the sum of a finite-rank and a
Hilbert--Schmidt operator, is compact self-adjoint.

\begin{proposition}[Finite matrices with a certified omitted remainder]
\label{f15v78:residual}
Let \(F_n\) be finite-rank orthogonal projections on
\(\mathcal H_\mu^0\), with \(F_n\to I\) strongly.
Suppose
\[
 \lambda_{\max}(F_nB_\chi F_n|_{\operatorname{ran}F_n})
                         \le b_n,\qquad
 \|B_\chi-F_nB_\chi F_n\|_{\rm HS}\le r_n .
\]
Then
\begin{equation}
 d[f]\ge
       \{\beta-\max(0,b_n)-r_n\}\|f\|^2
                         \qquad(f\in\mathcal H_\mu^0).
 \label{eq:f15v78:residual}
\end{equation}
For a fixed exact \(B_\chi\), its actual omitted
Hilbert--Schmidt remainder tends to zero.
\end{proposition}

\begin{proof}
The compression, extended by zero, is bounded above
by \(\max(0,b_n)I\). The operator norm of its
difference from \(B_\chi\) is at most the stated
Hilbert--Schmidt norm. Apply
\eqref{eq:f15v78:combinedbound}. Convergence of the
actual remainder follows by first approximating
\(B_\chi\) in Hilbert--Schmidt norm by a finite-rank
operator and then using strong convergence of
\(F_n\) on its finitely many vectors.
\end{proof}

This specifies the direction of numerical certification.
Eigenvalues of finite compressions alone are not upper
bounds for the complete operator. Enclosures of the
actual vacuum, local spectral spaces, all Gram entries,
and the omitted kernel are still required.
No finite table computed for a surrogate probability
law can supply those enclosures.

\subsection{Revised direction and the physical clock}

The immediate target is now one positive certified
complete-channel margin for the native transfer in
\eqref{eq:f15v78:residual}, with its scale dependence
recorded. V75--V77 can suggest configurations for
channels, but their gradient estimates do not yet
prove \eqref{eq:f15v78:local-low} for those physical
channels. The exact compact-kernel construction
above is an alternative to assuming that identification.

If subsequent reduction is used, its forms must
be exact relaxations of the identified physical
form; the shell bounds must have the full-vector
scope illustrated in Proposition~\ref{f15v78:vertical}.
A static marginal, a compressed energy and a relaxed
energy are not interchangeable. Likewise the
paper's vacuum-score comparison cannot be applied
to retained action densities without a same-carrier,
same-electric-symbol bridge.

At regulator \(j\), let the eventual certified
dimensionless margin be \(\theta_j\).
To prove a physical lower bound \(m_*>0\), an
independently calibrated time step must satisfy
\begin{equation}
        -\frac{\log(1-\theta_j)}{a_{t,j}}\ge m_* .
 \label{eq:f15v78:clock}
\end{equation}
A sufficient target is
\(\theta_j\ge1-e^{-m_*a_{t,j}}\).
V63's one-step gap tends to zero along its specified
sequence. Thus a fixed positive one-step margin
along that same sequence would contradict the
inherited result; the target is a correctly scaled
margin, not an arbitrarily reset clock.

Even that finite-regulator spectral estimate would
not finish the continuum construction. A total
renormalized probe family, spectral tightness,
compatible limiting dynamics, locality and
relativistic reconstruction remain to be proved.
The paper's conditional criteria organize these
requirements but do not verify them for this programme.
No completion gate is declared closed by this audit.

\subsection{Diagnostics and source preservation}

The new diagnostic is
\begin{center}\small
\path{scripts/verify_ym_f15_v78_transfer_channel_certificate.py}.
\end{center}
It reruns V77 and its inherited checks. New rational
fixtures verify four positive transfer eigenpairs,
243 jump/localization identities and bounds, eight
local channel complement inequalities, 48 centered
leakage and combined-error checks, and three exact
positive finite-model certificates. Sixteen
nonconstant-vacuum pair-measure identities retain
the Doob weights. The restricted-block counterexample,
the nonembeddable positive transfer, and 27
Hilbert--Schmidt remainder bounds are checked as well.
These fixtures validate algebra; they supply no
numerical Yang--Mills margin.

The script has 7,969 bytes and SHA--256
\begin{center}\scriptsize\ttfamily
ED2BF7ACD9B9882C5F96FF941DD6A8F64A008670D87BFBD2CFE1574E9BBB958D.
\end{center}
The supplied paper has 584,030 bytes and SHA--256
\begin{center}\scriptsize\ttfamily
F21A15CF26B916F008C3CF41B576682F2A81C8C53DB860CEC0C0F5A265A42A95.
\end{center}
The V77 predecessor has 1,743,096 bytes and SHA--256
\begin{center}\scriptsize\ttfamily
BABC2053553FC257779E55A85DA350BDA58E807DBE236B25950927B91E832685.
\end{center}
Removing this final section and restoring the title
version recovers V77 exactly. V78 becomes the sole
active main note; protected archives and companions
remain unchanged. Only \path{.tex} sources are kept
in \path{latex_notes}. Compilation products are kept
outside that folder, in \path{artifacts/ym_f15_v78_texcheck/}.

\begin{firewall}
V78 audits the supplied paper, corrects the scope
of its possible transfers, and proves a native
nonlocal complete-channel certificate for the
existing physical Wilson transfer, including the
centered low-mode Gram matrix and a full omitted-space
remainder. Its required positive uniform physical
margin is not established. The interacting continuum
and full Yang--Mills mass gap remain unproved.
The original goal remains active.
\end{firewall}

\section{V79: vacuum-independent Haar certificates for the literal Wilson transfer}
\label{f15v79:section}

\subsection{Context and the upstream change}

V78 transferred the supplied paper's complete-channel idea
to the actual nonlocal one-step Wilson form. Its implementation
still required the ground-state transform and hence quantities
involving the unknown vacuum. This section, dated
16 September 2026, removes that requirement from a
finite-regulator certification route. It keeps the physical
Wilson transfer and works in product Haar space.

The new ingredients are an explicit positive finite-rank
approximation to that transfer, a full omitted-space operator
bound, a finite matrix with the same nonzero spectrum, and
a vacuum-independent version of the channel certificate.
The first ingredient truncates the kinetic factor, not the
magnetic action and not the transfer's unknown eigenbasis.
Polynomial Haar integration then gives a convergent way to
enclose the finite matrix entries. None of these steps supplies
a uniform continuum spectral margin.

The one-link character expansion is already used in
Proposition~\ref{f15v57:transfer}. The positive Bessel series
and its fixed-coupling far-tail asymptotic were also proved
in f5 V42, Lemma \emph{Fixed-coupling Bessel tail}.
They are reused, not counted as new results here.
The additional statement below is an all-mode
operator-order bound for the literal magnetic sandwich,
followed by a full-complement spectral certificate.
The paper's general channel strategy remains credited
to the V78 source audit; no new claim is made about the
paper's unverified continuum hypotheses.

\subsection{Normalize only by a known scalar}

Fix a finite spatial lattice with \(E\) links and
\(P_{\rm sp}\) spatial plaquettes. Write \(\gamma>0\)
for the coupling denoted by \(\beta\) in V57.
Let \(m\) be normalized product Haar measure and let
\(\mathcal H=L^2(m)^{\rm gauge}\).
The magnetic action and multiplication factor are
\[
 S_{\rm sp}(U)=\sum_{p\ {\rm spatial}}
       \left(1-\tfrac12\Tr U_p\right),\qquad
 b(U)=e^{-\gamma S_{\rm sp}(U)/2},\qquad
 0\le S_{\rm sp}\le2P_{\rm sp}.
\]
Thus \(0<b\le1\), and multiplication \(M_b\) is bounded
and invertible at this fixed regulator. Let \(C_\gamma\)
be one-link convolution with
\(w_\gamma(u)=e^{-\gamma(1-\frac12\Tr u)}\), and put
\(c_0=\int w_\gamma\,dm_{\SU(2)}>0\).
V57 gives, on the physical subspace,
\begin{equation}
 \overline T=c_0^{-E}\mathsf T
       =M_b C M_b,\qquad
 C=\bigotimes_{e=1}^E(C_\gamma/c_0).
 \label{eq:f15v79:transfer}
\end{equation}
The gauge average acts as the identity on this subspace.
All the factors in this formula preserve it.
The scalar \(c_0^{-E}\) changes neither
\(\lambda_1/\lambda_0\) nor the vacuum-normalized
physical Hamiltonian. It is not division by an
unknown eigenvalue.

Label irreducible \(\SU(2)\) representations by
\(n=2j\in\mathbb N_0\), of dimension \(n+1\).
One-link Peter--Weyl matrix coefficients are eigenvectors
of \(C_\gamma/c_0\) with eigenvalues
\begin{equation}
 r_n=\frac{I_{n+1}(\gamma)}{I_1(\gamma)},\qquad
 r_0=1,\qquad 0<r_n\le1.
 \label{eq:f15v79:ratios}
\end{equation}
For completeness, the Weyl integral for the character
\(\vartheta_n(\theta)=\sin((n+1)\theta)/\sin\theta\)
gives its convolution eigenvalue
\[
 c_n=\frac{e^{-\gamma}}{n+1}
       \bigl(I_n(\gamma)-I_{n+2}(\gamma)\bigr)
     =e^{-\gamma}\frac{2I_{n+1}(\gamma)}{\gamma}.
\]
The Fourier coefficients of
\(e^{\gamma\cos\theta}
 =e^{(\gamma/2)e^{i\theta}}e^{(\gamma/2)e^{-i\theta}}\)
give the absolutely convergent positive series
\begin{equation}
 I_j(\gamma)=\sum_{k=0}^{\infty}
       \frac{(\gamma/2)^{2k+j}}{k!(k+j)!}.
 \label{eq:f15v79:bessel}
\end{equation}
Comparing coefficients gives
\(I_n-I_{n+2}=2(n+1)I_{n+1}/\gamma\).
These standard formulas agree with
\href{https://dlmf.nist.gov/10.25.E2}{NIST DLMF 10.25.2}
and \href{https://dlmf.nist.gov/10.29}{DLMF 10.29}.
Positivity follows from the series; the upper bound
in \eqref{eq:f15v79:ratios} follows because
\(w_\gamma/c_0\) is a probability density, so its
convolution is an \(L^2\) contraction. In particular
\(0\le C\le I\) and \(0\le\overline T\le I\).

\subsection{A positive approximation with a full-space remainder}

Let \(\Pi_N\) be the orthogonal projection onto the
Peter--Weyl sectors with \(n_e\le N\) on every link,
restricted to \(\mathcal H\). These projections commute
with the vertex gauge action and with \(C\). They
increase strongly to the identity on \(\mathcal H\).
Their ranges \(\mathcal F_N\) are finite dimensional;
before imposing gauge invariance the dimension is
\[
 \left(\sum_{n=0}^N(n+1)^2\right)^E
   =\left(\frac{(N+1)(N+2)(2N+3)}6\right)^E.
\]
Set \(C_N=\Pi_N C\Pi_N\) and \(T_N=M_b C_N M_b\).

\begin{theorem}[Literal-transfer positive tail bound]
\label{f15v79:tail}
If \(N+2\ge\gamma/2\), define
\[
 \epsilon_N=\min\left\{1,\,
                \frac{(\gamma/2)^{N+1}}{(N+2)!}\right\}.
\]
Then \(T_N\) has finite rank, preserves the physical
subspace, and
\begin{equation}
 0\le\overline T-T_N
       \le\epsilon_N M_{b^2}\le\epsilon_N I .
 \label{eq:f15v79:tail}
\end{equation}
Also \(T_N\le T_{N+1}\le\overline T\), and
\(T_N\to\overline T\) in operator norm at fixed
volume and coupling.
There is no factor \(E\) in this operator-norm
tail bound.
\end{theorem}

\begin{proof}
Termwise comparison with the series for \(I_1\) gives
\[
 \frac{(\gamma/2)^{2k+n+1}/[k!(k+n+1)!]}
      {(\gamma/2)^{2k+1}/[k!(k+1)!]}
 =(\gamma/2)^n\frac{(k+1)!}{(k+n+1)!}
 \le\frac{(\gamma/2)^n}{(n+1)!}.
\]
Consequently \(r_n\) is bounded by this last expression.
For \(n\ge N+1\), these factorial majorants decrease,
since their successive ratio is
\((\gamma/2)/(n+2)\le1\).
Together with \(r_n\le1\), this proves
\(r_n\le\epsilon_N\) for every omitted one-link label.

On the product Peter--Weyl sector with labels
\((n_e)\), the eigenvalue of \(C\) is
\(\prod_e r_{n_e}\). An omitted sector has at least
one label exceeding \(N\); that factor is at most
\(\epsilon_N\), while every other factor is at most
one. Thus \(0\le C-C_N\le\epsilon_N I\).
This operator inequality restricts to the gauge
subspace because all involved operators commute
with its projection. Sandwiching by \(M_b\)
proves \eqref{eq:f15v79:tail}.
The increasing spectral cutoffs give monotonicity.
At fixed \(\gamma\), the factorial bound tends
to zero, proving norm convergence.
\end{proof}

This is an operator-norm estimate, not a trace-norm
tail estimate. It does not discard representation
multiplicities in a trace calculation. Its absence
of a link-count factor is therefore not a
dimension-free assertion about partition functions.

It is essential that
\[
 T_N=M_b\Pi_N C\Pi_N M_b
       \quad\hbox{need not equal}\quad
       \Pi_N\overline T\Pi_N .
\]
The latter compression does not in general lie
below \(\overline T\) in operator order, because
\(\Pi_N\) need not commute with \(M_b\).
Also positive operator order does not mean that
the truncated kernel is pointwise nonnegative.
Already the one-link \(N=1\) kernel has value
\(1-4r_1\) at \(-I\), which is negative at
\(\gamma=2\): \(I_2(2)>1/2\), whereas
\(I_1(2)=\sum_k1/[k!(k+1)!]<2\).
No Markov or positivity-improving property of
the truncated kernel is used below.

\subsection{An exact finite matrix retaining the magnetic action}

Choose an orthonormal physical Peter--Weyl basis
\((\phi_\alpha)\) of \(\mathcal F_N\), with each
vector in a fixed link-label sector, and write
\(d_\alpha=\prod_e r_{n_e(\alpha)}^{1/2}\).
Define
\[
 W_N:\mathcal F_N\longrightarrow\mathcal H,\qquad
 W_N=M_b C_N^{1/2}|_{\mathcal F_N},\qquad
 G_N=W_N^*W_N.
\]
Then
\begin{equation}
 T_N=W_NW_N^*,\qquad
 (G_N)_{\alpha\beta}
   =d_\alpha d_\beta\int b(U)^2
                  \overline{\phi_\alpha(U)}\phi_\beta(U)\,dm(U).
 \label{eq:f15v79:gram}
\end{equation}
Since \(d_\alpha>0\) and \(b>0\), \(W_N\) is
injective and \(G_N\) is positive definite.
The nonzero eigenvalues of \(T_N\), including
multiplicities, are exactly those of \(G_N\).
Indeed \(T_NW_N=W_NG_N\), and every nonzero
eigenvector of \(T_N\) belongs to its range.
Equivalently the polar decomposition of \(W_N\)
identifies \(G_N\) with the restriction of \(T_N\)
to that range.

Let \(\nu_0(G_N)\ge\nu_1(G_N)\) denote the top
two eigenvalues after adjoining zero eigenvalues
when necessary. If \(\ell_N>0\) and \(u_N\ge0\)
are certified bounds
\[
 \nu_0(G_N)\ge\ell_N,\qquad \nu_1(G_N)\le u_N,
\]
the min--max principle and \eqref{eq:f15v79:tail}
give
\begin{equation}
 \lambda_0(\overline T)\ge\ell_N,\qquad
 \lambda_1(\overline T)\le u_N+\epsilon_N .
 \label{eq:f15v79:top-two}
\end{equation}
Here \(\lambda_1\) is the second eigenvalue with
multiplicity, not the largest eigenvalue seen
by a selected observable. This is a full physical
space assertion.

\begin{proposition}[Vacuum-independent physical certificate]
\label{f15v79:certificate}
If \(0<U_N:=u_N+\epsilon_N<\ell_N\), then at a
specified physical time step \(a_t>0\),
\begin{equation}
 \operatorname{gap}\left[
   -a_t^{-1}\log\bigl(\mathsf T/\lambda_0(\mathsf T)\bigr)
                    \right]
 \ge\frac1{a_t}\log\frac{\ell_N}{U_N}>0.
 \label{eq:f15v79:certificate}
\end{equation}
No value, lower pointwise bound, or derivative
of the vacuum wavefunction is needed for this
certificate.
\end{proposition}

\begin{proof}
For compact positive operators, each ordered
eigenvalue is increasing in operator order and
is Lipschitz with constant one in operator norm.
These facts follow directly from the Rayleigh
quotient min--max formulas. They give
\eqref{eq:f15v79:top-two}, including the zero
complement of \(T_N\). Hence
\(\lambda_1(\overline T)/\lambda_0(\overline T)
 \le U_N/\ell_N<1\).
The known scalar normalization cancels in this
ratio. The spectral calculus used in
Proposition~\ref{f15v78:conversion} gives the
claimed logarithmic gap.
\end{proof}

At \(N=0\), \(\mathcal F_0\) consists of the
constant function, so \(T_0=|b\rangle\langle b|\)
and \(G_0=\int b^2\,dm\). This checks that the
magnetic weight and its normalization have not
been dropped. It is not a uniform gap certificate.

\subsection{Certified matrix entries without a vacuum density}

The finite entries in \eqref{eq:f15v79:gram}
are actual Haar integrals with \(b^2=e^{-\gamma S_{\rm sp}}\);
they are not averages under a retained marginal
or a Gaussian replacement. There is also a
direct convergent enclosure for these integrals.

Set \(A=2\gamma P_{\rm sp}\) and, for \(K\ge0\), put
\[
 p_K(s)=\sum_{k=0}^{K}\frac{(-\gamma s)^k}{k!},
 \qquad
 \eta_K=\frac{A^{K+1}}{(K+1)!},\qquad
 B_K=\Pi_N M_{p_K(S_{\rm sp})}\Pi_N|_{\mathcal F_N}.
\]
Taylor's theorem on \(0\le\gamma s\le A\),
whose next exponential derivative has absolute
value at most one, proves
\[
 \|B_K-\Pi_NM_{b^2}\Pi_N|_{\mathcal F_N}\|
        \le\eta_K .
\]
This is again an operator norm estimate, not
an entrywise bound multiplied by the matrix
dimension.

\begin{proposition}[Action and kinetic-weight error ledger]
\label{f15v79:entries}
Write \(D=C_N^{1/2}|_{\mathcal F_N}\). Suppose a
Hermitian diagonal approximation \(\widetilde D\)
satisfies \(\|\widetilde D-D\|\le\delta\).
If a Hermitian matrix \(\widetilde G\) satisfies
\(\|\widetilde G-\widetilde D B_K\widetilde D\|\le\kappa\),
then
\begin{equation}
 \|\widetilde G-G_N\|\le
 \zeta:=\delta(2+\delta)+(1+\delta)^2\eta_K+\kappa .
 \label{eq:f15v79:entry-error}
\end{equation}
In particular a certified lower bound \(L\) for
the largest eigenvalue of \(\widetilde G\), and
an upper bound \(U\ge0\) for its second eigenvalue
including the zero complement, give
\[
 \ell_N=L-\zeta,\qquad u_N=U+\zeta
\]
as valid inputs to \eqref{eq:f15v79:top-two},
provided \(\ell_N>0\).
\end{proposition}

\begin{proof}
Let \(B=\Pi_NM_{b^2}\Pi_N|_{\mathcal F_N}\).
Then \(\|B\|\le1\), \(\|D\|\le1\), and
\(\|\widetilde D\|\le1+\delta\). Decompose
\[
 DBD-\widetilde D B_K\widetilde D
  =(D-\widetilde D)BD
     +\widetilde D B(D-\widetilde D)
     +\widetilde D(B-B_K)\widetilde D .
\]
The norms of the three terms are bounded by
\(\delta\), \((1+\delta)\delta\), and
\((1+\delta)^2\eta_K\). Add \(\kappa\).
The eigenvalue bounds follow from min--max after
adjoining a zero subspace to both finite matrices.
\end{proof}

These enclosures can be made finite and explicit.
Identify each link with a unit quaternion
\(q=(q_0,q_1,q_2,q_3)\); Haar measure is uniform
on \(S^3\). Odd coordinate monomials integrate
to zero, and for \(a_i\in\mathbb N_0\),
\begin{equation}
 \int_{S^3}\prod_{i=0}^3q_i^{2a_i}\,dm
  =\frac1{(1+\sum_i a_i)!}
          \prod_{i=0}^3\frac{(2a_i)!}{4^{a_i}a_i!}.
 \label{eq:f15v79:moments}
\end{equation}
To prove this formula, integrate
\(e^{-|x|^2}\prod_i x_i^{2a_i}\) first by four
one-dimensional Gaussian integrals and then in
radial coordinates in \(\mathbb R^4\); divide
by the same identity for all \(a_i=0\).
Reflection proves the assertion for odd powers.

The Wilson action is a polynomial in the link
quaternion coordinates and their conjugates.
Fixed Peter--Weyl matrix coefficients are
polynomials as well. Gauge averaging introduces
only further finite polynomial Haar integrals.
One may therefore obtain a finite physical basis
by gauge averaging a finite Peter--Weyl spanning
set, removing its exact null Gram directions,
and orthonormalizing. In standard algebraic
representation coordinates this uses algebraic
coefficients. Formula \eqref{eq:f15v79:moments},
link by link and gauge variable by gauge variable,
then evaluates every entry of \(B_K\).

For rational \(\gamma\), those polynomial integrals
have algebraic coefficients; for a coupling
supplied by certified approximations, their
additional rounding errors must be included
in \(\kappa\). Positive partial sums of
\eqref{eq:f15v79:bessel} have geometric tail
bounds once the successive term ratio is below
one. Dividing their positive intervals and
enclosing square roots bounds every \(d_\alpha\).
The maximum diagonal error is \(\delta\).
Any further finite matrix-entry errors can,
for example, be converted into \(\kappa\)
by a Frobenius norm bound.

At fixed regulator, \(\eta_K\to0\), and
\(\delta,\kappa\) may be reduced by certified
arithmetic. This establishes a convergent
entry-enclosure route, not an efficient
large-volume algorithm. In particular it
does not claim to have evaluated the required
large physical matrices in this iteration.

\subsection{Haar channels using any trial vector}

There is a channel version that also avoids
the exact vacuum. For a normalized
\(h\in\mathcal H\), write
\(Q_h=I-|h\rangle\langle h|\).
The min--max principle gives
\begin{equation}
 \lambda_1(\overline T)
 \le\sup_{\substack{f\perp h\\\|f\|=1}}
                      \langle f,\overline T f\rangle .
 \label{eq:f15v79:minmax}
\end{equation}
This uses the entire complement of \(h\).
An upper Rayleigh bound on a few vectors in
that complement would not suffice.

Fix \(N\), abbreviate \(W=W_N\), and let
\(\chi_1,\ldots,\chi_s\) be bounded nonnegative
gauge-invariant functions with \(\sum_i\chi_i^2=1\)
almost everywhere. Put \(D_i=\{\chi_i>0\}\),
let \(E_i:L^2(D_i,m)^{\rm gauge}\to\mathcal H\)
be extension by zero, and let \(F_i=E_i^*W\).
The local operator is \(F_iF_i^*=E_i^*T_NE_i\).
Its finite Gram matrix is
\[
 A_i=F_i^*F_i=W^*1_{D_i}W .
\]
For \(\sigma>0\) define the continuous function
on \([0,\infty)\)
\[
 k_\sigma(c)=
 \begin{cases}
   1-\sigma/c,&c>\sigma,\\
   0,&0\le c\le\sigma .
 \end{cases}
\]
Spectral decomposition of \(F_i^*F_i\) gives
\[
 (F_iF_i^*-\sigma I)_+
       =F_i k_\sigma(A_i)F_i^* .
\]
This formula includes the null space correctly
and requires no inverse of a singular Gram matrix.

Define \(X_i=M_{\chi_i}W\), and the finite-column
operator and Hermitian matrix
\[
 Z=[\,W,\ X_1,\ldots,X_s\,],\qquad
 J=\operatorname{diag}
       \bigl(I,k_\sigma(A_1)-I,\ldots,k_\sigma(A_s)-I\bigr).
\]
Finally set \(G_h=Z^*Q_h Z\).
Its entries are Haar inner products of known
finite families, with the trial-vector
projection retained.

\begin{theorem}[Finite signed-Gram channel certificate]
\label{f15v79:channels}
In the preceding notation, let
\[
 \mathfrak b_{N,h}=
   \max\{0,\lambda_{\max}(G_h^{1/2}JG_h^{1/2})\}.
\]
Then
\begin{equation}
 \lambda_1(\overline T)
    \le\sigma+\mathfrak b_{N,h}+\epsilon_N,\qquad
 \lambda_0(\overline T)
    \ge a_{N,h}:=\langle h,T_Nh\rangle=\|W^*h\|^2 .
 \label{eq:f15v79:channel}
\end{equation}
If certified bounds \(a_{N,h}\ge a_*>0\) and
\(\mathfrak b_{N,h}\le b_*\) satisfy
\(0<\sigma+b_*+\epsilon_N<a_*\), the physical
gap is at least
\[
          a_t^{-1}\log
            \frac{a_*}{\sigma+b_*+\epsilon_N}.
\]
All projections in this certificate are known
trial-vector or finite Gram projections;
none is the unknown vacuum projection.
\end{theorem}

\begin{proof}
Define the exact localization correction for
the truncated transfer by
\[
 R_N=T_N-\sum_i M_{\chi_i}T_NM_{\chi_i}.
\]
For each local operator,
\[
 F_iF_i^*\le\sigma I+(F_iF_i^*-\sigma I)_+ .
\]
Evaluate this inequality at \(E_i^*(\chi_i f)\),
sum over \(i\), and use \(\sum_i\|\chi_i f\|^2=\|f\|^2\).
The identity for the positive excess and
\(M_{\chi_i}1_{D_i}=M_{\chi_i}\) give
\[
 \begin{split}
 T_N
 &\le \sigma I+R_N+\sum_iX_i k_\sigma(A_i)X_i^*\\
 &=\sigma I+WW^*+\sum_iX_i(k_\sigma(A_i)-I)X_i^*
   =\sigma I+ZJZ^* .
 \end{split}
\]
This is an inequality on the full physical
Hilbert space, not just the span of the columns.

Let \(Z_h=Q_hZ\) and use its polar decomposition
\(Z_h=V G_h^{1/2}\). On the support of \(G_h\),
\(V\) is an isometry, so the nonzero spectrum of
\(Z_hJZ_h^*\) is exactly that of
\(G_h^{1/2}JG_h^{1/2}\).
Outside the finite range the operator is zero.
Including zero therefore gives the valid upper
bound \(\mathfrak b_{N,h}\) on all of \(h^\perp\),
even when a spanning Gram matrix is singular.
Compress the preceding operator inequality
to \(h^\perp\), add the full remainder
\(\epsilon_N I\), and apply
\eqref{eq:f15v79:minmax}. The lower bound is the
Rayleigh quotient and \(T_N\le\overline T\).
The logarithmic conversion is the same as in
Proposition~\ref{f15v79:certificate}.
\end{proof}

The signed matrix \(J\) must not be replaced
by a positive matrix or by selected favorable
channels. Its negative blocks carry the exact
localization correction. Nor may a singular
\(G_h\) be inverted without passing to its
support. The square-root formulation above
needs no such inverse.

Channel indicators and partition functions
introduce additional Haar integrals involving
\(1_{D_i}\), \(\chi_i\), and \(h\).
The polynomial-action enclosure alone does
not automatically evaluate those integrals.
They need their own certified bounds for the
chosen channels. The direct global matrix
certificate remains available without this
extra choice. Channel decomposition is an
organizational option, not an assertion that
it necessarily improves the optimal global
eigenvalue estimate.

\subsection{The remaining relative-error and continuum obligations}

For a fixed finite regulator, positivity and
compactness already give an isolated largest
eigenvalue. Norm convergence of \(T_N\) makes
increasingly accurate top-two enclosures
possible. This does not create a new
uniform physical gap from the previously
known finite-volume Perron separation.

The full tail bound is absolute. Its useful
relative size is \(\epsilon_N/\ell_N\), not
\(\epsilon_N\) alone. For example, the
rank-one lower approximation gives only
\[
 \lambda_0(\overline T)\ge\int b^2\,dm
                \ge e^{-2\gamma P_{\rm sp}}.
\]
The leading eigenvalue can be very small as
volume grows. The absence of \(E\) from the
absolute kinetic tail is not a proof that a
fixed \(N\) resolves the needed spectral ratio.
The matrix dimension and the polynomial-action
scale \(A=2\gamma P_{\rm sp}\) also grow.

Along a continuum sequence, one must actually
produce certified \(0<U_j<L_j\) and use the
pre-existing physical time calibration to prove
\begin{equation}
       \log(L_j/U_j)\ge m_*a_{t,j}
       \quad\hbox{with a single }m_*>0 .
 \label{eq:f15v79:physical-target}
\end{equation}
Here \(L_j,U_j\) include all entry, arithmetic,
channel and omitted-space errors appropriate
to the chosen certificate. No such sequence
is obtained in this section.
V63's vanishing one-step gap on its specified
sequence is still binding; a uniform positive
one-step margin there would be inconsistent
with that result. Recalibrating time after
seeing a spectral estimate is not allowed.

The concrete upstream task is now to exploit
the magnetic action and suitable physical
channels to obtain relative top-two estimates
for these Haar matrices at growing regulators.
Unknown pointwise vacuum bounds are no longer
a prerequisite for this route. Uniform matrix
estimates, a total limiting probe family,
spectral tightness, compatible limiting
dynamics, locality and relativistic
reconstruction remain separate obligations.

\subsection{Diagnostics, nonduplication and preservation}

The diagnostic
\begin{center}\small
\path{scripts/verify_ym_f15_v79_haar_transfer_enclosures.py}
\end{center}
reruns V78 and its inherited checks. New checks
include 150 outward Bessel-ratio bounds,
120 entire-tail majorant fixtures,
1,302 omitted product-mode bounds and 81
full-vector transfer-tail inequalities.
They distinguish the positive kinetic
sandwich from the potentially invalid
operator-order compression. A non-vacuum
trial vector gives an exact finite-model
certificate. Four local positive-excess
blocks, a signed centered Gram certificate
and five spectral power identities verify
the channel algebra. Thirty exponential
polynomial bounds, nine normalized
quaternion Haar identities and nine combined
entry-error budgets are also checked.
The all-mode analytic theorems above do not
follow from finite sampling, and these
fixtures give no numerical Yang--Mills margin.

The script has 8,661 bytes and SHA--256
\begin{center}\scriptsize\ttfamily
83BC80AA4271E12044EEAFE57628CA215B943B547AC2297ABB54BD391CF6F63B.
\end{center}
The V78 predecessor has 1,765,414 bytes and SHA--256
\begin{center}\scriptsize\ttfamily
0728D35F67A036585BA177B8DCA03C9B75633DFC4E8F8CD828B5E242A9111E25.
\end{center}
Removing this final section and restoring the
title version recovers V78 byte for byte.
The protected sources were unchanged before
this append. The last periodic companion
checkpoint was V75; the next is V80.
V78's separate audit of the newly supplied
paper is preserved and is not presented
again as new work.

V79 is the sole active main note.
Archives and companion notes remain in place.
Only \path{.tex} sources are stored in
\path{latex_notes}; compilation products go
to \path{artifacts/ym_f15_v79_texcheck/}.

\begin{firewall}
V79 proves an explicit positive finite-rank
approximation and full-space tail bound for
the native physical Wilson transfer, derives
vacuum-independent global and channel
certificates, and supplies convergent
matrix-entry enclosures. A positive
regulator-uniform physical margin has not
been established. The interacting continuum
and full Yang--Mills mass gap remain unproved.
The original goal remains active.
\end{firewall}

\section{V80: relative transfer tails and the large-volume reference obstruction}
\label{f15v80:section}

\subsection{The scheduled companion checkpoint}

This section, dated 16 September 2026, first performs
the requested five-version checkpoint. Fresh inventories
of \path{latex_notes} and the supplied Downloads sources
showed no newer relevant versions. The current terminal
conclusions of the seven companion sources below were
reread, together with \path{suggestions.txt}. Their hashes
are unchanged from the preceding checkpoint.
The newly supplied transfer paper remains covered by
V78's separate audit.

\begin{center}\small
\begin{tabular}{@{}lrl@{}}
\toprule
Source & Bytes & SHA--256 prefix\\
\midrule
SM continuum V72 & 607372 & \texttt{F44F9745D43379E1}\\
NS stability V26 & 278283 & \texttt{82C887F8C2C69E4F}\\
Shell mixing V16 & 623261 & \texttt{BF138E63BA826BC0}\\
Parallel YM V5 & 54174 & \texttt{306F1BB84CFA8A8D}\\
SM source selection V31 & 1095849 & \texttt{34312755DFB63E2E}\\
Native stationarity V34 & 1175242 & \texttt{08BC4B7DE876D011}\\
Exact-action Einstein V24 & 836365 & \texttt{A9B80D1A2689489C}\\
Suggestions & 11319 & \texttt{4F00C5864C9B7A90}\\
\bottomrule
\end{tabular}
\end{center}

SM V72 concerns a fixed-time auxiliary many-copy limit,
not a cutoff-uniform physical finite-species continuum.
NS V26 improves specific rank-one kernel norms, not
global regularity. Shell V16 keeps the measured local
one-loop factors together, but its flat-source result
has no joint nonperturbative volume/coupling remainder.
Parallel YM V5 rules out its extensive-global-charge
tightness target. Source-selection V31 still requires
word-native typing. Stationarity V34 has not evaluated
the 330 moving-fibre derivative entries, and Einstein
V24 proves a corrected first-order tangent rather than
a nonlinear upstream root. None supplies the missing
uniform physical transfer estimate.

Their useful methodological constraint remains to keep
the actual operator, carrier, and normalization together.
Accordingly the next step addresses V79's relative-error
problem directly on its native Wilson transfer.
All documents remain mathematical source material,
not instructions to merge files or adopt a proposed
conclusion. The next periodic checkpoint is V85.

\subsection{Exact spatial-Gibbs coordinates, not a vacuum identification}

Use the fixed-regulator notation of V79:
\[
 \overline T=M_b C M_b,\qquad
 b=e^{-\gamma S_{\rm sp}/2},\qquad
 C=\bigotimes_{e=1}^E(C_\gamma/c_0),\qquad
 r_n=\frac{I_{n+1}(\gamma)}{I_1(\gamma)} .
\]
All operators in the physical problem are restricted
to the gauge-invariant subspace.
Define the known spatial normalizer and probability law
\begin{equation}
 Z_b=\int b^2\,dm>0,\qquad
 d\nu_b=Z_b^{-1}b^2\,dm,\qquad
 \mathcal U_b f=Z_b^{-1/2}bf .
 \label{eq:f15v80:carrier}
\end{equation}
Then \(\mathcal U_b:L^2(\nu_b)^{\rm gauge}
\to L^2(m)^{\rm gauge}\) is unitary. Define
\[
 \mathcal K_b=\mathcal U_b^*(\overline T/Z_b)\mathcal U_b,
 \qquad
 \mathcal K_{b,N}=\mathcal U_b^*(T_N/Z_b)\mathcal U_b .
\]
These are not ground-state transforms. In particular,
\(\nu_b\) is not being identified with the physical
vacuum probability law.

If \(c(U,V)\) is the normalized product convolution
kernel of \(C\), direct substitution gives
\begin{equation}
       (\mathcal K_b f)(U)=\int c(U,V)f(V)\,d\nu_b(V).
 \label{eq:f15v80:kernel}
\end{equation}
A gauge average of this kernel gives the same operator
on invariant functions. The kernel is symmetric and
positive, and the operator is positive and compact,
but its row integral need not equal one.
No division by that row integral is made.
The eigenvalue ratios of \(\mathcal K_b\),
\(\overline T\), and the original physical transfer
\(\mathsf T\) agree exactly.

Since \(C_0=|1\rangle\langle1|\) in Haar space,
\[
       \mathcal K_{b,0}=|1\rangle\langle1|
           \quad\hbox{in }L^2(\nu_b).
\]
Consequently \(\mathcal K_b\ge\mathcal K_{b,N}
\ge|1\rangle\langle1|\), and their largest
eigenvalues are at least one. The small factor
\(Z_b\) has been accounted for exactly, not
estimated by its potentially exponentially small
lower bound.

\subsection{A magnetic-weight-uniform relative trace remainder}

Introduce the one-link kinetic traces
\begin{equation}
 s=\sum_{n\ge0}(n+1)^2r_n=\frac1{c_0},\qquad
 s_N=\sum_{n=0}^N(n+1)^2r_n,\qquad
 \tau_N=s^E-s_N^E .
 \label{eq:f15v80:trace-scalars}
\end{equation}
The representation multiplicity \((n+1)^2\)
is essential here. The identity \(s=1/c_0\)
follows by evaluating the absolutely convergent
character kernel at the identity, where
\(w_\gamma(I)=1\). Absolute convergence also
follows directly from V79's factorial majorant.

\begin{theorem}[Relative native-transfer trace tail]
\label{f15v80:relative}
At every finite regulator and for every \(N\ge0\),
\begin{equation}
 \begin{split}
 0&\le\mathcal K_b-\mathcal K_{b,N},\\
 \|\mathcal K_b-\mathcal K_{b,N}\|
 &\le
 \|\mathcal K_b-\mathcal K_{b,N}\|_{\mathfrak S_1}
 \le\tau_N,\\
 \frac{\|\overline T-T_N\|}
      {\lambda_0(\overline T)}
 &\le\tau_N .
 \end{split}
 \label{eq:f15v80:relative}
\end{equation}
The bound is uniform over bounded strictly positive
gauge-invariant magnetic multipliers \(b\) on this
same finite slice. Its dependence on the kinetic
coupling and link count is retained in \(\tau_N\).
\end{theorem}

\begin{proof}
First work on the full, ungauged product Haar space.
On one link use the normalized matrix coefficients
\(\sqrt{n+1}\,D^{(n)}_{ab}(U)\).
Unitarity of \(D^{(n)}(U)\) gives
\[
 \sum_{a,b}
    \left|\sqrt{n+1}\,D^{(n)}_{ab}(U)\right|^2
       =(n+1)^2
\]
at every \(U\). Hence the diagonal sum of the
product kinetic feature kernel is the constant
\(s^E\); restricting every link label to \(n\le N\)
gives the constant \(s_N^E\).

Expand the positive operator \(C-C_N\) in that
orthonormal product basis. Positivity permits
Tonelli's theorem, and the trace of a positive
rank-one operator is the squared norm of its
generating vector. Thus
\[
 \begin{split}
 \Tr_{\rm full}\bigl(M_b(C-C_N)M_b\bigr)
 &=\int b(U)^2
       \sum_{\alpha\ {\rm omitted}}
             r_\alpha|\phi_\alpha(U)|^2\,dm(U)\\
 &=Z_b(s^E-s_N^E)=Z_b\tau_N .
 \end{split}
\]
This is an exact full-space trace identity.
It also proves trace-class convergence of the
positive series.

The physical projection commutes with \(b,C,C_N\).
The trace of the positive operator restricted
to that subspace is no larger than its full
trace. It need not equal the full trace.
Divide by \(Z_b\) and conjugate by
\(\mathcal U_b\). Since the trace norm of a
positive operator equals its trace and dominates
its operator norm, the first two bounds follow.
Finally
\(\lambda_0(\overline T)\ge\lambda_0(T_0)=Z_b\),
which proves the last bound. The argument used
no estimate on \(Z_b\), no vacuum wavefunction,
and no restriction on the particular magnetic
weight beyond the stated scope.
\end{proof}

This is different from V79's link-count-free
\emph{absolute operator-norm} tail. That bound
remains valid. The new relative bound pays a
kinetic representation-counting cost:
\(s>1\) when \(\gamma>0\), and the factor \(s^E\)
can be very large. Removing an unknown magnetic
denominator does not remove volume dependence.

\begin{proposition}[Explicit selection of a relative cutoff]
\label{f15v80:cutoff}
Put \(x=\gamma/2\) and
\[
 q_N=x\,\frac{N+3}{(N+2)^2}.
\]
If \(q_N<1\), then
\begin{equation}
 \begin{split}
 s-s_N&\le
      \frac{(N+2)x^{N+1}}{(N+1)!(1-q_N)}
       =:\Delta_N^{\rm maj},\\
 \tau_N&\le E\,s^{E-1}\Delta_N^{\rm maj}.
 \end{split}
 \label{eq:f15v80:cutoff}
\end{equation}
For fixed \(E,\gamma\), the right side tends to
zero as \(N\to\infty\). A desired positive
relative tolerance can therefore be met without
knowing \(Z_b\) or any transfer eigenvalue.
\end{proposition}

\begin{proof}
V79's termwise positive-series comparison gives
\[
 (n+1)^2r_n
      \le(n+1)\frac{x^n}{n!}=:t_n .
\]
The ratio \(t_{n+1}/t_n=x(n+2)/(n+1)^2\)
decreases in \(n\). The first omitted term is
\(t_{N+1}=(N+2)x^{N+1}/(N+1)!\).
A geometric majorant with ratio \(q_N\) proves
the first inequality. The identity
\[
 s^E-s_N^E=(s-s_N)
             \sum_{j=0}^{E-1}s^{E-1-j}s_N^j
\]
proves the second. Factorial decay finishes
the proof.
\end{proof}

The scalar \(s=e^\gamma\gamma/(2I_1(\gamma))\)
can be enclosed by the positive Bessel series
already used in V79, consistent with
\href{https://dlmf.nist.gov/10.25.E2}{NIST DLMF 10.25.2}.
A more conservative bound requiring no cancellation is
\(s\le(1+x)e^x\), obtained by summing \(t_n\).
For computably specified regulator parameters,
any certified upper bound for \(s\) makes
\eqref{eq:f15v80:cutoff} a finite cutoff-selection
procedure. No efficiency or manageable matrix
dimension is asserted.

\subsection{A spectral-ratio enclosure with no unknown leading scale}

In V79's physical basis let
\(\psi_\alpha=d_\alpha\phi_\alpha\), and write
\begin{equation}
 (\Gamma_N)_{\alpha\beta}
     =\mathbb E_{\nu_b}
             [\,\overline{\psi_\alpha}\psi_\beta\,]
     =(G_N)_{\alpha\beta}/Z_b .
 \label{eq:f15v80:gram}
\end{equation}
The nonzero spectrum of \(\Gamma_N\) is that
of \(\mathcal K_{b,N}\). This follows either
from V79's factorization or by using the
finite-column operator with columns
\(\psi_\alpha\) in \(L^2(\nu_b)^{\rm gauge}\).
The trivial representation gives
\((\Gamma_N)_{00}=1\), so its top eigenvalue
\(\nu_{0,N}\) is at least one.
Include the zero complement when defining
its second eigenvalue \(\nu_{1,N}\), and put
\[
 \rho=\frac{\lambda_1(\mathsf T)}{\lambda_0(\mathsf T)},
 \qquad
 \rho_N=\frac{\nu_{1,N}}{\nu_{0,N}} .
\]

\begin{proposition}[Full-space ratio error]
\label{f15v80:ratio}
The native spectral ratio satisfies
\begin{equation}
 \frac{\nu_{1,N}}{\nu_{0,N}+\tau_N}
   \le\rho\le
 \frac{\nu_{1,N}+\tau_N}{\nu_{0,N}},
 \qquad
 |\rho-\rho_N|\le\frac{\tau_N}{\nu_{0,N}}
                         \le\tau_N .
 \label{eq:f15v80:ratio}
\end{equation}
In particular \(\rho_N+\tau_N<1\) is a
vacuum-independent sufficient certificate, with
physical gap at least
\[
             -a_t^{-1}\log(\rho_N+\tau_N).
\]
\end{proposition}

\begin{proof}
The positive remainder in
\eqref{eq:f15v80:relative} and min--max imply
\[
 \nu_{i,N}\le\lambda_i(\mathcal K_b)
                          \le\nu_{i,N}+\tau_N,
 \qquad i=0,1 .
\]
Apply the lower numerator with the upper
denominator and vice versa. The upper-ratio
excess over \(\rho_N\) is at most
\(\tau_N/\nu_{0,N}\). The lower-ratio deficit is
\[
 \frac{\nu_{1,N}\tau_N}
        {\nu_{0,N}(\nu_{0,N}+\tau_N)}
       \le\frac{\tau_N}{\nu_{0,N}},
\]
because \(0\le\nu_{1,N}\le\nu_{0,N}\).
The final statement is the inherited
logarithmic spectral conversion.
\end{proof}

For approximate entries, the matrix error must
still be paid. If a Hermitian approximation to
\(\Gamma_N\) has operator error at most \(z_N\),
a top-eigenvalue lower bound \(L_N\) and a
nonnegative second-eigenvalue upper bound \(U_N\)
give the valid inputs
\[
 \ell_N=\max\{1,L_N-z_N\},\qquad
 u_N=U_N+z_N,\qquad
 \rho\le(u_N+\tau_N)/\ell_N .
\]
The entry error here is for the \emph{normalized}
matrix \(\Gamma_N\), not the unnormalized matrix
\(G_N\). V79's polynomial integration method can
enclose numerator and positive normalizer, but
a small unnormalized entry error cannot simply
be reused after division by \(Z_b\).

\begin{proposition}[Scalar physical-rate approximation]
\label{f15v80:clock}
On a specified regulator sequence suppose
\(\rho_j\to1\), the pre-existing physical time
steps satisfy \(a_{t,j}\to0\), and cutoffs are
chosen with \(\tau_{N_j,j}/a_{t,j}\to0\).
Then, eventually, \(\rho_{N_j,j}>0\) and
\begin{equation}
 \left|
   -\frac{\log\rho_j}{a_{t,j}}
      +\frac{\log\rho_{N_j,j}}{a_{t,j}}
 \right|
 \le\frac{2\tau_{N_j,j}}{a_{t,j}}\longrightarrow0 .
 \label{eq:f15v80:clock}
\end{equation}
\end{proposition}

\begin{proof}
The hypotheses and \eqref{eq:f15v80:ratio} give
\(\rho_{N_j,j}\to1\). Both ratios are eventually
in \([1/2,1]\), on which the derivative of
\(\log\) has absolute value at most two.
Apply the mean-value theorem.
\end{proof}

V63 supplies the ratio convergence only on
its stated sequence; it is not silently
extended to other couplings or aspect ratios.
The physical clock is an independent input,
not something chosen to make
\eqref{eq:f15v80:clock} favorable. Proposition
\ref{f15v80:cutoff} allows the required scalar
tail tolerance to be prescribed on any given
sequence, at an explicitly retained dimensional
cost. Proposition~\ref{f15v80:clock} does not
construct continuum observables, identify a
limiting Hilbert space, or prove a positive
limit for either rate.

\subsection{What the exact covariance formulation does and does not give}

Let \(W_N^\nu\) have columns \(\psi_\alpha\)
in \(L^2(\nu_b)^{\rm gauge}\), and put
\[
 m_N=(W_N^\nu)^*1,\qquad
 a_N=\|m_N\|^2,\qquad
 \Sigma_N=\Gamma_N-m_Nm_N^*,\qquad
 Q_1=I-|1\rangle\langle1|.
\]
Thus \(\Sigma_N\) is the covariance matrix of
the weighted physical Peter--Weyl features
under the exact spatial Wilson law, and
\(a_N\ge1\). Factorization gives
\[
 \mathcal K_{b,N}=W_N^\nu(W_N^\nu)^*,\qquad
 \Sigma_N=(W_N^\nu)^*Q_1W_N^\nu .
\]
Therefore the nonzero spectrum of
\(Q_1\mathcal K_{b,N}Q_1\) is that of
\(\Sigma_N\).

\begin{proposition}[A sufficient mean-centered test]
\label{f15v80:covariance}
One has
\[
 \lambda_0(\mathcal K_b)\ge a_N,\qquad
 \lambda_1(\mathcal K_b)\le\|\Sigma_N\|+\tau_N .
\]
In particular \(\|\Sigma_N\|+\tau_N<a_N\)
is sufficient for a positive finite-regulator
gap. These are native transfer estimates,
not a Poincare inequality for an auxiliary
spatial diffusion.
\end{proposition}

\begin{proof}
Use the Rayleigh quotient at the normalized
constant function for the lower bound.
For the upper bound, min--max on the entire
orthogonal complement of that constant gives
\(\lambda_1(\mathcal K_b)\le\|Q_1\mathcal K_bQ_1\|\).
Compress the positive configuration-space
remainder, whose norm is at most \(\tau_N\),
and use the preceding spectral factorization.
\end{proof}

The common space in this proof is the
configuration Hilbert space \(L^2(\nu_b)\).
A finite feature covariance matrix padded by
zeros need not lie below the larger feature
covariance matrix: their cross blocks can make
the difference indefinite. Such a padded
matrix difference is not the positive
remainder used in the proof.

The complete kinetic weights \(d_\alpha\)
remain part of these covariances. This is
consistent with f6 V54's static-law
non-identifiability theorem: a spatial law
alone does not determine temporal dynamics.
Here the original temporal kernel is fixed
and retained exactly. Nor do the available
single-action-probe or retained-layer
covariance estimates automatically bound
this complete weighted physical matrix.

\subsection{Two exact warnings against a spurious simplification}

First, row normalization changes the target.
On two points with uniform reference Haar
weights let
\[
 c=\begin{pmatrix}3/2&1/2\\1/2&3/2\end{pmatrix},
 \qquad \nu=(1/4,3/4).
\]
The Haar convolution associated with \(c\)
has eigenvalues \(1,1/2\). The probability
\(\nu\) is obtained from a positive magnetic
weight \(b^2=(1/3,1)\). In Gibbs coordinates,
\[
 \mathcal K=c\,\operatorname{diag}(\nu)
   =\begin{pmatrix}3/8&3/8\\1/8&9/8\end{pmatrix}.
\]
It is self-adjoint in \(L^2(\nu)\), with
positive eigenvalues
\[
       \lambda_\pm=(3\pm\sqrt3)/4,
       \qquad \lambda_-/\lambda_+=2-\sqrt3.
\]
Its row sums are \(3/4,5/4\).
Dividing the rows by these sums instead gives
\[
 P_{\rm row}=
       \begin{pmatrix}1/2&1/2\\1/10&9/10\end{pmatrix},
\]
whose nontrivial eigenvalue is \(2/5\).
Since \(2/5\ne2-\sqrt3\), this Markov kernel
has a different gap. It is not the native
vacuum Doob transform.

Second, even without row normalization,
fixing the constant trial function can fail
badly with volume.

\begin{theorem}[Tensor-volume obstruction to the fixed-mean test]
\label{f15v80:overlap}
Let \(K\) be a compact positive operator on
\(L^2(\nu)\), with simple top eigenvalue
\(\lambda_0>\lambda_1>0\) and normalized
top eigenvector \(\omega\). Suppose the
normalized constant \(1\) is not a top
eigenvector and has nonzero overlap with
\(\omega\). Write
\[
 a=\langle1,K1\rangle<\lambda_0,\qquad
 q=|\langle1,\omega\rangle|^2\in(0,1).
\]
For \(K^{\otimes M}\) and
\(Q_M=I-|1^{\otimes M}\rangle\langle1^{\otimes M}|\),
\begin{equation}
 \langle1^{\otimes M},K^{\otimes M}1^{\otimes M}\rangle=a^M,
 \qquad
 \|Q_MK^{\otimes M}Q_M\|
       \ge\lambda_0^M(1-q^M).
 \label{eq:f15v80:overlap}
\end{equation}
Thus the untruncated mean-centered sufficient
test fails for all sufficiently large \(M\),
although the true normalized spectral ratio
of \(K^{\otimes M}\) stays equal to
\(\lambda_1/\lambda_0<1\).
\end{theorem}

\begin{proof}
The first identity is tensor factorization.
Positivity and the top spectral projection give
\[
 K^{\otimes M}\ge\lambda_0^M
       |\omega^{\otimes M}\rangle
       \langle\omega^{\otimes M}|.
\]
Compress by \(Q_M\); the resulting rank-one
operator has norm
\(\lambda_0^M\|Q_M\omega^{\otimes M}\|^2
 =\lambda_0^M(1-q^M)\).
Its ratio to \(a^M\) tends to infinity.
The second largest tensor eigenvalue is
\(\lambda_1\lambda_0^{M-1}\), proving the
last assertion.
\end{proof}

The preceding two-point kernel gives an explicit
finite-state example of this obstruction:
\[
a=9/8,\qquad \lambda_0=(3+\sqrt3)/4,\qquad
q=(2+\sqrt3)/4.
\]
The overlap formula follows from
\(a=\lambda_0q+\lambda_1(1-q)\).
For \(M=20\), already
\(\lambda_0^M(1-q^M)>a^M\).
For example the rational bounds
\(1.732<\sqrt3<1.733\) prove this strict
inequality without floating-point eigenvalues.
The true product gap remains
\(a_t^{-1}\log(2+\sqrt3)\), independent of \(M\).

This is a finite-state test of the proposed
method, not a Yang--Mills counterexample.
Its consequence is methodological and exact:
the spatial-mean hyperplane cannot be assumed
to remain a useful sufficient certificate at
large volume. Improving the high-mode cutoff
does not repair that loss of overlap, since
it persists with no truncation at all.
The full top-two matrix criterion does not
have this particular defect; it recovers the
correct product ratio in the same example.

\subsection{Revised direction and the remaining physical estimate}

The relative omitted-space error can now be
specified using \(E,\gamma,N\), without an
unknown vacuum eigenvalue or an unevaluated
magnetic normalizer in the bound. This resolves
one concrete error-accounting obstacle raised
in V79, with its representation-counting cost
explicit. It does not evaluate the low
physical spectrum.

The next analysis should therefore keep the
full normalized matrix \(\Gamma_N\), or use
regulator-adapted physical trial vectors and
complete channels whose overlap loss is
controlled. A fixed spatial-mean covariance
bound must not become a replacement goal.
The exact normalized expectations in
\eqref{eq:f15v80:gram} still need estimates
sharp enough to separate the first two
eigenvalues at the scale \(a_t\); neither
a trace bound nor a single-observable
susceptibility estimate does that.

Specifically, after all finite-entry errors
are paid, one still needs a regulator sequence
of certified ratios bounded by
\(e^{-m_*a_{t,j}}\) for one \(m_*>0\).
The relative trace tail can be chosen smaller
than that scale, but the retained-matrix
separation has not been proved.
A total limiting probe family, nontriviality,
tightness, compatible limiting dynamics,
locality and relativistic reconstruction
remain further obligations. No new companion
result closes them.

\subsection{Diagnostics and append-only preservation}

The script
\begin{center}\small
\path{scripts/verify_ym_f15_v80_relative_transfer_trace.py}
\end{center}
reruns V79 and its inherited checks.
It checks 15 dimension-weighted Bessel trace
enclosures, 45 relative product-tail budgets,
18 prescribed-error cutoff selections and
12 multiplicity-retaining tensor trace identities.
Three magnetic-weight fixtures verify
the exact Gibbs conjugation and normalized
trace identity, and three verify the weighted
feature mean/covariance relations. Two
counterexamples distinguish covariance-block
padding from the valid positive
configuration-space tail. A nontrivial gauge
projection checks the physical/full-trace
distinction. The row-normalization counterexample,
231 spectral-ratio error brackets, and three
tensor-volume overlap failures are checked
in exact rational arithmetic.
These are diagnostics of the written proofs,
not numerical certificates for a large
physical Yang--Mills matrix.

The script has 9,144 bytes and SHA--256
\begin{center}\scriptsize\ttfamily
7F6AAA751BE1B3B99F569B22A5EBDD378A858ACEF8647535B96EF6E3EDD1F903.
\end{center}
The V79 predecessor has 1,788,476 bytes and SHA--256
\begin{center}\scriptsize\ttfamily
EEC720C468B3D49F2976011C88256A902833E3CA47172D099361F35F270CAA8E.
\end{center}
Removing this final section and restoring the
title version recovers V79 exactly.
All 139 previously protected files were unchanged
before this append. V80 is the sole active main
note; archives and companions remain preserved.
Only \path{.tex} sources belong in
\path{latex_notes}; build products go to
\path{artifacts/ym_f15_v80_texcheck/}.

\begin{firewall}
V80 completes the companion checkpoint, proves
a magnetic-weight-uniform relative transfer
trace bound and a full spectral-ratio error
estimate, and identifies the tensor-volume
failure of an unadapted spatial-mean
certificate. It does not establish the
required uniform retained-matrix spectral
separation. The interacting continuum and
full Yang--Mills mass gap remain unproved.
The original goal remains active.
\end{firewall}


\section{V81: antisymmetric transfer and gauge-boundary returns}
\label{f15v81:section}

\subsection{Context and nonduplication}

V80 gave relative full-space truncation bounds but
also showed that centering at a fixed spatial mean
can fail at large volume even for a gapped product
system. The next estimate must retain the actual
first two transfer eigenvalues without relying on
that reference vector. This section, dated
16 September 2026, tests a tensor-stable
antisymmetric two-copy formulation.

The exterior-square identities used below are
standard linear algebra; their proofs fix the
normalization and omitted-space accounting.
The programme-specific calculation is what
happens when one tries to localize the two
copies by exchanging only some lattice links.
We compute the resulting Gauss-law defect,
the exact return on physical cycle characters,
and the actual Wilson magnetic boundary factor.
The defect is not small merely because both
replicas have zero plaquette action.

The archive search found extensive Schur/Feshbach
machinery, f7 V63--V74's glued-replica density
estimates, and f12 V54's endpoint cross-ratios.
Those results are not being re-proved or
identified with the present spectral exterior
square. In particular a glued marginal-density
norm is not automatically the operator norm
that occurs below. The last companion checkpoint
was V80; the next remains V85.

\subsection{The antisymmetric sector isolates the first two eigenvalues}

Let \(K\) be a compact positive operator on
a Hilbert space \(\mathcal H\), with
\(\lambda_0(K)>0\), and list its eigenvalues
in decreasing order with multiplicity.
Write \(F(f\otimes g)=g\otimes f\), let
\(P_-=(I-F)/2\), and identify
\(\bigwedge^2\mathcal H=\operatorname{Ran}P_-\).
Define
\[
 K^{\wedge2}=(K\otimes K)|_{\bigwedge^2\mathcal H}.
\]
The symmetric-square operator \(K^{\odot2}\)
is the analogous restriction to \(\ker(F-I)\).
If necessary, second eigenvalues and zero
exterior spaces are interpreted as zero.

\begin{proposition}[Reference-free two-copy ratio]
\label{f15v81:exterior}
One has
\begin{equation}
 \|K^{\wedge2}\|=\lambda_0(K)\lambda_1(K),
 \qquad
 \|K^{\odot2}\|=\lambda_0(K)^2,
 \qquad
 \frac{\lambda_1(K)}{\lambda_0(K)}
       =\frac{\|K^{\wedge2}\|}{\|K\|^2}.
 \label{eq:f15v81:exterior}
\end{equation}
For compact positive \(K,L\) with positive
top eigenvalues,
\begin{equation}
 \frac{\lambda_1(K\otimes L)}{\lambda_0(K\otimes L)}
  =\max\left\{
     \frac{\lambda_1(K)}{\lambda_0(K)},
     \frac{\lambda_1(L)}{\lambda_0(L)}\right\}.
 \label{eq:f15v81:tensor}
\end{equation}
Thus this target does not suffer V80's fixed-mean
overlap loss under independent tensor products.
\end{proposition}

\begin{proof}
In an orthonormal eigenbasis \(e_i\) of \(K\),
the normalized vectors
\((e_i\otimes e_j-e_j\otimes e_i)/\sqrt2\),
\(i<j\), have eigenvalues \(\lambda_i\lambda_j\).
The symmetric sector includes \(e_0\otimes e_0\).
These facts prove \eqref{eq:f15v81:exterior},
including zero eigenspaces.

Regroup two copies of \(\mathcal H\otimes\mathcal G\)
as \(\mathcal H^{\otimes2}\otimes\mathcal G^{\otimes2}\).
The global flip becomes \(F_{\mathcal H}F_{\mathcal G}\).
Its odd sector is the orthogonal sum of
the odd-even and even-odd sectors. Hence
\[
 (K\otimes L)^{\wedge2}
 \simeq(K^{\wedge2}\otimes L^{\odot2})
       \oplus(K^{\odot2}\otimes L^{\wedge2}).
\]
Take norms and use the first identities.
\end{proof}

This applies directly to V80's physical
\(\mathcal K_b\), whose spectral ratio equals
that of the original Wilson transfer.
It also applies to genuinely disconnected
physical graph components when their magnetic
weights factor. It does not assert that the
gauge-invariant Hilbert space of a connected
lattice factors over its individual links.

If \(k\) is a physical gauge-averaged integral
kernel of \(\mathcal K_b\) in \(L^2(\nu_b)\),
the extended odd operator \(P_-(\mathcal K_b\otimes
\mathcal K_b)P_-\) has kernel
\begin{equation}
 \frac12\left[
  k(U_1,V_1)k(U_2,V_2)
       -k(U_1,V_2)k(U_2,V_1)\right].
 \label{eq:f15v81:detkernel}
\end{equation}
The factor \(1/2\) comes from \(P_-\).
The kernel changes sign on exchanging its two
output arguments. Its operator positivity
does not make it a positive Markov kernel.
Discarding the second term would discard
the spectral cancellation being measured.

\subsection{A finite physical compound matrix with the V80 tail retained}

Use V80's normalized finite physical Gram
matrix \(\Gamma_N\), with top eigenvalue
\(\nu_{0,N}\ge1\), second eigenvalue
\(\nu_{1,N}\), and full ratio remainder
\(\tau_N\). For ordered pairs \(\alpha<\beta\)
and \(\eta<\zeta\), its second compound matrix is
\begin{equation}
 (\Gamma_N^{\wedge2})_{\alpha\beta,\eta\zeta}
   =(\Gamma_N)_{\alpha\eta}(\Gamma_N)_{\beta\zeta}
      -(\Gamma_N)_{\alpha\zeta}(\Gamma_N)_{\beta\eta}.
 \label{eq:f15v81:compound}
\end{equation}
There is no additional factor \(1/2\) in this
orthonormal exterior-basis matrix.

If \(W_N^\nu\) is the finite-column map used
in V80, functoriality of the exterior square
or direct expansion of its two-column
determinants gives
\[
 \mathcal K_{b,N}^{\wedge2}
   =(W_N^\nu)^{\wedge2}((W_N^\nu)^{\wedge2})^*,
 \qquad
 \Gamma_N^{\wedge2}
   =((W_N^\nu)^{\wedge2})^*(W_N^\nu)^{\wedge2}.
\]
Thus the nonzero spectra agree, and
\(\|\Gamma_N^{\wedge2}\|=\nu_{0,N}\nu_{1,N}\).

\begin{proposition}[Complete antisymmetric certificate]
\label{f15v81:certificate}
Suppose certified bounds satisfy
\(\nu_{0,N}\ge\ell_N>0\) and
\(\|\Gamma_N^{\wedge2}\|\le\omega_N\). Then
\begin{equation}
 \frac{\lambda_1(\mathsf T)}{\lambda_0(\mathsf T)}
       \le\frac{\omega_N}{\ell_N^2}+\tau_N .
 \label{eq:f15v81:certificate}
\end{equation}
If the right side is strictly between zero
and one, the corresponding physical gap is
at least its negative logarithm divided by
the independently specified time step.
\end{proposition}

\begin{proof}
The finite ratio is
\(\nu_{1,N}/\nu_{0,N}
 =\|\Gamma_N^{\wedge2}\|/\nu_{0,N}^2\).
Apply Proposition~\ref{f15v80:ratio}.
\end{proof}

For a Hermitian approximate matrix \(A_N\)
with \(\|\Gamma_N-A_N\|\le z_N\),
\begin{equation}
 \|\Gamma_N^{\wedge2}-A_N^{\wedge2}\|
      \le z_N(2\|A_N\|+z_N).
 \label{eq:f15v81:compound-error}
\end{equation}
Indeed, the difference of the tensor squares
is
\((\Gamma_N-A_N)\otimes\Gamma_N+
 A_N\otimes(\Gamma_N-A_N)\); restrict to the
odd subspace and use
\(\|\Gamma_N\|\le\|A_N\|+z_N\).
If \(A_N\) is not positive, its full operator
norm, not merely its largest signed eigenvalue,
must be used. This accounts for matrix errors;
it does not supply a small physical value
of \(\omega_N/\ell_N^2\).

This reformulation is an alternative target
for local analysis, not automatically a
cheaper numerical method. The compound
matrix is larger, and an unproved upper
bound on its norm would simply restate
the missing spectral separation.

\subsection{Put partial swaps on the correct carrier}

For the locality calculation work on the
full two-copy Haar space
\[
 \mathcal H_{\rm dbl}=L^2(G^E\times G^E,dm\,dm),
       \qquad G=\SU(2).
\]
For a link set \(A\), let \(s_A(U,V)=(U^A,V^A)\)
exchange \(U_e,V_e\) for \(e\in A\), and define
the unitary self-adjoint pullback
\(S_Af=f\circ s_A\).
The global flip is \(S_E\).
Every \(S_A\) commutes with \(S_E\) and with
the free doubled kinetic operator \(C\otimes C\).

Let \(P_v^{(r)}\) be the Haar gauge projection
at vertex \(v\) in replica \(r\), and put
\[
 \Pi_v=P_v^{(1)}P_v^{(2)},\qquad
 \mathbb P=\prod_v\Pi_v,\qquad
 B(U,V)=b(U)b(V).
\]
Vertex gauge actions commute at different
vertices because left and right actions on
a shared link commute. Thus these projections
commute. The full extension of the doubled
physical transfer is
\begin{equation}
       \mathbb P M_B(C\otimes C)M_B\mathbb P .
 \label{eq:f15v81:doubled}
\end{equation}
Its physical odd restriction is the exterior
square of \(\overline T\).
Both \(\mathbb P\) and \(M_B\) commute with
the global flip. They need not commute
with a partial swap.

Define the vertex boundary
\[
 \partial_{\rm v}A
  =\{v:\text{some incident links lie in }A
                      \text{ and some outside }A\}.
\]
Let \(\mathbb P_\partial\) be the product of
\(\Pi_v\) over this boundary, and
\(\mathbb P_{\rm int}\) the product over all
remaining vertices.

\begin{proposition}[Exact boundary localization of the Gauss defect]
\label{f15v81:gauss-boundary}
One has
\begin{equation}
 \begin{split}
 [\mathbb P,S_A]&=[\mathbb P_\partial,S_A]\mathbb P_{\rm int},\\
 \mathbb P S_A\mathbb P
   &=\mathbb P_\partial S_A
                    \mathbb P_\partial\mathbb P_{\rm int}.
 \end{split}
 \label{eq:f15v81:gauss-boundary}
\end{equation}
\end{proposition}

\begin{proof}
If no incident link at \(v\) is swapped,
\(S_A\) commutes with both replica gauge
actions there. If every incident link is
swapped, it interchanges the two replica
gauge actions and therefore commutes with
their product Haar projection \(\Pi_v\).
Consequently \(S_A\) commutes with
\(\mathbb P_{\rm int}\). Use
\(\mathbb P=\mathbb P_\partial\mathbb P_{\rm int}\)
and the commutativity of the projections.
\end{proof}

This identity isolates the boundary carrier;
it does not estimate the defect by a small
constant. The next calculation shows that
its norm can be exactly one on a genuine
physical excitation.

\subsection{A maximal gauge defect on a physical plaquette mode}

Let \(p\) be a simple plaquette, or more generally
a simple cycle, and let
\(\chi_1(U_p)=\Tr_{\mathbb C^2}U_p\).
In the notation \(n=2j\), this is the \(n=1\)
character. Product Haar integration gives
\(\int\chi_1(U_p)\,dm=0\) and
\(\int\chi_1(U_p)^2\,dm=1\).
Indeed its holonomy is Haar distributed,
and these are the constant and quadratic
unit-quaternion moments already used in V79.
Hence
\[
 \Phi_p(U,V)=
        \frac{\chi_1(U_p)-\chi_1(V_p)}{\sqrt2}
\]
is a normalized physical odd vector.

\begin{theorem}[A proper cycle cut is not a physical involution]
\label{f15v81:maximal-defect}
If \(A\) contains a nonempty proper subset
of the links of \(p\), then
\begin{equation}
       \mathbb P S_A\Phi_p=0,\qquad
       \|[\mathbb P,S_A]\|=1.
 \label{eq:f15v81:maximal-defect}
\end{equation}
The norm is still one when restricted to
the global odd Haar sector.
On the physical space the two effects
\[
       E_A^\pm=
          \tfrac12(I_{\rm phys}\pm\mathbb P S_A|_{\rm phys})
\]
are not orthogonal projections:
\(E_A^\pm\Phi_p=\Phi_p/2\), and their
idempotence defect on this vector is
\(-\Phi_p/4\).
\end{theorem}

\begin{proof}
Along the cycle there is a vertex where one
of its two incident cycle links is in \(A\)
and the other is not. In each mixed holonomy
appearing in \(S_A\Phi_p\), exactly one of
these two links comes from the first replica.
Apply the first-replica gauge transformation
\(-I\) at that vertex. Its central sign
changes each mixed trace by \(-1\), so
\(S_A\Phi_p\) changes sign. Haar gauge
averaging annihilates such a function:
the average is invariant under that same
central transformation. Thus
\(\mathbb P S_A\Phi_p=0\).

Since \(\mathbb P\Phi_p=\Phi_p\),
\([\mathbb P,S_A]\Phi_p=-S_A\Phi_p\), of
norm one. For an orthogonal projection
\(P\) and a self-adjoint unitary \(S\),
the commutator has off-diagonal blocks
\(PS(I-P)\) and \(-(I-P)SP\); its norm
is at most \(\|S\|=1\).
This proves equality. Both \(\mathbb P\)
and \(S_A\) commute with the global flip,
and the test vector is odd, proving the
restricted assertion. The effect identities
follow by applying their definitions.
\end{proof}

For a four-link plaquette the doubled free
kinetic operator also satisfies
\((C\otimes C)\Phi_p=r_1^4\Phi_p\).
Each fundamental matrix coefficient on its
four distinct links acquires one factor
\(r_1\). This identifies a real physical
electric mode, not a gauge-variant test
outside the target space. Magnetic
multiplication is not being claimed to
preserve this eigenvector.

The theorem does not say that gauge
projection worsens the mass gap. It says
that compressed link swaps cannot be
treated as a family of local orthogonal
parity projections. In particular their
norm defect is not made small by increasing
the coupling or by refining the cutoff.

\subsection{Exact physical cycle-character returns}

The preceding annihilation can be extended
to all representation labels on one cycle.
Let \(p\) have \(L\) distinct links and
write
\[
 e_n(U)=\chi_n(U_p),\qquad d_n=n+1 .
\]
These normalized functions span the
cycle-character subspace, with every other
link in the trivial representation.
Let \(b_A(p)\) be the number of cycle
vertices where the membership of the
two incident cycle links in \(A\) differs.
For a nonempty proper cut it is positive
and even.

\begin{theorem}[Gauge-averaged swap on a cycle sector]
\label{f15v81:cycle-return}
For a nonempty proper \(A\cap p\), the
physical Haar compression satisfies
\begin{equation}
 \mathbb P S_A\mathbb P(e_m\otimes e_n)
   =
 \begin{cases}
   0,&m\ne n,\\
   d_n^{-b_A(p)}\,e_n\otimes e_n,&m=n .
 \end{cases}
 \label{eq:f15v81:cycle-return}
\end{equation}
In particular it vanishes on the entire
antisymmetric square of the cycle-character
subspace. For a connected proper arc,
\(b_A(p)=2\), so the equal-label return
is \(d_n^{-2}\), at most \(1/4\) for
\(n\ge1\).
\end{theorem}

\begin{proof}
Gauge averaging preserves the representation
label on each link. After a partial swap
of \(e_m\otimes e_n\), the first replica
has label \(n\) on swapped cycle links
and \(m\) on the others; the second
replica has the opposite pattern.
At a boundary cycle vertex, all other
incident links are trivial. The invariant
space between its two nontrivial half-link
representations is
\(\operatorname{Hom}_G(V_m,V_n)\).
It is zero when \(m\ne n\), by irreducibility
and inequivalence. This proves the first
case.

For \(m=n\), the physical subspace with
these labels on the cycle and trivial
labels off it is one dimensional in each
replica: every degree-two intertwiner is
the scalar identity. The compression must
therefore be a scalar on \(e_n\otimes e_n\).

To compute it, orient the cycle consistently,
inverting link variables where necessary.
In the orthonormal one-link basis
\(\sqrt{d_n}D^{(n)}_{ab}\), its character
is the normalized vector
\[
 e_n=d_n^{-L/2}
    \sum_{a_1,\ldots,a_L}
      |a_1,a_2\rangle\otimes|a_2,a_3\rangle
         \otimes\cdots\otimes|a_L,a_1\rangle .
\]
Here each \(a_i\) ranges over \(d_n\)
values; product Haar orthogonality
makes the displayed product vectors
orthonormal. Regrouping half-link
indices by vertices writes this vector
as a product of normalized maximally
entangled pairs. The link partition
\(A\cap p\) cuts precisely \(b_A(p)\)
of those pairs. Its Schmidt coefficients
are therefore all \(d_n^{-b_A(p)/2}\),
with multiplicity \(d_n^{b_A(p)}\).

For a normalized Schmidt expansion
\(f=\sum_j s_j u_j\otimes v_j\), direct
expansion shows that the expectation
of the partial swap in \(f\otimes f\)
is \(\sum_j s_j^4\).
Apply this identity to the preceding
coefficients: the scalar is
\(d_n^{b_A(p)}d_n^{-2b_A(p)}
 =d_n^{-b_A(p)}\).
This proves the second case.
Antisymmetric cycle vectors are spanned
by \(e_m\wedge e_n\), \(m<n\), and hence
are annihilated by the first case.
\end{proof}

The formula also holds for an embedded
simple cycle of the full lattice when
all off-cycle link labels are trivial.
It is a Haar Gauss-projection identity.
It is not an evaluation of the interacting
magnetic Gram matrix \(\Gamma_N\).
Magnetic multiplication can generate
other loops and nontrivial boundary
intertwiners, so this cycle subspace
need not remain invariant under the
full transfer. The factor \(1/4\)
must not be relabelled as a physical
Hamiltonian gap.

\subsection{The exact magnetic boundary factor and its sharp native size}

Define the plaquette boundary
\[
 \partial_{\rm p}A
    =\{p:0<|p\cap A|<|p|\}.
\]
Writing
\(s_p(U)=1-\tfrac12\Tr U_p\), the swap
density factor is exactly
\begin{equation}
 \begin{split}
 R_A(U,V)&=\frac{B(U^A,V^A)}{B(U,V)},\\
 \log R_A
   &=-\frac{\gamma}{2}
       \sum_{p\in\partial_{\rm p}A}
       \bigl[s_p(U^A)+s_p(V^A)-s_p(U)-s_p(V)\bigr].
 \end{split}
 \label{eq:f15v81:magnetic}
\end{equation}
Plaquettes with all or none of their
links swapped cancel from the sum.
Since each \(s_p\in[0,2]\),
\begin{equation}
       |\log R_A|\le2\gamma|\partial_{\rm p}A|.
 \label{eq:f15v81:magnetic-bound}
\end{equation}
This localizes the magnetic defect to
mixed plaquettes. It is not a
weak-coupling smallness estimate.

\begin{proposition}[Sharp defect starting from two pure gauges]
\label{f15v81:sharp-magnetic}
On a three-dimensional periodic spatial
lattice of side at least three, let
\(A=\{e\}\) consist of one link.
There are two pure-gauge configurations
with zero spatial Wilson action for
which
\begin{equation}
               \log R_A=-8\gamma
                   =-2\gamma|\partial_{\rm p}A|.
 \label{eq:f15v81:sharp-magnetic}
\end{equation}
The reverse swap attains \(+8\gamma\).
\end{proposition}

\begin{proof}
Let every link of \(U\) be \(I\).
Choose an endpoint \(x\) of \(e\).
Let \(V\) be its gauge transform by the
central value \(-I\) at \(x\), with
the identity at all other vertices.
Every link incident to \(x\) changes
sign, but all plaquette holonomies of
\(V\) remain \(I\).

Swap only \(e\). In \(U^A\) this changes
one link from \(I\) to \(-I\); in \(V^A\)
it reverses one of the pure-gauge signs.
Each of the four plaquettes containing
\(e\) now has holonomy \(-I\) in both
replicas. All other plaquettes remain
flat. Each affected plaquette contributes
\(2+2-0-0=4\) to the action difference
in \eqref{eq:f15v81:magnetic}.
The total difference is \(16\), proving
\(-8\gamma\). Involutivity reverses its
sign for the opposite direction.
\end{proof}

The starting configurations represent the
same physical gauge orbit. This is not
a tunnelling barrier between inequivalent
vacua. It shows why an unframed link
exchange can have a large raw magnetic
factor even on perfectly flat input
replicas. Small curvature of each copy
alone cannot justify treating that
factor as a small perturbation.

In Gibbs coordinates there is a second
carrier issue. The plain pullback
\(f\mapsto f\circ s_A\) need not be
unitary for \(\nu_b\otimes\nu_b\).
The correctly transported unitary is
\begin{equation}
  \widetilde S_A f=R_A\,f\circ s_A
    =(\mathcal U_b\otimes\mathcal U_b)^*
          S_A(\mathcal U_b\otimes\mathcal U_b)f .
 \label{eq:f15v81:weighted-swap}
\end{equation}
Indeed the product Gibbs density is
\(B^2/Z_b^2\); multiplying by \(R_A^2\)
changes it to its swapped density.
A change of variables by \(s_A\) proves
the norm identity, and
\(R_A(R_A\circ s_A)=1\) proves
involutivity. Thus a large multiplier
ratio is not itself a large norm of
the transported swap: that operator
is still unitary.

Gauge averaging intertwines with
\(\mathcal U_b\otimes\mathcal U_b\),
because \(b\) is gauge invariant.
Consequently the maximal Gauss defect
in Theorem~\ref{f15v81:maximal-defect}
survives this exact change of carrier.
Correct density transport does not
turn a partial swap into a physical
local involution.

\subsection{The revised local target}

The antisymmetric physical norm is an
exact reference-free target and is
stable under genuine independent
tensor products. The attempted
single-link parity localization has
two explicitly identified terms:
boundary Gauss constraints and mixed
Wilson plaquettes. The former cannot
be dropped or replaced by commuting
compressed parity projections; the
latter cannot be declared small from
flatness of the two input replicas.

The cycle-return formula supplies a
concrete boundary building block.
A local continuation must retain the
open representation indices at a cut,
integrate the correct boundary
intertwiners, and include the magnetic
mixed-plaquette factors in the same
operator. Gauge framing or charged
boundary spaces may organize that
calculation, but their cost and
completeness must be proved.
The full physical compound matrix
\eqref{eq:f15v81:compound} already
retains these constraints and remains
a valid finite target.

No estimate here bounds its complete
interacting odd-sector norm below
\(e^{-m_*a_t}\ell_N^2\) with a
regulator-uniform \(m_*>0\).
V80's prescribed relative tail and
the finite entry errors must still
be added to any such estimate.
The interacting continuum, a total
limiting observable family, tightness,
locality and reconstruction remain
open as before.

\subsection{Diagnostics and preservation}

The new script is
\begin{center}\small
\path{scripts/verify_ym_f15_v81_antisymmetric_gauge_boundary.py}.
\end{center}
It reruns V80 and its inherited checks.
New exact-rational fixtures verify six
exterior eigenpairs, the determinant
kernel's factor \(1/2\), 49 Haar
character orthogonality identities
and 72 independently enumerated
cycle-cut returns. Two Cauchy--Binet
identities and 729 full-vector
positivity checks test the compound
Gram factorization. Nine tensor
decompositions and nine perturbation
budgets are also checked.

For the native boundary calculation,
168 noncommuting quaternion fixtures
verify the charged mixed-loop sign,
and 288 verify interior gauge-swap
covariance. There are 3,584 central
boundary-cancellation checks, 3,584
exact density-cocycle checks and
14 weighted isometry checks.
Finally 273 single-link cuts on
three-dimensional side-three and
side-four tori attain the exact
pure-gauge action difference \(16\).
The analytic proofs, not these finite
fixtures, establish the Haar projection
and full operator statements.

The script has 11,204 bytes and SHA--256
\begin{center}\scriptsize\ttfamily
C803D4F34D063AFFFA6BFD94B2C05EE34A66A53B6DA0FDD19983E581D6EC2720.
\end{center}
The V80 predecessor has 1,810,122 bytes
and SHA--256
\begin{center}\scriptsize\ttfamily
2A9E53966FFF95A33856F869F375AA656B2A14FF1282CC3F1148307D875634D8.
\end{center}
Removing this final section and
restoring the title version recovers
V80 byte for byte. All 140 previously
protected files were unchanged before
this append. V81 is the sole active
main note. Archives and companions
remain in place; only \path{.tex}
sources are kept in \path{latex_notes}.
Build products are stored outside
that folder, in
\path{artifacts/ym_f15_v81_texcheck/}.

\begin{firewall}
V81 fixes a tensor-stable physical
spectral target and computes exact
Gauss-boundary returns and a sharp
native magnetic swap defect. These
results rule out an unjustified
local parity decomposition; they do
not prove a uniform interacting
antisymmetric contraction.
The full continuum Yang--Mills
mass gap remains unproved.
The original goal remains active.
\end{firewall}


\section{V82: boundary-charge transfer and magnetic charge walks}
\label{f15v82:section}

\subsection{Context and nonduplication audit}

V81 computed the return of a partially swapped cycle after
Gauss projection. The next question is what replaces that
one-dimensional cycle calculation when a cut carries general
boundary charges and bulk loop multiplicities. This question
keeps the literal interacting transfer as the target of the
paper-inspired continuation begun in V78. No additional
continuum claim is imported from that paper.

There are important prior results not to reprove as new.
The archive f7 V80, in its theorems labelled
\texttt{thm:s7v80-projector} and
\texttt{thm:s7v80-mixture}, already gives the compact-group
boundary decomposition, normalized singlet gluing, and the
absence of a representation-dimension factor in the physical
sector trace. The archive f3 V91 computes the partial transpose
of a normalized singlet as a flip divided by its dimension.
Those are existing representation-theoretic inputs, not new
mass-gap estimates. The general distinction between charged
boundary carriers and unconstrained interior multiplicities
also appears in W.~Donnelly,
\emph{Entanglement entropy and nonabelian gauge symmetry},
Section~2,
\href{https://arxiv.org/html/1406.7304}{arXiv:1406.7304}.
That reference is background, not a spectral-gap input.

The new work here is the full multiplicity-space formula for
the compressed partial swap, its neutral-versus-charged
consequences, and its connection to the native Wilson cut
transfer. Mixed plaquette multiplication then gives an explicit
graph of allowed charge changes and an operator-norm
approximation by bounded-length charge walks. All statements
below are at a fixed finite spatial regulator unless an
algebraic bound is explicitly stated otherwise. The next
companion checkpoint remains V85.

\subsection{Cut notation and inherited singlet gluing}

Let the finite spatial graph have compact link variables
\(G=\SU(2)\), no matter fields and no external charged sources.
Partition its links into \(A\) and \(B=A^c\).
A boundary vertex meets links from both parts.
Impose Gauss invariance at vertices incident only to \(A\),
respectively only to \(B\), and write the resulting Hilbert
spaces as \(\mathcal H_A^{\rm int}\) and
\(\mathcal H_B^{\rm int}\).
Isolated vertices have no effect.
The boundary group is
\(K=\prod_{v\in\partial A}\SU(2)\).
The inherited compact-group decomposition is
\begin{align}
 \mathcal H_A^{\rm int}
 &=\bigoplus_r V_r\otimes\mathcal M_{A,r},
 &
 \mathcal H_B^{\rm int}
 &=\bigoplus_r V_r^*\otimes\mathcal M_{B,r},
 \label{eq:f15v82:regional}\\
 V_r&=\bigotimes_{v\in\partial A}V_{j_v},
 &
 D_r&=\prod_{v\in\partial A}(2j_v+1).
 \notag
\end{align}
The second decomposition uses a dual relabelling; this avoids
any choice of phases for self-dual \(\SU(2)\) representations.
Zero multiplicity spaces may be omitted. Define
\[
 \mathcal N_r=\mathcal M_{A,r}\otimes\mathcal M_{B,r},
 \qquad
 \Omega_r=D_r^{-1/2}\sum_{\alpha=1}^{D_r}
                   e_\alpha\otimes e^\alpha .
\]
The physical Haar space and its embedding are
\begin{equation}
 \mathcal H_{\rm phys}\simeq\bigoplus_r\mathcal N_r,
 \qquad
 J_r(a\otimes b)=\Omega_r\otimes a\otimes b,
 \label{eq:f15v82:gluing}
\end{equation}
with tensor factors reordered into the regional spaces.
Let \(Q_r\) be its orthogonal sector projections and
let \(0\) denote the all-trivial boundary label.
These are electric boundary representation labels, not
labels of an interacting spectral decomposition.
The countable Hilbert direct sums in
\eqref{eq:f15v82:regional}--\eqref{eq:f15v82:gluing}
are the full compact-group decompositions, not a cutoff
assumption. Finite Peter--Weyl spans form a dense subspace.

\subsection{Exact partial swap on multiplicity spaces}

On two replicas of the regional tensor product, let \(S_A\)
exchange the two \(A\) factors. Let \(P\) impose the boundary
Gauss constraint independently in both replicas, after the
interior constraints have been imposed. Thus \(P S_A P\),
restricted to the physical two-copy space, is the same
compressed partial swap as in V81. Denote it by \(R_A\)
in this subsection only; it is an operator, not V81's
multiplicative density cocycle.

\begin{theorem}[Complete compressed-swap formula]
\label{f15v82:swap}
The operator \(R_A\) kills
\(\mathcal N_r\otimes\mathcal N_s\) when \(r\ne s\).
On \(\mathcal N_r^{\otimes2}\), regrouped as
\(\mathcal M_{A,r}^{\otimes2}\otimes
  \mathcal M_{B,r}^{\otimes2}\), it is
\begin{equation}
 R_A=D_r^{-1} F_{A,r}\otimes I_{B,r}^{\otimes2},
 \label{eq:f15v82:swap}
\end{equation}
where \(F_{A,r}\) exchanges the \(A\) multiplicities.
On the globally antisymmetric physical subspace, its
matched-sector restrictions are
\begin{align}
 -D_r^{-1}I
 &\quad\hbox{on }\;
 \bigwedge\nolimits^2\mathcal M_{A,r}
       \otimes\operatorname{Sym}^2\mathcal M_{B,r},
 \notag\\
 +D_r^{-1}I
 &\quad\hbox{on }\;
 \operatorname{Sym}^2\mathcal M_{A,r}
       \otimes\bigwedge\nolimits^2\mathcal M_{B,r}.
 \label{eq:f15v82:odd}
\end{align}
The antisymmetrizations of unequal-sector pairs are killed.
Empty summands contribute nothing.
\end{theorem}

\begin{proof}
Start with elementary multiplicity tensors. After the \(A\)
swap, the first replica carries \(V_s\otimes V_r^*\)
and the second carries \(V_r\otimes V_s^*\).
Schur's lemma makes both Gauss projections zero unless
\(r=s\). For equal labels the carrier matrix element is
\[
 \langle\Omega_r\otimes\Omega_r,\,
           S_A(\Omega_r\otimes\Omega_r)\rangle
 =D_r^{-2}\sum_{\alpha,\beta=1}^{D_r}
                          {\bf1}_{\alpha=\beta}
 =D_r^{-1}.
\]
The multiplicities undergo exactly
\(a_1\otimes b_1\otimes a_2\otimes b_2
 \mapsto a_2\otimes b_1\otimes a_1\otimes b_2\).
This proves \eqref{eq:f15v82:swap}.
The global replica flip in a matched sector is
\(F_{A,r}\otimes F_{B,r}\).
Its eigenvalue \(-1\) consists of the two opposite-parity
summands in \eqref{eq:f15v82:odd}, giving the displayed signs.
The global flip exchanges unequal-sector summands, on both
of which the compression is zero.
All operators are bounded contractions; density extends
the elementary-tensor calculation to the full Hilbert space.
\end{proof}

For a single cycle with fixed edge representation \(n\),
the multiplicities in V81's cycle-only sector are
one-dimensional. Its boundary carrier dimension is
\(D_r=(n+1)^{b_A(p)}\). Thus
\eqref{eq:f15v82:swap} specializes to
\eqref{eq:f15v81:cycle-return}; it does not add a second
independent cycle damping factor.

\begin{proposition}[Charged return and neutral exceptions]
\label{f15v82:charged}
Under the no-matter assumptions above, every allowed
nontrivial boundary label satisfies \(D_r\geq3\). Consequently,
with \(Q_{\rm ch}=I-Q_0\otimes Q_0\),
\begin{equation}
 \|R_A Q_{\rm ch}\|\leq\tfrac13 .
 \label{eq:f15v82:third}
\end{equation}
The same bound holds after global antisymmetrization.
In contrast, if \(A\) contains a simple closed loop then
\(R_A\) has a neutral physical odd eigenvector of eigenvalue
\(-1\). If \(B\) contains such a loop it has one of
eigenvalue \(+1\).
\end{proposition}

\begin{proof}
For each connected component of the \(A\)-edge graph, the
gauge transformation \(-I\) at every vertex of that component
acts identically on every \(A\) link: its two endpoint
factors cancel. Interior Gauss invariance therefore implies
that the product of the boundary central transformations
in that component acts as the identity. On a carrier \(r\)
this says
\[
 \sum_{v\text{ on that component's boundary}}2j_v
                         \quad\hbox{is even}.
\]
The corresponding condition also holds for \(B\).
If any nonzero \(j_v\) is integral, its dimension is at
least three. Otherwise a nontrivial label has at least
two half-integral nonintegral spins, so its dimension is
at least four. This proves \(D_r\geq3\).
Unequal replica sectors vanish by
Theorem~\ref{f15v82:swap}; equal nontrivial sectors have
norm at most \(D_r^{-1}\). This proves
\eqref{eq:f15v82:third}, including the cross sectors with
one neutral replica.

The constant and the fundamental character of a simple
loop in \(A\) are orthogonal nonzero vectors invariant
under all \(A\)-vertex transformations, including those
at the boundary. Orthogonality follows by integrating
one loop link against Haar measure. They give two
independent vectors in \(\mathcal M_{A,0}\).
Use their antisymmetric product in the two \(A\)
multiplicities and constants in both \(B\) multiplicities.
Equation~\eqref{eq:f15v82:odd} gives eigenvalue \(-1\).
Interchanging \(A\) and \(B\) gives \(+1\).
\end{proof}

The constant \(1/3\) is attainable in this class of finite
pure-gauge graphs. Take two simple cycles meeting at exactly
one vertex, and put one cycle in each part. Interior gauge
reduction on a cycle leaves \(L^2(\SU(2))\), with conjugation
at the common vertex. Its edge-holonomy Peter--Weyl block
of spin \(k\) is \(\operatorname{End}(V_k)\) and decomposes
under conjugation into spins \(J=0,1,\ldots,2k\), once each.
The blocks \(k=1/2\) and \(k=1\) therefore give two independent
adjoint (\(J=1\)) multiplicities in \(A\), and \(k=1/2\)
gives one in \(B\). For this matched boundary sector
\(D_r=3\), and the first odd summand of
\eqref{eq:f15v82:odd} is nonzero. Its eigenvalue is
\(-1/3\). This is a kinematic sharpness example, not
an assertion about the Wilson vacuum or an interacting
excitation energy.

\subsection{The literal magnetic cut transfer}

Use V79's normalized native Haar transfer
\(\bar T=M_b C M_b\), where
\(C=C_A\otimes C_B\) is the product of normalized
one-link kinetic convolutions. Separate the spatial
plaquettes into those wholly in \(A\), wholly in \(B\),
and the mixed set \(\mathcal P_\times\).
For \(s_p=1-\tfrac12\Tr U_p\in[0,2]\), define
\begin{align}
 b_A&=\exp\!\left(-\frac{\gamma}{2}
                      \sum_{p\subset A}s_p\right),
 &
 b_B&=\exp\!\left(-\frac{\gamma}{2}
                      \sum_{p\subset B}s_p\right),
 \notag\\
 b_\times&=\exp\!\left(-\frac{\gamma}{2}
                      \sum_{p\in\mathcal P_\times}s_p\right),
 &
 T_A&=M_{b_A}C_A M_{b_A},
 \qquad T_B=M_{b_B}C_B M_{b_B}.
 \label{eq:f15v82:cut-factors}
\end{align}
All statements concern \(\gamma>0\) and the full kinetic
convolution. A Peter--Weyl truncation need not have a
pointwise strictly positive kernel and is not used in
the following Perron argument.

\begin{theorem}[Native cut spectrum]
\label{f15v82:cut}
On the physical space,
\begin{align}
 T_{\rm cut}
 &:=(T_A\otimes T_B)|_{\mathcal H_{\rm phys}}
   =\bigoplus_r T_{A,r}\otimes T_{B,r},
 \notag\\
 \bar T&=M_{b_\times}T_{\rm cut}M_{b_\times}.
 \label{eq:f15v82:cut}
\end{align}
Here \(T_{A,r}\) and \(T_{B,r}\) are the multiplicity
operators in the regional decomposition. No factor \(D_r\)
occurs. Write
\(\alpha=\|T_{A,0}\|\), \(\beta=\|T_{B,0}\|\).
The cut transfer has a simple top eigenvalue
\(\alpha\beta\), in the neutral sector. Its exact second
eigenvalue ratio is
\begin{equation}
 \rho(T_{\rm cut})
 =\max\left\{\rho(T_{A,0}),\rho(T_{B,0}),
    \sup_{r\ne0}
       \frac{\|T_{A,r}\|\|T_{B,r}\|}{\alpha\beta}
        \right\}<1.
 \label{eq:f15v82:cut-ratio}
\end{equation}
Only sectors present on both sides enter the supremum.
An absent charged part or a one-dimensional neutral
factor contributes zero to the corresponding term.
The strict inequality is only a fixed-regulator assertion.
\end{theorem}

\begin{proof}
A wholly regional plaquette is invariant under every
vertex gauge transformation in that region, including
boundary transformations. Each normalized central
one-link convolution commutes with independent left
and right translations. Thus \(T_A,T_B\) preserve the
interior constraints and commute with their boundary
group representations. Schur's lemma gives
\(I_{V_r}\otimes T_{A,r}\) and
\(I_{V_r^*}\otimes T_{B,r}\).
Applying the isometric singlet gluing
\eqref{eq:f15v82:gluing} proves the direct sum in
\eqref{eq:f15v82:cut}. Its carrier factor is the identity
on a single normalized singlet, not a trace over \(D_r\)
physical copies. This is the native-transfer
specialization of the previously established gluing,
not a new generic factorization principle.
The equality \(b=b_A b_B b_\times\) proves the sandwich
formula directly; no commutation of a magnetic
multiplier with \(C\) is asserted.

On the full regional Haar space the kernel of \(T_A\)
is continuous and strictly positive on a compact
configuration space. It defines a compact positive
self-adjoint operator. Its top eigenvalue is simple
with a strictly positive eigenfunction: replacement
of a maximizing vector by its absolute value increases
the Rayleigh quotient unless its phase is constant,
and two strictly positive top eigenvectors cannot be
orthogonal. Positivity of the kernel makes the
nonnegative top eigenfunction strictly positive.
Every regional gauge transformation commutes with
\(T_A\) and preserves positivity and normalization.
Uniqueness therefore makes that eigenfunction invariant
under all these transformations. It survives the
interior projection and belongs to the neutral
boundary carrier. The same proof applies to \(T_B\).
The positivity as an operator used here also follows
from the positive kinetic convolution sandwich.

The neutral tensor product has eigenvalues given by
products of its factor eigenvalues; its second ratio
is the maximum of the two neutral factor ratios.
Every nonneutral direct summand has norm
\(\|T_{A,r}\|\|T_{B,r}\|\).
Taking the largest competitor gives
\eqref{eq:f15v82:cut-ratio}.
Finally the compact operator \(T_A\otimes T_B\) has a
simple top eigenvalue, and its physical restriction
contains that eigenvector. Compactness gives a strict
separation from the rest of its spectrum at this fixed
regulator. It does not give a common separation for
growing graphs or varying coupling.
\end{proof}

This separates three dynamical tasks: neutral excitations
on each side, charged-sector energy costs, and the effect
of the actual mixed magnetic factors. The dimension
estimate \eqref{eq:f15v82:third} is not an estimate for
any of the spectral quantities in
\eqref{eq:f15v82:cut-ratio}.

\subsection{Mixed plaquettes and bounded charge walks}

Assume now a cubic spatial lattice with simple four-edge
plaquettes (for example a periodic lattice of side at
least three). For a mixed plaquette \(p\), let
\(\partial_A p\) be the vertices of \(p\) at which exactly
one of its two incident plaquette links lies in \(A\).
Write \(t_p=\tfrac12\Tr U_p\), a real bounded function.

\begin{proposition}[Necessary magnetic charge selection]
\label{f15v82:selection}
The physical block \(Q_s M_{t_p}Q_r\) vanishes unless
\begin{equation}
 \begin{cases}
  s_v\in\{j_v-\tfrac12,j_v+\tfrac12\}\cap[0,\infty),
       &v\in\partial_A p,\\
  s_v=j_v,&v\notin\partial_A p,
 \end{cases}
 \qquad r=(j_v)_v .
 \label{eq:f15v82:selection}
\end{equation}
Here \(s_v\) denotes the spin at vertex \(v\) of label \(s\).
These are necessary conditions only, with all regional
multiplicity and intertwiner constraints retained.
\end{proposition}

\begin{proof}
Temporarily apply a boundary gauge transformation only
to the \(A\) links. At a vertex with two \(A\) plaquette
links, the adjacent fundamental matrices and their
inverses cancel in the trace. With no \(A\) plaquette
link there is no action. With exactly one, the trace
depends on one fundamental matrix element of the
boundary transformation or its dual. Thus the finite
orbit span of \(t_p\), under this partial boundary
action, is a sum of copies of the tensor product of
fundamental carriers at \(\partial_A p\), with trivial
carriers elsewhere. Both fundamental and dual have
spin \(1/2\) for \(\SU(2)\).

Multiplication by this orbit span is an equivariant map
from its tensor product with any input carrier space
to the regional tensor Hilbert space. The
Clebsch--Gordan decomposition
\[
 V_j\otimes V_{1/2}
   =V_{j+1/2}\oplus V_{j-1/2}
\]
omits the negative-spin summand when \(j=0\).
It gives \eqref{eq:f15v82:selection} for the \(A\)
isotypic projections. Although partial boundary
transformations need not preserve the physical space,
these isotypic projections do commute with diagonal
Gauss averaging. Restriction to the physical singlet
subspace is consequently legitimate. Multiplication
by \(t_p\) itself is fully gauge invariant, so input
and output are physical; further constraints can
annihilate a permitted block but cannot create an
excluded one.
\end{proof}

Define a graph on the occurring matched labels \(r\):
join two labels whenever they satisfy the necessary
rule \eqref{eq:f15v82:selection} for at least one mixed
plaquette. Let \(d_\times(r,s)\) be its graph distance,
infinite between disconnected components. This graph
may include edges whose actual operator blocks vanish.
That enlargement makes the following support statement
weaker, not invalid.

\begin{theorem}[A norm-controlled charge-walk approximation]
\label{f15v82:walk}
Let \(m=|\mathcal P_\times|\),
\(x=\gamma m/2\), and
\(L_\times=\sum_{p\in\mathcal P_\times}M_{t_p}\)
on the physical Haar space. For an integer \(k\geq0\), put
\begin{align}
 D_k&=e^{-x}\sum_{h=0}^k
                  \frac{(\gamma L_\times/2)^h}{h!},
 &
 \delta_k(x)&=e^{-x}\sum_{h=k+1}^{\infty}\frac{x^h}{h!}.
 \label{eq:f15v82:walk-def}
\end{align}
Then
\begin{align}
 \|M_{b_\times}-D_k\|&\leq\delta_k(x),&
 \|D_k\|&\leq1,
 \label{eq:f15v82:walk-error}\\
 Q_sD_k Q_r&=0
          &&\text{if }d_\times(r,s)>k .
 \notag
\end{align}
Consequently \(Q_s M_{b_\times}Q_r\) has norm at most
\(\delta_k(x)\) for those pairs, and the positive compact
operator
\[
 T^{[k]}=D_k T_{\rm cut}D_k
\]
satisfies
\begin{align}
 \|\bar T-T^{[k]}\|
       &\leq2\delta_k(x)\|T_{\rm cut}\|,\\
 Q_sT^{[k]}Q_r&=0
       \quad\text{if }d_\times(r,s)>2k .
 \label{eq:f15v82:walk-transfer}
\end{align}
The approximation is banded in charge distance, not
finite rank. For \(k+2>x\), a scalar upper bound is
\begin{equation}
 \delta_k(x)\leq
 e^{-x}\frac{x^{k+1}}{(k+1)!}
               \frac{1}{1-x/(k+2)}.
 \label{eq:f15v82:poisson-tail}
\end{equation}
At \(m=0\), \(D_k=I\) and the error is zero.
\end{theorem}

\begin{proof}
Each \(M_{t_p}\) is self-adjoint of norm at most one.
Thus \(\|L_\times\|\leq m\), and the exact identity
\[
 M_{b_\times}=\exp(-xI+\gamma L_\times/2)
\]
has a norm-convergent power series. Its discarded
series has norm at most \(\delta_k(x)\).
The retained series has norm at most
\(e^{-x}\sum_{h=0}^k x^h/h!\leq1\).
Also \(0<b_\times\leq1\), so its multiplication norm
is at most one.

Proposition~\ref{f15v82:selection} implies that each
application of \(L_\times\) moves the charge support
by at most one edge of the graph. Induction gives
support within distance \(h\) for its \(h\)-th power.
This argument is valid for countably many sectors:
orthogonal projection onto a union of sectors is
bounded, and finite partial sums converge strongly.
It proves the polynomial support statement and the
off-band error bound.

The polynomial \(D_k\) is self-adjoint, but need not
be a positive multiplier. Nevertheless its sandwich
with the positive compact \(T_{\rm cut}\) is positive
and compact. With \(M=M_{b_\times}\),
\[
 MT_{\rm cut}M-D_kT_{\rm cut}D_k
 =(M-D_k)T_{\rm cut}M
       +D_kT_{\rm cut}(M-D_k).
\]
The two norm bounds give
\eqref{eq:f15v82:walk-transfer}.
The cut transfer preserves each charge sector, so
the two polynomials permit at most \(2k\) charge
steps in total. Finally, successive terms in the
discarded scalar series have ratio at most
\(x/(k+2)<1\), which gives
\eqref{eq:f15v82:poisson-tail}.
\end{proof}

The entire exponential is not asserted to have finite
charge range. Nor is the approximation asserted to
preserve a Markov kernel. Its useful properties are
the stated support, positivity of the transfer
sandwich, and a full operator-norm error. That error
is measured relative to \(\|T_{\rm cut}\|\), not
automatically relative to the interacting
\(\|\bar T\|\). For large \(\gamma m\), keeping \(k\)
fixed does not make the scalar tail small.

\subsection{Direction after this calculation}

A charged boundary dimension gain survives exact Gauss
gluing, but neutral multiplicities can saturate the
partial-swap norm. Moreover the genuine mixed magnetic
factor connects the electric charge sectors. Dropping
either fact would turn a kinematic \(1/3\) into an
unsupported physical contraction.

The next useful dynamical target is therefore the complete
neutral-plus-charged operator in
\eqref{eq:f15v82:cut}, with its mixed-plaquette charge
walks included. A local estimate must control neutral
loop excitations and charged-sector energies together,
then pay the magnetic coupling and approximation
errors in the actual top-eigenvalue normalization.
The exact finite physical Gram and compound certificates
of V79--V81 remain available for that task; they do not
supply those estimates by themselves. Choosing a charge
cutoff without controlling its multiplicity spaces
would again omit physical states.

In particular no bound proved here establishes
\(\rho(\bar T)\leq e^{-m_*a_t}\) with a positive common
\(m_*\) on a prescribed physical time scale. Continuum
tightness, a total limiting observable family,
reconstruction and nontrivial interaction remain
separate open obligations. The paper review in V78
and its claim boundaries are unchanged.

\subsection{Diagnostics and preservation}

The new script is
\begin{center}\small
\path{scripts/verify_ym_f15_v82_boundary_charge_transfer.py}.
\end{center}
It reruns V81 and the inherited chain; all 6,629 new
exact-rational fixtures passed. They independently contract
two singlets with nontrivial multiplicities, test
the surviving odd signs and unmatched-sector zeros,
enumerate necessary central-parity dimension bounds,
and check the three terms of the cut spectral ratio.
Further fixtures check the fundamental character
identity, native noncommuting plaquette dependence
on a partial boundary transformation, polynomial
charge-walk support, positive sandwiches and scalar
tail majorants. These finite diagnostics do not
replace the analytic proofs or certify a large
interacting matrix estimate.

The script has 8,206 bytes and SHA--256
\begin{center}\scriptsize\ttfamily
787DBC3B6B3B0C495F39E554EE2F07DC314F00D16F3A2BA33B39A92A0D5C4F5E.
\end{center}
The predecessor V81 has 1,832,730 bytes and SHA--256
\begin{center}\scriptsize\ttfamily
7F44432A15E69721D98AF941DABC8DC9D426297042E4566565DDD18C246112A0.
\end{center}
Removing this final section and restoring the title
version recovers that predecessor byte for byte.
All 141 previously protected source and diagnostic
files were unchanged before this append. V82 replaces
V81 as the sole active main note; archives and companion
sources remain in place. Only \path{.tex} sources
are kept in \path{latex_notes}. Build products belong
in \path{artifacts/ym_f15_v82_texcheck/}.

\begin{firewall}
V82 proves the exact multiplicity-space boundary return,
a charged \(1/3\) bound with explicit neutral exceptions,
the native cut-transfer spectral decomposition, and a
controlled mixed-plaquette charge-walk approximation.
It does not prove uniform neutral or charged energy
floors, stability of a gap under magnetic gluing, or
the four-dimensional continuum Yang--Mills mass gap.
The original goal remains active.
\end{firewall}


\section{V83: native charged-sector domination and charge tails}
\label{f15v83:section}

\subsection{The dynamical question left by V82}

V82 separated neutral bulk excitations, nontrivial boundary
charges, and the mixed magnetic plaquettes. Its \(1/3\)
compressed-swap bound was kinematic and did not bound a
charged transfer eigenvalue. This continuation supplies a
different estimate: a charged-sector norm bound relative to
the \emph{same regional interacting transfer's} Perron
eigenvalue. It holds on arbitrary finite regional volumes,
at fixed coupling, and retains every regional magnetic
plaquette. It is not a one-loop model.

The calculation uses a boundary frame on the full regional
Haar space, before interior Gauss projection. An exact
Fourier integral in that frame is dominated by its neutral
scalar integral. Several nonadjacent boundary vertices can
be treated together. This also yields a relative-error
boundary-charge truncation of the full magnetically glued
transfer, without knowing its vacuum vector.

The archive search found no earlier native charged-sector
comparison of this form. V79's one-link Bessel coefficients,
V82's singlet gluing and cut decomposition, and the existing
normalized-star Haar integrals are reused. No new claim is
made for those inputs. The charge labels here are boundary
gauge representations, not the charged colour coordinates
of the older gauge-fixed Gaussian calculations. No companion
checkpoint is due before V85.

\subsection{A monotone scalar coefficient}

Retain the \(\SU(2)\) convention
\(t(u)=\tfrac12\Tr u\).
For \(s>0\) set
\begin{equation}
 z(s)=\int_{\SU(2)}e^{s t(u)}\,du
       =\frac{2I_1(s)}s,
 \qquad
 r_n(s)=\frac{I_{n+1}(s)}{I_1(s)}
       \quad(n=0,1,\ldots).
 \label{eq:f15v83:coefficients}
\end{equation}
The continuous endpoint values are \(z(0)=r_0(0)=1\)
and \(r_n(0)=0\) for \(n\geq1\).
The irreducible representation indexed by \(n\) has
spin \(n/2\) and dimension \(n+1\).

\begin{proposition}[Coefficient properties used in the comparison]
\label{f15v83:scalar}
For \(s>0\), \(n\geq1\), the function \(r_n(s)\) is
strictly increasing in \(s\), lies in \((0,1)\), and
is decreasing in \(n\). In particular \(r_n(s)\leq r_1(s)\).
For fixed \(n\), as \(s\) tends to infinity,
\begin{equation}
 r_n(s)=1-\frac{n(n+2)}{2s}+O_n(s^{-2}).
 \label{eq:f15v83:asymptotic}
\end{equation}
For a unitary irreducible matrix \(D_n\) and \(u\in\SU(2)\),
\begin{equation}
 \int_{\SU(2)}e^{s t(q u)}D_n(q)\,dq
       =z(s)r_n(s)D_n(u^{-1}).
 \label{eq:f15v83:fourier}
\end{equation}
There is no extra factor \(n+1\) in this matrix integral.
\end{proposition}

\begin{proof}
For each integer \(j\geq1\), the adjacent ratio
\(I_{j+1}(s)/I_j(s)\) is the mean of \(t\) under the
probability density proportional to
\(e^{st}(1-t^2)^{j-1/2}\) on \((-1,1)\).
Its derivative is the strictly positive variance of \(t\);
its value lies in \((0,1)\). This follows by differentiating
the integral representation of \(s^{-j}I_j(s)\).
The standard identities used here are recorded in
\href{https://dlmf.nist.gov/10.32.E2}{DLMF 10.32.2}
and
\href{https://dlmf.nist.gov/10.29.E2}{DLMF 10.29.2}.
The product
\[
 r_n(s)=\prod_{j=1}^n\frac{I_{j+1}(s)}{I_j(s)}
\]
proves both monotonicity assertions and the bounds.
The fixed-order expansion
\[
 I_\nu(s)=\frac{e^s}{\sqrt{2\pi s}}
 \left[1-\frac{4\nu^2-1}{8s}+O_\nu(s^{-2})\right],
\]
as in
\href{https://dlmf.nist.gov/10.40.E1}{DLMF 10.40.1},
gives \eqref{eq:f15v83:asymptotic} by division.

For completeness, the central matrix integral is a
scalar by Schur's lemma. Its trace divided by \(n+1\)
is
\[
 \frac{1}{n+1}\int e^{s t(q)}\chi_n(q)\,dq
   =\frac{2I_{n+1}(s)}s=z(s)r_n(s),
\]
the character integral already used in V79.
Changing variables \(p=q u\) gives
\eqref{eq:f15v83:fourier}.
The trace division cancels the representation dimension;
retaining it would give a wrong norm comparison.
\end{proof}

\subsection{An exact regional frame, before gauge reduction}

Let \(R\) be a finite regional link graph with distinct
endpoints on every link. Put
\[
 C_R=\bigotimes_{e\in R}(C_\gamma/c_0),
 \qquad T_R=M_{b_R}C_RM_{b_R},
\]
on its full product Haar space, where \(b_R\) is the
positive Wilson multiplier from the plaquettes wholly
in \(R\). In fact the following argument permits any
strictly positive continuous \(b_R\) invariant under
every regional vertex gauge transformation.
The kernel is
\begin{equation}
 k_R(U,V)=b_R(U)b_R(V)z(\gamma)^{-|R|}
             \exp\!\left(\gamma\sum_{e\in R}
                                  t(U_eV_e^{-1})\right).
 \label{eq:f15v83:native-kernel}
\end{equation}
This is the actual kinetic normalization of V79, since
\(c_0=e^{-\gamma}z(\gamma)\).

Choose a set \(S\) of nonadjacent regional vertices,
each incident to at least one regional link, and write
\(d_R(v)\) for its regional degree. Reverse link
orientations if necessary so that every link incident
to \(S\) points away from its vertex in \(S\).
This is possible because no link joins two vertices
of \(S\). Orientation reversal preserves the Haar
measure and the \(\SU(2)\) kernel in
\eqref{eq:f15v83:native-kernel}.

At each \(v\in S\), choose one incident anchor link.
Write \(U_e=g_v X_e\) on that star, with the anchored
\(X_e=I\); use the unchanged \(U_e\) on all other links.
Distinct stars have disjoint link sets. The resulting
coordinates have exactly product measure
\[
 dU=\left(\prod_{v\in S}dg_v\right)d\xi,
 \qquad \xi\in \SU(2)^{|R|-|S|}.
\]
Indeed, \(g_v\) is the old anchor variable and each
remaining star coordinate is obtained by a
Haar-preserving left translation. A gauge
transformation at \(v\) acts only as left translation
of \(g_v\). In particular \(b_R(U)=b_R(\xi)\).

Use \(V=(h,\eta)\) and write its star links as \(h_vY_e\).
Define the quaternion-valued field
\begin{equation}
 F_v(\xi,\eta)=\sum_{e\ni v}Y_eX_e^{-1},
 \qquad
 s_v=\gamma|F_v|\leq\gamma d_R(v).
 \label{eq:f15v83:field}
\end{equation}
The norm is the Euclidean norm of a real quaternion.
Every summand is a unit quaternion. When \(F_v\ne0\),
write \(F_v=|F_v|u_v\) with \(u_v\in\SU(2)\).
At \(F_v=0\), choose \(u_v=I\).
Cyclicity and invariance of the real \(\SU(2)\) trace
under inversion give
\[
 \sum_{e\ni v}t(g_vX_eY_e^{-1}h_v^{-1})
                  =t(g_v^{-1}h_v F_v).
\]
Consequently the complete kernel has the exact form
\begin{equation}
 k_R((g,\xi),(h,\eta))
   =a(\xi,\eta)
       \prod_{v\in S}e^{s_v t(g_v^{-1}h_v u_v)},
 \qquad a(\xi,\eta)>0.
 \label{eq:f15v83:fibre-kernel}
\end{equation}
The factor \(a\) includes both magnetic multipliers,
the normalization, and every link outside the selected
stars. None is dropped or replaced by a free-field law.

This frame construction is made \emph{before} imposing
interior Gauss constraints. A reduced loop variable can
have only a conjugation action, which is not free.
No free boundary action on that already reduced
configuration space is assumed here.

\begin{theorem}[Normalized native charged-sector domination]
\label{f15v83:domination}
Let \(\lambda_R=\|T_R\|\) on the full regional Haar
space. On the joint isotypic sector of spins \(n_v/2\)
at \(v\in S\),
\begin{equation}
 \|T_R|_{\mathbf n}\|
       \leq \lambda_R
                     \prod_{v\in S}r_{n_v}(\gamma d_R(v)).
 \label{eq:f15v83:domination}
\end{equation}
Now impose interior Gauss invariance and retain the
boundary multiplicity operators of V82. Their neutral
top eigenvalue is still \(\lambda_R=\|T_{R,0}\|\).
For any occurring boundary label \(r=(n_v/2)_v\) and
any nonadjacent subset \(S\) of boundary vertices,
\begin{equation}
 \frac{\|T_{R,r}\|}{\|T_{R,0}\|}
       \leq\prod_{v\in S}r_{n_v}(\gamma d_R(v)).
 \label{eq:f15v83:regional-charge}
\end{equation}
These bounds have no regional-volume or
representation-dimension prefactor.
\end{theorem}

\begin{proof}
Decompose the group coordinates \(g_v\) in their
orthonormal Peter--Weyl matrix bases. Formula
\eqref{eq:f15v83:fourier} gives the joint Fourier
kernel, on each carrier's multiplicity variables, as
\[
 \kappa_0(\xi,\eta)
      \left(\prod_{v\in S}r_{n_v}(s_v)\right)
                       \mathcal U_{\mathbf n}(\xi,\eta),
 \qquad
 \kappa_0=a\prod_{v\in S}z(s_v),
\]
where \(\mathcal U_{\mathbf n}\) is unitary, a tensor
product of the matrices \(D_{n_v}(u_v^{-1})\), up to
the fixed Fourier index convention. The zero-spin
sector has the scalar kernel \(\kappa_0\).
Repeated carrier indices give identical blocks, not
a factor multiplying their norms.

Let \(K_0\) be that scalar integral operator.
For a vector-valued coefficient function \(f\), the
pointwise triangle inequality, unitarity and
Proposition~\ref{f15v83:scalar} imply
\[
 |(T_{R,\mathbf n}f)(\xi)|
 \leq
 \left(\prod_{v\in S}r_{n_v}(\gamma d_R(v))\right)
                         (K_0|f|)(\xi).
\]
The same statement holds for matrix coefficients
with their Hilbert--Schmidt norm. Taking the \(L^2\)
norm bounds \(\|T_{R,\mathbf n}\|\) by the displayed
product times \(\|K_0\|\).

The full native kernel is strictly positive and
continuous on a compact space, and \(T_R\) is a
positive operator. The Perron argument in
Theorem~\ref{f15v82:cut} gives a simple strictly positive
top eigenfunction invariant under every regional
vertex gauge transformation. It belongs to the
zero-spin sector at \(S\), so
\(\|K_0\|=\lambda_R\), not just an unrelated
normalizer. This proves
\eqref{eq:f15v83:domination}.

Interior Gauss averaging commutes with the boundary
group actions and with \(T_R\). Each specified
boundary block of the interior-invariant subspace
is a subspace of the full joint isotypic sector
just bounded. Restriction cannot increase the norm.
The same full Perron eigenfunction survives all
interior constraints and is neutral at every
boundary vertex. Thus its eigenvalue equals
\(\|T_{R,0}\|\), proving
\eqref{eq:f15v83:regional-charge}.
\end{proof}

\subsection{Charged costs for the physically glued cut}

Use the link partition \(A,B\), charge sectors \(Q_r\),
cut transfer \(T_{\rm cut}\), and factors
\(b_\times\) from V82. Write
\[
 \alpha=\|T_{A,0}\|,\qquad
 \beta=\|T_{B,0}\|,\qquad
 n_v=2j_v.
\]
At each boundary vertex, both \(d_A(v)\) and \(d_B(v)\)
are positive. Define nonnegative dimensionless costs
\begin{align}
 e_v(r)&=-\log\!\left[
       r_{n_v}(\gamma d_A(v))r_{n_v}(\gamma d_B(v))\right],
 \notag\\
 \mathcal E(r)&=\sum_{v\in\partial A}e_v(r),
 &
 \mathcal W(r)&=\max_{\substack{S\subset\partial A\\
                              S\ {\rm independent}}}
                                      \sum_{v\in S}e_v(r).
 \label{eq:f15v83:costs}
\end{align}
Here independence refers to the full spatial graph
induced on the boundary vertices, so it is sufficient
on both regional graphs. For \(n_v=0\) the cost is zero.
These are lower-bound coefficients for transfer
attenuation, not measured physical excitation energies.

\begin{theorem}[A charged cut bound uniform in spatial volume]
\label{f15v83:cut-cost}
If the induced boundary graph has a proper colouring
with \(\chi\) colours, then every occurring label satisfies
\begin{equation}
 \frac{\|T_{A,r}\|\,\|T_{B,r}\|}{\alpha\beta}
       \leq e^{-\mathcal W(r)}
       \leq e^{-\mathcal E(r)/\chi}.
 \label{eq:f15v83:cost-bound}
\end{equation}
For a cubic three-dimensional spatial lattice, set
\begin{equation}
 q_{\partial}(\gamma)
  =\max_{\substack{a,b\in\mathbb N,\ a,b\geq1\\a+b\leq6}}
                     r_1(a\gamma)r_1(b\gamma)<1.
 \label{eq:f15v83:boundary-ratio}
\end{equation}
Then
\begin{align}
 \sup_{r\ne0}
   \frac{\|T_{A,r}\|\,\|T_{B,r}\|}{\alpha\beta}
        &\leq q_{\partial}(\gamma),
 \notag\\
 \rho(T_{\rm cut})
   &\leq
     \max\{\rho(T_{A,0}),\rho(T_{B,0}),q_{\partial}(\gamma)\}.
 \label{eq:f15v83:cut-ratio}
\end{align}
A valid universal colouring choice is \(\chi\leq7\);
on a bipartite spatial graph \(\chi\leq2\).
At fixed \(\gamma>0\), all displayed charged bounds
are independent of spatial volume.
As \(\gamma\) tends to infinity,
\begin{equation}
 q_{\partial}(\gamma)=1-\gamma^{-1}+O(\gamma^{-2}),
 \qquad
 -\log q_{\partial}(\gamma)=\gamma^{-1}+O(\gamma^{-2}).
 \label{eq:f15v83:boundary-asymptotic}
\end{equation}
No constant positive lower bound in lattice time units
is asserted as \(\gamma\) tends to infinity.
\end{theorem}

\begin{proof}
Apply \eqref{eq:f15v83:regional-charge} to both regions
with the same independent set \(S\), and multiply.
The result is \(\exp(-\sum_{v\in S}e_v(r))\).
Taking the strongest of these bounds proves the first
inequality in \eqref{eq:f15v83:cost-bound}.
The colour classes are independent and partition
the vertices. One class has at least
\(\mathcal E(r)/\chi\) of the nonnegative cost,
proving the second inequality.

If \(r\ne0\), choose one vertex with \(n_v\geq1\).
For that singleton, \(r_{n_v}\leq r_1\) on both sides.
Its degrees are positive integers with sum at most
six. This proves the first line of
\eqref{eq:f15v83:cut-ratio}. The second follows from
the exact physical sector formula
\eqref{eq:f15v82:cut-ratio}. There are fifteen pairs
in \eqref{eq:f15v83:boundary-ratio}; each product is
strictly less than one. The degree bound also gives
a greedy colouring with at most seven colours,
and bipartiteness gives two. No count of charge
labels is inserted.

For each of the fifteen pairs, the fixed-order
expansion \eqref{eq:f15v83:asymptotic} gives
\[
 r_1(a\gamma)r_1(b\gamma)
 =1-\frac{3}{2\gamma}\left(\frac1a+\frac1b\right)
                                      +O(\gamma^{-2}).
\]
The minimum of its leading cost over the finite
set is \(1\), attained at \(a=b=3\).
The remainders are uniform over this finite set.
Taking the maximum and then the logarithm proves
\eqref{eq:f15v83:boundary-asymptotic}.
\end{proof}

The no-matter parity constraints of V82 may exclude
some labels, which only improves these upper bounds.
No assertion that the envelope
\(q_{\partial}(\gamma)\) is attained by a physical
cut eigenvector is needed.

\subsection{A positive relative-error boundary-charge truncation}

The previous theorem can be applied before the mixed
magnetic factors are glued back. It is essential to
keep their sandwich, because they need not commute
with charge projections.

\begin{theorem}[Charge-tail certificate with actual normalization]
\label{f15v83:tail}
For a cost cutoff \(L\geq0\), define
\begin{align}
 Q_{\leq L}&=\sum_{\mathcal E(r)\leq L}Q_r,
 \notag\\
 \bar T_{\leq L}
   &=M_{b_\times}Q_{\leq L}T_{\rm cut}
                         Q_{\leq L}M_{b_\times}.
 \label{eq:f15v83:charge-cut}
\end{align}
Let \(m=|\mathcal P_\times|\).
At every fixed finite lattice and coupling,
\(Q_{\leq L}\) contains only finitely many boundary
labels, but it need not have finite rank.
One has
\begin{align}
 0\leq\bar T_{\leq L}&\leq\bar T,
 &
 \|\bar T-\bar T_{\leq L}\|
       &\leq\alpha\beta\,e^{-L/\chi},
 \label{eq:f15v83:absolute-tail}\\
 \lambda_0(\bar T_{\leq L})
       &\geq\alpha\beta\,e^{-2\gamma m},
 &
 \frac{\|\bar T-\bar T_{\leq L}\|}
        {\lambda_0(\bar T_{\leq L})}
       &\leq e^{\,2\gamma m-L/\chi}.
 \label{eq:f15v83:relative-tail}
\end{align}
The relative bound also holds with
\(\lambda_0(\bar T)\) in the denominator.
Thus any \(0<\eta<1\) is guaranteed by choosing
\begin{equation}
 L\geq\chi\bigl(2\gamma m+\log(1/\eta)\bigr).
 \label{eq:f15v83:tail-choice}
\end{equation}
Every neutral multiplicity state of \(T_{\rm cut}\)
is retained by \(Q_{\leq L}\).
\end{theorem}

\begin{proof}
The cut transfer is block diagonal in \(r\), with
nonnegative blocks. Equation
\eqref{eq:f15v83:cost-bound} implies
\[
 0\leq T_{\rm cut}-Q_{\leq L}T_{\rm cut}Q_{\leq L},
 \qquad
 \|T_{\rm cut}-Q_{\leq L}T_{\rm cut}Q_{\leq L}\|
                         \leq\alpha\beta e^{-L/\chi}.
\]
The exact magnetic sandwich
\eqref{eq:f15v82:cut} and \(0<b_\times\leq1\)
prove \eqref{eq:f15v83:absolute-tail}.

Let \(\omega_{\rm cut}\) be the normalized cut
Perron vector. It is neutral, so it belongs to
the retained subspace for every \(L\geq0\).
Because \(0\leq s_p\leq2\),
\[
 e^{-\gamma m}\leq b_\times\leq1.
\]
The physical test vector
\(f=b_\times^{-1}\omega_{\rm cut}\) is therefore
well defined, with \(\|f\|^2\leq e^{2\gamma m}\).
Its numerator for the retained transfer is exactly
\[
 \langle f,\bar T_{\leq L}f\rangle
     =\langle\omega_{\rm cut},T_{\rm cut}
                                \omega_{\rm cut}\rangle
     =\alpha\beta .
\]
Rayleigh's principle proves the first inequality
in \eqref{eq:f15v83:relative-tail}; division gives
the second and \eqref{eq:f15v83:tail-choice}.
The order bound implies
\(\lambda_0(\bar T)\geq\lambda_0(\bar T_{\leq L})\).

At any fixed positive argument, \(r_n\) tends to
zero as \(n\) tends to infinity, for example by
V79's bound \(r_n(s)\leq(s/2)^n/(n+1)!\).
There are finitely many boundary vertices, and
each has positive degree on both sides.
Thus the condition \(\mathcal E(r)\leq L\)
bounds every individual \(n_v\), leaving only
finitely many possible labels. Their bulk
multiplicities are not bounded by this argument.
Since \(\mathcal E(0)=0\), the entire neutral
space \(Q_0\mathcal H_{\rm phys}\) is retained.
\end{proof}

The boundary magnetic factor is still present in
\(\bar T_{\leq L}\) on both sides, with no expansion.
The finite set of labels constrains the middle cut
operator; its magnetically weighted output vectors
need not have only those boundary charges.
The theorem is not the generally invalid claim
\(Q_{\leq L}\bar T Q_{\leq L}\leq\bar T\).
Its relative error depends on the actual
mixed-plaquette count and coupling, even though it
does not divide by an unknown vacuum eigenvalue
or sum a dimension factor over omitted labels.
The possible growth of the required cutoff and
of retained multiplicity spaces is not hidden.

\subsection{What is now established, and the next bottleneck}

There is now a native, normalized and spatial-volume
uniform charged-sector comparison for the cut transfer.
The all-charge tail can also be removed to prescribed
relative operator accuracy after exact magnetic gluing,
at the explicit cost
\eqref{eq:f15v83:tail-choice}.
This addresses the charged-energy and high-charge
parts left open in V82 at a fixed coupling.

It does not control the neutral bulk transfer ratios
in \eqref{eq:f15v83:cut-ratio}. Those spaces contain
the loop excitations isolated in V82, and the charge
cost vanishes on all of them. Applying the same
argument to the complete physical lattice without
a cut gives only the trivial vertex charges:
Gauss invariance makes every such \(n_v=0\).
One cannot relabel a neutral physical excitation
as a charged one to use this estimate.

The next task is consequently the low-charge,
infinite-multiplicity transfer with its magnetic
couplings, especially its neutral loop part.
Its full physical spectral ratio, not a local
boundary return, must be controlled on the
prescribed physical time scale. The finite Gram
and antisymmetric certificates of V79--V81 and
the charge-walk organization of V82 can be
combined with this charge-tail estimate, but
their remaining interacting norm inequalities
are not proved by that combination alone.

In particular, the \(1/\gamma\) charged attenuation
in \eqref{eq:f15v83:boundary-asymptotic} is not
declared to be a continuum mass. No temporal
lattice scale is chosen from that attenuation.
The interacting continuum and its positive
physical mass gap remain the unchanged objective.

\subsection{Diagnostics and preservation}

The new script is
\begin{center}\small
\path{scripts/verify_ym_f15_v83_charged_transfer_domination.py}.
\end{center}
It reruns V82 and the inherited checks; all 3,438
new exact-rational fixtures passed. These tests verify the native anchored-star
kernel identity and field magnitude, then independently
integrate the scalar, fundamental and adjoint shifted
Fourier coefficients with quaternion Haar moments.
Outward Bessel-series intervals test the scalar
orderings and the finite cubic-degree envelope.
Additional fixtures test full-vector unitary-kernel
norm domination, independent-set colouring costs,
and the positive charge-tail sandwich with a
noncommuting magnetic multiplier. These tests
check finite identities and bounds; the preceding
proofs establish the full operator statements.

For illustration, outward rational series enclosures
give the following intervals for the scalar envelope,
not for an actual physical excited eigenvalue:
\[
 \begin{array}{c|c}
 \gamma&q_{\partial}(\gamma)\\ \hline
 1&[0.322537,\ 0.322538]\\
 2&[0.581751,\ 0.581752]\\
 8&[0.880181,\ 0.880182]
 \end{array}
\]
Every displayed decimal endpoint is an exact rational
with denominator \(10^6\).
The script has 11,004 bytes and SHA--256
\begin{center}\scriptsize\ttfamily
FD870677A2170CF9F252EE6A037E0D24B51084D22C1603C4ACB22096FB0A87E9.
\end{center}
The predecessor V82 has 1,856,070 bytes and SHA--256
\begin{center}\scriptsize\ttfamily
34FAD1C6D30FEF23157D0D9CC4A51DEB229E022105FB156F831121408FC160EA.
\end{center}
Removing this final section and restoring the title
version recovers V82 byte for byte.
All 142 previously protected source and diagnostic
files were unchanged before this append.
V83 is the sole active main version; archives and
companion sources remain in place.
Only \path{.tex} sources are kept in
\path{latex_notes}; build files belong in
\path{artifacts/ym_f15_v83_texcheck/}.

\begin{firewall}
V83 proves normalized native charged-sector bounds
and a relative boundary-charge tail certificate.
The charged bound is uniform in spatial volume
at fixed coupling, not a uniform continuum
physical mass gap. Neutral loop multiplicities,
their interacting magnetic coupling and the
four-dimensional continuum construction remain
open. The original goal remains active.
\end{firewall}


\section{V84: one-link spectral reduction and vacuum charge escape}
\label{f15v84:section}

\subsection{Context and the upstream adjustment}

V83 controlled charged blocks of the cut transfer relative
to their own interacting Perron normalization. Magnetic
gluing still mixes those blocks. Before using a neutral
reference as a spectral approximation, one must distinguish
an unweighted neutral subspace from a magnetically weighted
transfer image.

This distinction is substantive on the actual lattice
vacuum. We prove that its projection onto any fixed
one-link Peter--Weyl band tends to zero along the existing
selected couplings. Nevertheless a kinetic truncation
\emph{inside the magnetic sandwich} can have a controlled
relative operator error. A new one-link comparison proves
such an error without a factor from the number of other
links. A positive operator union bound then improves
V80's whole-lattice product-trace remainder to a sum of
one-link tails. Finally we identify the precise neutral
one-link reduction and the magnetic weight it generates.

The projection estimate below is an elementary weighted
Hilbert--Schmidt argument, not a new general uncertainty
principle. Its programme-specific application uses the
actual vacuum exponential moment of V62, not a trial
wavefunction. The archive's f1 V31 one-link character
debits and f14 V81 gauge-matrix spectral concentration
concern different operators. V79's Gram factorization
and V66's normalized Haar star integral are reused;
the new spectral normalization and omitted-space
estimates are proved here. The next companion
checkpoint remains V85.

\subsection{The actual vacuum escapes a fixed one-link band}

Fix a finite spatial regulator, a link \(e\), and a
simple spatial plaquette \(p\) containing \(e\) exactly
once. Work first on the full product Haar Hilbert
space. Let \(\Pi_{e,N}\) retain all one-link irreducibles
with labels \(0\leq n\leq N\), including both matrix
indices. Its one-link dimension is
\begin{equation}
 D_N=\sum_{n=0}^N(n+1)^2
       =\frac{(N+1)(N+2)(2N+3)}6.
 \label{eq:f15v84:dimension}
\end{equation}
Recall \(s_p=1-\tfrac12\Tr U_p\), and put
\begin{equation}
 c_0(a)=\int_{\SU(2)}e^{-a(1-t(u))}\,du
            =e^{-a}\frac{2I_1(a)}a
            \quad(a>0).
 \label{eq:f15v84:well}
\end{equation}
This is the same one-link scalar used to normalize
the native kinetic convolution, now evaluated at
a possibly different positive argument \(a\).

\begin{theorem}[A volume-free weighted projection estimate]
\label{f15v84:uncertainty}
For every \(\psi\) in the full Haar space and \(a>0\),
\begin{equation}
 \|\Pi_{e,N}\psi\|^2
       \leq D_N c_0(a)
                       \langle\psi,e^{a s_p}\psi\rangle .
 \label{eq:f15v84:uncertainty}
\end{equation}
Moreover
\begin{equation}
 c_0(a)\leq\sqrt{\frac2\pi}\,a^{-3/2}.
 \label{eq:f15v84:well-bound}
\end{equation}
No conditional moment bound under the law of \(\psi\)
is assumed.
\end{theorem}

\begin{proof}
Fix all link variables other than \(U_e\).
As a function of \(U_e\), the plaquette holonomy is
a Haar-preserving left/right translate of \(U_e\)
or \(U_e^{-1}\). The diagonal of the one-link
projection kernel is the constant \(D_N\):
in the orthonormal basis
\(\sqrt{n+1}D_n(u)_{ij}\), unitarity gives
\[
 \sum_{n=0}^N\sum_{i,j}
            (n+1)|D_n(u)_{ij}|^2=D_N.
\]
Consequently the fibrewise Hilbert--Schmidt norm obeys
\[
 \|\Pi_{e,N}M_{e^{-a s_p/2}}\|_{\rm HS}^2
       =D_N\int e^{-a s_p}\,dU_e
       =D_Nc_0(a).
\]
The operator norm is no larger. Apply this to
\(e^{a s_p/2}\psi\) on that fibre, then integrate
over the remaining Haar coordinates.
At a fixed regulator \(s_p\) is bounded, so these
multiplications are bounded. This proves
\eqref{eq:f15v84:uncertainty} without a conditional-law
assertion or a factor from other links.

For Haar \(\SU(2)\), \(s=1-t(u)\) has density
\((2/\pi)\sqrt{s(2-s)}\) on \([0,2]\). Therefore
\[
 c_0(a)
 \leq\frac{2\sqrt2}{\pi}\int_0^\infty
                                s^{1/2}e^{-as}\,ds
 =\sqrt{\frac2\pi}\,a^{-3/2},
\]
using \(\Gamma(3/2)=\sqrt\pi/2\).
\end{proof}

\begin{proposition}[A native-vacuum boundary-charge consequence]
\label{f15v84:vacuum}
Let \(\Omega_M\) be the normalized positive physical
vacuum eigenvector on V60's selected spatial regulator,
with \(\gamma_M=\widetilde\beta_M\to\infty\).
For all sufficiently large \(M\),
\begin{align}
 \|\Pi_{e,N}\Omega_M\|^2
  &\leq e^3D_Nc_0(\gamma_M/16)
 \notag\\
  &\leq64e^3\sqrt{\frac2\pi}\,
                         D_N\gamma_M^{-3/2}.
 \label{eq:f15v84:vacuum-band}
\end{align}
In particular this projection tends to zero whenever
\(N=N_M=o(\sqrt{\gamma_M})\).
For the cut consisting of the single link \(A=\{e\}\),
its raw neutral projection \(Q_0\) satisfies
\(\|Q_0\Omega_M\|^2\to0\).

If \(F_M\) is a uniformly bounded gauge-invariant
multiplication observable, the first line of
\eqref{eq:f15v84:vacuum-band} also bounds
\(\|\Pi_{e,N}F_M\Omega_M\|^2\) after multiplication
by \(\|F_M\|_\infty^2\).
For normalized observable states, the corresponding
norm denominator must be retained.
\end{proposition}

\begin{proof}
Integrating the second endpoint and temporal links
in the vacuum slab measure
\eqref{eq:f15v59:vacuum-slab} gives the equal-time
marginal \(\Omega_M(U)^2\,dH(U)\), by the Perron
eigenvalue equation. Proposition~\ref{f15v62:moments}
therefore gives
\[
 \langle\Omega_M,e^{\gamma_M s_p/16}\Omega_M\rangle
                                                     \leq e^3
\]
for this spatial plaquette. Apply
Theorem~\ref{f15v84:uncertainty}.
Since \(D_N=O((N+1)^3)\), the claimed limiting
statement follows.

The one-link Peter--Weyl projection commutes with
both endpoint gauge actions, and hence with the
full Gauss projection. For a one-link cut, the
regional carrier labels at its two endpoints are
\((n/2,n/2)\), with the dual convention of V82.
The neutral label is exactly \(n=0\), so
\(Q_0=\Pi_{e,0}\) on the physical space.
Finally \(|F_M|^2\leq\|F_M\|_\infty^2\) in the
same weighted expectation proves the observable
version. Its normalized form requires division
by \(\|F_M\Omega_M\|^2\), which need not have a
uniform lower bound.
\end{proof}

This rules out a fixed unweighted electric band
as a norm approximation to the actual selected
vacuum. It does \emph{not} rule out a magnetic
image \(M_b\Pi_{e,N}\mathcal H\). Multiplication
by \(b\) can generate arbitrarily high electric
labels. The next result deliberately uses that
weighted image, not the projection just ruled out.

\subsection{A local relative remainder with no spectator-volume factor}

Use the normalized native physical transfer
\[
 \bar T=M_b(C_e\otimes C_{\widehat e})M_b,
 \qquad C_e=C_\gamma/c_0,
\]
where \(C_{\widehat e}\) is the untruncated product
of the other one-link convolutions. Write
\[
 C_{e,N}=\Pi_{e,N}C_e\Pi_{e,N},\qquad
 \bar T_{e,N}=M_b(C_{e,N}\otimes C_{\widehat e})M_b.
\]
The tensor notation refers to the full Haar space
before Gauss restriction. All these operators
commute with the Gauss projection.
In particular \(C_{e,0}=P_{e,0}\), the Haar
average on the selected link, with kernel one.
Define the scalar tail
\begin{equation}
 \tau_N(\gamma)
      =\sum_{n>N}(n+1)^2 r_n(\gamma).
 \label{eq:f15v84:local-tail}
\end{equation}

\begin{theorem}[Native one-link relative tail]
\label{f15v84:local-tail}
On the physical Hilbert space,
\begin{align}
 0\leq\bar T-\bar T_{e,N},\qquad
 \|\bar T-\bar T_{e,N}\|
      &\leq\tau_N(\gamma)\lambda_0(\bar T_{e,0}),
 \label{eq:f15v84:local-relative}\\
 \frac{\|\bar T-\bar T_{e,N}\|}
      {\lambda_0(\bar T_{e,N})}
      &\leq\tau_N(\gamma).
 \notag
\end{align}
The same relative bound holds with
\(\lambda_0(\bar T)\) in the denominator.
It is uniform over spatial volume and over the
native positive gauge-invariant magnetic multiplier.
The other links remain untruncated; this is not
yet a finite-rank approximation.
\end{theorem}

\begin{proof}
The one-link omitted kernel is the uniformly
absolutely convergent character series
\[
 R_N(u,v)=\sum_{n>N}(n+1)r_n(\gamma)
                                         \chi_n(uv^{-1}).
\]
Convergence follows, for example, from V79's
factorial coefficient bound. Since
\(|\chi_n|\leq n+1\), it satisfies
\(|R_N(u,v)|\leq\tau_N(\gamma)\).
Every other full kinetic factor has a strictly
positive continuous kernel. Since \(b>0\),
the full Haar kernels consequently obey
\[
 |(\bar T-\bar T_{e,N})(U,V)|
                       \leq\tau_N(\gamma)\bar T_{e,0}(U,V).
\]
Thus for any complex \(f\),
\[
 |(\bar T-\bar T_{e,N})f|
                    \leq\tau_N(\gamma)\bar T_{e,0}|f|,
\]
and taking the \(L^2\) norm gives the corresponding
full-space operator bound.

The operator \(\bar T_{e,0}\) is positive and compact
and has a strictly positive continuous kernel.
It commutes with all gauge transformations, so its
unique positive Perron vector is gauge invariant,
by the same argument as in V83.
Its full-space norm therefore equals its largest
physical eigenvalue. Restricting the preceding
bound to the physical space proves
\eqref{eq:f15v84:local-relative}.
The operator inequalities
\[
 0\leq\bar T_{e,0}\leq\bar T_{e,N}\leq\bar T
\]
follow from their commuting positive kinetic
factors and the same magnetic sandwich.
Monotonicity of the largest eigenvalue proves
both relative conclusions.
\end{proof}

This argument does not estimate a full trace over
the other links, and it does not divide by a
magnetic partition function. The positive kernel
used for comparison is that of \(\bar T_{e,0}\);
a finite nonzero Peter--Weyl cutoff need not
itself have a positive kernel.

\begin{theorem}[A simultaneous positive union bound]
\label{f15v84:union}
Choose a cutoff \(N_e\geq0\) for every spatial link
and put
\[
 C_{\boldsymbol N}=\bigotimes_e C_{e,N_e},
 \qquad
 \bar T_{\boldsymbol N}=M_bC_{\boldsymbol N}M_b,
 \qquad
 \eta_{\boldsymbol N}=\sum_e\tau_{N_e}(\gamma).
\]
Then \(\bar T_{\boldsymbol N}\) has finite rank
on the physical space, is positive, and
\begin{equation}
 0\leq\bar T-\bar T_{\boldsymbol N},\qquad
 \|\bar T-\bar T_{\boldsymbol N}\|
        \leq\eta_{\boldsymbol N}\lambda_0(\bar T).
 \label{eq:f15v84:union}
\end{equation}
If \(\eta_{\boldsymbol N}<1\), then
\begin{align}
 \lambda_0(\bar T_{\boldsymbol N})
       &\geq(1-\eta_{\boldsymbol N})\lambda_0(\bar T),
 \notag\\
 \frac{\|\bar T-\bar T_{\boldsymbol N}\|}
      {\lambda_0(\bar T_{\boldsymbol N})}
       &\leq\frac{\eta_{\boldsymbol N}}
                    {1-\eta_{\boldsymbol N}} .
 \label{eq:f15v84:union-relative}
\end{align}
In particular equal cutoffs cost \(E\tau_N(\gamma)\),
where \(E\) is the number of spatial links, rather
than the full product-trace tail of V80.
\end{theorem}

\begin{proof}
All kinetic factors and their label projections
commute. In the joint Peter--Weyl basis the scalar
inequality
\[
 1-\prod_e{\bf1}_{n_e\leq N_e}
       \leq\sum_e{\bf1}_{n_e>N_e}
\]
gives the positive operator inequality
\[
 C-C_{\boldsymbol N}
  \leq\sum_e
       \left[(C_e-C_{e,N_e})\otimes C_{\widehat e}\right].
\]
Sandwich by \(M_b\), restrict to the physical
space, and apply Theorem~\ref{f15v84:local-tail}
to each term, with its \emph{original untruncated}
other-link factors. Each
\(\lambda_0(\bar T_{e,0})\leq\lambda_0(\bar T)\).
This proves \eqref{eq:f15v84:union}.
Min--max gives
\(\lambda_0(\bar T_{\boldsymbol N})
 \geq\lambda_0(\bar T)-\|\bar T-\bar T_{\boldsymbol N}\|\),
and division gives
\eqref{eq:f15v84:union-relative}.
Finite rank is preserved by the magnetic
sandwich of the finite tensor-product kinetic range.
\end{proof}

For comparison, if
\(s=\sum_{n\geq0}(n+1)^2r_n\) and
\(s_N=\sum_{n=0}^N(n+1)^2r_n\), then
\[
 s^E-s_N^E
    =(s-s_N)\sum_{j=0}^{E-1}s^{E-1-j}s_N^j
    \geq E\tau_N,
\]
because \(s,s_N\geq1\).
Thus the new error coefficient is no larger
than V80's product-trace coefficient.
This improves the remainder, not the dimension
of the retained full physical Gram matrix.
It also avoids an invalid iteration of pointwise
positivity after successive signed-kernel cutoffs.

\subsection{A better large-coupling scale for the one-link tail}

The scalar \(\tau_N\) can already be enclosed by
V79--V80's outward positive series. The following
bound gives a more informative large-coupling
scale than a factorial estimate alone.

\begin{proposition}[Gaussian-to-exponential coefficient control]
\label{f15v84:bernstein}
For \(\gamma>0\) and \(n\geq1\),
\begin{equation}
 r_n(\gamma)
 \leq\exp\!\left[-\frac{n(n+2)}{2\gamma+2n+1}\right].
 \label{eq:f15v84:bernstein}
\end{equation}
For \(\gamma\geq1\), this implies
\[
 r_n(\gamma)\leq
 \begin{cases}
  e^{-n^2/(5\gamma)},&1\leq n\leq\gamma,\\
  e^{-n/5},&n>\gamma .
 \end{cases}
\]
Let \(q=e^{-1/5}\) and define
\[
 H(K)=q^K\left[
    \frac{(K+1)^2}{1-q}
   +\frac{2(K+1)q}{(1-q)^2}
   +\frac{q(1+q)}{(1-q)^3}\right].
\]
For integer \(\sqrt{5\gamma}\leq N\leq\gamma\),
\begin{equation}
 \tau_N(\gamma)
 \leq\left(10\gamma N+\frac{25\gamma^2}{N}\right)
                     e^{-N^2/(5\gamma)}
       +H(\lfloor\gamma\rfloor+1).
 \label{eq:f15v84:gaussian-tail}
\end{equation}
For \(N>\gamma\), one instead has
\(\tau_N(\gamma)\leq H(N+1)\).
In particular the choice
\begin{equation}
 N(\gamma)=\left\lceil\sqrt{10\gamma\log\gamma}\right\rceil
 \quad\hbox{for sufficiently large }\gamma
 \label{eq:f15v84:local-scale}
\end{equation}
gives
\(\tau_{N(\gamma)}(\gamma)
 =O(\gamma^{-1/2}\sqrt{\log\gamma})\to0\).
\end{proposition}

\begin{proof}
For \(j\geq1\) set \(u_j=I_{j+1}/I_j\).
The Bessel recurrences, as recorded in
\href{https://dlmf.nist.gov/10.29}{DLMF 10.29},
give
\[
 u_j'(\gamma)=1-u_j(\gamma)^2
                    -\frac{2j+1}{\gamma}u_j(\gamma).
\]
V83 identified the derivative as a strictly positive
variance. Solving the resulting quadratic
inequality yields
\[
 u_j(\gamma)
 <\sqrt{1+\left(\frac{j+1/2}{\gamma}\right)^2}
                        -\frac{j+1/2}{\gamma}
 =\exp\!\left[-\operatorname{arsinh}
                       \frac{j+1/2}{\gamma}\right].
\]
For \(x\geq0\),
\(\operatorname{arsinh}x\geq x/(1+x)\).
Both sides vanish at zero, and the derivative
inequality follows from
\((1+x)^4\geq1+x^2\).
Multiplying the adjacent-ratio estimates gives
\[
 -\log r_n(\gamma)
 \geq\sum_{j=1}^n\frac{j+1/2}{\gamma+j+1/2}
 \geq\frac{n(n+2)}{2\gamma+2n+1}.
\]
The two simpler exponentials follow by bounding
the denominator by \(5\gamma\) when \(n\leq\gamma\),
and by \(5n\) when \(n>\gamma\).

For \(N\geq\sqrt{5\gamma}\),
\(x^2e^{-x^2/(5\gamma)}\) is decreasing on
\([N,\infty)\). Since \((n+1)^2\leq4n^2\),
the part \(N<n\leq\gamma\) of the tail is at most
\[
 4\int_N^\infty x^2e^{-x^2/(5\gamma)}\,dx
 \leq\left(10\gamma N+\frac{25\gamma^2}{N}\right)
                                        e^{-N^2/(5\gamma)}.
\]
Integration by parts and
\(\int_N^\infty e^{-x^2/a}\,dx
 \leq(a/(2N))e^{-N^2/a}\) give the last inequality.
The remaining geometric sum is exactly
\(\sum_{n\geq K}(n+1)^2q^n=H(K)\).
This proves the tail bounds.
For \eqref{eq:f15v84:local-scale}, its Gaussian
exponential is at most \(\gamma^{-2}\);
the prefactor is \(O(\gamma^{3/2}\sqrt{\log\gamma})\).
The geometric contribution is
\(O(\gamma^2e^{-\gamma/5})\).
The displayed scale lies in the required interval
for all sufficiently large \(\gamma\), proving
the last assertion.
\end{proof}

The scale \eqref{eq:f15v84:local-scale} concerns one
link or a fixed number of links. On a growing lattice,
the simultaneous certificate still requires
\(\sum_e\tau_{N_e}\) to meet the chosen error budget.
If that budget must be \(o(a_t)\), the prescribed
physical \(a_t\) must be used in selecting it.
No physical clock is defined by the cutoff.
The raw-vacuum projection and the magnetically
weighted kinetic truncation are different objects;
their scales are not asserted to be matching
necessary and sufficient estimates for one problem.

\subsection{The exact local Gram operator and its neutral member}

The inherited Gram construction becomes particularly
explicit when only one link is truncated.
Use an orthonormal one-link basis
\[
 \phi_a(u)=\sqrt{n_a+1}\,D_{n_a}(u)_{i_a j_a},
 \qquad 0\leq n_a\leq N,
 \qquad d_a=\sqrt{r_{n_a}(\gamma)} .
\]
Write the remaining links as \(V\). Before Gauss
restriction, define the matrix-valued magnetic weight
\begin{equation}
 \mathcal H^{(N)}_{ab}(V)
   =\int_{\SU(2)}b(u,V)^2
                        \overline{\phi_a(u)}\phi_b(u)\,du .
 \label{eq:f15v84:local-magnetic-matrix}
\end{equation}
It is a strictly positive matrix at each fixed \(V\),
since \(b>0\) and the retained functions are independent.

\begin{proposition}[A local Gram identity with all other links retained]
\label{f15v84:local-gram}
The nonzero physical spectrum of \(\bar T_{e,N}\)
is the nonzero spectrum of the operator matrix
\begin{equation}
 (\mathcal G_N)_{ab}
  =d_a d_b\,C_{\widehat e}^{1/2}
             M_{\mathcal H^{(N)}_{ab}}
                         C_{\widehat e}^{1/2},
 \label{eq:f15v84:local-gram}
\end{equation}
restricted to the combined Gauss-invariant subspace
of the retained edge carriers and the other-link
space. The other-link space remains infinite
dimensional. Its components cannot individually
be declared gauge invariant when \(N>0\).
\end{proposition}

\begin{proof}
Let \(\mathcal F_N=\Pi_{e,N}\mathcal H\) before
Gauss restriction, and define
\[
 (W_N f)(u,V)
     =b(u,V)\sum_a d_a\phi_a(u)
                         (C_{\widehat e}^{1/2}f_a)(V).
\]
Direct multiplication gives
\(W_NW_N^*=\bar T_{e,N}\), while integration in
\(u\) gives \(W_N^*W_N=\mathcal G_N\).
The magnetic multiplier, kinetic factors and
complete Peter--Weyl projection are gauge equivariant.
Thus the same factorization restricts to the full
combined Gauss-invariant spaces. The positive
nonzero eigenvalues of \(WW^*\) and \(W^*W\) agree,
including multiplicities. This is V79's Gram
principle in a local factorization, not an
additional abstract spectral theorem.
\end{proof}

For \(N=0\), the retained edge function is exactly
one. On a cubic three-dimensional spatial slice,
the deleted edge belongs to \(m=4\) plaquettes.
Separate the spatial magnetic multiplier as
\[
 b(u,V)=b_{\widehat e}(V)b_\times(u,V).
\]
With the appropriate oriented three-link staples,
there is a quaternion \(S_e(V)\), the sum of
four unit quaternions, such that
\[
 \sum_{p\ni e}t(U_p)=t(uS_e(V)),\qquad |S_e(V)|\leq4 .
\]
Consequently the exact scalar weight is
\begin{align}
 Z_e(V)
    &:=\int b_\times(u,V)^2\,du
       =e^{-m\gamma}z(\gamma|S_e(V)|),
 \label{eq:f15v84:neutral-weight}\\
 \mathcal H^{(0)}_{00}(V)
    &=b_{\widehat e}(V)^2Z_e(V).
 \notag
\end{align}
At \(S_e=0\) this uses \(z(0)=1\).
The staple norm, and hence \(Z_e\), is invariant
under all vertex gauge transformations on the
graph with the edge removed.

\begin{proposition}[The neutral reduction retains a generated magnetic weight]
\label{f15v84:neutral}
The nonzero physical spectrum of \(\bar T_{e,0}\)
is exactly that of
\begin{equation}
 T_{{\rm eff},e}
   =M_{\,b_{\widehat e}\sqrt{Z_e}}\,
       C_{\widehat e}\,
                   M_{\,b_{\widehat e}\sqrt{Z_e}}
 \label{eq:f15v84:neutral-effective}
\end{equation}
on the fully gauge-invariant other-link space.
This is an exact statement about the positive
part \(\bar T_{e,0}\), not about the entire
\(\bar T\). The discarded charged modes must
still be paid.
\end{proposition}

\begin{proof}
For \(N=0\), formula
\eqref{eq:f15v84:local-gram} is
\[
 C_{\widehat e}^{1/2}
        M_{b_{\widehat e}^2Z_e}
                         C_{\widehat e}^{1/2}.
\]
A second \(W^*W,WW^*\) factorization with
\(W=M_{b_{\widehat e}\sqrt{Z_e}}
 C_{\widehat e}^{1/2}\) gives
\eqref{eq:f15v84:neutral-effective}.
The constant edge function carries no charge,
so combined Gauss invariance is precisely
invariance of the other-link function at
all its vertices. Gauge invariance of the
weight permits this restriction.
The integral \eqref{eq:f15v84:neutral-weight}
follows from quaternion rotation invariance
and the definition of \(z\).
\end{proof}

The scalar Haar integral is familiar from the
earlier star calculations; its spectral role
here is the new point. It involves four
\emph{spatial} plaquettes, not the six
spacetime plaquettes of V66's conditional
star. It is also not the conditional
density of the actual vacuum, whose
equal-time law is \(\Omega^2dH\).
It occurs because the Gram operator
integrates the square of the magnetic
multiplier.

This weight is not generally constant:
\[
 Z_e=e^{-\gamma(m-|S_e|)}c_0(\gamma|S_e|).
\]
Thus even the neutral reduction contains
a generated interaction depending on the
entire staple sum. Dropping \(Z_e\) would
change the operator whose spectrum is
being estimated. For \(N>0\), the full
matrix weight
\eqref{eq:f15v84:local-magnetic-matrix}
retains the charged mixing and the
combined Gauss constraints.

\subsection{The next spectral target}

The new remainder bound makes a local
magnetically weighted reduction quantitatively
legitimate when its computed \(\tau_N\) is small.
Its simultaneous version improves the global
finite-matrix error budget without changing
the target theory. What remains is a
nonperturbative bound for the leading
physical eigenvalue ratio of the retained
scalar or matrix-weighted operator.

In particular, setting \(N=0\) is not justified
at large coupling merely by calling the
retained sector neutral:
\(\tau_0=c_0(\gamma)^{-1}-1\to\infty\),
by V80's one-link trace identity and
\eqref{eq:f15v84:well-bound}.
Conversely, the raw-vacuum escape theorem
does not invalidate the weighted Gram
construction; its image includes the
magnetic multiplier and need not lie in
the raw band.

The next step should therefore keep the
generated magnetic weight and its physical
matrix-valued extensions while seeking
neutral-loop spectral control. A covariance
or auxiliary diffusion estimate for a
retained marginal is still not a bound
for this native transfer. No positive
regulator-uniform physical mass or
interacting continuum construction has
been obtained.

\subsection{Diagnostics and preservation}

The new script is
\begin{center}\small
\path{scripts/verify_ym_f15_v84_local_transfer_reduction.py}.
\end{center}
It reruns V83 and the inherited chain;
all 3,690 new exact-rational fixtures
passed. They check
projection dimensions and reproducing
kernels, weighted full-vector projection
bounds, outward one-link normalization,
Riccati and representation-tail enclosures,
and both Gaussian and geometric tail
budgets. Further fixtures check the local
kernel domination with coupled magnetic
weights and spectator links, the positive
simultaneous-cutoff union inequality,
local Gram matrices, and literal
four-plaquette staple integrals on a
three-dimensional torus.

These diagnostics test finite identities
and scalar enclosures. The preceding
proofs establish the full operator
estimates and the selected-vacuum
consequence; no large physical
eigenvalue gap is inferred from sampling.

As concrete outward-series choices for
one link, with all other links untruncated,
the script certifies \(\tau_N\leq1/100\) at
\[
 \begin{array}{c|ccc}
  \gamma&1&4&16\\ \hline
  N&4&8&18
 \end{array}
\]
These are relative operator-error
certificates, not mass-gap estimates.
The script has 11,672 bytes and SHA--256
\begin{center}\scriptsize\ttfamily
FF6CF514C98D9C563EE89265D389803AC925F5FC780D8C3DB380D55A46C08515.
\end{center}
The predecessor V83 has 1,877,741 bytes
and SHA--256
\begin{center}\scriptsize\ttfamily
3E458A3A5C4CC30AC57B371C6857FE6D34542FA22A693F8F0A6B60B57F260238.
\end{center}
Removing this final section and restoring
the title version recovers V83 byte for
byte. All 143 previously protected source
and diagnostic files were unchanged
before this append. V84 replaces V83
as the sole active main note; archives
and companions remain in place.
Only \path{.tex} sources are kept in
\path{latex_notes}. Build products belong
in \path{artifacts/ym_f15_v84_texcheck/}.

\begin{firewall}
V84 proves native local and simultaneous
relative truncation bounds, a quantitative
escape of the actual vacuum from fixed
unweighted one-link bands, and the exact
generated weight in the neutral reduction.
It does not replace the interacting
retained operator by a scalar mixing
model, discard charged errors, or prove
the continuum Yang--Mills mass gap.
The original goal remains active.
\end{firewall}

\section{V85: exact charged magnetic blocks and their retained transport}
\label{f15v85:context}

V84 replaced a product-trace truncation cost by a sum of one-link
tails and identified the exact generated neutral weight. The next
unresolved object is the retained, interacting physical transfer,
not the discarded tail. This section calculates its entire
one-link magnetic matrix, including charged carriers. It proves
an exact axial block reduction, explicit polynomial ellipticity,
and sharp fixed-band conditioning. It then keeps the associated
moving-frame factors in the native transfer kernel.

This is not another proof of V79's Gram factorization, V66's
scalar star integral, or the one-link diffusion decompositions
in f5 V20 and f4 V78. Those inputs are credited at their original
scope. The new object is V84's finite matrix-valued magnetic
weight inside the native transfer, with all other links retained.

\subsection{The scheduled companion checkpoint}

A fresh inventory found the same active companions as at V80.
The current SM continuum V72, NS stability V26, shell-mixing V16
and parallel YM V5 conclusions were reread, together with the
supplied source-selection V31, stationarity V34, Einstein-bridge
V24 and suggestions. All 144 previously protected files,
including these sources, are unchanged by hash.

The SM result is a fixed-time auxiliary many-copy limit, not
cutoff-uniform physical dynamics. The NS rank-one refinement
does not prove regularity. Shell mixing retains a measured flat
one-loop calculation but explicitly lacks the joint
nonperturbative remainder. The parallel YM total-charge
tightness route remains obstructed. Source selection still
requires word-native typing; stationarity has not populated its
330 moving-fibre derivatives; and the Einstein bridge has a
nonzero upstream tangent, not a nonlinear root. These facts
supply no new physical YM spectral estimate.

The useful methodological transfer is to retain the actual
carrier, generated density and moving coordinates together.
Here that means computing the charged matrix and transporting
its frame through the same kinetic kernel. No auxiliary
stationarity or diffusion theorem is substituted for the native
transfer. The next companion checkpoint is V90.

\subsection{The full one-link matrix in Cartesian coordinates}

Keep V84's three-dimensional spatial lattice and a simple edge
\(e\) with four incident plaquettes. Write \(V\) for all other
links, \(q\in S^3=\SU(2)\) for the deleted edge, and
\[
 b(q,V)^2=b_{\widehat e}(V)^2 e^{-4\gamma}
                 \exp\!\bigl(\gamma t(qS_e(V))\bigr).
\]
The quaternion \(S_e\) is the sum of four oriented staples.
Put \(F=\gamma\overline{S_e}\in\mathbb R^4\), \(s=|F|\leq4\gamma\),
and \(v=F/s\) when \(s>0\). Thus
\(\gamma t(qS_e)=q\cdot F\).
As before,
\[
 z(s)=\int_{S^3}e^{s q_0}\,dH(q)=\frac{2I_1(s)}s,\qquad z(0)=1,
 \qquad Z_e=e^{-4\gamma}z(s).
\]
Let \(\mathcal F_N\) be the complete Peter--Weyl band \(0\leq n\leq N\),
of dimension
\[
 D_N=\sum_{n=0}^N(n+1)^2.
\]
In any Haar-orthonormal basis \(\phi_a\) of that band define
\begin{equation}
 \widehat H_N(F)_{ab}
   =\frac1{z(s)}\int_{S^3}e^{q\cdot F}
                       \overline{\phi_a(q)}\phi_b(q)\,dH(q).
 \label{eq:f15v85:normalized-matrix}
\end{equation}
Then the matrix in V84 is exactly
\begin{equation}
 \mathcal H^{(N)}(V)
    =b_{\widehat e}(V)^2 Z_e(V)\widehat H_N(F(V)).
 \label{eq:f15v85:full-weight}
\end{equation}
In particular, division by \(z(s)\) here is an algebraic
factorization of the magnetic matrix. It neither row-normalizes
the native transfer nor asserts a formula for its vacuum law.

The entries of \(\widehat H_N\) are real-analytic functions of the
four Cartesian components of \(F\). Indeed the numerator is an
integral of a fixed polynomial on a compact sphere against an
entire exponential; its Taylor series converges locally uniformly,
and \(z(|F|)>0\) on real \(F\). This also covers \(F=0\), where
\(\widehat H_N(0)=I\). Directional formulas below introduce no
singularity into the original matrix.

\subsection{An exact axial reduction for every band}

Use normalized probability measure on \(S^2\), and choose
orthonormal spherical harmonics \(Y_{\ell m}\), \(1\leq m\leq2\ell+1\).
Write \(q=(t,\sqrt{1-t^2}\,\omega)\). Haar measure is
\[
 dH(q)=\frac2\pi\sqrt{1-t^2}\,dt\,d\sigma(\omega).
\]
For \(k\geq0\), set
\[
 h_{\ell k}
   =\frac{(k+2\ell+1)!}
          {4^\ell(k+\ell+1)(\ell!)^2 k!},
 \qquad
 \psi_{\ell km}(q)
   =h_{\ell k}^{-1/2}(1-t^2)^{\ell/2}
                C_k^{(\ell+1)}(t)Y_{\ell m}(\omega).
\]
Here the conventional Gegenbauer normalization is the one in
\href{https://dlmf.nist.gov/18.3}{DLMF Table 18.3.1}.

\begin{theorem}[Finite axial blocks of the native magnetic weight]
\label{f15v85:axial}
The functions with \(0\leq\ell\leq N\), \(0\leq k\leq N-\ell\)
form a Haar-orthonormal basis of \(\mathcal F_N\).
At \(F=s(1,0,0,0)\), the normalized matrix is
\begin{equation}
 A_N(s)=\bigoplus_{\ell=0}^N
                      A_{\ell,N}(s)\otimes I_{2\ell+1},
 \label{eq:f15v85:axial-block}
\end{equation}
where the block of size \(N-\ell+1\) has entries
\begin{equation}
 (A_{\ell,N}(s))_{kj}
  =\frac{2}{\pi z(s)\sqrt{h_{\ell k}h_{\ell j}}}
     \int_{-1}^1 e^{st}(1-t^2)^{\ell+1/2}
        C_k^{(\ell+1)}(t)C_j^{(\ell+1)}(t)\,dt.
 \label{eq:f15v85:radial-entry}
\end{equation}
For general \(F=sv\), a unitary \(U_v\), preserving every
harmonic degree \(n=\ell+k\), gives
\[
 \widehat H_N(F)=U_v A_N(s)U_v^*.
\]
Consequently the kinetic diagonal
\(D=\operatorname{diag}(\sqrt{r_{n_a}(\gamma)})\)
commutes with \(U_v\).
\end{theorem}

\begin{proof}
Separation of the spherical Laplacian in \((t,\omega)\) reduces
the degree-\(n\) equation to
\[
 (1-t^2)C''-(2\ell+3)tC'
                       +k(k+2\ell+2)C=0,\qquad n=\ell+k.
\]
Its polynomial solution is \(C_k^{(\ell+1)}\).
The factor \((1-t^2)^{\ell/2}Y_{\ell m}\) is a homogeneous
harmonic polynomial in the three transverse coordinates, so
the displayed functions extend smoothly across the poles.
The degree-\(n\) spherical harmonics on \(S^3\) are the
Peter--Weyl matrix-coefficient space of label \(n\); both have
Laplacian eigenvalue \(n(n+2)\) in the unit-sphere normalization
and dimension \((n+1)^2\).
Angular orthogonality and the Gegenbauer squared norm give
exactly \(h_{\ell k}\). Moreover
\[
 \sum_{\ell=0}^N(2\ell+1)(N-\ell+1)=D_N,
\]
so no carrier is missing.
The density \(e^{st}\) is angularly invariant, proving
\eqref{eq:f15v85:axial-block}--\eqref{eq:f15v85:radial-entry}.

Left multiplication by the unit quaternion \(v\) maps the pole
to \(v\), preserves Haar measure, and preserves each harmonic
degree. Its finite-dimensional unitary action, with the
appropriate inverse convention on coefficient vectors, gives
the stated conjugation. Since \(D\) is scalar on each complete
degree, it commutes with this action.
\end{proof}

This is a pointwise decomposition of a matrix weight, not a
decomposition of the entire interacting transfer into
independent angular sectors.

\subsection{A global, explicit polynomial ellipticity bound}

For \(d=N-\ell\) define rational positive matrices of size \(d+1\):
\[
 (K_{\ell,d})_{ij}=\frac{2}{2\ell+2i+2j+3},\qquad
 (R_{\ell,d})_{ij}
   =\int_{S^3}(1-q_0^2)^\ell(1-q_0)^{i+j}\,dH(q).
\]
Rationality of \(R\) follows from the even Haar moments.
These are the Gram matrices of independent monomials in
positive measures. Set
\begin{equation}
 c_N=\min\left\{\frac19,\
  \min_{0\leq\ell\leq N}
  \frac{\det K_{\ell,N-\ell}}
  {6(\operatorname{Tr}K_{\ell,N-\ell})^{N-\ell}
                       \operatorname{Tr}R_{\ell,N-\ell}}\right\}>0.
 \label{eq:f15v85:constant}
\end{equation}

\begin{proposition}[All-field ellipticity with an explicit band cost]
\label{f15v85:ellipticity}
For every real field \(F\),
\begin{equation}
 c_N(1+|F|)^{-2N} I
       \ \leq\ \widehat H_N(F)\ \leq\ D_N I,
 \qquad \operatorname{Tr}\widehat H_N(F)=D_N.
 \label{eq:f15v85:ellipticity}
\end{equation}
For \(N=1\), the displayed construction gives \(c_1=8/3375\).
The constants depend on \(N\); no uniform-in-band assertion is made.
\end{proposition}

\begin{proof}
In an angular sector write the radial polynomial as
\(p(1-t)=\sum_{j=0}^d a_j(1-t)^j\).
Its Haar squared norm is \(a^*Ra\). For \(s\geq1\), set
\(y=s(1-t)\) and
\[
 J_s=\int_0^{2s}e^{-y}y^{1/2}(2-y/s)^{1/2}\,dy.
\]
The weighted squared norm, normalized by \(z(s)\), is
\[
 \frac{s^{-\ell}}{J_s}
 \int_0^{2s}e^{-y}y^{\ell+1/2}(2-y/s)^{\ell+1/2}
                     |p(y/s)|^2\,dy.
\]
We have \(J_s\leq\sqrt2\,\Gamma(3/2)<2\).
On \(0\leq y\leq1\), \(2-y/s\geq1\) and \(e^{-y}>1/3\).
With \(E_s=\operatorname{diag}(1,s^{-1},\ldots,s^{-d})\), this is
at least
\[
 \frac16s^{-\ell}a^*E_s K E_s a
 \geq\frac{\det K}{6(\operatorname{Tr}K)^d}
                      s^{-\ell-2d}|a|^2.
\]
The last inequality uses
\(\lambda_{\min}(K)\geq\det K/(\operatorname{Tr}K)^d\).
Since \(a^*Ra\leq\operatorname{Tr}(R)|a|^2\) and
\(\ell+2d=2N-\ell\leq2N\), the lower bound follows for \(s\geq1\).
For \(0\leq s\leq1\), the density \(e^{st}/z(s)\geq e^{-2}>1/9\);
this gives the remaining range and the factor \((1+s)^{-2N}\).
Sum the orthogonal angular sectors and conjugate by \(U_v\).

The reproducing-kernel diagonal of the complete band is \(D_N\).
Integrating it against the normalized positive density proves
the trace identity. Positivity then gives the upper bound.
Direct evaluation of the two rational candidates at \(N=1\)
gives \(8/3375\) and \(4/45\), proving the stated value.
\end{proof}

This has an immediate native-operator consequence. Define
\(w(V)=b_{\widehat e}(V)\sqrt{Z_e(V)}\) and
\[
 \mathcal S_N
   =D\bigl(I_{\mathcal F_N}\otimes
        C_{\widehat e}^{1/2}M_{w^2}C_{\widehat e}^{1/2}\bigr)D .
\]
V84's local Gram matrix satisfies, on the same combined
Gauss-invariant space,
\begin{equation}
 \frac{c_N}{(1+4\gamma)^{2N}}\mathcal S_N
       \leq\mathcal G_N\leq D_N\mathcal S_N.
 \label{eq:f15v85:operator-comparison}
\end{equation}
Indeed sandwich \eqref{eq:f15v85:full-weight} and
\eqref{eq:f15v85:ellipticity} by \(DC_{\widehat e}^{1/2}\).
This is an actual operator inequality, but its deteriorating
comparison factor is not a uniform spectral-ratio estimate.

\subsection{Sharp fixed-band strong-field conditioning}

\begin{theorem}[Orders of every axial eigenvalue]
\label{f15v85:orders}
Fix \(N\). List the eigenvalues of \(A_{\ell,N}(s)\) in
decreasing order as \(\lambda_{\ell,k}(s)\),
\(0\leq k\leq N-\ell\). Then
\begin{equation}
 \lambda_{\ell,k}(s)\asymp_N s^{-(\ell+2k)}
                       \quad(s\longrightarrow\infty).
 \label{eq:f15v85:orders}
\end{equation}
Every such eigenvalue is repeated \(2\ell+1\) times in \(A_N\).
In particular,
\begin{align}
 \lambda_{\max}(A_N(s))&\longrightarrow D_N, \notag\\
 \lambda_{\min}(A_N(s))
    &\sim \kappa_N s^{-2N}, \qquad
 \kappa_N=\frac{N!(3/2)_N(3)_{2N}}
                     {2^{2N}(3/2)_{2N}}>0, \label{eq:f15v85:sharp}\\
 \det A_N(s)&\sim a_N s^{-ND_N}\quad\hbox{for some }a_N>0. \notag
\end{align}
Here \((a)_j=a(a+1)\cdots(a+j-1)\), with \((a)_0=1\).
Thus the condition number is asymptotic to
\((D_N/\kappa_N)s^{2N}\). These are fixed-\(N\) statements.
\end{theorem}

\begin{proof}
In the monomial basis used above the weighted radial Gram
matrix is
\[
 Q_{\ell,d}(s)=s^{-\ell} E_s M_{\ell,d}(s)E_s .
\]
After the same substitution \(y=s(1-t)\), dominated convergence
gives
\[
 (M_{\ell,d}(s))_{ij}
  \longrightarrow
  2^\ell\frac{\Gamma(\ell+i+j+3/2)}{\Gamma(3/2)}.
\]
For domination use
\((2-y/s)^{\ell+1/2}\leq2^{\ell+1/2}\) on its support and
the integrable Gamma density with each fixed polynomial power.
The limiting matrix is positive definite, being a monomial
Gram matrix in a strictly positive Gamma measure.
Its eigenvalues are therefore bounded above and below away
from zero for large \(s\).

The Haar metric on these coefficients is the fixed positive
matrix \(R_{\ell,d}\). Congruence by its inverse square root
identifies the axial block's eigenvalues. Min--max comparison
with \(E_s^2\), and the fixed invertible change of basis, prove
\eqref{eq:f15v85:orders}. Determinants give an asymptotic
positive constant in each block. The total exponent is
\[
 \sum_{\ell=0}^N(2\ell+1)
           \sum_{k=0}^{N-\ell}(\ell+2k)
 =N\sum_{\ell=0}^N(2\ell+1)(N-\ell+1)=ND_N.
\]

For the sharp least eigenvalue only \(\ell=0,d=N\) contributes:
all other blocks have their smallest exponent strictly below
\(2N\). In that block let \(M_\infty\) be the limiting Gamma
moment matrix and \(e_N\) the last coordinate. Then
\[
 s^{-2N}Q_{0,N}(s)^{-1}
       \longrightarrow (M_\infty^{-1})_{NN}e_Ne_N^* .
\]
The inverse of the Haar-orthonormal form is conjugate to
\(R^{1/2}Q^{-1}R^{1/2}\), whose largest eigenvalue after this
rescaling tends to \((M_\infty^{-1})_{NN}R_{NN}\).
The reciprocal of \((M_\infty^{-1})_{NN}\) is the squared norm
of the degree-\(N\) monic orthogonal polynomial for the
normalized Gamma law of shape \(3/2\). It equals
\(N!(3/2)_N\): the Laguerre norm and leading coefficient in
DLMF Table 18.3.1 give this directly. The elementary beta
integral gives
\[
 R_{NN}=\mathbb E_H(1-q_0)^{2N}
                    =2^{2N}(3/2)_{2N}/(3)_{2N}.
\]
Taking the reciprocal proves \eqref{eq:f15v85:sharp}.

Finally the normalized density concentrates at the pole.
Since the finite basis is continuous, \(A_N(s)\) tends to the
rank-one matrix of evaluations at that pole. Its trace is
\(D_N\) by the reproducing-kernel identity. This proves the
largest-eigenvalue limit and completes the proof.
\end{proof}

For example \(\kappa_1=6/5\) and \(\kappa_2=20/7\).
The small pointwise eigenvalues describe nearly vanishing
polynomials in the sharply concentrated one-link weight.
They are not small eigenvalues of the complete native transfer,
nor a counterexample to a physical mass gap. The other-link
kinetic operator acts on the coefficients of such polynomials.

The local magnetic comparison is already nonuniform for any
fixed \(N\geq1\). Moreover V84 requires a growing band for its
vanishing weak-coupling truncation error. Substitution of
\(N=N(\gamma)\) into the fixed-band asymptotic theorem is not
valid without new uniform estimates. The explicit global
constant \eqref{eq:f15v85:constant} remains valid for every \(N\),
but its band dependence must be retained.

\subsection{The first charged block in closed form}

For \(N=1\) use the orthonormal real basis
\((1,2q_0,2q_1,2q_2,2q_3)\).
Put \(a(s)=I_2(s)/I_1(s)\), \(b(s)=1-3a(s)/s\).
The continuous values at zero are understood.

\begin{proposition}[Radial and transverse channels of the five-dimensional weight]
\label{f15v85:first}
For \(F=sv\ne0\),
\begin{equation}
 \widehat H_1(F)=
 \begin{pmatrix}
  1&2a(s)v^{\mathsf T}\\
  2a(s)v&
       4b(s)vv^{\mathsf T}
         +4a(s)s^{-1}(I-vv^{\mathsf T})
 \end{pmatrix}.
 \label{eq:f15v85:first}
\end{equation}
Its transverse eigenvalue is \(4a(s)/s\), with multiplicity
three; its other two eigenvalues are those of
\[
 \begin{pmatrix}1&2a(s)\\2a(s)&4b(s)\end{pmatrix}.
\]
The determinant of this last matrix is \(4a'(s)>0\).
As \(s\to\infty\), the three transverse eigenvalues are
\(4s^{-1}(1+o(1))\), and the radial pair consists of
\(5+o(1)\) and \((6/5)s^{-2}(1+o(1))\).
\end{proposition}

\begin{proof}
Differentiate \(z(|F|)\) once and twice in Cartesian components.
The normalized mean is \(a(s)v\); the second moment is
\(b(s)vv^{\mathsf T}+(a(s)/s)(I-vv^{\mathsf T})\).
The Bessel recurrence gives
\(a'=1-a^2-3a/s=b-a^2\), also the strictly positive variance
of \(q\cdot v\). These formulas give the displayed matrix
and its determinant without suppressing the neutral--charged
off-diagonal entries.
The Gamma rescaling in the preceding theorem yields
\(a\to1\) and \(a'\sim(3/2)s^{-2}\).
The trace of the radial block tends to five, giving its smaller
eigenvalue by determinant divided by its larger eigenvalue.
At zero, \(a(s)/s\to1/4\), and the full matrix tends to \(I\).
\end{proof}

Thus retaining only the scalar \(Z_e\) misses both a
neutral--fundamental coupling and parametrically different
radial and transverse weights. Restoring the fixed-band kinetic
diagonal \(D\) does not remove these orders when
\(s\to\infty\) with \(s\leq4\gamma\): then \(\gamma\to\infty\)
and every fixed \(r_n(\gamma)\to1\).
This last limit follows from concentration of the normalized
one-link convolution at the identity. Congruence by \(D\)
changes each ordered positive eigenvalue by a factor \(1+o(1)\),
by comparing \(D^2\) to \(I\). Again this is not uniform in \(N\).

\subsection{Why the moving frame must remain in the transfer}

Let \(\mathsf C=I_{\mathcal F_N}\otimes C_{\widehat e}\) and
write \(H(V)=\mathcal H^{(N)}(V)\).
V84 gives \(\mathcal G_N=D\mathsf C^{1/2}M_H\mathsf C^{1/2}D\).
A second use of the already established Gram principle, with
\(W=M_{H^{1/2}}\mathsf C^{1/2}D\), identifies its nonzero
spectrum with
\[
 \mathcal L_N=M_{H^{1/2}}D^2\mathsf C M_{H^{1/2}}.
\]
All statements restrict to the combined Gauss-invariant space:
these factors are equivariant, and the restricted \(W,W^*\)
give the same positive nonzero eigenvalues.

Choose the measurable unitary field \(U(V)=U_{v(V)}\) above,
with \(U=I\) at \(F=0\), and let \(\mathcal U\) be its
multiplication operator. If \(c_{\widehat e}(V,V')\) is the
original scalar kinetic kernel on the other links, then
\(\mathcal U^*\mathcal L_N\mathcal U\) has the exact matrix kernel
\begin{equation}
 \begin{split}
 &w(V)A_N(s(V))^{1/2}\,
     c_{\widehat e}(V,V')\,D^2 U(V)^*U(V')\\
 &\hspace{35mm}\cdot A_N(s(V'))^{1/2}w(V').
 \end{split}
 \label{eq:f15v85:transport-kernel}
\end{equation}
The physical space after this conjugation is
\(\mathcal U^*\mathcal H_{\rm combined}\), not a collection of
individually gauge-invariant scalar components.

Formula \eqref{eq:f15v85:transport-kernel} follows simply by
inserting \(H=w^2UA_NU^*\) and using \([D,U]=0\).
In general \(U(V)^*U(V')\ne I\); it cannot be deleted.
One may further diagonalize each radial block, but then its
field-dependent diagonalizing matrices also enter the
two-point kernel and need not commute with \(D\).
At \(N=0\) there is only the trivial carrier, and this formula
reduces exactly to V84's scalar neutral operator.

The apparent loss of smoothness of \(U\) at zero field belongs
to the coordinate choice; it is absent from
\eqref{eq:f15v85:normalized-matrix}. No differentiated moving
frame has been used. For derivative estimates near zero one
must use the Cartesian matrix or control a genuine chart
transition, not divide by \(s\) without a bound.

\subsection{What changes in the attack and what is still missing}

The all-band magnetic matrix is now explicitly populated by
one-dimensional integrals, with exact carrier multiplicities.
It has a rigorous all-field floor and a sharp fixed-band
conditioning law. The scalar neutral multiplier alone cannot
stand in for these matrix weights. The corresponding native
kernel is \eqref{eq:f15v85:transport-kernel}, including its
orientation transport and physical subspace.

The next target is an estimate for this matrix-weighted
interacting operator, or a quantitatively equivalent
multi-link retained operator, at a band chosen from V84's
actual tail budget. Pointwise diagonalization and pointwise
ellipticity do not bound its leading physical eigenvalue ratio.
One must control the spatial variation and coupling of the
retained channels together with the charged tail, at the
prescribed physical clock. The independent continuum
construction and nontriviality requirements also remain open.

\subsection{Verification and append record}

The checker
\begin{center}\small
\path{scripts/verify_ym_f15_v85_axial_magnetic_gram.py}
\end{center}
reruns V84's diagnostic chain and adds exact rational tests of
the separated harmonic equations, Haar normalizations, band
dimensions, positive Gamma limits, polynomial floor constants,
sharp Gamma--Schur constants, and Cartesian magnetic moments
in nonaxial quaternion fields. Outward Bessel intervals check
the first charged block. Finite matrix analogues check the
second Gram factorization and show a changed spectrum if the
moving-frame transport is omitted. These last analogues are
not substitutes for a calculation of the actual lattice gap.

All 1,926 new exact-rational fixtures passed, together with the
inherited diagnostic chain. The final checker has 10,243 bytes
and SHA--256
\begin{center}\scriptsize\ttfamily
2C8FAB535793DAC573C1E54EBDB7460D6E9C11BB77C1E95F4BDD4F3E5AB2EA59.
\end{center}

The analytic proofs, rather than finite tests, establish the
all-band results and their fixed-band asymptotics. No continuum
limit or positive physical spectral margin is certified by
the tests. The predecessor V84 has 1,901,998 bytes and SHA--256
\begin{center}\scriptsize\ttfamily
617C94F2C2B503B398C9EB2F0DB8DC60C9CCE8FEA114746EA9A62C12EEF7C697.
\end{center}
Removing this section and restoring the title version recovers
that predecessor byte for byte. V85 replaces only the active
main note; all archives and companions are preserved.
Only \path{.tex} sources are stored in \path{latex_notes};
build products belong in \path{artifacts/ym_f15_v85_texcheck/}.

\begin{firewall}
V85 does not solve the retained interacting spectral problem.
It calculates the charged magnetic weight and its exact native
placement, identifies its nonuniform conditioning, and retains
the transport lost by a scalar replacement.
The full interacting continuum Yang--Mills mass gap remains
unproved, and the original goal remains active.
\end{firewall}

\section{V86: compact local transfer and two-endpoint channels}
\label{f15v86:context}

V85 computed the magnetic matrix but did not estimate the
interacting retained spectrum. This continuation returns one
step upstream: keep the complete compact one-link kinetic
operator, multiply by the actual magnetic factors at its two
endpoints, and compute the resulting local transfer block.
This avoids using a fixed-band asymptotic at a growing band.

The results below concern literal blocks of the native
unprojected kernel, before the combined Gauss restriction.
They prove a uniform Hilbert--Schmidt scaling limit on a
specified nonzero-staple region, its complete limiting
singular spectrum, and an all-angle decay estimate.
The limiting leading singular line depends on both endpoints.
Thus local fast-mode control is acquired, but its assembly
into a physical global gap remains an additional problem.

The prior harmonic-oscillator word bounds in f7 V75, the
auxiliary Mehler formula in f8, and the Gaussian comparison
of V64 concern different objects. No such theorem is
relabelled as a proof about this compact Wilson block.
The last companion checkpoint was V85; the next remains V90.

\subsection{The literal local block and its place in the native kernel}

Identify \(\SU(2)\) with the unit quaternions \(S^3\), with
probability Haar measure \(dH\). For \(\gamma>0\), write
\[
 c_0(\gamma)=\int_{S^3}e^{-\gamma(1-q_0)}\,dH(q),
 \qquad
 C_\gamma(q,q')=\frac{e^{-\gamma(1-q\cdot q')}}{c_0(\gamma)} .
\]
This is the original normalized one-link Wilson convolution.
For a unit quaternion \(v\) and \(0\leq\rho\leq4\), put
\[
 m_{\gamma,\rho,v}(q)
       =e^{-\gamma\rho(1-q\cdot v)/2},\qquad
 K_\gamma(\rho,v;\sigma,w)
       =M_{m_{\gamma,\rho,v}}C_\gamma M_{m_{\gamma,\sigma,w}} .
 \tag{86.1}
\]
For equal parameters this is a positive compact self-adjoint
transfer. For unequal parameters it is generally not
self-adjoint; its singular values, not its eigenvalues, will
be used.

To verify its native placement, use the deleted-edge staples
of V84--V85. If
\[
 \rho(V)=|S_e(V)|,\qquad
 v(V)=\overline{S_e(V)}/|S_e(V)|,
\]
then
\[
 b(q,V)=a(V)m_{\gamma,\rho(V),v(V)}(q),\qquad
 a(V)=b_{\widehat e}(V)e^{-\gamma(4-\rho(V))/2}.
 \tag{86.2}
\]
At \(\rho=0\) the multiplier is independent of \(v\).
The full unprojected native transfer kernel, viewed as an
operator-valued kernel in the other links, is exactly
\begin{equation}
 c_{\widehat e}(V,V')a(V)a(V')\,
 K_\gamma\bigl(\rho(V),v(V);\rho(V'),v(V')\bigr).
 \label{eq:f15v86:native-placement}
\end{equation}
The global physical operator is the restriction of that
equivariant operator to the combined Gauss-invariant space.
Freezing \(V,V'\) does not itself define a physical one-link
Hilbert space. In particular, three local oscillator
directions below are not being declared three physical
particle states.

\subsection{A norm limit over the entire compact link space}

Fix \(0<\rho_*\leq4\) and \(B<\infty\).
Take \(\rho,\sigma\in[\rho_*,4]\) and
\(\xi\in\mathbb R^3\), \(|\xi|\leq B\).
A common quaternion rotation reduces \(v\) to \(1\);
consider the other direction
\[
 w_\gamma=\operatorname{Exp}(\xi/\sqrt\gamma),\qquad
 \operatorname{Exp}(x)
     =\left(\cos|x|,\frac{\sin|x|}{|x|}x\right).
\]
For sufficiently large \(\gamma\), this parametrizes every
pair of directions at geodesic distance at most
\(B/\sqrt\gamma\).

Let \(B_\gamma=\{x\in\mathbb R^3:|x|<\pi\sqrt\gamma\}\).
The exponential chart covers the whole sphere except the
antipode, a Haar-null point, and its density is
\begin{equation}
 J_\gamma(x)
  =\frac1{2\pi^2\gamma^{3/2}}
       \left(\frac{\sin(|x|/\sqrt\gamma)}
                       {|x|/\sqrt\gamma}\right)^2.
 \label{eq:f15v86:jacobian}
\end{equation}
Define an isometry \(U_\gamma:L^2(S^3,dH)\to L^2(\mathbb R^3,dx)\)
by multiplication by \(J_\gamma^{1/2}\) in this chart and
extension by zero off \(B_\gamma\). It is not onto the whole
Euclidean space. Conjugation by \(U_\gamma,U_\gamma^*\) adds
only a zero subspace to the local block.

\begin{theorem}[Uniform full-compact Gaussian limit]
\label{f15v86:compact-limit}
Let \(\widetilde K_\gamma
 =U_\gamma K_\gamma(\rho,1;\sigma,w_\gamma)U_\gamma^*\).
On \(L^2(\mathbb R^3)\) define the integral operator with kernel
\begin{equation}
 G_{\rho,\sigma,\xi}(x,y)
  =(2\pi)^{-3/2}
     \exp\left[-\frac{\rho|x|^2}{4}
               -\frac{|x-y|^2}{2}
               -\frac{\sigma|y-\xi|^2}{4}\right].
 \label{eq:f15v86:gaussian-kernel}
\end{equation}
There are finite constants \(C,\gamma_0\), depending only on
\(\rho_*,B\), such that
\begin{equation}
 \sup_{\substack{\rho,\sigma\in[\rho_*,4]\\|\xi|\leq B}}
 \|\widetilde K_\gamma-G_{\rho,\sigma,\xi}\|_{\rm HS}
                  \leq C/\gamma,\qquad \gamma\geq\gamma_0.
 \label{eq:f15v86:hs-limit}
\end{equation}
In particular this is an operator-norm limit, with no
unestimated compact complement.
\end{theorem}

\begin{proof}
For \(q_x=\operatorname{Exp}(x/\sqrt\gamma)\), the transformed
kernel on \(B_\gamma^2\) is
\[
 \frac{\sqrt{J_\gamma(x)J_\gamma(y)}}{c_0(\gamma)}
                         e^{-A_\gamma(x,y)},
\]
where
\[
 A_\gamma=\gamma\left[
 \frac\rho2(1-(q_x)_0)+(1-q_x\cdot q_y)
              +\frac\sigma2(1-q_y\cdot w_\gamma)\right].
\]
All three summands are nonnegative.
The Haar integral with \(u=\gamma(1-q_0)\) gives
\[
 c_0(\gamma)=\frac2{\pi\gamma^{3/2}}
       \int_0^{2\gamma}e^{-u}u^{1/2}(2-u/\gamma)^{1/2}\,du
 =\sqrt{\frac2\pi}\gamma^{-3/2}(1+O(\gamma^{-1})).
 \tag{86.3}
\]
For the error estimate, subtract \(\sqrt2\) inside the
integral. On its support the absolute difference is at most
\(u/(\sqrt2\gamma)\), and the omitted Gamma tail is
exponentially small. Thus the remainder is bounded as stated.
A restriction to \(0\leq u\leq1\) also gives a positive
uniform lower bound for \(\gamma^{3/2}c_0(\gamma)\) when
\(\gamma\geq1\).

Put
\[
 A(x,y)=\rho|x|^2/4+|x-y|^2/2+\sigma|y-\xi|^2/4.
\]
The elementary Taylor inequalities for cosine and
\(\sin t/t\), applied to the explicit quaternion dot products,
give on \(B_\gamma^2\)
\begin{equation}
 |A_\gamma-A|
   \leq \frac{C_B}{\gamma}(1+|x|+|y|)^4 .
 \label{eq:f15v86:action-error}
\end{equation}
For example,
\[
 \gamma(1-q_x\cdot q_y)
  =\frac{|x-y|^2}{2}
       +O\!\left(\frac{(|x|+|y|)^4}{\gamma}\right).
\]
This bound is valid throughout the chart, not only on a
fixed-radius subset. To see this directly, expand each
cosine to quadratic order, each sine quotient to constant
order, use \(|\cos t|\leq1\), \(|\sin t/t|\leq1\), and their
global quadratic and quartic remainders. Products of the
quadratic errors are bounded by the displayed fourth-degree
polynomial. The same calculation with \(y,\xi\) treats the
last summand.

We also need a uniform integrable bound, not just this
expansion. For \(0\leq t\leq\pi\),
\(1-\cos t\geq2t^2/\pi^2\).
The triangle inequality for spherical distance gives
\[
 d(q_y,w_\gamma)\geq
                (|y|-B)_+/\sqrt\gamma.
\]
Using \((r-B)_+^2\geq r^2/2-B^2\), the two magnetic summands
alone imply
\[
 A_\gamma(x,y)\geq
                c_{\rho_*}(|x|^2+|y|^2)-C_B.
\]
The same type of bound holds for \(A\).
Consequently
\[
 |e^{-A_\gamma}-e^{-A}|
 \leq |A_\gamma-A|(e^{-A_\gamma}+e^{-A})
 \leq \frac{C}{\gamma}(1+|x|+|y|)^4
                         e^{-c(|x|^2+|y|^2)} .
\]
The constants are uniform in the declared parameter set.

Finally \eqref{eq:f15v86:jacobian} and (86.3) imply
\[
 \frac{\sqrt{J_\gamma(x)J_\gamma(y)}}{c_0(\gamma)}
  =(2\pi)^{-3/2}
       +O\!\left(\frac{1+|x|^2+|y|^2}{\gamma}\right)
\]
on \(B_\gamma^2\); the prefactor is also uniformly bounded
there. Combine the last two estimates and square-integrate
the polynomial Gaussian majorant. Off \(B_\gamma^2\) the
transformed compact kernel is zero, and the Euclidean kernel
has an exponentially small \(L^2\) tail. This proves
\eqref{eq:f15v86:hs-limit}.
\end{proof}

This proof explicitly controls configurations near the
antipode. It does not assume that a small-field chart already
has probability one under a different measure.

\subsection{The complete limiting singular spectrum}

Set
\[
 \begin{aligned}
 &a=1+\rho/2,\quad d=1+\sigma/2,\quad \Delta=ad-1,\\
 &Q=\sqrt{ad}-\sqrt{ad-1}\in(0,1),\qquad
 \kappa=\frac{\rho\sigma}{8\Delta}
             =\left(\frac4\rho+2+\frac4\sigma\right)^{-1}.
 \end{aligned}
 \tag{86.4}
\]

\begin{theorem}[Two-endpoint Gaussian singular channels]
\label{f15v86:mehler}
The singular values of \(G_{\rho,\sigma,\xi}\), grouped by
nonnegative integer level \(n\), are
\begin{equation}
 e^{-\kappa|\xi|^2}Q^{n+3/2},
          \qquad\hbox{multiplicity }\binom{n+2}{2}.
 \label{eq:f15v86:singular-spectrum}
\end{equation}
Its leading left and right singular vectors are normalized
Gaussians with respective centers and widths
\begin{align}
 x_0&=\frac{\sigma}{2\Delta}\xi,&
 y_0&=\frac{a\sigma}{2\Delta}\xi, \notag\\
 \omega_L&=\sqrt{a/d}\sqrt{\Delta},&
 \omega_R&=\sqrt{d/a}\sqrt{\Delta}. \label{eq:f15v86:ground-lines}
\end{align}
Here width \(\omega\) means the vector
\((\omega/\pi)^{3/4}e^{-\omega|x-x_0|^2/2}\).
\end{theorem}

\begin{proof}
Completing the square in \eqref{eq:f15v86:gaussian-kernel}
gives the exact identity
\[
 A(x,y)=\kappa|\xi|^2+
       \frac a2|x-x_0|^2+\frac d2|y-y_0|^2
                      -(x-x_0)\cdot(y-y_0).
\]
Separate unitary translations on the input and output remove
the centers. Opposite dilations
\[
 x-x_0=(d/a)^{1/4}X,\qquad
 y-y_0=(a/d)^{1/4}Y
\]
then replace both diagonal coefficients by \(\sqrt{ad}\).
The product of their Jacobian factors is one. These are
left and right unitary equivalences, so they preserve
singular values; they are not asserted to preserve the
eigenvalues of a nonsymmetric block.

It remains to diagonalize
\[
 (2\pi)^{-3/2}
       e^{-A_0(|X|^2+|Y|^2)/2+X\cdot Y},
 \qquad A_0=\sqrt{ad}>1.
\]
Write \(\omega=\sqrt{A_0^2-1}\).
A one-dimensional Gaussian integration sends
\(e^{-\omega Y^2/2}\) to
\((A_0+\omega)^{-1/2}e^{-\omega X^2/2}\).
The Hermite generating function, or
\href{https://dlmf.nist.gov/18.18.E28}{DLMF 18.18.28},
shows that each successive oscillator polynomial adds the
factor \(Q=(A_0+\omega)^{-1}\).
Completeness of the Hermite basis gives all eigenvalues,
not only trial Rayleigh quotients.
Taking three tensor factors gives \(Q^{n+3/2}\) and the
number \(\binom{n+2}{2}\) of triples with total degree \(n\).
Restore the scalar \(e^{-\kappa|\xi|^2}\), translations and
dilations. The ground widths become
\(\omega/(d/a)^{1/2}\) and \(\omega/(a/d)^{1/2}\), as stated.
\end{proof}

\begin{proposition}[Consequences for the full compact local block]
\label{f15v86:local-spectrum}
In the parameter region of Theorem~\ref{f15v86:compact-limit},
the ordered compact singular values differ from those in
\eqref{eq:f15v86:singular-spectrum} by at most \(C/\gamma\).
In particular
\begin{equation}
 \frac{s_1(K_\gamma)}{s_0(K_\gamma)}
             =Q+O(\gamma^{-1})
 \label{eq:f15v86:compact-ratio}
\end{equation}
uniformly in that region. For \(\rho=\sigma\), \(v=w\),
these are positive eigenvalues and
\[
 \frac{\lambda_1(K_\gamma(\rho,v;\rho,v))}
      {\lambda_0(K_\gamma(\rho,v;\rho,v))}
   =1+\frac\rho2-\sqrt{(1+\rho/2)^2-1}+O(\gamma^{-1}).
 \tag{86.5}
\]
The normalized leading left and right singular projectors,
after the same isometric embedding, converge in operator norm
at rate \(O(\gamma^{-1})\) to those in
\eqref{eq:f15v86:ground-lines}.
\end{proposition}

\begin{proof}
Ordered singular values are Lipschitz with respect to the
operator norm, by their min--max characterization.
The limiting top value in \eqref{eq:f15v86:singular-spectrum}
has a positive lower bound on the compact parameter set, so
division proves \eqref{eq:f15v86:compact-ratio}.
For equal endpoints the compact and limiting kernels are
positive self-adjoint, giving (86.5).

For the projectors apply the same perturbation estimate to
\(\widetilde K_\gamma\widetilde K_\gamma^*\) and
\(\widetilde K_\gamma^*\widetilde K_\gamma\).
Their errors are \(O(\gamma^{-1})\), and each limiting top
eigenvalue is separated from the rest by a uniform positive
amount. For completeness, if \(u_\gamma\) is its normalized
perturbed top eigenvector, project its residual equation
onto the complement of the limiting top line. The inverse
of the limiting operator minus the perturbed eigenvalue on
that complement has norm at most twice the reciprocal gap.
Thus \(\|(I-P_0)u_\gamma\|=O(\gamma^{-1})\), which is the
rank-one projector estimate.
\end{proof}

The statement is a local native-transfer result, not an
auxiliary diffusion estimate. Nevertheless no physical time
unit is inferred from (86.5), and no full lattice spectral
ratio follows by multiplying these local ratios.

\subsection{Uniform decay for arbitrary direction separation}

The near-diagonal theorem has bounded \(\xi\). The following
bound handles all pairs of directions without expanding them.

\begin{proposition}[Global angular suppression of the native block]
\label{f15v86:far-frames}
For \(0<\rho,\sigma\leq4\), let \(\theta=d(v,w)\in[0,\pi]\) and
\[
 \eta_{\rho,\sigma}
   =\frac1{\pi^2(1/\rho+1/2+1/\sigma)}.
\]
For every \(\gamma>0\),
\begin{equation}
 \|K_\gamma(\rho,v;\sigma,w)\|
 \leq
 \frac{\sqrt{c_0(\gamma\rho/2)c_0(\gamma\sigma/2)}}{c_0(\gamma)}
                     e^{-\gamma\eta_{\rho,\sigma}\theta^2/2}.
 \label{eq:f15v86:far-bound}
\end{equation}
For \(\rho,\sigma\in[\rho_*,4]\) and \(\gamma\geq1\), the
prefactor is uniformly bounded. Thus these local blocks
have a uniform Gaussian off-diagonal bound on the angular
scale \(\gamma^{-1/2}\).
\end{proposition}

\begin{proof}
Put \(r_1=d(v,q)\), \(r_2=d(q,q')\), \(r_3=d(q',w)\).
The native exponent divided by \(\gamma\), denoted \(E\),
satisfies
\[
 E\geq\pi^{-2}(\rho r_1^2+2r_2^2+\sigma r_3^2)
   \geq\eta_{\rho,\sigma}(r_1+r_2+r_3)^2
   \geq\eta_{\rho,\sigma}\theta^2.
\]
The middle inequality is weighted Cauchy--Schwarz, and the
last is the spherical triangle inequality.
Consequently \(2E\geq E+\eta_{\rho,\sigma}\theta^2\).
In the squared Hilbert--Schmidt integral use this inequality
and then discard only the nonnegative kinetic term in the
remaining \(E\). The two Haar integrals are
\(c_0(\gamma\rho/2)\) and \(c_0(\gamma\sigma/2)\).
Taking square roots proves \eqref{eq:f15v86:far-bound}.
The two-sided bound
\(c_0(t)\asymp(1+t)^{-3/2}\), from (86.3), its lower-bound
argument, and compactness for \(0\leq t\leq1\), gives the
uniform prefactor assertion.
\end{proof}

The scalar factors in \eqref{eq:f15v86:native-placement}
must still be restored when this estimate is used in the
full operator. They are not omitted normalizers.

\subsection{Why one endpoint does not determine the leading channel}

Even in the controlled near-diagonal region, the leading
singular vectors are genuinely two-endpoint objects.

\begin{proposition}[Order-one variation of the limiting leading line]
\label{f15v86:pair-dependence}
Fix \(\rho=\sigma>0\), \(a=1+\rho/2\),
\(\omega=\sqrt{a^2-1}\), and \(\xi\ne0\).
Keep the left endpoint fixed at \(v=1\), and compare right
directions \(1\) and \(\operatorname{Exp}(\xi/\sqrt\gamma)\).
The norm distance between the two leading left singular
projectors tends to
\begin{equation}
 \left(1-\exp\left[
       -\frac{\omega|\xi|^2}{2(a+1)^2}\right]\right)^{1/2}>0.
 \label{eq:f15v86:line-separation}
\end{equation}
In particular a rank-one projector chosen solely from the
left endpoint cannot approximate both leading lines with
an error tending to zero.
\end{proposition}

\begin{proof}
Formula \eqref{eq:f15v86:ground-lines} gives equal width
\(\omega\) and centers \(0\) and \(\xi/(a+1)\).
The overlap of normalized equal-width Gaussians with center
difference \(h\) is \(e^{-\omega|h|^2/4}\), by completing the
square. The norm distance of their rank-one projectors is
the square root of one minus the squared overlap.
Use the projector convergence in
Proposition~\ref{f15v86:local-spectrum}. The last claim follows
from the triangle inequality: any single candidate has error
at least half their positive limiting separation for one
of the two pairs.
\end{proof}

This is not a failure of the local spectral limit. It shows
why a collection of separately frozen local ground states
does not already supply a compatible global vacuum bundle.
Larger oscillator channels, or an explicitly two-endpoint
effective kernel, remain possible. Their kinetic transport
and Gauss constraints must be calculated, not suppressed.

\subsection{Zero staples and band removal retain their real costs}

The lower staple bound is essential. At \(\rho=\sigma=0\),
\[
 K_\gamma(0,v;0,w)=C_\gamma,\qquad
 \lambda_0=1,\qquad \lambda_1=r_1(\gamma)=I_2(\gamma)/I_1(\gamma)
                                      \longrightarrow1.
 \tag{86.6}
\]
Zero spatial staples are actual compact configurations.
On a cubic torus of side at least four, the four plaquettes
incident to an edge have four distinct opposite parallel
edges. Set all other links to the identity, two of these
four opposite edges to \(-I\), and the other two to \(I\).
Their staples sum to zero. This construction places no
probability estimate on such configurations and is not a
counterexample to the full physical mass gap. It does
exclude a uniform frozen-block ratio bound below one over
every exterior configuration at large \(\gamma\).

There is no hidden fixed Peter--Weyl cutoff in the
full-compact limit. If desired, the kinetic band can be
restored with an explicit operator error. Since
\(0<r_{n+1}(\gamma)\leq r_n(\gamma)\), the two-sided truncated
block satisfies
\[
 \|K_\gamma-M_{m_{\gamma,\rho,v}}
          C_{\gamma,N}M_{m_{\gamma,\sigma,w}}\|
 \leq r_{N+1}(\gamma)
 \leq
 \exp\left[-\frac{(N+1)(N+3)}{2\gamma+2N+3}\right].
 \tag{86.7}
\]
The first inequality uses both magnetic sup norms at most
one and the exact diagonal kinetic tail norm. The second
is V84's coefficient bound.
For example \(N/\sqrt\gamma\to\infty\) makes this error
vanish, without inserting a growing \(N\) into V85's
fixed-band asymptotic. On the controlled near-diagonal
parameter set the actual local leading singular value is
bounded below, so this also gives a relative local error.
It is not a lower bound for the full volume-dependent
leading eigenvalue.

\subsection{The next interacting calculation}

We now have, on a declared good-staple region, the actual
compact local transfer's limiting singular spectrum,
its leading two-endpoint modes, and a bound for large
relative direction. The proof is independent of lattice
volume because it is a statement about a frozen one-link
block. That independence does not make it a full-volume
gap theorem.

The next useful calculation is a covariant finite oscillator
channel for \eqref{eq:f15v86:native-placement}, retaining the
two-endpoint translation and width changes. The Gaussian
kernel makes those channel matrix elements explicit.
To reach the physical goal, such control must also survive
integration over the other links, the regions with small
staple magnitude, and the combined Gauss projection,
with errors measured against the actual global leading
eigenvalue and the prescribed physical clock. Neither
pointwise local ratios nor the good-region assumption
prove those remaining estimates. The interacting
continuum construction and its nontriviality remain
separate unmet requirements.

\subsection{Diagnostics and preservation}

The script
\begin{center}\small
\path{scripts/verify_ym_f15_v86_compact_local_transfer.py}
\end{center}
reruns V85's inherited chain. New exact-rational fixtures
check the native exponential-chart action jets, Haar
density and normalization coefficients, Hermite
eigenpolynomials, multiplicities and trace identities,
two-endpoint energy completions, singular-vector widths
and overlaps, all-angle path-energy bounds, and literal
four-staple zero-field configurations.
The optional numerical mode computes radial quadrature
illustrations; they are explicitly non-certified and are
not used in the proof of any uniform estimate.

All 3,509 new exact-rational fixtures passed, as did the
inherited chain. Six non-certified quadrature illustrations
were also run, separately from those exact checks.
The final script has 10,061 bytes and SHA--256
\begin{center}\scriptsize\ttfamily
0F55039B8C575F586CA58CC8C66499F1B73A9C1ED74950E5D30CF2F522A9E7BF.
\end{center}

The predecessor V85 has 1,923,533 bytes and SHA--256
\begin{center}\scriptsize\ttfamily
E24D0BDF7F7797B27A3967D46FB3831C141FCEC09D4D16AE7507BC0D1C4F8FE6.
\end{center}
Its body is retained byte for byte; removing this section
and restoring the title version recovers that source.
V86 replaces only the active main note. Archives,
companions and all 145 previously protected files remain
unchanged. Only \path{.tex} sources belong in
\path{latex_notes}; build products are written under
\path{artifacts/ym_f15_v86_texcheck/}.

\begin{firewall}
V86 proves a full-compact native local transfer theorem,
not a physical global mass gap. Its fast local channels
have nonseparable endpoint dependence, and its uniform
ratio estimate requires nonzero staples. No averaging,
Gauss restriction, physical-time calibration or
interacting continuum limit is inferred without the
remaining estimates. The full goal remains active.
\end{firewall}

\section{V87: covariant endpoint channels and a relative near-field error}
\label{f15v87:context}

V86 proved a full-compact local transfer limit but also showed
that its leading singular line depends on both endpoints.
This section constructs larger channels that depend on each
endpoint separately. Their matrix elements retain the
two-endpoint translation, unequal widths, and the first
noncommuting correction. The omitted oscillator levels have
an explicit exponential bound. On a declared good, nearby-field
part of the full native kernel, this becomes an error relative
to the actual physical leading eigenvalue, without a
spectator-volume factor.

The last statement is restricted to that part of the kernel.
No estimate for its complementary small-staple or distant-field
part, and no global physical spectral gap, is asserted.
This is a continuation of V86, not a reuse of the affine
Gaussian conditional channel in f14 V69 or the oscillator
word bounds in f7 V75. The last companion review was V85;
the next is V90.

\subsection{Separate endpoint frames expose a noncommuting term}

Keep V86's local block
\(K_\gamma(\rho,v;\sigma,w)\), and fix
\(\rho,\sigma\in[\rho_*,4]\), \(\rho_*>0\).
Write
\[
 \epsilon=\gamma^{-1/2},\qquad
 w=v\operatorname{Exp}(\epsilon\xi),\qquad |\xi|\leq B.
\]
Use separate full exponential charts
\[
 q=v\operatorname{Exp}(\epsilon x),\qquad
 q'=w\operatorname{Exp}(\epsilon z).
\]
Both have V86's density \(J_\gamma\) on
\(\{|x|<\pi\sqrt\gamma\}\).
Let \(U_{\gamma,v}\) denote the corresponding Haar isometry
into \(L^2(\mathbb R^3)\), extended by zero outside that ball.

Define the Euclidean kernels
\begin{align}
 G(x,z)&=(2\pi)^{-3/2}e^{-A(x,z)},&
 A(x,z)&=\frac{\rho|x|^2}{4}
       +\frac{|x-z-\xi|^2}{2}+\frac{\sigma|z|^2}{4}, \notag\\
 G^{[1]}(x,z)&=\bigl[x\cdot(\xi\times z)\bigr]G(x,z).
 \label{eq:f15v87:gaussian-and-cubic}
\end{align}

\begin{theorem}[First nonabelian endpoint correction in compact norm]
\label{f15v87:compact-jet}
Uniformly on the declared parameter set,
\begin{equation}
 \left\|U_{\gamma,v}K_\gamma U_{\gamma,w}^*
             -G-\gamma^{-1/2}G^{[1]}\right\|_{\rm HS}
                         \leq C_{\rho_*,B}/\gamma
 \label{eq:f15v87:compact-jet}
\end{equation}
for sufficiently large \(\gamma\).
The sign of \(G^{[1]}\) corresponds to the displayed order
\(w=v\operatorname{Exp}(\epsilon\xi)\) and the quaternion
product in \(q'\).
\end{theorem}

\begin{proof}
A common rotation removes \(v\). Quaternion multiplication gives
\[
 \operatorname{Exp}(\epsilon x)\cdot
       \bigl[\operatorname{Exp}(\epsilon\xi)
                         \operatorname{Exp}(\epsilon z)\bigr]
 =1-\frac{\epsilon^2}{2}|x-z-\xi|^2
                  +\epsilon^3x\cdot(\xi\times z)
                  +O\bigl(\epsilon^4(1+|x|+|z|)^4\bigr).
\]
The right magnetic term is exactly aligned with its own
endpoint:
\[
 q'\cdot w=\cos(\epsilon|z|).
\]
The left one is \(\cos(\epsilon|x|)\). Consequently the
native exponent is
\[
 A_\gamma=A-\epsilon L+R_\gamma,\qquad
 L=x\cdot(\xi\times z),\qquad
 |R_\gamma|\leq C_B\gamma^{-1}(1+|x|+|z|)^4
\]
throughout the full charts. Uniform Taylor remainders follow
from bounded fourth derivatives of the quaternion exponential
on the relevant compact parameter ball.

The magnetic terms give
\(A_\gamma\geq c_{\rho_*}(|x|^2+|z|^2)\).
The Euclidean exponent has a similar bound up to a constant
depending on \(B\). For nonnegative real \(a,b\),
\[
 |e^{-a}-e^{-b}+(a-b)e^{-b}|
             \leq \tfrac12(a-b)^2(e^{-a}+e^{-b}).
\]
Apply this with \(a=A_\gamma,b=A\), and separate \(R_\gamma\).
It bounds the error after \(e^{-A}(1+\epsilon L)\) by
\(\gamma^{-1}\) times a fixed polynomial Gaussian.
The two Haar square-root densities and \(c_0(\gamma)^{-1}\)
have no term of order \(\epsilon\); their error from the
Euclidean prefactor is \(O(\gamma^{-1})\) with the same
integrable control, as in V86. The off-chart Euclidean tails
are exponentially small. Square integration proves the claim.
\end{proof}

Replacing the product of endpoint exponentials by the
exponential of their sum would erase \(G^{[1]}\).
In the common chart used by V86 the same information is
carried by the change of the right endpoint basis.
There is no contradiction with V86's \(O(\gamma^{-1})\)
common-chart approximation.

\subsection{The complete Gaussian channel matrix}

Set
\[
 \begin{aligned}
 &a=1+\rho/2,\quad d=1+\sigma/2,\quad
 u=\sqrt{a^2-1},\quad v_0=\sqrt{d^2-1},\\
 &A_0=a+u,\quad B_0=d+v_0,\quad D=A_0B_0-1.
 \end{aligned}
\]
The symbol \(v_0\) is a positive oscillator width, not the
unit quaternion \(v\).
Let \(h_\alpha^u\), \(\alpha\in\mathbb N^3\), be the normalized
Hermite functions of width \(u\). Their generating identity is
\begin{equation}
 \sum_\alpha h_\alpha^u(x)\frac{t^\alpha}{\sqrt{\alpha!}}
 =\left(\frac u\pi\right)^{3/4}
      e^{-u|x|^2/2+\sqrt{2u}\,t\cdot x-t\cdot t/2}.
 \label{eq:f15v87:hermite-generating}
\end{equation}
This is the conventional Hermite identity
\href{https://dlmf.nist.gov/18.12.E15}{DLMF 18.12.15}
with its \(L^2\) normalization. Squares of complex generating
variables below are bilinear, not absolute squares.

Define
\[
 \alpha_0=\frac{2uB_0}{D}-1,\quad
 \beta_0=\frac{2v_0A_0}{D}-1,\quad
 \eta=\frac{2\sqrt{uv_0}}{D},\quad
 k_0=\frac{(A_0-1)(B_0-1)}{2D},
\]
\[
 L_0=\frac{\sqrt{2u}(B_0-1)}D,\qquad
 R_0=-\frac{\sqrt{2v_0}(A_0-1)}D,\qquad
 M_0=\begin{pmatrix}\alpha_0&\eta\\\eta&\beta_0\end{pmatrix}.
\]
Here \(\alpha_0,\beta_0\) are scalar coefficients, distinct
from the multi-indices of the Hermite basis.

\begin{theorem}[Populated endpoint matrix and its chiral correction]
\label{f15v87:generating}
Let
\[
 \Gamma_{\alpha\beta}
      =\langle h_\alpha^u,Gh_\beta^{v_0}\rangle,\qquad
 \Gamma^{[1]}_{\alpha\beta}
      =\langle h_\alpha^u,G^{[1]}h_\beta^{v_0}\rangle.
\]
Their complete generating functions are
\begin{align}
 \mathcal F(t,s)
 &=\sum_{\alpha,\beta}
       \Gamma_{\alpha\beta}\frac{t^\alpha s^\beta}
                                  {\sqrt{\alpha!\beta!}} \notag\\
 &=\eta^{3/2}e^{-k_0|\xi|^2}
   \exp\left[
     \frac{\alpha_0}{2}t\cdot t+\frac{\beta_0}{2}s\cdot s
       +\eta t\cdot s+L_0\xi\cdot t+R_0\xi\cdot s\right],
 \label{eq:f15v87:matrix-generating}\\
 \mathcal F^{[1]}(t,s)
 &=\eta\,t\cdot(\xi\times s)\,\mathcal F(t,s).
 \label{eq:f15v87:chiral-generating}
\end{align}
Thus every finite channel matrix is given by derivatives of
a specified exponential quadratic, with no missing matrix
elements.
\end{theorem}

\begin{proof}
Multiply the two generating functions
\eqref{eq:f15v87:hermite-generating} by \(G\) and integrate in
\(x,z\). The quadratic precision in each coordinate pair is
\[
 \begin{pmatrix}A_0&-1\\-1&B_0\end{pmatrix},
 \qquad
 \begin{pmatrix}A_0&-1\\-1&B_0\end{pmatrix}^{-1}
                  =D^{-1}\begin{pmatrix}B_0&1\\1&A_0\end{pmatrix}.
\]
The source vectors are
\(\ell_x=\xi+\sqrt{2u}\,t\) and
\(\ell_z=-\xi+\sqrt{2v_0}\,s\).
The Gaussian determinant gives \(\eta^{3/2}\).
Expanding
\[
 -|\xi|^2/2-(t\cdot t+s\cdot s)/2+
       \frac{B_0\ell_x\cdot\ell_x+
                     2\ell_x\cdot\ell_z+A_0\ell_z\cdot\ell_z}{2D}
\]
gives \eqref{eq:f15v87:matrix-generating}.

For the inserted scalar triple product the Gaussian means
are \(\mu_x=(B_0\ell_x+\ell_z)/D\) and
\(\mu_z=(\ell_x+A_0\ell_z)/D\).
Their cross covariance is \(D^{-1}I\), whose contraction
with the alternating tensor is zero. The insertion is
therefore \(\mu_x\cdot(\xi\times\mu_z)\).
The terms containing two copies of \(\xi\) vanish, and
the remaining terms equal
\(\eta\,t\cdot(\xi\times s)\).
This proves \eqref{eq:f15v87:chiral-generating}.
\end{proof}

In particular the first correction has zero entire row or
column whenever the corresponding oscillator multi-index
is zero. At level one the entries are
\[
 \Gamma_{e_i e_j}
 =\Gamma_{00}\bigl(\eta\delta_{ij}+L_0R_0\xi_i\xi_j\bigr),
 \qquad
 \Gamma^{[1]}_{e_i e_j}
 =\Gamma_{00}\eta\,e_i\cdot(\xi\times e_j),
 \tag{87.1}
\]
where \(\Gamma_{00}=\eta^{3/2}e^{-k_0|\xi|^2}>0\).
For equal widths and \(\xi=0\), one has
\(\alpha_0=\beta_0=0\), \(\eta=Q\), and
\(\mathcal F=Q^{3/2}e^{Q t\cdot s}\), recovering V86's
diagonal Mehler spectrum. Unequal widths and nonzero
displacement must not be dropped from the general formula.

\subsection{Explicit exponential bounds for all omitted levels}

For \(r>1\) such that \(r\|M_0\|<1\), put
\begin{align}
 D_r&=\det(I-r^2M_0^2),\qquad \ell=(L_0,R_0)^{\mathsf T},
 \notag\\
 \mathcal B_r
 &=\eta^3D_r^{-3/2}
     \exp\left[-2k_0|\xi|^2+
        r|\xi|^2\ell^{\mathsf T}(I-rM_0)^{-1}\ell\right],
 \label{eq:f15v87:weighted-norm}\\
 \mathcal B_r^{[1]}
 &=\frac{2\eta^2r^2|\xi|^2}{D_r}\,\mathcal B_r .
 \label{eq:f15v87:weighted-chiral-norm}
\end{align}

\begin{theorem}[Exact weighted coefficient norms and finite-channel tails]
\label{f15v87:tails}
For the displayed range of \(r\),
\begin{align}
 \sum_{\alpha,\beta}r^{|\alpha|+|\beta|}
                           |\Gamma_{\alpha\beta}|^2&=\mathcal B_r,
 \notag\\
 \sum_{\alpha,\beta}r^{|\alpha|+|\beta|}
                     |\Gamma^{[1]}_{\alpha\beta}|^2&=\mathcal B_r^{[1]}.
 \label{eq:f15v87:exact-weighted-sums}
\end{align}
Let \(P_m^u\) project onto all three-dimensional oscillator
levels \(|\alpha|\leq m\), of dimension \(\binom{m+3}{3}\).
Then
\begin{align}
 \|G-P_m^uGP_m^{v_0}\|_{\rm HS}
        &\leq r^{-(m+1)/2}\sqrt{\mathcal B_r}, \notag\\
 \|G^{[1]}-P_m^uG^{[1]}P_m^{v_0}\|_{\rm HS}
        &\leq r^{-(m+1)/2}\sqrt{\mathcal B_r^{[1]}} .
 \label{eq:f15v87:tails}
\end{align}
There is an explicit \(r>1\), independent of the endpoint
parameters in \([\rho_*,4]\), for which both right-hand
constants are bounded uniformly when \(|\xi|\leq B\).
\end{theorem}

\begin{proof}
Monomials \(z^\nu/\sqrt{\nu!}\) are orthonormal in
\(L^2(\mathbb C^6,\pi^{-6}e^{-|z|^2}dz)\) within the space
of entire functions. This follows by polar integration in
each complex coordinate. Hence the first sum in
\eqref{eq:f15v87:exact-weighted-sums} is the squared norm of
\(\mathcal F(\sqrt r\,t,\sqrt r\,s)\) in that Gaussian measure.

The real and imaginary parts of each coordinate pair have
precisions \(I-rM_0\) and \(I+rM_0\).
Integrating them gives the determinant \(D_r^{-3/2}\) and
the exponential in \eqref{eq:f15v87:weighted-norm}.
The condition \(r\|M_0\|<1\) makes both real Gaussian
integrals convergent.

For the second sum insert the squared modulus of
\(\eta r\,t\cdot(\xi\times s)\).
Rotate \(\xi\) to the third axis. Only the first two
coordinate pairs enter this alternating product. They
are independent, centered, and have complex covariance
\[
 R_r=\tfrac12\bigl((I-rM_0)^{-1}+(I+rM_0)^{-1}\bigr)
                         =(I-r^2M_0^2)^{-1}.
\]
Their expectation of \(|t_1s_2-t_2s_1|^2\) is
\(2\det R_r=2/D_r\). The nonzero Gaussian means lie in
the third coordinate and do not enter this product.
This proves \eqref{eq:f15v87:weighted-chiral-norm}.

If \(|\alpha|>m\) or \(|\beta|>m\), then
\(r^{|\alpha|+|\beta|}\geq r^{m+1}\).
Parseval and the two weighted sums prove the tail bounds.

To give a uniform admissible \(r\), put
\[
 u_*=\sqrt{(1+\rho_*/2)^2-1},\quad u^*=\sqrt8,\qquad
 \delta_*=\min\left\{
       \frac{2u_*}{4+\sqrt8},
       \frac{\rho_*}{\sqrt8+\rho_*/2}\right\},\qquad
 r=1+\delta_*/2.
 \tag{87.2}
\]
Indeed with \(S=\operatorname{diag}(\sqrt u,\sqrt{v_0})\) and
\(K=\left(\begin{smallmatrix}A_0&-1\\-1&B_0\end{smallmatrix}\right)\),
\[
 M_0=2SK^{-1}S-I.
\]
The lower bound on \(I+M_0\) is
\(2u_*/(4+\sqrt8)\).
Also \(K=S^2+Q_0\), where
\(Q_0=\left(\begin{smallmatrix}a&-1\\-1&d\end{smallmatrix}\right)
 \geq(\rho_*/2)I\).
Since \(S^2\leq\sqrt8\,I\), inversion gives
\(I-M_0\geq\rho_*/(\sqrt8+\rho_*/2)\).
Thus \(\|M_0\|\leq1-\delta_*\), and
\(r(1-\delta_*)<1\).
The explicit formulas are continuous on the compact
parameter set with denominators bounded away from zero,
so their suprema are finite.
\end{proof}

These are two-sided channel tails, not just decay of a
selected row. The endpoint channel depends on \(\rho,v\)
alone; its entries still depend on both endpoints.

\subsection{Genuine compact projections and gauge covariance}

Restrict \(h_\alpha^u\), \(|\alpha|\leq m\), to the ball
\(\{|x|<\pi\sqrt\gamma\}\). Their Gram matrix is positive
definite, since an analytic Gaussian polynomial cannot
vanish on an open ball unless it is zero. Let
\(P_{\gamma,m}^u\) be the orthogonal projection onto their
span in \(L^2(\mathbb R^3)\), after zero extension.
This is a projection, not a raw unnormalized feature matrix.

For fixed \(m\), uniformly in the declared widths,
\begin{equation}
 \|P_{\gamma,m}^u-P_m^u\|
               \leq C_{m,\rho_*}e^{-c_{\rho_*}\gamma}.
 \label{eq:f15v87:projector-limit}
\end{equation}
To prove it, bound the finitely many Hermite tails by a
polynomial times \(e^{-u_*|x|^2}\).
The restricted Gram matrix differs from \(I\) by the sum
of these tails; for large \(\gamma\) its inverse square
root is bounded and differs from \(I\) by the same order.
Comparing the resulting orthonormal basis maps gives
\eqref{eq:f15v87:projector-limit}, with a reduced positive
exponential constant absorbing the fixed polynomial.

Pull back by \(U_{\gamma,v}\) to obtain the native
finite-rank fibre projection
\[
 \Pi_{\gamma,m}(\rho,v)
     =U_{\gamma,v}^*P_{\gamma,m}^u U_{\gamma,v}.
 \tag{87.3}
\]
Its rank is \(\binom{m+3}{3}\).
The inverse chart functions are only required to be in
\(L^2(S^3)\); no statement about a differential form domain
at the antipode is being made.

Under an endpoint gauge change \(q\mapsto gqh^{-1}\),
\(v\mapsto gvh^{-1}\), the relative coordinate \(v^{-1}q\)
is conjugated by \(h\). Its logarithm rotates by
\(\operatorname{Ad}_h\in SO(3)\).
The ball, Haar density and complete oscillator levels are
rotation invariant. Their span and its orthogonal
projection therefore transform equivariantly.
Selecting only a preferred Cartesian oscillator vector
would not have this property. Both the Gaussian and the
scalar triple product in
\eqref{eq:f15v87:gaussian-and-cubic} obey the same simultaneous
rotation covariance.

Combining Theorem~\ref{f15v87:compact-jet} with
\eqref{eq:f15v87:tails}--\eqref{eq:f15v87:projector-limit}
gives, uniformly on good nearby pairs,
\begin{equation}
 \begin{split}
 &\|K_\gamma-\Pi_{\gamma,m}(\rho,v)
                          K_\gamma\Pi_{\gamma,m}(\sigma,w)\|\\
 &\quad\leq \varepsilon_{m,\gamma}:=
 r^{-(m+1)/2}\bigl(B_r^{1/2}
                     +\gamma^{-1/2}(B_r^{[1]})^{1/2}\bigr)
       +C/\gamma+C_{m,\rho_*}e^{-c_{\rho_*}\gamma},
 \end{split}
 \label{eq:f15v87:compact-channel-error}
\end{equation}
where \(B_r,B_r^{[1]}\) are suprema of the explicit
\(\mathcal B\)'s on that parameter set.
The remainder minus its double compression costs at most
twice its norm, absorbed in \(C\).
In orthonormalized endpoint bases the retained matrix is
\(\Gamma+\gamma^{-1/2}\Gamma^{[1]}+O_{m,\rho_*,B}(\gamma^{-1})\).

For every prescribed fixed tolerance, first choose a
finite \(m\) and then a sufficiently large \(\gamma\).
This gives a uniform finite-fibre approximation on that
region. The constants in
\eqref{eq:f15v87:projector-limit} depend on \(m\); no
undeclared growing-\(m\) estimate or physical-clock error
budget follows from the order of choices.

\subsection{Paying the actual global leading eigenvalue on the near region}

Fix \(0<\rho_*<4\), and let
\[
 \mathcal A=\{V:|S_e(V)|\geq\rho_*\},\qquad
 \mathcal N=\{(V,V')\in\mathcal A^2:
            d(v(V),v(V'))\leq B/\sqrt\gamma\}.
\]
The set \(\mathcal A\) is gauge invariant. The pair condition
is invariant under the simultaneous gauge action, as
required for an equivariant kernel; independent rotations
of its two arguments need not preserve it. Use
\(\Pi(V)=\Pi_{\gamma,m}(\rho(V),v(V))\) on \(\mathcal A\)
and the identity off \(\mathcal A\).
The resulting decomposable projection on the full link
space preserves the combined physical subspace.
It is finite dimensional only in the selected fibre on
\(\mathcal A\), not finite rank on the full lattice.

Let \(E_{\mathcal N}\) be the part of the kernel of
\(\bar T-\Pi\bar T\Pi\) restricted to \(\mathcal N\).
Restore V86's scalar factor
\[
 a(V)c_{\widehat e}(V,V')a(V').
\]
The compression is positive, but
\(\bar T-\Pi\bar T\Pi\) need not be positive; no Loewner
ordering of that difference is assumed.

\begin{theorem}[Volume-independent relative error for the good-near part]
\label{f15v87:relative-near}
For fixed \(\rho_*,B\), put
\[
 h_*=
 \left(\frac{2u_*}{(3+\sqrt8)^2-1}\right)^{3/2}e^{-B^2/2}>0.
 \tag{87.4}
\]
For sufficiently large \(\gamma\), independently of the
spectator volume,
\begin{equation}
 \|E_{\mathcal N}|_{\rm phys}\|
       \leq \frac{2\varepsilon_{m,\gamma}}{h_*}
                                \lambda_0(\bar T|_{\rm phys}).
 \label{eq:f15v87:relative-near}
\end{equation}
All factors refer to the actual native transfer.
No unknown vacuum density or volume-dependent leading
normalizer occurs on the right except the leading
eigenvalue to which the error is compared.
\end{theorem}

\begin{proof}
Let \(\psi_V\) be the normalized native pullback of the
restricted oscillator ground function on \(\mathcal A\).
It is positive almost everywhere and defines an isometry
\[
 (J_0 f)(q,V)=\psi_V(q)f(V),\qquad f\in L^2(\mathcal A).
\]
It is gauge equivariant. Its exact scalar compression
\(S_{00}=J_0^*\bar T J_0\) has nonnegative kernel
\[
 a(V)c_{\widehat e}(V,V')a(V')\,g_{00,\gamma}(V,V'),\qquad
 g_{00,\gamma}=\langle\psi_V,K_\gamma\psi_{V'}\rangle.
\]
The first chiral row and column vanish in
\eqref{eq:f15v87:chiral-generating}; hence on \(\mathcal N\)
\[
 g_{00,\gamma}=\eta^{3/2}e^{-k_0|\xi|^2}
                                  +O_{\rho_*,B}(\gamma^{-1}).
\]
Here \(\eta\geq2u_*/((3+\sqrt8)^2-1)\) and
\(0<k_0<1/2\), so the leading expression is at least \(h_*\).
Eventually \(g_{00,\gamma}\geq h_*/2\).

Define \(S_{\mathcal N}\) to have scalar kernel
\[
 {\bf1}_{\mathcal N}(V,V')
                       a(V)c_{\widehat e}(V,V')a(V').
\]
Pointwise nonnegative-kernel domination gives
\[
 \|S_{\mathcal N}\|\leq\frac2{h_*}\|S_{00}\|
                           \leq\frac2{h_*}\|\bar T\|.
\]
This does not require \(S_{\mathcal N}\) to be positive
semidefinite: apply its kernel to the absolute value of a
function and dominate by \(S_{00}\).
The original full transfer has a strictly positive
compact kernel and commutes with gauge transformations.
Its unique positive Perron vector is gauge invariant, so
\(\|\bar T\|=\lambda_0(\bar T|_{\rm phys})\), as in the
earlier native Perron arguments.

For an arbitrary full vector-valued function \(f(V)\), the
local bound \eqref{eq:f15v87:compact-channel-error} gives
\[
 \|(E_{\mathcal N}f)(V)\|_{L^2_q}
 \leq \varepsilon_{m,\gamma}
            (S_{\mathcal N}\|f(\cdot)\|_{L^2_q})(V).
\]
Take the outer \(L^2\) norm and then restrict to the
physical space. The two preceding estimates prove
\eqref{eq:f15v87:relative-near}.
\end{proof}

The new ingredient is the computed ground overlap
\(\Gamma_{00}\), with a positive lower bound independent
of volume. Kernel domination and the Perron principle
are inherited results. The estimate uses the actual
physical top eigenvalue, even when spectator volume
makes that eigenvalue small.

\subsection{What the remaining operator problem now is}

The good-near truncation error is controlled. The retained
matrix is explicit, equivariant, and includes its first
noncommuting correction. This is more than a pointwise
conditioning statement or a list of local spectral ratios.

It is not yet a bound for the complete remainder. The
complement of \(\mathcal N\) includes mixed small-staple
pairs and pairs whose endpoint directions are farther
apart. V86 supplies an absolute far-direction estimate,
but it has not been converted into a sufficiently small
error relative to the same global leading eigenvalue.
Nor has the physical leading ratio of the retained
matrix-valued other-link operator been bounded.
The projection is the identity on the bad region, so
the bad--bad block has not been discarded or solved.

The next step is to examine the joint native kinetic
cost of staple-direction changes and the scalar ground
compression used above. This may control more than an
arbitrary pairwise norm comparison; it must be derived
from the actual other-link convolution and generated
weights. The two-endpoint oscillator matrix cannot be
replaced by the product of separately frozen local gaps.
The prescribed physical clock, small-staple contributions,
and interacting continuum construction remain outstanding.

\subsection{Diagnostics and append record}

The new script
\begin{center}\small
\path{scripts/verify_ym_f15_v87_endpoint_oscillator_channels.py}
\end{center}
reruns V86's inherited exact chain. Its new rational
fixtures compare independent Hermite moments with the
coherent generating coefficients; verify quaternion
endpoint jets, the alternating correction, determinant
and Hilbert--Schmidt identities, rotation covariance,
complete-level counting and exponential-tail weights;
and check the scalar-compression relative-error argument
on positive finite operator kernels. These checks do not
supply a numerical physical spectral margin.

All 6,038 new exact-rational fixtures passed, together with
the inherited diagnostic chain. The checker has 9,297 bytes
and SHA--256
\begin{center}\scriptsize\ttfamily
0A58B2379EAA3A83280C89C8A29BC0DB9C9BA9AFCAA09B568F13662BD0479078.
\end{center}

The predecessor V86 has 1,944,734 bytes and SHA--256
\begin{center}\scriptsize\ttfamily
F3BD482E3234B23F957FA2A4B07AF49E153D44D8F004004523A67EB8560DEE25.
\end{center}
Its body is preserved byte for byte. Removing this final
section and restoring the title version recovers V86.
V87 replaces only the active main source; archives,
companions and all 146 previously protected files remain
unchanged. Only \path{.tex} files belong in
\path{latex_notes}; compilation products belong in
\path{artifacts/ym_f15_v87_texcheck/}.

\begin{firewall}
V87 proves a controlled finite-fibre native approximation
and a relative error estimate on the good-near region.
It does not silently discard the complementary kernel,
prove contraction of the retained interacting operator,
or construct a nontrivial continuum theory.
The full Yang--Mills mass gap remains unproved and the
original goal remains active.
\end{firewall}

\section{V88: kinetic-source domination, hybrid coverage, and an undecidability audit}
\label{f15v88:context}

This iteration closes the complete approximation-error problem
left by V87, for a hybrid projection at one selected spatial
link. It does not close the spectral-gap problem for the
retained operator. The new projection keeps every small-staple
fibre, and its error is bounded relative to the actual physical
top eigenvalue, uniformly in spectator volume. No
small-probability set of physical states is removed without
an operator estimate.

The new input is an exact monotonicity argument for the
\emph{native kinetic convolution}. It converts a large change
in the twelve links of a spatial staple into a relative
whole-operator bound. The compact endpoint expansions,
oscillator channels, Bessel-ratio monotonicity and Perron
normalization are inherited from V83 and V86--V87; their
proofs are not repeated as new results.

During this iteration the user supplied
\path{spectral_gap.pdf} and a proposed undecidability audit.
They are source material, not instructions to change the
mathematical target. The audit below is integrated into the
existing spectral route, not a separate claim that
Yang--Mills is decidable or undecidable. The last companion
checkpoint remains V85; the next is V90.

\subsection{What the undecidability paper does and does not transfer}

The supplied paper is T.~Cubitt, D.~Perez-Garcia and
M.~M.~Wolf, \emph{Undecidability of the Spectral Gap},
published version, 2022,
\href{https://arxiv.org/abs/1502.04573v5}{arXiv:1502.04573v5},
\href{https://doi.org/10.1017/fmp.2021.15}
{Forum of Mathematics, Pi 10, e14}.
Definitions 1--2, Theorem 3, the physical-scope discussion,
the proof overview and the final assembly in Section 6
are the relevant parts for this transfer audit. The detailed
quantum-machine implementation is not being re-proved here.

Theorem 3 constructs a computable family of
translation-invariant, nearest-neighbour two-dimensional
spin Hamiltonians, with one fixed finite local dimension.
Halting selects a uniformly gapped thermodynamic system;
nonhalting selects a system with densely filling low-energy
spectrum. The promise is stronger than merely failure of a
uniform lower gap. Arbitrarily small local quantum
perturbations of suitable classical interactions are already
enough for the construction. The hierarchy of tilings and
computational history states transmits information to scales
not controlled by a fixed finite-volume inspection.

This is not a theorem about the particular Wilson action,
the combined Gauss-invariant Hilbert space, or the continuum
Yang--Mills construction in these notes. In particular, the
paper explicitly distinguishes its theorem from proving
axiomatic independence of the Yang--Mills mass-gap conjecture
(pp.~9--10). Conversely, gauge symmetry or positivity by
themselves are not an exclusion theorem for such mechanisms.
No reduction identifying the paper's constructed family
with this native YM family has been supplied in either
direction.

The legitimate methodological import is to require
complete spectral coverage and uniform constants, while
testing every proposed simplification against hidden
retained modes. The suggested implication from summable RG
errors to absence of unbounded memory is not automatic:
summability alone says nothing about the dimension or
contraction of retained variables. Likewise, a supported
defect-energy floor is not yet coercivity for an arbitrary
superposition of dilute, separated or hierarchical defects.
Those claims require their own theorems. A full operator
norm bound remains a valid spectral perturbation tool;
a small norm of each local interaction is not the same
hypothesis when the number of interactions grows.

\subsection{A sharp kinetic-source form bound}

Fix a finite spatial graph and the native normalized
unprojected transfer
\begin{equation}
 \bar T=M_b C_\gamma^{\otimes E}M_b,\qquad
 C_s(U,V)=\frac{e^{-s(1-U\cdot V)}}{c_0(s)},\qquad
 c_0(s)=\int_{\SU(2)}e^{-s(1-q_0)}\,dH(q).
 \label{eq:f15v88:native}
\end{equation}
Haar measure has mass one. The positive continuous
magnetic multiplier \(b\) is the actual one at the chosen
regulator; it is never changed when a kinetic source is
introduced below. At finite volume \(b^{-1}\) is bounded.
All tensor products and congruences in this subsection
are on the full Haar Hilbert space before restriction
to the combined physical subspace.

By Proposition~\ref{f15v83:scalar}, the eigenvalues
\(r_n(s)=I_{n+1}(s)/I_1(s)\) of \(C_s\), in its common
complete Peter--Weyl basis, are nondecreasing in \(s\).
Consequently
\begin{equation}
 0\leq C_{\alpha s}\leq C_s,\qquad 0<\alpha\leq1,
 \label{eq:f15v88:kinetic-order}
\end{equation}
in positive-operator order. This is not pointwise order
between their kernels.

\begin{proposition}[Native kinetic source and its exact normalization]
\label{f15v88:source}
Let \(J\subset E\) be finite and nonempty, and
\(0\leq\theta_e<1\) for \(e\in J\). Multiply the kernel
of \(\bar T\) by
\[
 e^{X_\theta(U,V)},\qquad
 X_\theta=\gamma\sum_{e\in J}\theta_e(1-U_e\cdot V_e),
\]
and denote the resulting operator by \(\bar T_\theta\).
Then
\begin{equation}
 0\leq\bar T_\theta\leq R_\theta\bar T,\qquad
 R_\theta=
 \prod_{e\in J}\frac{c_0((1-\theta_e)\gamma)}{c_0(\gamma)}
 \leq\prod_{e\in J}(1-\theta_e)^{-3/2}.
 \label{eq:f15v88:source-order}
\end{equation}
The first upper constant \(R_\theta\) is sharp as a
quadratic-form comparison, including on the physical
subspace. The estimate has no spectator-volume factor.
\end{proposition}

\begin{proof}
The exact kernel identity is
\[
 \bar T_\theta=R_\theta M_b
 \left(\bigotimes_{e\in J}C_{(1-\theta_e)\gamma}
       \otimes\bigotimes_{e\notin J}C_\gamma\right)M_b.
\]
Tensoring positive-operator inequalities in
\eqref{eq:f15v88:kinetic-order}, and then taking the
congruence by \(M_b\), proves the first assertion. All
operators commute with the simultaneous gauge action,
so the inequality restricts to the physical space.

The one-dimensional integral already used in V86 gives
\[
 s^{3/2}c_0(s)=\frac2\pi\int_0^{2s}
       e^{-u}u^{1/2}(2-u/s)^{1/2}\,du.
\]
The right side is nondecreasing in \(s\): both its
integration domain and its nonnegative integrand increase.
Thus \(c_0(\alpha s)/c_0(s)\leq\alpha^{-3/2}\).
This is an exact all-\(s>0\) inequality, not an
asymptotic estimate. For fixed \(\alpha\) the bound is
asymptotically saturated as \(s\to\infty\).

Finally take \(f=b^{-1}\). It is physical, since \(b\)
is gauge invariant, and \(M_bf=1\). Each normalized
convolution leaves the constant function fixed, whence
\[
 \langle f,\bar T_\theta f\rangle
       =R_\theta\langle f,\bar T f\rangle.
\]
This proves sharpness of the form constant; no equality
of operator norms is asserted.
\end{proof}

The full \(\bar T\) has a strictly positive compact kernel.
Its unique positive Perron vector is gauge invariant,
so its full norm equals
\[
 \Lambda:=\|\bar T\|
         =\lambda_0(\bar T|_{\rm phys}).
\]
This inherited normalization is essential. A lower bound
for \(\Lambda\) which deteriorates with volume is not
inserted anywhere below.

\subsection{Relative operator tails for temporal link increments}

Put
\[
 D_J(U,V)=\sum_{e\in J}(1-U_e\cdot V_e),\qquad
 \nu_J=\frac{3|J|}{2},
\]
and define the scalar function
\begin{equation}
 \Psi_\nu(L)=
 \begin{cases}
 1,&0\leq L\leq\nu,\\
 e^{\nu-L}(L/\nu)^\nu,&L>\nu.
 \end{cases}
 \label{eq:f15v88:psi}
\end{equation}

\begin{theorem}[Native relative Chernoff bound]
\label{f15v88:tail}
Let \(\mathcal A\) be a measurable pair event contained
in \(\{\gamma D_J\geq L\}\), and let \(\bar T_{\mathcal A}\)
be the integral operator with kernel
\({\bf1}_{\mathcal A}\bar T\). Then
\begin{equation}
 \|\bar T_{\mathcal A}\|_{\rm full}
       \leq\Psi_{\nu_J}(L)\Lambda
       \leq \min\{1,\,2^{\nu_J}e^{-L/2}\}\Lambda.
 \label{eq:f15v88:relative-tail}
\end{equation}
The first inequality also holds for any measurable
mask taking values in \([0,1]\) supported on that event.
No positive-semidefiniteness of the masked operator is
required. For a simultaneous-gauge-invariant mask the
physical restriction is well defined and obeys the
same estimate.
\end{theorem}

\begin{proof}
For \(0\leq\theta<1\), pointwise nonnegative-kernel
domination gives
\[
 |\,\bar T_{\mathcal A}f\,|
       \leq e^{-\theta L}\bar T_\theta|f|,
 \qquad \theta_e=\theta\quad(e\in J).
\]
Thus Proposition~\ref{f15v88:source} yields
\[
 \|\bar T_{\mathcal A}\|
       \leq e^{-\theta L}(1-\theta)^{-\nu_J}\Lambda.
\]
The logarithm of this last scalar bound is convex;
its derivative is \(-L+\nu_J/(1-\theta)\).
Its minimum occurs at \(\theta=0\) if \(L\leq\nu_J\)
and at \(\theta=1-\nu_J/L\) otherwise. This proves
the first estimate. Comparing the minimum with
\(\theta=0\) and \(\theta=1/2\) proves the second.
\end{proof}

There is also a vacuum-pair consequence, but it is weaker
than the operator statement needed here. For the positive
unit Perron vector \(\Omega\), define
\[
 d\pi_\Omega(U,V)=\Lambda^{-1}
   \Omega(U)\bar T(U,V)\Omega(V)\,dH(U)\,dH(V).
\]
This is a probability measure in the specified unprojected
Haar representation, invariant under simultaneous gauge
transformations. It is not a claim about independently
gauge-fixed representatives at its two endpoints.
The source form bound implies
\[
 \mathbb E_{\pi_\Omega}e^{\theta\gamma D_J}
           \leq(1-\theta)^{-\nu_J},
 \qquad
 \mathbb E_{\pi_\Omega}[\gamma D_J]\leq\nu_J.
\]
For the second inequality use equality at \(\theta=0\)
and right derivatives. No coefficientwise bound on all
higher moments or full stochastic domination follows
just from this moment-generating-function inequality.
In particular, the proof of
\eqref{eq:f15v88:relative-tail} has not replaced an
operator norm by a vacuum probability.

\subsection{The twelve native links control a staple change}

Work now on a three-dimensional cubic spatial torus of
side at least four. Fix an oriented link \(e\), and write
\(U=(q,V)\), with \(q=U_e\). The four plaquettes containing
\(e\) determine the native staple sum \(S_e(V)\).
Each staple word has length three. Their other-link
supports are disjoint, so their union \(J_e\) contains
twelve distinct links and does not contain \(e\).
Write
\[
 \rho=|S_e(V)|,\quad \sigma=|S_e(W)|,\quad
 v=\overline{S_e(V)}/\rho,\quad
 w=\overline{S_e(W)}/\sigma,\quad
 \vartheta=d_{\mathbb S^3}(v,w)
\]
when \(\rho\sigma>0\). Here \(0\leq\rho,\sigma\leq4\).

\begin{proposition}[Literal staple-increment bound]
\label{f15v88:staple}
For every pair of other-link configurations,
\begin{equation}
 |S_e(V)-S_e(W)|^2
       \leq24D_{J_e}(V,W).
 \label{eq:f15v88:staple-lipschitz}
\end{equation}
If \(\rho,\sigma>0\), the left side is exactly
\begin{equation}
 (\rho-\sigma)^2+2\rho\sigma(1-\cos\vartheta).
 \label{eq:f15v88:radial-angular}
\end{equation}
Consequently, for fixed \(0<l<h<4\) and \(B>0\),
the angular and deep radial events
\[
 \begin{split}
 \mathcal F&=\{\rho,\sigma\geq l,\
                         \vartheta>B/\sqrt\gamma\},\\
 \mathcal D&=\{\max(\rho,\sigma)\geq h,\
                              \min(\rho,\sigma)<l\}
 \end{split}
\]
satisfy
\begin{align}
 \|\bar T_{\mathcal F}\|/\Lambda
   &\leq\Psi_{18}\left(\frac{l^2B^2}{6\pi^2}\right),
       \label{eq:f15v88:angular-tail}\\
 \|\bar T_{\mathcal D}\|/\Lambda
   &\leq\Psi_{18}\left(\frac{\gamma(h-l)^2}{24}\right).
       \label{eq:f15v88:deep-tail}
\end{align}
The estimates remain valid for smaller pair events.
\end{proposition}

\begin{proof}
The chord distance on unit quaternions is bi-invariant,
and inversion is an isometry. Telescoping a product of
three factors therefore bounds the difference of each
staple word by the sum of its three link differences.
The triangle inequality and Cauchy--Schwarz yield
\[
 |S_e(V)-S_e(W)|^2
   \leq\left(\sum_{j\in J_e}|V_j-W_j|\right)^2
   \leq12\sum_{j\in J_e}|V_j-W_j|^2
   =24D_{J_e}.
\]
The twelve-link count follows directly by listing the
positive and negative transverse plaquettes in the two
directions perpendicular to \(e\); side at least four
prevents their distinct links from being identified by
periodicity. Formula \eqref{eq:f15v88:radial-angular}
is the Euclidean polarization identity in \(\mathbb R^4\),
unchanged by quaternion conjugation.

On \(\mathcal F\), the inequality
\(1-\cos\vartheta\geq2\vartheta^2/\pi^2\) for
\(0\leq\vartheta\leq\pi\) gives
\[
 \gamma D_{J_e}\geq
       \frac{\gamma l^2\vartheta^2}{6\pi^2}
       >\frac{l^2B^2}{6\pi^2}.
\]
On \(\mathcal D\), the radial term instead gives
\(\gamma D_{J_e}>\gamma(h-l)^2/24\).
Apply Theorem~\ref{f15v88:tail}, with \(|J_e|=12\).
\end{proof}

These are estimates for the entire native transfer with
all its magnetic weights and other links still present,
not the absolute one-link bound from V86. The pair
conditions are invariant under simultaneous gauge
transformations. Independent rotations of the two
endpoint configurations need not preserve the angular
condition; no such invariance has been used.

\subsection{A complete relative error bound for hybrid compression}

Choose the V87 compact endpoint oscillator projections
\(\Pi_{\gamma,m}(\rho,v)\), using complete levels up to
\(m\). Their rank in the selected fibre is
\(\binom{m+3}{3}\). With \(0<l<h<4\) as above define
the decomposable orthogonal projection
\begin{equation}
 (Pf)(q,V)=P(V)f(q,V),\qquad
 P(V)=
 \begin{cases}
 \Pi_{\gamma,m}(\rho(V),v(V)),&\rho(V)\geq h,\\
 I,&\rho(V)<h.
 \end{cases}
 \label{eq:f15v88:hybrid}
\end{equation}
It is gauge equivariant and preserves the combined physical
space. The low-staple region is kept exactly, including
its infinite-dimensional fibre. The full operator
\(\bar T^{\rm hyb}=P\bar T P\) is positive and compact,
but the difference \(\bar T-\bar T^{\rm hyb}\) need
not be positive.

Use the V87 error
\(\varepsilon_{m,\gamma}(l,B)\) from
\eqref{eq:f15v87:compact-channel-error}, uniformly for
\(\rho,\sigma\in[l,4]\) and \(|\xi|\leq B\).
Let
\[
 u_l=\sqrt{(1+l/2)^2-1},\qquad
 h_*(l,B)=
   \left(\frac{2u_l}{(3+\sqrt8)^2-1}\right)^{3/2}
             e^{-B^2/2}.
\]

\begin{theorem}[Whole-operator hybrid approximation]
\label{f15v88:hybrid}
For fixed \(l,h,B,m\) and sufficiently large \(\gamma\),
\begin{equation}
 \frac{\|\bar T-P\bar T P\|_{\rm full}}{\Lambda}
       \leq\delta_{m,\gamma}(l,h,B),
 \label{eq:f15v88:hybrid-error}
\end{equation}
where
\begin{equation}
 \begin{split}
 \delta_{m,\gamma}(l,h,B)
    ={}&\frac{2\varepsilon_{m,\gamma}(l,B)}{h_*(l,B)}
       +2\Psi_{18}\left(\frac{l^2B^2}{6\pi^2}\right)\\
       &+2\Psi_{18}\left(\frac{\gamma(h-l)^2}{24}\right).
 \end{split}
 \label{eq:f15v88:delta}
\end{equation}
All constants are independent of spectator volume.
The same estimate holds after physical restriction.
For any fixed tolerance \(\tau>0\), fixed \(0<l<h<4\),
there exist finite \(B,m\) and \(\gamma_0\), independent
of volume, such that the right side is below \(\tau\)
for every \(\gamma\geq\gamma_0\).
\end{theorem}

\begin{proof}
If both \(\rho,\sigma<h\), the error kernel is zero
because both projections are the identity. On the
remaining pairs, where \(\max(\rho,\sigma)\geq h\),
partition the error into the following three disjoint parts:
\[
 \begin{array}{ll}
 \mathcal N:&\min(\rho,\sigma)\geq l,\
                        \vartheta\leq B/\sqrt\gamma,\\
 \mathcal F':&\min(\rho,\sigma)\geq l,\
                        \vartheta>B/\sqrt\gamma,\\
 \mathcal D:&\min(\rho,\sigma)<l.
 \end{array}
\]
The angular part \(\mathcal F'\) is contained in
\(\mathcal F\) of Proposition~\ref{f15v88:staple}.

First consider \(\mathcal N\). If both endpoints have
\(\rho,\sigma\geq h\), the local error bound is exactly
the V87 two-sided bound. If only one endpoint is above
\(h\), the other projection is the identity; the same
bound still holds. For example, with \(Q_\sigma\) the
hypothetical V87 finite projection at that other endpoint,
\[
 (I-\Pi_\rho)K_\gamma
   =(I-\Pi_\rho)
           (K_\gamma-\Pi_\rho K_\gamma Q_\sigma).
\]
The right-endpoint case uses the analogous multiplication
on the right. Since both field magnitudes are at least
\(l\), the two-sided bound is available even though
\(Q_\sigma\) is not used in the actual hybrid operator.

For completeness, the relative-normalization step is
precisely the one from
Theorem~\ref{f15v87:relative-near}, now on
\(\{\rho\geq l\}\). The positive fibre ground isometry
\(J_0\) gives a scalar compression
\(S_{00}=J_0^*\bar T J_0\). Its kernel contains the
ground overlap \(g_{00,\gamma}\geq h_*(l,B)/2\)
on \(\mathcal N\), for sufficiently large \(\gamma\).
If \(S_{\mathcal N}\) denotes the native scalar
other-link kernel restricted to this event, then
\[
 \|S_{\mathcal N}\|
      \leq2h_*^{-1}\|S_{00}\|
      \leq2h_*^{-1}\Lambda.
\]
Pointwise fibre-norm domination by
\(\varepsilon_{m,\gamma}S_{\mathcal N}\) proves
\(\|E_{\mathcal N}\|\leq2\varepsilon_{m,\gamma}h_*^{-1}\Lambda\)
on the full Hilbert space.

For either other-link pair event \(\mathcal A=\mathcal F'\)
or \(\mathcal D\), decomposability of \(P\) gives the exact
error identity
\[
 E_{\mathcal A}=\bar T_{\mathcal A}
                              -P\bar T_{\mathcal A}P,
 \qquad
 \|E_{\mathcal A}\|\leq2\|\bar T_{\mathcal A}\|.
\]
The event masks do not depend on the selected fibre
variables \(q,q'\), which is needed for that identity.
Proposition~\ref{f15v88:staple} bounds the two norms
relative to \(\Lambda\). Adding the three estimates
proves \eqref{eq:f15v88:delta}.

Finally choose \(B\) large enough to make the angular
term smaller than \(\tau/3\). With that \(B\) fixed,
V87's exponential level tail allows a finite \(m\)
making the limiting near term smaller than \(\tau/3\).
Then increase \(\gamma_0\) to control the compact
remainders, the ground-overlap threshold and the deep
radial term. A slight strict margin in the first two
choices gives the stated bound for all
\(\gamma\geq\gamma_0\). None of these choices depends
on spectator volume.
\end{proof}

This finishes the complementary-error task explicitly
left open in V87. The buffer \(l<h\) is substantive:
a jump between high and genuinely small staple magnitude
has a fixed kinetic cost, while the intermediate region
still admits the V87 local comparison. A hard truncation
on every nonzero staple, without such a buffer or a
replacement estimate, is not covered.

The theorem is a one-selected-link hybrid result.
Compressions at different links depend on overlapping
other-link configurations and need not commute.
No multiplication of these bounds into an all-link
finite-dimensional operator, or volume-free bound for
such a multiplication, is asserted.

\subsection{Complete near-top coverage, but not a gap certificate}

The following is an elementary consequence of the full
norm estimate; its value here is that it applies to all
physical states, not only to an enumerated set of
configuration patterns.

\begin{proposition}[No missing near-top state outside the hybrid range]
\label{f15v88:coverage}
Suppose \(\|\bar T-P\bar T P\|\leq\delta\Lambda\).
For \(0<\alpha\leq1\), let
\[
 E_\alpha={\bf1}_{[\alpha\Lambda,\Lambda]}
                         (\bar T|_{\rm phys}).
\]
Then
\begin{equation}
 \|(I-P)E_\alpha\|\leq\delta/\alpha.
 \label{eq:f15v88:coverage}
\end{equation}
Thus any unit physical eigenvector with eigenvalue
\(\lambda\geq\alpha\Lambda\) has at most
\((\delta/\alpha)^2\) of its squared norm outside
the retained range, irrespective of its localization,
defect pattern or spectator-volume dependence.
\end{proposition}

\begin{proof}
Let \(R=\bar T-P\bar T P\), restricted to the physical
space. Since \((I-P)P=0\), one has
\((I-P)\bar T=(I-P)R\). The spectral inverse of
\(\bar T\) on the range of \(E_\alpha\) is bounded
by \(1/(\alpha\Lambda)\). Therefore
\[
 (I-P)E_\alpha=(I-P)R
               (\bar T|_{\operatorname{Ran}E_\alpha})^{-1}
                    E_\alpha,
\]
which proves the bound. No assumed overlap with a
reference product vector is used.
\end{proof}

In Hamiltonian language, transfer eigenvalues close to
the top describe low excitation energies. Equation
\eqref{eq:f15v88:coverage} provides approximate
completeness for that whole window. It does not imply
that the retained component is far from the retained
vacuum, or that it has a positive excitation energy.

For clarity, specialize the earlier V79--V80 min--max
argument to this hybrid compression. Put
\[
 r_{\rm full}=\lambda_1(\bar T|_{\rm phys})/\Lambda,
 \qquad
 r_{\rm hyb}=\lambda_1(\bar T^{\rm hyb}|_{\rm phys})
                  /\lambda_0(\bar T^{\rm hyb}|_{\rm phys}).
\]
For \(\delta<1\), positivity, compression interlacing
and the norm error imply
\begin{equation}
 (1-\delta)r_{\rm hyb}\leq r_{\rm full}
               \leq\min\{1,r_{\rm hyb}+\delta\}.
 \label{eq:f15v88:ratio}
\end{equation}
Indeed the compressed top lies in
\([(1-\delta)\Lambda,\Lambda]\); each compressed
eigenvalue is at most the corresponding full eigenvalue,
and their differences are at most \(\delta\Lambda\).
These statements give both inequalities directly.
This is a handoff to the retained spectral problem,
not a new abstract min--max theorem.

\begin{proposition}[Adversarial positive-kernel retained slow mode]
\label{f15v88:countermodel}
An arbitrarily accurate full relative compression,
a uniformly contracting discarded fibre, and complete
near-top coverage can coexist with a closing spectral gap.
\end{proposition}

\begin{proof}
For \(0\leq r<1\), define the strictly positive
stochastic matrix
\[
 A_r=\frac12
   \begin{pmatrix}1+r&1-r\\1-r&1+r\end{pmatrix}.
\]
Its eigenvalues are \(1,r\). For \(n\geq2\), take
\[
 T_n=A_{n^{-2}}\otimes A_{1-n^{-1}},\qquad
 P_n=A_0\otimes I.
\]
Here \(A_0\) is the orthogonal projection onto constants.
The respective eigenvalue lists of \(T_n\) and
\(P_nT_nP_n\) are
\[
 (1,\ 1-n^{-1},\ n^{-2},\ n^{-2}(1-n^{-1})),
 \qquad (1,\ 1-n^{-1},\ 0,\ 0).
\]
Consequently
\[
 \|T_n-P_nT_nP_n\|=n^{-2},\qquad
 r_{\rm full}=r_{\rm hyb}=1-n^{-1}.
\]
Every matrix \(T_n\) is positive semidefinite and has
strictly positive entries, with a unique positive
Perron vector. The fast factor has gap \(1-n^{-2}\),
and the approximation error vanishes; nevertheless the
retained slow factor has gap \(n^{-1}\).
All conclusions follow by the constant/alternating
eigenbasis in each tensor factor.
\end{proof}

This is not a Wilson-action counterexample or a
construction of an undecidable model. It pinpoints
the logical distinction that the audit requires:
controlled discarded error cannot replace a bound
on the entire retained interacting operator. In the
actual hybrid operator, all low-staple fibres and all
other-link collective motion are still retained.

\subsection{Uniformity ledger and the next upstream target}

The constants proved here are uniform in spectator
volume but not uniformly small under every refinement
choice. In particular, \(h_*(l,B)^{-1}\) grows rapidly
with \(B\), and V87's compact projection constants
depend on \(m\). The order of choosing a fixed tolerance,
then \(B\), then \(m\), then sufficiently large \(\gamma\),
is not a quantitative theorem for arbitrary
\(m(\gamma),B(\gamma)\), or a prescribed vanishing
physical time step.

For the actual refinement sequence and independently
specified clock \(a_{t,j}\), a sufficient spectral
handoff would be
\begin{equation}
 r_{{\rm hyb},j}+\delta_j
             \leq e^{-m_*a_{t,j}},\qquad m_*>0,
 \label{eq:f15v88:clock-target}
\end{equation}
uniformly in the required volumes. An error
\(\delta_j=o(a_{t,j})\), together with a suitably
strong retained gap, is one possible way to meet this
budget. Neither component is supplied by the fixed
tolerance statement alone. Finite-volume positivity
and successful finite diagnostics do not replace the
uniform estimate.

The direction now moves upstream to the physical
spectral ratio of \(P\bar T P\), especially its unchanged
low-staple sector and collective other-link modes.
A sparse or hierarchical state cannot be lost outside
the hybrid range without paying
\eqref{eq:f15v88:coverage}; it can still be a slow
\emph{retained} state. Charged-sector estimates do not
automatically control every neutral retained state.
The next iteration must exploit additional native
structure there, not re-prove the kinetic tail or
rename the approximation theorem a contraction theorem.

The continuum construction, nontriviality, physical
clock calibration and a uniform positive physical mass
remain outstanding. No claim of decidability, exclusion
of all computational encodings, or independence from
any formal axiom system is made.

\subsection{Diagnostics and append record}

The script
\begin{center}\small
\path{scripts/verify_ym_f15_v88_kinetic_source_hybrid.py}
\end{center}
reruns the V87 inherited chain. Its new exact-rational
checks cover outward Bessel and normalization bounds,
tensor order under magnetic congruence and sharpness,
literal twelve-link incidence, noncommuting quaternion
staple differences, radial/angular decomposition,
Chernoff minimization, exhaustive pair partitions,
one-sided projection identities and the positive-kernel
closing-gap countermodels. These are diagnostics of
the proved formulas, not numerical certification of a
physical gap.

All 2,557 new exact-rational fixtures passed, along
with the inherited diagnostic chain. The checker has
8,355 bytes and SHA--256
\begin{center}\scriptsize\ttfamily
7186CFB52A2435BFF5BF520A4520DD6C0284F27DFE092C832E830DC3247C0FCF.
\end{center}

The predecessor V87 has 1,967,023 bytes and SHA--256
\begin{center}\scriptsize\ttfamily
70475C4CD14048B8C053C2E0D2DEF149AC738C1B76A223F36DF122F6A2DE9B07.
\end{center}
Its body is preserved byte for byte. Removing this final
section and restoring the title version recovers V87.
V88 replaces only the active main source; the prior
archives, companion notes and protected diagnostic
files are unchanged. The two newly supplied source
files are read-only inputs. Only \path{.tex} files belong
in \path{latex_notes}; compilation products belong in
\path{artifacts/ym_f15_v88_texcheck/}.

\begin{firewall}
V88 proves a volume-independent relative approximation
for the complete one-link hybrid transfer, and complete
near-top coverage up to its stated error. It does not
prove a gap for the retained hybrid operator, discard
small-staple states, or construct the interacting
continuum Yang--Mills theory. The undecidability audit
strengthens these distinctions; it does not settle the
Yang--Mills conjecture in either direction. The original
mass-gap goal remains active.
\end{firewall}

\section{V89: temporal persistence controls the small-staple operator}
\label{f15v89:context}

V88 left every low-staple fibre unchanged because a
one-time small probability is not an operator-norm bound.
This iteration goes upstream to the nonzero-source
vacuum estimate in V62. Its scope includes every finite
temporal support, with the error paid per source cell.
That stronger input controls persistence for arbitrarily
many consecutive slices. At each fixed spatial regulator,
the long-time persistence rate detects the entire positive
transfer restricted to a specified bad region.

This closes the small-staple block problem for the
selected native Wilson sequence. It gives an exponential
relative bound, a full-vector low-energy exclusion estimate,
and a justified replacement of V88's hybrid projection by
a projection of bounded rank in the selected fibre
\emph{at every other-link configuration}. The other links
remain infinite dimensional and interacting. No gap for
their retained collective motion is asserted.

Theorems~\ref{f15v62:source} and \ref{f15v62:window},
V87's endpoint channels and V88's kinetic angular tail
are reused at their proved scopes.
The new bridge is from temporal persistence to a
restricted physical spectral norm. In particular, no
sampling Dirichlet form is identified with the physical
transfer energy. The compulsory companion checkpoint
is V90, not this iteration.

\subsection{The fixed-regulator persistence bridge}

Let \(\bar T\) be the normalized native transfer of
\eqref{eq:f15v88:native}, and put
\[
 \Lambda=\|\bar T\|
       =\lambda_0(\bar T|_{\rm phys}),\qquad
 K=\bar T/\Lambda.
\]
Let \(\Omega>0\) be its unit Perron vector.
The full kernel is strictly positive and compact,
\(K\geq0\), and \(K\Omega=\Omega\).
These are inherited finite-regulator facts.

For a measurable gauge-invariant slice event \(A\), let
\(Q_A\) denote multiplication by its indicator, and put
\[
 K_A=Q_AKQ_A,\qquad u_A=Q_A\Omega.
\]
Regard \(K_A\) as an operator on \(L^2(A,dH)\), or
extend it by zero. For \(n\geq1\), define
\[
 p_A(n)=\mu_{\rm vac}\{U_0,\ldots,U_{n-1}\in A\}.
\]
These are probabilities in the actual stationary Wilson
vacuum, not an independently sampled sequence of slices.

\begin{proposition}[A temporal tube controls a whole restricted operator]
\label{f15v89:persistence}
At a fixed spatial regulator and coupling,
\begin{equation}
 p_A(n)=\langle u_A,K_A^{\,n-1}u_A\rangle,\qquad
 \lim_{n\to\infty}p_A(n)^{1/n}=\|K_A\|.
 \label{eq:f15v89:persistence}
\end{equation}
For a Haar-null event both sides of the limit are zero.
Otherwise the full restricted norm equals its physical
restricted norm. If, for every \(n\),
\[
 p_A(n)\leq C e^{-an}
\]
with finite \(C\) independent of \(n\), then
\begin{equation}
 \|K_A\|\leq\min\{1,e^{-a}\}.
 \label{eq:f15v89:tube-norm}
\end{equation}
No lower bound on a Perron overlap uniform in spatial
volume is needed.
\end{proposition}

\begin{proof}
The literal vacuum insertion formula, or the ground-state
Markov transform in \eqref{eq:f15v78:ground}, gives
\[
 p_A(n)=\langle\Omega,Q_AKQ_AK\cdots KQ_A\Omega\rangle.
\]
There are \(n\) indicators and \(n-1\) transfers.
Gauge-invariant insertions preserve the physical
subspace, so using the full unprojected transfer in this
formula agrees with the physical Wilson vacuum.

If \(A\) has positive Haar measure, the restricted
strictly positive compact kernel is positivity improving.
Its top eigenvalue \(r_A>0\) has a unit eigenvector
\(\phi_A>0\). Gauge invariance of the kernel and uniqueness
of that positive eigenvector imply that \(\phi_A\) is
physical. In particular its eigenvalue is the full norm
and the physical norm. Since \(u_A>0\) almost everywhere
on \(A\), the coefficient
\(c_A=|\langle\phi_A,u_A\rangle|^2\) is strictly positive.
Positivity in the operator sense gives
\[
 c_A r_A^{\,n-1}
       \leq p_A(n)\leq\|u_A\|^2r_A^{\,n-1}.
\]
Take \(n\)-th roots at this fixed regulator.
The same operation on the assumed temporal bound proves
\eqref{eq:f15v89:tube-norm}; also \(\|K_A\|\leq1\).
The coefficient \(c_A\) may be arbitrarily small as
volume grows. It disappears before that limit is taken.
\end{proof}

This statement uses all temporal lengths, not just
\(p_A(1)\). For example, for a positive unit vector
\(\omega\) and \(0<\eta<1\), the matrix
\[
 K=(1-\eta)I+\eta|\omega\rangle\langle\omega|
\]
is a positive contraction with strictly positive entries
and Perron vector \(\omega\). If \(A\) is one coordinate
of mass \(\epsilon=|\omega_A|^2\), then
\[
 p_A(1)=\epsilon,\qquad
 \|K_A\|=1-\eta(1-\epsilon),\qquad
 p_A(n)=\epsilon[1-\eta(1-\epsilon)]^{n-1}.
\]
Both \(\epsilon\) and \(\eta\) can tend to zero.
Thus rarity and a nearly invariant bad sector can coexist.
The temporal source input below rules out this persistence
mechanism for the declared native small-staple events.

\subsection{The actual source cost of a small spatial staple}

Fix a spatial link \(e=(x,i)\), on an even spatial torus
of side at least four. Let \(\mathcal P_e\) be the four
spatial plaquettes incident to \(e\), and write
\[
 \rho_e(U)=|S_e(U_{\widehat e})|,\qquad
 A_e(h)=\{\rho_e\leq h\},\qquad 0<h<4.
\]
Although \(\rho_e\) does not depend on \(U_e\), the
following inequality holds for every value of that link:
\begin{equation}
 \sum_{p\in\mathcal P_e}X_p
   =\gamma\{4-\operatorname{Re}(U_eS_e)\}
   \geq\gamma(4-\rho_e),\qquad X_p=\gamma s_p.
 \label{eq:f15v89:energy}
\end{equation}
All quantities are the literal spatial plaquette energies.

The bases of the four positive plaquettes have
multiplicities \(2,1,1\): two at \(x\), and one at
each of \(x-\widehat j\) for the two directions
\(j\ne i\). These are three distinct spatial cells.
By the original positive cell assignment,
\(X_p\leq4T_{\operatorname{base}(p)}\). Consequently,
with \(s=1/32\),
\begin{equation}
 s\sum_{t=0}^{n-1}\sum_{p\in\mathcal P_e(t)}X_p
       \leq\sum_y b_yT_y,
 \label{eq:f15v89:single-source}
\end{equation}
where on each time slice the three nonzero source
values are \(1/4,1/8,1/8\).
Different times give distinct source cells.
In particular
\[
 \|b\|_\infty=\frac14,\qquad
 \sum_yb_y=\frac n2,\qquad
 |\operatorname{supp}b|=3n.
\]

There is an exact fixed-regulator version. Define
\[
 \Delta_{L,\gamma}(b)
       =G_L(\gamma(1-b))-G_L(\gamma),\qquad
 C_{L,\gamma}=\Delta_{L,\gamma}(1/4)
                            +2\Delta_{L,\gamma}(1/8).
\]
Here \(G_L\) is \emph{V62's pressure of the literal
unnormalized Wilson transfer}. It must not be replaced
by the pressure after division by the kinetic Haar
normalizer. That scalar division leaves normalized
spectral ratios and vacuum states unchanged, but not
the pressure differences in a physical action source.

V62's finite-source theorem gives
\begin{equation}
 p_{A_e(h)}(n)
       \leq\exp\left[-n\left\{
          \frac{\gamma(4-h)}{32}-C_{L,\gamma}\right\}\right].
 \label{eq:f15v89:exact-tube}
\end{equation}
Indeed, on the event in question, the left side of
\eqref{eq:f15v89:single-source} is at least
\(n\gamma(4-h)/32\). Exponential Markov followed by
the source theorem bounds its exponential expectation
by \(\exp(nC_{L,\gamma})\).

\begin{theorem}[Small-staple physical operator bound on the selected sequence]
\label{f15v89:single}
Use the selected even spatial sides \(L=M\) and
couplings \(\gamma_M=\widetilde\beta_M\) of V60--V62.
For all sufficiently large \(M\), independently of
\(e\) and temporal length,
\begin{equation}
 \|Q_{A_e(h)}K_MQ_{A_e(h)}\|_{\rm full}
   =\|Q_{A_e(h)}K_MQ_{A_e(h)}|_{\rm phys}\|
   \leq q_M(h),
 \label{eq:f15v89:single-norm}
\end{equation}
where
\begin{equation}
 q_M(h)=\min\left\{1,\,
       \exp\left(3-\frac{\gamma_M(4-h)}{32}\right)\right\}.
 \label{eq:f15v89:q}
\end{equation}
For every fixed \(h<4\), this bound tends to zero
exponentially in \(\gamma_M\).
\end{theorem}

\begin{proof}
Choose the common sequence threshold so that
\[
 e_{\rm vac}\leq5,\qquad k_M\leq5,\qquad
 {\cal E}_M(1/4)\leq\frac1{36}.
\]
It exists by the inherited V62 estimates and is
independent of the source support. For \(0\leq b\leq1/4\),
\[
 J(b)\leq\frac{b^2}{2(1-b)}\leq\frac b6.
\]
The source in \eqref{eq:f15v89:single-source}
therefore obeys
\[
 \log\mathbb E_{\rm vac}e^{\sum_yb_yT_y}
  \leq\left(5+\frac56\right)\frac n2+\frac{3n}{36}
  =3n.
\]
This proves the temporal bound
\(p_{A_e(h)}(n)\leq q_M(h)^n\) for every \(n\).
Apply Proposition~\ref{f15v89:persistence}.
\end{proof}

The order of quantifiers is consequential. First choose
a sufficiently large selected \(M\); then V62 applies
to every finite array at that \(M\), however long its
temporal support. Its error is kept as \(3n/36\),
not sent to zero at fixed \(M\). The limit
\(n\to\infty\) extracts the restricted norm at that
same \(M\). Only after this do we use
\(\gamma_M\to\infty\).
There is no exchange of a growing support with the
infinite-time construction of the vacuum.

The bound is uniform along the selected sequence.
It is not a statement that \(C_{L,\gamma}\leq3\)
for arbitrary independently varied \(L,\gamma\).
At such parameters \eqref{eq:f15v89:exact-tube} is
the available estimate and may be uninformative.

\subsection{Simultaneous, moving and sparse bad patterns}

The same mechanism applies without independence of bad
labels. For a finite nonempty set \(\mathcal S\) of
spatial links, set
\[
 A_{\mathcal S}(h)=\bigcap_{e\in\mathcal S}A_e(h).
\]
Allow a different deterministic set \(\mathcal S_t\)
at each of finitely many successive time slices.
Let \(N_*=\sum_t|\mathcal S_t|\).
Summing \eqref{eq:f15v89:energy} on the corresponding
event gives an energy debit at least
\(\gamma_M(4-h)N_*\).

A spatial plaquette has four boundary links, so its
multiplicity in such a sum at one time is at most four.
There are three positive spatial plaquette orientations
at one base. Thus the summed base multiplicity is at
most twelve. Take \(s_*=1/192\), and define
\[
 b_{(t,y)}=4s_*
       \sum_{\substack{p:\operatorname{base}(p)=(t,y)}}
             \#\{e\in\mathcal S_t:p\in\mathcal P_e\}.
\]
The same literal geometry gives
\begin{equation}
 \|b\|_\infty\leq\frac14,\qquad
 \sum b=\frac{N_*}{12},\qquad
 |\operatorname{supp}b|\leq3N_*.
 \label{eq:f15v89:many-source}
\end{equation}
There is no coloring or disjointness premise.

\begin{theorem}[Operator cost of specified bad patterns and their unions]
\label{f15v89:patterns}
At the same sufficiently large selected \(M\), put
\[
 a_M(h)=\frac{\gamma_M(4-h)}{192}-\frac{41}{72}.
\]
Then the moving-pattern probability satisfies
\begin{equation}
 \mu_{\rm vac}\{U_t\in A_{\mathcal S_t}(h)
                         \text{ for every }t\}
     \leq \exp\left(-a_M(h)\sum_t|\mathcal S_t|\right).
 \label{eq:f15v89:moving}
\end{equation}
In particular,
\begin{equation}
 \|Q_{A_{\mathcal S}(h)}K_MQ_{A_{\mathcal S}(h)}\|
     \leq\min\{1,e^{-a_M(h)|\mathcal S|}\}.
 \label{eq:f15v89:pattern-norm}
\end{equation}
For any finite family \(\mathfrak F\) of nonempty
spatial-link sets and
\(A_{\mathfrak F}=\bigcup_{\mathcal S\in\mathfrak F}
 A_{\mathcal S}(h)\),
\begin{equation}
 \|Q_{A_{\mathfrak F}}K_MQ_{A_{\mathfrak F}}\|
   \leq\min\left\{1,\
        \sum_{\mathcal S\in\mathfrak F}e^{-a_M(h)|\mathcal S|}
              \right\}.
 \label{eq:f15v89:union-norm}
\end{equation}
All these norms can be taken on the full or physical space.
\end{theorem}

\begin{proof}
The source window and \eqref{eq:f15v89:many-source}
bound the logarithmic exponential expectation by
\[
 \left(5+\frac56\right)\frac{N_*}{12}
          +\frac{3N_*}{36}=\frac{41}{72}N_*.
\]
Exponential Markov proves \eqref{eq:f15v89:moving}.
For a fixed pattern use it at every time and apply
Proposition~\ref{f15v89:persistence}.

For the union event at every time, choose a pattern
\(\mathcal S_t\in\mathfrak F\) at that time, and take
the union over all such histories. Equation
\eqref{eq:f15v89:moving} bounds the sum by
\[
 \sum_{(\mathcal S_0,\ldots,\mathcal S_{n-1})
                        \in\mathfrak F^n}
   e^{-a_M(h)\sum_t|\mathcal S_t|}
   =\left(\sum_{\mathcal S\in\mathfrak F}
                   e^{-a_M(h)|\mathcal S|}\right)^n.
\]
Apply the persistence proposition again. This is a
union bound on actual dependent events, not a product
law assumption.
\end{proof}

For example, fix a deterministic graph on spatial links
of maximum degree \(\Delta\geq2\), and a specified root.
If the bad component of that root has at least \(k\)
vertices, it contains a connected bad set of size \(k\)
through that root. The already-used depth-first count
gives at most \(\Delta^{2(k-1)}\) such sets. Thus its
\emph{restricted transfer norm}, not just its one-time
probability, is at most
\[
 \min\{1,\Delta^{2(k-1)}e^{-a_M(h)k}\}.
\]
The combinatorial count is inherited from V71; the
operator upgrade is new.

Separated or hierarchical locations in a specified
\(\mathcal S\) cause no additional loss. Nevertheless,
a union over all possible locations has entropy.
For example, the event that at least \(k\) of \(N\)
specified links are bad has the bound
\[
 \min\{1,\binom Nk e^{-a_M(h)k}\}.
\]
For one bad link anywhere in that set, the sharper
single-link source also gives
\(\min\{1,Nq_M(h)\}\): the proof of the single-link
source bound allows its chosen edge to move at every
time. No uniform suppression of every possible sparse
defect in an arbitrarily large volume follows after
discarding this factor \(N\).

\subsection{From the bad-block norm to full-vector removal}

The following elementary positive-block estimates are
applied to the actual operator bound just proved.
They make explicit what is controlled beyond supported
vectors.

\begin{proposition}[Supported floor, low-energy mass and deletion error]
\label{f15v89:deletion}
Let \(K\) be a positive contraction, \(Q\) an orthogonal
projection and \(G=I-Q\). If \(\|QKQ\|\leq q\leq1\), then
\begin{align}
 \langle f,(I-K)f\rangle&\geq(1-q)\|f\|^2
                     &&(f\in\operatorname{Ran}Q),
             \label{eq:f15v89:supported}\\
 \|Q{\bf1}_{[\alpha,1]}(K)\|&\leq\sqrt{q/\alpha}
                     &&(0<\alpha\leq1),\label{eq:f15v89:mass}\\
 \|K-GKG\|&\leq d(q):=\frac{q+\sqrt{q^2+4q}}2
                      \leq\sqrt q+q. \label{eq:f15v89:delete}
\end{align}
The deletion estimate can also be capped at the
trivial bound one.
\end{proposition}

\begin{proof}
The supported estimate is immediate.
Since
\(\|QK^{1/2}\|^2=\|QKQ\|\leq q\), bounded spectral
inversion on the range of
\(E_\alpha={\bf1}_{[\alpha,1]}(K)\) gives
\[
 QE_\alpha=QK^{1/2}
          (K|_{\operatorname{Ran}E_\alpha})^{-1/2}E_\alpha,
\]
proving \eqref{eq:f15v89:mass}.

In the decomposition \(\operatorname{Ran}G\oplus
\operatorname{Ran}Q\), the difference has blocks
\[
 K-GKG=
   \begin{pmatrix}0&C\\C^*&B\end{pmatrix},
 \qquad B=QKQ,\quad C=GKQ.
\]
Factorization through \(K^{1/2}\) gives
\(\|C\|\leq\sqrt{\|GKG\|\|QKQ\|}\leq\sqrt q\).
For a unit vector with component norms \(x,y\), its
absolute quadratic form is at most
\(2\sqrt q\,xy+qy^2\). The top eigenvalue of the
scalar matrix
\(\left(\begin{smallmatrix}0&\sqrt q\\\sqrt q&q\end{smallmatrix}\right)\)
is \(d(q)\). Self-adjointness proves the norm bound.
Finally both \(K\) and \(GKG\) lie between zero and
the identity, so their difference has norm at most one.
\end{proof}

Apply this proposition with \(Q=Q_{A_e(h)}\) and
\(q=q_M(h)\). It proves a genuine native one-step
supported floor, approaching one, on this bad region.
More strongly, every physical near-top spectral window
has squared operator leakage into that region at most
\(q_M(h)/\alpha\). This is not an estimate only for
the vacuum or for a chosen local observable.

The square-root deletion loss cannot generally be
replaced by \(q\). For the rank-one positive contraction
\[
 K=|(\sqrt{1-q},\sqrt q)\rangle
             \langle(\sqrt{1-q},\sqrt q)|,
 \qquad Q=\begin{pmatrix}0&0\\0&1\end{pmatrix},
\]
the off-diagonal block has norm \(\sqrt{q(1-q)}\).
Its presence has been paid in \eqref{eq:f15v89:delete}.
For the exponentially small native \(q_M(h)\),
this loss remains exponentially small.

\subsection{A bounded-rank selected fibre everywhere}

Fix \(0<h<4\), and use the exact V87 compact oscillator
projection on the good region. Define
\begin{equation}
 R(V)=
 \begin{cases}
 \Pi_{\gamma_M,m}(\rho_e(V),v_e(V)),&\rho_e(V)>h,\\
 0,&\rho_e(V)\leq h.
 \end{cases}
 \label{eq:f15v89:finite-fibre}
\end{equation}
This is an orthogonal, gauge-equivariant decomposable
projection. Its rank in the selected link is at most
\(\binom{m+3}{3}\) everywhere, unlike the identity
retained by V88 on the bad region.

Let \(\varepsilon_{m,\gamma}(h,B)\) and \(h_*(h,B)\)
have exactly the meanings in V87--V88. Let
\(\Psi_{18}\) be \eqref{eq:f15v88:psi}.
No new Gaussian or growing-band estimate is assumed.

\begin{theorem}[Native finite-fibre reduction after paid bad-sector removal]
\label{f15v89:finite-fibre}
At all sufficiently large selected \(M\),
\begin{equation}
 \frac{\|\bar T_M-R\bar T_MR\|_{\rm full}}{\Lambda_M}
       \leq\eta_{m,M}(h,B),
 \label{eq:f15v89:finite-error}
\end{equation}
where
\begin{equation}
 \eta_{m,M}(h,B)
    =d(q_M(h))
       +\frac{2\varepsilon_{m,\gamma_M}(h,B)}{h_*(h,B)}
       +2\Psi_{18}\left(\frac{h^2B^2}{6\pi^2}\right).
 \label{eq:f15v89:eta}
\end{equation}
The same estimate holds after physical restriction.
For each fixed tolerance \(\tau>0\) and fixed \(h\),
there are finite \(B,m\) and a sequence threshold
\(M_0\) such that \(\eta_{m,M}<\tau\) for all
selected \(M\geq M_0\).
\end{theorem}

\begin{proof}
Set \(G=I-Q_{A_e(h)}\). Since \(R=GR=RG\),
\[
 \bar T-R\bar T R
   =(\bar T-G\bar T G)
                 +(G\bar T G-R\bar T R).
\]
Proposition~\ref{f15v89:deletion} pays the first term
by \(d(q_M(h))\Lambda_M\).

Both endpoints of the second term have staple magnitude
greater than \(h\). Split its pair kernel into
\(\vartheta\leq B/\sqrt{\gamma_M}\) and its complement.
V87's scalar ground compression and two-sided fibre
error give the near contribution
\(2\varepsilon_{m,\gamma_M}h_*^{-1}\Lambda_M\).
On the far part the error is exactly a masked transfer
minus its double compression. Its norm is at most
twice the norm of the masked transfer; V88's
native angular bound, now with threshold \(h\), gives
the final term of \eqref{eq:f15v89:eta}.
The masks depend only on other-link variables, as
required for this compression identity.

For the final statement, choose \(B\) large, then \(m\)
large but fixed, using V87's level tail. Finally let
\(M\) grow so that its compact remainders are small and
\(q_M(h)\) is exponentially small. The source threshold
and all error constants have the stated uniform scope.
\end{proof}

Thus the retained operator is an actual matrix-valued
integral operator on the good other-link region, with
fibre dimension at most \(\binom{m+3}{3}\). Its entries
are the exact native matrix elements
\[
 a(V)c_{\widehat e}(V,W)a(W)
 \langle\psi_{\alpha,V},K_{\gamma_M}(V,W)
                               \psi_{\beta,W}\rangle.
\]
The moving endpoint bases and their gauge transport
are retained. It is not a product of independent
one-link spectral ratios.

For \(h=2\), the newly paid bad factor is explicitly
\[
 q_M(2)\leq \exp(3-\gamma_M/16)
\]
once that expression is below one. The \(O(\gamma_M^{-1})\)
compact channel error and its dependence on \(B,m\)
have not been improved in this iteration.

\subsection{What is still upstream of a physical mass gap}

V89 removes a specific infinite-dimensional obstruction
in V88: small-staple fibres at the selected edge need
no longer remain unchanged. It also rules out a specified
persistent sparse or hierarchical small-staple pattern
by an operator bound, not a probability slogan.
However, the retained operator still includes all
collective neutral motion of the other links.

The min--max and spectral-window estimates of V88 apply
with \(\eta_{m,M}\) in place of its hybrid error.
They do not supply a bound for the retained leading
ratio. The target remains
\[
 r_{{\rm retained},M}+\eta_{m,M}
             \leq e^{-m_*a_{t,M}},\qquad m_*>0,
\]
on the independently specified physical clock.
A fixed-tolerance approximation is not that estimate.
Nor can noncommuting projections at overlapping links
be combined without further work. An all-location
union must pay its explicit spatial entropy.

The next upstream task is to analyze the coupled
good-region matrix-valued operator and its neutral
collective low modes. The paid bad-pattern estimates
can be used when forming that analysis, but the remaining
good sector cannot be assigned a gap merely because its
individual conditional fibres are confined. The
interacting continuum construction and its positive
physical mass gap remain open.

\subsection{Diagnostics and append record}

The new script
\begin{center}\small
\path{scripts/verify_ym_f15_v89_persistence_bad_sector.py}
\end{center}
reruns V88 and the inherited chain. Its independent
rational fixtures check literal source carriers in all
spatial orientations, arbitrary overlapping and moving
patterns, the extensive per-cell budgets, stationary
Markov path probabilities against restricted transfer
moments, rare-but-sticky countermodels, positive-block
deletion and its square-root loss, and temporal union
entropy. Finite tests do not replace the analytic
all-length persistence proof or certify a physical gap.

All 1,833 new exact-rational fixtures passed, along
with the inherited diagnostic chain. The checker has
7,610 bytes and SHA--256
\begin{center}\scriptsize\ttfamily
5C79353E01861EDBE4E6814C61904A33543C0B8F2934D796C02EEEBDD120D55B.
\end{center}

The predecessor V88 has 1,993,559 bytes and SHA--256
\begin{center}\scriptsize\ttfamily
BDCA0C72BA7D68E76A22393A5AA44E1F8E0A23D3456F06FFEF8E8683890A3479.
\end{center}
Its body is preserved byte for byte. Removing this final
section and restoring the title version recovers V88.
V89 replaces only the active main source. Protected
archives, companions, supplied papers and prior
diagnostics remain unchanged. Only \path{.tex} files
belong in \path{latex_notes}; builds are directed to
\path{artifacts/ym_f15_v89_texcheck/}.

\begin{firewall}
V89 proves actual small-staple restricted-transfer
bounds along the selected vacuum sequence, full-vector
near-top exclusion from those regions, and their
controlled removal in a native finite-fibre reduction.
It does not prove a gap for the remaining neutral
collective operator, calibrate the physical clock,
or construct a nontrivial interacting continuum
Yang--Mills theory. The original goal remains active.
\end{firewall}

\section{V90: source-weighted transfer bounds and local physical energy}
\label{f15v90:context}

V89 upgraded a vacuum probability estimate to a norm
bound for an operator killed on a small-staple region.
This iteration uses the same all-length source input
without replacing the source by an indicator. It
bounds the entire source-weighted native transfer.
Taking an operator logarithm then controls local
magnetic energy by the \emph{actual logarithmic transfer
energy}, rather than by an auxiliary sampling gradient.

The result also supplies exponential local-energy
moments uniformly for every vector in a declared
near-top spectral window. Thus the estimate is not
restricted to the vacuum or to vectors created by
a preselected set of probes. The continuum construction,
physical clock and collective gap are not supplied
by these bounds. The source theorem of V62 and the
native finite-regulator positivity facts are reused,
not claimed anew.

\subsection{The V90 companion checkpoint}

The current companion files were checked again at this
required five-version checkpoint. Their latest available
versions remain SM continuum V72, NS stability V26,
YM shell mixing V16 and parallel YM V5. The downloaded
exact-action Einstein V24, regenerative stationarity V34
and SM source-selection V31 inputs were also revisited.
Their protected hashes are unchanged.

The SM result transports fixed physical mode sets at
fixed time; it does not supply a growing-cutoff theorem.
The shell note controls a measured local one-loop
coefficient and a finite-volume coefficient remainder,
not the full joint nonperturbative transfer remainder.
The parallel YM note's distinction between a fixed
window and a global, spatially extensive activity is
particularly relevant below: source costs grow with
the number of plaquettes. The NS pointwise kernel
calculation supplies no missing YM spectral estimate.
The other three notes retain, respectively, a nonlinear
root problem beyond tangent correction, a stationarity
closure beyond a first-escape reduction, and native
source-typing requirements beyond a conditional compiler.

None of these statements is imported as a new proof
of a physical gap. The useful direction remains a
native all-vector transfer estimate with explicit
support and clock dependence. V88's undecidability
audit is retained: no finite local calculation, source
bound or summable truncation error is taken to exclude
all slow retained modes. The next companion checkpoint
is V95.

\subsection{The positive-source spectral-radius bridge}

At a fixed finite spatial regulator and positive
coupling, use the full native operator
\(\bar T\) of \eqref{eq:f15v88:native} and set
\[
 \Lambda=\|\bar T\|=\lambda_0(\bar T|_{\rm phys}),\qquad
 K=\bar T/\Lambda,\qquad K\Omega=\Omega,\qquad \|\Omega\|=1.
\]
The kernel is strictly positive, compact and
gauge invariant, with \(\Omega>0\). Moreover \(K\)
is a positive injective contraction. For injectivity,
every Peter--Weyl eigenvalue of each \(C_\gamma\)
in \eqref{eq:f15v88:native} is strictly positive,
and the finite-volume multiplier and its inverse
are bounded. The same facts give injectivity on
the physical subspace.

Let \(F\) be a bounded, nonnegative, gauge-invariant
real slice multiplication at this fixed regulator.
Its bound need not be uniform in coupling or spatial
volume. Define
\[
 W=e^{F/2},\qquad L_F=WKW,\qquad
 Z_F(n)=\mathbb E_{\rm vac}
              \exp\left\{\sum_{t=0}^{n-1}F(U_t)\right\}.
\]
The expectation is in the actual stationary Wilson
vacuum. No independence of time slices is assumed.

\begin{proposition}[All-length source cost controls a whole weighted transfer]
\label{f15v90:weighted}
At the fixed regulator,
\begin{equation}
 Z_F(n)=\langle W\Omega,L_F^{\,n-1}W\Omega\rangle,\qquad
 \lim_{n\to\infty}Z_F(n)^{1/n}=\|L_F\|.
 \label{eq:f15v90:cylinder}
\end{equation}
The norm on the right is the same on the full and
physical spaces. If a finite \(C\) satisfies
\begin{equation}
 Z_F(n)\leq e^{nC}\qquad(n\geq1),
 \label{eq:f15v90:source-premise}
\end{equation}
then
\begin{equation}
 \|e^{F/2}Ke^{F/2}\|\leq e^C,\qquad
 0\leq K\leq e^C e^{-F},\qquad
 \|e^{F/2}K^{1/2}\|^2\leq e^C .
 \label{eq:f15v90:weighted-order}
\end{equation}
An additional finite prefactor independent of \(n\)
in \eqref{eq:f15v90:source-premise} does not change
these conclusions.
\end{proposition}

\begin{proof}
The literal vacuum insertion formula is
\[
 Z_F(n)=\langle\Omega,e^F K e^F K\cdots K e^F\Omega\rangle,
\]
with \(n\) insertions and \(n-1\) transfers.
Collecting the half-source factors proves the first
identity in \eqref{eq:f15v90:cylinder}.
Gauge-invariant insertions preserve the physical
subspace, so the full operator computes the same
physical vacuum cylinder.

The operator \(L_F\) is positive and compact with
a strictly positive kernel. Its largest eigenvalue
\(r_F=\|L_F\|>0\) has a unit positive eigenvector
\(\phi_F\). Uniqueness and gauge invariance imply
that \(\phi_F\) is physical. Since \(W\Omega>0\),
\(c_F=|\langle\phi_F,W\Omega\rangle|^2>0\). Therefore
\[
 c_F r_F^{\,n-1}\leq Z_F(n)
                 \leq \|W\Omega\|^2r_F^{\,n-1}.
\]
Take \(n\)-th roots at this fixed regulator. There
is no need for a volume-uniform lower bound on \(c_F\).
This proves the limit and the first inequality in
\eqref{eq:f15v90:weighted-order}.
The congruence of \(WKW\leq e^C I\) by \(W^{-1}\)
proves the operator order. Finally, for
\(A=WK^{1/2}\), \(AA^*=WKW\), giving the norm bound.
No commutation of \(F\) with \(K\) has been used.
\end{proof}

This is the positive exponential-source extension of
V89's indicator argument. In particular, a source
estimate at one time alone would not prove
\eqref{eq:f15v90:weighted-order}. The per-cell errors
in the repeated source must be paid at all lengths
before taking the spectral-radius limit.

\subsection{Logarithmic physical energy and entire spectral windows}

Define the nonnegative self-adjoint transfer energy
\[
 H=-\log K
\]
by spectral calculus. This is dimensionless. If an
independent physical time spacing \(a_t>0\) is supplied,
then \(H_{\rm phys}=H/a_t\); that choice is not derived
in this iteration.

\begin{theorem}[Local energy is bounded by actual transfer energy]
\label{f15v90:energy}
Under the hypotheses of Proposition~\ref{f15v90:weighted},
\begin{equation}
 H\geq F-CI
 \label{eq:f15v90:log-order}
\end{equation}
as quadratic forms on \(D(H^{1/2})\).
In particular, for a unit vector in that domain,
\begin{equation}
 \langle\psi,F\psi\rangle\leq C+\|H^{1/2}\psi\|^2.
 \label{eq:f15v90:mean}
\end{equation}
For \(E_\alpha={\bf1}_{[\alpha,1]}(K)\), \(0<\alpha\leq1\),
one also has
\begin{equation}
 E_\alpha e^F E_\alpha\leq
       \frac{e^C}{\alpha}E_\alpha,\qquad
 \|e^{F/2}E_\alpha\|^2\leq\frac{e^C}{\alpha}.
 \label{eq:f15v90:window}
\end{equation}
These are full-subspace bounds. They apply, in particular,
to every physical unit vector and every density matrix
supported in the indicated spectral window.
\end{theorem}

\begin{proof}
We give the operator-logarithm step explicitly to
avoid any scalar or commuting substitution.
For bounded strictly positive invertible operators
\(A\leq B\), inversion reverses order. Indeed
\(D=A^{-1/2}BA^{-1/2}\geq I\), so
\(D^{-1}\leq I\) by spectral calculus and hence
\(B^{-1}\leq A^{-1}\). Applying this to \(A+tI,B+tI\)
and using the norm-convergent difference integral gives
\[
 \log B-\log A
   =\int_0^\infty
          \big\{(A+tI)^{-1}-(B+tI)^{-1}\big\}\,dt
   \geq0 .
\]
The formula follows from the scalar logarithm
representation by continuous functional calculus;
the integrand is bounded near zero and is \(O(t^{-2})\)
in norm at infinity.

In \eqref{eq:f15v90:weighted-order} put
\(B=e^Ce^{-F}\), which is bounded and bounded below
at this fixed regulator. For every \(\varepsilon>0\),
\(K+\varepsilon I\leq B+\varepsilon I\). Thus
\[
 -\log(K+\varepsilon I)\geq-\log(B+\varepsilon I).
\]
The right side converges in norm to \(F-CI\).
For a vector in \(D(H^{1/2})\), the left quadratic
form converges to that of \(H\) by the spectral
theorem and monotone convergence (with a fixed
lower bound for \(0<\varepsilon\leq1\)).
This proves \eqref{eq:f15v90:log-order} and
\eqref{eq:f15v90:mean}.

Bounded inversion on the spectral window gives
\[
 e^{F/2}E_\alpha
     =e^{F/2}K^{1/2}K^{-1/2}E_\alpha .
\]
The last factor has norm at most \(\alpha^{-1/2}\).
Use \eqref{eq:f15v90:weighted-order}, then square
the norm. The equivalent compressed positive-operator
inequality proves \eqref{eq:f15v90:window}.
Taking its trace against a positive trace-one
operator supported in the window gives the mixed-state
assertion. The mean-energy bound likewise extends
to density matrices with finite \(\operatorname{Tr}(DH)\)
by their positive spectral decomposition.
\end{proof}

More generally, for \(\psi\in D(K^{-1/2})\),
\begin{equation}
 \|e^{F/2}\psi\|^2\leq e^C\|K^{-1/2}\psi\|^2.
 \label{eq:f15v90:inverse}
\end{equation}
The right side is an exponential moment of transfer
energy, not its first moment. For a hard physical
energy window \({\bf1}_{[0,E]}(H_{\rm phys})\),
\eqref{eq:f15v90:window} reads
\[
 \langle e^F\rangle\leq e^{C+a_tE},
 \qquad \langle F\rangle\leq C+a_tE.
\]
The second inequality also follows directly from
\eqref{eq:f15v90:mean}. Neither assertion prescribes
\(a_t\).

\subsection{A fixed mean energy is not a hard energy window}

The distinction just made is necessary even for
strictly positive finite-state transfer kernels.
Let \(N\geq2\), \(q=e^{-N}\), and
\[
 \omega=(\sqrt{1-q},\sqrt q),\quad
 v=(-\sqrt q,\sqrt{1-q}),\quad
 K_N=|\omega\rangle\langle\omega|
                         +q|v\rangle\langle v|,\quad
 F_N=\operatorname{diag}(0,N/2).
\]
Then \(K_N\) is a positive contraction with strictly
positive entries, Perron vector \(\omega\), and
\(H_N=-\log K_N=N|v\rangle\langle v|\).
Writing \(W_N=e^{F_N/2}\), positivity and the trace
bound give
\[
 \begin{split}
 \|W_NK_NW_N\|
 &\leq\operatorname{Tr}(W_NK_NW_N)\\
 &=1-q+q^2+e^{N/2}(2q-q^2)
 \leq1+2e^{-N/2}\leq3,\\
 \|W_N\omega\|^2&=1-q+qe^{N/2}\leq2 .
 \end{split}
\]
Consequently its genuine stationary cylinders obey
\(Z_{F_N}(n)\leq2\cdot3^{n-1}\leq3^n\).
The source premise therefore holds uniformly with
\(C=\log3\), in the same strictly positive-kernel
class as the abstract proposition.

Nevertheless the density matrix
\[
 D_N=(1-N^{-1})|\omega\rangle\langle\omega|
                         +N^{-1}|v\rangle\langle v|
\]
satisfies
\[
 \operatorname{Tr}(D_NH_N)=1,\qquad
 \operatorname{Tr}(D_Ne^{F_N})
       \geq\frac{1-e^{-N}}{N}e^{N/2}\longrightarrow\infty .
\]
It has a small weight at arbitrarily high energy,
so it is not supported in a fixed spectral window.
This is a counterexample to a stronger inference
from the abstract source bound, not a counterexample
to the Yang--Mills mass-gap conjecture.

\subsection{Literal spatial Wilson sources and their costs}

Let \(X_p=\gamma s_p\geq0\) denote the actual spatial
plaquette energy in the existing \(\SU(2)\) convention.
For finitely many nonnegative spatial plaquette
weights \(w_p\), put
\[
 F_w=\sum_p w_pX_p,\qquad
 b_x=4\sum_{\operatorname{base}(p)=x}w_p.
\]
Assume \(\max_xb_x<1\). The positive cell assignment
\(X_p\leq4T_{\operatorname{base}(p)}\) implies
\[
 \sum_{t=0}^{n-1}F_w(U_t)
       \leq\sum_{t=0}^{n-1}\sum_x b_xT_{(t,x)} .
\]
Cells at different times are distinct source
coordinates. Shared plaquettes in their cell
energies cause no problem: all contributions
are nonnegative and the cell-source comparison
allows arbitrary overlaps.

At fixed \(L,\gamma\), define the exact source cost
\begin{equation}
 C_{L,\gamma}(w)=\sum_{x:b_x\ne0}
       [G_L(\gamma(1-b_x))-G_L(\gamma)] .
 \label{eq:f15v90:exact-cost}
\end{equation}
As in V89, \(G_L\) is the pressure of the literal
\emph{unnormalized} Wilson transfer used in V62.
Replacing it by the pressure after kinetic Haar
normalization would change this cost and is not
allowed. Theorem~\ref{f15v62:source}, with its centering
terms canceled, proves
\[
 Z_{F_w}(n)\leq \exp\{nC_{L,\gamma}(w)\}
                       \qquad(n\geq1).
\]
Thus all conclusions above apply at the exact cost
\eqref{eq:f15v90:exact-cost}.

For an explicit bound use the same selected sequence
\(L=M,\gamma=\gamma_M\) and common threshold as in
V89. In particular,
\[
 e_{\rm vac}\leq5,\quad k_M\leq5,\quad
 {\cal E}_M(1/4)\leq1/36.
\]
If \(0\leq b_x\leq1/4\), then
\(J(b_x)\leq b_x^2/[2(1-b_x)]\leq b_x/6\), so
\begin{equation}
 C_M(w)\leq
 5\sum_xb_x+5\sum_xJ(b_x)
                         +\frac{|\operatorname{supp}b|}{36}
 \leq\frac{35}{6}\sum_xb_x
                         +\frac{|\operatorname{supp}b|}{36}.
 \label{eq:f15v90:cost}
\end{equation}
Here \(C_M(w)=C_{M,\gamma_M}(w)\).
The first inequality follows from V62's scalar
pressure-window bound, not from differentiation
of a source estimate. In particular the additive
per-cell floor in this bound cannot be differentiated
at zero. At a zero source the exact cost is zero.

This threshold is independent of source cardinality
and temporal length. For every finite repeated
source the error is \(n|\operatorname{supp}b|/36\).
Its division by \(n\), not its omission, yields
\eqref{eq:f15v90:cost}. The limit in
\eqref{eq:f15v90:cylinder} is always taken at a
fixed selected \(M\).

\begin{theorem}[Explicit local magnetic bounds for every physical low-energy state]
\label{f15v90:wilson}
For all sufficiently large selected \(M\), uniformly
in spatial location, the following statements hold.

For one spatial plaquette,
\begin{equation}
 H_M\geq \frac{X_p}{16}-\frac{107}{72}I,\qquad
 E_\alpha e^{X_p/16}E_\alpha
          \leq\frac{e^{107/72}}{\alpha}E_\alpha .
 \label{eq:f15v90:one}
\end{equation}
For a finite set \(\mathcal P\) of \(m\geq1\)
distinct spatial plaquettes,
\begin{equation}
 \begin{split}
 H_M&\geq \frac1{48}\sum_{p\in\mathcal P}X_p
                                      -\frac{37m}{72}I,\\
 E_\alpha\exp\left(\frac1{48}
                        \sum_{p\in\mathcal P}X_p\right)E_\alpha
      &\leq\frac{e^{37m/72}}{\alpha}E_\alpha .
 \end{split}
 \label{eq:f15v90:patch}
\end{equation}
The form inequalities use \(D(H_M^{1/2})\), and
\(E_\alpha={\bf1}_{[\alpha,1]}(K_M)\).
No separation or independence of the plaquettes
is required.
\end{theorem}

\begin{proof}
For \(F=X_p/16\) the single nonzero cell source is
\(b=1/4\). Formula \eqref{eq:f15v90:cost} gives
\[
 \frac54+\frac5{24}+\frac1{36}=\frac{107}{72}.
\]
Apply Theorem~\ref{f15v90:energy}.

For the finite set, use
\(F=\sum_{\mathcal P}X_p/48\).
There are three positive spatial orientations at
a base, hence
\[
 b_x=\frac{\#\{p\in\mathcal P:\operatorname{base}(p)=x\}}{12}
       \leq\frac14,\qquad
 \sum_xb_x=\frac m{12},\qquad
 |\operatorname{supp}b|\leq m .
\]
The cost is therefore at most
\[
 \frac{35}{6}\frac m{12}+\frac m{36}
                         =\frac{37m}{72}.
\]
Apply the same theorem. These are comparisons for
the interacting Wilson transfer, not for a product
of single-plaquette laws.
\end{proof}

For example, the squared norm of the entire spectral
window leaking into a large local-energy event obeys
\begin{equation}
 \left\|{\bf1}_{\{\sum_{\mathcal P}X_p\geq R\}}E_\alpha\right\|^2
       \leq\min\left\{1,\,
          \alpha^{-1}\exp\left(\frac{37m}{72}-\frac R{48}\right)
                    \right\}.
 \label{eq:f15v90:patch-tail}
\end{equation}
Indeed the indicator is bounded pointwise by
\(e^{-R/48}\exp(\sum_{\mathcal P}X_p/48)\);
compress and use \eqref{eq:f15v90:patch}.
The inequality is for arbitrary superpositions in
the window, including neutral collective states.

For a spatial link define only in this subsection
the \emph{four-spatial-staple} deficit
\[
 {\mathcal F}^{\rm sp}_e=\gamma_M(4-\rho_e)\geq0 .
\]
It is not the six-plaquette spacetime deficit of V66.
Equation \eqref{eq:f15v89:energy} gives
\({\mathcal F}^{\rm sp}_e\leq\sum_{p\in\mathcal P_e}X_p\).
Use \(F={\mathcal F}^{\rm sp}_e/32\). V89's three-cell
source \(1/4,1/8,1/8\) bounds every repeated source
with cost at most three. Thus
\begin{equation}
 H_M\geq{\mathcal F}^{\rm sp}_e/32-3I,\qquad
 E_\alpha e^{{\mathcal F}^{\rm sp}_e/32}E_\alpha
                       \leq(e^3/\alpha)E_\alpha .
 \label{eq:f15v90:staple}
\end{equation}
In particular, for \(r\geq0\),
\[
 \|{\bf1}_{\{\rho_e\leq4-r/\gamma_M\}}E_\alpha\|^2
       \leq\min\{1,\alpha^{-1}e^{3-r/32}\}.
\]
When the threshold is negative the event is empty.
At \(r=\gamma_M(4-h)\) this recovers the earlier
fixed-threshold estimate; that specialization is not
a new result. The new statement controls the entire
thermal deficit distribution and its transfer-energy
form.

\subsection{Thermal curvature tightness without a collective gap claim}

Write the spatial plaquette holonomy as
\(U_p=(u_{p,0},u_{p,\mathrm{im}})\) and let
\(Y_p=\sqrt{\gamma_M}\,u_{p,\mathrm{im}}\),
using any of the previously fixed root transports
when a vector frame is needed. Its norm is invariant
under such transports, and the exact compact identity is
\[
 |Y_p|^2=2X_p-X_p^2/\gamma_M\leq2X_p .
\]
For every density matrix \(D\) supported in the
physical \(E_\alpha\) window, \eqref{eq:f15v90:one}
therefore gives
\begin{equation}
 \operatorname{Tr}(D e^{|Y_p|^2/32})\leq
                         \alpha^{-1}e^{107/72},\qquad
 \begin{cases}
 \operatorname{Tr}(D X_p^k)
      \leq\alpha^{-1}e^{107/72}16^k k!,\\
 \operatorname{Tr}(D |Y_p|^{2k})
      \leq\alpha^{-1}e^{107/72}32^k k!,
 \end{cases}
 \label{eq:f15v90:moments}
\end{equation}
for every integer \(k\geq0\).
The moment bounds use \(z^k\leq k!e^z\) for \(z\geq0\).
For a finite patch, likewise
\[
 \operatorname{Tr}\left(D
       e^{\sum_{p\in\mathcal P}|Y_p|^2/96}\right)
            \leq\alpha^{-1}e^{37m/72}.
\]
Thus fixed-patch configuration distributions are
uniformly tight at the thermal curvature scale
when \(m\) is fixed and \(\alpha\geq\alpha_0>0\).
The estimates are uniform over all such physical
states, not just over a vacuum subsequence.
They do not identify the limiting distribution
of an arbitrary excited state with the previously
identified vacuum law.

For a state with merely finite mean transfer energy,
the form bound, but not a uniform exponential bound,
still applies:
\[
 \operatorname{Tr}\left(D\sum_{\mathcal P}X_p\right)
       \leq48\operatorname{Tr}(D H_M)+\frac{74m}{3}.
\]
The extensive term is indispensable here.
A fixed physical region generally contains a growing
number of lattice plaquettes; this note does not
silently replace its \(m\) by a constant.

The advance is an actual-transfer local-energy
bound for the complete physical spectral subspace.
It is not Hilbert-space compactness of the states,
a bound on their electric derivatives, contraction
of a spatially coupled retained operator, or
long-distance clustering. In particular a local
multiplication bound remains valid after tensoring
with an unobserved positive slow factor; V88's
retained-sector counterexample already demonstrates
why such factors cannot be dismissed by a local
truncation estimate. Here they must instead be
controlled through native collective couplings.

The next missing estimate is therefore still a
coercivity or leading-ratio bound for the retained
neutral collective operator, with error small on
the independently supplied physical time scale.
V90's local-energy form can be used in that
analysis, but it is not a positive lower bound
on \(H_M\) orthogonal to its vacuum. No conclusion
about decidability or undecidability of this
particular YM family follows.

\subsection{Exact diagnostics and append record}

The checker
\begin{center}\small
\path{scripts/verify_ym_f15_v90_weighted_transfer_energy.py}
\end{center}
reruns the inherited V89 chain. New rational fixtures
compare stationary weighted cylinders with their
half-source transfer expression; check noncommuting
congruences and regularized resolvent order; test
entire, including proper multidimensional, spectral
windows; distinguish a mean energy from a hard
window; and verify literal plaquette, patch and
four-spatial-staple source costs.
The strictly positive finite-state example's
uniform tilted trace and endpoint bounds are
also checked. These tests do not replace the
analytic infinite-time or operator-logarithm
arguments, and do not certify a physical gap.

All 1,854 new exact-rational fixtures passed, together
with the inherited chain. The checker has 8,470 bytes
and SHA--256
\begin{center}\scriptsize\ttfamily
E9BB303A3351B62E2555E53DBC23C75E478AAA1DFC00BD1BBB885314E7556E58.
\end{center}

The predecessor V89 has 2,016,083 bytes and SHA--256
\begin{center}\scriptsize\ttfamily
D8366970C85F494A20096362301773E28BE925B68243CF04EC0E3AA52F125C4B.
\end{center}
Its body is preserved byte for byte. Removing this
section and restoring the title version recovers V89.
V90 replaces only the active main source. The
protected archives, companion notes, supplied papers
and previous diagnostics remain unchanged.
Only \path{.tex} files belong in \path{latex_notes};
builds are directed to
\path{artifacts/ym_f15_v90_texcheck/}.

\begin{firewall}
V90 proves native source-weighted operator bounds,
local magnetic energy controlled by actual logarithmic
transfer energy, and exponential thermal-curvature
bounds for every state in the declared physical
spectral windows, along the selected sequence.
It does not prove the retained collective gap,
calibrate physical time, or construct an interacting
four-dimensional continuum theory. The full
Yang--Mills mass gap remains unproved.
\end{firewall}

\section{V91: native electric tails and a targeted literature audit}
\label{f15v91:context}

This iteration incorporates the newly supplied literature
memorandum after the completed V90 companion checkpoint.
Its recommendations are source material, not instructions
or established theorems. In particular, a targeted search
does not establish novelty of our approach. The objective
is still a nontrivial interacting four-dimensional
continuum Yang--Mills theory with a positive physical
Hamiltonian gap; that objective remains unproved.

V90 supplied local magnetic-energy bounds for the actual
transfer spectral windows. The new positive result below
controls their electric representation labels and gives
a finite-rank approximation of the complete native
transfer, with explicit coupling and volume dependence.
It uses complex analytic continuation of its literal
Wilson kernel, not a replacement by a heat kernel.
The estimate holds at every finite regulator specified
below and every positive coupling, not only along the
selected pressure sequence used in V90.

Two proposed literature imports receive direct tests.
The global Hilbert-projective diameter cannot detect
a volume-uniform gap here, even on the gauge-invariant
positive cone. Separately, particular positivity and
physical-clock steps in a recent claimed construction
do not follow from their stated estimates. These are
limitations of specified arguments, not proofs that
Yang--Mills is gapless. Fusion recoupling and comparisons
on an actual integrated block remain useful candidates;
neither is assumed to contract the retained neutral
collective degrees of freedom.

\subsection{Literature ledger and the scope of the checks}
\label{f15v91:literature}

The following primary sources were checked to different
depths. The depth is part of the ledger.

\begin{enumerate}
\item Lemmens--Nussbaum,
\emph{Birkhoff's version of Hilbert's metric and its
applications in analysis},
\href{https://arxiv.org/abs/1304.7921}{arXiv:1304.7921},
pp.~9--10, Theorem~2.9 and its preceding cone discussion,
was read in text and rendered form. We use its classical
identity between projective diameter and the optimal
positive-cone contraction coefficient. The new native
diameter calculations are proved below.

\item Delcamp, Dittrich and Riello,
\href{https://arxiv.org/abs/1607.08881}{arXiv:1607.08881},
and Cunningham, Dittrich and Steinhaus,
\href{https://arxiv.org/abs/2002.10472}{arXiv:2002.10472},
were checked at the level of their declared setting.
The former develops a fusion basis for the specified
\(2+1\)-dimensional setting; the latter studies
quantum-group \(\SU(2)_k\) tensor-network coarse graining
in three spacetime dimensions. They motivate a
recoupling of V82's retained boundary intertwiners.
Their dimension, representation category and dynamics
must not silently be substituted for the present
four-dimensional ordinary-\(\SU(2)\) Wilson law.
No spectral inequality from either paper is imported.

\item Lawler--Sokal's general-state-space Cheeger framework
and Diaconis--Saloff-Coste's reversible-chain comparison
framework are possible tools for V78's exact Doob kernel.
The primary-source entry points are
\href{https://www.ams.org/journals/tran/1988-309-02/S0002-9947-1988-0930082-9/}
{the AMS article} and
\href{https://pi.math.cornell.edu/~lsc/papers/comp-rev.pdf}
{the authors' comparison paper}.
No numerical constant or applicability theorem is being
imported from an abstract or bibliographic record.
A comparison requires an actual same-space quadratic
form bound and compatible stationary measures;
conductance requires control of every measurable cut,
not just typical local plaquette events.

\item Shen--Zhu--Zhu,
\href{https://arxiv.org/abs/2204.12737}{arXiv:2204.12737},
explicitly works at strong coupling; its stated
\(\SU(N)\) sufficient region includes
\(\lvert\beta\rvert<1/[16(d-1)]\).
This is useful evidence of a rigorous functional-inequality
method in its own regime. It is not a uniform estimate
on our \(\gamma\to\infty\) trajectory, and an auxiliary
stochastic generator is not identified with the
physical logarithmic transfer Hamiltonian.

\item Faizal--Shabir,
\href{https://arxiv.org/abs/2606.19362v1}{arXiv:2606.19362v1},
\href{https://doi.org/10.1002/prop.70097}
{doi:10.1002/prop.70097},
is a 593-page combined file including the main paper
and supporting papers. We audited the main paper's
Lemma~5.3, Theorems~5.4 and 5.6, and the relevant
remainder argument of Theorem~10.4.
The corresponding printed pages are 50--53, 55,
and 114--117; selected formula pages were also rendered
and inspected. This is not a theorem-by-theorem audit
of all 593 pages. The specific failed implications
below are sufficient reason not to import the
claimed complete gap through those arguments.

\item The gapped-reference hypotheses in
\href{https://arxiv.org/abs/math-ph/0412040}{Yarotsky's
relatively bounded perturbation work} and the
commuting-projector/topological-order setting of
\href{https://arxiv.org/abs/1001.0344}
{Bravyi--Hastings--Michalakis} do not provide our
missing native reference gap. Balaban's RG work and
the Magnen--Rivasseau--Seneor ultraviolet construction
remain prospective sources from the memorandum;
their detailed hypotheses have not been audited
in this iteration and no continuum conclusion is
imported from them.
\end{enumerate}

\subsection{The literal kernel and its electric generators}

Fix a spatial cubic torus of even side \(L\geq4\).
Let \(\mathcal E\) be its oriented positive links,
\(E=|\mathcal E|\), and
\(\mathcal X=\SU(2)^{\mathcal E}\), with product
probability Haar measure. Each spatial link occurs
once in each of four spatial plaquettes. For
\(\gamma>0\), write
\begin{equation}
 b(U)=\exp\left[-\frac{\gamma}{2}
            \sum_{p\ {\rm spatial}}s_p(U)\right],
 \qquad
 \bar T=M_b C_\gamma^{\otimes E}M_b,\qquad
 K=\bar T/\Lambda,\quad \Lambda=\|\bar T\|.
 \label{eq:f15v91:native}
\end{equation}
Here \(s_p=1-\frac12\Tr U_p\), and \(C_\gamma\)
is exactly \eqref{eq:f15v88:native}.
At this finite regulator the kernel is continuous,
strictly positive, and positive semidefinite as an
operator. Its unique positive Perron vector is gauge
invariant, because the operator commutes with gauge
transformations. Thus its full norm is also its
physical leading eigenvalue, as in V88--V90.

Identify \(\SU(2)\) with the unit quaternions.
For a unit imaginary quaternion \(X\), \(X^2=-1\),
let \(D_{e,X}\) generate the unitary group
\[
 f(U)\longmapsto f(U_1,\ldots,e^{tX}U_e,\ldots),
 \qquad J_{e,X}=-iD_{e,X}.
\]
The self-adjoint \(J_{e,X}\) has integral weights
\(-n,-n+2,\ldots,n\) on the label-\(n\) irreducible
left representation. Fix the round unit-\(S^3\)
normalization
\[
 \mathcal C_e=-\Delta_e,\qquad
 \mathcal C_e|_n=n(n+2),\qquad
 \mathcal N_e=\sqrt{1+\mathcal C_e}-1.
\]
Thus \(\mathcal N_e|_n=n\). This normalization is
specified to avoid importing a different factor
of four in the Casimir. The spectral projections
\(\Pi_{e,N}={\bf1}_{[0,N]}(\mathcal N_e)\) commute
with both left and right link translations and
therefore with the combined Gauss projection.
The individual \(J_{e,X}\) need not preserve
the physical subspace; the estimates are first
proved on the full Haar space.

\begin{theorem}[One-link native imaginary-shift bound]
\label{f15v91:shift}
For every real \(\sigma\), the range of \(K\) lies
in the domain of \(e^{\sigma J_{e,X}}\), and
\begin{equation}
 \|e^{\sigma J_{e,X}}K\|_{\rm full}
       \leq \exp[3\gamma(\cosh\sigma-1)].
 \label{eq:f15v91:shift}
\end{equation}
The constant has no spectator-volume dependence.
\end{theorem}

\begin{proof}
Extend the kernel holomorphically by replacing \(U_e\)
with \(e^{-i\sigma X}U_e\) in its output argument.
Inverses in a plaquette word are holomorphic matrix
inverses, not complex conjugates. Since
\(e^{-i\sigma X}=\cosh\sigma-iX\sinh\sigma\),
and quaternionic half-traces of products of real
unit quaternions and \(X\) are real, every affected
oriented plaquette satisfies
\[
 \Re\left(\tfrac12\Tr U_p^{(-i\sigma)}\right)
       =\cosh\sigma\,\tfrac12\Tr U_p .
\]
For an inverse occurrence, the factor is
\(U_e^{-1}e^{i\sigma X}\); this changes the sign
of the imaginary term, not the displayed real part.
The same identity holds for the single kinetic
half-trace \(\frac12\Tr(U_eV_e^{-1})=U_e\cdot V_e\).

Every real half-trace is at most one. The four
output magnetic plaquettes each have exponent
weight \(\gamma/2\); the kinetic factor has weight
\(\gamma\). The input multiplier \(b(V)\) and all
spectator factors are unchanged. Hence, if
\(\mathsf k\) denotes the normalized kernel of \(K\),
\begin{equation}
 |\mathsf k(U^{(-i\sigma)},V)|
 \leq \exp[3\gamma(\cosh\sigma-1)]\mathsf k(U,V).
 \label{eq:f15v91:kernel-shift}
\end{equation}
The scalar normalizers are fixed and cancel.
Pointwise domination gives a bounded integral
operator \(A_\sigma\) with
\(\|A_\sigma f\|_2\leq
 e^{3\gamma(\cosh\sigma-1)}\|K|f|\|_2\).
Since \(\|K\|=1\), this proves the asserted bound
for \(A_\sigma\).

It remains to identify the unbounded exponential.
For any \(f\in L^2(\mathcal X)\), compactness of
\(\mathcal X\) and uniform kernel bounds on compact
complex strips make the shifted integral an entire
Hilbert-valued function of the flow parameter \(z\).
On the real axis it is
\(e^{izJ_{e,X}}Kf\), and it is \(2\pi\)-periodic.
Let \(P_m\) be the Fourier projection of this circle
action onto weight \(m\). Contour shifting in its
Fourier integral, or the uniqueness of analytic
continuation after applying \(P_m\), gives
\[
 P_m A_\sigma f=e^{\sigma m}P_mKf .
\]
Parseval applied to the bounded vector \(A_\sigma f\)
shows that
\(\sum_m e^{2\sigma m}\|P_mKf\|_2^2<\infty\).
This is precisely the spectral domain condition
for \(e^{\sigma J_{e,X}}\), with value \(A_\sigma f\).
It proves \eqref{eq:f15v91:shift} including its
domain assertion. No probability measure on the
complexified group has been introduced.
\end{proof}

\subsection{Averaging axes and electric spectral tails}

For \(\sigma\geq0\), define
\begin{equation}
 h_\sigma(n)=\frac{1}{n+1}
       \sum_{m=-n,-n+2,\ldots,n}e^{2\sigma m}
 =\frac{\sinh(2\sigma(n+1))}
        {(n+1)\sinh(2\sigma)},\qquad h_0(n)=1.
 \label{eq:f15v91:character}
\end{equation}

\begin{proposition}[A continued character as a positive form]
\label{f15v91:character}
On the full space, in quadratic-form notation,
\begin{equation}
 K h_\sigma(\mathcal N_e)K
       \leq e^{6\gamma(\cosh\sigma-1)}I.
 \label{eq:f15v91:character-form}
\end{equation}
Moreover,
\begin{equation}
 h_\sigma(n)\geq\tfrac14 e^{\sigma n},
 \qquad
 h_\sigma(n)\geq1+\frac{2\sigma^2}{3}n(n+2).
 \label{eq:f15v91:character-lower}
\end{equation}
\end{proposition}

\begin{proof}
First insert the finite output label projection
\(\Pi_{e,M}\) into the norm square in
Theorem~\ref{f15v91:shift}. It commutes with every
\(J_{e,X}\), so the same norm bound holds.
Average over \(X=gX_0g^{-1}\) with normalized Haar
measure \(dg\). On the label-\(n\) left representation,
Schur's lemma makes the averaged \(e^{2\sigma J}\)
a scalar, equal to its normalized trace
\(h_\sigma(n)\). Right multiplicities and all other
links carry the identity. Letting \(M\) increase
proves the form bound by monotone convergence.
This finite-cutoff argument avoids an unjustified
average of unbounded operators.

There are \(\lfloor n/4\rfloor+1\) weights
\(m\geq n/2\), at least \((n+1)/4\).
Their contribution gives the first lower bound.
The weights are symmetric and have second moment
\[
 \frac{1}{n+1}\sum_m m^2=\frac{n(n+2)}3 .
\]
Averaging \(\cosh(2\sigma m)\geq1+2\sigma^2m^2\)
gives the second bound.
\end{proof}

For \(R\geq0\), put
\begin{align}
 I_\gamma(R)
 &=R\,\operatorname{arsinh}\!\left(\frac{R}{6\gamma}\right)
   -6\gamma\left(\sqrt{1+\frac{R^2}{36\gamma^2}}-1\right),
 \label{eq:f15v91:rate}\\
 B_\gamma(R)&=\min\left\{\frac{R^2}{16\gamma},
                                  \frac R2\right\}.
 \label{eq:f15v91:rate-simple}
\end{align}

\begin{theorem}[Native electric tails and complete-window energy]
\label{f15v91:tails}
For every integer \(N\geq0\), with \(Q_{e,N}=I-\Pi_{e,N}\),
\begin{equation}
 \|Q_{e,N}K\|^2\leq
        \min\{1,4e^{-I_\gamma(N+1)}\},
 \qquad I_\gamma(R)\geq B_\gamma(R).
 \label{eq:f15v91:tail}
\end{equation}
Let \(0<\alpha\leq1\) and
\(E_\alpha={\bf1}_{[\alpha,1]}(K)\), on either the
full or physical space. Then
\begin{equation}
 \|Q_{e,N}E_\alpha\|^2
       \leq\min\{1,4\alpha^{-2}e^{-I_\gamma(N+1)}\}.
 \label{eq:f15v91:window-tail}
\end{equation}
Every unit vector \(\psi\in\operatorname{Ran}E_\alpha\)
satisfies
\begin{equation}
 \langle e^{\sigma\mathcal N_e}\rangle_\psi
       \leq4\alpha^{-2}e^{6\gamma(\cosh\sigma-1)}.
 \label{eq:f15v91:moment}
\end{equation}
For \(\gamma\geq1\), it also satisfies
\begin{equation}
 \langle\mathcal C_e\rangle_\psi
       \leq6\gamma\bigl(e\,\alpha^{-2}-1\bigr).
 \label{eq:f15v91:casimir}
\end{equation}
These expectation bounds extend to every density
matrix supported in the same window.
\end{theorem}

\begin{proof}
On labels \(n>N\), the first inequality in
\eqref{eq:f15v91:character-lower} bounds
\(h_\sigma(n)\) below by \(e^{\sigma(N+1)}/4\).
Consequently
\[
 \|Q_{e,N}K\|^2
 \leq4\exp\{6\gamma(\cosh\sigma-1)-\sigma(N+1)\}.
\]
Optimize over \(\sigma\geq0\).
The stationary point is
\(\sigma=\operatorname{arsinh}(R/(6\gamma))\);
the exponent's negative infimum is \(I_\gamma(R)\).
The cap by one follows from \(\|K\|=1\).

For completeness, if \(0\leq\sigma\leq1\), then
\[
 \cosh\sigma-1
 \leq \sigma^2\sum_{k\geq1}\frac1{2\cdot12^{k-1}}
 =\frac6{11}\sigma^2\leq\frac23\sigma^2 .
\]
Here \((2k)!\geq2\cdot12^{k-1}\).
If \(R\leq8\gamma\), choose
\(\sigma=R/(8\gamma)\); the variational rate is at
least \(R^2/(16\gamma)\).
If \(R\geq8\gamma\), choose \(\sigma=1\);
it is at least \(R-4\gamma\geq R/2\).
This proves \eqref{eq:f15v91:rate-simple}.

On the window, write \(\psi=Kf\) with
\(f=K^{-1}\psi\) and \(\|f\|\leq\alpha^{-1}\).
Both the tail and the exponential moment follow
from the already proved full-space bounds.
For the Casimir, retain the constant term in
the second inequality of
\eqref{eq:f15v91:character-lower}:
\[
 1+\frac{2\sigma^2}{3}
       \langle\mathcal C_e\rangle_\psi
 \leq \alpha^{-2}e^{6\gamma(\cosh\sigma-1)}.
\]
For \(\gamma\geq1\), take
\(\sigma=1/(2\sqrt\gamma)\); the exponent is at
most one. Rearrangement proves
\eqref{eq:f15v91:casimir}. Finite label truncations
justify the unbounded expectations and their limits.
Decomposing a supported density matrix into unit
vectors in its support, followed by monotone
convergence, gives the final assertion.
\end{proof}

Thus an entire near-top physical spectral subspace,
not just its vacuum or selected eigenvectors,
has native electric regularity of scale \(\gamma\).
This complements V90's local magnetic estimates.
It does not prove a global form inequality between
\(\mathcal C_e\) and \(H_M=-\log K\), nor does it
identify the latter with a diffusion generator.

\subsection{A volume-explicit raw-band outer compression}
\label{f15v91:compression}

Let \(\Pi=\prod_{e\in\mathcal E}\Pi_{e,N_e}\);
the projections commute. It preserves the physical
subspace. Unlike a fibre projection that depends
on neighboring links, it has an explicitly finite
full-space rank.

\begin{proposition}[Complete native finite-rank error]
\label{f15v91:compression-bound}
With \(Q=I-\Pi\),
\begin{align}
 \|QK\|^2&\leq4\sum_{e\in\mathcal E}
                         e^{-I_\gamma(N_e+1)},\\
 \|K-\Pi K\Pi\|&\leq4
       \left(\sum_{e\in\mathcal E}
                         e^{-I_\gamma(N_e+1)}\right)^{1/2}.
 \label{eq:f15v91:compression}
\end{align}
For a common cutoff \(N\), \(0<\eta<1\), set
\(\ell_\eta=\log(16E/\eta^2)\).
It suffices to choose an integer
\begin{equation}
 N\geq\max\{4\sqrt{\gamma\ell_\eta},\,2\ell_\eta\}
 \label{eq:f15v91:cutoff}
\end{equation}
to make the second norm at most \(\eta\).
Its retained rank is at most
\begin{equation}
 d_N^E,\qquad
 d_N=\sum_{n=0}^N(n+1)^2
       =\frac{(N+1)(N+2)(2N+3)}6 .
 \label{eq:f15v91:rank}
\end{equation}
\end{proposition}

\begin{proof}
The commuting-projection union inequality gives
\(Q\leq\sum_e Q_{e,N_e}\). Therefore
\[
 \|QKf\|^2=\langle Kf,QKf\rangle
 \leq\sum_e\|Q_{e,N_e}Kf\|^2.
\]
Apply Theorem~\ref{f15v91:tails}.
The exact identity
\(K-\Pi K\Pi=QK+\Pi KQ\), self-adjointness of \(K\),
and \(\|\Pi\|\leq1\) give
\(\|K-\Pi K\Pi\|\leq2\|QK\|\).
For \eqref{eq:f15v91:cutoff}, both terms in
\(B_\gamma(N+1)\) are at least \(\ell_\eta\).
Thus the sum of tails is at most
\(Ee^{-\ell_\eta}=\eta^2/16\).
The Peter--Weyl dimension sum gives the rank.
Restriction to the invariant physical subspace
cannot increase any of these norms or ranks.
\end{proof}

For fixed \(E\) and tolerance, this sufficient cutoff
is \(O(\sqrt\gamma)\), without an additional
\(\sqrt{\log\gamma}\) factor from a dimension-weighted
trace-tail estimate. The volume enters explicitly
through \(\ell_\eta\), and the retained dimension
is still very large. There is no assertion of a
fixed cutoff in a continuum limit: V84's escape
from fixed raw representation bands is consistent
with this growing cutoff.

The matrix to retain is precisely
\(\Pi M_b C_\gamma^{\otimes E}M_b\Pi/\Lambda\).
One must not insert further projections between
its magnetic and kinetic factors or replace its
magnetic exponential by the exponential of a
compressed action. Nor does this proof claim the
false operator-order statement \(\Pi K\Pi\leq K\).
After this outer compression, a unitary recoupling
of all retained representation and boundary
multiplicity indices changes neither its spectrum
nor its certified error. This is a well-defined
place to test the fusion literature without
discarding the mixed plaquette matrix elements.

The earlier min--max handoff
\eqref{eq:f15v88:ratio} applies with error \(\eta\).
If \(r_N\) is the retained physical top-two ratio,
it gives \(r_{\rm full}\leq\min\{1,r_N+\eta\}\).
The difficult estimate is still a bound on \(r_N\)
for the coupled neutral collective operator,
with \(r_N+\eta\leq e^{-m_*a_t}\) on an independently
specified physical time step. The present theorem
does not supply that bound or the clock.

\subsection{Testing the global Hilbert-projective route}
\label{f15v91:projective}

For strictly positive continuous functions on
\(\mathcal X\), or on its gauge quotient, define
\[
 d_H(f,g)=\log\sup_U\frac{f(U)}{g(U)}
                -\log\inf_U\frac{f(U)}{g(U)}.
\]
For a positive operator \(A\), let
\(\Delta(A)=\sup_{f,g>0}d_H(Af,Ag)\).
The classical theorem cited above states
\[
 \kappa(A)=\tanh(\Delta(A)/4),
\]
where \(\kappa(A)\) is the optimal Hilbert-metric
Lipschitz constant, with the usual infinite-diameter
convention. We distinguish this cone contraction
coefficient from an actual Hilbert-space eigenvalue
ratio.

\begin{proposition}[Endpoint multiplication does not change diameter]
\label{f15v91:diameter}
At each finite regulator, on both the full and
gauge-invariant positive cones,
\begin{equation}
 \Delta(M_b C_\gamma^{\otimes E}M_b)
           =\Delta(C_\gamma^{\otimes E}).
 \label{eq:f15v91:diameter-cancel}
\end{equation}
On the full cone the common value is \(4\gamma E\).
\end{proposition}

\begin{proof}
Multiplication by the strictly positive continuous
\(b\) is an onto cone automorphism, since \(b^{-1}\)
is bounded at finite volume. It is an isometry
for \(d_H\), and \(b\) is gauge invariant.
The output multiplier cancels from the metric;
the input multiplier merely permutes the pairs
over which its supremum is taken.
Scalar normalization by \(\Lambda\) also cancels.

For a strictly positive continuous kernel \(k\)
on a compact space with full-support measure,
\[
 \Delta(A)=\sup_{U,U',V,V'}
 \log\frac{k(U,V)k(U',V')}{k(U,V')k(U',V)}.
\]
Indeed a bound on every such cross ratio integrates
against \(f(V)g(V')\,dV\,dV'\) to bound the output
ratio for arbitrary positive \(f,g\). Conversely,
continuous nonnegative bumps concentrating near
\(V,V'\), with a vanishing positive constant added,
recover every cross ratio by continuity.

For the product kinetic kernel its logarithmic
cross ratio is
\[
 \gamma\sum_{e\in\mathcal E}
        (U_e-U'_e)\cdot(V_e-V'_e).
\]
Every term is at most \(4\gamma\) by Cauchy--Schwarz.
Take opposite unit quaternions at both pairs of
endpoints to attain \(4\gamma E\).
This last full-cone argument does not itself
provide a lower bound on the smaller physical cone.
\end{proof}

In particular the magnetic endpoint oscillation
suggested in the memorandum is not the cause of
the global diameter obstruction: the endpoint
weights cancel exactly. The next proposition
checks that Gauss restriction does not remove
the volume obstruction.

\begin{proposition}[A physical many-loop diameter obstruction]
\label{f15v91:physical-diameter}
Let \(N_0\) elementary spatial plaquettes have
pairwise disjoint edge sets. Set
\[
 r=\left(\frac{I_2(\gamma)}{I_1(\gamma)}\right)^4
       \in(0,1).
\]
Then the native physical-cone diameter satisfies
\begin{equation}
 \Delta_{\rm phys}(\bar T)
        \geq4N_0\operatorname{artanh}r,\qquad
 \kappa_{\rm phys}(\bar T)
        \geq\tanh\bigl(N_0\operatorname{artanh}r\bigr).
 \label{eq:f15v91:physical-diameter}
\end{equation}
On the even spatial torus one can take \(N_0=L^3/4\).
Thus this global-diameter certificate tends to one
as \(L\to\infty\), at every fixed \(\gamma>0\).
\end{proposition}

\begin{proof}
For each selected plaquette let
\(h_i(U)=\frac12\Tr U_{p_i}\in[-1,1]\).
It is gauge invariant and belongs to label \(n=1\)
on each of its four links. The exact kinetic
eigenvalue on label one is \(I_2(\gamma)/I_1(\gamma)\);
hence \(C_\gamma^{\otimes E}h_i=rh_i\).
Disjoint edge supports imply, for \(0<s<1\),
\[
 f_\pm=\prod_{i=1}^{N_0}(1\pm s h_i)>0,\qquad
 C_\gamma^{\otimes E}f_\pm
             =\prod_{i=1}^{N_0}(1\pm srh_i).
\]
Because the loops are edge-disjoint, all their
half-traces can be simultaneously \(+1\), or
simultaneously \(-1\): assign a central link sign
on one distinct edge of each loop and take the
other links to be the identity.
The ratio of these two outputs therefore ranges
from \(((1-sr)/(1+sr))^{N_0}\) to its reciprocal.
Its projective distance is
\(4N_0\operatorname{artanh}(sr)\).
Let \(s\) increase to one and use
\eqref{eq:f15v91:diameter-cancel}.

For the geometry, choose plaquettes in directions
1 and 2 based at \((2i,2j,z)\), with
\(0\leq i,j<L/2\) and \(0\leq z<L\).
Their edges are disjoint, including across periodic
boundaries, and their number is \(L^3/4\).
The formula for \(\kappa\) completes the proof.
\end{proof}

For clarity about what fails, a positive continuous
kernel's secondary eigenvalue ratio is bounded
\emph{above} by its projective contraction coefficient.
In this setting it can also be seen directly:
if \(v\) is a real secondary eigenfunction,
the strictly positive continuous Perron vector
\(\Omega\) permits the positive perturbation
\(\Omega+\varepsilon v\). Expanding the projective
distance to first order gives the oscillation
seminorm of \(v/\Omega\); after applying the
operator the seminorm is multiplied by the
absolute eigenvalue ratio. The definition of
\(\kappa\) gives the stated upper bound.

A lower bound tending to one on this \emph{upper
certificate} does not give a lower bound on the
actual spectral ratio. It only shows that the
optimal global Hilbert-metric coefficient cannot
certify a volume-uniform positive physical mass
at a time step independent of volume. It does
not show that the actual gap closes.
A genuinely integrated finite block, an adapted
invariant cone, or a different comparison metric
is not ruled out. A fusion-basis matrix is not
automatically entrywise positive; a cone argument
there requires a specified proper preserved cone.

\subsection{Targeted audit of the proposed RG imports}
\label{f15v91:rg-audit}

The following counterexamples are not native
Wilson counterexamples and are not an exhaustive
assessment of the cited series. They isolate
logical implications used in the identified
main-paper arguments. Consequently the claimed
complete mass-gap theorem is not adopted as an
input to these notes.

\begin{proposition}[A mixed-cell mask is not a positive compression]
\label{f15v91:gram-audit}
Positivity of a Gram matrix does not imply positivity
of the matrix retaining only entries joining
different cells.
\end{proposition}

\begin{proof}
For \(0<r<1\),
\[
 G=\begin{pmatrix}1&r\\r&1\end{pmatrix}>0.
\]
Assign one index to each of two cells. The
different-cell mask is
\[
 D=\begin{pmatrix}0&r\\r&0\end{pmatrix},
\]
whose eigenvalues are \(r,-r\). In particular
\(\langle(1,-1),D(1,-1)\rangle=-2r\).
More generally a positive semidefinite matrix
with all diagonal entries zero is zero, since
each \(2\)-by-\(2\) principal minor forces its
off-diagonal entry to vanish.
\end{proof}

In the main paper's Lemma~5.3, printed pp.~50--51,
the passage from the full Gram kernel to its
different-cell pair selection is described as
a positive principal block. Selecting pairs of
indices by their relation is not selecting the
same subset of rows and columns: the latter
retains its diagonal. The general Gram argument
therefore does not prove positivity of the
mixed-cell operator. To use such an operator
here requires an actual positive factorization
for the specific native expression, or a bound
on its complete signed operator error.
The counterexample does not assert that every
particular realization of that expression is
indefinite; it disproves the stated automatic
positivity inference.

\begin{proposition}[The correct physical-clock loss]
\label{f15v91:clock-loss}
Suppose positive contractions \(T,T'\) have a common
vacuum under a vacuum-preserving isometric embedding
\(V\), and on the relevant spaces
\[
 T'=V^*T^bV-D+F,\qquad D\geq0,\quad\|F\|\leq\epsilon.
\]
Let \(r,r'\) be their norms on vacuum orthogonal
complements, and let the physical time steps obey
\(a'=ba>0\). If \(0<r,r'\leq1\), then
\begin{equation}
 m'=-\frac{\log r'}{ba}
 \geq m-\frac1{ba}\log(1+\epsilon e^{bam}),
 \qquad m=-\frac{\log r}{a}.
 \label{eq:f15v91:clock-loss}
\end{equation}
\end{proposition}

\begin{proof}
The embedding sends the coarse vacuum orthogonal
complement into the fine one. For a unit vector
in the former, the quadratic form of \(V^*T^bV\)
is at most \(r^b\), and that of \(-D\) is
nonpositive. Since \(T'\) is positive, its norm
on that complement is its supremum Rayleigh
quotient. Thus \(r'\leq r^b+\epsilon\).
Taking logarithms and using \(r^b=e^{-bam}\)
gives \eqref{eq:f15v91:clock-loss}.
This is a quadratic-form argument, not the
generally false vector-norm assertion
\(\|(A-D)f\|\leq\|Af\|\) for \(D\geq0\).
\end{proof}

Theorem~5.4 of the cited main paper changes between
a dimensionless gap and \(-\log\lambda_1/a_k\).
Its printed p.~53 argument does not establish a
scale-independent physical loss from the displayed
absolute transfer error: \(a_k=b^ka_0\) is not
uniformly comparable to one over unbounded \(k\)
at fixed \(a_0>0\), and nonnegativity of a gap
alone is not an upper bound of one on that gap.
The following example makes the missing clock
control explicit, even with identical carriers
and exactly preserved vacua.

\begin{proposition}[Summable transfer errors can coexist with vanishing mass]
\label{f15v91:clock-counterexample}
There are strictly positive two-state kernels,
positive semidefinite normalized transfers \(T_k\),
a common vacuum, and errors of summable norm such that
\[
 T_{k+1}=T_k^2+F_k,\qquad
 \sum_{k\geq0}\|F_k\|<1-\lambda_0,
\]
but \(-\log\lambda_k/a_k\to0\) when \(a_k=2^k\).
Here \(\lambda_k\) denotes the secondary eigenvalue,
not the Perron eigenvalue.
\end{proposition}

\begin{proof}
Let
\[
 P=\tfrac12\begin{pmatrix}1&1\\1&1\end{pmatrix},
 \quad Q=I-P,\quad
 T_k=P+\lambda_kQ,\quad
 \lambda_0=\tfrac14,\quad
 \lambda_{k+1}=\lambda_k^2+2^{-k-4}.
\]
Inductively \(0<\lambda_k\leq1/2\), since
\(\lambda_k^2+2^{-k-4}\leq1/4+1/16<1/2\).
Thus each matrix has strictly positive entries,
spectrum \(\{1,\lambda_k\}\), and the same
normalized vacuum \((1,1)/\sqrt2\).
The identity holds with
\(F_k=2^{-k-4}Q\), and
\(\sum_k\|F_k\|=1/8<3/4=1-\lambda_0\).
For \(k\geq1\),
\(\lambda_k\geq2^{-k-3}\). Consequently
\[
 0\leq-\frac{\log\lambda_k}{2^k}
      \leq\frac{(k+3)\log2}{2^k}\longrightarrow0.
\]
All claims follow. This example does not model
YM; it disproves an abstract implication from
summable absolute transfer errors to a uniform
physical mass.
\end{proof}

Theorem~5.6's use of a telescoping gap-loss bound
also needs a strictly positive final margin,
not just convergence of the loss series.
For example \(\delta_k=2^{-k}\) and
\(\epsilon_k=2^{-k-1}\) satisfy
\(\delta_{k+1}=\delta_k-\epsilon_k\) and
\(\sum_k\epsilon_k=\delta_0=1\), yet
\(\inf_k\delta_k=0\).
If the strict inequality is imposed as an
additional condition, it must be verified on the
actual trajectory. This observation does not
disprove a theorem conditional on that condition.

Finally, the main paper's Theorem~10.4 uses a
remainder estimate of the form
\begin{equation}
 g'=g-\beta_0g^3+R,\qquad
 |R|\leq c_2g^5+c_3g\|K_{\rm pol}\|
                         +c_3\|K_{\rm pol}\|^2
 \label{eq:f15v91:rg-remainder}
\end{equation}
to infer cubic contraction on a small fixed
rectangular polydisc. Here \(K_{\rm pol}\) denotes
a polymer variable, not our normalized transfer.

\begin{proposition}[A rectangular remainder bound is insufficient]
\label{f15v91:polydisc}
Estimate \eqref{eq:f15v91:rg-remainder} alone does
not force \(g'<g\) for all sufficiently small
positive \(g\) and small independent
\(\|K_{\rm pol}\|\).
A sufficient size condition for the error part
is instead a suitably narrow scaled tube
\(\|K_{\rm pol}\|\leq\kappa g^2\), with additional
coefficient control.
\end{proposition}

\begin{proof}
Set \(\beta_0=c_2=c_3=1\),
\(\|K_{\rm pol}\|=g>0\), and \(R=2g^2\).
The remainder bound holds, but
\(g'=g-g^3+2g^2>g\) for \(0<g<2\).
This choice lies in every positive fixed small
rectangle for sufficiently small \(g\).
It is a counterexample to the inference from
the estimate, not an identification of the
actual RG remainder.

In the proposed tube, the upper error divided
by \(g^3\) is at most
\[
 c_2g^2+c_3\kappa+c_3\kappa^2g .
\]
If this is at most \(\beta_0/2\), the one-step
bound \(g'\leq g-(\beta_0/2)g^3\) follows.
Continuing the argument requires proving
invariance of the tube under the full coupled
RG map, as well as all its carrier, positivity
and physical-clock hypotheses. None is supplied
by choosing a fixed polydisc.
\end{proof}

\subsection{Direction after the import tests and verification record}

The new upstream gain is electric compactness
at a controlled growing representation scale
for the literal transfer. It gives a complete
finite raw-band carrier on which fusion
recoupling and exact mixed-plaquette charge
transitions can be investigated.
We do not repeat V81's exterior-square identity,
V82's Gauss decomposition, or V88's min--max
handoff as new results.

The global-cone route has now failed a precise
volume-uniform certificate test. A block-cone
experiment remains eligible only after its
actual integration and invariant cone are
specified. Cheeger/comparison methods remain
eligible only with a quantitative same-law
flow or quadratic-form estimate. Generic
positivity, local regularity, or summable
unnormalized RG errors cannot stand in for
either estimate. The next unresolved target
is the retained neutral collective spectral
ratio on the physical clock.

The checker
\begin{center}\small
\path{scripts/verify_ym_f15_v91_complex_shift_and_projective_audit.py}
\end{center}
passed all 4,849 new exact-rational fixtures and
the inherited V90 chain. Its new checks cover
holomorphic quaternion inverses and both link
orientations, the four magnetic half-weights,
complete representation-weight averages,
the rate and cutoff inequalities, two-sided
compression, projective cross-ratio invariance,
edge-disjoint torus geometry, and the explicit
Gram, clock and RG-remainder counterexamples.
Finite checks do not replace the analytic
domain, all-volume or infinite-sequence proofs.
The script has 10,177 bytes and SHA--256
\begin{center}\scriptsize\ttfamily
8C403833C8BAB1E5D9F410EF45AE97C72C1B4605B4AE21C227E0DC238DE2A536.
\end{center}

The supplied memorandum has SHA--256
\begin{center}\scriptsize\ttfamily
1995031F017BCE1DCDF350B658615EB66487F7E83EDB926F39265CED9E496D15.
\end{center}
The read-only Lemmens--Nussbaum PDF has SHA--256
\begin{center}\scriptsize\ttfamily
97AF8D0D527E029D38C8A695FA895662A2C8CE68F2A373F0911D0BAAAF66535C.
\end{center}
The read-only Faizal--Shabir combined PDF has SHA--256
\begin{center}\scriptsize\ttfamily
0648F30F67719A6C83976F1DA53131FC6F9A9DDC9BA1C406345F5ED959B32939.
\end{center}
The PDFs and their inspection products are outside
\path{latex_notes}, in
\path{artifacts/ym_f15_v91_literature/}.

The predecessor V90 has 2,037,550 bytes and SHA--256
\begin{center}\scriptsize\ttfamily
43B8F4F6B8E61879D20C3AEC530EB865D624C3A3678AEC203C6D7147F11A7E02.
\end{center}
Its body is preserved byte for byte: remove this
section and restore the title version to recover
V90 exactly. V91 replaces only the active main
source; archives, companions and previous
diagnostics remain unchanged. The most recent
five-version companion checkpoint is V90 and
the next is V95. Only \path{.tex} files belong
in \path{latex_notes}; compilation products go
to \path{artifacts/ym_f15_v91_texcheck/}.

\begin{firewall}
V91 proves native electric-label tails, whole-window
Casimir bounds, a volume-explicit finite-rank
outer compression, and a physical-cone obstruction
to one proposed global projective certificate.
It also isolates specific invalid implications
in proposed literature imports without claiming
an exhaustive audit of those works.
It proves neither a coupled neutral-sector gap,
a physical continuum clock, nor an interacting
four-dimensional continuum construction.
The full Yang--Mills mass gap remains unproved.
\end{firewall}

\section{V92: a coupled neutral two-plaquette transfer and its complete gap}
\label{f15v92:context}

V91 gave a full-space electric cutoff, but did not
bound the retained neutral collective spectrum.
This iteration takes a genuinely coupled neutral
block: two adjacent spatial plaquettes sharing
one link. All seven links fluctuate, all vertices
obey Gauss invariance, and both plaquette weights
and the shared-link kinetic convolution are retained.
There is no prescribed staple, external field,
or independent-loop substitution.

The result is a complete fixed-block theorem:
the native physical spectral ratio has a strict
upper bound independent of the positive coupling.
At large coupling its limit is
\[
 \left(\frac{5-\sqrt{21}}2\right)^2
       =\frac{23-5\sqrt{21}}2<\frac1{20}.
\]
The proof establishes Hilbert--Schmidt convergence
of the whole compact transfer after an exact
isometric change of coordinates, then identifies
the entire limiting Gauss-invariant spectrum.
It is not a Rayleigh calculation on a few states.
The all-coupling bound is qualitative; no explicit
numerical entry coupling or worst-case ratio
over intermediate couplings is claimed.

This is a starting block, not a replacement for
the objective. The seven-link graph is not the
three-dimensional spatial torus, and its fully
neutral Gauss space is not an arbitrary charged
regional boundary space. A comparison with the
coupled many-block transfer is still missing.
The full interacting four-dimensional continuum
Yang--Mills mass gap remains unproved.

There are existing ingredients not to recount
as new. V82 gives general boundary selection
rules; V74 uses the absence of a linear colour
singlet in an invariant Gaussian sector; V86
proves a full-compact single-link Gaussian limit
with fixed endpoint staples and gives the
Mehler diagonalization. The new work here is
the exact shared-link reduction, actual fusion
amplitudes and an omitted triplet, followed by
a whole-operator limit and complete physical
gap for this coupled, unfrozen two-loop block.
A search of the available YM notes found no
earlier proof of this specified block theorem.
This is a local nonduplication check, not a
claim of literature novelty.

\subsection{Exact gauge reduction of the seven-link block}
\label{f15v92:reduction}

Let the graph consist of three internally disjoint
paths from \(v_0\) to \(v_1\). Paths \(R_1,R_2\)
have length three, and \(R_0\) is their common
one-link chord. They form two adjacent elementary
squares. Orient the paths from \(v_0\) to \(v_1\),
reversing link variables when necessary, and
denote their ordered holonomies by \(H_i\).
Set
\[
 A=H_0^{-1}H_1,\qquad B=H_0^{-1}H_2,\qquad
 h(A)=\tfrac12\Tr A,\qquad s(A)=1-h(A).
\]
The two Wilson plaquette actions are \(s(A)\)
and \(s(B)\). Taking the inverse orientation
of either square leaves them unchanged.
On the seven-link Haar space the native
normalized transfer is
\begin{equation}
 \bar T_{\Box,\gamma}
   =M_b C_\gamma^{\otimes7}M_b,\qquad
 b=\exp[-\gamma(s(A)+s(B))/2].
 \label{eq:f15v92:seven-link}
\end{equation}
The kinetic normalization is exactly
\eqref{eq:f15v88:native}, a known scalar
normalization, not division by a guessed vacuum
eigenvalue. Write
\[
 k_\gamma(g)=\frac{e^{-\gamma(1-h(g))}}{c_0(\gamma)},
 \qquad
 (C_\gamma f)(A)=\int k_\gamma(g)f(g^{-1}A)\,dH(g).
\]
At \(\gamma=0\), use \(k_0=1\), \(b=1\).

\begin{proposition}[Exact neutral carrier and shared-link convolution]
\label{f15v92:exact-reduction}
The physical seven-link space is isometric to
\begin{equation}
 \mathcal H_\Box
     =L^2(G^2,dH(A)dH(B))^{\operatorname{Ad}G},
 \qquad G=\SU(2),
 \label{eq:f15v92:carrier}
\end{equation}
where \(g\) conjugates \(A,B\) simultaneously.
On this space the native transfer becomes
\begin{equation}
 T_{\Box,\gamma}=M_b\mathcal C_\gamma M_b,\qquad
 \mathcal C_\gamma=(C_\gamma^3\otimes C_\gamma^3)D_\gamma,
 \label{eq:f15v92:reduced-transfer}
\end{equation}
where
\[
 (D_\gamma f)(A,B)
     =\int k_\gamma(g)f(g^{-1}A,g^{-1}B)\,dH(g).
\]
These formulas also define a positive self-adjoint
compact extension with strictly positive kernel
on full \(L^2(G^2)\). Its simple Perron vector
belongs to \(\mathcal H_\Box\), so its full and
physical top eigenvalues agree.
\end{proposition}

\begin{proof}
Gauss invariance at the degree-two interior
vertices makes a function depend on each path
only through \(H_i\). Product probability Haar
measure on the links pushes forward to
probability Haar measure on each \(H_i\);
thus this reduction is isometric. At the two
remaining vertices the gauge action is common
left and common right translation of all
three holonomies. Changing variables from
\((H_0,H_1,H_2)\) to \((H_0,A,B)\) preserves
product Haar measure. Common left invariance
removes \(H_0\); common right invariance becomes
simultaneous conjugation of \(A,B\). Integration
of the removed Haar variable contributes one.

Central convolution on an individual path link
acts as that same central convolution on its
path holonomy: intervening left and right
translations can be moved through a central
kernel. The three links on \(R_1,R_2\) therefore
give \(C_\gamma^3\) in the respective holonomies.
For the common link, evaluate the invariant
function in the representative \(H_0=I\).
Changing \(H_0\) by \(g\) changes both \(A,B\)
by common left multiplication by \(g^{-1}\);
this is \(D_\gamma\). The two individual
central convolutions commute with simultaneous
left translation, hence with \(D_\gamma\).
The magnetic multiplier is already a function
of \(A,B\). This proves the exact reduction.

The Fourier coefficient of \(k_\gamma\) on
label \(n\) is
\(r_n(\gamma)=I_{n+1}(\gamma)/I_1(\gamma)>0\).
Consequently \(C_\gamma,D_\gamma\) are positive
operators; for the latter decompose simultaneous
left translation into irreducible representations.
Their commuting product is positive.
Its kernel is
\begin{equation}
 \mathsf c_\gamma(A,B;A',B')
  =\int k_\gamma(g)
       k_\gamma^{*3}(g^{-1}AA'^{-1})
       k_\gamma^{*3}(g^{-1}BB'^{-1})\,dH(g).
 \label{eq:f15v92:kernel}
\end{equation}
It is continuous and strictly positive.
The magnetic sandwich has the same properties
at finite coupling and commutes with simultaneous
conjugation. Compactness and strict kernel
positivity give a unique positive Perron vector,
which is invariant by uniqueness.
At \(\gamma=0\) the operator is the projection
onto constants, with the same Perron conclusion.
\end{proof}

The complete physical spin-network basis is
indexed by triples \((n_1,n_2,n_0)\) satisfying
\[
 |n_1-n_2|\leq n_0\leq n_1+n_2,\qquad
 n_1+n_2+n_0\ \hbox{even}.
\]
The trivalent invariant tensor has multiplicity
one at each endpoint, by the Clebsch--Gordan
rule already used in V82. There is one physical
basis vector for each admissible triple.
Its kinetic eigenvalue is exactly
\begin{equation}
 r_{n_1}(\gamma)^3r_{n_2}(\gamma)^3r_{n_0}(\gamma).
 \label{eq:f15v92:kinetic-labels}
\end{equation}
The common-link label records the fusion of the
two path carriers; it cannot be omitted.
Neither the kinetic operator nor the full
transfer is a product of two independent
plaquette transfers.

\subsection{Actual fusion amplitudes and an omitted triplet}
\label{f15v92:fusion}

Let \(\Pi_1\) retain labels \(0,1\) on every
original link and impose all Gauss constraints.
Define
\[
 \phi_0=1,\quad \phi_p=\Tr A,\quad
 \phi_q=\Tr B,\quad \phi_o=\Tr(AB^{-1}).
\]

\begin{proposition}[The complete first raw band]
\label{f15v92:first-band}
These functions are an orthonormal basis of
\(\operatorname{Ran}\Pi_1\). In this order,
the compressed magnetic multipliers are
\begin{equation}
 A_p=\Pi_1M_{h(A)}\Pi_1
   =\begin{pmatrix}
 0&1/2&0&0\\1/2&0&0&0\\0&0&0&1/4\\0&0&1/4&0
 \end{pmatrix},
 \quad
 A_q=\Pi_1M_{h(B)}\Pi_1
   =\begin{pmatrix}
 0&0&1/2&0\\0&0&0&1/4\\1/2&0&0&0\\0&1/4&0&0
 \end{pmatrix}.
 \label{eq:f15v92:plaquette-matrices}
\end{equation}
Their commutator is
\begin{equation}
 [A_p,A_q]=\frac3{16}
   \bigl(|\phi_p\rangle\langle\phi_q|
                -|\phi_q\rangle\langle\phi_p|\bigr),
 \qquad \|[A_p,A_q]\|=\frac3{16}.
 \label{eq:f15v92:commutator}
\end{equation}
\end{proposition}

\begin{proof}
The central Gauss transformation at a vertex
forces an even number of incident label-one
edges. On this degree-at-most-three graph,
each occupied vertex therefore has degree two.
The binary cycle space has dimension two:
its four elements are the empty graph, the
two squares, and their outer six-edge perimeter.
The fundamental singlet at a degree-two vertex
has multiplicity one. These are precisely
the four vectors above.

Write \(A=(a_0,\boldsymbol a)\),
\(B=(b_0,\boldsymbol b)\) as unit quaternions.
The reduced Haar variables are independent
uniform points of \(S^3\), and
\[
 \phi_p=2a_0,\quad \phi_q=2b_0,\quad
 \phi_o=2(a_0b_0+\boldsymbol a\cdot\boldsymbol b).
\]
Their Haar first moments vanish and
\(\mathbb E a_\mu a_\nu=\delta_{\mu\nu}/4\).
These identities prove orthonormality and
give the two nonzero independent amplitudes
\[
 \langle\phi_0,h(A)\phi_p\rangle=\frac12,\qquad
 \langle\phi_q,h(A)\phi_o\rangle=\frac14.
\]
The other entries follow by interchange of
\(A,B\), symmetry, or a vanishing centered
Haar coordinate. Matrix multiplication
then gives \eqref{eq:f15v92:commutator}.
\end{proof}

The original multiplication operators commute:
their compressed commutator is caused by a
missing intermediate channel. In fact put
\begin{equation}
 \tau=\phi_p\phi_q-\frac12\phi_o,\qquad
 \phi_t=\frac2{\sqrt3}\tau .
 \label{eq:f15v92:triplet}
\end{equation}
Haar integration gives \(\|\tau\|^2=3/4\)
and orthogonality to all four retained vectors.
The individual left Casimirs in \(A,B\) are
\(3,3\), but the simultaneous-left Casimir is
\(8\). Indeed \(\phi_o\) is invariant under
common left multiplication, and the remaining
summand of \(V_1\otimes V_1=V_0\oplus V_2\)
has label two. Direct quaternion differentiation
also gives
\[
 -\sum_{i=1}^3(L_i^A+L_i^B)^2\tau=8\tau.
\]
Thus \(\phi_t\) is the normalized physical
state with labels \((1,1,2)\), and
\begin{equation}
 h(A)\phi_q=\frac14\phi_o+\frac{\sqrt3}{4}\phi_t,
 \qquad
 \|(I-\Pi_1)M_{h(A)}\phi_q\|^2=\frac3{16}.
 \label{eq:f15v92:triplet-leak}
\end{equation}
It is globally a physical singlet although its
shared link carries a triplet. Neutrality does
not make all internal representations trivial.

Inserting \(I=\Pi_1+(I-\Pi_1)\) between two
commuting multipliers gives exactly
\[
 [\Pi_1M_p\Pi_1,\Pi_1M_q\Pi_1]
   =-\Pi_1M_p(I-\Pi_1)M_q\Pi_1
     +\Pi_1M_q(I-\Pi_1)M_p\Pi_1.
\]
The omitted channels cancel
\eqref{eq:f15v92:commutator}. Also
\[
 \Pi_1M_p^2\Pi_1-(\Pi_1M_p\Pi_1)^2
          =\Pi_1M_p(I-\Pi_1)M_p\Pi_1\geq0,
\]
whose \(\phi_q\) diagonal entry is \(3/16\)
for \(M_p=M_{h(A)}\).
This explicitly prevents replacing the
outer-compressed magnetic exponential by an
exponential of compressed plaquette matrices.
The gap proof below uses the full transfer,
not this four-dimensional approximation.

\subsection{Fixed-step convolution on the full compact chart}
\label{f15v92:convolution}

Put \(j_0=1/(2\pi^2)\),
\(q_\gamma(x)=\operatorname{Exp}(x/\sqrt\gamma)\),
and \(\mathbb B_\gamma=\{|x|<\pi\sqrt\gamma\}\).
The complete chart covers the group except
the Haar-null antipode and has density
\[
 J_\gamma(x)=j_0\gamma^{-3/2}
       \left(\frac{\sin(|x|/\sqrt\gamma)}
                         {|x|/\sqrt\gamma}\right)^2.
\]
This is \eqref{eq:f15v86:jacobian}.
Write
\(\varphi_m(z)=(2\pi m)^{-3/2}e^{-|z|^2/(2m)}\).

\begin{proposition}[Native fixed-step convolution limit]
\label{f15v92:convolution-limit}
For each fixed integer \(m\geq1\),
\begin{equation}
 \gamma^{-3/2}k_\gamma^{*m}(q_\gamma(z))
       \longrightarrow j_0^{-1}\varphi_m(z)
 \label{eq:f15v92:convolution-limit}
\end{equation}
uniformly on compact \(z\)-sets. For \(\gamma\geq1\),
\[
 \|k_\gamma^{*m}\|_\infty\leq C\gamma^{3/2},
 \qquad
 k_\gamma(q_\gamma(w))J_\gamma(w)\leq Ce^{-c|w|^2}
                  \quad(w\in\mathbb B_\gamma).
\]
\end{proposition}

\begin{proof}
The elementary Laplace normalization proved
in V86 is
\[
 c_0(\gamma)=j_0(2\pi)^{3/2}\gamma^{-3/2}
                                      (1+O(\gamma^{-1})).
\]
It follows from the Gamma integral there,
or from this chart and dominated convergence.
The bound \(1-\cos t\geq2t^2/\pi^2\) for
\(0\leq t\leq\pi\) supplies domination;
a small chart ball gives a positive uniform
lower bound for \(\gamma^{3/2}c_0(\gamma)\).
This proves the displayed bounds for \(m=1\)
and, by Taylor expansion, the locally uniform
limit. Convolution by a probability density
does not increase the supremum norm, giving
that bound for all \(m\).

For induction use
\[
 \begin{split}
 &\gamma^{-3/2}k_\gamma^{*(m+1)}(q_\gamma(z))\\
 &\quad=\int_{\mathbb B_\gamma}
 k_\gamma(q_\gamma(w))J_\gamma(w)
 \bigl[\gamma^{-3/2}k_\gamma^{*m}
                  (q_\gamma(w)^{-1}q_\gamma(z))\bigr]\,dw.
 \end{split}
\]
On compact \(w,z\)-sets the smooth local
group product gives
\(q_\gamma(w)^{-1}q_\gamma(z)
       =q_\gamma(z-w+o(1))\), uniformly.
The inductive limit applies to the bracket,
which is uniformly bounded. The other factor
has the stated Gaussian bound in \(w\).
Dominated convergence yields
\(j_0^{-1}(\varphi_1*\varphi_m)(z)
       =j_0^{-1}\varphi_{m+1}(z)\).
For uniformity on a compact \(z\)-set,
first discard a large \(w\)-tail with the
common Gaussian bound and then use uniform
convergence on the remaining compact set.
No finite-\(\gamma\) kernel has been
replaced by a heat kernel.
\end{proof}

\subsection{A whole-operator limit retaining the shared covariance}
\label{f15v92:operator-limit}

Use the isometry
\[
 (\mathcal U_\gamma f)(x,y)
   =\sqrt{J_\gamma(x)J_\gamma(y)}
                   f(q_\gamma(x),q_\gamma(y))
\]
on \(\mathbb B_\gamma^2\), extended by zero
to \(L^2(\mathbb R^6)\). It exactly intertwines
simultaneous group conjugation with simultaneous
rotations of the two three-vectors, and embeds
the physical space in
\(L^2(\mathbb R^6)^{\mathrm{SO}(3)}\).
The zero complement does not alter nonzero
eigenvalues.

For \(X=(x,y)\), set
\[
 \mathsf M=\begin{pmatrix}4&1\\1&4\end{pmatrix}\otimes I_3,
 \qquad
 \Phi_{\mathsf M}(Z)
   =\frac{\exp[-Z^{\mathsf T}\mathsf M^{-1}Z/2]}
                  {(2\pi)^3\,15^{3/2}}.
\]
Let \(\mathcal G\) have kernel
\begin{equation}
 \mathcal G(X,X')
   =e^{-|X|^2/4}\Phi_{\mathsf M}(X-X')e^{-|X'|^2/4}.
 \label{eq:f15v92:gaussian}
\end{equation}

\begin{theorem}[Full-compact coupled-block convergence]
\label{f15v92:hs-limit}
For the full-space extension in
Proposition~\ref{f15v92:exact-reduction},
\begin{equation}
 \|\mathcal U_\gamma T_{\Box,\gamma}\mathcal U_\gamma^*
                                     -\mathcal G\|_{\rm HS}
                  \longrightarrow0\quad(\gamma\to\infty).
 \label{eq:f15v92:hs-limit}
\end{equation}
The same convergence holds on the complete
physical invariant space with the indicated
zero complement. No compact large-field
region or physical representation sector
is discarded.
\end{theorem}

\begin{proof}
In \eqref{eq:f15v92:kernel}, put \(g=q_\gamma(w)\),
\(A=q_\gamma(x)\), \(B=q_\gamma(y)\), and
similarly at the primed endpoint.
Proposition~\ref{f15v92:convolution-limit}
with \(m=3\) gives, for each fixed \(X,X'\),
\[
 \begin{split}
 &\gamma^{-3}\mathsf c_\gamma(A,B;A',B')\\
 &\quad\longrightarrow j_0^{-2}\int_{\mathbb R^3}
    \varphi_1(w)\varphi_3(x-x'-w)
                          \varphi_3(y-y'-w)\,dw\\
 &\quad=j_0^{-2}\Phi_{\mathsf M}(X-X').
 \end{split}
\]
The Gaussian bound for the \(w\)-probability
density and the two rescaled supremum bounds
justify the limiting integral. Its covariance
is exact: if \(W,U,V\) are independent centered
three-dimensional Gaussians with covariances
\(I_3,3I_3,3I_3\), then \((W+U,W+V)\) has
covariance \(\mathsf M\). The off-diagonal
\(I_3\) comes from the shared link.

The two output and two input Haar square-root
factors have product asymptotic
\(j_0^2\gamma^{-3}\), and
\[
 b(q_\gamma(x),q_\gamma(y))
       \longrightarrow e^{-(|x|^2+|y|^2)/4}.
\]
This proves pointwise convergence to
\eqref{eq:f15v92:gaussian}.

For the required global bound,
\(\|k_\gamma^{*3}\|_\infty\leq C\gamma^{3/2}\)
and \(\int k_\gamma=1\) imply
\(\mathsf c_\gamma\leq C\gamma^3\).
Since \(J_\gamma\leq j_0\gamma^{-3/2}\),
the rescaled kinetic kernel is uniformly
bounded. Globally throughout the chart,
\[
 b(q_\gamma(x),q_\gamma(y))
           \leq e^{-(|x|^2+|y|^2)/\pi^2}.
\]
Thus the rescaled transfer kernel, including
its zero extension, is bounded by
\(C e^{-c(|X|^2+|X'|^2)}\). The limiting kernel
has a bound of the same form. Squaring the
difference and applying dominated convergence
proves \eqref{eq:f15v92:hs-limit}.
This controls the compact complement rather
than assuming concentration is a spectral
estimate.

Finally \(\mathcal U_\gamma\) intertwines
the group actions, hence their averaged
orthogonal invariant projections. Both
operators preserve the fixed Euclidean
invariant subspace. Compressing their
difference to it cannot increase its
Hilbert--Schmidt norm. This proves the
physical-space conclusion.
\end{proof}

This argument gives convergence, not an
explicit rate or threshold coupling.
Neither is needed for the qualitative
fixed-block theorem, but the limiting
constant is not a numerical certificate
at an arbitrary finite coupling.

\subsection{The complete physical Gaussian spectrum}
\label{f15v92:spectrum}

For \(s>0\), define
\begin{equation}
 q_s=\frac{s+2-\sqrt{s^2+4s}}2\in(0,1),\qquad
 \omega_s=\frac{\sqrt{s+4}}{2\sqrt s}.
 \label{eq:f15v92:roots}
\end{equation}
The orthogonal coordinates
\(z_+=(x+y)/\sqrt2\), \(z_-=(x-y)/\sqrt2\)
diagonalize \(\mathsf M\) into \(5I_3,3I_3\)
and preserve the magnetic quadratic form.

\begin{proposition}[Normal modes and the first allowed singlet]
\label{f15v92:gaussian-spectrum}
The full-space eigenvalues of \(\mathcal G\)
before residual Gauss projection are
\begin{equation}
 \ell_0q_5^{n_+}q_3^{n_-},\qquad
 \ell_0=(q_5q_3)^{3/2},\qquad n_+,n_-\geq0,
 \label{eq:f15v92:full-spectrum}
\end{equation}
with oscillator multiplicities.
The ground state is an invariant Gaussian.
On the entire invariant subspace,
\begin{equation}
 \frac{\lambda_1(\mathcal G|_{\rm inv})}
      {\lambda_0(\mathcal G|_{\rm inv})}
       =q_3^2=\frac{23-5\sqrt{21}}2<\frac1{20}.
 \label{eq:f15v92:singlet-ratio}
\end{equation}
\end{proposition}

\begin{proof}
For one coordinate of variance \(s\), the
kernel is
\[
 g_s(u,v)=(2\pi s)^{-1/2}
      e^{-u^2/4-(u-v)^2/(2s)-v^2/4}.
\]
Put \(a=(s+2)/(4s)\), \(c=1/s\).
The exponent is \(-a(u^2+v^2)+cuv\).
Gaussian integration sends
\(e^{-\omega_sv^2/2}\) to
\(q_s^{1/2}e^{-\omega_su^2/2}\).
The constants obey
\[
 a-\frac{c^2}{4(a+\omega_s/2)}=\frac{\omega_s}{2},
 \quad 2s(a+\omega_s/2)=q_s^{-1},
 \quad \frac{c}{2(a+\omega_s/2)}=q_s.
\]
The Hermite generating function therefore
gives eigenvalues \(q_s^{n+1/2}\).
Equivalently use the standard
\href{https://dlmf.nist.gov/18.18.E28}
{Mehler formula, DLMF 18.18.28},
as in Theorem~\ref{f15v86:mehler}.
Completeness of the Hermite functions gives
all eigenvalues. Three coordinates for each
of \(s=5,3\) yield
\eqref{eq:f15v92:full-spectrum}.

The ground Gaussian \(\Omega_{\mathcal G}\)
depends only on \(|z_+|^2,|z_-|^2\).
The total degree-one space consists of two
copies of the defining rotation representation,
neither of which has an invariant vector.
All nonconstant singlets have total degree
at least two. As \(0<q_5<q_3<1\), their
eigenvalue ratios are at most \(q_3^2\).
This is attained by the nonzero invariant
Hermite vector
\[
 \left(|z_-|^2-\frac3{2\omega_3}\right)\Omega_{\mathcal G},
\]
which is orthogonal to the vacuum.
Thus the result bounds the entire spectrum,
not only a selected Rayleigh quotient.
Finally \(\sqrt{21}>229/50\) proves the
strict upper bound \(q_3^2<1/20\).
\end{proof}

Without residual Gauss projection the first
ratio is \(q_3\), not \(q_3^2\).
The projection does not make every even
oscillator state physical; the proof uses
the absence of linear singlets and exhibits
the quadratic scalar attaining the maximum.

The shared-link coupling also changes the
fixed-block gap. An isolated square has
covariance \(4I_3\); two independent squares
would give first singlet ratio \(q_4^2\).
Here \(q_3^2>q_4^2\), so coupling weakens
that independent-loop separation. It has
not been discarded or assumed beneficial.

\subsection{A uniform gap on this fixed native block}
\label{f15v92:block-gap}

Define
\[
 \Lambda_\Box(\gamma)
  =\lambda_0(T_{\Box,\gamma}|_{\mathcal H_\Box}),
 \qquad
 r_\Box(\gamma)
  =\frac{\lambda_1(T_{\Box,\gamma}|_{\mathcal H_\Box})}
              {\Lambda_\Box(\gamma)}.
\]
The ordered eigenvalues include multiplicity
on the whole physical space.

\begin{theorem}[Coupling-uniform fixed-block neutral transfer gap]
\label{f15v92:uniform-gap}
For this seven-link two-plaquette graph,
\begin{equation}
 \Lambda_\Box(\gamma)\longrightarrow(q_5q_3)^{3/2}>0,
 \qquad r_\Box(\gamma)\longrightarrow q_3^2
                    \quad(\gamma\to\infty).
 \label{eq:f15v92:limits}
\end{equation}
Also \(r_\Box(\gamma)\to0\) as \(\gamma\downarrow0\).
There is a constant \(q_\Box<1\), depending
on this fixed graph but not on \(\gamma\geq0\),
such that
\begin{equation}
 r_\Box(\gamma)\leq q_\Box.
 \label{eq:f15v92:uniform-block}
\end{equation}
Thus the dimensionless physical block
Hamiltonian
\(-\log(T_{\Box,\gamma}/\Lambda_\Box(\gamma))\)
has gap at least \(-\log q_\Box>0\) for
every \(\gamma>0\).
\end{theorem}

\begin{proof}
Theorem~\ref{f15v92:hs-limit} implies operator-norm
convergence on the entire invariant space.
Ordered eigenvalues of compact positive
operators are Lipschitz in operator norm by
min--max; the isometry's zero complement
does not change them.
Proposition~\ref{f15v92:gaussian-spectrum}
therefore gives both limits in
\eqref{eq:f15v92:limits}. The positive
limiting top value justifies the ratio.

On every compact interval of
\(\gamma\geq0\), the original reduced kernel
is continuous in \(\gamma\) and configuration,
and the scalar normalizer stays positive.
The operators on fixed \(L^2(G^2)\) are
therefore Hilbert--Schmidt-norm continuous,
as are their ordered eigenvalues and ratio.
The top eigenvalue does not vanish.
The simple Perron property gives
\(r_\Box(\gamma)<1\) for each finite coupling.
At zero the rank-one projection gives
\(r_\Box(0)=0\).

Choose a sufficiently large finite
\(\gamma_*\) that \(r_\Box(\gamma)<1/10\)
for \(\gamma\geq\gamma_*\), using the
large-coupling limit. On \([0,\gamma_*]\),
continuity and pointwise strictness give
a maximum strictly less than one.
The maximum of that value and \(1/10\)
is a valid \(q_\Box<1\).
This proves existence, not a numerical
evaluation of \(q_\Box\) or \(\gamma_*\).
Spectral calculus gives the logarithmic
gap assertion.
\end{proof}

This proof does not assume that the vacuum
stays in the first raw band. The triplet
in \eqref{eq:f15v92:triplet-leak} and every
higher representation generated by the
magnetic weight remain in the transfer of
Theorem~\ref{f15v92:uniform-gap}.
The four-dimensional matrices above diagnose
coupling structure; they do not supply
the gap theorem by themselves.

\subsection{What this supplies to the global attack}

There is now a native gapped neutral reference
block containing two coupled plaquettes at
all couplings, with an explicit asymptotic
ratio. The shared link changes both fusion
amplitudes and kinetic normal modes.

Embedding these blocks in the spatial torus
introduces extra incident links, boundary
representations and plaquettes joining
blocks. The fully neutral closed-block
Gauss projection cannot be imposed
independently at such boundaries without
changing the physical carrier.
V82 describes the required matched charged
multiplicities and V91 supplies a complete
growing raw cutoff; neither provides a
coercivity estimate for the assembled
interblock operator.

The interblock Wilson multipliers are not
small in operator norm at large coupling.
Thus a gapped fixed reference is not enough
to invoke a stability theorem with unverified
relative-perturbation hypotheses.
The next substantive estimate must be a
boundary-carrying extension or a same-law
block-form comparison paying all couplings.
A tensor product of closed-block gaps is
not that estimate.

No volume-uniform conclusion follows from
\eqref{eq:f15v92:uniform-block}; graph size
was fixed throughout. The collective
comparison, independently calibrated
physical time, continuum tightness and
nontrivial interacting reconstruction
remain open. The original full YM
objective is unchanged.

\subsection{Exact diagnostics and append record}

The checker
\begin{center}\small
\path{scripts/verify_ym_f15_v92_coupled_neutral_block.py}
\end{center}
uses rational Haar moments and quaternionic
differential operators, not floating-point
diagonalization. It enumerates the complete
first raw-band cycle space; computes the
two magnetic matrices, their commutator and
the omitted triplet norm; and checks the
three path Casimirs independently.
It also checks the common-link Gaussian
covariance, its integrated cross term,
Mehler constants, rational isolating
intervals for the roots, and the linear
and quadratic residual-Gauss spaces.

All 454 new exact-rational fixtures and
the inherited V91 chain passed. They do
not replace the full-compact convergence,
complete Hermite expansion, or compact
coupling-interval argument.
The checker has 10,524 bytes and SHA--256
\begin{center}\scriptsize\ttfamily
4B2060FC5EC108A19A37AA134ED2347CBCBD58D5406875DA6812480A47015849.
\end{center}

The predecessor V91 has 2,071,529 bytes
and SHA--256
\begin{center}\scriptsize\ttfamily
0F2839929B040E134AD8BFB0722F0B0F2396261A775A07FF262E6DA4E62F7B1B.
\end{center}
Its body is preserved byte for byte:
remove this section and restore the
title version to recover V91.
V92 replaces only the active main source;
archives, companions, supplied papers
and previous diagnostic scripts remain
unchanged. The last five-version companion
checkpoint is V90 and the next is V95.
Only \path{.tex} files belong in
\path{latex_notes}; compilation products
go to \path{artifacts/ym_f15_v92_texcheck/}.

\begin{firewall}
V92 proves an exact coupled-block reduction,
native fusion amplitudes and a coupling-uniform
gap on the entire physical space of the
specified fixed two-plaquette block.
This is not a volume-uniform gap for the
spatial torus, not a bound under arbitrary
boundary fields, and not a four-dimensional
continuum construction. The full Yang--Mills
mass gap remains unproved.
\end{firewall}

\section{V93: native boundary carriers, open-block scale, and magnetic gluing}
\label{f15v93:context}

V92 proved a coupling-uniform gap on the fully
neutral physical space of a fixed two-plaquette
block. Gluing such blocks requires a larger
regional carrier: Gauss invariance is imposed
at interior vertices, while boundary
representations remain available for matching.
The neutral block theorem cannot simply be
applied to that larger space.

This continuation obtains quantitative bounds
on the missing boundary modes from the literal
native transfer. Open Wilson lines multiplied
by the actual regional Perron vector give
orthogonal charged trial spaces, with complete
matrix-element control by Schur's lemma.
Their energy is bounded by a native temporal
increment estimate. Together with V83's charged
upper bounds and V92's neutral gap, this proves
that the complete open two-plaquette block has
dimensionless gap of order \(1/\gamma\), rather
than a gap bounded below independently of
\(\gamma\).

The trial spaces also yield actual physical
states after a cut is Gauss-glued. We calculate
the exact changed Gram matrix required when
the mixed magnetic multiplier is restored.
For one mixed plaquette this matrix is explicit
and is not a small perturbation of the identity.
Thus the slow cut states are not being declared
slow states of the uncut YM transfer.

The new content is the native variational
lower bound and spectral counting, the
two-sided complete open-block scale, and the
explicit interface metric on these particular
states. V78's ground-state form, V82's singlet
gluing, V83's charged domination and V88's
kinetic source inequality are used as existing
inputs, not presented as new abstract theorems.
The full volume-uniform physical YM mass gap
and interacting continuum construction remain
unproved.

\subsection{A native Perron-pair increment bound}

Let \(R\) be a finite regional graph with
distinct link endpoints and probability
product Haar measure \(dm_R\). Take
\[
 T_R=M_{b_R}C_\gamma^{\otimes |R|}M_{b_R},
 \qquad \Lambda_R=\|T_R\|,\qquad K_R=T_R/\Lambda_R,
\]
where \(b_R>0\) is continuous and invariant
under every regional vertex gauge
transformation. The intended \(b_R\) is the
literal Wilson multiplier of the plaquettes
wholly in \(R\). The estimate below also
allows any such gauge-invariant multiplier.

The full positive continuous kernel has a
unique normalized positive Perron vector
\(\Omega_R\), invariant under every regional
gauge transformation. It is the same Perron
vector after imposing any subset of the
Gauss constraints. Define
\begin{equation}
 \begin{split}
 d\mu_R(U)&=\Omega_R(U)^2dm_R(U),\\
 d\mathcal J_R(U,V)
   &=\Lambda_R^{-1}\Omega_R(U)t_R(U,V)\Omega_R(V)
                                      dm_R(U)dm_R(V).
 \end{split}
 \label{eq:f15v93:pair}
\end{equation}
The pair law is symmetric with both marginals
\(\mu_R\), as in \eqref{eq:f15v78:ground}.
For arbitrary \(f\in L^2(\mu_R;\mathbb C)\),
not necessarily gauge invariant,
\begin{equation}
 \langle\Omega_R f,(I-K_R)\Omega_R f\rangle
       =\frac12\mathbb E_{\mathcal J_R}
                            |f(U)-f(V)|^2 .
 \label{eq:f15v93:jump}
\end{equation}
Expansion of the square and the two marginals
prove this full-space version of V78's form.

Write
\[
 s_e(U,V)=1-U_e\cdot V_e,\qquad
 a_\gamma=1-r_1(\gamma),\qquad
 r_n(\gamma)=I_{n+1}(\gamma)/I_1(\gamma).
\]

\begin{proposition}[Exact one-link Perron increment cost]
\label{f15v93:increment}
For every link \(e\in R\),
\begin{equation}
 \mathbb E_{\mathcal J_R}s_e
       \leq a_\gamma
       \leq\min\left\{1,\frac3{2\gamma}\right\}.
 \label{eq:f15v93:increment}
\end{equation}
There is no spectator-volume or magnetic
oscillation factor in this estimate.
\end{proposition}

\begin{proof}
Hold the magnetic multiplier fixed at the
chosen regulator. V88's one-link kinetic
source inequality
\eqref{eq:f15v88:source-order} gives
\[
 \mathbb E_{\mathcal J_R}e^{\theta\gamma s_e}
 =\Lambda_R^{-1}
       \langle\Omega_R,T_{R,\theta}\Omega_R\rangle
 \leq\frac{c_0((1-\theta)\gamma)}{c_0(\gamma)}
 \leq(1-\theta)^{-3/2},\qquad 0\leq\theta<1 .
\]
At zero both sides equal one. Differentiation
from the right is legitimate because
\(s_e\in[0,2]\) at this finite regulator.
It gives
\[
 \mathbb E_{\mathcal J_R}s_e
       \leq-\frac{c_0'(\gamma)}{c_0(\gamma)}.
\]
Since \(c_0(\gamma)=e^{-\gamma}z(\gamma)\) and
\(z'(\gamma)/z(\gamma)=r_1(\gamma)\) by the
fundamental character integral, this derivative
is \(a_\gamma\). Differentiating the scalar
inequality
\(c_0((1-\theta)\gamma)/c_0(\gamma)
 \leq(1-\theta)^{-3/2}\) at zero gives
\(\gamma a_\gamma\leq3/2\).
Finally \(0<r_1(\gamma)<1\) implies
\(0<a_\gamma<1\). The argument differentiates
a kinetic source with fixed \(b_R\), not
the full coupling-dependent vacuum.
\end{proof}

\subsection{Open Wilson lines as entire charged trial carriers}
\label{f15v93:carriers}

Choose a boundary set \(\partial R\) and
impose Gauss invariance only at its complement.
Let \(P\) be a simple oriented path of length
\(\ell\) joining two distinct vertices
\(v,w\in\partial R\). Let \(H_P(U)\) be its
ordered holonomy, with inverses when the
path opposes a chosen link orientation.
For the label-\(n\) unitary representation,
put
\[
 d_n=n+1,\qquad c_n=n(n+2),\qquad
 \Phi_{n,ab}(U)
      =\sqrt{d_n}\,\Omega_R(U)D_n(H_P(U))_{ab},
 \quad 1\leq a,b\leq d_n .
\]
Here \(c_n\) is the unit-\(S^3\) Casimir,
with the same normalization as V91--V92.

\begin{theorem}[Native Wilson-line carrier bound]
\label{f15v93:carrier}
For each \(n\), the vectors \(\Phi_{n,ab}\)
are orthonormal. Their span \(\mathcal W_{n,P}\)
is one irreducible boundary carrier of
dimension \(d_n^2\), with label \(n/2\) at
\(v,w\) and trivial labels at all other
boundary vertices. Distinct \(n\)'s are
orthogonal. If \(P_{n,P}\) is its projection,
then
\begin{equation}
 P_{n,P}K_RP_{n,P}=t_{n,P}P_{n,P},
 \qquad
 t_{n,P}=\frac1{d_n}
    \mathbb E_{\mathcal J_R}
          \chi_n(H_P(U)^{-1}H_P(V)),
 \label{eq:f15v93:carrier-compression}
\end{equation}
and
\begin{equation}
 0\leq1-t_{n,P}
      \leq\frac{c_n\ell^2}{3}a_\gamma
      \leq\frac{c_n\ell^2}{2\gamma}.
 \label{eq:f15v93:carrier-loss}
\end{equation}
In particular, for a single link and \(n=1\),
\(t_{1,P}\geq r_1(\gamma)\).

The trial space need not be invariant under
\(K_R\). The scalar in
\eqref{eq:f15v93:carrier-compression} is a
compressed Rayleigh value, not an asserted
exact transfer eigenvalue.
\end{theorem}

\begin{proof}
The marginal law of \(H_P\) under \(\mu_R\)
is exactly Haar. Indeed \(\mu_R\) is invariant
under the gauge transformation at \(v\),
which left-translates \(H_P\) arbitrarily.
Uniqueness of probability Haar measure
proves the assertion. Schur orthogonality
then gives
\[
 \int\overline{D_n(H_P)_{ab}}D_m(H_P)_{cd}\,d\mu_R
 =\frac{\delta_{nm}\delta_{ac}\delta_{bd}}{d_n}.
\]
All intermediate gauge transformations cancel
along the path, so only the two endpoint
carriers remain. This gives the stated
irreducible \(G_v\times G_w\) representation
and interior Gauss invariance.

The operator \(K_R\) commutes with the
boundary group. Its compression to this
single irreducible copy is therefore a
scalar by Schur's lemma. Cross matrix
elements between different \(n\)'s vanish.
Take the trace over the \(d_n^2\) orthonormal
vectors. The matrix identity
\[
 \sum_{a,b}\overline{D_n(H_U)_{ab}}D_n(H_V)_{ab}
                      =\chi_n(H_U^{-1}H_V)
\]
gives \eqref{eq:f15v93:carrier-compression}.
Positivity and normalization of \(K_R\)
give \(0\leq t_{n,P}\leq1\).

For \(g\in\SU(2)\) with eigenangles
\(\pm\theta\), the representation weights
are \(-n,-n+2,\ldots,n\). The elementary
inequality
\(\lvert e^{im\theta}-1\rvert
        \leq |m|\lvert e^{i\theta}-1\rvert\)
and their second moment
\(d_n^{-1}\sum_m m^2=c_n/3\) give
\begin{equation}
 1-\frac{\chi_n(g)}{d_n}
       \leq\frac{c_n}{3}(1-\tfrac12\Tr g).
 \label{eq:f15v93:character-chord}
\end{equation}
Equivalently this bounds the normalized
Hilbert--Schmidt square distance of the
two representation matrices.

Quaternion multiplication preserves Euclidean
norm. Telescoping the two ordered path
products and applying Cauchy--Schwarz yields
\[
 |H_P(U)-H_P(V)|^2
        \leq\ell\sum_{e\in P}|U_e-V_e|^2
        =2\ell\sum_{e\in P}s_e(U,V).
\]
The same bound holds for inverse-oriented
links because inversion preserves the
chord distance. Use
\(1-\frac12\Tr(H_U^{-1}H_V)=|H_U-H_V|^2/2\),
then \eqref{eq:f15v93:character-chord} and
Proposition~\ref{f15v93:increment}.
This gives
\[
 1-t_{n,P}\leq
       \frac{c_n}{3}\ell\sum_{e\in P}
                            \mathbb E_{\mathcal J_R}s_e
       \leq\frac{c_n\ell^2}{3}a_\gamma .
\]
For \(n=1,\ell=1\), \(c_1=3\), so the
last inequality reads \(t_{1,P}\geq1-a_\gamma
=r_1(\gamma)\).
\end{proof}

Let \(\rho_{R,n,P}\) denote the norm of the
whole corresponding boundary multiplicity
operator, divided by \(\Lambda_R\).
The preceding finite carrier is a trial
copy within that sector, hence
\begin{equation}
 \rho_{R,n,P}\geq
       \max\left\{0,1-\frac{c_n\ell^2}{3}a_\gamma\right\}.
 \label{eq:f15v93:sector-lower}
\end{equation}
This is a lower bound on a spectral ratio,
or an upper bound on its excitation energy.
It complements, rather than replaces,
V83's upper ratio bound.

\subsection{Spectral crowding before boundary Gauss gluing}

Let \(\mathcal W_{N,P}=\bigoplus_{n=0}^N\mathcal W_{n,P}\).
Cross terms vanish by boundary symmetry, so
the preceding bound holds for every vector
in this direct sum, not just for its chosen
basis elements.

\begin{proposition}[An entire near-top trial space]
\label{f15v93:regional-count}
Let \(\lambda_j(T_R^{\rm int})\) be the ordered
eigenvalues after interior Gauss constraints,
indexed from \(j=0\). Put
\[
 D_N=\sum_{n=0}^N(n+1)^2
       =\frac{(N+1)(N+2)(2N+3)}6.
\]
Then
\begin{equation}
 \frac{\lambda_{D_N-1}(T_R^{\rm int})}{\Lambda_R}
       \geq
       \max\left\{0,1-\frac{N(N+2)\ell^2}{3}a_\gamma\right\}.
 \label{eq:f15v93:regional-count}
\end{equation}
For \(0<\varepsilon<1\), every integer \(N\geq0\)
with \(N(N+2)\ell^2\leq2\varepsilon\gamma\)
therefore certifies at least \(D_N\)
eigenvalues in
\([(1-\varepsilon)\Lambda_R,\Lambda_R]\).
\end{proposition}

\begin{proof}
Every unit vector in the \(D_N\)-dimensional
trial subspace has Rayleigh quotient at least
\(1-N(N+2)\ell^2a_\gamma/3\) for \(K_R\).
The variational min--max principle gives
\eqref{eq:f15v93:regional-count}. Positivity
gives the cap by zero, and
\eqref{eq:f15v93:increment} gives the
simpler sufficient cutoff.
\end{proof}

For fixed path length and fixed tolerance,
one can take \(N\) of order \(\sqrt\gamma\);
the certified regional count is of order
\(\gamma^{3/2}\). This does not assert that
\(\mathcal W_{N,P}\) itself is an exact
spectral subspace. It proves a spectral
count through a whole-subspace variational
bound. The estimate with \(N\) growing uses
the exact Casimir inequality, not the
fixed-\(n\) Bessel asymptotic.

\subsection{The complete open two-plaquette block has a different scale}
\label{f15v93:open-block}

Return to V92's seven-link graph. Let the
boundary set include the two endpoints
of its common link; any or all of the
other four vertices may also be boundary
vertices. Impose Gauss invariance at the
remaining vertices only. Thus this includes
the fully open regional block, not merely
the case with two exposed vertices.

Let \(r_{\rm open}(\gamma)\) be the second
eigenvalue ratio on this complete regional
space. Its top eigenvector is still the
fully gauge-invariant Perron vector.
Let \(q_\Box<1\) be the coupling-uniform
neutral bound of
Theorem~\ref{f15v92:uniform-gap}.

\begin{theorem}[Two-sided native open-block scale]
\label{f15v93:open-scale}
For every \(\gamma>0\),
\begin{equation}
 r_1(\gamma)\leq r_{\rm open}(\gamma)
            \leq\max\{q_\Box,r_1(3\gamma)\}.
 \label{eq:f15v93:open-squeeze}
\end{equation}
Consequently
\begin{equation}
 \frac12\leq
 \liminf_{\gamma\to\infty}\gamma(1-r_{\rm open}(\gamma))
 \leq
 \limsup_{\gamma\to\infty}\gamma(1-r_{\rm open}(\gamma))
 \leq\frac32 .
 \label{eq:f15v93:open-scale}
\end{equation}
The same bounds hold with
\(1-r_{\rm open}\) replaced by
\(-\log r_{\rm open}\).
\end{theorem}

\begin{proof}
Use the common link as a path of length
one. Its fundamental carrier is orthogonal
to the vacuum, and
Theorem~\ref{f15v93:carrier} bounds its
Rayleigh value below by \(r_1(\gamma)\).
This gives the left inequality.

Decompose the full regional space into
boundary isotypic sectors. Its all-trivial
sector is exactly the fully neutral
physical space of V92, whose nonvacuum
ratio is at most \(q_\Box\).
In every nontrivial sector choose a
boundary vertex carrying a nonzero label
\(n\). All regional vertex degrees are
at most three. V83's single-vertex bound
\eqref{eq:f15v83:regional-charge}, together
with monotonicity in the parameter and
decrease in representation label, gives
\[
 \|T_{R,r}\|/\Lambda_R
       \leq r_n(\gamma d_R(v))
       \leq r_1(3\gamma).
\]
This controls every charged sector,
including its entire multiplicity space.
Taking the maximum proves the upper bound
in \eqref{eq:f15v93:open-squeeze}.

For sufficiently large \(\gamma\),
\(r_1(3\gamma)>q_\Box\).
The fixed-order expansion
\eqref{eq:f15v83:asymptotic} gives
\[
 1-r_1(\gamma)=\frac3{2\gamma}+O(\gamma^{-2}),
 \qquad
 1-r_1(3\gamma)=\frac1{2\gamma}+O(\gamma^{-2}).
\]
Substitution proves \eqref{eq:f15v93:open-scale}.
The squeeze also gives \(1-r_{\rm open}=O(1/\gamma)\);
hence
\(-\log r_{\rm open}=(1-r_{\rm open})+O(\gamma^{-2})\).
\end{proof}

For completeness, this expansion is the
standard fixed-order modified-Bessel
expansion, also checked against
\href{https://dlmf.nist.gov/10.40.E1}
{DLMF 10.40.1}; it is not a new asymptotic
claim or a formula uniform in growing \(n\).

There is a label-resolved version. For the
sector with label \(n/2\) at both common-link
endpoints and trivial labels at all other
boundary vertices, write its normalized
norm as \(\rho_{\Box,n}(\gamma)\).
For \(n\geq1\),
\begin{equation}
 \max\left\{0,1-\frac{n(n+2)}3a_\gamma\right\}
       \leq\rho_{\Box,n}(\gamma)\leq r_n(3\gamma).
 \label{eq:f15v93:fixed-charge}
\end{equation}
The lower bound is the common-link trial
carrier and the upper bound uses one
degree-three endpoint. For each fixed \(n\),
\[
 \frac{n(n+2)}6
  \leq\liminf_{\gamma\to\infty}\gamma(1-\rho_{\Box,n})
  \leq\limsup_{\gamma\to\infty}\gamma(1-\rho_{\Box,n})
  \leq\frac{n(n+2)}2 .
\]
This establishes the Casimir scale of actual
native charged sectors, but not convergence
to a particular effective boundary generator.

There is no contradiction with V92.
The fully neutral carrier removes these
open endpoint charges; the regional carrier
must keep them for subsequent matching.
A coupling-independent neutral block gap
is not a coupling-independent open-block
gap. Nor is the latter's \(1/\gamma\)
dimensionless scale a statement that a
physical mass vanishes: division by an
independently supplied \(a_t(\gamma)\)
still has to be justified.

\subsection{Gauss gluing makes physical cut states, with no carrier multiplicity}
\label{f15v93:cut-states}

Let \(A,B\) be disjoint link regions, with
two common boundary vertices \(v,w\).
Suppose each region contains a path between
them, of lengths \(\ell_A,\ell_B\).
Other shared vertices are allowed; path
holonomies transform only at their endpoints.
Use the actual regional Perron vectors
\(\Omega_A,\Omega_B\), and the cut transfer
\[
 T_{\rm cut}=(T_A\otimes T_B)|_{\rm phys},
 \qquad\Lambda_{\rm cut}=\Lambda_A\Lambda_B.
\]
The leading eigenvalue identity is V82's
native cut theorem.

\begin{theorem}[A native physical tower for the cut transfer]
\label{f15v93:cut-tower}
The vectors
\begin{equation}
 \Gamma_n(U_A,U_B)=\Omega_A(U_A)\Omega_B(U_B)
      \chi_n(H_A(U_A)H_B(U_B)^{-1}),\qquad n\geq0,
 \label{eq:f15v93:closed-trials}
\end{equation}
are normalized, mutually orthogonal and
fully gauge invariant. The compression of
\(T_{\rm cut}/\Lambda_{\rm cut}\) to their
span is diagonal, with entries
\(t_{A,n}t_{B,n}\), where \(t_{R,n}\) is
\eqref{eq:f15v93:carrier-compression} for
the selected regional path. In particular
\begin{equation}
 \frac{\lambda_N(T_{\rm cut})}{\Lambda_{\rm cut}}
 \geq
 \max\left\{0,1-
  \frac{N(N+2)(\ell_A^2+\ell_B^2)}3a_\gamma\right\}.
 \label{eq:f15v93:cut-count}
\end{equation}
\end{theorem}

\begin{proof}
Both holonomies transform as
\(H_R\mapsto g_vH_Rg_w^{-1}\), so their
relative product transforms by conjugation.
The character is fully invariant, including
at all other vertices. Under the product
regional vacuum law, the two holonomies
are independent Haar variables by the
argument in Theorem~\ref{f15v93:carrier}.
Their relative product is Haar, proving
normalization and character orthogonality.

More explicitly the two endpoint carrier
has dimension \(D=(n+1)^2=d_n^2\).
The normalized singlet gluing is
\[
 \Gamma_n=\frac1{d_n}\sum_{a,b=1}^{d_n}
       \Phi_{A,n,ab}\otimes\overline{\Phi_{B,n,ab}}.
\]
There are \(d_n^2\) terms, each with
coefficient \(1/d_n\); the squared norm
is one. The regional compressions are
scalar on those copies, so their product
has expectation \(t_{A,n}t_{B,n}\).
Reality of the native kernel makes the
conjugate copy have the same scalar.
Different \(n\)'s lie in different matched
boundary sectors and have zero cross
matrix elements. No factor \(D\) occurs.

Let
\(\delta_{R,n}=n(n+2)\ell_R^2a_\gamma/3\).
When \(\delta_{A,n}+\delta_{B,n}<1\), both
individual lower bounds are nonnegative,
and
\[
 t_{A,n}t_{B,n}
 \geq(1-\delta_{A,n})(1-\delta_{B,n})
 \geq1-\delta_{A,n}-\delta_{B,n}.
\]
If their sum is at least one, positivity
provides the weaker zero bound.
Apply this to every \(n\leq N\) and use
min--max on the \(N+1\) dimensional span.
\end{proof}

For fixed path lengths and tolerance,
\(N\) can again grow as \(\sqrt\gamma\).
But the physical count is now \(N+1\),
not \(\sum_{n\leq N}(n+1)^2\).
Gauss gluing contracts a single normalized
singlet in each matched carrier; it does
not produce \(D\) physical copies.
The state \(\Gamma_0=\Omega_A\Omega_B\)
is the cut vacuum and every \(n>0\) state
is orthogonal to it.

These are genuine physical states of
the cut operator, with the regional
magnetic weights still present.
They are not yet trial states with the
same energy for the uncut transfer.
The following identity retains the
precise interface cost.

\subsection{The exact magnetic interface metric}

Write the full native physical transfer
as in V82,
\[
 T=M T_{\rm cut}M,\qquad M=M_{b_\times},
 \quad b_\times=\exp[-\gamma\sum_{p\ {\rm mixed}}s_p/2].
\]
At each finite regulator \(M^{-1}\) is
bounded and gauge invariant.
For \(0\leq n\leq N\), put
\(\Psi_n=M^{-1}\Gamma_n\) and define
\[
 B_{nm}=\langle\Gamma_n,M^{-2}\Gamma_m\rangle,
 \qquad
 D_{nm}=\delta_{nm}t_{A,n}t_{B,n}.
\]

\begin{proposition}[Exact interface-weighted trial pencil]
\label{f15v93:interface-pencil}
The Gram matrix of the \(\Psi_n\)'s is
\(B\geq I\), and their full-transfer
quadratic-form matrix is
\begin{equation}
 \Lambda_{\rm cut}^{-1}
       \langle\Psi_n,T\Psi_m\rangle=D_{nm}.
 \label{eq:f15v93:interface-pencil}
\end{equation}
For \(0\leq j\leq N\),
\begin{equation}
 \frac{\lambda_j(T)}{\Lambda_{\rm cut}}
       \geq\lambda_j(B^{-1/2}DB^{-1/2}).
 \label{eq:f15v93:interface-minmax}
\end{equation}
These are lower spectral bounds from a
trial space, not a lower bound on the
full Hamiltonian gap.
\end{proposition}

\begin{proof}
The Gram identity is the definition.
Since \(0<b_\times\leq1\), \(M^{-2}\geq I\)
as a multiplication operator, and the
\(\Gamma_n\)'s are orthonormal, giving
\(B\geq I\). The exact congruence yields
\[
 \langle M^{-1}\Gamma_n,
           MT_{\rm cut}M M^{-1}\Gamma_m\rangle
       =\langle\Gamma_n,T_{\rm cut}\Gamma_m\rangle.
\]
Theorem~\ref{f15v93:cut-tower} gives
\eqref{eq:f15v93:interface-pencil}.
A coefficient vector \(c\) has norm square
\(c^*Bc\) and transfer energy
\(\Lambda_{\rm cut}c^*Dc\). Orthonormalizing
this finite trial space gives
\(B^{-1/2}DB^{-1/2}\); min--max proves
\eqref{eq:f15v93:interface-minmax}.
No commutation of \(M\) with \(T_{\rm cut}\)
or false operator-order comparison between
these congruences is used.
\end{proof}

The metric is explicit for a single mixed
plaquette formed by the two paths, with
\(\ell_A+\ell_B=4\), and with no other
mixed plaquettes in this calculation.
Set \(W=H_AH_B^{-1}\) and
\(M=e^{-\gamma(1-\frac12\Tr W)/2}\).
Then the cut vacuum makes \(W\) exactly
Haar, regardless of the other wholly
regional plaquette weights. Hence
\begin{equation}
 B_{nm}=\int_G e^{\gamma(1-\frac12\Tr W)}
                              \chi_n(W)\chi_m(W)\,dH(W).
 \label{eq:f15v93:interface-gram}
\end{equation}

\begin{proposition}[One-plaquette metric and its nonperturbative size]
\label{f15v93:interface-explicit}
For the preceding single interface,
\begin{equation}
 B_{nm}=e^{2\gamma}c_0(\gamma)(-1)^{n+m}
       \sum_{\substack{k=|n-m|\\k\equiv n+m\ (2)}}^{n+m}
                                    (k+1)r_k(\gamma).
 \label{eq:f15v93:interface-explicit}
\end{equation}
In particular
\begin{equation}
 B_{00}=e^{2\gamma}c_0(\gamma)
       =\sqrt{\frac2\pi}\,
          e^{2\gamma}\gamma^{-3/2}(1+O(\gamma^{-1})).
 \label{eq:f15v93:interface-size}
\end{equation}
For each fixed finite \(N\),
\[
 B/B_{00}\longrightarrow vv^*,\qquad
 v_n=(-1)^n(n+1),\quad 0\leq n\leq N.
\]
\end{proposition}

\begin{proof}
The character fusion identity is
\[
 \chi_n\chi_m
    =\sum_{k=|n-m|,\,|n-m|+2,\ldots,n+m}\chi_k .
\]
Changing \(W\) to \(-W\) in each term gives
\[
 \int e^{-\gamma h(W)}\chi_k(W)\,dH(W)
       =(-1)^k z(\gamma)(k+1)r_k(\gamma),
 \qquad h(W)=\tfrac12\Tr W.
\]
All labels in the sum have parity \(n+m\).
Use \(z(\gamma)=e^\gamma c_0(\gamma)\)
to obtain \eqref{eq:f15v93:interface-explicit}.
Its \(00\) entry gives
\eqref{eq:f15v93:interface-size}, using the
already proved scalar Laplace normalization.
Finally \(r_k(\gamma)\to1\) for each fixed
label, and
\[
 \sum_{k=|n-m|,\,|n-m|+2,\ldots,n+m}(k+1)
                                   =(n+1)(m+1).
\]
Entrywise convergence in this fixed finite
matrix proves the last assertion.
\end{proof}

Thus \(\|B\|\geq B_{00}\) grows exponentially.
The inverse-multiplier transport is not
near-isometric and cannot justify carrying
the cut energy estimate unchanged into
the full transfer. The fixed-band rank-one
limit is not a uniform estimate for growing
\(N\), nor a lower bound of that exponential
size in every direction. In particular
one cannot replace the matrix \(B\) by
\(B_{00}I\).

This example makes the cost explicit without
asserting that all gluing destroys a gap.
Magnetic gluing can remove these slow
relative boundary orientations. To determine
the actual full spectrum one must control
the resulting interacting boundary
operator, its true normalization and the
remaining collective modes.

\subsection{Direction and verification record}

The closed-block result now has a precise
open-boundary counterpart. The native
neutral internal sector can stay uniformly
separated while an increasing family of
boundary representations lies within
\(O(c_n/\gamma)\) of the regional top.
Their scale matches the need for growing
representation cutoffs, rather than a
fixed set of charged channels.

The next global step must retain these
boundary variables and the actual mixed
plaquette matrix elements. It cannot
start from a coupling-independent
open-block gap, count singlet carriers
as independent physical states, or
transfer the cut Rayleigh bounds through
an unpaid magnetic multiplier.
The available routes are an exact
boundary effective operator or a
same-law comparison that proves its
collective coercivity. Such a comparison,
physical time calibration and interacting
continuum reconstruction are still missing.
No conclusion of physical gaplessness
has been drawn from regional or cut
dimensionless gaps.

The checker
\begin{center}\small
\path{scripts/verify_ym_f15_v93_boundary_carriers_and_gluing.py}
\end{center}
passed all 4,778 new exact-rational fixtures
and the inherited V92 chain. Its tests
cover the Casimir character bound,
noncommuting path telescoping in both
orientations, Perron-pair trace algebra,
regional and glued counting budgets,
normalized singlet contractions,
coefficient-by-coefficient interface
character fusion, and the exact
generalized-eigenvalue congruence.
The finite positive kernels used to
test algebra are not called Wilson
laws. These diagnostics supplement,
but do not replace, the all-region
analytic and variational proofs.
The script has 8,460 bytes and SHA--256
\begin{center}\scriptsize\ttfamily
CEDCF7F7FABD57BFC7078C215534E207946F0393D4FCF4B4EB3FB409408B0F4A.
\end{center}

The predecessor V92 has 2,098,034 bytes
and SHA--256
\begin{center}\scriptsize\ttfamily
128E1F1F6421A72A8494A6861EAA61D6643EC48F3B17D2450E4D86428B95C2C6.
\end{center}
Its body is preserved byte for byte;
removing this section and restoring
the title version recovers V92.
V93 replaces only the active main
source. Archives, companion notes
and earlier diagnostic scripts remain
unchanged. The last companion checkpoint
is V90 and the next is V95.
Only \path{.tex} files belong in
\path{latex_notes}; builds are directed
to \path{artifacts/ym_f15_v93_texcheck/}.

\begin{firewall}
V93 proves native charged trial-space
bounds, spectral counting, and the
complete \(1/\gamma\) scale of the
open two-plaquette block. It also
constructs physical cut states and
their exact magnetic interface metric.
It does not prove that the cut states
remain slow after gluing, a volume-uniform
full-transfer gap, or a nontrivial
four-dimensional continuum theory.
The full Yang--Mills mass gap remains
unproved.
\end{firewall}

\section{V94: forward magnetic fusion and a uniform one-interface normalizer}
\label{f15v94:context}

V93 showed that inverse-multiplier transport of cut
states carries an exponentially large Gram matrix.
That is a statement about a particular transported
trial space, not an exponential upper bound on the
actual glued Perron eigenvalue. This section works
with the forward multiplier and keeps every charge
label. It proves that restoring one specified mixed
plaquette has a coupling-uniform, spectator-independent
Perron cost. This is an actual native-transfer
normalization result, but not a gap theorem.

The forward multiplier becomes an exact killed
nearest-neighbour walk on the nonnegative charge
labels. We use its full infinite matrix, derive a
growing-band error bound, and keep the coupling to
the orthogonal physical complement. The latter is
essential: an exact compression is not an invariant
effective theory.

V85 already calculated one-link magnetic Gram
blocks, V84 already proved the Bessel coefficient
tail, and V93 already constructed the cut Wilson-line
carriers. Those results are not repeated as new
discoveries. The additions here are their native
forward interface combination, the uniform Perron
bound, the resulting relative charge truncation, and
an explicit bound on the missing complementary block.

\subsection{A single mixed plaquette and its exact forward compression}
\label{f15v94:setup}

Use V93's two disjoint link regions \(A,B\), actual
normalized positive regional Perron vectors
\(\Omega_A,\Omega_B\), and paths from \(v\) to \(w\).
Assume the paths form a simple mixed plaquette:
\[
 W=H_AH_B^{-1},\qquad \ell_A+\ell_B=4.
\]
There are no other mixed magnetic factors in this
section. Arbitrarily many wholly regional
plaquettes are allowed. More generally the regional
multipliers may be any positive continuous functions
invariant under all regional vertex gauge
transformations. The path and gauge assumptions of
V93 are retained.

On the full physical space put
\[
 C=T_{\rm cut}/\Lambda_{\rm cut},\qquad
 M=M_{b_\gamma(W)},\qquad
 b_\gamma(W)=e^{-\gamma(1-\frac12\Tr W)/2},\qquad
 \widehat T=T/\Lambda_{\rm cut}=MCM.
\]
Thus \(0\leq C\leq I\), \(\|C\|=1\), and \(0<M\leq I\).
The normalizer is still
\(\Lambda_{\rm cut}=\Lambda_A\Lambda_B\), not the
unknown glued eigenvalue.

Let \(\mathcal H_\chi=L^2(G,dH)^{\operatorname{Ad}G}\),
with orthonormal character basis
\(\{\chi_n:n\geq0\}\). V93 gives the isometry
\[
 Jf=\Omega_A\Omega_B f(W),\qquad
 J\chi_n=\Gamma_n,\qquad Q=JJ^*.
\]
Write \(F_\gamma\) for multiplication by \(b_\gamma\)
on \(\mathcal H_\chi\), and define
\[
 w_n=t_{A,n}t_{B,n},\qquad D\chi_n=w_n\chi_n,
 \qquad S=\ell_A^2+\ell_B^2.
\]
Here \(w_0=1\) and \(0\leq w_n\leq1\). The notation
\(w_n\) is not a representation dimension.

\begin{proposition}[Exact native forward charge chain]
\label{f15v94:forward}
The closed subspace \(Q\mathcal H_{\rm phys}\) reduces
\(M\), and
\begin{equation}
 MJ=JF_\gamma,\qquad J^*CJ=D,\qquad
 E_\gamma:=J^*\widehat T J=F_\gamma D F_\gamma.
 \label{eq:f15v94:forward}
\end{equation}
In particular \(E_\gamma\) is the exact compression
on this infinite-dimensional physical trial space.
No invariance of that space under \(C\) or
\(\widehat T\) is asserted.
\end{proposition}

\begin{proof}
Under the cut vacuum law, \(W\) is exactly Haar,
by V93's endpoint gauge argument and independence
of the two regional laws. Thus \(J\) is an isometry
on the whole central \(L^2\) space, not just on
finite character sums. Multiplication gives
\(MJf=J(b_\gamma f)\). Since \(M\) is bounded
self-adjoint, invariance of the range of \(J\)
also implies invariance of its orthogonal
complement. V93's diagonal cut compression gives
\(J^*CJ=D\), first on finite sums and then by
boundedness. Taking adjoints of \(MJ=JF_\gamma\)
proves the final identity.
\end{proof}

There are useful all-label bounds already available.
For \(a_\gamma=1-r_1(\gamma)\), V93 gives
\begin{equation}
 w_n\geq
 \max\left\{0,1-\frac{n(n+2)S}{3}a_\gamma\right\}
 \geq\max\left\{0,1-\frac{n(n+2)S}{2\gamma}\right\}.
 \label{eq:f15v94:diagonal-lower}
\end{equation}
If \(s_A=\gamma d_A(v)\), \(s_B=\gamma d_B(v)\),
V83's single-vertex charged domination gives
\begin{equation}
 w_n\leq r_n(s_A)r_n(s_B).
 \label{eq:f15v94:diagonal-upper}
\end{equation}
Indeed each \(t_{R,n}\) is a nonnegative Rayleigh
value in the endpoint charge sector, bounded above
by its normalized sector norm. The degrees here
belong to the actual regional graph.

\subsection{The forward multiplier is a killed charge walk}
\label{f15v94:walk}

Identify \(\mathcal H_\chi\) with
\(\ell^2(\mathbb N_0)\) by the character basis.
Let \(A_+\) be the half-line adjacency operator:
\[
 A_+e_0=e_1,\qquad A_+e_n=e_{n-1}+e_{n+1}
 \quad(n\geq1).
\]
The fusion rule \(\chi_1\chi_n=\chi_{n-1}+\chi_{n+1}\),
with \(\chi_{-1}=0\), says that multiplication by
\(\frac12\Tr W\) is \(A_+/2\). Thus
\begin{equation}
 F_\gamma=\exp\left[\frac{\gamma}{4}(A_+-2I)\right].
 \label{eq:f15v94:walk}
\end{equation}

\begin{proposition}[Exact infinite charge-walk matrix]
\label{f15v94:walk-matrix}
For all \(n,m\geq0\),
\begin{equation}
 (F_\gamma)_{nm}
 =e^{-\gamma/2}
   \left[I_{|n-m|}(\gamma/2)-I_{n+m+2}(\gamma/2)\right].
 \label{eq:f15v94:walk-matrix}
\end{equation}
Its entries are nonnegative and its row and column
sums are at most one. It is the transition
semigroup, at time \(\gamma/4\), of the continuous-time
walk on \(\mathbb N_0\) that jumps at rate one
in either direction and is killed upon reaching
\(-1\). This is a representation of the magnetic
multiplier, not an auxiliary dynamics substituted
for the physical transfer.
\end{proposition}

\begin{proof}
Equation~\eqref{eq:f15v94:walk} follows by bounded
functional calculus. With \(h=\cos\theta\),
central Haar measure is
\((2/\pi)\sin^2\theta\,d\theta\), and
\(\chi_n=\sin((n+1)\theta)/\sin\theta\).
Consequently the matrix entry is
\[
 \frac2\pi e^{-\gamma/2}\int_0^\pi
  e^{(\gamma/2)\cos\theta}
          \sin((n+1)\theta)\sin((m+1)\theta)\,d\theta .
\]
The product-to-sum formula and the Fourier
coefficients of \(e^{z\cos\theta}\) give
\eqref{eq:f15v94:walk-matrix}. The Fourier identity
is also recorded in
\href{https://dlmf.nist.gov/10.35.E2}{DLMF 10.35.2}.

For a direct positivity proof, the coefficient
\((A_+^j)_{nm}\) counts length-\(j\) nearest-neighbour
paths from \(m\) to \(n\) which stay nonnegative.
Expanding \(e^{-2u}e^{uA_+}\), \(u=\gamma/4\),
is the rate-two Poisson jump expansion, with each
allowed direction weighted by one half.
Paths exiting at \(-1\) are removed. This proves
the transition interpretation and the row-sum
bound; symmetry gives the column-sum bound.
Reflection at the first visit to \(-1\) gives
the difference of full-line kernels in
\eqref{eq:f15v94:walk-matrix}.
\end{proof}

Let \(P_N\) project onto charges \(0,\ldots,N\).
A finite exponential
\(\exp[(\gamma/4)P_N(A_+-2I)P_N]\) on that finite
space is generally not \(P_NF_\gamma P_N\):
paths can leave the retained interval and return.
Already for \(N=0\), the full \(A_+^2\) has
\(00\) entry one, whereas \((P_0A_+P_0)^2=0\).
The finite spectral construction below therefore
uses the exact infinite-walk matrix, not a
silently absorbing upper charge boundary.

\begin{proposition}[A spectator-independent magnetic charge buffer]
\label{f15v94:buffer}
For integers \(N,L\geq0\), set
\[
 \mathcal I_\gamma(x)=
 x\,\operatorname{arsinh}(2x/\gamma)
 -\frac{\gamma}{2}
          \left(\sqrt{1+4x^2/\gamma^2}-1\right).
\]
Then
\begin{equation}
 \|(I-P_{N+L})F_\gamma P_N\|
       \leq e^{-\mathcal I_\gamma(L+1)},\qquad
 \mathcal I_\gamma(x)\geq
             \min\left\{\frac{x^2}{2\gamma},\frac x2\right\}.
 \label{eq:f15v94:buffer}
\end{equation}
\end{proposition}

\begin{proof}
The killed kernel is bounded entrywise by the
full-line walk kernel
\(p_k=e^{-\gamma/2}I_{|k|}(\gamma/2)\).
The latter is the distribution of \(X=N_+-N_-\),
where the two independent Poisson variables have
mean \(\gamma/4\). Hence
\[
 \mathbb E e^{\theta X}
       =\exp\left[\frac{\gamma}{2}(\cosh\theta-1)\right].
\]
Every row and every column sum of the nonnegative
off-band block is bounded by
\(\mathbb P(X\geq L+1)\). The Schur test bounds its
\(\ell^2\) norm by that same number. Exponential
Markov inequality, optimized at
\(\theta=\operatorname{arsinh}(2x/\gamma)\),
gives the first assertion. For the simpler bound,
use \(\cosh\theta-1\leq\theta^2\) on \([0,1]\)
and the admissible choice
\(\theta=\min\{1,x/\gamma\}\).
\end{proof}

Thus fixed-error magnetic buffering costs
\(O(\sqrt\gamma)\) charges at large \(\gamma\);
for a varying error its logarithmic budget must
be retained. This does not use a fixed-order
Bessel asymptotic in a growing-order regime.

\subsection{A coupling-uniform native Perron cost for one interface}
\label{f15v94:normalizer-section}

The exponential inverse metric of V93 does not
prevent an appropriate forward trial space from
having a uniformly positive Rayleigh quotient.
Define the explicit, deliberately nonoptimal constant
\begin{equation}
 \kappa_S=\frac{e^{-1/2}}{384\pi^3S^{3/2}}>0.
 \label{eq:f15v94:kappa}
\end{equation}

\begin{theorem}[One-interface Perron normalization without a coupling loss]
\label{f15v94:normalizer}
For every \(\gamma>0\) in the preceding single-interface
native setting,
\begin{equation}
 \kappa_S\leq\|E_\gamma\|
       \leq\frac{\lambda_0(T)}{\Lambda_A\Lambda_B}
       \leq1.
 \label{eq:f15v94:normalizer}
\end{equation}
The constant depends only on the two path lengths,
not on coupling, regional volume, or oscillation
of the wholly regional magnetic multipliers.
\end{theorem}

\begin{proof}
First suppose \(\gamma\geq16S\), and put
\[
 N=\left\lfloor\sqrt{\frac{\gamma}{4S}}\right\rfloor.
\]
Then \(N\geq2\),
\(N\geq\sqrt\gamma/(4\sqrt S)\), and
\(N(N+2)\leq2N^2\leq\gamma/(2S)\).
Equation~\eqref{eq:f15v94:diagonal-lower} gives
\(D\geq(3/4)P_N\), whence
\[
 \|E_\gamma\|\geq\frac34\|F_\gamma P_NF_\gamma\|
                   =\frac34\|P_NF_\gamma^2P_N\|.
\]
For a unit vector in this last finite space take
\[
 \psi_N=Z_N^{-1/2}\sum_{n=0}^N(n+1)\chi_n,\qquad
 Z_N=\sum_{n=0}^N(n+1)^2\geq N^3/3.
\]
Here \(Z_N\) is the squared norm of a weighted
character sum. The physical trial space still
has dimension \(N+1\), not \(Z_N\).

For \(0\leq\theta\leq\gamma^{-1/2}\), V93's
character-chord bound and \(1-\cos\theta\leq\theta^2/2\)
imply, for all \(n\leq N\),
\[
 \frac{\chi_n(\theta)}{n+1}
 \geq1-\frac{n(n+2)\theta^2}{6}
 \geq1-\frac1{12S}\geq\frac12.
\]
Thus \(\psi_N(\theta)\geq\sqrt{Z_N}/2\).
On the same ball \(b_\gamma^2\geq e^{-1/2}\).
Since \(\gamma\geq16S\) and \(S\geq2\), its angular
radius is less than \(\pi/2\). The inequality
\(\sin\theta\geq2\theta/\pi\) gives
\[
 H\{\theta\leq\gamma^{-1/2}\}
    =\frac2\pi\int_0^{\gamma^{-1/2}}\sin^2\theta\,d\theta
    \geq\frac8{3\pi^3}\gamma^{-3/2}.
\]
Consequently
\[
 \begin{split}
 \langle\psi_N,F_\gamma^2\psi_N\rangle
 &\geq\frac{Z_N}{4}e^{-1/2}
                         \frac8{3\pi^3}\gamma^{-3/2}\\
 &\geq\frac{e^{-1/2}}{288\pi^3S^{3/2}}.
 \end{split}
\]
Multiplication by \(3/4\) proves the lower bound.

For \(0<\gamma\leq16S\), use \(D\geq P_0\).
The rank-one operator \(F_\gamma P_0F_\gamma\)
has norm
\[
 \|F_\gamma\chi_0\|^2
       =c_0(\gamma)\geq c_0(16S).
\]
The same Haar-ball estimate, now with parameter
\(16S\), gives
\[
 c_0(16S)\geq
   e^{-1/2}\frac8{3\pi^3}(16S)^{-3/2}
   =\frac{e^{-1/2}}{24\pi^3S^{3/2}}
   \geq\kappa_S.
\]
Finally a compression cannot increase operator
norm, while \(\|MCM\|\leq\|C\|=1\). The physical
Perron eigenvalue is the norm of \(T\), so these
two observations prove the remaining inequalities.
\end{proof}

This is a bound on the actual glued leading
eigenvalue, not merely on a selected character
matrix. It uses a \(\gamma\)-dependent coherent
sum of charges. It neither asserts that the
glued vacuum equals that sum nor bounds the
second eigenvalue away from the first.
For a square, \(S=8\) or \(S=10\).
Arbitrary simultaneous interfaces do not satisfy
this single-relative-holonomy hypothesis;
no iteration over an extensive interface set
is implicit in the theorem.

\subsection{A finite charge matrix with a relative error}
\label{f15v94:finite-section}

Define
\[
 D_N=P_NDP_N,\qquad E_{\gamma,N}=F_\gamma D_NF_\gamma,
 \qquad
 \eta_N=r_{N+1}(s_A)r_{N+1}(s_B).
\]
These are kinetic charge cutoffs followed by
the full magnetic multiplier, not raw projection
of all output charge labels.

\begin{theorem}[Relative approximation of the exact boundary compression]
\label{f15v94:finite}
At every fixed regulator,
\begin{equation}
 0\leq E_\gamma-E_{\gamma,N},\qquad
 \|E_\gamma-E_{\gamma,N}\|\leq\eta_N,\qquad
 \frac{\|E_\gamma-E_{\gamma,N}\|}{\|E_\gamma\|}
                       \leq\frac{\eta_N}{\kappa_S}.
 \label{eq:f15v94:finite}
\end{equation}
The nonzero spectrum of \(E_{\gamma,N}\) is that of
the \((N+1)\)-by-\((N+1)\) positive matrix
\begin{equation}
 (H_N)_{nm}=
 \sqrt{w_nw_m}\,e^{-\gamma}
       [I_{|n-m|}(\gamma)-I_{n+m+2}(\gamma)],
 \qquad 0\leq n,m\leq N.
 \label{eq:f15v94:finite-matrix}
\end{equation}
Ordered eigenvalues, padded by zeros, differ
from those of \(E_\gamma\) by at most \(\eta_N\).
The conclusion concerns \(E_\gamma\), not the
entire full physical transfer.
\end{theorem}

\begin{proof}
The upper bound
\eqref{eq:f15v94:diagonal-upper} and decrease
of \(r_n\) with label give
\(0\leq D-D_N\leq\eta_N I\). Congruence by
the positive contraction \(F_\gamma\) proves
the first two bounds. Theorem~\ref{f15v94:normalizer}
gives the relative denominator.

Set \(V=F_\gamma P_ND_N^{1/2}\), from the finite
charge space to the full character space.
Then \(VV^*=E_{\gamma,N}\) and
\(V^*V=D_N^{1/2}P_NF_\gamma^2P_ND_N^{1/2}\).
Their nonzero eigenvalues agree. Since
\(F_\gamma^2=F_{2\gamma}\), the exact walk formula
proves \eqref{eq:f15v94:finite-matrix}.
Finally \(\eta_N\to0\) at fixed \(s_A,s_B\);
alternatively compactness follows by compression
of the compact native operator. Min--max gives
the asserted eigenvalue perturbation bound.
\end{proof}

All unknown regional dynamics remain in the
actual scalars \(w_n\). Equations
\eqref{eq:f15v94:diagonal-lower}--\eqref{eq:f15v94:diagonal-upper}
also give rigorous diagonal upper and lower
operator comparisons; they are not a license
to choose \(w_n\) freely as a model.

There is an explicit growing cutoff. V84's
already proved bound \eqref{eq:f15v84:bernstein}
gives, for \(k=N+1\),
\begin{equation}
 \eta_N\leq
 \exp\left[-\frac{k(k+2)}{2s_A+2k+1}
           -\frac{k(k+2)}{2s_B+2k+1}\right].
 \label{eq:f15v94:kinetic-tail}
\end{equation}
On the spatial cubic graph, \(d_A(v),d_B(v)\leq6\).
Bounding both parameters by \(6\gamma\) yields
\begin{equation}
 \eta_N\leq
       \exp\left[-\min\left\{\frac{k^2}{15\gamma},
                                      \frac{2k}{5}\right\}\right].
 \label{eq:f15v94:simple-tail}
\end{equation}
For clarity, if \(k\leq6\gamma\), then \(6\gamma\geq1\)
and \(12\gamma+2k+1\leq30\gamma\).
If \(k>6\gamma\), the same denominator is at most
\(5k\). These two cases prove the displayed
simplification for every \(\gamma>0\).

For \(0<\varepsilon<1\), put
\(b_\varepsilon=\log(1/(\kappa_S\varepsilon))\).
The sufficient choice
\[
 N+1\geq
 \max\left\{\sqrt{15\gamma b_\varepsilon},
                          \frac52 b_\varepsilon\right\}
\]
makes the relative error in
\eqref{eq:f15v94:finite} at most \(\varepsilon\).
At fixed error this retains \(O(\sqrt\gamma)\)
physical charge states. There is no regional-volume
count or representation-dimension multiplier in
this estimate. It is not V91's finite-rank
approximation of the entire spatial transfer.

\subsection{The complementary coupling is retained explicitly}
\label{f15v94:leakage-section}

The exact forward compression is not a closure
of the full physical Hilbert space. The missing
block has a useful paid bound. Put
\[
 R=(I-Q)CJ,\qquad
 B=(I-Q)\widehat T J,\qquad
 M_\perp=M|_{(I-Q)\mathcal H_{\rm phys}} .
\]

\begin{proposition}[Native forward complementary leakage]
\label{f15v94:leakage}
The following inequalities hold on the complete
character space:
\begin{equation}
 R^*R\leq D-D^2,\qquad
 B=M_\perp R F_\gamma,\qquad
 B^*B\leq F_\gamma(D-D^2)F_\gamma.
 \label{eq:f15v94:leakage}
\end{equation}
In particular, with \(N,L\geq0\),
\begin{equation}
 \|BP_N\|\leq
 \sqrt{\min\left\{\frac14,
            \frac{(N+L)(N+L+2)S}{2\gamma}\right\}}
       +\frac12 e^{-\mathcal I_\gamma(L+1)}.
 \label{eq:f15v94:buffered-leakage}
\end{equation}
\end{proposition}

\begin{proof}
The orthogonal decomposition \(CJ=JD+R\) gives
\[
 D^2+R^*R=J^*C^2J\leq J^*CJ=D,
\]
because \(C\) is a positive contraction.
Since \(Q\) reduces \(M\), the formula for \(B\)
follows from \(MJ=JF_\gamma\). Also
\(M_\perp^2\leq I\), proving its form bound.
For every unit \(f\),
\[
 \|Bf\|\leq\|(D-D^2)^{1/2}F_\gamma f\|.
\]
The square-root diagonal has norm at most \(1/2\).
For \(n\leq K\), use
\(w_n(1-w_n)\leq1-w_n\) and
\eqref{eq:f15v94:diagonal-lower} to bound its
squared norm on \(P_K\) by
\(\min\{1/4,K(K+2)S/(2\gamma)\}\).
Split \(F_\gamma P_N\) at \(K=N+L\) and apply
Proposition~\ref{f15v94:buffer} to the remaining
block. This proves
\eqref{eq:f15v94:buffered-leakage}.
\end{proof}

The unmixed low-charge defect is small for fixed
charge. Forward magnetic multiplication, however,
spreads over the \(\sqrt\gamma\) charge scale.
The last bound does not tend to zero uniformly
on the retained growing band. No small-leakage
conclusion is inferred from it. Also there is
as yet no uniform spectral estimate for
\((I-Q)\widehat T(I-Q)\).
An eventual Schur-complement or same-law comparison
must pay both quantities; a gap of \(H_N\) alone
would not be a gap of the full YM transfer.

\subsection{What changes in the attack}
\label{f15v94:direction}

The single-interface normalization no longer
requires an exponentially bad inverse-multiplier
estimate. The exact forward chain and its relative
finite approximation retain native regional
Perron data and the whole magnetic charge walk.
This is a more useful interface object than a
fixed set of raw charges.

The next obstruction is consequently specific:
control the other physical multiplicities and
their coupling to this boundary chain, then
handle overlapping mixed plaquettes without
losing uniformity over the growing spatial graph.
The one-interface theorem does not supply that
collective coercivity. The independent physical
clock and nontrivial continuum reconstruction
also remain open. Neither a volume-uniform
mass gap nor physical gaplessness follows here.

The checker
\begin{center}\small
\path{scripts/verify_ym_f15_v94_forward_magnetic_interface.py}
\end{center}
checks the reflected half-line path counts,
the Bessel-difference coefficients, Haar fusion
coefficients, coherent-character normalization,
the Perron-bound constants, finite Gram spectra,
positive truncation errors, and the complementary
leakage inequality. Its finite positive
contractions are algebra fixtures, not substituted
Wilson laws. The analytic all-coupling arguments
above are the proofs; the checks are supplementary.
All 7,087 new exact-rational fixtures and the
inherited V93 chain passed. The checker has
8,772 bytes and SHA--256
\begin{center}\scriptsize\ttfamily
A55AB706014CD85A713FEE674DA52C78D7D71664C1AC704BB7C89A14536BA766.
\end{center}

The predecessor V93 has 2,123,725 bytes and SHA--256
\begin{center}\scriptsize\ttfamily
88826AD4E009ABC13B4A44FBE504FFE6E0BCD8482E66C1BBA8EED32FD135636F.
\end{center}
Its body is preserved byte for byte: remove this
section and restore the title version to recover
V93. Only the active main source is replaced.
Archives, supplied papers, companion notes and
earlier diagnostic scripts are unchanged.
The last five-version companion checkpoint is
V90; the next is V95. Only \path{.tex} files
belong in \path{latex_notes}; build products go
to \path{artifacts/ym_f15_v94_texcheck/}.

\begin{firewall}
V94 proves a coupling-uniform Perron normalization
bound for one native mixed plaquette, an exact
forward charge compression, its controlled
finite approximation, and an explicit
complementary leakage bound. It does not
identify that compression with the full physical
transfer, control an extensive mixed interface,
prove a volume-uniform physical gap, or
construct the interacting four-dimensional
continuum theory. The full Yang--Mills mass
gap remains unproved.
\end{firewall}

\section{V95: a growing charged ground band and the complete regional complement}
\label{f15v95:context}

V94 calculated a native forward interface compression,
but correctly left its coupling to other physical
multiplicities unresolved. We now go upstream to a
regional operator for which those multiplicities can
be controlled. The region is V92's two-plaquette
seven-link graph, with its two common-link endpoints
open and Gauss invariance imposed at the other four
vertices. Both endpoint charges are retained.

The result is a uniform large-coupling description
for every label \(0\leq n\leq R\sqrt\gamma\), with
arbitrary fixed \(R<\infty\). Each such left-endpoint
charge has exactly one leading matched right-endpoint
multiplicity. A charge-dependent change of frame
separates its boundary energy from the complete
internal oscillator spectrum. Retaining these true
charged ground carriers up to \(6\sqrt\gamma\) leaves
an entire regional complement whose normalized norm
tends to \(q_3=(5-\sqrt{21})/2<1/4\).

This is not a gap above the full regional vacuum:
the retained boundary charges still have energies
approaching zero in lattice units, as V93 proved.
It is a genuine internal-complement estimate with
a growing retained carrier. The general many-boundary
region, interacting gluing on the spatial lattice,
physical time calibration and continuum construction
remain open.

\subsection{The scheduled V95 companion checkpoint}
\label{f15v95:checkpoint}

The current folder was inventoried again. The latest
companions remain SM continuum V72, NS stability V26,
shell mixing V16 and parallel YM V5. Their concluding
mathematical sections were reread, along with the
supplied SM source-selection V31, native stationarity
V34 and exact-action Einstein bridge V24.
All 159 previously protected files were unchanged
by SHA--256.

SM V72 proves a fixed-mode, bounded-time auxiliary
many-copy limit, not a cutoff-uniform physical
dynamics. NS V26 concerns pointwise rank-one kernel
norms, not a YM spectral estimate. Shell V16 controls
a complete measured flat one-loop coefficient but
does not supply the joint noncommuting
nonperturbative remainder. Parallel YM V5 rules out
the specified extensive global-charge tightness
mechanism; it does not provide a substitute gap.
The three supplied bridge notes still distinguish
word-native typing from external completion, an
unpopulated moving-fibre derivative calculation,
and a nonzero first-jet tangent from a nonlinear
root.

None supplies the missing collective YM coercivity.
Their useful methodological constraint is to retain
the actual operator, changing frame, multiplicity
and error norm together. That is implemented below;
no companion's auxiliary limit is identified with
the physical transfer. The next checkpoint is V100,
at the requested lineage rollover.

\subsection{Exact native Fourier reduction with two open endpoints}
\label{f15v95:reduction-section}

Let \(H_0,H_1,H_2\) be the three path holonomies of
V92, of lengths \(1,3,3\), respectively. After Gauss
invariance at the four interior vertices, the space
is \(L^2(G^3)\) with probability product Haar measure.
Set
\[
 g=H_0,\qquad A=g^{-1}H_1,\qquad B=g^{-1}H_2,\qquad
 b_\gamma(A,B)=e^{-\gamma(s(A)+s(B))/2}.
\]
The path kinetic operator before the coordinate
change is \(C_\gamma\otimes C_\gamma^3\otimes C_\gamma^3\).
All its factors and both plaquette multipliers are
the literal native ones of V92.

Put \(d_n=n+1\) and let \(D_n\) be the label-\(n\)
unitary representation. For a fixed left matrix
index \(a\), expand a function in the form
\[
 f_a(g,A,B)=\sqrt{d_n}\sum_{b=1}^{d_n}
                          D_n(g)_{ab}\psi_b(A,B).
\]
Peter--Weyl orthogonality identifies its multiplicity
space with
\(\mathcal H_n=L^2(G^2)\otimes V_n\).
Write \(T_{n,\gamma}\) for its native operator.

\begin{proposition}[The complete charged multiplicity kernel]
\label{f15v95:reduction}
On the two-endpoint-open regional space,
\begin{equation}
 T_{{\rm op},\gamma}
       \simeq\bigoplus_{n\geq0}
                     I_{V_n^{\rm left}}\otimes T_{n,\gamma}.
 \label{eq:f15v95:direct-sum}
\end{equation}
For \(X=(A,B)\), \(Y=(A',B')\), the matrix kernel is
\begin{equation}
 \begin{split}
 K_{n,\gamma}(X,Y)
  ={}&b_\gamma(X)b_\gamma(Y)
       \int_G k_\gamma(h)D_n(h)\\
     &\qquad{}\cdot k_\gamma^{*3}(h^{-1}AA'^{-1})
                    k_\gamma^{*3}(h^{-1}BB'^{-1})\,dH(h).
 \end{split}
 \label{eq:f15v95:kernel}
\end{equation}
It defines a positive self-adjoint compact operator.
The right-endpoint representation on \(\mathcal H_n\) is
\begin{equation}
 (\mathcal R_u\psi)(A,B)
       =D_n(u)\psi(u^{-1}Au,u^{-1}Bu),
 \label{eq:f15v95:right-action}
\end{equation}
and commutes with \(T_{n,\gamma}\).
The \(n=0\) operator is the full \(L^2(G^2)\)
extension of V92, before its residual Gauss
projection.
\end{proposition}

\begin{proof}
The change \((H_0,H_1,H_2)\mapsto(g,A,B)\) preserves
product Haar measure. In the integral for the input
coordinate write \(g'=gh\). The common-link
central kernel is \(k_\gamma(h)\), while the first
length-three path gives
\[
 k_\gamma^{*3}\!\left(
       gAA'^{-1}h^{-1}g^{-1}\right)
       =k_\gamma^{*3}(h^{-1}AA'^{-1}).
\]
The second gives the analogous factor.
Finally \(D_n(gh)=D_n(g)D_n(h)\), giving the matrix
kernel with exactly the displayed orientation.
The left matrix index is unchanged. This proves
the full orthogonal direct sum.

The original native operator is a positive
self-adjoint compact operator, so its Fourier
blocks have these properties; pointwise positivity
of a matrix entry is neither required nor asserted.
The right endpoint sends
\((g,A,B)\) to \((gu,u^{-1}Au,u^{-1}Bu)\),
which gives \eqref{eq:f15v95:right-action}.
The original gauge commutation proves the
intertwining claim. At \(n=0\) the matrix
factor is one and the kinetic kernel is
\eqref{eq:f15v92:kernel}, with the same magnetic
sandwich as V92.
\end{proof}

Let
\[
 \Lambda_\gamma=\|T_{{\rm op},\gamma}\|
                    =\|T_{0,\gamma}\|.
\]
The equality follows from the full positive Perron
vector's gauge invariance, as in V82--V93.
By V92,
\(\Lambda_\gamma\to\ell_0=(q_5q_3)^{3/2}>0\).

\subsection{The moving frame forced by the shared link}
\label{f15v95:frame-section}

Use the complete exponential chart and isometry
\(\mathcal U_\gamma\) of V92, tensoring with \(I_{V_n}\).
Write \(X=(x,y)\in\mathbb R^6\), and let
\(\widetilde T_{n,\gamma}\) denote the resulting
operator, extended by zero outside
\(\mathbb B_\gamma^2\).
The scalar coupled Gaussian \(\mathcal G\) is
\eqref{eq:f15v92:gaussian}.

There is an exact Gaussian square completion.
For \(d_x=x-x'\), \(d_y=y-y'\), put
\[
 m=\frac{d_x+d_y}{5},\qquad \sigma=\frac35.
\]
With \(\varphi_\sigma\) the centered three-dimensional
Gaussian density of covariance \(\sigma I_3\),
\begin{equation}
 \varphi_1(z)\varphi_3(d_x-z)\varphi_3(d_y-z)
      =\Phi_{\mathsf M}(d_x,d_y)\varphi_\sigma(z-m),
 \quad
 \mathsf M=\begin{pmatrix}4&1\\1&4\end{pmatrix}\otimes I_3.
 \label{eq:f15v95:completion}
\end{equation}
For one Cartesian coordinate, its exponent identity is
\[
 \frac{z^2}{2}+\frac{(d_x-z)^2+(d_y-z)^2}{6}
 =\frac{4d_x^2-2d_xd_y+4d_y^2}{30}
            +\frac56\left(z-\frac{d_x+d_y}{5}\right)^2.
\]
The normalizations agree because
\(15(3/5)=1\cdot3\cdot3\).

Define the pointwise unitary moving frame
\begin{equation}
 (V_{n,\gamma}\psi)(x,y)
    =D_n\!\left(\operatorname{Exp}
                   \frac{x+y}{5\sqrt\gamma}\right)\psi(x,y),
 \qquad
 \alpha_{n,\gamma}
     =\exp\left[-\frac{3n(n+2)}{10\gamma}\right].
 \label{eq:f15v95:frame}
\end{equation}
It commutes with the right action, now represented
by simultaneous rotations of \(x,y\) together with
\(D_n(u)\). This follows directly from equivariance
of the group exponential under conjugation.

The frame is not negligible throughout the growing
band: if \(n\) is of order \(\sqrt\gamma\), the
representation can resolve an order-one phase.
Replacing it by the identity is not part of the
argument.

\begin{proposition}[Dimension-independent representation product control]
\label{f15v95:representation}
For \(g,h\in\SU(2)\), in the fundamental matrix norm,
\[
 \|D_n(g)-D_n(h)\|\leq n\|g-h\|.
\]
For \(u,v\) in any fixed bounded subset of the
imaginary quaternions,
\begin{equation}
 \left\|D_n\!\left(\operatorname{Exp}
                         \frac{u+v}{\sqrt\gamma}\right)
       -D_n\!\left(\operatorname{Exp}\frac u{\sqrt\gamma}\right)
        D_n\!\left(\operatorname{Exp}\frac v{\sqrt\gamma}\right)
 \right\|\leq C\,\frac n\gamma.
 \label{eq:f15v95:product}
\end{equation}
The constant does not depend on representation
dimension.
\end{proposition}

\begin{proof}
Realize \(D_n\) on the symmetric subspace of the
\(n\)-fold tensor power of the fundamental
representation. Telescoping the difference of the
two tensor powers gives a sum of \(n\) terms of
norm at most \(\|g-h\|\); restriction cannot
increase the norm. At \(n=0\) the difference is
zero. Taylor expansion of the two fundamental
matrix exponentials gives
\[
 \left\|\operatorname{Exp}((u+v)/\sqrt\gamma)
    -\operatorname{Exp}(u/\sqrt\gamma)
                         \operatorname{Exp}(v/\sqrt\gamma)\right\|
        \leq C/\gamma
\]
uniformly on each bounded set. Combine the two
bounds. This argument does not commute \(u\) and \(v\).
\end{proof}

\begin{proposition}[Central Gaussian average for growing charges]
\label{f15v95:central}
For \(Z\sim N(0,\sigma I_3)\), \(\sigma=3/5\),
\[
 \mathbb E D_n(\operatorname{Exp}(Z/\sqrt\gamma))
                         =a_{n,\gamma}I_{V_n}.
\]
For every fixed \(R<\infty\),
\begin{equation}
 \sup_{0\leq n\leq R\sqrt\gamma}
           |a_{n,\gamma}-\alpha_{n,\gamma}|
                    =O_R(\gamma^{-1}).
 \label{eq:f15v95:central}
\end{equation}
\end{proposition}

\begin{proof}
Rotational invariance and Schur's lemma give the
scalar. Taking the normalized character trace
and using the weights \(-n,-n+2,\ldots,n\) gives
\[
 a_{n,\gamma}=\frac1{n+1}\sum_{j=0}^n
 f_\sigma\!\left(\frac{-n+2j}{\sqrt\gamma}\right),
 \qquad
 f_\sigma(t)=(1-\sigma t^2)e^{-\sigma t^2/2}.
\]
Indeed a radial Gaussian integration gives
\(\mathbb E\cos(t|Z|)
       =(1-\sigma t^2)e^{-\sigma t^2/2}\):
differentiate the ordinary one-dimensional
Gaussian Fourier integral twice and use the
three-dimensional radial density.

These sample points are the midpoints of \(n+1\)
equal intervals filling \([-L,L]\), where
\(L=(n+1)/\sqrt\gamma\), with mesh \(2/\sqrt\gamma\).
The normalized midpoint error is bounded by
\[
 \frac1{6\gamma}
           \sup_{|t|\leq L}|f_\sigma''(t)|.
\]
For \(n\leq R\sqrt\gamma\) and \(\gamma\geq1\)
this is \(O_R(\gamma^{-1})\), independently of
the number of weights. The integral average is
\[
 \frac1{2L}\int_{-L}^L f_\sigma(t)\,dt
                    =e^{-\sigma L^2/2},
\]
because \(f_\sigma(t)\) is the derivative of
\(t e^{-\sigma t^2/2}\).
Finally \((n+1)^2=n(n+2)+1\), so replacing this
last exponential by \(\alpha_{n,\gamma}\) costs
another \(O_R(\gamma^{-1})\).
\end{proof}

\subsection{A full-compact operator limit uniform in the growing carrier}
\label{f15v95:limit-section}

\begin{theorem}[Native charged block with its moving frame]
\label{f15v95:limit}
For every fixed \(R<\infty\),
\begin{equation}
 \delta_R(\gamma):=
 \sup_{\substack{n\in\mathbb N_0\\n\leq R\sqrt\gamma}}
 \left\|\widetilde T_{n,\gamma}
   -\alpha_{n,\gamma}V_{n,\gamma}
                (\mathcal G\otimes I_{V_n})V_{n,\gamma}^*
 \right\|\longrightarrow0.
 \label{eq:f15v95:limit}
\end{equation}
The norm is the operator norm on the indicated
vector-valued \(L^2(\mathbb R^6)\) space.
No representation-dimension factor is hidden
in the convergence assertion.
\end{theorem}

\begin{proof}
In the exact kernel put \(h=\operatorname{Exp}(z/\sqrt\gamma)\)
and use the complete chart for \(A,B,A',B'\).
V92's convolution limit and Haar factors give
the scalar limiting integrand
\[
 e^{-|X|^2/4-|X'|^2/4}
       \varphi_1(z)\varphi_3(d_x-z)\varphi_3(d_y-z),
\]
now multiplied by the exact unitary
\(D_n(\operatorname{Exp}(z/\sqrt\gamma))\).
For \(X,X',z\) in fixed compact sets the scalar
errors converge uniformly, independently of \(n\).
For \(z\) outside a large fixed ball the
Gaussian bound on
\(k_\gamma(\operatorname{Exp}(z/\sqrt\gamma))J_\gamma(z)\),
and the two rescaled convolution supremum bounds
from V92, make the error uniformly small.
The matrix factor has norm one and does not
alter either assertion.

Use \eqref{eq:f15v95:completion} in the limiting
scalar integrand and write \(z=m+Z\). On bounded
\(m,Z\), Proposition~\ref{f15v95:representation} has error
at most \(C_R/\sqrt\gamma\) throughout
\(n\leq R\sqrt\gamma\). Outside a large \(Z\)-ball,
unitarity bounds the matrix difference by two
and the fixed Gaussian tail is uniformly small
for bounded \(m\). It follows that, uniformly on
compact \(X,X'\)-sets and throughout this charge band,
the matrix integral differs by \(o_R(1)\) from
\[
 \Phi_{\mathsf M}(X-X')\,a_{n,\gamma}
       D_n(\operatorname{Exp}(m/\sqrt\gamma)).
\]
Proposition~\ref{f15v95:central} replaces \(a_{n,\gamma}\)
by \(\alpha_{n,\gamma}\). A second application of
\eqref{eq:f15v95:product}, with
\(m=(x+y)/5-(x'+y')/5\), replaces the last matrix by
\(V_{n,\gamma}(X)V_{n,\gamma}(X')^*\).
This proves uniform compact convergence of the
matrix-kernel difference in its matrix operator norm.

For the global step, unitarity bounds the norm
of the exact charged kernel by the scalar \(n=0\)
kernel. V92's complete-chart estimate therefore
bounds it, after scaling and zero extension, by
\[
 C e^{-c(|X|^2+|X'|^2)}
\]
with constants independent of \(n,\gamma\geq1\).
The comparison kernel has a bound of the same
form, since \(0<\alpha_{n,\gamma}\leq1\) and the
frames are unitary. The supremum over the growing
charge band of the matrix-norm difference thus
converges to zero in scalar \(L^2(dX\,dX')\)
by dominated convergence.

If a vector-valued integral kernel \(L(X,X')\)
obeys \(\|L(X,X')\|\leq h(X,X')\), then
\[
 |Lf(X)|\leq\int h(X,X')|f(X')|\,dX',\qquad
                      \|L\|\leq\|h\|_{L^2(dX\,dX')}.
\]
Apply this dimension-independent scalar
majorant estimate to the difference kernel.
It proves \eqref{eq:f15v95:limit}.
Taking a matrix Hilbert--Schmidt norm instead
would insert a factor \(\sqrt{d_n}\); that is
not the estimate used here.
\end{proof}

The compact chart complement is controlled by the
global bound, not discarded by a typical-field
argument. The conclusion is uniform convergence
for fixed \(R\), with no explicit rate or entry
coupling for the full native operator. Only the
central Gaussian average has the separately stated
\(O_R(\gamma^{-1})\) error.

\subsection{True charged ground multiplicities and their internal gap}
\label{f15v95:bands-section}

Recall from V92 the complete scalar spectrum
\[
 \operatorname{spec}_{\ne0}(\mathcal G)
       =\{\ell_0 q_5^{j_+}q_3^{j_-}:j_+,j_-\geq0\},
 \quad
 q_s=\frac{s+2-\sqrt{s^2+4s}}2,
 \quad \ell_0=(q_5q_3)^{3/2}.
\]
It follows from the full Hermite expansion, not
just radial trial states. The underlying Mehler
formula is also recorded in
\href{https://dlmf.nist.gov/18.18.E28}{DLMF 18.18.28}.
The Gaussian ground vector \(\Omega_{\mathcal G}\)
is a rotation singlet, while the next full-space
eigenvalue is \(\ell_0q_3\).

\begin{theorem}[A unique matched ground copy throughout a growing charge band]
\label{f15v95:ground-band}
Fix \(R<\infty\). For all sufficiently large
\(\gamma\), simultaneously for
\(0\leq n\leq R\sqrt\gamma\), the top eigenvalue
\(\mu_{n,\gamma}\) of \(T_{n,\gamma}\) has
multiplicity exactly \(d_n=n+1\).
Its eigenspace carries the right-endpoint
irreducible representation \(V_n\).
Thus the full regional endpoint sector \((n/2,n/2)\)
has exactly one leading multiplicity copy.

With eigenvalues indexed from zero, uniformly
throughout that band,
\begin{align}
 \frac{\mu_{n,\gamma}}{\Lambda_\gamma}
                   -\alpha_{n,\gamma}&\longrightarrow0,
 \label{eq:f15v95:ground-energy}\\
 \frac{\lambda_{d_n}(T_{n,\gamma})}{\mu_{n,\gamma}}
                                      &\longrightarrow q_3.
 \label{eq:f15v95:internal-edge}
\end{align}
The exact leading projection, after the chart
embedding, converges in operator norm to
\begin{equation}
 \pi_{n,\gamma}
   =V_{n,\gamma}
      (|\Omega_{\mathcal G}\rangle\langle\Omega_{\mathcal G}|
                                      \otimes I_{V_n})
                         V_{n,\gamma}^*,
 \label{eq:f15v95:ground-projector}
\end{equation}
uniformly for these growing labels.
\end{theorem}

\begin{proof}
For fixed \(R\), \(\alpha_{n,\gamma}\) has a
strictly positive common lower bound for
\(\gamma\geq1\), \(n\leq R\sqrt\gamma\).
The comparison operator in
\eqref{eq:f15v95:limit} has top eigenvalue
\(a=\alpha_{n,\gamma}\ell_0\), exactly \(d_n\)
times, and next eigenvalue \(a q_3\).
Its gap \(a(1-q_3)\) is therefore uniformly
positive on the specified band.

Let \(\delta=\delta_R(\gamma)\).
Min--max puts the first \(d_n\) exact eigenvalues
within \(\delta\) of \(a\), and all subsequent
ones at most \(a q_3+\delta\).
For \(\delta<a(1-q_3)/2\), the spectral interval
above \(a(1+q_3)/2\) defines an exact projection
\(P_{n,\gamma}\) of rank \(d_n\).
For every unit vector \(\psi\) in its range,
the Rayleigh value of the comparison operator
is at least \(a-2\delta\). Since that operator
equals \(a\) on \(\pi_{n,\gamma}\) and is at most
\(a q_3\) on its complement,
\[
 \|(I-\pi_{n,\gamma})\psi\|^2
                      \leq\frac{2\delta}{a(1-q_3)}.
\]
For two orthogonal projections of the same finite
rank, their norm distance equals the maximal sine
of the principal angles, equivalently the norm
of either cross-complement map. Thus the two
projections converge in norm, uniformly here.

Both projections commute with the right-endpoint
action. The frame commutes with that action and
\(\Omega_{\mathcal G}\) is invariant, so
\(\pi_{n,\gamma}\)'s range is precisely one \(V_n\).
When their norm distance is less than one,
\(\pi_{n,\gamma}\) restricted to
\(\operatorname{ran}P_{n,\gamma}\) is an invertible
intertwiner between the two equal-dimensional
spaces. The exact range is consequently the
same irreducible representation. Schur's lemma
makes \(T_{n,\gamma}\) scalar on it.
The strict separation from the rest proves
the asserted exact multiplicity of its top
eigenvalue.

Weyl's min--max bounds now give
\[
 |\mu_{n,\gamma}-a|\leq\delta,\qquad
 |\lambda_{d_n}(T_{n,\gamma})-a q_3|\leq\delta .
\]
Use the common positive lower bound for \(a\),
and \(\Lambda_\gamma\to\ell_0\), to obtain
\eqref{eq:f15v95:ground-energy}--\eqref{eq:f15v95:internal-edge}.
Reintroducing the left carrier gives one
\(V_n^{\rm left}\otimes V_n^{\rm right}\)
copy, not \(d_n\) additional physical
multiplicity copies.
\end{proof}

For a sequence with \(n_\gamma/\sqrt\gamma\to p\),
the actual native charged ratio therefore obeys
\[
 \frac{\mu_{n_\gamma,\gamma}}{\Lambda_\gamma}
                         \longrightarrow e^{-3p^2/10}.
\]
This identifies the boundary dispersion at the
growing-charge scale. It does \emph{not} prove
\(\gamma(1-\mu_{n,\gamma}/\Lambda_\gamma)
       \to3n(n+2)/10\) for a fixed \(n\):
the error in the full operator theorem is only
\(o(1)\), not \(o(\gamma^{-1})\).
V93's fixed-label two-sided bounds retain
their stated scope.

Other right-endpoint labels, and all additional
multiplicities of the matched label, lie below
the internal edge in
\eqref{eq:f15v95:internal-edge}.
The value \(q_3\), rather than \(q_3^2\), is
appropriate because the right endpoint has
not been projected to a singlet.
For \(n=0\), imposing that extra singlet
constraint recovers V92's different neutral
edge \(q_3^2\).

\subsection{The entire regional complement after retaining the boundary band}
\label{f15v95:complement-section}

Let \(N_\gamma=\lfloor6\sqrt\gamma\rfloor\).
For sufficiently large \(\gamma\), define on the
actual open regional space
\begin{equation}
 \Pi_\gamma
   =\bigoplus_{n=0}^{N_\gamma}
           I_{V_n^{\rm left}}\otimes P_{n,\gamma},
 \qquad
 \operatorname{rank}\Pi_\gamma
       =\sum_{n=0}^{N_\gamma}(n+1)^2
       =\frac{(N_\gamma+1)(N_\gamma+2)(2N_\gamma+3)}6.
 \label{eq:f15v95:retained}
\end{equation}
Here \(P_{n,\gamma}\) is understood on the original
compact multiplicity space through the isometry.
It is an exact spectral projection, not the
Gaussian projection used to approximate it.

\begin{theorem}[Complete fast-complement bound]
\label{f15v95:complement}
The projection \(\Pi_\gamma\) commutes with the
native transfer and both endpoint gauge groups,
contains the actual vacuum, and satisfies
\begin{equation}
 \frac{\|(I-\Pi_\gamma)T_{{\rm op},\gamma}
                          (I-\Pi_\gamma)\|}{\Lambda_\gamma}
                  \longrightarrow q_3<\frac14.
 \label{eq:f15v95:complement}
\end{equation}
In particular this ratio is at most \(1/3\) for
all sufficiently large \(\gamma\).
This controls every state in the discarded
regional complement.
\end{theorem}

\begin{proof}
The exact block decomposition and spectral
projections give commutation. Their right-endpoint
invariance was proved above. The \(n=0\) top
eigenvector is the actual Perron vacuum.
For \(n\leq N_\gamma\), Theorem~\ref{f15v95:ground-band}
bounds the omitted edge, divided by
\(\Lambda_\gamma\), by \(q_3+o(1)\), uniformly
in all these labels, since
\(\alpha_{n,\gamma}\leq1\).

For all remaining labels, V83's domination at
the degree-three left endpoint and V84's exact
Bessel bound give, with \(k=N_\gamma+1\),
\[
 \sup_{n>N_\gamma}\frac{\|T_{n,\gamma}\|}{\Lambda_\gamma}
 \leq r_k(3\gamma)
 \leq\exp\left[-\frac{k(k+2)}{6\gamma+2k+1}\right]
                         \longrightarrow e^{-6}<q_3.
\]
This is an upper bound on every omitted high-charge
multiplicity, not a count of labels.
Taking the supremum over the full direct sum
proves the limsup bound. Conversely, the \(n=0\)
block contributes its first nonvacuum full-space
eigenvalue, whose ratio tends to \(q_3\) by V92's
unprojected operator limit. It is orthogonal to
the retained \(n=0\) ground vector and to all
other left-charge sectors. This proves the
matching liminf.
\end{proof}

The retained dimension is \(O(\gamma^{3/2})\)
before physical boundary gluing. These are
actual boundary carrier dimensions; subsequent
singlet matching must still use V82's normalized
intertwiners. The complete open-block gap still
shrinks, because the retained band contains
the charged modes of V93. No contradiction
or positive physical mass is inferred.

\subsection{Arbitrarily accurate regional truncation by retaining internal bands}
\label{f15v95:precision-section}

The preceding \(1/3\) bound need not be adequate
after dividing by a small glued normalization.
There is an arbitrary-precision version that
keeps more actual internal eigenvectors.

\begin{proposition}[Relative block approximation with controlled growing rank]
\label{f15v95:precision}
For every \(0<\varepsilon<1\), there are constants
\(C_\varepsilon,\gamma_\varepsilon<\infty\) and
exact finite-rank projections
\(\Pi_{\varepsilon,\gamma}\), commuting with
\(T_{{\rm op},\gamma}\) and both endpoint groups,
such that for \(\gamma\geq\gamma_\varepsilon\),
\begin{equation}
 \operatorname{rank}\Pi_{\varepsilon,\gamma}
                   \leq C_\varepsilon\gamma^{3/2},
 \qquad
 \|T_{{\rm op},\gamma}
       -\Pi_{\varepsilon,\gamma}T_{{\rm op},\gamma}
                                 \Pi_{\varepsilon,\gamma}\|
                   \leq\varepsilon\Lambda_\gamma .
 \label{eq:f15v95:precision}
\end{equation}
The constants may depend on the requested accuracy.
No explicit numerical entry coupling is claimed.
\end{proposition}

\begin{proof}
Choose \(R\) so that \(e^{-R^2/6}<\varepsilon/4\).
Choose \(0<\tau<\varepsilon/4\) outside the
discrete nonzero spectrum of \(\mathcal G/\ell_0\).
Its spectral subspace above \(\tau\) has a finite
rank \(m_\tau\), and there is a strictly positive
gap around \(\tau\). This subspace is invariant
under rotations.

For every \(n\leq R\sqrt\gamma\), the uniform
operator approximation, and the positive lower
bound for \(\alpha_{n,\gamma}\), identify an exact
spectral projection of \(T_{n,\gamma}\) above
\(\alpha_{n,\gamma}\ell_0\tau\) of rank
\(m_\tau d_n\), for all sufficiently large
\(\gamma\), simultaneously in \(n\).
It commutes with the right-endpoint group.
The omitted eigenvalues are at most
\(\alpha_{n,\gamma}\ell_0\tau+\delta_R(\gamma)\);
dividing by \(\Lambda_\gamma\to\ell_0\) gives a
limsup at most \(\tau\).
Labels above \(R\sqrt\gamma\) are bounded as in
the previous proof, now with limsup
\(e^{-R^2/6}<\varepsilon/4\).
For sufficiently large \(\gamma\), both bounds
are less than \(\varepsilon\).

Include the left carriers and sum these exact
finite spectral projections. Its rank is
\[
 m_\tau\sum_{n=0}^{\lfloor R\sqrt\gamma\rfloor}(n+1)^2
                            =O_\varepsilon(\gamma^{3/2}).
\]
The projection commutes with the transfer, so
the norm of the removed part is exactly the
norm on its orthogonal complement. This proves
\eqref{eq:f15v95:precision}.
\end{proof}

This improves the regional internal reduction,
not V91's theorem for the whole spatial lattice.
It retains exact native spectral subspaces but
does not yet provide closed formulas for their
mixed-plaquette matrix elements. Multiplication
by an external interface weight need not commute
with these projections.

\subsection{Direction, verification and preserved scope}
\label{f15v95:direction}

The missing regional complement is now controlled
for this specified two-endpoint-open block,
uniformly over the necessary growing charges.
The common-link frame gives a nontrivial charged
boundary energy rather than a replacement by
independent neutral plaquettes. Higher internal
bands can be retained to meet any fixed relative
block error.

The next assembly must compute or bound the
actual interface multiplier between these
retained native spaces, paying the omitted
parts under the correct normalization. Blocks
with additional exposed vertices require their
own boundary analysis. Overlapping interfaces
and collective neutral excitations on the
growing three-dimensional graph are not handled
by the present fixed-region theorem.
The independent physical clock and nontrivial
four-dimensional continuum reconstruction also
remain missing. A gapped internal complement
is not the full Yang--Mills mass gap.

The checker
\begin{center}\small
\path{scripts/verify_ym_f15_v95_growing_charge_block.py}
\end{center}
tests the literal quaternion kernel orientation,
right-endpoint covariance, exact conditional
Gaussian square completion, noncommuting frame
jets, radial Gaussian character coefficients,
the normalized midpoint error, projection-angle
budgets, carrier ranks and the high-label
exponent. Finite positive matrices in these
checks are algebra fixtures, not native Wilson
laws. The full-compact uniform convergence
and spectral results are analytic proofs,
not conclusions from sampling.

The checker passes 3,699 new exact-rational
fixtures together with the inherited V94 chain.
It has 7,665 bytes and SHA--256
\begin{center}\scriptsize\ttfamily
F092FB82177A08049FCC2BC6549A9B2ADB17EDF7BBEC7A7270614CFB37D99F71.
\end{center}

The predecessor V94 has 2,144,131 bytes and SHA--256
\begin{center}\scriptsize\ttfamily
D7B750042307EA96DE9DA8D29977E7FDFE4B33D9B0BDAE1BDFC64755CAF89B77.
\end{center}
Its body is retained byte for byte; removing
this section and restoring the title version
recovers V94. Only the active main source is
replaced. Archives, companions, supplied papers
and prior diagnostic scripts remain unchanged.
The companion checkpoint is now V95, with V100
next. Only \path{.tex} files are stored in
\path{latex_notes}; builds go to
\path{artifacts/ym_f15_v95_texcheck/}.

\begin{firewall}
V95 proves a growing-charge native operator
limit, a unique matched charged ground
multiplicity, and complete internal-complement
control on the specified two-endpoint-open
two-plaquette block. It also gives an
arbitrary-precision relative regional truncation.
It does not establish the many-block interacting
gap, handle arbitrary exposed boundaries,
derive a physical clock, or construct the
interacting continuum. The full Yang--Mills
mass gap remains unproved.
\end{firewall}

% V96 append begins; widen section numbers from 100 onward in the TOC.
\makeatletter
% The temporary ! parameter character survives both auxiliary-file writes.
\addtocontents{toc}{\protect\catcode33=6\protect\relax
 \protect\renewcommand*{\protect\l@section}[2]{%
 \protect\ifnum\protect\c@tocdepth>\protect\z@
 \protect\addpenalty\protect\@secpenalty
 \protect\addvspace{1.0em \protect\@plus\protect\p@}%
 \protect\setlength\protect\@tempdima{2.2em}%
 \protect\begingroup
 \protect\parindent\protect\z@
 \protect\rightskip\protect\@pnumwidth
 \protect\parfillskip-\protect\@pnumwidth
 \protect\leavevmode\protect\bfseries
 \protect\advance\protect\leftskip\protect\@tempdima
 \protect\hskip-\protect\leftskip
 !1\protect\nobreak\protect\hfil
 \protect\nobreak\protect\hb@xt@\protect\@pnumwidth{%
 \protect\hss !2\protect\kern-\protect\p@\protect\kern\protect\p@}%
 \protect\par\protect\endgroup\protect\fi}%
 \protect\catcode33=12\protect\relax
 \protect\renewcommand*{\protect\l@subsection}{%
 \protect\@dottedtocline{2}{2.2em}{3.1em}}}
\makeatother
\section{V96: complete Gauss-matched gluing across one actual mixed plaquette}
\label{f15v96:context}

V94 controlled a forward interface compression,
but not the entire physical transfer.
V95 controlled all omitted internal states of
a particular two-endpoint-open regional block.
The next step is to combine these results on an
actual link partition, keeping the internal
multiplicities through the magnetic sandwich.

We attach a three-link path to the common-link
endpoints of V95's two-plaquette region and restore
the square formed by that path and the common link.
The union is a three-page plaquette book with ten
links, embedded in the cubic spatial lattice.
For this specified region, the complete interacting
physical transfer has a relative-error finite-rank
approximation of rank \(O_\varepsilon(\sqrt\gamma)\).
The rank counts Gauss-matched multiplicities,
not the dimensions of uncontracted carriers.
An exact fusion formula integrates the mixed
holonomy while retaining native internal eigenfunctions.

We also check the full glued operator directly.
Its complete compact-chart limit is a nine-dimensional
coupled Gaussian, and its physical spectral ratio
tends to \(q_3^2<1/20\). The Gaussian convergence
method and Mehler diagonalization are inherited
from V92, not claimed as new methods.
The new interface conclusion is a full-space
approximation, rather than a gap of a selected
trial compression. None of these fixed-region
results is a volume-uniform theorem on the
three-dimensional spatial lattice.

\subsection{A literal cubic-lattice cut}
\label{f15v96:geometry-section}

Take the common edge from \(v=(0,0,0)\) to
\(w=(1,0,0)\). Choose the three length-three paths
through the respective pairs of vertices
\[
 ((0,1,0),(1,1,0)),\qquad
 ((0,-1,0),(1,-1,0)),\qquad
 ((0,0,1),(1,0,1)).
\]
Together with the common edge these are three
unit squares. Their interiors are disjoint.
There are ten links and eight vertices.

Let \(\mathsf A\) consist of the common edge
and the first two paths; it is exactly V95's
seven-link block. Let \(\mathsf B\) be the third
path. The link sets are disjoint and their
only shared vertices are \(v,w\).
At the six other vertices impose Gauss invariance,
and then match both shared endpoints to singlets.
All vertices of the union satisfy Gauss invariance.
This is a finite graph with that specified
boundary condition; it does not represent the
open boundary of an arbitrary region in a larger
lattice.

Write the four path holonomies as \(g,H_1,H_2,h\)
and put
\[
 A=g^{-1}H_1,\qquad B=g^{-1}H_2,\qquad C=g^{-1}h.
\]
The fully physical space is
\begin{equation}
 \mathcal H_{\rm book}
       =L^2(G^3,dH(A)dH(B)dH(C))^{\operatorname{Ad}G}.
 \label{eq:f15v96:physical}
\end{equation}
The path region \(\mathsf B\), before endpoint
matching, has native transfer \(C_\gamma^3\).
Its Peter--Weyl sector
\(V_n^{\rm left}\otimes V_n^{\rm right}\) occurs
once and has eigenvalue \(r_n(\gamma)^3\);
its Perron eigenvalue is one.

Let \(T_{\rm cut,\gamma}\) be the physical cut
transfer. Its Perron eigenvalue is
\(\Lambda_{\mathsf A,\gamma}\), the regional
Perron eigenvalue in V95. Restore the third
plaquette by
\begin{equation}
 M_\gamma f=e^{-\gamma s(C)/2}f,\qquad
 T_{\rm book,\gamma}
             =M_\gamma T_{\rm cut,\gamma}M_\gamma.
 \label{eq:f15v96:sandwich}
\end{equation}
The mixed Wilson loop may equivalently be
written \(gh^{-1}\), since its trace is the
trace of \(C\). There are no additional mixed
plaquettes in this graph.

V94 applies with path lengths \(1,3\), hence
\(S=10\). In particular, with its explicit
constant
\[
 \kappa_{10}=\frac{e^{-1/2}}{384\pi^3\,10^{3/2}},
\]
the \emph{actual full physical} Perron value obeys
\begin{equation}
 \kappa_{10}
 \leq\frac{\lambda_0(T_{\rm book,\gamma})}
                  {\Lambda_{\mathsf A,\gamma}}\leq1
 \qquad(\gamma>0).
 \label{eq:f15v96:normalizer}
\end{equation}
This denominator control will pay the complete
discarded space, not just a charged trial family.

\subsection{The exact rotational content of a retained internal band}
\label{f15v96:content-section}

Use V95's notation
\[
 \ell_0=(q_5q_3)^{3/2},\qquad
 \alpha_{n,\gamma}=e^{-3n(n+2)/(10\gamma)},\qquad
 \delta_R(\gamma)\longrightarrow0 .
\]
Fix \(R<\infty\), and choose \(0<\tau<1\) outside
the nonzero spectrum of \(\mathcal G/\ell_0\).
Let
\begin{equation}
 W_\tau=\operatorname{Ran}
             {\bf1}_{(\ell_0\tau,\infty)}(\mathcal G).
 \label{eq:f15v96:gaussian-band}
\end{equation}
It is finite dimensional and invariant under
simultaneous rotations.
For \(n\leq R\sqrt\gamma\), let \(E_{n,\gamma,\tau}\)
be the exact spectral subspace of V95's
\(T_{n,\gamma}\) above
\(\alpha_{n,\gamma}\ell_0\tau\).

\begin{proposition}[Equivariant identification of the whole retained band]
\label{f15v96:equivariant}
For all sufficiently large \(\gamma\), uniformly
for \(0\leq n\leq R\sqrt\gamma\),
\[
 E_{n,\gamma,\tau}\simeq W_\tau\otimes V_n
\]
as right-endpoint representations. This identifies
representation content, not equality of the
native eigenvalues with Gaussian eigenvalues.
\end{proposition}

\begin{proof}
After V95's complete-chart embedding, the comparison
operator is
\[
 A_{n,\gamma}
       =\alpha_{n,\gamma}V_{n,\gamma}
                   (\mathcal G\otimes I_{V_n})V_{n,\gamma}^*.
\]
For fixed \(R\), the coefficients \(\alpha_{n,\gamma}\)
have a common positive lower bound. Since
\(\tau>0\) lies in a spectral gap, the threshold
\(\alpha_{n,\gamma}\ell_0\tau\) is separated
from this comparison spectrum uniformly in \(n\).

For completeness, put a rectangular contour
around the part of the spectrum above the threshold,
with right edge \(2\ell_0\) and horizontal edges
at imaginary parts \(\pm\ell_0\).
Its left edge is the threshold.
The contour lengths are uniformly bounded,
and their distances to the comparison spectra
have a positive common lower bound.
When \(\delta_R(\gamma)\) is smaller than half
that bound, the exact resolvents on the contours
are also uniformly bounded. The resolvent identity
and the contour integral for the projections give
\[
 \left\|{\bf1}_{(\alpha_{n,\gamma}\ell_0\tau,\infty)}
                (\widetilde T_{n,\gamma})
       -V_{n,\gamma}
             ({\bf1}_{(\ell_0\tau,\infty)}(\mathcal G)
                                      \otimes I_{V_n})
                         V_{n,\gamma}^*\right\|
                    \leq C_{R,\tau}\delta_R(\gamma).
\]
The contours may depend on \(n,\gamma\); the
displayed bounds do not. No rank factor enters
the resolvent estimate.

Both projections commute with the right action.
For norm distance less than one, restriction of
one projection to the range of the other is an
invertible equivariant map, as in V95.
Finally the moving frame commutes with that
action. Its comparison range consequently has
representation \(W_\tau\otimes V_n\), proving
the assertion.
\end{proof}

The rotational content of \(W_\tau\) can be
counted explicitly. The degree-\(r\) excitations
of a three-dimensional isotropic oscillator
transform as \(\operatorname{Sym}^r(V_2)\),
where \(V_2\) has spin one.
Indeed its three creation operators transform
as a vector, and their commuting degree-\(r\)
products give this symmetric power.
Thus
\begin{equation}
 W_\tau\simeq
 \bigoplus_{\substack{r,s\geq0\\q_5^r q_3^s>\tau}}
          \operatorname{Sym}^r(V_2)
                   \otimes\operatorname{Sym}^s(V_2)
       \simeq\bigoplus_{j\geq0}m_{\tau,j}V_{2j}.
 \label{eq:f15v96:content}
\end{equation}
Only finitely many \(m_{\tau,j}\) are nonzero.
They are integer multiplicities, not carrier
dimensions. A finite formula is
\begin{equation}
 m_{\tau,j}=
 \sum_{\substack{r,s\geq0\\q_5^r q_3^s>\tau}}
 \ \sum_{a=0}^{\lfloor r/2\rfloor}
 \ \sum_{b=0}^{\lfloor s/2\rfloor}
 {\bf1}\{|(r-2a)-(s-2b)|\leq j\leq(r-2a)+(s-2b)\}.
 \label{eq:f15v96:multiplicity}
\end{equation}
To verify it, homogeneous polynomials of degree
\(r\) decompose into \(|x|^{2a}\) times harmonic
homogeneous polynomials of degree \(r-2a\).
The latter carry spin \(r-2a\).
This decomposition follows by successively
subtracting Laplacian traces; its dimensions
sum to \(\binom{r+2}{2}\).
Apply this to both symmetric powers and then
use the multiplicity-one Clebsch--Gordan rule.

\subsection{Gauss matching reduces the retained rank}
\label{f15v96:rank-section}

Let \(N_\gamma=\lfloor R\sqrt\gamma\rfloor\).
On the open \(\mathsf A\) space retain
\(I_{V_n^{\rm left}}\otimes E_{n,\gamma,\tau}\)
for \(n\leq N_\gamma\), and retain nothing at
higher left charge. Tensor with the entire
path space of \(\mathsf B\) and impose both
endpoint Gauss constraints.
Denote the resulting physical projection by
\(Q_{\gamma,R,\tau}\).
It commutes with \(T_{\rm cut,\gamma}\), but
need not commute with \(M_\gamma\).

\begin{theorem}[The complete matched rank, without carrier overcounting]
\label{f15v96:rank}
For sufficiently large \(\gamma\), the number
of retained physical copies at charge \(n\) is
\[
 \nu_{\tau,n}=\sum_{0\leq j\leq n}m_{\tau,j},
\]
and
\begin{equation}
 \operatorname{rank}Q_{\gamma,R,\tau}
   =\sum_{n=0}^{N_\gamma}\nu_{\tau,n}
   =\sum_{j\geq0}m_{\tau,j}\max\{0,N_\gamma-j+1\}.
 \label{eq:f15v96:rank}
\end{equation}
In particular it is \(O_{R,\tau}(\sqrt\gamma)\).
For \(N_\gamma\) beyond the largest occurring \(j\),
it is exactly
\[
 (N_\gamma+1)\sum_jm_{\tau,j}-\sum_jj\,m_{\tau,j}.
\]
\end{theorem}

\begin{proof}
Left-endpoint matching forces the path label
to equal the \(\mathsf A\) left label \(n\).
That singlet occurs once. The path's other
endpoint has the same label \(n\), again with
no multiplicity. Consequently the remaining
physical multiplicity is
\[
 \dim\operatorname{Hom}_G(V_n,E_{n,\gamma,\tau}).
\]
By Proposition~\ref{f15v96:equivariant}, this is
the multiplicity of \(V_n\) in
\(\bigoplus_jm_{\tau,j}(V_{2j}\otimes V_n)\).
The Clebsch--Gordan interval contains label \(n\)
exactly when \(j\leq n\), and then once.
This proves the formula for \(\nu_{\tau,n}\).
Summing and interchanging the two finite sums
gives \eqref{eq:f15v96:rank}.
No factor \(n+1\) survives either normalized
endpoint singlet contraction.
\end{proof}

This is stronger than the unglued
\(O(\gamma^{3/2})\) carrier count in V95.
It is not obtained by deleting actual internal
states: every multiplicity in the selected
internal spectral band that can match the
path has been counted.

\subsection{A relative approximation of the entire interacting transfer}
\label{f15v96:approximation-section}

Set \(k=N_\gamma+1\) and define the native error
budget
\begin{equation}
 e_{\gamma,R,\tau}
   =\max\left\{
       \frac{\ell_0\tau}{\Lambda_{\mathsf A,\gamma}},
       r_k(3\gamma)r_k(\gamma)^3
          \right\}.
 \label{eq:f15v96:error-budget}
\end{equation}
The subspace identification requires sufficiently
large coupling, but this expression is in the
actual native normalizations.

\begin{theorem}[Full physical forward truncation]
\label{f15v96:approximation}
For those sufficiently large couplings, put
\[
 T_{\gamma,R,\tau}^{\rm ret}
   =M_\gamma Q_{\gamma,R,\tau}T_{\rm cut,\gamma}
                         Q_{\gamma,R,\tau}M_\gamma .
\]
Then
\begin{align}
 0&\leq T_{\rm book,\gamma}-T_{\gamma,R,\tau}^{\rm ret},
 \notag\\
 \|T_{\rm book,\gamma}-T_{\gamma,R,\tau}^{\rm ret}\|
     &\leq \Lambda_{\mathsf A,\gamma}\,e_{\gamma,R,\tau},
 \label{eq:f15v96:absolute-error}\\
 \frac{\|T_{\rm book,\gamma}-T_{\gamma,R,\tau}^{\rm ret}\|}
             {\lambda_0(T_{\rm book,\gamma})}
     &\leq \frac{e_{\gamma,R,\tau}}{\kappa_{10}}.
 \label{eq:f15v96:relative-error}
\end{align}
For every \(0<\varepsilon<1\), suitable fixed
\(R,\tau\) give relative error at most
\(\varepsilon\), eventually in \(\gamma\), with
\[
 \operatorname{rank}T_{\gamma,R,\tau}^{\rm ret}
                              =O_\varepsilon(\sqrt\gamma).
\]
This approximates the full physical operator
on \(\mathcal H_{\rm book}\), not its compression
on a vacuum-times-character trial space.
\end{theorem}

\begin{proof}
For \(n\leq N_\gamma\), the omitted exact
\(\mathsf A\) eigenvalues are at most
\(\alpha_{n,\gamma}\ell_0\tau\).
Their cut products acquire the path factor
\(r_n(\gamma)^3\leq1\), so their normalized
sizes are at most
\(\ell_0\tau/\Lambda_{\mathsf A,\gamma}\).
There is no additional \(\delta_R\) in this
tail estimate: the projection is defined by
an exact spectral threshold. The V95 error
was needed to identify its rank and representation
content, not to redefine its cutoff.

For \(n>N_\gamma\), V83 bounds the whole
\(\mathsf A\) charge sector by
\(\Lambda_{\mathsf A,\gamma}r_n(3\gamma)\).
The path factor is exactly \(r_n(\gamma)^3\).
Monotonicity in \(n\) therefore bounds every
remaining physical multiplicity by the
second term in \eqref{eq:f15v96:error-budget}.
The direct-sum norm is a supremum, not a
sum over charges.

Since \(Q_{\gamma,R,\tau}\) is a commuting
cut projection, its removed cut part is positive
with norm at most
\(\Lambda_{\mathsf A,\gamma}e_{\gamma,R,\tau}\).
Sandwiching by the contraction \(M_\gamma\)
preserves positivity and cannot increase this
norm. This proves \eqref{eq:f15v96:absolute-error};
\eqref{eq:f15v96:normalizer} proves
\eqref{eq:f15v96:relative-error}.

V84's Bessel bound gives
\[
 r_k(3\gamma)r_k(\gamma)^3
 \leq\exp\left[
 -\frac{k(k+2)}{6\gamma+2k+1}
 -\frac{3k(k+2)}{2\gamma+2k+1}\right]
                 \longrightarrow e^{-5R^2/3}.
\]
Also \(\Lambda_{\mathsf A,\gamma}\to\ell_0\).
Choose \(R\) with
\(e^{-5R^2/3}<\varepsilon\kappa_{10}/4\)
and choose a permissible
\(0<\tau<\varepsilon\kappa_{10}/4\).
Then eventually \(e_{\gamma,R,\tau}
                  <\varepsilon\kappa_{10}\).
The rank bound follows from
Theorem~\ref{f15v96:rank} and bounded multiplication.
At finite coupling \(M_\gamma\) is invertible,
and the retained cut eigenvalues are positive,
so the retained rank is in fact unchanged
by this magnetic sandwich.
\end{proof}

This argument does not assert
\(Q M_\gamma=M_\gamma Q\), or replace the
forward sandwich by \(Q T_{\rm book,\gamma}Q\).
The output vectors \(M_\gamma\Phi\) include all
magnetically generated charges and internal
components. Their finite-dimensional span
is different from the original cut subspace.
The entry coupling is still qualitative;
the proof does not give a numerical finite-coupling
certificate without further estimates.

\subsection{Native matrix elements with the mixed holonomy integrated exactly}
\label{f15v96:matrix-section}

Choose orthogonal isometric embeddings of
\(V_n\) into the retained right-\(V_n\)
isotypic part of \(E_{n,\gamma,\tau}\) which
diagonalize its multiplicity operator.
Write their columns as the matrix-valued
functions \(\Psi_{n,a}(A,B)\), where
\(1\leq a\leq\nu_{\tau,n}\).
Thus
\begin{align}
 \int_{G^2}\Psi_{n,a}^*\Psi_{n,b}\,dH(A)dH(B)
      &=\delta_{ab}I_{V_n},
 \label{eq:f15v96:psi-normalization}\\
 D_n(u)\Psi_{n,a}(u^{-1}Au,u^{-1}Bu)
      &=\Psi_{n,a}(A,B)D_n(u).
 \label{eq:f15v96:psi-covariance}
\end{align}
Let their native \(\mathsf A\) eigenvalues be
\(\lambda^{\mathsf A}_{n,a,\gamma}\).

\begin{proposition}[Normalized physical basis and exact forward Gram matrix]
\label{f15v96:gram}
An orthonormal cut basis for
\(\operatorname{Ran}Q_{\gamma,R,\tau}\) is
\begin{equation}
 \Phi_{n,a}(A,B,C)
    =\Tr\bigl(D_n(C)^*\Psi_{n,a}(A,B)\bigr).
 \label{eq:f15v96:phi}
\end{equation}
Its normalized cut eigenvalues are
\[
 c_{n,a,\gamma}
    =\frac{\lambda^{\mathsf A}_{n,a,\gamma}}
                 {\Lambda_{\mathsf A,\gamma}}r_n(\gamma)^3.
\]
The nonzero spectrum of
\(T_{\gamma,R,\tau}^{\rm ret}/\Lambda_{\mathsf A,\gamma}\)
is exactly that of the positive finite matrix
\begin{equation}
 H_{(n,a),(m,b)}
   =\sqrt{c_{n,a,\gamma}c_{m,b,\gamma}}
      \int_{G^3} e^{-\gamma s(C)}
                     \overline{\Phi_{n,a}}\Phi_{m,b}\,dH^3 .
 \label{eq:f15v96:gram}
\end{equation}
\end{proposition}

\begin{proof}
The covariance identity makes
\eqref{eq:f15v96:phi} invariant under simultaneous
conjugation of \(A,B,C\).
Peter--Weyl orthogonality in \(C\) gives zero
between different labels and, at label \(n\),
\[
 \int_G
   \overline{\Tr(D_n(C)^*X)}\Tr(D_n(C)^*Y)\,dH(C)
                        =\frac1{d_n}\Tr(X^*Y).
\]
Together with \eqref{eq:f15v96:psi-normalization},
this proves orthonormality.
Equivalently, contracting the two endpoint
singlets in the normalized regional matrix
bases gives precisely \eqref{eq:f15v96:phi}:
the two Peter--Weyl factors \(\sqrt{d_n}\)
cancel the two singlet factors \(1/\sqrt{d_n}\).

The cut transfer is the regional tensor product,
so its eigenvalue is the stated product.
If \(Z\) sends the finite basis vector indexed
by \((n,a)\) to
\(\sqrt{c_{n,a,\gamma}}M_\gamma\Phi_{n,a}\),
then \(ZZ^*\) is the retained normalized
operator and \(Z^*Z\) is \eqref{eq:f15v96:gram}.
Their nonzero eigenvalues agree, with multiplicity.
\end{proof}

The integral over \(C\) is an explicit finite
fusion expression. Let \(P_\ell^{n,\bar m}\)
be the orthogonal projection onto label \(\ell\)
in \(V_n\otimes\overline{V_m}\).
For indices \(i,j\) of \(V_n\) and \(k,l\) of
\(V_m\),
\begin{equation}
 \begin{split}
 &\int_G e^{-\gamma s(C)}
             D_n(C)_{ij}\overline{D_m(C)_{kl}}\,dH(C)\\
 &\quad=c_0(\gamma)
       \sum_{\substack{|n-m|\leq\ell\leq n+m\\
                        \ell\equiv n+m\ ({\rm mod}\ 2)}}
           r_\ell(\gamma)
                 (P_\ell^{n,\bar m})_{(i,k),(j,l)} .
 \end{split}
 \label{eq:f15v96:fusion}
\end{equation}
Here \(c_0(\gamma)=\int e^{-\gamma s(C)}\,dH(C)\).
To prove the identity, decompose the tensor
representation by Clebsch--Gordan. The central
integral on each irreducible summand is
\(c_0(\gamma)r_\ell(\gamma)I\), as in V83.
There is no additional dimension factor.

In \eqref{eq:f15v96:gram}, multiply the right
side of \eqref{eq:f15v96:fusion} by
\(\overline{\Psi_{n,a}(A,B)_{ij}}\,
                    \Psi_{m,b}(A,B)_{kl}\),
sum the four indices, and integrate over \(A,B\).
This integrates the mixed holonomy exactly.
The displayed paired index order matters;
reshuffling a positive tensor is not a license
to assert an unrelated matrix order.
Positivity of \(H\) follows from its original
Haar Gram representation.

For example, at \(n=m=1\),
\(P_0=|\operatorname{vec}I\rangle
      \langle\operatorname{vec}I|/2\) and
\(P_2=I-P_0\). The magnetic Gram entry,
before the two square-root cut weights, is
\begin{equation}
 c_0(\gamma)\left[
 (1-r_2(\gamma))\delta_{ab}
 +r_2(\gamma)\int_{G^2}
          \overline{\Tr\Psi_{1,a}}\Tr\Psi_{1,b}\,dH^2
                    \right].
 \label{eq:f15v96:fundamental}
\end{equation}
The trace overlap retains actual internal
wavefunctions. Replacing it by one or by a
freely chosen channel weight would alter the
native transfer.

If \(h_0\geq h_1\geq\cdots\) are the eigenvalues
of \(H\), with zero padding, the positive
remainder gives
\[
 h_j\leq
 \frac{\lambda_j(T_{\rm book,\gamma})}
                        {\Lambda_{\mathsf A,\gamma}}
       \leq h_j+e_{\gamma,R,\tau}.
\]
In particular an evaluated matrix satisfying
\begin{equation}
 h_1+e_{\gamma,R,\tau}<h_0
 \quad\Longrightarrow\quad
 \frac{\lambda_1(T_{\rm book,\gamma})}
                    {\lambda_0(T_{\rm book,\gamma})}
       \leq\frac{h_1+e_{\gamma,R,\tau}}{h_0}<1
 \label{eq:f15v96:certificate}
\end{equation}
certifies the \emph{full} graph's ratio.
This is the earlier one-sided spectral comparison
applied to the newly controlled full gluing.
The internal \(\Psi_{n,a}\) have not been numerically
enclosed here. Formula \eqref{eq:f15v96:fusion}
is not a claim that every entry of \(H\) has
already been explicitly evaluated.

\subsection{A direct check of the full three-page operator}
\label{f15v96:full-limit-section}

There is an independent analytic way to check
this specified glued graph.
Put \(A_1=A,A_2=B,A_3=C\) and
\(b_3=\exp[-\gamma\sum_{i=1}^3s(A_i)/2]\).
The exact full-space kernel before residual
conjugation projection is
\begin{equation}
 K_{3,\gamma}(A;A')
   =b_3(A)b_3(A')\int_G k_\gamma(h)
         \prod_{i=1}^3
              k_\gamma^{*3}(h^{-1}A_iA_i'^{-1})\,dH(h).
 \label{eq:f15v96:book-kernel}
\end{equation}
This follows by the same Haar-preserving path
reduction as V92, now retaining the third path.
Equivalently its kinetic operator is
\((C_\gamma^3)^{\otimes3}D_\gamma\), with
\(D_\gamma\) simultaneous left convolution.
Its commuting factors are positive. The
magnetic sandwich is positive, compact and
has a strictly positive continuous kernel.

In the complete scaled chart \(X=(x_1,x_2,x_3)\in
\mathbb R^9\), let
\[
 \mathsf M_3=(3I_3+\boldsymbol1\boldsymbol1^{\mathsf T})
                       \otimes I_3,\qquad
 \boldsymbol1=(1,1,1)^{\mathsf T},
\]
and define \(\mathcal G_3\) by kernel
\begin{equation}
 \mathcal G_3(X,X')
       =e^{-|X|^2/4}\Phi_{\mathsf M_3}(X-X')
                                      e^{-|X'|^2/4}.
 \label{eq:f15v96:book-gaussian}
\end{equation}
The covariance has off-diagonal \(I_3\) blocks.
These record the common edge and cannot be
dropped when restoring the interface.

\begin{theorem}[Full interacting fixed-region spectral check]
\label{f15v96:book-gap}
The complete-chart embedding of
\(T_{\rm book,\gamma}\), with zero extension,
converges in Hilbert--Schmidt norm to
\(\mathcal G_3\), both before and after the
simultaneous-conjugation projection.
Writing \(q_s=(s+2-\sqrt{s^2+4s})/2\), one has
\begin{align}
 \lambda_0(T_{\rm book,\gamma})
       &\longrightarrow (q_6q_3^2)^{3/2},
 \label{eq:f15v96:book-top}\\
 \frac{\lambda_0(T_{\rm book,\gamma})}
                    {\Lambda_{\mathsf A,\gamma}}
       &\longrightarrow
                       \left(\frac{q_6q_3}{q_5}\right)^{3/2}>0,
 \label{eq:f15v96:book-normalizer}\\
 \frac{\lambda_1(T_{\rm book,\gamma})}
                    {\lambda_0(T_{\rm book,\gamma})}
       &\longrightarrow q_3^2<\frac1{20}.
 \label{eq:f15v96:book-edge}
\end{align}
There is also a constant \(q_{\rm book}<1\)
bounding this physical ratio for every
\(\gamma>0\) on this fixed graph.
No explicit optimal value or finite-coupling
entry threshold is asserted.
\end{theorem}

\begin{proof}
Apply V92's native fixed-step convolution limit
to the three factors in \eqref{eq:f15v96:book-kernel}.
The common-link integration becomes the law of
\[
 (Z+U_1,Z+U_2,Z+U_3),
\]
where the four vectors are independent centered
three-dimensional Gaussians. The covariance of
\(Z\) is \(I_3\), and each \(U_i\) has covariance
\(3I_3\). Their joint covariance
is exactly \(\mathsf M_3\).
For one Cartesian coordinate,
\(M_3=3I_3+\boldsymbol1\boldsymbol1^{\mathsf T}\)
has determinant \(54\) and inverse
\(I_3/3-\boldsymbol1\boldsymbol1^{\mathsf T}/18\).
For differences \(d_1,d_2,d_3\), the exact
square completion is
\[
 \frac{z^2}{2}+\sum_{i=1}^3\frac{(d_i-z)^2}{6}
 =\frac{\sum_i d_i^2}{6}
       -\frac{(\sum_i d_i)^2}{36}
       +\left(z-\frac{\sum_i d_i}{6}\right)^2 .
\]
The conditional common-link variance is \(1/2\),
and the normalizations agree since \(54/2=27\).

The three convolution factors are bounded by
\(C\gamma^{9/2}\) in product and
\(\int k_\gamma=1\). The six Haar square-root
factors contribute at most \(C\gamma^{-9/2}\).
Globally in the complete chart the magnetic
factor is at most
\(\exp[-\sum_i|x_i|^2/\pi^2]\).
The rescaled kernels, including the zero
complement, are therefore bounded by
\(Ce^{-c(|X|^2+|X'|^2)}\).
The native local limit and this integrable
square majorant prove Hilbert--Schmidt convergence.
All arguments retain the compact chart complement.
Equivariance of the chart allows restriction
to the entire physical invariant space.

An orthogonal change of the three path indices
diagonalizes \(M_3\) into \(6,3,3\).
It preserves the magnetic quadratic form and
commutes with simultaneous spatial rotations.
The complete one-dimensional Hermite
diagonalization of V92, based on the
\href{https://dlmf.nist.gov/18.18.E28}{Mehler formula},
therefore gives full-space eigenvalues
\[
 (q_6q_3^2)^{3/2}\,q_6^{r}q_3^{s+t},
                   \qquad r,s,t\in\mathbb N_0,
\]
with oscillator multiplicities.
Its ground state is a rotation singlet.
The degree-one space consists of three spin-one
representations, so has no singlet.
Every nonconstant physical state has degree
at least two, and \(q_6<q_3\); its ratio is
at most \(q_3^2\).
A radial degree-two excitation in either
relative normal mode is a nonzero singlet
attaining this value. Completeness, rather
than a variational lower bound alone, identifies
the physical second eigenvalue.

Min--max continuity under operator-norm
convergence proves
\eqref{eq:f15v96:book-top} and
\eqref{eq:f15v96:book-edge}.
V95's
\(\Lambda_{\mathsf A,\gamma}\to(q_5q_3)^{3/2}\)
then proves \eqref{eq:f15v96:book-normalizer}.
The elementary inequality
\(21>(229/50)^2\) verifies \(q_3^2<1/20\).

Finally the original compact kernel is
Hilbert--Schmidt continuous on each compact
coupling interval including zero.
At zero it is the constant projection.
At each positive coupling its full Perron
eigenvector is simple, strictly positive and
gauge invariant, so the physical ratio is
strictly below one. Continuity on a compact
interval bounds its maximum strictly below
one; the limiting ratio controls the remaining
large-coupling interval. This gives
\(q_{\rm book}<1\) for the fixed graph.
\end{proof}

Thus this actual magnetic gluing confines the
retained boundary band as part of a complete
interacting fixed-region operator. A shrinking
cut charged gap did not force a shrinking
fully Gauss-matched glued gap.
Neither direction can be inferred from a cut
spectrum alone. This check does not establish
stability under an extensive family of mixed
plaquettes, nor does it turn the regional
lattice-time gap into a positive physical
continuum mass.

\subsection{What remains after this interface test}
\label{f15v96:scope-section}

The distinction between the V94 compression
and the full physical operator has been resolved
for this specified link partition: all discarded
states have a relative bound, all allowed
retained multiplicities are counted, and the
mixed holonomy has an exact native fusion formula.
The internal eigenfunctions remain part of the
data. They have not been replaced by freely
adjustable channel parameters.

The next obstruction is geometric and collective.
When these plaquette blocks sit inside the
full spatial lattice, more than their two
common-link endpoints are exposed to other
links. Those extra boundary charges require
their own analysis. Simultaneous mixed plaquettes
share variables; their forward sandwiches and
normalizations cannot be multiplied as if
independent. A uniform gap for the neutral
collective modes of the growing graph, an
independently calibrated physical time, and
the nontrivial interacting four-dimensional
continuum construction remain unproved.

The current source and prior result ledger were
checked for overlap before this continuation.
The last companion checkpoint remains V95;
the next is V100, at the requested archive/restart.
No companion or archived body is modified.

The checker
\begin{center}\small
\path{scripts/verify_ym_f15_v96_full_native_gluing.py}
\end{center}
verifies the literal cubic-lattice geometry,
independent character and harmonic descriptions
of oscillator rotation content, Gauss-matched
multiplicities and ranks, the fundamental fusion
tensor with its paired index order, the shared-link
Gaussian square completion, and the native
combined charge-tail exponent. Small positive
matrix sandwiches are algebra fixtures, not
substituted Wilson measures. The full-compact
limits and all-coupling fixed-region assertion
are proved analytically above.

All 2,284 new exact-rational fixtures and
the inherited V95 chain pass.
The checker has 9,043 bytes and SHA--256
\begin{center}\scriptsize\ttfamily
57CA4E0A55CAE1A5C1394452BA30BCC52AED908DB128F94460F7D7095474A9CC.
\end{center}

The predecessor V95 has 2,172,363 bytes and SHA--256
\begin{center}\scriptsize\ttfamily
343291585D2E23F3377CB9B8FBB4D0A6A01D2A7F9D454578C8006164EDD56E13.
\end{center}
Its body is retained byte for byte; removing
this V96 append, including its TOC formatting,
and restoring the title version recovers V95.
Only the active main source is
replaced. All 160 previously protected files
remain unchanged. Only \path{.tex} files belong
in \path{latex_notes}; compilation products go
to \path{artifacts/ym_f15_v96_texcheck/}.

\begin{firewall}
V96 proves a complete Gauss-matched finite-rank
approximation of one specified full interacting
three-page region, its exact native mixed-holonomy
fusion matrix, and a fixed-region physical
spectral check after restoring that interface.
It does not prove a many-region mass gap,
treat all exposed boundaries, derive physical
time scaling, or construct interacting continuum
Yang--Mills theory. The full mass gap remains
unproved.
\end{firewall}

\section{V97: the fully open cube and its coupled boundary heat operator}
\label{f15v97:context}

V96's two-endpoint gluing theorem did not retain
the other vertices exposed when a region is placed
inside the spatial lattice. We now treat a literal
elementary cube with \emph{all eight vertices open}.
Throughout this note \(G=\SU(2)\), as in the
preceding local-transfer calculations.
No vertex Gauss constraint is imposed within this
regional Hilbert space. All six native plaquette
multipliers and all twelve native link kernels
are retained.

The main result is a uniform operator-norm
factorization for boundary labels growing as
\(\sqrt\gamma\), with error \(O_R(\gamma^{-1/2})\).
The fast factor has fifteen real coordinates.
The slow factor is a positive, generally
noncommuting spin operator determined by the
inverse reduced graph Laplacian. It is not a
product of independently chosen scalar charge
weights. The eighth, root vertex acts on the
retained tensor product and is not removed.

This gives a nearly invariant native boundary
subspace and a bound on its \emph{entire}
orthogonal complement. It also gives arbitrary
relative-error finite-rank approximations on the
full open cube. The dimension grows with coupling,
and no estimate uniform in a growing spatial
volume is obtained. In particular, the internal
fast bound is not a mass gap above the full
open cube's vacuum, still less the continuum
Yang--Mills mass gap.

\subsection{A Haar-exact tree coordinate system with no Gauss projection}
\label{f15v97:coordinates-section}

Let \(V=\{0,1\}^3\), with root \(o=(0,0,0)\),
and orient the twelve edges in the positive
coordinate directions. Use the seven-edge tree
consisting of all four \(x\)-edges, the two
\(y\)-edges at \(x=0\), and the \(z\)-edge at
\(x=y=0\). Its five oriented chords are
\[
 a=U_y(1,0,0),\quad b=U_y(1,0,1),\quad
 c=U_z(0,1,0),\quad d=U_z(1,0,0),\quad
 e=U_z(1,1,0)
\]
in the tree-fixed representative.

Each original configuration has unique coordinates
\(g_v\in G\) for \(v\ne o\), together with \(X\in G^5\).
Set \(g_o=I\) and write \(X=(a,b,c,d,e)\). They satisfy
\begin{equation}
 U_f=g_{s(f)}X_f g_{t(f)}^{-1},
 \qquad X_f=I\text{ on the tree}.
 \label{eq:f15v97:coordinates}
\end{equation}
Here chord \(X_f\) means the corresponding entry
of \(X\). Solve successively along the rooted
tree for the \(g_v\), then solve for each chord.
Every step is a Haar translation or inversion,
so product probability Haar measure becomes
\(dH^{7}(g)dH^5(X)\).
These are coordinates on the \emph{full}
\(L^2(G^{12})\), not a gauge quotient.

The six plaquette words, up to conjugation and
orientation which leave \(s\) unchanged, are
\begin{equation}
 a,\quad b,\quad d,\quad ec^{-1},\quad
 c,\quad aeb^{-1}d^{-1}.
 \label{eq:f15v97:faces}
\end{equation}
Put
\[
 S_\Box(X)=s(a)+s(b)+s(d)+s(ec^{-1})+s(c)
                              +s(aeb^{-1}d^{-1}),\qquad
 b_{\Box,\gamma}=e^{-\gamma S_\Box/2}.
\]
The native open transfer is
\(T_{\Box,\gamma}=M_b C_\gamma^{\otimes12}M_b\).
The one-link normalization is the inherited
\(k_\gamma(U)=e^{-\gamma s(U)}/c_0(\gamma)\).

\begin{proposition}[A global compact magnetic bound]
\label{f15v97:magnetic}
On all of \(G^5\),
\begin{equation}
 \sum_{X_i\in\{a,b,c,d,e\}}s(X_i)\leq3S_\Box(X).
 \label{eq:f15v97:coercivity}
\end{equation}
In particular the simultaneous zero of the
six plaquette actions in these coordinates is
\(a=b=c=d=e=I\).
\end{proposition}

\begin{proof}
For unit quaternions, \(s(U)=\frac12|I-U|^2\).
The chord identity \(e=(ec^{-1})c\), the triangle
inequality and \((x+y)^2\leq2x^2+2y^2\) give
\(s(e)\leq2s(ec^{-1})+2s(c)\).
Add the other four chord actions. Each term
on the resulting right side occurs in
\(3S_\Box\) with at least its required coefficient.
The zero-set assertion follows at once, without
a small-field restriction.
\end{proof}

\subsection{The exact seven-label kernel and the retained root action}
\label{f15v97:kernel-section}

Decompose the seven \(g_v\) variables in their
normalized Peter--Weyl bases. Write
\[
 \mathbf n=(n_v)_{v\ne o},\qquad
 \mathcal V_{\mathbf n}=\bigotimes_{v\ne o}V_{n_v},
 \qquad d_{\mathbf n}=\prod_{v\ne o}(n_v+1).
\]
For each fixed collection of left indices, the
multiplicity space is
\(\mathcal H_{\mathbf n}=L^2(G^5)\otimes\mathcal V_{\mathbf n}\).
Let \(\mathsf D_{\mathbf n}(h)=\bigotimes_vD_{n_v}(h_v)\)
and set \(h_o=I\).

\begin{proposition}[Full-boundary native reduction]
\label{f15v97:kernel}
There is an exact decomposition
\[
 T_{\Box,\gamma}\simeq
 \bigoplus_{\mathbf n}
        I_{\mathcal V_{\mathbf n}^{\rm left}}\otimes
                                      T_{\mathbf n,\gamma},
\]
where
\begin{equation}
 \begin{split}
 K_{\mathbf n,\gamma}(X,Y)
   ={}&b_{\Box,\gamma}(X)b_{\Box,\gamma}(Y)\\
   &{}\cdot\int_{G^7}
       \prod_{f\in E}
          k_\gamma(h_{s(f)}^{-1}X_fh_{t(f)}Y_f^{-1})
                        \mathsf D_{\mathbf n}(h)\,dH^7(h).
 \end{split}
 \label{eq:f15v97:kernel}
\end{equation}
Each block is positive, self-adjoint and compact.
The root action on its multiplicity space is
\begin{equation}
 (\mathcal R_u\psi)(X)
       =\left(\bigotimes_{v\ne o}D_{n_v}(u)\right)
                                \psi(u^{-1}Xu),
 \label{eq:f15v97:root-action}
\end{equation}
where all five chords are conjugated simultaneously.
It commutes with the block transfer.
\end{proposition}

\begin{proof}
In the input tree variables set \(g'_v=g_vh_v\).
For each edge the original kinetic argument is,
up to conjugation,
\(X_fh_{t(f)}Y_f^{-1}h_{s(f)}^{-1}\).
Centrality moves the final factor to the front,
giving \eqref{eq:f15v97:kernel}.
The Fourier factors obey
\(D_{n_v}(g_vh_v)=D_{n_v}(g_v)D_{n_v}(h_v)\);
hence all left indices remain unchanged.
Positivity and compactness follow from the
original native operator, not from entrywise
positivity of this matrix kernel.

A nonroot gauge transformation is left
translation of its own \(g_v\).
A root transformation, with the inverse
parameter convention, sends
\(g_v\) to \(g_vu\) and \(X\) to \(u^{-1}Xu\).
This proves \eqref{eq:f15v97:root-action}.
Original gauge commutation gives the last claim.
Thus the root charge still acts on the tensor
carrier and the internal coordinates. No eighth
independent Peter--Weyl factor should be inserted.
\end{proof}

Let \(\Lambda_\Box(\gamma)=\|T_{\Box,\gamma}\|\).
The original full kernel is strictly positive
on a compact configuration space. Its simple
positive Perron vector is invariant under all
eight gauge groups, so
\begin{equation}
 \Lambda_\Box(\gamma)=\|T_{\mathbf0,\gamma}\|.
 \label{eq:f15v97:perron}
\end{equation}
This does not project the other regional states
to gauge singlets.

\subsection{The graph covariance and the internal quadratic form}
\label{f15v97:graph-section}

Order the nonroot vertices lexicographically.
Let \(\mathsf B\) be the \(12\)-by-\(7\)
root-deleted incidence matrix, with row
\(z_{t(f)}-z_{s(f)}\), and let \(\mathsf J\)
inject the five chord coordinates in the edge
space. Define
\begin{equation}
 \mathsf L=\mathsf B^{\mathsf T}\mathsf B,\qquad
 \Sigma=\mathsf L^{-1},\qquad
 \mathsf F=-\Sigma\mathsf B^{\mathsf T}\mathsf J .
 \label{eq:f15v97:laplacian}
\end{equation}
The positive definiteness of \(\mathsf L\)
follows from connectedness and the fixed root.

Let \(\mathsf Z\) be the fundamental cycle matrix
specified by
\(\mathsf Z\mathsf B=0\), \(\mathsf Z\mathsf J=I_5\).
Its tree columns are
\(-\mathsf B_{\rm chord}\mathsf B_{\rm tree}^{-1}\).
The chord covariance is
\begin{equation}
 \mathsf M=\mathsf Z\mathsf Z^{\mathsf T}
   =\begin{pmatrix}
 4&0&-1&1&-2\\
 0&4&1&-1&2\\
 -1&1&4&1&3\\
 1&-1&1&4&1\\
 -2&2&3&1&6
 \end{pmatrix}.
 \label{eq:f15v97:covariance}
\end{equation}
Both \(\det\mathsf L\) and \(\det\mathsf M\)
are \(384\).

For one Cartesian component of \(d=x-x'\),
there is the exact Gaussian identity
\begin{equation}
 \Phi_{I_{12}}(\mathsf Bz+\mathsf Jd)
       =\Phi_{\mathsf M}(d)\Phi_\Sigma(z-\mathsf Fd).
 \label{eq:f15v97:completion}
\end{equation}
For the group Lie algebra take three copies
of this identity, equivalently tensor every
covariance with \(I_3\).
To prove it, the map
\((z,d)\mapsto\mathsf Bz+\mathsf Jd\)
has determinant of absolute value one:
its tree block has determinant \(\pm1\)
and its chord block is triangular with \(I_5\).
The covariance of \(d=\mathsf Z(\mathsf Bz+\mathsf Jd)\)
is \(\mathsf M\). Completing the \(z\) square
gives precision \(\mathsf L\) and mean
\(\mathsf Fd\), with no missing determinant factor.

Linearizing the six words
\eqref{eq:f15v97:faces} gives
\begin{equation}
 \mathsf A=
 \begin{pmatrix}
 1&0&0&0&0\\0&1&0&0&0\\0&0&0&1&0\\
 0&0&-1&0&1\\0&0&1&0&0\\1&-1&0&-1&1
 \end{pmatrix},
 \qquad \mathsf H=\mathsf A^{\mathsf T}\mathsf A.
 \label{eq:f15v97:hessian}
\end{equation}
Here \(\det\mathsf H=6>0\).
The positive matrix
\(\mathsf M^{1/2}\mathsf H\mathsf M^{1/2}\)
has eigenvalues
\begin{equation}
 4,\quad4,\quad4,\quad6,\quad6 .
 \label{eq:f15v97:frequencies}
\end{equation}
One exact verification is
\[
 (\mathsf M\mathsf H-4I)(\mathsf M\mathsf H-6I)=0,
 \qquad \Tr(\mathsf M\mathsf H)=24.
\]
Similarity to a positive symmetric matrix gives
diagonalizability; the trace determines the
multiplicities. These are full internal
frequencies, not a restriction to radial trial
functions.

\subsection{A dimension-independent noncommuting Gaussian estimate}
\label{f15v97:wick-section}

The conditional Gaussian in
\eqref{eq:f15v97:completion} is correlated across
vertices. Its matrix average must not be
replaced by seven scalar Casimir factors.

\begin{proposition}[Gaussian exponential versus a quadratic operator]
\label{f15v97:wick}
Let \(S_1,\ldots,S_d\) be self-adjoint matrices
on any finite-dimensional Hilbert space.
Assume \(\|S_i\|\leq L\), \(L\geq1\), and
\(\|[S_i,S_j]\|\leq\eta\).
For independent standard real Gaussians \(\xi_i\),
put \(Q=\frac12\sum_iS_i^2\) and \(a=dL^2/2\).
Then
\begin{equation}
 \left\|\mathbb E e^{\,i\sum_i\xi_iS_i}-e^{-Q}\right\|
               \leq4\eta(a^2+a)e^a.
 \label{eq:f15v97:wick}
\end{equation}
There is no matrix-dimension factor.
\end{proposition}

\begin{proof}
The exponential series may be integrated term
by term in norm because
\(\mathbb E e^{L\sum_i|\xi_i|}<\infty\).
Odd moments vanish. The Gaussian moment
generating function \(e^{|t|^2/2}\), differentiated
at zero, expresses each even moment as the
sum over \((2k-1)!!\) pairings.

For each pairing and each of its \(d^k\)
assignments of indices, move the two occurrences
of every paired index next to one another,
ordering pairs by their first occurrence.
At most \((2k)^2\) adjacent interchanges suffice.
Each interchange changes the matrix word by
norm at most \(\eta L^{2k-2}\).
After the reordering, summing the assignments
gives \((\sum_iS_i^2)^k\).
Consequently
\[
 \left\|\mathbb E\Big(\sum_i\xi_iS_i\Big)^{2k}
       -(2k-1)!!\Big(\sum_iS_i^2\Big)^k\right\|
 \leq(2k-1)!!\,d^k(2k)^2\eta L^{2k-2}.
\]
Since \((2k-1)!!/(2k)!=1/(2^k k!)\) and \(L\geq1\),
summing the exponential coefficients bounds
the difference by
\[
 4\eta\sum_{k\geq1}\frac{k^2a^k}{k!}
                         =4\eta(a^2+a)e^a .
\]
This also bounds the absolutely convergent
series of errors and completes the justification.
\end{proof}

Let \(J_{v,\alpha}\), \(\alpha=1,2,3\), be the
Hermitian generators on factor \(V_{n_v}\). Our convention is
\[
 D_{n_v}(\operatorname{Exp}(t e_\alpha))=e^{itJ_{v,\alpha}}.
\]
They satisfy
\[
 \|J_{v,\alpha}\|=n_v,\qquad
 \sum_\alpha J_{v,\alpha}^2=n_v(n_v+2)I,\qquad
 \|[J_{v,\alpha},J_{v,\beta}]\|\leq2n_v.
\]
Different vertex factors commute.
Define the coupled boundary operator
\begin{equation}
 Q_{\mathbf n,\gamma}
   =\frac1{2\gamma}\sum_{v,w\ne o}\Sigma_{vw}
                      \sum_{\alpha=1}^3J_{v,\alpha}J_{w,\alpha},
 \qquad W_{\mathbf n,\gamma}=e^{-Q_{\mathbf n,\gamma}}.
 \label{eq:f15v97:boundary-heat}
\end{equation}
This is positive in the sense \(Q_{\mathbf n,\gamma}\geq0\):
factor the positive covariance \(\Sigma\) and
write the quadratic expression as a sum of
squares of self-adjoint linear combinations.
It commutes with the diagonal root representation.
For \(\max_v n_v\leq R\sqrt\gamma\),
\[
 \|Q_{\mathbf n,\gamma}\|
       \leq \frac32R^2\sum_{v,w}|\Sigma_{vw}|=:K_R,
 \qquad e^{-K_R}I\leq W_{\mathbf n,\gamma}\leq I.
\]

\begin{proposition}[The correlated boundary average]
\label{f15v97:average}
If \(Z\) is the centered \(21\)-dimensional
Gaussian of covariance \(\Sigma\otimes I_3\),
then, uniformly for \(\max_v n_v\leq R\sqrt\gamma\),
\begin{equation}
 \left\|\mathbb E\bigotimes_{v\ne o}
       D_{n_v}(\operatorname{Exp}(Z_v/\sqrt\gamma))
                     -W_{\mathbf n,\gamma}\right\|
                         \leq C_R\gamma^{-1/2}.
 \label{eq:f15v97:average}
\end{equation}
\end{proposition}

\begin{proof}
Write \(Z=(\Sigma^{1/2}\otimes I_3)\xi\).
The tensor product of exponentials is a
single exponential, since different vertex
factors commute. Its exponent has the form
\(i\sum_{i=1}^{21}\xi_iS_i\), with \(S_i\)
linear combinations of \(J_{v,\alpha}/\sqrt\gamma\).
Their norms have a common bound depending
only on \(R,\Sigma\). Their commutators have
norm at most \(C_R/\sqrt\gamma\), by the
displayed generator bound. Their quadratic
sum is exactly \(Q_{\mathbf n,\gamma}\).
Apply Proposition~\ref{f15v97:wick} with
\(d=21\) and \(L\) increased to at least one.
\end{proof}

The off-diagonal entries of \(\Sigma\) produce
operators \(J_v\cdot J_w\).
The root tensor representation is generally
reducible, so Schur's lemma does not turn
\(W_{\mathbf n,\gamma}\) into one scalar.
Its multiplicity action is part of the
boundary dynamics, not a dispensable dimension
correction.

\subsection{A full-compact kernel estimate with a rate}
\label{f15v97:rate-section}

Use the complete exponential chart
\(q_\gamma(x)=\operatorname{Exp}(x/\sqrt\gamma)\),
\(|x|<\pi\sqrt\gamma\), and its inherited density
\[
 J_\gamma(x)=j_0\gamma^{-3/2}
       \left(\frac{\sin(|x|/\sqrt\gamma)}
                         {|x|/\sqrt\gamma}\right)^2,
 \qquad j_0=(2\pi^2)^{-1}.
\]
Apply this chart to the five output chords
\(x\), the five input chords \(x'\), and all
seven integration variables \(z\).
Let \(\mathcal U_\gamma\) denote the
Haar-square-root isometry on the five chords,
extended by zero to \(L^2(\mathbb R^{15})\).

Set \(\mathcal W=(x,x',z)\in\mathbb R^{51}\).
The exact scalar integrand, after this scaling
but before the \(z\) integration, is
\[
 A_\gamma P_\gamma(\mathcal W)e^{-E_\gamma(\mathcal W)},
 \qquad
 A_\gamma=(j_0\gamma^{-3/2}/c_0(\gamma))^{12}.
\]
Here \(P_\gamma\) is the product of the ten
chord sinc factors and the seven squared
vertex sinc factors, and
\[
 E_\gamma
   =\frac\gamma2S_\Box(q_\gamma(x))
      +\frac\gamma2S_\Box(q_\gamma(x'))
      +\gamma\sum_{f\in E}
       s(q_\gamma(z_{s(f)})^{-1}X_f
                    q_\gamma(z_{t(f)})Y_f^{-1}).
\]
The root \(z_o\) is zero. Its quadratic limit is
\begin{equation}
 E_0=\frac14x^{\mathsf T}(\mathsf H\otimes I_3)x
       +\frac14{x'}^{\mathsf T}(\mathsf H\otimes I_3)x'
       +\frac12|(\mathsf B\otimes I_3)z
                         +(\mathsf J\otimes I_3)(x-x')|^2.
 \label{eq:f15v97:joint-quadratic}
\end{equation}

\begin{proposition}[Global domination and scalar remainder]
\label{f15v97:scalar-rate}
Throughout the joint complete chart, for
\(\gamma\geq1\),
\begin{align}
 E_\gamma(\mathcal W),\,E_0(\mathcal W)
       &\geq\frac{|\mathcal W|^2}{6\pi^2},
 \label{eq:f15v97:global-domination}\\
 |E_\gamma-E_0|&\leq24\gamma^{-1/2}|\mathcal W|^3,
 \label{eq:f15v97:action-remainder}\\
 0\leq P_\gamma\leq1,\qquad
 1-P_\gamma&\leq\frac{|\mathcal W|^2}{3\gamma}.
 \label{eq:f15v97:haar-remainder}
\end{align}
Also
\(A_\gamma=(2\pi)^{-18}(1+O(\gamma^{-1}))\),
with a uniform bound for \(\gamma\geq1\).
Replacing this scalar integrand by
\((2\pi)^{-18}e^{-E_0}\), and extending the
latter to all \(\mathbb R^{51}\), costs
\(O(\gamma^{-1/2})\) after \(z\) integration
in scalar \(L^2(dx\,dx')\).
The same estimate holds when both integrands
are multiplied by any common matrix factor
of operator norm at most one.
\end{proposition}

\begin{proof}
The magnetic bound
\eqref{eq:f15v97:coercivity} and
\(1-\cos t\geq2t^2/\pi^2\) on \([0,\pi]\)
give the lower bound
\((|x|^2+|x'|^2)/(3\pi^2)\) from the two
magnetic ends.

For a tree path of length \(d_v\),
the chord triangle inequality gives
\[
 s(h_v)\leq d_v\sum_{f\text{ on the path}}s(h_{s(f)}^{-1}h_{t(f)}).
\]
The sum of the seven tree depths is \(12\).
Summing over vertices thus bounds
\(\sum_vs(h_v)\) by twelve times the tree-edge
action. The latter is part of the kinetic
action, and supplies \(|z|^2/(6\pi^2)\).
This proves the bound for \(E_\gamma\).
Taking its pointwise large-\(\gamma\) limit
proves the same bound for \(E_0\) on all of
Euclidean space.

To bound the remainder, expand each word
along the line \(t\mathcal W/\sqrt\gamma\),
\(0\leq t\leq1\), in fundamental unitary matrices.
Every word has at most four distinct coordinate
factors. Its third derivative has norm at most
\((2|\mathcal W|/\sqrt\gamma)^3\), because all
other factors are unitary.
The real normalized trace has norm at most one.
There are twelve kinetic words of coefficient
\(\gamma\), and twelve magnetic words of
coefficient \(\gamma/2\).
The total Taylor remainder is consequently
at most
\((12+6)8|\mathcal W|^3/(6\sqrt\gamma)\),
which is \eqref{eq:f15v97:action-remainder}.
The first derivatives vanish and the second
derivatives give exactly
\eqref{eq:f15v97:joint-quadratic}.
This argument is valid throughout the chart,
not merely in a fixed small ball.

On \([0,\pi]\), sinc is between zero and one,
and \(1-\operatorname{sinc}t\leq t^2/6\).
For its square the bound is \(t^2/3\).
Use \(1-\prod_i a_i\leq\sum_i(1-a_i)\)
for factors in \([0,1]\) to get
\eqref{eq:f15v97:haar-remainder}.
The normalization estimate follows from
V92's proved
\(c_0(\gamma)=j_0(2\pi)^{3/2}\gamma^{-3/2}
                          (1+O(\gamma^{-1}))\);
continuity extends its uniform bounds down to one.

Finally
\[
 |e^{-E_\gamma}-e^{-E_0}|
 \leq|E_\gamma-E_0|e^{-\min(E_\gamma,E_0)}
 \leq24\gamma^{-1/2}|\mathcal W|^3
                          e^{-|\mathcal W|^2/(6\pi^2)}.
\]
Together with the other two remainders this
gives an integrable polynomial times a
Gaussian majorant. Integrating \(z\) leaves
a polynomial times a Gaussian in \(x,x'\),
whose \(L^2\) norm is \(O(\gamma^{-1/2})\).
The excluded Euclidean chart complement has
at least one coordinate norm \(\pi\sqrt\gamma\);
the same Gaussian bound makes its contribution
exponentially small.
Multiplication by a common norm-one matrix
does not change these scalar estimates.
\end{proof}

\subsection{The uniform coupled boundary factorization}
\label{f15v97:factor-section}

On \(L^2(\mathbb R^{15})\) define
\[
 b_0(x)=e^{-x^{\mathsf T}(\mathsf H\otimes I_3)x/4},
 \qquad
 \mathcal G_\Box(x,x')
         =b_0(x)\Phi_{\mathsf M\otimes I_3}(x-x')b_0(x').
\]
Define the charge-dependent unitary multiplication
operator
\begin{equation}
 V_{\mathbf n,\gamma}(x)
    =\bigotimes_{v\ne o}
          D_{n_v}\!\left(\operatorname{Exp}
                       \frac{(\mathsf Fx)_v}{\sqrt\gamma}\right).
 \label{eq:f15v97:frame}
\end{equation}
It commutes with the root action, by equivariance
of the exponential and the real linearity of
\(\mathsf F\) in the spatial vectors.

\begin{theorem}[All-boundary native factorization with a growing-label rate]
\label{f15v97:factor}
For every fixed \(R<\infty\), there is a finite
constant \(C_R\) such that, for \(\gamma\geq1\),
\begin{equation}
 \sup_{\max_vn_v\leq R\sqrt\gamma}
 \left\|\mathcal U_\gamma T_{\mathbf n,\gamma}\mathcal U_\gamma^*
  -V_{\mathbf n,\gamma}
       (\mathcal G_\Box\otimes W_{\mathbf n,\gamma})
                                      V_{\mathbf n,\gamma}^*\right\|
                      \leq C_R\gamma^{-1/2}.
 \label{eq:f15v97:factor}
\end{equation}
The norm is the operator norm on each indicated
vector-valued Euclidean \(L^2\) space.
The constant does not depend on its representation
dimension. It does depend on the fixed cube and
on \(R\).
\end{theorem}

\begin{proof}
Proposition~\ref{f15v97:scalar-rate} replaces
the scalar part of the exact matrix kernel
by its Gaussian integrand, while keeping
\(\mathsf D_{\mathbf n}(q_\gamma(z))\) intact.
The scalar \(L^2\) error is already
\(O(\gamma^{-1/2})\), uniformly in all labels.

Use \eqref{eq:f15v97:completion} and put
\(z=\mathsf F(x-x')+Z\).
For anti-Hermitian matrices \(A,B\), Duhamel's
formula gives the global unitary estimate
\[
 \|e^{A+B}-e^Ae^B\|\leq\tfrac12\|[A,B]\|.
\]
For example write the difference as a single
Duhamel integral; its integrand contains
\([B,e^{tA}]\), whose norm is at most
\(t\|[A,B]\|\), and integrate \(t\) from zero
to one. In label \(n_v\), this bounds the
product defect for Lie algebra vectors \(u,v\)
by \(n_v|u||v|/\gamma\).
Telescoping the seven tensor factors sums
these bounds, rather than multiplying by
their dimensions.

The conditional mean is linear in \(x,x'\).
Gaussian integration of this product defect
therefore has a polynomial times Gaussian
bound of size \(O_R(\gamma^{-1/2})\).
Proposition~\ref{f15v97:average} replaces the
remaining centered matrix average by
\(W_{\mathbf n,\gamma}\) with the same order
of error. A second application of the unitary
product estimate replaces the mean factor by
\(V_{\mathbf n,\gamma}(x)V_{\mathbf n,\gamma}(x')^*\).
The order of factors is preserved: the
conditional average first gives the mean
matrix times \(W_{\mathbf n,\gamma}\).

There is one additional noncommuting point.
The desired kernel has
\(V(x)W_{\mathbf n,\gamma}V(x')^*\),
not \(V(x)V(x')^*W_{\mathbf n,\gamma}\).
For \(A(x')=\sum_{v,\alpha}(\mathsf Fx')_{v,\alpha}
                                      J_{v,\alpha}/\sqrt\gamma\),
the frame is \(V(x')=e^{iA(x')}\).
The generator commutator bounds imply
\[
 \|[Q_{\mathbf n,\gamma},A(x')]\|
                    \leq C_R|x'|\gamma^{-1/2}.
\]
Indeed each quadratic term contributes one
commutator of size \(O_R(\gamma^{-1/2})\)
times one bounded scaled generator.
The integral formulas for the commutators
with \(e^{-Q}\) and \(e^{iA}\), using
\(Q\geq0\), give
\[
 \|[W_{\mathbf n,\gamma},V(x')^*]\|
                    \leq C_R|x'|\gamma^{-1/2}.
\]
This pays the necessary reordering.
It is not exact commutation.
Again its polynomial factor is integrable
against the magnetic Gaussian ends.

All resulting matrix-norm errors are bounded
by scalar kernels of \(L^2(dx\,dx')\) norm
at most \(C_R/\sqrt\gamma\).
A vector-valued kernel bounded in matrix
operator norm by \(h(x,x')\) defines an
operator of norm at most \(\|h\|_{L^2}\),
as proved in V95. This yields
\eqref{eq:f15v97:factor}.
No matrix Hilbert--Schmidt norm, which would
insert a growing carrier dimension, is used.
\end{proof}

The extra commutator in this proof is essential.
A correlated boundary average does not, in
general, commute exactly with its moving frame.

\subsection{The complete fast spectrum and the correct boundary subspace}
\label{f15v97:fast-section}

The linear change \(x=\mathsf M^{1/2}y\),
with its \(L^2\) Jacobian, turns the covariance
into the identity and the magnetic matrix into
\(\mathsf M^{1/2}\mathsf H\mathsf M^{1/2}\).
An orthogonal change in the five path indices
then uses \eqref{eq:f15v97:frequencies}.
Both changes commute with simultaneous spatial
rotations. The complete Hermite expansion,
as in V92 and the
\href{https://dlmf.nist.gov/18.18.E28}{Mehler formula},
gives
\begin{equation}
 \operatorname{spec}_{\ne0}(\mathcal G_\Box)
    =\{\ell_\Box q_4^r q_6^s:r,s\geq0\},
 \qquad
 \ell_\Box=(q_4^3q_6^2)^{3/2},
 \quad q_t=\frac{t+2-\sqrt{t^2+4t}}2,
 \label{eq:f15v97:fast-spectrum}
\end{equation}
with all oscillator multiplicities.
The normalized ground vector \(\Omega_\Box\)
is a rotation singlet. The first full-space
excited eigenvalue is \(\ell_\Box q_4\), where
\(q_4=3-\sqrt8<1/5\).
The root is open, so degree-one spin-one
excitations must not be discarded.

At \(\mathbf n=\mathbf0\) the frame and boundary
heat factor are identities. Equations
\eqref{eq:f15v97:perron} and
\eqref{eq:f15v97:factor} therefore imply
\begin{equation}
 \Lambda_\Box(\gamma)=\ell_\Box+O(\gamma^{-1/2}).
 \label{eq:f15v97:top-limit}
\end{equation}

Let \(\chi_\gamma\) be the indicator of the
five-coordinate complete chart, and put
\(a_\gamma=\|\chi_\gamma\Omega_\Box\|^2\).
For \(\mathbf n\ne\mathbf0\), define an exact
isometry \(J_{\mathbf n,\gamma}:\mathcal V_{\mathbf n}
                                      \to\mathcal H_{\mathbf n}\) by
\begin{equation}
 \mathcal U_\gamma J_{\mathbf n,\gamma}v
      =a_\gamma^{-1/2}\chi_\gamma(x)\Omega_\Box(x)
                                      V_{\mathbf n,\gamma}(x)v .
 \label{eq:f15v97:ground-isometry}
\end{equation}
The scalar Gaussian tail makes
\(1-a_\gamma\) exponentially small in \(\gamma\).
For \(\mathbf n=\mathbf0\), instead choose the
actual normalized native Perron vector as
\(J_{\mathbf0,\gamma}1\).
This small correction ensures exact vacuum
inclusion, without assuming that the Gaussian
ground vector is the native one.

For \(N_\gamma=\lfloor6\sqrt\gamma\rfloor\), let
\begin{equation}
 \Pi_\gamma
   =\bigoplus_{\substack{\mathbf n\\0\leq n_v\leq N_\gamma}}
       I_{\mathcal V_{\mathbf n}^{\rm left}}\otimes
                   J_{\mathbf n,\gamma}J_{\mathbf n,\gamma}^* .
 \label{eq:f15v97:projection}
\end{equation}

\begin{theorem}[All eight boundary actions retained, with a complete fast bound]
\label{f15v97:complement}
The projection \(\Pi_\gamma\) commutes with all
eight regional gauge actions and contains the
actual vacuum. Its rank is
\begin{equation}
 \operatorname{rank}\Pi_\gamma
    =\left[\frac{(N_\gamma+1)(N_\gamma+2)(2N_\gamma+3)}6\right]^7
                              =O(\gamma^{21/2}).
 \label{eq:f15v97:rank}
\end{equation}
Under the displayed exact isometries, uniformly
over the retained label box,
\begin{align}
 \|J_{\mathbf n,\gamma}^*T_{\mathbf n,\gamma}
       J_{\mathbf n,\gamma}
              -\ell_\Box W_{\mathbf n,\gamma}\|
                                  &\leq C\gamma^{-1/2},
 \label{eq:f15v97:slow-compression}\\
 \frac{\|\Pi_\gamma T_{\Box,\gamma}(I-\Pi_\gamma)\|}
                         {\Lambda_\Box(\gamma)}
                                  &=O(\gamma^{-1/2}),
 \label{eq:f15v97:offdiagonal}\\
 \frac{\|(I-\Pi_\gamma)T_{\Box,\gamma}(I-\Pi_\gamma)\|}
                         {\Lambda_\Box(\gamma)}
                                  &\longrightarrow q_4<\frac15.
 \label{eq:f15v97:complement}
\end{align}
In particular the last ratio is at most \(1/4\)
for sufficiently large coupling.
In general \(\Pi_\gamma\) does \emph{not} commute
with the transfer: this is a nearly invariant
subspace, not an exact spectral projection.
\end{theorem}

\begin{proof}
The frame and chart are root-equivariant,
and \(\Omega_\Box\) is invariant.
Thus each ground isometry intertwines the
root tensor representation with the actual
root action. The left factors are untouched.
The native Perron replacement at \(\mathbf0\)
is invariant as well. This proves gauge
commutation and vacuum inclusion.

For each label there are \(d_{\mathbf n}\)
left states and \(d_{\mathbf n}\) retained
right-tensor states. The rank sum factorizes:
\[
 \sum_{0\leq n_v\leq N_\gamma}d_{\mathbf n}^2
                      =\left(\sum_{n=0}^{N_\gamma}(n+1)^2\right)^7.
\]
The root representation can be decomposed
into its irreducibles, but is already counted
in these right-tensor dimensions.
There is no additional root factor.

In the Euclidean comparison the ground
projection is
\(V(|\Omega_\Box\rangle\langle\Omega_\Box|\otimes I)V^*\).
It commutes with
\(V(\mathcal G_\Box\otimes W)V^*\).
The chart-normalized projection differs from
it in norm by an exponentially small amount,
uniformly in representation dimension:
multiplication by \(V\) is unitary.
Theorem~\ref{f15v97:factor} therefore gives
\eqref{eq:f15v97:slow-compression} and the
low-label contribution to
\eqref{eq:f15v97:offdiagonal}.
At \(\mathbf0\), the native Perron projection
has zero off-diagonal block, and its eigenvalue
is covered by \eqref{eq:f15v97:top-limit}.
No off-diagonal block connects distinct
nonroot label sectors.

On the remaining internal part of any retained
label sector, the comparison norm is at most
\(\ell_\Box q_4\), because \(0<W\leq I\).
The norm error is \(O(\gamma^{-1/2})\).
For the corrected \(\mathbf0\) sector, ordinary
min--max comparison gives the same bound
directly on its exact nonvacuum eigenvalues.

If any nonroot label exceeds \(N_\gamma\),
V83's single-vertex domination applies at that
degree-three vertex. The whole sector, including
every root representation and every internal
multiplicity, is bounded by
\[
 \frac{\|T_{\mathbf n,\gamma}\|}{\Lambda_\Box(\gamma)}
 \leq r_{N_\gamma+1}(3\gamma)
 \leq\exp\left[-\frac{k(k+2)}{6\gamma+2k+1}\right]
                    \longrightarrow e^{-6}<q_4,
 \qquad k=N_\gamma+1.
\]
Taking the supremum over sectors proves the
limsup in \eqref{eq:f15v97:complement}.
Conversely the first excited eigenvalue of
the \(\mathbf0\) block tends to
\(\ell_\Box q_4\). Its root spin-one states
are allowed here, and it is exactly orthogonal
to the native Perron replacement.
This proves the matching liminf.
The lower bound on \(\Lambda_\Box\) from
\eqref{eq:f15v97:top-limit} justifies all ratios.
\end{proof}

Boundary heat eigenvalues at different fusion
multiplicities can overlap internal excited
bands. No separated exact ground cluster
for the whole growing label box is inferred.
The noncommuting matrix \(W_{\mathbf n,\gamma}\)
and the off-diagonal error must travel together
in any subsequent gluing argument.

\subsection{An arbitrary-accuracy approximation on the full open region}
\label{f15v97:precision-section}

\begin{proposition}[Complete open-cube finite-rank control]
\label{f15v97:precision}
For each \(0<\varepsilon<1\), the native open
cube transfer has, for sufficiently large
\(\gamma\), a positive finite-rank approximation
with relative operator-norm error at most
\(\varepsilon\) and rank \(O_\varepsilon(\gamma^{21/2})\).
It can alternatively be chosen as
\(P_{\varepsilon,\gamma}T_{\Box,\gamma}
 P_{\varepsilon,\gamma}\), with
\(P_{\varepsilon,\gamma}\) an exact finite-rank
spectral projection commuting with all eight
gauge groups, the same order of rank bound,
and positive removed part.
\end{proposition}

\begin{proof}
Choose a fixed \(R\) with
\(e^{-R^2/6}<\varepsilon/8\).
Retain a finite Gaussian spectral subspace
\(W_\tau\) whose complementary norm for
\(\mathcal G_\Box\) is at most
\(\ell_\Box\tau\), with \(0<\tau<\varepsilon/8\).
For all seven labels at most \(R\sqrt\gamma\),
replace \(\mathcal G_\Box\) by this finite
positive truncation in the comparison operator
of \eqref{eq:f15v97:factor}, and pull it back
by \(\mathcal U_\gamma^*\).
Set the approximation to zero on other labels.

Its low-label error is at most
\(C_R\gamma^{-1/2}+\ell_\Box\tau\).
Its high-label error, relative to
\(\Lambda_\Box(\gamma)\), is bounded by
\(r_{\lfloor R\sqrt\gamma\rfloor+1}(3\gamma)\),
whose limsup is at most \(e^{-R^2/6}\).
Equation~\eqref{eq:f15v97:top-limit} gives the
claimed relative tolerance eventually.
Its rank is at most
\[
 \dim W_\tau\left(
       \sum_{n=0}^{\lfloor R\sqrt\gamma\rfloor}(n+1)^2
                         \right)^7
                  =O_\varepsilon(\gamma^{21/2}).
\]
The approximation is positive and gauge-equivariant,
but need not commute with the exact transfer;
nor is it asserted to lie below that transfer.

For the alternative native choice, apply this
construction with error strictly smaller than
\(\varepsilon\Lambda_\Box(\gamma)\), and call
its rank bound \(K\).
Min--max implies that the exact eigenvalue
with zero-based index \(K\) is at most
\(\varepsilon\Lambda_\Box(\gamma)\).
The exact spectral projection onto eigenvalues
strictly greater than this threshold has rank
at most \(K\), contains the vacuum and commutes
with the transfer and all gauge groups.
Its removed part is positive with norm at
most the threshold, proving the final assertion.
\end{proof}

These ranks are deliberately not presented
as economical many-cube algorithms.
The constants and required entry coupling
have not been numerically enclosed.
The gain is full exposed-boundary control
for an actual cube, with a specified coupled
charge operator and a rate, not a
volume-independent complexity claim.

\subsection{Direction, verification and limits}
\label{f15v97:scope-section}

The boundary geometry excluded from V96 has
now been addressed on one elementary cube:
every vertex remains open, the root action
is retained, and all internal states outside
the growing boundary subspace have a native
norm bound. The moving frame is determined by
the conditional graph Gaussian, while the
correlated average is compared to a positive
coupled boundary operator with an explicitly
paid noncommuting error.

This does not make different cubes independent.
A link partition in the spatial lattice must
still assign shared edges, retain its actual
boundary intertwiners, and restore all mixed
plaquettes with the correct normalization.
A spectral bound for the collective neutral
modes after that assembly remains missing.
The constants above are for one fixed cube;
no uniformity in a growing graph has been proved.
Independent physical time calibration and
nontrivial interacting four-dimensional
continuum reconstruction also remain open.

Even for a fixed boundary label, the present
\(O(\gamma^{-1/2})\) error is larger than a
putative first-order \(1/\gamma\) charged
energy. Thus it does not identify that
fixed-label coefficient. The shrinking
open charged gap discussed in V93 is not
removed by the fast-complement bound.

The checker
\begin{center}\small
\path{scripts/verify_ym_f15_v97_all_boundary_cube.py}
\end{center}
tests the exact incidence, cycle, covariance
and magnetic matrices; literal noncommuting
cube words and kernel orientation; the
global tree-depth budget; independent Wick
pairing and matrix-word reorderings; full
boundary carrier ranks; and the degree-three
high-charge tail. These finite algebra checks
are supplementary. The global compact bounds,
dimension-independent operator estimates and
spectral conclusions are analytic proofs above.
It passed all 3,192 new exact-rational fixtures
and reran the V96 checker and its inherited
chain successfully. The checker has 10,050 bytes
and SHA--256
\begin{center}\scriptsize\ttfamily
4E27FC3F2575E4ACD9BE9F68C2CEC08405F9C23C9DD184585EE0992E0673DE14.
\end{center}

The predecessor V96 has 2,202,122 bytes and SHA--256
\begin{center}\scriptsize\ttfamily
40C339981853E17061726F58220DD783705E513FF50DA3FBB0B5E9AC11D0D45A.
\end{center}
Its body, including its TOC formatting, is
retained byte for byte; deleting this V97
section and restoring the title recovers V96.
Only the active main source is replaced.
All 161 previously protected files remain
unchanged. The companion checkpoint remains
V95 and the next is V100, at the planned
archive/restart. Only \path{.tex} files are
stored in \path{latex_notes}; builds go to
\path{artifacts/ym_f15_v97_texcheck/}.

\begin{firewall}
V97 proves a full-boundary elementary-cube
operator approximation with a growing-charge
rate, a positive coupled boundary heat factor,
and complete internal-complement and finite-rank
control. The retained ground subspace is only
nearly invariant in general. No many-region
spectral gap, physical time calibration or
interacting continuum Yang--Mills construction
has been proved. The full mass-gap goal
remains unachieved.
\end{firewall}


\section{V98: an actual spatial assembly and its coarse boundary kinetic form}
\label{f15v98:context}

V97 retained every boundary action of one cube.
The next issue is not another isolated region:
elementary cubes in the usual cell decomposition
share links, so their transfers cannot be tensored
as independent factors. We use a disjoint
\emph{vertex-block} partition instead. It yields
native cube transfers and separate bridge-link
convolutions, with an exact list of the remaining
plaquettes.

The resulting Gauss-matched ground carrier is
a coarse multigraph Hilbert space. Its kinetic
comparison is a coupled electric operator, not
a product of independently damped boundary spins.
We derive this operator, prove volume-independent
quadratic-form bounds and compute its lowest
nonzero energy. These statements concern the
cut comparison, before its magnetic interfaces
are restored. A separate theorem controls the
full transfer at each fixed spatial volume;
its constants are explicitly not volume uniform.

The boundary-charge organization is consistent
with the motivation of the
\href{https://arxiv.org/abs/1607.08881}{Delcamp--Dittrich--Riello
fusion-basis work}. That paper addresses a
\(2+1\)-dimensional setting. No spectral theorem
from it is being imported into four-dimensional
Yang--Mills theory. The carrier identifications
and bounds below are proved for our native
\(\SU(2)\) spatial transfer.

\subsection{A link partition of the entire spatial torus}
\label{f15v98:partition-section}

Let the spatial lattice be
\((\mathbb Z/(2m)\mathbb Z)^3\), with \(m\geq4\).
Write \(B=m^3\) for the number of vertex blocks.
For \(b\in(\mathbb Z/m\mathbb Z)^3\), the block
has vertices
\[
 2b+u,\qquad u\in\{0,1\}^3.
\]
An oriented positive-\(\mu\) link is internal
when its base has even \(\mu\)-coordinate;
it is a bridge when that coordinate is odd.
The internal links of each block are exactly
the twelve edges of V97's cube. All internal
link sets and all bridge singletons are disjoint.

\begin{proposition}[Exact incidence and plaquette inventory]
\label{f15v98:inventory}
There are \(12B\) internal links and \(12B\)
bridges. Every original vertex has three internal
and three bridge incidences.
The \(24B\) spatial plaquettes divide into
\[
 6B\text{ internal},\qquad
 12B\text{ one-direction-crossing},\qquad
 6B\text{ two-direction-crossing}.
\]
Collapsing each block to a coarse vertex gives
four parallel bridges along each positive
nearest-neighbour coarse direction. This
multigraph has \(B\) vertices and \(12B\) links.
Its underlying simple graph has no triangles
when \(m\geq4\).
\end{proposition}

\begin{proof}
A link in direction \(\mu\) has exactly one
of the two possible base parities. Within a
block, there are four links of each type
for each \(\mu\), proving the link counts.
At a port \(u\), the bridge is outgoing in
direction \(\mu\) if \(u_\mu=1\), and incoming
from the previous block if \(u_\mu=0\).
This proves the incidence statement.

For a plaquette in directions \(\mu,\nu\),
classify the base by the parities of those
two coordinates. The even--even class is
internal. One odd coordinate gives two
bridges and two internal edges in neighbouring
blocks. Two odd coordinates give four bridges.
For each direction pair, the third coordinate
has two choices within a block; hence the
counts per block are \(2,4,2\), respectively.
Sum over the three direction pairs.
The four bridges between coarse neighbours
are indexed by the two transverse port bits.
A three-step closed walk in the underlying
coarse graph would require a length-three
periodic coordinate cycle or a triangle in
the cubical lattice; neither occurs for \(m\geq4\).
\end{proof}

Use V97's rooted axial tree separately in
each block. In its coordinates an internal
link is \(g_sX_eg_t^{-1}\), with \(X_e=I\)
on that tree. A bridge is written
\(g_sY_eg_t^{-1}\), using the local \(g\)
at each endpoint, and \(g=I\) at every block
root. The coordinate changes are Haar
translations after the independent internal
tree changes. Thus the exact full product
measure is
\[
 dH^{7B}(g)\,dH^{5B}(X)\,dH^{12B}(Y).
\]
Nonroot Gauss invariance removes the \(g\)
coordinates. The remaining block-root action
conjugates that block's five \(X\)'s and acts
by the usual left/right transformations on
its incident \(Y\)'s. Consequently the
\emph{entire} physical Hilbert space is exactly
\begin{equation}
 \mathcal H_{\rm phys}\simeq
 L^2(G^{5B}\times G^{12B})^{G^B}.
 \label{eq:f15v98:full-carrier}
\end{equation}
No internal oscillator approximation is used
in this identity.

For a one-direction-crossing plaquette, with
orientation chosen along its first bridge,
the exact reduced word is
\begin{equation}
 Y_0 X_{\rm right}Y_1^{-1}X_{\rm left}^{-1}.
 \label{eq:f15v98:digon}
\end{equation}
Each displayed \(X\) is the corresponding
internal edge, possibly \(I\), in the local
tree representative. The intermediate \(g\)'s
cancel consecutively in the original plaquette
word. A two-direction-crossing plaquette
similarly becomes the word of its four
oriented \(Y\)'s. These are decorated digons
and coarse quadrilaterals; none is dropped.

\subsection{The exact native magnetic sandwich}
\label{f15v98:sandwich-section}

On the full link Haar space define
\[
 T_{\rm cut,\gamma}
   =\left(\bigotimes_{b=1}^B T_{\Box,\gamma}^{(b)}\right)
          \otimes\left(\bigotimes_{e\ {\rm bridge}}C_\gamma^{(e)}\right),
\]
using V97's normalized \(C_\gamma\), whose
Peter--Weyl eigenvalues are \(r_n(\gamma)\).
All factors commute with the original gauge
actions. Let \(S_\times(X,Y)\) be the sum of
the \(18B\) mixed plaquette actions just listed,
and put \(M_\times=e^{-\gamma S_\times/2}\).
After restriction to \(\mathcal H_{\rm phys}\),
the full native transfer is exactly
\begin{equation}
 T_{\rm full,\gamma}
       =M_\times T_{\rm cut,\gamma}M_\times,\qquad
 0<M_\times\leq1.
 \label{eq:f15v98:sandwich}
\end{equation}
The cut Perron value is
\begin{equation}
 \|T_{\rm cut,\gamma}|_{\rm phys}\|
                      =\Lambda_\Box(\gamma)^B.
 \label{eq:f15v98:cut-top}
\end{equation}
Indeed the product of the cube Perron vectors
and constant bridge vectors is invariant
under every original vertex action and attains
the full tensor-product norm.
Both cut and full strictly positive kernels
have gauge-invariant Perron vectors, so their
full Haar norms equal their physical norms.

\subsection{Exact Gauss matching of the cube ground carriers}
\label{f15v98:carrier-section}

Define the coarse bridge Hilbert space
\[
 \mathcal K=L^2(G^{12B})^{G^B},
\]
where all bridge ends attached to one block
transform under the same coarse gauge group.
Let \(\mathcal K_N\) retain bridge labels
\(0\leq n_e\leq N\), including all their matrix
indices before the coarse Gauss projection.
It is finite dimensional.

At each of the eight ports \(u\) of a block,
the three bridge ends have a total
\(\SU(2)\) representation, with spin label
\(n_{b,u}\leq3N\). All fusion multiplicities
are retained. At the seven nonroot ports,
match this total representation to the
corresponding left carrier of V97's
\(J_{\mathbf n,\gamma}\). At the root use
the remaining original Gauss constraint.

\begin{proposition}[A native all-volume ground-carrier isometry]
\label{f15v98:isometry}
For every \(N,\gamma\) there is an isometry
\[
 S_{\gamma,N}:\mathcal K_N\longrightarrow\mathcal H_{\rm phys}
\]
whose image uses V97's exact cube ground
isometries at every block and no bridge label
above \(N\). In particular,
\[
 \operatorname{rank}S_{\gamma,N}
        =\dim\mathcal K_N
        \leq D_N^{12B},\qquad
 D_N=\sum_{n=0}^N(n+1)^2.
\]
This is an identification of a selected native
subspace, not yet an approximation to the
entire physical transfer.
\end{proposition}

\begin{proof}
For each block the retained cube representation,
before original Gauss matching, has the form
\[
 \bigoplus_{\mathbf n}
   \left(\bigotimes_{u\ne o}V_{n_u}^{\rm left}\right)
       \otimes
   \left(\bigotimes_{u\ne o}V_{n_u}^{\rm root}\right),
\]
with a dual placed on one side according to
the endpoint convention. This is precisely
the representation content of the regular
functions on seven group variables. V97's
ground isometries intertwine the actual
root action with the displayed tensor product;
they preserve norms, also in the all-zero
sector where the native Perron vector is used.

Decompose each three-end bridge tensor at a
nonroot port into irreducibles and their
multiplicity spaces. Its contraction with
the cube left carrier uses the normalized
invariant vector in \(V_n\otimes V_n^*\).
Such contraction transfers that port carrier
to the block root without changing its
multiplicity. Schur orthogonality, or directly
the normalized Peter--Weyl bases, shows that
this is a unitary identification of the
invariant tensor spaces. Apply it to all
seven ports. The last Gauss constraint is
exactly invariance under the simultaneous
action on all eight transported port tensors.
This is the coarse \(G^B\) invariance defining
\(\mathcal K_N\).

Conversely every original Gauss-invariant vector
in the stated ground carriers has this
decomposition. Thus no further multiplicity
or representation-dimension weight is missing.
The rank bound follows by forgetting the
coarse Gauss restriction.
\end{proof}

For example its exact finite rank can also be
written without a choice of fusion tree:
\[
 \dim\mathcal K_N
  =\int_{G^B}\prod_{e\ {\rm bridge}}
       \left(\sum_{n=0}^N
          \chi_n(h_{s(e)})\overline{\chi_n(h_{t(e)})}\right)
                                            \,dH^B(h).
\]
This is the trace of the coarse Gauss projector
on the full truncated bridge space; character
orthogonality retains every intertwiner
multiplicity. It is not a product of independent
singlet probabilities.

This construction is the representation
version of contracting each internal spanning
tree. Its use of \(J_{\mathbf n,\gamma}\)
adds the native, charge-dependent internal
wavefunction. It does not identify that
wavefunction with a delta function at \(X=I\).

\subsection{A root-independent coarse electric operator}
\label{f15v98:electric-section}

Use Hermitian endpoint generators in the
same normalization as V97:
\(\sum_\alpha J_{e,\alpha}^2=n_e(n_e+2)\).
Opposite endpoints of a bridge carry commuting
left and right actions with the appropriate
dual convention. At a port define
\[
 \mathbf J_{b,u}
       =\sum_{e\ {\rm incident\ at}\ (b,u)}\mathbf J_{e,b}.
\]
There are exactly three terms. Different
ports have commuting gauge actions.
On \(\mathcal K\) the coarse Gauss constraint is
\begin{equation}
 \sum_{u\in\{0,1\}^3}\mathbf J_{b,u}=0
                      \quad\hbox{for every block }b.
 \label{eq:f15v98:coarse-gauss}
\end{equation}
Initially all formulas below are on finite
Peter--Weyl sums.

Let \(L_\Box\) be the full \(8\)-by-\(8\)
combinatorial Laplacian of the cube, and
\(L_\Box^+\) its inverse on zero-sum vectors,
extended by zero on constants. Set
\begin{align}
 \mathcal E&=\sum_{e\ {\rm bridge}}n_e(n_e+2),\notag\\
 \mathcal Q_b&=\sum_{u,v}(L_\Box^+)_{uv}
                          \mathbf J_{b,u}\cdot\mathbf J_{b,v},\notag\\
 \mathcal D&=\mathcal E+\sum_b\mathcal Q_b.
 \label{eq:f15v98:kinetic}
\end{align}
Here the first line denotes the bridge
Casimir operator, not a fixed scalar on
the whole Hilbert space.

\begin{proposition}[Exact origin of the coupled port form]
\label{f15v98:root}
On the coarse physical space, the port form
\(\mathcal Q_b\) equals
\[
 \sum_{u,v\ne o}\Sigma_{uv}
                         \mathbf J_{b,u}\cdot\mathbf J_{b,v},
\]
where \(\Sigma\) is V97's inverse root-deleted
cube Laplacian. Thus it is independent of
the chosen root after Gauss matching.
The operators \(\mathcal Q_b\) for distinct
blocks commute, and each commutes with every
bridge Casimir. They preserve the coarse
Gauss constraints and every bridge-label cutoff.
\end{proposition}

\begin{proof}
For ordinary real zero-sum charges \(j\), solve
\(L_\Box\phi=j\), choosing \(\phi_o=0\).
Then \(\phi_{u\ne o}=\Sigma j_{u\ne o}\)
and
\[
 j^{\mathsf T}L_\Box^+j
       =\sum_u j_u\phi_u
       =j_{u\ne o}^{\mathsf T}\Sigma j_{u\ne o}.
\]
Polarization gives the corresponding bilinear
identity. The difference of the two full
quadratic coefficient matrices has the form
\({\bf1}a^{\mathsf T}+a{\bf1}^{\mathsf T}\).
For each Cartesian component, replacing \(j\)
by endpoint operators produces terms containing
the total port generator. That generator
commutes with the other linear combination
of the \emph{same} Cartesian component, and
annihilates physical vectors by
\eqref{eq:f15v98:coarse-gauss}.
This proves the operator identity without
commuting different Cartesian components.

Under the normalized carrier contractions
above, V97's right generators become the
transported port generators, up to simultaneous
duality signs that cancel in these dot products.
It is therefore this quadratic form, not
separate port Casimirs, that its boundary
heat factor induces.

Each dot product is invariant under simultaneous
rotation at the coarse vertex. Different
blocks act on distinct endpoint indices;
even on a shared bridge its left and right
actions commute. Casimirs commute with both
endpoint actions. All remaining assertions follow.
\end{proof}

\begin{theorem}[Volume-independent ellipticity of the cut kinetic form]
\label{f15v98:ellipticity}
On \(\mathcal K\), as quadratic forms,
\begin{equation}
 \mathcal E\leq\mathcal D\leq4\mathcal E.
 \label{eq:f15v98:ellipticity}
\end{equation}
The nonnegative self-adjoint realization
obtained from the finite Peter--Weyl blocks
has compact resolvent and a unique zero
eigenvector, the constant function.
The constants in \eqref{eq:f15v98:ellipticity}
do not depend on \(m\), \(B\) or the cutoff.
\end{theorem}

\begin{proof}
For a binary mask \(\eta\in\{0,1\}^3\), the
Walsh vector \((-1)^{\eta\cdot u}\) is a
cube-Laplacian eigenvector of eigenvalue
\(2|\eta|\). Thus the nonzero eigenvalues
of \(L_\Box\) are \(2,4,6\), and on total-zero
charges
\[
 \frac16\sum_u|\mathbf J_{b,u}|^2
       \leq\mathcal Q_b
       \leq\frac12\sum_u|\mathbf J_{b,u}|^2.
\]
These form inequalities follow by factoring
the positive scalar coefficient matrices and
summing squared norms; they remain valid
for the noncommuting Cartesian generators.

For each component, Cauchy--Schwarz on the
three endpoint vectors at a port gives
\[
 \|\mathbf J_{b,u}\psi\|^2
       \leq3\sum_{e\ni(b,u)}\|\mathbf J_{e,b}\psi\|^2,
\]
where vector norms include the three components.
Every bridge has two endpoints. Summing gives
\[
 0\leq\sum_b\mathcal Q_b
       \leq\tfrac32\sum_{b,u,e\ni(b,u)}
                                |\mathbf J_{e,b}|^2
       =3\mathcal E.
\]
This proves the form comparison.

The operators preserve each complete bridge
Peter--Weyl label sector, which is finite
dimensional. Define their self-adjoint direct
sum on those sectors. Below any fixed energy,
\(\mathcal E\) permits only finitely many
labels, and \(\mathcal D\geq\mathcal E\).
This proves compact resolvent.
Finally \(\mathcal E=0\) forces every bridge
label to be zero, leaving just the constant.
The upper bound also shows equality of the
closed form domains.
\end{proof}

\subsection{The exact lowest coarse kinetic energy}
\label{f15v98:edge-section}

The preceding form comparison can be sharpened
at the spectral bottom for this actual
multigraph.

\begin{theorem}[Sharp edge of the assembled cut comparison]
\label{f15v98:edge}
For every \(m\geq4\),
\begin{equation}
 \inf\operatorname{spec}
       (\mathcal D|_{\mathbf1^\perp})=\frac{19}{2}.
 \label{eq:f15v98:edge}
\end{equation}
Consequently the top-two ratio of
\(e^{-\mathcal D/(2\gamma)}\) on \(\mathcal K\)
is exactly \(e^{-19/(4\gamma)}\).
This is not asserted to be the ratio of the
full native Wilson transfer.
\end{theorem}

\begin{proof}
The cube effective resistance between two
ports at Hamming distance \(h\) is
\[
 R_h=(e_u-e_v)^{\mathsf T}L_\Box^+(e_u-e_v)
       =\begin{cases}
          7/12,&h=1,\\
          3/4,&h=2,\\
          5/6,&h=3.
         \end{cases}
\]
For an explicit verification use the Walsh
diagonalization above:
\[
 R_h=\frac18\sum_{\eta\ne0}
       \frac{\big((-1)^{\eta\cdot u}
                    -(-1)^{\eta\cdot v}\big)^2}{2|\eta|}.
\]
The finite sum gives the three displayed
rational values.

Every bridge Casimir commutes with \(\mathcal D\).
In a sector of electric energy exceeding \(19/2\),
the form bound already suffices.
The only nonempty label lists of electric
energy at most \(19/2\) are one label \(1\),
one label \(2\), two labels \(1,1\), or three
labels \(1,1,1\): their energies are \(3,8,6,9\).
A single nontrivial link cannot satisfy Gauss
invariance at its endpoints.
For three spin-half links, the central gauge
element \(-I\) requires even incidence at every
coarse vertex. With three links and no
self-loops this would be a triangle, absent
by Proposition~\ref{f15v98:inventory}.
For two spin-half links, Gauss invariance
therefore requires parallel links between
the same two blocks.

That sector is one dimensional: the singlet
at each of its two occupied coarse vertices
is unique. It is represented by
\(\chi_1(Y_eY_f^{-1})\).
The two endpoint port charges at either block
sum to zero and have Casimir \(3\), so
\(\mathcal Q_b=3R_h\) there.
The ports lie in one face; distinct bridges
have transverse Hamming distance \(h=1\) or \(2\),
the same at both blocks. The total eigenvalue is
\[
 6+6R_h=\frac{19}{2}\quad\hbox{or}\quad\frac{21}{2}.
\]
Adjacent transverse ports exist at every coarse
nearest-neighbour pair, giving the first value.
All other physical sectors have already been
bounded from below. This proves both the edge
and its attainment. The heat ratio follows
by functional calculus.
\end{proof}

This positive auxiliary edge is in the
unscaled operator \(\mathcal D\). Its heat
deficit is \(19/(4\gamma)+O(\gamma^{-2})\);
it does not by itself provide a positive
physical mass or determine a physical clock.

\subsection{A native compression estimate, with its volume cost exposed}
\label{f15v98:comparison-section}

Fix \(R>0\), let \(N_\gamma=\lfloor R\sqrt\gamma\rfloor\),
and abbreviate \(S_\gamma=S_{\gamma,N_\gamma}\).
The bridge cutoff is preserved by \(\mathcal D\).

\begin{theorem}[From native cut transfer to the coarse coupled heat]
\label{f15v98:comparison}
There are constants \(C_R,\gamma_R\) independent
of the number of blocks such that, for
\(\gamma\geq\gamma_R\),
\begin{equation}
 \left\|
 S_\gamma^*\frac{T_{\rm cut,\gamma}}{\Lambda_\Box(\gamma)^B}S_\gamma
       -e^{-\mathcal D/(2\gamma)}|_{\mathcal K_{N_\gamma}}
 \right\|
                \leq\frac{B C_R}{\sqrt\gamma}.
 \label{eq:f15v98:comparison}
\end{equation}
No estimate on the entire orthogonal complement
of \(S_\gamma\mathcal K_{N_\gamma}\) is claimed here.
\end{theorem}

\begin{proof}
Each nonroot cube-port label is at most
\(3N_\gamma\), by three-end fusion.
V97's native compression theorem and its
Perron asymptotic imply, uniformly on these
carriers,
\[
 \left\|
 \frac{J_{\mathbf n,\gamma}^*T_{\mathbf n,\gamma}
                         J_{\mathbf n,\gamma}}{\Lambda_\Box(\gamma)}
             -e^{-Q_{\mathbf n,\gamma}}\right\|
                         \leq C_R\gamma^{-1/2}.
\]
Both operators are positive contractions.
In the all-zero sector the native compression
is exactly one. Direct sums over labels do
not multiply the error by their dimensions.

For each bridge, the normalized Fourier
integral and V97's noncommuting Gaussian
estimate give
\[
 \sup_{0\leq n\leq R\sqrt\gamma}
 \left|r_n(\gamma)-e^{-n(n+2)/(2\gamma)}\right|
                           \leq C_R\gamma^{-1/2}.
\]
To see the required uniformity without a
fixed-order Bessel asymptotic, scale the full
one-link exponential chart. The normalized
scalar density differs in \(L^1\) from the
standard three-dimensional Gaussian by
\(O(\gamma^{-1})\): the Taylor remainder for
\(1-\cos t\), the sinc-squared Jacobian and
the complete-chart Gaussian tail give this
bound, dominated by a polynomial Gaussian.
Keep the norm-one matrix \(D_n\) during that
replacement. Then apply
Proposition~\ref{f15v97:wick} to
\(J_\alpha/\sqrt\gamma\), whose commutators
are \(O_R(\gamma^{-1/2})\).
Their quadratic sum is the scalar
\(n(n+2)/(2\gamma)\). This proves the display.

Tensor telescoping for contractions gives
the sum of the \(B\) cube errors and \(12B\)
bridge errors. Restriction to invariant
tensors and the carrier isometry cannot
increase the operator norm.
Under the exact matching, the cube factors
become \(e^{-\mathcal Q_b/(2\gamma)}\) by
Proposition~\ref{f15v98:root}. Bridge factors
become their Casimir heats.
They commute with one another for the
endpoint reasons proved there, so their
product is exactly \(e^{-\mathcal D/(2\gamma)}\).
This proves \eqref{eq:f15v98:comparison}.
\end{proof}

For fixed volume this is a convergent
growing-carrier estimate. It is not uniform
when \(B/\sqrt\gamma\) fails to tend to zero.
Even at fixed \(B\), its error is larger than
the \(1/\gamma\) heat deficit in
Theorem~\ref{f15v98:edge}. Thus that theorem
does not, through this estimate, identify
the native cut transfer's first-order gap.
Restoring \(M_\times\) is an additional
problem, not included in this comparison.

\subsection{The full transfer can be approximated, but the denominator costs volume}
\label{f15v98:full-approx-section}

Here we retain enough \emph{internal} cube
states, instead of only the ground carriers.
The tensor spectral-tail observation used
in V96 applies to the actual torus partition.

\begin{proposition}[An explicit native denominator and fixed-volume approximation]
\label{f15v98:full-approx}
Put
\[
 c=\frac{e^{-5/2}}{3\sqrt2\,\pi^{5/2}},\qquad
 \kappa=c^{24}>0.
\]
For \(\gamma\geq1\),
\begin{equation}
 \kappa^B\leq
 \frac{\|T_{\rm full,\gamma}|_{\rm phys}\|}
                       {\Lambda_\Box(\gamma)^B}
 \leq1.
 \label{eq:f15v98:denominator}
\end{equation}
For fixed \(B\) and \(0<\delta<1\), the full
physical transfer has a positive finite-rank
approximation with positive removed part,
relative norm error at most \(\delta\), and
rank
\begin{equation}
 O_{B,\delta}\big(\gamma^{57B/2}\big)
                     \quad(\gamma\longrightarrow\infty).
 \label{eq:f15v98:full-rank}
\end{equation}
The constants and entry coupling are not
uniform in \(B\).
\end{proposition}

\begin{proof}
For the lower bound take the normalized
indicator of the product of link balls
\(|\log U_e|\leq1/(2\sqrt\gamma)\) in the
full Haar space. Write \(E=P=24B\) for
its link and plaquette counts, and \(v_\gamma\)
for the one-link ball measure.
The complete-chart Haar formula gives
\[
 v_\gamma\geq\frac{\gamma^{-3/2}}{3\pi^3},
\]
using \(\sin t/t\geq2/\pi\) on
\(0\leq t\leq\pi/2\).
V84's scalar normalization bound is
\(c_0(\gamma)\leq\sqrt{2/\pi}\,\gamma^{-3/2}\).

For two configurations in those balls,
each kinetic action is at most \(1/(2\gamma)\).
Each plaquette action is at most \(2/\gamma\),
by the four-factor chord or geodesic bound.
The two magnetic ends together consequently
cost at most \(e^{-2P}\); the kinetic
exponents cost at most \(e^{-E/2}\).
The Rayleigh quotient of the indicator is
at least
\[
 e^{-2P-E/2}
       \left(\frac{v_\gamma}{c_0(\gamma)}\right)^E
                           \geq c^{24B}.
\]
Although this test vector is not gauge
invariant, the full strictly positive
transfer has an invariant Perron vector.
Its full norm equals its physical norm,
so the test is a valid lower bound for that
physical norm as well.
Also \(\Lambda_\Box\leq1\).
These facts and the contraction sandwich give
\eqref{eq:f15v98:denominator}.

Set \(\rho=\delta\kappa^B\).
In each cube choose V97's exact spectral
projection with omitted norm at most
\(\rho\Lambda_\Box\); in each bridge choose
the exact Peter--Weyl cutoff with omitted
norm at most \(\rho\).
Their tensor product \(P_\rho\) commutes with
the cut transfer and all gauge actions.
In the complete tensor eigenbasis, any omitted
state has at least one omitted factor.
All other factor eigenvalues are bounded
by their own top values. Hence
\[
 0\leq T_{\rm cut,\gamma}
          -P_\rho T_{\rm cut,\gamma}P_\rho,\qquad
 \|T_{\rm cut,\gamma}
          -P_\rho T_{\rm cut,\gamma}P_\rho\|
                     \leq\rho\Lambda_\Box^B.
\]
There is no factor counting regions in this
maximum-tail estimate.
Sandwich by \(M_\times\), restrict to the
physical space and divide by
\eqref{eq:f15v98:denominator}.
This gives the required relative error
and positivity.

For fixed \(\rho\), each cube rank is
\(O_\rho(\gamma^{21/2})\), by V97.
A bridge cutoff of size \(O_\rho(\sqrt\gamma)\)
suffices by V83's coefficient tail bound,
and has rank \(O_\rho(\gamma^{3/2})\).
There are \(B\) cubes and \(12B\) bridges.
Multiplication by \(M_\times\) and Gauss
restriction do not increase rank, giving
\(21B/2+18B=57B/2\).
\end{proof}

This approximation uses the magnetically
weighted image of \(P_\rho\), not an
unweighted coarse ground projection.
Its denominator lower bound is intentionally
crude but noncircular and independent of
coupling. The exponential volume loss
\(\kappa^B\) prevents it from being the
needed volume-uniform spectral certificate.

\subsection{Why the mixed plaquettes are not a small perturbation}
\label{f15v98:nonfreeze-section}

Consider a one-direction-crossing \(xy\)
plaquette across the \(x\)-face at height
\(z=0\). With the chosen axial trees, its
right \(y\)-edge is a tree edge \(I\), and
its left \(y\)-edge is the chord \(a\).
When its two bridge coordinates equal \(I\),
\eqref{eq:f15v98:digon} becomes \(a^{-1}\).
Let \(\omega_\gamma\) be the normalized
native Perron vector of V97's \(\mathbf n=0\)
cube block. Then
\begin{equation}
 \int_{G^5}|\omega_\gamma(X)|^2
                      e^{-\gamma s(a)/2}\,dH^5(X)
       \longrightarrow
 \int_{\mathbb R^{15}}|\Omega_\Box(x)|^2
                              e^{-|x_a|^2/4}\,dx
                         \quad\in(0,1).
 \label{eq:f15v98:nonfreeze}
\end{equation}

To justify this, V97's norm convergence and
the simple separated Gaussian top eigenvalue
give convergence of the normalized Perron
vectors under the complete-chart embedding.
For example the spectral projection follows
from a fixed contour around that top eigenvalue;
the resolvent identity gives norm convergence,
and positivity fixes the phase.
Their squared densities therefore converge
in \(L^1\).
The displayed bounded multiplier converges
pointwise to the Gaussian multiplier, so
dominated convergence proves
\eqref{eq:f15v98:nonfreeze}.
The Gaussian ground density is strictly positive,
and its \(a\)-coordinate is nondegenerate;
the limit is strictly below one and above zero.

Freezing \(a=I\) would instead give exactly
one. Thus the natural \(1/\sqrt\gamma\)
internal fluctuations already change this
mixed plaquette average by order one.
This is a diagnostic for one actual native
fast-state average, not an assertion that
it equals the full charge-dependent
compression \(S_\gamma^*M_\times S_\gamma\).
It rules out justifying that magnetic
replacement solely by small coordinate size.
The associated internal response and the
moving frames must be retained in a subsequent
effective magnetic calculation.

\subsection{What has moved, and what remains}
\label{f15v98:scope-section}

We now have a literal whole-spatial-lattice
link partition, an exact mixed-plaquette
inventory and Gauss-matched coarse carrier,
a native cut compression estimate, and
volume-independent bounds on its derived
coupled kinetic form. The sharp auxiliary
edge \(19/2\) is proved on the full coarse
physical Hilbert space, not on a chosen
finite list of trial states.

None of those facts licenses discarding
the \(18B\) magnetic interfaces.
The full-transfer approximation pays them
exactly, but its relative-error certificate
loses an exponential in volume.
The ground-carrier comparison instead has
an explicit \(B/\sqrt\gamma\) error and does
not control all discarded internal states.
These are different estimates, and their
advantages must not be combined by silently
forgetting their respective costs.

The next substantive problem is the actual
magnetic response on the Gauss-matched coarse
carrier, including virtual internal excitations
and the neutral collective modes.
A volume-uniform coercive estimate for that
response, followed by independent physical
time calibration and nontrivial continuum
reconstruction, is still missing.
The known strong-coupling or auxiliary
electric gaps cannot be substituted for it.

The checker
\begin{center}\small
\path{scripts/verify_ym_f15_v98_parity_assembly.py}
\end{center}
checks the complete torus link and plaquette
partition at several sizes, noncommuting
mixed words, exact cube pseudoinverses and
resistances, root-independent quadratic
forms, three-end fusion multiplicities and
Casimir bounds, the low-energy label inventory,
and the tensor maximum-tail identity.
Its finite fixtures supplement the analytic
proofs; they do not certify a Yang--Mills
mass gap.
All 29,644 new exact-rational and combinatorial
fixtures passed, together with the inherited
V97 chain. The checker has 7,739 bytes and SHA--256
\begin{center}\scriptsize\ttfamily
FE676DFF3C71BF97355A174F9CF92FBD9D286BE3B2EEA9C921C7BFBAFA58BDF5.
\end{center}

The predecessor V97 has 2,237,664 bytes and SHA--256
\begin{center}\scriptsize\ttfamily
881F12B9BEE4AE7E885E42C6CCBCCA3EAD32057CE7661088064E2568A9EF6EBA.
\end{center}
Its body is retained byte for byte; removing
this V98 section and restoring the title
recovers V97. Only the active main source
is replaced. All 162 previously protected
files remain unchanged. The last companion
checkpoint is V95 and the next is V100,
at the planned archive/restart.
Only \path{.tex} files are stored in
\path{latex_notes}; compilation products go to
\path{artifacts/ym_f15_v98_texcheck/}.

\begin{firewall}
V98 gives an exact spatial assembly, an
explicit coupled coarse kinetic comparison
with a sharp auxiliary edge, and a full
fixed-volume native approximation.
The magnetic response has not been shown
to preserve that edge. No volume-uniform
native physical gap, independent physical
clock or interacting continuum Yang--Mills
construction has been proved. The full
mass-gap goal remains unachieved.
\end{firewall}


\section{V99: moving-frame magnetic response and dressed carriers}
\label{f15v99:context}

V98's spatial assembly did not justify treating
the mixed plaquettes as a small perturbation
of its coarse electric comparison. We now
retain the moving frame during the magnetic
calculation. It gives an exact conditional
Gram formula for the native compact magnetic
multiplier, and a quantitative failure of
ground-only magnetic projection. The appropriate
replacement is a normalized magnetically
dressed carrier.

The leading weak-field response of this
carrier can be integrated explicitly.
Its internal conditional covariance has
volume-independent exponential locality;
its determinant normalization has a uniformly
summable local expansion. These are Gaussian
response results with a proved fixed-volume
connection to the native compact integral.
They are not a uniform approximation of the
entire native transfer or its low spectrum.
In particular, the coarse magnetic quadratic
form retains its long-wavelength directions.

Throughout, use V98's side-\(2m\) spatial
torus, \(m\geq4\), with \(B=m^3\) blocks,
\(12B\) bridge links and \(18B\) mixed
plaquettes. The group is \(G=\SU(2)\).
All original links, plaquettes and Gauss
constraints remain present.

\subsection{The moving frame as an exact Haar-preserving bridge map}
\label{f15v99:frame-section}

Let \(\Omega_\Box\) be V97's normalized
Gaussian ground vector on \(\mathbb R^{15}\).
For the complete five-chord chart put
\[
 a_\gamma=\int_{\{|x_i|<\pi\sqrt\gamma,\ i=1,\ldots,5\}}
                         |\Omega_\Box(x)|^2\,dx,
 \qquad
 \varphi_\gamma(X)
   =a_\gamma^{-1/2}
       \frac{\Omega_\Box(x)}
            {\prod_{i=1}^5J_\gamma(x_i)^{1/2}}.
\]
The chart is one-to-one off a Haar-null set,
on which values may be assigned arbitrarily.
Thus \(\varphi_\gamma\) is an exact normalized
Haar \(L^2\) function, and its squared density
in \(x\) is the normalized truncated Gaussian.
It is invariant under simultaneous conjugation.

Use V97's real matrix \(\mathsf F\) and define,
at each nonroot port \(u\) of block \(b\),
\[
 h_{b,u}(X)=\operatorname{Exp}
                   ((\mathsf F x_b)_u/\sqrt\gamma),
 \qquad h_{b,o}=I.
\]
For a bridge from port \(s\) to port \(t\), set
\begin{equation}
 (\Theta_XY)_e=h_s(X)^{-1}Y_eh_t(X).
 \label{eq:f15v99:theta}
\end{equation}
At fixed \(X\) this is a product of Haar
translations. Its inverse is
\(Y_e=h_sZ_eh_t^{-1}\).
Both maps respect the coarse root gauge action.
Write
\(\Phi_\gamma(X)=\prod_b\varphi_\gamma(X_b)\),
and define on V98's coarse physical space
\(\mathcal K=L^2(G^{12B})^{G^B}\)
\begin{equation}
 (\widehat S_\gamma f)(X,Y)
             =\Phi_\gamma(X)f(\Theta_XY).
 \label{eq:f15v99:isometry}
\end{equation}

\begin{proposition}[Exact realization of the growing-charge frame]
\label{f15v99:frame}
The map \(\widehat S_\gamma:\mathcal K\to
\mathcal H_{\rm phys}\) is an isometry.
Its Peter--Weyl decomposition is V97's moving
Gaussian ground carrier at every nonroot
port, with all charge and fusion multiplicities
retained. It uses the Gaussian ground vector
also in the zero-label cube sector.
For the bridge cutoff \(\mathcal K_N\),
\[
 \|\widehat S_\gamma|_{\mathcal K_N}-S_{\gamma,N}\|
                   \leq C B\gamma^{-1/2}
\]
for sufficiently large \(\gamma\), where
\(S_{\gamma,N}\) is V98's isometry, and \(C\)
does not depend on \(N\) or \(B\).
\end{proposition}

\begin{proof}
For fixed \(X\), change \(Y\) to \(Z=\Theta_XY\).
The product Haar measure is unchanged, and
\(\|\Phi_\gamma\|_2=1\), proving isometry.
Conjugation equivariance of the exponential,
the real linearity of \(\mathsf F\) and the
invariance of \(\Phi_\gamma\) prove the
Gauss assertion.

The cube Fourier factor is \(D_n(g)\).
Replacing its right carrier by the V97 frame
multiplies it by \(D_n(h(X))\), giving
\(D_n(gh(X))\).
After matching a bridge end at that port,
the transported bridge argument is therefore
\(h_s^{-1}Y_eh_t\).
This verifies \eqref{eq:f15v99:isometry},
including the inverse convention, first on
finite Peter--Weyl sums and then by density.
Normalized singlet contractions are those of
V98 and introduce no dimension factor.

The two isometries differ only where V98
replaced a cube's zero-label Gaussian vector
by its actual Perron vector.
V97's \(O(\gamma^{-1/2})\) operator convergence
and the simple separated Gaussian top eigenvalue
give the same order of convergence of their
rank-one projections by a fixed resolvent
contour. Positivity fixes the phase of the
normalized vectors. Hence each local isometry
changes in norm by \(O(\gamma^{-1/2})\),
uniformly in the other labels.
Tensor telescoping and restriction to invariant
tensors give the stated \(B\)-factor.
\end{proof}

This exact carrier is defined on all of
\(\mathcal K\), without a finite electric
cutoff. It is not claimed to consist of
native eigenvectors. Combining the last
bound with V98 preserves its cut-compression
estimate on \(\mathcal K_{\lfloor R\sqrt\gamma\rfloor}\),
with an \(O_R(B/\sqrt\gamma)\) error.
It does not remove that error or control
all omitted internal states.

\subsection{An exact compact magnetic Gram and leakage identity}
\label{f15v99:gram-section}

Retain the actual multiplier
\(M_\times=e^{-\gamma S_\times(X,Y)/2}\).
For \(t>0\) define the scalar function
\begin{equation}
 F_{\gamma,t}(Z)=
  \int|\Phi_\gamma(X)|^2
   \exp\!\left[-\frac{t\gamma}{2}
        S_\times(X,\Theta_X^{-1}Z)\right]\,dH^{5B}(X).
 \label{eq:f15v99:Ft}
\end{equation}
No mixed plaquette is replaced by its quadratic
part in this definition.
The functions are continuous, coarse-gauge
invariant, and satisfy
\[
 e^{-18t\gamma B}\leq F_{\gamma,t}\leq1.
\]
Continuity follows by dominated convergence;
gauge invariance follows by a simultaneous
change of the internal variables.

\begin{theorem}[Native magnetic compression with the moving frame retained]
\label{f15v99:gram}
Writing \(M_F\) for multiplication by \(F\),
one has exactly
\begin{align}
 \widehat S_\gamma^*M_\times^t\widehat S_\gamma
                         &=M_{F_{\gamma,t}},\notag\\
 \widehat S_\gamma^*M_\times
       (I-\widehat P_\gamma)M_\times\widehat S_\gamma
                         &=M_{F_{\gamma,2}-F_{\gamma,1}^2},
 \quad
 \widehat P_\gamma=\widehat S_\gamma\widehat S_\gamma^*.
 \label{eq:f15v99:gram}
\end{align}
In particular \(F_{\gamma,2}\geq F_{\gamma,1}^2\).
\end{theorem}

\begin{proof}
Apply the Haar-preserving change of variables
in \eqref{eq:f15v99:theta} to the quadratic
form of the first operator. The two coarse
test functions then depend only on \(Z\);
their remaining internal integral is
\eqref{eq:f15v99:Ft}.
Polarization proves the operator identity.
Subtract the compressed product
\((\widehat S_\gamma^*M_\times\widehat S_\gamma)^2\)
from
\(\widehat S_\gamma^*M_\times^2\widehat S_\gamma\)
to prove the second. Its positivity also
follows from the conditional variance
inequality in \eqref{eq:f15v99:Ft}.
\end{proof}

The square of the compressed magnetic factor
is thus not its compressed square.
If a bridge Peter--Weyl projection \(P_N\)
is imposed, its Gram is
\(P_NM_{F_{\gamma,t}}P_N\) on \(\mathcal K_N\).
Multiplication by \(F_{\gamma,t}\) need not
preserve that cutoff; it must not be treated
as its diagonal restriction to a list of
electric labels.

\subsection{The balanced internal embedding and its ground covariance}
\label{f15v99:balanced-section}

The transformation in \eqref{eq:f15v99:theta}
can equally be viewed as a gauge change on
the whole original graph: bridges become
\(Z\), and internal edges become
\[
 \overline X_e=h_s^{-1}X_eh_t,
\]
including the tree edges, for which the old
\(X_e\) was \(I\). Plaquette traces are unchanged.
This is a pointwise rewriting: the internal
integration variables and their density remain
\(X\). No Haar Jacobian for \(X\mapsto\overline X\)
is being asserted.
For one Cartesian component its linear
internal-edge embedding is
\begin{equation}
 \mathsf U=\mathsf J+\mathsf B\mathsf F
     =\mathsf Z^{\mathsf T}\mathsf M^{-1},\qquad
 \mathsf U^{\mathsf T}\mathsf U=\mathsf M^{-1}.
 \label{eq:f15v99:balanced}
\end{equation}
Here \(\mathsf B,\mathsf J,\mathsf Z,\mathsf M\)
are the single-cube matrices of V97.
Indeed \(\mathsf B^{\mathsf T}\mathsf U=0\)
and \(\mathsf Z\mathsf U=I\), which uniquely
specify the cycle-space representative.
Thus \(\mathsf U\mathsf M^{1/2}\) is an
isometry from the five chord coordinates
into the twelve-edge space.

Let \(\mathsf C\) be the one-Cartesian-component
covariance of \(|\Omega_\Box|^2\).
With
\[
 \mathsf T=\mathsf M^{1/2}\mathsf H\mathsf M^{1/2},
 \qquad a(\mathsf T)=\sqrt{\mathsf T+\mathsf T^2/4},
\]
it is
\begin{equation}
 \mathsf C=\mathsf M^{1/2}
                [2a(\mathsf T)]^{-1}\mathsf M^{1/2}.
 \label{eq:f15v99:covariance}
\end{equation}
For a scalar oscillator of magnetic parameter
\(q\), inserting \(e^{-a x^2/2}\) in its
Gaussian kernel gives \(a^2=q+q^2/4\).
Its squared normalized ground vector has
variance \(1/(2a)\). Diagonalize \(\mathsf T\)
as in V97 to prove
\eqref{eq:f15v99:covariance}; this is also the
ground term of the
\href{https://dlmf.nist.gov/18.18.E28}{Mehler expansion}.
The five parameters are \(4,4,4,6,6\), so
\begin{align}
 \mathsf U\mathsf C\mathsf U^{\mathsf T}
       &\leq\frac1{4\sqrt2}I_{12},\notag\\
 \Tr(\mathsf M^{-1}\mathsf C)
       &=\frac3{4\sqrt2}+\frac1{\sqrt{15}}.
 \label{eq:f15v99:covariance-bounds}
\end{align}
All matrices in these displays act on path
indices. Tensoring by \(I_3\) restores the
three Cartesian components.

Let \(\mathsf C_{\rm int}\) and \(\mathsf D\)
be the signed mixed-plaquette incidence
matrices on, respectively, the \(12B\)
internal edges and \(12B\) bridges.
Define
\begin{equation}
 \mathsf A=\mathsf C_{\rm int}
                     \left(\bigoplus_b\mathsf U\right),
 \qquad
 \mathsf C_B=\bigoplus_b\mathsf C,\qquad
 \mathsf K=\mathsf A\mathsf C_B\mathsf A^{\mathsf T}.
 \label{eq:f15v99:response-matrices}
\end{equation}
Thus the linear mixed curvature is
\(\mathsf A x+\mathsf D y\), not the
unbalanced expression obtained by forgetting
the moving frame.

\begin{proposition}[Uniform size and exact trace of the magnetic response]
\label{f15v99:matrix-bounds}
One has
\begin{equation}
 0\leq\mathsf K,\qquad
 \|\mathsf K\|\leq\rho:=\frac1{\sqrt2}<1,\qquad
 \operatorname{rank}\mathsf K\leq5B,
 \quad
 \Tr\mathsf K
   =B\left(\frac3{2\sqrt2}+\frac2{\sqrt{15}}\right)>B.
 \label{eq:f15v99:matrix-bounds}
\end{equation}
The block diagonal of
\(\mathsf A^{\mathsf T}\mathsf A\) is
\(2\mathsf M^{-1}\) in every block; its
off-diagonal blocks only join coarse neighbours.
\end{proposition}

\begin{proof}
A plaquette crossing one block interface has
two internal edges, one in each neighbouring
block. A plaquette crossing in both directions
has no internal edge. Every internal edge
belongs to exactly two mixed plaquettes.
Thus the maximum absolute row and column
sums of \(\mathsf C_{\rm int}\) are both two,
and \(\|\mathsf C_{\rm int}\|\leq2\).
Use \eqref{eq:f15v99:covariance-bounds} to
obtain \(\|\mathsf K\|\leq4/(4\sqrt2)\).

Within one block each such row contains only
one internal edge. Consequently its diagonal
Gram block before insertion of \(\mathsf U\)
is \(2I_{12}\); afterwards it is
\(2\mathsf U^{\mathsf T}\mathsf U=2\mathsf M^{-1}\).
The support statement follows from the two
neighbouring blocks of the plaquette.
Since \(\mathsf C_B\) is block diagonal,
taking the trace gives exactly twice \(B\)
times \eqref{eq:f15v99:covariance-bounds}.
The strict lower bound follows already from
\(3/(2\sqrt2)>1\). Positivity and rank are
immediate from the factorization.
\end{proof}

\subsection{The full compact average has this Gaussian response at fixed volume}
\label{f15v99:limit-section}

In bridge coordinates set
\(Z_e=\operatorname{Exp}(y_e/\sqrt\gamma)\).
For \(t>0\) define
\begin{equation}
 \begin{split}
 F_t(y)={}&\det(I+(t/2)\mathsf K)^{-3/2}\\
 &{}\cdot\exp\!\left[
 -\frac t4\,y^{\mathsf T}
   \big(\mathsf D^{\mathsf T}(I+(t/2)\mathsf K)^{-1}
                              \mathsf D\otimes I_3\big)y\right].
 \end{split}
 \label{eq:f15v99:gaussian-response}
\end{equation}

\begin{theorem}[Native-to-Gaussian magnetic response, with absolute-error scope]
\label{f15v99:limit}
For \(t=1,2\), fixed \(R<\infty\), and all
sufficiently large \(\gamma\),
\begin{equation}
 \sup_{\max_e|y_e|\leq R}
 \left|F_{\gamma,t}(\operatorname{Exp}(y/\sqrt\gamma))
                              -F_t(y)\right|
 \leq\frac{C_t B(1+R^3)}{\sqrt\gamma}
                          +C_t B e^{-c\gamma}.
 \label{eq:f15v99:limit}
\end{equation}
Constants may depend on the fixed single-cube
geometry but not on \(B\).
This is an absolute error. No volume-uniform
relative estimate or global bridge-field
approximation is asserted.
\end{theorem}

\begin{proof}
In the internal chart, the probability
\(|\Phi_\gamma|^2dH^{5B}\) is a product of
truncated Gaussian ground densities exactly.
Replacing it by the untruncated product
changes a bounded expectation by at most
\(C B e^{-c\gamma}\), using the one-block
Gaussian tail and the product total-variation
bound.

On the complete internal chart, write the
mixed words after the exact balanced gauge
change. Each word contains a bounded number
of exponential factors, with arguments
linear in the internal coordinates of at
most finitely many adjacent cubes and in
its bridge coordinates. All coefficient
matrices are fixed single-cube matrices.
Taylor's theorem along a radial line and
unitarity of the factors give
\[
 \left|\frac{t\gamma}{2}s(\hbox{word})
       -\frac t4|\hbox{its linear curvature}|^2\right|
 \leq \frac{C_t}{\sqrt\gamma}
               (|x_{\rm local}|+|y_{\rm local}|)^3.
\]
First derivatives vanish. The quadratic
term is the squared norm of the linear
Lie algebra word, yielding precisely
\(\mathsf A x+\mathsf D y\).
The derivative bound is valid on the
whole internal chart, not only on a fixed ball.

Both the exact and quadratic plaquette
exponents are nonnegative.
The exponential is one-Lipschitz on
\([0,\infty)\), and telescoping the \(18B\)
factors bounds their product error by the
sum of these local errors.
Their expectations under either normalized
truncated or untruncated one-block Gaussian
have bounded third moments, independently
of \(B\) and sufficiently large \(\gamma\).
For \(\max_e|y_e|\leq R\) this gives the first
term in \eqref{eq:f15v99:limit}.
The Gaussian average left over is
\[
 \mathbb E_{x\sim N(0,\mathsf C_B\otimes I_3)}
               e^{-t|\mathsf A x+\mathsf D y|^2/4}.
\]
Complete the square, or integrate the
Gaussian image \(\mathsf A x\).
The determinant lemma and quadratic
completion give \eqref{eq:f15v99:gaussian-response},
including its power \(3/2\).
\end{proof}

For each fixed \(B\), the limit is strictly
positive on compact \(y\)-sets, so the same integrand estimate, divided by
this fixed-volume lower bound, proves convergence
of normalized conditional densities there.
The lower bound deteriorates with \(B\);
dividing the displayed error by \(F_t\)
without retaining that cost would be invalid.

\subsection{An extensive relative leakage survives the corrected frame}
\label{f15v99:leakage-section}

\begin{theorem}[Failure of a uniform small-leakage ground-only projection]
\label{f15v99:leakage}
The Gaussian response satisfies
\begin{equation}
 \frac{F_2(0)}{F_1(0)^2}\geq e^{B/30},
              \qquad F_1(0)\leq e^{-B/2}.
 \label{eq:f15v99:leakage-determinant}
\end{equation}
For each fixed \(B\), the exact native
magnetic multiplier consequently obeys
\begin{equation}
 \liminf_{\gamma\to\infty}
 \sup_{0\ne f\in\mathcal K}
 \frac{\|(I-\widehat P_\gamma)
                         M_\times\widehat S_\gamma f\|^2}
      {\|\widehat P_\gamma
                         M_\times\widehat S_\gamma f\|^2}
                       \geq e^{B/30}-1.
 \label{eq:f15v99:leakage-native}
\end{equation}
This concerns the full moving ground carrier
before a bridge-label cutoff. It is not a
lower bound on a Yang--Mills excitation energy
or a counterexample to the mass-gap goal.
\end{theorem}

\begin{proof}
Let \(\lambda_j\) be the nonzero eigenvalues
of \(\mathsf K\), all at most \(\rho<1\).
By \eqref{eq:f15v99:gaussian-response},
\[
 \log\frac{F_2(0)}{F_1(0)^2}
     =\frac32\sum_j
        \left(2\log(1+\lambda_j/2)-\log(1+\lambda_j)\right).
\]
For \(0\leq\lambda\leq1\),
\[
 \begin{split}
 2\log(1+\lambda/2)-\log(1+\lambda)
   &=\log\left(1+\frac{\lambda^2}{4+4\lambda}\right)\\
   &\geq\frac{\lambda^2}{4+4\lambda+\lambda^2}
    \geq\frac{\lambda^2}{9}.
 \end{split}
\]
The first inequality follows by integrating
\(1/(1+s)\), or from
\(\log(1+u)\geq u/(1+u)\).
Cauchy--Schwarz and
Proposition~\ref{f15v99:matrix-bounds} give
\[
 \sum_j\lambda_j^2
       \geq\frac{(\Tr\mathsf K)^2}{5B}\geq\frac B5.
\]
This proves the first bound.
Also
\(\log(1+\lambda/2)\geq\lambda/(2+\lambda)
\geq\lambda/3\), proving the second.

For the native assertion,
Theorem~\ref{f15v99:gram} writes the quotient
as the ratio of the integrals of
\(F_{\gamma,2}-F_{\gamma,1}^2\) and
\(F_{\gamma,1}^2\) against \(|f|^2dH\).
Its supremum is at least the pointwise
ratio at \(Z=I\): use coarse-gauge-invariant
functions supported in successively smaller
invariant neighbourhoods of its pure-gauge
orbit. Continuity, invariance, positivity of
the denominator and full support of Haar
measure justify this localization.
At fixed \(B\), Theorem~\ref{f15v99:limit}
gives convergence of the pointwise ratio
to \(F_2(0)/F_1(0)^2-1\).
Apply the first bound.
\end{proof}

The determinant factors are independent of
\(y\) in the Gaussian response; a common
scalar normalization may cancel from a
normalized transfer spectrum.
The conclusion is not that such an extensive
free-energy term precludes a gap.
It is that \(F_1^2\) cannot replace \(F_2\),
and that discarding the corresponding
magnetic image is not a uniformly small
relative projection error.

\subsection{Normalize the magnetic image instead}
\label{f15v99:dressing-section}

Since \(F_{\gamma,2}\) is bounded below at
each finite regulator and coupling, define
\begin{equation}
 \mathcal I_\gamma
     =M_\times\widehat S_\gamma
                           M_{F_{\gamma,2}^{-1/2}}.
 \label{eq:f15v99:dressed}
\end{equation}
The exact Gram identity gives
\(\mathcal I_\gamma^*\mathcal I_\gamma=I\).
Thus this is an isometry on the same coarse
physical carrier, incorporating all the
mixed plaquettes and their normalization.

For any bounded positive coarse operator \(C\),
\begin{equation}
 \begin{split}
 M_\times\widehat S_\gamma C
                       \widehat S_\gamma^*M_\times
    =\mathcal I_\gamma\big(
        M_{F_{\gamma,2}^{1/2}}C
        M_{F_{\gamma,2}^{1/2}}\big)\mathcal I_\gamma^*.
 \end{split}
 \label{eq:f15v99:dressed-gram}
\end{equation}
This follows by substitution, and identifies
the nonzero spectrum of this specified
magnetically weighted ground-carrier operator.
It does \emph{not} identify it with
\(T_{\rm full,\gamma}\): that would additionally
require a paid approximation of
\(T_{\rm cut,\gamma}\) by
\(\widehat S_\gamma C\widehat S_\gamma^*\),
including all internal complements.

For a finite bridge cutoff the correct
metric is instead
\[
 G_{\gamma,N}
     =P_NM_{F_{\gamma,2}}P_N|_{\mathcal K_N}.
\]
It is positive definite. The normalized
finite-carrier isometry is
\(M_\times\widehat S_\gamma|_{\mathcal K_N}
G_{\gamma,N}^{-1/2}\), and the corresponding
coarse operator is \(G_{\gamma,N}^{1/2}
C_NG_{\gamma,N}^{1/2}\).
One cannot move \(P_N\) through the multiplier
or substitute \(P_NM_{F_{\gamma,1}}P_N\)
squared for this Gram.

\subsection{The dressed conditional Gaussian is uniformly local}
\label{f15v99:posterior-section}

At fixed bridge \(y\), the squared density
of the dressed fast vector in the limit is
proportional to
\[
 \exp\left[-\frac12x^{\mathsf T}
                     (\mathsf C_B^{-1}\otimes I_3)x
           -\frac12|\mathsf A x+\mathsf D y|^2\right].
\]
Completing the square gives its path-index
covariance and mean:
\begin{equation}
 \mathsf C_{\rm post}
        =(\mathsf C_B^{-1}+\mathsf A^{\mathsf T}\mathsf A)^{-1},
 \qquad
 \mu(y)=-(\mathsf C_{\rm post}\mathsf A^{\mathsf T}
                                    \mathsf D\otimes I_3)y.
 \label{eq:f15v99:posterior}
\end{equation}
These are the normalized \(t=2\) response,
not the unperturbed independent cube covariance.

Put
\[
 \mathsf L=\mathsf C_B^{1/2}\mathsf A^{\mathsf T}
                              \mathsf A\mathsf C_B^{1/2}.
\]
It has the same nonzero eigenvalues as
\(\mathsf K\), has only nearest-neighbour
off-diagonal block entries, and
\(\|\mathsf L\|\leq\rho<1\).
Hence
\begin{equation}
 \mathsf C_{\rm post}
  =\mathsf C_B^{1/2}
        \sum_{r=0}^{\infty}(-\mathsf L)^r
                                  \mathsf C_B^{1/2}
 \label{eq:f15v99:neumann}
\end{equation}
converges in operator norm, uniformly in \(B\).
If \(E,F\) are block sets of coarse graph
distance \(d\), their coordinate projections
satisfy
\begin{equation}
 \|P_E\mathsf C_{\rm post}P_F\|
       \leq\frac{\|\mathsf C\|\,\rho^d}{1-\rho}.
 \label{eq:f15v99:locality}
\end{equation}
Indeed \(\mathsf L^r\) cannot connect sets
at distance greater than \(r\); sum the
remaining geometric norm bound. This estimate
has no factor from the number of blocks
in either set.

The bridge quadratic form is
\begin{equation}
 \mathsf H_{\rm eff}
      =\mathsf D^{\mathsf T}(I+\mathsf K)^{-1}\mathsf D,
 \qquad
 (1+\rho)^{-1}\mathsf D^{\mathsf T}\mathsf D
       \leq\mathsf H_{\rm eff}
       \leq\mathsf D^{\mathsf T}\mathsf D.
 \label{eq:f15v99:hessian}
\end{equation}
The inequalities are direct spectral bounds
on \(I+\mathsf K\).
Woodbury's identity gives
\[
 (I+\mathsf K)^{-1}
     =I-\mathsf A\mathsf C_{\rm post}\mathsf A^{\mathsf T}.
\]
Together with \eqref{eq:f15v99:locality}
and the bounded, finite-range incidence
factors, this proves exponential locality
of the effective Hessian entries and of
the conditional mean kernel, with a fixed
change in the distance and constants.

There is also a uniformly summable
normalizer expansion:
\begin{equation}
 \frac1B\log F_2(0)
     =-\frac3{2B}\sum_{r=1}^{\infty}
                   \frac{(-1)^{r+1}}r\Tr(\mathsf L^r).
 \label{eq:f15v99:log-normalizer}
\end{equation}
Each trace is a sum of closed nearest-block
walks, with the five-dimensional internal
matrices retained. After order \(q\), the
absolute remainder per block is at most
\[
 \frac{15}{2}\,
       \frac{\rho^{q+1}}{(q+1)(1-\rho)}.
\]
Use \(|\Tr\mathsf L^r|\leq5B\rho^r\) to
prove this bound and absolute convergence.
This gives a local accounting of the leading
vacuum normalization, not a native non-Gaussian
polymer expansion.

\subsection{The magnetic response does not manufacture a coarse mass}
\label{f15v99:soft-section}

\begin{proposition}[Exact Gaussian null space and long-wave stiffness]
\label{f15v99:soft}
Per Cartesian component,
\(\ker\mathsf H_{\rm eff}=\ker\mathsf D\)
has dimension \(B+2\).
It consists of coarse gradients and three
constant directional holonomies, represented
by equal values on each set of four parallel
bridges. There are vectors perpendicular
to this kernel with Rayleigh quotient
\begin{equation}
 \frac{\langle y,\mathsf H_{\rm eff}y\rangle}{\|y\|^2}
                           \leq2\sin^2(\pi/m).
 \label{eq:f15v99:soft}
\end{equation}
\end{proposition}

\begin{proof}
The equality of kernels follows from the
strict positivity of \((I+\mathsf K)^{-1}\).
A decorated-digon row of \(\mathsf D\)
is a difference between two parallel
bridges at adjacent transverse ports.
Those four ports form a connected square,
so vanishing of all such rows forces
a common value \(a_\mu(b)\) on the four
bridges of each coarse edge.

The quadrilateral rows then give the
ordinary coarse lattice curl of \(a\);
there are two copies for each direction
pair. A curl-free periodic edge field is
a gradient plus its three periods:
subtract constant directional fields to
remove the periods, and integrate the
remaining field along lattice paths.
Vanishing plaquette curls and periods
make this path-independent.
Gradients have dimension \(B-1\), and
the three periods are independent.
This proves the null-space description
and its dimension.

For a test outside that kernel set
\(a_2(b)=\cos(2\pi b_1/m)\) and the other
two components to zero; repeat each value
on its four parallel bridges.
The field has zero mean and zero discrete
divergence, hence is orthogonal to periods
and gradients. All digon rows vanish.
There are two copies of the nonzero coarse
curl but four copies of the edge field, so
\[
 \frac{\|\mathsf D y\|^2}{\|y\|^2}
      =\frac12|e^{2\pi i/m}-1|^2
      =2\sin^2(\pi/m).
\]
The real cosine has the same quotient.
Apply the upper bound in
\eqref{eq:f15v99:hessian}.
\end{proof}

This is a statement about the linearized
static magnetic Hessian on its gauge
quotient, not a construction of gapless
native physical Hamiltonian states.
Nonlinear compact holonomies and the
coupled transfer dynamics are not replaced
by this Gaussian calculation.
It does show why uniform locality of the
fast conditional response alone is not
the missing infrared mass estimate.

\subsection{Direction, verification and preserved scope}
\label{f15v99:scope-section}

The upstream change is now explicit:
use the magnetically dressed Gram
\(F_{\gamma,2}\), retaining the moving frame
and the full conditional fast response,
rather than iterating a ground-only
\(F_{\gamma,1}\) compression.
The leading determinant and response have
a local, volume-independent organization,
while the exact compact identities remain
valid for every bridge configuration.

The unresolved step is to control the
\emph{native} dressed transfer, including
excited internal bands, its bridge-field
tails and its nonlinear neutral collective
modes, with a relative estimate uniform
in spatial volume. The absolute
\(B/\sqrt\gamma\) comparison cannot be
divided by exponentially small normalizers
without an additional argument.
Nor does the positive conditional fast
precision control the long-wave directions
of Proposition~\ref{f15v99:soft}.
Independent physical time calibration
and interacting four-dimensional continuum
reconstruction remain separate open tasks.

The checker
\begin{center}\small
\path{scripts/verify_ym_f15_v99_magnetic_response.py}
\end{center}
checks the balanced cycle embedding;
the exact radical coefficients of the
Gaussian covariance; all-block mixed
incidence and trace identities; coarse
gradient and holonomy null vectors;
noncommuting moving-frame words; conditional
Gram and leakage algebra; Gaussian
completion; determinant inequalities;
Neumann remainders and finite propagation;
and explicit transverse stiffness quotients.
The numerical fixtures use exact rational
arithmetic and supplement the analytic
proofs; they do not prove the full mass gap.
All 36,245 new exact-rational and combinatorial
fixtures passed, together with the inherited
V98 chain. The checker has 9,172 bytes and SHA--256
\begin{center}\scriptsize\ttfamily
81F677AB3D591F40ACB7CE8F0E0763454F5E332A8E5202D27C959AC4A9E89872.
\end{center}

The predecessor V98 has 2,267,905 bytes and SHA--256
\begin{center}\scriptsize\ttfamily
7FF35943F74A91FE1BBFDDB3321F7836D8454C92BBC9D1898C946AE46FBAE50E.
\end{center}
Its body is retained byte for byte; removing
this V99 section and restoring the title
recovers V98. Only the active main source
is replaced. All 163 previously protected
files remain unchanged. The companion
checkpoint and the planned archive/restart
are both due at V100.
Only \path{.tex} files belong in
\path{latex_notes}; builds are placed in
\path{artifacts/ym_f15_v99_texcheck/}.

\begin{firewall}
V99 proves exact moving-frame magnetic Gram
and dressing identities, a native fixed-volume
Gaussian response limit, an extensive
relative leakage bound for the undressed
ground carrier, and uniform locality of
the leading dressed fast covariance.
These results do not establish a volume-uniform
native transfer gap, a physical clock or
an interacting continuum Yang--Mills theory.
The full mass-gap goal remains unachieved.
\end{firewall}


\section{V100: the excitation density of the magnetic image}
\label{f15v100:context}

The V99 magnetic image is normalized, but it is not close to the old
product ground carrier uniformly in volume. We now quantify which
internal oscillator levels it actually occupies. This determines a
usable band budget and rules out every subextensive total-degree
budget in the leading response, not just the vacuum band.
It does not prove invariance of that band under the full transfer.
The Hermite generating method is already used in V87; the new object
is the full V98 spatial assembly with V99's correlated magnetic image.

\subsection{The required companion checkpoint}

The current inventory still contains SM continuum V72, NS stability V26,
shell-mixing V16 and parallel YM V5. The supplied newer related notes
are SM source-selection V31, regenerative stationarity V34 and
exact-action Einstein bridge V24. Their terminal mathematical sections
were reread; all 164 previously protected source and diagnostic files
have unchanged hashes at this checkpoint.

SM V72 proves a fixed-mode, fixed-time many-copy limit, not a
cutoff-uniform physical finite-species limit. NS V26 concerns radial
rank-one kernel contractions; its numerical/global comparisons refer
onward to V27 certification and provide no YM coercivity.
Shell V16 gives a complete measured flat-holonomy one-loop coefficient,
not the noncommuting remainder or a common volume/coupling estimate.
Parallel YM V5 excludes a volume-uniform tail for an extensive global
topological count. It recommends density or local-window control.
That distinction is directly useful here: the internal excitation
count should also be allowed to be extensive.

Source-selection V31 still requires word-native typing operations;
stationarity V34 leaves the moving-fibre derivative entries unevaluated;
Einstein bridge V24 supplies a nonzero first-jet derivative, not an
upstream nonlinear root. None closes the missing native YM transfer
estimate. The supplied literature suggestions remain a source of
methods, not additional instructions or established theorems.
The scope-checked literature audit in V91 and the exact boundary
fusion work in V82--V99 are not repeated.

\subsection{The complete internal oscillator counting problem}

Use the matrices \(A,D,C_B,K,L\) of V99, where
\[
 K=AC_BA^{\mathsf T},\qquad
 L=C_B^{1/2}A^{\mathsf T}AC_B^{1/2},\qquad
 0\le L,\quad \|L\|\le\rho=1/\sqrt2,\quad
 \Tr L>B,\quad \operatorname{rank}L\le5B .
\]
The matrices here first act on path indices. A tensor product by
\(I_3\) restores colour components. There are \(n=15B\) internal real
coordinates. Write \(x=(C_B^{1/2}\otimes I_3)u\) and set
\[
 \mathsf L=L\otimes I_3,\quad Q=I+\mathsf L,\quad
 b=(C_B^{1/2}A^{\mathsf T}D\otimes I_3)y,\quad
 \mu=-Q^{-1}b.
\]
In these coordinates the product ground vector and the normalized
magnetic image from V99 are respectively
\begin{equation}
 \Omega(u)=(2\pi)^{-n/4}e^{-|u|^2/4},\qquad
 \psi_y(u)=(2\pi)^{-n/4}\det(Q)^{1/4}
       e^{-(u-\mu)^{\mathsf T}Q(u-\mu)/4}.
 \label{f15v100:state}
\end{equation}
No independence of the dressed blocks is asserted.

Let \(\mathrm{He}_{\alpha}\) be the product probabilists' Hermite
polynomials and let
\(h_\alpha=\Omega\,\mathrm{He}_\alpha/\sqrt{\alpha!}\).
They are an orthonormal basis of \(L^2(\mathbb R^n)\). Define
\[
 \mathcal N h_\alpha=|\alpha|h_\alpha,\qquad
 P_{\le k}=1_{[0,k]}(\mathcal N).
\]
Equivalently \(\mathcal N=-\Delta+|u|^2/4-n/2\), initially on the
Hermite span and then by its nonnegative self-adjoint closure.
This is the total degree relative to the old product oscillator,
not a physical Hamiltonian and not a count of spatial topological charge.

Put
\begin{equation}
 S=\mathsf L(2I+\mathsf L)^{-1},\qquad
 d=-(2I+\mathsf L)^{-1}b,\qquad
 s_*=\frac{\rho}{2+\rho}<\frac13 .
 \label{f15v100:squeeze}
\end{equation}
Thus \(0\le S\le s_*I\).

\begin{theorem}[Exact normalized occupation generating function]
\label{f15v100:pgf}
For \(0\le z<s_*^{-1}\),
\begin{align}
 \langle\psi_y,z^{\mathcal N}\psi_y\rangle
 &=\left(\frac{\det(I-S^2)}{\det(I-z^2S^2)}\right)^{1/2}
 \notag\\
 &\quad\cdot
 \exp\!\left\{(z-1)d^{\mathsf T}
          (I+zS)^{-1}(I+S)^{-1}d\right\}.
 \label{eq:f15v100:pgf}
\end{align}
At \(z=0\) this denotes the ground projection. All levels, including
all colour and block multiplicities, are counted.
\end{theorem}

\begin{proof}
An orthogonal change of coordinates diagonalizes \(\mathsf L\) and
preserves \(\mathcal N\). In one eigen-coordinate write \(q=1+\lambda\),
\(s=(q-1)/(q+1)\), and \(d=q\mu/(q+1)\).
Multiplying the Hermite generating function
\[
 \sum_{j\ge0}h_j(u)\frac{w^j}{\sqrt{j!}}
          =\Omega(u)e^{uw-w^2/2}
\]
by the Gaussian in \eqref{f15v100:state} and integrating gives
\[
 \sum_{j\ge0}\langle h_j,\psi\rangle\frac{w^j}{\sqrt{j!}}
          =k\,e^{-sw^2/2+dw}.
\]
The standard generating identity is the same Hermite input used in
V87 and the
\href{https://dlmf.nist.gov/18.18.E28}{Mehler formula};
the integral here fixes its width and displacement.

The monomials \(w^j/\sqrt{j!}\) are orthonormal for
\(\pi^{-1}e^{-|w|^2}d^2w\) on \(\mathbb C\).
At argument \(\sqrt z\,w\), the real and imaginary Gaussian
precisions are \(1+zs\) and \(1-zs\). Hence
\[
 \sum_{j\ge0}z^j|\langle h_j,\psi\rangle|^2
      =k^2(1-z^2s^2)^{-1/2}
                    \exp\{zd^2/(1+zs)\}.
\]
At \(z=1\) the left side is one, which determines \(k^2\).
Divide by this normalization and use
\[
 \frac z{1+zs}-\frac1{1+s}
                   =\frac{z-1}{(1+zs)(1+s)}.
\]
The product over all \(n\) coordinates gives the assertion.
Convergence of the displayed complex Gaussian integrals for
\(z<s_*^{-1}\) also justifies the coefficient sums.
\end{proof}

At \(y=0\) the mean occupation is
\begin{equation}
 \langle\mathcal N\rangle_{\psi_0}
       =\frac34\Tr\!\left[L^2(I+L)^{-1}\right]
       \ge \frac{3B}{20(1+\rho)}.
 \label{f15v100:mean}
\end{equation}
Indeed \(\Tr L^2\ge(\Tr L)^2/(5B)\ge B/5\).
This bound alone would not rule out a large low-degree mass.
The next lower-tail estimate does.

\subsection{Subextensive old bands lose the state; an extensive band suffices}

\begin{theorem}[Two-sided scale of the total excitation budget]
\label{f15v100:budget}
Define
\[
 c_-=\frac{9}{40(2+\rho)^2},\qquad
 c_+=\frac{15}{2}
       \log\frac{1-s_*^2}{1-4s_*^2}.
\]
For every integer \(k\ge0\),
\begin{equation}
 \|P_{\le k}\psi_0\|^2
            \le \exp\{k\log2-c_-B\},
 \label{f15v100:lower-tail}
\end{equation}
and for all bridge vectors \(y\),
\begin{equation}
 \|(I-P_{\le k})\psi_y\|^2
 \le \exp\left\{c_+B+\frac{\rho}{4}\|Dy\|^2-k\log2\right\}.
 \label{f15v100:upper-tail}
\end{equation}
Here \(D\) in a colour expression means \(D\otimes I_3\).
In particular any \(k=o(B)\) captures a vanishing fraction at \(y=0\).
Conversely an integer
\begin{equation}
 k\ge\frac{c_+B+(\rho/4)\|Dy\|^2+\log(\delta^{-1})}{\log2}
 \label{f15v100:choose-k}
\end{equation}
loses at most \(\delta\), for \(0<\delta<1\).
\end{theorem}

\begin{proof}
For \(0<z<1\) and an eigenvalue \(s\) of \(S\),
\[
 \log\frac{1-s^2}{1-z^2s^2}
 =-\log\left(1+\frac{(1-z^2)s^2}{1-s^2}\right)
 \le -(1-z^2)s^2 .
\]
Use \(\log(1+v)\ge v/(1+v)\) for \(v\ge0\).
Moreover
\[
 \Tr S^2
   =3\Tr[L^2(2I+L)^{-2}]
   \ge\frac{3B}{5(2+\rho)^2}.
\]
At \(y=0\), \eqref{eq:f15v100:pgf} and \(z=1/2\) therefore give
\(\langle2^{-\mathcal N}\rangle\le e^{-c_-B}\).
On degrees at most \(k\), \(2^{-\mathcal N}\ge2^{-k}\).
This proves \eqref{f15v100:lower-tail}.

For \(z=2\), the logarithm of the determinant factor is at most
\(c_+B\), since \(n=15B\) and
\((1-s^2)/(1-4s^2)\) increases with \(s^2\).
Both resolvents in the exponent are contractions, so its exponent
is at most \(\|d\|^2\). From \eqref{f15v100:squeeze},
\[
 \|d\|^2\le\frac14\|b\|^2
       \le\frac{\rho}{4}\|Dy\|^2.
\]
Apply the upper Markov bound for \(2^{\mathcal N}\).
Using \(k\) rather than \(k+1\) only weakens the result.
\end{proof}

There are \(12B\) decorated-digon rows of \(D\), each with two
bridge terms, and \(6B\) quadrilateral rows with four terms.
Consequently \(\max_e|y_e|\le R\) implies
\(\|Dy\|^2\le144BR^2\).
Thus one budget \(k=O(B(1+R^2)+\log\delta^{-1})\) works on
that entire bridge window. The old total-degree space has dimension
\[
                    \operatorname{rank}P_{\le k}
                          =\binom{15B+k}{k}.
\]
This is an extensive band, not a uniformly bounded number of old
oscillator excitations. It is also not a lower bound on the rank
of an adaptive carrier: the dressed vector itself spans one line
at each fixed \(y\). The result specifies the cost of using the
old product oscillator basis, not an intrinsic Hilbert-space
dimension obstruction.

\subsection{A genuine compact projection and the order of limits}

A Euclidean projection restricted to a compact chart is generally
not a projection. We state the correct finite-volume handoff
explicitly. Let \(E_\gamma\subset\mathbb R^{15B}\) be the image in
\(u\)-coordinates of the full five-chord chart in every block.
It exhausts \(\mathbb R^{15B}\) as \(\gamma\to\infty\).

For fixed \(B,k\), restrict all \(h_\alpha\), \(|\alpha|\le k\),
to \(E_\gamma\). Write \(V_{\gamma,k}\) for the map from their
coefficient space to \(L^2(E_\gamma)\), and put
\begin{equation}
 G_{\gamma,k}=V_{\gamma,k}^*V_{\gamma,k},\qquad
 \Pi_{\gamma,k}
       =V_{\gamma,k}G_{\gamma,k}^{-1}V_{\gamma,k}^*.
 \label{f15v100:chart-projection}
\end{equation}
The Gram is positive definite: a polynomial times a nonzero
Gaussian cannot vanish on the nonempty open chart unless all
its coefficients vanish. Thus this is an orthogonal projection.
The Haar chart isometry and the exact moving-frame map carry
it to a genuine internal projection in compact variables.

\begin{proposition}[Fixed-volume native meaning of the band estimates]
\label{f15v100:native-band}
Let \(\psi_{\gamma,y}\), expressed in \(u\)-coordinates and extended
by zero, be the normalized conditional fast vector of V99's exact
dressed isometry at \(Z_e=\operatorname{Exp}(y_e/\sqrt\gamma)\).
For fixed \(B,k,R\),
\begin{equation}
 \sup_{\max_e|y_e|\le R}
 \left|
   \|\Pi_{\gamma,k}\psi_{\gamma,y}\|^2
                  -\|P_{\le k}\psi_y\|^2
 \right|\longrightarrow0 .
 \label{f15v100:native-limit}
\end{equation}
The compact projection respects the remaining root-gauge action
when used with all Hermite multiplicities and then restricted
to the physical subspace.
\end{proposition}

\begin{proof}
V99's proof gives \(L^1\) convergence of the normalized conditional
densities on fixed bridge windows at fixed \(B\): the integrand
error is divided by a strictly positive fixed-volume lower bound.
All amplitudes are nonnegative, and
\((\sqrt a-\sqrt b)^2\le|a-b|\), so
\(\psi_{\gamma,y}\to\psi_y\) in \(L^2\), uniformly on that window.
Each of the finitely many restricted Hermite vectors converges
in \(L^2\) to the full one. Hence \(G_{\gamma,k}\to I\) in matrix
norm, and the projections in \eqref{f15v100:chart-projection},
extended by zero, converge in operator norm to \(P_{\le k}\).
The triangle inequality proves the formula.

At a block, a root gauge rotation acts orthogonally on all three
colour coordinates. The chart, Gaussian, and total-degree
polynomial space are invariant under this action.
Their Gram and its inverse therefore commute with it.
The moving-frame map is gauge equivariant by V99.
The transported projection consequently preserves Gauss
invariance, with no removal of nontrivial internal multiplets.
\end{proof}

For the extensive choice \(k=k(B,R,\delta)\), the proposition
is applied separately for each fixed \(B\). It does not supply
a rate uniform in \(B\) or permit an exchange of limits.
It concerns the magnetic image, not the norm of the whole
native transfer complement. The volume-independent local
separation of V97's internal oscillator is compatible with
an extensive number of its excitations in a differently
weighted vector.

\subsection{Archive decision and the next actual transfer calculation}

This checkpoint closes the fifteenth lineage at V100.
It is archived as
\begin{center}\small
\path{Renewal_Yang_Mills_Mass_Gap_Research_Notes_V100_f15.tex}.
\end{center}
The new active V1 will summarize the mathematical interfaces
and continue with the overlap of conditional dressed vectors
under the actual bridge kinetic kernel. This is needed because
a normalized family is not automatically preserved by one
kinetic step. Counting its old excitations is not a substitute
for computing that overlap or for controlling the internal
excited-state transfer.

The checker
\path{scripts/verify_ym_f15_v100_excitation_density.py}
compares Hermite coefficients with an independent generating-function
differential recurrence, verifies occupation moments and matrix
resolvent identities, the algebraic Chernoff inputs, complete
degree ranks, and the difference between a true Gram-normalized
projection and a raw chart restriction. These exact finite
checks supplement the analytic proofs and do not certify a
continuum theorem. All 2,433 new exact fixtures passed, together
with the inherited V99 chain. The checker has 4,530 bytes and SHA--256
\begin{center}\scriptsize\ttfamily
3F61190A0EBD2BF8E1B739B2DA908D365C857AA6EF2915F0F81B09B803FC0777.
\end{center}

The predecessor V99 has 2,295,961 bytes and SHA--256
\begin{center}\scriptsize\ttfamily
7259F37FBBA1B4ED00DDE8071B81117CBF8FCC19F90D9F802C0B98488A781325.
\end{center}
Its complete body is preserved; deleting this append and restoring
the title recovers those bytes. The archive is retained, not
replaced by the compact summary. Only \texttt{.tex} files belong
in \texttt{latex\_notes}; compilation and visual QA remain outside it.

\begin{firewall}
V100 gives the exact excitation generating function of the
leading correlated magnetic image, a necessary extensive old-band
budget, a sufficient budget with an exponential tail, and a
fixed-volume compact projection limit. It proves no
volume-uniform native transfer gap, physical-time calibration,
or interacting four-dimensional continuum Yang--Mills construction.
The full mass-gap goal remains unachieved.
\end{firewall}

\end{document}
